\documentclass[10pt]{article}
\usepackage[letterpaper,margin=0.92in,footskip=0.46in]{geometry}
\usepackage{amsmath,amssymb,amsthm,mathtools}
\usepackage[unicode,colorlinks=true,linkcolor=blue,urlcolor=blue,bookmarksnumbered=true]{hyperref}
\hypersetup{pdftitle={Ordinary RH through fourth-power samples: a working exposition}}
\usepackage{xurl}
\setcounter{tocdepth}{1}
\setlength{\emergencystretch}{2em}
\allowdisplaybreaks[2]
\DeclareRobustCommand{\srcid}[1]{\ifmmode\text{\nolinkurl{#1}}\else\nolinkurl{#1}\fi}
\newtheorem{srtheorem}{Theorem}[section]
\newtheorem{srproposition}[srtheorem]{Proposition}
\newtheorem{srlemma}[srtheorem]{Lemma}
\newtheorem{srcorollary}[srtheorem]{Corollary}
\theoremstyle{definition}
\newtheorem{srdefinition}[srtheorem]{Definition}
\theoremstyle{plain}
\title{Ordinary RH through fourth-power samples\\[4pt]
\large Proved reductions and the remaining arithmetic--geometric inputs}
\author{A working exposition of Alpöge's formal development\\
with reductions from Principia contributors}
\date{21 November 2026}
\begin{document}
\maketitle
\begin{abstract}
This working exposition organizes the conductor-one argument around
fourth-power prime sampling, two uniform character estimates, and a
compatible geometric construction on common windows.
The analytic ending, signed finite comparison, minus and dyadic estimates,
and a global high-twist bound are proved in prose.
The latter uses published arithmetic K-theory comparison and Hurewicz
theorems, with the integral counts printed here.
We construct classical ordinary Eisenstein seeds, their rank-two and
rank-three unitary realizations, an arithmetic frame with its cohomology
bound chosen before cyclotomic depth, and the actual entry algebra,
flat local charts, selected curves and sharp ideal.
The ordinary modules and finite-level Selmer estimates are given;
rank-four descent and the conditional rank-three endpoints are matched
with their full local hypotheses.
On the theta side we construct every finite matching map, prove the
exact good-place coinvariant condition and both ordinary-convergence
and decay clauses, and retain the actual measures, vacuum vectors and
bar-model tests. The canonical cusp-weight mean and quadratic CM
period are computed.
The two uniform character estimates still require the remaining
common-support and local original-level producers.
The archimedean matching, global theta comparison and rank recursions,
and uniform integral-model, slope and metric estimates remain to be
completed on the same windows. Part II preserves the supporting
longer-route translations with their explicit hypotheses.
This is not yet a completed standalone proof of RH.
\end{abstract}
\begingroup\small
\tableofcontents
\endgroup
\clearpage
\section*{Scope, sources and verification status}
\addcontentsline{toc}{section}{Scope, sources and verification status}
This is a consolidated working exposition, not a completed unconditional
proof of RH. The prime-only analytic ending and signed finite comparison
are printed first. The elementary arithmetic packet, minus and dyadic
bounds, global high-twist bound and common-integer finiteness arguments
are proved next. Both uniform character estimates are proved in
Theorem~\ref{thm:sr-generic-character-estimate} and
Proposition~\ref{prop:sr-nontrivial-final-assembly}.
Corollary~\ref{cor:sr-principal-plus-height-proved} assembles
their polynomial plus-height bound.

The deformation argument now includes actual classical ordinary seeds,
their same-module unitary insertion, the cyclotomic arithmetic frame,
its uniform class-rank and strict-cohomology bounds, and the actual entry
space with flat charts, selected curves and sharp ideal.
The rank-four tower witnesses and exact descent are proved in
Section~\ref{sec:sr-rank-four-actual-endpoint}.
Section~\ref{sec:sr-rank-three-host-endpoints} gives finite host algebras
and exact automorphic endpoints from the stated common support output.
Section~\ref{sec:actual-seed-two-propagations} proves the common
rank-two/rank-three propagation output. The actual local crystalline
component and original-level comparison
are proved, including the nearest-weight case. The rank-four common space, good residual lattice, actual modular
ideal, residual paths and nonsplit block extension are constructed.
Its final finite-weight and no-block arithmetic endpoint is proved
in Proposition~\ref{prop:sr-four-actual-final-endpoint}.
Proposition~\ref{prop:sr-gaussian-actual-g10} specializes the
same rank-two/rank-three construction to the actual Gaussian
frame, including coefficient prime three and the boundary weight. No uniform character estimate is inferred
from a seed alone.

The geometric chapters construct the archimedean common Fock quotient
and all finite local matching maps,
including inert places in either rank parity, and prove the exact
good-place coinvariant condition.
Subsection~\ref{subsec:sr-theta-decay-closed} proves both ordinary
convergence and the small-side rapid-decay and moderate-growth clauses
with the specified bar-model test.
The canonical cusp-weight mean has constant four and the quadratic
CM period is evaluated. Section~\ref{sec:sr-compact-global-comparison}
also proves the nonzero global comparison at the compact base.
The finite cone vacuum, absolutely convergent adelic Euler
factorization and four global constant-term producers are also proved.
The all-signature global comparison, quotient descent, rank recursions
and the general integral-model, saturation and covolume bound
remain to be completed. The universal principal Lie-bundle
slope estimate is proved in
Theorem~\ref{thm:sr-general-principal-slope}.

The term \emph{connected} in the formal family predicate means one
compatible eventual window witness, as explained in
Section~\ref{sec:sr-latest-community-check}.
The required local inequalities and quantifier order are retained in
Definition~\ref{def:sr-supply}.
The second part preserves the earlier translations of the longer route.
Their explicit hypotheses remain in force, with the newer first-part
proofs discharging only the clauses identified in the closing ledger.

The shorter route follows Anton's Principia files
\href{https://principia-math.com/projects/d16ab288ed353bc5b10e292ace0ddbc1/editor?file=d59903a6498c7db021b728233ee6bdc2}{\srcid{FourthSamplesRH}},
\href{https://principia-math.com/projects/d16ab288ed353bc5b10e292ace0ddbc1/editor?file=dfeee1c759dca3c2b8e378a2dec2c780}{\srcid{PrimeFourthSamples}},
\href{https://principia-math.com/projects/d16ab288ed353bc5b10e292ace0ddbc1/editor?file=57674e7030ca8ce9bc2b4113ee37ab82}{\srcid{SignedSignal}},
\href{https://principia-math.com/projects/d16ab288ed353bc5b10e292ace0ddbc1/editor?file=150a1ccfbf0ccf9a0b4297260c71c05d}{\srcid{PointwiseGeometry}},
\href{https://principia-math.com/projects/d16ab288ed353bc5b10e292ace0ddbc1/editor?file=116c16d3654fa8b487602e67ead50b96}{\srcid{SignedGeometry}},
and
\href{https://principia-math.com/projects/d16ab288ed353bc5b10e292ace0ddbc1/editor?file=41bb998820fe961dbf8e023c9469578b}{\srcid{ArchiveThin}}.
The separate
\href{https://principia-math.com/projects/d16ab288ed353bc5b10e292ace0ddbc1/editor?file=c36fb16dbf85efd1520d208bf47edf3a}{\srcid{PrincipalWindows}}
interface isolates the conductor-one connected-family route.
Its explicit geometric premises must still be proved in a standalone text.
The 552-line specialization, its 630-line cubic-logarithmic revision,
the subsequent two-character height reduction, and the historical
five-PDF review are recorded in
Section~\ref{sec:sr-source-review}; the implications it proves have
not been mistaken for proofs of those premises.

The
\href{https://github.com/leanprover-community/riemann-hypothesis/pull/12}{signed-signal upstream supplement}
at commit \srcid{9430017510358b5ed1ab12c8bb900fd706e60a32}
has an independent successful Lean 4.33.1 source rebuild of its 41 project
dependencies. That verifies the supplement's implications with their
stated hypotheses. It does not verify this new prose edition or prove
its supply assertion.
The simplification of the real Mellin identity by
\href{https://principia-math.com/projects/d16ab288ed353bc5b10e292ace0ddbc1/editor?file=58b04d0292b512962984e43972698e31}{uimspeedrun}
preserves its statement, deriving it from the complex identity.
We use the same specialization principle, with the actual integral
calculation printed in the text.

Recent geometric translations by Baelson, reviewed by lowrank, include
\href{https://principia-math.com/api/rh/step?name=Development.AbelianVarietyTheory.Proofs.DivisorProofs.ArCalcOfPrincipalisation}{arithmetic intersection calculus from principalisation},
\href{https://principia-math.com/api/rh/step?name=Development.AbelianVarietyTheory.Proofs.DivisorProofs.SemiPullOfPrincipalisation}{semipositivity under pullback},
\href{https://principia-math.com/api/rh/step?name=Development.AbelianVarietyTheory.Proofs.DivisorProofs.RigidBoundaryChartsForFanUnitaryDatumStatement}{rigid toroidal boundary charts}, and
\href{https://principia-math.com/api/rh/step?name=Development.AbelianVarietyTheory.Proofs.DivisorProofs.ExistsFramesFockMapPolynomialLieActionsStatement}{Fock frames and Lie actions}.
They locate the required geometric dependencies.
The short re-export modules do not themselves contain the proofs of
all their imported results.

The ordinary-zeta oscillation argument in this edition excludes positive
real zeros with the alternating zeta series; it consequently omits the
general real-pole case in the Lean argument.
The elementary bounds here also give the constant four in
\eqref{eq:sr-Elip}. These prose variations, including the dyadic pro-\(2\) group proof and
spectral uniqueness, dyadic matrix and integral-homology arguments,
have not been
separately formalized or checked by Lean.
The class-number proof prints the one-ramified-prime argument and the
strict cyclotomic descent explicitly, using the classical background
listed in Section~\ref{sec:sr-arithmetic}.

The full formal export has also been typeset as an indexed record edition.
That is a record transcription, not a human exposition.
Its reverse-conversion audit establishes fidelity to the input records,
not completion of this paper.
The
\href{https://github.com/ChristianKudeba/alpoge-rh-verification/releases/tag/rh-literal-audit-v1}{record edition and audit}
and
\href{https://math.mit.edu/~drew/grh/}{independent formal verification report}
are separate evidence from the mathematical arguments printed here.


\subsection*{The logical route through the paper}
The analytic target is Theorem~\ref{thm:sr-prime}: a
fourth-power prime-sampling estimate implies the Riemann
hypothesis for ordinary zeta. Theorem~\ref{thm:sr-window}
permits that estimate to be supplied on the admissible
windows. The signed finite comparison in
Theorem~\ref{thm:sr-supply} reduces it to the arithmetic
and geometric witnesses in Definition~\ref{def:sr-supply}.
These are the two substantive branches of the argument.

On the arithmetic branch, the minus and dyadic parts are
already bounded. The plus part is reduced, with its prime
weights and quantifiers intact, to
Definition~\ref{def:sr-character-inputs}. A single exponent
must bound the sum over every finite generic prime set;
a second exponent must bound each nontrivial exceptional
or resonant character. Both exponents precede the varying
host, character order and twist.
Finiteness for one host is useful but does not establish
these uniform estimates.

The deformation chapters build the mechanism for those
estimates: actual ordinary seeds, an arithmetic frame,
curve and link presentations, transport in one ordinary
module, and automorphic descent. The class-group constant
in Section~\ref{sec:arithmetic-frame-uniform-depth} is
chosen before cyclotomic depth. It need not itself be
polynomial in the host. Its role is to make the auxiliary
frame sufficiently deep; the final numerical character
count must be made after descent, with its stated uniform
conductor and weight bounds.

On the geometric branch, one must choose compatible data
on the same windows, prove the global theta comparison
and its rank recursions, and control the integral lattices
and metrics of those data. The local maps, convergence
estimates and fixed-model identities in the theta chapters
are components of this construction. They do not separately
choose the compatible family or prove the uniform
covolume bound. Here \emph{connected} refers to the
source's compatible eventual window witness, as explained
in Section~\ref{sec:sr-latest-community-check}.

The closing ledger in Section~\ref{sec:sr-source-review}
records the exact frontier in these two branches.
Until both supply obligations there have been discharged,
the final RH conclusion remains conditional in this prose
edition. Part II preserves the earlier longer-route
translations with their hypotheses and makes their
supporting arguments available to the first part.

\part{The shorter route for ordinary zeta}
\section{A prime-only criterion for ordinary RH}
\label{sec:sr-prime}
Write \(\Lambda\) for the von Mangoldt function and
\(\gamma=\lim_{m\to\infty}(\sum_{r=1}^m r^{-1}-\log m)\).
For \(x\geq1\) put
\[
 T(x)=\sum_{n\leq x}\Lambda(n)(x/n-1),\qquad
 E(x)=T(x)-x\log x+(1+\gamma)x.
\]
All sums over \(n\) in this part are sums over positive integers. For an
integer \(m\geq1\) define
\[
 R(m)=\sum_{n\leq m}\Lambda(n)\left(1-\frac{m\bmod n}{n}\right),\quad
 P(m)=\sum_{\substack{p\leq m\\p\ {\rm prime}}}
 \log p\left(1-\frac{m\bmod p}{p}\right).
\]
The notation \(m\bmod n\) denotes the integer in \([0,n-1]\) congruent to
\(m\) modulo \(n\).

\begin{srtheorem}[Fourth-power prime criterion]\label{thm:sr-prime}
Suppose that for every \(\epsilon>0\) there is a constant
\(C_\epsilon\geq0\), independent of the subsequent interval, such that
arbitrarily large \(Y\) satisfy
\begin{equation}\label{eq:sr-prime}
 P(j^4)\geq\gamma j^4-C_\epsilon(j^4)^{1/2+\epsilon}
 \quad\hbox{whenever }j\in\mathbb N,\quad Y\leq j^4\leq Y^{3/2}.
\end{equation}
Then every nontrivial zero of the Riemann zeta function has real part
\(1/2\).
\end{srtheorem}
The hypothesis in this theorem is an arithmetic input. The theorem does not
assert it unconditionally. This section and the next two sections prove
the entire implication without importing a theorem from the Lean archive.
The criterion follows Anton's \texttt{FourthSamplesRH} and
\texttt{PrimeFourthSamples}; the exclusion of positive real zeros below
allows a shorter ordinary-zeta oscillation proof.

\subsection{Elementary bounds and the exact factorial identity}
Put \(\psi(x)=\sum_{n\leq x}\Lambda(n)\) and
\(\vartheta(x)=\sum_{p\leq x}\log p\). We first record a bound used several
times below:
\begin{equation}\label{eq:sr-cheb}
 0\leq\vartheta(x)\leq\psi(x)\leq3x\qquad(x\geq0).
\end{equation}
Indeed, for any integer \(a\geq1\), Legendre's formula gives
\[
 v_p\binom{2a}{a}=\sum_{r\geq1}
 \bigl(\lfloor2a/p^r\rfloor-2\lfloor a/p^r\rfloor\bigr).
\]
Every summand is nonnegative, and any \(p^r\in(a,2a]\) contributes one.
Consequently
\(\psi(2a)-\psi(a)\leq\log\binom{2a}{a}\leq2a\log2\).
Summing with \(a=1,2,\ldots\) along successive powers of two gives
\(\psi(2^k)\leq(2^{k+1}-2)\log2\). If \(x>1\), choose the least \(k\)
with \(x\leq2^k\); then \(2^k<2x\), and hence
\(\psi(x)<4x\log2<3x\). For \(0\leq x\leq1\) the sum is zero.
The elementary inequality \(4\log2<3\) follows, for example, by retaining
four positive terms in the series for \(\exp(3/4)>2\).

Since \(\log r=\sum_{n\mid r}\Lambda(n)\), interchanging two finite sums gives
\[
 \log(m!)=\sum_{n\leq m}\Lambda(n)\lfloor m/n\rfloor .
\]
Using \(m/n=\lfloor m/n\rfloor+(m\bmod n)/n\) now yields the exact identity
\begin{equation}\label{eq:sr-factorial}
 E(m)=\gamma m-R(m)+q(m),\qquad
 q(m)=\log(m!)-m\log m+m.
\end{equation}
Monotonicity of \(\log t\), applied on each unit interval, gives
\[
 m\log m-m+1\leq\log(m!)\leq m\log m-m+1+\log m.
\]
Thus \(1\leq q(m)\leq1+\log m\leq2\sqrt m\).
For the last inequality integrate \(1/t\leq1/\sqrt t\) from \(1\) to \(m\).
In particular,
\begin{equation}\label{eq:sr-integer}
 R(m)\geq\gamma m-E(m).
\end{equation}
A lower bound \(R(m)\geq\gamma m-Cm^{1/2+\epsilon}\) implies
\(E(m)\leq(C+2)m^{1/2+\epsilon}\).

Higher prime powers have a uniformly bounded cost:
\begin{equation}\label{eq:sr-discard}
 0\leq R(m)-P(m)\leq K\sqrt m,\qquad
 K=3+\frac{18}{\mathrm e\log2}.
\end{equation}
The first difference is nonnegative because every weight is in \([0,1]\).
For \(m\geq2\) it is at most
\[
 \psi(m)-\vartheta(m)
 =\sum_{r=2}^{\lfloor\log m/\log2\rfloor}\vartheta(m^{1/r})
 \leq3\sqrt m+\frac{3\log m}{\log2}m^{1/3}.
\]
The function \(v\exp(-v/6)\), for \(v\geq0\), has maximum \(6/\mathrm e\).
Taking \(v=\log m\) proves \eqref{eq:sr-discard}. The case \(m=1\) is zero.
It follows that the fourth-power criterion with \(R\) is equivalent to
\eqref{eq:sr-prime}, allowing a change of \(C_\epsilon\) by \(K\).

\subsection{Interpolation between fourth powers}
For any bounded interval of positive \(x\), the expression
\(T(x)=\sum_n\Lambda(n)(x/n-1)_+\) has only finitely many nonzero terms.
It is a convex continuous function. For \(0<a\leq x\leq b\), \(a<b\), put
\(\lambda=(b-x)/(b-a)\). Convexity and the identity defining \(E\) imply
\[
 E(x)\leq\lambda E(a)+(1-\lambda)E(b)
 +\lambda a\log a+(1-\lambda)b\log b-x\log x.
\]
The last three terms are at most
\((x-a)(b-x)/(2a)\). To see this, subtract that quadratic from the chord
of \(f(t)=t\log t\) minus \(f(x)\). The resulting function of \(x\) has
second derivative \(1/a-1/x\geq0\), and vanishes at both endpoints.
A convex function vanishing at both endpoints is nonpositive between them.
Therefore
\begin{equation}\label{eq:sr-interpolate}
 E(x)\leq\max(E(a),E(b))+\frac{(b-a)^2}{8a}.
\end{equation}
Take \(a=j^4,b=(j+1)^4\), with \(j\geq1\).
Then \(b\leq16a\) and
\(b-a=4j^3+6j^2+4j+1\leq15j^3\). The interpolation loss is at most
\(225\sqrt a/8\).

Fix \(\epsilon>0\) and an interval from \eqref{eq:sr-prime}.
Equation \eqref{eq:sr-factorial} gives
\(E(j^4)\leq(C_\epsilon+2)(j^4)^{1/2+\epsilon}\) at every supplied sample.
Set \(X=16Y\). For sufficiently large \(Y\),
\[
 16(16Y)^{5/4}\leq Y^{3/2}.
\]
If \(X\leq x\leq X^{5/4}\), its bracketing fourth powers satisfy
\(Y\leq a\leq x<b\leq Y^{3/2}\), since \(a\geq x/16\) and \(b\leq16x\).
Consequently \eqref{eq:sr-interpolate} gives throughout that interval
\[
 E(x)\leq
 \left((C_\epsilon+2)16^{1/2+\epsilon}+225/8\right)x^{1/2+\epsilon}.
\]
Increasing the lower threshold for \(Y\) absorbs the displayed constant
into \(x^\epsilon\). Thus for every \(\epsilon>0\) there are arbitrarily
late intervals
\begin{equation}\label{eq:sr-continuous}
 [X,X^{5/4}]\quad\hbox{on which}\quad E(x)\leq x^{1/2+2\epsilon}.
\end{equation}
Both the uniformity of \(C_\epsilon\) and the shortening of the interval
have been used explicitly.

\section{The zeta transform and its poles}\label{sec:sr-transform}
We give the analytic facts needed in the window argument. The underlying
tools are integration, Fourier series for a Gaussian, and Cauchy's
integral formula.

\subsection{Continuation, absence of positive real zeros, and reflection}
For \(\Re s>1\), the absolutely convergent series and Euler product give
\[
 \zeta(s)=\sum_{n\geq1}n^{-s}=\prod_p(1-p^{-s})^{-1},\qquad
 -\frac{\zeta'(s)}{\zeta(s)}=\sum_{n\geq1}\Lambda(n)n^{-s}.
\]
The product follows by expanding finitely many geometric series and then
using absolute convergence to pass to all primes. Its logarithm converges
absolutely, so the product is nonzero. Termwise differentiation of the
logarithm is locally uniform in this half-plane.

Partial summation gives
\begin{equation}\label{eq:sr-continuation}
 \zeta(s)=\frac{s}{s-1}
  -s\int_1^\infty \{t\}t^{-s-1}\,dt .
\end{equation}
The integral and all its derivatives converge locally uniformly for
\(\Re s>0\); its \(r\)-th derivative is dominated on
\(\Re s\geq a>0\) by a constant times
\(t^{-a-1}(1+\log t)^r\), which is integrable. This proves meromorphic
continuation there, with a simple pole of residue one at \(1\).
Furthermore
\[
 \int_1^m\{t\}t^{-2}\,dt=1+\log m-\sum_{r=1}^m1/r
\]
for integers \(m\geq1\). Taking the limit in this identity and then
expanding \eqref{eq:sr-continuation} at \(1\) shows
\begin{equation}\label{eq:sr-laurent}
 \zeta(s)=\frac1{s-1}+\gamma+O(s-1).
\end{equation}

The alternating series \(\eta(s)=\sum_{n\geq1}(-1)^{n-1}n^{-s}\) is
holomorphic for \(\Re s>0\). Summation by parts, using bounded partial
sums of the signs, bounds its tails uniformly on every compact subset
by a constant times \(M^{-a}\), where \(a\) is the minimum real part and
\(M\) the first tail index. In \(\Re s>1\), separating even and odd terms
gives \(\eta(s)=(1-2^{1-s})\zeta(s)\), and continuation gives the same
identity throughout \(\Re s>0\).
For real \(s>0\), pairing successive terms expresses \(\eta(s)\) as a
sum of positive terms, the first of which is \(1-2^{-s}>0\).
Thus \(\zeta(s)<0\) for \(0<s<1\); there are no positive real zeros.

There are no zeros on \(\Re s=1\) either. For \(\sigma>1\) and real \(t\),
the logarithmic Euler product and
\(3+4\cos u+\cos(2u)=2(1+\cos u)^2\geq0\) give
\[
 \zeta(\sigma)^3|\zeta(\sigma+it)|^4
 |\zeta(\sigma+2it)|\geq1.
\]
If \(t\ne0\) and \(\zeta(1+it)\) vanished to order \(m\geq1\), the left
side would be \(O((\sigma-1)^{4m-3})\) as \(\sigma\downarrow1\), since
\(\zeta\) is holomorphic at \(1+2it\). This tends to zero, a contradiction.
At \(t=0\) there is a pole. Hence every zero in \(\Re s>0\) lies in
\(0<\Re s<1\), is nonreal, and has finite positive multiplicity.
The zeros form a locally finite set: a nonzero holomorphic function has
isolated zeros, and a compact subset of its domain cannot contain an
infinite discrete zero set without an accumulation point.

For completeness, reflection can be recovered as follows.
The Gaussian Fourier transform is
\[
 \int_{\mathbb R}e^{-\pi t x^2}e^{-2\pi i\xi x}\,dx
 =t^{-1/2}e^{-\pi\xi^2/t}\qquad(t>0).
\]
The value at zero is the Gaussian integral, obtained by squaring and
integrating in polar coordinates. Differentiation in \(\xi\) and
integration by parts give the differential equation
\(F'(\xi)=-(2\pi\xi/t)F(\xi)\), proving the formula.
The smooth periodization \(\sum_{n\in\mathbb Z}e^{-\pi t(x+n)^2}\)
has these values at integral \(\xi\) as its Fourier coefficients.
Both the periodization and the resulting Fourier series converge
uniformly with all needed derivatives. Evaluation at \(x=0\) gives
\[
 \Theta(t)=\sum_{n\in\mathbb Z}e^{-\pi n^2t}
 =t^{-1/2}\Theta(1/t).
\]
Mellin integration, initially for \(\Re s>1\), and a split at \(t=1\)
now give
\begin{equation}\label{eq:sr-reflection}
 \pi^{-s/2}\Gamma(s/2)\zeta(s)
 =-\frac1s+\frac1{s-1}
 +\frac12\int_1^\infty(\Theta(t)-1)
 \left(t^{s/2-1}+t^{(1-s)/2-1}\right)\,dt .
\end{equation}
The integral is entire because of Gaussian decay, locally uniform with
every derivative. Its right side is invariant under \(s\mapsto1-s\).
This proves the functional equation and continuation to the whole plane.
The standard Euler product
\[
 \Gamma(z)^{-1}=z e^{\gamma z}
       \prod_{n=1}^{\infty}(1+z/n)e^{-z/n}
\]
shows that \(\Gamma\) has no zeros and has poles only at the nonpositive
integers. Its product converges locally uniformly by
\(\log(1+z/n)-z/n=O(n^{-2})\); the identity follows from Euler's
limit formula for the gamma integral. These elementary gamma facts and
\eqref{eq:sr-reflection} imply that a zero away from the negative even
integers and \(0\) reflects to a zero at \(1-s\).
At \(0\), comparison of the residues in \eqref{eq:sr-reflection} gives
\(\zeta(0)=-1/2\). Thus proving zero-freeness in \(\Re s>1/2\)
suffices for RH.

\subsection{The smoothed Mellin identity}
For \(\Re s>1\) and every positive integer \(n\), direct integration gives
\[
 \int_1^\infty(x/n-1)_+x^{-s-1}\,dx
 =\frac{n^{-s}}{s(s-1)}.
\]
Indeed the integral starts at \(n\), and integrating
\(n^{-1}x^{-s}-x^{-s-1}\) gives the displayed answer.
Absolute convergence justifies summing this identity after multiplication
by \(\Lambda(n)\): the sum of integrals of absolute values is
\(\sum_n\Lambda(n)n^{-\sigma}/(\sigma(\sigma-1))<\infty\),
where \(\sigma=\Re s>1\). Also
\(\int_1^\infty x\log x\,x^{-s-1}\,dx=(s-1)^{-2}\) and
\(\int_1^\infty x\,x^{-s-1}\,dx=(s-1)^{-1}\). Consequently
\begin{equation}\label{eq:sr-mellin}
 \Phi(s):=\int_1^\infty E(x)x^{-s-1}\,dx
 =-\frac{\zeta'(s)}{s(s-1)\zeta(s)}
       -\frac1{(s-1)^2}+\frac{1+\gamma}{s-1}.
\end{equation}
Initially this holds in \(\Re s>1\).
The expansion \eqref{eq:sr-laurent} cancels both apparent polar terms at
\(1\). Therefore \(\Phi\) is holomorphic on \(\Re s>0\) except at the
zeros of \(\zeta\). If \(\rho\) is a zero of multiplicity \(m\), factor
\(\zeta(s)=(s-\rho)^m h(s)\), \(h(\rho)\ne0\), locally. Then
\[
 \Phi(s)=\frac{c_\rho}{s-\rho}+H_\rho(s),\qquad
 c_\rho=-\frac{m}{\rho(\rho-1)}\ne0,
\]
where \(H_\rho\) is holomorphic near \(\rho\).
Every such pole is nonreal.

Set \(g(u)=E(e^u)\), \(u\geq0\). This is a continuous real function.
The trivial estimates \(\Lambda(n)\leq\log n\) and
\(\sum_{n\leq x}1/n\leq1+\log x\) give
\(T(x)\leq x\log x(1+\log x)\).
Thus \(|g(u)|\leq C_0 e^{2u}\) for a fixed \(C_0\):
each of \(1,u,u^2\), multiplied by \(e^{-u}\), is bounded on
\([0,\infty)\). The change of variables \(x=e^u\) gives
\[
 \Phi(s)=\int_0^\infty g(u)e^{-su}\,du\qquad(\Re s>2).
\]
For every \(k\geq0\), differentiation under the integral yields
\begin{equation}\label{eq:sr-deriv}
 A_k(s):=\frac{(-1)^k\Phi^{(k)}(s)}{k!}
 =\int_0^\infty g(u)\frac{u^k e^{-su}}{k!}\,du .
\end{equation}
On any closed sub-half-plane \(\Re s\geq2+a\), the differentiated
integrands are dominated by \(C_0u^k e^{-au}\), an integrable function.
This supplies the justification for every derivative in \eqref{eq:sr-deriv}.

\section{The window oscillation argument in full}\label{sec:sr-landau}
\begin{srtheorem}[One-sided windows exclude right-hand poles]\label{thm:sr-window}
Let \(0<\theta<1\) and \(\delta>0\). If arbitrarily large \(U\) satisfy
\[
 g(u)\leq e^{\theta u}\qquad(U\leq u\leq(1+\delta)U),
\]
then \(\zeta\) has no zero with real part greater than \(\theta\).
\end{srtheorem}

\subsection{A localized gamma kernel}
For \(a>0\), integration by parts \(k\) times gives
\[
 \int_0^\infty\frac{u^k e^{-au}}{k!}\,du=a^{-k-1},\qquad
 \frac{u^k e^{-au}}{k!}\leq a^{-k}.
\]
The second inequality follows by retaining the \(k\)-th term of the
series for \(e^{au}\). The boundary terms in the integration vanish
because exponential decay dominates each fixed power.

We will choose \(\sigma>2\), depending only on \(\delta\), and, for each
\(U\geq1\), an integer \(k\geq U\) such that
\begin{equation}\label{eq:sr-tail}
 \int_{(0,\infty)\setminus[U,(1+\delta)U]}
       \frac{u^k e^{(2-\sigma)u}}{k!}\,du
 \leq C_1\sigma^{-k}.
\end{equation}
Here is an explicit existence argument, including both tails.
Put \(\alpha=1+\delta/2\) and \(l=2+6/\delta\).
For sufficiently large \(\sigma>2+l\), both
\[
 -l(1+\delta)+\alpha\sigma
       \log\frac{\sigma}{\sigma-2-l}<-\tfrac12,\qquad
 l+\alpha\sigma\log\frac{\sigma}{\sigma-2+l}<-\tfrac12
\]
hold. The limits of their left sides as \(\sigma\to\infty\) are both
\(-1\), since \(\sigma\log(\sigma/(\sigma-a))\to a\).
Choose such a \(\sigma\), also with \(\alpha\sigma\geq1\), and set
\(k=\lceil\alpha\sigma U\rceil\).
For \(u\geq(1+\delta)U\), insert \(e^{l(u-(1+\delta)U)}\geq1\)
into the integrand. For \(0<u<U\), insert \(e^{l(U-u)}\geq1\).
Enlarging each resulting integral to \((0,\infty)\) gives the tail bound
\[
 \frac{e^{-l(1+\delta)U}}{(\sigma-2-l)^{k+1}}
 +\frac{e^{lU}}{(\sigma-2+l)^{k+1}}.
\]
Multiply by \(\sigma^k\).
In the first term, use \(k<\alpha\sigma U+1\), since its logarithmic
ratio is positive. In the second, use \(k\geq\alpha\sigma U\), since
its logarithmic ratio is negative. The sum is at most
\[
 e^{-U/2}\left(
 \frac{\sigma}{(\sigma-2-l)^2}+\frac1{\sigma-2+l}\right).
\]
This proves \eqref{eq:sr-tail} with the parenthesized expression as
\(C_1\). The same choice has \(k\geq U\).

\subsection{An inequality for the derivatives of the transform}
For a good interval \(I=[U,(1+\delta)U]\), let
\(w(u)=u^k e^{-\sigma u}/k!\).
For any real \(v\), \eqref{eq:sr-deriv} and the triangle inequality give
\[
 |A_k(\sigma+iv)|+A_k(\sigma)
 \leq\int_0^\infty(|g(u)|+g(u))w(u)\,du .
\]
The quantity \(A_k(\sigma)\) is real. On \(I\),
\(|g|+g=2\max(g,0)\leq2e^{\theta u}\); outside \(I\), it is at most
\(2C_0e^{2u}\). The kernel maximum and \eqref{eq:sr-tail} therefore give
\begin{equation}\label{eq:sr-key}
 |A_k(\sigma+iv)|+A_k(\sigma)
 \leq\frac{2\delta U}{(\sigma-\theta)^k}
       +\frac{2C_0C_1}{\sigma^k}.
\end{equation}
For every \(\eta>0\), all sufficiently large \(U\) satisfy
\(\delta U\leq e^{\eta U}\). Since \(\sigma-\theta<\sigma\),
\begin{equation}\label{eq:sr-key2}
 |A_k(\sigma+iv)|+A_k(\sigma)
 \leq C_2 e^{\eta U}(\sigma-\theta)^{-k},
 \qquad C_2=2+2C_0C_1.
\end{equation}
For every \(\eta>0\), this holds for arbitrarily large good \(U\),
with a corresponding \(k\geq U\), simultaneously for all \(v\).

\subsection{Isolating one pole}
Suppose for a contradiction that \(\rho_0\) is a zero with
\(\Re\rho_0>\theta\). Write \(Z\) for the zeros in the positive half-plane.
Consider the closed disk centered at the real point \(\sigma\), of radius
\(\sigma-\theta\). It lies in \(\Re s\geq\theta>0\) and contains finitely
many elements of \(Z\).

There exist a number \(D\leq\sigma-\theta\) and a zero \(\rho_1\) such that
the open disk \(B(\sigma,D)\) is pole-free and
\begin{equation}\label{eq:sr-horizontal}
 \sigma-\Re\rho_1<D.
\end{equation}
If the closed disk has no zeros, take \(D=\sigma-\theta,\rho_1=\rho_0\).
Otherwise choose a zero \(\rho_*\) in it closest to \(\sigma\) and take
\(D=|\sigma-\rho_*|,\rho_1=\rho_*\).
All positive-half-plane zeros are nonreal, so
\(\sigma-\Re\rho_*<|\sigma-\rho_*|\).
This proves \eqref{eq:sr-horizontal} in the second case as well.

We next find a center \(z=\sigma+iv\) and a radius \(r\leq D\) whose open
disk contains exactly one zero \(\rho\), with \(d=|z-\rho|<r\).
Start with \(z_0=\sigma+i\Im\rho_1\), so
\(|z_0-\rho_1|=\sigma-\Re\rho_1<D\).
Choose \(h>0\) so small that
\[
 |z_0-\rho_1|+2h<D,\qquad D+2h<\sigma .
\]
The set \(F=Z\cap\overline B(z_0,D+h)\) is finite. As \(v\) varies in
\((\Im\rho_1-h,\Im\rho_1+h)\), two distinct points \(a,b\) of \(F\)
can have equal distances from \(\sigma+iv\) for at most one value of \(v\).
Indeed equality of their squared distances is affine in \(v\).
If their imaginary parts coincide, equality would force equal real
parts too, since both real parts are less than \(1<\sigma\).
Exclude the finitely many possible tie values and choose \(v\) in the
remaining interval. There is a unique nearest point \(\rho\) of \(F\);
its distance is at most \(|z-\rho_1|<D\).
Every zero outside \(F\) has distance greater than \(D\) from \(z\),
by the triangle inequality.
Choose \(r\) strictly between \(d\) and the smaller of \(D\) and the
distances to all other points of \(F\). This proves the claim.
Both disks under discussion are contained in \(\Re s>0\).

\subsection{Cauchy estimates and the contradiction}
Choose \(b_0\) with \(d<b_0<r\).
Since \(\Phi\) is holomorphic in \(B(\sigma,D)\), Cauchy's formula
on the circle of radius \(b_0<D\) gives
\[
 |A_k(\sigma)|\leq M_0 b_0^{-k}
\]
with \(M_0\) independent of \(k\).
In \(B(z,r)\), subtract the single principal part
\(c_\rho/(s-\rho)\). The resulting function is holomorphic throughout
the disk, by the local factorization at \(\rho\). Cauchy's formula gives
\[
 \left|A_k(z)-\frac{c_\rho}{(z-\rho)^{k+1}}\right|
 \leq M_1 b_0^{-k}.
\]
Combining these estimates with \eqref{eq:sr-key2} yields, for arbitrarily
large good \(U\) and the corresponding \(k\geq U\),
\[
 \frac{|c_\rho|}{d^{k+1}}
 \leq(M_0+M_1)b_0^{-k}
       +C_2e^{\eta U}(\sigma-\theta)^{-k}.
\]
Here \(d>0\), since all zeros have real part less than \(1\), whereas
\(\Re z=\sigma>2\). Also \(d<b_0<r\leq D\leq\sigma-\theta\).
Multiply by \(d^k\), put \(\lambda=\log(b_0/d)>0\), and choose
\(0<\eta<\lambda\).
The right side is at most
\[
 (M_0+M_1)e^{-\lambda U}
   +C_2e^{-(\lambda-\eta)U},
\]
which tends to zero. The left side is the fixed positive number
\(|c_\rho|/d\). This contradiction proves
Theorem~\ref{thm:sr-window}.

\subsection{Completion of the conditional analytic deduction}
Suppose \(\zeta(\rho)=0\) with \(\Re\rho>1/2\).
The previous analytic facts give \(\Re\rho<1\).
Choose \(\epsilon>0\) so small that
\(\theta=1/2+2\epsilon<\Re\rho\).
Equation \eqref{eq:sr-continuous}, after \(x=e^u\), supplies arbitrarily
late intervals \([\log X,(5/4)\log X]\) with
\(g(u)\leq e^{\theta u}\).
Theorem~\ref{thm:sr-window}, with \(\delta=1/4\), rules out \(\rho\).
Thus there are no zeros in \(\Re s>1/2\).
Reflection from \eqref{eq:sr-reflection} proves
Theorem~\ref{thm:sr-prime}.
\qed

\section{The prime distribution shared by the two analytic routes}
\label{sec:sr-shared-distribution}
The arithmetic-geometric approach controls the smoothed error \(E\).
A spectral approach seeks a trace identity for its prime-power
distribution. The following calculation identifies the common object,
including its normalization, without assuming RH or a trace formula.

\subsection{An exact identity on the logarithmic line}
Set \(H(u)=e^{-u/2}E(e^u)\) for \(u\geq0\), and let \(D_u=d/du\).
For \(u>0\), in distributions,
\begin{equation}\label{eq:sr-prime-distribution}
 (D_u^2-\tfrac14)H
  =\sum_{n\geq2}\frac{\Lambda(n)}{\sqrt n}\delta_{\log n}
          -e^{u/2}\,du .
\end{equation}
This sum is locally finite. In particular, for every
\(\varphi\in C_c^\infty(0,\infty)\),
\begin{equation}\label{eq:sr-prime-test}
 \int_0^\infty H(u)(\varphi''(u)-\tfrac14\varphi(u))\,du
 =\sum_{n\geq2}\frac{\Lambda(n)}{\sqrt n}\varphi(\log n)
       -\int_0^\infty e^{u/2}\varphi(u)\,du.
\end{equation}

Here is a direct proof. Put \(a=\log n\). The contribution of \(n\)
to \(e^{-u/2}T(e^u)\) is
\[
 h_n(u)=\frac{\Lambda(n)}{\sqrt n}
 \bigl(e^{(u-a)/2}-e^{-(u-a)/2}\bigr){\bf1}_{u\geq a}.
\]
It is continuous, is zero before \(a\), and solves
\(h_n''-\tfrac14 h_n=0\) away from \(a\).
Its first derivative jumps at \(a\) by \(\Lambda(n)/\sqrt n\).
Splitting the integral at \(a\) and integrating by parts twice
therefore gives
\[
 \int_0^\infty h_n(u)(\varphi''(u)-\tfrac14\varphi(u))\,du
       =\frac{\Lambda(n)}{\sqrt n}\varphi(a).
\]
Only finitely many \(h_n\) contribute on the support of \(\varphi\).
The remaining term in \(H\) is
\[
 (1+\gamma-u)e^{u/2},\qquad
 (D_u^2-\tfrac14)((1+\gamma-u)e^{u/2})=-e^{u/2}.
\]
This proves \eqref{eq:sr-prime-test} and
\eqref{eq:sr-prime-distribution}. There is no boundary term at zero,
since the test functions vanish in a neighborhood of zero.

The values at zero, useful for a one-sided Laplace transform, are
\[
 H(0)=1+\gamma,\qquad H'(0)=\frac{\gamma-1}{2}.
\]
For \(\Re z>3/2\) the elementary growth bound already proved for
\(E(e^u)\) justifies integration by parts at infinity.
Changing variables gives
\(\mathcal L H(z)=\Phi(z+\tfrac12)\), and
\eqref{eq:sr-prime-distribution} yields
\begin{align}
 (z^2-\tfrac14)\Phi(z+\tfrac12)
  &=-\frac{\zeta'(z+\tfrac12)}{\zeta(z+\tfrac12)}
       -\frac1{z-\tfrac12}\notag\\
  &\qquad +(1+\gamma)z+\frac{\gamma-1}{2}.
 \label{eq:sr-centered-transform}
\end{align}
The boundary terms have been displayed, so they cannot be lost by
formally transforming a differential operator on a half-line.
Equation~\eqref{eq:sr-mellin} gives the same identity algebraically.
The poles associated to nontrivial zeros are located at
\(z=\rho-\tfrac12\). Thus a bound on \(E\) and a valid global spectral
identity address the same zeros, by different additional arguments.
The local prime-power calculation itself proves neither of them.

\subsection{The growth estimate actually needed for spectral sums}
The following global zero count is enough for smooth compactly
supported tests; a universal logarithmic bound in each unit-height
box for arbitrary Mellin transforms is not needed.
\begin{srlemma}\label{lem:sr-global-zero-count}
The nontrivial zeros of \(\zeta\), counted with multiplicity, satisfy
\[
 \#\{\rho:|\rho|\leq R\}=O(R\log(R+2)).
\]
\end{srlemma}
\begin{proof}
From \eqref{eq:sr-reflection},
\[
 \xi(s)=\frac{s(s-1)}2\pi^{-s/2}\Gamma(s/2)\zeta(s)
  =\frac12+\frac{s(s-1)}4J(s),
\]
where
\[
 J(s)=\int_1^\infty(\Theta(t)-1)
           (t^{s/2-1}+t^{(1-s)/2-1})\,dt
\]
is entire and \(\xi(0)=1/2\).
For \(t\geq1\),
\[
 0\leq\Theta(t)-1
 \leq\frac{2e^{-\pi t}}{1-e^{-\pi}}.
\]
Indeed \(n^2\geq n\) bounds the positive-index sum by a geometric
series. For \(|s|\leq R\), \(R\geq2\), both powers in \(J\)
are bounded by \(t^m\), with \(m=\lceil(R+1)/2\rceil\).
Consequently
\[
 |J(s)|\leq C\int_1^\infty e^{-\pi t}t^m\,dt
       \leq C\,\frac{m!}{\pi^{m+1}}
       \leq C(R+2)^{R+2}.
\]
The integral over \([0,\infty)\) equals \(m!/\pi^{m+1}\), by
\(m\) integrations by parts.
Hence
\(\log\max_{|s|\leq R}|\xi(s)|\leq C'(R+2)\log(R+2)\).

For completeness the needed Jensen inequality follows by factoring
the finitely many zeros in a closed disk and applying the mean-value
property to the logarithm of the remaining nonvanishing holomorphic
factor. Its contribution from a zero \(\rho\) in the disk of radius
\(R_1\) is \(\log(R_1/|\rho|)\); this follows from
the mean of \(\log|R_1e^{i\theta}-\rho|\), which is \(\log R_1\).
Choose \(2R<R_1<2R+1\) with no zero on its boundary.
Each zero with \(|\rho|\leq R\) then contributes at least \(\log2\).
The sum of these contributions is bounded by
\(\log\max_{|s|\leq R_1}|\xi(s)|-\log|\xi(0)|\).
This proves the assertion for zeros of \(\xi\).
Every nontrivial zeta zero occurs in \(\xi\) with the same multiplicity,
since the other factors are holomorphic and nonzero there.
\end{proof}

\subsection{What an exact spectral identity on shrinking strips would prove}
\begin{srlemma}[Uniqueness for a vertical discrete spectrum]
\label{lem:sr-spectrum-unique}
Let \(\mathcal Z\) be a closed discrete set in
\(\{z:|\Re z|\leq B\}\).
Give its points nonzero multiplicities \(m_z\) such that
\(\sum_{z\in\mathcal Z}|m_z|(1+|\Im z|)^{-M}<\infty\)
for some \(M\geq0\).
For \(\varphi\in C_c^\infty(0,\infty)\) put
\(F_\varphi(z)=\int_0^\infty\varphi(u)e^{zu}\,du\).
If
\[
 \sum_{z\in\mathcal Z}m_zF_\varphi(z)=0
       \quad\hbox{for every such }\varphi,
\]
then \(\mathcal Z\) is empty.
\end{srlemma}
\begin{proof}
Integration by parts an arbitrary number of times in \(u\)
shows rapid decay of \(F_\varphi(z)\) as \(|\Im z|\to\infty\),
uniformly for \(|\Re z|\leq B\).
There are no boundary terms, because the support is compact in
\((0,\infty)\). Thus
\(a_z=m_zF_\varphi(z)\) is absolutely summable.

For \(T\geq0\), define
\(\varphi_T(u)=\varphi(u-T)\), with zero extension.
It is another permitted test, and
\(F_{\varphi_T}(z)=e^{zT}F_\varphi(z)\).
The hypothesis becomes \(\sum_z a_ze^{zT}=0\).
For \(\Re w>B\), multiply by \(e^{-wT}\) and integrate \(T\)
from zero to infinity. The sum of the absolute integrals is at
most \(\sum_z|a_z|/(\Re w-B)\), so interchange is justified.
We obtain
\[
 \sum_{z\in\mathcal Z}\frac{a_z}{w-z}=0\qquad(\Re w>B).
\]
This series converges normally on every compact set disjoint from
\(\mathcal Z\): only finitely many \(z\) are near that compact set,
and \(|w-z|\geq1\) uniformly for all remaining sufficiently large
\(z\). It therefore defines a meromorphic function, with residue
\(a_z\) at \(z\).
The plane minus a closed discrete set is connected. One can join
two points by a polygonal path and detour around the finitely many
points on its compact image.
The identity theorem makes this meromorphic function zero off
\(\mathcal Z\), so all its residues vanish.

If \(z_0=\sigma+i\tau\) belonged to \(\mathcal Z\), choose a nonzero
nonnegative smooth bump \(b\) supported in \((0,\infty)\) and set
\(\varphi(u)=b(u)e^{-i\tau u}\).
Then \(F_\varphi(z_0)=\int b(u)e^{\sigma u}\,du>0\).
This contradicts \(a_{z_0}=0\), since \(m_{z_0}\ne0\).
\end{proof}

Suppose, as an additional hypothesis, that for every \(\delta>0\)
one had an identity
\begin{equation}\label{eq:sr-strip-trace}
 \mathcal A(\varphi)=
  \sum_{\substack{\rho\\|\Re\rho-1/2|<\delta}}
        m_\rho F_\varphi(\rho-\tfrac12)
       \qquad(\varphi\in C_c^\infty(0,\infty)),
\end{equation}
with exactly the same functional \(\mathcal A\) at every width.
The zero count in Lemma~\ref{lem:sr-global-zero-count} makes these
sums absolutely convergent, by integration by parts as above.
Subtracting the identities at \(0<\delta_1<\delta_2\) gives the
zero distribution for
\[
 \{\,\rho-\tfrac12:
             \delta_1\leq|\Re\rho-\tfrac12|<\delta_2\,\}.
\]
This is closed discrete, because zeros have no finite accumulation
point. Lemma~\ref{lem:sr-spectrum-unique} makes it empty.
Any off-line zero would belong to such a set for some two widths,
so \eqref{eq:sr-strip-trace} would imply RH.

This implication explains the burden of a width-independent trace
claim. Local calculations of the prime and gamma terms do not prove
\eqref{eq:sr-strip-trace}. A construction must prove the spectral
identification, convergence and equality of the global traces,
including any terms introduced by changing the strip.
This paper does not assume \eqref{eq:sr-strip-trace} in its
arithmetic-geometric argument.

\subsection{Two checks on proposed replacements for that burden}
For \(\Re A>0\) put
\[
 g_A(x)=(4\pi A)^{-1/2}\exp(-x^2/(4A)),
 \qquad h(x)=\frac{g_{1+i}(x)-g_{1-i}(x)}{2i}.
\]
Use the square root continuous in the right half-plane and positive
on the positive real axis. Then \(h\) is real-valued.
For every real \(c\), every derivative of \(e^{cx}h(x)\) is a
polynomial times a Gaussian with strictly negative quadratic real
part. It is therefore Schwartz.
Completing the square for positive real \(A\) gives
\(\int_{\mathbb R}g_A(x)e^{zx}\,dx=e^{Az^2}\).
Both sides are holomorphic for \(\Re A>0\), locally uniformly by
Gaussian domination; continuation in \(A\) proves the same formula
there. Thus, for \(\varphi(t)=h(\log t)\),
\[
 \int_0^\infty\varphi(t)t^z\,\frac{dt}{t}
                  =e^{z^2}\sin(z^2).
\]
This entire transform has simple zeros at
\(z=i\sqrt{n\pi}\), \(n\geq1\).
The number in \([-1,1]+i[T,T+1]\), for large \(T\), equals
\[
 \left\lfloor\frac{(T+1)^2}{\pi}\right\rfloor
 -\left\lceil\frac{T^2}{\pi}\right\rceil+1
                  =\frac{2T}{\pi}+O(1).
\]
It cannot satisfy an \(O(\log T)\) unit-box bound.
Rapid vertical decay alone consequently cannot justify that bound
for arbitrary exponential-Schwartz Mellin transforms.
In Jensen's inequality an upper bound on a rapidly decaying value
does not supply the lower bound on that value needed to bound its
negative logarithm.

There is also an elementary projection check.
Let \(\mathcal H\) be a Hilbert space and \(\lambda_0\) a nonzero
continuous linear functional.
If the range of \(P_T\) lies in \(\ker\lambda_0\) for every
\(T\geq T_0\), then \(P_T\) cannot converge strongly to the identity:
otherwise, for every \(x\),
\(\lambda_0(x)=\lim_T\lambda_0(P_Tx)=0\).
More generally intersections of an increasing collection of
evaluation kernels decrease as the collection grows.
Passing to complementary projectors changes the limiting subspace;
it does not prove density in the full space.

These two checks apply to the general zero-count and kernel-projector
claims in Masullo's current draft, Lemma~9.3.2.4 and
Theorem~9.3.2.6 (printed pages 425--426).
The affected statements need replacement before their downstream
spectral argument can supply an input here.
This identifies particular obstructions; it is not a verdict on every
other result in that draft.
The source PDF and the separate, longer audit are listed in
Section~\ref{sec:sr-source-review}.

\section{The signed finite-group comparison}\label{sec:sr-finite}
This section gives the finite comparison underlying Anton's shorter route.
Every estimate in this section is proved; the existence of data satisfying
them on late intervals is stated separately in
Section~\ref{sec:sr-construction}.

Let \(H\) be a finite abelian group, \(n=|H|\), and put
\[
 \Gamma=(\mathbb Z/2)^2\times H,\quad
 \Phi=\{(0,b,h):b\in\mathbb Z/2,\ h\in H\},\quad f=|\Phi|=2n.
\]
Use additive notation for \(\Gamma\). Define
\(c=(1,1,0)\), \(t^\circ=(0,0,t)\),
\(\chi(a,b,h)=(-1)^a\), and \(\omega(a,b,h)=(-1)^b\).
The dual group \(\widehat\Gamma\) consists of its complex unitary characters;
its identity character is denoted \(1\). It has \(4n=2f\) elements.

For each rational prime \(p\), choose a subgroup \(I_p\subset\Gamma\) and
a Frobenius representative \(F_p\in\Gamma\).
Fix distinct odd primes \(q,d\). The host conditions are:
\[
 I_p=0\ (p\notin\{2,q,d\}),\qquad
 I_p\subset\Phi,\ F_p\in\Phi\ (p=2,q),\qquad
 I_d=\{0,(1,0,0)\}.
\]
For \(p^r\), \(r\geq1\), set
\[
 c_\psi(p^r)=\frac1{|I_p|}\sum_{h\in I_p}\psi(rF_p+h).
\]
Extend these coefficients multiplicatively to all positive integers.
Their absolute values are at most one, and \(c_1(m)=1\).
Choose arbitrary constants \(\ell_\psi\in\mathbb C\) for \(\psi\ne1\).
For \(y>0\), define
\[
 G_1(y)=E(y),\qquad
 G_\psi(y)=\sum_{r\leq y}\Lambda(r)c_\psi(r)(y/r-1)+y\ell_\psi
 \quad(\psi\ne1),
\]
using the same formula defining \(E\) also when \(0<y<1\).
Put
\[
 {\cal C}(y)=\Re(G_\chi(y)+G_{\chi\omega}(y)-G_\omega(y)),
 \qquad {\cal S}(y)=G_1(y)+G_\omega(y)-G_\chi(y)-G_{\chi\omega}(y).
\]
Thus \(\Re{\cal S}(y)=E(y)-{\cal C}(y)\).

\subsection{Orthogonality, including the ramified coefficients}
For every \(g\in\Gamma\), orthogonality on the finite abelian group gives
\begin{equation}\label{eq:sr-coset}
 \frac1f\sum_{\psi\in\widehat\Gamma}\psi(-g)c_\psi(p^r)
 =\begin{cases}2/|I_p|&g-rF_p\in I_p,\\0&\text{otherwise}.\end{cases}
\end{equation}
To verify this without any representation theory, interchange the
two finite sums. For each \(h\in I_p\), the inner sum is
\(\sum_\psi\psi(rF_p+h-g)\), equal to \(|\Gamma|\) if its argument is zero
and zero otherwise. The latter assertion follows by multiplying the sum
by a character nontrivial on the argument; such a character exists by
the decomposition of a finite abelian group into cyclic groups.
There is at most one \(h\) giving zero. Division by \(f|I_p|\) proves
\eqref{eq:sr-coset}.

In particular the raw four-character prime-power coefficient is
\[
 1+c_\omega(p^r)-c_\chi(p^r)-c_{\chi\omega}(p^r)
 =\frac4{|I_p|}
   \#\{h\in I_p:\chi(rF_p+h)=-1,\ \omega(rF_p+h)=1\}.
\]
This follows pointwise from
\(1+\omega-\chi-\chi\omega=(1-\chi)(1+\omega)\).
It is real and lies in \([0,4]\), even at ramified primes.
This positivity statement concerns the coefficients; it does not by itself
bound the complete signal, which also includes the normalizing terms in
the functions \(G_\psi\).

\subsection{The exact colour balance}
For \(t\in H\) write \(\rho_t(g)=c+g+t^\circ\).
It is a bijection from \(\Gamma\setminus\Phi\) to \(\Phi\).
Let \(a>0\), let \(w\) be positive on \(\Phi\), and choose an integer
\(N\geq a/w(\rho)\) for every \(\rho\in\Phi\). Define
\[
 \begin{split}
 {\cal A}(a,t,w)={}&
 2\sum_{\substack{p\leq N\ {\rm prime}\\p\notin\{2,q,d\}\\F_p\notin\Phi}}
  \log p
  \sum_{\substack{1\leq r\leq N\\rF_p\notin\Phi}}
       (a/p^r-w(\rho_t(rF_p)))_+\\
 &+\log d\sum_{r=1}^N
   \sum_{\substack{g-rF_d\in I_d\\g\notin\Phi}}
       (a/d^r-w(\rho_t(g)))_+ ,
 \end{split}
\]
and
\[
 M(a,w)=a\left(\log a-1-\gamma+\Re\ell_\chi
                    -\frac1f\sum_{\rho\in\Phi}\log w(\rho)\right).
\]
Finally put
\[
 B(a,t,w)=\Re\left[
 \frac1f\sum_{g\notin\Phi}\sum_{\psi\in\widehat\Gamma}
    \psi(-g)w(\rho_t(g))G_\psi(a/w(\rho_t(g)))\right].
\]
Then
\begin{equation}\label{eq:sr-balance}
 {\cal A}(a,t,w)-M(a,w)=B(a,t,w).
\end{equation}
Here is a direct verification. First expand the ramp sums in \(B\) and
use \eqref{eq:sr-coset}. Terms beyond \(N\) vanish by the choice of \(N\).
The terms at \(2,q\) vanish, because their Frobenius cosets lie in \(\Phi\).
For every other unramified prime the coefficient is two at the single
point \(rF_p\). At \(d\), its two-point coset has exactly one point outside
\(\Phi\), with coefficient one. These are precisely the two sums in
\({\cal A}\). No prime with \(p>N\), or exponent \(r>N\), contributes,
because \(p^r>N\) then and every relevant \(a/w\leq N\).

For the remaining terms, sum each character on the complement of
the index-two subgroup \(\Phi\):
\[
 \frac1f\sum_{g\notin\Phi}\psi(-g)
 =\begin{cases}1&\psi=1,\\-1&\psi=\chi,\\0&\text{otherwise}.\end{cases}
\]
This follows from the same orthogonality, first on \(\Phi\) and then
on its coset. The nontrivial-character linear terms therefore contribute
\(-a\ell_\chi\).
The normalizing terms in \(G_1\) contribute
\(-a\log a+(1+\gamma)a+(a/f)\sum_{\rho\in\Phi}\log w(\rho)\).
Their real parts give \(-M\), proving \eqref{eq:sr-balance}.

\subsection{Two uniform Lipschitz estimates}
For all \(A,w,w'>0\) and every character,
\begin{equation}\label{eq:sr-weightlip}
 |wG_\psi(A/w)-w'G_\psi(A/w')|
 \leq\frac{3A}{\min(w,w')}|w-w'|.
\end{equation}
For \(\psi\ne1\), write
\[
 wG_\psi(A/w)=\sum_r\Lambda(r)c_\psi(r)(A/r-w)_++A\ell_\psi .
\]
The linear term cancels. Each positive-part function is Lipschitz with
constant one, and only \(r\leq A/\min(w,w')\) can contribute.
Equation \eqref{eq:sr-cheb} proves \eqref{eq:sr-weightlip}.
For \(\psi=1\), the expression is the real continuous piecewise smooth
function
\[
 \sum_r\Lambda(r)(A/r-w)_+-A\log A+A\log w+(1+\gamma)A.
\]
Off its finitely many breakpoints in any compact interval of \(w>0\),
its derivative is \(A/w-\psi(A/w)\).
By \eqref{eq:sr-cheb}, its absolute value is at most \(2A/w\).
Integrating across the finitely many pieces proves the stronger bound
with constant two, and hence \eqref{eq:sr-weightlip}.

Also,
\begin{equation}\label{eq:sr-Elip}
 |E(x)-E(y)|\leq4|x-y|\qquad(x,y\geq1).
\end{equation}
For an integer \(m\geq1\), set \(L(m)=\sum_{r\leq m}\Lambda(r)/r\).
The factorial identity gives
\[
 \log(m!)=mL(m)-\sum_{r\leq m}\Lambda(r)\{m/r\},
\]
where the second sum lies between zero and \(\psi(m)\leq3m\).
The factorial integral bounds above imply
\(-1\leq L(m)-\log m\leq3\), using
\(1+\log m\leq m\).
If \(m\leq x<m+1\), then \(L(x)=L(m)\) and
\(0\leq\log x-\log m\leq\log2\). Therefore
\[
 -1-\log2\leq L(x)-\log x\leq3.
\]
The elementary harmonic integral inequalities give \(0\leq\gamma\leq1\).
Away from its integer breakpoints,
\(E'(x)=L(x)-\log x+\gamma\), whose absolute value is at most four.
Continuity and integration over the finitely many pieces between \(x,y\)
prove \eqref{eq:sr-Elip}.

\subsection{Transporting one signed bound}
Suppose \(x\geq3\), \(|y-x|\leq1\).
Apply \eqref{eq:sr-weightlip} with \(A=xy,w=x,w'=y\).
Taking real parts gives, for every \(\psi\),
\[
 |x\Re G_\psi(y)-y\Re G_\psi(x)|
 \leq3\max(x,y)|x-y|\leq4x.
\]
Using the upper inequalities for \(\chi,\chi\omega\) and the lower
one for \(\omega\), we obtain
\(x{\cal C}(y)\leq y{\cal C}(x)+12x\).
If \({\cal C}(x)\leq Bx^\theta\), \(B\geq0\), \(\theta\leq1\), then
\[
 {\cal C}(y)\leq Bx^\theta+B+12,
\]
because \(0\leq y\leq x+1\) and \(x^{\theta-1}\leq1\).
Together with \eqref{eq:sr-Elip} this proves
\begin{equation}\label{eq:sr-signaltransport}
 \Re{\cal S}(y)\geq E(x)-Bx^\theta-(B+16).
\end{equation}
The additive loss here uses the elementary constant four proved above.

\subsection{Averaging the thick part and bounding the repair}
Choose subsets \(T,T_2\subset H\), an integer \(m\geq1\),
numbers \(a_t>0\), a scale \(b>0\), and weights
\(w_{u,j}(\rho)>0\), \(u\in H,1\leq j\leq m,\rho\in\Phi\).
Call a point \(\rho\) thick for \(u\) if
\(\rho=(0,1,h+u)\) for some \(h\in T_2\).
Replace its weight by \(b\), keeping all other weights unchanged,
to obtain \(\widetilde w_{u,j}\).
Write \(\delta_2=|T_2|/n\). Define the thin mean
\[
 {\cal R}=\frac1n\sum_{u\in H}\sum_{t\in T}\sum_{j=1}^m
 \Re\left[\frac1f
 \sum_{\substack{g\notin\Phi\\\rho_t(g)\ {\rm not\ thick\ for}\ u}}
 \sum_{\psi\in\widehat\Gamma}
 \psi(-g)w_{u,j}(\rho_t(g))
 G_\psi(a_t/w_{u,j}(\rho_t(g)))\right].
\]
Then
\begin{equation}\label{eq:sr-thick}
 \frac1n\sum_{u,t,j}B(a_t,t,\widetilde w_{u,j})
 =\frac{\delta_2b}{2}\sum_{t\in T}\sum_{j=1}^m
       \Re{\cal S}(a_t/b)+{\cal R}.
\end{equation}
Indeed thickness requires that \(g=(1,0,h)\).
Averaging the translates \(T_2+u\) gives multiplicity \(|T_2|\)
at each \(h\in H\). Summing over \(h\) kills every character nontrivial
on \(H\). The four survivors \(1,\omega,\chi,\chi\omega\) have values
\(1,1,-1,-1\), respectively, at \((1,0,0)\).
Division by \(nf=2n^2\) leaves \(\delta_2/2\).
The complementary part is unchanged by the replacement, proving the identity.

Suppose moreover that \(|a_t/b-x|\leq1\), and the original thick weights
satisfy \(b\leq w_{u,j}(\rho)\leq b+\eta+1\), where \(\eta\geq0\).
For each fixed \(u,t,j\), \eqref{eq:sr-weightlip} bounds the change in
each weighted \(G_\psi\) by
\(3(x+1)(\eta+1)\).
There are \(f\) possible \(g\), \(2f\) characters, and a factor \(1/f\).
The triangle inequality therefore gives
\[
 |B(a_t,t,\widetilde w_{u,j})-B(a_t,t,w_{u,j})|
 \leq6f(x+1)(\eta+1).
\]
Consequently, if \(Q=|T|mb\) and
\begin{equation}\label{eq:sr-repair}
 6f(x+1)(\eta+1)\leq b/x,
\end{equation}
the total change in the averaged balance is at most \(Q/x\).

The thin mean can also be recovered from the original stronger
pointwise input. If at every thin node and for every \(\psi\),
\[
 |w_{u,j}(\rho)G_\psi(a_t/w_{u,j}(\rho))|\leq L,\qquad L\geq0,
\]
then the same count of terms gives
\begin{equation}\label{eq:sr-thin}
 {\cal R}\geq-2f|T|mL.
\end{equation}
If, more specifically, \(a_t\leq(x+1)b\),
\(a_t/w_{u,j}(\rho)\geq X'/3\) at thin nodes and
\(|G_\psi(z)|\leq z\vartheta\) for \(z\geq X'/3\), take
\(L=(x+1)b\vartheta\).
The inequality \(2f(x+1)\vartheta\leq1/x\) then yields
\({\cal R}\geq-Q/x\). No averaging cancellation is assumed in this argument.

\section{The exact geometric construction requirement}
\label{sec:sr-construction}
We now specify what must be supplied by the arithmetic geometry.
The distinction between a local implication and the existence statement
in this section is essential.

\begin{srdefinition}[A signed datum at \(x\)]
Fix \(C\geq0\) and \(0<\epsilon\leq1/2\). A signed datum consists of the
host and node data above, a cutoff \(N\), numbers \(\eta\geq0,k>0,A_E\geq0\)
and \({\rm Err}\), satisfying the following conditions.
Write \(t_0=|T|,t_2=|T_2|,Q=t_0mb\).
The host conditions hold, \(x\geq3\),
\(|a_t/b-x|\leq1\), and every \(a_t/w_{u,j}(\rho)\leq N\).
Thick weights lie in \([b,b+\eta+1]\), all weights are positive, and
\eqref{eq:sr-repair} holds. The three substantive inequalities are
\begin{align}
 \frac1n\sum_{u,t,j}\bigl({\cal A}(a_t,t,w_{u,j})-M(a_t,w_{u,j})\bigr)
       &\leq{\rm Err},\label{eq:sr-top}\\
 {\cal C}(x)&\leq3x^{1/2+\epsilon},\label{eq:sr-correction}\\
 {\cal R}&\geq-Q/x.\label{eq:sr-rest}
\end{align}
Finally the numerical shape requirements are
\begin{gather}
 A_E\leq2xb,\quad
 t_0\geq n x^{-\epsilon/8}/2,\quad
 x^{-\epsilon}\leq t_2/n\leq1,\label{eq:sr-shape1}\\
 x^{1/2-\epsilon}\leq m\leq x^{1/2+\epsilon/2},\quad
 xm\leq k\leq n^2x^{1/2-\epsilon},\quad
 26\leq x^{3\epsilon/8},\label{eq:sr-shape2}\\
 {\rm Err}\leq x^{C\epsilon}\left[
 t_0A_E(1+m^2/k)+m\left(2nmbx^\epsilon+
                 \frac{t_2kb x^\epsilon}{n^2x^{-\epsilon}}\right)\right].
 \label{eq:sr-error}
\end{gather}
\end{srdefinition}

\begin{srproposition}[Pointwise consequence]\label{prop:sr-point}
A signed datum at \(x\) gives \(E(x)\leq x^{1/2+(C+3)\epsilon}\).
\end{srproposition}
\begin{proof}
The condition \(26\leq x^{3\epsilon/8}\) implies
\(19\leq x^{1/2+\epsilon}\).
Apply \eqref{eq:sr-signaltransport} with \(B=3,\theta=1/2+\epsilon\)
and \(y=a_t/b\). It gives
\(\Re{\cal S}(a_t/b)\geq E(x)-4x^{1/2+\epsilon}\).
Equation \eqref{eq:sr-thick}, the thin lower bound, the repair bound,
the exact balance, and \eqref{eq:sr-top} therefore give
\begin{equation}\label{eq:sr-sizecomparison}
 Q\frac{\delta_2}{2}(E(x)-4x^{1/2+\epsilon})
       \leq{\rm Err}+2Q/x.
\end{equation}
All quantities divided by below are positive: \(n,b,m,t_0,k,\delta_2\)
are positive by the shape requirements.

Put \(Y=x^{1/2+13\epsilon/8}\) and \(P_0=x^{C\epsilon}\).
Divide the bracket in \eqref{eq:sr-error} by \(Q\).
Its four terms satisfy, respectively,
\[
 \frac{A_E}{mb}\leq2x^{1/2+\epsilon},\qquad
 \frac{A_Em}{bk}\leq2,\qquad
 \frac{2nmx^\epsilon}{t_0}\leq4x^{1/2+13\epsilon/8},\qquad
 \frac{t_2kx^{2\epsilon}}{n^2t_0}\leq2x^{1/2+9\epsilon/8}.
\]
These follow in order from \(A_E\leq2xb\) and the lower bound on \(m\);
from \(k\geq xm\); from the lower bound on \(t_0\) and upper bound on \(m\);
and from the upper bounds on \(t_2,k\) and lower bound on \(t_0\).
Each right side is at most its numerical coefficient times \(Y\), so
\({\rm Err}\leq10P_0YQ\).
Dividing \eqref{eq:sr-sizecomparison} by \(Q\delta_2/2\) gives
\[
 E(x)\leq4x^{1/2+\epsilon}
       +20P_0Yx^\epsilon+4x^{\epsilon-1}
 \leq4x^{1/2+\epsilon}+22P_0Yx^\epsilon .
\]
For the last step use \(x\geq3\) and \(P_0Y\geq1\).
Since \(C\geq0\),
\(x^{1/2+\epsilon}\leq P_0Yx^\epsilon
=x^{1/2+(C+21/8)\epsilon}\).
The result is at most \(26x^{1/2+(C+21/8)\epsilon}\).
Absorb 26 by \(x^{3\epsilon/8}\), obtaining the asserted bound.
\end{proof}

\begin{srdefinition}[Uniform supply]\label{def:sr-supply}
The supply assertion is the following precise statement:
there exist \(C\geq0,\epsilon_0>0\), with \(\epsilon_0\leq1/2\), such that
for every \(0<\epsilon<\epsilon_0\) and every \(Y_0\), there is
\(Y\geq\max(Y_0,1)\) for which a signed datum with the same \(C,\epsilon\)
exists at every fourth power \(j^4\in[Y,Y^{3/2}]\).
The host, group, field, weights and auxiliary parameters may vary with
the sample. The constant \(C\) and the order of quantifiers may not.
\end{srdefinition}

\begin{srtheorem}[Signed supply implies RH]\label{thm:sr-supply}
The assertion of Definition~\ref{def:sr-supply} implies ordinary RH.
\end{srtheorem}
\begin{proof}
Given any \(\alpha>0\), choose \(0<\epsilon<\epsilon_0\) with
\((C+3)\epsilon\leq\alpha\).
Proposition~\ref{prop:sr-point} bounds \(E(j^4)\) by
\((j^4)^{1/2+\alpha}\) at all the supplied samples.
By \eqref{eq:sr-integer} and \eqref{eq:sr-discard},
\[
 P(j^4)\geq\gamma j^4-(1+K)(j^4)^{1/2+\alpha}.
\]
The constant \(1+K\) is uniform across the chosen interval.
Theorem~\ref{thm:sr-prime} now applies.
\end{proof}

\paragraph{Status of the construction.}
The preceding theorem proves an implication. This edition does not prove
Definition~\ref{def:sr-supply}. In particular it does not construct
the geometric witnesses giving \eqref{eq:sr-top} with the error
\eqref{eq:sr-error} on arbitrarily late intervals. It also does not give
an independent proof that the signed correction bound
\eqref{eq:sr-correction} and the complementary bound
\eqref{eq:sr-rest} can be achieved simultaneously with those witnesses.
The archive-to-signed-data map transfers these requirements from stronger
archive inputs; it does not remove the requirement to prove their existence.
Consequently Theorem~\ref{thm:sr-supply} cannot be described as an
unconditional proof supplied by this paper.

\section{The conductor-one arithmetic packet}\label{sec:sr-arithmetic}
The arithmetic specialization in this section follows the contributors'
\texttt{PrincipalHostField}, \texttt{PrincipalCohomology},
\texttt{LengthBudget}, \texttt{PrincipalMinus},
\texttt{PrincipalPlus}, and \texttt{PowerWindow} files.
We prove the elementary field, class-number, descent, and numerical
comparisons here. The remaining uniform height and rank inputs are
distinguished from these consequences at the end of the section.
Classical background used here consists of finite Galois theory,
the Dirichlet unit theorem, Hilbert theorem 90, the Minkowski ideal-class
bound, and the unramified Artin reciprocity theorem.
None is supplied by a reference to a Lean module.

\subsection{An explicit field and its degree}
Let \(q\equiv3\pmod8,d\equiv7\pmod{16}\) be primes and \(e\geq0\).
They are distinct odd primes. Set \(r=e+2\), \(N=2^rqd\), and
\(L=\mathbb Q(\zeta_N)\). Let \(A\) be the subgroup of
\((\mathbb Z/N)^\times=\operatorname{Gal}(L/\mathbb Q)\) defined by
\[
 a\bmod2^r\in\{1,-1\},\qquad
 \left(\frac aq\right)=\left(\frac ad\right)=1,
\]
and put \(E=L^A\).
The Chinese remainder theorem and surjectivity of each quadratic character
give
\[
 \operatorname{Gal}(E/\mathbb Q)\simeq
 ((\mathbb Z/2^r)^\times/\{\pm1\})\times(\mathbb Z/2)^2,\qquad
 [E:\mathbb Q]=2^r=2^{e+2}.
\]
Indeed the first factor has order \(2^{r-2}\), including \(r=2\),
and each of the other factors has order two.
In particular no odd prime divides the degree.

Complex conjugation is the class of \(-1\).
It is trivial on the first factor and nontrivial on both quadratic
factors, since \(q,d\equiv3\pmod4\).
It is therefore a nonidentity central involution.
Its fixed field \(E^+\) lies in \(\mathbb R\), and
\([E:E^+]=2\). Every embedding of the Galois field \(E\) has the same
nontrivial conjugation action. Thus \(E\) is a CM field.
The least positive conductor \(f_E\) for a cyclotomic embedding
satisfies \(f_E\leq N\).

This construction is independent of the ambient cyclotomic field used
to describe it. In a larger cyclotomic field of conductor divisible by
\(N\), the subgroup specified by these same three tests is the inverse
image of \(A\). Restriction of automorphisms is surjective, and Galois
correspondence identifies its fixed field with \(E\).
Changing the chosen primitive root changes an exponent description by
multiplication by a unit and leaves the automorphism's exponent itself
unchanged. Therefore the fixed subgroup, and hence its fixed field,
are unchanged. This proves the canonical-field reduction without
retaining any choice of a realization.

\subsection{Strict lengths and a single integer}
For a prime \(\ell\), write \(S_\ell\) for the primes of \(E\) above \(\ell\).
For \(a\geq1\), define the strict group
\[
 T_{\ell,a}(E)=\ker\left[
 H^1(G_{E,S_\ell},(\mathbb Q_\ell/\mathbb Z_\ell)(-a))
 \longrightarrow
 \prod_{v\mid\ell}H^1(E_v,(\mathbb Q_\ell/\mathbb Z_\ell)(-a))\right].
\]
Cohomology is continuous, and the local maps are restrictions after
inflation from the Galois group unramified outside \(S_\ell\).
Let \(s_{\ell,E}(a)\) be its \(\mathbb Z_\ell\)-length, with value infinity
allowed until finiteness has been proved.
The group and length are independent of choices of closures and primes
above local places: field isomorphisms induce isomorphisms of global
and local groups, commute with restriction, and preserve these kernels.

Here is a useful exact reformulation, following \texttt{LengthBudget}.
For nonnegative integers \(L_\ell\) and a real \(B\),
\begin{equation}\label{eq:sr-integerbudget}
 \left(\forall S\ {\rm finite},\ \sum_{\ell\in S}L_\ell\log\ell\leq B\right)
 \quad\Longleftrightarrow\quad
 \left(\exists M\in\mathbb Z_{\geq1},\
 \log M\leq B,\ \ell^{L_\ell}\mid M\ {\rm for\ every}\ \ell\right).
\end{equation}
For the forward implication, the empty set first gives \(B\geq0\).
If \(L_\ell>0\), the singleton inequality gives \(\ell\leq e^B\).
There are only finitely many such primes. Take \(M\) to be the product
of \(\ell^{L_\ell}\) over this finite support; its logarithm is at most
\(B\) by the hypothesis. Conversely, distinct prime powers are coprime,
so each finite product divides \(M\), and its logarithm is at most
\(\log M\leq B\).
Apply this with \(L_\ell=\sum_{1\leq a<t}s_{\ell,E}(a)\), after proving
every length in that sum finite.
Finiteness cannot be dropped: replacing infinity by zero is not a
mathematical interpretation of a length.

\subsection{The one-ramified-prime class-number lemma}
\begin{srlemma}\label{lem:sr-oneprime}
Let \(L/K\) be a cyclic extension of degree \(p\), with \(p\) prime,
and both fields totally imaginary. Suppose exactly one finite prime
of \(K\) ramifies, and it is totally ramified. If
\(p\nmid h_K\), then \(p\nmid h_L\).
\end{srlemma}
\begin{proof}
Write \(G=\operatorname{Gal}(L/K)\), \(U_L\) for the units,
\(I_L\) for fractional ideals, and \(P_L\) for principal fractional ideals.
The group \(I_L\) is a direct sum of permutation modules on the primes.
Its first cohomology is zero: for each orbit this reduces to
\(H^1(J,\mathbb Z)=\operatorname{Hom}(J,\mathbb Z)=0\), with \(J\) finite.
Taking invariants in \(0\to P_L\to I_L\to\operatorname{Cl}_L\to0\)
therefore gives
\[
 |\operatorname{Cl}_L^G|
   =|I_L^G/P_L^G|\ |H^1(G,P_L)|.
\]
Extension of ideals identifies \(I_K\) with a subgroup of \(I_L^G\).
Invariants on each prime orbit show
\([I_L^G:I_K]=\prod_{v\ {\rm finite}}e_v=p\).
The sequence \(1\to U_L\to L^\times\to P_L\to1\), together with
Hilbert theorem 90, gives
\(P_L^G/P_K\simeq H^1(G,U_L)\). Thus
\[
 |I_L^G/P_L^G|=h_Kp/|H^1(G,U_L)|.
\]
In the same sequence,
\[
 H^1(G,P_L)\simeq
 \ker\bigl(H^2(G,U_L)\to H^2(G,L^\times)\bigr)
 = (U_K\cap N_{L/K}L^\times)/N_{L/K}U_L.
\]
For a cyclic group, \(H^2(G,U_L)=U_K/N_{L/K}U_L\).
Consequently
\[
 |\operatorname{Cl}_L^G|
 =\frac{h_Kp}{[U_K:U_K\cap N_{L/K}L^\times]}\,
       \frac{|H^2(G,U_L)|}{|H^1(G,U_L)|}.
\]

For clarity, the last ratio is \(1/p\).
Finite modules have Herbrand quotient one, as is seen by counting the
kernel and image of \(g-1\) and of the norm.
The logarithmic unit lattice, modulo its finite torsion, spans the
hyperplane of sum zero in the real vector space on the complex places.
By the Dirichlet unit theorem its rational representation is
\(\mathbb Q[G]^{r_2(K)}-\mathbb Q\).
Herbrand quotient is unchanged under passage to a commensurable
integral lattice, since its finite quotient has Herbrand quotient one.
For \(\mathbb Z[G]\) that quotient is one; for the trivial lattice
\(\mathbb Z\) it is \(p\).
The exact sequence from the sum-zero lattice to
\(\mathbb Z[G]^{r_2(K)}\) and then to \(\mathbb Z\) therefore gives \(1/p\).
These facts also show that the cohomology groups counted above are finite.
We obtain the ambiguous-class formula, in precisely this case,
\[
 |\operatorname{Cl}_L^G|
 =\frac{h_K}{[U_K:U_K\cap N_{L/K}L^\times]}.
\]
Its right side is prime to \(p\).
If the \(p\)-Sylow subgroup \(A\) of \(\operatorname{Cl}_L\) were
nontrivial, the action of \(G\) on \(A\) would have a nontrivial fixed
point. Indeed its orbits have sizes one or \(p\), so
\(|A^G|\equiv|A|\equiv0\pmod p\), while the identity is fixed.
This would make \(p\) divide \(|\operatorname{Cl}_L^G|\), a contradiction.
\end{proof}

The units argument uses the usual rational representation form of the
Dirichlet unit theorem. One can see the rational representation directly
from its real logarithmic embedding: the invariants for every subgroup
have dimension \(r_2(L^J)-1\); these determine the rational character
of the cyclic group and give the stated representation.
Thus the passage between the logarithmic real space and integral
lattices does not assume a rational-valued logarithm.

In \(\mathbb Q(\zeta_{p^{n+1}})/\mathbb Q(\zeta_{p^n})\), the degree
is \(p\), the unique prime above \(p\) is totally ramified, and no other
finite prime ramifies. These facts follow from the Eisenstein
cyclotomic polynomial for \(\zeta_{p^n}-1\) and its discriminant,
which is a power of \(p\).
Starting with \(\mathbb Q(i)\), whose Gaussian integer ring is Euclidean
by nearest-integer division in each coordinate, Lemma~\ref{lem:sr-oneprime}
therefore proves
\[
 h_{\mathbb Q(\zeta_{2^n})}\ \hbox{is odd}\qquad(n\geq2).
\]
This proves the Weber input used in the contributors' dyadic reduction.

The same argument applies to the primes 3 and 7 once the first class
numbers are checked. For \(\mathbb Q(\zeta_3)\), the Minkowski bound is
\((2/\pi)\sqrt3<2\), so every class has an ideal of norm one.
For \(\mathbb Q(\zeta_7)\), the degree is six and the discriminant has
absolute value \(7^5\); its Minkowski bound is
\[
 (4/\pi)^3\frac{6!}{6^6}7^{5/2}<5.
\]
Primes above 2 have residue degree three, and primes above 3 have
residue degree six, because those are the orders of 2 and 3 modulo 7.
Thus no ideal has norm 2, 3 or 4.
Every ideal class has a representative of norm at most four, so the
class number is one.
Here the discriminant is obtained by taking the product of the
derivatives of \(\Phi_p=(X^p-1)/(X-1)\) at its roots;
its absolute value is \(p^{p-2}\).
We have proved
\[
 p\nmid h_{\mathbb Q(\zeta_{p^n})}
 \qquad(p=3,7,\ n\geq1).
\]

\subsection{Strict descent in the regular cyclotomic towers}
\begin{srproposition}\label{prop:sr-regularsha}
For \(p=3,7\), \(F=\mathbb Q(\zeta_{p^n})\), \(n\geq1\), and every
\(a\geq1\), the strict group \(T_{p,a}(F)\) is zero.
\end{srproposition}
\begin{proof}
Let \(D=(\mathbb Q_p/\mathbb Z_p)(-a)\) and represent a strict class
by a continuous cocycle \(z:G_F\to D\), inflated from the group
unramified outside \(p\).
Its image is finite because \(G_F\) is compact and \(D\) is discrete.
Hence all its values belong to \(D[p^b]\) for some \(b\).
At the unique prime above \(p\), strictness says that the cocycle is a
coboundary, \(z(g)=(g-1)v\), on the decomposition group, for some \(v\in D\).
Choose a finite higher cyclotomic layer \(F_m\) in which the cyclotomic
action is trivial on both \(D[p^b]\) and \(v\).

The restriction \(z|_{G_{F_m}}\) is then a homomorphism into \(D[p^b]\).
It is zero on decomposition at \(p\).
At every other finite prime, it is zero on inertia: the coefficient
action is unramified there and the original cocycle was inflated from
the group unramified outside \(p\).
All higher primes above \(p\) are conjugate, so the same vanishing
holds at each choice of a local prime.
Unramified Artin reciprocity makes this homomorphism factor through
\(\operatorname{Cl}_{F_m}\).
The latter group has order prime to \(p\), by the preceding result.
A homomorphism from it into a finite \(p\)-group is zero.
Thus \(z|_{G_{F_m}}=0\).

Let \(F_\infty=\bigcup_mF_m\).
The action on \(D\) is trivial on \(G_{F_\infty}\), and the preceding
vanishing lets \(z\) descend to a continuous cocycle on
\(\Gamma=\operatorname{Gal}(F_\infty/F)\simeq\mathbb Z_p\).
Take a topological generator \(g\) of \(\Gamma\).
The scalar \(g-1\) on \(D\) is
\(\chi_{\rm cyc}(g)^{-a}-1\ne0\).
Multiplication by every nonzero element of \(\mathbb Z_p\) is surjective
on \(\mathbb Q_p/\mathbb Z_p\). Choose \(v_0\in D\) such that
\((g-1)v_0=z(g)\).
After subtracting this coboundary the cocycle vanishes on \(g\), hence
on all its integral powers by the cocycle identity, and hence on
\(\Gamma\) by continuity and density.
The original class is zero.
\end{proof}
This proves the zero bottom-layer budgets for the two exceptional
auxiliary primes, rather than merely naming an unramified-character input.

\subsection{The odd-prime numerical budget and its remaining premise}
For odd \(\ell\), complex conjugation on \(T_{\ell,a}(E)\) has complementary
idempotents \((1+(-1)^ac)/2\) and \((1-(-1)^ac)/2\).
Call their lengths \(s^+_{\ell,E}(a)\) and \(s^-_{\ell,E}(a)\);
the minus eigenvalue is \((-1)^{a+1}\).
When finite, they sum to \(s_{\ell,E}(a)\).

Define \(b(f,k)=16f^{k+1}k!\).
The precise prime-to-degree estimate used in the contributors' odd
packet is the following assertion, here denoted \(\mathrm{PB}(E)\):
all the minus lengths are finite, and, for every \(a\geq1\) and finite
set \(S\) of odd primes not dividing \([E:\mathbb Q]\),
\begin{equation}\label{eq:sr-PB}
 \sum_{\ell\in S}s^-_{\ell,E}(a)\log\ell
       \leq[E:\mathbb Q]\log b(f_E,a+1).
\end{equation}
For the canonical conductor-one field, every odd prime is eligible.
Assuming \eqref{eq:sr-PB}, for \(t\geq2\) we get the fully explicit bound
\begin{equation}\label{eq:sr-minus}
 \sum_{\substack{\ell\in S\\\ell\ne2}}\log\ell
       \sum_{a=1}^{t-1}s^-_{\ell,E}(a)
 \leq3[E:\mathbb Q]t^2\log(tf_E).
\end{equation}
To check the entire numerical estimate, set \(u=tf_E\geq2\).
For \(a<t\), \(16\leq u^4\), \(f_E^{a+2}\leq u^{a+2}\), and
\((a+1)!\leq(a+1)^{a+1}\leq u^{a+1}\).
Thus \(b(f_E,a+1)\leq u^{2a+7}\).
Consequently
\[
 \sum_{a=1}^{t-1}\log b(f_E,a+1)
 \leq(t-1)(2t+5)\log u\leq3t^2\log u.
\]
The last inequality follows from
\(3t^2-(t-1)(2t+5)=t^2-3t+5>0\).
Interchange the two finite sums in \eqref{eq:sr-PB} to obtain
\eqref{eq:sr-minus}.
This proves the polynomial budget from \(\mathrm{PB}(E)\);
it does not prove \(\mathrm{PB}(E)\).
The characterwise Selmer estimate behind that premise must still be
verified in a completed standalone argument.

\subsection{Selecting a whole interval away from finitely many exceptions}
Let \(s,B\geq1\), \(U>0\), and let \(A\) be a set of at most \(R\)
real numbers. There exists \(y\) such that
\[
 \frac{U}{s(sB)^R}\leq y,\qquad sy\leq U,
\]
and, for \(y\leq u\leq sy\) and \(a\in A\),
\begin{equation}\label{eq:sr-logavoid}
 a\leq u\quad\Longrightarrow\quad Ba\leq u.
\end{equation}
Here is the complete induction on \(R\).
For \(R=0\), use \(y=U/s\).
If some \(b\in A\) lies in \((U/(sB),U]\), remove it and apply the
induction hypothesis to \(U'=U/(sB)\) and \(A\setminus\{b\}\).
The resulting lower bound is exactly \(U/[s(sB)^R]\), and
\(sy\leq U'<b\), so \(b\leq u\) never occurs on the resulting interval.
The other points satisfy \eqref{eq:sr-logavoid} by induction.
If there is no such \(b\), choose \(y=U/s\).
Any \(a\leq u\leq U\) must satisfy \(a\leq U/(sB)\), hence
\(Ba\leq U/s=y\leq u\).
The required lower bound follows from \(sB\geq1\).

Now let \(0<\eta\leq C_t\), \(X>1\), and let \(F\) be a set of at most
\(R\) nonnegative integers.
Apply the result with
\(U=\log X,B=C_t/\eta\), and
\(A=\{\log a/C_t:a\in F,\ a>0\}\).
Exponentiating gives \(Y\) with
\[
 X^{1/J}\leq Y,\qquad Y^s\leq X,\qquad
 J=s(sC_t/\eta)^R,
\]
such that for every \(Y\leq x\leq Y^s\), every \(a\in F\) with
\(a\leq x^{C_t}\) already satisfies \(a\leq x^\eta\).
The value \(a=0\) satisfies this automatically.
For \(s=2,R=24,\eta=\epsilon/(2C_K)\), this is exactly the explicit
interval exponent
\begin{equation}\label{eq:sr-J}
 J=2^{25}(2C_KC_t/\epsilon)^{24}.
\end{equation}
The exceptional set must be fixed before \(x\) varies on the whole
interval. A new exceptional set at each \(x\) would not justify this argument.

\subsection{Absorbing the conductor-one budgets uniformly}
We finish the numerical passage needed after the difficult bounds have
been obtained. Fix \(\epsilon>0,C\geq0\), take
\(t=\lfloor x^{401}\rfloor\), and suppose
\[
 x^{2+8\epsilon}\leq2^e\leq2x^{2+8\epsilon},
 \qquad \log(qd)\leq x^{\epsilon/4}.
\]
For the explicit field,
\[
 [E:\mathbb Q]\leq8x^{2+8\epsilon},\qquad
 f_E\leq N=2^{e+2}qd,\qquad
 \log(Nt)\leq(403+8\epsilon)\log x+\log8+x^{\epsilon/4}.
\]
In particular \(\log(Nt)\leq2x^{\epsilon/4}\) for all sufficiently large
\(x\), with a threshold independent of \(q,d,e\) satisfying these inequalities.

Assume the dyadic strict length sum is at most
\(B_2=\exp(x^{4\epsilon/5})\), the odd minus sum is bounded by
\eqref{eq:sr-minus}, and the odd plus sum is at most \((Nt)^C\).
For every finite prime set, the total strict sum is then at most
\begin{equation}\label{eq:sr-budgetsum}
 (\log2)e^{x^{4\epsilon/5}}
 +3[E:\mathbb Q]t^2\log(tf_E)+(Nt)^C.
\end{equation}
The degree and conductor estimates bound its middle term, eventually, by
\(48x^{804+8\epsilon+\epsilon/4}\).
Its last term is at most \(\exp(2Cx^{\epsilon/4})\).
Each of the three displayed terms divided by
\(\exp(x^\epsilon/2)\) tends to zero:
take logarithms, divide by \(x^\epsilon\), and use respectively
\(x^{-\epsilon/5}\to0\), \((\log x)/x^\epsilon\to0\), and
\(x^{-3\epsilon/4}\to0\).
Thus \eqref{eq:sr-budgetsum} is at most \(\exp(x^\epsilon/2)\)
for all sufficiently large \(x\), uniformly in the auxiliary data.
Equation \eqref{eq:sr-integerbudget} supplies a common integer \(M\)
with \(\log M\leq\exp(x^\epsilon/2)\), if the stipulated finiteness
and the three length estimates have been proved.

\paragraph{What is still required here.}
The class-number and strict descent arguments above prove the auxiliary
regularity and zero bottom-layer inputs.
They do not prove the uniform plus estimate \((Nt)^C\) furnished by
the archive's pointwise height theorem.
The odd estimate \eqref{eq:sr-PB}, the quantitative dyadic rank bound
that produces at most 24 exceptional values, and the cohomological
size estimate called Lemma \(K\) in the source also need their full
standalone proofs before this packet constructs the required window.
The exceptional-interval selection and final absorption are proved here
conditional on those precise inputs.

\section{Dyadic cokernel bounds and a common exceptional window}
\label{sec:sr-dyadic-bounds}
The generator and descent results of
Sections~\ref{sec:sr-dyadic-generators}--\ref{sec:sr-dyadic-descent}
permit a quantitative argument with fixed constants.
We prove that argument here. Its last step is conditional on a
single high-twist cohomology bound, stated explicitly in
Definition~\ref{def:sr-K2}; that bound is not proved by the
linear algebra below.

Fix the principal primes \(q,d\), the field
\(L=\mathbb Q(i,\sqrt{-q},\sqrt{-d})\), its cyclotomic extension
\(L_\infty\), and the split unramified module \(M=X_2\).
We have a continuous surjection \(\mathbb Z_2^{24}\to M\).
The generator \(5\) of the cyclotomic group induces the continuous
automorphism \(\tau\) of \(M\).
For \(a\geq1,m\geq0\), write
\[
 c_m(a)=\operatorname{length}_{\mathbb Z_2}
           \bigl(M/((5^a\tau)^{2^m}-1)M\bigr).
\]
Lengths may be infinite until finiteness has been established.
Proposition~\ref{prop:sr-dyadic-descent} identifies this with
\(s_{2,L_m}(a)\), and \eqref{eq:sr-dyadic-return} bounds
\(s_{2,E_h}(a)\) by \(1+c_h(a)\).
In this section the fixed depth is \(j_0=7\).
It is distinct from the exponent \(J\) controlling the analytic window.

\subsection{A free lift and its uniform congruence}
Put \(r=\dim_{\mathbb F_2}(M/2M)\leq24\).
Lifts of an \(\mathbb F_2\)-basis generate \(M\) over
\(\mathbb Z_2\). Indeed their closed span is the compact image of
\(\mathbb Z_2^r\); the quotient \(Q\) satisfies \(Q=2Q\).
Then \(Q=2^nQ\) for every \(n\), while
\(\bigcap_n2^nQ=0\) by the finite-quotient argument in
Lemma~\ref{lem:sr-dyadic-bidual}. Thus \(Q=0\).
If \(r=0\), \(M=0\) and the assertions below are immediate.
Assume henceforth \(r>0\), and choose a surjection
\(\pi:\mathbb Z_2^r\to M\) lifting such a basis.

Choose a matrix \(A\in M_r(\mathbb Z_2)\) such that
\(\pi A=\tau\pi\), by lifting the images of the basis vectors.
Modulo two, \(\pi\) is an isomorphism, so \(\overline A\)
is the automorphism induced by \(\tau\).
This finite action is an image of the pro-\(2\) group
\(\Gamma\simeq\mathbb Z_2\), and consequently has \(2\)-power order.
In characteristic two it is unipotent:
\((\overline A-1)^{2^b}=0\) for some \(b\);
its minimal polynomial divides \((T-1)^r\).
Thus \((\overline A-1)^r=0\).
Since \(r\leq24<32\), we have
\(\overline A^{32}=1\).
Also \(\det\overline A\ne0\), so \(A\in\mathrm{GL}_r(\mathbb Z_2)\).

Write \(A^{32}=1+2B\).
Squaring gives \(A^{64}=1+4(B+B^2)\), and squaring again gives
\begin{equation}\label{eq:sr-A128}
                   A^{128}\equiv1\pmod8.
\end{equation}
Applying \(\pi\) proves the same image congruence for \(\tau\):
\((\tau^{128}-1)M\subseteq8M\).
For \(T_a=5^a\tau\) it follows that
\begin{equation}\label{eq:sr-twist128}
                  (T_a^{128}-1)M\subseteq8M ,
\end{equation}
because \(5^{128a}-1\) is divisible by eight.
The lift is used to prove these congruences; division by two in
an endomorphism of the possibly torsion module \(M\) has not been
assumed.

\subsection{Growth from the fixed layer}
\begin{srproposition}\label{prop:sr-coker-growth}
For every \(a\geq1,m\geq0\),
\[
 c_m(a)\leq c_7(a)+24(m-7)_+ .
\]
If \(c_7(a)\) is finite, every \(c_m(a)\) is finite.
\end{srproposition}
\begin{proof}
For \(m<7\), the factorization of \(T_a^{128}-1\) by
\(T_a^{2^m}-1\) shows that the latter cokernel is a quotient
of the former. Hence \(c_m(a)\leq c_7(a)\).

For \(m=7+k\), put \(V=T_a^{128}\) and
\(S_k=1+V+\cdots+V^{2^k-1}\).
Then \(T_a^{2^m}-1=(V-1)S_k\).
For \(k=0\), \(S_k=1\).
For \(k\geq1\), let \(N=V-1\); by
\eqref{eq:sr-twist128}, \(N^jM\subseteq2^{3j}M\).
The finite binomial identity gives
\[
 S_k=\sum_{j=0}^{2^k-1}{2^k\choose j+1}N^j .
\]
For \(j\geq1\),
\[
 v_2\binom{2^k}{j+1}\geq k-v_2(j+1),\qquad
                       v_2(j+1)\leq j .
\]
The first inequality follows from
\(\binom{2^k}{j+1}=\frac{2^k}{j+1}\binom{2^k-1}{j}\);
the second follows from \(j+1\leq2^j\).
Every nonconstant term consequently maps \(M\) into
\(2^{k+1}M\). Thus
\((S_k-2^k)M\subseteq2^{k+1}M\).

The compact quotient \(Q=M/S_kM\) satisfies
\(2^kQ\subseteq2^{k+1}Q\).
The reverse inclusion is automatic, so iteration and
\(\bigcap_n2^nQ=0\) give \(2^kQ=0\).
Since \(Q\) has at most 24 generators, it is a quotient of
\((\mathbb Z/2^k)^{24}\), and its length is at most \(24k\).

Finally, for any two endomorphisms \(B,C\) of a module, there is an
exact sequence
\[
 0\longrightarrow B(M)/BC(M)
 \longrightarrow M/BC(M)\longrightarrow M/B(M)\longrightarrow0,
\]
and \(B(M)/BC(M)\) is a quotient of \(M/C(M)\).
Apply this with \(B=V-1,C=S_k\).
Length additivity yields \(c_m(a)\leq c_7(a)+24k\),
including in extended lengths, and proves the finiteness assertion.
\end{proof}

\subsection{At most 24 centers govern all the other twists}
Use \(v_2(2)=1\) on finite extensions of \(\mathbb Q_2\).
For a matrix with entries in \(8\mathbb Z_2\), its logarithm and
exponential converge, and the usual inverse identities hold.
Here are the estimates needed to use these functions.
For \(v_2(z)\geq2\),
\[
 v_2(\exp z-1)=v_2(z).
\]
Indeed \((\exp z-1)/z=1+\sum_{n\geq2}z^{n-1}/n!\);
every other term has valuation at least
\(2(n-1)-v_2(n!)\geq n-1\geq1\).
The identical valuation estimates prove convergence for matrices
in the maximum-entry norm.
The formal identities for logarithm and exponential can then be
substituted, because all the relevant power series and their
products converge in this complete matrix algebra.
In particular \(\exp(\log U)=U\) for \(U\equiv1\pmod8\),
and exponentials of commuting matrices multiply.
Also \(v_2(\log5)=2\), since the first term of
\(\log(1+4)\) is four and every subsequent term is divisible by eight.

Take \(U=A^{128}\) and the matrix
\[
              S=-\frac{\log U}{128\log5}\in M_r(\mathbb Q_2).
\]
Choose a finite extension \(K/\mathbb Q_2\) splitting its
characteristic polynomial, and list the eigenvalues
\(s_1,\ldots,s_r\in K\) with multiplicity.
They depend on \(q,d\) and the one chosen lift \(A\), not on the
twist \(a\), the layer \(m\), or a subsequent window.

\begin{srproposition}\label{prop:sr-centers}
If \(B\geq0\) is an integer and
\(|a-s_i|_2\geq2^{-B}\) for every \(i\), then
\begin{equation}\label{eq:sr-center-bound}
                  c_m(a)\leq24(\max(m,7)+2+B)
                         \qquad(m\geq0).
\end{equation}
In particular these cokernels are finite.
\end{srproposition}
\begin{proof}
First let \(m\geq7\).
On the free module, the relevant endomorphism is
\[
 (5^aA)^{2^m}-1
 =\exp\left(2^{m-7}(128a\log5+\log U)\right)-1 .
\]
The matrix in parentheses has entries divisible by eight.
Its eigenvalues, after triangularization over \(K\), are
\(2^m\log5\,(a-s_i)\), and have valuation at least three.
One can check this last assertion without making an isometric
change of basis: a matrix divisible by eight has characteristic
roots divisible by eight, because the characteristic roots of
an integral matrix are integral.
The scalar exponential estimate therefore gives
\[
 v_2\det((5^aA)^{2^m}-1)
       =\sum_{i=1}^r\bigl(m+2+v_2(a-s_i)\bigr)
       \leq r(m+2+B).
\]
The assumed lower distances make every eigenfactor nonzero.
Over the discrete valuation ring \(\mathbb Z_2\), row and column
elimination at an entry of least valuation diagonalizes a matrix
of nonzero determinant.
Its cokernel is finite and has length the valuation of that
determinant: the diagonal factors \(2^{b_i}\) contribute
\(\sum b_i=v_2\det\).
The surjection \(\pi\) induces a surjection from this free
cokernel onto the cokernel defining \(c_m(a)\).
This proves the required bound for \(m\geq7\).
For \(m<7\), use the quotient comparison with \(c_7(a)\)
from the preceding proposition.
\end{proof}

Fix an integer horizon \(Z\geq1\), and put
\(B_Z=\lfloor\log_2 Z\rfloor\).
Define
\begin{equation}\label{eq:sr-exception-set}
 \mathcal E_Z=
 \bigcup_{i=1}^r\{a\in\mathbb Z:1\leq a\leq Z,\
                              |a-s_i|_2<2^{-B_Z}\}.
\end{equation}
Each set in the union has at most one member.
If distinct \(a,b\in[1,Z]\) belonged to it, ultrametricity would give
\(|a-b|_2<2^{-B_Z}\).
But \(0<|a-b|<2^{B_Z+1}\) as ordinary integers, so
\(v_2(a-b)\leq B_Z\), a contradiction.
Thus \(|\mathcal E_Z|\leq r\leq24\), and
\begin{equation}\label{eq:sr-outside-exception}
 c_m(a)\leq24(\max(m,7)+2+B_Z)
       \quad(1\leq a\leq Z,\ a\notin\mathcal E_Z).
\end{equation}
The set is chosen once for the whole horizon.
No new exceptional set is chosen when the real scale varies.

\subsection{The fixed-depth discriminant}
At depth seven,
\[
 L_7=\mathbb Q(\zeta_{512},\sqrt{-q},\sqrt{-d}),\qquad
                         [L_7:\mathbb Q]=1024.
\]
The quadratic radicals are independent over
\(\mathbb Q(\zeta_{512})\) by their odd valuations at \(q,d\).
For later use we prove the explicit bound
\begin{equation}\label{eq:sr-L7-discriminant}
            |\operatorname{disc}L_7|
                     \leq2^{8192}(qd)^{512}.
\end{equation}
The power basis for \(\zeta_{512}\) has discriminant of absolute
value \(2^{2048}\): its polynomial is \(X^{256}+1\), and the
product of the derivatives \(256\zeta^{255}\) at all roots has
absolute value \(256^{256}\).
The quadratic integral bases for the other two fields have
discriminants \(-q,-d\).
Products of these three bases form an integral basis of a full-rank
order in \(L_7\), since the three fields are linearly disjoint.
The trace matrix is the tensor product of the three trace matrices.
Its determinant is
\[
       (2^{2048})^4 q^{512}d^{512}
                          =2^{8192}(qd)^{512}.
\]
For an integral full-rank order of index \(I\) in the maximal
order, the discriminants differ by \(I^2\), by the determinant
of the basis change.
Dropping this index proves \eqref{eq:sr-L7-discriminant}.
In particular, if \(x\geq1\) and \(\log(qd)\leq x^{\epsilon/4}\),
\begin{equation}\label{eq:sr-L7-logdisc}
 \log(2|\operatorname{disc}L_7|)
 \leq (8193\log2+512)x^{\epsilon/4}.
\end{equation}

\subsection{The one high-twist input still required}
Let \(S_2(F)\) be the primes above two in a number field \(F\).
Write
\[
 \mathcal H_2(F,a)=
 H^2(G_{F,S_2(F)},\mathbb Z_2(a+1)).
\]
We use continuous cohomology for the compact coefficient module;
equivalently one uses the inverse finite-coefficient system with
its continuous transition maps.
The classical Poitou--Tate duality background gives, for totally
imaginary \(F\), the pairing
\[
 T_{2,a}(F)\ \times\
 \ker\left(\mathcal H_2(F,a)\longrightarrow
       \prod_{v\mid2}H^2(F_v,\mathbb Z_2(a+1))\right)
                      \longrightarrow\mathbb Q_2/\mathbb Z_2 .
\]
It is a perfect Pontryagin pairing.
This is the same global and local Tate duality input used in
Section~\ref{sec:minus-characterwise}, with mutually dual
coefficients of twists \(-a\) and \(a+1\).
There are no real-place corrections for these fields.
The finite-coefficient pairings pass to the stated limits:
the finite-level groups are finite, hence their inverse systems
satisfy the Mittag--Leffler condition; direct limits are exact,
and only finitely many local places occur.
Consequently finiteness of \(\mathcal H_2(F,a)\) gives
\begin{equation}\label{eq:sr-Sha-H2}
 s_{2,F}(a)\log2\leq\log\#\mathcal H_2(F,a).
\end{equation}
This inequality uses the inclusion of the displayed kernel into
\(\mathcal H_2(F,a)\), not an equality with the whole \(H^2\).

\subsection{The local cohomology terms throughout the tower}
For these particular fields the local correction in global duality can
be computed exactly. This removes a further imported estimate without
proving the remaining global class-counting bound.

\begin{srproposition}\label{prop:sr-H2-exact}
For every \(m\geq0\), the field \(F=L_m\) has exactly two primes above two.
For every
\(a\geq1\), the group \(\mathcal H_2(F,a)\) is finite if and only if
the strict group \(T_{2,a}(F)\) is finite. When they are finite,
\begin{equation}\label{eq:sr-H2-exact}
 \operatorname{length}_{\mathbb Z_2}\mathcal H_2(F,a)
                 =s_{2,F}(a)+m+2+v_2(a).
\end{equation}
\end{srproposition}
\begin{proof}
First determine the completions. Put \(b=m+2\geq2\) and \(C=\mathbb Q(\zeta_{2^b})\).
The polynomial of \(\zeta_{2^b}-1\) is
\((Y+1)^{2^{b-1}}+1\). Its constant term is two, every other
nonleading coefficient is divisible by two, and the leading coefficient
is one. The divisibility of the intermediate binomial coefficients
follows by reducing \((1+Y)^{2^{b-1}}\) modulo two, where it is
\(1+Y^{2^{b-1}}\). Thus the polynomial is Eisenstein at two.
Consequently \(C\) has a single dyadic completion
\(C_2=\mathbb Q_2(\zeta_{2^b})\), totally ramified of degree \(2^{b-1}\).

An odd \(2\)-adic unit congruent to one modulo eight is a square in
\(\mathbb Q_2\). To see sufficiency directly, start with an odd square
root modulo \(2^3\). Given an odd \(b_n\) with
\(b_n^2\equiv u\pmod {2^n}\), \(n\geq3\), choose
\(b_{n+1}=b_n+\delta 2^{n-1}\), \(\delta\in\{0,1\}\).
Modulo \(2^{n+1}\), the square changes by \(\delta b_n2^n\).
Exactly one choice gives the next congruence. These choices are Cauchy
and their limit squares to \(u\).
Necessity follows by squaring an odd integer modulo eight.
Since \(d\equiv7\pmod {16}\), \(-d\equiv1\pmod8\).
Hence adjoining \(\sqrt{-d}\) splits over \(C_2\).

On the other hand \(-q\equiv5\pmod8\).
For \(u=-q\), the polynomial
\[
                 Z^2-Z+(1-u)/4
\]
reduces to \(Z^2+Z+1\) over \(\mathbb F_2\).
It is irreducible and separable. Its root generates the unramified
quadratic extension \(U/\mathbb Q_2\).
For completeness, the ring obtained by adjoining this root to
\(\mathbb Z_2\) is free of rank two and its reduction is
\(\mathbb F_4\). An element whose reduction is nonzero is invertible:
lift its inverse modulo two and correct by the convergent geometric
series. Every nonzero element is a power of two times such a unit,
by factoring the least valuation of its two coefficients.
This ring is therefore a discrete valuation ring with uniformizer two
and residue field \(\mathbb F_4\). It is integrally closed, hence is
the full ring of integers of its fraction field. This proves the
unramified assertion.

A totally ramified extension and an unramified extension have trivial
intersection; after unramified base change \(C_2U/C_2\) has degree
two. Since \(F=C(\sqrt{-q},\sqrt{-d})\), its dyadic algebra is
\[
              F\otimes_{\mathbb Q}\mathbb Q_2
                     \simeq C_2U\times C_2U .
\]
Thus there are exactly two dyadic primes and their completions are
both \(C_2U\), of degree \(2^b\) over \(\mathbb Q_2\).
Here \([F:\mathbb Q]=2^{b+1}\), by the independence of the two
odd radicals already used above.

We next determine the cyclotomic images. Every finite extension in
\(\mathbb Q_2(\mu_{2^\infty})/C_2\) is totally ramified, whereas
\(C_2U/C_2\) is unramified. Their intersection is \(C_2\).
It follows that the cyclotomic image of each local absolute Galois
group is
\[
                    1+2^b\mathbb Z_2 .
\]
The global image is the same. Indeed \(F/C\) is biquadratic.
Each of its three quadratic intermediate fields is ramified at
\(q\) or \(d\), since the corresponding radical
\(-q,-d\), or \(qd\) has odd valuation there and \(C/\mathbb Q\)
is unramified at these odd primes.
No such field lies in the cyclotomic \(2\)-power tower, which is
unramified outside two.
As the cyclotomic tower over \(C\) is procyclic, a nontrivial
intersection with \(F\) would contain a quadratic intermediate
field. Thus \(F\cap\mathbb Q(\mu_{2^\infty})=C\), proving the
global image assertion.

The group \(1+2^b\mathbb Z_2\) is topologically generated by
\(1+2^b\). One direct verification uses logarithm: it identifies
this group with \(2^b\mathbb Z_2\), and
\(\log(1+2^b)/2^b\) is a unit.
For \(a\geq1\),
\[
                 v_2((1+2^b)^a-1)=b+v_2(a).
\]
For odd \(a\) the first binomial term has valuation \(b\) and all
subsequent terms have larger valuation.
Replacing an odd exponent by twice that exponent raises the valuation
by one, using \(z^2-1=(z-1)(z+1)\) with \(z\equiv1\pmod4\).
Iteration proves the formula for every \(a\).
The same valuation holds for the inverse power minus one.

Let \(D_a=\mathbb Q_2/\mathbb Z_2(-a)\).
Its invariants under any of the global or local groups just
considered are therefore cyclic of order
\(2^{b+v_2(a)}\). A class in \(\mathbb Q_2/\mathbb Z_2\)
is fixed precisely when it is killed by the action of the dense
generator minus one; continuity then gives invariance under the
whole group. Local Tate duality identifies
\(H^2(F_v,\mathbb Z_2(a+1))\) with the Pontryagin dual of these
local invariants, so it has length \(b+v_2(a)\).

Finally use the following tail of global Poitou--Tate duality:
\[
 0\longrightarrow T_{2,a}(F)^\vee
 \longrightarrow \mathcal H_2(F,a)
 \longrightarrow \bigoplus_{v\mid2}
                    H^2(F_v,\mathbb Z_2(a+1))
 \longrightarrow H^0(F,D_a)^\vee
 \longrightarrow0 .
\]
Here \({}^\vee\) denotes Pontryagin duality.
The leftmost identification is the strict pairing printed above.
The final map is dual to the diagonal map on invariants, which is
injective. The totally imaginary hypothesis removes archimedean
terms. As for the preceding pairing, this is the compact-coefficient
tail obtained from the finite-coefficient duality sequence.
The other three groups in this sequence are finite.
It follows at once that its first two groups are finite
simultaneously. In that case length additivity gives
\[
 s_{2,F}(a)+2(b+v_2(a))-(b+v_2(a))
                  =s_{2,F}(a)+m+2+v_2(a),
\]
as asserted.
\end{proof}

\begin{srproposition}\label{prop:sr-H2-growth}
If \(\mathcal H_2(L_7,a)\) is finite, then every
\(\mathcal H_2(L_m,a)\) is finite, and
\[
 \operatorname{length}_{\mathbb Z_2}\mathcal H_2(L_m,a)
 \leq
 \operatorname{length}_{\mathbb Z_2}\mathcal H_2(L_7,a)
                                      +25(m-7)_+ .
\]
\end{srproposition}
\begin{proof}
Finiteness at depth seven gives finiteness of \(c_7(a)\) by
Proposition~\ref{prop:sr-H2-exact}. Proposition~\ref{prop:sr-coker-growth}
then gives finiteness of every \(c_m(a)\), and the exact local
calculation gives finiteness of the corresponding \(H^2\).
For \(m\geq7\), add \(m+2+v_2(a)\) to
\(c_m(a)\leq c_7(a)+24(m-7)\), and apply
\eqref{eq:sr-H2-exact} at both depths. This gives \(25(m-7)\).
For \(m<7\), \(c_m(a)\leq c_7(a)\) and \(m+2\leq9\),
giving the asserted bound with zero extra term.
\end{proof}

In particular, the missing size theorem is a global strict-group
estimate, rather than an additional local correction.
More precisely, suppose a uniform \(C\geq1\) gave
\[
 s_{2,L_7}(a)\log2\leq
 \exp\bigl((2048a)^C\log(2|\operatorname{disc}L_7|)\bigr)
\]
together with finiteness, for every pair and every \(a\geq1\).
Then \(K_2(C+1)\) would follow. Set
\(B=(2048a)^C\log(2|\operatorname{disc}L_7|)\).
We have \(B\geq2048a\log2\geq1\), and
\((9+v_2(a))\log2\leq B\), since \(v_2(a)\leq a\).
By \eqref{eq:sr-H2-exact},
\[
 \log\#\mathcal H_2(L_7,a)\leq e^B+B
        \leq e^{2B}
        \leq
 \exp\bigl((2048a)^{C+1}\log(2|\operatorname{disc}L_7|)\bigr).
\]
For the middle inequality, \(B\leq e^B\) and
\(2e^B\leq e^{2B}\) when \(B\geq1\).
Conversely \(K_2(C_K)\) bounds the strict length by
\eqref{eq:sr-Sha-H2}. Neither direction proves the global strict
estimate itself.

\begin{srdefinition}[The specialized high-twist bound]
\label{def:sr-K2}
The assertion \(K_2(C_K)\), for a fixed real \(C_K\geq1\), says
that for every principal pair \(q,d\) and every integer \(a\geq1\),
the group \(\mathcal H_2(L_7,a)\) is finite and
\begin{equation}\label{eq:sr-K2}
 \log\#\mathcal H_2(L_7,a)
 \leq
 \exp\bigl((2048a)^{C_K}
                \log(2|\operatorname{disc}L_7|)\bigr).
\end{equation}
\end{srdefinition}
The archive's general Lemma \(K\) implies this assertion by the
degree calculation above and \(\mu_4\subset L_7\).
Its class-counting and geometric proof has not been replaced by
a standalone proof here. Thus \(K_2(C_K)\) remains an explicit
input to the next theorem.

\subsection{The dyadic estimate on the same common windows}
\begin{srtheorem}\label{thm:sr-dyadic-window}
Assume \(K_2(C_K)\).
Fix \(C_t=401\), \(0<\epsilon<1/52\), and put
\[
 \eta=\frac{\epsilon}{2C_K},\qquad
 J=2^{26}(2C_KC_t/\epsilon)^{24}.
\]
For sufficiently large \(X\), uniformly in the principal pair
\(q,d\) satisfying
\(\log(qd)\leq X^{\epsilon/(4J)}\), there is \(Y\) with
\[
 X^{1/J}\leq Y,\qquad Y^2\leq X
\]
such that for every \(Y\leq x\leq Y^2\), every natural \(h=h(x)\)
with
\[
 x^{2+8\epsilon}\leq2^h\leq2x^{2+8\epsilon},
                  \qquad t=\lfloor x^{401}\rfloor,
\]
all \(s_{2,E_h}(a)\), \(1\leq a<t\), are finite and
\begin{equation}\label{eq:sr-full-dyadic-window}
                \sum_{a=1}^{t-1}s_{2,E_h}(a)
                          \leq\exp(x^{4\epsilon/5}).
\end{equation}
The threshold depends on \(C_K,\epsilon\), not on \(X,q,d\)
or the subsequent choices of \(x,h\).
An admissible depth always exists: take
\(h=\lceil(2+8\epsilon)\log_2x\rceil\).
\end{srtheorem}
\begin{proof}
For the fixed pair choose the one module, lift and list of centers
above. Take \(Z=\lfloor X^{401}\rfloor\) and form
\(\mathcal E_Z\) from \eqref{eq:sr-exception-set}.
Apply the finite-exception interval selection proved in
\eqref{eq:sr-logavoid}--\eqref{eq:sr-J}, with \(s=2\),
\(R=24\) and this \(\eta\).
Its exponent is \(J/2\); therefore its stronger lower bound
\(Y\geq X^{2/J}\) implies the required \(Y\geq X^{1/J}\).
On the resulting interval,
\[
 a\in\mathcal E_Z,\quad a\leq x^{401}
                         \quad\Longrightarrow\quad a\leq x^\eta .
\]
This one exceptional set is used for the entire interval.

For every \(x\) there, \(x\leq X\), so \(t\leq Z\).
Also
\[
 \log(qd)\leq x^{\epsilon/4},\quad
 B_Z\leq401\log_2 X\leq401J\log_2 x,\quad
 h\leq1+(2+8\epsilon)\log_2 x.
\]
For a twist \(a\notin\mathcal E_Z\), combine
\eqref{eq:sr-outside-exception} with
\eqref{eq:sr-dyadic-return}.
Summing over fewer than \(x^{401}\) such twists contributes at most
\[
                     C_1 x^{401}(1+\log x),
\]
where \(C_1\) depends only on \(\epsilon,C_K\) through \(J\).

For an exceptional twist in the range, \(a\leq x^\eta\).
By \eqref{eq:sr-K2} and~\eqref{eq:sr-L7-logdisc},
\[
 \log\#\mathcal H_2(L_7,a)
 \leq\exp(D_0x^{3\epsilon/4}),\qquad
 D_0=2048^{C_K}(8193\log2+512).
\]
Indeed \(a^{C_K}\leq x^{\epsilon/2}\).
Once \(x^{\epsilon/40}\geq D_0\), this is at most
\(\exp(x^{31\epsilon/40})\).
Equations~\eqref{eq:sr-Sha-H2} and
\eqref{eq:sr-dyadic-coker}, followed by
Proposition~\ref{prop:sr-coker-growth} and
\eqref{eq:sr-dyadic-return}, now give
\[
 s_{2,E_h}(a)
 \leq1+\frac{\exp(x^{31\epsilon/40})}{\log2}
                          +24(h-7)_+ .
\]
There are at most 24 exceptional twists.
All these lengths are finite by \(K_2(C_K)\).

The total is therefore at most
\[
 C_1x^{401}(1+\log x)
 +24+\frac{24}{\log2}\exp(x^{31\epsilon/40})
                              +576h .
\]
Dividing this expression by \(\exp(x^{4\epsilon/5})\) gives a
quantity tending to zero:
the logarithm of the polynomial terms is \(O(\log x)\), and
\(31\epsilon/40<4\epsilon/5\).
Thus \eqref{eq:sr-full-dyadic-window} holds for all sufficiently
large \(x\), uniformly in the pair.
Increasing the lower threshold for \(X\) ensures that every
\(x\geq X^{1/J}\) meets all those thresholds.
This proves the theorem and all its quantifiers.
\end{proof}

This proves the dyadic tower growth, finite-center estimate,
24-element exceptional-set construction, discriminant comparison
and common-window deduction.
The high-twist bound \(K_2(C_K)\) is still required.
The uniform plus estimates and the geometric supply assertion
are separate inputs; none follows from the dyadic calculation.


\section{The global high-twist bound from integral homology}
\label{sec:sr-global-cohomology}
We prove the quantitative cohomology estimate needed at the exceptional
dyadic twists. The argument uses integral linear-group homology, retaining
the finite stabilizers and hence the small coefficient primes.
The alternative route in Sections 6.76--6.83 of the
\href{https://principia-math.com/api/ws/files/8c527eb7a1dbba0b1dc69c2ba0e60db0}{777-page reorganized working draft}
suggested this approach. We give the geometric count, the integral
torsion calculation and all numerical absorptions explicitly.
The standard comparison theorems used are stated below with their
published sources; no theorem of the Lean archive supplies the estimate.

For a totally imaginary number field \(F\), put
\[
 n=[F:\mathbb Q],\qquad D=|\operatorname{disc}F|,\qquad
 S_\ell=\{v:v\mid\ell\},\qquad
 X_{\ell,a}(F)=H^2(G_{F,S_\ell},\mathbb Z_\ell(a+1)).
\]
Here \(G_{F,S_\ell}\) is the full Galois group unramified outside
\(S_\ell\), and cohomology uses continuous cochains with the compact
coefficient module.

\begin{srtheorem}\label{thm:sr-global-high-twist}
Let \(a\geq1\) and let \(\ell\) be prime.
Suppose \(\mu_\ell\subset F\) if \(\ell\) is odd, and
\(\mu_4\subset F\) if \(\ell=2\).
Then \(X_{\ell,a}(F)\) is finite and
\begin{equation}\label{eq:sr-global-high-twist}
 \log\#X_{\ell,a}(F)
       \leq \exp\bigl((2an)^{24}\log(2D)\bigr).
\end{equation}
The exponent is independent of \(F,\ell,a\).
\end{srtheorem}
The proof occupies the following subsections.

\subsection{The standard comparison theorems}
We use the following established results, in these precise forms.
First, the even group \(K_{2a}(\mathcal O_F)\), \(a\geq1\),
is finite by Quillen's finite generation and Borel's rank theorem.
The arithmetic Quillen--Lichtenbaum comparison gives
\begin{equation}\label{eq:sr-K-comparison}
 K_{2a}(\mathcal O_F)\otimes\mathbb Z_\ell
                      \simeq X_{\ell,a}(F).
\end{equation}
For odd \(\ell\), see
\href{https://www.sas.rochester.edu/mth/sites/doug-ravenel/otherpapers/weibel-ints.pdf}{Weibel, Theorem 70 and Corollary 71}.
For two, the totally imaginary comparison is
\href{https://sites.math.rutgers.edu/~weibel/archive/papers-dir/RognesWeibel.pdf}{Rognes--Weibel, Theorem 0.4}.
Inverting \(\ell\) does not change the \(\ell\)-primary group in
positive even degree: in Quillen's localization sequence the adjacent
positive-degree groups of the finite residue fields have orders prime
to \(\ell\).
Finite-coefficient restricted Galois cohomology agrees with the
corresponding etale cohomology of \(\mathcal O_F[1/\ell]\).
The hypotheses for this comparison are \(S_\ell\) containing all
places above \(\ell\), and no real-place correction at two; see
\href{https://www.mathi.uni-heidelberg.de/~schmidt/papers/kpi1.pdf}{Schmidt, Proposition 2.3 and its proof}.
The finite-coefficient groups are finite, so inverse limits have the
Mittag--Leffler property. The degree-two compact-coefficient comparison
therefore has no \(\lim^1 H^1\) term. These observations specify both
the ramification set and the twist in \eqref{eq:sr-K-comparison}.
In particular this comparison proves finiteness before we estimate
the order.

Second, the plus construction gives
\(K_m(\mathcal O_F)=\pi_m(B\mathrm{GL}(\mathcal O_F)^+)\) for \(m>0\).
The indicated connected component is an infinite loop space, and the
plus map preserves integral homology.
Integral homology stability for number rings identifies its \(m\)-th
homology with
\(H_m(\mathrm{GL}_N(\mathcal O_F),\mathbb Z)\) when
\(N\geq2m+1\); we use the more generous choice
\(N=2m+4\).
This is the number-ring stability theorem of van der Kallen, also used
in
\href{https://www.maths.tcd.ie/EMIS/journals/DMJDMV/vol-kato/soule.dm.pdf}{Soule, Section 4.2}.

Third,
\href{https://msp.org/agt/2016/16-2/agt-v16-n2-p12-p.pdf}{Mathew, Theorem 6.3}
bounds the kernel of the unstable Hurewicz map of a connected infinite
loop space. In degree \(2a\), its \(\ell\)-primary kernel is killed by
\(\ell^{b_\ell}\), where one may take
\[
 b_2=a+3,\qquad
 b_\ell=\left\lceil\frac{2a+2}{2\ell-2}\right\rceil+1
                                             \quad(\ell>2).
\]
Both are at most \(a+3\leq4a\).
The connectivity parameter in that theorem is one; no simple
connectivity of \(B\mathrm{GL}^+\) is required.
We will also use Minkowski's convex body and ideal-class theorems,
the Kummer exact sequence, and the local-invariant form of
Brauer--Hasse--Noether. Their applications and estimates are printed
below.

\subsection{A positive-definite region of forms}
Fix \(N\geq1\), set \(d=nN\), and choose one embedding from each
complex pair. In
\[
 V=\prod_{\sigma:F\hookrightarrow\mathbb C/\mathrm{conj}}\mathbb C^N
                       \simeq\mathbb R^d
\]
use the Euclidean norm
\(\|x\|^2=\sum_\sigma\sum_j|x_{\sigma,j}|^2\), and let
\(\Lambda=\mathcal O_F^N\) with its diagonal embedding.
It is a full lattice with
\[
           \operatorname{covol}(\Lambda)=2^{-d/2}D^{N/2}.
\]
Indeed the determinant of the embedding matrix of an integral basis
squares to the discriminant; replacing each pair of conjugate
coordinates by its real and imaginary parts contributes a factor
\(1/2\) to the absolute determinant.

Every nonzero element of \(\Lambda\) has norm at least one.
A nonzero coordinate \(z\in\mathcal O_F\) satisfies
\(\prod_\sigma|\sigma z|^2=|\operatorname{Norm}_{F/\mathbb Q}z|\geq1\),
so at least one of its archimedean absolute values is at least one.
Let
\[
 W=\prod_\sigma\operatorname{Herm}_N(\mathbb C),\qquad
 q_h(x)=\sum_\sigma x_\sigma^*h_\sigma x_\sigma,\qquad
 P=\{h\in W:q_h(y)\geq1\ (0\ne y\in\Lambda)\}.
\]
The set \(P\) is closed, convex and nonempty: the tuple of identity
matrices belongs to it. Let \(T\subset W\) be the closed cone of
positive semidefinite tuples, and put
\[
 I_0=D^N d^d .
\]

\begin{srlemma}\label{lem:sr-form-region}
Every \(h\in P\) is positive definite and
\begin{equation}\label{eq:sr-form-det}
                       \det_{\mathbb R}h\geq I_0^{-2}.
\end{equation}
Near each point of \(P\), only finitely many lattice inequalities
are needed to define \(P\). A point is a vertex exactly when its
active functionals \(h\mapsto q_h(y)\) span \(W^*\).
\end{srlemma}
\begin{proof}
The rational span \(\Lambda\otimes\mathbb Q\) is dense in \(V\).
If \(q_h\) were negative somewhere, it would be negative on a rational
lattice direction. Clearing the denominator would contradict
\(q_h(y)\geq1\). Thus \(h\in T\).

For \(\delta>0\), the ellipsoid
\(q_{h+\delta\,1}(x)<1\) contains no nonzero lattice point.
If \(v_d\) is the Euclidean unit-ball volume, Minkowski gives
\[
 \frac{v_d}{\sqrt{\det_{\mathbb R}(h+\delta\,1)}}
                      \leq2^d\operatorname{covol}(\Lambda).
\]
The cube \([-1/\sqrt d,1/\sqrt d]^d\) is contained in the unit
ball, so \(v_d\geq2^dd^{-d/2}\).
Dropping the factor \(2^{-d/2}\leq1\) in the covolume formula,
\[
 \sqrt{\det_{\mathbb R}(h+\delta\,1)}
       \geq D^{-N/2}d^{-d/2}=I_0^{-1/2}\geq I_0^{-1}.
\]
Let \(\delta\downarrow0\). The determinant stays positive, hence
the semidefinite tuple \(h\) is positive definite and
\eqref{eq:sr-form-det} follows.

In a sufficiently small neighborhood of such a tuple, all quadratic
forms are bounded below by a common positive multiple of \(\|x\|^2\).
Only a bounded set of lattice vectors can violate an inequality
there, and that set is finite.
Thus \(P\) is locally cut out by finitely many linear inequalities.
If the active functionals at \(h\) do not span \(W^*\), a nonzero
annihilating direction gives a short segment through \(h\) inside
\(P\); the finitely many inactive nearby inequalities have positive
slack. The point is then not extreme.
If they span, any expression of \(h\) as the midpoint of two points
of \(P\) forces all the active inequalities to be equalities at both
points, and their span forces both points to equal \(h\).
This proves the vertex characterization.
\end{proof}

\subsection{The normal-cone complex is integrally acyclic}
Let \(U\subset W^*\) be the open cone of functionals strictly positive
on \(T\setminus\{0\}\).
For each vertex \(v\) of \(P\), set
\[
 C_v=\{u\in U:u(h)\geq u(v)\text{ for all }h\in P\}.
\]
These are convex, relatively closed subsets of \(U\).

\begin{srlemma}\label{lem:sr-normal-acyclic}
The \(C_v\) cover \(U\).
The augmented chain complex freely generated in degree \(p\) by all
ordered tuples \((v_0,\ldots,v_p)\) with
\(C_{v_0}\cap\cdots\cap C_{v_p}\ne\varnothing\), repetitions allowed,
is exact. Its boundary is the alternating sum deleting one entry.
\end{srlemma}
\begin{proof}
For \(u\in U\), compactness of the unit sphere in \(T\) gives
\(u(h)\geq c\|h\|\) on \(T\) for some \(c>0\).
Consequently a sublevel of \(u\) on \(P\) is compact, so its minimum
is attained. The compact minimizing face has an extreme point:
successively minimize independent linear coordinates on that face
until the remaining face is a point. An extreme point of a
minimizing face is extreme in \(P\), hence is a vertex by
Lemma~\ref{lem:sr-form-region}. This proves the cover.

Only finitely many \(C_v\) meet a fixed compact \(K\subset U\).
Fix \(h_0\in P\). On \(K\) the constants \(c\) above have a common
positive lower bound and \(u(h_0)\) has a finite upper bound.
If \(u\in C_v\cap K\), then \(u(v)\leq u(h_0)\); all such \(v\)
therefore lie in a fixed compact subset of \(P\).
Local finite presentation, followed by a finite covering of that
compact set, gives a finite list of possible active inequalities.
Each vertex is specified by a spanning sublist, so only finitely
many occur.

Consider a finite cycle in the ordinary oriented nerve of the
\(C_v\). Choose one witness in the common intersection for every
simplex in its support, and let \(K\) be the convex hull of those
witnesses. This is a compact convex subset of \(U\).
The nonempty \(C_v\cap K\) form a finite closed convex cover of \(K\),
and its nerve contains the cycle.

Here is the finite closed-convex nerve argument used at this point.
For a finite family of compact sets, every subfamily with empty
intersection still has empty intersection after sufficiently small
open thickening. Otherwise a sequence of almost common points would
have a convergent subsequence with limit in the common intersection.
Choose one thickness working for all subfamilies.
The thickenings are open and convex, their intersection pattern is
unchanged, and their union is the corresponding open neighborhood
of \(K\), which is convex. The open convex nerve is acyclic:
the double complex of singular chains on intersections computes
both the nerve and the union, since the intersections are
contractible. Exactness in the covering direction follows after
subdividing each finite singular chain until its simplices lie in
members of the open cover. Thus the original finite cycle bounds.
Every finite cycle in the full ordinary nerve consequently bounds.

Finally compare ordered tuples with the ordinary oriented nerve.
Send a tuple with repeated entries to zero and a tuple with distinct
entries to its sorted oriented simplex. This is a chain map:
the two faces that remove a repeated entry cancel, and on distinct
entries the assertion is the usual orientation rule.
Increasing tuples give a section.
The composite is chain homotopic to the identity on the ordered
complex. To construct the homotopy inductively on a basis tuple,
subtract the already defined homotopy on its boundary from the
difference of the two chain maps.
The result is a cycle supported on the same finite vertex set.
That entire tuple subcomplex contracts by prepending one chosen
vertex, because all its vertices have a common witness.
Use this contraction to fill the cycle and define the next homotopy.
This proves the induction in every degree.
Hence the augmented ordered complex has the same, zero, homology
as the ordinary nerve.
\end{proof}

\subsection{A common frame and a uniform orbit count}
The group \(\Gamma=\mathrm{GL}_N(\mathcal O_F)\) acts by change of
variables on \(P\), preserving the preceding complex.

\begin{srlemma}\label{lem:sr-common-frame}
Every tuple in that complex has \(N\) \(F\)-independent lattice
vectors of value one for every member of the tuple.
If they are the columns of \(A\in M_N(\mathcal O_F)\), then
\[
                   1\leq|\operatorname{Norm}\det A|\leq I_0.
\]
\end{srlemma}
\begin{proof}
Let \(u\) be a common minimizing functional and \(h\) the average
of the tuple. Its active vectors are precisely those active for every
member, since all their individual values are at least one.
If their \(F\)-span were proper, choose a nonzero \(F\)-linear
functional annihilating them.
At one complex place its squared absolute value, and zero at the
other places, defines a nonzero element \(b\in T\).
It vanishes on all active vectors.
Local finite presentation gives \(h-\delta b\in P\) for sufficiently
small \(\delta>0\).
But \(u(b)>0\), contradicting minimality of \(h\).
Thus the common active vectors span \(F^N\), and \(N\) independent
ones can be chosen.

Put \(I=|\operatorname{Norm}\det A|\), a positive integer.
For a member \(h_i\) of the tuple, the pulled-back blocks
\(A_\sigma^*h_{i,\sigma}A_\sigma\) have diagonal entries at most
one: for each column, their sum over the places is one.
Hadamard therefore gives
\(\det_{\mathbb R}(A^*h_iA)\leq1\).
Its real determinant is \(I^2\det_{\mathbb R}h_i\).
Apply \eqref{eq:sr-form-det} to get \(I\leq I_0\).
\end{proof}

Define the positive integers
\[
 R=I_0^2N^d,\qquad V_0=(2R+1)^d,\qquad
 B_p=I_0^{d^2+d+1}(V_0+1)^{dN(p+1)}.
\]
\begin{srproposition}\label{prop:sr-orbit-bound}
The number of \(\Gamma\)-orbits of ordered tuples of \(p+1\) vertices
with a common witness is at most \(B_p\).
Equivalently, \(B_p\) bounds the degree-\(p\) simplex orbits.
For \(L_0=\log(2D)\), and \(m\geq1\), one common bound
\(B=B_{m+1}\) in degrees \(0,\ldots,m+1\) satisfies
\begin{equation}\label{eq:sr-orbit-log}
                   \log B\leq22(m+2)N d^4L_0 .
\end{equation}
\end{srproposition}
\begin{proof}
The row module of the matrix \(A\) in the preceding lemma is an
additive sublattice of \(\mathcal O_F^N\) of index \(I\leq I_0\).
Matrices with the same row module differ by left multiplication by
\(\mathrm{GL}_N(\mathcal O_F)\), since their rows are bases of that
same free module.
Thus these row modules distinguish the possible left orbits of frames.

A fixed integral basis identifies the ambient additive group with
\(\mathbb Z^d\). Hermite normal form bounds the number of its
sublattices of index \(j\) by \(j^{d(d+1)/2}\):
there are at most \(j^d\) diagonal choices and at most
\(j^{d(d-1)/2}\) reduced off-diagonal choices.
Summing over \(1\leq j\leq I_0\) gives at most
\(I_0^{d^2+d+1}\) possibilities.
Counting all additive sublattices only enlarges the number of
\(\mathcal O_F\)-row modules.
Choose one frame representative per orbit and move the tuple so
that its common frame is this representative.

For a fixed representative \(A\), put \(h'=A^*hA\).
All its real eigenvalues are at most \(N\), by the diagonal bound
and the trace bound in each positive Hermitian block.
Their product is at least \(I_0^{-2}\), since \(I\geq1\).
Thus every eigenvalue is at least \(I_0^{-2}N^{-(d-1)}\).
If \(y\in\mathcal O_F^N\) is active for one such vertex and
\(z=A^{-1}y\), then
\[
                    \|z\|\leq I_0N^{(d-1)/2}.
\]
Moreover \(Iz\in\mathcal O_F^N\): the finite additive group
\(\mathcal O_F^N/A\mathcal O_F^N\) has order \(I\), hence is killed
by \(I\). Therefore
\(\|Iz\|\leq I_0^2N^{(d-1)/2}\leq R\).

Distinct points of \(\Lambda\) have distance at least one.
Disjoint open balls of radius \(1/2\) about points in its radius-\(R\)
ball fit inside the radius-\(R+1/2\) ball. Comparing Euclidean volumes
bounds their number by \(V_0\).
The injection \(y\mapsto IA^{-1}y\) therefore leaves at most \(V_0\)
possible active vectors.
A vertex is determined by independent active equations numbering
\(\dim_{\mathbb R}W=nN^2/2\leq dN\).
Its number is at most \((V_0+1)^{dN}\), allowing a dummy entry in
an encoding shorter than \(dN\).
Taking \(p+1\) vertices and multiplying by the number of frames proves
the formula for \(B_p\).

For clarity, the logarithmic estimate also has fixed constants.
Use \(L_0\geq1/2\), \(N\leq d\), and
\(\log t\leq t\) for \(t\geq1\). Then
\[
 \log I_0=N\log D+d\log d\leq3d^2L_0,
 \qquad
 \log R\leq2NL_0+3d^2.
\]
As \(R\geq1\), these give
\[
 \log V_0\leq d(2+2NL_0+3d^2)\leq12d^3L_0,
 \qquad
 \log(V_0+1)\leq13d^3L_0.
\]
Consequently
\[
 \log B\leq
       [9+13(m+2)N]d^4L_0
                 \leq22(m+2)Nd^4L_0,
\]
which is \eqref{eq:sr-orbit-log}.
\end{proof}

\subsection{Integral homology with the stabilizers retained}
\begin{srlemma}\label{lem:sr-stabilizer-size}
Every stabilizer of an ordered simplex is finite of order at most
\(J=3^{d^2}\).
\end{srlemma}
\begin{proof}
The stabilizer of a vertex preserves its positive definite real form.
The image of each vector of a fixed integral basis of \(\Lambda\)
therefore lies in a finite set of lattice vectors of fixed norm.
There are finitely many possible matrices, so this stabilizer is finite.
An ordered-tuple stabilizer is a subgroup of the stabilizer of its
first vertex.

A finite subgroup of \(\mathrm{GL}_d(\mathbb Z)\) injects under
reduction modulo three. Indeed, otherwise its congruence kernel
would contain an element \(g\) of prime order \(p\).
Write \(g=1+3^sA\), where \(s\geq1\) and some entry of \(A\) is
not divisible by three.
If \(p\ne3\), expansion of \(g^p\) modulo \(3^{s+1}\) gives
\(pA\equiv0\pmod3\), a contradiction.
If \(p=3\), expansion modulo \(3^{s+2}\) gives the same contradiction:
\[
 (1+3^sA)^3-1
      =3^{s+1}A+3^{2s+1}A^2+3^{3s}A^3,
\]
and the last two terms are divisible by \(3^{s+2}\).
The action on the additive lattice embeds the finite vertex stabilizer
in \(\mathrm{GL}_d(\mathbb Z)\).
Its reduction consequently embeds in
\(\mathrm{GL}_d(\mathbb F_3)\), of order at most \(3^{d^2}\).
\end{proof}

\begin{srlemma}\label{lem:sr-integral-action-bound}
Let a group \(G\) act on an exact augmented complex whose degree-\(p\)
basis consists of ordered simplices with the alternating deletion
boundary. Suppose that in degrees \(0,\ldots,m+1\) there are at most
\(B\) orbits and stabilizers of order at most \(J\), with \(J\geq2\).
Then \(H_m(G,\mathbb Z)\) is finitely generated, its torsion is finite,
and
\begin{equation}\label{eq:sr-integral-action-bound}
 \log\#H_m(G,\mathbb Z)_{\rm tors}
                 \leq B\log(m+2)+mBJ^m\log J .
\end{equation}
\end{srlemma}
\begin{proof}
Apply a free group resolution to the augmented complex.
Exactness of the augmentation identifies the total homology with
\(H_*(G,\mathbb Z)\). Taking group homology first and using the
permutation bases gives the first-quadrant spectral sequence
\[
 E^1_{p,q}=\bigoplus_{[\sigma]\in G\backslash X_p}
                    H_q(G_\sigma,\mathbb Z)
                  \ \Longrightarrow\ H_{p+q}(G,\mathbb Z).
\]
There is no orientation-character twist: fixing an ordered basis
simplex fixes each of its entries.
This description also follows directly by writing its permutation
module as a sum of induced trivial modules and using the free
resolution on each summand.

For a finite group \(H\) of order at most \(J\), the integral
bar group in degree \(q\) is free of rank at most \(J^q\).
Its cycles have at most that rank, so its homology is generated by
at most \(J^q\) elements.
For \(q>0\), restriction and transfer to the trivial subgroup show
that \(|H|\) kills \(H_q(H,\mathbb Z)\).
Thus
\[
                  \#H_q(H,\mathbb Z)\leq J^{J^q}\quad(q>0).
\]
Every positive-row term in total degree \(m\) consequently has order
at most \(J^{BJ^q}\), where \(1\leq q\leq m\).
The same bound applies to its eventual subquotient.

The row \(q=0\) is the orbit chain complex, free of rank at most
\(B\) in each relevant degree.
Each column of its boundary from degree \(m+1\) has Euclidean norm
at most \(m+2\).
If the boundary has rank \(r\leq B\), the torsion order of its
cokernel is the greatest common divisor of its nonzero
\(r\)-by-\(r\) minors, by Smith normal form.
Hadamard bounds each such minor by \((m+2)^r\).
If \(r=0\), the cokernel is free.
The row-zero homology is a subgroup of that cokernel, so its torsion
order is at most \((m+2)^B\).

For \(r\geq2\), the differential has bidegree \((-r,r-1)\).
There is no incoming differential to \((m,0)\), since its source
would have second coordinate \(1-r<0\).
Thus \(E^\infty_{m,0}\) is a subgroup of \(E^2_{m,0}\), and its
torsion has the same upper bound.
This direction avoids any assertion that a quotient of a free
abelian group must be torsion-free.

All the relevant terms are finitely generated.
The abutment has a finite filtration with these subquotients.
For an exact sequence of abelian groups, torsion in the middle
maps into torsion of the last group with kernel contained in
torsion of the first; its order is at most the product of the
two torsion orders.
Apply this successively to the filtration.
The positive-row terms contribute at most
\(B\sum_{q=1}^mJ^q\log J\leq mBJ^m\log J\).
Together with the row-zero bound this proves
\eqref{eq:sr-integral-action-bound}, finiteness of torsion and finite
generation.
\end{proof}

\begin{srproposition}\label{prop:sr-all-prime-homology}
For \(a\geq1\), \(m=2a\), \(N=4a+4\), and \(z=2an\),
\begin{equation}\label{eq:sr-homology-final}
 \log\#H_{2a}(\mathrm{GL}_{4a+4}(\mathcal O_F),\mathbb Z)_{\rm tors}
                           \leq\exp(z^{23}\log(2D)).
\end{equation}
\end{srproposition}
\begin{proof}
Use the ordered normal-cone complex, with
\(B=B_{m+1}\), \(J=3^{d^2}\), \(d=nN\).
It is exact by Lemma~\ref{lem:sr-normal-acyclic}.
Its orbits and stabilizers have the bounds just proved.
Write \(L_0=\log(2D)\geq1/2\).
Since \(\log J\leq2d^2\) and \(m\geq1\),
\[
 \log(m+2)+mJ^m\log J
       \leq5md^2e^{2md^2}\leq e^{7md^2}.
\]
Here \(\log(m+2)\leq3m\), and \(5md^2\leq e^{5md^2}\).
Therefore \eqref{eq:sr-integral-action-bound} and
\eqref{eq:sr-orbit-log} give
\[
 \log\#H_m(\Gamma,\mathbb Z)_{\rm tors}
 \leq
 \exp\bigl([22(m+2)Nd^4+14md^2]L_0\bigr)
 \leq \exp\bigl(36(m+2)Nd^4L_0\bigr).
\]
Finally \(m+2\leq4a\), \(N\leq8a\), \(d\leq8an\), so
\[
 36(m+2)Nd^4
       \leq 4{,}718{,}592\,a^6n^4
       \leq (2an)^{23}.
\]
The last inequality uses \(4{,}718{,}592<2^{23}\) and \(a,n\geq1\).
This proves \eqref{eq:sr-homology-final}, including at the small
torsion primes.
\end{proof}


\subsection{Arithmetic-group cohomology and finite levels}
\label{sec:sr-arithmetic-cohomology}
The same complex gives the coefficient-uniform cohomology input used
in Section~\ref{sec:spectrum-deformation}.
This application is over \(\mathbb C\); the earlier integral argument
was needed to retain torsion at the small primes.

Keep \(F,N,n,d,D,I_0,V_0\) as above, without any roots-of-unity
hypothesis, and set
\[
 s=nN^2,\qquad T_0=(V_0+1)^{dN},\qquad
 C_0=I_0^{d^2+d+1}T_0^{T_0+1}.
\]
These integers depend only on \(F,N\).

\begin{srproposition}\label{prop:sr-arithmetic-cohomology}
For every finite-dimensional complex representation \(V\) of
\(\Gamma_0=\mathrm{GL}_N(\mathcal O_F)\),
\begin{equation}\label{eq:sr-arithmetic-cohomology}
 \sum_{j\geq0}\dim_{\mathbb C}H^j(\Gamma_0,V)
                                      \leq C_0\dim_{\mathbb C}V.
\end{equation}
The cohomology vanishes for \(j\geq T_0\).
If \(H\subset\mathrm{GL}_N(F)\) is commensurable with \(\Gamma_0\),
the corresponding sum for any finite-dimensional \(H\)-module \(V\)
is at most
\[
 C_0[\Gamma_0:H\cap\Gamma_0]\dim_{\mathbb C}V.
\]
\end{srproposition}
\begin{proof}
Use the ordinary oriented nerve in the proof of
Lemma~\ref{lem:sr-normal-acyclic}.
A simplex has a common frame by Lemma~\ref{lem:sr-common-frame}.
For a fixed frame, the proof of Proposition~\ref{prop:sr-orbit-bound}
bounds all possible vertices by \(T_0\).
Its vertices are distinct, so every simplex has at most \(T_0\)
vertices. The augmented oriented complex is therefore exact and
zero in degrees \(p\geq T_0\).
The number of its degree-\(p\) simplex orbits is at most
\[
               I_0^{d^2+d+1}T_0^{p+1}.
\]
Indeed an ordering of its vertices is an ordered simplex in the
previous count, and forgetting the order maps orbit sets onto orbit
sets.

A stabilizer of an unordered simplex permutes its finite vertex set.
The kernel fixes every vertex and is contained in a finite vertex
stabilizer. Thus the full stabilizer is finite.
It can act on the oriented generator by its permutation sign.
Over \(\mathbb C\), invariants of a finite group are exact: average
a lift over the group to lift any invariant vector through a
surjection. All its positive cohomology consequently vanishes.
Applying a group resolution to the exact oriented complex therefore
computes \(H^*(\Gamma_0,V)\) by a cochain complex whose degree-\(p\)
term has dimension at most
\[
           I_0^{d^2+d+1}T_0^{p+1}\dim V.
\]
More explicitly each orbit contributes the stabilizer invariants of
\(V\) tensored with the orientation character; their dimension is
at most \(\dim V\).
There are finitely many orbits in each degree.
This is the cohomological version of the resolution argument in
Lemma~\ref{lem:sr-integral-action-bound}.
It also proves the asserted vanishing.

A cohomology group has dimension at most its cochain group.
Since \(T_0\geq2\),
\[
 \sum_{p=0}^{T_0-1}T_0^{p+1}
       =\frac{T_0^{T_0+1}-T_0}{T_0-1}
       \leq T_0^{T_0+1}.
\]
Summing the preceding cochain dimensions proves
\eqref{eq:sr-arithmetic-cohomology}.

For the last assertion put \(J=H\cap\Gamma_0\).
Restriction from \(H\) to \(J\) is injective in every degree:
the composition of restriction and the transfer back to \(H\)
is multiplication by \([H:J]\), a nonzero complex scalar.
Shapiro's lemma identifies \(H^*(J,V)\) with
\(H^*(\Gamma_0,\operatorname{Coind}_J^{\Gamma_0}V)\).
The latter coefficient module has dimension
\([\Gamma_0:J]\dim V\).
Apply \eqref{eq:sr-arithmetic-cohomology} and the injectivity.
No action of \(H\) on the original integral lattice is being assumed.
\end{proof}

We now give the finite-orbit and level factors rather than leaving
them as an arithmetic-group hypothesis.
Let \(\widehat{\mathcal O}_F\) be the product of the finite local
integer rings, put
\[
 G_f=\mathrm{GL}_N(\mathbb A_{F,f}),\quad
 G_r=\mathrm{GL}_N(F),\quad
 U=\mathrm{GL}_N(\widehat{\mathcal O}_F),\quad
 U(q)=\ker\bigl(U\longrightarrow
        \mathrm{GL}_N(\widehat{\mathcal O}_F/q\widehat{\mathcal O}_F)\bigr)
\]
for positive integers \(q\).
In particular \(U(1)=U\).
Put \(M_0=2^nD\) and
\[
                  C_*=C_0 M_0^{n+1}(9M_0)^s.
\]
For a complex representation \(V\) of \(G_r\), the action on
functions is
\((\gamma f)(x)=\gamma\cdot f(\gamma^{-1}x)\).

\begin{srproposition}\label{prop:sr-level-cohomology}
For every positive integer \(M\) and every finite-dimensional
complex \(G_r\)-module \(V\),
\begin{equation}\label{eq:sr-level-cohomology}
 \sum_{j\geq0}\dim_{\mathbb C}
 H^j\bigl(G_r,\operatorname{Map}(G_f/U(3M),V)\bigr)
                                   \leq C_* M^s\dim_{\mathbb C}V.
\end{equation}
The constant is chosen before \(M\) and \(V\).
\end{srproposition}
\begin{proof}
First identify the \(G_r\)-orbits in \(G_f/U\).
For \(g\in G_f\), the module
\[
                P_g=F^N\cap g\widehat{\mathcal O}_F^N
\]
is a full \(\mathcal O_F\)-lattice and its completion is
\(g\widehat{\mathcal O}_F^N\).
Here is the elementary approximation in that assertion.
Choose an integer \(b>0\) with
\[
 b\widehat{\mathcal O}_F^N\subset
 g\widehat{\mathcal O}_F^N\subset
 b^{-1}\widehat{\mathcal O}_F^N.
\]
Only finitely many primes occur in these inclusions.
The Chinese remainder theorem identifies
\(b^{-1}\mathcal O_F^N/b\mathcal O_F^N\) with the product of its
local quotients at those primes.
Taking the inverse image of the given subgroups in that finite
quotient constructs \(P_g\), proves that it lies between the two
global lattices, and proves the completion assertion.
Conversely the completion of any full lattice is locally free over
each discrete valuation ring, hence has the form
\(g\widehat{\mathcal O}_F^N\).
Right multiplication by \(U\) leaves this lattice unchanged.
Two such lattices are in the same rational matrix orbit exactly
when they are isomorphic as \(\mathcal O_F\)-modules: an isomorphism
extends uniquely to their \(F\)-spans.

For completeness, the classification needed here is already
elementary over a Dedekind domain.
Project a nonzero coordinate of a full lattice onto its image,
a nonzero fractional ideal \(I\).
An invertible ideal is projective.
Indeed choose finitely many \(x_i\in I,y_i\in I^{-1}\) with
\(\sum_i x_iy_i=1\).
Given a surjection onto \(I\) and lifts \(e_i\) of \(x_i\), the
map \(z\mapsto\sum_i(y_i z)e_i\) splits it.
The coordinate projection therefore splits.
Its kernel is a lattice of rank one less, so induction expresses
the original lattice as a direct sum of fractional ideals.
Proposition~\ref{prop:sr-ideal-lattices} converts that sum to
\(\mathcal O_F^{N-1}\oplus I\), where \(I\) is their product.
Two such normalized modules are isomorphic precisely when their
determinant ideals have the same ideal class.
The necessity follows by taking the top exterior power.
For sufficiency, multiply the final coordinate by a generator of
the ratio of the ideals.
Thus the orbit set is the ideal class group.

Choose integral ideal representatives \(I_\lambda\) with
\(s_\lambda=\#(\mathcal O_F/I_\lambda)\leq M_0\), using
Minkowski, and write
\(P_\lambda=\mathcal O_F^{N-1}\oplus I_\lambda\).
The ideal count in Lemma~\ref{lem:sr-global-generator} gives
\(h_F\leq M_0^{n+1}\); that argument uses no roots of unity.
Choose finite adelic representatives \(g_\lambda\) completing
these lattices.
The fiber over an orbit in \(G_f/U\) has at most
\[
 [U:U(3)]
       \leq\#M_N(\mathcal O_F/3\mathcal O_F)=3^s
\]
orbits at level three.
The ring has \(3^n\) elements because \(\mathcal O_F\) is free of
rank \(n\) over \(\mathbb Z\).
The stabilizer in \(G_r\) of any such orbit representative is
\[
 H_\lambda=G_r\cap g_\lambda U(3)g_\lambda^{-1}.
\]
It is independent of the representative inside \(g_\lambda U\),
since \(U(3)\) is normal in \(U\).
It consists of the automorphisms of \(P_\lambda\) acting identically
on \(P_\lambda/3P_\lambda\).
The equality follows by the completion assertion and the same
Chinese remainder theorem.

The group \(\Gamma_0(3s_\lambda)\) is contained in
\(H_\lambda\cap\Gamma_0\).
To verify this, the finite group \(\mathcal O_F/I_\lambda\) has
order \(s_\lambda\), so
\(s_\lambda\mathcal O_F^N\subset P_\lambda\subset\mathcal O_F^N\).
If \(\gamma\equiv1\pmod{3s_\lambda}\), then
\[
       (\gamma-1)P_\lambda\subset
                    3s_\lambda\mathcal O_F^N\subset3P_\lambda.
\]
The inverse has the same congruence, which proves preservation in
both directions as well as the level-three condition.
The automorphism group of \(P_\lambda\) is commensurable with
\(\Gamma_0\), as both lattices lie between rational scalar multiples
of one another. Explicitly its level-\(3s_\lambda\) kernel on
\(P_\lambda\) is contained in \(\Gamma_0\), and has finite index;
the preceding inclusion gives the finite index on the other side.
Therefore
\[
 [\Gamma_0:H_\lambda\cap\Gamma_0]
                             \leq(3s_\lambda)^s\leq(3M_0)^s.
\]
Proposition~\ref{prop:sr-arithmetic-cohomology} bounds the total
cohomology of every finite-dimensional \(H_\lambda\)-module \(W\)
by \(C_0(3M_0)^s\dim W\).
This proves the required coefficient-uniform stabilizer estimate.

Finally consider \(G_f/U(3M)\to G_f/U(3)\).
Every fiber has \(D_M=[U(3):U(3M)]\) elements.
In fact \(D_M\leq M^s\): write \(u=1+3A\) with
\(A\in M_N(\widehat{\mathcal O}_F)\).
Distinct cosets in this quotient have distinct images of \(A\)
modulo \(M\), since equality of these images is equivalent to
\(u^{-1}u'\equiv1\pmod{3M}\).
There are \(M^s\) possible images.

Over a \(G_r\)-orbit at level three, the function module on its
inverse image is coinduced from its stabilizer, with coefficient
module the functions on this \(D_M\)-point fiber with values in
\(V\). This module has dimension \(D_M\dim V\).
This assertion follows by restricting a function to the fiber at
the chosen orbit representative; its values on other fibers are
then the coinduction function on rational translates.
The diagonal action stated before the proposition is exactly the
action transported by this restriction.
Shapiro's lemma applies with no finiteness condition on the index
in \(G_r\).
There are at most \(h_F3^s\) orbits, so summing their bounds gives
\[
 C_0 h_F3^s(3M_0)^sD_M\dim V
                   \leq C_*M^s\dim V.
\]
The product decomposition over the orbit set is finite.
This proves \eqref{eq:sr-level-cohomology}.
\end{proof}

For the rank-three spectrum over the fixed imaginary quadratic
field used in the character estimate, \(s=18\).
Thus this proves the coefficient-uniform arithmetic cohomology and
finite-level inputs in Section~\ref{sec:spectrum-deformation},
with constants independent of the weight and level.
The separate infinitesimal-character criterion, cuspidal cohomology
injection, and the realization and rigidity of deformation points
must still be checked on the representations to which they are applied.
A dimension estimate alone does not provide that realization.

\subsection{The number of cohomology generators}
Put \(X=X_{\ell,a}(F)\), already finite by
\eqref{eq:sr-K-comparison}, and
\(r=\dim_{\mathbb F_\ell}(X/\ell X)\).

\begin{srlemma}\label{lem:sr-global-generator}
With \(L_0=\log(2D)\),
\begin{equation}\label{eq:sr-global-generator}
                         r\log\ell\leq6n^2L_0 .
\end{equation}
\end{srlemma}
\begin{proof}
Reduction of continuous cochains gives an injection
\[
 X/\ell X\ \hookrightarrow\
                  H^2(G_{F,S_\ell},\mathbb F_\ell(a+1)).
\]
One can verify its kernel without any finiteness assertion.
If an integral two-cocycle reduces to a coboundary, choose a
continuous lift of the finite-valued one-cochain, subtract its
coboundary and divide by \(\ell\).
The result is still a continuous two-cocycle, since multiplication
by \(\ell\) is injective on the coefficient lattice.
Thus the original class was in \(\ell X\).

The roots-of-unity hypothesis makes the mod-\(\ell\) cyclotomic
character trivial. The reduced coefficient module is isomorphic
to \(\mu_\ell\).
On \(R=\mathcal O_F[1/\ell]\), the
\href{https://stacks.math.columbia.edu/tag/03PK}{Kummer exact sequence}
therefore gives
\[
 0\longrightarrow\operatorname{Pic}(R)/\ell
 \longrightarrow H^2_{\rm et}(R,\mu_\ell)
 \longrightarrow\operatorname{Br}(R)[\ell]\longrightarrow0.
\]
The Picard group is a quotient of the ideal class group.
Brauer--Hasse--Noether embeds the last group into
\(\bigoplus_{v\mid\ell}\operatorname{Br}(F_v)[\ell]\):
the local invariants vanish outside \(S_\ell\), and all infinite
places are complex.
Each local group has order \(\ell\); there are at most \(n\)
such places. Consequently
\[
                         r\log\ell\leq\log h_F+n\log\ell.
\]

We include a crude sufficient class-number bound.
Minkowski gives an integral representative in every ideal class
with norm at most \(2^nD\).
The number of ideals of norm \(k\) is at most the coefficient
\(d_n(k)\) of \(\zeta(s)^n\).
To check this locally, for primes above a rational prime \(p\)
with residue degrees \(f_i\), send an ideal exponent vector
\((e_i)\), \(\sum f_ie_i=b\), to an \(n\)-tuple of nonnegative
integers by repeating \(e_i\) in \(f_i\) slots and filling the
remaining slots with zeros. This is injective, because
\(\sum f_i\leq n\).
Multiplying over rational primes gives the assertion.
Also \(d_n(k)\leq k^n\), by counting ordered \(n\)-tuples of positive
factors of \(k\).
Summing up to the integer \(2^nD\) proves
\[
 h_F\leq(2^nD)^{n+1},\qquad
 \log h_F\leq(n+1)^2L_0\leq4n^2L_0.
\]
For odd \(\ell\), \(\mathbb Q(\mu_\ell)\subset F\) implies
\(\ell-1\leq n\). For two, \(\ell\leq n+1\) is automatic.
Hence
\(n\log\ell\leq n\log(n+1)\leq n^2\leq2n^2L_0\).
Combining these inequalities proves
\eqref{eq:sr-global-generator}.
\end{proof}

\subsection{Completion of the global bound}
\begin{proof}[Proof of Theorem~\ref{thm:sr-global-high-twist}]
Put \(m=2a,N=4a+4,z=2an,L_0=\log(2D)\).
The comparison \eqref{eq:sr-K-comparison}, followed by Hurewicz
and homology stability, gives a map
\[
 X\longrightarrow
 H_{2a}(\mathrm{GL}_{4a+4}(\mathcal O_F),\mathbb Z)
                                      \otimes\mathbb Z_\ell .
\]
Flatness of \(\mathbb Z_\ell\) over \(\mathbb Z\) preserves the
kernel description. Mathew's bound makes this kernel a subgroup
of \(X[\ell^{b_\ell}]\), with \(b_\ell\leq4a\).
A finite \(\ell\)-group with \(r\) cyclic factors has at most
\(\ell^{br}\) elements killed by \(\ell^b\), as follows immediately
by writing it as a direct sum of cyclic groups.
The kernel therefore has logarithmic order at most
\(4ar\log\ell\).

The image is finite and lies in the torsion of the target.
By Proposition~\ref{prop:sr-all-prime-homology} its logarithmic
order is at most \(\exp(z^{23}L_0)\).
Lemma~\ref{lem:sr-global-generator} consequently gives
\[
 \log\#X\leq24an^2L_0+\exp(z^{23}L_0).
\]
As \(a,n\geq1\), \(24an^2\leq z^5\leq z^{23}\).
For \(t=z^{23}L_0\geq1\),
\[
 \log\#X\leq t+e^t\leq e^{2t}
                    \leq\exp(z^{24}L_0),
\]
using \(z\geq2\). This is \eqref{eq:sr-global-high-twist}.
Every bound precedes the field, the coefficient prime and the twist.
\end{proof}

\begin{srproposition}\label{prop:sr-K2-proved}
The specialized assertion \(K_2(24)\) of
Definition~\ref{def:sr-K2} holds.
The dyadic common-window theorem
\ref{thm:sr-dyadic-window} is consequently unconditional with
\[
 C_K=24,\qquad
 \eta=\epsilon/48,\qquad
 J=2^{26}(19248/\epsilon)^{24}.
\]
\end{srproposition}
\begin{proof}
The field \(L_7\) is totally imaginary, contains \(\mu_4\), and
has degree 1024. Apply Theorem~\ref{thm:sr-global-high-twist}
with \(\ell=2\), \(F=L_7\), and \(2an=2048a\).
This gives exactly the finiteness and inequality in
Definition~\ref{def:sr-K2}.
The remaining conclusions follow by substitution into
Theorem~\ref{thm:sr-dyadic-window}.
\end{proof}

The \(H^2\) group here is also the one in the archive's
\srcid{LemmaK}: its ramification set is precisely the primes above
\(\ell\), and its compact coefficient is the cyclotomic
\((a+1)\)-st twist.
Changing the algebraic closure identifies these Galois groups;
the factorization of the cyclotomic character through the unramified
quotient is unique. The local embedding choices do not enter this
global \(H^2\).
Thus this bound applies uniformly to the allowed localization data.
The character-counting bounds and connected geometric supply remain
separate tasks. This new prose proof has not been separately
formalized in Lean.

\section{The latest reduction of the plus height to two character estimates}
\label{sec:sr-two-counts}
Anton's subsequent
\href{https://github.com/antoshashakov/riemann-hypothesis/blob/76ffabc62ab67b126afc128013ab6efa6f7d2a7c/verification/TwoCharacterCounts.lean}{\srcid{TwoCharacterCounts}}
reduces the polynomial plus-height input to two rational-character
bounds. We give the precise estimates and their assembly here.
Neither estimate is asserted unconditionally in this section.
We first supply the trivial-character vanishing needed in that assembly.

\subsection{Two vanishing arguments over \(\mathbb Q\)}
\begin{srproposition}\label{prop:sr-rational-regular}
For \(\ell=3,7\) and every \(a\geq1\),
\(T_{\ell,a}(\mathbb Q)=0\).
\end{srproposition}
\begin{proof}
Restriction takes this strict group into
\(T_{\ell,a}(\mathbb Q(\zeta_\ell))\), which is zero by
Proposition~\ref{prop:sr-regularsha}.
The extension has degree \(\ell-1\), a unit in \(\mathbb Z_\ell\).
Restriction is injective: restriction followed by corestriction is
multiplication by \(\ell-1\), or equivalently the finite-quotient
cohomology kernel vanishes by averaging over that quotient.
The same argument applies to the local restriction conditions.
Thus the original group is zero.
\end{proof}

\begin{srproposition}[The resonant trivial character]\label{prop:sr-rational-resonant}
For every odd prime \(\ell\) and every positive \(a\) with
\(\ell-1\mid2a\), one has \(T_{\ell,a}(\mathbb Q)=0\).
\end{srproposition}
\begin{proof}
Let \(\omega:\operatorname{Gal}(\mathbb Q(\mu_{\ell^\infty})/\mathbb Q)
\to\mu_{\ell-1}\subset\mathbb Z_\ell^\times\) be the finite
Teichmuller part of the cyclotomic character.
The condition implies that \(\omega^a\) has order at most two.
Its fixed field \(K\) in \(\mathbb Q(\zeta_\ell)\) is either
\(\mathbb Q\) or the unique quadratic subfield.
In the quadratic case its discriminant is
\((-1)^{(\ell-1)/2}\ell\).
Indeed the Gauss-sum calculation in Section~\ref{sec:sr-dyadic-dual}
works for every odd prime: it gives
\(\tau_\ell^2=(-1)^{(\ell-1)/2}\ell\), and this squarefree radicand
is \(1\bmod4\).
The integral quadratic basis \(1,(1+\sqrt D)/2\) has discriminant
\(D\); the trace and norm integrality conditions give the usual
full ring of integers for squarefree \(D\equiv1\bmod4\).

In both cases \(\ell\nmid h_K\). For the quadratic case we include
an elementary bound. Minkowski gives an integral ideal in each class
of norm at most \(B\leq\sqrt\ell\).
There are at most \(d(n)\) ideals of norm \(n\): at a split rational
prime the count is the exponent plus one, and the inert and ramified
counts are smaller. Here \(d(n)\) is the divisor function.
For \(B\geq1\),
\[
 h_K\leq\sum_{n\leq B}d(n)
 =\sum_{u\leq B}\lfloor B/u\rfloor
 \leq B(1+\log B)
 \leq\sqrt\ell(1+\log\sqrt\ell)<\ell .
\]
The last inequality is \(\log z<z-1\) for \(z>1\).
For \(K=\mathbb Q\) the class number is one.

Let \(K_\infty\) be the subfield of
\(\mathbb Q(\mu_{\ell^\infty})\) fixed by
\(\ker(\omega^a)\) in its finite factor \(\mu_{\ell-1}\).
Then \(K_\infty/K\) is a \(\mathbb Z_\ell\)-extension, and
\(\chi_{\rm cyc}^{-a}\) is trivial on \(G_{K_\infty}\).
There is a unique prime above \(\ell\) in \(K\), totally ramified
through this tower, and no other finite prime ramifies.
The one-ramified-prime class-number lemma also applies to these
cyclic degree-\(\ell\) steps if \(K\) is real.
For odd \(\ell\) no real place can complexify in such a Galois
extension. In its proof the unit representation is then
\(\mathbb Q[G]^{r_1(K)+r_2(K)}-\mathbb Q\);
the same Herbrand quotient \(1/\ell\) and ideal calculation apply.
Thus \(\ell\) is prime to every finite-layer class number.

The strict descent of Proposition~\ref{prop:sr-regularsha} now applies
to this tower, as can also be seen directly.
A strict cocycle has finite image. In a sufficiently high layer the
coefficient action fixes that image and the chosen local primitive
at the unique prime above \(\ell\).
Its restriction is a homomorphism, zero at that decomposition group
and unramified everywhere else.
Artin reciprocity and the prime-to-\(\ell\) class number kill it.
The cocycle therefore descends to
\(\operatorname{Gal}(K_\infty/K)\simeq\mathbb Z_\ell\).
The generator acts on \(\mathbb Q_\ell/\mathbb Z_\ell\) by a scalar
different from one, and multiplication by the difference is
surjective. Subtracting a coboundary kills its value at that
generator, hence kills the cocycle by continuity.
We obtain \(T_{\ell,a}(K)=0\).
Finally \([K:\mathbb Q]\leq2\) is prime to \(\ell\), so the
restriction argument of Proposition~\ref{prop:sr-rational-regular}
gives the result over \(\mathbb Q\).
\end{proof}

\subsection{The rational groups and their normalizations}
For an odd prime \(\ell\), choose a finite splitting extension of
\(\mathbb Q_\ell\), with ring \(\mathcal O\), fraction field \(Q\),
and \(r_{\mathcal O}=v_{\mathcal O}(\ell)>0\).
For a Dirichlet character \(\psi\) of level \(N\), let
\(\mathcal V'_{\ell,N,a}(\psi)\) be the subgroup of
\[
 H^1(\mathbb Q,(Q/\mathcal O)(\psi^{-1}\chi_{\rm cyc}^{-a}))
\]
whose restrictions to all finite inertia groups and to the full
local groups at primes dividing \(\ell N\) are zero.
These are restrictions of cohomology classes, not assertions that
every representative cocycle is pointwise zero.
Put
\[
 v_{\ell,a}(\psi)=
 \frac{\operatorname{length}_{\mathcal O}
              \mathcal V'_{\ell,N,a}(\psi)}{r_{\mathcal O}}
                  \quad\hbox{in }[0,\infty].
\]
This is the normalized length called \srcid{Vn} when its finiteness
has been proved; \srcid{VLe} includes that finiteness as well as
the weighted inequality. An infinite length is never set to zero.

For a principal host let \(H\subseteq(\mathbb Z/N)^\times\)
be its fixing subgroup, and let
\[
 \mathcal X_a=\{\psi:\psi|_H=1,\ \psi(-1)=(-1)^a\}.
\]
All these characters have \(2\)-power order and
\(|\mathcal X_a|\leq[E:\mathbb Q]\leq N\).
The plus parity here is different from the minus parity in
Proposition~\ref{prop:sr-PB}.

The restriction and averaging proof in that proposition, applied
to the other parity, gives
\begin{equation}\label{eq:sr-plus-character}
 s^+_{\ell,E}(a)\leq
             \sum_{\psi\in\mathcal X_a}v_{\ell,a}(\psi).
\end{equation}
There is no index loss: every odd \(\ell\) is prime to
\([E:\mathbb Q]\).
We check the extra local conditions used in \(\mathcal V'\).
At \(p\ne\ell\), if \(\psi\) is ramified, a nontrivial inertia
value differs from one by an \(\mathcal O\)-unit, because its
order is prime to \(\ell\).
The argument in Section~\ref{sec:minus-characterwise} turns an
inertia coboundary into a full local coboundary.
If \(\psi\) is unramified, an unramified local cocycle factors
through the procyclic unramified quotient.
Its Frobenius scalar is \(\psi(\operatorname{Frob}_p)^{-1}p^{-a}\ne1\):
a positive rational power \(p^a>1\) cannot be a root of unity.
Subtract a coboundary using surjectivity of this nonzero scalar
minus one on \(Q/\mathcal O\). The cocycle is then zero at
Frobenius, hence zero by continuity.
Thus every relevant unramified class is already locally zero
at these additional primes. This proves that imposing the
\(N\)-prime conditions introduces no missing local term in
\eqref{eq:sr-plus-character}.

For \(\psi=1\), the rational vanishing propositions imply
\[
 v_{\ell,a}(1)=0\quad
 \hbox{if }\ell\in\{3,7\}\ \hbox{or }\ell-1\mid2a .
\]
The statements remain true over \(\mathcal O\), since its scalar
extension is finite free over \(\mathbb Z_\ell\), and continuous
cohomology and the local kernels commute with finite direct sums.
Adding the further local conditions only takes a subgroup.

\subsection{The two remaining estimates}
Write \(M=Nt\), with \(t\geq2\).
All the following quantifiers concern the same sufficiently large
principal hosts as \eqref{eq:sr-principal-height}.
The fixed thresholds and the exponents precede the host and \(t\).
The splitting coefficient ring may vary with \(\ell\).

\begin{srdefinition}\label{def:sr-character-inputs}
The generic estimate asks for one natural exponent \(A\) such that,
for each \(1\leq a<t\) and every finite prime set \(S\) with
\(\ell\notin\{2,3,7\}\) and \(\ell-1\nmid2a\),
\begin{equation}\label{eq:sr-generic-count}
 \sum_{\ell\in S}(\log\ell)
       \sum_{\psi\in\mathcal X_a}v_{\ell,a}(\psi)\leq M^A .
\end{equation}
The nontrivial estimate asks for one natural exponent \(B\) such that
\begin{equation}\label{eq:sr-nontrivial-count}
 (\log\ell)v_{\ell,a}(\psi)\leq(\ell N a)^B
\end{equation}
whenever \(\psi\in\mathcal X_a,\ \psi\ne1,\ \ell\) is odd, and
\[
 \ell\in\{3,7\}\quad\hbox{or}\quad
                       (a\geq2\ \hbox{and }\ell-1\mid2a).
\]
Finiteness is included in both estimates.
\end{srdefinition}
The first is a simultaneous sum over arbitrarily large finite
prime sets, rather than a pointwise bound which could not be
summed over all primes. These two estimates are still assumptions
of the next proposition.

\begin{srproposition}\label{prop:sr-two-counts}
Definition~\ref{def:sr-character-inputs} implies the polynomial
plus-height bound \eqref{eq:sr-principal-height}.
One may use the uniform exponent \(A+3B+5\).
At the exceptional pair \(3,7\) alone, exponent \(2B+2\) suffices.
\end{srproposition}
\begin{proof}
All sums below may first be taken over an arbitrary fixed finite
prime set. The inputs and the vanishing results make each rational
length finite; \eqref{eq:sr-plus-character} then proves finiteness
of each plus length.

For \(\ell\notin\{2,3,7\}\) and a fixed \(a\), the generic primes
contribute at most \(M^A\).
Every other prime satisfies \(\ell-1\mid2a\), so
\(\ell\leq2a+1\leq2t\), and necessarily \(a\geq2\).
There are at most \(2t\) such primes.
The trivial character contributes zero.
For every other character,
\[
 \ell N a\leq2Nt^2=2Mt\leq M^3,
\]
since \(t\leq M\) and \(M\geq2\).
Its weighted contribution is therefore at most \(M^{3B}\).
There are at most \(N\) characters.
Summing the fixed-\(a\) estimate over \(a<t\), and using \(tN=M\),
bounds the entire outside contribution by
\[
 t(M^A+2tN M^{3B})
 \leq M^{A+1}+2M^{3B+2}
 \leq3M^{A+3B+2}\leq M^{A+3B+4}.
\]
All exponents are nonnegative; the last inequality uses \(M^2\geq4\).

At \(\ell=3,7\), the trivial character again contributes zero for
every positive \(a\).
Here \(M\geq7\) for the principal hosts, so
\(\ell N a\leq M^2\).
Each remaining weighted contribution is at most \(M^{2B}\).
The two primes, at most \(N\) characters and fewer than \(t\)
twists give at most
\[
          2tN M^{2B}=2M^{2B+1}\leq M^{2B+2}.
\]
Add the inside and outside estimates.
Since \(A+3B+4\geq2B+2\) and \(M\geq2\), their sum is at most
\(M^{A+3B+5}\).
Equation~\eqref{eq:sr-plus-character} proves the asserted plus bound.
The constants are independent of the chosen finite prime set,
host, splitting rings and twist range, as required.
\end{proof}

In the source's notation the exceptional set at conductor one is
\(B_\chi(1)=\{3,7\}\).
Its quadratic auxiliary field is \(\mathbb Q(\sqrt{-7})\), whose
Minkowski bound \((2/\pi)\sqrt7<2\) gives class number one.
The first eligible quadratic auxiliary prime is seven and the
other auxiliary prime is three; the base character has order one.
The odd prime factors in the defining exceptional set therefore
reduce to those of \(7\cdot3\), together with three itself.
An odd-order nontrivial host character cannot occur, because its
order divides a power of two.
This is the elementary content of the exceptional-set and
Case III simplifications in the contributors' files.
The proof above avoids retaining the case labels: the generic,
resonant and exceptional conditions are stated directly.

The arithmetic content not supplied here is
\eqref{eq:sr-generic-count} and~\eqref{eq:sr-nontrivial-count}.
Their source proofs involve the global counting and geometric
comparisons in Part II. The vanishing and the assembly above
are not substitutes for those two estimates.


\section{Integer bounds for the principal character packet}
\label{sec:sr-character-integers}
The two target estimates of Definition~\ref{def:sr-character-inputs}
concern the same principal host and the same normalized strict lengths.
This section proves their numerical and character-orbit reductions.
The later support arguments make several underlying implications
explicit; the automorphic producers still required are identified at
each use.

\subsection{A single Frobenius prime for the generic character sum}
\label{sec:sr-fixed-seven-generic}
Throughout this subsection the principal host is the field defined in
Section~\ref{sec:sr-arithmetic}, with its fixing subgroup $H_{\rm host}$.
Only its abelian $2$-group character packet, its degree bound, and its
local quadratic characters at the two odd principal primes are used.
We increase the fixed lower threshold for the principal primes to ensure
$q,d>7$. This does not change the order of the uniform quantifiers.

\begin{srproposition}[The exact orbit-integer implication]
\label{prop:sr-orbit-integer-implication}
Fix an imaginary quadratic field $K$. Assume the common trace-integer
construction and the one-character deformation inequality of
Sections~\ref{sec:common-trace-integer}
and~\ref{sec:single-character}, including their stated arithmetic inputs.
There are $c,C\in\mathbb N$, depending only on $K$, with the following
property. For every $a\ge1$, positive modulus $N$, subgroup
$J\le(\mathbb Z/N)^\times$, and prime
$p\nmid30N|\operatorname{disc}K|$, there is a nonzero integer
$Z=Z_{K,J,a,N,p}$ such that
\begin{equation}\label{eq:sr-orbit-integer-size}
 \log|Z|\le c\,\varphi([(\mathbb Z/N)^\times:J])
                    [N(a+1)]^C\log p.
\end{equation}
The integer is chosen before the coefficient prime and coefficient ring.
For every prime $\ell\ne p$ and finite splitting coefficient ring
$\mathcal O$ at $\ell$, if one character of kernel $J$ satisfies the
good-data conditions in Section~\ref{sec:single-character}, all such
characters have finite strict groups and
\begin{equation}\label{eq:sr-orbit-integer-valuations}
       \sum_{\ker\psi=J}v_{\ell,a}(\psi)
                              \le2\operatorname{ord}_\ell Z.
\end{equation}
\end{srproposition}
\begin{proof}
Choose $w\mid p$ and put $Q=Nw$. The complex orbit
$\mathcal O_J$ consists of all characters modulo $N$ with kernel $J$.
Its indexed roots $z_\chi=\chi(Q)^{-1}$ are permuted by every
automorphism of $\mathbb C$, since an injective field homomorphism
preserves a character kernel. If the orbit is nonempty, its quotient
$(\mathbb Z/N)^\times/J$ is cyclic; its faithful characters form one
orbit of cardinality $\varphi([(\mathbb Z/N)^\times:J])$.
All of them are powers $\chi^u$ with $u$ prime to the order.
If the orbit is empty, take $Z=1$.

The trace-integer construction applied to this indexed family gives
a nonzero integer $Z$ satisfying \eqref{eq:sr-orbit-integer-size}.
For every $\ell\ne p$ it controls the product of all the nonzero
differences between admissible traces at inverse Frobenius and the
corresponding Eisenstein traces, for every finite list having pairwise
nonisomorphic members within each orbit position.

Fix such an $\ell$ and $\mathcal O$, and choose an embedding of its
fraction field in $\overline{\mathbb Q}_\ell$ and a field isomorphism
$\iota:\overline{\mathbb Q}_\ell\to\mathbb C$.
Composition preserves kernels and identifies the coefficient-valued
orbit with $\mathcal O_J$. This identification is onto: every character
with a given kernel is a power, prime to the order, of one character
already in the coefficient ring. Good data transfer to every character
of this orbit. This is checked by the same powers: the kernels coincide,
the quadratic local component is preserved, and all the remaining
conditions depend only on the kernel, the order, or the prescribed parity.

Arithmetic Frobenius at $w$ has cyclotomic value $Q$ and Dirichlet
value $\psi(Q)$. Hence its inverse gives the reducible trace
\[
 T_\psi=Q^{-a}\psi(Q)^{-1}+1+Q^a\psi(Q).
\]
The separation supplied by the trace-integer construction permits the
one-character inequality at this inverse Frobenius. It proves finiteness
and supplies a finite set $X_\psi$ of pairwise nonisomorphic admissible
representations with
\[
 v_{\ell,a}(\psi)\le
       2\sum_{\rho\in X_\psi}
          v_\ell\bigl(\operatorname{tr}\rho(\operatorname{Frob}_w^{-1})
                                            -T_\psi\bigr).
\]
Here $v_\ell(\ell)=1$. The product inequality for $Z$ says
\[
 |Z|_\ell\le
       \prod_{\psi}\prod_{\rho\in X_\psi}
       \left|\operatorname{tr}\rho(\operatorname{Frob}_w^{-1})
                                            -T_\psi\right|_\ell .
\]
All factors are nonzero. Taking $-\log_\ell$ bounds the sum of their
valuations by $\operatorname{ord}_\ell Z$. Summing the preceding
one-character inequalities proves
\eqref{eq:sr-orbit-integer-valuations}.
No distinctness is required between different orbit positions: the
integer product retains their separate multiplicities.
The coefficient normalization is exactly
$\operatorname{length}_{\mathcal O}/v_{\mathcal O}(\ell)$, as in
Section~\ref{sec:single-character}.
\end{proof}

\begin{srproposition}[The fixed-seven specialization]
\label{prop:sr-fixed-seven-generic}
Under the arithmetic inputs of
Proposition~\ref{prop:sr-orbit-integer-implication} for the single field
$K=\mathbb Q(\sqrt{-3})$, the generic estimate
\eqref{eq:sr-generic-count} holds with one natural exponent fixed
before the host, twist range, coefficient rings and finite prime set.
More precisely, if $c,C$ are the constants above, choose
$k\in\mathbb N$ once and for all with
\[
                2^k\ge\max(1,2c\log7).
\]
Then one may use $A=C+1+k$.
\end{srproposition}
\begin{proof}
Fix $t\ge2$, a principal host with $q,d>7$, and $1\le a<t$.
Let $\mathcal J_a$ be the set of kernels of its complex characters
in $\mathcal X_a$. The entire kernel-$J$ orbit is contained in
$\mathcal X_a$ for every $J\in\mathcal J_a$.
Indeed, a character of the same kernel is $\psi^u$ for $u$ prime to
$r=\operatorname{ord}\psi$. It is trivial on $H_{\rm host}$.
If $r>1$ it is a power of two, so $u$ is odd and
$\psi^u(-1)=\psi(-1)$; if $r=1$ the orbit is a singleton.
Consequently these orbits partition $\mathcal X_a$, and
\begin{equation}\label{eq:sr-orbit-totient-sum}
 \sum_{J\in\mathcal J_a}
           \varphi([(\mathbb Z/N)^\times:J])
             =|\mathcal X_a|\le[E:\mathbb Q]\le N.
\end{equation}
This equality concerns character kernels in the full group of units
modulo $N$, with their common host fixing subgroup retained.

Put $p=7$. The choice $q,d>7$ gives $7\nmid N$, while
$|\operatorname{disc}K|=3$, so $7\nmid30N|\operatorname{disc}K|$.
Proposition~\ref{prop:sr-three-auxiliary} verifies the good-data
conditions for every odd nonresonant coefficient prime and every
character in $\mathcal X_a$. For the source's character $\psi'$ with
the $\ell$-local component removed, the final kernel condition follows
from inertia at three as well: its conductor still divides $N$.
Use Proposition~\ref{prop:sr-orbit-integer-implication} once for
each $J\in\mathcal J_a$, obtaining $Z_J=Z_{K,J,a,N,7}$.
Every $\ell$ in the prescribed generic set $S$ differs from seven.
Thus all the following strict lengths are finite and
\begin{align*}
 \sum_{\ell\in S}\log\ell
                   \sum_{\psi\in\mathcal X_a}v_{\ell,a}(\psi)
 &\le2\sum_{J\in\mathcal J_a}
                     \sum_{\ell\in S}
                      (\log\ell)\operatorname{ord}_\ell Z_J\\
 &\le2\sum_{J\in\mathcal J_a}\log|Z_J|\\
 &\le2c\log7\,N[N(a+1)]^C .
\end{align*}
The middle inequality is prime factorization of the nonzero integers
$Z_J$; it does not contain $|S|$.
Put $M=Nt$. Since $a<t$, we have $N(a+1)\le M$, $N\le M$,
and $M\ge2$. Our fixed choice of $k$ gives
\[
        2c\log7\,N[N(a+1)]^C
                         \le M^{C+1+k}.
\]
This proves the asserted uniform quantifier order.
\end{proof}

\paragraph{Source verification and remaining scope.}
The exact proposition $\srcid{PropGK}$ is defined in
\href{https://github.com/ChristianKudeba/alpoge-rh-verification/blob/alpoge-source-7cbf7a0c/src/Development/AutomorphicArithmetic/AutomorphicVarieties/UnitaryThreeGroups/KatoUse.lean}
{\srcid{UnitaryThreeGroups.KatoUse}}.
Its producer is
\href{https://github.com/ChristianKudeba/alpoge-rh-verification/blob/alpoge-source-7cbf7a0c/src/Development/AutomorphicArithmetic/AutomorphicVarieties/ModularityArgument/ThresholdArguments/DenseThresholdCases/FPropGKU.lean}
{\srcid{FPropGKU.propGK_of_numberU3U_of_oneCharU}}.
The two arguments above prove the orbit assembly and the entire
coefficient-prime summation. Their stated arithmetic input still
includes the unramified automorphy used by the one-character inequality.
In the archive its producer runs through
\srcid{DnEndsTop.ends0_alt}, \srcid{LinkedEndE},
\srcid{PrLinkedEnd}, \srcid{AbsMain}, and the characteristic-$\ell$
component-propagation theorem \srcid{PropagationTheoremGL3}.
The fixed-seven specialization does not replace that producer.

\subsection{Arithmetic absorption in the principal exceptional packet}
\label{sec:sr-exceptional-absorption}

Throughout this subsection the principal host is the actual fixed field
used in Section~\ref{sec:sr-arithmetic}:
\[
 E=\mathbb Q\bigl(\zeta_{2^{h+2}}+\zeta_{2^{h+2}}^{-1},
                   \sqrt{-q},\sqrt{-d}\bigr),\qquad N=2^{h+2}qd.
\]
Equivalently, its fixing subgroup at the dyadic factor is generated by
\(-1\). We use precisely its character packet \(\mathcal X_a\).
In particular the dyadic component of an admitted character is even.
Enlarging this field to the full dyadic cyclotomic compositum would
enlarge that packet and is not a step in the following argument.
The primes \(q,d\) satisfy the principal congruences, exceed \(2t+1\),
and exceed any fixed thresholds stated below.

For clarity, the assertions below concerning a deformation comparison
are implications from explicitly stated comparisons. They establish the
uniform scalar absorptions, not the construction of the deformation
points or their automorphic realization.

\begin{srproposition}\label{prop:sr-principal-case-II-arithmetic}
Let \(1\leq a<t\), let \(\psi\in\mathcal X_a\) be nontrivial, and suppose
\[
 \ell\in\{3,7\}\quad\hbox{or}\quad
                    (a\geq2,\ \ell-1\mid2a).
\]
Assume \(q,d>7\). Then \(\ell\nmid N\), the order of \(\psi\) is a
nontrivial power of two, and its local components at odd primes have
order at most two. Consequently, when \(a\geq2\), this is Case II
of the source; neither Case III nor Case IV is needed for this packet.

Write \(c=c_2(\psi)\) for the exponent in the conductor \(2^c\) of its
dyadic component. Remove that component from \(\psi\), and call the
order of the remaining character \(b\). Then \(b\mid2\), and
\[
 f=\operatorname{ord}_{\operatorname{lcm}(4,b)}(\ell)
       =\operatorname{ord}_4(\ell)\leq2,\qquad
 w=v_2(\ell^f-1),\qquad 4^w<\ell^4 .
\]
\end{srproposition}
\begin{proof}
The odd ramification of \(E/\mathbb Q\) comes from the two quadratic
radicals, so the components at \(q,d\) have exponent two; there are no
other odd ramified components. The dyadic quotient is a \(2\)-group.
Thus every host character has \(2\)-power order, and its order is even
when it is nontrivial.

For an exceptional prime, \(\ell<q,d\) follows from \(q,d>7\).
For a resonant prime, \(\ell\leq2a+1<q,d\).
Since \(\ell\) is odd, in either case \(\ell\nmid2^{h+2}qd=N\).
The \(\ell\)-power-order part of \(\psi\) and its local component at
\(\ell\) are therefore both trivial. Hence the character denoted
\(\mathrm{chiP}\) in the source is exactly \(\psi\), and its order is
even. Together with \(a\geq2\) and the displayed exceptional or resonant
alternative, these are exactly the three conditions defining Case II.
Case III requires a nontrivial character of odd order, which cannot
occur here. Case IV also requires odd order after taking \(\mathrm{chiP}\).

Removing the dyadic component leaves the product of the two quadratic
components, whose order divides two. Thus the modulus defining \(f\)
is four. Finally \(2^w\leq\ell^f-1<\ell^2\), which gives \(4^w<\ell^4\).
\end{proof}

\begin{srproposition}[The small-conductor twist-one branch]
\label{prop:sr-small-conductor-exceptional}
At \(a=1\), the threshold
\[
                    R_1(\ell,\psi)=16\ell^2\,4^w
\]
equals \(9216\) for \(\ell=3\) and \(200704\) for \(\ell=7\).
If \(2^{c_2(\psi)}<R_1(\ell,\psi)\), then
\[
 \operatorname{ord}\psi\leq
 \begin{cases}2^{11},&\ell=3,\\2^{15},&\ell=7.\end{cases}
\]
Consequently this entire branch satisfies
\[
 (\log\ell)v_{\ell,1}(\psi)\leq(\ell N)^{B_{\mathrm{small}}},
 \qquad
 B_{\mathrm{small}}=196609\cdot393216^{24}.
\]
The exponent is fixed before the host, character and splitting ring.
\end{srproposition}
\begin{proof}
Both three and seven are \(3\bmod4\), so \(f=2\).
The valuations \(v_2(3^2-1)=3\) and \(v_2(7^2-1)=4\) give the displayed
thresholds. Since
\[
 2^{13}<9216<2^{14},\qquad
 2^{17}<200704<2^{18},
\]
the strict threshold bounds \(c_2(\psi)\) by thirteen or seventeen.

For \(c\geq3\), the exponent of \((\mathbb Z/2^c)^\times\) is
\(2^{c-2}\). Indeed the units are generated by \(-1\) and \(5\), and
\(v_2(5^{2^j}-1)=j+2\), by induction using the factorization
\(x^2-1=(x-1)(x+1)\). For \(c=0,1,2\), the unit-group exponent is at
most two. The odd local components also have exponent at most two.
The order of their product is the least common multiple of these
\(2\)-power orders, proving the two order bounds.

Apply Proposition~\ref{prop:sr-bounded-order-one} with \(R=2^{15}=32768\).
Its exponent \((6R+1)(12R)^{24}\) is the displayed integer.
This also proves finiteness in this branch.
\end{proof}

\paragraph{Normalization under a change of splitting ring.}
If \(\mathcal O\subset\mathcal O'\) is a finite extension of splitting
discrete valuation rings and \(\psi'\) is the same character in
\(\mathcal O'\), then
\[
                         v_{\ell,a}(\psi')=v_{\ell,a}(\psi).
\]
Indeed \(\mathcal O'\) is finite free over \(\mathcal O\), and
\[
 (Q/\mathcal O)\otimes_{\mathcal O}\mathcal O'
                       \simeq Q'/\mathcal O'.
\]
Choosing an \(\mathcal O\)-basis identifies the continuous cochain
complex, and every local restriction complex, with a finite direct
sum of the old complex. Thus cohomology and the simultaneous local
kernel commute with this scalar extension. If
\(e'=v_{\mathcal O'}(\varpi_{\mathcal O})\), the cyclic decomposition
of a finite torsion module \(T\) gives
\[
 \operatorname{length}_{\mathcal O'}(T\otimes\mathcal O')
       =e'\operatorname{length}_{\mathcal O}T,\qquad
 v_{\mathcal O'}(\ell)=e'v_{\mathcal O}(\ell).
\]
Dividing gives the asserted identity. Finiteness descends by faithful
flatness as well as ascending by scalar extension. In particular,
adjoining the auxiliary dyadic roots of unity creates no factor in
the normalized estimate.

\begin{srproposition}[Uniform absorption of the rank-four comparison]
\label{prop:sr-rank-four-absorption}
Fix natural constants \(C,A\) in the rank-four common-integer
comparison, before the principal host is chosen.
For each Case II instance with \(a\geq2\), set
\[
 u=\max\{\ell N4^w,\ 2^{c_2(\psi)+6}\ell^3\},\quad
 n=\lceil\log_2u\rceil,\quad
 M=\ell\,2^nN,\quad k=2(\ell-1)2^n.
\]
Suppose the following comparisons have been constructed for these
same data:
\begin{enumerate}
\item At an eligible prime \(p\geq5\), a finite pairwise nonisomorphic
family of admissible rank-four representations \(\mathcal R\), with
nonzero trace differences \(d_\rho\), satisfies
\[
 (\log\ell)v_{\ell,a}(\psi)
                  \leq2\sum_{\rho\in\mathcal R}-\log|d_\rho|_\ell .
\]
\item A nonzero integer \(J\), for the common root-of-unity order \(k\),
satisfies
\[
 |J|_\ell\leq\prod_{\rho\in\mathcal R}|d_\rho|_\ell,\qquad
 \log|J|\leq C\varphi(k)\{M(a+1)\}^{A}\log p.
\]
\end{enumerate}
These are the normalizations in
\(\mathrm{Raw.Ineq}\), \(\mathrm{NumberSp4Alt}\), and
\(\mathrm{CountS.exists\_frob\_nonvanishing\_with\_sha\_bound}\).
There is no additional factor \(a+1\) outside the indicated \(A\)-th
power.

Choose once and for all \(b_4\in\mathbb N\) with \(2C\leq2^{b_4}\).
If the two comparisons hold at every eligible prime, they imply
\[
 (\log\ell)v_{\ell,a}(\psi)\leq(\ell Na)^{B_4},\qquad
                         B_4=10A+12+b_4 .
\]
\end{srproposition}
\begin{proof}
Put \(x=\ell Na\). The conductor \(2^{c_2(\psi)}\) divides \(N\).
By Proposition~\ref{prop:sr-principal-case-II-arithmetic},
\[
 u\leq\max\{\ell^5N,\ 64N\ell^3\}\leq\ell^7N.
\]
Here \(64\leq\ell^4\) since \(\ell\geq3\).
For every positive \(u\), \(2^{\lceil\log_2u\rceil}\leq2u\);
hence
\[
                    2^n\leq2\ell^7N\leq\ell^8N.
\]
It follows that
\[
 \begin{split}
 k&=2(\ell-1)2^n\leq2\ell^9N\leq x^{10},\\
 M(a+1)&\leq\ell^9N^2(2a)
                    \leq\ell^{10}N^2a\leq x^{10}.
 \end{split}
\]

Choose a prime divisor \(p\) of the positive integer \(6\ell N-1\).
It is different from two and three, so \(p\geq5\), and
\(p\nmid\ell N\). Since the only added factor in \(M\) is a power
of two, \(p\nmid M\). It is also outside the Case II ramification
set, whose prime divisors are those of \(2\ell N\).
Moreover
\[
       p\leq6\ell N-1<6\ell N\leq x^2,\qquad \log p\leq x^2,
\]
where \(a\geq2\), \(\ell\geq3\) ensure \(\ell Na^2\geq6\).

Take negative logarithms of the product comparison. Since the
differences are nonzero,
\[
 (\log\ell)v_{\ell,a}(\psi)
       \leq2v_\ell(J)\log\ell
       \leq2\log|J|
       \leq2Ck\{M(a+1)\}^{A}\log p
       \leq2C x^{10A+12}.
\]
The integer inequality \(v_\ell(J)\log\ell\leq\log|J|\) is immediate
from \(\ell^{v_\ell(J)}\mid J\). Finally
\(2C\leq2^{b_4}\leq x^{b_4}\), proving the stated exponent.
All constants precede the host, \(\ell,a,\psi\), and coefficient
extension.
\end{proof}

\begin{srproposition}[A fixed Gaussian test prime]
\label{prop:sr-gaussian-fixed-prime}
Raise the fixed principal-host threshold to \(q,d>11\).
Then \(p=11\) is one eligible inert test prime for every Gaussian
twist-one instance at \(\ell=3,7\). Its unique Gaussian prime \(w\)
has norm \(121\), and \(11\nmid\ell N\).

Suppose the Gaussian common-integer comparison has fixed constants
\(C_G,A_G\in\mathbb N\), so that, with the Galois-stable list of
all complex characters modulo \(N\), it supplies \(J_G\ne0\) with
\[
              \log|J_G|
                 \leq C_G\varphi(N)(2N)^{A_G}\log11 .
\]
Suppose also that the marked universal ring, its finite weight-zero
part, the automorphy and rigidity of its type-I points, and the
cotangent comparison have supplied the source's inequality
\[
 v_{\ell,1}(\psi)
     \leq2\sum_{\rho} -\log_\ell|d_\rho|_\ell ,
\]
with the nonisomorphism conditions needed for the common-integer
comparison after restriction to \(G_{\mathbb Q(i)}\).
Then
\[
 (\log\ell)v_{\ell,1}(\psi)
        \leq4C_G\varphi(N)(2N)^{A_G}\log11
        \leq(\ell N)^{B_G},
 \qquad B_G=A_G+1+b_G,
\]
where \(b_G\) is any fixed natural number satisfying
\(4C_G2^{A_G}\log11\leq3^{b_G}\).
\end{srproposition}
\begin{proof}
Since \(11\equiv3\bmod4\), the polynomial \(X^2+1\) is irreducible
over \(\mathbb F_{11}\): a root would give an element of order four
in the group \(\mathbb F_{11}^{\times}\), whose order is ten.
Thus \(11\) is inert in \(\mathbb Q(i)\) and its residue-field norm
is \(11^2\). It divides neither the discriminant four nor \(N\);
also \(\ell\ne11\).

A point's automorphic realization may choose either of the two
embeddings \(\mathbb Q(i)\hookrightarrow\Omega\). Partition the
finite representative family accordingly. For each part the
common-integer comparison bounds the sum of the negative
\(\ell\)-adic logarithms by \(v_\ell(J_G)\).
The same integer works for either embedding because its choice
precedes the algebraic closure, embedding and coefficient prime
in that comparison. Adding the two inequalities and then applying
the cotangent factor two gives
\[
             (\log\ell)v_{\ell,1}(\psi)
                         \leq4v_\ell(J_G)\log\ell
                         \leq4\log|J_G|.
\]
This argument uses the source's distinctness after restriction;
pairwise distinct rational representations alone would not
justify this step.

Finally \(\varphi(N)\leq N\), so the right side is at most
\(4C_G2^{A_G}\log11\cdot N^{A_G+1}\).
Since \(\ell\geq3\), \(N\geq1\), and \(b_G\geq0\), this is at most
\((\ell N)^{A_G+1+b_G}\).
\end{proof}

Combining the preceding three branches would permit the single
exponent
\[
                  B=\max\{B_{\mathrm{small}},B_4,B_G\}.
\]
The small-conductor twist-one branch is proved above without
deformation inputs. The other two displayed absorptions require
their stipulated deformation and common-integer comparisons.
They do not assert that the finite weight-zero, no-block,
automorphic-point or marked-ring constructions have been supplied.

\paragraph{Frozen-source locations.}
The precise source definitions are
\(\mathrm{HeightAndVanishing.DirichletPrimeToLPartMulLPowerPart}\)
(\(\mathrm{Host.H}\), \(\mathrm{R1}\), \(\mathrm{nII}\), and the cases),
\(\mathrm{DenseThresholdCases.AfStmt}\)
(\(\mathrm{NumberSp4Alt}\)),
\(\mathrm{DenseThresholdCases.CountS}\),
\(\mathrm{DenseThresholdCases.ConductorAtMapInjectiveCoeff}\),
\(\mathrm{DenseThresholdCases.EHostNumberU}\), and
\(\mathrm{DenseThresholdCases.CountVOf}\), at tag
\texttt{alpoge-source-7cbf7a0c}.
The finite weight-zero/no-block/automorphic rank-four producer is
\(\mathrm{Sp4In0.sp4In0\_alt}\), through
\(\mathrm{CuFinal.main4\_alt}\).
The Gaussian marked-ring producer is
\(\mathrm{DnEndsTop.ends0\_alt}\), second component
\(\mathrm{ThmS6Core\_s}\).

\section{Restriction and infinitesimal rigidity}
\label{sec:sr-restriction-rigidity}
We now prove the restriction and local tangent statements needed when
the integer comparisons are applied to their actual automorphic points.
The existence and automorphy of those points are separate inputs.


\subsection{Restriction to the infinite cyclotomic kernel}
\label{sec:sr-cyclotomic-restriction}
The ordinary type-I argument in
Section~\ref{sec:spectrum-open-irreducibility} proves irreducibility
on every open subgroup of the relevant absolute Galois group.
The cyclotomic kernel is closed and usually not open.
We give the additional group-theoretic argument needed to pass to it.

\begin{srlemma}\label{lem:sr-procyclic-restriction}
Let \(G\) be a profinite group, \(H\) a closed normal subgroup
with procyclic quotient, and \(L\) an algebraically closed
Hausdorff topological field of characteristic zero.
Let \(V\) be a finite-dimensional continuous \(L\)-representation
of \(G\).
Suppose \(V\) is irreducible on every open subgroup of \(G\).
Then \(V|_H\) is irreducible.
Finite-dimensional linear subspaces are understood with their
usual coordinate topology, so they are closed.
\end{srlemma}
\begin{proof}
Choose a nonzero \(H\)-submodule of smallest dimension.
It is simple. Every \(G\)-translate is again a simple
\(H\)-submodule, since \(H\) is normal.
Their sum is nonzero and \(G\)-invariant, hence is \(V\).
A finite subcollection already spans \(V\).
A sum of simple submodules is semisimple: add them successively,
omitting a summand if it is contained in the sum already chosen;
otherwise its intersection with that sum is zero by simplicity.
Thus \(V|_H\) is semisimple.

Group its simple summands into their isotypic components
\(V_1,\ldots,V_j\). These are intrinsic \(H\)-submodules and \(G\)
permutes them.
The stabilizer of \(V_i\) is closed: in a basis adapted to
\(V_i\), invariance is the vanishing of a fixed matrix block.
It has finite index, since there are finitely many components.
A closed finite-index subgroup of a profinite group is open,
as its complement is a finite union of closed cosets.
If there were more than one component, this open stabilizer
would preserve the proper nonzero subspace \(V_i\), contradicting
the hypothesis. Hence there is only one isotypic component.

Choose an \(H\)-module identification \(V=W\otimes_L M\), with
\(W\) simple and \(H\) acting trivially on \(M\).
Schur's lemma here is elementary: an \(H\)-equivariant endomorphism
of \(W\) has an eigenvalue over \(L\); the kernel after subtracting
that scalar is a nonzero \(H\)-submodule, so the endomorphism is scalar.
It follows, by writing a matrix on the multiplicity factors, that
\[
                 \operatorname{End}_H(V)
                           =1_W\otimes\operatorname{End}_L(M).
\]

Choose \(g\in G\) lifting a topological generator of \(G/H\).
Conjugation by \(g\) preserves the simple \(H\)-type \(W\).
Choose an intertwining isomorphism \(T\) on \(W\) for that conjugation.
After multiplying the action of \(g\) by \(T^{-1}\) on the first
factor, the resulting operator commutes with \(H\).
Thus the action of \(g\) is \(T\otimes B\) for some
\(B\in\mathrm{GL}(M)\).
Over the algebraically closed field \(L\), \(B\) has an eigenline
\(M_1\). The subspace \(W\otimes M_1\) is invariant under \(H\)
and \(g\), in both directions.

The subgroup generated by \(H\) and \(g\) is dense in \(G\).
Indeed its closure is compact, its image in \(G/H\) is closed and
contains a topological generator, and it contains \(H\);
its image is therefore all of \(G/H\), so the closure is all of \(G\).
The stabilizer of \(W\otimes M_1\) is closed by the same block-matrix
argument. It consequently equals \(G\).
Irreducibility of \(V\) now implies \(W\otimes M_1=V\), so
\(\dim M=1\). This proves irreducibility on \(H\).
\end{proof}

\begin{srproposition}\label{prop:sr-cyclotomic-restriction}
Let \(F\) be a number field and let \(\ell\) be odd.
A continuous representation over \(\overline{\mathbb Q}_\ell\)
which is irreducible on every open subgroup of \(G_F\)
is irreducible on \(G_{F(\mu_{\ell^\infty})}\).
This conclusion remains true after a finite extension of \(F\)
and after any scalar character twist.
\end{srproposition}
\begin{proof}
The quotient by the indicated kernel is the closed cyclotomic
image in \(\mathbb Z_\ell^\times\).
For odd \(\ell\),
\(\mathbb Z_\ell^\times\simeq\mu_{\ell-1}\times\mathbb Z_\ell\)
is procyclic; one element generates every finite quotient by
the Chinese remainder theorem, since the finite factor has
order prime to \(\ell\).
Every closed subgroup of a procyclic profinite group is procyclic.
Apply Lemma~\ref{lem:sr-procyclic-restriction}.
Every open subgroup of \(G_{F'}\), for \(F'/F\) finite, is also
open in \(G_F\), so the same hypothesis holds after this base change.
A scalar twist has precisely the same invariant subspaces as the
original representation on every subgroup.
\end{proof}

For the normalized ordinary type-I points, the open-subgroup
hypothesis is proved in
Section~\ref{sec:spectrum-open-irreducibility}; normalization by
the inverse cyclotomic character does not change it.
The proposition therefore verifies the infinite-extension
irreducibility condition required by
\href{https://arxiv.org/pdf/2207.04925}{Thorne, Theorem 6.1},
once the specified regular algebraic polarized cuspidal
automorphic realization has been constructed.
That theorem concerns the adjoint descended to the maximal
totally real subfield.
Its application still requires the actual tangent classes to
satisfy the Bloch--Kato local conditions.
The proposition supplies the irreducibility condition, rather
than the automorphic realization or those local conditions.

\subsection{A fixed finite catalogue of imaginary quadratic fields}
\label{sec:sr-finite-CM-catalogue}
The rank-four automorphy predicate used by the source provides
automorphy over imaginary quadratic fields where the restriction
is irreducible. It does not itself assert automorphy over
\(\mathbb Q\). The following elementary argument selects such a
field from a fixed finite list.

\begin{srlemma}\label{lem:sr-quadratic-self-twists}
Let \(G\) be a group and \(V\) an irreducible representation of
dimension \(n\) over an algebraically closed field of
characteristic zero.
There are at most \(n^2\) characters \(\eta:G\to L^\times\)
with \(V\simeq V\otimes\eta\), including the trivial character.
If \(H\) is normal of index two with quadratic character
\(\eta_H\), then \(V|_H\) is reducible exactly when
\(V\simeq V\otimes\eta_H\).
An irreducible representation of \(H\) has at most two isomorphism
classes of extensions to \(G\).
\end{srlemma}
\begin{proof}
A self-twist gives a nonzero intertwiner in \(\operatorname{End}_L(V)\);
it is invertible by irreducibility.
It spans an eigenline for conjugation by \(G\), with character
\(\eta\) or its inverse according to the intertwiner convention.
Different characters give linearly independent eigenlines.
To prove this directly, take a linear relation of minimal
nonzero length among chosen eigenvectors. Conjugate by an element
separating two of its characters and subtract the appropriate
scalar multiple of the original relation. This gives a nonzero
relation of smaller length, a contradiction.
Their number is at most \(\dim\operatorname{End}_L(V)=n^2\).

The restriction to a normal subgroup is semisimple by the
simple-submodule orbit argument in the preceding proof.
For an index-two subgroup its isotypic components form one orbit
of size one or two.
If there is one type with multiplicity greater than one,
the \(T\otimes B\) argument for a lift of the nontrivial coset
gives a proper \(G\)-invariant subspace from an eigenline of \(B\).
This contradicts irreducibility.
Thus reducibility of \(V|_H\) forces two isotypic components.
The nontrivial coset swaps them.
The operator equal to \(1\) on the first and \(-1\) on the second
commutes with \(H\) and anticommutes with that coset.
It intertwines \(V\) with \(V\otimes\eta_H\).

Conversely, let \(T\) be such an intertwiner.
Its square commutes with \(G\), hence is a nonzero scalar.
Rescale \(T\) so that \(T^2=1\).
The operator \(T\) cannot be scalar, since it anticommutes with an
invertible operator in the nontrivial coset and the characteristic
is not two.
Both its \(1\) and \(-1\) eigenspaces are therefore nonzero,
and both are \(H\)-invariant. This proves the converse.

Finally identify the restrictions of two extensions to an
irreducible \(H\)-representation.
For a fixed element \(g\) outside \(H\), their two operators
implement the same conjugation action on \(H\).
Their ratio commutes with \(H\), and so is scalar by Schur's lemma.
Their squares agree because \(g^2\in H\), so that scalar is \(1\)
or \(-1\). These are precisely the two twists by the quotient
characters. This proves the last assertion.
\end{proof}

Choose \(n^2+1\) distinct rational primes \(p_i\equiv3\pmod4\)
and put \(K_i=\mathbb Q(\sqrt{-p_i})\).
These choices are made before the conductor, twist and coefficient
prime. Only finitely many primes are needed; their existence follows,
for example, from Dirichlet's theorem.
For any continuous irreducible rank-\(n\) representation of
\(G_{\mathbb Q}\), at most \(n^2-1\) of these restrictions are
reducible, by the lemma.
For an odd coefficient prime \(\ell\), at most one \(K_i\) is
ramified at \(\ell\), since its discriminant is \(-p_i\).
Thus at least one fixed \(K_i\) has irreducible restriction and
is unramified at \(\ell\).
For rank four a list of seventeen fields suffices.

There is a corresponding finite-family counting statement.
Assign each member of a finite set of pairwise nonisomorphic
irreducible \(G_{\mathbb Q}\)-representations to one field from
the fixed list where it restricts irreducibly.
If the restrictions assigned to \(K_i\) have at most \(B_i\)
isomorphism classes, the last assertion of the lemma gives
\[
                         |\mathscr R|\leq2\sum_i B_i.
\]
In particular constants in established counts over the fixed
fields can be maximized before choosing the varying host.
This is a transfer of an already verified restriction count.
Applying an automorphic count also requires its precise
weight and level conditions for these restrictions.

\subsection{A rank-four cyclotomic irreducibility criterion}
\label{sec:sr-rank-four-cyclotomic}

The following criterion supplies the cyclotomic irreducibility part of
the rank-four rigidity argument. It uses the residual quadratic
self-twist exclusion, rather than an irreducibility assertion on all
open subgroups.

\begin{srlemma}\label{lem:sr-rank-four-no-self-twist}
Let \(\ell\) be odd, let \(\Gamma\) be profinite, and let
\(\epsilon:\Gamma\to\mathbb Z_\ell^\times\) be continuous.
Let \(V\) be a continuous absolutely irreducible representation of
dimension four over a finite extension of \(\mathbb Q_\ell\).
Suppose a stable lattice has residual semisimplification
\(\chi_1\oplus\cdots\oplus\chi_4\), over a finite splitting residue
field \(k\). Write \(\bar\epsilon\) for the residual cyclotomic
character. Assume
\[
 \left.
 \begin{gathered}
 \eta:\Gamma\to k^\times,\quad \eta^2=1,\quad
            \ker\bar\epsilon\subseteq\ker\eta,\\
 \{\eta\chi_1,\ldots,\eta\chi_4\}
                     =\{\chi_1,\ldots,\chi_4\}
                   \quad\hbox{as multisets}
 \end{gathered}\right\}
                  \quad\Longrightarrow\quad\eta=1.
 \tag{C}
\]
Then \(V|_{\ker\epsilon}\) is absolutely irreducible.
\end{srlemma}
\begin{proof}
Extend scalars to an algebraic closure \(L\) of the coefficient
field, and put \(H=\ker\epsilon\).
The quotient \(\Gamma/H\) is a closed subgroup of
\(\mathbb Z_\ell^\times\), hence is procyclic for odd \(\ell\).
Choose \(g\in\Gamma\) lifting a topological generator.

The restriction to \(H\) is semisimple.
Indeed choose an \(H\)-submodule \(W\ne0\) of smallest dimension.
It is simple, and its \(\Gamma\)-translates are simple \(H\)-submodules
because \(H\) is normal. Their sum is nonzero and
\(\Gamma\)-invariant, so is all of \(V_L\).
A finite subcollection spans \(V_L\).
A finite sum of simple modules is semisimple: when a further simple
module is added, its intersection with the preceding sum is either
zero or that whole simple module.

Let \(V_0,\ldots,V_{d-1}\) be the isotypic components.
The action of \(\Gamma\) on them is transitive; otherwise the sum
over one orbit is a proper \(\Gamma\)-submodule.
Each stabilizer is closed, by the vanishing of a block of matrix
entries in a basis adapted to that component.
It has finite index and is therefore open.
The permutation action consequently factors continuously through
the procyclic quotient \(\Gamma/H\).
Thus \(g\) cyclically permutes the components.
They have equal dimension, so \(d\mid4\).

If \(d>1\), then \(d=2\) or \(4\).
Let \(S\) be the stabilizer of \(V_0\). Its image in
\(\Gamma/H\) has cyclic quotient of order \(d\).
The character of this quotient taking \(g\) to \(-1\) defines a
continuous nontrivial quadratic character
\(\eta_0:\Gamma\to\{1,-1\}\).
The operator \(T\) which is \((-1)^i\) on \(V_i=g^iV_0\)
commutes with \(S\) and anticommutes with \(g\).
Therefore \(V_L\simeq V_L\otimes\eta_0\).

Any continuous quadratic character of a closed subgroup of
\(\mathbb Z_\ell^\times\) kills its intersection with
\(1+\ell\mathbb Z_\ell\): that intersection is pro-\(\ell\),
whereas its image has order at most two.
Hence \(\eta_0\) factors through \(\bar\epsilon\).
Reduce its values to \(k\), obtaining \(\bar\eta_0\).
The self-twist gives equality of characteristic polynomials on
every element; the residual semisimplifications therefore have
the same Jordan--Hölder constituents, by the characteristic-polynomial
form of Brauer--Nesbitt. Consequently \(\bar\eta_0\) permutes the
multiset of the four residual characters.
Condition (C) gives \(\bar\eta_0=1\).
Reduction is injective on \(\{1,-1\}\) because \(\ell\ne2\),
contradicting the choice of \(\eta_0\).
Thus \(d=1\).

Write the single isotypic restriction as
\(V_L=W\otimes_L M\), with \(W\) simple for \(H\).
Conjugation by \(g\) preserves the isomorphism class of \(W\).
Choose an intertwiner \(A\) for this conjugation.
After multiplying the action of \(g\) by \(A^{-1}\) on \(W\),
the resulting operator commutes with \(H\).
Schur's lemma gives
\(\operatorname{End}_H(V_L)=1_W\otimes\operatorname{End}_L(M)\);
hence \(g\) acts as \(A\otimes B\), for some invertible \(B\) on \(M\).
An eigenline \(M_1\) of \(B\) makes \(W\otimes M_1\) invariant
under \(H\) and \(g\).

The subgroup generated by \(H,g\) is dense in \(\Gamma\).
The stabilizer of the linear subspace \(W\otimes M_1\) is closed,
again by block-matrix equations, so this subspace is
\(\Gamma\)-invariant. Absolute irreducibility forces
\(W\otimes M_1=V_L\). Thus \(\dim M=1\), proving the result.
\end{proof}

The criterion (C) is exactly the residual exclusion in
\(\mathrm{TNiceFour.nice0four\_alt}\), called
\(\mathrm{P1.condC}\) at the linked point.
The latter is proved in
\(\mathrm{WeakShimuraTheory.ProjectiveLineFlags.CondC}\).
The four characters in the actual linked packet are pairwise distinct
by the tower's separation condition at its distinguished element.
Thus the source's inclusion condition on the four characters is
equivalent here to permutation of the multiset.
Its transport to the linked point involves no estimate:
the level quotient map is surjective, its residual cyclotomic
character is the usual one by \(\mathrm{P1.levEps}\), and its
four residual characters have a common scalar-character factor.
That common factor cancels from a proposed permutation.
A quadratic character takes only the values \(1,-1\), so it descends
through an injective residue-field extension.
Pulling it back, cancelling the common factor, and using the
tower's residual self-twist exclusion gives (C).
This explains exactly which residual property of the constructed
tower is being used. It does not construct that tower.

Applied to the linked rank-four point, the lemma proves
irreducibility over the full cyclotomic kernel.
Scalar twists do not change invariant subspaces.
Surjectivity of the level quotient identifies its cyclotomic
kernel with the image of \(G_{L(\mu_{\ell^\infty})}\).
Thus the same assertion holds for the underlying restricted
Galois representation. The additional condition called
\(\mathrm{P1.condD}\), and the source's regular-coset property
\(\mathrm{N3}\), are not needed for this particular deduction.
They have separate uses in the source's patching argument.

\subsection{The local condition for an ordinary first-order lift}
\label{sec:sr-rank-four-local-tangent}

We give the local reason for the Bloch--Kato condition; ordinary
triangularity by itself would not suffice at an exceptional adjacent
character ratio.

\begin{srlemma}\label{lem:sr-positive-ordinary-extension}
Let \(F/\mathbb Q_\ell\) be finite and \(E/\mathbb Q_\ell\) a finite
coefficient extension. Let
\(\widetilde r:G_F\to\operatorname{GL}_n(E[\tau]/\tau^2)\)
be continuous and upper triangular.
Its diagonal characters are \(\widetilde\kappa_i\), with reductions
\(\kappa_i\).
Suppose there are strictly decreasing integers
\(w_1>\cdots>w_n\) such that on inertia
\[
                \widetilde\kappa_i=\epsilon^{w_i}
\]
as \(E[\tau]/\tau^2\)-valued characters.
Assume also that
\[
               \kappa_i\kappa_j^{-1}\ne\epsilon
                  \quad\hbox{on }G_F\quad (i<j).
 \tag{N}
\]
Then \(\widetilde r\), regarded as a \(2n\)-dimensional
\(E\)-representation, is crystalline.
Its reduction \(r\) is crystalline, and its first-order class
belongs to \(H^1_f(F,\operatorname{ad}r)\).
\end{srlemma}
\begin{proof}
The \(i\)-th diagonal block, regarded over \(E\), is a
two-dimensional unramified self-extension of
\(\kappa_i=\epsilon^{w_i}\alpha_i\), with \(\alpha_i\) unramified.
It is crystalline; unramified representations are crystalline
and tensoring by an integral cyclotomic power preserves that property.

Build the upper triangular representation one diagonal block at a time.
Suppose the preceding blocks form a crystalline representation \(U\),
and the next block is \(W\), of weight \(w_j\).
The representation \(T=\operatorname{Hom}_E(W,U)\) is crystalline.
Its semisimplification consists of the characters
\(\kappa_i\kappa_j^{-1}\), for \(i<j\), with their multiplicities.
Their cyclotomic exponents \(w_i-w_j\) are positive.
Thus \(H^0(F,T)=0\) and
\[
      \dim_E D_{\rm dR}(T)/\mathrm{Fil}^0
                          =[F:\mathbb Q_\ell]\dim_E T.
\]
Moreover \(H^0(F,T^*(1))=0\): if \(w_i-w_j>1\), its corresponding
character has nonzero cyclotomic exponent; if \(w_i-w_j=1\),
a fixed vector would contradict (N).
A representation with a finite filtration and no trivial
Jordan--Hölder constituent has no fixed vector, so this check on
the constituents suffices.

Local Tate duality gives \(H^2(F,T)=0\).
The local Euler characteristic and the crystalline fundamental
exact sequence give respectively
\[
 \dim_E H^1(F,T)=[F:\mathbb Q_\ell]\dim_E T,\qquad
 \dim_E H^1_f(F,T)
       =\dim_E D_{\rm dR}(T)/\mathrm{Fil}^0+\dim_E H^0(F,T).
\]
The dimensions agree. Since \(H^1_f(F,T)\subseteq H^1(F,T)\),
every extension class is crystalline.
Therefore the extension of \(W\) by \(U\) in \(\widetilde r\)
is crystalline. Induction proves the assertion.
Finally the sequence
\(0\to r\to\widetilde r\to r\to0\), regarded over \(E\),
identifies its class with the defining crystalline local
condition in \(H^1(F,\operatorname{ad}r)\).
\end{proof}

Here is the concrete check for the rank-four packet.
Write the ordinary cyclotomic exponents as
\[
                         (a+1,a,1,0),\qquad a\geq2.
\]
Choose the source's Frobenius lift \(\sigma\) at \(\ell\) with
\(\epsilon(\sigma)=1\).
On the ordered residual flag its eigencharacters at \(\sigma\)
have values
\[
                              uv,\quad u,\quad v,\quad1,
\]
where \(u=\bar\psi(\sigma)\) and \(v=\bar\xi(\sigma)\).
The good-data condition is
\[
                              vu^m\ne1\qquad(m\in\mathbb Z).
 \tag{D}
\]
For weight difference one, the first and second positions, and the
third and fourth positions, have residual ratio \(v\ne1\).
When \(a=2\), the middle positions also have weight difference one,
and their residual ratio is \(u/v\ne1\), by (D) with \(m=-1\).
These are all the pairs with weight difference one.

The remaining finite inertial characters can be killed over a
finite local extension whose Galois subgroup still contains \(\sigma\).
To see this, collect those characters into one finite-image character
\(\chi\). Divide \(\chi\) by the unramified character taking
\(\sigma\) to \(\chi(\sigma)\), and take its kernel.
This is open and normal, contains \(\sigma\), and kills \(\chi\)
on inertia. Over this extension every ordinary diagonal has the
indicated pure cyclotomic restriction on inertia.
A ratio with weight difference other than one cannot become
\(\epsilon\) on an open group, since the cyclotomic image on inertia
is infinite. For difference one, the residual check at the retained
\(\sigma\) excludes that equality.
Thus Lemma~\ref{lem:sr-positive-ordinary-extension} applies.

Crystalline local vanishing descends from this finite extension.
For the period-module cohomology defining \(H^1_f\), restriction
is injective after replacing the extension by a finite normal one:
the inflation--restriction kernel is a degree-one cohomology group
of a finite group acting on a characteristic-zero vector space,
and averaging kills it.
The argument therefore proves the original Bloch--Kato condition.

Away from \(\ell\), the source's local conditions make the lift
trivial on an open subgroup of inertia.
The tangent cocycle is then zero on that subgroup.
An explicit finite-index averaging proves it is an inertia
coboundary: after a local coboundary is subtracted, let the cocycle
be \(z\), zero on \(J\leq I\) of finite index \(m\).
It factors as a function on left cosets, and if
\(S=\sum_{gJ\in I/J}z(g)\), the cocycle identity gives
\(S=mz(h)+hS\). Hence \(z(h)=(h-1)(-S/m)\).

For the formal point problem, the ordered flags and their residual
labels are the ones specified by the local deformation conditions.
One first descends the point and its dual-number lift to a finite
coefficient field, as in the source's
\(\mathrm{TbTriDiag.tb\_desc}\).
The preceding proof then supplies the local content of
\(\mathrm{TbTriDiag.tanBK}\), including the adjacent-ratio check.
It does not infer the Bloch--Kato condition solely from ordinary
triangularity.

\subsection{The correct polarized adjoint and the symplectic projection}
\label{sec:sr-rank-four-adjoint-projection}

\begin{srproposition}\label{prop:sr-polarized-symplectic-rigidity}
Let \(r:G_{\mathbb Q}\to\operatorname{GSp}_4(E)\) be continuous,
with fixed similitude character \(\theta\), and let \(K\) be an
imaginary quadratic field with character \(\eta\).
Fix a complex conjugation \(c\) with \(\theta(c)=-1\).
Let \(J\) be the alternating matrix preserved up to \(\theta\), and set
\[
 X^\#=J^{-1}X^{\mathsf T}J,\qquad
 P(X)=\tfrac12(X-X^\#),\qquad
 A(g)X=r(g)Xr(g)^{-1}.
\]
Define the representation on \(M_4(E)\) by
\[
 R_\eta(g)X=A(g)P(X)+\eta(g)A(g)(X-P(X)).
 \tag{R}
\]
It is the polarized adjoint descended from \(r|_{G_K}\), with
polarizing matrix \(B=Jr(c)^{-1}\):
\[
 R_\eta(h)X=A(h)X\quad(h\in G_K),\qquad
 R_\eta(c)X=-B^{-1}X^{\mathsf T}B.
\]
Suppose its Bloch--Kato Selmer group over \(\mathbb Q\) vanishes.
Then every fixed-similitude first-order lift of \(r\) whose tangent
class satisfies the local conditions proved above is conjugate
to the constant lift by \(1+\tau Y\), with \(Y\in\mathfrak{sp}_4(E)\).
\end{srproposition}
\begin{proof}
The symplectic adjoint is an involution and commutes with
conjugation by a symplectic similitude. Thus \(P^2=P\),
its image is \(\mathfrak{sp}_4(E)\), and
\(P A(g)=A(g)P\).
The two invariant summands \(\operatorname{im}P\) and \(\ker P\)
show directly that (R) is a representation.
On \(G_K\), \(\eta=1\), giving the ordinary adjoint action.
At \(c\), \(\eta=-1\), so
\[
 R_\eta(c)X=A(c)(2P(X)-X)=-A(c)X^\#
                                      =-B^{-1}X^{\mathsf T}B.
\]
For \(h'=chc^{-1}\), the similitude identity gives
\[
                  r(h)^{\mathsf T}B\,r(h')=\theta(h)B,
\]
which identifies the asserted polarization and its descent.
Since \(r(c)^2=1\) and \(\theta(c)=-1\), the same identity gives
\(B^{\mathsf T}=B\), the required odd polarized descent in this
rank-four setting.

Write a first-order lift as
\(\widetilde r(g)=r(g)+\tau\,\dot r(g)\), and put
\(z(g)=\dot r(g)r(g)^{-1}\).
The homomorphism identity says
\(z(gh)=z(g)+A(g)z(h)\).
Differentiating the fixed-similitude identity
\(\widetilde r(g)^{\mathsf T}J\widetilde r(g)=\theta(g)J\)
gives \(z(g)^\#=-z(g)\). Hence \(Pz=z\).

The map \(P\) intertwines \(A\) with \(R_\eta\):
\(P A(g)=R_\eta(g)P\).
It carries inertia-coboundary witnesses to inertia-coboundary
witnesses, and its tensor product with the identity on
\(B_{\rm cris}\) does the same for the crystalline witnesses.
Thus \(z\) is a Bloch--Kato cocycle for \(R_\eta\).
By the assumed vanishing it is a coboundary,
\(z(g)=R_\eta(g)X-X\).
Projecting once more gives
\[
                    z(g)=A(g)P(X)-P(X).
\]
Set \(Y=-P(X)\). Then \(Y\in\mathfrak{sp}_4(E)\) and
\[
                    \dot r(g)=Yr(g)-r(g)Y.
\]
This is exactly conjugation by \(1+\tau Y\).
\end{proof}

The applicable published vanishing statement is
\href{https://arxiv.org/pdf/2207.04925}{Thorne, Theorem 6.1}:
for an actual regular algebraic polarized cuspidal realization
over a CM field, cyclotomic-kernel irreducibility implies
vanishing of its polarized adjoint descended to the maximal
totally real subfield.
This is not a statement about the full adjoint Selmer group
over the CM field.

For a realization over a larger CM field \(L\supset K\), the
same application is valid after the polarization is matched.
Since \(K\cap L^+=\mathbb Q\), one has \(L=KL^+\).
Consequently \(R_\eta|_{G_{L^+}}\) is precisely the descended
adjoint of \(r|_{G_L}\) with the displayed polarization.
An algebraic scalar twist does not change its adjoint on \(G_L\);
its polarizing form must still be identified, up to scalar,
with the form defining this descent.
This is the purpose of the scalar norm identity in
\(\mathrm{AdDown}\), not a consequence of adjoint
twist-invariance alone.
Once this matching and the actual polarized cuspidal realization
are supplied, the vanishing over \(L^+\) implies the vanishing
over \(\mathbb Q\): restriction preserves the local conditions
and is injective in degree one by finite-index averaging in
characteristic zero.

The rank-four automorphy predicate \(\mathrm{IsAut4Alt}\) is
\(\mathrm{IsAutQA}\). Its exact definition provides a regular
algebraic cuspidal realization over every imaginary quadratic
field where the restriction is irreducible, together with
\(\mathrm{WitA}\). It does not assert a cuspidal
\(\operatorname{GL}_4/\mathbb Q\) realization.
The field selection, local level compatibility and polarization
matching in the preceding paragraph are therefore essential.

\paragraph{What these arguments close.}
They prove the residual criterion for cyclotomic irreducibility,
the ordinary tangent's Bloch--Kato local condition with its
nonexceptionality check, and the symplectic projection from
the correctly descended polarized adjoint.
They do not prove the quantitative exceptional estimate without
the actual rank-four finite/no-block/automorphic-point producer.
That source producer is
\(\mathrm{Sp4In0.sp4In0\_alt}\), through
\(\mathrm{CuFinal.main4\_alt}\) and
\(\mathrm{ModularityRankFourOfPromodularity}\).
Its Case II input is
\(\mathrm{Gsp4PromodularityCaseTwo.proMod\_of\_pmEx\_of\_ctx2\_of\_kPair\_of\_star}\),
which constructs the deep residue chain and propagates the modular
support from an automorphic seed.
For the present principal packet, no Case III count is needed.
The Gaussian twist-one branch still needs the marked universal-ring
and type-I automorphy statement \(\mathrm{ThmS6Core\_s}\), supplied
in the formal dependency chain by \(\mathrm{DnEndsTop.ends0\_alt}\).
All paths in this paragraph are in the frozen tag
\texttt{alpoge-source-7cbf7a0c}.

\section{Regular cosets on the actual curve branches}
\label{sec:sr-curve-regular-cosets}
This section proves the regular-element density needed for the
rank-three Taylor--Wiles construction.  Its hypotheses are the
local separation and residual character conditions of the selected
curve points.  Cyclotomic irreducibility by itself is not used as a
substitute for them.

Write \(A=k[[T]]\), \(F=k((T))\), and let \(\overline F\) be an
algebraic closure.  Here \(k\) is finite of characteristic
\(\ell\ge5\), as in the generic principal-character branch.
Let \(G\) be the compact profinite global group and let
\(G'\) be its closed cyclotomic kernel.  A curve representation has
entries in \(A\), continuous modulo every \(T^j\).
The local ordinary characters are denoted by
\(\psi_{v,i}:D_v\to A^\times\).
The residual global determinant is the product of the characters
\(\chi_i:G\to k^\times\).  All residual characters in this section
are characters of the full group \(G\), not merely of \(G'\).

\subsection{Local separation on every open subgroup}
\begin{srlemma}[No torsion in the principal units]
\label{lem:sr-curve-principal-units}
The group \(1+Tk[[T]]\) is torsion free.
Suppose a compact local subgroup \(N_v\subset D_v\) maps into \(G'\),
the ratios \(\psi_{v,i}/\psi_{v,j}\) have residue one on \(N_v\),
and no such ratio is identically one there.
For every open subgroup \(V\subset G\), there is a
\(g\in N_v\) mapping into \(V\cap G'\) such that
\[
       \psi_{v,i}(g)^m\ne\psi_{v,j}(g)^m
       \qquad(i\ne j,\ m\ge1).
\]
\end{srlemma}
\begin{proof}
If \(z\in Tk[[T]]\) is nonzero and
\(m=\ell^r s\), with \(\ell\nmid s\), then
\[
 (1+z)^m=(1+z^{\ell^r})^s
       =1+s z^{\ell^r}+\text{terms of higher \(T\)-order}.
\]
The displayed leading term is nonzero.  This proves the first claim.

Put \(W=N_v\cap \iota_v^{-1}(V)\), where
\(\iota_v:D_v\to G\) is the local map.
It is open and of finite index in \(N_v\).
A character ratio nontrivial on \(N_v\) stays nontrivial on \(W\):
a positive power of any chosen element of \(N_v\) lies in \(W\),
and the principal units have no torsion.
Its image on \(W\) is therefore infinite.
The kernel of that ratio is closed and has infinite index in \(W\).
For completeness, such a subgroup has Haar measure zero:
for every positive integer \(M\) there are \(M\) distinct cosets,
all of its measure, and their disjoint union has measure at most one.
There are only finitely many pairs \(i\ne j\).
Their kernels cannot cover \(W\), because a finite union of
measure-zero sets has measure zero.
Choose \(g\) outside this union.
Torsion freeness now proves every asserted inequality.
\end{proof}

At a selected nice curve point, these hypotheses hold with
\(N_v\) the prescribed local subgroup on which both the
residual character and the cyclotomic character are trivial.
The equal residual values follow from that definition.
The condition usually denoted \(\mathrm{Nice.sep}\) says precisely
that the differences of the two universal local characters do
not vanish on this subgroup after specialization.
A scalar twisting character cancels from every ratio.
Thus the preceding lemma applies to the actual twisted point,
on every open subgroup, before any regular detector is selected.
This deduction still requires a selected curve with that
nonvanishing property; it does not prove existence of such curves.

\subsection{The prime-rank restriction argument}
We give the elementary irreducibility argument that strengthens
the global type condition on these curves.

\begin{srlemma}[Characters and reduction]
\label{lem:sr-curve-character-matching}
Let \(r:G\to\mathrm{GL}_d(A)\), \(d<\ell\), have residual
characteristic polynomial
\(\prod_{i=1}^d(X-\chi_i(g))\) at every \(g\).
Suppose on a subgroup \(P\) its extension to \(\overline F\)
is a sum of characters \(\lambda_1,\ldots,\lambda_d\).
There is a single permutation of the indexed multisets such that
the reductions of the \(\lambda_i\) equal
\(\chi_i|_P\), after that permutation.
The same assertion holds over \(\overline F\) for two sums of
characters with equal characteristic polynomials.
\end{srlemma}
\begin{proof}
Extend the \(T\)-adic valuation to \(\overline F\).
Every \(\lambda_i(g)\) is integral because it is a root of a
monic characteristic polynomial over \(A\).
Applying the same argument to \(g^{-1}\) shows its inverse is
integral.  It is a unit of this valuation ring and has a
nonzero residue in an algebraic closure of \(k\).
Reduction therefore gives genuine characters of \(P\).

We recall the independence argument for distinct characters.
If a nonzero relation \(\sum_i a_i\theta_i(g)=0\) holds for
every \(g\), choose one with the fewest nonzero coefficients.
For \(\theta_1\ne\theta_j\), choose \(h\) with
\(\theta_1(h)\ne\theta_j(h)\).
Subtract \(\theta_j(h)\) times the relation at \(g\) from the
relation at \(hg\).  This gives a shorter nonzero relation,
a contradiction.  Thus distinct characters are linearly
independent.

Take traces of the two characteristic-polynomial identities.
Group together equal characters and apply this independence.
Their multiplicities are integers between zero and \(d<\ell\),
so equality in characteristic \(\ell\) is equality of those
integers.  The two indexed multisets therefore agree.
Choosing a bijection of equal entries gives one fixed
permutation, valid on the whole subgroup.
\end{proof}

\begin{srproposition}[Irreducibility on the required restrictions]
\label{prop:sr-curve-open-irreducibility}
Let \(d=2\) or \(3\), and suppose \(r:G\to\mathrm{GL}_d(A)\)
is absolutely irreducible over \(\overline F\).
Assume that, for one local group \(D\), its semisimplification
is the sum of \(d\) ordinary characters \(\psi_i\).
Assume also:
\begin{enumerate}
\item for every open subgroup \(V\subset G\), some element of
\(D\) maps into \(V\cap G'\) with all \(\psi_i\)-values distinct;
\item the residual determinant is a sum of global characters
\(\chi_i\), and two of their restrictions to \(D\) are unequal.
\end{enumerate}
Then \(r|_{V\cap G'}\) is absolutely irreducible for every
open subgroup \(V\subset G\).
\end{srproposition}
\begin{proof}
Replace \(V\) by its open normal core and put \(H=V\cap G'\).
It suffices to prove irreducibility on \(H\).

Here is the prime-dimension Clifford argument in the needed form.
Choose a minimal nonzero \(H\)-invariant subspace \(W\).
Every translate \(r(g)W\) is simple for \(H\), because \(H\)
is normal.  The sum of all these translates is a nonzero
\(G\)-invariant subspace, hence the whole representation.
Only finitely many translates are needed to span it.
A sum of simple modules is semisimple: adjoining a simple
module either has zero intersection with the sum already
formed or is contained in it.
Grouping the resulting simple modules into their isotypic
parts, \(G\) acts transitively on these parts.
Their simple dimension, their common multiplicity and
their number have product \(d\).
Since \(d\) is prime, either the restriction is irreducible,
or \(H\) acts by one scalar character, or it has exactly
\(d\) distinct one-dimensional character spaces, permuted
transitively by \(G\).
In the homogeneous one-dimensional case \(H\) acts by scalars
on the entire space; this also follows directly from
semisimplicity.

The scalar possibility contradicts the element in (1):
its eigenvalues are the distinct values of the ordinary
characters.  Suppose then that there are \(d\) distinct
lines, and write
\[
        \pi:G\to S_d,\qquad P=\ker\pi .
\]
The action of \(P\) on these lines is by characters
\(\lambda_1,\ldots,\lambda_d\).
On \(D\cap H\) the line characters are the restrictions of
the \(\psi_i\), in one fixed order, by
Lemma~\ref{lem:sr-curve-character-matching}.
They are distinct there, again by (1).
For \(t\in D\), each character \(\psi_i\) is unchanged by
conjugation by \(t\).
Conjugation therefore cannot permute the distinct characters
on \(D\cap H\).  It fixes every one of the \(d\) lines.
Thus \(D\subset P\).

Reduce the \(\lambda_i\) using
Lemma~\ref{lem:sr-curve-character-matching}.
They become the restrictions of the global \(\chi_i\), in
one fixed order.  A global character is invariant under
conjugation by every element of \(G\).
But the permutation action of \(G\) is transitive on the
lines, so conjugation transports any \(\lambda_i\) to any
\(\lambda_j\).
Their reductions are consequently all the same character
on \(P\).
Since \(D\subset P\), all the \(\chi_i|_D\) are equal,
contradicting (2).
Only the irreducible case remains.
\end{proof}

For the principal rank-three stable branch, the local residual
characters are a common scalar times
\(\overline\Psi,1,\overline\Psi^{-1}\).
The host condition supplies a local element with
\(\overline\Psi^2\ne1\), so the residual distinction in (2)
holds.  Lemma~\ref{lem:sr-curve-principal-units} supplies (1).
The same argument applies to the rank-two block on the
two-plus-one branch.  In the actual integral coordinates,
this block uses the outer indices \(0,2\), and its residual
characteristic polynomial is
\[
        (X-\chi_0(g))(X-\chi_2(g)).
\]
Indeed one of the two opposite off-diagonal entries is zero
after reduction, so their product in the two-by-two
determinant vanishes.  The residual distinction is now
\(\overline\Psi\ne\overline\Psi^{-1}\).
The middle line is a direct summand of the ordinary local
representation; its character is one of the three ordinary
Jordan--Hölder characters.  The two-dimensional summand has
the other two.  Their restrictions to every \(V\cap G'\)
are distinct by the preceding local separation.
Proposition~\ref{prop:sr-curve-open-irreducibility} therefore
makes the outer block absolutely irreducible on every such
restriction.  No irreducibility is asserted for the full
two-plus-one representation.

\subsection{Closed zero sets and principal minors}
\begin{srlemma}[Closedness for algebraic-coefficient linear forms]
\label{lem:sr-curve-linear-closedness}
For every \(\overline F\)-linear functional \(L\) on
\(M_d(\overline F)\), the set
\(\{g:L(r(g))=0\}\) is closed in \(G\).
\end{srlemma}
\begin{proof}
The finitely many coefficients of \(L\) lie in a
finite-dimensional \(F\)-subspace of \(\overline F\).
Choose an \(F\)-basis \(b_1,\ldots,b_s\) for their span and
write
\[
                  L(r(g))=\sum_{\nu=1}^s b_\nu L_\nu(r(g)),
\]
where the \(L_\nu\) have coefficients in \(F\).
The maps on the right take values in \(F\) and are continuous
for its \(T\)-adic topology: each is a finite linear
combination of the continuous entries in \(A\).
Linear independence of the \(b_\nu\) says \(L(r(g))=0\)
exactly when all \(L_\nu(r(g))=0\).
This is a finite intersection of closed sets.
\end{proof}

\begin{srlemma}[A matrix in the coset with nonzero minors]
\label{lem:sr-curve-coset-minors}
Let \(Q\) be a compact profinite group with a representation
\(s:Q\to\mathrm{GL}_d(\overline F)\).
Assume its restriction to every open subgroup is irreducible,
and all linear-functional zero sets in the preceding lemma
are closed.
For every \(y\in\mathrm{GL}_d(\overline F)\), there exists
\(b\in Q\) such that every diagonal entry of \(ys(b)\)
and of its inverse is nonzero.
If \(d\le3\), every principal minor of \(ys(b)\) is nonzero.
\end{srlemma}
\begin{proof}
For a nonzero vector \(v\) and a nonzero linear form \(f\),
the zero set of \(f(s(g)v)\) has empty interior.
Otherwise it contains a coset \(hW\), with \(W\) open in \(Q\).
The span of \(s(W)v\) is a nonzero \(W\)-invariant subspace,
hence the whole space.  The relation
\(f(s(h)s(w)v)=0\) for every \(w\in W\) would make
the nonzero form \(f\circ s(h)\) zero on the whole space.
This is impossible.
The zero set is closed by the linear-functional hypothesis.
Its complement is therefore open and dense.

Apply this to the \(d\) diagonal coefficients of \(ys(g)\).
For the inverse coefficients use \(g\mapsto g^{-1}\), a
homeomorphism of \(Q\), and the diagonal coefficients of
\(s(g)y^{-1}\).
There are only \(2d\) closed zero sets.
The intersection of their open dense complements is nonempty:
intersect the complements one at a time, keeping a nonempty
open set at each step.
A point \(b\) in that intersection has the required property.

For a one-by-one principal minor this is the diagonal assertion.
For \(d=3\), a two-by-two principal minor complementary to
index \(i\) equals
\(\det(ys(b))(ys(b)^{-1})_{ii}\), by the cofactor formula.
It is nonzero.  The full determinant is nonzero because
\(ys(b)\) is invertible.  The empty minor equals one.
The cases \(d\le2\) are included in the same argument.
\end{proof}

The hypotheses here hold after every change of basis over
\(\overline F\), since a matrix coefficient in the new basis is
still a linear functional in the old entries.
They also hold on \(V\cap G'\): every open subgroup of this
closed subgroup contains a subgroup of the form
\(V_1\cap G'\), with \(V_1\) open in \(G\).
This follows by choosing an open ambient identity neighborhood
inside the given relative open neighborhood and then a
profinite open subgroup inside it.
The preceding restriction proposition gives irreducibility.

\subsection{The quadratic and cubic discriminants}
\begin{srlemma}[Separating a small matrix by diagonal powers]
\label{lem:sr-curve-discriminant}
Let \(d\le3\), \(M\in\mathrm{GL}_d(K)\), for a field \(K\),
have all principal minors nonzero.
Let \(\alpha_1,\ldots,\alpha_d\in K^\times\) satisfy
\((\alpha_i/\alpha_j)^m\ne1\) for \(i\ne j\) and \(m\ge1\).
Then the characteristic polynomial of
\(M\operatorname{diag}(\alpha_i)^m\) is separable for some
integer \(m\ge0\).
\end{srlemma}
\begin{proof}
Let \(B\) be the finitely generated multiplicative subgroup
generated by the \(\alpha_i\).
Its quotient by its finite torsion subgroup is a free
abelian group.  Each \(\alpha_i/\alpha_j\) is nonzero in
this quotient.  Choose a rational linear functional on
\(B\otimes_{\mathbb Z}\mathbb Q\) which is nonzero on every
one of these finitely many differences.
Such a functional exists because finitely many proper
hyperplanes cannot cover a vector space over the infinite
field \(\mathbb Q\). In a fixed basis, evaluate the finitely
many nonzero defining linear forms at \((1,t,t^2,\ldots)\).
Each becomes a nonzero rational polynomial in \(t\); an integer
outside their finitely many root sets gives the required functional.  It gives an additive weight
\(w:B\to\mathbb Q\) with all \(w(\alpha_i)\) distinct.

We shall use the following elementary fact.
An exponential sum \(\sum_{\beta}c_\beta\beta^m\), with
distinct nonzero bases \(\beta\), cannot vanish for every
\(m\ge0\) unless all \(c_\beta\) are zero.
For \(s\) bases, the equations for \(m=0,\ldots,s-1\)
have Vandermonde determinant
\(\prod_{i<j}(\beta_j-\beta_i)\ne0\).
In particular a term with uniquely largest \(w\)-weight
cannot cancel with the other bases.

For \(d=0,1\), separability is immediate.
For \(d=2\), relabel so \(w_1>w_2\), and put
\(d_i=M_{ii}\), \(D=\det M\).
The discriminant is the exponential expression
\[
 (d_1\alpha_1^m+d_2\alpha_2^m)^2
                      -4D(\alpha_1\alpha_2)^m .
\]
Its uniquely highest-weight term is
\(d_1^2(\alpha_1^2)^m\), with nonzero coefficient.
It is nonzero for some \(m\), proving the claim.

For \(d=3\), relabel so \(w_1>w_2>w_3\).
Write \(p_{ij}\) for the principal two-by-two minors.
The elementary symmetric coefficients of the characteristic
polynomial are
\[
 e_1=\sum_i d_i\alpha_i^m,\qquad
 e_2=\sum_{i<j}p_{ij}(\alpha_i\alpha_j)^m,\qquad
 e_3=D(\alpha_1\alpha_2\alpha_3)^m .
\]
The discriminant of \(X^3-e_1X^2+e_2X-e_3\) is
\[
 \Delta=e_1^2e_2^2-4e_2^3-4e_1^3e_3
                       -27e_3^2+18e_1e_2e_3 .
\]
The largest possible weights in these five terms, in order,
are
\[
 \begin{gathered}
 4w_1+2w_2,\quad 3w_1+3w_2,\quad
 4w_1+w_2+w_3,\\
 2w_1+2w_2+2w_3,\quad 3w_1+2w_2+w_3 .
 \end{gathered}
\]
The first is strictly largest.
Within its term the unique largest monomial is
\[
              d_1^2p_{12}^2(\alpha_1^4\alpha_2^2)^m .
\]
Its coefficient is nonzero.  The Vandermonde fact therefore
makes \(\Delta\ne0\) for some \(m\).
This reasoning is valid also when some of the integers
\(4,18,27\) become zero in \(K\); the leading coefficient
has coefficient one.
Nonzero discriminant is equivalent to separability for a
monic polynomial, completing the proof.
\end{proof}

\subsection{Density on both actual rank-three branches}
\begin{srtheorem}[The curve-point regularity condition]
\label{thm:sr-curve-N3}
Consider a selected nice curve point of the generic principal
rank-three construction.  On its stable branch assume the
global representation is absolutely irreducible.
On its two-plus-one branch assume the actual outer
rank-two block is absolutely irreducible and the middle
coordinate is a character.
Then, on either branch, every open subset of \(G\)
meeting \(G'\) contains a \(\sigma\in G'\) for which
\(\det(X-r(\sigma))\) is separable.
\end{srtheorem}
\begin{proof}
Let \(O\) be such an open subset, choose \(t\in O\cap G'\),
and choose an open subgroup \(V\subset G\) with \(tV\subset O\).
Local separation supplies \(u\in V\cap G'\) whose full
ordinary eigenvalue tuple \(\alpha\) has torsion-free
pairwise ratios.  It has distinct eigenvalues, so
\(r(u)\) is diagonalizable over \(\overline F\).
Work in that eigenbasis.

On the stable branch,
Proposition~\ref{prop:sr-curve-open-irreducibility} and
Lemma~\ref{lem:sr-curve-coset-minors}, applied to
\(Q=V\cap G'\) and \(y=r(t)\), give \(b\in Q\)
such that \(M=r(t)r(b)\) has every principal minor nonzero.
Lemma~\ref{lem:sr-curve-discriminant} now gives an \(m\)
for which \(Mr(u)^m\) has separable characteristic polynomial.
It is \(r(tb u^m)\), and \(tb u^m\in O\cap G'\).

On the two-plus-one branch diagonalize the rank-two block
of \(r(u)\); its remaining scalar is the middle-line value.
The rank-two restriction is irreducible on every open subgroup
of \(Q\), as proved above.
Choose \(b\in Q\) so that the two diagonal entries of
the rank-two block of \(r(t)r(b)\) are nonzero.
The scalar block is nonzero because it is a character.
Every principal minor of the resulting three-by-three
block matrix is nonzero: the mixed two-by-two minors are
that scalar times a diagonal entry, the remaining
two-by-two minor is the invertible rank-two determinant,
and the full determinant is their product.
Apply Lemma~\ref{lem:sr-curve-discriminant} to the full
three-by-three matrix and its full free-ratio tuple.
The same element \(tb u^m\) works.

Separable characteristic polynomials descend from
\(\overline F\) to \(F\), since extension of fields preserves
the nonvanishing of their discriminant.
This proves exactly the regular-density condition
usually denoted \(N3\), before taking finite lists of
simultaneously regular detectors.
\end{proof}

The arithmetic content used here is the actual nonvanishing
of the local character differences on the selected curve,
the residual local distinction, and the representation type
of that curve.  Selecting curves with these properties and
placing them in the required common support remain the
separate curve-selection and support arguments.
The elementary density proof introduces no additional
automorphy assumption or fixed positive-characteristic
valuation on \(\overline F\).

For correspondence with the frozen development, the relevant
paths under \srcid{CurveDVRSpecialization} are
\srcid{NiceSep}, \srcid{IrrNormal}, \srcid{CliffPrime},
\srcid{NiceIrr}, \srcid{NiceIrrE},
\srcid{LinClosedOfAdic}, \srcid{RegCoset},
\srcid{RegCosetE}, and
\srcid{NiceSpecializationTwistSatisfiesThirdCondition}.
The closed-zero-set argument above replaces the more general
abstract exponential-sum/full-matrix construction by a direct
argument on the actual profinite curve points.


\subsection{The remaining two local nice-point conditions}
\begin{srproposition}[The distinguished local invariant space]
\label{prop:sr-curve-N4}
Let \(r:G\to\mathrm{GL}_n(A)\) have the prescribed residual
determinant. Suppose a distinguished local image \(D_0\)
has the following two properties:
all residual \(\chi_i|_{D_0}\) agree, and some \(g\in D_0\)
has \(\overline\epsilon(g)\ne1\).
Then
\[
 \{Y\in M_n(F):
      \epsilon(g)r(g)Y=Yr(g)\text{ for every }g\in D_0\}=0.
\]
The assertion holds for every point of that same residual
determinant and in every rank.
\end{srproposition}
\begin{proof}
The zero-rank case is immediate.
Use just the distinguished \(g\).
Write \(C=r(g)\), \(e=\epsilon(g)\), and let
\(c\in k^\times\subset A^\times\) be its common residual
character value.
Cayley--Hamilton applied after reduction gives
\((C-cI)^n\in T M_n(A)\).
Put \(N=C-cI\) and \(V=eC-cI\).
Modulo \(T\),
\[
                  V=(\overline e-1)cI+\overline e\,N.
\]
This is invertible: a nonzero scalar plus a commuting
nilpotent matrix has inverse given by a finite geometric
series. Its determinant is a unit of \(A\), so \(V\) is
invertible over \(A\).

Clear a common denominator of \(Y\), so \(Y\in M_n(A)\).
The assumed equation is \(VY=YN\).
Induction gives \(V^nY=YN^n\).
Write \(N^n=TQ\), with \(Q\in M_n(A)\); then
\[
                         Y=T V^{-n}YQ.
\]
Iterating this equality shows every entry of \(Y\) belongs
to \(T^jA\) for every \(j\).
Their intersection is zero, so \(Y=0\).
Clearing denominators was injective, which proves the
assertion over \(F\).
\end{proof}

For the generic principal packet over
\(K=\mathbb Q(\sqrt{-3})\), this local condition can itself
be arranged concretely.
Here \(3\nmid\ell N\), \(\ell\) is odd, and
\(\psi(-1)=(-1)^a\).
Choose a prime \(p\), outside any prescribed finite set, with
\[
                       p\equiv-1\pmod{\ell N},
                       \qquad p\equiv1\pmod3.
\]
The Chinese remainder theorem gives a reduced residue
class, and Dirichlet's theorem supplies infinitely many
such primes.
It splits in \(K\).
Its residual arithmetic Frobenius has
\(\overline\epsilon(p)=-1\) and
\[
           \overline\epsilon(p)^a\overline\psi(p)
                     =(-1)^a\overline\psi(-1)=1.
\]
The three residual characters are therefore all equal
on its full decomposition group, after the common scalar
link twist; inertia is trivial in the residual characters
at this prime.
At the initial imaginary-quadratic packet, taking that local
group for \(D_0\) gives exactly the hypotheses of
Proposition~\ref{prop:sr-curve-N4}.  When the arithmetic link
passes to a further CM tower, it must retain these two
local properties; the initial congruence calculation alone
does not establish that transfer, since a residue-degree
power can change the cyclotomic value.
The local level introduced to obtain this distinguished
place must subsequently be removed by the two-link argument.
No bound for the least such prime is used here.
For other residual packets the two displayed local
properties must be checked separately.

Finally, the ordered ordinary condition \(N1\) follows from
the same local separation used above.
Choose its free-ratio element in the full cyclotomic kernel.
The \(\psi_i\) are pairwise unequal there.
If \(\psi_i=\psi_j\epsilon\) with \(i<j\), their values
would also be equal on this element, because
\(\epsilon=1\) there.  This is a contradiction.
Consequently the selected rank-three curves satisfy \(N1\)
and \(N3\).  Whenever the selected link supplies the
explicit distinguished-place condition \(\mathrm{V0ok}\),
they also satisfy \(N4\) by the full matrix argument above. Their stable and two-plus-one type conditions are
the ones retained throughout the preceding arguments.

The general \(N4\) calculation corresponds to
\srcid{WeakShimuraTheory/ProjectiveLineFlags/ContinuousLinearClosed.lean},
\srcid{P1.n4V0}.  Its hypothesis is the explicit residual
local condition \srcid{V0ok}; a global no-self-twist or
no-Speh condition is not substituted for it.

\section{Finite-level arithmetic Selmer estimates and auxiliary limits}
\label{sec:sr-arithmetic-selmer-producers}
The statements below construct the finite trace detectors and the
auxiliary arithmetic tower, and prove the first- and second-Selmer
bounds needed for the common-support argument.  They apply to the
actual nice paired points with their displayed local conditions.
Their constants may depend on that fixed finite family of points;
the uniform conductor estimates later absorb this dependence through
the common support and automorphic counting argument.


\subsection{The arithmetic Taylor--Wiles tower and its common limit}
\label{sec:arithmetic-tw-tower}

We construct the actual auxiliary arithmetic groups and the limit
used by the completed-support argument.
The number of auxiliary primes is fixed by the selected trace
detectors; it need not equal a conjectural dual Selmer dimension.

Let \(L/L^+\) be the CM extension in the unitary problem, let
\(\ell\) be odd, and let \(S\) contain all places above \(\ell\)
and all ramification of the fixed residual packet.
Write \(\Gamma_0\) for the arithmetic quotient over \(L^+\),
\(G_0=\ker\eta\) for its CM subgroup, and
\(G_0'=\ker(\varepsilon|_{G_0})\), using the \emph{full}
\(\ell\)-adic cyclotomic character.
Keep the actual residual characters \(\bar\chi_1,\ldots,\bar\chi_n\)
and their polarized determinant.
Suppose a finite set
\[
                   s_1,\ldots,s_r\in G_0'
 \tag{TW1}
\]
has been chosen.
For each of the finitely many nice paired points under
consideration, require the characteristic polynomials
of \(\rho(s_j)\) to be separable over its fraction field.
The trace-detection proof chooses the \(s_j\) once, before the
tower.  Regularity is used in the Selmer estimates, not in
Chebotarev's construction below.
The source of the regular-coset property for the actual curve
points must therefore be established at the nice-point step;
absolute irreducibility on the full cyclotomic kernel alone is
not substituted for it.

We use only the following standard arithmetic finiteness input:
for every open subgroup \(U\) of the global arithmetic group,
and of each fixed local group, the set
\(\operatorname{Hom}_{\mathrm{cts}}(U,\mathbb F_\ell)\) is finite.
Globally, an open subgroup corresponds to a finite number-field
extension with a fixed finite ramification set.
Kummer theory, finite generation of its \(S\)-unit group, and
finiteness of its class group give this assertion after adjoining
\(\mu_\ell\); restriction and corestriction, or the
inflation--restriction sequence, give it before adjoining
\(\mu_\ell\).
Locally it follows from local class field theory.
This is the actual \srcid{PhiCond} input; it is independent of
the coefficient field and of the auxiliary tower level.

\paragraph{A residual-kernel subgroup and a cofinal sequence.}
Set
\[
 H=\{g\in\Gamma_0:\eta(g)=1,\ 
          \bar\chi_i(g)=1\ (1\leq i\leq n),\
          \bar\varepsilon(g)=1\}.
 \tag{TW2}
\]
It is open and normal.
Normality under \(G_0\) follows because the \(\bar\chi_i\) are
characters.  Under complex conjugation, polarization identifies
the multiset of conjugate characters with their inverses times
the cyclotomic multiplier; on the displayed kernel the multiplier
is \(1\), so that kernel is preserved.

For an open normal \(J\subset\Gamma_0\), put
\[
 \Phi_\ell(J)=
       \bigcap_{\alpha\in
          \operatorname{Hom}_{\mathrm{cts}}(J,\mathbb F_\ell)}
                         \ker\alpha .
\]
The intersection is finite, so \(\Phi_\ell(J)\) is open;
it is characteristic in \(J\) and hence normal in \(\Gamma_0\).
Take \(W_0=H\) and \(W_{N+1}=\Phi_\ell(W_N)\).
Define, exactly as in the source,
\[
 \mathrm{PL}_\ell(H)=
 \bigcap_{\substack{V\triangleleft_{\mathrm{open}}\Gamma_0\\
     \forall h\in H\ \exists e\geq0:\ h^{\ell^e}\in V}}V.
 \tag{TW3}
\]
Thus a finite quotient used in this intersection has an
\(\ell\)-group image of \(H\).
We have \(W_N\subset H\), and
\[
 \forall V\triangleleft_{\mathrm{open}}\Gamma_0,\
 \mathrm{PL}_\ell(H)\subset V
 \quad\Longrightarrow\quad
 W_N\subset V\quad(N\text{ sufficiently large}).
 \tag{TW4}
\]
To check this, the image \(P\) of \(H\) in \(\Gamma_0/V\)
is a finite \(\ell\)-group.
Indeed the maximal pro-\(\ell\) quotient of \(H\) is the quotient
by (TW3): normal cores show that using open normal subgroups
of \(\Gamma_0\), rather than all open normal subgroups of \(H\),
does not change it.
The images of \(W_N\) are contained in the iterated Frattini
subgroups of \(P\), which reach \(1\).
Equivalently, every nontrivial finite \(\ell\)-group has a
nonzero homomorphism to \(\mathbb F_\ell\), so each nontrivial
Frattini step strictly decreases its order.
This proves (TW4).

\paragraph{Choosing the primes, with the Frobenius convention fixed.}
Lift each \(s_j\) to \(g_j\in G_L\).
Since \(s_j\in G_0'\), its lift has
\(\varepsilon(g_j)=1\).
Let
\[
 U_N=\operatorname{pr}^{-1}(W_N)
                 \cap\ker(\varepsilon\bmod\ell^N)\subset G_L .
\]
This is open and normal.
Chebotarev, applied in the finite quotient by \(U_N\), chooses
distinct auxiliary places \(v_{Nj}\) of \(L^+\), outside \(S\),
split in \(L/L^+\), and places \(w_{Nj}\) of \(L\) above them,
with arithmetic Frobenius representatives \(\sigma_{Nj}\)
satisfying
\[
                         \sigma_{Nj}g_j\in U_N .
 \tag{TW5}
\]
Here one may choose degree-one primes of \(L\).
The usual avoidance of the finitely many rational primes
ramifying in \(L/L^+\), together with degree one, makes the
underlying \(L^+\)-place split; it also avoids every prescribed
finite set.
This is a finite Chebotarev problem at each \(N\), so simultaneous
choice of the \(r\) distinct places follows by finite induction.

Enlarge \(S\) by these \(v_{Nj}\) and let
\(h_N:\Gamma_N\twoheadrightarrow\Gamma_0\) be the resulting
arithmetic quotient map.
At \(w_{Nj}\), choose a topological generator \(\tau_{Nj}\)
of tame pro-\(\ell\) inertia and use the source convention
\[
             \phi_{Nj}=\operatorname{pr}_N(\sigma_{Nj}^{-1}),
 \qquad Q_{Nj}=\mathrm Nw_{Nj}.
 \tag{TW6}
\]
The inverse in this definition matters.
Cyclotomic compatibility with arithmetic Frobenius and (TW5)
give
\[
                           Q_{Nj}\equiv1\pmod{\ell^N}.
 \tag{TW7}
\]
Normality of \(W_N\), followed by inversion of (TW5), gives
\[
             h_N(\phi_{Nj})s_j^{-1}\in W_N .
 \tag{TW8}
\]
Consequently these products lie in \(H\) for every \(N\), and
converge to \(1\) modulo every open normal subgroup containing
\(\mathrm{PL}_\ell(H)\).

Put \(H_N=h_N^{-1}(H)\).
The tame relation, in the inverse-Frobenius convention (TW6),
is
\[
 \phi_{Nj}^{-1}\tau_{Nj}\phi_{Nj}
       =\tau_{Nj}^{Q_{Nj}}
                        \pmod{\mathrm{PL}_\ell(H_N)}.
 \tag{TW9}
\]
Wild inertia at this place has pro-residue-characteristic order,
hence has trivial image in a pro-\(\ell\) quotient.
The prime-to-\(\ell\) part of tame inertia also disappears.
The local kernel of the tame pro-\(\ell\) quotient maps into
\(\mathrm{PL}_\ell(H_N)\): every finite quotient in (TW3) has
an \(\ell\)-group image on \(H_N\), so it kills these local
prime-to-\(\ell\) subgroups.
The arithmetic kernel \( \ker h_N\) is the closed normal subgroup
generated by the new inertia groups.
It follows that
\[
 \begin{gathered}
 h_N(\tau_{Nj})=1,\\
 \ker h_N\subset
 \overline{\langle\!\langle
       \tau_{N1},\ldots,\tau_{Nr},
                         \mathrm{PL}_\ell(H_N)\rangle\!\rangle}.
 \tag{TW10}
 \end{gathered}
\]
The old local maps and their cyclotomic and polarization
identities are preserved by \(h_N\).

\paragraph{Uniform numbers of generators, without a Selmer assumption.}
There is a constant
\[
 b_H=
 \#\operatorname{Hom}_{\mathrm{cts}}(H,\mathbb F_\ell)\,
                   \ell^{\,r[\Gamma_0:H]}
 \tag{TW11}
\]
such that
\[
             \#\operatorname{Hom}_{\mathrm{cts}}
                         (H_N,\mathbb F_\ell)\leq b_H
                     \quad\hbox{for every }N.
 \tag{TW12}
\]
Here is the arithmetic kernel count.
Modulo homomorphisms inflated from \(H\), a character of \(H_N\)
is determined by its values at the \(r[\Gamma_0:H]\) elements
\(x\tau_{Nj}x^{-1}\), using one lift \(x\) of each coset
of \(H\) in \(\Gamma_0\).
Indeed conjugation by \(H_N\) does not affect a homomorphism
from \(H_N\) to an abelian group.
If all those values vanish, it kills the closed normal subgroup
generated by the new \(\tau_{Nj}\).
It also kills \(\mathrm{PL}_\ell(H_N)\), by (TW3).
Equation (TW10) then makes it descend to \(H\).
Each value has \(\ell\) choices, proving (TW11)--(TW12).
Thus the bound is obtained before either of the first- and
second-Selmer estimates.

\paragraph{Why the required points and first-order data descend.}
For a point of the fixed residual determinant, every
\(h\in H\) has residual characteristic polynomial \((X-1)^n\).
Its residual representation need not be the identity matrix.
Nevertheless its image on \(H\) modulo \(T^m\) is a finite
\(\ell\)-group: all residual elements are unipotent, so their
orders are powers of \(\ell\), and Cauchy's theorem shows that
their finite group has \(\ell\)-power order.
The congruence kernel is an \(\ell\)-group too.
It follows that the point kills \(\mathrm{PL}_\ell(H)\).
One may also see this directly from Cayley--Hamilton:
\[
 \bar\rho(h)\text{ unipotent}\quad\Longrightarrow\quad
     (\rho(h)^{\,\ell^{nm}}-1)\in T^mM_n(k[[T]]).
 \tag{TW13}
\]
Apply the same calculation to all conjugates of \(h\), take the
normal core of the finite representation kernel, and use (TW3).

This argument includes the odd generators and the actual
first- and second-order cocycles.
Their semidirect-product or square-zero-matrix representations
at finite coefficient precision have an additive
\(\ell\)-group kernel and the same residual unipotent image
on \(H\) or \(H_N\).
They therefore kill the corresponding \(\mathrm{PL}\).
The old local conditions are tested on the original local groups,
so composing their maps with the quotient does not change those
conditions.
Inflation consequently identifies the relevant finite Selmer
cocycles and boundaries on the original groups and on the
\(\mathrm{PL}\)-quotients.
In particular the first- and second-Selmer
bounds proved below transport with their original constants.

There is a corresponding statement for the completed determinant
and its paired Cayley--Hamilton algebra.
At any finite Artinian coefficient quotient the algebra is finite.
The residual characteristic polynomial of \(h\in H\) is
\((X-1)^n\), so \(h-1\) is nilpotent modulo the maximal ideal.
Nilpotence of that ideal makes \(h-1\) nilpotent in the finite
algebra.  A unit \(1+x\) with \(x\) nilpotent over a finite
\(\ell\)-power coefficient ring has \(\ell\)-power order,
by the binomial theorem.
Thus the finite \(H\)-image is again an \(\ell\)-group and
kills \(\mathrm{PL}_\ell(H)\).
Taking all finite coefficient quotients proves descent in the
completed topology, including the \emph{closed}
Cayley--Hamilton relations.
This is the same determinant and the same completed support
ideal used in the support-propagation subsection.
No new determinant ring is inserted at the limit step.

We may therefore pass to the actual quotient tower and henceforth
assume
\[
             \mathrm{PL}_\ell(H)=1,\qquad
             \mathrm{PL}_\ell(H_N)=1 .
 \tag{TW14}
\]
The sign and full cyclotomic character descend:
\(\eta|_H=1\), and the image of \(H\) in
\(1+\ell\mathbb Z_\ell\) is pro-\(\ell\).
Every map and every marked element above descends as well.
The image of \(H_N\) is now pro-\(\ell\), and (TW12) bounds its
minimal number of topological generators.
This is the Burnside basis theorem for pro-\(\ell\) groups:
generators of its finite Frattini quotient lift to topological
generators.  The coarse bound \(b_H\) is more than sufficient.

\paragraph{One finitely generated cover for every level.}
Let \(a_0=[\Gamma_0:H]\).
Choose \(a_0\) coset representatives of \(H\) and one additional
marked representative equal to complex conjugation \(c\).
Lift the coset representatives to each \(\Gamma_N\), using
\(c_N\) for the marked one, and add at most \(b_H\) generators
of \(H_N\).
The resulting \(a_0+1+b_H\) elements generate \(\Gamma_N\).
Let \(P\) be the free profinite group on this fixed number of
generators.
It has continuous surjections
\[
                       \pi_N:P\twoheadrightarrow\Gamma_N
 \tag{TW15}
\]
with one fixed \(c_P\in P\) mapping to \(c_N\), and with
\(h_N\pi_N\) independent of \(N\) after projection to
\(\Gamma_0/H\).
The kernel \(H_N\) is pro-\(\ell\) and finitely generated;
the finite coset representatives make its closure together with
these lifts the whole group.
This verifies surjectivity in (TW15), rather than merely
mapping a free group to each level.

\paragraph{The kernel limit and the base map.}
Choose an ultrafilter \(\mathcal U\) on \(\mathbb N\) containing
every cofinite set, and put \(K_N=\ker\pi_N\).
For every open normal \(V\subset P\), the images \(K_NV/V\)
range over the finite set of subgroups of \(P/V\).
Exactly one such image occurs on a set in \(\mathcal U\).
These images are compatible as \(V\) varies, and their transition
maps are surjective: this holds for the images of every \(K_N\)
and hence holds on the intersection of the finitely many
relevant sets in \(\mathcal U\).
Compactness produces a closed normal subgroup \(K_\infty\subset P\)
with exactly those finite images.
Set
\[
                        \Gamma_\infty=P/K_\infty .
\]
Then
\[
                K_NV=K_\infty V
                   \quad\text{for \(\mathcal U\)-almost all }N
 \tag{TW16}
\]
for every fixed \(V\).

Write \(f_N=h_N\pi_N\).
For any finite quotient \(\Gamma_0/V_0\), there are only finitely
many continuous homomorphisms from \(P\) to that quotient,
because \(P\) has finitely many topological generators.
Their \(\mathcal U\)-limits are compatible and define
a continuous \(f_\infty:P\to\Gamma_0\).
For each \(V_0\), the equality \(f_N=f_\infty\) modulo \(V_0\)
holds \emph{uniformly on all of \(P\)} for almost all \(N\).
This proves \(K_\infty\subset\ker f_\infty\), using (TW16).
It also proves surjectivity, since every \(f_N\) is surjective
onto every finite quotient and the image of a compact group
is closed.
We obtain
\[
                      h_\infty:\Gamma_\infty
                                      \twoheadrightarrow\Gamma_0 .
 \tag{TW17}
\]
The sign and cyclotomic characters and complex conjugation are
the ones obtained from \(f_\infty\).
In particular \(c_\infty^2=1\), since \(c_N^2=1\) at every level.

\paragraph{Uniform local limits; no lifting of local homomorphisms.}
Here an arbitrary pointwise limit is insufficient.
For each fixed old local group \(Y\), let
\(H_Y\) be the inverse image of \(H\) under its base map, and put
\(P_Y=\mathrm{PL}_\ell(H_Y)\), defined using \(Y\).
The local finiteness input and the pro-\(\ell\) Burnside basis
argument show that \(Y/P_Y\) is topologically finitely generated.
Every local map \(i_N:Y\to\Gamma_N\) kills \(P_Y\), by (TW14):
a finite quotient of \(H_N\) is an \(\ell\)-group, so the
pullback of each such quotient appears in the defining
intersection for \(P_Y\).

Fix \(V\triangleleft_{\mathrm{open}}P\).
For almost all \(N\), the subgroup \(S=K_NV=K_\infty V\)
is fixed.  At those levels there is a canonical quotient
\[
                         \Gamma_N\longrightarrow P/S ,
\]
because \(\ker\pi_N=K_N\subset S\).
Its composite with \(i_N\) is a continuous homomorphism
\(Y/P_Y\to P/S\).
There are finitely many such homomorphisms.
Thus it is a single \emph{whole local map}, simultaneously for
all \(y\in Y\), for almost all \(N\).
These stabilized maps are compatible as \(V\) varies, and give
a continuous local homomorphism
\(i_\infty:Y\to\Gamma_\infty\).
Its composite with \(h_\infty\) is the original base local map.

Equivalently, for every \(V\), the local elements at level \(N\)
and at the limit have lifts to \(P\) whose ratio lies in \(V\),
uniformly for all \(y\in Y\), for almost all \(N\).
This is the precise \srcid{Close} condition in the limit
definition.
The argument uses no continuous homomorphic lift
\(Y\to P\), which need not exist.
If one chooses arbitrary set sections of \(\pi_N\), their
reductions modulo \(S\) are the canonical quotient map just
described and are continuous; that is all that is required.
There are only finitely many old local places, so all these
conditions may be imposed simultaneously.

\paragraph{Marked Frobenius, inertia and the commuting limit.}
Choose set-theoretic lifts of \(\phi_{Nj}\) and \(\tau_{Nj}\)
to \(P\), and take their compact \(\mathcal U\)-limits.
Changing these lifts changes the limit only by \(K_\infty\),
by (TW16); hence they define elements
\(\phi_{\infty j},\tau_{\infty j}\in\Gamma_\infty\).
The finite set of marked elements has the same uniform
\srcid{Close} property.
Equations (TW4), (TW8), and (TW14) give
\[
 h_\infty(\phi_{\infty j})=s_j,\qquad
 h_\infty(\tau_{\infty j})=1 .
 \tag{TW18}
\]

Finally these two marked elements commute.
In the quotient tower, (TW9) is an exact identity.
At any fixed finite quotient arising from \(P/V\), the image of
\(\tau_{Nj}\) belongs to an \(\ell\)-group of bounded order,
because \(H_N\) is pro-\(\ell\).
Its exponent therefore divides some \(\ell^e\), independent of
\(N\) for that fixed finite quotient.
For \(N\geq e\), (TW7) makes
\(\tau_{Nj}^{Q_{Nj}}=\tau_{Nj}\) in that quotient.
The commutators consequently vanish in every finite quotient
at the limit:
\[
                        [\phi_{\infty j},\tau_{\infty j}]=1 .
 \tag{TW19}
\]
Together (TW15)--(TW19), the uniformly stabilized local maps,
and the cofinal ultrafilter are the actual limit data used in
the cohomology-limit and completed-support constructions.

\paragraph{Scope and exact frozen sources.}
This proves the arithmetic tower and the common group limit
from a fixed finite list (TW1).
The separable eigenvalues may be split by a finite extension of
the complete coefficient DVR.
For finitely many points, choose a common multiple of the residue
degrees, then extend their finite residue fields to the same
cardinality.  The resulting DVR extensions are finite, injective
and local; all point equations and odd generators extend with
them.
Thus the eigenvalue and common-residue-cardinality requirements
at the cohomology-limit interface are obtained by actual finite
coefficient extensions.

All paths below are under
\srcid{Development.AutomorphicArithmetic.AutomorphicVarieties}
at the frozen tag \srcid{alpoge-source-7cbf7a0c}.
The residual kernel is
\srcid{PatchingModules.ChebotarevFamilies.StarSingle.hEx};
its cofinal sequence is \srcid{FrSNormal.frat}.
The prime choice is \srcid{ChebOne} and
\srcid{ChebFam.chebFam}; the enlarged arithmetic quotient and
kernel are \srcid{LevUp.levUp} and \srcid{kerLe}.
The uniform character count is \srcid{HomBound.homBound},
and the residual-unipotent descent is \srcid{UnipPow.rhoPL}.
The assembled arithmetic construction is
\srcid{PatchingModules.AutomorphicDetection.ArPl.arPl}.
The quotient and cover are
\srcid{PatchingModules.SeparableMatrixStability.QuotEx.quot_ex},
\srcid{PlGen.pl_gen}, and \srcid{SysCover.sys_cover}.
The kernel limit uses \srcid{AlgebraTools.LimKer} and
\srcid{LimElt}; the uniform local argument is
\srcid{SeparableMatrixStability.Loc_fin.loc_fin} and
\srcid{InLim.in_lim}; the commuting limit is
\srcid{CommLim.comm_lim}.
Coefficient splitting and common residue degree are
\srcid{PatchingModules.DeterminantPadicFree.CharPOfPadicAlt.onePt}.
These sources are in
\url{https://github.com/ChristianKudeba/alpoge-rh-verification/tree/alpoge-source-7cbf7a0c/src/Development/AutomorphicArithmetic/AutomorphicVarieties}.



\subsection{The arithmetic second-Selmer bound}
\label{sec:arithmetic-second-selmer}

Here we prove the finite-level assertion \srcid{TWGrp.B2}, rather
than assuming it in the limit argument.
Let \(A=k[[T]]\), \(F=k((T))\), \(q=\#k\), with
\(\operatorname{char}k=\ell\ne2\).
Keep a nice paired point \(\rho\) of rank \(n\), with block sizes
\(a,b=n-a\), and the original odd Selmer generators
\(\kappa_1,\ldots,\kappa_g\).
The hypotheses on these objects used below are precisely:
\[
 \begin{gathered}
 \rho(G_0)\subset
   \mathrm{GL}_a(A)\times\mathrm{GL}_b(A),\qquad
 F[\rho(G_0')]=M_a(F)\oplus M_b(F),\\
 G_0'=\ker(\varepsilon:G_0\to\mathbb Z_\ell^\times),\\
 \psi_{v,i}\ne\psi_{v,j}\quad(i\ne j),\qquad
 T^s z-\sum_jc_j\kappa_j\text{ is an odd boundary}
 \text{ for every integral odd Selmer }z,
 \tag{B2a}
 \end{gathered}
\]
where \(s\) and \(c_j\in A\) in the last statement may depend on \(z\).
We also keep the actual polarized and unramified equations:
\(\rho(cgc^{-1})=\varepsilon(g)^{1-n}\rho(g)^{-\mathsf t}\),
\(\rho(I_u)=1\),
\(\kappa_j(cgc^{-1})=-\kappa_j(g)^{\mathsf t}\), and
\(\kappa_j(I_u)=0\).
Each \(\kappa_j\) is a continuous odd adjoint cocycle with its
integral ordinary local correcting matrices and prescribed
zero diagonal on the local \(\ell\)-primary torsion.
The ordinary bases are integral and respect the block type;
the characters \(\psi_{v,i}\) are trivial on that specified
local \(\ell\)-primary torsion.
The finite-coefficient Galois cohomology groups are finite.
We use the global and local Euler characteristics, local duality,
the global localization bound in degree two, and vanishing in degree
three for \(\ell\)-primary coefficients.  For the actual CM
arithmetic groups these are the cohomological statements proved
in the global-cohomology subsection: the group over the totally
real base has odd \(\ell\), and all primes above \(\ell\) are in the
ramification set.

Let \(G_N\to G_0\) be the actual auxiliary arithmetic groups.
For clarity, the Taylor--Wiles properties used in this proof are:
there is one fixed open normal \(H\subset\Gamma_0\) on which
\(\eta\), \(\bar\varepsilon\), and the residual character semisimplification
are trivial; there are fixed \(s_1,\ldots,s_r\in G_0'\), each
\(\rho(s_j)\) having distinct eigenvalues over \(\overline F\);
and there are auxiliary Frobenius and inertia elements
\(\phi_{Nj},\tau_{Nj}\) satisfying
\[
 \begin{gathered}
 h_N(\tau_{Nj})=1,\qquad
 \rho(h_N(\phi_{Nj}))\equiv\rho(s_j)\pmod{T^m}
       \quad(N\text{ sufficiently large for }m),\\
 \mathrm Nw_{Nj}\equiv1\pmod{\ell^N},\qquad
 \phi_{Nj}^{-1}\tau_{Nj}\phi_{Nj}
           =\tau_{Nj}^{\,\mathrm Nw_{Nj}}
               \pmod{\mathrm{PL}(H_N)},\\
 \ker h_N\subset
 \overline{\langle\!\langle\tau_{N1},\ldots,\tau_{Nr},
                      \mathrm{PL}(H_N)\rangle\!\rangle}.
 \tag{B2b}
 \end{gathered}
\]
Here \(\mathrm{PL}(H_N)\) is the kernel used to retain the maximal
pro-\(\ell\) quotient of the residual-kernel subgroup \(H_N\).
These are arithmetic prime-selection and tame inertia statements;
their construction is separate from the bound proved here.
The point kills \(\mathrm{PL}(H)\).
It suffices that every residual characteristic polynomial on \(H\)
is \((X-1)^n\); one need not assume the residual matrices are identity.
Their finite image consists of unipotent matrices and is therefore
an \(\ell\)-group: every element has \(\ell\)-power order and
Cauchy's theorem excludes other prime divisors of the group order.
The kernel of reduction modulo \(T\) is an \(\ell\)-group as well.
Thus every reduction modulo \(T^m\) has an \(\ell\)-group image on
\(H\), which kills \(\mathrm{PL}(H)\); separatedness proves the
integral assertion.

\paragraph{One uniform approximation lemma.}
We first record the complete-module argument responsible for uniformity
in \(m\).  Suppose \(C,C'\) are \(T\)-torsion-free \(A\)-modules,
\(C\) is complete and separated, \(C'\) is separated, and
\(\Phi:C\to C'\) is linear.  Put
\(S_m=\Phi^{-1}(T^mC')\).
Suppose \(S_1/TC\) is finite, and an \(A\)-submodule \(B\subset\ker\Phi\) has the
property that every \(z\in\ker\Phi\) has some \(T\)-power in \(B\).
Then, for one constant \(C_\Phi\) before all \(m\), the reduction
of \(S_m\) modulo \(T^mC\) is covered by at most \(q^{C_\Phi}\)
cosets of the reduction of \(B\).
The same is true after any linear projection of these reductions.

Here is a proof that does not assume \(C\) finitely generated.
The images of \(S_m\) in \(C/TC\) form a decreasing sequence of
subgroups of the finite image of \(S_1\); they stabilize, say from
\(m_0\) onwards.
Division by \(T\) is legitimate in \(C'\).
Here are the stabilization and limit steps explicitly.
For every \(j\geq1\), the images of \(S_m\) in \(C/T^jC\)
are constant for \(m\geq m_0+j-1\).
For the induction, let \(x\in S_{m_0+j}\) and choose a solution
\(s_1\) at any larger prescribed precision, with
\(x-s_1=T c_1\), using stabilization modulo \(T\).
Take that precision at least \(m_0+j\).
Then \(c_1\in S_{m_0+j-1}\), because
\(T\Phi(c_1)=\Phi(x-s_1)\) and \(C'\) is torsion free.
Apply the induction hypothesis to \(c_1\), approximating it
modulo \(T^j\) by a solution \(s_2\) at the prescribed precision.
The solution \(s_1+Ts_2\) approximates \(x\) modulo \(T^{j+1}\).
This proves the induction and, in particular, the convenient
slightly weaker assertion
\[
 \forall j,M,\quad
 x\in S_{m_0+j}\Longrightarrow
 \exists s\in S_M,\quad x-s\in T^jC.
 \tag{B2c1}
\]
Starting with \(r_0=x\in S_{m_0+j}\), apply (B2c1) successively
to choose
\[
 r_n\in S_{m_0+j+n},\qquad
 r_n-r_{n+1}=T^{j+n}d_n .
\]
Completeness gives \(L=\sum_{n\geq0}T^nd_n\in C\).
Set \(k=x-T^jL\).
For each \(n\), \(k-r_n\in T^{j+n}C\), whereas
\(\Phi(r_n)\in T^{m_0+j+n}C'\).
Separatedness of \(C'\) gives \(\Phi(k)=0\).
Moreover \(T^j\Phi(L)=\Phi(x)\in T^{m_0+j}C'\);
torsion freeness implies \(L\in S_{m_0}\).
Thus
\[
 S_{m_0+j}\subset\ker\Phi+T^jS_{m_0}\qquad(j\geq0).
 \tag{B2c}
\]
The kernel is complete and
\(\ker\Phi\cap TC=T\ker\Phi\), again by torsion freeness of \(C'\).
Thus its reduction modulo \(T\) is finite; successive approximation
over \(A\) proves it is a finite \(A\)-module.
The pointwise \(T\)-power condition on \(B\) is therefore uniform:
\(T^s\ker\Phi\subset B\) for some \(s\).
Choose finite representatives for \(\ker\Phi/B\) and for
\(S_{m_0}\) modulo \(T^{m_0}C\).
Formula (B2c) covers every \(S_{m_0+j}\) modulo \(T^{m_0+j}C\) by
the sums of one representative of the first kind and \(T^j\)
times one of the second.
There are uniformly finitely many such sums.
The finitely many \(m<m_0\) are absorbed into the constant.
Their finiteness follows inductively from finiteness at \(m=1\):
a difference divisible by \(T^r\) can be divided, and its new
residue is a solution modulo \(T\).
This proves the lemma.

Apply it to continuous integral global cochains together with their
finitely many local correcting matrices.
Let \(\Phi\) list the cocycle, polarization, inertia, ordinary
lower-entry and prescribed zero-diagonal defects.
Its target is a product of finite free \(A\)-modules.
The source is complete in the uniform \(T\)-adic topology:
compatible locally constant reductions define a continuous integral
cochain.  The solution set modulo \(T\) is finite by finite Galois
\(H^1\), finite boundaries, and finitely many local matrices.
Every finite-coefficient cochain lifts continuously, by using a
section of the finite coefficient reduction.
Take \(B\) to consist of exact solutions whose global cochain is
in the span of the \(\kappa_j\) and integral boundaries; leave the
compatible local correcting matrices unrestricted.
The last condition in (B2a) is exactly the \(T\)-power condition
on this \(B\).  Projecting to global cochains in the lemma yields
\[
 \#\operatorname{Sel}_{\mathrm{odd},0}(A/T^e)
                  \leq q^{ge+C_0}\qquad(e\geq0).
 \tag{B2d}
\]
This is a bound on classes, since the image of \(B\) has at most
\(q^{ge}\) classes.  It is not obtained merely by asserting that
the integral rank is \(g\); the approximation lemma controls
classes which exist only at finite precision.

\paragraph{The auxiliary primes add only a bounded odd Selmer factor.}
Let \(M_e=\operatorname{Odd}_a(A/T^e)\), with adjoint \(G_0\)-action,
extended to the full polarized group by minus transpose at complex
conjugation.
For each \(j\), the endomorphism
\[
 L_j:\operatorname{Odd}_a(A)\to\operatorname{Odd}_a(A),\qquad
       Y\longmapsto \rho(s_j)Y\rho(s_j)^{-1}-Y
\]
is injective after inverting \(T\).
Indeed the centralizer of a matrix with distinct eigenvalues is a
commutative algebra.  It contains the even block idempotent, so an
odd matrix in that centralizer commutes with the idempotent and is
zero.  Smith normal form consequently gives \(c_j\) with
\[
       \#\ker(L_j\bmod T^e)\leq q^{c_j}\quad\hbox{for every }e.
 \tag{B2e}
\]

For fixed \(m\) and sufficiently large \(N\), take an odd cocycle
modulo \(T^e\), \(e\leq m\).
Its value at \(\tau_{Nj}\) is unchanged by a coboundary, because
\(\rho(\tau_{Nj})=1\).
Its finite semidirect-product representation kills
\(\mathrm{PL}(H_N)\): the image of \(H_N\) in the residual kernel,
including the cocycle coordinate, is an \(\ell\)-group.
The tame relation (B2b) and
\(\mathrm Nw_{Nj}\equiv1\pmod\ell\) therefore show that its value
at \(\tau_{Nj}\) centralizes \(\rho(h_N(\phi_{Nj}))\).
This is the same centralizer as in (B2e) at precision \(e\).
If all these inertia values vanish, (B2b) makes the cocycle descend
to \(G_0\), with the same old local conditions.
Inflation on \(H^1\) is injective, so the inertia-value map gives
\[
 \#\operatorname{Sel}_{\mathrm{odd},N}(A/T^e)
 \leq \#\operatorname{Sel}_{\mathrm{odd},0}(A/T^e)
                                  \prod_jq^{c_j}
 \leq q^{ge+C_0+C_{\mathrm{aux}}},
 \quad C_{\mathrm{aux}}=\sum_jc_j.
 \tag{B2f}
\]
The constant \(C_{\mathrm{aux}}\) is fixed for the tower.

\paragraph{The actual relative obstruction group.}
At an ordinary place \(v\mid\ell\), let \(B_v\subset M_e\) be
the upper-triangular odd submodule in the chosen ordinary basis,
and put \(Q_v=M_e/B_v\).
At a prescribed non-\(\ell\) place put \(Q_u=M_e\), imposing
unramified classes on the specified inertia subgroup.
The ordinary type condition makes \(Q_v\) free of rank \(ab\)
over \(A/T^e\).
Define \(\mathcal H^2_e\) by pairs
\[
 \left(u\in Z^2(G_N,M_e),\
       b_w\in C^1(G_w,Q_w):\
       db_w=\pi_w\operatorname{res}_w u\right),
\]
modulo pairs coming from a global \(1\)-cochain and local
unramified \(1\)-cocycles.
At ordinary places unramified here means locally a boundary,
because the specified local subgroup is the full group.
This is precisely the group
\srcid{Dat.H2LocalConditions}; its defining cochains retain the
ordinary correcting matrices.

Write \(\mathcal C_e\) for the cokernel of
\(H^2(G_N,M_e)\to\prod_wH^2(G_w,Q_w)\).
The relative-cochain exact sequence, followed by Euler
characteristics, gives
\[
 \#\mathcal H^2_e
 \leq
 \#\operatorname{Sel}_{\mathrm{odd},N}(A/T^e)\,
 \left(\prod_{v\mid\ell}\#H^0(G_v,Q_v)\right)\#\mathcal C_e.
 \tag{B2g}
\]
Here is the cancellation explicitly.
The relative exact sequence has outer terms the cokernel of
\(H^1(G_N,M_e)\to\prod_wH^1(G_w,Q_w)/\mathcal L_w\)
and the kernel of localization in degree two.
Their cardinalities multiply to \(\#\mathcal H^2_e\).
The global odd module has rank \(2ab\), while minus transpose
has a fixed space of rank \(ab\) at each real place.
If \(d=[L^+:\mathbb Q]\), global Euler characteristic contributes
\(q^{-eabd}/\#H^0(G_N,M_e)\).
Local Euler characteristic at \(v\mid\ell\) contributes
\(q^{eab[L^+_v:\mathbb Q_\ell]}\#H^0(G_v,Q_v)\).
The degree sums are \(d\), so these powers cancel.
At the other local places, the unramified \(H^1\) cardinality
equals the \(H^0\) cardinality and contributes no remaining factor.
The finite prescribed inertia types have order prime to \(\ell\);
passing through them does not change that identity.
Discarding the positive global \(H^0\) denominator proves (B2g).

There are constants \(C_H,C_D\), independent of \(e,N\), with
\[
 \#H^0(G_v,Q_v)\leq q^{C_H},\qquad
 \#\mathcal C_e\leq q^{C_D}.
 \tag{B2h}
\]
For the first, choose for each distinct pair \(i,j\) an element
where \(\psi_{v,i}/\psi_{v,j}\ne1\).
In the lower-entry quotient, conjugation is triangular when entries
are ordered by their distance below the diagonal.
The diagonal coefficients of the invariance equations are these
nonzero character differences.
Successively eliminating the lower entries gives one nonzero
\(a_v\in A\) annihilating all invariants in every \(Q_v\).
Their cardinality is at most
\(q^{ab\,v_T(a_v)}\), uniformly in \(e\).
There are finitely many ordinary places.

For the second, global duality bounds \(\mathcal C_e\) by
\[
 \operatorname{Hom}_{\mathbb F_\ell[G_N]}(M_e,\mathbb F_\ell(1)).
\]
Localization first in the full module and then in the quotients
is legitimate because the local cohomological dimension is two.
Choose \(s\) with \(T^sE\in A[\rho(G_0')]\), where
\(E=\operatorname{diag}(1_a,0_b)\); this follows from the span
condition in (B2a).
An equivariant additive functional \(f\) satisfies
\(f(\rho(g)Y)=f(Y\rho(g))\) for \(g\in G_0'\), since the
cyclotomic character is \(1\) there.
The same holds with any \(A\)-multiple of \(Y\), so it extends to
the displayed \(A\)-linear span.
Thus
\[
                   f(T^s(EY-YE))=0 .
\]
On the upper and lower odd blocks the commutator is respectively
\(Y\) and \(-Y\).  Splitting the blocks proves
\(f(T^sM_e)=0\).
All such functionals are determined by \(M_e/T^sM_e\), and their
number is at most \(q^{2abs}\).
This proves the second assertion of (B2h), for example with
\(C_D=2abs\).

\paragraph{The obstruction quotient at every precision.}
A bound for \(\mathcal H^2_m\) alone is not enough.
For every \(e\leq m\), the two coefficient exact sequences
\[
 0\to A/T^{m-e}\xrightarrow{T^e}A/T^m\to A/T^e\to0,\qquad
 0\to A/T^e\xrightarrow{T^{m-e}}A/T^m\to A/T^{m-e}\to0
\]
give
\[
 \#(\mathcal H^2_m/T^e\mathcal H^2_m)
                         \leq\#\mathcal H^2_e\,\#\mathcal C_e.
 \tag{B2i}
\]
We spell out the maps and the defect argument.
Use the notation \(\mathcal E_r=\mathcal H^2_r\), and write
\[
 \begin{array}{c}
 f:\mathcal E_{m-e}\longrightarrow\mathcal E_m,\quad
 g:\mathcal E_m\longrightarrow\mathcal E_e,\\
 k:\mathcal E_e\longrightarrow\mathcal E_m,\quad
 r:\mathcal E_m\longrightarrow\mathcal E_{m-e}.
 \end{array}
\]
The first row is induced by the first coefficient sequence, the
second row by the second sequence.  The injections multiply
coefficients by \(T^e\) and \(T^{m-e}\), respectively, and
\[
                         f r=T^e .
 \tag{B2i1}
\]
The same coefficient maps act on the free lower-entry quotients
\(Q_w\), where both sequences are exact.
Local unramified cocycles lift under the coefficient surjections:
at an ordinary place they are boundaries, which lift by lifting
their correcting matrix; at another prescribed place first restrict to the finite
prime-to-\(\ell\) inertia quotient.
Averaging makes this restriction a boundary.
Lift its correcting matrix, subtract that boundary, and lift
the remaining Frobenius value in the inertia-invariant module;
averaging makes invariants exact under the coefficient
surjection.
A cocycle on the procyclic unramified quotient is specified
by that Frobenius value.
It is continuous for finite \(\ell\)-primary coefficients:
after an exponent of the finite Frobenius action, a further
\(\ell\)-power makes the norm sum zero.
Add back the lifted boundary.
This gives an actual lift of the original unramified cocycle,
including its finite inertial part.

It follows directly on representatives that
\(\ker g\subset\operatorname{im}f\).
Indeed, if a relative pair \(p\) maps to a relative coboundary,
lift its global correcting \(1\)-cochain and its local unramified
\(1\)-cocycles.
Subtract that lifted relative coboundary from \(p\).
Both its global and local coordinates now map to zero pointwise.
Exactness of the coefficient sequences gives unique preimages
under the first injection.  Injectivity makes those preimages
a global \(2\)-cocycle and compatible local \(1\)-cochains.
They define a class of \(\mathcal E_{m-e}\) mapping to \([p]\).
Consequently
\[
                 \#(\mathcal E_m/\operatorname{im}f)
                                    \leq\#\mathcal E_e.
 \tag{B2i2}
\]

For the other factor we give the actual defect construction.
Take a relative representative
\(p'=(u',b'_w)\) over \(A/T^{m-e}\).
Global degree-three vanishing makes
\(H^2(M_m)\to H^2(M_{m-e})\) surjective.
Hence \(u'\) lifts to an actual global \(2\)-cocycle \(u\):
first lift its cohomology class, then lift the correcting
\(1\)-cochain and subtract its differential.
Lift the local \(1\)-cochains \(b'_w\) arbitrarily to \(b_w\).
The local defects have unique expressions
\[
 \pi_w\operatorname{res}_w u-db_w
               =T^{m-e}c_w,\qquad c_w\in Z^2(G_w,Q_{w,e}).
 \tag{B2i3}
\]
The right side is interpreted using the coefficient injection
\(Q_{w,e}\hookrightarrow Q_{w,m}\).
Its injectivity and \(d^2=0\) prove that \(c_w\) is a cocycle.
The defect is measured in the \emph{actual} cokernel
\[
 \mathcal C_e=
  \left(\prod_wH^2(G_w,Q_{w,e})\right)/
       \operatorname{im}\!\left(H^2(G_N,M_e)
                       \longrightarrow\prod_wH^2(G_w,Q_{w,e})\right).
\]

The zero criterion is explicit.
If \(([c_w])_w\) is the localization of a global class represented
by \(v\in Z^2(G_N,M_e)\), choose local \(a_w\) such that
\(c_w=\pi_w\operatorname{res}_wv-da_w\).
Then
\[
       (u-T^{m-e}v,\ b_w-T^{m-e}a_w)
\]
is a relative cocycle over \(A/T^m\) mapping to \(p'\).
The same calculation applies to a difference of two relative
representatives.  Choose one representative and one system of
lifts for each element of
\(\mathcal E_{m-e}/\operatorname{im}r\).
Sending it to its defect class in \(\mathcal C_e\) is an
injective map of sets: equality of two defect classes makes
their difference a lift by the displayed correction, so the
two quotient classes agree.
This formulation needs no unproved independence of arbitrary
choices of relative representatives.  In particular,
\[
           \#(\mathcal E_{m-e}/\operatorname{im}r)
                                            \leq\#\mathcal C_e.
 \tag{B2i4}
\]
By (B2i1), \(f\) induces a surjection from this quotient onto
\(\operatorname{im}f/T^e\mathcal E_m\).
Multiplying the bounds for
\(\mathcal E_m/\operatorname{im}f\) and
\(\operatorname{im}f/T^e\mathcal E_m\), using (B2i2)--(B2i4),
proves (B2i).

Combining (B2f)--(B2i), with \(r_\ell\) the number of ordinary
places, yields, simultaneously for every \(e\leq m\) and all
sufficiently large \(N\),
\[
 \#(\mathcal H^2_m/T^e\mathcal H^2_m)
           \leq q^{ge+C},\qquad
 C=C_0+C_{\mathrm{aux}}+r_\ell C_H+2C_D.
 \tag{B2j}
\]
The constant is fixed before \(m,N\).

\paragraph{Primitive relations and the explicit second-order lift.}
Write the finite \(A\)-module \(\mathcal H^2_m\) in elementary
divisors, with lengths
\(\lambda_1\geq\lambda_2\geq\cdots\).
Equation (B2j) says
\(\sum_i\min(e,\lambda_i)\leq ge+C\).
Taking \(e=\lambda_{g+1}\), when it is nonzero, proves
\(\lambda_{g+1}\leq C\).
Thus \(T^C\mathcal H^2_m\) is generated by at most \(g\) elements.
For any \(g+1\) elements \(\xi_i\) of \(\mathcal H^2_m\), lift
their coordinates in such a generating set to a \(g\) by \(g+1\)
matrix over \(A\).  Its kernel over \(F\) has a nonzero vector.
Clear denominators and divide by its common \(T\)-power.
The resulting \(c_i\in A\) include a unit and satisfy
\[
                         T^C\sum_i c_i\xi_i=0 .
 \tag{B2k}
\]

For the even Selmer tuples \(z_{ij}\) in \srcid{B2}, use the
actual square-zero coefficient algebra
\[
 \mathbb D_g(A/T^m)
  =(A/T^m)\oplus\bigoplus_j\epsilon_j(A/T^m)
    \oplus\bigoplus_j\epsilon'_j(A/T^m)\oplus\eta(A/T^m),
 \quad\epsilon_i\epsilon'_j=\delta_{ij}\eta.
\]
The fixed first-order representation is
\((1+\sum_j\epsilon_j\kappa_j)\rho\).
The extension of coefficient modules governing the next step is
\[
 0\to\eta\,\operatorname{Odd}_a(A/T^m)
       \longrightarrow W_2
       \longrightarrow
       \bigoplus_j\epsilon'_j\operatorname{Even}_a(A/T^m)\to0.
\]
For the \(i\)-th tuple its connecting obstruction is the relative
class \(\xi_i\) just defined.
Its global two-cocycle is, up to the sign convention for \(d^1\),
\[
 (g,h)\longmapsto
        \sum_j[\kappa_j(g),\operatorname{Ad}(\rho(g))z_{ij}(h)].
 \tag{B2l}
\]
This formula follows immediately by multiplying the two
\(\mathbb D_g\)-valued matrices.
The local component is obtained by applying the same differential
to the chosen even local correcting matrices, using the fixed
odd correcting matrices for the \(\kappa_j\).
Its lower-entry projection has differential equal to the
restriction of (B2l), so it is a pair defining \(\mathcal H^2_m\).
This is a cochain construction, not an assumed obstruction
vanishing theorem.

When a linear combination of the relative classes vanishes,
a global odd one-cochain and local odd correcting matrices
trivialize the combination.
Putting that cochain in the \(\eta\)-coordinate gives
\[
 \Psi(g)=
    \sum_j\epsilon'_j\!\left(T^C\sum_i c_i z_{ij}(g)\right)
                           +\eta\,y(g).
\]
The cocycle equation is exactly cancellation of (B2l).
Complex-conjugation compatibility is obtained by averaging with
its polarized transform; \(2\) is a unit.
At ordinary places the same local correcting matrices make every
lower entry zero.  The required diagonal entries on the specified
inertia also remain zero: in the ordinary basis the relevant
commutators are commutators of upper-triangular matrices and have
zero diagonal.
At other prescribed inertia the relative unramified condition
makes the cocycle a boundary there.  Its action is trivial, by
\(\rho(I_u)=1\) and \(\kappa_j(I_u)=0\), so that boundary is zero.
Thus \(\Psi\in\operatorname{Sel}_{2,N}(A/T^m)\) and
\[
       \epsilon'_j(\Psi)=T^C\sum_i c_i z_{ij},\qquad
                    \text{some }c_i\in A^\times .
\]
This proves \srcid{TWGrp.B2} with one \(C\) before every
coefficient precision and every tower level.

\paragraph{Scope and exact sources.}
The proof applies in every rank and to both paired blocks; for
\(a=n\) the odd module is zero and the same argument simplifies.
It proves the actual finite-level arithmetic second-Selmer bound
from the displayed nice-point, original-generator, standard
arithmetic cohomology, and Taylor--Wiles group properties.
It does not silently manufacture the regular cosets or the
Taylor--Wiles tower; these are the separate prime-selection
inputs in (B2b).

All source paths below are at the frozen tag
\srcid{alpoge-source-7cbf7a0c}, under
\srcid{Development.AutomorphicArithmetic.AutomorphicVarieties}.
The uniform approximation proof is
\srcid{PatchingModules.DegreeCountCohomology.SolCnt}
(\srcid{scS\_stab}, \srcid{scS\_lim}, \srcid{sc\_kerfin},
\srcid{sc\_unif}, \srcid{solCnt}).
Its projected original odd-Selmer application is
\srcid{ArchimedeanH0Cards.ImgCard.sel\_cnt} and \srcid{sel0}.
The inertia comparison is
\srcid{ArchimedeanH0Cards.ExistsCore.odd\_kill},
\srcid{fix\_count}, and \srcid{selN}.
The local invariant bound and dual invariant bound are
\srcid{H0LocalCardBoundedOfClsKerLeKer} and \srcid{DualBound}.
The relative group is defined in
\srcid{PushforwardPreservesSelmerCocycles.Dat.H2LocalConditions};
its cardinal formula is \srcid{OdCount.odCount}.
The two exact-sequence chase is
\srcid{ObstructionBoundOfDuality.step}, using
\srcid{ExactMid} and \srcid{CokerBound}; primitive relations are
\srcid{GenTrick.genTrick}.
The explicit lift from those relations is
\srcid{Transl.transl}; the final arithmetic specialization is
\srcid{AutomorphicDetection.TaylorWilesObstructionBound.b2Ar}.
These file paths are below
\url{https://github.com/ChristianKudeba/alpoge-rh-verification/tree/alpoge-source-7cbf7a0c/src/Development/AutomorphicArithmetic/AutomorphicVarieties}.

\subsection{Trace detection and the arithmetic first-Selmer bound}
\label{sec:arithmetic-first-selmer}

We prove the actual first-order bound \srcid{TWGrp.B1}.
In particular the finite list of Taylor--Wiles elements is
constructed, rather than postulated to detect all dual classes.

Write \(A=k[[T]]\), \(F=k((T))\), \(Q=\#k\), and
\(\operatorname{char}k=\ell\ne2\).
Let \(\rho:G_0\to\mathrm{GL}_n(A)\) be one of the nice paired
points, with block sizes \(a,b=n-a\).
Put
\[
 E=M_a(F)\oplus M_b(F),\qquad
 O=\operatorname{Hom}(F^b,F^a)\oplus
                         \operatorname{Hom}(F^a,F^b).
\]
The precise global hypotheses used here are
\[
 F[\rho(G_0')]=E,\qquad
 G_0'=\ker(\varepsilon|_{G_0}),\qquad
 \exists g_*\in G_0:\bar\varepsilon(g_*)\ne1,
 \tag{B1a}
\]
together with polarization and the original odd generators
\(\kappa_1,\ldots,\kappa_g\).
Their independence means that a linear combination which is
the boundary of an odd skew matrix has all coefficients zero.
Their original spanning property is the one already used in
the second-Selmer proof.
At \(v\mid\ell\), the integral ordinary bases respect the block
type, and the ordered characters satisfy
\[
 \psi_{v,i}\ne\psi_{v,j}\ (i\ne j),\qquad
 \psi_{v,i}\ne\psi_{v,j}\bar\varepsilon_v\ (i<j).
 \tag{B1b}
\]
The prescribed local torsion is killed by the diagonal
characters.  At the other prescribed inertia groups,
\(\rho=1\) and \(\kappa_j=0\).
We use the finite arithmetic cohomology, local and global duality,
and Euler characteristics established previously.

Let \(D_0\subset G_0\) be the distinguished local image.
The remaining nice-point hypothesis used in this bound is exactly
\[
 \{Y\in M_n(F):
    \bar\varepsilon(g)\rho(g)Y\rho(g)^{-1}=Y
                         \text{ for all }g\in D_0\}=0 .
 \tag{B1c}
\]
This is \srcid{N4}, an explicitly specified local condition.
We do not replace it by irreducibility of \(\rho\).

\paragraph{The actual dual coefficient module and its traces.}
Use the square-zero coefficient algebra
\[
 {\mathbb D}_g=A\oplus\bigoplus_j\epsilon_jA
                   \oplus\bigoplus_j\epsilon'_jA\oplus\eta A,
 \qquad \epsilon_i\epsilon'_j=\delta_{ij}\eta,
\]
with all remaining products of positive-degree generators zero.
The fixed first-order representation is
\[
                   \rho_1=(1+\sum_j\epsilon_j\kappa_j)\rho .
\]
The primal and dual modules are
\[
 W_1=E_A\oplus\bigoplus_j\epsilon_jO_A,\qquad
 W_1^\vee=\eta E_A\oplus\bigoplus_j\epsilon'_jO_A.
 \tag{B1d}
\]
Their pairing is the \(\eta\)-coefficient of \(\operatorname{tr}(XY)\).
It is perfect over \(A\).
The \(G_0\)-action on the dual module is
\(\bar\varepsilon\,\operatorname{Ad}\rho_1\); its extension across
complex conjugation is minus the cyclotomic multiplier times
transpose.
For a continuous polarized twisted cocycle \(z\), put
\[
             \ell_s(z)=
       \eta\operatorname{tr}\bigl(z(s)\rho_1(s)\bigr)
                         \quad(s\in G_0').
 \tag{B1e}
\]
These are the actual trace functionals used by the construction.

We first prove that their common zero space over \(F\) consists
of boundaries.
Write the coordinates of \(z\) as
\[
                 z(g)=\eta\varphi(g)
                              +\sum_j\epsilon'_j\lambda_j(g),
 \qquad \chi=\bar\varepsilon .
\]
Then
\[
 \begin{split}
 \lambda_j(gh)&=\lambda_j(g)+
                  \chi(g)\operatorname{Ad}\rho(g)\lambda_j(h),\\
 \varphi(gh)&=\varphi(g)+\chi(g)\left(
       \operatorname{Ad}\rho(g)\varphi(h)+
       \sum_j[\kappa_j(g),
                     \operatorname{Ad}\rho(g)\lambda_j(h)]\right).
 \end{split}
 \tag{B1f}
\]
The odd \(\lambda_j\) and even \(\varphi\) have the same transpose
polarization sign \(s\), with \(s^2=1\).
For \(g\in G_0'\), vanishing of (B1e) is exactly
\[
       \operatorname{tr}\left(
        \varphi(g)\rho(g)+
                \sum_j\lambda_j(g)\kappa_j(g)\rho(g)\right)=0 .
 \tag{B1g}
\]

\paragraph{A relation calculation that detects the odd coordinates.}
Here are the algebraic steps behind this assertion.
Let \(R:F[G_0']\to E\) be induced by \(\rho\), and for
\(c=\sum c_g[g]\) write
\[
 K_j(c)=\sum_gc_g\kappa_j(g)\rho(g),\qquad
 L_j(c)=\sum_gc_g\lambda_j(g)\rho(g).
\]
The cocycle equations make \(K_j\) and \(L_j\) derivations
relative to \(R\) on this group algebra.
Define also
\[
 P(c)=\sum_gc_g\left(\varphi(g)\rho(g)+
                     \sum_j\lambda_j(g)\kappa_j(g)\rho(g)\right).
\]
Multiplying (B1f) gives the explicit product formula
\[
 P(cd)=P(c)R(d)+R(c)P(d)+
                  \sum_j\bigl(K_j(c)L_j(d)+L_j(c)K_j(d)\bigr).
 \tag{B1h}
\]
We have \(\operatorname{tr}P(c)=0\) by (B1g).
For \(c,d\in\ker R\) and opposite-block indices \(i,j\), insert
preimages under \(R\) of the two diagonal matrix units
\(e_{jj}\) and \(e_{ii}\) in \(\operatorname{tr}P(c e_{jj}d e_{ii})\).
Surjectivity of \(R\), which is (B1a), permits these choices.
Every term involving \(R(c)\) or \(R(d)\) vanishes.
Formula (B1h) leaves precisely
\[
 \sum_\nu\left(
      K_\nu(c)_{ij}L_\nu(d)_{ji}
                 +L_\nu(c)_{ij}K_\nu(d)_{ji}\right)=0 .
 \tag{B1i}
\]

We need the following elementary consequence of relation vanishing.
Suppose \(K\) is an untwisted odd cocycle and \(L\) is a
\(\chi\)-twisted odd cocycle.  If
\[
       \left(\sum_gc_g(K(g)+L(g))\rho(g)\right)_{ij}=0
                                  \quad(c\in\ker R)
 \tag{B1j}
\]
for one opposite-block pair \(i,j\), their corresponding entire
corner is principal on \(G_0\).
First, multiply the relation on the left and right by images
of elements of \(G_0'\).
The cocycle equation shows that the extra terms contain
\(R(c)=0\); hence (B1j) is preserved.
The matrix units in \(E\) move the indices to any entry in that
corner.
Thus the corner of \((K+L)(g)\rho(g)\) factors through a derivation
of \(E\) into that corner bimodule.

Every derivation of \(E=M_a(F)\oplus M_b(F)\) into \(M_n(F)\)
is inner.  An explicit formula is useful here.
Choose in each nonempty block an anchor index \(\beta(i)\),
put \(u_i=e_{i,\beta(i)}\), \(v_i=e_{\beta(i),i}\), and for a
derivation \(d\) take \(Y=\sum_i d(u_i)v_i\).
The identities
\(\sum_i u_iv_i=1\) and the Leibniz rule give
\(d(e_{ij})=Ye_{ij}-e_{ij}Y\) whenever \(i,j\) are in one block.
Linearity proves the assertion for all of \(E\).
After changing the sign of \(Y\), we obtain
\[
                  K+L=\delta Y\quad\hbox{on }G_0'
\]
for the chosen corner.

The two summands separate because \(\chi(g_*)\ne1\).
For completeness, if \(c_*=\chi(g_*)\), direct expansion of
the cocycle equations at \(g g_*\) and \(g_*g_*^{-1}gg_*\)
gives
\[
 \begin{split}
 Y_1=(c_*-1)^{-1}\bigl(
  c_*\operatorname{Ad}\rho(g_*)Y
             -c_*K(g_*)-L(g_*)-Y\bigr),\\
             K=\delta Y_1,\qquad L=\delta(Y-Y_1)
                                \quad\hbox{on }G_0'.
 \end{split}
 \tag{B1k}
\]
Here \(K,L\) denote their chosen corner projections.
Project the correcting matrices to that corner.
The errors from these equations, extended to \(G_0\) with their
respective twists, vanish on \(G_0'\).
Their values therefore centralize \(\rho(G_0')\).
An odd matrix centralizing \(E\) is zero, because it commutes
with the two block idempotents.
Thus the corner equations hold on all of \(G_0\).

Polarization reconstructs the opposite corner.
If \(K(\sigma g)=-K(g)^{\mathsf t}\), replace its corner
correcting matrix \(Y\) by \(Y-Y^{\mathsf t}\).
If \(L(\sigma g)=sL(g)^{\mathsf t}\), replace its corner
correcting matrix by \(Y+sY^{\mathsf t}\).
The identity
\[
        \delta_\chi(Y)(\sigma g)
                        =\delta_\chi(Y^{\mathsf t})(g)^{\mathsf t}
\]
follows from the polarized equation for \(\rho\).
These replacements give the full odd boundaries with the
required signs.

Apply this consequence to
\(K=\sum_\nu\alpha_\nu\kappa_\nu\),
\(L=\sum_\nu\beta_\nu\lambda_\nu\).
If (B1j) holds, the preceding argument makes \(K\) an odd skew
boundary; independence of the \(\kappa_\nu\) implies
\(\alpha_\nu=0\) for every \(\nu\).
In (B1i), fix \(d\in\ker R\), and put
\[
                 \alpha_\nu=L_\nu(d)_{ji},\qquad
                 \beta_\nu=K_\nu(d)_{ji}.
\]
Equation (B1i) is exactly (B1j), so all
\(L_\nu(d)_{ji}\) vanish.
Now each \(\lambda_\nu\) kills relations in the opposite entry.
Applying the same argument with \(\alpha=0\) and
\(\beta_\nu=1\), the other \(\beta\)'s zero, proves
\[
        \lambda_\nu=\delta_\chi Y_\nu,\qquad
                Y_\nu\in O,\quad sY_\nu^{\mathsf t}=Y_\nu .
 \tag{B1l}
\]
If \(a=n\), the odd module is zero and this paragraph is vacuous.

\paragraph{The even coordinate and the full dual boundary.}
Subtract from \(\varphi(g)\) the explicit cross term
\[
             \chi(g)\sum_\nu
                  [\kappa_\nu(g),\operatorname{Ad}\rho(g)Y_\nu].
\]
Call the result \(\varphi_0(g)\).
Equations (B1f) and (B1l) show, by expansion, that it is an
ordinary \(\chi\)-twisted even cocycle.
Cyclicity of trace and (B1g) give
\(\operatorname{tr}(\varphi_0(g)\rho(g))=0\) on \(G_0'\).
For a relation \(c\in\ker R\), put
\(M=\sum c_g\varphi_0(g)\rho(g)\in E\).
Apply the trace equation to \(gh\) and sum over \(c\).
It gives \(\operatorname{tr}(M\rho(h))=0\) for all \(h\in G_0'\).
The trace pairing on \(E\) is nondegenerate, so \(M=0\).
Thus \(\varphi_0(g)\rho(g)\) factors through an even derivation
of \(E\), which is inner by the same matrix-unit calculation.

Subtract the resulting even boundary.
The error vanishes on \(G_0'\), so its values centralize \(E\)
and therefore commute with every \(\rho(g)\).
The quotient \(G_0/G_0'\) is abelian.
Its two cocycle expansions at \(g g_*\) and \(g_*g\) yield
\[
      (c_*-1)\varphi_{\rm err}(g)
                =(\chi(g)-1)\varphi_{\rm err}(g_*).
\]
This is the boundary of
\((c_*-1)^{-1}\varphi_{\rm err}(g_*)\).
Averaging the correcting matrix with \(s\) times its transpose
makes it have the required polarization, since \(2\) is a unit.
Together with (B1l), this produces
\[
 Y=\eta Y_0+\sum_\nu\epsilon'_\nu Y_\nu\in W_1^\vee\otimes_A F,
 \qquad
             z(g)=\chi(g)\operatorname{Ad}\rho_1(g)Y-Y .
 \tag{B1m}
\]
Clearing the finitely many matrix denominators proves that an
integral \(z\) annihilated by all \(\ell_s\), \(s\in G_0'\),
has a \(T\)-power in the integral polarized boundary module.
This is the precise assertion called \srcid{Detects}.

\paragraph{Choosing finitely many simultaneous regular detectors.}
The module of continuous twisted cocycles with values in a fixed
finite free \(A\)-module is finite over \(A\).
Indeed its reduction modulo \(T\) injects into the finite set
of residual continuous cocycles: the kernel of reduction is
\(T\) times the integral cocycle module.
Dividing a cocycle by \(T\) preserves continuity by testing
one higher coefficient precision.
The module is complete in the uniform \(T\)-adic topology,
so successive approximation, or topological Nakayama, proves
finite generation.
Residual cocycle finiteness follows from \srcid{PhiCond}:
on the open kernel of the residual action a cocycle is a
homomorphism into a finite elementary \(\ell\)-group, and
there are finitely many choices on its finite complement.
This argument applies to the full matrix coefficient module
over \(\mathbb D_g\), and hence to all its submodules.

For finitely many nice points, let \(\mathcal S\subset G_0'\)
be the set where every point has separable characteristic
polynomial.
For each point this regular set is open: the discriminant is
a continuous \(A\)-valued function, and its nonvanishing is
open by separatedness of \(A\).
The actual regular-coset property \srcid{N3} says this open
set meets every nonempty relatively open part of \(G_0'\).
Finite induction, retaining the regular open sets already
obtained, shows that \(\mathcal S\) is dense in \(G_0'\).

For each integral cocycle \(z\), the function
\(s\mapsto\ell_s(z)\) is continuous.
Hence vanishing on \(\mathcal S\) implies vanishing on all
of \(G_0'\), and (B1m) applies.
On the finite \(A\)-module of integral cocycles, the functionals
\(\ell_s\), \(s\in\mathcal S\), span a finite submodule of its
\(A\)-dual.
Select finitely many which generate that span.
Their common kernel is therefore the common kernel of all
the functionals.
Do this for every point and take the finite union.
It produces one list \(s_1,\ldots,s_r\in\mathcal S\) which
detects every point, all of whose characteristic polynomials
are separable at every point.
The finite DVR extensions described in the tower subsection
split all these polynomials and give the same residue-field
cardinality across the finite family.
This proves the actual choice input.

\paragraph{A uniform bound on the dual Selmer factor.}
The trace calculation over \(F\) alone does not bound classes
which exist only modulo \(T^m\).
Apply the uniform approximation lemma from the second-Selmer
proof to continuous cochains, with defect map listing
their cocycle, polarization, and the finitely many
\(\ell_{s_j}\) equations.
Take \(B\) to be the \(A\)-submodule of polarized boundaries.
The hypothesis that every exact solution has a \(T\)-power
in \(B\) is (B1m); residual solutions are finite by
\srcid{PhiCond}.
The lemma consequently gives one \(C_{\rm tr}\), before all \(m\),
such that
\[
 \#\{z\in\operatorname{TwZ}(A/T^m):
                     \ell_{s_j}(z)=0\ \forall j\}
       \leq \#\operatorname{TwB}(A/T^m)\,Q^{C_{\rm tr}} .
 \tag{B1n}
\]

Use the actual arithmetic tower constructed in this section.
For fixed \(m\) and all sufficiently large \(N\), a full
polarized dual cocycle locally a boundary at each new place
has value zero at \(\tau_{Nj}\), because \(\tau_{Nj}\) acts
trivially on the coefficient module.
Its finite semidirect-product image kills
\(\mathrm{PL}_\ell(H_N)\), as proved in the tower construction.
The normal-kernel assertion (TW10) makes it descend to
\(\Gamma_0\).
Normalize its value at complex conjugation to zero by a
boundary, using \(2^{-1}\), and restrict to \(G_0\).

At finite precision there are only finitely many such
base cocycles, not just finitely many point matrices.
Intersect their finite semidirect-product kernels, along
with the kernels of \(\rho_1\) and the cyclotomic character
at that precision.
This open subgroup contains \(\mathrm{PL}_\ell(H)\).
Frobenius convergence (TW8) therefore identifies their
values at \(h_N(\phi_{Nj})\) and at \(s_j\), simultaneously.
The boundary equation at \(\phi_{Nj}\), multiplied by
\(\rho_1(s_j)\) and traced, gives
\[
 \eta\operatorname{tr}\bigl(
   (\operatorname{Ad}\rho_1(s_j)Y-Y)\rho_1(s_j)\bigr)=0 .
\]
Thus (B1n) applies.
Normalization at complex conjugation has exactly the same
fibers for cocycles and boundaries; canceling those fibers
gives the actual dual bound
\[
       \#H^1_{\mathcal L^\perp,N}
                    (W_1^\vee/T^m)\leq Q^{C_{\rm tr}},
 \tag{B1o}
\]
when at the auxiliary places the dual local condition is
being a boundary.
For the global counting inequality it suffices to impose
only these auxiliary local conditions, so any further
orthogonality conditions can only decrease this cardinality.

\paragraph{The local ordinary and invariant counts.}
Put
\[
 \nu=\frac{a(a-1)}2+\frac{b(b-1)}2+gab .
\]
Then \(W_1\) has rank \(n+2\nu\), and its upper-triangular
submodule in the actual ordinary basis has rank \(n+\nu\).
For an ordinary place \(v\), let
\(d_v=[L^+_v:\mathbb Q_\ell]\), and let \(d'_v\) be the free
\(\mathbb Z_\ell\)-rank of the local unit quotient used by the
ordinary weight algebra.
Local class field theory gives
\[
                         \sum_{v\mid\ell}d_v
                        =\sum_{v\mid\ell}d'_v=d,
             \qquad d=[L^+:\mathbb Q].
 \tag{B1p}
\]
The local abelian quotient, with the prescribed finite torsion
removed, has rank \(d'_v+1\), including its unramified factor.

Let \(Z_{b,v,m}\) be the actual upper-triangular cocycles,
with diagonal zero on that prescribed torsion.
Their diagonal coordinates are \(n\) continuous characters
of this abelian quotient.
They have at most \(Q^{mn(d'_v+1)}\) choices.
Filter the strictly upper-triangular coordinates by their
distance from the diagonal, retaining the even block entries
and the \(g\) odd copies.
There are \(\nu\) root coordinates.
Each associated rank-one character is
\(\psi_{v,i}/\psi_{v,j}\), \(i<j\).
Its degree-two local cohomology is uniformly bounded in \(m\):
local duality and (B1b) identify it with the dual of invariants
of a nontrivial character.
A fixed nonzero character difference annihilates those
invariants at every precision.
Local Euler characteristic, together with the elementary
boundary count, therefore bounds its cocycles by
\(Q^{m(d_v+1)+C_{ij}}\).
Successively choosing the filtered coordinates never gives
more choices than the product of these cocycle counts.
Thus, with \(C_v\) fixed before \(m\),
\[
 \#Z_{b,v,m}\leq
       Q^{\,m[n(d'_v+1)+\nu(d_v+1)]+C_v}.
 \tag{B1q}
\]
This counts the actual ordinary cocycles; the fixed local
correctors in \(\rho_1\) enter by the integral basis change.

There are also constants \(C_G,C_*,C_j\) such that
\[
 \begin{array}{rl}
 \#H^0(\Gamma_N,W_1/T^m)&\leq Q^{C_G},\\
 \#H^0(D_0,W_1^\vee/T^m)&\leq Q^{C_*},\\
 \#H^0(\langle s_j\rangle,W_1^\vee/T^m)
                                      &\leq Q^{mn+C_j}.
 \end{array}
 \tag{B1r}
\]
These assertions follow by finite matrix equations and Smith
normal form, as follows.
A \(G_0'\)-invariant constant even matrix centralizes \(E\),
hence is scalar on each block.
The minus-transpose equation at complex conjugation kills
both scalars.
The remaining odd components centralize \(E\) and vanish.
Thus the first invariant space is zero after inverting \(T\).
For the second, first apply (B1c) to every \(\epsilon'\)-coordinate,
then to the remaining \(\eta\)-coordinate; the cross terms
disappear, so its invariant space also vanishes over \(F\).
For the third, the odd components centralizing the regular
\(\rho(s_j)\) vanish.
Indeed the centralizer of a matrix with distinct eigenvalues
is commutative and contains the block idempotent.
The remaining even centralizer has dimension \(n\).
A fixed finite set of the defining linear equations has these
same kernels over \(F\).
Their Smith factors bound finite-precision kernels by
\(Q^C\) times \(Q^{m}\) for each surviving generic dimension,
which gives exactly (B1r).
The Frobenius approximation in the tower transports these
bounds to the actual new local groups for large \(N\).

\paragraph{The finite duality count and the exponent cancellation.}
Let \(\operatorname{Sel}_{1,N,m}\) and
\(\operatorname{Bd}_{1,N,m}\) be the actual primal cocycle and
boundary sets.
We give the count with every term visible.
At an ordinary place take local cocycles equal to a boundary
plus an element of \(Z_{b,v,m}\).
The map from a correcting matrix and an upper cocycle onto
this sum has fibers containing the upper-triangular module,
of cardinality \(Q^{m(n+\nu)}\).
Thus, writing \(M_m=W_1/T^m\),
\[
       \#\mathcal L^Z_v
               \leq \#M_m\,
                       \#Z_{b,v,m}/Q^{m(n+\nu)} .
 \tag{B1s}
\]
At an old non-\(\ell\) place the unramified cocycle set has
cardinality \(\#M_m\), since inertia acts trivially.
At the distinguished and the new places take the full local
cocycle set.
Local duality at a non-\(\ell\) place gives its cardinality
as \(\#M_m\) times the cardinality of the dual invariants.
These are bounded by (B1r).

Finally each real place has at least \(Q^{m\nu}\) invariants
in \(M_m\).
For the chosen complex conjugation these are its skew matrices,
of rank \(\nu\).
For another real conjugation, the polarized representation
identifies its fixed space with that skew space by multiplication
by the corresponding symmetric invertible matrix.
This works integrally, so the same lower bound holds at
finite precision.

Apply the finite global-duality cardinal formula, using the
perfect pairing in (B1d), the local subgroups just specified,
and (B1o).
In cocycle notation the local factors are
\(\#\mathcal L^Z_w/\#M_m\); the real factors occur in the
denominator, and the global primal invariants in the numerator.
The denominator of global dual invariants can be discarded
for an upper bound.
Normalization at complex conjugation identifies the primal
class quotient with
\(\operatorname{Sel}_{1,N,m}/\operatorname{Bd}_{1,N,m}\),
because \(2\) is invertible.
Using (B1q)--(B1s) therefore gives
\[
 \begin{split}
 \#\operatorname{Sel}_{1,N,m}\,
   Q^{m[(n+\nu)\#\{v\mid\ell\}+\nu d]}
 \ \leq\ 
 \#\operatorname{Bd}_{1,N,m}\,
 Q^{C_G+C_{\rm tr}+C_*+\sum_vC_v+\sum_jC_j}\\
 {}\cdot
 Q^{\,m\sum_{v\mid\ell}
        [n(d'_v+1)+\nu(d_v+1)]+mnr}.
 \end{split}
 \tag{B1t}
\]
All cohomology and cocycle sets here are finite by the
arithmetic finiteness input.
The auxiliary level \(N\) only needs to exceed a threshold
depending on \(m\).

The identity (B1p) gives the exact cancellation
\[
 \sum_v[n(d'_v+1)+\nu(d_v+1)]+nr
       -[(n+\nu)\#\{v\mid\ell\}+\nu d]=nd+nr .
\]
Hence one constant \(C\), fixed before every precision and
tower level, satisfies
\[
 \boxed{\ 
 \#\operatorname{Sel}_{1,N,m}
 \leq\#\operatorname{Bd}_{1,N,m}\,
                       Q^{\,m(nd+nr)+C}
 \quad\text{for all sufficiently large \(N\), for each \(m\)}.
       \ }
 \tag{B1u}
\]
This is \srcid{TWGrp.B1} with its actual parameter
\(D_0^{\rm bound}=nd+nr\).
The point-dependent bounded constants affect only the
precision threshold and additive constant; the displayed
linear slope is explicit.

\paragraph{Exact sources and what this closes.}
All paths in this paragraph are under
\srcid{Development.AutomorphicArithmetic.AutomorphicVarieties}
at \srcid{alpoge-source-7cbf7a0c}.
The odd relation calculation is
\srcid{PatchingModules.DeterminantPadicFree.RelPair.relPair},
\srcid{RelMove.relMove}, \srcid{DTw.der},
\srcid{DerEv.derEv}, \srcid{DTwAlt.eigen},
\srcid{Infres.infres}, \srcid{Shapiro.shapiro},
and \srcid{BlockPrinc.ann}; these are assembled in
\srcid{OddPart.oddPart}, \srcid{EvenPart.evenPart},
and \srcid{DetAlg.detAlg}.
Integral denominator clearing is \srcid{DetAsm.detAsm}.
The actual dense and finite choices are
\srcid{AutomorphicDetection.Detection.detection},
\srcid{DeterminantPadicFree.MemPowSmul.z1Fg},
\srcid{AcConst.matFg}, \srcid{FinSub.finSub},
\srcid{DetFin.detFin}, \srcid{OmDense.omDense},
and \srcid{AutomorphicDetection.Choice.choice}.
Finite-precision detection is
\srcid{EpsilonEtaMatrices.ExistsNcardAnnihilatedCocyclesLeCoboundariesMul};
descent and normalization are
\srcid{Descend} and \srcid{DualN.dualN}.
The local and invariant estimates are
\srcid{EntryZero}, \srcid{Rho1MapReductionLevel},
\srcid{FinrankCentralizerW1Le}, and the local ordinary-count
modules imported by \srcid{OneLevel}.
The finite duality count is \srcid{OneLevel.oneLevel};
its arithmetic assembly is
\srcid{W1BoundArithmeticRankTwo.W1_bound_arithmetic_rank_two}
and \srcid{AutomorphicDetection.BoundW1Arithmetic.b1Ar}.
These sources are in
\url{https://github.com/ChristianKudeba/alpoge-rh-verification/tree/alpoge-source-7cbf7a0c/src/Development/AutomorphicArithmetic/AutomorphicVarieties}.

Thus the first- and second-Selmer estimates and the arithmetic
Taylor--Wiles group tower have explicit proofs once the actual
nice points, their regular-coset property, and the established
arithmetic cohomology are in place.
This result supplies those numerical and group-theoretic
producers to the completed-support construction.
It does not identify an arbitrary polarized representation
with an automorphic point; that realization remains the
separate classical and seed/support part of the argument.

\section{A constructed ordinary Eisenstein seed}
\label{sec:constructed-eisenstein-seed}

We construct the classical seed used by the rank-three branch.
The construction below produces actual cuspidal eigenforms.
Its last paragraph distinguishes this existence statement from
placing the seed in the required localized unitary support.

Let \(\ell\ge5\), fix an embedding of the algebraic coefficient
fields into \(\overline{\mathbb Q}_\ell\), and write
\(\mathfrak m_\ell\) for the ideal of elements of positive
\(\ell\)-adic valuation.
Take a primitive nontrivial Dirichlet character \(\varphi\)
of conductor \(M>1\), with \(\ell\nmid M\).
Choose a prime \(w\nmid M\) and a nontrivial character
\(\theta\) of conductor \(w\) such that
\[
       w\equiv1\pmod\ell,\qquad
       \varphi(w)=1,\qquad \theta^\ell=1 .
 \tag{ES1}
\]
Set \(R=Mw\) and \(\chi=\varphi\theta\), interpreted at level
\(R\) by the Chinese remainder theorem.
It is primitive: the two conductors are coprime and each
factor has its stated conductor.
In particular \(\ell\nmid R\).

We shall prove that, for every sufficiently large \(k\) with
\(\chi(-1)=(-1)^k\), there is a normalized cuspidal eigenform
\(f\) of weight \(k\), level \(R\), and character \(\chi\), with
\[
          a_p(f)\equiv1+\chi(p)p^{k-1}\pmod{\mathfrak m_\ell}
                       \quad\hbox{for every prime }p .
 \tag{ES2}
\]
This includes \(p\mid R\) and \(p=\ell\).
The latter congruence makes \(f\) ordinary.

\subsection{Integral forms with a prescribed constant term}
For \(r\ge1\), let \(P_r\in\mathbb Z[X]\) be defined by
\[
 P_1=X,\qquad
 P_{r+1}=X((1-X)P'_r+rP_r),\qquad
                    L_r(x)=\frac{P_r(x)}{(1-x)^r}.
\]
For \(|x|<1\), this is
\(\sum_{n\ge1}n^{r-1}x^n\); differentiation proves the recursion.
In particular
\[
 L_3(x)=\frac{x(1+x)}{(1-x)^3},\qquad
 L_5(x)=L_3(x)\left(1+12\frac{x}{(1-x)^2}\right).
 \tag{ES3}
\]
If \(x\ne1\) has finite order prime to \(\ell\), then
\(1-x\) is an \(\ell\)-adic unit.
For example, reduction preserves all roots of unity of
prime-to-\(\ell\) order: the geometric-sum identity shows
that two such roots with equal residue are equal.
Thus every \(L_r(x)\) is integral.
If also \(x^2\ne1\), then \(L_3(x)\) is a unit.

We record a separation calculation for these constant terms.
If \(x,y\) are \(R\)-th roots, both primitive, and
\(y\ne x,x^{-1}\), put
\[
                    u=\frac{x}{(1-x)^2},\qquad
                    v=\frac{y}{(1-y)^2}.
\]
Both are units, and
\[
 u-v=\frac{x(1-y/x)(1-xy)}{(1-x)^2(1-y)^2}
\]
is a unit as well.
Let \(A=L_3(x)^2\), \(B=L_3(y)^2\).
If \(A-B\) is a unit it already separates \(x,y\).
Otherwise (ES3) gives
\[
 L_3(x)L_5(x)-L_3(y)L_5(y)
       =12A(u-v)+(A-B)(1+12v).
\]
Its first term is a unit and its second term lies in
\(\mathfrak m_\ell\), since \(\ell\ge5\).
Their sum is a unit.
Consequently one of the constant-term functions
\(L_3^2\), \(L_3L_5\), of weights \(6,8\), separates any
two primitive roots not equal up to inversion.
This is the generic-prime specialization of the source's
weight-\(4,6,8\) calculation.

Here are the actual modular forms to which we apply it.
For \(a\in(\mathbb Z/R)^\times\) and \(r\ge3\), set
\[
 {\cal E}_{r,a}(z)=
 \frac{R^r}{(-2\pi i)^r/(r-1)!}
       \sum_{\substack{(c,d)\in\mathbb Z^2\\
                 c\equiv0\ (R),\ d\equiv a\ (R)}}
                                    (cz+d)^{-r}.
 \tag{ES4}
\]
The sum is absolutely and locally uniformly convergent.
Changing the summation variables proves its
\(\Gamma_1(R)\)-modularity, and for
\(\gamma\in\Gamma_0(R)\),
\[
             {\cal E}_{r,a}|_r\gamma
                           ={\cal E}_{r,a\gamma_{22}}.
 \tag{ES5}
\]
The same convergence estimate after a cusp change proves
holomorphy at all cusps.
If \(R\nmid\gamma_{21}\), the constant term of
\({\cal E}_{r,a}|_r\gamma\) is zero.
Indeed a constant term would require a lattice vector whose
transformed first coordinate is zero; the two congruences in
(ES4), with \(a\) a unit, would then force
\(R\mid\gamma_{21}\).
Equivalently these forms vanish at every cusp not equivalent
to infinity under \(\Gamma_0(R)\).

Writing \(\zeta=e^{2\pi i/R}\), the elementary Lipschitz
expansion of the inner sum gives
\[
 \begin{split}
 a_0({\cal E}_{r,a})&=L_r(\zeta^a),\\
 a_n({\cal E}_{r,a})&=
   \sum_{d\mid n}d^{r-1}
       \bigl(\zeta^{ad}+(-1)^r\zeta^{-ad}\bigr),\qquad n\ge1 .
 \end{split}
 \tag{ES6}
\]
One obtains this by summing first over the second lattice
coordinate, expanding the resulting geometric series for
positive first coordinate, and pairing negative first
coordinates with positive ones.
Absolute convergence for \(r>2\) justifies these operations.
These formulas prove integrality of every Fourier coefficient.
They also prove stability under every automorphism of
\(\mathbb C\): such an automorphism replaces \(\zeta\) by
\(\zeta^s\), and hence replaces the family by
\({\cal E}_{r,sa}\).
No invariance of the transcendental normalizing factor in
(ES4) is asserted; (ES6) is the coefficient calculation
which proves the required conjugation statement.

For every \(d\in(\mathbb Z/R)^\times\setminus\{1,-1\}\), use
the separation calculation to choose
\[
 H_a={\cal E}_{3,a}^2\quad\hbox{or}\quad
 H_a={\cal E}_{3,a}{\cal E}_{5,a}
\]
whose constant terms at \(a=1,d\) differ by a unit.
Let \(r_d=6\) or \(8\), accordingly.
There is an integral level-one form \(J_{r_d}\) with constant
term one: take \(E_6\) or \(E_4^2\).
The family
\[
                 H_a-a_0(H_d)J_{r_d}
 \tag{ES7}
\]
has constant term zero at \(a=d\) and a unit at \(a=1\).
It is integral and has the family transformation law (ES5).
Take the product of (ES7) over all \(d\ne\pm1\).
It is an even-weight family, of some weight \(m_0\), whose
constant term is zero except possibly at \(a=\pm1\),
and is a unit at \(a=1\).
All its coefficients lie in a fixed cyclotomic number field.
One may take \(m_0\le8(\phi(R)-2)\).

To reach every sufficiently large weight, use
\({\cal E}_{3,a}\), whose constant at \(a=1\) is a unit.
For odd \(j\ge10\), multiply it by an integral level-one
form of weight \(j-3\) and constant one.
For even \(j\ge10\), use \({\cal E}_{3,a}^2\) and a level-one
form of weight \(j-6\).
Such level-one forms exist in every nonnegative even weight
other than \(2\), as products of \(E_4,E_6\).
Multiplying by the previous product gives, for every
\(k\ge m_0+10\), an integral weight-\(k\) family \(F_a\)
which vanishes at the other cusps, satisfies (ES5), and has
\[
  u=a_0(F_1)\in\mathcal O_\ell^\times,\qquad
  a_0(F_a)=0\quad(a\ne\pm1),\qquad
                   a_0(F_{-1})=(-1)^ku .
\]
For any character \(\chi\) with the weight parity, take
\[
             G_\chi=\frac1{2u}
               \sum_{a\in(\mathbb Z/R)^\times}\chi(a^{-1})F_a .
 \tag{ES8}
\]
The transformation law makes its character \(\chi\).
It is integral, has constant term one, and vanishes at all
the other cusps.
Only the unit \(2u\) was divided out.
In particular this projection remains integral even if
\(\ell\mid\phi(R)\).
The coefficient formulas and the construction give modular
conjugates of \(G_\chi\) for every automorphism of
\(\mathbb C\).

\subsection{The small Eisenstein constant, proved by distribution}
The character projection of (ES4), divided by its Gauss sum,
gives a form \(E_\chi\) of weight \(k\), character \(\chi\),
vanishing at the other cusps, with
\[
 \begin{split}
 a_n(E_\chi)&=\sum_{d\mid n}\chi(d)d^{k-1}\quad(n\ge1),\\
 2\tau(\chi^{-1})a_0(E_\chi)
        &=\sum_{a\bmod R}\chi^{-1}(a)L_k(\zeta^a),
            \qquad \tau(\chi^{-1})=\sum_{a\bmod R}
                                      \chi^{-1}(a)\zeta^a .
 \end{split}
 \tag{ES9}
\]
The usual Gauss identity
\(\tau(\chi)\tau(\chi^{-1})\chi(-1)=R\) shows that both
Gauss sums are \(\ell\)-adic units, since \(\ell\nmid R\).
This identity follows by applying the finite Fourier
transform twice to the primitive character.
All the nonconstant coefficients in (ES9) are integral.

We prove \(c=a_0(E_\chi)\in\mathfrak m_\ell\) directly.
For \(x^w\ne1\), the geometric-series expansion and then
differentiation give
\[
             \sum_{j=0}^{w-1}L_k(x\zeta_w^j)
                                       =w^kL_k(x^w).
 \tag{ES10}
\]
Both sides are rational functions, so equality initially
proved for \(|x|<1\) proves it wherever their denominators
are nonzero.
Let
\[
          T_\varphi=\sum_{b\bmod M}
                         \varphi^{-1}(b)L_k(\zeta^{wb}).
\]
This is integral, since all nonzero contributing roots have
order dividing \(M\), prime to \(\ell\).
Split the sum on the right of (ES9) into multiples of \(w\)
and their complement, and use (ES10) on the full sum with
the \(\theta\)-factor omitted.
Because \(\varphi(w)=1\), the result is the exact identity
\[
 \begin{split}
 \sum_{a\bmod R}\chi^{-1}(a)L_k(\zeta^a)
  =\sum_{\substack{0\le n<R\\w\nmid n}}
        (\theta^{-1}(n)-1)\varphi^{-1}(n)L_k(\zeta^n)\\
                 +(w^k-1)T_\varphi .
 \end{split}
 \tag{ES11}
\]
The terms with zero character value contribute zero and may
be omitted before evaluating any rational function at \(1\).
Every \(\ell\)-th root of unity is congruent to \(1\) modulo
\(\mathfrak m_\ell\), and \(w^k\equiv1\pmod\ell\).
Each term of (ES11) therefore lies in
\(\mathfrak m_\ell\).
Since \(2\tau(\chi^{-1})\) is a unit, (ES9) proves
\[
                               c\in\mathfrak m_\ell .
 \tag{ES12}
\]

Consequently
\[
                             h=E_\chi-cG_\chi
\]
is a cusp form: it vanishes at the other cusps, and its
constant term at infinity is zero.
The character law transports this vanishing to the remaining
cusps equivalent to infinity under \(\Gamma_0(R)\).
Its coefficients are integral and
\[
 a_n(h)\equiv\sum_{d\mid n}\chi(d)d^{k-1}
                    \pmod{\mathfrak m_\ell}\quad(n\ge1),
                   \qquad a_1(h)\equiv1 .
 \tag{ES13}
\]
Thus its residual form is nonzero.

\subsection{Lifting the simultaneous eigenvalues}
The divisor-sum identity at every prime, including those
dividing \(R\), is
\[
 \sigma_\chi(pn)+\chi(p)p^{k-1}
       {\boldsymbol1}_{p\mid n}\sigma_\chi(n/p)
        =(1+\chi(p)p^{k-1})\sigma_\chi(n).
\]
It follows either by splitting the divisors according to
their \(p\)-valuation or by their geometric-series factor.
Applying it to (ES13) makes \(h\) a simultaneous residual
eigenvector for all the prime Hecke operators, with the
eigenvalues in (ES2).

Here is the integral eigenvalue lifting argument.
Choose one finite coefficient field containing every
coefficient in the explicit construction of \(h\), and let
\(\mathcal O\) be its \(\ell\)-adic integer ring.
Let \(V\) be the span over its fraction field of all Hecke
translates of \(h\).
It is finite dimensional.
Indeed all these translates have coefficients in that one
field, and linear independence of finitely many such
\(q\)-series is detected by finitely many coefficients.
Their rank over the coefficient field equals their rank
over \(\mathbb C\); the complex cusp-form space is finite
dimensional.
The lattice of elements of \(V\) with all coefficients
in \(\mathcal O\) embeds in \(\mathcal O^B\) using a finite
injective list of coefficient evaluations.
It is therefore finite and free.
It spans \(V\), since the integral Hecke translates of \(h\)
span \(V\).
The Hecke operators preserve it by their coefficient formula.

Their image \(\mathbb T\) in its endomorphism ring is a
finite, torsion-free, commutative \(\mathcal O\)-algebra.
The nonzero residual eigenvector (ES13) gives a maximal
ideal \(\mathfrak n\) and the specified residual eigenvalues.
Choose a minimal prime of \(\mathbb T\) below \(\mathfrak n\).
It does not contain the uniformizer, since the latter is a
non-zero-divisor in \(\mathbb T\).
The resulting generic character takes values in a finite
extension of the coefficient field and has a place reducing
to \(\mathfrak n\).
It occurs on \(V\): the faithful finite module has support
all of \(\operatorname{Spec}\mathbb T\), and the socle of
the corresponding local Artinian generic algebra gives a
common eigenvector.
This is the elementary Deligne--Serre lifting argument.

For a nonzero common eigenvector the first coefficient is
nonzero.
Otherwise the Hecke recurrence and induction on \(n\), using
a prime divisor of \(n\), make every positive coefficient
zero; cuspidality makes the zeroth coefficient zero as well.
Normalize that first coefficient to one.
The resulting eigenform satisfies (ES2).
Since \(\chi\) has conductor exactly \(R\), it cannot be
old from a proper divisor of \(R\): the central character
of such an old constituent would have smaller conductor.
Thus the form is a newform of the stated primitive level.
This completes the construction of the ordinary cuspidal
Eisenstein congruence.

\subsection{Its Galois representation and its rank-three adjoint}
We specify the classical inputs at this point, so that their
fields and hypotheses are visible.
For a normalized holomorphic newform of weight \(k\ge3\)
over \(\mathbb Q\), Deligne's construction supplies its
continuous absolutely irreducible representation \(\rho_f\),
with
\[
 \operatorname{tr}\rho_f(\operatorname{Frob}_p)=a_p(f),
 \qquad
 \det\rho_f=\varepsilon^{k-1}\chi .
 \tag{ES14}
\]
Here we use arithmetic Frobenius and its corresponding
standard dual convention; replacing every representation
by the dual gives the inverse-Frobenius convention used
elsewhere in the source.
Its Hodge--Tate weights in the convention
\(\mathrm{HT}(\varepsilon)=1\) are \(0,k-1\).
The prime \(\ell\) does not divide its level, and (ES2)
makes \(a_\ell(f)\) a unit.
The ordinary local theorem therefore gives an upper
triangular restriction at \(\ell\), whose unramified
quotient takes Frobenius to the unit root of
\[
                    X^2-a_\ell(f)X+\chi(\ell)\ell^{k-1}.
\]
At primes different from \(\ell\), local--global compatibility
identifies its Weil--Deligne representation with that of
the newform.
These are applications to a classical newform over
\(\mathbb Q\), of weight at least three, prime-to-\(\ell\)
level, and ordinary unit root; no Hilbert even-degree or
residual irreducibility hypothesis is inserted.

Equations (ES2) and (ES14), followed by density of Frobenius
and character independence in dimension two, give
\[
                  \bar\rho_f^{\,\mathrm{ss}}
                        =1\oplus\overline{\varepsilon^{k-1}
                                                      \varphi}.
 \tag{ES15}
\]
The \(\theta\)-factor disappears modulo
\(\mathfrak m_\ell\).
At every \(p\mid R\), primitive conductor forces the local
representation to be a principal series with one
unramified character and the ramified character prescribed
by \(\chi\).
Indeed the conductor of the central character already
equals the entire newform conductor; a special or
supercuspidal constituent would have larger representation
conductor than central-character conductor.
Thus its inertia type is exactly \(1\oplus\chi|_{I_p}\)
and its monodromy operator is zero.

Put \(m=k-1\), and suppose, as in the principal seed data,
\[
                  \overline{\varepsilon^m\varphi}^{\,2}\ne1
                         \quad\hbox{on }G_{\mathbb Q}.
 \tag{ES16}
\]
Then \(\rho_f\) remains absolutely irreducible on every
open subgroup.
Here is a proof avoiding any image-size assertion.
If its restriction to an open normal subgroup were
reducible, prime-dimensional Clifford theory gives either
a scalar restriction or two lines interchanged transitively.
A scalar restriction would make its projective image finite.
On an open local subgroup at \(\ell\) its two Hodge--Tate
weights would then coincide, contrary to \(0,m\), \(m>0\).
Two interchanged lines give a quadratic self-twist
\(\alpha\ne1\).
Reducing that twist in (ES15), in odd characteristic, gives
\[
        \{1,\bar\Phi\}=\{\bar\alpha,\bar\alpha\bar\Phi\},
                         \qquad \Phi=\varepsilon^m\varphi.
\]
Since the nontrivial quadratic character stays nontrivial
modulo \(\ell\), these multisets force
\(\bar\alpha=\bar\Phi\) and \(\bar\Phi^2=1\), contradicting
(ES16).

In characteristic zero this also makes
\(\operatorname{ad}^0\rho_f\) irreducible on every open
subgroup.
The connected Zariski closure of \(\rho_f\) is reductive:
its unipotent radical has nonzero fixed vectors, and
normality and irreducibility make it act trivially.
It acts irreducibly in dimension two and cannot be a torus.
Its nonzero semisimple derived Lie algebra is therefore
\(\mathfrak{sl}_2\); the connected group contains
\(\mathrm{SL}_2\).
The adjoint action on its three-dimensional simple Lie
algebra is irreducible.
Every open subgroup has the same connected Zariski closure,
which proves the assertion.

Now let \(L/\mathbb Q\) be the finite abelian CM extension
in the seed packet, with
\[
          \eta|_{G_L}=\theta|_{G_L}=1,\qquad
                         \varphi=\chi'\eta .
\]
For a prescribed residue weight \(j\), choose
\[
                         m\equiv j\pmod{\ell-1}
\]
as large as desired and \(k=m+1\) above the constructed
threshold.
The seed parity is
\(\varphi(-1)=(-1)^{j+1}\), so it gives the required
weight parity.
On \(G_L\), (ES15) and the adjoint formula give the exact
rank-three residual semisimplification
\[
 \overline{\operatorname{ad}^0\rho_f|_{G_L}}^{\,\mathrm{ss}}
       =\bar\Psi\oplus1\oplus\bar\Psi^{-1},
                          \qquad \Psi=\varepsilon^j\chi' .
 \tag{ES17}
\]
The ordinary flag has Hodge--Tate weights \(m,0,-m\).
Its ordered characters are the ratio of the two ordinary
characters of \(\rho_f\), the trivial character, and the
inverse ratio.
Their reductions are exactly the ordered tuple in (ES17).

The trace form on trace-zero two-by-two matrices is
nondegenerate, symmetric, and invariant under conjugation.
For a complex conjugation \(c\), multiply its Gram matrix
by \(\operatorname{ad}^0\rho_f(c)\).
The result is symmetric, since \(c^2=1\), and gives
\[
       r(g)^{\mathsf t}P\,r(cgc^{-1})=P,
                  \qquad r=\operatorname{ad}^0\rho_f|_{G_L}.
 \tag{ES18}
\]
Thus \(r\varepsilon^{-1}\) has the standard unitary
multiplier \(\varepsilon^{-2}\).
Away from \(\ell\) its finite inertia type is
\(\chi'\oplus1\oplus\chi'^{-1}\); the \(\eta,\theta\)
factors disappear on \(G_L\).
In particular it is unramified away from the ramification
set allowed in the seed packet.
All these assertions concern the actual adjoint
representation, rather than a determinant with the same
generic roots.

\subsection{The automorphic transfer and local root bounds}
There is a direct classical realization of this seed.
Iterate cyclic base change for \(\mathrm{GL}_2\) along a
prime-degree tower for the abelian extension \(L/\mathbb Q\).
The base change is cuspidal: a noncuspidal step would be
the dihedral/self-twist case, making the attached
two-dimensional representation reducible on a finite-index
subgroup, already excluded above.
At each step the unramified parameters match restriction,
so the Galois realization is \(\rho_f|_{G_L}\).
It has no nontrivial self-twist, by the same open-subgroup
irreducibility.
The adjoint lift of Gelbart--Jacquet, Theorem~9.3, therefore
gives a cuspidal \(\mathrm{GL}_3(\mathbb A_L)\)-representation
with the local adjoint parameters.
It realizes \(r\), and after the indicated algebraic scalar
normalization realizes \(r\varepsilon^{-1}\).
The weights are regular, since \(m>0\), and (ES18) gives
the required conjugate self-duality.

The same theorem supplies the unlinked local property of this
seed without a new Selmer hypothesis.
At a finite place, unitarize the \(\mathrm{GL}_2\)-representation.
For a complementary principal-series component its two
Weil eigenvalue absolute values have ratio \(q^{2t}\).
Gelbart--Jacquet, Theorem~9.3(4), gives \(|t|<1/4\).
For a tempered component take \(t=0\).
The primitive-level ramified components above are principal
series with ramified ratio; their unitarity makes both
inducing characters unitary, so they too have \(t=0\).
For the adjoint representation, the three absolute values
are consequently
\[
                              q^{2t},\quad1,\quad q^{-2t}.
\]
The absolute value of a ratio between any two is strictly
between \(q^{-1}\) and \(q\).
After local base change of residue degree \(f\), the exponents
and the logarithm of the residue norm both multiply by \(f\),
so the same strict inequality holds with the new norm.
Finite inertial character factors have absolute value one.
Thus no two adjoint Frobenius roots have ratio equal to the
local residue norm or its inverse.
The monodromy is zero at the finite ramified places under
consideration.
This proves the actual unlinked condition for this seed;
it is a property of the constructed local parameters,
rather than a condition imposed on an arbitrary deformation.

The classical sources just used are Deligne,
\emph{Formes modulaires et repr\'esentations \(\ell\)-adiques},
Bourbaki~355, for the two-dimensional realization;
the ordinary local theorem in Wiles,
\emph{On ordinary \(\lambda\)-adic representations associated
to modular forms}, \emph{Invent.\ Math.}~94 (1988), 529--573;
Carayol, \emph{Sur les repr\'esentations \(\ell\)-adiques
associ\'ees aux formes modulaires de Hilbert}, Theorem~A,
whose base field here is \(\mathbb Q\), of odd degree;
Langlands, \emph{Base Change for \(\mathrm{GL}(2)\)},
Chapter~2, for the cyclic transfers;
and Gelbart--Jacquet, \emph{A relation between automorphic
representations of \(\mathrm{GL}(2)\) and \(\mathrm{GL}(3)\)},
Theorem~9.3, for the adjoint lift.
Primary copies are respectively
\url{https://archive.numdam.org/item/SB_1968-1969__11__139_0/},
\url{https://doi.org/10.1007/BF01394275},
\url{https://archive.numdam.org/articles/10.24033/asens.1512/},
\url{https://publications.ias.edu/sites/default/files/base-change-for-gl2_rpl.pdf},
and
\url{https://www.numdam.org/articles/10.24033/asens.1355/}.
The specific ordinary matrix statement, including the
unit-root quotient, is matched with the duality convention
in Lemma~\ref{lem:sr-stable-three-seed-insertion}.

The constructive part is the actual
\srcid{UnlinkedHeckeRoots.PropEC.propEC} producer.
Its ingredients at the frozen tag are
\srcid{IdeleGenerators.HolomorphicIrreducibility.EisSlash},
\srcid{EisVan}, \srcid{EisCoef},
\srcid{EulerianNumeratorThree.lkSep},
\srcid{LevelOne.levelOne}, and \srcid{ConstSmall.constSmall};
the conjugation-preserving version is
\srcid{UnlinkedHeckeRoots.UnitFamC.unitFamC},
\srcid{UnitFormC.unitFormC}, and \srcid{EigLiftC.eigLiftC}.
They are under
\srcid{Development.AutomorphicArithmetic.AutomorphicVarieties}.
The subsequent source calls are
\srcid{ThresholdArguments.DenseThresholdCases.WEGCore},
\srcid{TPropEGW}, \srcid{TSeedInW}, and
\srcid{AdjointLiftFromEisensteinSeedOfHyps}.

What has been produced is an actual ordinary classical
eigenform and an actual strongly irreducible polarized
rank-three automorphic adjoint with the required residual
and finite inertia packet, at arbitrarily large compatible
weights.
The field construction below chooses the definite unitary
pure inner form and retains the scalar Frobenius on the
selected tower. Lemma~\ref{lem:sr-stable-three-seed-insertion}
then supplies its ordinary local vector and localization
in the same Hecke module and determinant ring. That
insertion uses the actual characteristic-zero seed and
its local types, as well as its residual determinant.


\subsection{The Eisenstein seed at coefficient prime \(3\)}
\label{subsec:sr-eisenstein-three}
The constructed Eisenstein seed also works at \(\ell=3\)
on the actual prime-to-\(3\) level.
Only the constant-term separation calculation changes;
the absolutely convergent Eisenstein series have weight
at least three throughout.

Use the notation of
Section~\ref{sec:constructed-eisenstein-seed}.
For \(x\ne1\), put \(u=x/(1-x)^2\).
The Eulerian recurrence gives
\[
        L_4(x)=\frac{x(1+4x+x^2)}{(1-x)^4}
                                     =u+6u^2 .
 \tag{E3.1}
\]
If \(x,y\) are primitive roots of order \(R\), \(3\nmid R\),
and \(y\ne x,x^{-1}\), the same factorization as before gives
\[
 u-v=\frac{x(1-y/x)(1-xy)}
                         {(1-x)^2(1-y)^2}
                      \in\mathcal O_3^\times .
\]
All denominators and numerator factors are units, since
reduction preserves roots of unity of order prime to \(3\).
Consequently
\[
          L_4(x)-L_4(y)
               =(u-v)\bigl(1+6(u+v)\bigr)
                          \in\mathcal O_3^\times .
 \tag{E3.2}
\]
The second factor is congruent to \(1\) modulo \(3\).

Thus, for each \(d\ne\pm1\), replace the weight-\(6\)
or weight-\(8\) separating family in (ES7) by
\[
                H_a={\cal E}_{4,a},\qquad
                H_a-L_4(\zeta^d)E_4 .
 \tag{E3.3}
\]
The forms \({\cal E}_{4,a}\) are exactly the absolutely
convergent series (ES4); their constant terms and
coefficients are given by (ES6).
The integral level-one form \(E_4\) has constant term one.
Equation (E3.2) gives the required zero at \(a=d\)
and unit at \(a=1\).
The product has weight
\(m_0=4(\phi(R)-2)\).

Here \(R=Mw\), with \(M>1\) and a fresh prime
\(w\equiv1\pmod3\), so primitive \(R\)-th roots are not
\(1\) or \(-1\).
Their \(L_3\)-values are units.
Multiply this product by \({\cal E}_{3,a}\) in odd
remaining weight and by \({\cal E}_{3,a}^2\) in even
remaining weight, followed by a suitable product of
\(E_4,E_6\).
As in the original argument this reaches every
\(k\ge m_0+10\), with the chosen parity.
These final \({\cal E}_{3,a}\)-factors give vanishing
at every cusp not equivalent to infinity under
\(\Gamma_0(R)\).
The projection divides only by the unit \(2u\).
It therefore gives the integral \(G_\chi\) with
constant term one and the prescribed character.

The distribution identity (ES10)--(ES11) applies unchanged
to an auxiliary character \(\theta\) of exact order \(3\).
Its values are congruent to \(1\) modulo the maximal ideal;
also \(w^k-1\) is divisible by \(3\).
The Gauss sum is a unit because \(3\nmid R\).
Thus the Eisenstein constant is in the maximal ideal,
and \(E_\chi-cG_\chi\) is the same nonzero residual
cuspidal eigenform.
The finite integral Hecke-algebra lifting argument,
the primitive-level newform argument, Deligne realization
and ordinary local theorem impose no restriction
\(\ell\ge5\).
All of them now apply at \(3\).
In particular the all-odd-prime Eisenstein seed, its
rank-two normalization and its rank-three adjoint seed
are constructed at the prime required by the Gaussian
branch.
There is no weight-two regularization or inverted
factor \(3\) in this extension.


\section{The seed packet at every generic coefficient prime}
\label{sec:seed-all-generic-primes}

The prime-to-\(\ell\) hypothesis in the classical
Eisenstein construction concerns its finite character
conductor, not the original host conductor.
We make the necessary adjustment explicitly.
This also constructs a common CM field in which the
distinguished scalar Frobenius retains its required
cyclotomic value.

Let \(N=2^{h+2}qd\) be the principal host conductor, with
\(q,d>7\), and let \(\psi\) be a host character of parity
\(\psi(-1)=(-1)^a\).
Fix
\[
       \ell\notin\{2,3,7\},\qquad
                         \ell-1\nmid2a .
 \tag{SP1}
\]
Every odd-primary conductor component of a host character
is quadratic: the odd ramification comes from the two
quadratic radicals.
Consequently
\[
        \psi=\psi^{(\ell)}\widetilde\omega_\ell^{\,b},
       \qquad b\in\{0,(\ell-1)/2\},
 \tag{SP2}
\]
where \(\widetilde\omega_\ell\) is the global finite
Teichm\"uller character obtained from the cyclotomic character,
and \(\psi^{(\ell)}\) has conductor prime to \(\ell\).
For \(\ell\nmid N\), take \(b=0\).
For \(\ell=q\) or \(d\), the nonzero choice of \(b\) is
exactly the possible quadratic component of conductor \(\ell\).

Choose \(j\) with
\[
                       j\equiv a+b\pmod{\ell-1},
 \qquad
              \chi'=\psi^{(\ell)},\quad
              K=\mathbb Q(\sqrt{-3}),\quad\eta=\chi_K .
 \tag{SP3}
\]
Then
\[
 \chi'(-1)=(-1)^j,\qquad
  2j\equiv2a\pmod{\ell-1},\qquad
       \bar\varepsilon^{\,j}\bar\chi'
                       =\bar\varepsilon^{\,a}\bar\psi .
 \tag{SP4}
\]
The finite character
\(\varphi=\chi'\eta\) has conductor \(M>1\) prime to
\(\ell\).
It is nontrivial because \(\eta\) is ramified at \(3\),
whereas the sufficiently large principal host has no
ramification at \(3\).
This verifies the nontriviality assumption of the classical
seed construction, including the cases \(\ell=q,d\).

Let \(N_\ell\) be a prime-to-\(\ell\) common conductor of
\(\chi'\) and \(\eta\), for example
\(3N/\ell^{v_\ell(N)}\).
Choose a fresh prime
\[
                       w\equiv1\pmod{\ell N_\ell},
 \tag{SP5}
\]
and a nontrivial order-\(\ell\) Dirichlet character
\(\theta\) modulo \(w\).
Such a character exists because \(\ell\mid w-1\).
It has conductor \(w\), is even, and cuts out a totally
real cyclic field \(F_\theta/\mathbb Q\) of degree \(\ell\).
Equation (SP5) gives
\(\chi'(w)=\eta(w)=\varphi(w)=1\).
Thus all conductor, split-at-\(w\), parity, and nontriviality
requirements of the classical Eisenstein construction hold.
Its weights may be taken
\[
                           m\equiv j\pmod{\ell-1},
                         \qquad k=m+1
 \tag{SP6}
\]
arbitrarily large.

\subsection{The full local packet after the adjustment}
The representation constructed at weight \(m\) has global
residual character
\(\bar\varepsilon^{\,m}\bar\chi'\bar\eta\).
On the CM field containing \(K\) and \(F_\theta\), its
rank-three adjoint therefore has residual tuple
\[
        \bar\varepsilon^{\,a}\bar\psi,\quad1,\quad
                      (\bar\varepsilon^{\,a}\bar\psi)^{-1},
 \tag{SP7}
\]
by (SP4).
This is an equality of global residual characters.

The finite inertia types away from \(\ell\) agree
\emph{before reduction}.
Indeed \(\widetilde\omega_\ell\) is unramified at every
finite prime different from \(\ell\).
Consequently \(\psi|_{I_p}=\chi'|_{I_p}\) there, while
\(\eta,\theta\) vanish after the indicated field restriction.
The adjoint's type
\(\chi'\oplus1\oplus\chi'^{-1}\) is exactly the original
host finite type at those places.

At a place above \(\ell\), local class field theory identifies
the torsion of the local abelianization with the torsion
in the local multiplicative group.
The cyclotomic character on it takes values in
\(\mathbb Z_\ell^\times\).
It kills its \(\ell\)-primary torsion, because
\(\mathbb Z_\ell^\times\) has no nontrivial finite
\(\ell\)-power subgroup for odd \(\ell\).
The finite host characters have \(2\)-power order and
therefore kill the same \(\ell\)-primary torsion.
This is precisely the subgroup denoted
\srcid{GFrame.mu} in the source.

On the remaining finite torsion, the cyclotomic values are
Teichm\"uller units, so
\[
           \varepsilon^m
              =\varepsilon^a\widetilde\omega_\ell^{\,b}
 \tag{SP8}
\]
there, by (SP6).
The ordinary unit-root quotient of the seed is unramified
and kills these unit torsion elements.
Thus the ordered ordinary characters of its adjoint have
exactly the same finite-torsion restrictions as the original
packet.
Their principal-unit powers vary with the arithmetic
weight \(m\), as allowed in the common ordinary weight
algebra.
No equality between the distinct principal-unit powers
\(\varepsilon^m\) and \(\varepsilon^a\) is being claimed.
Equations (SP7)--(SP8), the fixed away-\(\ell\) types, and
the common polarization put the seed in the same prescribed
ordinary torsion component.
Applying the same fixed algebraic scalar normalization to
both packets preserves all these equalities.

\subsection{A finite-unramified CM tower with a scalar prime}
Choose a fresh prime \(p_0\), different from \(w\), with
\[
                 p_0\equiv-1\pmod{\ell N},\qquad
                              p_0\equiv1\pmod3 .
 \tag{SP9}
\]
The congruences are compatible because \(3\nmid\ell N\).
Dirichlet's theorem provides arbitrarily large such primes.
It splits in \(K\).
At this prime,
\[
 \bar\varepsilon(\operatorname{Frob}_{p_0})=-1,\qquad
 \bigl(\bar\varepsilon^{\,a}\bar\psi\bigr)
                      (\operatorname{Frob}_{p_0})=1.
 \tag{SP10}
\]
The second equation is the host parity:
\((-1)^a\psi(-1)=1\).
Thus all three residual adjoint eigencharacters coincide
at this initial Frobenius.

Let \(\mathcal P\) be any prescribed finite set of rational
primes at which the eventual CM extension must split
over its real part; include \(2,3,\ell\), and all primes
of the intended tame level.
The following construction also permits arbitrary prescribed
lower bounds for its real degree.

Choose a fresh positive prime \(D\) satisfying
\[
 \begin{gathered}
 D\equiv1\pmod4,\qquad D\equiv2\pmod3,\qquad
                      \left(\frac{3D}{p_0}\right)=1,\\
 \left(\frac Dp\right)=\left(\frac{-1}p\right)
 \quad\left(p\in\mathcal P,\ p\ge5,\
                     \left(\frac{-3}p\right)=-1\right).
 \end{gathered}
 \tag{SP11}
\]
It avoids \(\mathcal P\cup\{p_0,w\}\).
Each displayed Legendre condition is a nonempty choice of
unit residue classes at its respective odd prime.
There is no extra condition at \(p_0\) in the second line:
it splits in \(K\).
The Chinese remainder theorem and then Dirichlet's theorem
give a positive prime \(D\) in one of these simultaneous
unit classes.

Put
\[
            F_0^+=\mathbb Q(\sqrt{3D}),\qquad
                     L_0=F_0^+K.
\]
The third quadratic subfield of \(L_0\) is
\(\mathbb Q(\sqrt{-D})\).
Since \(D\equiv1\pmod4\), their discriminants are
\[
       -3,\qquad 12D,\qquad-4D .
\]
The discriminant formula for a biquadratic field gives
\[
        \operatorname{disc}(L_0)=144D^2
                          =\operatorname{disc}(F_0^+)^2 .
 \tag{SP12}
\]
The relative discriminant ideal therefore has norm one.
Thus \(L_0/F_0^+\) is unramified at \emph{every finite place}.

For \(p\in\mathcal P,\ p\ge5\), if \(p\) splits in \(K\),
then every place of \(F_0^+\) above it splits in \(L_0\).
If \(p\) is inert in \(K\), (SP11) gives
\[
                 \left(\frac{3D}p\right)
                    =\left(\frac{-3}p\right)=-1 .
\]
The two local quadratic fields \(K_p\) and \(F_{0,p}^+\)
are then the same unramified quadratic extension of
\(\mathbb Q_p\).
Their compositum splits over \(F_{0,p}^+\), so again the
CM extension splits at the desired place.
At \(3\), the ratio of the two quadratic radicands is
\((3D)/(-3)=-D\).
It is a square in \(\mathbb Q_3^\times\), since
\(D\equiv2\pmod3\).
Thus \(F_{0,3}^+=K_3\) and the CM extension splits there.

To arrange splitting at \(2\), choose another fresh positive
prime \(B\) with
\[
                      B\equiv5\pmod8,\qquad
                           \left(\frac B{p_0}\right)=1 .
 \tag{SP13}
\]
Again CRT and Dirichlet provide it, outside every already
specified finite set.
The field \(\mathbb Q_2(\sqrt B)\) is the unramified
quadratic extension of \(\mathbb Q_2\).
The real field \(F_0^+\) is ramified quadratic at \(2\),
since \(3D\equiv3\pmod4\).
The relative extension in (SP12) is unramified there,
and is killed by adjoining this unramified quadratic
field.  If it was already split there, splitting is
preserved.
Hence, after adjoining the real field \(\mathbb Q(\sqrt B)\),
the CM extension is split at \(2\) as well.
All previously obtained splitting conditions are preserved
by base change.

For each depth \(r\ge0\), let \(C_r\) be the unique
cyclic degree-\(\ell^r\) subfield of
\(\mathbb Q(\zeta_{\ell^{r+1}})\).
It is the field fixed by the Teichm\"uller subgroup of
order \(\ell-1\).
In particular \(C_r\) is totally real and totally ramified
at \(\ell\).
Its conductor is an \(\ell\)-power, whereas the conductors
of the preceding real fields are prime to \(\ell\).
This is the cyclotomic depth choice required for the
uniform class-group argument below.
Take
\[
 \begin{split}
 F^+&=F_0^+\mathbb Q(\sqrt B)F_\theta C_r,\\
                         L&=F^+K .
 \end{split}
 \tag{SP14}
\]
The fixed real biquadratic field has degree four.
The order-\(\ell\) field \(F_\theta\), of conductor \(w\),
is disjoint from the \(\ell\)-power conductor field \(C_r\):
their intersection would be an extension of \(\mathbb Q\)
unramified at every finite prime.
It is also disjoint from the biquadratic field.
Consequently
\[
                 [F^+:\mathbb Q]=4\ell^{r+1},
                        \qquad [L:F^+]=2 .
 \tag{SP15}
\]
The real degree is even.
At \(\ell\) the fixed real fields are unramified and
\(C_r\) is totally ramified of degree \(\ell^r\).
Every ordinary local degree is therefore at least
\(\ell^r\).
The global and all ordinary local degrees can exceed
any prescribed bound.
The uniform arithmetic bound for the three Selmer
dimensions denoted \(C_0\) is proved separately, before
this depth is chosen.
The fields are abelian over \(\mathbb Q\), \(L\) is CM,
and \(F^+\) is its maximal real subfield.
Also \(\eta|_{G_L}=\theta|_{G_L}=1\).

The finite-unramified property (SP12) survives every real
base change: an unramified local extension remains
unramified or split after finite base change.
It follows that \(L/F^+\) is unramified at every finite
place and split at every prime in \(\mathcal P\).
Thus it satisfies the exact \srcid{UnrSplit} condition.
The fixed fields in (SP14) have conductor prime to
\(\ell\); adjoining \(C_r\) gives ramification index
\(\ell^r\) at \(\ell\).
It follows that \(L\) is locally disjoint from the
tamely ramified extension
\(\mathbb Q_\ell(\mu_\ell)/\mathbb Q_\ell\),
whose degree is \(\ell-1\).
The image of \(\bar\varepsilon\) on inertia over \(L\)
is therefore still all of \(\mathbb F_\ell^\times\).
In particular no completion of \(L\) at \(\ell\)
contains \(\mu_\ell\) or nontrivial \(\ell\)-power roots
of unity.
Absolute unramifiedness at \(\ell\) is not needed.
Equations (SP1) and (SP4) show that
\[
               (\bar\varepsilon^{\,j}\bar\chi')^2
                                      \ne1\quad\hbox{on }G_L .
 \tag{SP16}
\]
This verifies the actual square-separation field condition,
not just the corresponding condition over \(\mathbb Q\).

Finally \(p_0\) splits in both real quadratic factors in
(SP14), by (SP11) and (SP13), and splits in \(K\).
It is unramified in \(F_\theta C_r\).
All the remaining residue degrees divide
\(\ell^{r+1}\), hence are odd.
For every \(u\mid p_0\) in \(L\), with residue degree \(f_u\),
restriction of unramified Frobenius gives
\[
 \bar\varepsilon(\operatorname{Frob}_u)
                       =(-1)^{f_u}=-1,\qquad
 (\bar\varepsilon^{\,a}\bar\psi)(\operatorname{Frob}_u)=1.
 \tag{SP17}
\]
Thus the scalar residual Frobenius and nontrivial
cyclotomic value survive the \emph{actual selected tower}.
The same conclusion holds for the adjusted character in
(SP4).
Any residual lattice need only have this scalar
characteristic polynomial; it need not have a scalar
matrix.  The matrix-minus-scalar is nilpotent after
reduction, which is exactly the local input used in the
separate \(N4\) matrix argument.
It also supplies the element with nontrivial
\(\bar\varepsilon\) required by first-Selmer trace detection.

\subsection{The definite rank-three group}
There is now no separate existence or parity problem for
the required definite unitary pure inner form.
Use the hermitian form
\[
                             h=1_3
\]
on \(L^3\).
It is positive definite at every real place of \(F^+\).
At a finite split place its unitary group is
\(\mathrm{GL}_3(F_v^+)\).
At a finite nonsplit place \(L_v/F_v^+\) is unramified.
Norms of units are surjective there, so \(-1\) is a norm.
The hermitian forms \(1_3\) and
\(\operatorname{diag}(1,-1,1)\) have the same determinant
class and hence are isometric over that local field.
The latter contains a hyperbolic plane.
Thus \(U(h)\) is quasi-split at every finite place.
Both forms are global hermitian forms over the same
quadratic extension, so their comparison defines a
pure inner twist, rather than merely an abstract inner
form.
The even degree in (SP15) is consistent with the global
product condition; it has been produced by the field
construction itself.

The constructed adjoint seed has a simple cuspidal
conjugate self-dual rank-three parameter.
With this group, the prescribed split finite local
members, the spherical members at the nonsplit finite
places, and the compact cohomological members at infinity
are the relevant local choices.
The separate unitary packet/vector realization must place
those choices in the ordinary module at its fixed type
and weight specialization.
The field construction and torsion matching above supply
its actual arithmetic hypotheses.

\paragraph{Correspondence and exact scope.}
The prime-to-\(\ell\) finite-character requirement is
\srcid{FieldPairCharacterTripleConditions.aux.cop} together
with \srcid{chi\_split} in
\srcid{IdeleGenerators.SlopeTestRemainingChoices}.
The local torsion subgroup is
\srcid{WeilGroupCharacters.UniversalOrdinaryDeterminantRing.GFrame.mu}.
The finite-unramified and split-at-level requirements are
\srcid{ModularityArgument.SymmetricPowerTransfer.UnrSplit}.
All these paths are under
\srcid{Development.AutomorphicArithmetic.AutomorphicVarieties}
at \srcid{alpoge-source-7cbf7a0c}.

This proves the packet adjustment at every coefficient
prime in (SP1), including \(\ell\mid N\), and gives explicit
common fields satisfying the seed's character, square,
finite-unramified, splitting, degree, and distinguished
Frobenius requirements.
It proves these assertions for the tower (SP14).
An independently chosen CM tower is not entitled to
(SP17) without an analogous residue-degree check.


\subsection{Inserting a stable cuspidal seed in the ordinary module}
\label{subsec:sr-stable-seed-insertion}
The local lattice and ordinary normalizations in
Section~\ref{subsec:sr-case-two-hecke} apply in arbitrary fixed rank.
Here is the precise specialization needed for the rank-three seed.
It uses the same finite function module at every auxiliary level.

\begin{srlemma}\label{lem:sr-stable-three-seed-insertion}
Let \(L/F\) be the actual CM tower, with \(F\) totally real,
\(L/F\) unramified at finite places and split at all places in
the prescribed finite set \(S\), including those above \(\ell\)
and the distinguished auxiliary place \(v_0\).
Suppose the definite hermitian group \(U(3)\) is a pure inner
form of its quasi-split unitary group and is quasi-split at
every finite place.
Let \(\Pi\) be a regular algebraic conjugate-self-dual
cuspidal representation of \(\mathrm{GL}_3(\mathbb A_L)\)
with the following actual local data:
\[
 \begin{gathered}
 r_{\ell,\iota}(\Pi)=r_0,\qquad
 r_0\bmod\mathfrak m\simeq
                  \bar\chi_0\oplus\bar\chi_1\oplus\bar\chi_2,\\
 \mathrm{WD}(r_0|_{G_{L_w}})^{F\text{-ss}}
       =\bigoplus_{i=0}^2 \lambda_{wi},\quad N=0
                    \quad(w\nmid\ell),\\
 \lambda_{wi}|_I=[\bar\chi_i]|_I
        \quad\hbox{at the prescribed split tame places},\\
 r_0|_{G_{L_w}}\ \hbox{ordinary with the prescribed ordered
 rank-one characters }\psi_{wi}
                    \quad(w\mid\ell).
 \end{gathered}
 \tag{S1}
\]
Assume at \(\ell\) that the local factor of \(\Pi\) is a
principal series with the ordinary refinement specified by
these characters. At nonsplit finite places outside \(S\),
assume it is unramified.
Assume also that \(\Pi\) is unramified at \(v_0\).
Then its seed point is supported by the residual ordinary
function module with these tame types and these weight
and ordinary eigenvalues.
In particular the kernel \(J_0\) of the actual ordinary
Hecke quotient vanishes at that point.
\end{srlemma}
\begin{proof}
Use
\href{https://arxiv.org/pdf/1409.3731v3}
{Kaletha--M\'inguez--Shin--White, Theorem 1.7.1}
only in the proved generic-parameter, pure-inner-twist
scope specified immediately after its statement.
The simple cuspidal parameter of \(\Pi\) has trivial global
component group modulo the center.
Its compact cohomological archimedean members therefore
occur with the prescribed split local members and
spherical members at every nonsplit finite place.
The hypotheses on the actual hermitian group supply
the pure-inner and global existence conditions.
This gives an actual automorphic representation \(\pi\)
of the chosen definite unitary group, with base change \(\Pi\).

At a split tame place use the coefficient representation
\[
                   \sigma^{\mathrm{crys}}(\tau_w)^\vee,
       \qquad \tau_w=\bigoplus_{i=0}^2[\bar\chi_i]|_I .
 \tag{S2}
\]
The local type criterion
\href{https://math.uchicago.edu/~emerton/pdffiles/bonn-lectures-survey-article.pdf}
{Emerton--Gee, Theorem 9.2.4}
and (S1) give
\(\operatorname{Hom}(\sigma^{\mathrm{crys}}(\tau_w),\Pi_w)
\ne0\).
The dual in (S2) is exactly the coefficient convention for
the function module. Choose its stable finite-free integral
lattice as in (H3).

At \(w\mid\ell\), the chosen principal-series refinement
has a one-dimensional Jacquet eigenspace.
After taking a sufficiently deep ordinary level, a
nonzero vector on this eigenspace is represented by a
local vector of the principal series.
This follows directly from the induced-function model:
on the open Bruhat cell prescribe its value at the identity,
extend by the inducing character and a sufficiently small
right compact open subgroup, and sum over the finite
unipotent quotient for the ordinary projector.
The normalized Jacquet action gives
\[
 \begin{aligned}
 \psi_{wi}(g)\epsilon_\ell(g)^i\,
                      \theta(e_{wi}(y))&=1
                  &&(\operatorname{Art}_w(g)=y\in\mathcal O_{L_w}^\times),\\
 U_{wj}&=\prod_{i=0}^j
                \psi_{wi}(g_w)\epsilon_\ell(g_w)^i
                  &&(\operatorname{Art}_w(g_w)=\varpi_w).
 \end{aligned}
 \tag{S3}
\]
These are (H5)--(H6) in rank three, so both the function
coefficient character \(\theta^{-1}\) and the operator
normalization agree with the common ordinary module.
The right-hand sides in the second line are units:
the \(\psi_{wi}\) are actual continuous Galois characters,
and the chosen local reciprocity lift has cyclotomic
value a unit.
The normalized ordinary projector thus acts as the
identity on the refined line.

Choose spherical vectors at the other unramified places.
At the distinguished auxiliary place use a spherical vector
inside the prescribed smaller level.
Tensor these with the local vectors above and the
cohomological archimedean factor.
The resulting automorphic vector is a nonzero classical
ordinary vector in precisely the finite function space used
to construct the common module.
All its coefficients, finite type matrices and eigenvalues
are defined over a finite coefficient extension.
There are finitely many double cosets.
Clear their finitely many denominators and scale to obtain
a primitive integral vector.
This operation retains every Hecke eigenvalue and
ordinary eigenvalue in (S3).

The spherical residual polynomial and (S3) are the
specified residual Hecke system.
The residual direct-summand projector therefore acts as
the identity on this vector. It survives in that
localized ordinary module.
The rank-independent coefficient-change and classical-control
isomorphisms identify its classical specialization with the
corresponding specialization of the same module over
\(\Lambda\).
Thus its Hecke eigensystem factors through the actual
closed Hecke algebra in \(\operatorname{End}_\Lambda(M)\).
The universal ordinary determinant maps to this algebra
with kernel \(J_0\). Evaluation at the constructed vector
annihilates \(J_0\), as claimed.
\end{proof}

For the rank-three adjoint seed, the preceding hypotheses are
checked on the actual seed, rather than imposed on an
arbitrary ordinary deformation.
Its construction over \(\mathbb Q\), cyclic CM base change
and Gelbart--Jacquet lift give the stable cuspidal \(\Pi\).
The local Galois representation after the inverse-cyclotomic
normalization is the actual \(r_0\) in (S1).
The local \(\ell\)-factor is unramified; its ordered Satake
parameters are those of the three crystalline ordinary
characters.
Their Jacquet refinement has the unit eigenvalues (S3).
The strict local adjoint root-ratio bound excludes linked
inducing characters, so this is an irreducible principal
series. The stated seed inertia at tame places gives (S2)
with \(N=0\).
The field construction gives a definite pure inner
\(U(3)\) quasi-split at finite places.
Consequently Lemma~\ref{lem:sr-stable-three-seed-insertion}
inserts this seed in the common module; no additional
existence statement for a modular ideal is used.

% Rank-independent part of THidaU, TTypeEx, TThetaEx,
% UnitaryHeckeData and ExistsHeckeIdeal, specialized to the
% stable rank-three cuspidal seed. The characteristic-zero
% eigenvector and residual support are in the SAME module.


\subsection{The actual rank-two companion seed on the same packet}
\label{subsec:generic-rank-two-companion}

The rank-three propagation uses the rank-two space
for the paired outer block.
Its initial point can be constructed on the same
fields and with the same auxiliary order-\(\ell\)
character as the rank-three seed.
We give that construction explicitly.

Retain
\(\Psi=\bar\varepsilon^j\bar\chi'
 =\bar\varepsilon^a\bar\psi\),
\(\eta=\chi_{\mathbb Q(\sqrt{-3})}\),
and \(\ell\notin\{2,3,7\}\).
Run the proved Eisenstein construction with
\[
                    \varphi_2=(\chi')^2\eta,
                         \qquad k_2=2m+1,\quad
                            m\equiv j\pmod{\ell-1}.           \tag{RS1}
\]
The finite character \(\varphi_2\) is primitive at
its own conductor, nontrivial, and of conductor
prime to \(\ell\).
Nontriviality is supplied by its factor at three.
It is odd, as required for the odd weight \(k_2\).
The same fresh \(w\) satisfies
\(\varphi_2(w)=1\), and the same \(\theta\)
has order \(\ell\), conductor \(w\), and trivial
restriction to \(G_L\).
Thus there are ordinary newforms \(f_2\) at all
sufficiently large weights (RS1).
Their actual Galois representations satisfy
\[
 \det\rho_{f_2}=\varepsilon^{2m}(\chi')^2\eta\theta,
 \qquad
 \bar\rho_{f_2}^{\rm ss}
                =1\oplus\bar\varepsilon^{2m}(\bar\chi')^2\eta .
                                                               \tag{RS2}
\]
At \(\ell\) they are crystalline ordinary of
weights \(2m,0\), with rank-one unramified
quotient; away from \(\ell\) their inertia type
is \(1\oplus(\chi')^2\eta\theta\), with \(N=0\).

These representations are irreducible on every
open subgroup, even in a case where
\(\Psi^4=1\).
Indeed failure of strong irreducibility in
dimension two gives either finite projective
image or a nontrivial quadratic self-twist
\(\alpha\).
Finite projective image contradicts the two
distinct Hodge--Tate weights on an open local
subgroup.
In the self-twist case, reduction of (RS2)
gives
\[
 \{1,\bar\Phi_2\}
          =\{\bar\alpha,\bar\alpha\bar\Phi_2\},
       \qquad\bar\Phi_2=\Psi^2\eta .
\]
Since \(\alpha\) is nontrivial quadratic in
odd characteristic, this forces
\(\bar\alpha=\bar\Phi_2\).

An isomorphism of the local ordinary
representation with its \(\alpha\)-twist
preserves its unique higher-weight stable line.
Uniqueness follows because any map from a
rank-one character of Hodge--Tate weight \(2m\)
to the rank-one weight-zero quotient is zero.
It consequently identifies the weight-zero
quotient with its \(\alpha\)-twist, which
forces \(\alpha|_{G_{\mathbb Q_\ell}}=1\).
But the restriction of
\(\bar\Phi_2=\Psi^2\eta\) to a place above
\(\ell\) of the selected real base is nontrivial:
\(\eta\) is trivial on that local group and
\(\Psi^2\ne1\) there.
This is a contradiction.
This argument uses the actual full local
non-equality and does not assume that
\(\Psi^4\ne1\) globally.

On \(G_L\), put
\[
             s_2=\rho_{f_2}\varepsilon^{-m}(\chi')^{-1}.
                                                               \tag{RS3}
\]
Its determinant is one.
Its residual ordered tuple is
\((\Psi,\Psi^{-1})\), its ordinary weights
are \(m,-m\), and its away-\(\ell\) inertia
type is \((\chi',(\chi')^{-1})\).
The \(\eta,\theta\) factors disappear on \(G_L\).
These are the same outer characters, in the
same order, as for the rank-three adjoint seed.

The pairing is symmetric conjugate self-dual,
rather than only an alternating self-dual
pairing.
Let \(A_2\) be the standard alternating
matrix on the two-dimensional space and use
the rational representation in (RS3)
before restriction.
For a complex conjugation \(c\),
\(\det s_2(c)=-1\) and \(s_2(c)^2=1\).
Hence
\[
                      P_2=A_2s_2(c),\qquad P_2^t=P_2 ,
\]
and, on \(G_L\),
\[
                    s_2(g)^tP_2s_2(cgc^{-1})=P_2 .
                                                               \tag{RS4}
\]
Both equations follow from
\(s_2(c)^tA_2s_2(c)=-A_2\) and
\(s_2(g)^tA_2s_2(g)=A_2\).
This supplies the exact paired rank-two
polarization.

For standard unitary normalization use a
fixed concrete CM character.
The elliptic curve
\[
                               E:\ y^2=x^3+1
\]
has CM by \(\mathbb Z[\zeta_3]\), with all
endomorphisms defined over
\(K=\mathbb Q(\sqrt{-3})\).
Over a finite coefficient extension its
Tate module restricted to \(G_K\) splits
into the characters \(\lambda,\lambda^c\).
The Weil pairing gives
\(\lambda\lambda^c=\varepsilon\).
Take
\[
                              \kappa=\lambda^{-1}|_{G_L}.
                                                               \tag{RS5}
\]
Then
\(\kappa\kappa^c=\varepsilon^{-1}\),
so \(s_2\kappa\) has the standard rank-two
unitary multiplier.
The displayed integral equation of \(E\)
has discriminant \(-432\), so all its
bad finite primes lie above \(2,3\),
already included in the level packet.
At the generic coefficient primes here it
has good reduction, and its characters are
crystalline.
Its finite inertia characters away from
\(\ell\) are killed by \(6\).
To see this directly, pass locally to a finite
extension over which \(E\) has good reduction.
The finite inertia action on the good special
fibre commutes with the specialized
\(\mathbb Z[\zeta_3]\)-action, since these
endomorphisms are defined over \(K\).
In the endomorphism algebra its centralizer is
\(K\); intersecting with integral automorphisms
gives \(\mathcal O_K^\times=\mu_6\).
This proves the assertion also at \(2,3\),
where the full good-fibre automorphism group
can be larger.
Since \(\ell\ge5\), this is prime-to-\(\ell\)
inertia.
These characters are retained in the common
residual and tame coefficient packet for the
scalar twist (RS5).
Thus (RS5) does not add an uncontrolled
finite ramification set.

The transfer identity in \srcid{IsTwist.ver}
also holds with its exact sign.
For \(t\in G_{\mathbb Q}\setminus G_K\), the
Tate representation interchanges the two CM
lines, so in their basis it has matrix
\(\left(\begin{smallmatrix}0&u\\v&0\end{smallmatrix}\right)\).
Its determinant is \(\varepsilon(t)=-uv\).
The square acts on either line by \(uv\);
hence
\[
        \lambda(t^2)=-\varepsilon(t),\qquad
        \kappa(t^2)=-\varepsilon(t)^{-1}.                     \tag{RS6}
\]
After restriction this is the same identity
for the selected CM packet and every allowed
element of its nontrivial real coset.


Strong irreducibility makes cyclic base
change of \(f_2\) to \(L\) cuspidal.
Classical CM theory identifies \(\lambda\)
with its algebraic Hecke character.
Twisting that base change by the algebraic
characters in (RS3)--(RS5) therefore gives
an actual regular algebraic conjugate
self-dual cuspidal rank-two representation
with Galois realization \(s_2\kappa\).

Its local inducing characters are unlinked.
For the unramified complementary components,
the same Gelbart--Jacquet bound
\(|t|<1/4\) used in the rank-three seed
bounds the ratio of the two roots by
\(q^{2|t|}<q\).
At a ramified primitive-level component the
ramified ratio and unitarity give \(t=0\).
Residue-degree base change preserves the
strict ratio inequality.
The scalar twists (RS3)--(RS5) cancel
from every ratio.

Finally \(h=I_2\) is the actual definite
hermitian form over \(L/F^+\).
It is positive at every real place.
At a nonsplit finite place the CM
extension is unramified, and \(-1\)
is a norm of a unit; thus \(I_2\)
is isometric there to the hyperbolic
hermitian plane.
At split places its group is
\(\mathrm{GL}_2\).
The even degree of \(F^+\) is supplied
by the field construction.
This is the required pure inner
unitary group.
The rank-independent stable-parameter
descent and local ordinary-vector
construction, with the exact
\(\theta,U_j\) normalizations already
proved for the common module, insert
(RS3)--(RS5) in its rank-two localized
ordinary factor.
The residual tuple and finite types
are exactly its prescribed tuple,
including the scalar \(\kappa\).
Thus the actual rank-two Hecke ideal
vanishes at this seed, at arbitrarily
large compatible \(m\).

This is the rank-two initial-point
producer called \srcid{Seed2Alt}.
Its classical constituent is the same
\srcid{BirthEG.propEGP\_of\_printsH}
construction, with finite character
\((\chi')^2\eta\).
The symmetric pairing and scalar
normalization above are the data used
by
\srcid{CurveDVRSpecialization.SeedPtOf}
and \srcid{PrLinkedEnd}.
The construction supplies a stable
rank-two seed; it does not replace
the paired rank-three support
argument by individually polarized
residual constituents.

\section{The arithmetic frame with a bound chosen before depth}
\label{sec:arithmetic-frame-uniform-depth}

The constant in the curve argument is chosen before the
cyclotomic depth.  Finiteness of cohomology at one varying
number field would not justify this order.
We prove the required uniformity and the two strict
classes on the same fields as the seed.

Fix an odd prime \(\ell\), a finite field \(k\) of
characteristic \(\ell\), an imaginary quadratic character
\(\eta\), and a continuous finite-order character
\(\Psi:G_{\mathbb Q}\to k^\times\) with \(\Psi(c)=1\).
Its order is prime to \(\ell\).
In the principal application
\(\eta=\chi_{\mathbb Q(\sqrt{-3})}\) and
\(\Psi=\bar\varepsilon^a\bar\psi
      =\bar\varepsilon^j\bar\chi'\).
The fixed totally real abelian field \(F_b\) has
conductor prime to \(\ell\).
For the principal generic packet it is the
biquadratic field constructed in
Section~\ref{sec:seed-all-generic-primes}.
The argument also applies to a fixed real quadratic
base after the local conditions below have been checked.
It has the following checked properties:
\begin{enumerate}
\item \(F_bK/F_b\), where \(K=\mathbb Q^{\ker\eta}\),
is unramified at every finite place and split at the
prescribed level primes, including \(\ell\).
\item At \(\ell\), \(F_b/\mathbb Q\) is unramified,
\(\Psi|_{I_\ell}=\bar\varepsilon^j\), and at every
place \(v\mid\ell\) of \(F_b\) the character
\(\Psi^2|_{G_{F_{b,v}}}\) is nontrivial.
The generic packet has the stronger property
\(2j\not\equiv0\pmod{\ell-1}\), which verifies this
condition on inertia itself.
For an exceptional coefficient prime the stated
condition may instead be verified on the full
decomposition group.
\item There is a prime \(p_0\), split in \(F_bK\), with
\(\Psi(\operatorname{Frob}_{p_0})=1\) and
\(\bar\varepsilon(\operatorname{Frob}_{p_0})=-1\).
\end{enumerate}
All finite ramification of \(\eta,\Psi\) is included in
a fixed integer \(M\), divisible by \(2\ell\).

Choose the fresh prime \(w\) used for the Eisenstein
seed, with \(w\equiv1\pmod{\ell}\) and with
\(\eta,\Psi,\bar\varepsilon\) trivial on its decomposition
group.  Let \(F_\theta/\mathbb Q\) be its cyclic
degree-\(\ell\) field.
Put
\[
 \begin{gathered}
 C_r=\mathbb Q(\zeta_{\ell^{r+1}})^{\mu_{\ell-1}},
 \qquad F_r^+=F_bF_\theta C_r,\qquad L_r=F_r^+K,\\
 E_0=F_bF_\theta K\,\mathbb Q^{\ker\Psi}\,
                                  \mathbb Q(\mu_\ell),
 \qquad E_r=E_0C_r.
 \end{gathered}                                             \tag{AF1}
\]
The finite abelian field \(E_0\) is fixed before \(r\).
The degrees of \(E_r/F_r^+\) and \(E_r/L_r\) are prime to
\(\ell\).

\subsection{One class-group bound throughout the cyclotomic tower}
\begin{srlemma}\label{lem:frame-class-rank-uniform}
There is \(C_0\ge0\), depending on the fixed field \(E_0\)
and on \(\ell\), such that
\[
       \dim_{\mathbb F_\ell}
             \operatorname{Cl}(E_r)/\ell\operatorname{Cl}(E_r)
                                      \le C_0
                         \qquad(r\ge0).                       \tag{AF2}
\]
No bound on the numerical size of \(C_0\) is needed here.
\end{srlemma}
\begin{proof}
We use the classical cyclotomic Iwasawa module and the
\href{https://doi.org/10.2307/1971116}
{Ferrero--Washington theorem} in its actual scope:
the base field \(E_0\) is a finite abelian extension of
\(\mathbb Q\), and the tower is its cyclotomic
\(\mathbb Z_\ell\)-extension.
Write
\[
 A_r=\operatorname{Cl}(E_r)[\ell^\infty],
       \qquad X=\varprojlim_r A_r
\]
with transition maps given by norms.
The Iwasawa module theorem makes \(X\) a finitely
generated torsion \(\mathbb Z_\ell[[T]]\)-module.
Ferrero--Washington gives \(\mu(X)=0\).
The structure theorem therefore makes \(X\) a finite
\(\mathbb Z_\ell\)-module in the sense of finite generation.
Indeed the elementary modules
\(\mathbb Z_\ell[[T]]/(f(T)^e)\), with \(f\) distinguished,
are finite free over \(\mathbb Z_\ell\), and the finite
kernels and cokernels in the structure theorem preserve
finite generation.

For completeness, projections to finite class groups
are surjective after discarding finitely many initial
layers.
In a cyclotomic tower all primes above \(\ell\) are
totally ramified beyond a fixed layer.
Thus no nontrivial subextension of \(E_t/E_r\), for
\(t\ge r\) beyond that layer, is unramified.
Its intersection with the Hilbert class field of \(E_r\)
is \(E_r\).
Given any ideal class of \(E_r\), Chebotarev supplies a
prime in that class which splits in \(E_t/E_r\).
A prime above it has that class as its norm.
The norm on class groups is consequently surjective,
and so is the norm on their \(\ell\)-primary parts.
Compactness of the finite inverse system now gives
\(X\twoheadrightarrow A_r\).
Hence
\(\dim_{\mathbb F_\ell}A_r/\ell A_r
 \le\dim_{\mathbb F_\ell}X/\ell X\)
for all those layers.
Enlarge this bound to cover the finitely many initial
layers.  This proves (AF2).
\end{proof}

This is the actual argument of
\srcid{FrobeniusSmoothVectors.RankLeOfCardLe.lev\_rk}
and \srcid{b0}.
The source fixes
\[
 V=H_0\cap H_{\rm th}(w)\cap H_b\cap
                         \ker\eta\cap\ker\Psi\cap\ker\bar\varepsilon
\]
and uses the exact equality
\[
 U_r\cap\ker\eta\cap\ker\Psi\cap\ker\bar\varepsilon
              =V\cap\ker(\varepsilon\bmod\ell^{r+1}).           \tag{AF3}
\]
The fixed field of \(V\) is the fixed abelian base in
that proof.  Formula (AF1) permits an additional fixed
real abelian factor; the same proof applies verbatim.
A tower of unrelated fresh odd conductors would not
give (AF2) from Ferrero--Washington.

\subsection{The three cohomology bounds defining the frame}
Let \(\Gamma_r\) be the quotient of \(G_{F_r^+}\) by the
closed normal subgroup generated, away from \(\ell\), by
the inertia elements on which both \(\eta\) and \(\Psi\)
are trivial.
Let \(\delta\) denote its induced quadratic character,
\(\Gamma_{r,0}=\ker\delta\), and use the actual local
decomposition groups above \(\ell\) as local maps.
For the CM subgroup use both members of each conjugate
pair of places.

For a representation \(W\) and stable local subspaces
\(b_v\), denote by \(H_C^i(W,b)\) the mapping-cone
cohomology of
\[
 C^\bullet(\Gamma_r,W)\longrightarrow
          \prod_{v\mid\ell}C^\bullet(G_{F_{r,v}^+},W/b_v),
\]
with the cone shifted so that its degree-one elements
are cocycles \(z\) together with local correcting
vectors \(Y_v\), modulo global coboundaries.
When \(b_v=0\) it is the strict local condition.

\begin{srlemma}\label{lem:frame-three-bounds}
After enlarging the same \(C_0\) if necessary,
simultaneously for every depth \(r\),
\[
 \begin{split}
 \dim_kH_C^1(\Gamma_r,k(\Psi^{-2}\delta),0)&\le C_0,\\
 \dim_kH_C^1(\Gamma_{r,0},k(\Psi^{-1}),0)&\le C_0,\\
 \dim_kH^1(\Gamma_r,k)&\le1+C_0 .
 \end{split}                                                 \tag{AF4}
\]
These are precisely \srcid{Frame.Ok.b1p},
\srcid{b1}, and \srcid{h1}.
\end{srlemma}
\begin{proof}
At every ordinary local group, \(\Psi^2\ne1\).
Indeed this was checked over \(F_b\), and the subsequent
local extensions have \(\ell\)-power degree; restriction
through them cannot decrease the nontrivial image of a
character of order prime to \(\ell\).
The local invariant spaces in the first two lines are
therefore zero.
Consequently forgetting the local correcting vectors
injects their mapping-cone \(H^1\) into global \(H^1\):
two local primitives of a global boundary differ by a
local invariant, which is zero.

Restrict further to \(G_{E_r}\).
The intermediate quotient has order prime to \(\ell\),
so averaging kills the kernel in inflation--restriction.
Both coefficient characters are now trivial.
The quotient defining \(\Gamma_r\) kills inertia away
from \(\ell\) once \(\eta,\Psi\) are trivial, and the
strict condition kills inertia above \(\ell\).
Thus the restrictions are everywhere-unramified
homomorphisms \(G_{E_r}\to k\).
Class field theory identifies them with subspaces of
\(\operatorname{Hom}(\operatorname{Cl}(E_r),k)\).
Its dimension over \(k\) is the
\(\mathbb F_\ell\)-dimension in (AF2).
This proves the first two bounds.

For the third line, restriction injects into
\(H^1(G_{E_r,S_\ell},k)\).
Kummer theory, followed by the finite prime-to-\(\ell\)
quotient invariants, identifies its possible image
with the cyclotomic-character part of
\[
             E_r^\times/(E_r^\times)^\ell
\]
which is unramified away from \(\ell\).
There is the exact Kummer sequence with terms
\[
 \mathcal O_{E_r}[1/\ell]^\times/
                       (\mathcal O_{E_r}[1/\ell]^\times)^\ell,
 \qquad
 \operatorname{Cl}(\mathcal O_{E_r}[1/\ell])[\ell].
                                                               \tag{AF5}
\]
The class-group term has dimension at most \(C_0\).

The cyclotomic part of the valuation module at
\(\ell\) is zero.
Indeed the tame cyclotomic inertia is nontrivial at
each such place over \(F_r^+\); it fixes that place
but acts nontrivially through \(\bar\varepsilon\).
The coefficient of an eigenvector at that place must
therefore vanish.
Hence the relevant \(S_\ell\)-unit part in (AF5)
comes from ordinary units.

The field \(E_r\) is CM.
For any unit \(u\), the unit \(u/\overline u\) has
absolute value one at every archimedean embedding;
it is a root of unity.
Since two is invertible modulo \(\ell\), the
minus part of units modulo \(\ell\) comes from
roots of unity.
The character \(\bar\varepsilon\) is odd, and the
\(\ell\)-primary roots of unity form a cyclic group.
Its contribution has dimension at most one.
This proves the final bound in (AF4).
\end{proof}

\subsection{The two actual nonzero strict classes}
The two nontriviality fields of the frame do not require
assuming the existence of an unexplained deformation.
They are forced by one additional ramified place.

\begin{srlemma}\label{lem:frame-strict-classes}
For \(a'=1,2\), put \(\kappa_{a'}=\Psi^{-a'}\eta\).
The quotient \(\Gamma_r\) has
\[
                H_C^1(\Gamma_r,k(\kappa_{a'}),0)\ne0
                                  \qquad(r\ge0).              \tag{AF6}
\]
\end{srlemma}
\begin{proof}
The character \(\kappa_{a'}\) is odd.
It is nontrivial on every full local decomposition
group at \(\ell\), by the displayed \(\Psi^2\ne1\)
condition and splitting of the CM extension there.
Nontriviality on inertia is not required.
It differs from \(\bar\varepsilon\): at the selected
scalar prime \(p_0\) it has value one, whereas
\(\bar\varepsilon\) has value minus one.

Over \(\mathbb Q\), take the finite ramification set
consisting of \(\ell,w\), and the primes ramifying in
\(\kappa_{a'}\).
Use zero local conditions at \(\ell\) and at the
ramified primes, and the full local \(H^1\) at \(w\).
The Selmer dimension formula from finite
Poitou--Tate duality gives
\[
 \dim_k H^1_{\mathcal L}(\mathbb Q,k(\kappa_{a'}))
 -\dim_k H^1_{\mathcal L^\perp}
           (\mathbb Q,k(\kappa_{a'}^{-1}\bar\varepsilon))
                                                        =1 . \tag{AF7}
\]
Here are all correction terms.
Both global invariant spaces are zero because
\(\kappa_{a'}\ne1,\bar\varepsilon\).
The archimedean invariant space is zero because
\(\kappa_{a'}\) is odd.
At \(\ell\), and at the primes ramifying in
\(\kappa_{a'}\), local invariants are zero.
At \(w\), both \(\kappa_{a'}\) and
\(\bar\varepsilon\) are trivial; local duality and the
local Euler formula give \(h^1=2\), \(h^0=1\).
Thus \(w\) contributes exactly one.
At unramified primes outside the set, the unramified
condition has dimension \(h^0\) and contributes zero.
These are exactly the terms in (AF7).
In particular there is a nonzero Selmer class \(z\).

The class stays nonzero on \(G_{F_r^+}\).
The extension \(F_r^+/\mathbb Q\) is Galois, and
\(\kappa_{a'}\) is nontrivial on a full ordinary local
decomposition group over \(F_r^+\), as just proved.
Its invariant space over \(G_{F_r^+}\) is therefore zero.
Inflation--restriction therefore makes restriction
on \(H^1\) injective.

At \(w\), the character is trivial, and a cocycle
restricts on inertia to a homomorphism into the
additive group of \(k\).
Wild inertia has pro-\(w\) image and is killed.
The remaining \(\ell\)-part of tame inertia is
\(\mathbb Z_\ell(1)\).
The degree-\(\ell\) field \(F_\theta\) is totally
ramified at \(w\); inertia after restriction to it
maps into \(\ell\mathbb Z_\ell(1)\).
Our homomorphism is zero there.
At every other finite prime different from \(\ell\),
a local boundary vanishes on inertia whenever its
coefficient character is one; at primes outside
the ramification set the class is unramified.
Thus \(z\) vanishes on exactly the inertia elements
normally killed in defining \(\Gamma_r\).
Products, conjugates and continuity give vanishing
on their closed normal closure.
The class descends to \(\Gamma_r\).
It is locally a boundary at every place above
\(\ell\), so it supplies (AF6), including its actual
local correcting vectors.
No nonzero target character group, large Selmer
length, or desired character estimate was assumed
in this construction.
\end{proof}

\subsection{Completing the frame and then choosing depth}
We record the remaining arithmetic assertions, to
make the use of \srcid{Frame.Ok} explicit.

The group \(\Gamma_r\) is profinite.
Its quadratic character \(\delta\) has open kernel,
complex conjugation is an involution with
\(\delta(c)=-1\), and the induced characters
\(\Psi,\varepsilon\) are continuous.
There are \(d=[F_r^+:\mathbb Q]\) archimedean places,
and the sum of the ordinary local degrees is \(d\).
The marked element \(g_0\) with \(\Psi(g_0)^2\ne1\)
may be selected in a full ordinary local decomposition
group at \(\ell\).  In the generic packet it may be
selected in inertia.

Every residue degree at \(p_0\) in \(F_r^+K\) is odd:
the fixed real quadratic factors split there, and
the remaining extensions have \(\ell\)-power degree.
Thus a local Frobenius \(s_0\in\ker\delta\) has
\[
                         \Psi(s_0)=1,\qquad
                  \bar\varepsilon(s_0)=-1 .                  \tag{AF8}
\]
It proves all three nonvanishing requirements
\(\bar\varepsilon\ne1\),
\(\bar\varepsilon\Psi^{\pm1}\ne1\), and
\(\bar\varepsilon\Psi^{\pm2}\delta\ne1\).
It also excludes
\(\bar\varepsilon=\Psi^b\) on \(\ker\delta\)
for every integer \(b\).
Local square separation holds at every ordinary place
by the full-decomposition condition above and its
preservation under the \(\ell\)-power layers.
On the inertia kernel of \(\varepsilon\), \(\Psi=1\).

A completion of \(F_r^+\) at \(\ell\) contains no
\(\mu_\ell\): its ramification index is an
\(\ell\)-power.
Its pro-\(\ell\) unit group is therefore a free
\(\mathbb Z_\ell\)-module of rank \(d_v\).
Choose a basis adapted to the cyclotomic norm,
whose kernel has rank \(d_v-1\).
All basis elements lie in the kernel of \(\Psi\),
and all but one lie in the kernel of \(\varepsilon\).
Any continuous character into a finite abelian
\(\ell\)-group vanishing on this basis vanishes on
inertia.
This is the exact \srcid{Frame.Ok.gen} assertion.
A pro-\(\ell\) character of \(\ker\delta\) unramified
at every ordinary place is unramified at all finite
places: the other allowed inertia quotients have
order prime to \(\ell\).
It factors through the finite class group of \(L_r\).
Its exponent supplies \srcid{Frame.Ok.cl}.

The finite-cohomology condition \(\Phi_\ell\) follows
from the ordinary arithmetic finiteness theorem on
every open subgroup: its finite representations
are unramified outside a finite set, and restriction
to a finite splitting field reduces degree-one
finiteness to the unit and class-group calculation.
This is a finiteness assertion, not a uniform
bound in \(r\).

Finally the Euler--Poitou--Tate comparison for the
same global quotient and these local maps gives,
for a finite \(k'\)-representation \(W\) and stable
local subspaces \(b_v\), with
\(\operatorname{Hom}_{\Gamma_r}
       (W,k'(\bar\varepsilon))=0\),
\[
 h_C^2(W,b)+h_C^0(W,b)+\sum_v d_v\dim b_v
       \le h_C^1(W,b)+\sum_{u\mid\infty}\dim W^{c_u}.            \tag{AF9}
\]
This is the precise \srcid{Frame.Ok.ec} statement.
It uses the standard finite global and local duality
already proved in the cohomology part of the manuscript.
The local unramified conditions away from \(\ell\)
have Euler contribution zero.
The last possible cone degree is the global dual
invariant term, which vanishes by the displayed
\(\operatorname{Hom}\) hypothesis.
Thus the degree-two cone is finite and its Euler
formula gives (AF9).
For the normalized residual matrix representations
used below, the dual invariant hypothesis is checked
by triangular entries: the diagonal differences are
\(0,\pm1,\pm2\); (AF8) excludes the first two sorts,
and complex conjugation gives the required
\(\delta\)-factor on the opposite-corner
\(\pm2\) entries.
This is the matrix proof
\srcid{FrobeniusSmoothVectors.PrDual.pr\_dual}.

Equations (AF4), (AF6), and the preceding assertions
verify every field of \srcid{Frame.Ok}\((C_0)\).
Only now choose \(r\) with
\[
                        \ell^r>3C_0+7 .                       \tag{AF10}
\]
Since every \(d_v\ge\ell^r\) and \(d\ge d_v\),
\[
                         d>3C_0+7,\qquad d_v>5 .
\]
These are exactly \srcid{Frame.Big}\((C_0)\).
Any additional prescribed global or local depth
bound may be imposed by further increasing \(r\);
the same \(C_0\) still applies.

\paragraph{Exact source correspondence and scope.}
The actual producers are
\srcid{RankLeOfCardLe.b0},
\srcid{FrameEx.ok\_link},
\srcid{FrameEx.frame\_ex},
\srcid{Strict.st\_char}, \srcid{Strict.strict},
and
\srcid{FrobeniusGaloisFieldFourCompatibility}.
They are under
\srcid{Development.AutomorphicArithmetic.AutomorphicVarieties.%
FrobeniusSmoothVectors}
at the frozen tag \srcid{alpoge-source-7cbf7a0c}.
The three comparison bounds are assembled by
\srcid{GaloisCochains.B1p},
\srcid{B1}, and \srcid{H1}.
The additional fixed real abelian factors in (AF1)
do not require a new Iwasawa theorem:
their compositum is absorbed into \(E_0\) before
the cyclotomic depth is chosen.

This proves the actual arithmetic frame used by
the curve and sharp-ideal arguments on the
principal nonresonant packet, and on any exceptional
packet for which the same displayed fixed-base and
full-decomposition conditions have been verified.
It supplies auxiliary existence and depth; it
does not assert a polynomial numerical bound for
\(C_0\), nor use such a bound in the final
host-uniform character estimate.

\section{The actual entry space and its flat local deformation charts}
\label{sec:entry-space-flat}

Fix the arithmetic frame of
Section~\ref{sec:arithmetic-frame-uniform-depth},
with \(\operatorname{Frame.Ok}(C_0)\) and
\(\operatorname{Frame.Big}(C_0)\).
Write \(d=\sum_vd_v\), let \(\mathcal O\) be a complete
coefficient DVR with residue field \(k\), and use ranks
\(n=2,3\).
The weights are \(w=(1,-1)\) and \(w=(1,0,-1)\),
respectively; \(J\) is the reversal matrix, and
\(g^*=cg^{-1}c^{-1}\).

The normalized representations used here consist of
actual matrices and ordered local characters satisfying
\[
 \begin{split}
 &\rho(gh)=\rho(g)\rho(h),\qquad
                 \rho(g^*)=J\,{}^t\rho(g)J,\\
 &(\rho(g_0)-\rho(g_0^*))_{ij}=0\quad(i\ne j),\\
 &\prod_{i=0}^{n-1}
       \bigl(\rho(\operatorname{loc}_v(g_i))
                         -\psi_{vi}(g_i)I_n\bigr)=0
                                  \quad\hbox{for every }(g_i).
 \end{split}                                                  \tag{LS1}
\]
The product is in the displayed order.
The residual diagonal and local characters are
\(\Psi^{w_i}\).
The ideals \(\mathfrak t_+,\mathfrak t_-\) are generated
by entries below and above the diagonal.
In rank three \(\mathfrak e\) is generated by the four
outer--middle entries.
Type I means off both triangular loci and off
\(V(\mathfrak e)\); Type E means off both triangular
loci and on \(V(\mathfrak e)\).

\subsection{Construction and algebraization of the one entry space}
There is a complete Noetherian local ring
\(R^{\rm diag}\) representing continuous data (LS1)
with the diagonal residual representation.
Here is the actual construction.
For the continuous framed representation functor
with its fixed diagonal residue, the condition
\(\Phi_\ell\) makes the relative tangent space finite.
The functor preserves fibre products, including the
fibre products used in the infinitesimal uniqueness
condition: a pair of compatible matrices has a
unique matrix lift to the fibre product.
The standard successive small-extension construction
therefore gives its universal complete local ring.
There are no automorphisms to identify, since the
basis is part of the data.
Lifting a tangent basis gives a surjection from a
finite-variable \(\mathcal O\)-power-series ring.

Adjoin the independent ordered local characters.
For each \(v\), their deformation ring is the
completed group algebra of the maximal pro-\(\ell\)
quotient of \(G_v^{\rm ab}\), translated by their
fixed Teichm\"uller residue characters.
Local class field theory makes that quotient a
finite \(\ell\)-group times \(\mathbb Z_\ell^{d_v+1}\).
Thus the completed tensor product of these finitely
many character rings with the framed ring is again
complete Noetherian.
In this ring take the ideal generated by every
entry of the polarization, marking, and ordered
product equations in (LS1).
This ideal is closed: every ideal of a complete
Noetherian local ring is closed.
Its quotient \(R^{\rm diag}\) represents precisely
the continuous data (LS1).
Indeed a continuous map kills the ideal exactly
when it kills all its displayed generators, and
all equations are preserved under the inverse
limit of the finite coefficient quotients.

For each ordinary place choose \(s_v\) with
\(\Psi(\operatorname{loc}_v s_v)^2\ne1\).
Hensel diagonalization at its distinct residual
eigenvalues gives inverse matrices \(P_v,Q_v\),
congruent to one, diagonalizing \(\rho(s_v)\).
Impose the polarization and marked-element
equations and the entries below the diagonal of
\(Q_v\rho(\operatorname{loc}_v s)P_v\), for all \(s\).
Its diagonal entries give the ordered characters.
Conversely (LS1) implies these same equations.
Set every \(g_i=s_v\) except one.
The other factors are diagonal.
The \((a,b)\)-entry for \(b\le a\) isolates
\[
       \bigl(Q_v\rho(\operatorname{loc}_v s)P_v
                            -\psi_{va}(s)I_n\bigr)_{ab}
\]
times products of unit eigenvalue differences.
For \(n\le3\), direct multiplication of the at
most three factors proves that it is zero.
Thus this quotient represents the stated ordinary
functor, including its ordered characters.

Give \(\rho(g)_{ij}\) weight \(w_i-w_j\) and each
local character weight zero.
Conjugation by \(\operatorname{diag}(t^{w_i})\)
preserves (LS1) and gives a completed torus grading.
Define the actual source rings
\[
 \begin{split}
 B&=\overline{\mathcal O[
       \text{weight-zero entry monomials},\,\psi_{vi}(s)]}
                                \subset R^{\rm diag},\\
 A&=B[\rho(g)_{ij}:g,i,j]\subset R^{\rm diag}.
 \end{split}                                                  \tag{LS2}
\]
The closure is in the maximal-ideal topology of
\(R^{\rm diag}\).

We record their finiteness and topology.
The entries and characters topologically generate
\(R^{\rm diag}\), by the representing property.
Choose finitely many homogeneous such elements
spanning its relative cotangent space.
They give a graded surjection from
\(\mathcal O[[X_1,\ldots,X_h]]\).
The exponent semigroup of weight-zero monomials
is finitely generated.
For each fixed weight the exponent vectors
have finitely many coordinatewise minimal
elements, by Dickson's lemma; every other
exponent is one of these plus a weight-zero
exponent.
Decomposing a formal series accordingly proves
that the degree-zero part is complete Noetherian
and each fixed-weight part is finite over it.
These assertions pass to the homogeneous quotient,
by projecting coefficients to their weight.
Thus the degree-zero part is exactly \(B\),
each \(A_j\) is finite over \(B\), and \(A\) is
a finitely generated graded \(B\)-algebra.
In particular \(A\) is Noetherian.
On each \(A_j\) the induced topology is cofinal
with the \(\mathfrak m_B\)-adic topology: a long
fixed-weight monomial contains arbitrarily
many nonconstant weight-zero factors.

This one algebra also represents the points
with nondiagonal triangular residue.
For clarity we give the actual algebraization
argument.
Suppose \(Q\) is complete local Noetherian with
finite residue extension of \(k\), and \(\rho_Q\)
satisfies (LS1), continuously, with residual
diagonal entries and local characters
\(\Psi^{w_i}\).
Its residual representation is upper or lower
triangular.
Here are the coordinate equations for this assertion.
Write \(a_{ij}(g)=\bar\rho_Q(g)_{ij}\).
Each diagonal function is the character
\(\Psi^{w_i}\), so for every \(g,h\)
\[
 \sum_{j\ne i}a_{ij}(g)a_{ji}(h)=0.
\]
The marked matrix
\[
 X=\bar\rho_Q(g_0)-\bar\rho_Q(g_0^*)
                         =\operatorname{diag}(\zeta,0,-\zeta)
\]
has \(\zeta\ne0\).
Comparing the \(i\)-th diagonal entries of
\(\bar\rho_Q(g)X\bar\rho_Q(h)\) gives the second equation
\[
 \sum_{j\ne i}X_{jj}a_{ij}(g)a_{ji}(h)=0.
\]
Since \(2\zeta\ne0\), these two equations imply
\[
 a_{01}(g)a_{10}(h)=
 a_{02}(g)a_{20}(h)=
 a_{12}(g)a_{21}(h)=0
                              \quad\text{for all }g,h.
\]
Polarization gives
\(a_{01}(g^*)=a_{12}(g)\) and
\(a_{10}(g^*)=a_{21}(g)\).
Consequently either \(a_{01}=a_{12}=0\) identically,
or \(a_{10}=a_{21}=0\) identically.
Also either \(a_{02}=0\) or \(a_{20}=0\) identically.

Only the two mixed cases require another equation.
If \(a_{01}=a_{12}=a_{20}=0\), the \((0,1)\)-entry
of \(\bar\rho_Q(gh)=\bar\rho_Q(g)\bar\rho_Q(h)\) gives
\(a_{02}(g)a_{21}(h)=0\).
Either \(a_{02}=0\), giving lower triangularity,
or \(a_{21}=a_{10}=0\), giving upper triangularity.
If \(a_{10}=a_{21}=a_{02}=0\), the \((2,1)\)-entry
instead gives \(a_{20}(g)a_{01}(h)=0\);
its two alternatives give upper or lower
triangularity in the same way.
The other cases are already triangular.
In rank two the one diagonal equation is
\(a_{01}(g)a_{10}(h)=0\), which proves the assertion.
For points of the fibre in (LS2), the same
conclusion follows immediately: nonzero entries
of opposite weights would give a nonzero
weight-zero product in the residual augmentation.

Put \(Q_0=\{q\in Q:\bar q\in k\}\).
It is complete local with residue \(k\).
The ring \(Q\) is finite over \(Q_0\):
lift a residue basis including \(1\); subtract
its residue combination from any element.
The remainder belongs to
\(\mathfrak m_Q\subset Q_0\).
These lifts generate \(Q\) over \(Q_0\).
Noetherianity descends along this finite extension,
and the maximal ideals of \(Q_0,Q\) have the
same underlying set and the same powers.

Let \(P_j\subset Q\) be the \(Q_0\)-span of
entry monomials of weight \(j\), and form
\[
                          D=\bigoplus_jP_jt^j
                                    \subset Q[t,t^{-1}].      \tag{LS3}
\]
Every weight-zero off-diagonal monomial belongs
to \(\mathfrak m_Q\subset Q_0\); diagonal and
local-character entries belong to \(Q_0\).
Thus \(P_0=Q_0\), and the polynomials whose
degree-zero coefficient lies in
\(\mathfrak m_{Q_0}\) form a maximal ideal
\(\mathfrak n\) of \(D\).
Products of opposite nonzero degrees have that
property.
Each entry span is finite over \(Q_0\), so
\(D\) is a finitely generated \(Q_0\)-algebra.
Each \(P_j\) is finite over \(Q_0\) by the same
fixed-weight monomial decomposition.
Its induced \(\mathfrak n\)-adic topology is
cofinal with its
\(\mathfrak m_{Q_0}\)-adic topology.
Indeed remove primitive zero-weight factors
from a long fixed-weight monomial.
There are finitely many such factors, of
bounded length; each lies in
\(\mathfrak m_{Q_0}\), and the remaining
minimal monomial has bounded length.

In \(\widehat D_{\mathfrak n}\), the entries
\(t^{w_i-w_j}\rho_Q(g)_{ij}\), together with
the degree-zero characters, give continuous
residually diagonal data (LS1).
The universal map from \(R^{\rm diag}\) exists.
Its image on \(B\) lies in the closed
degree-zero part \(Q_0\), and its image on
\(A\) lies in (LS3).
Evaluation at \(t=1\) gives
\[
                              A\longrightarrow Q              \tag{LS4}
\]
with exactly the original entries and characters,
and with \(\mathfrak m_B\) mapping into
\(\mathfrak m_Q\).
Continuity on each entry lattice follows from
Artin--Rees inside the finite \(Q_0\)-module
\(Q\).
Uniqueness holds on entries and characters,
then on \(B\) by fixed-weight adic approximation,
then on \(A\).

For locally constant lower data over an arbitrary
\(k\)-field or over \(k[T]\), their entries have
finite image and hence finite \(k\)-span.
All nonconstant weight-zero monomials vanish.
Use \(Q_0=k\) in (LS3).
The same vertex completion and evaluation give
the actual field and polynomial points of \(A\).
Consequently the previously constructed lower
interpolation families are families in this
one entry space.
Their connected core and dense nonexceptional
locus require no further point-identification
hypothesis.

\subsection{Flat local charts and their dimension formula}
Let \(x\) be a closed point of \(A\) above
\(\mathfrak m_B\), with residue field \(k'\).
Choose the finite unramified coefficient extension
\(\mathcal O'/\mathcal O\) with residue field \(k'\).
For its specified residue map let
\[
 R_x=\widehat{(A\otimes_{\mathcal O}\mathcal O')_{\mathfrak n_x}}.
                                                               \tag{LS5}
\]
The map \(A\to R_x\) is flat: coefficient extension,
localization, and Noetherian local completion
are flat.
After localizing \(A\) at \(x\), it is faithfully
flat.
By (LS4) and its uniqueness, (LS5) represents
all normalized deformations of the actual
residue representation at \(x\).
This proves the source's \srcid{DefIso} statement;
it does not assume a new unrelated deformation
ring.

Minimal primes of \(R_x\) map to generic points
of components through \(x\), and every such
component has a minimal-prime lift.
Going down and faithful flatness prove these
two assertions.
Thus (LS5) gives the required flat chart.

For a prime \(P\subset R_x\), put \(p=P\cap A\),
\(p_B=p\cap B\) and
\(t=\operatorname{trdeg}_{B/p_B}(A/p)\).
Then
\[
                       \dim R_x/P\le\dim B/p_B+t.             \tag{LS6}
\]
For each \(p\subseteq x\), equality is attained
at some prime of \(R_x\) above \(p\).
The finite coefficient extension preserves
transcendence degree.
The finite-type dimension formula at the closed
point gives
\(\dim(A_x/pA_x)=\dim B/p_B+t\).
Completion preserves this quotient dimension;
excellence makes the completion of the local
domain \(A_x/pA_x\) equidimensional.
Any prime over it has at most that dimension,
and its minimal primes have exactly that
dimension.
This proves (LS6) with its existence assertion.

\subsection{The actual tangent and obstruction maps}
Write \(R=R_x\).
Extend its residue representation to the unitary
semidirect product on the real group, with
complex conjugation of matrix part \(I_n\).
The adjoint action on the CM subgroup is
conjugation by \(\bar\rho\); on the other coset
it includes \(\vartheta(X)=-J\,{}^tXJ\).
Choose the actual ordinary flags and denote
their Borel Lie algebras by \(b_v\).
All cone groups below include local correcting
vectors as in (AF9).

For a small extension \(C'\twoheadrightarrow C\)
with kernel \(I\) and \(\mathfrak m_{C'}I=0\),
choose continuous setwise global matrix lifts
\(\widetilde r(s)\), flag lifts
\(\widetilde g_v\), and upper-triangular
local lifts \(\widetilde t_v(s)\).
They exist by finite sections on coefficient
quotients.
The defects are the matrices with entries in \(I\)
\[
 \begin{split}
 \Omega(s,t)&=\widetilde r(s)
    \operatorname{sgnAct}_{\delta(s)}(\widetilde r(t))
                                  \widetilde r(st)^{-1}-I_n,\\
 a_v(s)&=\widetilde r(\operatorname{loc}_v s)
   \bigl(\widetilde g_v^{-1}\widetilde t_v(s)
                                      \widetilde g_v\bigr)^{-1}
                                                   -I_n .
 \end{split}                                                  \tag{LS7}
\]
Associativity gives \(d\Omega=0\), and local
multiplication gives
\(\Omega|_{G_v}-da_v\in b_v\).
Applying \(u:I\to k'\) gives a class in
\(H_C^2(M_n,b)\).

If this class is zero, write \(u\Omega=d\beta\)
and
\(ua_v-\beta|_{G_v}
       +(\operatorname{Ad}(\bar r)-1)Y_v\in b_v\).
In the one-dimensional pushout selected by \(u\),
replace the global lifts by
\((1-\beta(s))\widetilde r(s)\), and the flags by
\(\widetilde g_v(1-Y_v)\).
Square-zero multiplication proves the corrected
representation and the ordinary condition.
For complex conjugation normalize \(1+X\),
where \(J\,{}^tXJ=X\), by \(1-X/2\).
For the marked difference use
\(Y_{ij}=-E_{ij}/(\lambda_j-\lambda_i)\)
on its off-diagonal error.
The denominators are units and its
antidiagonal symmetry makes \(\vartheta(Y)=Y\);
this conjugation preserves the first normalization.
Thus zero obstruction gives an actual
normalized ordinary lift.

Choose a minimal relative presentation
\(P=\mathcal O'[[X_1,\ldots,X_t]]\twoheadrightarrow R\),
with kernel \(I_R\).
The obstruction gives an injection
\[
              (I_R/\mathfrak m_PI_R)^*
                   \lhook\joinrel\longrightarrow H_C^2(M_n,b).
                                                               \tag{LS8}
\]
Indeed a nonzero kernel functional would give
a nonzero one-dimensional quotient of
\(P/\mathfrak m_PI_R\to R\) with a section.
A section is impossible for a minimal
presentation.
For a relation write
\(f=f(0)+\sum_iX_i h_i\), with
\(h_i\in\mathfrak m_P\).
Subtract its value under the section.
Every kernel element is then in the maximal
ideal times that kernel; Nakayama kills it.
This also treats a coefficient-uniformizer
relation.

A cone \(H^1\)-class gives first-order matrices
\((1+\epsilon z(s))\bar r(s)\) with their
corrected local flags.
Normalize them as above.
This shows that every class comes from a
relative tangent vector.
The torus matrix \(T_0=\operatorname{diag}(w_i)\)
is nonzero and its first-order conjugation
preserves both normalizations.
If its cocycle is nonzero it is a nonzero
tangent direction killed in \(H_C^1\);
if zero it is a nonzero cone invariant.
Consequently
\[
                    h_C^1+1\le t+h_C^0 .                     \tag{LS9}
\]
The local boundary rigidity in this comparison
is elementary: if
\((\operatorname{Ad}(\bar r(s))-1)Y\in b_v\)
for every local \(s\), induction on lower
entries, using the separated characters,
forces \(Y\in b_v\).

At each real place the unitary-adjoint fixed
space has dimension \(n(n-1)/2\):
multiplication by the pairing matrix identifies
it with skew-symmetric matrices.
The ordinary Borel has dimension \(n(n+1)/2\).
Equations (AF9), (LS8), and (LS9) give
\(t-\dim H_C^2\ge nd+1\).
Pad a generating list of \(I_R\) by zeros.
We obtain the actual ordinary presentation
\[
 R\simeq\mathcal O'[[X_1,\ldots,X_t]]/(f_1,\ldots,f_b),
                         \qquad t=b+nd+1 .                   \tag{LS10}
\]
The dual-invariant hypothesis in (AF9) is
satisfied here.
The residual matrices are triangular with
diagonal \(\Psi^{w_i}\).
Solving a twisted invariant by entries reduces
to \(\bar\varepsilon\Psi^j=1\), \(j=0,\pm1\),
or \(\bar\varepsilon\Psi^{\pm2}\delta=1\);
the latter sign comes from opposite-corner
transpose at complex conjugation.
The scalar element (AF8) excludes every case.

Suppose now \(n=3\) and \(\mathfrak e\subseteq x\).
The sign action extends to an involution of
the same \(R\), fixing \(\mathcal O'\) and \(B\).
Use \((1\pm\sigma)/2\) on tangent variables,
minimal relations, and the lifts in (LS7).
The obstruction maps split into the even and
odd cone groups.
Write \(a,g\) for the tangent dimensions,
and \(r_e,r_o\) for the relation numbers.
Equation (LS8) gives \(r_e\le h_{C,e}^2\),
\(r_o\le h_{C,o}^2\).
The tangent comparison gives
\[
 h_{C,e}^1+1\le a+h_{C,e}^0,\qquad h_{C,o}^1\le g .
\]
The extra direction is the even matrix
\(\operatorname{diag}(1,0,-1)\).
The local flags can be chosen compatible
with parity by projecting their separated
eigenvectors.
Their upper-triangular matrix units give
\(\dim b_{v,e}\ge4,\ \dim b_{v,o}\ge2\).
At a real involution the skew-matrix
identification bounds the two fixed spaces
by \(1\) and \(2\), respectively.
Apply (AF9) in each parity:
\[
 h_{C,e}^1-h_{C,e}^2\ge h_{C,e}^0+3d,\qquad
 h_{C,o}^1-h_{C,o}^2\ge h_{C,o}^0 .
\]
Hence \(a-r_e\ge3d+1\), \(g\ge r_o\).
Pad the relation lists separately to get
\[
 R\simeq
 \mathcal O'[[x_1,\ldots,x_a,y_1,\ldots,y_g]]
   /(F_1,\ldots,F_b,G_1,\ldots,G_g),\qquad a=b+3d+1 .           \tag{LS11}
\]
Here \(F_i\) are even, \(G_i\) are odd, every
linear relation coefficient is nonunit,
all entries have their prescribed parity,
\(B\) is even, and \(\mathfrak e=(y_j)\).
Assign each odd monomial to its first
occurring odd variable.
This writes \(G=My\), with every constant
coefficient of \(M\) in
\(\mathfrak m_{\mathcal O'}\).
All fields of the source's \srcid{EPres}
object have now been constructed.

\paragraph{Sources and scope.}
The actual producers under
\srcid{Development.AutomorphicArithmetic.AutomorphicVarieties}
at \srcid{alpoge-source-7cbf7a0c} are
\srcid{FrobeniusSmoothVectors.SpaceEx},
\srcid{AlgModel}, \srcid{AlgFamily},
\srcid{AlgLimit}, \srcid{PointLoc},
\srcid{DefIso}, \srcid{LocCpl},
\srcid{LocFlat}, and \srcid{LocDim}.
The obstruction and presentation calculations
expand \srcid{HasObs.pres},
\srcid{PrSetlift}, \srcid{PrCorrect},
\srcid{PrNorm}, \srcid{PrCount},
\srcid{PrMazur}, \srcid{PrDual},
\srcid{LieTMem.ep\_tan},
\srcid{EpHasobs}, \srcid{EpMazur},
and \srcid{SumLe.epres}.
This proves these specialized producers for
the arithmetic frame already constructed;
no automorphic-support implication has been
assumed in proving them.

\section{Small loci, selected curves, and the actual sharp ideal}
\label{sec:entry-nice-sharp}
We continue with the single entry algebra and its
completed charts in Section~\ref{sec:entry-space-flat}.
All dimensions and cohomology constants below
belong to the arithmetic frame already constructed.

\subsection{The two small loci}
The lower-triangular characteristic fibre has
dimension at most \(2C_0\).
Here is a generator proof, valid over rings.
In rank three its matrices are determined by
the lower functions \((y,a)\) used in the
explicit interpolation.
The linear space of possible \(y\)'s injects
into the strict \(H^1\) for \(\Psi^{-1}\)
over the CM subgroup.
The marked equality at \(g_0,g_0^*\) kills a
possible coboundary: it gives
\((\Psi(g_0)-\Psi(g_0)^{-1})h=0\).
Thus its dimension is at most \(C_0\).
The space of strict linear differences of
the \(a\)'s has dimension at most \(C_0\),
by normalization at \(c\) and the strict
\(H^1\) for \(\Psi^{-2}\delta\).

Choose a basis of the \(y\)-space.
Finitely many evaluations recover its
coefficients, so all values of \(y\) belong
to an algebra generated by at most \(C_0\)
elements.
Modulo this algebra, the quadratic term
\(y(g^*)y(h)\) in the equation for \(a\)
vanishes.
Its local primitive products vanish there
too: separated local characters recover
each primitive from one \(y\)-value.
The resulting additive quotient function
is in the strict linear \(a\)-space.
Choose representatives for its at most
\(C_0\) coefficients.
Together these elements generate every
matrix entry.
The universal fibre is generated by those
entries, so has dimension at most \(2C_0\).
In rank two only the strict linear space
occurs.
The fixed-weight topology and (LS6) give
the same bound for the completed fibre
over \(\mathfrak m_B\).

For the diagonal characters outside that
fibre, use the following finite detection
of abelian characters.
There are \(C_0+1+d\) elements \(s_i\)
in \(\ker\Psi\), and \(d\) local elements
\(s_i'\) in \(\ker\Psi\), with one fixed
exponent \(h\ne0\), such that:
\begin{enumerate}
\item a continuous character of the CM group
into a finite abelian \(\ell\)-group,
trivial on the \(s_i\), has exponent dividing \(h\);
\item if additionally
\(\chi(g^*)=\chi(g)\), its triviality on
the \(s_i'\) has the same consequence.
\end{enumerate}
To construct them, choose \(C_0+1\)
evaluations detecting
\(\operatorname{Hom}(\Gamma,k)\), by (AF4),
and append the \(d\) ordinary unit generators.
Split a character into its \(c\)-even and
\(c\)-odd parts, using that two is invertible.
The even part extends to \(\Gamma\);
the first evaluations detect it.
The odd part, after its local inertia
values vanish, factors through the finite
CM class group.
Its exponent is \(h\).
Concretely the first evaluations may be
replaced by
\((u\,cuc^{-1})^{\#k^\times}\);
they then belong to \(\ker\Psi\), without
changing their detection of an
\(\ell\)-primary character.
The \(s_i'\) are the \(d\) local unit
generators alone.
If a finite abelian \(\ell\)-group character
escaped these tests, its nonzero quotient
of order \(\ell\) would contradict the
corresponding \(k\)-linear evaluation
property.
This proves the assertion for every such
finite target, not just \(k\).

Now cut the lower-triangular locus by
the coefficient uniformizer, by
\(\chi_0(s_i)-1\), and in rank three by
\(\chi_1(s_i')-1\).
In each prime quotient containing these
equations, the diagonal characters equal
their residue characters.
Indeed their principal-unit parts have
finite exponent by the preceding tests.
In a characteristic-\(\ell\) domain there
is no nontrivial \(\ell\)-power torsion;
the remaining prime-to-\(\ell\) torsion is
uniquely fixed by its residue.
Polarization determines the final
diagonal character from the first.
The middle character in rank three
satisfies the star condition in item 2.
The ordinary product identities identify
the local characters with these same
ordered diagonal characters.
All off-diagonal monomials of weight
zero vanish on a triangular matrix.
Therefore this cut kills the entire
maximal ideal \(\mathfrak m_B\).

There are at most
\(2+C_0+(n-1)d\) cutting equations.
The height theorem and the characteristic
fibre bound prove, in the relaxed
constants of the frozen source,
\[
     \dim R_x/\mathfrak t_-
                    \le(n-1)d+3C_0+7<nd .                   \tag{SH1}
\]
On a chart with
\(\mathfrak t_+\not\subseteq x\), the
other triangular ideal is the unit ideal.

Consider next a supported local coincidence.
The actual finite-weight ideal \(I\subset B\)
has
\[
       \mathfrak m_B\subseteq\sqrt{I+\mathfrak l},
\]
where \(\mathfrak l\) is generated by the
coefficient uniformizer and ordinary
inertia augmentation on
\(\ker\varepsilon\).
Suppose \(\psi_i=\psi_j\) on
\(N_v=I_v\cap\ker\varepsilon\).
The local unit basis has \(d_v-1\)
members in \(N_v\).
In the system cutting out ordinary
augmentation omit the \(j\)-th character
on precisely those members.
Its values are supplied by the \(i\)-th
character and the coincidence equations.
Keep its value on the one remaining
generator.
Together with the coefficient uniformizer
this uses at most \(nd-d_v+2\) equations.
Their radical is maximal by the
finite-weight property.
Hence the corresponding quotient of \(B\)
has dimension at most \(nd-d_v+2\).

At an off-triangular prime the relative
transcendence degree of \(A/B\) is at most
one.
Choose nonzero entries \(u,v\) of opposite
weights.
For any other entry \(z\), a positive
power of \(z\) times a power of \(u\) or
\(v\) has weight zero.
These products lie in \(B\), and are
nonzero in the prime quotient.
Thus \(z\) is algebraic over
\(\operatorname{Frac}(B)[u]\).
Apply (LS6).
For a characteristic-\(\ell\) prime on
the coincidence locus, the relaxed
source bound is
\[
                         \dim R_x/P+d_v\le nd+5 .             \tag{SH2}
\]
Since \(d_v>5\), its dimension is
strictly less than \(nd\).
For triangular primes use (SH1).
Taking the supremum over prime quotients
proves the same strict bound for each
full coincidence ideal.

\subsection{The selected characteristic-prime curves}
Suppose \(P\subset R_x\) contains the
supported ideal and has
\(\dim R_x/P\ge nd+1\).
Apply the one-dimensional prime-selection
lemma already proved, with threshold \(nd\),
the coefficient uniformizer, and the
finitely many small ideals in (SH1)--(SH2).
It gives an actual prime \(p\supseteq P\)
satisfying
\[
 \ell\in p,\quad \dim R_x/p=1,\quad
 \mathfrak t_\pm\not\subseteq p,\quad
       \psi_i/\psi_j\ne1\text{ on every }N_v.                 \tag{SH3}
\]
These are every field of the source
\srcid{Nice} predicate.

The normalization of the complete
one-dimensional local domain \(R_x/p\)
is finite.
It is a complete local DVR with finite
residue field \(K\).
Since it has equal characteristic \(\ell\)
and \(K\) is perfect, it contains its unique
coefficient field lifting \(K\): its elements
are the roots of \(X^{|K|}-X\).
Successive expansion in a uniformizer identifies
this DVR with \(K[[T]]\).
Thus it gives the actual \srcid{W.Point}
coefficient ring, with kernel \(p\).
The local character product identities
give its ordinary flags by the separated
eigenvalue argument of (LS1).
The characteristic-\(\ell\) regular-coset
proof from these separated points supplies
the previously established \(N1,N3\);
the distinguished scalar Frobenius
supplies \(N4\).
Thus the selected curves give actual
nice points in the same determinant
space, rather than hypothetical points
of an unrelated ring.

\subsection{Connectedness after a determinantal section}
We use and prove the needed complete-local
connectedness statement.
If \(C\) is an excellent complete local
domain of dimension \(q\), and \(h\)
elements of its maximal ideal generate
\(J_0\), every minimal prime of \(J_0\)
has quotient dimension at least \(q-h\),
and the graph of these minimal primes
with edges
\[
                        \dim C/(P+Q)\ge q-h-1                \tag{SH4}
\]
is connected.
The component bound is the height theorem
and catenarity.

For the graph, normalize \(C\), obtaining
a finite complete normal local domain.
A disconnected graph would partition its
components after the same \(h\) equations
into two nonempty unions with intersection
dimension less than \(q-h-1\).
This reduction is legitimate under the
finite normalization map: dimensions are
preserved, each minimal component upstairs
has dimension at least \(q-h\), and hence
none maps entirely into that smaller
intersection.
Regroup the minimal components upstairs
according to the two unions downstairs.

Localize at a prime minimal over the
intersection of the two unions.
Its height is at least \(h+2\).
The localization is excellent and normal,
so its completion is a domain of that
dimension \(q'\ge h+2\).
The two unions give a disconnection of
the punctured spectrum after the same
\(h\) equations.
Let their ideals be \(J_1,J_2\).
Neither is maximal-primary and
\(\sqrt{J_1+J_2}=\mathfrak m\),
\(\sqrt{J_1\cap J_2}=\sqrt{J_0}\).
The Cech complex on \(h\) generators
gives \(H^{q'-1}_{J_0}=0\).
The
\href{https://stacks.math.columbia.edu/tag/0EB0}
{Hartshorne--Lichtenbaum theorem}
gives \(H^{q'}_{J_1}=H^{q'}_{J_2}=0\).
Mayer--Vietoris would give
\(H^{q'}_{\mathfrak m}=0\), contrary
to top local-cohomology nonvanishing.
This proves (SH4), including the case
of a nonequidimensional quotient.

Put \(D_0=3d+1\), and consider the even/odd
presentation (LS11), with \(G=My\).
Introduce independent variables
\(Z_{ij},T_j\) over its regular ambient
power-series ring.
By the explicit incidence parametrization
already proved, the radical of
\[
                               (ZT,\det Z)
\]
is prime of height at most \(g\).
Its quotient domain has dimension at least
\(D_0+1+b+g^2+g\).
Cut by \(Z_{ij}-M_{ij}\), by \(T_j-y_j\),
and by the \(b\) even relations.
All \(g^2+g+b\) equations belong to the
maximal ideal.
Statement (SH4) shows that
\[
             I^\sharp=(\det M)\subset R_x
            \text{ has its minimal-prime graph at threshold }D_0:
                                                               \tag{SH5}
\]
every minimal prime has quotient dimension
at least \(D_0+1\), and adjacent ones may
be joined with
\(\dim R_x/(P+Q)\ge D_0\).
Radicals and the surjection from the
ambient ring preserve this graph and
its dimensions.

In the actual exceptional chart \(g>0\).
Otherwise (LS11) would give
\(\mathfrak eR_x=0\).
Faithful flatness after localizing
\(A\) at \(x\) then gives a single
\(s\notin x\) annihilating the finitely
generated ideal \(\mathfrak e\subset A\).
The actual polynomial families after
(LS4) prove that the nonexceptional
locus is dense in the characteristic
core.
It therefore meets its basic
neighborhood \(D(s)\).
But \(s\mathfrak e=0\) makes every
prime in \(D(s)\) contain \(\mathfrak e\),
a contradiction.

The adjugate identity gives
\((\det M)y_j=0\) in \(R_x\).
Every prime avoiding
\(\mathfrak e=(y_j)\) consequently
contains \(I^\sharp\).
Every minimal component in (SH5)
is off both triangular loci, by
its dimension lower bound and (SH1).
If it avoids \(\mathfrak e\), it is
Type I.

\subsection{The actual thick representation on an exceptional component}
Let \(P\) be a minimal prime in (SH5)
containing \(\mathfrak e\).
Over \(R_x/P\), the singular matrix \(M\)
has a nonzero kernel vector \(\nu\).
Choose it over the fraction field and
clear its finitely many denominators.
The operator
\[
          D(f)=\sum_j\nu_j\frac{\partial f}{\partial y_j}
                                               \pmod P        \tag{SH6}
\]
kills the even relations: their odd
derivatives belong to \((y_j)\).
It kills \(My\) because \(M\nu=0\).
Thus it is a derivation \(R_x\to R_x/P\),
zero on the even subring and on \(B\),
and nonzero on at least one \(y_j\).

Define actual matrices by
\[
 s_D(g)=
 \begin{pmatrix}
 \rho_{00}(g)&D\rho_{01}(g)&\rho_{02}(g)\\
 0&\rho_{11}(g)&0\\
 \rho_{20}(g)&D\rho_{21}(g)&\rho_{22}(g)
 \end{pmatrix}
                    \quad\text{over }R_x/P .                 \tag{SH7}
\]
Let \(\pi:R_x\to R_x/P\) be the quotient map.
For a matrix \(a\), require that its four odd
entries vanish under \(\pi\) and that the
derivatives of its five even entries vanish.
The matrices with these properties are closed
under multiplication.
For two such matrices, the Leibniz rule gives
\[
 \operatorname{th}(ab)
       =\operatorname{th}(a)\operatorname{th}(b),
 \qquad
 \operatorname{th}(a-cI)
       =\operatorname{th}(a)-\pi(c)I
                          \quad\text{if }D(c)=0,
\]
where \(\operatorname{th}\) retains the even
entries, differentiates entries \((0,1),(2,1)\),
and sets entries \((1,0),(1,2)\) to zero.
In particular \(\operatorname{th}(I)=I\), and
\(s_D=\operatorname{th}\circ\rho\) is a group
representation.

The original polarization and its ordered
ordinary identity imply
\[
 \prod_{i=0}^{2}
   \bigl(\rho((\operatorname{loc}_v g_i)^*)
                    -\psi_{v,2-i}(g_i)I\bigr)=0.             \tag{SH8}
\]
Indeed apply \(a\mapsto J\,{}^taJ\) to the original
identity with the arguments in reversed order;
transpose reverses the product.
Every factor in the ordinary identity and
(SH8) has the two properties just specified,
since \(D\psi_{vi}=0\).
Applying the multiplicative map
\(\operatorname{th}\) therefore proves exactly
\[
 \begin{split}
 \prod_{i=0}^{2}
   \bigl(s_D(\operatorname{loc}_v g_i)
                    -\pi(\psi_{vi}(g_i))I\bigr)&=0,\\
 \prod_{i=0}^{2}
   \bigl(s_D((\operatorname{loc}_v g_i)^*)
                    -\pi(\psi_{v,2-i}(g_i))I\bigr)&=0.
 \end{split}                                                \tag{SH9}
\]
These are the two ordered local identities in
\srcid{IsThickMod}.
Polarization is used here for the original
\(\rho\); \srcid{IsThickMod} requires the
identities (SH9) for \(s_D\).
Continuity is retained because a
derivation sends
\(\mathfrak m_{R_x}^{N+1}\) into
\(\mathfrak m_{R_x/P}^N\).

This middle-character extension is
nonsplit after every injective embedding
of \(R_x/P\) into a field.
A splitting vector may be normalized
by \(v_1=1\).
Subtract its eigenvector equations
at \(g_0\) and \(g_0^*\).
The marked normalization makes the
off-diagonal entries equal.
Their outer diagonal difference is
a unit and their middle difference
is zero.
Therefore \(v_0=v_2=0\).
The remaining eigenvector equations
give
\(D\rho_{01}=D\rho_{21}=0\).
Polarization gives the other odd
derivatives zero, and the even
derivatives already vanish.
The ideal of odd entries is
\((y_j)\); the derivation rule would
make \(D(y_j)=0\) for every \(j\),
contrary to (SH6).
This proves all parts of the actual
\srcid{IsThickMod} predicate:
the same even part, the specified
zero middle row, both ordered local
identities, continuity, and
nonsplitting after scalar extension.

Off the exceptional chart use
\(I^\sharp=0\).
The ordinary presentation (LS10)
is \(b\) equations in \(b+D_0\)
variables, so (SH4) gives its graph
at the same threshold.
All minimal components are off
both triangular loci by (SH1),
and none contains \(\mathfrak e\),
which is the unit ideal on this
chart.
They are all Type I.

Thus the specialized sharp ideal
has actually been constructed.
It is zero off the exceptional
chart, lies in every Type-I prime,
has the required connected
minimal-prime graph, and every
minimal component is Type I or
has the explicit thick model
(SH7).

\paragraph{Exact source correspondence.}
The small-fibre generators are
\srcid{FinrankStructMapYRange},
\srcid{FinrankStructMapSRange},
\srcid{FinrankFrameSpaceEq},
and \srcid{FibDim}.
Abelian detection and the two small
loci are
\srcid{AbEval.ab\_gen},
\srcid{CharpolyCoeffsInSubalgebra.l51},
\srcid{TameInertiaCondition.l52},
and
\srcid{KrullDimLtOfNotLeTriangular}.
The curve producer is
\srcid{CurveDVRSpecialization.%
ExistsNicePrimeAboveOfLargeDim}.
The determinant section and thick model
expand \srcid{Detsec.detsec},
\srcid{Oddm.oddm},
\srcid{AnnOfFlat.gpos},
\srcid{ThickOfDer.thick}, and
\srcid{ChainedOfComap.sharp3}.
Except for the indicated curve module,
these paths are under
\srcid{FrobeniusSmoothVectors}, within
\srcid{Development.AutomorphicArithmetic.AutomorphicVarieties}
at \srcid{alpoge-source-7cbf7a0c}.

These producers supply the actual
flat-chart, sharp-ideal, dichotomy
and curve-selection hypotheses in
the already proved local/global
propagation argument.
The remaining insertion into its
supported locus uses the common
Hecke module and its completed
support inclusions; it is not a
consequence of connectedness alone.

\section{Automorphic support and local presentations}
\label{sec:sr-support-propagation}
The arithmetic endpoint requires one automorphic support to propagate
through a connected deformation space. We first give the specialization
and connectedness argument, then construct a single patched module
from explicit finite-level inputs, and finally prove the determinant
algebra used at the paired locus. These proofs do not replace the
construction of the finite-level automorphic modules or the Galois
presentations to which they apply.

\subsection{The characteristic-prime propagation mechanism}
\label{sec:sr-characteristic-propagation}
We isolate and prove the prime-selection and connectedness arguments in
the rank-three endpoint. These arguments do not invoke an automorphy
lifting theorem. Their application still requires the specific
deformation presentations and the Hecke-support implications stated
at the end of this section.

\subsubsection{A one-dimensional characteristic-prime specialization}
\begin{srlemma}\label{lem:sr-nice-prime}
Let $(R,\mathfrak m)$ be a Noetherian local ring, let
$\varpi\in\mathfrak m$, and let $D\ge1$.
Suppose that $I,\mathfrak a_1,\ldots,\mathfrak a_s$ are ideals with
\[
 \dim R/I\ge D+1,\qquad
                  \dim R/\mathfrak a_i<D\quad(1\le i\le s).
\]
There is a prime $\mathfrak p$ such that
\[
 I+(\varpi)\subseteq\mathfrak p,\qquad
 \dim R/\mathfrak p=1,\qquad
                 \mathfrak a_i\not\subseteq\mathfrak p
                              \quad(1\le i\le s).
\]
\end{srlemma}
\begin{proof}
The local dimension inequality for one equation gives
$\dim R/(I,\varpi)\ge D$.
One proof of that inequality is to lift a system of parameters of
$R/(I,\varpi)$ to $R/I$ and adjoin $\varpi$; the resulting maximal-primary
ideal has at most one additional generator, and Krull's height theorem
bounds the dimension. Since a Noetherian ring has finitely many minimal
primes, choose a minimal prime $\mathfrak q$ of $(I,\varpi)$ with
$\dim R/\mathfrak q\ge D$.
Then $\mathfrak a_i\not\subseteq\mathfrak q$ for every $i$.
Choose $x_i\in\mathfrak a_i\setminus\mathfrak q$ and
$g\in\mathfrak m\setminus\mathfrak q$, and put $f=g\prod_i x_i$.
Thus $f\in\mathfrak m\setminus\mathfrak q$.

Choose a prime $\mathfrak p$ maximal among the primes containing
$\mathfrak q$ and avoiding $f$; the ascending-chain condition supplies
a maximal element. Put $A=R/\mathfrak p$, with maximal ideal
$\mathfrak n$ and image $\bar f$ of $f$.
This is a local domain, $0\ne\bar f\in\mathfrak n$, and every nonzero
prime of $A$ contains $\bar f$, by maximality of $\mathfrak p$.
We claim $\dim A=1$.

For any $0\ne x\in\mathfrak n$, choose a prime $P$ minimal over $(x)$.
It is nonzero and has height at most one by the principal ideal theorem.
It contains $\bar f$. Choose $Q$ minimal over $(\bar f)$ inside $P$.
Since $\bar f\ne0$, $Q\ne0$. A strict inclusion $Q\subsetneq P$ would
give the chain $0\subsetneq Q\subsetneq P$, contrary to the height
bound for $P$. Hence $P=Q$.
It follows that every element of $\mathfrak n$ lies in the union of
the finitely many primes minimal over $(\bar f)$.
Finite prime avoidance places $\mathfrak n$ inside one of them, so
$\mathfrak n$ itself has height at most one.
The nonzero element $\bar f\in\mathfrak n$ gives
$0\subsetneq\mathfrak n$, so the dimension is exactly one.
Finally $f\notin\mathfrak p$ forces
$x_i\notin\mathfrak p$ for every $i$, proving the required avoidance.
\end{proof}
For $\varpi=\ell$, this produces a characteristic-$\ell$ prime.
The conclusion concerns a one-dimensional quotient; it does not assert
that the selected prime has height one.
The source proof is
\href{https://github.com/ChristianKudeba/alpoge-rh-verification/blob/alpoge-source-7cbf7a0c/src/Development/AutomorphicArithmetic/AutomorphicVarieties/CurveDVRSpecialization/NicePrimeEx.lean}
{\srcid{CurveDVRSpecialization.NicePrimeEx}}.

\subsubsection{Local propagation across ordinary and exceptional components}
For a Noetherian ring $R$ and a prime $P$, write
$d(P)=\dim R/P$. Fix an integer $D$, an ideal $\mathfrak s$, and five
sets of primes $\mathcal T,\mathcal E,\mathcal G,\mathcal N,\mathcal W$.
Think of these as the ordinary locus, exceptional locus, supported
locus, selected specialization primes, and thick exceptional components.
The following lemma states every hypothesis used by the local argument.

Let $\mathcal C=\operatorname{Min}(\mathfrak s)$.
Two members $P,Q\in\mathcal C$ are adjacent if
$\dim R/(P+Q)\ge D$.
Two selected exceptional primes $p,q\in\mathcal N\cap\mathcal E$ are
linked if there is $W\in\mathcal C\cap\mathcal W$ with $W\subseteq p,q$.
A chain in either relation means a finite chain, allowing length zero.

\begin{srlemma}\label{lem:sr-local-propagation}
Assume:
\begin{enumerate}
\item $\mathcal G$ is closed under specialization, and $\mathcal T$
is closed under generalization.
\item Every $p\in\mathcal N$ belongs to $\mathcal T\cup\mathcal E$.
\item If $I\in\mathcal G$ and $d(I)\ge D$, there is
$p\in\mathcal N$ containing $I$.
\item If $W\in\mathcal C\cap\mathcal W$, $W\subseteq I$, and
$d(I)\ge D$, there is $p\in\mathcal N\cap\mathcal E$ containing $I$.
\item Every minimal prime of $R$ belonging to $\mathcal T$ is in
$\mathcal C$. Every two members of $\mathcal C$ are connected by an
adjacency chain. Every member of $\mathcal C$ is in
$\mathcal T\cup\mathcal W$.
\item If $p\in\mathcal N\cap\mathcal T\cap\mathcal G$, every prime
$P\subseteq p$ belongs to $\mathcal G$.
\item Suppose $p,q\in\mathcal N\cap\mathcal E$ are linked by a chain,
and some $P\subseteq p$ belongs to $\mathcal T\cap\mathcal G$.
Then every $P'\subseteq q$ belonging to $\mathcal T$ belongs to
$\mathcal G$.
\end{enumerate}
If $I\in\mathcal T\cap\mathcal G$ and $d(I)\ge D$, then every minimal
prime of $R$ in $\mathcal T$ belongs to $\mathcal G$.
\end{srlemma}
\begin{proof}
First observe a consequence valid for every $p\in\mathcal N$:
if some $Q\subseteq p$ belongs to $\mathcal T\cap\mathcal G$,
then every $Q'\subseteq p$ in $\mathcal T$ belongs to $\mathcal G$.
If $p\in\mathcal T$, specialization gives $p\in\mathcal G$ and
hypothesis 6 applies. If $p\in\mathcal E$, hypothesis 7 applies to
the chain of length zero.

Choose a minimal prime $P_0$ of $R$ inside $I$.
Generalization gives $P_0\in\mathcal T$, hence $P_0\in\mathcal C$.
Hypothesis 3 supplies $p\in\mathcal N$ above $I$; the observation
with witness $I$ proves $P_0\in\mathcal G$.

We propagate the following information along an adjacency chain.
A component $Q\in\mathcal C$ is informed if either
$Q\in\mathcal T\cap\mathcal G$, or $Q\in\mathcal W$ and there are
selected exceptional primes $p_0,p$ joined by a link chain such that
some ordinary supported prime lies below $p_0$ and $Q\subseteq p$.
Informed ordinary components are supported, by definition.

Suppose $Q$ is informed and adjacent to $Q'$.
Choose a prime $M$ above $Q+Q'$ with $d(M)\ge D$; a minimal prime
of this quotient with maximal quotient dimension suffices.
If $Q\in\mathcal T\cap\mathcal G$, specialization makes $M$ supported.
Hypothesis 3 gives a selected $p\supseteq M$, and the first observation
shows that every ordinary prime below $p$ is supported.
If $Q$ is informed through an exceptional link chain, hypothesis 4
gives a selected exceptional $p\supseteq M$.
The old endpoint and $p$ are linked through $Q$; append this link to
the old chain and apply hypothesis 7.
Again every ordinary prime below $p$ is supported.
Thus if $Q'\in\mathcal T$, then $Q'\in\mathcal G$.

If $Q'\notin\mathcal T$, hypothesis 5 gives $Q'\in\mathcal W$.
Hypothesis 4 gives a selected exceptional $p'\supseteq M$.
When $Q$ is ordinary supported, $Q\subseteq p'$ is the initial
ordinary supported witness, so $Q'$ is informed with a chain of
length zero. When $Q$ is thick and informed, append the link from
the previous endpoint to $p'$ through $Q$.
This informs $Q'$ as well.

Starting at the supported ordinary component $P_0$, hypothesis 5
therefore informs every component in $\mathcal C$.
Every minimal prime of $R$ in $\mathcal T$ lies in $\mathcal C$
and is supported, which proves the conclusion.
\end{proof}
This is the entire argument of
\href{https://github.com/ChristianKudeba/alpoge-rh-verification/blob/alpoge-source-7cbf7a0c/src/Development/AutomorphicArithmetic/AutomorphicVarieties/CurveDVRSpecialization/LocProp.lean}
{\srcid{CurveDVRSpecialization.LocProp}}.
The separate propagation through thick exceptional primes is essential:
ordinary propagation at a single specialization does not cross those
components by itself.

\subsubsection{Global propagation from a connected characteristic-prime core}
\begin{srlemma}\label{lem:sr-global-propagation}
Let $X$ be a Noetherian topological space. Fix subsets
$T,G,C,D,E\subseteq X$.
Call an irreducible component relevant if it meets $T$.
Suppose:
\begin{enumerate}
\item $D\subseteq C$ and $C$ is connected.
\item Every closed subset meeting $C$ meets $D$.
\item $C\subseteq\overline{D\setminus E}$.
\item Every relevant irreducible component meets $C$.
\item If relevant components $Y,Y'$ meet at $x\in D$ and $Y\subseteq G$,
then $Y'\subseteq G$.
\item Every $x\in D\setminus E$ lies on a relevant component.
\item Some relevant component is contained in $G$.
\end{enumerate}
Then $T\subseteq G$.
\end{srlemma}
\begin{proof}
There are finitely many irreducible components.
Let $A$ be the union of all relevant components contained in $G$, and
let $B$ be the union of the other relevant components.
These are closed sets. Hypotheses 3 and 6 give
$C\subseteq A\cup B$.
The seed and hypothesis 4 imply $C\cap A\ne\varnothing$.
If a relevant component were not contained in $G$, hypothesis 4
would also give $C\cap B\ne\varnothing$.
Connectedness would then give $C\cap A\cap B\ne\varnothing$.
Hypothesis 2, applied to the closed set $A\cap B$, would give
$x\in D\cap A\cap B$.
Choose a component in each of the two unions containing $x$.
Hypothesis 5 would make the component from $B$ supported,
contradicting its definition.
Thus every relevant component is contained in $G$.
Every point of $T$ lies on a relevant component, proving the claim.
\end{proof}
This is
\href{https://github.com/ChristianKudeba/alpoge-rh-verification/blob/alpoge-source-7cbf7a0c/src/Development/AutomorphicArithmetic/AutomorphicVarieties/CurveDVRSpecialization/GlobProp.lean}
{\srcid{CurveDVRSpecialization.GlobProp}}.
The proof works with preconnected $C$; no nonemptiness is needed beyond
the seed and the hypotheses that the seed component meets $C$.

Here is how local charts give hypothesis 5. Let $T$ be open and $G$
closed. At $x\in D$, suppose that a flat local chart
$f:\operatorname{Spec}R\to X$ maps its closed point to $x$.
Its minimal primes map to the generic points of components through $x$,
and every such component is obtained this way.
For an affine neighborhood this follows by factoring through
$\mathcal O_{X,x}\to R$: the map is flat and local, hence faithfully
flat; lying over and going down give the two assertions.
Suppose every minimal prime of $R$ has quotient dimension at least
the threshold $D_0$ in Lemma~\ref{lem:sr-local-propagation}.
A relevant component $Y$ through $x$ has a minimal-prime lift $P$.
Its generic point is in the open set $T$, so $P\in f^{-1}T$.
If $Y\subseteq G$, then $P\in f^{-1}G$.
The local lemma now places every ordinary minimal-prime lift in
$f^{-1}G$. A lift of the generic point of any other relevant
component $Y'$ through $x$ is one of these primes.
Closedness of $G$ then gives $Y'\subseteq G$.
The same argument starts a supported component from any local ordinary
supported prime of sufficiently large quotient dimension.
These are the chart comparisons called
\srcid{ChartOfFlat} and \srcid{GlobOfCharts} in the source.

\subsubsection{An explicit connected family in characteristic $\ell$}
We prove the interpolation used for the characteristic-$\ell$ core.
Let $k$ be a field of odd characteristic, $\Gamma$ a profinite group,
$\delta:\Gamma\to\{\pm1\}$ a continuous surjection, and
$c\in\Gamma$ an involution with $\delta(c)=-1$.
Put $\Gamma_0=\ker\delta$.
Let $\phi:\Gamma\to k^\times$ be a locally constant character with
$\phi(c)=1$. On $\Gamma_0$ put $g^*=cg^{-1}c^{-1}$ and
$P(g)=\phi(g)$, so $P(g^*)=P(g)^{-1}$.
Fix $g_0\in\Gamma_0$ with $P(g_0)^2\ne1$ and local homomorphisms
$\Gamma_v\to\Gamma_0$, on each of which $P^2$ is nontrivial.

For a $k$-algebra $Q$, call a locally constant function
$b:\Gamma_0\to Q$ a strict linear datum if
\[
 \begin{split}
 b(gh)&=P(h)b(g)+P(g)^{-1}b(h),\\
 b(g^*)&=b(g),\\
 b(s)&=(P(s)-P(s)^{-1})\gamma_v
                \quad(s\in\Gamma_v)
 \end{split}
\]
for some $\gamma_v\in Q$.
Call a pair of locally constant functions $(y,a)$ a lower datum if
\begin{align*}
 y(gh)&=P(h)y(g)+y(h),\\
 a(gh)&=P(h)a(g)+y(g^*)y(h)+P(g)^{-1}a(h),\\
 a(g^*)&=a(g),\qquad y(g_0)=y(g_0^*),
\end{align*}
and, on each local group, there are $\alpha_v,\beta_v,\gamma_v\in Q$
with
\[
 \begin{split}
 y(s)&=(P(s)-1)\alpha_v,\\
 y(s^*)&=(1-P(s)^{-1})\beta_v,\\
 a(s)&=(P(s)-P(s)^{-1})\gamma_v
                  -(1-P(s)^{-1})\alpha_v\beta_v.
 \end{split}
\]
These formulas give the normalized lower-triangular representation
\[
 R_{y,a}(g)=
 \begin{pmatrix}
 P(g)&0&0\\
 y(g)&1&0\\
 a(g)&y(g^*)&P(g)^{-1}
 \end{pmatrix}.
\]
The two multiplication identities are precisely the matrix
multiplication identity. The involution gives
$R_{y,a}(g^*)=J\,{}^tR_{y,a}(g)J$ for the reversal matrix $J$,
and the equality at $g_0$ is the marked off-diagonal normalization.
The displayed local formulas imply the ordered three-character
ordinary determinant identity by multiplying the three triangular
factors. This checks the actual normalized representation conditions.

\begin{srlemma}\label{lem:sr-low-interpolation}
The lower data are stable under scalar extension and under
\[
             (y,a)\longmapsto(ty,t^2a+b),
\]
where $t\in Q$ and $b$ is a strict linear datum.
Suppose there exists a nonzero strict linear datum over $k$ and a lower
datum $(y_*,a_*)$ over $k$ with $y_*\ne0$.
Then the nonzero lower-data core is connected, and its subset
$y\ne0$ is dense, in every representing characteristic-$\ell$ core
having the field-point and polynomial-point universal properties
for these formulas.
\end{srlemma}
\begin{proof}
The operation preserves the multiplication and involution identities
by direct substitution. In the local formulas replace
\[
  (\alpha_v,\beta_v,\gamma_v)
       \quad\hbox{by}\quad
  (t\alpha_v,t\beta_v,t^2\gamma_v+\eta_v),
\]
where $b(s)=(P(s)-P(s)^{-1})\eta_v$.

Fix a nonzero strict linear datum $b_0$.
The pair $(0,b_0)$ is a lower datum.
For any field-valued lower datum $(y,a)$, the polynomial-valued pair
\[
            (Ty,\ T^2a+(1-T)b_0)
\]
is a lower datum over $Q[T]$. Its evaluations at one and zero are
$(y,a)$ and $(0,b_0)$.
The generic prime of this polynomial family specializes to both
endpoint primes. It stays in the nonzero core: a zero polynomial
datum would have zero evaluation at zero, contradicting $b_0\ne0$.
Thus every point is joined to the fixed point $(0,b_0)$ by a connected
three-point subset consisting of the two endpoint primes and the
generic prime of the family. Their union is connected.

Now take an exceptional point $(0,a)$.
The defining identities show that $a$ is a strict linear datum.
The family
\[
                    (Ty_*,\ T^2a_*+a)
\]
specializes to $(0,a)$ at zero, and its generic point has nonzero
first coordinate. Hence every exceptional point is in the closure
of the nonexceptional locus. Points already outside the exceptional
locus require no argument.
\end{proof}

The two nonzero data have a direct cohomological origin.
For $m=1,2$, let $\xi_m=\phi^{-m}\delta$, and suppose that the strict
degree-one cohomology for $\xi_m$, with local coboundaries at the
chosen groups, is nonzero.
These are precisely the two nontriviality fields in the source frame.
The local nontriviality of $P^2$ ensures that the corresponding
mapping-cone cohomology injects into $H^1(\Gamma,k(\xi_m))$:
if a global cocycle is a coboundary, the difference of every local
primitive and the global primitive is locally invariant, hence zero.
Restriction to $\Gamma_0$ is injective by averaging over the quotient
of order two.

For $m=2$, choose a nonzero class and a cocycle $x$ with $x(c)=0$,
which is possible because $\xi_2(c)=-1$.
Then
\[
                         b_0(g)=P(g)x(g)
\]
is a nonzero strict linear datum.
The cocycle equation proves its multiplication formula.
Normalization at $c$ gives
$x(g^*)=P(g)^2x(g)$, proving its involution formula.
A local coboundary $x(s)=(P(s)^{-2}-1)h_v$ gives the local formula
with $\gamma_v=-h_v$.

For $m=1$, the same normalization gives a nonzero locally constant
function $u(g)=P(g)x(g)$ with
\[
 u(gh)=P(h)u(g)+u(h),\qquad
 u(g^*)=P(g)^{-1}u(g),\qquad
 u(s)=(P(s)-1)\alpha_v.
\]
Put $p_0=P(g_0)$ and choose
\[
                 \lambda=-\frac{u(g_0)}{p_0+1}.
\]
The denominator is nonzero because $p_0^2\ne1$.
Define
\[
 \begin{split}
 y_*(g)&=u(g)+\lambda(P(g)-1),\\
 a_*(g)&=\lambda^2
   -\frac{\lambda^2}{2}(P(g)+P(g)^{-1})
   -\lambda(u(g)+P(g)^{-1}u(g))
   +\frac{P(g)^{-1}u(g)^2}{2}.
 \end{split}
\]
Substitution of the three identities for $u$ proves the lower-data
identities. At $g_0$ the difference
$y_*(g_0)-y_*(g_0^*)$ is
\[
       \frac{p_0-1}{p_0}
                  \bigl(u(g_0)+\lambda(p_0+1)\bigr)=0.
\]
For the local formulas use
\[
 \alpha_v'=\lambda+\alpha_v,\quad
 \beta_v'=\alpha_v-\lambda,\quad
 \gamma_v'=\frac{\alpha_v^2}{2}
                         -\lambda\alpha_v-\frac{\lambda^2}{2}.
\]
Finally $y_*=0$ would give $u=-\lambda(P-1)$.
Comparison with the involution formula at $g_0$ forces
$2\lambda(p_0-1)=0$, hence $\lambda=0$ and $u=0$, a contradiction.
This proves the required nonzero lower datum.

The source formulas and their universal-property application are in
\href{https://github.com/ChristianKudeba/alpoge-rh-verification/blob/alpoge-source-7cbf7a0c/src/Development/AutomorphicArithmetic/AutomorphicVarieties/FrobeniusSmoothVectors/StrSMul.lean}
{\srcid{StrSMul}},
\href{https://github.com/ChristianKudeba/alpoge-rh-verification/blob/alpoge-source-7cbf7a0c/src/Development/AutomorphicArithmetic/AutomorphicVarieties/FrobeniusSmoothVectors/ExistsNonzeroLowCocycle.lean}
{\srcid{ExistsNonzeroLowCocycle}}, and
\href{https://github.com/ChristianKudeba/alpoge-rh-verification/blob/alpoge-source-7cbf7a0c/src/Development/AutomorphicArithmetic/AutomorphicVarieties/FrobeniusSmoothVectors/ExistsNonzeroLowCocycle.lean}
{\srcid{preconnected_rank_three}}.
These arguments produce connectedness and density once the source
field-point and polynomial-point identifications have been established;
those identifications must not be inferred merely from the existence
of a formal local universal ring.

\subsubsection{The remaining specialized interface}
For the source's rank-three universal framed space $S$, the generic
and exceptional strata are
\[
 T=\{P:\mathfrak t_+\not\subseteq P,\
          \mathfrak t_-\not\subseteq P,\
          \mathfrak e\not\subseteq P\},\qquad
 E=\{P:\mathfrak t_+\not\subseteq P,\
          \mathfrak t_-\not\subseteq P,\
          \mathfrak e\subseteq P\},
\]
where $\mathfrak t_\pm$ are generated by the two triangular directions,
and $\mathfrak e$ by entries of odd weight for weights $(1,0,-1)$.
The supported locus is the zero locus of the Hecke ideal pulled back
to the framed entry algebra.
A selected prime has quotient dimension one, contains $\ell$, lies
off both triangular loci, and has pairwise separated local characters.

The local propagation threshold is $3d+1$ in rank three and $2d+1$
in rank two. The preceding lemmas prove the abstract propagation
once the following specialized facts are supplied:
\begin{enumerate}
\item The field-point and polynomial-point descriptions of the
characteristic-$\ell$ core, the flat local charts, and the dimension
bounds on the triangular and coincident-character loci.
\item The sharp local ideal: outside $E$ it is zero. At a point of $E$
the source constructs an even/odd presentation
\[
 R=\mathcal O'[[x_1,\ldots,x_a,y_1,\ldots,y_g]]
                         /(F_1,\ldots,F_b,G_1,\ldots,G_g),
 \qquad a=b+3d+1,\qquad G=M y .
\]
The even equations are $F_i$, the odd equations are $G_i$, and
$\mathfrak e$ is generated by the images of the $y_j$.
The sharp ideal is generated by the image of $\det M$.
The adjugate identity $(\operatorname{adj}M)G=(\det M)y$ shows directly
that every prime outside $E$ contains this ideal.
The required connectedness of its minimal-prime graph and the
ordinary/thick dichotomy still require the determinant-section and
deformation-obstruction arguments of
\srcid{Detsec}, \srcid{SumLe}, \srcid{ChainedOfComap}, and
\srcid{ThickOfDer}.
\item The Hecke-support implications P1 and P23, corresponding exactly
to hypotheses 6 and 7 of Lemma~\ref{lem:sr-local-propagation}, and a
seed of full weight-space dimension. The two-dimensional propagation
must also supply the finite exceptional support ideal used in rank
three. These are the inputs assembled in \srcid{PrLinkedEnd} and
\srcid{AbsMain}.
\end{enumerate}
Thus the prime avoidance, interpolation and propagation arguments
above are proved statements. They do not assert the full arithmetic
endpoint before these specialized facts have been proved.
In particular a theorem assuming that every residual constituent
is individually polarized cannot supply this interface merely from
the paired outer constituents $\bar\Psi$ and $\bar\Psi^{-1}$.

\subsection{One patched module and its completed supports}
\label{sec:sr-common-patched-module}
This section proves the common-module and completed-support argument
used in the source's \srcid{Raw} construction. In particular, the
modules at different points of a prime chain arise from one module.
Choosing a new module independently at each point would not prove
the support-transport assertion.

\subsubsection{The exact finite-level hypotheses}
Let $\Lambda$ be a complete Noetherian local ring with finite residue
field, and let $(R,\mathfrak m_R)$ be a Noetherian local ring.
Fix a homomorphism
\[
       s:\Lambda[[Y_1,\ldots,Y_r]]\longrightarrow R.
\]
Write $y_i=s(Y_i)$, $Y=(y_1,\ldots,y_r)$ and
$\mathfrak m=\mathfrak m_\Lambda R$.
Let $f\in R$, let $I\subseteq R$ be an ideal, and let $X$ be an
$R$-module finite over $\Lambda$, with compatible scalar actions and
$YX=0$. Fix a nonprincipal ultrafilter $\mathcal U$ on $\mathbb N$.
Let $M_N$ be $R$-modules satisfying:
\begin{enumerate}
\item For every $k,j\ge1$, for $\mathcal U$-almost all $N$ the module
\[
                  M_N/(Y^k+\mathfrak m_\Lambda^j)M_N
\]
is finite free over
$\Lambda[[Y_1,\ldots,Y_r]]/((Y_1,\ldots,Y_r)^k,\mathfrak m_\Lambda^j)$.
No common bound on the ranks is assumed.
\item For every $j\ge1$, for $\mathcal U$-almost all $N$ there are
$R$-linear maps
\[
 X/\mathfrak m_\Lambda^jX
       \ \mathop{\rightleftarrows}^{\alpha_{N,j}}_{\beta_{N,j}}\
             M_N/(Y+\mathfrak m_\Lambda^j)M_N
\]
such that both composites are multiplication by $f$.
\item There are an integer $e\ge1$ and a set of generators
$\mathscr R$ of $I$ such that, for every $a\in\mathscr R$ and $c\ge1$,
for $\mathcal U$-almost all $N$ there is $b\in\operatorname{Ann}_R M_N$
with
\[
                            (fa)^e-b\in\mathfrak m_R^c.
\]
The exponent $e$ precedes the approximation depth, the level, and
the relation generator.
\end{enumerate}
These hypotheses separate finite-level automorphic control from
the commutative algebra proved below. In the application,
$\Lambda$ is the fixed ordinary weight algebra with its extra scalar
variable, the $Y_i$ are the Taylor--Wiles variables, $f$ is the
auxiliary eigenvalue discriminant, and $I$ contains the kernel,
polarization, ordinary, inertia and two auxiliary relations.

\begin{srproposition}\label{prop:sr-one-patched-module}
Put $S=R[1/f]$ and $C=\widehat S^{\,Y}$, the $Y$-adic completion.
Under the three finite-level hypotheses there is one finite
$C$-module $\mathcal M$ with
\begin{equation}\label{eq:sr-patched-fiber}
 \mathcal M/Y\mathcal M\simeq C\otimes_R X,\qquad
 I C\subseteq\sqrt{\operatorname{Ann}_C\mathcal M}.
\end{equation}
Moreover
\begin{equation}\label{eq:sr-patched-gr}
 \operatorname{gr}_Y\mathcal M
           \simeq(\mathcal M/Y\mathcal M)[T_1,\ldots,T_r].
\end{equation}
Thus the variables are a quasi-regular sequence on $\mathcal M$.
\end{srproposition}
\begin{proof}
For $k\ge0$, define the $R$-module
\[
 U_k=\prod_N M_N\big/
 \left\{(z_N):
    \text{for every }j,\ 
     z_N\in(Y^k+\mathfrak m_\Lambda^j)M_N
                  \text{ for }\mathcal U\text{-almost all }N\right\}.
\]
In particular $U_0=0$.
The identity on representatives gives surjective maps
$U_{k+1}\to U_k$. Their kernels are $Y^kU_{k+1}$.

Here is the needed diagonal argument for the last assertion.
Choose generators $a_1,\ldots,a_t$ of $Y^k$.
If a representative $(z_N)$ maps to zero in $U_k$, then for each
$j$ it admits, for almost all $N$, a decomposition
\[
       z_N=\sum_{\nu=1}^t a_\nu w_{\nu,N}+v_N,
                          \qquad v_N\in\mathfrak m_\Lambda^jM_N.
\]
Choose a function $j(N)\to\infty$ along $\mathcal U$ on which these
decompositions hold. Such a choice is obtained by imposing the first
$j$ conditions and $N\ge j$, and taking the largest admissible
$j\le N$. Use a decomposition at $j(N)$ for each $N$.
The error sequence represents zero in $U_{k+1}$, which proves that
the kernel is generated by the $a_\nu$.
The reverse inclusion is immediate.
The same argument shows
\[
 U_k/Y^aU_k\simeq U_a\quad(1\le a\le k),\qquad Y^kU_k=0.
\]

We next identify $U_1$ after inverting $f$.
Choose a function $j(N)\to\infty$ along $\mathcal U$ and control maps
from hypothesis 2 at precision $j(N)$.
Lifting the images of elements of $X$ gives an $R$-linear map
$\alpha:X\to U_1$; all failures of linearity of arbitrary lifts vanish
in the defining quotient.
For a representative $(z_N)$ and fixed $j$, reduce
$\beta_{N,j(N)}(z_N)$ modulo $\mathfrak m_\Lambda^jX$ when
$j(N)\ge j$. The finite set $X/\mathfrak m_\Lambda^jX$ has a unique
$\mathcal U$-limit. These limits are compatible in $j$.
Since the finite $\Lambda$-module $X$ is complete, they define an
$R$-linear map $\beta:U_1\to X$.
Changing a representative by a null sequence does not change the limit:
$\beta_{N,j(N)}$ kills the $Y$-part and preserves the
$\mathfrak m_\Lambda^j$-part.
At each fixed precision the control identities give
\[
                      \alpha\beta=f,\qquad\beta\alpha=f.
\]
Consequently
\[
                         S\otimes_R U_1\simeq S\otimes_R X.
\]

Put $P_k=S\otimes_R U_k$.
Localization is exact, so $P_{k+1}\to P_k$ is surjective with kernel
$Y^kP_{k+1}$, and $Y^kP_k=0$.
Also $P_1$ is finite over $S$.
The inverse limit
$\mathcal M=\varprojlim_kP_k$ is therefore finite over $C$,
and $\mathcal M/Y^k\mathcal M\simeq P_k$ for every $k$.
We give the finite-generation detail to make this use of completion
explicit. Lift a finite generating set of $P_1$ compatibly through the
surjective tower. If $K_k$ is the kernel of the resulting surjection
$S^t\to P_k$, then
\[
             Y^kS^t\subseteq K_k,\qquad
             K_k=K_{k+1}+Y^kS^t.
\]
Successively correcting coefficient vectors makes their differences
belong to $Y^kS^t$; they converge in $C^t$.
This proves $C^t\to\varprojlim P_k$ is onto.
The same correction applied to relation vectors identifies the
kernel of the $k$-th projection with $Y^k\mathcal M$.
It also proves separatedness.

For completeness, hypothesis 1 gives the associated graded
calculation. A degree-$a$ relation
\[
               \sum_{|\nu|=a}y^\nu c_\nu
                           \in(Y^{a+1}+\mathfrak m_\Lambda^j)M_N
\]
in a module free over the truncated power-series ring forces
$c_\nu\in(Y+\mathfrak m_\Lambda^j)M_N$ for every $\nu$:
in each free coordinate the degree-$a$ monomials are independent over
$\Lambda/\mathfrak m_\Lambda^j$.
Only finitely many monomials occur.
Pass to $\mathcal U$-classes and use the diagonal kernel argument.
For $k>a$ this proves the same coefficient implication in $U_k$.
It survives localization and passage to the separated inverse limit.
The usual map from the polynomial module on $\mathcal M/Y\mathcal M$
to the associated graded is always surjective, and this coefficient
implication proves injectivity. This gives
\eqref{eq:sr-patched-gr}.
The identification for $k=1$ gives the first assertion of
\eqref{eq:sr-patched-fiber}.

Finally fix $a\in\mathscr R$ and $k,j\ge1$.
The finite $R$-module $X/\mathfrak m_\Lambda^jX$ has finite length,
so some power $\mathfrak m_R^c$ kills it.
The control maps at precision $j$ therefore give, for almost all $N$,
\[
                   f\mathfrak m_R^cM_N
                            \subseteq(Y+\mathfrak m_\Lambda^j)M_N.
\]
Taking the $k$-fold product of the acting ideals gives
\[
              f^k\mathfrak m_R^{ck}M_N
                            \subseteq(Y^k+\mathfrak m_\Lambda^j)M_N.
\]
Use hypothesis 3 at precision $ck$ and subtract an actual annihilator.
It follows that
\[
              f^{k+e}a^eM_N
                            \subseteq(Y^k+\mathfrak m_\Lambda^j)M_N
\]
for almost all $N$. The power $f^{k+e}$ is independent of $j$.
Hence $f^{k+e}a^e$ kills $U_k$.
After inverting $f$, $a^e$ kills $P_k$ for every $k$.
Separatedness proves $a^e\mathcal M=0$.
The generators of $I$ therefore map to the radical of the annihilator,
which completes the proof.
\end{proof}

This reconstructs the common module produced in
\srcid{AnnihilatorBaseChange.PmodEx} from the actual finite-level
freeness, paired control maps, and uniform nilpotence hypotheses.
It does not posit its existence separately at each specialization.

\subsubsection{Dimension at a completed point}
\begin{srlemma}\label{lem:sr-patched-support-dimension}
Let $L$ be a complete regular local ring of dimension $d$, let $B$ be
a complete Noetherian local $L$-algebra with local structure map,
and let $M$ be a finite nonzero $B$-module.
Suppose $y_1,\ldots,y_r\in\mathfrak m_B$ satisfy
\[
 \operatorname{gr}_{(y)}M
                  \simeq(M/(y)M)[T_1,\ldots,T_r],
\]
and $M/(y)M$ is finite free over $L$.
Then every irreducible component of $\operatorname{Supp}_BM$ has
dimension exactly $d+r$.
\end{srlemma}
\begin{proof}
Evaluation gives a local map $A=L[[Y_1,\ldots,Y_r]]\to B$ with
$Y_i\mapsto y_i$.
Lift an $L$-basis of $M/(y)M$ to $M$.
Successive approximation in powers of $(y)$ gives a surjection
$A^t\to M$: at each stage choose the next homogeneous coefficient,
and use completeness of the finite $B$-module to take the limit.
The associated graded hypothesis proves injectivity:
a nonzero relation has a least total degree in the $Y_i$, whose
coefficient vector would give a nonzero relation in the corresponding
free homogeneous piece.
Thus $M$ is finite free over $A$.

The faithful action of $B/\operatorname{Ann}_BM$ embeds that ring in
$\operatorname{End}_A(M)$, a finite free $A$-module.
Consequently $B/\operatorname{Ann}_BM$ is a finite torsion-free
$A$-algebra.
Every minimal prime contracts to zero in the domain $A$, since a
nonzero element of $A$ acts injectively on this torsion-free algebra
and thus lies in no minimal associated prime.
For each such prime $P$, the domain
$(B/\operatorname{Ann}_BM)/P$ is finite integral over $A$.
Its dimension is $\dim A=d+r$, proving the claim.
\end{proof}
This is the support-dimension content of
\srcid{AnnihilatorBaseChange.PmDimQuot}
and \srcid{AlgebraTools.CaCmPatch}.

In the application, the regular local ring $L$ is the completion of
the ordinary weight algebra at the contracted point.
The relevant special fiber is finite free over $L$.
Here is why finite freeness is preserved in that step.
The action of $R$ on the finite free $\Lambda$-module $X$ factors
through the finite $\Lambda$-algebra
$R/\operatorname{Ann}_RX\subseteq\operatorname{End}_\Lambda X$.
After completing $\Lambda$ at the contracted prime, this finite algebra
is a product of its complete local factors.
Completion at the selected $R$-prime picks one factor.
The corresponding piece of $X\otimes_\Lambda L$ is an idempotent
direct summand of a finite free $L$-module, and is therefore finite
free because $L$ is local. Zero summands have empty support.

\subsubsection{Fiber support and transport between completions}
Let $q_i$ be any finite family of primes of $R$ satisfying
\[
                        I+Y\subseteq q_i,\qquad f\notin q_i.
\]
For such a prime, $q_iC$ is prime and contracts to $q_i$.
Indeed $C/YC\simeq R[1/f]/Y$, so primes containing $Y$ are determined
by the ordinary localization correspondence.
Completing $C_{q_iC}$ gives the same ring as completing $R_{q_i}$:
their quotients modulo every power of the selected maximal ideal
agree, because $Y\subseteq q_i$ and $f$ is a unit there.
Write this common completion as $B_i^0$, and put
\[
 \mathcal M_i=B_i^0\otimes_C\mathcal M,\quad
 B_i=B_i^0/IB_i^0,\quad
 S_i=(\operatorname{Ann}_{B_i^0}\mathcal M_i+IB_i^0)/IB_i^0.
\]
The radical assertion in \eqref{eq:sr-patched-fiber} ensures that
the minimal support components and their dimensions are unchanged
by this quotient.

\begin{srlemma}\label{lem:sr-completed-support}
The ideals $S_i$ satisfy:
\begin{enumerate}
\item If the completed special fiber is finite free over a regular
local weight ring of dimension $d$, then
$\dim B_i/P=d+r$ for every $P\in\operatorname{Min}(S_i)$.
\item For every prime $Q\subseteq B_i$ containing $Y$,
\[
         S_i\subseteq Q
          \quad\Longleftrightarrow\quad
          \operatorname{Ann}_RX\subseteq Q\cap R.
\]
\item Suppose $Z\subseteq Q$ are primes of $B_i$, $S_i\subseteq Z$,
$Y\subseteq Q$, and $Q\cap R\subseteq q_j$.
There is a prime $Z'\subseteq B_j$ with
\[
                       S_j\subseteq Z',
                  \qquad Z'\cap(R/I)=Z\cap(R/I).
\]
\end{enumerate}
\end{srlemma}
\begin{proof}
Flat base change to a completion preserves the associated graded
description and the special fiber. The first assertion therefore
follows from Lemma~\ref{lem:sr-patched-support-dimension}.

For the second, localize at the inverse image of $Q$ in $B_i^0$.
Since the $y_i$ are in that maximal ideal and $\mathcal M_i$ is finite,
Nakayama gives
\[
 (\mathcal M_i)_Q\ne0
        \quad\Longleftrightarrow\quad
      (\mathcal M_i/Y\mathcal M_i)_Q\ne0.
\]
Use \eqref{eq:sr-patched-fiber}.
For any ring map $A\to B$ and finite $A$-module $N$,
\[
     \operatorname{Supp}_B(B\otimes_AN)
              =(\operatorname{Spec}B\to\operatorname{Spec}A)^{-1}
                                      \operatorname{Supp}_AN.
\]
To verify this identity, tensor at a prime with its residue field.
The resulting tensor product is the extension of
$N\otimes_A\kappa(p)$ through the field inclusion
$\kappa(p)\hookrightarrow\kappa(Q)$, so it is nonzero exactly when
the former vector space is nonzero.
For finite modules, support is the zero locus of the annihilator.
This proves assertion 2.

For assertion 3, first work before quotienting by $I$.
Let $p=Z\cap C$.
The prime $p$ lies in $\operatorname{Supp}_C\mathcal M$, by the
same support identity.
Since $Q$ contains $Y$, its contraction to $C$ is determined by
its contraction to $R$ modulo $Y$; hence $Q\cap R\subseteq q_j$
implies $Q\cap C\subseteq q_jC$.
Thus $p\subseteq q_jC$.
The completion of $C_{q_jC}$ is faithfully flat, so there is a prime
$Z'$ above $p$. The support identity puts it in the support of the
base change of the same module $\mathcal M$.
Contraction to $R$ is unchanged.
The ideal $I$ lies in every support prime by the radical annihilation
assertion, so the primes descend to the displayed quotients.
This proves assertion 3.
\end{proof}
The source separates these three facts as
\srcid{PmDimQuot}, \srcid{PmFib}, and \srcid{PmMove}.
The last proof explicitly uses the common module: equality of numerical
dimensions of independently chosen modules would not give transport
of a support prime.

\subsubsection{The completed-support chain}
Here is the final algebraic form of the source's \srcid{ChainAbs}.
Let $q_0,\ldots,q_h$ be selected points as above.
Assume that the completed support components have dimension at least
$d_0$. Let $\mathcal T$ be a property of primes of $R/I$ stable under
generalization, and let $\mathcal E_i$ be a property of primes of $B_i$.
Suppose the following local inclusion statement holds:
\[
 \left.
 \begin{array}{c}
 Z,Q\in\operatorname{Spec}B_i,\quad
 \mathcal T(Z\cap(R/I)),\quad\dim B_i/Z\ge d_0,\\
 \mathcal T(Q\cap(R/I))\ \text{or}\ \mathcal E_i(Q)
 \end{array}\right\}
                           \ \Longrightarrow\ Z\subseteq Q.       \tag{*}
\]
Suppose bridge primes $E_i$ of $R/I$ contain $Y$, lie below both
$q_i/I$ and $q_{i+1}/I$, and every prime of $B_i$ above $E_i$ has
property $\mathcal E_i$.
Suppose there is a prime $p_0\subseteq q_0/I$ containing $Y$, with
$\mathcal T(p_0)$ and
$\operatorname{Ann}_RX\subseteq p_0\cap R$.

\begin{srproposition}\label{prop:sr-completed-support-chain}
If $p\subseteq q_h/I$ contains $Y$ and satisfies either
$\mathcal T(p)$ or
$\mathcal E_h(Q)$ for every prime $Q$ of $B_h$ above $p$, then
\[
                            \operatorname{Ann}_RX\subseteq p\cap R.
\]
\end{srproposition}
\begin{proof}
Faithful flatness lifts $p_0$ to a prime of $B_0$.
The fiber-support equivalence puts it in $V(S_0)$.
Choose a minimal support prime $Z_0$ below it.
Then $\mathcal T(Z_0\cap(R/I))$ holds by generalization, and its
dimension is at least $d_0$.

Inductively suppose there is a supported prime $Z_i$ of $B_i$
whose contraction has property $\mathcal T$.
Replace it by a minimal support prime underneath.
Lift the bridge $E_i$ to a prime $Q_i$ of $B_i$.
The bridge hypothesis gives $\mathcal E_i(Q_i)$, so (*) gives
$Z_i\subseteq Q_i$.
The support-transport lemma then moves $Z_i$ to a supported prime
of $B_{i+1}$ with the same contraction to $R/I$.
This preserves $\mathcal T$ and proves the induction.

At the last point, choose a minimal support prime beneath the
supported prime obtained by induction, and lift $p$ to $Q\in B_h$.
The final hypothesis and (*) put the minimal support prime below $Q$.
Hence $S_h\subseteq Q$.
The fiber-support equivalence now gives
$\operatorname{Ann}_RX\subseteq p\cap R$, as required.
\end{proof}

\paragraph{What this closes in the actual source route.}
The common-module construction, its uniform nilpotence passage,
its completed support dimensions, the fiber identity, support
transport, and the chain argument have now been proved from the
displayed finite-level hypotheses and (*).
The source modules are
\href{https://github.com/ChristianKudeba/alpoge-rh-verification/blob/alpoge-source-7cbf7a0c/src/Development/AutomorphicArithmetic/AutomorphicVarieties/PatchingModules/AnnihilatorBaseChange/PmodEx.lean}
{\srcid{PmodEx}},
\href{https://github.com/ChristianKudeba/alpoge-rh-verification/blob/alpoge-source-7cbf7a0c/src/Development/AutomorphicArithmetic/AutomorphicVarieties/PatchingModules/AutomorphicDetection/AnnBc.lean}
{\srcid{AnnBc}}, and
\href{https://github.com/ChristianKudeba/alpoge-rh-verification/blob/alpoge-source-7cbf7a0c/src/Development/AutomorphicArithmetic/AutomorphicVarieties/PatchingModules/AnnihilatorBaseChange/ChainAbs.lean}
{\srcid{ChainAbs}}.
The single universal ring and relation ideal are defined in
\srcid{SpecializationMapCompUniversalSection}; \srcid{LiftPt.core}
applies the preceding facts to them.
The arithmetic local inclusion (*) is the separate producer
\srcid{AutomorphicDetection.GalOut}. It uses \srcid{S6} to realize the
completed determinant over a finite quadratic coefficient algebra,
then the first- and second-Selmer bounds and the Laurent presentation
\srcid{SelmerScaling.PresOut0}.
The finite-level freeness, paired control maps, uniform relation
nilpotence, and that Galois inclusion have to be verified in this
same universal ring before the full \srcid{Raw} endpoint is asserted.

\subsection{The determinant algebra at paired components}
\label{sec:sr-paired-determinant-algebra}
We prove the last commutative-algebra step of the local inclusion
used in the completed-support chain. It treats the two paired blocks
through a finite coefficient algebra with an odd summand.

\subsubsection{The generic incidence ideal}
Let $g\ge1$, let $R$ be a Noetherian domain, and let
\[
 P_0=R[Z_{ij},T_j:1\le i,j\le g],\qquad
 J=((ZT)_1,\ldots,(ZT)_g,\det Z).
\]
\begin{srlemma}\label{lem:sr-generic-incidence}
The radical of $J$ is prime and contains no nonzero element of $R$.
It does not contain $T_1$. The same assertions hold after extending
the homogeneous ideal to $R[[Z_{ij},T_j]]$.
If the latter power-series ring is Noetherian, the height of this
prime is at most $g$.
\end{srlemma}
\begin{proof}
Put $E=\operatorname{diag}(0,1,\ldots,1)$.
Use two generic $g$-by-$g$ matrices $N,U$ and a scalar $t$ and define
a polynomial map on coordinate rings by
\[
                 Z\longmapsto NE\operatorname{adj}U,\qquad
                 T\longmapsto tUe_1.
\]
The image satisfies $ZT=0$ and $\det Z=0$.
Its coordinate ring is a subring of the polynomial domain
$R[N_{ij},U_{ij},t]$, so the kernel $K$ is prime.
The images of a nonzero constant and of $T_1=tU_{11}$ are nonzero.

Every field-valued pair $(M,v)$ with $Mv=0$ and $\det M=0$
factors through this parametrization.
If $v\ne0$, take the first column of an invertible $U$ to be $v$
and take $t=1$.
If $v=0$, take its first column to be any nonzero vector in $\ker M$
and take $t=0$.
In either case $MU$ has first column zero.
Thus
\[
         N=(\det U)^{-1}MU
       \quad\Longrightarrow\quad
         NE\operatorname{adj}U=M,\qquad tUe_1=v.
\]
Apply this factorization to the fraction field of every prime quotient
of $P_0$ containing $J$.
Every element of $K$ belongs to all such primes, so $K=\sqrt J$.

The ideal $J$ is homogeneous for total degree, and therefore so is $K$.
For a homogeneous finitely generated ideal $K\subseteq R[\mathbf Z]$,
a formal series belongs to $KR[[\mathbf Z]]$ exactly when each of its
homogeneous polynomial pieces belongs to $K$.
Indeed choose finite homogeneous generators and sum homogeneous
coefficient decompositions degree by degree.
If two series are outside this extended ideal, choose the least
degrees $a,b$ of homogeneous pieces outside $K$.
The degree-$a+b$ piece of their product is, modulo $K$, the product
of those two pieces, and is nonzero because $K$ is prime.
Thus the extended ideal is prime.
Its contraction is $K$, and it is the radical of the extension of $J$.

Write $\widehat K$ for that prime.
Choose a prime $Q$ minimal over the $g$ row relations $ZT$
inside $\widehat K$.
Since $T_1\notin\widehat K$, also $T_1\notin Q$.
The adjugate identity gives $(\det Z)T_1\in Q$, hence $\det Z\in Q$.
Therefore $\widehat K=Q$.
Krull's height theorem gives $\operatorname{ht}\widehat K\le g$.
\end{proof}
This is the explicit parametrization and homogeneous-completion
argument of \srcid{CommutativeAlgebraHeight.GDetrad},
\srcid{GHomprime} and \srcid{GDetps}.

Now let $\mathcal O$ be a complete DVR and work in
\[
 P=\mathcal O[[Z_{ij},T_j,W_a]].
\]
The $g$ row relations $ZT$ form a regular sequence.
One direct check gives diagonal $Z_{ii}$ weight one,
off-diagonal $Z_{ij}$ weight two, and every $T_j$ weight one.
The initial forms of the rows are the monomials $Z_{ii}T_i$.
They involve disjoint sets of variables and form a regular sequence.
Taking the least weighted degree of a relation then proves regularity
of the original rows, one row at a time.
Since $P$ is regular local and catenary, Lemma~\ref{lem:sr-generic-incidence}
therefore gives
\begin{equation}\label{eq:sr-incidence-dimension}
               \dim P/\sqrt{(ZT,\det Z)}=\dim P-g.
\end{equation}

The same conclusion applies to elements $z_{ij},t_j$ of a power-series
ring over $\mathcal O$ whose classes are linearly independent modulo
$(\varpi,\mathfrak m_P^2)$.
Complete the independent classes to a relative coordinate basis.
The homomorphism from the power-series ring in those coordinates is
an isomorphism: successive approximation gives surjectivity, and its
invertible linear term gives an inverse by solving each homogeneous
degree. Constants in $\varpi\mathcal O$ cause no problem because the
map is continuous for the maximal-ideal topology.
Thus the displayed incidence ideal is carried to the generic one.

\subsubsection{Stability detected by odd products}
Let $(B,\mathfrak m_B)$ be local and let $A$ be a finite commutative
$B$-algebra with an injective structure map and a decomposition
\[
                          A=B\oplus N,\qquad NN\subseteq\mathfrak m_B.
\]
Define the ideal
\[
  I_{\rm odd}=(b\in B:\ uv=b\text{ in }A
                              \text{ for some }u,v\in N).
\]
Consider a representation whose $B$-linear span is exactly
\[
 \begin{pmatrix}
 M_a(B)&M_{a,n-a}(N)\\
 M_{n-a,a}(N)&M_{n-a}(B)
 \end{pmatrix},\qquad0<a<n .
\]
\begin{srlemma}\label{lem:sr-odd-product-detector}
At a prime $q$ of $B$, this representation has full matrix span over
an algebraic closure of $\kappa(q)$ if and only if
$I_{\rm odd}\not\subseteq q$.
\end{srlemma}
\begin{proof}
Pass to a prime of the integral algebra $A$ over $q$ and then to its
residue field.
If $I_{\rm odd}\subseteq q$, every $u\in N$ has square zero in that
field, hence has zero image. All off-block entries vanish, so the span
is block diagonal and is not the full matrix algebra.

Otherwise choose $u,v\in N$ with $uv=b\notin q$.
The image of $u$ in any residue field above $q$ is nonzero.
The given span contains every same-block matrix unit and every
off-block matrix unit multiplied by $u$.
After extending scalars, division by the nonzero image of $u$ gives
every matrix unit. The span is the full matrix algebra.
\end{proof}
For a single block ($a=n$), the given span is already the full matrix
algebra; in the source the odd module is then zero.
The trace-Gram determinant criterion gives the same statement without
choosing an algebraic closure: the trace pairing on the full matrix
algebra is nondegenerate, so a nonzero $n^2$-by-$n^2$ Gram minor is
equivalent to full span. This is exactly the determinant predicate
\srcid{DiscSt} and the proof in \srcid{ClosureConjugation.UrepStable}.

\subsubsection{The ordinary presentation}
\begin{srlemma}\label{lem:sr-ordinary-presentation-prime}
Let $\mathcal O\to B$ be a local homomorphism from a complete DVR to
a complete Noetherian local ring, inducing an isomorphism on residue
fields. Suppose $x_1,\ldots,x_D\in\mathfrak m_B$ satisfy
\[
            \mathfrak m_B\subseteq
                 (x_1,\ldots,x_D)+\varpi B+\mathfrak m_B^2.
\]
If $Z$ is prime and $\dim B/Z\ge D+1$, then $Z=0$.
Consequently $Z\subseteq Q$ for every prime $Q$ of $B$.
\end{srlemma}
\begin{proof}
The displayed cotangent generation and completeness give a surjection
$\pi:\mathcal O[[X_1,\ldots,X_D]]\twoheadrightarrow B$,
$X_i\mapsto x_i$, by successive approximation.
The source is a regular local domain of dimension $D+1$.
If its kernel after composition with $B\to B/Z$ contained a nonzero
element, the principal ideal theorem and the regular local dimension
formula would give $\dim B/Z\le D$.
Hence the composite has zero kernel.
Surjectivity now forces both $\ker\pi=0$ and $Z=0$.
\end{proof}
This proves the complete algebraic content of \srcid{PresOut0}.

\subsubsection{The paired presentation}
The following lemma is the last step in the exceptional local
inclusion. It states the exact presentation data, including the
independence that the Selmer calculation must supply.

\begin{srlemma}\label{lem:sr-paired-presentation-prime}
Let $\mathcal O\to B$ have the hypotheses of
Lemma~\ref{lem:sr-ordinary-presentation-prime}, and let
$A=B\oplus N$ be a finite commutative $B$-algebra as above.
Let $m_1,\ldots,m_g$ generate $N$ over $B$, and let
$x_1,\ldots,x_D\in\mathfrak m_B$.
Suppose the vectors
\[
                 \bar c_{ij}=(c_{ij,1},\ldots,c_{ij,D})\bmod\varpi
                    \quad(1\le i,j\le g)
\]
are linearly independent over the residue field, and put
\[
                  H=\varpi B+\mathfrak m_B^2+I_{\rm odd}.
\]
Assume
\[
 \mathfrak m_B\subseteq(x_1,\ldots,x_D)+H,\qquad
        \sum_j\left(\sum_a c_{ij,a}x_a\right)m_j\in HN
                                       \quad(1\le i\le g).
\]
If $Z,Q$ are primes of $B$ with
\[
 \dim B/Z\ge D+1,\qquad I_{\rm odd}\not\subseteq Z,\qquad
 I_{\rm odd}\not\subseteq Q\ \text{or}\ \operatorname{Ann}_BN\subseteq Q,
\]
then $Z\subseteq Q$.
\end{srlemma}
\begin{proof}
If $g=0$, then $N=0$ and $I_{\rm odd}=0$, contradicting the hypothesis
at $Z$. Thus $g\ge1$.
Since the $m_j$ generate $N$, choose $w_{ij}\in H$ such that
\[
                \sum_j\zeta_{ij}m_j=0,\qquad
          \zeta_{ij}=\sum_a c_{ij,a}x_a-w_{ij}.
\]
Multiplication by the adjugate proves
$(\det\zeta)m_j=0$ for every $j$, hence
\begin{equation}\label{eq:sr-det-ann-odd}
                             \det\zeta\in\operatorname{Ann}_BN.
\end{equation}
If $I_{\rm odd}\not\subseteq Q$, choose an odd product
$uv=b\notin Q$. Equation~\eqref{eq:sr-det-ann-odd} gives
$(\det\zeta)b=0$, using the injection $B\hookrightarrow A$.
Primality then gives $\det\zeta\in Q$.
If instead $\operatorname{Ann}_BN\subseteq Q$, the same conclusion
follows directly. Thus both alternatives at $Q$ give
$\det\zeta\in Q$.

We make the formal coordinates explicit.
The ring $A$ is complete local with maximal ideal
$\mathfrak m_BA+N$, since the image of $N$ modulo $\mathfrak m_BA$
has square zero. Every element of $I_{\rm odd}$ maps into
$\mathfrak m_A^2$.
The cotangent-generation assumption and the generators $m_j$
therefore give a surjection
\[
        \pi:P=\mathcal O[[X_1,\ldots,X_D,T_1,\ldots,T_g]]
                         \twoheadrightarrow A,
          \qquad X_a\mapsto x_a,\quad T_j\mapsto m_j .
\]
Lift each $w_{ij}$ to
$\widetilde w_{ij}\in(\varpi,\mathfrak m_P^2)$ and put
\[
                    z_{ij}=\sum_a c_{ij,a}X_a-\widetilde w_{ij}.
\]
The vectors $\bar c_{ij}$ are independent, so the $z_{ij}$ together
with the $T_j$ are independent in the relative cotangent space.
They extend to formal coordinates, as above.
The row relations $zT$ belong to $\ker\pi$.

Choose primes $Z',Q'$ of $A$ over $Z,Q$ by integrality.
At least one $m_j$ is outside $Z'$; otherwise all of $N$ lies in
$Z'$, so $I_{\rm odd}\subseteq Z'\cap B=Z$, a contradiction.
Put $P_Z=\pi^{-1}Z'$.
Then $zT\subseteq P_Z$ and at least one $T_j\notin P_Z$.
The adjugate identity gives $\det z\in P_Z$.
Consequently the prime
\[
                           K=\sqrt{(zT,\det z)}
\]
is contained in $P_Z$.
Lemma~\ref{lem:sr-generic-incidence} and
\eqref{eq:sr-incidence-dimension} give
\[
                        \dim P/K=\dim P-g=D+1.
\]
Since $A/Z'$ is finite integral over $B/Z$,
\[
                         \dim P/P_Z=\dim B/Z\ge D+1 .
\]
A strict inclusion $K\subsetneq P_Z$ in the complete regular local
ambient ring would lower the quotient dimension by at least one.
Hence $P_Z=K$.

Finally $P_Q=\pi^{-1}Q'$ contains $zT$ and $\det z$, because
$\det\zeta\in Q$. Thus
$P_Z=K\subseteq P_Q$, giving $Z'\subseteq Q'$.
Contracting to $B$ gives $Z\subseteq Q$.
\end{proof}
The proof supplies the power-series version needed for the source's
\srcid{CotangentIndependentOfResidues.presOut1}.
In particular it displays why the coefficient vectors must be
independent before choosing the formal coordinates.

\subsubsection{The actual local inclusion in the patched chain}
For the completed determinant algebra used in
Section~\ref{sec:sr-common-patched-module}, the source constructs
a finite commutative coefficient algebra $A=B\oplus N$ with
$NN\subseteq\mathfrak m_B$ and a matrix representation whose span is
the two-block algebra above.
The stable locus is detected by $I_{\rm odd}$ through
Lemma~\ref{lem:sr-odd-product-detector}.
The thick exceptional locus lies in
$V(\operatorname{Ann}_BN)$.

The first- and second-Selmer bounds and the spanning by the chosen
odd cocycles produce the presentation in
Lemma~\ref{lem:sr-paired-presentation-prime}:
\srcid{AlCSmul.pres} supplies the $x_a,m_j,c_{ij,a}$ with the
displayed generation, independence and relation properties.
In the single-block case $N=0$ it supplies
Lemma~\ref{lem:sr-ordinary-presentation-prime}.
Thus, once that finite coefficient algebra and those explicit Selmer
presentation data are constructed, the local implication (*) in
Proposition~\ref{prop:sr-completed-support-chain} is a proved
consequence of the two presentation lemmas.
The Galois source performs these two cases in
\srcid{AutomorphicDetection.GalOut.ctx_out}.

The algebra $A=B\oplus N$ is not an assumption that the two residual
constituents are separately polarized.
In the source it is constructed from the completed
Cayley--Hamilton algebra with its transpose involution, lifting its
orthogonal rank-one idempotents and identifying the two opposite
off-block corner modules with the odd summand $N$.
Its multiplication is
\[
       (b,u)(c,v)=(bc+\beta(u,v),\,bv+cu),
                    \qquad\beta:N\times N\longrightarrow B.
\]
The pairing is symmetric, and associativity of the corner products
gives $\beta(u,v)w=\beta(v,w)u$.
Residual block diagonality gives $\beta(N,N)\subseteq\mathfrak m_B$.
The construction, including determinant and involution compatibility,
is in \srcid{ClosureConjugation.OBLt},
\srcid{ComplOddPart}, \srcid{GrI} and \srcid{UrepOf},
assembled in \srcid{AutomorphicDetection.S6}.
Those constructions and the Selmer presentation are the remaining
arithmetic producers; the ordinary/paired prime-inclusion algebra
above has been proved independently.

\paragraph{Source locations.}
The generic incidence sources are
\href{https://github.com/ChristianKudeba/alpoge-rh-verification/blob/alpoge-source-7cbf7a0c/src/Development/AutomorphicArithmetic/AutomorphicVarieties/CommutativeAlgebraHeight/GDetrad.lean}
{\srcid{GDetrad}} and
\href{https://github.com/ChristianKudeba/alpoge-rh-verification/blob/alpoge-source-7cbf7a0c/src/Development/AutomorphicArithmetic/AutomorphicVarieties/CommutativeAlgebraHeight/GDetps.lean}
{\srcid{GDetps}}.
The completed local inclusion is
\href{https://github.com/ChristianKudeba/alpoge-rh-verification/blob/alpoge-source-7cbf7a0c/src/Development/AutomorphicArithmetic/AutomorphicVarieties/PatchingModules/AutomorphicDetection/GalOut.lean}
{\srcid{GalOut}}.
The exact paired presentation is
\href{https://github.com/ChristianKudeba/alpoge-rh-verification/blob/alpoge-source-7cbf7a0c/src/Development/AutomorphicArithmetic/AutomorphicVarieties/PatchingModules/SelmerScaling/CotangentIndependentOfResidues.lean}
{\srcid{CotangentIndependentOfResidues}},
and its Selmer producer is
\href{https://github.com/ChristianKudeba/alpoge-rh-verification/blob/alpoge-source-7cbf7a0c/src/Development/AutomorphicArithmetic/AutomorphicVarieties/PatchingModules/SelmerScaling/AlCSmul.lean}
{\srcid{AlCSmul}}.

\section{Ordinary modules and integral auxiliary control}
\label{sec:sr-ordinary-control}
We construct the ordinary function modules, their coefficient change and
classical specialization maps, and their integral old/new operators.
These supply the module-theoretic part of common support patching.
The attached Galois representations, their local compatibility, the
arithmetic Selmer bounds and the seed remain separate requirements.
The same constructions apply to ranks three and four.


\subsection{Construction of the ordinary modules and their finite-level control}
\label{sec:ordinary-module-construction}

This subsection constructs the finite-projective modules used by the
auxiliary-level argument.  It does not take \srcid{LevelAlt} or
\srcid{LevMod} as a premise.  The arithmetic setting is the definite
unitary setting of \srcid{HidaSetting}: \(L\) is a CM field,
\(G/L^+\) is the unitary group of the standard positive definite
Hermitian form of rank \(n\geq1\), and all places used for the ordinary
operators and for varying the auxiliary level are split places.
At the distinguished fixed auxiliary place the level is contained in
the first principal congruence subgroup at an odd rational prime \(p\),
with \(p\ne\ell\) and \(p\) unramified in \(L\).  These are exactly the
smallness conditions used below.  The construction is valid in every
rank; in the generic application \(n=3\).

Let \(A\) be a complete Noetherian local coefficient algebra with finite
residue field and maximal ideal \(\mathfrak m\).  Let \(A^m\) be the
coefficient lattice, with the integral tame action denoted by \(\tau\).
Write \(G_{\mathrm a}\) for the factors away from \(\ell\), \(G_\ell\)
for the product of the split factors above \(\ell\), and
\(\Gamma=G(L^+)\) for the rational subgroup.  On \(G_\ell\) fix the
Iwahori data
\[
 I(c)=N\,T\,N^-_c,\qquad c\geq0,
\]
with their usual opposite-unipotent filtration \(N^-_c\), diagonal
congruence subgroups \(T_c\), and a strictly positive central element
\(z\) in the positive diagonal monoid.  Products and factorizations here
are interpreted in the Iwahori decomposition; in particular
\(I(c)\subset I(0)\), \([I(0):I(c)]<\infty\), and multiplication by \(z\)
contracts the opposite-unipotent filtration by at least one step.
The character \(\theta:T\to A^\times\) is continuous in the precise
sense
\[
 \forall k\ \exists c\quad
 \theta(t)-1\in\mathfrak m^k\quad(t\in T_c).
 \tag{OC1}
\]
The universal weight characters occurring in the application have this
property.

For a tame compact open subgroup \(U\), let \(S(U)\) be the \(A\)-module
of functions
\[
 f:G_{\mathrm a}\times G_\ell\longrightarrow A^m
\]
which are left \(\Gamma\)-invariant, satisfy the tame
\(\tau\)-equivariance, are right \(N\)-invariant and right \(T\)-equivariant
with character \(\theta^{-1}\), and are uniformly continuous along
\(N^-_c\) modulo every \(\mathfrak m^k\).
The inverse in \(\theta^{-1}\) is the convention in the function model.
Put
\[
 V_k=\{f\in S(U):f(g)\in(\mathfrak m^k)^m\ \hbox{for every }g\},
 \qquad T_z=U_z.
\]
Here \(U_z\) is the integral double-coset sum, with the normalization of
the source Iwahori data.  It preserves \(S(U)\) and every \(V_k\).
Define the ordinary module by
\[
 O(U)=
 \{f\in S(U):\ \forall k\ \exists r\geq1,\
                 T_z^r f-f\in V_k\}.
 \tag{OC2}
\]

\begin{srlemma}[Smallness of the actual arithmetic levels]
\label{lem:ordinary-neat-level}
For every collection of compact open subgroups at the varying
auxiliary places, the tame level in the preceding construction is
neat in the following explicit sense:
\[
 \gamma g(u,b)=g,\quad
 \gamma\in\Gamma,\ u\in U,\ b\in I(0)
 \quad\Longrightarrow\quad u=b=1.
 \tag{OC3}
\]
Consequently arbitrary values in \(A^m\) can be prescribed on distinct
double cosets for \(\Gamma\backslash
(G_{\mathrm a}\times G_\ell)/(U\times I(c))\).
\end{srlemma}

\begin{proof}
Represent \(\gamma\) by a matrix \(M\in\mathrm{GL}_n(L)\).
At every archimedean embedding the matrices \(M^r\) are unitary.
At the varying auxiliary places, compactness implies that some positive
power of the local matrix belongs to a conjugate of
\(\mathrm{GL}_n(\mathcal O_{L,w})\): its orbit on the discrete set of
lattices is finite.  Taking a common multiple for the finitely many
varying places, and using the defining integral levels at the remaining
places, we find a positive integer \(a\) such that \(M^a\) is locally
conjugate to an integral invertible matrix at a selected place above
each finite place of \(L^+\).

Thus the coefficients of the characteristic polynomial of \(M^a\)
are integral at those selected places, as are their conjugates:
unitarity gives \(\overline{M^a}{}^t=(M^a)^{-1}\), and the latter is
also locally conjugate to an integral invertible matrix.
For a coefficient \(x\), both
\(x+\overline x\) and \(x\overline x\) belong to \(\mathcal O_{L^+}\);
the equation
\(X^2-(x+\overline x)X+x\overline x=0\) proves that \(x\) is integral
over \(\mathbb Z\).  Cayley--Hamilton now puts all powers of \(M^a\)
in a finitely generated \(\mathcal O_L\)-module of matrices.
After one fixed denominator is cleared, every entry of every power
is an algebraic integer bounded at all archimedean embeddings.
The Minkowski lattice has only finitely many such elements.
The set of powers of \(M^a\), and hence of \(M\), is finite.

At the distinguished auxiliary place the matrix is conjugate to a
finite-order element of \(1+pM_n(\mathcal O_{L,w})\).
This subgroup is torsion free when \(p\) is odd and unramified.
Indeed, if \(1+C\ne1\) and the minimum valuation of its entries is
\(r\geq1\), raising to a power prime to \(p\) retains that minimum,
whereas raising to the \(p\)-th power gives minimum \(r+1\).
For the latter assertion the first term \(pC\) has minimum \(r+1\);
all remaining binomial terms have larger minimum, since \(p\geq3\).
No positive power can therefore equal \(1\).
It follows that \(M=1\); the stabilizer identity then gives \(u=b=1\).

The finite double-coset assertion needed here follows from the definite
unitary finite-adelic orbit argument already proved in the arithmetic
finiteness subsection.  Passing from \(I(0)\) to \(I(c)\) preserves
finiteness because its index is finite.  In view of (OC3), the values on
the resulting double-coset representatives have no stabilizer
compatibility condition.
\end{proof}

\begin{srlemma}[Ordinary projector and finite freeness]
\label{lem:ordinary-projector-free}
For every arithmetic level above there is an \(A\)-linear projection
\[
 e_U:S(U)\longrightarrow O(U),\qquad
 e_U f=\lim_{r\to\infty}T_z^{\,r!}f
 \tag{OC4}
\]
in the valuewise \(\mathfrak m\)-adic filtration.  The module \(O(U)\)
is finite free over \(A\), and
\[
 O(U)\cap V_k=\mathfrak m^kO(U).
 \tag{OC5}
\]
Every positive diagonal \(U\)-operator is invertible on \(O(U)\).
All assertions hold uniformly for arbitrary compact open choices at
the finitely many varying auxiliary places; the finite free rank may
depend on the level.
\end{srlemma}

\begin{proof}
We give the finite approximation and algebra arguments, since this
avoids importing an unspecified ordinary control theorem.

Fix \(k\), choose \(c\) as in (OC1), and choose representatives
\(x_1,\ldots,x_b\) for the finite double-coset set with \(I(c)\).
Evaluation gives
\[
 E_k:S(U)\longrightarrow A^{bm}.
\]
There is an \(A\)-linear interpolation
\(s_k:A^{bm}\to S(U)\) with \(E_ks_k=1\), and its values are
\(N^-_c\)-invariant modulo \(\mathfrak m^k\).
Here is the explicit construction.  For \(b\in I(0)\), let \(t(b)\)
be the diagonal factor in \(b=y\,t(b)\,x\), with
\(y\in N^-_0\), \(x\in N\).  Its uniqueness gives
\[
 t(sb)=s\,t(b),\quad t(bs)=t(b)s,\quad t(bx)=t(b).
\]
The Iwahori interchange relation further gives
\(t(by)=t(b)s\) for some \(s\in T_c\) whenever \(y\in N^-_c\).
For prescribed vectors \(v_i\), set
\[
 s_k(v)\bigl(\gamma x_i(u,b)\bigr)
       =\theta(t(b))^{-1}\tau(u^{-1})v_i,
       \qquad b\in I(c).
 \tag{OC6}
\]
Neatness makes this well-defined; the identities for \(t(b)\) give
the required equivariances.  The final identity, together with
\(\theta(s)\equiv1\pmod{\mathfrak m^k}\), proves the asserted
\(N^-_c\)-invariance.  Choosing a deeper congruence for any higher
precision proves that \(s_k(v)\) belongs to \(S(U)\), rather than
merely defining a function at precision \(k\).

The elementary \(U_z\) contraction is
\[
 \begin{split}
 v\ \hbox{fixed by }N,T_c,N^-_{d+1},\quad d\geq c
 \quad\Longrightarrow\quad
 U_zv\ \hbox{fixed by }N,T_c,N^-_d .
 \end{split}
 \tag{OC7}
\]
It follows by writing the double-coset sum and moving an element of
\(N^-_d\) through its representatives; the Iwahori interchange places
the error in \(T_cN^-_{d+1}\).
Apply (OC7) modulo \(\mathfrak m^k\).  For a given \(f\), uniform
continuity supplies an initial opposite-unipotent depth, and (OC1)
makes \(T_c\) act trivially.  Thus, for all sufficiently large \(a\),
\(T_z^af\) is \(N^-_c\)-invariant modulo \(\mathfrak m^k\).
The difference
\[
 T_z^af-s_kE_k(T_z^af)
 \tag{OC8}
\]
has zero values at the representatives and is invariant modulo
\(\mathfrak m^k\) under the full level.  Hence (OC8) belongs to \(V_k\).
Also \(E_k(V_j)\subset(\mathfrak m^j)^{bm}\).

The finite set \((A/\mathfrak m^k)^{bm}\) implies that two sufficiently
late evaluations \(E_k(T_z^af)\) coincide modulo \(\mathfrak m^k\).
Equation (OC8) then makes the \(T_z\)-orbit eventually periodic
modulo \(V_k\).  Since \(T_zV_k\subset V_k\), all subsequent iterates
have the same period.  Consequently \(T_z^{r!}f\) is Cauchy modulo
every \(V_k\).
The valuewise limit belongs to \(S(U)\): left and right equivariance
pass to limits, and at each fixed precision a sufficiently late
approximating function supplies the required continuity depth.
Separatedness and linearity give the map \(e_U\) in (OC4).
Its image is exactly (OC2), it fixes that image, and it preserves
every \(V_k\).

We next prove finite freeness, without assuming it while constructing
the projector.  Applying \(e_U\) to (OC8), and using ordinary
recurrence, gives for every \(k\) finite maps \(E_k,s_k\) with
\[
 f-e_Us_kE_kf\in V_k\qquad(f\in O(U)).
 \tag{OC9}
\]
For precision \(1\), write
\(\sigma=e_Us_1:A^b\to O(U)\) and
\(E=E_1|_{O(U)}\), enlarging \(b\) to include coefficient coordinates.
If \(f\in O(U)\cap V_j\), use (OC9) at precision \(j+1\).
Its coefficients lie in \(\mathfrak m^j\).
Apply (OC9) at precision \(1\) to each of the finitely many ordinary
vectors \(e_Us_{j+1}(e_i)\), multiply the errors by those coefficients,
and sum.  This gives
\[
 \exists a\in(\mathfrak m^j)^b,\qquad
 f-\sigma a\in V_{j+1}.
 \tag{OC10}
\]
Iteration starting at \(j=0\), followed by completeness of the finite
module \(A^b\), proves that \(\sigma\) is surjective.
Starting at a general \(j\) proves
\(O(U)\cap V_j\subset\mathfrak m^jO(U)\); the reverse inclusion follows
directly from values.  This proves (OC5).

Now \(C=\sigma E\) is an endomorphism of the finite \(A\)-module
\(O(U)\) congruent to the identity modulo \(\mathfrak mO(U)\).
Nakayama makes \(C\) surjective, and a surjective endomorphism of a
finite module over a Noetherian ring is bijective.
Thus \(EC^{-1}\) is a section of \(\sigma\), so \(O(U)\) is finite
projective over the local ring \(A\), hence finite free.

Finally, ordinary recurrence modulo \(V_1=\mathfrak mO(U)\) says that
\(T_z\) is surjective modulo \(\mathfrak m\), so the same argument makes
\(T_z\) invertible on \(O(U)\).
For any positive diagonal \(a\), choose \(r\) so that \(z^ra^{-1}\)
is positive.  The positive \(U\)-operators commute and multiply
according to the positive monoid.  Since their product
\(U_{z^ra^{-1}}U_a=U_z^r\) is invertible, \(U_a\) is invertible.
\end{proof}

\begin{srlemma}[Auxiliary group-ring projectivity and exact norm control]
\label{lem:ordinary-aux-projective}
Let \(V\triangleleft U\) have finite index and set \(\Delta=U/V\).
Then \(O(V)\) is finite projective over \(A[\Delta]\).  Moreover the
unnormalized norm
\[
 \operatorname{Nm}_{U/V}=\sum_{\delta\in\Delta}\delta
\]
has image \(O(U)\), and
\[
 \ker\bigl(\operatorname{Nm}_{U/V}:O(V)\to O(U)\bigr)
       =\sum_{\delta\in\Delta}(\delta-1)O(V).
 \tag{OC11}
\]
In particular \(O(V)_\Delta\simeq O(U)\).
No invertibility of \(|\Delta|\) in \(A\) is required.
\end{srlemma}

\begin{proof}
Choose one representative \(R(g)\) for each double coset at level
\(U\times I(0)\).  In
\(g=\gamma R(g)(u,b)\), neatness makes the \(U\)-coordinate \(u\)
unique.  Let \(\chi(g)\) be \(1\) when \(u\in V\), and \(0\) otherwise.
Then \(\chi\) is left \(\Gamma\)-invariant and right
\(V\times I(0)\)-invariant, and
\[
 \sum_{\delta\in\Delta}\chi(g\delta)=1.
 \tag{OC12}
\]
Indeed precisely one right coset sends the \(U\)-coordinate to \(V\).
Multiplication by \(\chi\) maps \(O(V)\) into \(S(V)\), since \(\chi\)
is \(I(0)\)-invariant.  Put
\(\Phi=e_V\circ M_\chi\in\operatorname{End}_A(O(V))\).
The ordinary projector commutes with the away translations: this
holds for each \(U_z^{r!}\), hence for their separated limits.
Taking \(e_V\) of the partition identity (OC12) yields
\[
 \sum_{\delta\in\Delta}\delta\Phi\delta^{-1}=1.
 \tag{OC13}
\]
Since \(O(V)\) is finite free over \(A\), the equivariant surjection
\[
 A[\Delta]\otimes_AO(V)\longrightarrow O(V),\qquad
 \delta\otimes v\longmapsto\delta v
\]
has the \(A[\Delta]\)-linear section
\[
 v\longmapsto
 \sum_{\delta\in\Delta}\delta\otimes\Phi(\delta^{-1}v).
\]
Linearity follows by reindexing and the composite is the identity by
(OC13).  This proves finite projectivity.

The fixed submodule \(O(V)^\Delta\) is exactly \(O(U)\), by the
function definitions and the level-independent recurrence (OC2).
For \(v\) fixed, (OC13) gives
\(\operatorname{Nm}(\Phi(v))=v\), proving norm surjectivity.
If \(\operatorname{Nm}(v)=0\), then
\[
 \sum_\delta\Phi(\delta^{-1}v)=\Phi(\operatorname{Nm}(v))=0
\]
and (OC13) gives
\[
 v=\sum_\delta
       \bigl(\delta\Phi(\delta^{-1}v)-\Phi(\delta^{-1}v)\bigr).
\]
This proves the nontrivial inclusion in (OC11); the reverse inclusion
is the norm identity.
\end{proof}

\paragraph{Hecke localization and the actual smooth auxiliary module.}
Use only the tame Hecke operators away from the varying auxiliary
places, together with the ordinary operators at \(\ell\).
They commute on \(O(U)\); their image
\(\mathbb T_U\subset\operatorname{End}_A(O(U))\) is a finite commutative
\(A\)-algebra.  Completeness of \(A\) decomposes \(\mathbb T_U\) into
its local factors.  The factor selected by the prescribed residual
Hecke character is a central-idempotent summand
\[
 M(U)=\epsilon_UO(U).
 \tag{OC14}
\]
It is finite free over \(A\).  At level \(V\triangleleft U\), the
projector \(\epsilon_V\) commutes with the \(\Delta\)-action, because
every selected Hecke operator does.  Thus \(M(V)\) is a projective
\(A[\Delta]\)-summand of \(O(V)\).

For clarity about invariants, view all \(M(K)\), for compact open
\(K\) in the product of the varying local groups, in the single
function space, and put
\[
 \mathcal M=\bigcup_KM(K).
 \tag{OC15}
\]
This union is directed and translation-stable, so it is an actual
smooth auxiliary representation.  It satisfies
\(\mathcal M^K=M(K)\).
Here the only point beyond the function invariance is compatibility
of the localized summand with intersection.  If \(Y\subset Y'\) are
the corresponding finite ordinary modules, then Artin--Rees gives a
constant \(c\) with
\(Y\cap\mathfrak m^{r+c}Y'\subset\mathfrak m^rY\).
Consequently the simultaneous generalized residual-eigenvalue
condition, expressed at every \(\mathfrak m\)-adic precision, cuts out
the same submodule of \(Y\) whether tested in \(Y\) or \(Y'\).
Intersecting two levels and using this observation proves the claimed
fixed-vector equality.
The quotient \(K/K'\) acting on \(M(K')\) is the actual auxiliary
quotient: the local embedding used to form the tame level has
preimage \(K'\) inside \(K\).  Lemma
\ref{lem:ordinary-aux-projective} and the commuting central idempotent
therefore prove that \(\mathcal M^{K'}\) is projective over
\(A[K/K']\).

These constructions supply finite freeness, group-ring projectivity,
the smooth module, and the norm/augmentation control in the
finite-level input.  They do not by themselves prove that the chosen
localized factor is nonzero, that its Hecke determinant is realized
by a Galois representation, or that every classical specialization
has the prescribed local-global compatibility.  Those are separate
automorphic arithmetic inputs.  In particular this construction must
not be used to declare the automorphic seed proved.

\paragraph{Exact source locations.}
The preceding construction follows the proofs, rather than just the
output interfaces, in the frozen source
\srcid{alpoge-source-7cbf7a0c}, under
\srcid{Development.AutomorphicArithmetic.AutomorphicVarieties}.
The arithmetic smallness proof is
\srcid{HeckeAlgebras.OrdinaryHeckeCharacters.NtExpand.neatU}.
The interpolation and contraction proofs are
\srcid{BuildingsResidueChars.ExistsLinearSectionAutomorphicFunctionsOfFree},
\srcid{ExistsIwahoriTorusProjection},
\srcid{ExistsSectionAutomorphicSpaceWithPrescribedValue}, and
\srcid{UOperatorPowEventuallyInvariantModPow}.
The algebra proof of (OC4)--(OC5) is
\srcid{HeckeAlgebras.OrdinaryHeckeCharacters.OrdFree.ordAbs},
\srcid{ordStruct}, and \srcid{ordFree}; (OC12)--(OC13) and
group-ring projectivity are proved in
\srcid{OrdProj.ordProj\_chi}, \srcid{ordProj\_phi},
\srcid{ordProj\_higman}, and \srcid{ordProj}.
Hecke localization is \srcid{HiComm.hiComm}; the fixed-level
comparison and localized projectivity are the explicit
\srcid{pb\_locInf} and \srcid{pb\_main} arguments in
\srcid{PiBasic}, culminating in \srcid{piBasic}.
All these file paths are below
\url{https://github.com/ChristianKudeba/alpoge-rh-verification/tree/alpoge-source-7cbf7a0c/src/Development/AutomorphicArithmetic/AutomorphicVarieties}.


\subsection{Base change and classical specialization of the constructed ordinary modules}
\label{sec:ordinary-classical-control}

We continue with the function modules of
\ref{sec:ordinary-module-construction}.  This subsection proves the
ordinary base-change and classical control statements used to obtain
separating specializations.  It does not assume an ordinary
automorphy-lifting theorem.

\begin{srlemma}[Ordinary base change]
\label{lem:ordinary-base-change}
Let \(A\to A'\) be a local homomorphism of complete Noetherian local
rings with finite residue fields.  Base-change the tame coefficient
action and the torus character.  At each of the neat finite levels,
coefficient change induces an isomorphism
\[
 A'\otimes_A O_A(U)\ \xrightarrow{\ \sim\ }\ O_{A'}(U).
 \tag{CC1}
\]
This assertion does not require \(A'\) to be flat over \(A\).
\end{srlemma}

\begin{proof}
Coefficient change carries value filtration \(V_k\) into \(V'_k\).
It intertwines every \(U_z\)-power and hence the factorial ordinary
projectors.  In particular it defines a map \(j:O_A(U)\to O_{A'}(U)\).

Choose one finite interpolation system at precision \(1\), with
evaluation functionals \(E_p\), interpolation sections \(s_p\), and
\(w_p=e_As_p\).  The proof of (OC9), using this fixed interpolation
system, gives
\[
 f-\sum_pE_p(f)w_p\in\mathfrak m_AO_A(U).
 \tag{CC2}
\]
The same sections, changed to \(A'\), give
\[
 f'-\sum_pE'_p(f')j(w_p)\in\mathfrak m_{A'}O_{A'}(U).
 \tag{CC3}
\]
To justify using precisely the same system, both characters are
trivial modulo their maximal ideals on the chosen \(T_c\).
For an ordinary \(f'\), a sufficiently large factorial power is
congruent to \(f'\) modulo \(V'_1\) and is \(N^-_c\)-invariant there,
by (OC7).  Its interpolation has the same representative values.
Subtract, apply that factorial power to the finite interpolation
sections, and use \(e_As_p=w_p\) modulo \(V_1\).
This proves (CC3) with the displayed \(j(w_p)\), rather than with an
independently chosen generating family.

Let \(C(f)=\sum_pE_p(f)w_p\).  By (CC2), Nakayama and the Noetherian
endomorphism lemma make \(C\) an automorphism of \(O_A(U)\).
Let
\[
 L:A'\otimes_AO_A(U)\to O_{A'}(U),\qquad
 R(f')=\sum_pE'_p(f')\otimes w_p.
\]
Evaluation is compatible with \(j\), so
\(RL=1\otimes C\).
The latter is invertible, proving injectivity of \(L\).
Equation (CC3) and Nakayama prove surjectivity of \(L\), since
\(O_{A'}(U)\) is finite over \(A'\).
\end{proof}

\paragraph{The elementary algebraic weight data.}
Let \(\mathcal O\) be a sufficiently large finite coefficient DVR
and let \(\Omega\) be an algebraically closed characteristic-zero
field containing its fraction field.
An algebraic dominant weight at the split \(\ell\)-adic places gives
a finite-dimensional representation \(\xi\) with a lowest-weight
coordinate numbered \(0\), diagonal character \(\chi\), and an
integral weight lattice.  Enlarge \(\mathcal O\) to contain the
finitely many coefficient-field embeddings used to define it.
The properties needed below are the following concrete properties
of this weight lattice:
\[
 \begin{array}{ll}
 \text{(i)}&
  e_0^\vee\xi(n)=e_0^\vee\quad(n\in N),\qquad
  e_0^\vee\xi(t)=\chi(t)e_0^\vee\quad(t\in T);\\[2mm]
 \text{(ii)}&
  \xi(N)\ \hbox{is integral},\qquad
  \chi(a)^{-1}\xi(a)\ \hbox{is integral for }a\hbox{ positive};\\[2mm]
 \text{(iii)}&
  \chi(z)^{-1}\xi(z)\ \hbox{has every column except column }0
                     \hbox{ in }\mathfrak m_{\mathcal O};\\[2mm]
 \text{(iv)}&
  \xi(y)\equiv1\pmod{\mathfrak m_{\mathcal O}^k}
   \quad(y\in N^-_c,\ c\hbox{ sufficiently large for }k).
 \end{array}
 \tag{CC4}
\]
One may construct the lattice using the integral Schur modules for
the successive differences of the highest weight, tensoring with
a determinant power.  The matrix coefficients of unipotents are
integral polynomials, which proves (i), (ii), and (iv).
In a weight basis the normalized action of the strictly positive
\(z\) is \(1\) at the extremal coordinate and has positive valuation
at every other weight; its extremal weight has multiplicity one.
This proves (iii).
These observations work for arbitrary \(\ell\); no
Fontaine--Laffaille range is used in this representation-theoretic
control argument.

Suppose now that the specialized weight character has the form
\[
 \theta(t)=\phi(t)\chi(t),\qquad \phi|_{T_c}=1
 \tag{CC5}
\]
for a finite-depth character \(\phi:T\to\mathcal O^\times\).
Choose \(c\) sufficiently deep for its finite part and the integral
weight representation.
Let \(C(U,c,\phi)\) be the classical algebraic function space:
its functions \(F\) transform on the left by \(\xi\) at the
\(\ell\)-adic component of \(\Gamma\), satisfy the tame coefficient
transformation, and have right Iwahori character \(\phi^{-1}\)
at depth \(c\).
For these functions set
\[
 qF(g)=\xi(g_\ell)^{-1}F(g),\qquad
 \operatorname{ctrl}(F)(g)=(qF(g))_0 .
 \tag{CC6}
\]
The coefficient index from the tame lattice is carried along in
this formula.  The classical ordinary operator is
\[
 U_z^{\mathrm{cl}}=
   \chi(z)^{-1}\sum_{n\in N/(N\cap zNz^{-1})}[nz].
 \tag{CC7}
\]
The row identities in (CC4)(i) give, by evaluating the finite sum,
\[
 \operatorname{ctrl}(U_a^{\mathrm{cl}}F)
        =U_a\,\operatorname{ctrl}(F)
 \tag{CC8}
\]
for every positive \(a\).  The control map also intertwines away
translations and finite away-Hecke sums.

Evaluation of all coordinates of \(qF\) at representatives for
\(\Gamma\backslash(G_{\mathrm a}\times G_\ell)/(U\times I(c))\)
identifies \(C(U,c,\phi)\) with a finite-dimensional coordinate
space over \(\Omega\).
Indeed equivariance recovers the function from these coordinates;
arbitrary coordinates give a function because the level is neat.
The integral coordinate lattice \(L_c\) consists exactly of those
functions for which every coordinate of \(qF(g)\) belongs to
\(\mathcal O\) for every \(g\).
The integral Iwahori action shows the equivalence with integrality
at the chosen representatives.
Consequently \(L_c\) is finite free and is preserved by (CC7).

\begin{srproposition}[Classical ordinary control]
\label{prop:classical-ordinary-control}
Let \(C(U,c,\phi)^{\mathrm{ord}}\) be the sum of generalized
eigenspaces of \(U_z^{\mathrm{cl}}\) whose eigenvalue is an integral
unit.  Then \(\operatorname{ctrl}\) is injective on this space and
\[
 \operatorname{ctrl}\bigl(C(U,c,\phi)^{\mathrm{ord}}\bigr)
  =\Omega\cdot O_{\mathcal O}(U)
 \tag{CC9}
\]
inside the common scalar-valued function space.
If \(A\to\mathcal O\to\Omega\) is an arithmetic specialization,
(CC1) and (CC9) therefore give a specialization map from the
constructed \(A\)-ordinary module to classical ordinary forms.
The images of any \(A\)-basis are linearly independent over
\(\Omega\).
The map commutes with all selected Hecke operators and all away
translations, including the auxiliary local group actions.
\end{srproposition}

\begin{proof}
The ordinary part of the finite integral lattice \(L_c\) is the
image of \(\lim_r(U_z^{\mathrm{cl}})^{r!}\).
This follows from the same finite-quotient periodicity proof as
(OC4), now on a finite free module.
After extension to \(\Omega\), its span is exactly the direct sum
of the generalized eigenspaces with integral-unit eigenvalue:
the complementary eigenvalues have positive valuation, and the
factorial projector is \(0\) there and \(1\) on the unit part.

The contraction controlling the kernel is particularly simple.
If every coordinate of \(qF\) belongs to \(\mathfrak m^j\), while
its coordinate \(0\) belongs to \(\mathfrak m^{j+1}\), then
\[
 q(U_z^{\mathrm{cl}}F)\ \hbox{has all coordinates in }
                   \mathfrak m^{j+1}.
 \tag{CC10}
\]
For an individual summand of (CC7), its untwisted value is
\[
 \xi(n)\bigl(\chi(z)^{-1}\xi(z)\bigr)
                 qF(g(nz)).
\]
Columns different from \(0\) contribute
\(\mathfrak m\mathfrak m^j\), by (CC4)(iii).
Column \(0\) contributes \(\mathfrak m^{j+1}\) by hypothesis.
The integral matrix \(\xi(n)\) and the finite sum preserve the
ideal, proving (CC10).

If \(F\) belongs to the integral ordinary lattice and
\(\operatorname{ctrl}(F)=0\), (CC10), ordinary recurrence modulo
every \(\mathfrak m^kL_c\), and stability of these lattices show
inductively that \(F\in\mathfrak m^jL_c\) for every \(j\).
Separatedness gives \(F=0\).
The same proves injectivity after scalar extension: for finitely
many integral independent functions, independence is preserved
when the fraction field is extended to \(\Omega\).
To verify this last elementary fact, clear denominators first over
\(\operatorname{Frac}\mathcal O\), then expand the finitely many
extension-field coefficients in a basis of their span over that
fraction field and compare coefficients in the function identity.
It follows that the control map is injective on the whole classical
ordinary space.

For an integral ordinary classical \(F\), the function
\(f=\operatorname{ctrl}(F)\) has integral values.
It satisfies the scalar equivariances with \(\theta\) by
(CC4)(i) and (CC5).
It is uniformly continuous along the opposite unipotent:
\(\xi(y)\equiv1\pmod{\mathfrak m^k}\) at sufficiently deep level,
by (CC4)(iv).
Equation (CC8) transfers ordinary recurrence to \(f\).
Thus the control image is contained in \(\Omega\cdot O_{\mathcal O}(U)\).

We prove the reverse inclusion by an explicit approximation.
Take \(f\in O_{\mathcal O}(U)\).
Choose \(c'\geq c\) so deep that both \(f\) and \(\xi\) are
\(N^-_{c'}\)-invariant modulo \(\mathfrak m\).
At the finite representatives at depth \(c'\), prescribe the
coordinates
\[
 qF(x)_0=f(x),\qquad qF(x)_i=0\quad(i\ne0).
\]
Neat interpolation gives an integral classical form at that depth.
Its scalar control differs from \(f\) by valuewise
\(\mathfrak m\), since the two controls have the same
representative values and the same congruence invariance.

A sufficiently large positive power of \(U_z^{\mathrm{cl}}\)
moves this form from depth \(c'\) to depth \(c\).
This is the same one-step Iwahori interchange as (OC7):
the opposite-unipotent depth drops by one at each application;
the finite character is already trivial on \(T_c\).
Choose the power also divisible by a recurrence period for \(f\)
modulo \(\mathfrak m\).
By (CC8), its scalar control still agrees with \(f\) modulo
\(\mathfrak m\).
Apply the classical factorial ordinary projector at depth \(c\).
The resulting ordinary classical form has control congruent to
\(f\), because (CC8) identifies its control limit with the scalar
ordinary projector, which preserves \(V_1\) and fixes \(f\).

Let
\[
 L=\bigl(\operatorname{ctrl}
       (C(U,c,\phi)^{\mathrm{ord}})\bigr)
               \cap O_{\mathcal O}(U).
\]
This is an \(\mathcal O\)-submodule of a finite
\(\mathcal O\)-module.  The preceding approximation proves
\(L+\mathfrak m O_{\mathcal O}(U)=O_{\mathcal O}(U)\);
here valuewise \(\mathfrak m\) is
\(\mathfrak mO_{\mathcal O}(U)\) by (OC5).
Nakayama gives \(L=O_{\mathcal O}(U)\), proving (CC9).

Finally apply (CC1).
A basis of \(O_A(U)\) becomes an \(\mathcal O\)-basis of
\(O_{\mathcal O}(U)\).
As integral scalar-valued functions these basis vectors remain
independent over \(\Omega\), by the coefficient comparison just
used.  Their unique inverse-control images are therefore
independent classical forms.
All constructions commute with the indicated Hecke operators:
coefficient change and control do so by their pointwise formulas,
and the ordinary projectors do so by their factorial limits.
The same formulas commute with away translations, which proves
the asserted auxiliary equivariance.
\end{proof}

\paragraph{Automatic continuity of arithmetic coefficient maps.}
The use of a local specialization in (CC1) is justified
intrinsically.  Let \(T\) be a complete Noetherian local
\(\mathcal O\)-algebra with the same finite residue field and let
\(x:T\to\overline{\mathbb Q}_\ell\) be an
\(\mathcal O\)-algebra map.
For \(a\in\mathfrak m_T\) one has \(|x(a)|_\ell<1\).
Otherwise \(x(a)^{-1}\) is integral over \(\mathbb Z_\ell\);
reversing its monic integral equation gives a polynomial
\(1+c_1X+\cdots+c_rX^r\) vanishing at \(x(a)\), with
\(c_i\in\mathbb Z_\ell\).
Its value at \(a\) is a unit of \(T\), a contradiction.
Writing an arbitrary element as its \(\mathcal O\)-residue
representative plus an element of \(\mathfrak m_T\) gives
\(|x(t)|_\ell\leq1\) for every \(t\).
For finitely many generators of \(\mathfrak m_T\), take the maximum
of their norms to obtain \(b<1\); then
\[
 |x(\mathfrak m_T^j)|_\ell\leq b^j.
\]
Thus \(x\) is continuous.  Choose finitely many topological
generators \(t_1,\ldots,t_s\) of \(T\) over \(\mathcal O\); these
exist by completeness, finite generation of \(\mathfrak m_T\), and
surjectivity on the residue field.  The finitely many algebraic
elements \(x(t_i)\), together with the coefficient field, belong to
one finite extension \(K/\mathbb Q_\ell\).
Every element of \(T\) is an adic limit of polynomials in the
\(t_i\); continuity and completeness of \(K\) imply
\(x(T)\subset\mathcal O_K\).
The map \(T\to\mathcal O_K\) is local, since maximal-ideal elements
have norm less than one.
This gives the finite coefficient DVR required in (CC1).
The residue-preserving complete local coefficient rings in the
application therefore have exactly the continuity and coefficient
field containment needed for ordinary specialization; continuity is
not silently imposed on an arbitrary field-valued algebra map.

\paragraph{Remaining arithmetic input and exact sources.}
This proves the module-theoretic classical specialization and
basis independence.  To pass from classical forms to attached
Galois representations, one still uses the definite-unitary
automorphic decomposition, its attached semisimple Galois
representations, and local-global compatibility.  To separate
Hecke-algebra elements by these specializations one additionally
uses density of the algebraic weight characters in the specified
weight algebra.  Neither the Galois-realization assertion nor
weight density is being inferred solely from the control
isomorphism.

The exact frozen source is under
\srcid{Development.AutomorphicArithmetic.AutomorphicVarieties},
tag \srcid{alpoge-source-7cbf7a0c}.
Base change is the finite evaluation argument
\srcid{HeckeAlgebras.OrdinaryHeckeCharacters.OrdBCMapOrd.ordBC\_abs},
\srcid{ordBC\_key}, and \srcid{ordBC}.
Control is
\srcid{CtrlOrd.ctrlOrd\_inj}, \srcid{ctrlOrd\_approx},
and \srcid{ctrlOrd}; the elementary lattice calculation (CC10) is
\srcid{BuildingsResidueChars.UpOperatorMapsLatticeOfControlValues}.
The other explicit construction helpers are
\srcid{ExistsIntegralFormLiftingControlMap},
\srcid{ExistsAutomorphicFormWithPrescribedIwahoriValues},
and \srcid{ExistsFreeLatticeSpanningOrdinaryPart}.
The coefficient-change independence lemma is
\srcid{HeckeAlgebras.OrdinaryHeckeCharacters.FnIndep.fnIndep}.
The simultaneous specialization and all Hecke/auxiliary
equivariances are assembled by the actual proof
\srcid{HeckeAlgebras.HeckeFixedLevelSpaces.FreeMap.absSpec},
then by \srcid{OrdinaryHeckeCharacters.SpecIntIff.spec}.
The continuity calculation is the proof
\srcid{DenseThresholdCases.TCoreDet.cd\_norm\_lt},
\srcid{cd\_pow}, and \srcid{cd\_cont}.
All files are below
\url{https://github.com/ChristianKudeba/alpoge-rh-verification/tree/alpoge-source-7cbf7a0c/src/Development/AutomorphicArithmetic/AutomorphicVarieties}.


\subsection{Explicit old/new maps from the integral Iwahori algebra}
\label{sec:explicit-old-new}

We construct the two maps required in the auxiliary control argument.
They are defined by finite sums of Hecke operators.  In particular the
construction never divides by the index of an Iwahori subgroup in a
maximal compact subgroup.  The only arithmetic coefficient comparison
used in this subsection is displayed in (ON8) below.

\paragraph{The integral calculation at one split place.}
Let \(k\) be a nonarchimedean local field, with residue cardinality
\(Q\), \(K=\mathrm{GL}_n(\mathcal O_k)\), and standard Iwahori \(I\).
Let \(Q\) be a unit in the coefficient ring.
Write \(X_i\), for \(0\leq i<n\), for the commuting Bernstein operators
with the normalization
\[
 X_i=Q^i\,\Theta(\operatorname{diag}_i(\varpi)).
 \tag{ON1}
\]
Thus, if \(t_r=\operatorname{diag}(\varpi,\ldots,\varpi,1,\ldots,1)\)
has \(r\) entries \(\varpi\), then
\[
 X_i[I t_i I]=Q^i[I t_{i+1}I].
\]
The positive diagonal Hecke operators are units when \(Q\) is a unit.
One can see this in the affine Iwahori algebra: the simple operators
satisfy
\(T_s^2=(Q-1)T_s+Q\), hence
\(T_s^{-1}=Q^{-1}(T_s-(Q-1))\), and translations are products of these
operators and length-zero units.  The commuting positive diagonal
action consequently extends to the torus action \(\Theta\) in (ON1).
This proves existence of the torus action used here directly from
double-coset operators.

Put, in the polynomial ring \(\mathcal P=\mathbb Z[Q,Q^{-1}]
[X_0,\ldots,X_{n-1}]\),
\[
 D(X)=\prod_{i<j}(X_j-X_i),\qquad
 N_Q(X)=\prod_{i<j}(X_j-QX_i),\qquad
 \partial(f)=
 \frac{\sum_{w\in S_n}\operatorname{sgn}(w)\,w(f)}{D(X)}.
 \tag{ON2}
\]
The numerator is divisible by \(D\) over this integral polynomial
ring: it vanishes after each substitution \(X_i=X_j\), so the distinct
prime linear factors divide it.  In particular no factorial is
inverted in \(\partial\).  Its values are symmetric polynomials and
\(\partial(sf)=s\partial(f)\) for symmetric \(s\).

Let \(E=[K:I]\) denote the trace operator from Iwahori-fixed to
\(K\)-fixed vectors, considered as an endomorphism of the
Iwahori-fixed module.  This notation denotes the double-coset sum,
not scalar multiplication by the integer \([K:I]\) on arbitrary
vectors.  In the Iwahori algebra \(E=\sum_{w\in S_n}T_w\).
Its integral identities are
\[
 \begin{split}
 Es(X)&=s(X)E &&(s\ \hbox{symmetric}),\\
 E f(X)E&=\partial(N_Qf)(X)\,E .
 \end{split}
 \tag{ON3}
\]
Here and below a polynomial on the right means its Bernstein action.

For completeness, these identities can be verified in the universal
Bernstein algebra.  Its generators satisfy the finite Hecke
quadratic and braid relations, and, for adjacent \(i,j=i+1\),
\[
 \begin{split}
 T_iX_i&=X_jT_i+(Q-1)X_i,\\
 T_iX_j&=X_iT_i-(Q-1)X_i,\qquad
 T_iX_a=X_aT_i\quad(a\ne i,j).
 \end{split}
 \tag{ON4}
\]
These are the rank-one double-coset calculation with
\(I\operatorname{diag}(\varpi,1)I\); operators at other coordinates
commute.  To check the universal identities without making a
specialization injective, use the fraction field of \(\mathcal P\)
and the crossed product by \(S_n\).  There
\[
 T_i=
 \frac{(1-Q)X_i}{X_j-X_i}
   +\frac{QX_j-X_i}{X_j-X_i}\,[s_i].
 \tag{ON5}
\]
These operators satisfy (ON4), the quadratic relation, and the braid
relation, by the two-variable and three-variable expansions,
respectively.  Straightening (ON4) spans the universal algebra by
\(\mathcal P T_w\).  In (ON5), the coefficient of \([w]\) in \(T_w\)
is a nonzero product over the inversions of \(w\); all remaining
terms have smaller Bruhat length.  Thus these spanning elements are
independent and the crossed-product calculation is faithful.

The equality \(T_iE=QE\) implies that if \(A_w\) is the coefficient of
\([w]\) in \(E\), then \(A_{s_iw}=s_i(A_w)\).
The coefficient at the longest permutation is obtained solely from
its longest reduced word, and is
\(w_0(N_Q/D)\).  Hence
\[
 E=\sum_w w(N_Q/D)[w],\qquad
 E(f)=\partial(N_Qf).
 \tag{ON6}
\]
It follows that \(E\) commutes with symmetric polynomials and that
\(E f E=\partial(N_Qf)E\), since \(E\) has symmetric image.
The result is an identity in the integral universal algebra, by
faithfulness; (ON5) is only a verification device.
Specializing \(Q\) and applying the actual Iwahori action proves
(ON3) over the original coefficient ring.

Here is an explicit finite divided-difference dual family.
Let \(\mathcal R=\{(i,j):0\leq i<j<n\}\).
For each subset \(H\subset\mathcal R\), put
\[
 a_H=\prod_{(i,j)\in H}X_j,\qquad
 c_H=\prod_{(i,j)\in\mathcal R\setminus H}(-X_i).
\]
Then
\[
 \sum_H a_H\,w(c_H)=
 \prod_{i<j}(X_j-X_{w(i)})
 =\begin{cases}D,&w=1,\\0,&w\ne1.\end{cases}
\]
For \(w\ne1\), choose \(i\) with \(i<w(i)\); the factor indexed by
\((i,w(i))\) is zero.  Expanding the antisymmetrizer therefore gives
\[
 \sum_Ha_H\partial(c_Hf)=f,\qquad
 \sum_Hc_H\partial(a_Hf)=f,\qquad
 \sum_H a_H(X)E c_H(X)=N_Q(X).
 \tag{ON7}
\]
The second equality follows by applying a permutation to the dual
identity for its inverse, which interchanges the two families.
For the last equality, insert (ON6); only the identity permutation
survives in its crossed-product expansion.
These formulas prove the precise integral identity used in the
old/new construction.

\paragraph{The global modules and the coefficient comparison.}
Let \(\mathcal Q=\{w_1,\ldots,w_q\}\) be the distinct split auxiliary
places, away from \(\ell\) and from the fixed ramification set.
Use the smooth module \(\mathcal M\) constructed in
\ref{sec:ordinary-module-construction}; at these places write
\[
 V=\mathcal M^{\prod_j I_j},\qquad
 M=\mathcal M^{\prod_j K_j}.
\]
The spherical module is identified with the fixed base module by the
actual maximal-level comparison.  More generally its coefficients
are related by the ring map \(\beta:R\to B\); the inclusion
\(\iota:M\to V\) is additive and satisfies
\(\iota(\beta(r)m)=r\iota(m)\).
Let
\[
 P_j(T)\in B[T]
\]
be monic of degree \(n\).
The comparison with the base Galois determinant is exactly the
following spherical coefficient identity:
\[
 e_k(U_{j,0},\ldots,U_{j,n-1})\iota(m)
   =\iota\bigl((-1)^k[P_j]_{n-k}\,m\bigr),
 \quad 0\leq k\leq n.
 \tag{ON8}
\]
The normalization of \(U_{j,i}\) is (ON1).
In the application
\[
 P_j(T)=D_{\mathrm{base}}(T-[\mathrm{Frob}_{w_j}]),
\]
with the inverse-Frobenius convention of the original
\srcid{Pair}; (ON8) uses that same convention.

We spell out the proof of (ON8) from the arithmetic realization
input, so that it is not an additional lifting assumption.
The actual classical specializations separate vectors of
\(\mathcal M\).  Under such a specialization the auxiliary action
is an isotypic sum of an irreducible smooth complex local
representation \(\sigma_j\).  A nonzero image of a spherical vector
forces \(\sigma_j^{K_j}\ne0\).
The normalized Bernstein eigenvalues \(\alpha_i\) of this spherical
representation satisfy
\[
 \operatorname{charpoly}(\rho(\mathrm{Frob}_{w_j}))
       =\prod_i(T-\alpha_i)
\]
by the unramified local-global comparison with the attached Galois
representation.  The symmetric Bernstein operator \(e_k\) acts on
the spherical line by \(e_k(\alpha)\); Vieta's formula identifies
this with the specialized right hand side of (ON8).
If the specialized spherical vector is zero both sides are zero.
Separation by all these specializations proves (ON8) integrally.
The argument needs precisely the separating arithmetic
specializations, attached determinant, and this unramified
coefficient comparison; it does not need a theorem lifting
individually polarized residual constituents.

Let \(E_j:V\to V\) be the unnormalized local Iwahori-to-spherical
trace, and put \(E=\prod_jE_j\).
The operators at distinct places commute and \(E(V)\subset M\).
Write \(\operatorname{tr}:V\to M\) for this product trace under the
base identification, so \(\iota\operatorname{tr}=E\).
The module \(\mathcal M\) has no nonzero integer torsion: for example
this follows from its separating complex specializations, or from
its finite free weight-algebra fixed levels.
For \(v\in V\) spherical at \(w_j\), the sandwich identity in (ON3)
gives
\[
 E_j f(U_j)v=\partial(N_{Q_j}f)(U_j)v.
 \tag{ON9}
\]
Indeed \(E_jv=[K_j:I_j]v\); first multiply the proposed identity by
this positive integer and use (ON3), then cancel the integer using
torsion freeness.  Thus no assertion that \([K_j:I_j]\) is a unit
is hidden in (ON9).
The product trace is \(\beta\)-semilinear, and, by (ON8) and
symmetric commutation, it intertwines every symmetric polynomial
with its image in the coefficients of \(P_j\).

\begin{srproposition}[The two concrete maps]
\label{prop:concrete-two-maps}
Let
\[
 S=B[u_{j,i}:j,i]\big/
 \left(\hbox{coefficients of }
          \prod_i(T-u_{j,i})-P_j(T)\ \hbox{for every }j\right)
\]
be the simultaneous ordered splitting algebra, and let
\(\overline f\) denote coefficient change by \(\beta\) followed by
the quotient map to \(S\).
There are additive maps, intertwining the polynomial actions,
\[
 \Phi:S\otimes_BM\longrightarrow V,\qquad
 \Psi:V\longrightarrow S\otimes_BM,
\]
such that, for
\[
 N_{\mathcal Q}(u)=
     \prod_{j=1}^q\prod_{i<i'}(u_{j,i'}-Q_ju_{j,i}),
\]
\[
 \Phi\Psi=N_{\mathcal Q}(U),\qquad
 \Psi\Phi=\overline{N_{\mathcal Q}}.
 \tag{ON10}
\]
They may be chosen so that the two composites are instead
multiplication by
\[
 f_{\mathcal Q}(u)=
 \prod_j\prod_i\prod_{i'\ne i}
       (u_{j,i}-u_{j,i'})(u_{j,i}-Q_ju_{j,i'}).
 \tag{ON11}
\]
These are the two maps appearing in the finite-level ordinary
control relation.
\end{srproposition}

\begin{proof}
Define \(\Phi\) first on polynomials by
\[
 \Phi\!\left(\overline{\sum_\nu b_\nu u^\nu}\otimes m\right)
       =\sum_\nu U^\nu\iota(b_\nu m).
 \tag{ON12}
\]
The coefficient relations defining \(S\) kill this expression by
(ON8) and Vieta's formula.  Thus (ON12) descends to the quotient
and the tensor product.  It satisfies
\(\Phi(1\otimes m)=\iota(m)\) and
\(\Phi(\overline h\,z)=h(U)\Phi(z)\).
This definition is meaningful even when \(\beta\) is not surjective.

For a tuple \(\mathbf H=(H_1,\ldots,H_q)\) of subsets of
\(\mathcal R\), put
\[
 A_{\mathbf H}=\prod_j a_{H_j}(u_j),\qquad
 C_{\mathbf H}=\prod_j c_{H_j}(u_j).
\]
The reverse map is the finite sum
\[
 \Psi(v)=
 \sum_{\mathbf H}
    \overline{A_{\mathbf H}}\otimes
       \operatorname{tr}\bigl(C_{\mathbf H}(U)v\bigr).
 \tag{ON13}
\]
We verify its polynomial equivariance, since this is essential for
the second composite.  In
\(\mathcal P\otimes_{\mathcal P^{S_n}}\mathcal P\), (ON7) gives
\[
 \sum_Ha_H\otimes c_Hx=\sum_Hxa_H\otimes c_H .
 \tag{ON14}
\]
Indeed expand \(c_Hx\) by the second formula of (ON7), transfer
the symmetric coefficient
\(\partial(a_Kc_Hx)\) across the tensor sign, and use the first
formula of (ON7) to sum over \(H\).
After taking the product of (ON14) over the auxiliary places,
symmetric balance is precisely the symmetric intertwining property
of \(\operatorname{tr}\) and (ON8).
Consequently
\(\Psi(h(U)v)=\overline h\,\Psi(v)\).

Applying (ON12) to (ON13) and using
\(\iota\operatorname{tr}=\prod_jE_j\) gives
\[
 \Phi\Psi
 =\sum_{\mathbf H}A_{\mathbf H}(U)
            \left(\prod_jE_j\right)C_{\mathbf H}(U)
 =\prod_jN_{Q_j}(U_j),
\]
where (ON7) and commutation between distinct places give the final
equality.
For a spherical vector, repeated use of (ON9) yields
\[
 \operatorname{tr}\bigl(C_{\mathbf H}(U)\iota(m)\bigr)
  =\left(\prod_j
        \partial(N_{Q_j}c_{H_j})\right)(u)\,m.
\]
The factor on the right is symmetric at every place and hence is a
well-defined \(B\)-coefficient under (ON8).
The first equality of (ON7), applied to \(f=N_{Q_j}\), now gives
\[
 \Psi(\iota(m))=\overline{N_{\mathcal Q}}(1\otimes m).
\]
Since \(S\otimes_BM\) is generated from these elements by the
polynomial action, polynomial equivariance proves the second
identity in (ON10).

Finally \(N_{\mathcal Q}\) divides \(f_{\mathcal Q}\) over the integral
polynomial ring: each factor \(u_{j,i'}-Q_ju_{j,i}\), \(i<i'\),
occurs among the ordered factors in (ON11).
Write \(f_{\mathcal Q}=N_{\mathcal Q}c\).
Keep \(\Phi\) and replace \(\Psi\) by \(\Psi\circ c(U)\).
Equivariance and (ON10) make both composites multiplication by
\(f_{\mathcal Q}\), as claimed.
\end{proof}

\paragraph{What this proves and what still comes from arithmetic.}
The maps (ON12)--(ON13) are now constructed, with exactly the
discriminant used by the auxiliary support argument.  Together with
the finite-projective modules and norm control constructed in the
preceding subsection, and the uniform auxiliary nilpotence proof,
this supplies the module-theoretic finite-level control mechanism.
It does not supply an automorphic seed or construct attached Galois
representations from nothing.
The only remaining arithmetic realization premise in this particular
old/new step is the explicitly described classical separation and
coefficient comparison producing (ON8).
In particular the presence of a two-map interface in the formal
source is not being substituted for its proof.

\paragraph{Exact source locations.}
Use the frozen tag \srcid{alpoge-source-7cbf7a0c}, under
\srcid{Development.AutomorphicArithmetic.AutomorphicVarieties}.
Definitions of the splitting algebra and of the exact two-map
conclusion are
\srcid{AlgebraTools.MSplitDefs} and
\srcid{PatchingModules.NilpotentLevelLowering.LevelModule.TwoMaps}.
The actual normalized local Bernstein operators, their faithfulness,
the trace identities, and the divided-difference formulas are
\srcid{AlgebraTools.IwBernHk.bernHk},
\srcid{IwBernFaith.bernFaith},
\srcid{HkSandwich.hkSandwich},
\srcid{HkSum.hkSum},
\srcid{HkDualEx.hkDualEx}, and
\srcid{HkDualExp.hkDualExp}.
The passage from actual local group actions to these identities,
including cancellation of the positive integer index, is the proof
\srcid{PatchingModules.NilpotentLevelLowering.HkHypOf.locA} and
\srcid{hkHypOf}.
The maps and their two composites are constructed in
\srcid{MultiPhi.multiPhi}, \srcid{MultiPsi.multiPsi},
\srcid{MultiTwo.multiTwo}, and \srcid{TwoMul.twoMul}.
Their spherical inclusion and product trace in the global application
are written explicitly in
\srcid{ModularityArgument.ThresholdArguments.DenseThresholdCases.LevTwoOf.two\_gen}.
The arithmetic coefficient comparison (ON8) is the proof
\srcid{DenseThresholdCases.SymOfVieta.symOf\_alt}:
its specialization comparison is \srcid{symOf\_sym}, its coefficient
calculation is \srcid{symOf\_vieta}, and nonzero spherical specialization
is \srcid{symOf\_sph}.
All files are below
\url{https://github.com/ChristianKudeba/alpoge-rh-verification/tree/alpoge-source-7cbf7a0c/src/Development/AutomorphicArithmetic/AutomorphicVarieties}.

\subsection{The uniform auxiliary exponent and the finite-level module}
\label{sec:uniform-aux-exponent}

The integer exponent used before patching can be chosen to be
$n!$.  In rank three it is therefore $6$, chosen before the
host, the coefficient prime, the auxiliary level, and the
number of auxiliary places.

\begin{srlemma}[A polynomial vanishing on the torus support]
\label{lem:torus-uniform-nil}
Let a commutative group $T$ act by invertible commuting operators
on a complex vector space $W$ of dimension at most $d$.
Let $t_1,\ldots,t_r\in T$ and $F\in\mathbb C[X_1,\ldots,X_r]$.
If
\[
            F(\chi(t_1),\ldots,\chi(t_r))=0
\]
for every character $\chi$ occurring in a simultaneous
generalized-character space of $W$, then
$F(t_1,\ldots,t_r)^d=0$ on $W$.
The same exponent works on any direct sum of copies of this
$T$-module.
\end{srlemma}
\begin{proof}
A commuting family over an algebraically closed field has a
common eigenvector.  Apply this observation successively to
quotients to obtain a complete invariant flag.
The diagonal functions on $T$ are multiplicative characters,
since every group element is invertible.
The hypothesis says that the polynomial in the operators has
zero diagonal in this flag.  It is strictly upper triangular,
so its $d$-th power is zero.
On a direct sum the same identity holds on each summand.
\end{proof}

We state explicitly the local representation-theoretic input
used to take $d=n!$.
For a nonarchimedean local field $K$, a subgroup
$C\subseteq\kappa_K^\times$, and an irreducible smooth complex
representation $\sigma$ of $\operatorname{GL}_n(K)$, the space
$\sigma^{\operatorname{Iw}_C}$ is finite dimensional and
\[
                   \dim_\mathbb C
                      \sigma^{\operatorname{Iw}_C}\leq n!.
                                                               \tag{35}
\]
This is the exact bound \srcid{AlgebraTools.IwJqLine.jqLine},
using the printed geometric lemma and Bernstein--Zelevinsky
finiteness hypotheses.
Its mechanism is worth recording.
The canonical map from these invariants to the Borel Jacquet
module is injective: a vector in its kernel is killed by a
sufficiently positive Hecke operator, whereas the torus action
extending the positive Hecke action makes that operator
invertible.  A nonzero simultaneous torus weight, together
with Frobenius reciprocity, embeds $\sigma$ in a principal
series.  The Borel Jacquet module of this principal series has
a filtration with one one-dimensional quotient for each of
the $n!$ Bruhat cells, by the geometric lemma.
Exactness of the Jacquet functor in characteristic zero bounds
the Jacquet module of $\sigma$ by $n!$.
The injective canonical map gives (35).
If the invariant space is zero, (35) is immediate.
Thus this input has no dependence on the residue cardinality
or on $C$.

For an isotypic sum of copies of $\sigma$, invariants under the
compact group $\operatorname{Iw}_C$ are spanned by images of
$\sigma^{\operatorname{Iw}_C}$ under intertwining maps.
Indeed any finite expression can be averaged over this compact
group; smoothness reduces this averaging to a finite sum.
The torus action commutes with these intertwining maps.
Lemma~\ref{lem:torus-uniform-nil} and (35) therefore give the
same exponent $n!$ on an isotypic sum, regardless of its
multiplicity.  Restriction to invariants at other places
does not increase the required exponent.
This is the actual uniform-exponent argument in
\srcid{NilpotentLevelLowering.NilKill}.

\paragraph{The two auxiliary polynomial relations.}
At one auxiliary place write $Q$ for the residue cardinality
and $P,T$ for Frobenius and tame inertia matrices.
The characteristic-zero local data supplies a semisimplification
by characters $c_i$, and
\[
 P^{-1}TP=T^Q,\qquad
 \det(X-R(g))=\prod_i(X-c_i(g)).
\]
On the locus
\[
         \alpha_i=c_i(\mathrm{Frob}),\quad
         \alpha_i\ne\alpha_j,\quad
         \alpha_i\ne Q\alpha_j\quad(i\ne j),                    \tag{36}
\]
the following matrix identities hold:
\[
 \prod_i(P-\alpha_i)=0,\qquad
 \prod_{k\ne i}(P-\alpha_k)
                    (T-c_i(\tau))=0.                          \tag{37}
\]
Here the second product may be written in any order.

We include the matrix argument.
The character relation gives $c_i(\tau)^{Q-1}=1$; hence
$T^{Q-1}$ is unipotent.
Set $N=\log(T^{Q-1})$, using its finite nilpotent logarithm.
Then $P^{-1}NP=QN$.
The eigenvalues of $\operatorname{Ad}(P^{-1})$ are the ratios
$\alpha_j/\alpha_i$.  Condition (36), and $Q\ne1$, exclude $Q$,
so $N=0$.  Thus $T^{Q-1}=1$, and the tame relation gives $PT=TP$.
The eigenspaces of $P$ are one dimensional.
The character identities for the traces of $P^rT$,
$0\leq r<n$, and the invertible Vandermonde matrix in the
distinct $\alpha_i$, show that $T$ acts on the $i$-th
eigenspace by $c_i(\tau)$.  This proves (37).
This is the content of the finite matrix argument
\srcid{LocData.rel\_of\_chars}, also used by
\srcid{TLocDataT.locDataT\_alt} with the printed local--global
compatibility statement.

For all auxiliary places introduce commuting Hecke variables
$U_{ji}$ and $Y_{ji}$ and put
\[
 f_N=\prod_j\prod_i\prod_{i'\ne i}
            (U_{ji}-U_{ji'})(U_{ji}-Q_{Nj}U_{ji'}).
                                                               \tag{38}
\]
Let $r$ be either trace of a group element times the first
polynomial in (37), or trace of a group element times
\[
       \prod_{k\ne i}(P-U_{jk})(T-1-Y_{ji}).                    \tag{39}
\]
At every supported complex character of the relevant local
torus, the polynomial $f_Nr$ vanishes.
If $f_N$ vanishes this is immediate; otherwise its factors
give (36) and (37) gives the assertion.
The torus-support lemma proves
\[
                       (f_Nr)^{\,n!}M_N=0.                   \tag{40}
\]
All additional factors in (38), at the other auxiliary places,
commute with the operator at this place, so they do not
multiply the exponent by the number of places.

The passage back to the integral ordinary module uses the
actual separating family of classical automorphic
specializations: an element is zero if all its classical
images are zero.  Each specialization sends the auxiliary
operators to the normalized torus operators
\[
       U_{ji}\longmapsto Q_{Nj}^{\,i}\Theta_j(\varpi_i),
       \qquad 1+Y_{ji}\longmapsto\Theta_j(\delta_i).
\]
Their polynomial evaluations intertwine the integral and
complex actions.  Applying the already uniform identity
at every specialization and then separation proves (40)
integrally, with the same exponent $n!$.
This is \srcid{LevRelOf.relOf\_alt} and
\srcid{LocOkLevelModuleOfThetaData.level\_relation\_one\_two}.
In particular the exponent is fixed before $N$, before the
chosen relation, and before the element of $M_N$.

\begin{srlemma}[Projective norm and the augmentation variables]
\label{lem:projective-norm-control}
Let $(\Lambda,\mathfrak m)$ be a complete Noetherian local
ring with residue characteristic $\ell$ and
$\ell\in\mathfrak m$.  Let $\Delta$ be a finite $\ell$-group
and $M$ a finite projective $\Lambda[\Delta]$-module.
The norm
\[
       \mathcal N:M\longrightarrow M^\Delta,\qquad
                         x\longmapsto\sum_{\delta\in\Delta}\delta x
\]
is onto, with kernel the augmentation submodule
$\sum_{\delta\in\Delta}(\delta-1)M$.
If $\Delta$ is abelian, $\Lambda[\Delta]$ is local and
$M$ is finite free over it.
\end{srlemma}
\begin{proof}
For the regular module $\Lambda[\Delta]$, the invariant
elements are exactly the scalar multiples of
$\sum_\delta\delta$.  Its norm is multiplication of this
element by the augmentation.  Its kernel is therefore the
augmentation ideal, which is generated by $\delta-1$.
The same assertions hold on finite free modules.
A projective module is a direct summand of such a module;
norm, invariants, and augmentation submodules respect this
decomposition, proving the norm assertion.
For abelian $\Delta$, the reduction
$(\Lambda/\mathfrak m)[\Delta]$ is local:
its augmentation ideal is nilpotent and its quotient is the
residue field.  Thus the finite complete $\Lambda$-algebra
$\Lambda[\Delta]$ is local, with maximal ideal generated by
$\mathfrak m$ and $\delta-1$.
A finite projective module over a local ring is free.
\end{proof}

This proves the trace/control assertion at the auxiliary
level without dividing by $\#\Delta$, which is usually not a
unit.  For a chosen generating family $\delta_p$ one can use
only the operators $\delta_p-1$ in the kernel formula:
$\delta\delta'-1=\delta(\delta'-1)+(\delta-1)$.

\begin{srlemma}[Fixed truncations of the auxiliary group ring]
\label{lem:aux-trunc-free}
Let
\[
 \Delta_N=\prod_{j,i}\mathbb Z/\ell^{a_{Nj}}\mathbb Z,
                     \qquad a_{Nj}\geq N.
\]
Put $S=\Lambda[[Y_{ji}]]$, and let $M_N$ be finite projective
over $\Lambda[\Delta_N]$.  Identify $1+Y_{ji}$ with the
corresponding group generator.  Then for every fixed $k,c$,
and all sufficiently large $N$,
\[
 M_N/(Y^k+\mathfrak m^c)M_N
       \text{ is finite free over }
                      S/(Y^k+\mathfrak m^c).
                                                               \tag{41}
\]
Its rank may depend on $N$.
\end{srlemma}
\begin{proof}
There is an isomorphism
\[
 \Lambda[\Delta_N]\simeq
 S/\bigl((1+Y_{ji})^{\ell^{a_{Nj}}}-1\bigr)_{j,i}.
\]
The polynomials are monic and distinguished, so successive
power-series division proves the isomorphism as well as
finiteness and completeness.
For $1\leq r<k$,
\[
 v_\ell\binom{\ell^a}{r}\geq a-v_\ell(r),
\]
because $r\binom{\ell^a}{r}=\ell^a
\binom{\ell^a-1}{r-1}$.
Thus if $N\geq c+\max_{1\leq r<k}v_\ell(r)$, every displayed
relation belongs to $Y^k+\mathfrak m^c$.
Terms of degree at least $k$ already belong to $Y^k$.
The preceding lemma makes $M_N$ free over its group ring;
taking the specified quotient proves (41).
\end{proof}

\paragraph{Matching the common-module hypotheses.}
The module in (41) is the module entering the common
ultraproduct construction.  At $Y=0$, the projective norm
lemma gives its lower-level invariant module.
The two old/new-level maps, after the auxiliary root-splitting
algebra is introduced, have both composites equal to the
single discriminator $f_N$.  On the fixed base module these
maps use the same annihilator ideal $J$; their scalar
extensions do not introduce a new automorphic support ideal.
The two maps and their discriminator are the
\srcid{TwoMaps} input of \srcid{OpToTW.opCore}, assembled by
\srcid{LevTwoOf.twoOf\_alt}.

The source construction then takes the simultaneous
topologically nilpotent summand for $U_{ji}-u_{ji}$ and
$Y_{ji}$.  The complementary idempotent decomposition
commutes with all Hecke operators.  It preserves the
projective norm calculation and the identities (40).
Passing to the root-splitting algebra and completing sends
$Y$ to zero on the base module, and the two control maps
remain linear for this common specialization.
These are the operations \srcid{nilPart},
\srcid{freePart}, \srcid{quotPart}, and \srcid{spPart}
inside \srcid{NilpotentLevelLowering.OpToTW}.

For the fixed limit auxiliary elements, their finite-level
trace relations converge to the relations defining the
common ideal $I$.
Also $Q_{Nj}\equiv1\pmod{\ell^N}$, so (38) converges to the
fixed discriminant $f$.
Consequently (40), with the fixed exponent $e=n!$, says
exactly: for every generator $r$ of $I$ and every precision
$c$, for all sufficiently large $N$, $(fr)^e$ is congruent
modulo $\mathfrak m_R^c$ to an element annihilating $M_N$.
This is the uniform approximate annihilator condition of
the common patched-module theorem.

The exact assembly in \srcid{THasAutOf} therefore chooses
\srcid{TWMaps} with exponent \srcid{n.factorial}.
For rank three, $e=6$ throughout.  The free truncations
(41), the norm/control statement, and the uniform exponent
(40) have all been proved here.
The ordinary function-module construction, projectivity and old/new
maps are supplied earlier in Section~\ref{sec:sr-ordinary-control}.
For the actual arithmetic specialization one must still construct
the attached Galois determinant and verify its classical separation,
local--global compatibility, and the same universal ring and limit
data. These requirements are not consequences of the existence of
an abstract deformation ring.

\subsection{The cohomological passage to the Taylor--Wiles limit}
\label{sec:coh-at-tw-limit}

This section proves the limit conversion \srcid{CohLim}.
It is separate from the arithmetic estimates producing the
finite-level hypotheses \srcid{B1} and \srcid{B2}.

Keep $A_0=k[[T]]$, $F=k((T))$, $q=\#k$, the paired point $\rho$,
and the chosen odd classes $\kappa_1,\ldots,\kappa_g$ on the
\emph{original} group $\Gamma_0$.
Their original-group generating property is an explicit premise:
for every integral odd Selmer cocycle $z_0$ on $\Gamma_0$ there
are $s\geq0$ and $a_j\in A_0$ such that
$T^sz_0-\sum_j a_j\kappa_j$ is an integral odd coboundary.
This is exactly \srcid{Kap.gen}; independence is
\srcid{Kap.indep}.  It is not a spanning assumption on
$\Gamma_\infty$, which is what the proof establishes.
For $m\geq1$ write $A_m=A_0/T^m$.
Use the even, odd, first-order and second-order Selmer conditions
specified above, including their local correcting matrices.
Let $\Gamma_N\to\Gamma_0$ be the auxiliary groups and
$\Gamma_\infty\to\Gamma_0$ the group denoted by
$\Lambda.G$ in \srcid{Lim}.
The limit data include the following facts; these are the
specific convergence properties used in the proof.
\begin{enumerate}
\item A topologically finitely generated profinite group $P$
has continuous surjections $\pi_N:P\to\Gamma_N$ and
$\pi_\infty:P\to\Gamma_\infty$.
They preserve the sign characters and a chosen complex
conjugation.  The index-two sign kernel $P^+$ is therefore
also topologically finitely generated.
\item For every open normal $V\subseteq P$, eventually along
the chosen nonprincipal ultrafilter,
\[
       (\ker\pi_N)V=(\ker\pi_\infty)V.
                                                               \tag{30}
\]
The maps $f_N:P\to\Gamma_0$ and $f_\infty:P\to\Gamma_0$
obtained by composition agree modulo every fixed finite
quotient of $\Gamma_0$, eventually and uniformly on $P$.
\item The local maps, the prescribed unramified local maps,
and the auxiliary elements $\phi_{Nj},\tau_{Nj}$ converge
in the same finite-quotient sense on $P$.
More precisely their values have lifts to $P$ which agree
modulo every prescribed open normal $V$ eventually,
uniformly over the local group.  This is the relation
\srcid{Close} and the assertions
$\Lambda.\mathrm{phi\_lim}$,
$\Lambda.\mathrm{phiu\_lim}$,
$\Lambda.\mathrm{ph\_lim}$,
$\Lambda.\mathrm{ta\_lim}$.
\item After the \srcid{ProL} quotient, each
$\ker(\Gamma_N\to\Gamma_0)$ is contained in the closed normal
subgroup generated by the $\tau_{Nj}$.
At the limit
\[
 h_\infty(\tau_j)=1,\quad h_\infty(\phi_j)=s_j,\quad
                         \phi_j\tau_j=\tau_j\phi_j.             \tag{31}
\]
The matrix $\rho(s_j)$ has $n$ distinct eigenvalues in $F$.
\end{enumerate}
All these are statements about the existing limit object.
For example (30) is $\Lambda.\mathrm{ker\_lim}$, and (31)
is $\Lambda.\mathrm{p\_ta}$,
$\Lambda.\mathrm{p\_ph}$ and $\Lambda.\mathrm{comm}$.

The exact finite-level inputs, with constants chosen before $m,N$,
are these:
\[
 \begin{split}
 \#\operatorname{Sel}_{1,N}(A_m)
 &\leq
 \#\operatorname{Bd}_{1,N}(A_m)\,q^{\,mD+C_1},\\
 &\hspace{5mm}\operatorname{Sel}_{1,N}(A_m)\text{ is finite},
 \end{split}                                                    \tag{32}
\]
eventually in $N$ for every fixed $m$; and, for every $g+1$
tuples of even Selmer cocycles $z_{ij}$ over $A_m$,
\[
 \begin{gathered}
 \exists\,c_1,\ldots,c_{g+1}\in A_0
       \text{ with at least one }c_i\in A_0^\times,\quad
 \exists\,\Xi\in\operatorname{Sel}_{2,N}(A_m),\\
       \varepsilon_j'(\Xi)=T^{C_2}\sum_i c_i z_{ij}
                             \quad(1\leq j\leq g).
 \end{gathered}                                                   \tag{33}
\]
For each fixed $m$, (33) holds for all such tuples at all
sufficiently large $N$.
These are \srcid{TWGrp.B1} and \srcid{TWGrp.B2}.
The uniform exponent $C_2$ and the primitive coefficient vector
in (33) are essential.

\begin{srproposition}[The three limit bounds]
\label{prop:tw-limit-coh}
Under these hypotheses the integral first-order Selmer rank
is at most $D$; every $g+1$ tuples of integral even cocycles
have a nonzero $A_0$-linear combination lifting to an integral
second-order cocycle; and every integral odd Selmer cocycle
on $\Gamma_\infty$ is an $A_0$-linear combination of the
$\kappa_j$ modulo a coboundary after multiplication by some
power of $T$.
These are exactly $K1(h_\infty,\kappa,D)$,
$K2(h_\infty,\kappa)$ and $K3(h_\infty,\kappa)$.
They imply the three Laurent hypotheses of
Theorem~\ref{thm:selmer-paired-pres}.
\end{srproposition}

\begin{proof}
\emph{Finite coefficient Selmer spaces transfer to the limit.}
For a fixed finite coefficient module and a fixed continuous
action, the continuous cocycles on a topologically finitely
generated profinite group form a finite set.
Indeed a cocycle is determined by its values on finitely many
topological generators: its values on the abstract subgroup
they generate follow from the cocycle identity, and continuity
gives the rest.
Equivalently use the homomorphism into the finite semidirect
product of the module and its action group.

A locally constant function $c:P^+\to Z$ factors through a
surjection $\pi:P^+\to H$ precisely when $\ker\pi$ lies in its
right stabilizer
\[
             \operatorname{Stab}(c)
             =\{u:c(xu)=c(x)\ \forall x\}.
\]
This stabilizer is open, since $P^+$ is compact.
Its normal core, viewed in $P$, is open normal.
Thus (30) says that $c$ factors through $\pi_N$ precisely
when it factors through $\pi_\infty$, eventually.
There are only finitely many cocycles on $P^+$ for the chosen
finite action, so one may impose this assertion for all of
them at once.
The actions themselves agree for large $N$ by the finite
quotient convergence of $f_N$.
The local convergence property in item 3 makes the descended
functions agree on all specified local and unramified groups.
The same correcting matrix therefore verifies the ordinary
condition on both groups.
Polarization is preserved by the fixed conjugation.
Consequently the finite coefficient Selmer sets on
$\Gamma_N$ and $\Gamma_\infty$ are identified by their pullbacks
to $P^+$, for sufficiently large $N$.
The coboundary sets are identified as well.
This proves (32)--(33) on $\Gamma_\infty$, with the same
constants $C_1,C_2$.

\emph{A fixed loss of precision for independent classes.}
Let $z_1,\ldots,z_e$ be integral first-order cocycles independent
over $A_0$ modulo integral first-order coboundaries.
There is $C\geq0$ such that
\[
 \sum_i a_i z_i\text{ is a coboundary modulo }T^m
       \quad\Longrightarrow\quad
                         a_i\in T^{m-C}A_0\quad(m\geq C).       \tag{34}
\]
Here is a finite-module proof of this assertion.
Let $V$ be the finite free $A_0$-module of allowed
skew-transpose first-order correcting matrices.
On $M=A_0^e\oplus V$ consider the linear operators
\[
        L_\gamma(a,Y)=\sum_i a_i z_i(\gamma)-dY(\gamma).
\]
Their span in the finite module
$\operatorname{Hom}_{A_0}(M,\operatorname{Mat}_n(\mathbb D_g(A_0)))$
is finitely generated.  Choose finite generators and assemble
them into one linear map $L:M\to M'$ with $M'$ finite free.
The common kernel is the kernel of all $L_\gamma$.
Independence of the classes says that each coordinate
functional $a\mapsto a_i$ vanishes on this kernel.
Smith normal form over the DVR $A_0$ supplies $C$ with
\[
       L(x)\in T^mM'
                 \ \Longrightarrow\
       x\in\ker L+T^{m-C}M.
\]
A coboundary matrix modulo $T^m$ lifts to $V$:
lift its upper-triangular entries and reflect them with the
required minus sign, preserving the first-order parity.
Applying the last implication to this lift proves (34).
This is the \srcid{CoefficientsBoundedOfSelmerBoundary}
argument, expressed over a DVR instead of its more general
Artin--Rees formulation.

Choose representatives in $A_0$ for $A_0/T^{m-C}$.
Their $e$-tuples give distinct classes modulo the first-order
coboundaries by (34).  Hence
\[
        q^{(m-C)e}
       \leq
       \frac{\#\operatorname{Sel}_{1,\infty}(A_m)}
            {\#\operatorname{Bd}_{1,\infty}(A_m)}
       \leq q^{\,mD+C_1}.
\]
Letting the integer $m$ grow proves $e\leq D$.  This is $K1$.

\emph{Second-order lifting survives with its local conditions.}
For integral tuples $z_{ij}$, (33) gives primitive coefficient
vectors $c_m$ and second-order Selmer cocycles $\Xi_m$ modulo
$T^m$.  The ring $A_0$ is compact.
Moreover the set of second-order Selmer cocycles modulo
$T^j$ is finite for every fixed $j$, by the finite-cocycle
argument above and finite generation of $\Gamma_\infty$.

Take an ultralimit of $c_m$ and of all the reductions
$\Xi_m\bmod T^j$.  The latter reductions stabilize as whole
functions, not merely pointwise, because they belong to a
finite set of continuous cocycles.
They are compatible and define an integral continuous
cocycle $\Xi$.
Its coefficient, cocycle, polarization and unramified identities
hold at every precision and hence integrally.

The local correcting matrices also survive.
At each local place, lift the correcting matrices for
$\Xi_m$ arbitrarily to the finite free $A_0$-module of all
matrices and take their ultralimit in this compact module.
Parity and the equations for zero lower entries and the
specified zero diagonal entries are closed conditions.
Thus the limit is an allowed correcting matrix for $\Xi$.
This proves the ordinary condition without discarding it
when passing to the limit.

Pass to an ultrafilter set on which one fixed coordinate of
$c_m$ is a unit; one exists because there are only $g+1$ coordinates.
The limiting vector $c$ still has a unit coordinate.
The defining equations become
\[
                    \varepsilon_j'(\Xi)
                        =T^{C_2}\sum_i c_i z_{ij}.
\]
The vector $T^{C_2}c$ is nonzero since $A_0$ is a domain.
This proves $K2$.

\emph{The kernel is generated by the auxiliary inertia.}
Let $K$ be the closed normal subgroup of $\Gamma_\infty$
generated by all $\tau_j$.
We claim $\ker h_\infty\subseteq K$.
It suffices to test a finite quotient
$Q=\Gamma_\infty/W$ with $K\subseteq W$ open normal.
Let $V=\pi_\infty^{-1}(W)$.
Equation (30) makes $\ker\pi_N\subseteq V$ eventually, giving
$\Gamma_N\to Q$ through the common cover.
The convergence of the $\tau_{Nj}$ makes their images in $Q$
trivial.
Item 4 therefore makes this homomorphism descend to a
homomorphism $\theta_N:\Gamma_0\to Q$.

There are only finitely many continuous homomorphisms
$\Gamma_0\to Q$, because $\Gamma_0$ is a quotient of the
topologically finitely generated group $P$.
Take the intersection $U$ of their kernels.
Finite quotient convergence of $f_N$ to $f_\infty$ makes them
agree modulo $U$ eventually, so
$\theta_Nf_N=\theta_Nf_\infty$.
But the first of these is precisely
$P\to\Gamma_\infty\to Q$.
For $x\in\ker h_\infty$, choose a lift in $P$; its image under
$f_\infty$ is $1$, and therefore its image in $Q$ is $1$.
This proves the claim.  It is the substance of
\srcid{KerOo.ker\_oo}.

\emph{Odd classes descend.}
Let $z$ be an integral odd Selmer cocycle on $\Gamma_\infty$.
The cocycle identities for $\phi_j\tau_j=\tau_j\phi_j$, and
$h_\infty(\tau_j)=1$, imply
\[
                       \rho(s_j)z(\tau_j)
                                  =z(\tau_j)\rho(s_j).
\]
The matrix $\rho(s_j)$ is even and has distinct eigenvalues
over $F$, while $z(\tau_j)$ is odd.
The centralizer of a matrix with distinct eigenvalues is
commutative.  Its even block idempotent commutes with
$\rho(s_j)$, hence also with $z(\tau_j)$.
An odd matrix commuting with this block idempotent is zero.
Thus $z(\tau_j)=0$.

The subgroup on which both $h_\infty$ and $z$ vanish is closed
and normal.  Normality follows from the cocycle identity;
its stability under complex conjugation follows from
$z(\overline g)=-z(g)^{\mathsf t}$.
The kernel-generation claim therefore shows that $z$ vanishes
on $\ker h_\infty$, and $z$ descends to $\Gamma_0$.
Its continuity and all Selmer conditions descend:
the group map is a continuous quotient, the local maps
commute with it, and the same local correcting matrices apply.
The generating property of the original $\kappa_j$ now gives
$T^s z-\sum_j a_j\kappa_j$ a coboundary for some $s\geq0$ and
$a_j\in A_0$.  This proves $K3$.

Finally these integral conclusions imply the three Laurent
inputs.  For a finite family of bounded Laurent cocycles,
clear one common power of $T$ in the functions and in their
local correcting matrices; the distinguished local set is
finite in the arithmetic application.
This gives integral Selmer cocycles.
Apply $K1$ to an independent integral family, $K2$ to the
scaled even tuples, and $K3$ to a scaled odd cocycle.
After inverting $T$, their conclusions are exactly the
Laurent rank bound, nonzero second-order combination, and
spanning by $\kappa_j$.
No dimension or local condition is changed by this scalar
normalization.
\end{proof}

\paragraph{Sources and the remaining arithmetic input.}
Finite coefficient transfer is
\srcid{LcOfAdic.t1}, with
\srcid{EventuallyFactorThroughTowerIffLimit} and
\srcid{EventuallyFactorizationCompatibleWithLocalMaps}.
The first-order argument is
\srcid{FirstCohomologyBoundOfBoundaryBound}.
The second-order argument is
\srcid{SecondCohomologyBoundOfBoundaryBound},
\srcid{SelmerSecondExistsUltrafilterLimit}, and
\srcid{SelmerSecondClosedUnderLimits}.
Odd descent is \srcid{KoszulThreeOfDet}, \srcid{OdComm},
\srcid{CocKer}, and \srcid{KerOo}.
The assembly is \srcid{AutomorphicDetection.CohLim.cohLim}.
The modules are under
\srcid{Development.AutomorphicArithmetic.AutomorphicVarieties.%
PatchingModules} at the frozen source tag.
Thus $K1,K2,K3$ require no additional theorem once the actual
limit object, \srcid{ProL}, distinct auxiliary eigenvalues,
and the uniform finite-level bounds (32)--(33) have been
produced.  Their arithmetic production is not asserted by
the present limit argument.

\section{The Case II rank-four support calculation}
\label{sec:sr-rank-four-support}
The specialization in Section~\ref{sec:sr-character-integers} eliminates
the other rank-four cases. We give the remaining core calculation
from its stated tower, common-space and modular-support inputs.
The existence of the linked automorphic seed and the finite-level
automorphic modules remains explicit.

\subsection{The rank-four characteristic-\(\ell\) support calculation}
\label{subsec:sr-rank-four-core}
This subsection supplies the rank-four inputs to the common-space
completion and component-propagation lemmas. It is restricted to the
actual Case II tower and requires no Case III construction.

Write \(F\) for its totally real abelian field,
\(d=[F:\mathbb Q]\), \(s=\#\{v\mid\ell\}\), and
\(\delta=[F_v:\mathbb Q_\ell]\), so \(s\delta=d\).
Let \(\Lambda\) be a finite coefficient DVR with residue field \(k\)
of odd characteristic \(\ell\).
The prescribed residual characters are \(\chi_0,\ldots,\chi_3\),
with \(\chi_i\chi_{3-i}=\nu\), pairwise distinct at \(g_0\).
Representations have fixed multiplier \(\nu\), are diagonal at
\(g_0\), satisfy the prescribed global ramification condition,
and have ordinary flags represented in the lower open cell by
lower unipotent symplectic matrices \(n_v\).
Their ordinary characters are trivial on the maximal
\(\ell\)-power torsion subgroup of the local abelianization.
The tower makes the residual representation trivial on inertia
at \(\ell\); its multiplier is trivial there in characteristic
\(\ell\).

Let \(B\subset A\) be the common space: \(B\) is complete,
\(A\) is of finite type over \(B\), the universal matrix and flag
entries generate \(A\) over \(B\), and \(B\) contains all
characteristic-polynomial coefficients and ordinary diagonal
characters \(\kappa_{v,i}\).
The completion \(R\) at a residual point represents the normalized
deformation problem. The common-space construction and its completion/dimension transport
are retained as hypotheses: the field-transcendence bounds below must
give the indicated complete-local dimension bounds. Their verification
for the actual common space is part of its construction, and is not
inferred just from the existence of a universal formal ring.

The property of a modular ideal \(J\subset B\) needed below is
\begin{equation}\label{eq:sr-pm-finite}
 \mathfrak m_B^e\subset J+\mathfrak m_\Lambda B+
 (\kappa_{v,i}(t)-1:v\mid\ell,\ 0\le i\le3,\ t\in I_v)
 \quad\text{for some }e.
\end{equation}
We also require its actual curve propagation property:
if \(A\to k'[[T]]\), with \(k'\) finite, is centered over
\(\mathfrak m_B\), kills \(J\), has absolutely irreducible
generic representation, and has pairwise distinct ordinary
characters on
\(U_v=I_v\cap\ker\epsilon_\ell\cap\ker\bar r\),
then every prime of \(B\) through which this curve factors
contains \(J\). These are properties of an actual ideal; no
existence of such an ideal is asserted in this subsection.

\begin{srlemma}[Rank-four trace fiber]\label{lem:sr-four-trace-fiber}
At an absolutely irreducible field-valued point, the matrix and
flag entries are algebraic over the image field of \(B\) and
at most \(2+4s\) additional elements. Consequently, for every
irreducible generic prime \(\mathfrak q\) of the representing
completion,
\[
 \dim R/\mathfrak q\le
       \dim B/\mathfrak q_B+2+4s.
\]
\end{srlemma}
\begin{proof}
Pass to an algebraic closure of the point field and put
\(r(g_0)=\operatorname{diag}(d_0,d_1,d_2,d_3)\).
These distinct eigenvalues are algebraic over the image of \(B\).
Let \(E\) be an algebraically closed intermediate field containing
them, all traces, and all multiplier values.

The four identities
\[
 \operatorname{tr}(r(g_0)^m r(h))
       =\sum_i d_i^m r(h)_{ii}\qquad(0\le m\le3)
\]
and Vandermonde inversion put every diagonal entry in \(E\).
Apply the same inversion to the \(00\) entry of
\(r(u)r(g_0)^m r(w)\), and then twice to the \(00\) entry of
\(r(u)r(g_0)^m r(g)r(g_0)^n r(w)\). They give
\begin{equation}\label{eq:sr-four-triples}
 r(u)_{0i}r(w)_{i0}\in E,\qquad
 r(u)_{0i}r(g)_{ij}r(w)_{j0}\in E .
\end{equation}
The matrix-algebra spanning theorem supplies
\(a_i=r(u_i)_{0i}\ne0\) and \(b_i=r(w_i)_{i0}\ne0\).
Adjoin only \(a_1,a_2\); \(a_0\) is already in \(E\).

Our adjoint convention is
\(X^\#_{ij}=\sigma_i\sigma_jX_{3-j,3-i}\),
\(\sigma=(1,1,-1,-1)\).
The symplectic relation gives
\(-r(g)_{02}=\nu(g)r(g^{-1})_{13}\).
Thus \(c=r(u_2^{-1})_{13}=-a_2/\nu(u_2)\) is nonzero in \(E\).
Now \(a_1cb_3\in E\) gives \(b_3\in E\);
\(a_3b_3\in E\) gives \(a_3\in E\).
Then \(a_ib_i\in E\) gives every \(b_i\), and the triple
products give every entry of \(r(g)\).

Every lower unipotent symplectic flag matrix is
\[
 \begin{pmatrix}
 1&0&0&0\\ A&1&0&0\\ B&D&1&0\\
 C&B-DA&-A&1
 \end{pmatrix}.
\]
Four additional coordinates per place therefore suffice.
Coordinate generation and the common-space dimension-transport
lemma prove the result.
\end{proof}

\begin{srlemma}[Reducible tangent bound]
\label{lem:sr-four-reducible-tangent}
Suppose the tower cohomology bounds give
\[
 s\le C,\quad \dim Z_w\le C\ (w=0,\psi,-\psi),\quad
 \dim Z_w\le d+C\ (w\text{ odd}),\quad
 \dim S_w\le C\ (w\text{ a root}),
\]
where \(Z_w\) denotes the globally allowed normalized cocycles
and \(S_w\subset Z_w\) the subspace locally a coboundary.
At a good residual point every proper reducible support quotient
has tangent dimension at most \(d+20C\) and dimension at most
\(d+20C+1\).
\end{srlemma}
\begin{proof}
Put
\[
 \lambda=(1,0,1,0),\quad\mu=(1,1,0,0),\quad
 w_{ij}=(\lambda_i-\lambda_j,\mu_i-\mu_j),
\]
and \(\alpha=(1,0),\beta=(-1,1),\psi=(0,1),\theta=(1,1)\).
The odd roots are \(\pm\alpha,\pm\beta,\pm\theta\).
A support \(S\) means \(r(g)_{ij}=0\) when
\(i\ne j,w_{ij}\notin S\).
The nine maximal proper supports and the pivots needed below are
\[
\begin{array}{c|l|c|c}
&S&\text{split}&\text{partially split}\\\hline
1&\{\beta,-\beta,-\alpha,-\theta,-\psi\}&12&12\\
2&\{\beta,-\beta,\alpha,\theta,\psi\}&-&-\\
3&\{\theta,-\theta,\psi,\beta,-\alpha\}&-&12\\
4&\{\theta,-\theta,-\psi,-\beta,\alpha\}&-&-\\
5&\{\alpha,-\alpha,-\theta,-\psi,-\beta\}&01&01\\
6&\{\alpha,-\alpha,\theta,\psi,\beta\}&-&01\\
7&\{\psi,-\psi,\beta,-\alpha,-\theta\}&12&12\\
8&\{\psi,-\psi,-\beta,\alpha,\theta\}&-&-\\
9&\{\beta,-\beta,\theta,-\theta\}&-&03
\end{array}
\]
This list covers reducibility because an invariant subspace for
the regular diagonal \(r(g_0)\) is a coordinate subspace.
There are four line stabilizers, four coordinate Lagrangian
plane stabilizers, and one pair of complementary nondegenerate
coordinate planes; three-dimensional subspaces are handled by
their symplectic perpendicular lines.

The split alternative applies when
\[
 (-\psi\in S\text{ and }(\beta\in S\text{ or }-\alpha\in S))
 \quad\text{or}\quad\{-\alpha,-\beta\}\subset S.
\]
The partially split alternative applies when
\(-\alpha\in S\) or \(\{\beta,-\theta\}\subset S\).
Goodness of the residual point guarantees one of these
alternatives, and in the second it gives the partially split
shape described next. Thus rows \(2,4,8\) cannot occur in the
corresponding alternative.

At a split point \(r=D=\operatorname{diag}(\chi_i)\), \(n_v=1\).
For a tangent vector \((z,(Y_v))\), each \(z_{ij}\) belongs to
\(Z_{w_{ij}}\). The linearized ordinary condition makes every
lower entry locally a coboundary, and off-support entries vanish.
Among the independent upper odd entries \(01,03,12\), each
applicable row allows only the displayed pivot. Killing it
therefore makes all independent odd entries strict.

At a partially split point \(r=D\) and
\(n_v=1+b_vX_-\), where \(b_v\ne0\) and
\(X_-=E_{20}+E_{31}\).
The residual flag equation, in its \(20\) entry, gives
\(\chi_0|_{G_v}=\chi_2|_{G_v}\).
Thus roots with the same first coordinate induce the same local
character. Direct multiplication in \(\mathfrak{sp}_4\) gives
\[
 Z_{10}=-b_vZ_{12},\quad Z_{30}=-b_v^2Z_{12},\quad
 Z_{21}=2b_vZ_{01}+b_v^2Z_{03}
\]
whenever \(n_v^\#Zn_v\) is upper triangular.
Apply this to
\(Z=z(s)+D(s)Y_vD(s)^{-1}-Y_v\), and pass modulo local
coboundaries. We obtain
\[
 [z_{12}]=0\Longleftrightarrow[z_{10}]=0
            \Longleftrightarrow[z_{30}]=0,\qquad
 [z_{21}]=2b_v[z_{01}]+b_v^2[z_{03}].
\]
Both \(2b_v\) and \(b_v^2\) are invertible.
In each indicated row, its pivot together with off-support
zeros therefore kills all six independent odd classes.
This proves the finite table rather than assuming it.

The pivot image has dimension at most \(d+C\).
On its kernel the ten independent symplectic entries
\(i\le3-j\) are strict odd cocycles or even/zero cocycles,
of total dimension at most \(10C\).
The six lower entries of \(n_v^\#Y_vn_v\) per place have
dimension at most \(6s\le6C\).
These coordinates jointly determine the vector, since symplectic
symmetry recovers the remaining global entries and the transformed
local matrices are strictly lower triangular.
Thus its dimension is at most \(d+17C\le d+20C\).

The reduction to the two simple points uses an elementary finite
weight filtration. The residual shape is either supported in
\(\{\beta,-\alpha,-\theta,-\psi\}\), with local flag support
\(\{-\alpha,-\psi,-\theta\}\), or in
\(\{-\alpha,-\beta,-\psi,-\theta\}\).
Use diagonal conjugation weights \((0,2,1,3)\) in the first
case, and \((0,1,2,3)\) in the second.
The limits are \(D,1\).
For the partially split shape the global support is
\(\{\beta,-\alpha,-\theta\}\), the same local flag support
has \(n_{v,20}\ne0\), and \((0,1,0,1)\) gives
\(D,1+n_{v,20}X_-\).
Filter tangent matrices by their entry weights \(m_i-m_j\).
Taking the lowest nonzero homogeneous part preserves every
linearized cocycle, support, symplectic and flag equation:
all structure matrices have nonnegative degree and the limits
are their degree-zero parts.
The associated graded embeds, degree by degree, into the tangent
space of the stated limit. There are finitely many weights,
and the degree images are independent, proving the dimension
inequality without a smoothness assertion.

Finally a derivation of the representing support quotient
injects into this matrix tangent space. If its derivatives on
all universal representation and flag entries vanish, its two
maps to \(k[\varepsilon]/\varepsilon^2\) give the same universal
deformation and hence coincide.
Nakayama then generates the maximal ideal by a uniformizer and
at most \(d+20C\) further elements, proving the dimension bound.
\end{proof}

\begin{srlemma}[Ordinary-character collision]\label{lem:sr-four-collision}
Assume \eqref{eq:sr-pm-finite}.
Let \(\mathfrak q\subset B\) be a prime containing
\(J+\mathfrak m_\Lambda B\). If
\(\kappa_{v,i}=\kappa_{v,j}\) on \(U_v\) modulo \(\mathfrak q\),
for some \(i\ne j\), then
\[
 \dim B/\mathfrak q+\delta\le2d+1.
\]
\end{srlemma}
\begin{proof}
In a local domain of characteristic \(\ell\),
\(1+\mathfrak m\) has no nontrivial finite-order elements.
For \(\ell\)-power order use
\((x-1)^{\ell^r}=x^{\ell^r}-1\).
For prime-to-\(\ell\) order \(n\), the geometric sum is a
unit because its residue is \(n\); then separate these factors
for general finite order.

Local class field theory now gives \(\delta\) pinning elements
of \(I_u\) at each \(u\mid\ell\): a continuous
\(G_u\to1+\mathfrak m\) character trivial at those elements is
trivial on \(I_u\).
Indeed the torsion-free pro-\(\ell\) quotient of the local unit
group is \(\mathbb Z_\ell^\delta\), and its finite torsion is
killed by the preceding observation.
Also choose \(t_0\in I_v\) such that
\(\{t_0\}\cup(I_v\cap\ker\epsilon_\ell)\) pins such characters
on \(I_v\).
The cyclotomic image is a closed subgroup of
\(1+\ell\mathbb Z_\ell\), hence procyclic.
The extra kernel of \(\bar r\) in \(U_v\) makes no difference,
since the tower kills the residual inertia action.

Choose \(a\) with \(\{a,3-a,i,3-i\}=\{0,1,2,3\}\) and either
\(a=j\) or \(j=3-i\). In \(B/\mathfrak q\) impose
\[
 \begin{array}{ll}
 \kappa_{u,a}(t)-1&
    \text{at all \(\delta\) pins for all \(u\)},\\
 \kappa_{u,i}(t)-1&
    \text{at those pins for all \(u\ne v\)},\\
 \kappa_{v,i}(t_0)-\kappa_{v,j}(t_0).&
 \end{array}
\]
There are \(d+(d-\delta)+1=2d-\delta+1\) parameters,
all in the maximal ideal.
In every prime quotient containing them,
\(\kappa_{u,a}|_{I_u}=1\) for all \(u\), and
\(\kappa_{u,i}|_{I_u}=1\) for \(u\ne v\).
The original collision and the last parameter give equality on
all of \(I_v\).
If \(a=j\), this kills \(\kappa_{v,i}\);
if \(j=3-i\), the pairing
\(\kappa_{v,i}\kappa_{v,3-i}=1\) gives
\(\kappa_{v,i}^2=1\), which again kills it.
Pairing kills the remaining characters on all inertia.
Equation \eqref{eq:sr-pm-finite} forces this prime to be maximal.
Thus the parameter ideal has maximal radical.
The parameter dimension theorem gives
\(\dim B/\mathfrak q\le2d-\delta+1\).
\end{proof}

\begin{srlemma}[Normalized presentation]\label{lem:sr-four-presentation}
Let \(h^0,h^2\) be the degree zero and degree two dimensions of
the ordinary global-to-local adjoint cone, and \(h^1_c\) its
degree one dimension. The arithmetic cone estimate
\[
 h^2+h^0+6d\le h^1_c+4d
\]
gives a presentation
\[
 R=\Lambda[[x_1,\ldots,x_m]]/(f_1,\ldots,f_n),
 \qquad n+2+2d\le m+2s.
\]
\end{srlemma}
\begin{proof}
The fixed-similitude adjoint \(\mathfrak{sp}_4(k)\) has dimension
\(10\); an ordinary symplectic Borel has dimension \(6\);
the diagonal symplectic torus has dimension \(2\).

Temporarily omit the local \(\ell\)-power torsion condition.
For a small extension with square-zero kernel, choose lifts of
the representation and flags.
Their multiplication and flag defects form respectively a
global \(2\)-cochain and local \(1\)-cochains modulo the Borel.
Associativity is exactly the cone \(2\)-cocycle identity.
Changing lifts changes the defects by its coboundary.
Conversely a coboundary corrects both defects by its global
\(1\)-cochain and local \(0\)-cochains.
All corrections lie in \(\mathfrak{sp}_4\), maintaining the
multiplier; \(2\) is invertible.
Thus cone \(H^2\) is an obstruction space.
The minimal-relator argument gives at most \(h^2\) relations:
minimal relators map injectively to their universal lifting
obstructions, since a zero obstruction corrects the universal
lift and removes the proposed relation.

The normalized tangent dimension is at least \(h^1_c+2-h^0\).
To check the count, represent ordinary degree one classes by
global cocycles and local corrections chosen in the lower
complement of each Borel.
The global conjugation parameters have dimension \(10\), with
kernel of dimension at most \(h^0\).
Normalizing at \(g_0\) imposes at most \(10-2=8\) conditions.
The resulting dimension inequality is
\(h^1_c+2\le\dim T+h^0\).
Dual numbers identify \(T\) with the tangent space of the
representing ring.

Restore the local torsion condition. A generator of the
\(\ell\)-power torsion in each local abelianization suffices.
Its multiplier is \(1\): the cyclotomic image has no
\(\ell\)-power torsion for odd \(\ell\), and the remaining
finite multiplier has order prime to \(\ell\), including the
finite \(2\)-power character in this packet.
Only the first two of the four paired ordinary characters need
to be set equal to \(1\). This adds \(2s\) equations.
Therefore
\[
 n\le h^2+2s,\qquad m\ge h^1_c+2-h^0,
\]
and the cone estimate gives \(m-n\ge2d+2-2s\).

For this cone estimate the local Borel dimension is \(6\),
and the complex-conjugation fixed adjoint dimension is \(4\):
an anti-symplectic involution has two isotropic eigenspaces of
dimension \(2\), and its fixed symplectic Lie algebra is
\(\mathfrak{gl}_2\).
The required absence of a cyclotomic quotient is checked on the
residual adjoint root characters. The actual tower supplies,
for each such character, an element where it differs from the
mod-\(\ell\) cyclotomic character; the zero character is also
separated by complex conjugation. This is exactly the
noncyclotomic condition in the tower construction, rather than
an inference from pairwise distinct values alone.
Thus the arithmetic cone estimate applies to this adjoint.
\end{proof}

\begin{srproposition}[Rank-four component propagation]
\label{prop:sr-four-core}
Assume the common-space construction, the tower cohomology
bounds, \eqref{eq:sr-pm-finite}, and the actual modular curve
propagation property stated above.
If \(d\ge30(C+1)^2\), then at a good residual point, one
component containing \(JA\) implies that every component
through that point contains \(JA\).
\end{srproposition}
\begin{proof}
Put \(N_0=2d+1-2s\).
As \(s\le C\), \(s\delta=d\), and \(d\ge30(C+1)^2\),
\[
 \delta\ge30(C+1),\quad
 d+20C+1<N_0,\quad
 2d+1+2+4s<N_0+\delta.
\]
For the last inequality it suffices that \(2+6s<\delta\).

Lemma~\ref{lem:sr-four-presentation} and the complete-local
component-chain lemma join minimal components by a chain whose
neighbors meet in dimension at least \(N_0+1\).
Lemma~\ref{lem:sr-four-reducible-tangent} makes every reducible
bad set have dimension less than \(N_0\).
Outside these sets, a bad prime where \(J\), the uniformizer,
and an ordinary collision vanish satisfies
\[
 \dim R/\mathfrak q\le
 \dim B/\mathfrak q_B+2+4s
 \le2d+1-\delta+2+4s<N_0
\]
by the trace-fiber and collision lemmas.

Suppose \(J\) vanishes on the first of two neighbors.
The rank-independent prime-avoidance and curve lemma, applied
to their intersection and the finite collection of these bad
ideals, gives a characteristic-\(\ell\) curve with finite
residue field. It kills \(J\), has irreducible generic
representation, and has pairwise distinct ordinary characters
on each \(U_v\).
The modular curve property implies that \(J\) is contained in
the contraction of the second component, since that prime
also lies under the curve kernel.
Thus \(J\) vanishes on the second component.
Induction along the chain proves the assertion in the
completion; faithfully flat transport gives it in \(A\).
\end{proof}

\begin{srlemma}[Characteristic-\(\ell\) consequence]
\label{lem:sr-four-characteristic-ell}
Suppose an actual principal Case II point factors through such
a modular ideal and has characteristic \(\ell\).
Its characteristic polynomial at every \(g\in G_{\mathbb Q}\)
is then the product of the four prescribed residual characters.
It is neither absolutely irreducible of rank \(4\) nor a block
extension with an irreducible rank-\(2\) constituent.
\end{srlemma}
\begin{proof}
Use the stable valuation lattice and the factorization on \(G_F\).
The prescribed ordinary characters become \(1\) on inertia:
the cyclotomic factor reduces to its mod-\(\ell\) value killed
by the tower, and the other finite characters are unramified.
Equation \eqref{eq:sr-pm-finite} therefore kills every element
of \(\mathfrak m_B\), because the target is a domain.
On \(G_F\), all characteristic-polynomial coefficients are the
given finite-field residual constants.

For \(g\in G_{\mathbb Q}\), some positive power \(g^M\) lies in
\(G_F\).
The eigenvalues of \(r(g)^M\), and therefore of \(r(g)\), are
algebraic over \(\mathbb F_\ell\).
Extend the valuation to an algebraic closure of the fraction
field. It is trivial on the algebraic closure of
\(\mathbb F_\ell\), on which reduction is injective.
Thus the characteristic-polynomial coefficients of \(r(g)\)
are constants, and their reduction identifies them with the
given residual polynomial.
Characteristic-polynomial Brauer--Nesbitt identifies its
semisimplification with the four diagonal characters and gives
both exclusions. Complete characteristic polynomials are used,
so this argument also applies when \(\ell=3\).
\end{proof}

\paragraph{Scope and actual remaining producers.}
For the principal packet with \(a\ge2\), the required coefficient
prime is prime to \(N\): in the resonant branch
\(\ell\le2a+1<q,d\), and in the \(3,7\) branch \(q,d>7\).
Thus \(\chi_P=\psi\), and the nontrivial datum is Case II.
The actual host uses the real dyadic cyclotomic subfield;
its dyadic characters are even and its finite character orders
are powers of \(2\).
That makes the prime-to-\(\ell\) finite-character condition in
the tower automatic. The sufficiently large cyclotomic layer
is still necessary for \(d\ge30(C+1)^2\).

The preceding lemmas reconstruct the independent rank-four
characteristic-\(\ell\) core.
The actual Case II endpoint additionally requires the
construction of its linked automorphic rank-four seed and the
construction of its modular ideal with finite-weight, curve,
seed, purity, and automorphy properties.
After those constructions, one chooses a tame Case II tower,
passes to the sufficiently high layer, propagates from the
seed along the residual links, and factors the actual point
through the modular ideal.
The characteristic-\(\ell\) lemma supplies the contradiction;
the actual purity and automorphy properties apply to
characteristic-zero points. The usual deformation endpoint
then gives finite weight zero, no full upper block, and
automorphy at every irreducible point.
No Case III seed or Case III promotion is needed.

% Source alignment, frozen at alpoge-source-7cbf7a0c:
% WeakShimuraTheory/FibreTraces/FibTTrace.lean (fibTorus);
% WeakShimuraTheory/CompactIntegralFields/FibLow.lean (fib);
% WeakShimuraTheory/RedIdeal.lean (redTypes, Bounds, Good, Res);
% WeakShimuraTheory/RedTanFlow.lean (tan_lowest, tan_flow);
% WeakShimuraTheory/CompactIntegralFields/RedTan.lean (redTan);
% WeakShimuraTheory/CompactIntegralFields/HomOfDer.lean (red);
% WeakShimuraTheory/RedDim.lean (redDim);
% WeakShimuraTheory/CompactIntegralFields/ImportsAlt.lean (nn);
% WeakShimuraTheory/AdjointCoefficients/AbsPres.lean;
% WeakShimuraTheory/AdjointCoefficients/MuQuot.lean;
% WeakShimuraTheory/CompactIntegralFields/IncLoc.lean (x1);
% WeakShimuraTheory/CompactIntegralFields/ChCore.lean (ch);
% ModularityArgument/ThresholdArguments/DenseThresholdCases/
% Gsp4PromodularityCaseTwo.lean
%   (proMod_of_pmEx_of_ctx2_of_kPair_of_star);
% PotentialModularityTower.lean;
% ModularityRankFourOfPromodularity.lean.

\subsection{Generic rank-four curves and regular Frobenius cosets}
\label{subsec:sr-four-generic-curves}
We supply the rank-four generic-curve calculation used by
ordinary support patching, including its regular-coset assertion.

\begin{srlemma}[One local element separates all powers]
\label{lem:sr-four-free-local-element}
Let \(k\) be a finite field of characteristic \(\ell\), and
\(\kappa_0,\ldots,\kappa_3:H\to k[[T]]^\times\) be characters.
Let \(U\subset H\) be a subgroup on which all their residual
characters agree.
If every ratio \(\kappa_i/\kappa_j\), \(i\ne j\), is nontrivial
on \(U\), then some \(u\in U\) has all six ratios
\(\kappa_i(u)/\kappa_j(u)\) of infinite order.
\end{srlemma}
\begin{proof}
The group \(1+Tk[[T]]\) is torsion-free:
separate the \(\ell\)-power factor, using
\((x-1)^{\ell^r}=x^{\ell^r}-1\), from the prime-to-\(\ell\)
factor, whose geometric sum is a unit.
Choose finitely many elements of \(U\) witnessing nontriviality
of the six ratios.
Tensor the torsion-free subgroup generated by their ratio
values with \(\mathbb Q\).
For each ratio, products of these elements to integer powers
are trivial only on a proper rational linear subspace of the
exponent lattice. Finitely many proper rational linear
subspaces do not cover that lattice.
A product with exponents outside their union gives \(u\).
Every nontrivial ratio at \(u\) belongs to the torsion-free
group just described and has infinite order.
\end{proof}

\begin{srlemma}[A large residual ratio excludes virtual commutativity]
\label{lem:sr-four-noncommutative}
Let \(\Gamma\) be profinite, \(Z\triangleleft\Gamma\) closed,
and let \(r:\Gamma\to\mathrm{GL}_4(K)\) be continuous and
absolutely irreducible, where \(K\) is an algebraically closed
valued field. Suppose its characteristic polynomials are
integral and reduce to
\(\prod_i(X-\chi_i(g))\) for four residual characters.
Assume some \(\chi_i/\chi_j\) has order greater than \(4\).
If \(Z\) contains \(u\) for which every positive power \(r(u)^m\)
has four distinct eigenvalues, then every open normal subgroup
of \(Z\) has noncommutative image.
\end{srlemma}
\begin{proof}
Suppose \(N\triangleleft Z\) has finite index \(m\) and
commutative image. Its normal core \(C\) in \(\Gamma\)
contains \(u^m\): every conjugate of \(u\) is in \(Z\), whose
\(m\)-th power is in \(N\).
The image of \(C\) commutes and centralizes the regular
semisimple \(r(u^m)\), hence is diagonal in its eigenbasis.
Because \(C\) is normal, \(\Gamma\) permutes its four common
eigenlines. The characters on these lines are distinct, as
they are distinct at \(u^m\).
Absolute irreducibility makes that permutation transitive.
Thus \(r\) is monomial, with permutation map
\(\pi:\Gamma\to S_4\).

For \(g\in\ker\pi\), write its four eigenvalues as \(a_i(g)\).
They are units since the characteristic polynomials of both
\(g\) and \(g^{-1}\) are integral.
If \(h\) sends the \(j\)-th eigenline to the \(i\)-th, the
commutator \(ghg^{-1}h^{-1}\) has eigenvalue
\(a_i(g)/a_j(g)\).
Its residual characteristic polynomial is \((X-1)^4\),
because all \(\chi_i\) are characters.
Hence the four \(a_i(g)\) have the same reduction.
Comparing characteristic polynomials at \(g\) shows
\(\chi_i(g)=\chi_j(g)\) for every \(i,j\).

Thus every residual ratio factors through the subgroup
\(\pi(\Gamma)\subset S_4\).
Its character image is finite cyclic.
Choose an element mapping to a generator of that cyclic image.
An element of \(S_4\) has order \(1,2,3\), or \(4\), so the
ratio has order at most \(4\), a contradiction.
\end{proof}

\begin{srlemma}[Rank-four regular-coset lemma]
\label{lem:sr-four-regular-coset}
Let \(G\) be a profinite group and
\(r:G\to\mathrm{GL}_4(K)\) an absolutely irreducible
representation over an algebraically closed field of
characteristic different from \(2\).
Assume polynomial zero sets of its matrix entries are closed.
Suppose \(r(G)\) contains a semisimple element with all
eigenvalue ratios of infinite order, and every open normal
subgroup of \(G\) has noncommutative image.
Then for every open subgroup \(U\subset G\) and every
\(\sigma\in G\), the coset \(\sigma U\) contains an element
whose characteristic polynomial is separable.
\end{srlemma}
\begin{proof}
Let \(\mathcal G\) be the Zariski closure of \(r(G)\) and
\(\mathcal H=\mathcal G^\circ\).
The preimage of \(\mathcal H\) is closed of finite index,
hence open.
A positive power of the specified semisimple element lies
in \(\mathcal H\), still with all ratios of infinite order.

The restriction of the natural representation to the normal
group \(\mathcal H\) is semisimple by the socle argument:
conjugates of a simple submodule are simple and their sum is
\(\mathcal G\)-invariant.
The isotypic components form one transitive \(\mathcal G\)-orbit.
There is no multiplicity greater than one, since the regular
element in \(\mathcal H\) has no repeated eigenvalues.
Their equal dimensions therefore give one component of
dimension \(4\), two components of dimension \(2\), or four
lines.
The last case would make the open preimage of
\(\mathcal H\) commutative, contrary to hypothesis.

The connected torus which is the identity component of the
closure of the cyclic semisimple subgroup has four distinct
characters in this eigenbasis. Otherwise some eigenvalue
ratio would have finite order.
Choose a cocharacter of that torus with four distinct integer
weights \(w_0,\ldots,w_3\), and write
\(D(z)=\operatorname{diag}(z^{w_i})\).

If \(\mathcal H\) acts irreducibly in dimension \(4\), or the
coset under consideration preserves its two irreducible
\(2\)-blocks, every diagonal entry and every diagonal cofactor
is a nonzero regular function on that coset of \(\mathcal H\).
For a diagonal entry this follows because the orbit of the
corresponding vector spans its irreducible block; for a
cofactor apply the same observation to the inverse matrix.
Connectedness makes \(\mathcal H\) irreducible as a variety.
Thus there is a matrix \(x\) in the coset with all these
entries and cofactors nonzero.

We check directly that \(xD(z)\) has separable characteristic
polynomial over \(K(z)\).
Relabel so \(w_0<w_1<w_2<w_3\).
Write its polynomial as
\(X^4+c_1X^3+c_2X^2+c_3X+c_4\), and use the \(z\)-valuation.
The unique lowest terms give
\[
 v(c_1)=w_0,\quad v(c_3)=w_0+w_1+w_2,\quad
 v(c_4)=w_0+w_1+w_2+w_3.
\]
If the \(01\) principal \(2\)-minor of \(x\) is nonzero,
\(v(c_2)=w_0+w_1\). The Newton polygon has four
length-one edges with distinct slopes, giving four distinct
roots.
Otherwise every possible term of \(c_2\) has valuation at least
\(w_0+w_2\), which is greater than
\((2w_0+w_1+w_2)/2\).
There are the two outer length-one edges and one middle edge
of length two.
After adjoining \(z^{1/2}\) if necessary, its residual
polynomial is \(Y^2+c\), with \(c\ne0\).
It is separable because \(2\ne0\), so Hensel's lemma again
gives four distinct roots.
This proves the first case without a quartic discriminant
expansion.

In the remaining case the coset interchanges the two
irreducible \(2\)-blocks.
Its matrices have the form
\(\begin{psmallmatrix}0&B\\C&0\end{psmallmatrix}\).
Irreducibility of each block and the same nonzero-function
argument allow all entries of \(B,C\) to be chosen nonzero;
both blocks are invertible.
Order the two cocharacter weights in each block as
\(u_0<u_1\) and \(v_0<v_1\).
The characteristic polynomial after multiplying by \(D(z)\)
has the form
\[
 X^4-A(z)X^2+\Delta(z).
\]
Here the unique lowest term of \(A(z)\) is a nonzero product
of two block entries times \(z^{u_0+v_0}\), and
\[
 v(\Delta)=u_0+u_1+v_0+v_1>2v(A).
\]
The quadratic in \(X^2\) has two distinct nonzero roots of
different valuations. Taking their square roots gives four
distinct roots since \(2\ne0\).

Thus the separable locus is a nonempty Zariski open in every
coset of \(\mathcal H\).
The Zariski closure of \(r(U)\), for a finite-index \(U\),
contains \(\mathcal H\): finitely many translates of that
closure cover \(\mathcal G\), so its identity component has
the same dimension.
Consequently \(r(\sigma U)\) meets the nonempty separable
open just found. This proves the lemma.
\end{proof}

\begin{srproposition}[The actual rank-four curve conditions]
\label{prop:sr-four-curve-genericity}
Let an actual ordinary rank-four curve over \(k[[T]]\) have
absolutely irreducible generic representation.
Suppose its four ordinary characters are pairwise distinct
on each \(U_v\), its residual packet satisfies condition
\(\mathrm{(C)}\) of Lemma~\ref{lem:sr-rank-four-no-self-twist},
and one residual character ratio has order greater than \(4\).
Then it has:
\[
 \begin{gathered}
 \kappa_{v,i}\ne\kappa_{v,j},\qquad
 \kappa_{v,i}\ne\epsilon_\ell\kappa_{v,j}\quad(i\ne j),\\
 r|_{\ker\epsilon_\ell}\text{ absolutely irreducible},\\
 \text{a regular semisimple element in every open coset of }
                         \ker\epsilon_\ell .
 \end{gathered}
\]
In particular these give the rank-four local genericity,
cyclotomic irreducibility, and regular-Frobenius-coset inputs
to the support patching theorem.
\end{srproposition}
\begin{proof}
Both local inequalities follow by restriction to \(U_v\),
where \(\epsilon_\ell=1\).
Lemma~\ref{lem:sr-four-free-local-element} gives an element
\(u\) in one such subgroup whose ordinary eigenvalue ratios
all have infinite order; its characteristic polynomial and
every power's characteristic polynomial are separable.

The proof of Lemma~\ref{lem:sr-rank-four-no-self-twist}
works equally over the algebraic closure of \(k((T))\):
normal-subgroup semisimplicity uses the socle argument, and
the quotient by the cyclotomic kernel is still procyclic.
A nontrivial isotypic orbit has even length dividing \(4\),
hence gives a quadratic self-twist factoring through the
mod-\(\ell\) cyclotomic character. Its values \(1,-1\)
remain distinct after reduction.
Condition \(\mathrm{(C)}\) excludes it.
A single isotypic component with multiplicity greater than
one is excluded by the tensor-factor eigenline argument
and closedness of subspace stabilizers.
This proves cyclotomic irreducibility in characteristic
\(\ell\), without assuming it from the characteristic-zero
case.

Apply Lemma~\ref{lem:sr-four-noncommutative} with
\(Z=\ker\epsilon_\ell\).
It shows that no open normal subgroup of \(Z\) has
commutative image.
The regular-coset lemma now applies to \(r|_Z\).
Polynomial zero sets are closed because matrix entries lie
in the valued power-series coefficient field; a finite
coefficient extension containing a polynomial's coefficients
suffices for this assertion.
This gives the stated regular-coset condition.
\end{proof}

% Frozen-source correspondence:
% WeakShimuraTheory/CompactIntegralFields/H4.lean;
% WeakShimuraTheory/StabilizerIndexBounds/
% ExistsMemPowerSeriesCharRatiosNotIsOfFinOrder.lean, NonComm.lean;
% AlgebraTools/MonomialRepresentations/
% MonomialOfComm.lean, MonomialRed.lean, PermChar.lean, RegCosetFour.lean,
% DiscSwap.lean, RelWeight.lean, TorusTransfer.lean.
% The Newton-polygon argument replaces the source's explicit
% quartic discriminant algebra; it proves the same separability claim.

\subsection{The Case II Hilbert seed and its induction}
\label{subsec:sr-case-two-seed}
We separate the mixed Hilbert seed from its elementary induction.
The latter is ordinary induction from a real quadratic field; the
source name for its matrix change of basis does not denote a
symmetric-square functor.

\begin{srlemma}[Mixed ordinary Hilbert seed]
\label{lem:sr-mixed-hilbert-seed}
Let \(k_0\) be a real quadratic field in which the odd prime
\(\ell\) splits. Let \(\xi:G_{k_0}\to\overline{\mathbb Q}_\ell^\times\)
be finite and totally odd, unramified at \(\ell\), of order prime
to \(\ell\), with \(\bar\xi^2\ne1\).
Suppose its conductor is supported above \(2\), with a ramified
prime of norm \(2\), and its Frobenius value at each prime above
\(\ell\) is nontrivial modulo \(\ell\).
Write these two primes as \(v_+,v_-\).

Using the classical Hilbert Eisenstein, integral q-expansion and
bounded-denominator, ordinary control, and local--global
compatibility theorems at this tame conductor,
there are \(k,t\ge2\), both divisible by \(\ell-1\), and a
cuspidal Hilbert modular representation with associated
\(r:G_{k_0}\to\mathrm{GL}_2(E)\), for a finite
\(E/\mathbb Q_\ell\), such that
\[
 \det r=\epsilon_\ell^k\xi,\qquad
 \bar r^{\,\mathrm{ss}}=1\oplus\bar\xi.
\]
It is irreducible on every open subgroup. Its ordinary weights
and residual ordering are
\[
 \begin{array}{c|c|c}
 &\text{weights}&\text{ordered residual characters}\\\hline
 v_+&(k+t,-t)&(\bar\xi,1)\\
 v_-&(k,0)&(1,\bar\xi).
 \end{array}
\]
Away from \(\ell\), its restriction to inertia is conjugate to
\(1\oplus\xi\).
The same construction permits arbitrary sufficiently large \(k,t\).
\end{srlemma}
\begin{proof}
Let \(\omega\) be the finite Hecke character corresponding to
\(\xi\), with Artin reciprocity convention chosen accordingly.
The parallel weight-one Eisenstein series \(E_1(1,\omega)\)
has normalized nonconstant ideal coefficients
\[
 a(\mathfrak m)=\sum_{\mathfrak a\mid\mathfrak m}
                         \omega(\mathfrak a),\qquad a(1)=1.
\]
In particular these coefficients are integral.
At each prime above \(\ell\), the two old-space ordinary roots
are \(1,\omega(v)\).
Stabilize with the second root at \(v_+\) and the first at
\(v_-\). Denote the resulting form by \(E\).
Equivalently apply \(1-V_{v_+}\) and
\(1-\omega(v_-)V_{v_-}\), with the usual ideal-shift
normalization of \(V\).
Its nonconstant coefficients remain integral and its first is
\(1\).

Choose a parallel weight \(k\), divisible by \(\ell-1\), and
a level-one Hasse lift \(H\) whose expansion at every cusp is
\(1\) to sufficiently high \(\ell\)-adic precision.
Such lifts exist with arbitrarily large weight and precision.
The finitely many constant coefficients of \(E\) have bounded
denominators. Choose the precision larger than those denominators.
Then the nonconstant coefficients of \(EH-E\) are in the
coefficient maximal ideal, and \(EH\) has weight \(k+1\).
No assertion that the weight-one constant coefficients themselves
are integral is needed.

The Hilbert weight-one Fourier formula is classical; a
real-quadratic explicit statement is
\href{https://webpages.leeu.edu/rmoy/files/2018/08/noncm.pdf}
{Moy--Specter, Proposition 2.1}.
The level-one Hasse lift and arbitrary precision follow from
\href{https://arxiv.org/pdf/1911.11196}
{Deo--Dimitrov--Wiese, Lemmas 2.2--2.3}.

For primes \(\mathfrak p\nmid\ell\), the Hecke coefficient
formula changes on passing from weight \(1\) to weight \(k+1\)
only by the factor \(\mathrm N\mathfrak p^k\), which is \(1\)
modulo the maximal ideal. At the Iwahori primes above \(\ell\),
the \(U\)-formula is the same coefficient shift.
Thus \(EH\) is a simultaneous eigenvector modulo that ideal
with the chosen two ordinary roots.
Use the Hecke-stable lattice defined by its nonconstant Fourier
coefficients in the finite-dimensional holomorphic space.
It is a genuine lattice: finitely many coefficient evaluations
are injective, so the integrality conditions give a bounded
module, and the integral q-expansion/bounded-denominator theorem for
classical Hilbert forms supplies a basis with bounded coefficient
denominators, hence a full sublattice. This integral theorem is
an explicit published input here, not a consequence of finite
dimensionality alone.
The commuting finite Hecke algebra and the elementary
Deligne--Serre lifting lemma give a characteristic-zero
holomorphic eigenform \(g\) with the same residual eigensystem.
Its two \(U\)-eigenvalues are units, with residues
\(\bar\omega(v_+)\) and \(1\).

We check cuspidality, which is essential here.
If \(g\) were Eisenstein, its unramified eigenvalues would come
from two global finite characters \(\eta_1,\eta_2\), with
\(\eta_1\eta_2=\omega\), one carrying the factor
\(\mathrm N\mathfrak p^k\).
Chebotarev and characteristic polynomials identify their
residual sum with \(1\oplus\bar\omega\).
Hence the globally unordered residual pair is
\(\{1,\bar\omega\}\).
Since \(k>0\), precisely one of its two local roots at every
\(\ell\)-prime is a unit: the character without the positive
\(\ell^k\) factor. It is the same globally chosen character at
both primes.
But at \(v_+\) it must reduce to \(\bar\omega(v_+)\ne1\),
whereas at \(v_-\) it must reduce to \(1\ne\bar\omega(v_-)\).
Neither global choice is possible. Thus \(g\) is cuspidal.

Apply nearly ordinary control and independence of algebraic
weight, in the component of \(g\), with fixed central character.
Choose the arithmetic weight giving the displayed two weight
pairs, with \(t\) divisible by \(\ell-1\).
Both pairs have sum \(k\), so this holds the determinant fixed.
Their gaps are positive and can be chosen arbitrarily large,
so ordinary classicality applies.
For this application the ordinary module is an actual
definite-quaternion module, not a formal instance of an
axiomatic control statement.
Take the quaternion algebra over the real quadratic field
\(k_0\) ramified at both real places and at no finite place.
It exists because the number of ramified places is even.
At finite places identify its maximal compact subgroups
with those of \(\mathrm{GL}_2(k_{0,v})\), using the same
tame level \(\Gamma_1(\mathfrak n)\) and the same Iwahori
and deeper ordinary levels above \(\ell\).
The coefficient module consists of functions on its finite
double-coset set with the fixed central character.
The local matrix translations, lower-unipotent averaging
and ordinary projector are the rank-two instance of the
ordinary function-module construction.
That construction verifies the projectivity and transition
axioms required for control; the finite central-character
projector is exact since its relevant diamond orders are
powers of \(2\).
The algebraic coefficient at an infinite place has highest
weight difference equal to the prescribed positive ordinary
gap minus one.
Evaluation on its highest-weight line, with the normalized
\(U(y)\) operators, gives weight independence.
At classical regular weights, Jacquet--Langlands identifies
the resulting eigensystems with the cuspidal Hilbert
eigensystems at precisely these finite levels.
Thus this actual module verifies the hypotheses of the
ordinary control theorem.
Its finite torsion-free weight algebra is faithful on its
nonzero localized ordinary space; lying over supplies an
arithmetic point above the chosen weight.
The normalized \(U\)-characters retain the two residual
orientations because \(\bar\xi(\operatorname{Frob}_{v_\pm})\ne1\).
The control and weight normalization are the actual
definite-quaternion instances of
\href{https://www.math.ucla.edu/~hida/Lum07.pdf}
{Hida, Theorems 1.2 and 2.4}.
Theorem 1.2 requires its module axioms; the preceding
function-module construction verifies them here.
The operator \(U(y)\) is divided by its lowest-weight factor
before applying control.

The classical Hilbert Galois representation theorem gives an
irreducible ordinary \(r\) with the stated determinant and
residual semisimplification.
The tame level throughout is exactly
\(\mathfrak n=\operatorname{cond}(\xi)\):
the Eisenstein series starts there, stabilization adds only
the two primes above \(\ell\), the Hasse lift has level one,
and ordinary weight variation keeps this tame level.
If one realizes the spaces at a neat level, it is
\(\Gamma_1(\mathfrak n)\) already within this conductor;
no extra tame prime is introduced.

Here is the local conductor check.
At a prime where \(\xi\) is ramified, write \(V\) for the
two-dimensional Weil--Deligne representation.
Its conductor is at most the level exponent, which is the
conductor exponent of \(\det V=\xi\).
If \(V^I=0\), then
\(a(\det V)\le a(V)-1\):
in the Swan integral the codimension for the determinant is
at most that for \(V\) at every ramification subgroup, while
the tame codimension is strictly smaller.
This is impossible. Consequently \(\dim V^I=1\), and the
inertial representation is \(1\oplus\xi\).
Monodromy is zero, since it commutes with inertia and no
nonzero nilpotent preserves these two distinct one-dimensional
characters. At the other tame primes the level exponent is
zero and the representation is unramified.
Local--global compatibility proves the asserted exact inertia
type.

We finish by proving the stronger open-subgroup irreducibility
rather than adding it as an input.
If \(r\) becomes reducible on a normal open subgroup, Clifford
theory gives either two distinct one-dimensional constituents
or a scalar restriction.
In the former case the two constituents are interchanged,
giving a nontrivial quadratic self-twist \(\eta\).
Since \(\ell\) is odd, reduction preserves its nontriviality.
Then
\[
 \bar\eta(1\oplus\bar\xi)\simeq1\oplus\bar\xi
\]
forces either \(\bar\eta=1\), or
\(\bar\eta=\bar\xi\) and \(\bar\xi^2=1\), both impossible.
In the scalar case the projective image is finite.
The ordinary inertia ratio at \(v_-\) is
\(\epsilon_\ell^k\), of infinite order, excluding this case.
An arbitrary open subgroup contains a normal open subgroup,
so every open restriction is irreducible.
\end{proof}

\paragraph{Published input normalization.}
The preceding lemma uses the classical Hilbert ordinary theory
at the original tame conductor, not an automorphy lifting
theorem for arbitrary reducible residual representations.
The weight-one construction and the two opposite unit roots
produce the initial cuspidal eigensystem explicitly.
The Hasse lift has level one, so it adds no tame prime.
At small \(\ell\), one uses the neat \(2\)-power level already
present in the principal Case II datum and then its finite
character projector; the relevant diamond-group orders are
powers of \(2\), hence units at odd \(\ell\).
The local type assertion must be read in the same
local--global convention as the Hilbert Galois representation.

\begin{srproposition}[The induced Case II seed]
\label{prop:sr-induced-case-two-seed}
Consider an actual Case II tower \(F,L=FK\), and its real
quadratic field \(k_0\), split at \(2\) and \(\ell\).
Let \(\varphi\) be the quadratic character of \(k_0/\mathbb Q\).
On \(G_F\), its residual value is the quadratic root character
in the four-character packet.
Assume the tower root conditions: this character is nontrivial
on \(G_L\), and the packet has no nontrivial quadratic
self-twist on \(G_L\) factoring through mod-\(\ell\)
cyclotomic image.
Let the finite character \(\xi\) be the restriction of the
auxiliary rational character \(\xi:G_{\mathbb Q}\to E^\times\)
of the Case II datum. Let \(r\) be its mixed Hilbert seed from
Lemma~\ref{lem:sr-mixed-hilbert-seed}, with the two ordinary
patterns interchanged by the nontrivial automorphism of
\(k_0/\mathbb Q\).

Then, after a finite coefficient extension, ordinary induction,
a specified symplectic change of basis, and a unique character
twist congruent to \(1\) give the linked rank-four Case II seed:
it is irreducible on \(G_L\cap\ker\epsilon_\ell\), has the
prescribed multiplier and ramification conditions, and its
residual point is in the partially split good stratum.
Its ordinary squared-weight gaps are at least \(4\).
After an explicitly removable scalar twist it is attached to a
regular algebraic cuspidal representation over \(L\).
\end{srproposition}
\begin{proof}
Put \(H=G_{k_0}\), let \(t_0\) be the tower's distinguished
element outside \(H\), and set
\[
 \mu=\epsilon_\ell^k\xi\varphi
       \quad\text{on }G_{\mathbb Q}.
\]
The determinant of \(r\) is \(\mu|_H\).

First choose a Ribet lattice on \(H\cap G_F\).
This restriction is irreducible, whereas its residual
semisimplification is \(1\oplus\bar\xi\).
The values of these characters at \(t_0^2\) are distinct.
Thus the stable-lattice fixed set in the Bruhat--Tits tree lies
in the apartment determined by the two eigenlines of
\(r(t_0^2)\).
It is a bounded interval: an unbounded interval would give an
invariant line over the fraction field.
At its endpoint with upper residual orientation, reduction is
an extension of \(\bar\xi\) by \(1\) that is nonsplit on
\(H\cap G_F\); otherwise the interval could be extended.
In the corresponding eigenbasis,
\[
 r(t_0^2)=\operatorname{diag}(\alpha^2,\beta^2),\qquad
 \alpha\beta=\mu(t_0),
\]
and all matrices on \(H\cap G_F\) are integral, reducing to
\(\begin{psmallmatrix}1&b(h)\\0&\bar\xi(h)\end{psmallmatrix}\),
with \(b(h)\ne0\) for some \(h\).
Choose the square root \(\alpha\) in the prescribed residue
class \(\pm1\); Hensel's lemma applies since \(\ell\ne2\).
Then put \(\beta=\mu(t_0)/\alpha\).
The residual packet identity fixes its required residue too.

In the coset basis, \(I=\operatorname{Ind}_H^{G_{\mathbb Q}}r\)
has matrices
\[
 I(h)=
 \begin{pmatrix}r(h)&0\\0&r(t_0^{-1}ht_0)\end{pmatrix},
 \qquad
 I(t_0)=
 \begin{pmatrix}0&r(t_0^2)\\1&0\end{pmatrix}.
\]
Write \(J_2=\begin{psmallmatrix}0&1\\-1&0\end{psmallmatrix}\).
The alternating form
\(\operatorname{diag}(J_2,\mu(t_0)J_2)\) has multiplier
\(\mu\) under both displayed kinds of matrices.
For the second matrix this follows from
\(\det r(t_0^2)=\mu(t_0)^2\).
Define
\[
 P=
 \begin{pmatrix}
 0&\alpha&0&\alpha\\
 \beta&0&-\beta&0\\
 0&-1&0&1\\
 1&0&1&0
 \end{pmatrix}.
\]
It is invertible because \(2\alpha\beta\ne0\), and
\[
 P^t\operatorname{diag}(J_2,\alpha\beta J_2)P
       =-2\alpha\beta J_4,\qquad
 J_4=\begin{pmatrix}
 0&0&0&1\\0&0&1&0\\0&-1&0&0\\-1&0&0&0
 \end{pmatrix}.
\]
Consequently \(R=P^{-1}IP\) is a symplectic similitude of
multiplier \(\mu\), and
\[
 R(t_0)=\operatorname{diag}(\beta,-\alpha,-\beta,\alpha).
\]
These equations specify both the normalization and the pairing.

We prove the required cyclotomic irreducibility.
Put \(V=G_L\) and \(U=V\cap\ker\epsilon_\ell\).
If \(\varphi|_U=1\), then \(\varphi|_V\) factors through the
closed cyclotomic image.
Its quadratic character kills the pro-\(\ell\) subgroup, so
it factors through the mod-\(\ell\) image.
But multiplication by \(\varphi\) permutes the residual packet
\(\{1,\bar\xi,\varphi,\bar\xi\varphi\}\).
The no-self-twist condition would force \(\varphi|_V=1\),
contrary to the other root condition.
Thus \(U\not\subset H\).

The restriction \(r|_{H\cap U}\) is irreducible:
apply the strong-open irreducibility and procyclic-restriction
lemma to the open group \(H\cap V\) and its cyclotomic kernel.
Suppose its conjugate by \(t_0\) were isomorphic to it, with
intertwiner \(M\).
For every \(h\in H\cap V\), normality of \(H\cap U\) and
Schur's lemma make
\[
 r(h)^{-1}M^{-1}r(t_0^{-1}ht_0)M
\]
a scalar. The determinants of the two representations agree,
so that scalar has square \(1\).
It follows that \(r(h^2)\) and \(r(t_0^{-1}h^2t_0)\) have
the same characteristic polynomial.

Choose \(h\) in inertia above \(\ell\) and in \(H\cap V\)
with \(\epsilon_\ell(h)\) of infinite order.
At the two conjugate primes, the eigenvalue pairs of these
squares are respectively
\[
 \{\epsilon_\ell(h)^{2(k+t)},\epsilon_\ell(h)^{-2t}\},
 \qquad
 \{\epsilon_\ell(h)^{2k},1\}.
\]
They differ: \(1\) belongs only to the second pair.
Thus the two irreducible constituents on \(H\cap U\) are
nonisomorphic.
The elementary index-two induction criterion now gives
\(R|_U\) irreducible.

Here is the residual good-point calculation.
For \(h\in H\cap G_F\), the two blocks reduce to
\[
 A=\begin{pmatrix}1&b\\0&x\end{pmatrix},\qquad
 B=\begin{pmatrix}1&b'\\0&x\end{pmatrix},
 \qquad x=\bar\xi(h).
\]
Put \(\gamma=\bar\beta/\bar\alpha\).
Direct multiplication by \(P^{-1}\) and \(P\) gives
\[
 \bar R(h)=
 \begin{pmatrix}
 x&0&0&0\\
 (\gamma b-b')/2&1&-(\gamma b+b')/2&0\\
 0&0&x&0\\
 (\gamma b+b')/2&0&-(\gamma b-b')/2&1
 \end{pmatrix}.
\]
Its off-diagonal support is
\(\{-\alpha_{\rm root},\beta_{\rm root},-\theta_{\rm root}\}\)
in the root notation of the preceding subsection.
At a point with \(b\ne0\), either \(\gamma b-b'\ne0\),
or \(\gamma b+b'=2\gamma b\ne0\).
Hence either the negative-\(\alpha_{\rm root}\) entry is
nonzero, or both the \(\beta_{\rm root}\) and
negative-\(\theta_{\rm root}\) entries are nonzero.
The outside coset differs by the normalized diagonal
\(R(t_0)\), so the same support and the prescribed diagonal
packet hold on all of \(G_F\).

We also check the local flag rather than inferring it merely
from induction.
At the prime with weights \((k+t,-t)\), the high line in
the first block has residual eigencharacter \(\bar\xi\);
write a primitive vector on it as \((u,1)\).
The high line in the other block, with weights \((k,0)\),
has residual character \(1\), so it reduces to its first
coordinate line.
In the mixed basis these two vectors have coordinates
\[
 (1/(2\beta),u/(2\alpha),-1/(2\beta),u/(2\alpha)),
 \qquad (0,-1/2,0,1/2)
\]
modulo the maximal ideal, up to units.
Their first two successive flag minors are units.
At the conjugate place the roles are reversed, giving
\[
 (1/2,-u/2,1/2,u/2),\qquad
 (0,1/(2\alpha),0,1/(2\alpha)).
\]
Again the two minors are units.
The symplectic flags are therefore in the integral lower open
cell. Their unique lower unipotent representatives satisfy
\[
 n_{v,20}\equiv-1\text{ or }1,\qquad n_{v,21}\equiv0.
\]
The third step is the symplectic perpendicular of the first,
so these two minors suffice for the complete symplectic flag.
Together with the displayed residual matrix this is exactly
the partially split good stratum.

Let \(\vartheta\) be the required tower multiplier.
The residual packet gives
\(\vartheta/\mu\equiv1\) on \(G_F\).
Since squaring is an automorphism of \(1+\mathfrak m_E\)
for odd \(\ell\), there is a unique continuous
\(\delta:G_F\to1+\mathfrak m_E\) with
\(\delta^2=\vartheta/\mu\), after enlarging \(E\) to contain
the finitely many character values.
Set \(Q=\delta R|_{G_F}\).
It has multiplier \(\vartheta\), the same residual point,
and the same ordinary flags.
Its local torsion condition holds because cyclotomic images
have no \(\ell\)-power torsion, the other local diagonal
characters are unramified, and \(\delta\equiv1\).

For the global ramification condition, take an allowed inertia
element away from \(\ell\) with trivial prescribed residual
representation. Prime-to-\(\ell\) reduction gives
\(\xi=\psi=\varphi=1\) there.
Both rank-two blocks are \(1\) by their actual inertia type.
Thus \(R=1\), \(\vartheta/\mu=1\), and the unique square root
near \(1\) is \(\delta=1\).
So \(Q\) kills the required normal subgroup.

Finally put \(M=Lk_0\).
Solvable base change gives the cuspidal rank-two representation
over \(M\); its irreducible Galois representation rules out
cuspidal splitting.
The unequal ordinary weight pairs also rule out invariance
under \(\operatorname{Gal}(M/L)\).
Quadratic automorphic induction is therefore cuspidal over \(L\).
The union of weights is
\[
 m=(k+t,k,0,-t),
\]
with each consecutive gap at least \(2\), so this induction
is regular algebraic after the standard integral norm twist.
For example, under
\(r_\ell(\pi)=\operatorname{rec}_\ell
(\pi\otimes|\det|^{(1-n)/2})\),
one takes
\(\operatorname{AI}_{M/L}(\pi_M)\otimes|\det|\);
its attached representation is exactly \(R|_{G_L}\).
This checks the rank \(2\) to rank \(4\) normalization.

In the actual principal datum
\(\vartheta=\epsilon_\ell^{a+1}\omega^{-1}\xi\psi\).
Hence
\[
 \delta^2=\epsilon_\ell^{b}\,\omega^{-1}\psi\varphi^{-1},
 \qquad b=a+1-k,
\]
whose second factor has finite order prime to \(\ell\).
The squared weights of \(Q\) are \(2m_i+b\);
their consecutive gaps are at least \(4\).
Twisting \(Q|_{G_L}\) by \(\delta^{-1}\) removes the correction
and gives the actual attached automorphic \(R|_{G_L}\).
Moreover
\((\delta^{-2}\epsilon_\ell^b)^N=1\) for a positive \(N\),
and \(-b+(2m_0+b)\) is even.
These are exactly the finite-order and parity conditions
required of the seed's removable scalar twist.

All matrices and characters used are defined over a finite
coefficient field: \(r\) is, \(\alpha,\beta\) require only a
finite extension, the unique square root near \(1\) remains
there, and the local flag vectors are obtained by solving
finite linear equations with unit flag minors.
Thus the integral seed has the required finite coefficient
DVR, rather than only a completed algebraic-closure lattice.
\end{proof}

% Frozen-source correspondence (alpoge-source-7cbf7a0c):
% ModularityArgument/ThresholdArguments/DenseThresholdCases/
% MPropMHol.lean, THolLift.lean, SVRPropMV.lean, BirthMix.lean,
% MixSeedOfHolWC.lean, MixSeedOf.lean, SeedIIZ.lean, BirthSeedII.lean;
% BuildingsResidueChars/SymmetricSquareMixedExists.lean,
% ExistsRibetLattice.lean, GaloisCharacterResidualDecomposition.lean,
% ResidualCharactersPsiNe.lean, MixSignTwist.lean,
% InducedIrreducible.lean, InducedRepresentationConstituentVanishing.lean;
% Foundations/GaloisRepresentations/AdmissibleResidues/IndLoc0.lean.

\subsection{Local vanishing of the modular ideal at the actual seed}
\label{subsec:sr-seed-local-zero}

\begin{srlemma}[A faithful finite quotient and zero relative tangent]
\label{lem:sr-faithful-quotient-tangent}
Let \(W\) be a regular weight algebra in characteristic zero,
\(D\) a Noetherian deformation algebra, and \(J\subset D\).
Suppose \(D/J\) is finite and torsion-free over \(W\).
Let \(x:D\to E\), with \(E\) a finite coefficient field,
kill \(J\).
If every first-order deformation of \(x\) fixing its weight
is constant, then
\[
 J D_{\ker x}=0.
\]
In particular every prime \(\mathfrak q\subset\ker x\)
contains \(J\).
\end{srlemma}
\begin{proof}
Localize at the weight prime and at \(\ker x\), and complete.
After a finite separable residue-field extension, lifted
uniquely by finite etale extension of the complete weight
ring, the two residue fields agree.
Call the resulting complete weight ring \(W'\) and the
completed deformation ring \(S\).
The zero relative tangent condition says
\(\mathfrak m_S=\mathfrak m_{W'}S+\mathfrak m_S^2\).
Formal Nakayama, or successive approximation modulo powers of
the maximal ideal, makes \(W'\to S\) surjective.

The corresponding local factor \(T\) of the completion of
\(D/J\) is a nonzero finite torsion-free \(W'\)-module.
It is nonzero because \(x\) kills \(J\).
Torsion-freeness persists under completion and the finite
etale extension, and under passage to the local direct factor.
Since \(W'\) is a domain, \(W'\to T\) is injective.
But it factors as the surjection \(W'\to S\) followed by
\(S\to T\). The first map is therefore also injective, and
the second has zero kernel.
Thus the localized completed ideal is zero.
Faithful flatness of completion gives \(J D_{\ker x}=0\).

For \(j\in J\), choose \(s\notin\ker x\) with \(sj=0\).
If \(\mathfrak q\subset\ker x\), then \(s\notin\mathfrak q\),
so primality forces \(j\in\mathfrak q\).
\end{proof}

\begin{srproposition}[Seed local vanishing]
\label{prop:sr-seed-local-zero}
Let \(J_0\) be the actual unitary ordinary Hecke ideal for a
linked tower over \(L=FK\).
Use its finite torsion-free weight quotient and its classical
arithmetic-point property.
Let \(x\) be the arithmetic point obtained from the induced
Case II seed, after the harmless scalar twist making its
weights integral in the unitary normalization.
Then every prime below \(\ker x\) contains \(J_0\).
This remains true after transporting \(J_0\) to the common
coefficient algebra and changing that harmless scalar twist.
\end{srproposition}
\begin{proof}
The seed is an actual ordinary automorphic representation
with regular weights, the prescribed tame inertial type,
and unramified representation at the extra link place.
The classical arithmetic-point property therefore puts it
in the Hecke quotient: \(J_0\subset\ker x\).
This uses its actual automorphy and level, not a proposed
lifting theorem.

Consider a first-order deformation of \(x\) with weight fixed.
At an absolutely irreducible point, the determinant
deformation is represented by a matrix representation over
\(E[\tau]/\tau^2\), unique up to strict conjugation.
One elementary construction lifts matrix units in the
Cayley--Hamilton algebra: their reductions generate \(M_4(E)\)
by absolute irreducibility, and their multiplication errors
are corrected across a square-zero ideal.
The polarization relation lifts with them.
Write the resulting representation as
\[
 r'(g)=(1+\tau z(g))r(g).
\]
Its multiplication law is the adjoint cocycle equation.
The lifted conjugate-self-dual pairing extends this cocycle
to the descended polarized adjoint on \(G_F\).

The restriction at every \(\ell\)-place has an actual lifted
ordinary flag. To check this for an ordinary determinant
deformation, choose an inertia element after killing the
finite characters, with cyclotomic value of infinite order.
The distinct ordinary weights make its four eigenvalues
distinct. Their differences are units over
\(E[\tau]/\tau^2\), so its spectral idempotents lift.
The ordinary Cayley--Hamilton relations kill the lower
matrix entries between these ordered idempotents and
identify each diagonal character; they give the required
free flag. This is why an ordinary determinant relation
suffices here.
Fixing weight makes its diagonal inertial characters constant.
The consecutive seed gaps are at least \(4\) for the squared
weights, or at least \(2\) after the integral unitary
normalization. In particular no upper character ratio is
the cyclotomic character.
The positive-ordinary-extension lemma therefore makes the
first-order local representation crystalline after killing
its finite inertial characters.
Finite-index restriction is injective on \(H^1\) in
characteristic zero by averaging.
Thus its local cocycle lies in the Bloch--Kato finite
subspace before that extension too.

Away from \(\ell\), the deformation relations kill the
inertia subgroup on which the prescribed residual characters
are all trivial.
The seed's actual inertia representation is a direct sum of
those finite characters.
After this subgroup is killed, averaging over the remaining
finite quotient makes the inertial cocycle a coboundary.
At the additional unramified link place the same statement
is immediate. Hence the global cocycle belongs to
\[
 H^1_f\bigl(F,\operatorname{ad}_{\rm pol}(r|_{G_L})\bigr),
\]
the adjoint descended by its actual polarization.

All hypotheses of the published adjoint theorem are now
present: the seed over \(L\) is regular algebraic cuspidal,
its cyclotomic-kernel restriction is irreducible, and the
polarization just used is the one from its automorphic
conjugate self-duality.
Thorne's Theorem 6.1 gives vanishing of this descended
adjoint group over the totally real field \(F=L^+\).
It does not assert full adjoint vanishing over \(L\).
If the seed is first written with the correction character
\(\delta\), remove it using \(\delta^{-1}\), which extends
to \(G_F\); this leaves the adjoint unchanged and matches
the actual automorphic representation and polarization.
We conclude that \(z\) is a coboundary, so \(r'\) is strictly
conjugate to \(r\).

The ordered ordinary characters also have zero derivatives.
They cannot vary inside a constant representation: locally
they are pairwise distinct as characters, so the successive
invariant lines with the prescribed ordering are unique.
Equivalently, the infinitesimal flag deformation maps between
distinct character constituents and has no invariants.
Thus all characteristic-polynomial coefficients and all
ordered ordinary characters are constant.
They generate the universal ordinary determinant algebra;
the first-order map \(x\) is constant.

Apply Lemma~\ref{lem:sr-faithful-quotient-tangent} to the
actual finite torsion-free Hecke quotient.
This gives the local vanishing assertion.

For transport to the common space, a prime below the
seed's common-space point pulls back to a prime below \(x\),
hence contains \(J_0\).
Its image ideal is therefore contained in the original prime.
For a scalar change of normalization, twist the universal
determinant and ordinary characters by the corresponding
finite/algebraic character.
The same twist is an automorphism of the ordinary Hecke
module, with inverse the inverse-character twist.
It preserves its annihilator. Arithmetic-point density and
adic closedness identify the transported closed Hecke
ideals, proving the final assertion.
\end{proof}

\paragraph{Assembly with the two-link construction.}
The actual modular ideal used for Case II is the sum of the
images of two such unitary Hecke ideals, formed from links
with the same scalar normalization and two different
auxiliary unramified places.
Ordinary finite-weight control gives
\eqref{eq:sr-pm-finite}.
The rank-independent Taylor--Wiles support theorem applies
to its characteristic-\(\ell\) curves:
Proposition~\ref{prop:sr-four-curve-genericity} provides the
rank-four genericity, cyclotomic irreducibility, and regular
coset hypotheses; the local link condition gives the
additional unramified-place condition.
Proposition~\ref{prop:sr-seed-local-zero} gives the seed
property for each ideal, hence for their sum.
Arithmetic ordinary control gives purity and automorphy away
from the corresponding link place.
The two different link places remove the omitted-place
qualification, and the usual solvable descent gives the
automorphy statement over every allowed imaginary quadratic
restriction where it remains irreducible.
These are the five properties of the actual modular ideal
used in the Case II propagation, with every rank-four
specific input supplied above.
The remaining shared input in this assembly is the
rank-independent support patching theorem and its actual
unitary Hecke-module realization.

% Frozen-source map:
% ModularityArgument/ThresholdArguments/DenseThresholdCases/
% EigenbasisRankFour.lean, BirthPm.lean, HidaModularIdealMap.lean,
% PmTwo.lean, PmExTh.lean, PointTheoremBaseObject.lean;
% WeakShimuraTheory/EigenbasisLifts/RegOfTan.lean, Factor.lean,
% LiftRep.lean, RelForm.lean, CocRel.lean, CocExt.lean, FlagDual.lean,
% BkAbEmb.lean, AwAvg.lean, ConjOf.lean, ConstOf.lean;
% PatchingModules/AutomorphicDetection/PfromAut.lean;
% ModularityArgument/ThresholdArguments/DenseThresholdCases/
% THidaU.lean, ExistsHeckeIdeal.lean, UnitaryHeckeData.lean.

\subsection{The arithmetic Hecke quotient on the Case II tower}
\label{subsec:sr-case-two-hecke}
We identify the automorphic algebra in the rank-four support
calculation with the ordinary modules of
Section~\ref{sec:sr-ordinary-control}.
The coefficient lattice, local types and operator
normalizations must be the same on both sides.

Keep an actual Case II tower. Write \(F\) for its totally real
field, \(L=FK\) for its CM field, and \(S\) for the finite
ramification set, including the places above \(\ell\) and the
distinguished auxiliary place \(v_0\).
The tower conditions give
\[
 L/F\ \hbox{unramified at every finite place},\qquad
 v\in S\Longrightarrow v\ \hbox{splits in }L,\qquad
 [F:\mathbb Q]>1 .
 \tag{H1}
\]
The last assertion follows from the deep layer.
The definite rank-four unitary group is quasi-split at every
finite place: its global existence parity is satisfied since
\(4[F:\mathbb Q]/2=2[F:\mathbb Q]\) is even.

We use classical automorphic theorems in the following scope.
Base change of a classical representation on this definite
unitary group is an isobaric sum of discrete
conjugate-self-dual general-linear representations.
This is Labesse's Corollary 5.3 in
\emph{Changement de base CM et s\'eries discr\`etes}; see
\href{https://arxiv.org/pdf/0906.4189}{Guerberoff, Theorem 2.2}.
Its field, degree and finite quasi-splitness hypotheses
are (H1).
For the cuspidal seed's descent with specified local factors,
one can use
\href{https://arxiv.org/pdf/1409.3731}
{Kaletha--M\'inguez--Shin--White, Theorem 1.7.1}
in its proved generic-parameter, pure-inner-twist case
(the qualification immediately following the theorem).
Our definite hermitian group is a pure inner form.
The simple cuspidal parameter has trivial global component
group modulo the center. Thus one may choose its prescribed
split local members, its spherical nonsplit members, and
its compact cohomological archimedean members.
There is no remaining global multiplicity-character condition.
The existence parity was checked above. No general
nongeneric starred case of that theorem is used.
Attached semisimple Galois representations are supplied by
Guerberoff, Theorem 2.3.
For their cuspidal constituents we use the polarization and
full local--global compatibility of
\href{https://annals.math.princeton.edu/wp-content/uploads/annals-v179-n2-p03-p.pdf}
{Barnet--Lamb--Gee--Geraghty--Taylor, Theorem 2.1.1}.
At \(\ell\), potential semistability and compatibility can be
checked after killing the finite inducing characters, where
the local representation is Iwahori-spherical.
The ordinary local flag has the characters determined by
the unit and uniformizer equations below.
These inputs concern already automorphic representations;
none asserts automorphy of an arbitrary deformation.

\subsubsection{The local lattice and the weight normalization}
Let \(\mathcal O\) be a finite coefficient DVR, enlarged to contain
the needed finite character values, and let
\(\bar\chi_0,\ldots,\bar\chi_3\) be the link's residual
characters on \(G_L\).
At a split tame place \(u\), put
\[
                 \tau_u=\bigoplus_{i=0}^3
                       [\bar\chi_i]|_{I_{L_w}},
 \tag{H2}
\]
where brackets denote finite Teichm\"uller lifts.
This is a finite-image inertial type extending to the Weil group.

The local type used here must detect zero monodromy.
The Schneider--Zink type \(\sigma^{\rm crys}(\tau_u)\)
is finite dimensional and smooth on
\(\mathrm{GL}_4(\mathcal O_{L_w})\), and occurs in a smooth
irreducible \(\pi_w\) exactly when
\(\operatorname{rec}(\pi_w)|_I=\tau_u\) and \(N_{\rm WD}=0\).
A precise formulation is
\href{https://math.uchicago.edu/~emerton/pdffiles/bonn-lectures-survey-article.pdf}
{Emerton--Gee, Theorem 9.2.4}.
Use the contragredient in the function coefficient module:
\[
              \tau_0=\bigotimes_u
                              \sigma^{\rm crys}(\tau_u)^\vee .
 \tag{H3}
\]
Invariants of \(\tau_0\otimes\pi_w\) are then
\(\operatorname{Hom}(\sigma^{\rm crys}(\tau_u),\pi_w)\),
which fixes the dual convention.

There is a finite-free stable integral lattice in (H3)
over a finite coefficient extension.
Each compact representation factors through a finite quotient.
Its finitely many matrices and inverses are defined over
a finite extension; the sum of the finite orbit of a full
lattice is stable.
Tensor these lattices over the finitely many places.
No order of a finite quotient is inverted.

After the prescribed finite local torsion is fixed, the
inertial weight algebra is an actual power-series algebra
\[
 \Lambda\simeq
 \mathcal O[[X_{vij}:v\mid\ell,\ 0\le i<4,\ 1\le j\le d_v]],
           \qquad d_v=[L_w:\mathbb Q_\ell].
 \tag{H4}
\]
It represents the free pro-\(\ell\) unit quotients at
the selected \(w\mid v\).
Write \(\kappa_{vi}\) for its universal inertial characters.
Use the source's local reciprocity convention and define
the torus character \(\theta\) by
\[
 \kappa_{vi}(g)\epsilon_\ell(g)^i
                  \theta(e_{vi}(y))=1,\qquad
            \operatorname{Art}_v(g)=y\in\mathcal O_{L_w}^{\times}.
 \tag{H5}
\]
Here \(e_{vi}(y)\) has \(y\) in coordinate \(i\).
Local reciprocity makes (H5) well-defined and continuous.
The function model uses \(\theta^{-1}\).
Thus the factor \(\epsilon_\ell^i\) is part of the
normalization, and is retained.

For the commuting normalized ordinary operators, indexed
by \(0\le j<4\), the same convention gives
\[
       U_{vj}=\prod_{i=0}^{j}
                  \psi_{vi}(g_v)\epsilon_\ell(g_v)^i,
                  \qquad\operatorname{Art}_v(g_v)=\varpi_v .
 \tag{H6}
\]
Away from \(S\), the residual spherical values are specified by
\[
 \begin{split}
 P_w(X)&=\prod_{i=0}^3
          (X-\bar\chi_i(\operatorname{Frob}_w^{-1}))\\
       &=\sum_{j=0}^{4}(-1)^j
                  q_w^{j(j-1)/2}\bar a_{wj}X^{4-j},
                        \qquad\bar a_{w0}=1 .
 \end{split}
 \tag{H7}
\]
Since \(q_w\) is a unit in the residue field, these
\(\bar a_{wj}\) are uniquely determined.
The residual ordinary values are those of (H6).
Choose arbitrary integral lifts; the simultaneous residual
localization depends only on their reductions.

\subsubsection{Separation by classical specializations}
Let \(M\) be the corresponding residual direct summand of
the constructed ordinary function module, with lattice (H3)
and character (H5).
It is finite free over \(\Lambda\), by the ordinary projector
and coefficient-change proofs, also after the direct-summand
projection. If the residual system does not occur, set \(M=0\).
Otherwise let
\[
              \mathbb T\subseteq\operatorname{End}_\Lambda(M)
 \tag{H8}
\]
be the closed algebra generated by the selected spherical,
ordinary and scalar operators.
It is commutative: this is spherical and ordinary Iwahori
commutativity, and commutativity at disjoint places.
It is finite and torsion free over \(\Lambda\), because
it is a submodule of the finite endomorphism module over
the Noetherian domain \(\Lambda\).

\begin{srlemma}\label{lem:sr-unitary-classical-separation}
The evaluations of \(\mathbb T\) at actual classical ordinary
eigensystems separate its elements.
One may use a single fixed regular algebraic weight,
varying finite-order characters above \(\ell\), and require
the four finite inertial characters at each selected
place to be pairwise distinct.
Consequently \(\mathbb T\) is reduced.
\end{srlemma}
\begin{proof}
In the free coordinates (H4), fix the algebraic part.
Finite-order twists assign roots of unity of \(\ell\)-power
order to a basis of the free pro-\(\ell\) quotient,
after multiplication by the fixed characters for that
algebraic part.
They respect the prescribed finite torsion.
The resulting coordinate values, after subtracting \(1\),
have absolute value less than one.
Each assignment lies in a finite coefficient extension.

These assignments are Zariski dense for integral power series.
In one variable a nonzero element of \(\mathcal O[[X]]\)
has only finitely many zeros in the open unit disc.
Divide out its coefficient-content power of the uniformizer
and apply Weierstrass preparation at the first unit coefficient;
the remaining power-series factor is a unit on that disc.
It cannot vanish at every allowed torsion value.
For several variables, specialize all but one and induct.
The same proof permits finitely many values to be excluded
in each free coordinate.
As \(d_v>0\), it permits all equalities between the four
finite inertial characters to be excluded at each place.

At such a weight, coefficient and classical control identify
\(M\) with actual classical ordinary forms; a specialized
\(\Lambda\)-basis remains linearly independent.
The definite automorphic space is semisimple.
Its unramified operators act by their Satake characters.
At \(\ell\), distinct finite inertial characters exclude
linked inducing characters and monodromy between them.
The ordered Jacquet line has the specified unit and
uniformizer characters.
The commuting normalized ordinary operators act on that
line by one eigencharacter, and act by the same scalars
on repeated automorphic copies.
The selected commuting algebra is therefore a sum of its
eigenspaces on this classical ordinary space.

If \(t\in\mathbb T\) vanishes at every such eigencharacter,
it acts as zero in every such specialization of \(M\).
Each entry of its action matrix in a \(\Lambda\)-basis
vanishes on the dense set just proved, hence is zero.
The faithful inclusion (H8) gives \(t=0\).
A nilpotent element vanishes at every field-valued character,
so \(\mathbb T\) is reduced.
\end{proof}

\subsubsection{Integral interpolation at the prime three}
\begin{srlemma}\label{lem:sr-compact-determinant-interpolation}
Let \(T\) be a compact Hausdorff commutative ring, and
\(x_\alpha:T\to E_\alpha\) a separating family of continuous
maps to characteristic-zero fields.
Let \(\Gamma\) be profinite and
\(r_\alpha:\Gamma\to\mathrm{GL}_n(E_\alpha)\) continuous.
Suppose that on a dense union of conjugacy classes their
characteristic polynomials are specializations of
polynomials in \(T[X]\).
There is then a unique continuous \(n\)-dimensional
determinant on \(T[\Gamma]\) with these specializations.
No integer up to \(n!\) is inverted.
\end{srlemma}
\begin{proof}
The injection \(T\to\prod_\alpha E_\alpha\) is a
homeomorphism onto a closed subring, by compactness and
the Hausdorff property.
For \(g\in\Gamma\), its tuple of characteristic coefficients
is a limit of the corresponding tuples on the dense set.
Conjugation leaves these unchanged.
Thus they belong to \(T\), and their maps \(\Gamma\to T\)
are continuous.

Characteristic polynomials suffice without dividing power
sums by \(1,\ldots,n\).
For finitely many matrices \(A_1,\ldots,A_s\), the integral
Amitsur identity is
\[
 \det(1-\textstyle\sum_i t_iA_i)
       =\prod_{[w]\ {\rm primitive}}\det(1-t_wA_w).
 \tag{H9}
\]
The product is over primitive cyclic words, and \(A_w,t_w\)
are the matrix word and its commuting monomial.
In bounded degree only finitely many factors contribute.
For generic matrices over \(\mathbb Z\), prove this over
\(\mathbb Q\) by formal logarithms and grouping traces
into cyclic words.
Both sides have integral coefficients, so this proves
an identity over \(\mathbb Z\), valid after arbitrary
scalar extension.

Apply (H9) to \(A_i=r_\alpha(g_i)\).
Every factor is a characteristic polynomial of a group
word, whose coefficients already belong to \(T\).
Hence every coefficient on the left belongs to \(T\).
This defines the degree-\(n\) polynomial law on finite
sums in \(T[\Gamma]\).
Multiplicativity, the unit and scalar-extension identities
hold after every separating \(x_\alpha\), and therefore
over \(T\). The same argument proves uniqueness.
The characteristic-coefficient continuity above gives
its profinite continuity.
\end{proof}

Apply this lemma to (H8) and
Lemma~\ref{lem:sr-unitary-classical-separation}.
The finite algebra \(\mathbb T\) over the compact
power-series ring \(\Lambda\) is compact.
Its arithmetic maps are continuous and lie in finite
coefficient extensions, by the automatic-continuity
argument following classical control.
The attached classical semisimple representations give
the \(r_\alpha\).
Chebotarev supplies the dense conjugacy classes and
(H7), with Hecke operators, supplies their polynomials.
This constructs an actual continuous determinant
\(D_{\mathbb T}\).

We record the normalization for arbitrary discrete constituents,
including odd-dimensional ones.
This is the auxiliary-character construction in the proof
of Guerberoff's realization theorem, made explicit in rank four.
Write the discrete base change as
\(\boxplus_i\Pi_i\), of ranks \(n_i\), with
\(\sum_i n_i=4\).
Choose a unitary anticyclotomic Hecke character \(\eta_i\),
unramified at any prescribed finite set including the
places above \(\ell\), whose infinity type is
\[
 \eta_{i,\tau}(z)=(z/|z|)^{\delta_{i,\tau}},
 \quad
 \delta_{i,\tau}=0\ \hbox{if }4-n_i\ \hbox{is even},
 \quad
 \delta_{i,\tau}\in\{1,-1\}\ \hbox{otherwise}.
 \tag{H10a}
\]
The signs at conjugate embeddings are opposite.
The CM Hecke-character construction used in that theorem
supplies these characters with the stated finite
unramifiedness; for even \(n_i\) take \(\eta_i=1\).
Then \(\Pi_i\eta_i\) is cohomological.
The discrete-spectrum classification writes
\[
       \Pi_i\eta_i
             =\boxplus_{j=0}^{b_i-1}
                         \pi_i^0|\det|^j,\qquad n_i=a_i b_i,
\]
where \(\pi_i^0\) is cuspidal of rank \(a_i\) and
\(\pi_i^0|\det|^{(b_i-1)/2}\) is conjugate self-dual.

Choose a second unitary anticyclotomic character \(\xi_i\),
again unramified at the prescribed finite set, with angular
infinity exponent \(0\) if \(a_i\) or \(b_i\) is odd,
and exponent \(1\) or \(-1\) if both are even.
Then
\[
       C_i=\pi_i^0\xi_i|\det|^{(b_i-1)/2}
\]
is cuspidal, cohomological and conjugate self-dual.
The two Hecke characters
\[
 A_i=\xi_i^{-1}|\cdot|^{(a_i-1)(b_i-1)/2},
       \qquad
 B_i=\eta_i^{-1}|\cdot|^{(n_i-4)/2}
 \tag{H10b}
\]
are \emph{algebraic}.
Indeed an angular factor \((z/|z|)^\delta\) has
exponents \(\delta/2,-\delta/2\), while the normalized
complex norm has exponents \(1,1\).
For \(A_i\), the integer
\((a_i-1)(b_i-1)\) is odd exactly when both \(a_i,b_i\)
are even, the case where the angular exponent is odd.
For \(B_i\), \(n_i-4\) is odd exactly when the angular
exponent in (H10a) is odd.
Both pairs of resulting exponents are therefore integers.
Only these algebraic products are converted into
\(\ell\)-adic characters.

Let \(r_{\ell,\iota}(C_i)\) use its rank-\(a_i\)
normalization, and set
\[
 R_i=r_{\ell,\iota}(C_i)\,
            \epsilon_\ell^{\,a_i-n_i}\,
            r_{\ell,\iota}(A_i)\,
            r_{\ell,\iota}(B_i).
 \tag{H10c}
\]
All powers of \(\epsilon_\ell\) in this formula are integers.
At a compatible finite place, the norm exponent in its
parameter is
\[
 \frac{b_i-a_i}{2}+(a_i-n_i)
       +\frac{(a_i-1)(b_i-1)}{2}
       +\frac{n_i-4}{2}=-\frac32 .
\]
The factors \(\xi_i\) cancel and the remaining character
is \(\eta_i^{-1}\).
Thus \(R_i\) has parameter
\(\operatorname{rec}(\pi_i^0\eta_i^{-1}|\det|^{-3/2})\).
Consequently the total attached representation is
\(\bigoplus_i\bigoplus_{j=0}^{b_i-1}R_i\epsilon_\ell^j\).
Putting \(\rho_i'=R_i\epsilon_\ell^{b_i-1}\) and reversing
the index gives exactly the form used below:
\[
             \bigoplus_i\ \bigoplus_{j=0}^{b_i-1}
                                  \rho_i'\epsilon_\ell^{-j}.
 \tag{H10}
\]
This proves the decomposition with actual continuous
finite-coefficient Galois representations.
It does not choose a square root of \(\epsilon_\ell\).
The polarization and Hodge numbers are those of the
definite-unitary realization; the auxiliary characters
are untwisted by the algebraic products (H10b).

The ordinary characters over \(\mathbb T\) are explicit.
Set \(U_{v,-1}=1\); prescribe them on units by (H5)
and on a uniformizer by
\[
        \psi_{vi}(g_v)\epsilon_\ell(g_v)^i
                                  =U_{vi}U_{v,i-1}^{-1}.
 \tag{H11}
\]
The ordinary operators are units in the complete local
algebra \(\mathbb T\).
Local reciprocity extends these data continuously:
the local multiplicative group is the product of its
units and the cyclic uniformizer group, whose map
extends to its profinite completion since
\(\mathbb T^\times\) is profinite.
At each separating point these are the actual ordered
Galois characters, by the two local equations and the
ordinary flag.

At the separating classical points the Galois local flags
used above do not follow from semisimplified compatibility
alone. Here is the additional check.
Distinct finite inertial characters make the smooth local
factor at \(\ell\) an irreducible principal series with
zero-monodromy parameter.
After a finite abelian extension killing those characters
its parameter is unramified.
Apply full compatibility and the crystalline assertion of
the cuspidal Galois theorem, after a compatible solvable
base change, and then restrict back.
This proves potential crystallinity and the asserted
finite inertial type, with zero monodromy.

Over that local extension, order the Hodge numbers
\(h_{\sigma0}<\cdots<h_{\sigma3}\) by the ordered ordinary
characters. The normalized unit equations for \(U_{vj}\)
say that the Newton number of the first \(j\) parameter
lines is
\(\sum_\sigma\sum_{i<j}h_{\sigma i}\).
For any \(j\)-dimensional filtered subspace this is the
minimum possible Hodge number.
Weak admissibility gives the opposite inequality, so
equality holds on these parameter subspaces.
They are weakly admissible and yield Galois subrepresentations.
Their successive rank-one quotients are the characters
of (H5), (H11).
Their distinct slopes make this flag canonical after the
chosen ordering, so it descends.
This supplies the actual ordinary flag, including its
unit and uniformizer values.
Both Galois-realization inputs used here apply at every
odd coefficient prime; in particular rank four at
\(\ell=3\) is included.
\subsubsection{The quotient and its arithmetic points}
Let \(B\) be the universal ordinary polarized determinant
algebra for the same residual link and local characters.

\begin{srproposition}\label{prop:sr-case-two-hecke-quotient}
There is a continuous surjection
\[
                             B\longrightarrow\mathbb T
 \tag{H12}
\]
compatible with \(\Lambda\), the Galois determinant and
the ordered ordinary characters.
Its kernel \(J_0\) is the actual automorphic ideal.
The quotient \(B/J_0\) is finite and torsion free over
\(\Lambda\), with dense classical arithmetic points.

The normalized Case II seed defines an arithmetic point
\(x\) of \(B\) with \(J_0\subseteq\ker x\).
Every regular arithmetic point of \(B/J_0\) is realized
by actual classical definite-unitary forms with the
specified semisimple Galois determinant.
If that Galois representation is irreducible, its
general-linear base change is cuspidal.
\end{srproposition}
\begin{proof}
At the separating classical points, the attached
representations have the specified unitary polarization
and ordered ordinary flags (H5), (H11).
At tame type places they have inertia (H2) and zero
monodromy, by (H3) and full local--global compatibility.
They are unramified outside the specified ramification set.
Thus the polarization, ordinary Cayley--Hamilton and
killed-inertia relations hold at every separating point.
They are polynomial determinant relations, so they hold
over \(\mathbb T\).
Equations (H7), (H6), Chebotarev and the integral identity
(H9) give the specified residual determinant and local
characters. Universality gives (H12).

Its image contains \(\Lambda\), all spherical operators
by (H7), and all ordinary operators by (H6).
These topologically generate \(\mathbb T\).
The image is closed by compactness of the source,
so (H12) is surjective.
Finiteness, torsion freeness and density follow from
(H8) and Lemma~\ref{lem:sr-unitary-classical-separation}.
If the residual module is zero, set \(J_0=B\).
The seed argument next shows that this alternative
does not occur for the actual seed link.

Use the seed's actual cuspidal automorphic representation
from Proposition~\ref{prop:sr-induced-case-two-seed},
with the normalization of
Proposition~\ref{prop:sr-seed-local-zero}.
Its prescribed regular weights, tame inertial characters
and nonexceptional local ordinary flag were checked there.
Its ordinary local Jacquet vector has exactly the unit
and uniformizer values (H5)--(H6).
The auxiliary \(v_0\) was chosen away from its ramification;
the seed and link-normalizing character are unramified
there. Its split local factor therefore has a vector fixed
by the first principal congruence subgroup.
At tame places (H3) supplies the required coefficient
vectors; at nonsplit places it is unramified.

The proved generic pure-inner-form descent described
after (H1) now gives a nonzero vector with this exact
weight, tame type and auxiliary level in the classical
definite-unitary space.
Classical control places its eigensystem in the
specialization of the constructed module.
Its residual values are (H7), (H6), so it belongs to \(M\).
It gives a character \(\mathbb T\to E_x\) whose composite
with (H12) equals \(x\):
the weight characters, Frobenius polynomials and ordered
ordinary characters agree, and they generate \(B\).
Hence \(J_0\subseteq\ker x\).

Finally \(M\) is a faithful finite \(\mathbb T\)-module.
Every prime of \(\mathbb T\) therefore belongs to its
support. Localize at a characteristic-zero point and
apply Nakayama to see that the corresponding fiber of
\(M\) is nonzero.
Classical control gives actual classical forms.
A common eigenvector for the commuting finite-dimensional
algebra in this fiber realizes the point.
Its attached semisimple Galois determinant equals the
given determinant by Frobenius compatibility and
Chebotarev. If this representation is irreducible,
(H10) has one summand and its length is one.
The base change is then cuspidal.
\end{proof}

The point property just proved uses actual automorphy,
the matched tame type, the nonexceptional ordered local
flag and unramifiedness at \(v_0\).
Those are the properties checked for the constructed seed.

\subsubsection{Speh strings and the required spectral bound}
The spectral assertion used at this endpoint is the
strict bound called \(\mathrm{PureW}\):
\[
               |\iota(\alpha)|<
                    \operatorname{Norm}(w)|\iota(\beta)|
 \tag{H13}
\]
for any two nonzero Frobenius roots, off finitely many
places. It is weaker than equality of their absolute
values and suffices for the later block exclusion.

\begin{srlemma}\label{lem:sr-case-two-speh-exclusion}
At an arithmetic point on the actual link quotient
\(B/J_0\), every length \(b_i\) in (H10) is one.
Bound (H13) follows.
Undoing the scalar link twist gives the
\(\mathrm{PureW}\) condition on the original rational
rank-four representation.
\end{srlemma}
\begin{proof}
The exact residual input is
\[
          \bar\chi_i\ne\bar\chi_j\bar\epsilon_\ell
                       \quad\hbox{on }G_L\quad(0\le i,j<4).
 \tag{H14}
\]
It includes \(i=j\).
On the tower this follows from
\[
 \forall w\in\mathbb Z^2\ \exists g\in G_L:\quad
       \overline{(\chi_0/\chi_1)^{w_1}
                 (\chi_0/\chi_2)^{w_2}}(g)
                              \ne\bar\epsilon_\ell(g).
 \tag{H15}
\]
The paired relation \(\chi_0\chi_3=\chi_1\chi_2\)
expresses every \(\chi_i/\chi_j\) by its root weight
\(w=w(i,j)\). A scalar link twist cancels.
The tower condition \(\mathrm{Ok.g2}\) is exactly
(H15). The subgroup of \(G_F\) fixing \(L\) is
canonically \(G_L\), so its witnesses belong to the
required group.

Choose stable lattices in the summands of (H10).
Their direct sum has residual semisimplification
\(\bigoplus_i\bar\chi_i\), independently of the lattice:
the characteristic polynomials agree with those of
the actual determinant, and characteristic-polynomial
Brauer--Nesbitt applies.
If \(b_i\ge2\), a nonzero Jordan--H\"older
constituent of the reduction of \(\rho_i'\) is
some \(\bar\chi_s\).
The next summand
\(\rho_i'\epsilon_\ell^{-1}\) makes
\(\bar\chi_s\bar\epsilon_\ell^{-1}\) another constituent,
say \(\bar\chi_t\).
This contradicts (H14). All \(b_i\) are therefore one.

The remaining constituents are unitary cuspidal
general-linear representations. Their unramified
local components are unitary and generic.
The generic part of the
\href{https://www.numdam.org/item/ASENS_1986_4_19_3_335_0/}
{Tadi\'c unitary classification} expresses them using
tempered factors and complementary pairs with
exponents \(s,-s\), \(0<s<1/2\).
At an unramified place all factors are spherical;
the Satake absolute values lie strictly between
\(q_w^{-1/2}\) and \(q_w^{1/2}\).
This holds for all constituents at once, so every
two parameters have ratio strictly less than \(q_w\).
The common rank-four Galois normalization multiplies
them all by the same scalar and preserves this ratio.
This proves (H13).

A scalar link twist also cancels from the ratio.
For a rational Frobenius power lying in \(G_L\),
let \(f_w\) be the residue degree of the place below it.
A rational exponent \(f\) is a multiple of \(f_w\);
its eigenvalues are the \(f/f_w\)-th powers of
the local Frobenius eigenvalues.
Thus the bound \(<q_w\) becomes \(<p^f\).
Exclude the finitely many ramified places and take the
open subgroup \(G_L\). This is the required
\(\mathrm{PureW}\) condition.
\end{proof}

\subsubsection{One module on the two links}
The shared arithmetic Taylor--Wiles calculation has an
exact role in this match.
On a selected nice characteristic-\(\ell\) curve, its
finite-level Selmer inputs are the first- and second-Selmer bounds
of Section~\ref{sec:sr-arithmetic-selmer-producers}, on the actual auxiliary
Galois groups with the actual ordinary and polarized
local conditions. Their constants precede precision
and auxiliary level.
That shared arithmetic calculation proves these estimates and
constructs the actual auxiliary groups and primes, under the nice-point
conditions displayed there. The link must supply those same conditions,
including its distinguished local \srcid{V0ok} condition.
The modules to which they apply are precisely the
ordinary function modules with lattice (H3),
character (H5), residual system (H7) and varying
split auxiliary local levels.
One module system is used along the prime chain.

The first finite-level hypothesis of
Proposition~\ref{prop:sr-one-patched-module}
is ordinary auxiliary projectivity and truncated
freeness.
Its second is given by the explicit old/new and
trace maps of
Proposition~\ref{prop:concrete-two-maps},
with the same normalized auxiliary discriminant.
For the third, use the actual Galois determinant
constructed above at each auxiliary level.
Its polarization, ordinary, inertia and auxiliary
relations are the relations of \(B\).
The explicit torus calculation gives nilpotence
exponent \(n!=24\) in rank four.
All maps act on the same free tame coefficient
lattice and commute with its tame action, so the
lattice introduces no new denominator.

Consequently the three finite-level hypotheses hold
for this actual system once the stated arithmetic
Selmer bounds and auxiliary construction are invoked.
Proposition~\ref{prop:tw-limit-coh} supplies the local
presentation estimates, and
Propositions~\ref{prop:sr-one-patched-module} and
\ref{prop:sr-completed-support-chain} give a common
module, its completed supports and their transport.
This is the rank-four application of those shared proofs.

Choose the two links with the same scalar character
\(\kappa\) and distinct auxiliary primes.
For their universal ordinary determinant algebras
\(B,B'\), take the actual kernels \(J_0,J'_0\)
in (H12). On their common paired space \(X\), put
\[
                        J=\phi(J_0)X+\phi'(J'_0)X.
 \tag{H16}
\]
The finite weight-zero property follows from
finiteness of \(B/J_0\) over \(\Lambda\).
Indeed take a characteristic-\(\ell\) domain
quotient after setting the fixed inertial weight
characters to \(1\).
The image of \(\Lambda\) is finite, so the link
Hecke quotient is algebraic over a finite field.
Thus all characteristic coefficients on \(G_L\)
are algebraic constants.
Every element of \(G_F\) has square in \(G_L\).
Its eigenvalues and characteristic coefficients
are algebraic constants as well.
The complete local domain has the original finite
residue field. Reduction is injective on every finite
subfield, so any algebraic constant in it is the unique
corresponding element of that residue field.
All characteristic coefficients thus equal their
residual values.
These coefficient differences and the inertial weight
differences topologically generate the maximal ideal
of the paired trace algebra. Formal Nakayama, using
Noetherianity, makes that maximal ideal zero in the
domain quotient. Every such prime is therefore maximal.
A power of the maximal ideal is consequently contained
in the fixed-weight ideal, as required in
Proposition~\ref{prop:sr-four-core}.

On each nice curve the common module, the local
presentations and
Proposition~\ref{prop:sr-four-curve-genericity}
give the curve-local support assertion.
At the seed,
Proposition~\ref{prop:sr-case-two-hecke-quotient}
gives \(J_0x=J'_0x=0\), and
Proposition~\ref{prop:sr-seed-local-zero}
gives local vanishing of their two mapped ideals.
The residual chain and
Proposition~\ref{prop:sr-four-core} then transport
this support to the original paired point.
The characteristic-\(\ell\) exclusion is
Lemma~\ref{lem:sr-four-characteristic-ell};
the characteristic-zero spectral assertion is
Lemma~\ref{lem:sr-case-two-speh-exclusion}.

At a regular irreducible arithmetic point, classical
control gives an automorphic representation with
the prescribed local factors away from the
distinguished auxiliary place.
Both links describe the same global general-linear
representation: their scalar twist is the same,
and their Frobenius polynomials agree away from
their two auxiliary places.
Strong multiplicity one identifies their cuspidal
representations. At any one finite place at
least one link has no auxiliary exception.
Its local--global comparison therefore gives
the desired local statement there, including
at the other link's auxiliary place.
The level introduced for smallness has thereby
been removed from the actual endpoint.

This identifies the Case II Hecke and support
objects with the shared ordinary construction.
The arithmetic production of those Selmer bounds and
of the actual auxiliary groups and prime-selection limit is supplied
by Section~\ref{sec:sr-arithmetic-selmer-producers}; the matching nice-point
and common-space conditions remain the ones explicitly used above.
The further passage from automorphy over \(L\)
to the source's \(\mathrm{IsAutQA}\), including
its \(\mathrm{WitA}\) local level data over an
imaginary quadratic field, is the separate
solvable-descent and local-compatibility step.
It is not inferred solely from the existence
of a Hecke eigensystem.

% Exact frozen-source correspondence:
% DenseThresholdCases/HdEx, TTypeEx, TThetaEx, TResEx,
% UnitaryHeckeData, TDenseU, TCoreSep, TCoreDet, OrdCoreH,
% ExistsHeckeIdeal, TJRing, TJDense, THidaU, TLevelU, THasAutOf;
% PointTheoremBaseObject, BirthPm, TPmFin, PmTwo,
% Gsp4PromodularityCaseTwo, VWinPm, PmAutO, DnPmDown.
% NoOmLift and MemMax in
% WeakShimuraTheory/CharacteristicPolynomialRoots.

\section{The actual rank-four tower and automorphic endpoint}
\label{sec:sr-rank-four-actual-endpoint}

\subsection{The character witnesses on the principal Case II tower}
\label{subsec:sr-case-two-tower-witnesses}
We align the three tower conditions used above with their
actual principal-packet producers.
Here \(\ell\) is odd, \(\ell\nmid N\), and the nontrivial
principal character \(\psi\) has \(2\)-power order and
\(\psi(-1)=(-1)^a\).
Write \(\omega\) for the Teichm\"uller cyclotomic character,
\(\xi\) for the chosen primitive odd auxiliary character
of conductor \(2^n\), and
\[
            \bar\Psi=\bar\omega^{\,a}\bar\psi .
 \tag{T1}
\]
The actual residual diagonal, before the possible
ordinary permutation, is
\[
                  (\bar\xi\bar\Psi,\bar\Psi,\bar\xi,1).
 \tag{T2}
\]
Reduction is injective on all \(2\)-power character values.

Choose an exponent \(s\) with \(\psi^{2^s}=1\), as in
the source's \(\mathrm{psi\_exp}\), and let \(2^c\)
be the conductor of its dyadic component.
The actual choice \(n=\mathrm{nII}(\ell,\psi)\)
satisfies the two numerical lower bounds
\[
       2^n\ge\ell^2 2^{s+2},
                   \qquad 2^n\ge2^{c+6}.
 \tag{T3}
\]
These are precisely the two inequalities extracted from
that definition in
\srcid{ExistsDecompositionElementAtLXiPowFourNeOne.lean}
and
\srcid{ExistsInertiaElementAtTwoXiPowFourNeOne.lean}.
In \(\mathrm{psi\_exp}\) the annihilator is written
\(2^s u\) with \(u\) odd; since our packet has \(2\)-power
order, \(\psi^{2^s}=1\) follows.
No \(\ell\)-primary part of \(\psi\) needs to be killed
in this specialization.

\begin{srlemma}\label{lem:sr-primitive-dyadic-frobenius}
Let \(\xi\) be primitive of conductor \(2^n\), \(n\ge3\).
For an odd integer \(z>1\), if
\(\xi(z)^{2^b u}=1\) with \(u\) odd, then
\[
                              2^n<z^2 2^b .
 \tag{T4}
\]
\end{srlemma}
\begin{proof}
The character is faithful on
\((1+4\mathbb Z_2)/(1+2^n\mathbb Z_2)\).
Its value has \(2\)-power order, so the odd exponent
\(u\) can be removed.
For \(b=0\), squaring gives
\(z^2\equiv1\pmod{2^n}\), hence
\(2^n\le z^2-1<z^2\).
For \(b\ge1\), faithfulness gives
\(z^{2^b}\equiv1\pmod{2^n}\).
The elementary lifting-the-exponent identity yields
\[
 v_2(z^{2^b}-1)=v_2(z^2-1)+b-1.
\]
Consequently
\(2^n\le2^{b-1}(z^2-1)<2^b z^2\).
\end{proof}

At \(\ell\), choose the local reciprocity lift
\(\sigma_\ell\) of a uniformizer with
\(\epsilon_\ell(\sigma_\ell)=1\).
The two finite characters are unramified there, and
their values are \(\psi(\ell),\xi(\ell)\).
Then
\[
       g_\ell=\sigma_\ell^{2^s}\in D_\ell,\qquad
       \omega(g_\ell)=1,\qquad\psi(g_\ell)=1,\qquad
                              \bar\xi(g_\ell)^4\ne1 .
 \tag{T5}
\]
The last assertion follows from (T3)--(T4), with \(z=\ell\)
and \(b=s+2\).
At \(2\), choose an inertia element whose cyclotomic
projection modulo \(2^n\) is \(3\).
Raise it to \(2^c\), obtaining an element \(u\) with
\[
       u\in I_2,\quad \omega(u)=1,\quad\psi(u)=1,\quad
                                   \bar\xi(u)^4\ne1 .
 \tag{T6}
\]
The dyadic part of \(\psi\) is killed, and its other
components are unramified at \(2\).
If the final assertion failed, (T4) with \(z=3\) would
give \(2^n<9\cdot2^{c+2}\), contradicting
\(2^n\ge16\cdot2^{c+2}\).
These arguments also work after an odd
\(\ell\)-power is inserted in the exponents.
They are the actual local witnesses used by the source.

\subsubsection{The finite abelian base field}
Choose three distinct fresh odd primes, outside the
original ramification set \(S\), and quadratic characters
\(\phi,\eta,\eta'\), respectively ramified only at those
primes, with
\[
 \phi(-1)=1,\qquad \eta(-1)=\eta'(-1)=-1 ,
\]
all split at \(S\), and with \(\eta,\eta'\) also split
at the earlier chosen fresh primes.
Such characters are obtained by choosing their prime
conductors in arithmetic progressions.
For the even character choose the conductor prime
\(1\bmod8\); for an odd one choose it \(7\bmod8\).
At each prescribed odd prime require the signed conductor
to be a square modulo that prime.
The Chinese remainder theorem gives a reduced progression
and Dirichlet's theorem supplies a fresh prime in it.
The quadratic characters then have exactly the asserted
ramification and splitting.

Put
\[
 H_0=\{g:\bar\Psi(g)=\phi(g),\
                   \eta(g)=\eta'(g)\},\quad
 F_0=\overline{\mathbb Q}^{\,H_0},\quad
 K=\overline{\mathbb Q}^{\,\ker\eta},\quad L_0=F_0K .
 \tag{T7}
\]
All characters in (T7) are finite and continuous.
Their common kernel contains the commutator, so \(F_0\)
is finite abelian.
Complex conjugation belongs to \(H_0\), by
\(\bar\Psi(c)=1\), \(\phi(c)=1\), and
\(\eta(c)=\eta'(c)=-1\).
Thus \(F_0\) is totally real; \(K\) is imaginary quadratic.
Moreover \(H_0\subset\ker(\eta\eta')\), and
\(\eta\eta'\) is a nontrivial real quadratic character
(the two conductors are distinct).
Hence \([F_0:\mathbb Q]\) is even.

The CM extension \(L_0/F_0\) is unramified at every
finite place.
Away from the conductor of \(\eta\), its inertia is
already killed by \(\eta\).
At that conductor, \(\eta'\) is unramified, and the
condition \(\eta=\eta'\) on \(H_0\) kills the remaining
inertia.
Every prime in \(S\) splits in \(L_0/F_0\), since
\(\eta\) is trivial on its rational decomposition group.
At \(\ell\), \(\xi\) is unramified and the relation
\(\bar\Psi=\phi\) on \(H_0\) makes \(\bar\Psi\)
trivial on inertia, since the fresh \(\phi\) is
unramified there.
Thus (T2) is trivial on inertia over \(F_0\).
The extension \(F_0/\mathbb Q\) is tame at \(\ell\):
its inertia factors through the Teichm\"uller character,
whose order divides \(\ell-1\).

Multiplication by inertia at the three fresh conductor
primes changes \(\phi,\eta,\eta'\) independently while
leaving \(\xi,\Psi,\omega,\psi\) unchanged.
Indeed the latter characters are unramified at those
primes, and each of the other two auxiliary quadratic
characters is unramified at the prime used for a change.
This gives the explicit adjustment operation
\(\mathrm{adj2}\) in the source.

The element \(u\) of (T6) already lies in \(G_{L_0}\):
all three fresh quadratic characters are unramified
at \(2\), and \(\bar\Psi(u)=1\).
Apply the adjustment operation to complex conjugation,
setting \(\phi=1,\eta=\eta'=1\).
It gives an element \(v\in G_{L_0}\) with
\[
       \bar\Psi(v)=1,\qquad\bar\xi(v)=-1,\qquad
                                      \bar\omega(v)=-1 .
 \tag{T8}
\]
Since \(\psi\) is nontrivial of \(2\)-power order,
its image contains \(-1\).
The conductor of \(\psi\) is prime to \(\ell\);
Chinese remaindering in the cyclotomic quotient
therefore supplies an element with
\(\bar\psi=-1,\bar\omega=1\), and hence
\(\bar\Psi=-1\).
Adjust it to \(\phi=-1,\eta=\eta'=1\).
It lies in \(G_{L_0}\).
If its \(\bar\xi\)-value has square \(1\), multiply
it by \(u\).
Thus there is \(t\in G_{L_0}\) with
\[
       \bar\Psi(t)=-1,\qquad\bar\xi(t)^2\ne1,\qquad
                                      \bar\omega(t)=1 .
 \tag{T9}
\]

The witnesses \(u,v,t\) give all the character-separation
conditions used in the preceding subsections.
Every root ratio is
\(\bar\xi^x\bar\Psi^y\), with
\[
 (x,y)\in\{(\pm1,0),(0,\pm1),(\pm1,\pm1)\}.
\]
If \(x\ne0\), its value at \(u\) is nontrivial.
If \(x=0\), its value at \(t\) is \(-1\).
Hence every root character is nontrivial on \(G_{L_0}\);
this is \(\mathrm{Ok.g3}\).
These same witnesses belong to \(\ker\bar\omega\), so
the root characters remain nontrivial on that kernel.

More strongly,
\(\bar\omega\) is not in the subgroup generated by
all root ratios.
If
\(\bar\omega=\bar\xi^x\bar\Psi^y\) on \(G_{L_0}\),
evaluation at \(u\) makes \(x\) even, because
\(\bar\xi(u)\) has \(2\)-power order greater than four.
Evaluation at \(v\) then gives \(1=-1\), a contradiction.
This is \(\mathrm{Ok.g2}\) and supplies the premise of
Proposition~\ref{prop:sr-four-actual-scalar-place}.

The element \(t\) also has four distinct values in (T2):
they are
\(-\bar\xi(t),-1,\bar\xi(t),1\).
It can therefore be chosen as the frame's distinguished
regular element.
Finally, a character trivial on \(\ker\bar\omega\)
and permuting the four residual characters must be
trivial.
Compare its product with the character \(1\) in (T2);
it is a ratio of two of those characters.
Every nontrivial such ratio is nontrivial on
\(\ker\bar\omega\), as proved using \(u,t\).
This proves in particular the quadratic self-twist
exclusion \(\mathrm{Ok.g3'}\).

\subsubsection{The \(\ell\)-adic local witness and higher layers}
All three auxiliary quadratic characters split at \(\ell\).
The element \(g_\ell\) of (T5) therefore belongs to
\(G_{F_0}\), and to \(G_{L_0}\), while retaining its
rational local decomposition group.
Its values give
\[
          \bar\omega(g_\ell)=1,\qquad
          \bar\Psi(g_\ell)=1,\qquad
                                  \bar\xi(g_\ell)\ne1 .
 \tag{T10}
\]
These are exactly the second alternative of (D11),
hence the actual \(\mathrm{LocL}\) producer.
They are stronger than just a generic residual
non-equality of two characters.

For the larger dimension required by the support
calculation, adjoin successive finite layers of the real
cyclotomic \(\mathbb Z_\ell\)-extension.
Write \(F_m\) for the resulting real fields and
\(L_m=F_mK\).
The degrees \([F_m:\mathbb Q]\) grow without bound;
the extensions \(F_m/F_0\) have \(\ell\)-power degree.
Their finite unramified CM property and their splitting
at \(S\) follow by base change from \(L_0/F_0\).
The degree remains even.

In any such layer an element of \(G_{F_0}\) has an
odd \(\ell\)-power lying in \(G_{F_m}\).
Apply a common such power to \(u,v,t\).
Their \(\bar\xi,\bar\Psi\)-values have \(2\)-power order,
so all the displayed distinctions, signs and
root-separation arguments survive.
Thus \(\mathrm{g2}\), \(\mathrm{g3}\),
\(\mathrm{g3'}\), and the regular frame element survive
on every selected higher layer.
For the local witness, either take the same power of
\(g_\ell\), or note that its cyclotomic value is \(1\),
so it already fixes the cyclotomic layers.
Equation (T10) survives.
The source's odd-power transport and
\(\mathrm{leFit}\) implement exactly these facts.

This supplies the actual tower hypotheses used by
the scalar-place producer, selected-tower irreducibility
and the ordinary local nonlinking calculation.
All these witnesses come from the principal packet and
the explicitly chosen auxiliary dyadic character.
No Case III tower is used.

% Exact frozen-source correspondence:
% FrobeniusFields/ExistsDecompositionElementAtLXiPowFourNeOne,
% ExistsInertiaElementAtTwoXiPowFourNeOne, Adj2, Elts2,
% CharacterSubgroupOpenEvenIndexTamePower, Ok2B;
% CompactIntegralFields/ExistsTameTowerOkTwo;
% UpperTriangularPlaces/BigOm and NormalValuationBounds/NeOfPair.
% The ell-primary character Kap is trivial for this 2-power
% principal packet; the source permits it for larger packets.


\subsection{Producing the rank-four auxiliary scalar place}
\label{subsec:sr-four-scalar-place}
The auxiliary local vanishing condition used in the Selmer
calculation has an actual place producer.
It uses the full root-character condition of the selected
tower, rather than a guessed congruence for a rational prime.

\begin{srlemma}\label{lem:sr-finite-character-separation}
Let \(k\) be a finite field, \(\Gamma\) a group, and
\(\chi_0,\ldots,\chi_3,\omega:\Gamma\to k^\times\)
characters with
\(\chi_0\chi_3=\chi_1\chi_2\).
Suppose
\[
       \omega\ne(\chi_0/\chi_1)^u(\chi_0/\chi_2)^v
                            \quad\hbox{for every }u,v\in\mathbb Z .
 \tag{A1}
\]
Then there exists \(g\in\Gamma\) such that
\[
                     \chi_0(g)=\cdots=\chi_3(g),
                            \qquad\omega(g)\ne1.
 \tag{A2}
\]
\end{srlemma}
\begin{proof}
Replace \(\Gamma\) by its finite image under these five
characters, a finite abelian group \(A\) of exponent
dividing \(\#k^\times\).
Its character pairing with
\(A^\vee=\operatorname{Hom}(A,k^\times)\) is perfect.
Indeed the decomposition of \(A\) into cyclic factors
and the roots of unity in the cyclic group \(k^\times\)
prove this on each factor.
For a subgroup \(S\subset A^\vee\), the same decomposition
or the order formula for annihilators gives
\[
                       \operatorname{Ann}
                         (\operatorname{Ann}(S))=S.
\]
Let \(S\) be generated by all \(\chi_i/\chi_j\).
The paired relation expresses all of them as powers of
\(\chi_0/\chi_1\) and \(\chi_0/\chi_2\).
Hypothesis (A1) says \(\omega\notin S\).
Double annihilation therefore gives an element of
\(\operatorname{Ann}(S)\) not killed by \(\omega\).
It is exactly (A2).
\end{proof}

\begin{srproposition}\label{prop:sr-four-actual-scalar-place}
Let \(L/F\) be the CM extension of an actual Case II
tower, and restrict its residual diagonal characters
\(\bar\chi_i\) and cyclotomic character
\(\bar\epsilon_\ell\) to \(G_L\).
The tower condition
\[
 \bar\epsilon_\ell\ne
       (\bar\chi_0/\bar\chi_1)^u
       (\bar\chi_0/\bar\chi_2)^v
                   \quad\hbox{on }G_L
                         \quad(u,v\in\mathbb Z)
 \tag{A3}
\]
produces a finite place \(v_0\) of \(F\), split in \(L\),
outside any specified finite set, such that
the residual characters are unramified there and
\[
 \bar\chi_0|_{D_{v_0}}=\cdots=\bar\chi_3|_{D_{v_0}},
              \qquad
            \bar\epsilon_\ell|_{D_{v_0}}\ne1 .
 \tag{A4}
\]
Here the decomposition group is identified with that
of either chosen split place of \(L\).
Multiplying all four characters by the residual scalar
link character preserves (A4).
\end{srproposition}
\begin{proof}
The paired relation holds for the actual residual tuple.
Apply Lemma~\ref{lem:sr-finite-character-separation}
on \(G_L\) with \(\omega=\bar\epsilon_\ell\);
obtain \(g_0\) as in (A2).
All characters have open kernel.

Choose a finite Galois extension \(J/F\) containing \(L\)
on which these characters become trivial.
They are restrictions of the given rational characters,
so their common kernel is normal in \(G_F\).
The image of \(g_0\) in \(\operatorname{Gal}(J/F)\)
fixes \(L\).
Chebotarev supplies infinitely many places of \(F\),
unramified in \(J\), with this Frobenius conjugacy class.
Choose one outside the prescribed finite set.
It splits in \(L/F\), because its Frobenius fixes \(L\).
All the characters in question are unramified there.

Their values on a Frobenius lift are those at \(g_0\);
conjugation does not change a character.
Their values on the entire local decomposition group
are powers of these values, since unramified characters
factor through the procyclic unramified quotient.
The four values are equal on \(g_0\), hence on this
whole decomposition group.
The cyclotomic value at that Frobenius is nontrivial.
This proves (A4).
An arbitrary finite avoidance set allows the two links'
distinguished auxiliary places to be different.
A common scalar multiplication changes none of the
ratios or the cyclotomic assertion, proving the last claim.
\end{proof}

Condition (A3) is exactly \(\mathrm{Tower.Ok.g2}\), after
identifying the two root coordinates with the root weights.
The frozen source implements the finite-duality argument
in \srcid{WeakShimuraTheory/UnipotentIdempotents/GExRtZero.lean},
theorem \(\mathrm{Lk.gEx}\), and the actual place selection
in \srcid{EqOfKer.lean}, theorem \(\mathrm{Lk.setEx}\).
The version with an additional finite avoidance condition
is \(\mathrm{setAv}\) in
\srcid{Foundations/GaloisRepresentations/PlainTorsionPairs/LinkTwo.lean}.
The residual tuple in the Hida datum is the scalar link
twist of this tuple.
Thus \(\mathrm{Lk.v0ok}\) supplies exactly
\(\mathrm{r.V0ok}\), not merely a genericity statement
about residual characters.
The local vanishing transfer in
\(\mathrm{P1.n4V0}\) can therefore be applied with its
actual scalar-place premise.
No condition on a Frobenius at another local place
is substituted for (A4).


\subsection{Irreducibility on the actual CM tower}
\label{subsec:sr-actual-tower-irreducibility}
The irreducibility premise at the start of
Section~\ref{subsec:sr-case-two-descent} is produced by
the actual residual tuple.
The following argument also explains why no arbitrary
finite restriction is being assumed irreducible.

\begin{srlemma}\label{lem:sr-residual-distinct-abelian-restriction}
Let \(G\) be a group and \(H\) a normal subgroup of finite
index with \(G/H\) abelian.
Let \(V\) be an irreducible representation over an
algebraically closed field of characteristic zero,
with a stable lattice over a DVR.
Suppose the residual semisimplification is
\(\bigoplus_{i=1}^n\bar\chi_i\), where
the characters \(\bar\chi_i|_H\) are pairwise distinct.
Then \(V|_H\) is irreducible.
\end{srlemma}
\begin{proof}
The restriction \(V|_H\) is semisimple.
Choose a nonzero simple \(H\)-submodule of minimal
positive dimension.
Its \(G\)-translates are simple because \(H\) is normal;
their sum is all of \(V\) by irreducibility.
Successively adding these simple submodules gives a
semisimple module, proving the claim.

If \(V|_H\) were reducible,
\(\operatorname{End}_H(V)\) would have dimension greater
than one.
Conjugation by \(G\) acts on this algebra through the
finite abelian group \(G/H\).
In characteristic zero this action is a sum of characters,
and its invariant subspace is
\(\operatorname{End}_G(V)\), of dimension one by Schur.
There is therefore a nonzero eigenvector \(A\) for a
nontrivial character \(\eta\) of \(G/H\):
\[
                    \rho(g)A\rho(g)^{-1}=\eta(g)A .
\]
This is an intertwiner between two irreducible
\(G\)-representations, hence is invertible.
Choose \(g\) with \(\eta(g)\ne1\).
Since \(\eta|_H=1\), taking traces gives
\[
                      \operatorname{tr}\rho(gh)=0
                                     \quad(h\in H).
\]
Reduce this exact equality using the stable lattice.
The residual semisimplification has the same traces,
so
\[
              \sum_i\bar\chi_i(g)\bar\chi_i(h)=0
                                      \quad(h\in H).
 \tag{I1}
\]
Distinct multiplicative characters of a group are
linearly independent over their coefficient field.
For reference, in a shortest nonzero relation, translate
the argument by an element separating two characters
and subtract a scalar multiple of the original relation.
The remaining nonzero relation is shorter, a contradiction.
Apply this to \(\bar\chi_i|_H\) in (I1).
It gives \(\bar\chi_i(g)=0\) for every \(i\), impossible
for characters valued in a field's multiplicative group.
\end{proof}

\begin{srproposition}\label{prop:sr-case-two-tower-irreducible}
Every irreducible actual characteristic-zero rational
point of the Case II problem remains irreducible on
the CM field \(L\) of the selected tower.
\end{srproposition}
\begin{proof}
The compositum \(L=FK\) is abelian over \(\mathbb Q\):
\(F/\mathbb Q\) is abelian and \(K/\mathbb Q\) quadratic.
Thus \(H=G_L\) is normal of finite index with abelian
quotient.

The stable lattice in the actual point has residual
semisimplification
\(\bigoplus_{i=0}^3\bar\chi_i\).
The tower's root-separation condition is
\[
       \forall w\in\{\hbox{symplectic roots}\}\quad
       \exists g\in G_L:\quad \bar r_w(g)\ne1 .
 \tag{I2}
\]
For \(i\ne j\), the paired relation
\(\bar\chi_0\bar\chi_3=\bar\chi_1\bar\chi_2\)
expresses \(\bar\chi_i/\bar\chi_j\) as
\(\bar r_{w(i,j)}\) for a symplectic root.
Equation (I2) therefore says exactly that their
restrictions to \(G_L\) are pairwise distinct.
The possible relabeling of the ordinary tuple is a
permutation and changes none of these conclusions.
Apply Lemma~\ref{lem:sr-residual-distinct-abelian-restriction}.

If the supported representative is a conjugate of the
original lattice representation, first apply the lemma
to the original representation and then conjugate.
Irreducibility is unchanged.
Hence the exact premise of
Proposition~\ref{prop:sr-all-quadratic-automorphy}
has been supplied for every actual irreducible rational
point, and that proposition's conclusion gives
\(\mathrm{IsAutQA}\) for all of them.
\end{proof}

% Frozen source:
% DenseThresholdCases/MapRepComp.lean,
% private irr_sub, dist_L, irr_L, and theorem aut_n_alt.
% (I2) is exactly Tower.Ok.g3; the finite-abelian quotient
% is proved there by comm_L.


\subsection{Removing the link character and descending the actual point}
\label{subsec:sr-case-two-descent}
We finish the automorphic endpoint of the Case II support
argument. Keep the actual tower and Hecke quotients of
Section~\ref{subsec:sr-case-two-hecke}.
Let
\[
 r:G_{\mathbb Q}\longrightarrow
                   \mathrm{GSp}_4(\overline{\mathbb Q}_\ell)
 \quad\hbox{with multiplier }\theta
 \tag{D1}
\]
be the original regular arithmetic point.
Proposition~\ref{prop:sr-case-two-tower-irreducible} supplies
\(r|_{G_L}\) irreducible for the actual selected tower.
This is the precise hypothesis
in the source's \(\mathrm{pmAutO}\) and
\(\mathrm{pmDown}\).
It is not inferred by restricting an arbitrary irreducible
representation to an arbitrary finite extension.
The preceding support argument gives a regular algebraic
cuspidal representation attached to
\[
                         r|_{G_L}\otimes\kappa .
 \tag{D2}
\]
The two links, their common scalar character, and strong
multiplicity one give the asserted local factors away from
either distinguished auxiliary place, as proved there.

\begin{srlemma}\label{lem:sr-link-character-algebraic}
The actual link character \(\kappa\) is the \(\ell\)-adic
avatar of an algebraic Hecke character of \(L\).
Consequently (D2) can be untwisted to a regular algebraic
cuspidal representation attached exactly to \(r|_{G_L}\).
\end{srlemma}
\begin{proof}
Write the original twist as \(a=2q+b\), \(b\in\{0,1\}\).
The actual character construction gives
\[
 \kappa=\kappa_0\gamma^{-(a+4)}\beta^b,\qquad
 \gamma^2\omega=\epsilon_\ell,\qquad \omega^{\ell-1}=1,
 \qquad \kappa_0^m=1
 \tag{D3}
\]
for some positive \(m\).
At each place above \(\ell\) there are an integer \(j\)
and a positive integer \(k\) such that
\[
                 (\beta^2\epsilon_\ell^{-(2j+1)})^k=1
                          \quad\hbox{on inertia}.
\]
Put \(A_v=bj-q-2\), an integer. Direct multiplication
in this commutative character group gives
\[
              (\kappa\epsilon_\ell^{-A_v})
                          ^{\,2km(\ell-1)}=1
                           \quad\hbox{on }I_v .
 \tag{D4}
\]
For example, the cyclotomic exponent contributed by (D3)
is \(-(2q+b+4)/2+b(2j+1)/2=bj-q-2\);
multiplying by \(2km(\ell-1)\) eliminates all finite factors.
Equation (D4), rather than a choice of a half-cyclotomic
character, proves the assertion about integral local
weights.
The finite coefficient and continuity properties of the
actual character imply that the finite-order quotient in
(D4) has open kernel.
Thus \(\kappa\) is locally algebraic on units at every
place above \(\ell\), with integral, parallel exponent \(A_v\).
Its construction makes it unramified outside a finite set.

For completeness this local statement globalizes by
idele class reciprocity.
Let \(P_\ell\) be the algebraic product of the embeddings
of \(L\) into \(\overline{\mathbb Q}_\ell\) with these
integer exponents, using the sign determined by the chosen
reciprocity convention. Let \(P_\infty\) be the same product
of complex embeddings, transported by \(\iota\).
On finite ideles divide \(\kappa\circ\operatorname{Art}_L\)
by \(P_\ell\), and at infinity multiply by \(P_\infty\).
This defines
\[
 \chi(s_f,s_\infty)=
 \iota\!\left(
   \frac{\kappa(\operatorname{Art}_L(s_f))}
                    {P_\ell(s_\ell)}\right)
                         P_\infty(s_\infty).
 \tag{D5}
\]
The finite factor is locally constant: this is (D4) above
\(\ell\); away from \(\ell\) a continuous rank-one
\(\ell\)-adic character has finite image on local units
at all ramified places.
For \(\alpha\in L^\times\), global reciprocity kills
the principal finite idele after its connected complex
infinite component is included, and
\(P_\infty(\alpha)=\iota(P_\ell(\alpha))\).
Hence (D5) is trivial on \(L^\times\).
It is an algebraic Hecke character and its \(\ell\)-adic
avatar is precisely \(\kappa\).

Twisting the actual representation in (D2) by
\(\chi^{-1}\circ\det\) preserves cuspidality and regular
algebraicity and removes \(\kappa\) from its attached
Galois representation.
This gives the representation asserted in the lemma.
\end{proof}

We next spell out the elementary Galois correction in
solvable descent. It ensures exact attachment, including
the finite character which a bare descent statement may
leave unspecified.

\begin{srlemma}\label{lem:sr-cyclic-exact-descent}
Let \(M/K\) be a cyclic extension of prime degree between
CM fields, and let \(r_K:G_K\to
\mathrm{GL}_n(\overline{\mathbb Q}_\ell)\) be continuous,
with \(r_K|_{G_M}\) irreducible.
If \(r_K|_{G_M}\) is attached to a regular algebraic
cuspidal representation \(\Pi_M\), then \(r_K\) is attached
exactly to a regular algebraic cuspidal representation
\(\Pi_K\).
Conversely, regular algebraic cuspidal automorphy of
\(r_K\) implies such automorphy of its restriction, under
the same irreducibility hypothesis.
\end{srlemma}
\begin{proof}
For descent, \(\Pi_M\) is invariant under
\(\operatorname{Gal}(M/K)\).
Indeed its conjugate has the same unramified Frobenius
polynomials, since \(r_K\) extends the attached
representation; Chebotarev and strong multiplicity one
give the invariance.
The cyclic base-change and descent theorem of
Arthur--Clozel therefore gives a cuspidal descent
\(\Pi_K^0\).
At infinity every local extension is \(\mathbb C/\mathbb C\),
so its infinitesimal characters are the same algebraic
regular ones. Thus \(\Pi_K^0\) is regular algebraic.

To identify the Galois representation of this descent,
use
\href{https://virtualmath1.stanford.edu/~rltaylor/rigcoh.pdf}
{Harris--Lan--Taylor--Thorne, Theorem A}.
Its hypotheses are precisely a CM field, a regular
algebraic cuspidal general-linear representation and
an arbitrary coefficient prime.
In particular no polarization of the uncorrected
descent is needed.
It supplies the continuous semisimple representation
\(s=r_{\ell,\iota}(\Pi_K^0)\).
Unramified base change and Chebotarev give
\(s|_{G_M}\simeq r_K|_{G_M}\).

Conjugate \(s\) so that these restrictions agree.
For \(g\in G_K\), the operator \(s(g)r_K(g)^{-1}\)
commutes with their common irreducible \(G_M\)-action:
this follows by applying both extensions to \(ghg^{-1}\),
\(h\in G_M\).
Schur's lemma makes this operator a scalar \(\delta(g)\).
Multiplication shows that \(\delta\) is a character,
trivial on \(G_M\); it therefore has finite order and
factors through \(\operatorname{Gal}(M/K)\).
Use global class field theory to twist
\(\Pi_K^0\) by the inverse finite Hecke character.
The resulting \(\Pi_K\) is attached exactly to \(r_K\).
A finite twist preserves regular algebraicity and
cuspidality.

For ascent, cyclic base change of \(\Pi_K\) exists.
If it were not cuspidal, the cyclic cuspidality criterion
would give a nontrivial character \(\eta\) of
\(\operatorname{Gal}(M/K)\) with
\(\Pi_K\simeq\Pi_K\otimes\eta\).
Chebotarev would give \(r_K\simeq r_K\otimes\eta_\ell\).
An intertwiner for this self-twist commutes with \(G_M\),
hence is scalar by irreducibility on \(G_M\);
the intertwining identity would then force \(\eta=1\).
This contradiction proves cuspidality of the base change.
Its regular algebraicity is preserved at infinity,
and Chebotarev identifies the attached representation
with \(r_K|_{G_M}\).
\end{proof}

Only the following published base-change input is used here:
the cyclic prime-degree base-change/descent theorem and its
cuspidality criterion, in Arthur--Clozel,
\emph{Simple algebras, base change, and the advanced theory
of the trace formula}, Chapter 3.
Iterating the lemma along a subnormal series with
cyclic prime quotients proves the same exact descent
for a solvable Galois CM extension.
In our tower all intermediate fields containing \(K\)
are CM: the totally real field is abelian over
\(\mathbb Q\), and complex conjugation is central in
\(\operatorname{Gal}(L/\mathbb Q)\).

\begin{srproposition}\label{prop:sr-all-quadratic-automorphy}
The actual point (D1)--(D2) is regular algebraic cuspidal
automorphic over every imaginary quadratic field \(K'\)
for which \(r|_{G_{K'}}\) is irreducible.
\end{srproposition}
\begin{proof}
Apply Lemma~\ref{lem:sr-link-character-algebraic}
and solvable exact descent from \(L\) to its chosen
imaginary quadratic field \(K\).
The irreducibility on \(G_L\) implies irreducibility
on all larger Galois groups along this descent.
This gives an actual regular algebraic cuspidal
representation attached to \(r|_{G_K}\).

Fix \(K'\) with \(r|_{G_{K'}}\) irreducible.
By Lemma~\ref{lem:sr-quadratic-self-twists}, an irreducible
rank-four representation of \(G_K\) can become reducible
on at most \(16\) distinct quadratic extensions of \(K\).
As \(K''\ne K\) ranges over imaginary quadratic fields,
the extensions \(KK''/K\) are distinct.
Indeed a biquadratic field with two imaginary quadratic
subfields has a real third quadratic subfield.
Apply the same argument over \(K'\).
Among any fixed \(35\) distinct imaginary quadratic fields
there is therefore a field \(K''\), distinct from \(K,K'\),
for which
\[
                r|_{G_{KK''}},\qquad r|_{G_{K'K''}}
                         \quad\hbox{are irreducible}.
 \tag{D6}
\]
There are at most \(16+16+2=34\) excluded fields.
Cyclic ascent from \(K\) to \(KK''\), exact descent to
\(K''\), ascent to \(K'K''\), and exact descent to \(K'\)
now prove the assertion.
At every use of Lemma~\ref{lem:sr-cyclic-exact-descent},
the needed restriction is one of (D6).
In particular this argument never assumes that
\(r|_{G_{LK'}}\) is irreducible.
\end{proof}

\subsubsection{The actual polarization and full local comparison}
Let \(\Pi_{K'}\) be a representation obtained above.
The symplectic relation in (D1) and the fact that \(r\)
is defined on \(G_{\mathbb Q}\) give
\[
              (r|_{G_{K'}})^c
                    \simeq(r|_{G_{K'}})^\vee\theta .
\]
With the Galois normalization
\(\operatorname{rec}(\Pi|\det|^{(1-n)/2})\),
one has
\[
 r_{\ell,\iota}(\Pi^\vee)
       =r_{\ell,\iota}(\Pi)^\vee\epsilon_\ell^{1-n}.
\]
The polarizing Hecke character of \(\mathbb Q\) is therefore
the algebraic character with avatar
\(\theta\epsilon_\ell^{n-1}\), here
\(\theta\epsilon_\ell^3\).
The actual multiplier is locally algebraic and odd.
Thus its product with \(\epsilon_\ell^3\) has value \(1\)
at complex conjugation, the required rank-four parity.
Chebotarev and strong multiplicity one give the
corresponding automorphic polarization of \(\Pi_{K'}\).

We use full Frobenius-semisimplified local--global
compatibility for this regular algebraic essentially
conjugate-self-dual cuspidal representation.
This combined published scope is explicitly recorded in the
introduction to
\href{https://annals.math.princeton.edu/wp-content/uploads/annals-v179-n2-p03-p.pdf}
{Barnet--Lamb--Gee--Geraghty--Taylor, p.~507}.
Away from \(\ell\) it is the polarized realization theorem
there; at \(\ell\) its unrestricted monodromy comparison is
\href{https://arxiv.org/pdf/1202.4683}
{Caraiani, Theorem 1.1}, together with that polarized
normalization.
The latter theorem has no Iwahori or Shin-regular
hypothesis in its statement.
The comparison identifies \(N\), not only the
semisimplified Weil representation.
The attached \(r|_{G_{K'}}\) is de Rham above \(\ell\)
with the specified Hodge numbers.
These realization and comparison theorems apply to the
already automorphic point at every odd coefficient prime,
including rank four at \(\ell=3\);
the lower bounds on coefficient primes in automorphy
lifting theorems are not hypotheses here.

\subsubsection{The local window predicate}
We verify the complete local assertion called
\(\mathrm{WitA}\) in the source.
It has two clauses.
First, at \(v\mid\ell\), an ordinary flag with characters
\(\kappa_i\) satisfying
\[
                         \kappa_i\ne\kappa_j\epsilon_\ell
                                \quad\hbox{for all }i,j
 \tag{D7}
\]
can be replaced by a flag for which
\[
                    \kappa_i^m\ne
                      (\kappa_j\epsilon_\ell)^m
                       \quad(i,j,\ m\ge1).
 \tag{D8}
\]
These are \(\mathrm{Unl1}\) and
\(\mathrm{PotUnl}\), respectively.
Second, for every regular algebraic cusp representation
attached to \(r|_{G_{K'}}\), there is a single real number \(t\)
such that, at every \(v\nmid\ell\) and at every \(v\mid\ell\)
satisfying (D8), its cuspidal support consists of rank-one
characters \(\xi\) with
\[
                 q_v^{\,t-1/2}<|\xi(\varpi_v)|
                                  <q_v^{\,t+1/2}.
 \tag{D9}
\]
The source denotes this local assertion by
\(\mathrm{WinAt}(\Pi,t,v)\).
The normalizing scalar in Galois--automorphic comparison
shifts \(t\) and never changes a ratio.

\begin{srlemma}\label{lem:sr-parallel-wd-characters}
Let \(k/\mathbb Q_\ell\) be finite.
Among continuous de Rham characters of \(G_k\)
whose cyclotomic exponents are parallel integers,
the rank-one Weil--Deligne character determines the
Galois character.
Consequently, if a de Rham representation with such
rank-one ordinary graded characters satisfies (D7),
then its Weil--Deligne monodromy is zero.
\end{srlemma}
\begin{proof}
Write such a character as
\(\alpha=\epsilon_\ell^h\eta u\), with \(h\in\mathbb Z\),
\(\eta\) finite order and \(u\) unramified.
This follows from local reciprocity after dividing out
the parallel algebraic character on an open subgroup
of units.
The unramified value of \(u\) is an \(\ell\)-adic unit,
since \(\alpha\), \(\epsilon_\ell\), and \(\eta\)
are continuous characters of the compact group \(G_k\).

At a geometric Weil Frobenius, the valuation of
\(\mathrm{WD}(\alpha)\) is \(h\) times the nonzero
valuation of \(\mathrm{WD}(\epsilon_\ell)\).
The finite-order factor and the value of \(u\)
have valuation zero.
Thus the Weil--Deligne character first recovers \(h\).
After dividing out that power of \(\epsilon_\ell\),
it recovers the finite inertia and the unramified
value by local reciprocity.
This proves the first assertion.

The Weil--Deligne functor is exact on potentially
semistable representations.
For the ordinary flag its semisimplified Weil
representation therefore has the rank-one characters
\(\lambda_i=\mathrm{WD}(\kappa_i)\).
If \(N\ne0\), the Weil equivariance of monodromy
gives two of these characters with
\(\lambda_i=\lambda_j\mathrm{WD}(\epsilon_\ell)\),
after interchanging \(i,j\) if required by the
Frobenius convention.
The first assertion gives
\(\kappa_i=\kappa_j\epsilon_\ell\).
This contradicts (D7).
\end{proof}

The restriction to parallel exponents is essential to this
short proof and is satisfied here.
The rational ordinary point has cyclotomic exponents
\((a+1,a,1,0)\) on an open inertial subgroup.
Every restriction to a finite local field has the same
parallel exponents.
The ordinary flag over the real tower descends to
\(G_{\mathbb Q_\ell}\) as follows.
Pass to an open normal inertial subgroup where all finite
characters disappear.
Choose an element on which the cyclotomic character has
infinite order; it has the four distinct eigenvalues
\(\epsilon_\ell(g)^{a+1},
  \epsilon_\ell(g)^a,\epsilon_\ell(g),1\).
The ordered flag is the unique flag made from the
successive eigenspaces in this order.
Every local conjugate flag has the same ordered
characters on that subgroup, so equals it.
Hence the rational decomposition group preserves this
flag. Its graded characters have the stated parallel
exponents.
Restriction supplies such a flag at every place of every
imaginary quadratic field.

Away from \(\ell\), the actual point has finite abelian
inertia extending to decomposition characters.
Here is the precise ramification argument.
Use the fixed torus residual representation
\[
 \bar r_{\rm tor}
       =\operatorname{diag}(\bar\chi_0,\ldots,\bar\chi_3)
       =D.\mathrm{rbar},\qquad
                       H=G_F\cap\ker\bar r_{\rm tor}.
 \tag{D10}
\]
This is literally a diagonal representation.
In the original framed rational problem
\(\mathrm{KPair.red}\) gives entrywise reduction to
\(D.\mathrm{r0}=\operatorname{diag}(D.\mathrm{ps})\);
it is stronger than prescribing a semisimplification.
After the lattice changes used in the residual chain,
the paired ramification set \(\mathrm{fr.N}\) still uses
this same fixed \(D.\mathrm{rbar}\), rather than the
possibly nonsplit residual representation at a chain point.
Thus \(H\) is open and contains the commutator of
\(G_{\mathbb Q}\): \(F/\mathbb Q\) is abelian and
\(\bar r_{\rm tor}\) is diagonal in global characters.
The actual paired ramification condition kills
\(I_p\cap H\) for every \(p\ne\ell\).
Conjugation of the characteristic-zero point preserves
this killing assertion.
Thus \(r(I_p)\) is a finite abelian group, and conjugation
by a decomposition element is trivial on its quotient
through \(G_{\mathbb Q}/H\).
Each inertial isotypic subspace is decomposition-stable.
Diagonalizing a Frobenius on each such subspace shows
that the Frobenius-semisimplified Weil representation is
a sum of characters extending the inertial characters.
Finite inertia gives \(N=0\).
These statements remain true on every finite local
restriction.

\begin{srproposition}\label{prop:sr-actual-wit}
For every \(K'\) in Proposition~\ref{prop:sr-all-quadratic-automorphy},
the actual representation \(r|_{G_{K'}}\) satisfies
\(\mathrm{WitA}\).
Thus the endpoint is exactly the source predicate
\(\mathrm{IsAutQA}\), including its local assertion.
\end{srproposition}
\begin{proof}
A cuspidal automorphic representation of a general-linear
group becomes unitary after a single real norm twist.
The same twist works at all finite places.
Every local component of a cuspidal representation is
generic. The
\href{https://www.numdam.org/item/ASENS_1986_4_19_3_335_0/}
{generic unitary classification of Tadi\'c}
has the following consequence.
If its Weil--Deligne parameter is a sum of characters
and has \(N=0\), it is a principal series with inducing
characters in the strict window (D9).
Indeed its unitary building blocks are tempered
characters and complementary pairs with exponents
\(s,-s\), \(0<s<1/2\).
A discrete-series block of length greater than one
would have nonzero monodromy and is excluded.
Undoing the one global norm twist gives the one
number \(t\) in (D9).

At \(v\nmid\ell\), the finite abelian inertia calculation
after (D10) and full local--global compatibility give
exactly this character parameter with \(N=0\).
Hence (D9) holds there.

At \(v\mid\ell\), assume (D7).
There exists a flag with parallel graded exponents,
as just proved.
Its graded character multiset is the intrinsic multiset
of Jordan--H\"older characters of the local representation.
Any other full invariant flag gives the same multiset.
Thus (D7), if it holds for a flag, holds also for the
parallel flag.
The representation is de Rham by its actual polarized
automorphic realization, so
Lemma~\ref{lem:sr-parallel-wd-characters}
gives \(N=0\).
Its semisimplified Weil parameter is a sum of characters,
by exactness along the ordinary flag.
Full compatibility and the generic unitary classification
therefore give (D9).

Suppose now that (D8) fails.
Then for some \(i,j,m\ge1\),
\((\kappa_i/\kappa_j)^m=\epsilon_\ell^m\).
The Weil characters would have complex Frobenius
absolute-value ratio \(q_v\) or \(q_v^{-1}\);
the finite-order ambiguity has absolute value one.
Under local--global compatibility the common rank
normalizing scalar cancels, so the inducing characters
in (D9) would have ratio \(q_v\) in one order.
Their strict window makes that ratio strictly less
than \(q_v\), a contradiction.
The same parallel flag satisfies (D8).
This proves \(\mathrm{Unl1}\Rightarrow\mathrm{PotUnl}\).

Conversely (D8) implies (D7) by taking \(m=1\);
the preceding argument supplies (D9) at every such
place. This is exactly the second clause of
\(\mathrm{WitA}\).
Any other regular algebraic cusp representation
attached to this Galois representation equals
\(\Pi_{K'}\) by Chebotarev and strong multiplicity one,
so the asserted quantifier over all attached
representations introduces no extra hypothesis.
Together with
Proposition~\ref{prop:sr-all-quadratic-automorphy}
this proves \(\mathrm{IsAutQA}\).
\end{proof}

\subsubsection{The ordinary nonlinking condition at the selected tower}
We also record why the actual selected tower gives (D7)
before invoking the preceding implication.
This is the source's direct \(\mathrm{unlF}\) calculation;
it verifies a local hypothesis rather than assuming it
from generic residual characters.
The tower has the following alternative at every
\(\ell\)-adic place:
\[
 \begin{array}{ll}
 \text{either}&
 \exists g\in I_F:\ \bar\omega(g)\ne1,
                         \quad\bar r|_{I_F}=1,\\
 \text{or}&
 \exists g\in D_F:\quad
      \bar\omega(g)=1,\quad\bar\Psi(g)=1,\quad
                                     \bar\xi(g)\ne1 .
 \end{array}
 \tag{D11}
\]
Here \(\bar\omega=\bar\epsilon_\ell\),
\(\bar\Psi\) is the actual root character, and on
\(\ker\bar\Psi\) the residual diagonal is
\[
                      (\bar\xi,1,\bar\xi,1).
 \tag{D12}
\]
These are precisely the local tower condition
\(\mathrm{LocL}\) and the identities in
\(\mathrm{barFacts}\).

If \(\kappa_i=\kappa_j\epsilon_\ell\), restrict to an
open inertial subgroup where their finite factors
disappear.
The infinite cyclotomic image gives
\(h_i=h_j+1\).
For \((h_0,h_1,h_2,h_3)=(a+1,a,1,0)\), \(a\ge2\),
the only possibilities are
\((i,j)=(0,1),(2,3)\), or \((1,2)\), the last only
when \(a=2\).
Reduce the same character equality modulo the
coefficient maximal ideal.
In the first case of (D11) its two residual sides
are \(1\) and \(\bar\omega(g)\), a contradiction.
In the second case, (D12) has unequal \(i,j\)
entries for every listed pair, while
\(\bar\omega(g)=1\), again a contradiction.
Thus (D7) holds on the selected real tower.
The CM extension is split at these places, so the
same local decomposition groups give (D7) over \(L\).

Apply the local argument in
Proposition~\ref{prop:sr-actual-wit} to the actual
automorphic point over \(L\), undoing its algebraic
scalar link twist, whose ratios cancel.
It gives (D8) on \(G_{L_w}\).
An equality prohibited by (D8) on the rational
decomposition group would restrict to such an equality
on this finite-index subgroup.
Therefore (D8) holds already for the rational ordinary
flag and after every finite local restriction.
Indeed if the ratio became finite order on a
finite-index subgroup, a further positive power would
be trivial on the whole local group.
The actual point is consequently potentially crystalline
at \(\ell\), and has the required unlinked principal
series factors after each imaginary quadratic restriction.
This conclusion is obtained from actual automorphy,
the checked tower condition and full comparison.

The local--global comparison at a distinguished auxiliary
place uses the other link, as explained after (H16).
Thus no auxiliary prime or neat level is retained in this
endpoint.
This completes the actual Case II match from the common
ordinary module and its transported support to the full
automorphic-and-local predicate used by the counting
argument.

% Exact frozen-source counterparts:
% WeakShimuraTheory/UnipotentIdempotents/KapCyc;
% DenseThresholdCases/AuKap, TAlgOfCyc, TwistInvarianceAutomorphic;
% BaseChangeUnits/SolCOfSolDescent, VWOfOne, TAcross;
% DenseThresholdCases/PmAutO, DnPmDown;
% UpperTriangularPlaces/PtLawsOf and
% NormalValuationBounds/NeOfPair, CommLeRbar;
% AutomorphyLifting/WeakOfAtt and IsAutCsdA.
% The selected-tower irreducibility premise is explicit in (D1).

\section{The finite rank-three host and exact automorphic endpoints}
\label{sec:sr-rank-three-host-endpoints}

\subsection{Constant characteristic polynomials and the Gaussian finite part}
\label{subsec:sr-constant-character-finite-part}
We make explicit the last algebraic step from
\(\mathrm{FinEnd}\) to a finite weight-zero algebra.
All groups below are the actual finitely ramified arithmetic
quotients, so Lemma~\ref{lem:sr-finite-cross-cocycles}
and Proposition~\ref{prop:sr-finite-CH} apply.
The latter proposition and the integral determinant identity
(H9) are valid in characteristic \(3\).

\begin{srlemma}\label{lem:sr-constant-character-triangular}
Let \(k\) be a finite field of odd characteristic,
\(A\) a local domain containing \(k\), and
\(r:\Gamma\to\mathrm{GL}_3(A)\rtimes\{\pm1\}\)
an actual normalized rank-three unitary point.
Write \(H=\ker\eta\).
Suppose its reduction is the torus tuple
\((\chi,1,\chi^{-1})\), with \(\chi(g_1)^2\ne1\).
Suppose on an open subgroup of \(H\) its matrix
characteristic polynomials equal those of that tuple.
Then the matrices on \(H\) are all upper triangular
or all lower triangular in the marked basis, and every
diagonal entry equals its prescribed constant character.
\end{srlemma}
\begin{proof}
For each \(g\in H\), some positive power \(g^s\)
belongs to the given open subgroup.
An eigenvalue \(\alpha\) of \(r(g)\) has \(\alpha^s\)
algebraic over \(k\), so \(\alpha\) is algebraic over \(k\).
Thus every characteristic coefficient of \(r(g)\)
is algebraic over \(k\).
An element of \(A\) algebraic over the finite field \(k\)
belongs to a finite subfield of \(A\).
Reduction is injective on that subfield.
Its prescribed residual value consequently determines it
exactly.
It follows that
\[
 \operatorname{char}_{r(g)}(T)
                   =(T-\chi(g))(T-1)(T-\chi(g)^{-1})
                                \quad(g\in H).
 \tag{CF1}
\]

Here is the exact triangularity argument in the marked basis.
Let \(c\) be the normalized conjugation, so its matrix
part is \(1\), and put \(g^*=cg^{-1}c\).
The unitary law gives
\[
 r(g^*)=J_3r(g)^{\mathsf t}J_3^{-1}\qquad(g\in H).
\]
The marked normalization kills the two adjacent entries
of \(r(g_1)-r(g_1^*)\).
The other two off-diagonal entries cancel automatically
under this transpose reflection.
Hence
\[
             r(g_1)-r(g_1^*)=\operatorname{diag}(\zeta,0,-\zeta),
 \tag{CF2}
\]
where \(\zeta\) reduces to
\(\chi(g_1)-\chi(g_1)^{-1}\) and is a unit.
The three diagonal values in (CF2) are distinct,
since \(2\ne0\).

We first exclude irreducibility over \(\operatorname{Frac}A\).
By (H9), the characteristic polynomials on group words
determine the full determinant on the group algebra
without division by \(3!\).
Equation (CF1) therefore factors this representation
through the constant-character Cayley--Hamilton algebra
of Proposition~\ref{prop:sr-finite-CH}, after scalar extension.
Indeed continuity modulo every power of the maximal
ideal gives a finite quotient representation at each
precision; its determinant kills the defining
Cayley--Hamilton relations, and separatedness then makes
the original representation kill their closed ideal.
Thus passage to the fraction field does not silently
discard the completed-algebra condition.
Its nilpotent radical has quotient \(k^3\), as in that
proposition's proof.
Every irreducible representation of this algebra is
one-dimensional: the radical annihilates a simple module,
and the three coordinate idempotents of \(k^3\) split it.
Thus the rank-three representation is reducible.

Any invariant subspace is invariant under (CF2).
Its spectral projectors, obtained by Lagrange interpolation
at \(\zeta,0,-\zeta\), show that it is spanned by
a nonempty proper subset of the three coordinate vectors.
If every matrix entry \(ij\) vanishes on \(H\), so does
entry \(2-j,2-i\), by the unitary reflection.
There are six coordinate subsets to check.
The subsets \(\{0\},\{0,1\}\) force upper triangularity;
\(\{2\},\{1,2\}\) force lower triangularity.
The subsets \(\{1\},\{0,2\}\) force a block decomposition
of the middle line and the outer two coordinates.

In this last case the middle entry is an eigenvalue,
hence algebraic over \(k\); its residual value is \(1\),
so it equals \(1\).
The outer block has characteristic polynomial
\((T-\chi(g))(T-\chi(g)^{-1})\) for every \(g\).
Apply the rank-two instance of the same
Cayley--Hamilton argument.
The outer representation is reducible.
Its invariant line is a coordinate line because
the restriction of (CF2) is
\(\operatorname{diag}(\zeta,-\zeta)\).
One of the two corner entries consequently vanishes
on every \(g\), and the whole rank-three representation
is upper or lower triangular.
Injection into \(\operatorname{Frac}A\) transports
these zero entries back to \(A\).

Each diagonal entry in a triangular matrix is an
eigenvalue, so the finite-field argument following
(CF1) fixes it to its prescribed residual character.
Finally the matrix part of \(r(gc)\) equals that of
\(r(g)\); thus the same diagonal assertion holds
outside \(H\) as well.
\end{proof}

\begin{srlemma}\label{lem:sr-graded-finite-part-criterion}
Let \(\mathcal O\) be a complete coefficient DVR with
finite residue field and uniformizer \(\varpi\).
Let \(R\) be the actual complete Noetherian local
homogeneously presented algebra, with finitely many
homogeneous generators \(z_i\in\mathfrak m_R\),
of integer weights.
Let \(B=R_0\) be its completed weight-zero algebra.
Suppose every local specialization \(R\to A\), with
\(A\) a complete characteristic-residue-prime DVR
with finite residue field, kills every zero-weight
generator and kills either all positive-weight or
all negative-weight generators.
Then \(B\) is finite over \(\mathcal O\).
\end{srlemma}
\begin{proof}
The zero-weight monoid of exponent vectors has finitely
many generators.
For example intersect the nonnegative rational cone
with the integral kernel of the weight homomorphism;
triangulating its rational polyhedral cones and taking
the finitely many lattice points in their bounded
fundamental parallelepipeds proves this assertion.
Consequently \(B\) is a complete Noetherian local
\(\mathcal O\)-algebra with maximal ideal generated by
\(\varpi\) and finitely many nonconstant balanced
monomials \(b_1,\ldots,b_t\) in the \(z_i\).
The continuous weight-zero projection is
\(\mathcal O\)-linear, and gives
\[
                       (\varpi R)\cap B=\varpi B.
 \tag{CF3}
\]

We claim \(b_j\in\sqrt{\varpi R}\).
Otherwise choose a prime \(P\supset(\varpi)\) avoiding
\(b_j\).
In the complete local domain \(R/P\), its image is
nonzero and belongs to the maximal ideal.
There is a one-dimensional domain quotient preserving
that nonzero image.
Here is a direct curve-selection argument.
If the dimension exceeds one, choose an element of
the maximal ideal outside the finitely many height-one
minimal primes above that image.
A minimal prime above the chosen element has height
one by the principal ideal theorem and avoids the
image; quotient by it.
Dimension decreases, and iteration ends in dimension one.
The normalization of this one-dimensional complete local
domain is finite, local and a complete DVR, with finite
residue field.
Thus there is a local map \(R\to A\) of the kind in
the hypothesis with \(b_j\) nonzero.

A nonconstant zero-weight monomial either contains a
zero-weight variable, or contains variables of both
signs.
Every such monomial maps to zero under the hypothesis,
a contradiction.
This proves the claim.
Use (CF3) on a power of \(b_j\) to get
\(b_j\in\sqrt{\varpi B}\).
Since these finitely many monomials and \(\varpi\)
generate the maximal ideal,
\(\mathfrak m_B^s\subset\varpi B\) for some \(s\).
Hence \(B/\varpi B\) is finite over the residue field.
The \(\varpi\)-adic and maximal-ideal topologies are
equivalent by this containment.
Completeness and the usual successive lifting of a
basis of \(B/\varpi B\) show that finitely many elements
generate \(B\) over \(\mathcal O\).
\end{proof}

\begin{srproposition}\label{prop:sr-host-finite-part-from-finend}
For the actual marked rank-three host algebra,
\(\mathrm{FinEnd}\) implies finiteness of its weight-zero
part over the coefficient DVR.
\end{srproposition}
\begin{proof}
Its adapted homogeneous generators are matrix entries
of \(r-r_{\mathrm{tor}}\); their weights are the
differences of \(1,0,-1\).
They all belong to the maximal ideal by residual
normalization.
A characteristic-\(\ell\) DVR specialization gives the
ordinary point to which \(\mathrm{FinEnd}\) applies.
That condition fixes its characteristic polynomials
on an open subgroup of \(\ker\eta\).
Lemma~\ref{lem:sr-constant-character-triangular} makes
it upper or lower triangular with the prescribed
constant diagonal.
The zero-weight generator images are therefore zero;
in one triangular case all positive-weight images
vanish, and in the other all negative-weight images
vanish.
Lemma~\ref{lem:sr-graded-finite-part-criterion}
gives the required finite weight-zero part.
\end{proof}

The exact frozen-source chain is
\srcid{CurveDVRSpecialization/ConstOfChar.lean},
\srcid{UnitaryThreeDeformations/UniversalDiagonalOfTriangularPoints.lean},
and
\srcid{CurveDVRSpecialization/FinCrit.lean}.
The preceding proofs give their needed arithmetic specialization.
The intermediate \srcid{DiagEq} theorem fixes diagonals for
general triangular ordinary points using the frame's class-one
condition; here characteristic-\(\ell\) finiteness needs only
the stronger constant-character argument (CF1).
No unspecified triangularity or diagonal-finiteness assertion
is retained in the implication from (G10) to the finite
Gaussian host algebra.


\subsection{The actual Gaussian seeds and their common support}
\label{subsec:sr-gaussian-producer}
We give the arithmetic specialization needed for the
large-conductor Gaussian branch, including coefficient
prime \(3\).
The shared sharp/flat propagation theorem has a precise
remaining role at the end of this subsection;
automorphy is not assumed as a Gaussian literature input.

Fix \(K=\mathbb Q(i)\), \(\eta=\eta_K\),
\(\ell\in\{3,7\}\), twist \(1\), and an actual principal
character \(\psi\) in the large-conductor branch.
Thus \(\ell\nmid N\), \(4\mid N\),
\(\psi(-1)=-1\), and its dyadic component \(\psi_2\)
is even and primitive of conductor \(2^c\).
Put \(\Psi=\epsilon_\ell\psi\).
The bounds in Lemma~\ref{lem:sr-dyadic-three-seven}
give \(c\ge14\) at \(3\), \(c\ge18\) at \(7\), and
\[
           \operatorname{ord}\psi(3)=2^{c-2},\qquad
           \operatorname{ord}\psi(7)=2^{c-3}.
 \tag{G1}
\]
In particular these orders exceed \(8\).
All finite character values have order prime to \(\ell\).

\subsubsection{A scalar prime and the real quadratic base}
Choose a dyadic unit \(b\equiv1\pmod4\) with
\(\psi_2(b)=-1\).
Such a unit exists because \(\psi_2\) is faithful on
the cyclic \(1+4\)-quotient.
Set the odd-prime residues equal to \(1\).
The Chinese remainder theorem and Dirichlet's theorem
give a fresh prime \(p_0\) with
\[
         p_0\equiv b\pmod{2^c},\quad
         p_0\equiv1\pmod{qd},\quad
         p_0\equiv-1\pmod\ell .
\]
The harmless larger dyadic level in \(N\) can be included
by lifting the chosen residue.
Hence
\[
       \eta(p_0)=1,\quad
       \bar\psi(p_0)=-1,\quad
       \bar\Psi(p_0)=1,\quad
       \bar\epsilon_\ell(p_0)=-1 .
 \tag{G2}
\]
Thus \(p_0\) splits in \(K\).
Both residual tuples
\((\bar\Psi,\bar\Psi^{-1})\) and
\((\bar\Psi,1,\bar\Psi^{-1})\) are scalar there.
We may avoid any given finite set.

Choose a fresh positive prime \(D\equiv7\pmod8\).
At each odd level prime \(p\), including \(\ell\),
require
\[
                   (D/p)=\eta(p),
\]
and also require \((D/p_0)=1\).
These conditions specify nonzero residues modulo the
distinct primes, so another reduced arithmetic progression
and Dirichlet's theorem produce such a \(D\).
Put
\[
                       F_0^+=\mathbb Q(\sqrt D).
\]
At \(2\), its completion is \(\mathbb Q_2(i)\);
at every odd level prime its quadratic local character
equals that of \(K/\mathbb Q\).
Thus every level prime splits in
\(F_0^+K/F_0^+\).
Moreover
\[
 \operatorname{disc}(F_0^+K)
        =(-4)(4D)(-D)=16D^2
                     =\operatorname{disc}(F_0^+)^2 .
\]
The relative discriminant formula shows that this CM
quadratic extension is unramified at every finite place.
This real quadratic field is the exact index-two
\(\mathrm{H}_0\) of the source's
\srcid{FrobeniusSplitPrimeDatum}.
The extra ramification at \(D\), outside the level,
supplies that datum's fresh ramified prime.
The prime \(p_0\) splits in \(F_0^+K\).
At \(\ell\), the completion of \(F_0^+\) is the
unramified quadratic extension, since \(\ell\) is
inert in \(K\).

Choose a fresh prime
\(w\equiv1\pmod{\ell N}\) and a character \(\vartheta\)
of exact order \(\ell\), conductor \(w\).
It is even and congruent to \(1\) modulo the coefficient
maximal ideal.
Let \(F_\vartheta/\mathbb Q\) be its real cyclic
degree-\(\ell\) field.
For the depth \(r\) required by the common frame put
\[
 F_r^+=F_0^+F_\vartheta
        \mathbb Q(\zeta_{\ell^{r+1}})^{\mu_{\ell-1}},
                  \qquad L_r=F_r^+K .
 \tag{G3}
\]
The CM extension remains finite-unramified and split
at all level primes.
Its real degree is even.
The new residue degrees at \(p_0\) are powers of the
odd prime \(\ell\), so (G2) survives on its decomposition
groups over \(F_r^+\).
This supplies the actual scalar-place condition for
both ranks and their scalar twists.

\subsubsection{The exact local conditions of the Gaussian frame}
There is no need to require
\(\bar\Psi^2\ne1\) on \(\ell\)-adic inertia.
That assertion would fail at \(\ell=3\).
The source's local separation condition concerns the
whole decomposition group, and it holds here.

Let \(f\) be a rational local uniformizer lift with
\(\epsilon_\ell(f)=1\).
Since \(K/\mathbb Q\) is inert at \(\ell\), \(f^2\)
fixes \(K\), and
\[
                 \bar\Psi(f^2)^2=\bar\psi(\ell)^4\ne1 .
 \tag{G4}
\]
Over a field in (G3) the unramified local degree is
\(2\ell^u\) for some \(u\ge0\);
the cyclotomic extra ramification has \(\ell\)-power degree.
Thus (G1) still gives (G4) on every actual local
decomposition group.
More explicitly, for \(j=\pm1,\pm2\), both
\(\bar\Psi^j\) and
\(\bar\Psi^j\bar\epsilon_\ell^{-1}\)
are nontrivial there.
Evaluate on \(f^{2\ell^u}\):
the cyclotomic value is \(1\) and the nontriviality
is detected by \(\bar\psi(\ell)^{2j\ell^u}\).
These are the full-local-group invariant vanishings
used in local Tate duality and the Euler formula.

The mod-\(\ell\) cyclotomic inertia image is still full:
the initial real base is unramified at \(\ell\) and
the later ramification has \(\ell\)-power index,
whereas its order is \(\ell-1\).
This is exactly the source's \(\mathrm{UFacts.mu}\).
The local fields contain no \(\mu_\ell\).
Their pro-\(\ell\) unit groups therefore have the
free ranks used in the ordinary module construction.
The character \(\bar\Psi\) is trivial on this pro-\(\ell\)
inertial part.

The other Gaussian frame separation has an explicit
global witness.
Take the dyadic residue \(-1\), odd-prime residues \(1\),
and residue \(1\pmod\ell\).
Evenness of \(\psi_2\) gives an element with
\[
              \bar\epsilon_\ell=1,\qquad
              \bar\Psi=1,\qquad \eta=-1 .
 \tag{G5}
\]
Thus
\(\ker\bar\epsilon_\ell\cap\ker\bar\Psi
 \not\subset\ker\eta\), the exact \(\mathrm{Gen.g4}\).
Equation (G4) supplies \(\mathrm{Gen.g5}\).
All finite inertial values away from \(\ell\) have
\(2\)-power order, so their Teichm\"uller lift is
unchanged and the source's character \(\mathrm{H}_b\)
may be the whole group.

Apply the shared arithmetic-frame construction with this
fixed real quadratic base and these witnesses.
Its Ferrero--Washington class-rank bound is chosen
before the cyclotomic depth; its local invariant
conditions are the full-decomposition-group assertions
just checked.
For the two strict classes
\(\bar\Psi^{-a'}\eta\), \(a'=1,2\), nontriviality over
\(F_r^+\) follows from (G4), including when inertia is
trivial at \(3\).
Consequently the restriction on \(H^1\) is injective
by inflation--restriction.
The local Euler/Poitou--Tate proof and the class-rank
bounds are unchanged.
Choose the depth only afterwards, as in that shared proof.
The marked element is chosen using (G4), in the local
decomposition group.
It is not required to belong to inertia.
These observations match the source's
\(\mathrm{Frame.Ok}\), \(\mathrm{Frame.Big}\) and
\(\mathrm{GalLink}\) conditions at \(3\) as well as \(7\).

\subsubsection{Two actual cuspidal seeds}
Use \(\chi'=\psi\), \(j=1\), and the same
\(\eta,\vartheta,w,L_r\) for both ranks.
These satisfy the source's
\srcid{FieldPairCharacterTripleConditions}:
\(\eta\) cuts out \(K\), its conductor divides \(N\),
\(\chi'(-1)=-1\), the prime \(w\) splits both finite
characters, and \(\vartheta|_{G_{L_r}}=1\).
The local witness (G4) is its nontrivial-square condition.
Also \(\chi'\ne\eta\) and its order is prime to \(\ell\),
by the large dyadic conductor.

Take \(m\equiv1\pmod{\ell-1}\) as large as needed.
For rank three, apply the constructed Eisenstein seed
to the primitive character underlying
\[
                         \varphi_3=\psi\eta
\]
and the auxiliary \(\vartheta\), at weight \(m+1\).
This even character is nontrivial: a conductor at least
\(2^{14}\) cannot be canceled by the quadratic
conductor-four factor.
Let \(\rho_3\) be the resulting two-dimensional
ordinary cuspidal Galois representation.
Set
\[
                            r_3=\operatorname{ad}^0\rho_3
                                                    |_{G_{L_r}} .
\]
Its residual tuple, ordinary inertial exponents and
finite inertia are exactly
\[
 \begin{array}{c|c|c}
 \hbox{residual tuple}&\hbox{cyclotomic exponents}&
                                     \hbox{finite inertia tuple}\\ \hline
 (\bar\Psi,1,\bar\Psi^{-1})&(m,0,-m)&(\psi,1,\psi^{-1}) .
 \end{array}
 \tag{G6}
\]
The scalar automorphic normalization is
\(k_3=\epsilon_\ell^{-1}\).

For rank two use instead
\(\varphi_2=\psi^2\eta\), an odd, hence nontrivial,
character, and classical weight \(2m+1\).
Let \(\rho_2\) be its ordinary cuspidal representation
and put
\[
                   \sigma_2=\rho_2\epsilon_\ell^{-m}\psi^{-1}.
\]
The determinant and ordered residual characters are
\[
 \det\sigma_2=\eta\vartheta,\qquad
 \bar\sigma_2^{\,\mathrm{ss}}
       =(\bar\epsilon_\ell^m\bar\psi\eta)
                     \oplus(\bar\epsilon_\ell^{-m}\bar\psi^{-1}).
 \tag{G7}
\]
Its restriction \(r_2=\sigma_2|_{G_{L_r}}\) therefore
has determinant one, residual tuple
\((\bar\Psi,\bar\Psi^{-1})\), ordinary exponents
\((m,-m)\), and finite inertia
\((\psi,\psi^{-1})\), with the indicated ordering.
These are the exact rank-two seed data,
not just a matching adjoint trace.

Both classical levels are their primitive character
conductors times \(w\), and divide \(Nw\).
Their minimal-level local calculation gives \(N=0\)
and the specified rank-one inertial summands.
At \(w\), restriction to \(L_r\) kills
\(\vartheta\), so both seeds become unramified there.
Thus \(w\) does not enlarge their tame packet on \(L_r\).
At \(\ell\) they are crystalline ordinary, with the
unit-root congruence supplied by the Eisenstein seed.
The coefficient-\(3\) construction is
Section~\ref{subsec:sr-eisenstein-three}; no
\(\ell\ge5\) assumption is retained.

The strong-open irreducibility proof in the constructed
seed applies to both \(\rho_i\).
For \(\rho_3\), the ratio of the residual characters
on a local element above \(\ell\) has nontrivial square
by (G1)--(G4).
For \(\rho_2\), its square is detected by
\(\psi(\ell)^8\ne1\).
Thus neither residual ratio can be quadratic.
A quadratic self-twist is excluded, and the distinct
Hodge numbers exclude finite projective image.
Consequently both two-dimensional representations and
the rank-three adjoint remain irreducible on every
open subgroup.
This proof is valid in coefficient characteristic
\(3\), since only characteristic zero and preservation
of quadratic characters modulo an odd prime are used.

\subsubsection{An explicit Gaussian rank-two scalar character}
There is a particularly concrete \(k_2\) in this case.
Let \(E:y^2=x^3-x\).
Its CM endomorphism \(i\) is defined over \(K\);
after splitting coefficients its Tate representation on
\(G_K\) is a sum of characters \(\lambda,\lambda^c\).
The Weil pairing gives
\[
                              \lambda\lambda^c=\epsilon_\ell .
\]
The character \(\lambda\) is the usual Gaussian
algebraic Hecke character: a prime-to-\(2\) ideal has
a unique generator congruent to \(1\) modulo
\((1+i)^3\), and this generator gives its ideal
character, with the reciprocity orientation fixed above.
This explicit description is recorded, for example, in
\href{https://www.college-de-france.fr/media/professeurs-honoraires/UPL6386850193290496047_cours_Taiwan_avril__10juin_2011__5.pdf}
{Serre, \emph{Lectures on \(N_X(p)\)}, Section 2.2.2}.
It shows that the finite conductor is supported at \(2\)
and its finite inertial values lie in \(\mu_4\).
Choose the CM embedding for which \(\lambda\) has
cyclotomic exponents \((1,0)\), and set
\[
                                  k_2=\lambda^{-1}|_{G_{L_r}} .
 \tag{G8}
\]
It is algebraic, has exponents \((-1,0)\), and
\(k_2k_2^c=\epsilon_\ell^{-1}\).

This character also satisfies the source's transfer
condition, which cannot be replaced just by that norm
identity.
For \(t\in G_{\mathbb Q}\setminus G_K\), the matrix of
\(t\) on the two CM lines is off-diagonal.
Its determinant is \(\epsilon_\ell(t)\).
Its square on either line is therefore
\(-\epsilon_\ell(t)\); hence
\[
              \lambda(t^2)=-\epsilon_\ell(t),\qquad
              \lambda^{-1}(t^2)=-\epsilon_\ell(t)^{-1}.
 \tag{G9}
\]
This is exactly \(\mathrm{IsTwist.ver}\).
Finite inertia away from \(\ell\) is killed by the
prime-to-\(\ell\) integer \(4\).
Good reduction of \(E\) at \(3,7\) makes its rank-one
characters crystalline there.
All fields of coefficients are finite.
Thus (G8) supplies every character requirement in
the actual rank-two seed normalization.

The trace-zero trace form gives the symmetric
conjugate pairing of \(r_3\).
For rank two, write \(J\) for the alternating
two-by-two form and use \(J\sigma_2(c)\).
Since \(\det\sigma_2(c)=-1\) and \(c^2=1\),
this matrix is symmetric and gives a conjugate pairing
of \(r_2\) with multiplier one.
Twisting by (G8) changes that multiplier to
\(\epsilon_\ell^{-1}\), the rank-two conjugate-self-dual
Galois normalization.
Similarly \(r_3\epsilon_\ell^{-1}\) has multiplier
\(\epsilon_\ell^{-2}\), the rank-three normalization.

Cyclic base change of the two classical cusp forms
to \(L_r\) is cuspidal by strong-open irreducibility.
Their algebraic twists remain cuspidal and regular.
The Gelbart--Jacquet adjoint lift gives the cuspidal
rank-three realization of (G6), with normalization
\(r_3\epsilon_\ell^{-1}\).
The actual local principal-series parameters have
\(N=0\).
The ordinary eigenvalues are those determined by
the unit and uniformizer equations (H5)--(H6), in
ranks two and three respectively.
The generic unitary window, or the adjoint root bound
proved in the constructed seed, supplies the
unlinked local condition.
This verifies the full automorphic-and-local seed
conditions \(\mathrm{IsSeedAlt}\).

\subsubsection{Insertion in the same two function modules}
The definite groups \(U(2)\) and \(U(3)\) may both be
taken from the positive definite identity hermitian forms.
At every finite place the CM extension is unramified,
and units are local norms; these forms are quasi-split
there.
Their real signatures are definite, and the even
real degree in (G3) gives the required existence parity.
They are actual pure inner forms.

Apply the stable cuspidal seed insertion argument of
Lemma~\ref{lem:sr-stable-three-seed-insertion} in ranks
three and two.
Only the rank of the local type and Jacquet equations
changes.
KMSW is used only for its proved generic-parameter
pure-inner case.
At tame places the coefficient lattice is
\(\bigotimes\sigma^{\mathrm{crys}}(\tau)^\vee\).
At \(\ell\) the ordered ordinary Jacquet vector has
exactly the \(\theta^{-1}\) coefficient character and
the normalized unit \(U_j\)-values in (H5)--(H6).
The finite integral vector, its residual localization
and classical control therefore occur in the same
ordinary modules and universal determinant quotients
used for the two runs.
The actual Hecke kernels vanish at these actual seeds.
This supplies the Gaussian \(\mathrm{seedPtOf}\)
hypotheses in both ranks.

The coefficient interpolation, ordinary projector,
projectivity, coefficient change, old/new maps and
nilpotence calculation of the common construction
are rank independent.
They divide neither by \(3\) nor by \(3!\);
the rank-three nilpotence exponent is \(3!=6\).
The normalized inertia types and ordinary characters
here have just been identified.
Accordingly the only shared support input still needed
at this interface is the completed sharp/flat
rank-two/rank-three propagation theorem, with its
actual arithmetic frame bound and its two strict
classes.
It is the theorem implemented by
\(\mathrm{PrLinkedEnd}\to\mathrm{AbsMain}
 \to\mathrm{PropagationTheoremGL3}\), using the
rank-two sharp/flat run.
Once that shared construction has been completed,
it gives, on this very frame,
\[
 \mathrm{FinPt}(\Phi,3),\qquad
 \mathrm{AutPt}(\Phi,3,\epsilon_\ell^{-1}),\qquad
 \mathrm{AutPt}(\Phi,2,k_2).
 \tag{G10}
\]
The two actual distinguished auxiliary places remove
each other's omitted-place qualification.
This is the precise use of the common theorem;
(G10) is not inferred merely from an Eisenstein
congruence or a regular algebraic realization of the seed.

\subsubsection{From the shared output to the Gaussian host endpoint}
We record the final implication so that, after (G10),
no further Gaussian automorphy premise remains.
The quotient map in \(\mathrm{GalLink}\) is by the
closed normal closure of the inertia elements on which
\(\eta\) and \(\bar\Psi\) are both trivial.
An actual marked Gaussian point therefore restricts
to the exact ordinary rank-three representation
of \(\Phi\), with inertial tuple
\((\Psi,1,\Psi^{-1})\).
The source's \(\mathrm{exists\_nrep3}\) is this
restriction and ordinary flag, with its stable lattice;
it does not posit a new representation.

First use \(\mathrm{FinPt}(\Phi,3)\) in (G10) to obtain
\(\mathrm{FinEnd}\) on the open group \(G_{L_r}\).
Characteristic-\(\ell\) fixed-inertia points then have
the residual characteristic polynomials there.
Lemma~\ref{lem:sr-constant-character-triangular}
gives upper or lower triangularity in the actual marked
basis and fixes each diagonal to its constant character.
The balanced-monomial argument in
Lemma~\ref{lem:sr-graded-finite-part-criterion},
applied exactly in
Proposition~\ref{prop:sr-host-finite-part-from-finend},
therefore makes the weight-zero universal part finite
over the coefficient ring.
These proofs reconstruct the needed
\srcid{ConstOfChar}--\srcid{UniversalDiagonalOfTriangularPoints}--%
\srcid{FinCrit} chain and are valid in characteristic three.
It is exactly the required specialization of the
finiteness clause of \(\mathrm{ThmS6Core\_s}\).
Every resulting characteristic-zero universal point
therefore admits the finite integral DVR model and
normalization used by the ordinary-point argument.
This finiteness is obtained before applying the
strong-open assertion to such a point.

For a type-I characteristic-zero host point,
Section~\ref{sec:spectrum-open-irreducibility} proves
irreducibility on every open subgroup of \(G_K\).
Thus its restriction to \(G_{L_r}\) is irreducible.
The rank-three assertion in (G10) gives actual
regular algebraic cuspidal automorphy of
\[
                             \widehat\rho|_{G_{L_r}},
                         \qquad \widehat\rho=\rho\epsilon_\ell^{-1}.
\]
Use exact cyclic descent
(Lemma~\ref{lem:sr-cyclic-exact-descent}) along
the abelian CM extension \(L_r/K\).
The finite character correction in that lemma is
included, so the descended representation is attached
exactly to \(\widehat\rho\).

The semidirect unitary law and the normalized conjugation
matrix give
\[
              \rho(g)^{\mathsf t}J_3\rho(cgc^{-1})=J_3
                           \quad(g\in G_K).
\]
Thus \(\widehat\rho\) has the exact conjugate pairing
with multiplier \(\epsilon_\ell^{-2}\).
Chebotarev and strong multiplicity one give genuine
conjugate self-duality of its rank-three automorphic
realization.
Full local--global compatibility therefore applies.
Its ordinary, finite-inertia, and full local group
conditions give the actual local window predicate,
as in Proposition~\ref{prop:sr-actual-wit};
the scalar normalization cancels from every ratio.
At \(\ell=3,7\) one may also use directly
Lemma~\ref{lem:sr-gaussian-local-ratios} and the
full comparison: there is no cyclotomic edge,
monodromy is zero, and the fixed inertial type is
unramified in Weil--Deligne normalization.
This yields the asserted crystalline condition.
No comparison of semisimplified Weil characters alone
is used to identify monodromy.

The type-I points have the automorphy just descended,
and their strong-open irreducibility was proved
independently in the cited ordinary-point argument.
Proposition~\ref{prop:sr-cyclotomic-restriction} supplies
the infinite-cyclotomic-kernel condition.
Together these give all three clauses of
\(\mathrm{BPrime\_s}\) at the actual Gaussian host,
once the explicitly stated common support input (G10)
is supplied.

This locates the remaining common dependency exactly.
The Gaussian prime-\(3\) seed, its character normalization,
its actual module insertion, the finite-unramified
real quadratic base, local frame witnesses and the
descent/local-comparison implication have been proved
here. Completion of the shared sharp/flat propagation
construction supplies (G10); none of the predicates
\(\mathrm{BPrime\_s.aut}\) or
\(\mathrm{ThmS6Core\_s}\) is used as a substitute for it.

% Frozen source:
% DnEndsTop.ends0_alt -> DnEnds.ends_u_alt -> SymTop.linkedEnd_alt
% -> LinkedEndE.linked_endomorphism_dense -> TwoRuns.twoRuns_alt.
% Seed2Alt: SdStmt, BirthTop, BirthEG, IsSeed.
% Host interface: DnHost, AutEnd3Alt, SymmetricCubeTheoremOfEndomorphisms.
% Actual local conditions: FrobeniusSmoothVectors/
% FrobeniusSplitPrimeDatum.Gen.g4,g5 and UFacts.mu,f3.


\subsection{The actual generic rank-three endpoint after propagation}
\label{subsec:sr-generic-three-endpoint}
We spell out the endpoint over the fixed field
\(K=\mathbb Q(\sqrt{-3})\).
Let \(a\ge1\), let \(\ell\) and \(\psi\) satisfy the
good-data hypotheses of Section~\ref{sec:single-character},
and put \(\Psi=\epsilon_\ell^a\psi\).
Let \(R\) be the actual marked rank-three point, let
\(\rho=R|_{G_K}\), and put
\(\widehat\rho=\rho\epsilon_\ell^{-1}\).
Assume the common sharp/flat construction has supplied,
on the actual frame and actual ordinary function modules,
\[
 \mathrm{FinPt}(\Phi,3),\qquad
 \mathrm{AutPt}(\Phi,3,\epsilon_\ell^{-1}).
 \tag{GE1}
\]
These are conclusions of the propagation construction,
rather than hypotheses about an unrelated automorphic
representation.

\begin{srproposition}\label{prop:sr-generic-three-exact-endpoint}
Under (GE1), the weight-zero part of the marked host
algebra is finite over the coefficient DVR.
Every type-I characteristic-zero point gives an actual
regular algebraic conjugate self-dual cuspidal representation
over \(K\) attached exactly to \(\widehat\rho\).
It satisfies the full local predicate \(\mathrm{WitA}\).
Moreover \(\widehat\rho\) is irreducible on
\(G_{K(\mu_{\ell^\infty})}\).
\end{srproposition}
\begin{proof}
The map from \(G_{\mathbb Q}\) to the frame group is the
actual Galois link.
Its restriction carries the marked point, its ordinary
flag and its finite inertia characters to the normalized
rank-three frame point.
Thus the first assertion of (GE1) fixes the characteristic
polynomials on the open subgroup \(G_L\subset G_K\).
For characteristic-\(\ell\) points,
Proposition~\ref{prop:sr-host-finite-part-from-finend}
proves the required finite weight-zero algebra.
Its proof uses the actual marked generators, so applies
to the generic host as well as the Gaussian host.
The integral normalization of an ordinary characteristic-zero
point can now be made as in
Section~\ref{sec:spectrum-deformation}.

For a type-I point the residual inertia separator is
\[
 \exists u\in I_\ell:\quad \overline\Psi(u)^2\ne1.
 \tag{GE2}
\]
Indeed the good-data condition \(\psi_\ell^2=1\)
gives
\[
 \overline\Psi|_{I_\ell}
      =\overline\epsilon_\ell^{\,a+b},
       \qquad b\in\{0,(\ell-1)/2\},
 \quad
 \overline\Psi^{\,2}|_{I_\ell}
      =\overline\epsilon_\ell^{\,2a}.
\]
This is the local quadratic character calculation in the
actual principal packet.
The good-data condition \(\ell-1\nmid2a\) therefore
gives (GE2), including when \(\ell\mid N\).
Also \(I_\ell\subset G_K\), since \(\ell\ne3\).
Section~\ref{sec:spectrum-open-irreducibility},
with exactly this separator and the ordinary flag,
proves irreducibility on every open subgroup of \(G_K\).
In particular \(\widehat\rho|_{G_L}\) is irreducible.
The second assertion of (GE1) therefore gives actual
regular algebraic cuspidal automorphy of
\(\widehat\rho|_{G_L}\).
The twist in (GE1) is already \(\epsilon_\ell^{-1}\):
the geometric-Frobenius convention gives
\[
 r_{\ell,\iota}(\operatorname{Ad}\pi_f)
       =\operatorname{ad}^0(r_{\ell,\iota}(\pi_f))
                                      \epsilon_\ell^{-1},
\]
which is the normalization of the rank-three seed.

The extension \(L/K\) is finite abelian and CM.
Apply Lemma~\ref{lem:sr-cyclic-exact-descent}
successively along its cyclic composition series.
The invariant representation at a descent step is the
restriction of the original \(\widehat\rho\), and its
restriction to \(G_L\) remains irreducible.
Consequently the lemma includes its finite-character
correction and gives a cusp representation attached
exactly to \(\widehat\rho\), with no undetermined
finite twist left at the endpoint.

The normalized semidirect-product law gives
\[
 \rho(cgc^{-1})=J_3\rho(g)^{-\mathsf t}J_3^{-1}.
\]
Since \(\epsilon_\ell(cgc^{-1})=\epsilon_\ell(g)\),
\[
 \widehat\rho(cgc^{-1})
       =\epsilon_\ell(g)^{-2}
                   J_3\widehat\rho(g)^{-\mathsf t}J_3^{-1}.
 \tag{GE3}
\]
The exponent is \(1-3=-2\), exactly the prescribed
rank-three Galois normalization.
The pairing is symmetric and its real multiplier is
\(\epsilon_\ell^{-2}\).
Chebotarev and strong multiplicity one therefore make
the descended regular algebraic cusp representation
genuinely conjugate self-dual.
Full local--global compatibility applies to that
actual polarized representation, including monodromy
at places above \(\ell\).

Here are the local clauses; they use no comparison of
semisimplified Weil characters to identify monodromy.
Away from \(\ell\), the actual ramification condition
kills inertia intersected with the kernel of the fixed
residual torus tuple and the quadratic sign.
This kernel contains the global commutator subgroup.
The resulting inertia image is finite abelian and its
isotypic subspaces are decomposition-stable.
Thus its Frobenius-semisimplified Weil representation
is a sum of characters and has \(N=0\).
At a place above \(\ell\), the ordinary flag has parallel
cyclotomic exponents
\[
                         (a-1,-1,-a-1)
 \tag{GE4}
\]
after the scalar twist.
The de Rham assertion comes from the actual automorphic
realization.
If the flag satisfies
\(\kappa_i\ne\kappa_j\epsilon_\ell\) for all \(i,j\),
Lemma~\ref{lem:sr-parallel-wd-characters} implies \(N=0\).

A single real norm twist makes the global cusp representation
unitary. At every place with this rank-one parameter
and \(N=0\), its generic unitary local representation
has all its inducing characters in the strict window
\[
                  q_v^{\,t-1/2}<|\xi(\varpi_v)|
                                         <q_v^{\,t+1/2}
 \tag{GE5}
\]
for the same real \(t\).
This is the character-parameter instance of the
unitary classification used in
Proposition~\ref{prop:sr-actual-wit}.
An equality
\((\kappa_i/\kappa_j)^m=\epsilon_\ell^m\)
would give a Frobenius absolute-value ratio \(q_v\)
in one order, contrary to (GE5).
Thus \(\mathrm{Unl1}\) implies \(\mathrm{PotUnl}\),
and the same \(t\) gives the window at every place
away from \(\ell\), and at every \(\ell\)-adic place
satisfying \(\mathrm{PotUnl}\).
Any other attached cusp representation has the same
unramified Hecke factors and equals this one by strong
multiplicity one.
This proves every quantifier in \(\mathrm{WitA}\),
without assuming \(\mathrm{Unl1}\) at a place where
it has not been established.

Finally the open-subgroup irreducibility already proved,
and Proposition~\ref{prop:sr-cyclotomic-restriction},
give irreducibility on \(G_{K(\mu_{\ell^\infty})}\).
The scalar inverse-cyclotomic twist changes no invariant
subspace.
\end{proof}

\paragraph{The separate original-level clause.}
The preceding proposition gives the actual
\(\mathrm{AutEnd3Alt}\) and \(\mathrm{IsAutCsdA}\)
endpoint, including its complete local window.
For \(a\ge2\), the differences of the three exponents
in (GE4) have absolute values \(a,a,2a\), none equal
to one. Thus \(\mathrm{Unl1}\) is automatic, and
(GE5) applies at every \(\ell\)-adic place.
For \(a=1\), this Hodge-number argument cannot be used:
nearest differences are one.

The source's stronger counting class
\(\mathrm{Adm3U_s}\) also requires
\[
 \mathrm{Th.Unl}\bigl(\ell,N,(a,0,-a),\rho\bigr).
 \tag{GE6}
\]
In the archive this is supplied by
\(\mathrm{UnlEndLogNorm}\).
An alternative sufficient producer in the present
argument is the actual local component assertion
\[
 \begin{gathered}
 I_Y\subset I_{\mathrm{cr}}\subset I_t,\qquad
 I_t\subset\ker x
 \quad\Longrightarrow\\[-2mm]
 \rho_x|_{\ker\eta\cap\ker\tau}
                         \text{ is crystalline}.
 \end{gathered}
 \tag{GE7}
\]
Here \(I_Y\) is the ordinary ideal saturated by
\(\ell\,dd_0dd_1\), \(I_{\mathrm{cr}}=\sqrt{I_Y}\),
and \(I_t\) is the selected component through the torus
point, all in the actual local universal ring.
Once (GE7), including the factorization of every
selected global point through \(I_t\), is proved,
its crystalline monodromy is zero.
Full local--global compatibility then yields (GE5)
without assuming \(\mathrm{Unl1}\).
That strict window prohibits
\(\kappa_i=\kappa_j\epsilon_\ell\) itself, proving
(GE6) also when \(a=1\).
The specified finite inertial type gives the original
local level.
When \(\ell\nmid N\), that type is trivial,
so the attached representation is crystalline and its
local automorphic component is unramified.
When \(\ell\mid N\), one retains the prescribed
finite type \((\psi_\ell,1,\psi_\ell^{-1})\):
the representation is crystalline after restriction to
\(\ker\eta\cap\ker\tau\), hence potentially crystalline
over the original local field, and its finite conductor
is accounted for by the existing \(\ell\)-part of \(N\).
No unramifiedness assertion is made in that case.
This is precisely what prevents an unwanted factor
of the varying coefficient prime from entering the
common integer's fixed spectrum.

Neither ordinarity alone nor the name
\(I_{\mathrm{cr}}\) proves (GE7).
The theorem
\srcid{CrystallineOrdinaryIdeals/RtPoints.lean}
uses the separate normal-form result
\srcid{ideal\_cuts\_crystalline\_locus};
its proof and the actual selected-chart factorization
are the remaining local producers for this alternative.
They are kept distinct here from the support output
(GE1).

\section{Geometric and lattice comparisons with independent proofs}
\label{sec:sr-geometric}
Recent translations by Baelson, reviewed by lowrank, identify the
principalisation, arithmetic intersection, toroidal-chart and Fock-model
branches of the original development. A statement that merely re-exports
an earlier theorem does not yet explain that theorem's proof.
We record geometric and lattice comparisons that can be proved directly here,
and specify their domains to avoid applying them beyond their hypotheses.

\subsection{An exact comparison of fractional-ideal lattices}
\label{sec:sr-ideal-lattices}
One ingredient of the high-twist class-counting argument is a
determinant-prescribed comparison of integral lattices.
Here it has a direct proof over a Dedekind domain.

\begin{srproposition}\label{prop:sr-ideal-lattices}
Let \(R\) be a Dedekind domain with fraction field \(F\).
Let \(I_1,\ldots,I_r\) and \(J_1,\ldots,J_r\) be nonzero
fractional ideals, and let \(\beta\in F^\times\), with
\[
                  \prod_{j=1}^rJ_j=\beta\prod_{j=1}^rI_j.
\]
If \(r=0\), assume also \(\beta=1\).
There exists \(A\in\mathrm{GL}_r(F)\) with
\[
 \det A=\beta,\qquad
             A\left(\bigoplus_{j=1}^rI_j\right)
                         =\bigoplus_{j=1}^rJ_j.
\]
\end{srproposition}
\begin{proof}
We first construct, for any two nonzero fractional ideals \(I,J\),
a determinant-one matrix sending \(I\oplus J\) onto \(R\oplus IJ\).
Choose \(0\ne a\in I^{-1}\).
Then \(aI\) is a nonzero integral ideal; let \(\mathcal P\) be its
finite set of prime divisors.

Choose \(b\in J^{-1}\) whose image in \(J^{-1}/PJ^{-1}\) is nonzero
for every \(P\in\mathcal P\).
Here is a direct justification of that simultaneous choice.
For any such \(P\), \(J^{-1}\ne PJ^{-1}\), since multiplication
by the invertible ideal \(J\) would otherwise give \(R=P\).
Choose \(b_P\) outside \(PJ^{-1}\).
For distinct maximal ideals the Chinese remainder construction gives
\(e_P\in R\) with
\(e_P\equiv1\pmod P\) and \(e_P\equiv0\pmod Q\) for every other
\(Q\in\mathcal P\).
Then \(b=\sum_P e_Pb_P\) has the required residues.
The Chinese remainder construction used here follows inductively
from \(P+Q=R\) for distinct maximal ideals.
If \(\mathcal P\) is empty, take \(b=0\).

Both \(aI\) and \(bJ\) are integral ideals and
\[
                         aI+bJ=R.
\]
Indeed a maximal ideal containing their sum would contain \(aI\),
so would belong to \(\mathcal P\); but \(bJ\subseteq P\) would
imply \(b\in PJ^{-1}\), contrary to the choice.
Choose \(x\in I,y\in J\) with \(ax+by=1\), and put
\[
       B=\begin{pmatrix}a&b\\-y&x\end{pmatrix},
       \qquad
       B^{-1}=\begin{pmatrix}x&-b\\y&a\end{pmatrix}.
\]
The displayed relation gives \(\det B=1\) and verifies the inverse.
For \(u\in I,v\in J\), the first coordinate \(au+bv\) belongs
to \(R\), and the second coordinate \(-yu+xv\) belongs to \(IJ\).
Conversely, for \(z\in R,w\in IJ\),
\[
                 xz-bw\in I,\qquad yz+aw\in J,
\]
using \(b\in J^{-1},a\in I^{-1}\).
Thus \(B\) induces the claimed exact lattice comparison.

For \(r\geq2\), apply these two-coordinate matrices successively
in coordinates \((1,2),(2,3),\ldots,(r-1,r)\).
All full matrices have determinant one, and their product \(B_I\)
sends
\[
       \bigoplus_j I_j\quad\hbox{onto}\quad
                             R^{r-1}\oplus\prod_j I_j .
\]
Construct \(B_J\) in the same manner.
The diagonal matrix
\(D=\operatorname{diag}(1,\ldots,1,\beta)\) sends the first
normalized lattice onto the second and has determinant \(\beta\).
Therefore \(A=B_J^{-1}DB_I\) has the required determinant and
image. For \(r=1\), take \(A=(\beta)\).
For \(r=0,\beta=1\), take the identity of the zero space,
whose determinant is one.
\end{proof}

In a number field \(F\), write
\[
 I_j=\prod_P P^{a_{j,P}},\qquad
 J_j=\prod_P P^{b_{j,P}},
\]
with integer exponents of finite support.
If
\[
              v_P(\beta)=\sum_j(b_{j,P}-a_{j,P})
                                      \quad\hbox{for every }P ,
\]
unique ideal factorization gives
\(\prod J_j=\beta\prod I_j\).
The proposition consequently produces the exact equivalence
\begin{equation}\label{eq:sr-ideal-valuations}
 \bigl(v_P(x_j)\geq a_{j,P}\ \hbox{for all }j,P\bigr)
 \quad\Longleftrightarrow\quad
 \bigl(v_P((Ax)_j)\geq b_{j,P}\ \hbox{for all }j,P\bigr).
\end{equation}
We use \(v_P(0)=+\infty\).
Membership in a fractional ideal is equivalent to these valuation
inequalities, by unique factorization and localization at the
discrete valuation rings \(R_P\).
This proves, in particular, the determinant and valuation comparison
called \srcid{SStein} in the class-counting chain leading to \(K_2\).

There is also an exact archimedean consequence.
Use the standard Minkowski embedding of \(F^r\) into
\(\mathbb R^{rr_1}\times\mathbb C^{rr_2}\), with fixed Lebesgue
measure, and let \(\operatorname{covol}\) denote lattice covolume.
The two direct sums of fractional ideals are full lattices, and
\begin{equation}\label{eq:sr-ideal-covolume}
 \operatorname{covol}\left(\bigoplus_jJ_j\right)
   =|\operatorname{Norm}_{F/\mathbb Q}\beta|\,
                  \operatorname{covol}\left(\bigoplus_jI_j\right).
\end{equation}
Indeed the real absolute Jacobian of \(A\) is
\[
 \prod_{\sigma\ {\rm real}}|\sigma(\det A)|
 \prod_{\sigma\ {\rm complex}/\mathrm{conj}}|\sigma(\det A)|^2
                 =|\operatorname{Norm}_{F/\mathbb Q}\beta|.
\]
A complex-linear matrix has real determinant equal to the square
of the absolute value of its complex determinant; this follows
by triangularization, or by writing its real and imaginary blocks.
Applying the change-of-variables formula to a lattice fundamental
parallelepiped proves \eqref{eq:sr-ideal-covolume}.
This comparison applies uniformly at every layer and twist.
The later metric estimates also require bounds on operator norms;
the exact determinant comparison supplies the covolume factor.

\subsection{Logarithmic norms after principalisation}
Let \(X\) be a connected smooth projective complex variety of dimension \(d\),
\(L\) a holomorphic line bundle with smooth positive Hermitian norm,
\(s\) a nonzero holomorphic section, and \(\alpha\) a smooth top-degree
real form on \(X\).
Then
\begin{equation}\label{eq:sr-logint}
 \int_X|\log\|s\||\,|\alpha|<\infty.
\end{equation}
In particular this applies when \(\alpha\) is a product of smooth Chern
forms of total degree \(d\).

Here is the proof, with principalisation as the classical resolution
theorem in characteristic zero. Take a proper modification
\(\pi:Y\to X\) with \(Y\) smooth, such that the zero divisor of
\(\pi^*s\) is a simple normal crossings divisor.
On a coordinate polydisk of \(Y\), with a nonvanishing local frame,
\(\pi^*s=u z_1^{a_1}\cdots z_k^{a_k}\), where \(u\) is holomorphic
and nowhere zero and the \(a_i\) are nonnegative integers.
On a slightly smaller closed polydisk, both the norm of the frame
and \(|u|\) are bounded above and below by positive constants.
Consequently
\[
 |\log\|\pi^*s\||\leq C+\sum_{i=1}^ka_i|\log|z_i||.
\]
The smooth form \(\pi^*\alpha\) has bounded coefficients there.
Fubini and polar coordinates give
\[
 \int_{|z|\leq r}|\log|z||\,dA(z)
       =2\pi\int_0^r\rho|\log\rho|\,d\rho<\infty.
\]
Thus the required absolute integral is finite on each such polydisk.
Compactness supplies a finite cover and a partition of unity.
The modification is a biholomorphism away from analytic sets of measure
zero. Change of variables there transfers the finite absolute integral
to \(X\), proving \eqref{eq:sr-logint}.

This proof covers smooth metrics. It does not assert the corresponding
uniform estimates for every singular cusp metric appearing in the
arithmetic construction. Bounds at those cusps must be checked with
their actual metric growth and boundary multiplicities.

\subsection{The shifted positive-definite lattice identity}
Let \(V\) be a \(d\)-dimensional Euclidean space and \(L\subset V\)
a full lattice of covolume \(v\). Define
\(L^*=\{\lambda:\langle\lambda,\ell\rangle\in\mathbb Z
\text{ for all }\ell\in L\}\).
For \(t>0,u\in V\),
\begin{equation}\label{eq:sr-latticetheta}
 \sum_{\ell\in L}e^{-\pi t\|\ell+u\|^2}
 =\frac{t^{-d/2}}v
   \sum_{\lambda\in L^*}
     e^{-\pi\|\lambda\|^2/t}e^{2\pi i\langle\lambda,u\rangle}.
\end{equation}
Both sides converge absolutely, locally uniformly in \(t>0,u\);
all derivatives in \(u\) and any fixed number in \(t\) do as well.

To prove this, choose a basis of \(L\) and write it as \(A\mathbb Z^d\).
The least singular value of \(A\) is positive, so
\(\|A n+u\|\geq c\|n\|-\|u\|\).
A Gaussian in \(\|n\|\) dominates the summands and every fixed derivative
on compact parameter sets; sums of these Gaussians converge by comparison
with the product of \(d\) one-dimensional Gaussian sums.
The periodization \(F(u)=\sum_{\ell\in L}e^{-\pi t\|\ell+u\|^2}\)
is smooth and \(L\)-periodic. Its Fourier coefficient at \(\lambda\in L^*\)
is, by integration over a fundamental domain and then all its translates,
\[
 \frac1v\int_Ve^{-\pi t\|x\|^2}e^{-2\pi i\langle\lambda,x\rangle}\,dx
 =\frac{t^{-d/2}}v e^{-\pi\|\lambda\|^2/t}.
\]
The integral is a product of the one-dimensional Gaussian transforms
proved in Section~\ref{sec:sr-transform}.
The coefficients form an absolutely convergent Fourier series, whose
sum therefore equals \(F\). This proves the identity and the derivative
claims.

The factor \(v\) is fixed by the Euclidean measure used in that integral.
If \(L'\subset L\) has finite index \(h\), then
\(\operatorname{covol}(L')=h\operatorname{covol}(L)\):
a fundamental domain for \(L'\) is the disjoint union, up to its boundary,
of \(h\) translates of a fundamental domain for \(L\).
For an invertible real linear map \(B\),
\(\operatorname{covol}(BL)=|\det B|\operatorname{covol}(L)\),
by the ordinary change-of-variables formula.
Thus no unspecified volume normalization is hidden in
\eqref{eq:sr-latticetheta}.

Positive definiteness is essential to this particular proof.
For an indefinite quadratic form the displayed unsmoothed Gaussian
need not even tend to zero along the lattice.
Nor does \eqref{eq:sr-latticetheta} establish the required matching of
ramified local Weil sections. The relevant Schwartz functions,
regularizations and local measures in the global theta construction
must be supplied separately.

\section{Compatible descent maps and the principal signed-data bridge}
\label{sec:sr-compatible-descent}

The geometric construction uses one arithmetic model, its integral
lattice, and its archimedean metric at a time.  We prove two additional
parts of that construction: the finite group descent maps, with their
integral and norm bounds, and the passage from the principal host's
analytic estimates to the signed correction and complementary mean.
The existence of all geometric models with uniform bounds is distinguished
from these implications at the end of the section.

\subsection{The principal host in the coordinates of the signed comparison}
\label{sec:sr-principal-coordinate-bridge}

Put \(D_h=2^{h+2}\), with \(h\geq2\), and
\[
 H=(\mathbb Z/D_h)^\times/\{1,-1\},\qquad n=|H|=2^h .
\]
Let \(q,d\) be distinct odd primes, with \(q\equiv3\pmod8\),
\(d\equiv7\pmod{16}\) and \(q\mid d+1\).  The principal host is
\[
 E=\mathbb Q(\zeta_{D_h})^+(\sqrt{-q},\sqrt{-d}).
\]
The divisibility condition is a property of the selected window pair.
In the frozen arithmetic construction,
\srcid{CharacterBoxes.ExistsPageN.WinPair} includes the field
\srcid{d_q : q \mid d+1}; \srcid{CharacterBoxes.PairGood} retains
a \srcid{WinPair} as its \srcid{pair} field.
The map \srcid{CharacterBoxes.HostOfField.hostOfField} consumes that
same \srcid{PairGood} when making the host.
Thus the principal window selection supplies \(q\mid d+1\) before
the host is passed to the geometric construction.
It is not a consequence of the weaker abstract \srcid{Host.pair}
condition.  Every application of this bridge uses a selected window
host with this retained provenance.

This is the fixed field of the subgroup specified by the dyadic
congruence in the arithmetic construction.  Its real cyclotomic
factor, rather than the full cyclotomic field, is used here.

Write \(\lambda_q,\lambda_d\) for the two quadratic residue
characters.  The cyclotomic Artin map at the modulus \(D_hqd\) is
\[
 a\longmapsto \bigl(\lambda_q(a),\lambda_d(a),[a]_{D_h}\bigr).
\]
The three coordinate maps are simultaneously surjective by the Chinese
remainder theorem.  Their fixed field is the displayed compositum;
in particular \([E:\mathbb Q]=4n\).  If
\(\lambda_q(a)=(-1)^{a_q}\) and
\(\lambda_d(a)=(-1)^{a_d}\), use the additive coordinates
\begin{equation}\label{eq:sr-host-coordinate-swap}
 \operatorname{Gal}(E/\mathbb Q)\ \longrightarrow\
       \Gamma=(\mathbb Z/2)^2\times H,\qquad
 (\lambda_q(a),\lambda_d(a),[a])\longmapsto(a_d,a_q,[a]).
\end{equation}
Thus the first sign in the signed comparison is the \(d\)-sign.
Complex conjugation is \(c=(1,1,0)\), because the dyadic factor is real.
The embeddings carrying the chosen root of \(-d\) to \(i\sqrt d\)
form exactly \(\Phi=\{(0,b,h)\}\).  The characters of
Section~\ref{sec:sr-finite} are consequently
\[
                   \chi=\lambda_d,\qquad \omega=\lambda_q .
\]

For completeness this coordinate choice gives exactly the inertia
conditions of that section.  Outside \(2,q,d\) the extension is
unramified.  At \(d\), cyclotomic inertia maps onto
\(\{0,(1,0,0)\}\).  At \(q\), the field
\(\mathbb Q(\sqrt{-d})\) is split: the congruence
\(-d\equiv1\pmod q\) lifts to a square in \(\mathbb Q_q\).
It is also split at two, because \(-d\equiv1\pmod8\).
For the latter assertion start with an odd root modulo \(8\);
if \(z^2+d=2^j u\), \(j\geq3\), retain \(z\) when \(u\)
is even and replace it by \(z+2^{j-1}\) when \(u\) is odd.
Starting from \(z_3=1\), this defines odd integers \(z_j\) with
\(z_j^2+d\equiv0\pmod{2^j}\) and
\(z_{j+1}-z_j\in\{0,2^{j-1}\}\).
Indeed the change of its square is
\(2^jz_j+2^{2j-2}\equiv2^j\pmod{2^{j+1}}\).
The sequence is Cauchy in \(\mathbb Z_2\), since the valuations of
its successive differences tend to infinity.
Its limit \(z\) satisfies \(z^2+d=0\), by continuity.
Literal compatibility of the intermediate roots modulo \(2^j\)
is not required.
The decomposition groups at \(2,q\) therefore act trivially on the
\(d\)-quadratic subfield.  They lie in \(\Phi\); both their inertia
groups and their Frobenius representatives lie in \(\Phi\).
These are precisely the host conditions used in the signed balance.

Let \(\theta\in\widehat\Gamma\), and let \(\theta^*\) be the primitive
Dirichlet character inducing its cyclotomic Artin character.
Use its actual normalization
\[
                    \ell_\theta=L'(1,\theta^*)/L(1,\theta^*)
                    \quad(\theta\ne1).
\]
Then the inertial coefficient in the signed comparison is exactly
\begin{equation}\label{eq:sr-primitive-inertial-coefficient}
 \frac1{|I_p|}\sum_{u\in I_p}\theta(rF_p+u)
                         =\theta^*(p^r)\qquad(r\geq1).
\end{equation}
Indeed a nontrivial character on \(I_p\) has sum zero, and
the primitive Dirichlet character is ramified at \(p\) in exactly
this case.  If the inertia restriction is trivial, the sum is
\(\theta(F_p)^r\); the primitive character is unramified and its
value at \(p\) is \(\theta(F_p)\).
These assertions follow directly in the cyclotomic realization:
the \(p\)-primary unit factor is the inertia group, a character
kills that factor exactly when its conductor is prime to \(p\),
and Frobenius on the prime-to-\(p\) cyclotomic factor acts by \(p\).
This argument includes primes dividing the larger host modulus but
not the conductor of \(\theta^*\).
Thus all the \(G_\theta\), at both thick and thin nodes, are the
same centered primitive-character ramps as those in the host
analytic estimates.

Here are the estimates needed from that host, stated without a
predicate name.  For the three nontrivial characters of
\(\Gamma/H\), assume
\begin{equation}\label{eq:sr-host-three-quadratics}
 |G_\chi(x)|,\ |G_{\chi\omega}(x)|,\ |G_\omega(x)|
                         \leq x^{1/2+\epsilon}.
\end{equation}
For every \(\theta\in\widehat\Gamma\), assume
\begin{equation}\label{eq:sr-host-thin-analytic}
 |G_\theta(y)|\leq y\exp(-c x^{\epsilon/8})
 \quad\left(y\geq \tfrac13\exp(x^{\epsilon/3})\right),
 \qquad c>0 .
\end{equation}
The first estimate is the principal specialization of the source's
bounded-character estimate, rather than an additional independence
assumption.  To check this, the retained primitive set for the
cyclic subgroup generated by the principal character is \(\{1\}\).
The source first requires triviality on \(H_0\), and here \(H_0=H\).
It then bounds the characters taking value \(-1\) at
\(e=(1,0,0)\), together with those taking value \(1\) there that are
not retained.  The first class contains \(\chi,\chi\omega\);
the second contains \(\omega\), since it is nontrivial.
The surjectivity of the Artin map ensures that none of these
three primitive characters is the principal character.
This unpacks the definitions
\srcid{HostRaw.Bounded}, \srcid{BoundedCharsPointwise} and
\srcid{EAtH.bdG}; it verifies which three bounds are being used.

\begin{srproposition}\label{prop:sr-host-signed-analytic-bridge}
On this principal host, \eqref{eq:sr-host-three-quadratics}
implies the signed correction inequality
\(\mathcal C(x)\leq3x^{1/2+\epsilon}\).
Suppose also that \eqref{eq:sr-host-thin-analytic} holds and that
the thin weights satisfy
\[
 0<w_{u,j}(\rho)\leq2b',\qquad
 b'\exp(x^{\epsilon/3})\leq xb,\qquad
 |a_t/b-x|\leq1,\qquad x\geq3 .
\]
If
\begin{equation}\label{eq:sr-thin-host-budget}
       2f(x+1)\exp(-c x^{\epsilon/8})\leq x^{-1},
\end{equation}
then the complementary mean for these same data satisfies
\(\mathcal R\geq-Q/x\).
\end{srproposition}
\begin{proof}
Taking real parts and applying the triangle inequality to the three
bounds gives the correction inequality.
Put \(X'=\exp(x^{\epsilon/3})\) and
\(\vartheta=\exp(-c x^{\epsilon/8})\).  At a thin node,
\[
 \frac{a_t}{w_{u,j}(\rho)}
   \geq \frac{(x-1)b}{2b'}
   \geq\frac{x-1}{2x}X'\geq X'/3 .
\]
Consequently the modulus of its weighted ramp is at most
\(a_t\vartheta\leq(x+1)b\vartheta\).
In the definition of \(\mathcal R\), for each \(u,t,j\)
there are at most \(f\) summation points, \(2f\) characters and
the prefactor \(1/f\).  Taking absolute values before averaging
over \(u\) gives
\[
 \mathcal R\geq-2f|T|m(x+1)b\vartheta\geq-|T|mb/x .
\]
The last step is \eqref{eq:sr-thin-host-budget}.
Every ramp and every normalization in this calculation belongs to the
one host just fixed, by \eqref{eq:sr-primitive-inertial-coefficient}.
\end{proof}

The threshold in this implication is uniform over the host fields
on the prescribed dyadic size range.  In fact
\(n\leq2x^{2+8\epsilon}\) gives \(f\leq4x^{2+8\epsilon}\).
It is enough to require
\[
       c x^{\epsilon/8}\geq\log16+(4+8\epsilon)\log x .
\]
This holds above a threshold depending on the fixed
\(\epsilon,c\), before choosing any field or node, since
\(x^{\epsilon/8}/\log x\longrightarrow\infty\).

The same observation deals with the repair inequality without a
hidden dependence on the number of nodes.  If
\(f\leq A_f x^{B_f}\), \(\eta\leq A_\eta x^{B_\eta}\),
where the nonnegative exponents and constants are fixed on the
family, and \(b=\exp(x^\epsilon)\), then
\[
 6f x(x+1)(\eta+1)
 \leq12A_f(A_\eta+1)x^{B_f+B_\eta+2}\leq b
\]
eventually, uniformly over that family.
Finally a cutoff can be chosen as the ceiling of the largest of the
finitely many positive ratios \(a_t/w_{u,j}(\rho)\), increased to
one if necessary.  This chooses one cutoff for all the nodes and
changes none of their weights.

\subsection{Integral twisted traces and the exact norm loss}
\label{sec:sr-twisted-trace-descent}

Let \(K\) be a number field, \(R\) a commutative \(K\)-algebra and
\(S\) a finite group of order \(s\).
Suppose \(\rho_g\) are \(K\)-algebra automorphisms satisfying the
right action convention
\[
             \rho_1=\operatorname{id},\qquad
             \rho_{hg}=\rho_g\circ\rho_h .
\]
Fix a character \(\xi:S\to K^\times\), and set
\[
 \operatorname{Tr}_\xi(F)=\sum_{g\in S}\xi(g)^{-1}\rho_g(F),
 \qquad
 A(u)_{h,g}=\rho_{hg}(u),\qquad \Delta(u)=\det A(u).
\]
Assume that multiplication by \(\Delta(u)\) is injective on \(R\).

Let \(V_-,V,V_+\) be finite dimensional \(K\)-subspaces of \(R\).
Assume that \(V\) is \(S\)-stable and that
\[
       V_-\,\rho_h(u)\subset V,\qquad V\,\rho_h(u)\subset V_+
                         \quad(h\in S).
\]
Put \(N=V\cap R_\xi\), where
\(R_\xi=\{F:\rho_g(F)=\xi(g)F\text{ for all }g\}\), and define
\[
             \varphi_h(f)=\operatorname{Tr}_\xi(f\rho_h(u)).
\]

\begin{srproposition}\label{prop:sr-twisted-trace-descent}
The maps \(\varphi_h:V_-\to N\) have zero common kernel, and
\begin{equation}\label{eq:sr-descent-ranks}
           \dim_K V_-\leq s\dim_K N\leq\dim_K V_+ .
\end{equation}
Suppose lattices in the three spaces are preserved by the indicated
actions and products, \(u\) is integral, and the lattice in \(N\)
is \(L_\xi=L\cap R_\xi\), for the chosen lattice \(L\subset V\).
Then every \(\varphi_h\) carries the source lattice into \(L_\xi\).

At each complex embedding let the target norm be Hermitian,
let every \(\rho_g\) act isometrically on \(V\), and assume
\[
                  \|f\rho_h(u)\|_V\leq B\|f\|_{V_-}
                  \quad(h\in S),\qquad B\geq0 .
\]
For these very same maps,
\begin{equation}\label{eq:sr-descent-norms}
                \sum_{h\in S}\|\varphi_h(f)\|_V^2
                              \leq s^3 B^2\|f\|_{V_-}^2 .
\end{equation}
If the integral lattice \(L_\xi\) does not span \(N\), its rank
still satisfies \(s\operatorname{rk}L_\xi\leq\dim_K V_+\).
\end{srproposition}
\begin{proof}
Reindexing \(g\mapsto gk\) gives
\(\rho_k(\operatorname{Tr}_\xi F)=\xi(k)\operatorname{Tr}_\xi F\).
The stability and product assumptions therefore put each
\(\varphi_h(f)\) in \(N\).

If every \(\varphi_h(f)\) is zero, put
\(v_g=\xi(g)^{-1}\rho_g(f)\).  Expanding the trace gives
\[
                    (A(u)v)_h=\varphi_h(f)=0 .
\]
Multiplication by the adjugate gives \(\Delta(u)v=0\).
The assumed injectivity of multiplication by \(\Delta(u)\)
forces \(v=0\); its identity coordinate is \(f\).  Thus
\(f\mapsto(\varphi_h(f))_h\) is an injection \(V_-\to N^S\),
proving the first rank inequality.

For the other inequality use
\(\Psi:N^S\to V_+\), \(\Psi(F)=\sum_h F_h\rho_h(u)\).
If \(\Psi(F)=0\), applying \(\rho_g\) and using the common
eigencharacter of the \(F_h\) yields
\[
                  \sum_h \rho_{hg}(u)F_h=0\qquad(g\in S).
\]
Thus \(A(u)^t F=0\).  Its determinant is again \(\Delta(u)\),
so the same adjugate argument proves that \(\Psi\) is injective.
This proves the second rank inequality.

Every value of \(\xi\) is a root of unity, since \(S\) is finite.
Its inverse is an algebraic integer and a unit, so the
unnormalized trace preserves integral lattices.  The preceding
eigencharacter computation proves integrality in \(L_\xi\),
without dividing by \(s\).

At a complex embedding the same roots of unity have modulus one.
For each \(h\), the triangle inequality, invariance of the target
norm and the multiplication bound give
\[
 \|\varphi_h(f)\|_V
 \leq\sum_g\|\rho_g(f\rho_h(u))\|_V
 \leq sB\|f\|_{V_-}.
\]
Squaring and summing over the \(s\) possible \(h\) proves
\eqref{eq:sr-descent-norms}.  Finally
\(\operatorname{rk}L_\xi=\dim_K(KL_\xi)\leq\dim_K N\),
which proves the last assertion.
\end{proof}

In the geometric application the three spaces are cusp sections at
weights \(G-W,G,G+W\) on one cover, \(S\) is its prescribed level
subgroup, and \(\rho\) is pullback by level changes.
The equality between the deep type lattice and \(L_\xi\), preservation
of the cusp condition, integrality of the pullbacks and isometry of
their Petersson metrics must all refer to this same cover.
Under those identifications the proposition constructs the maps,
rather than merely assuming their existence.
The source names are \srcid{mapsCore},
\srcid{twistedTrace\_equivariant\_and\_injective\_of\_groupDet}
and \srcid{finrank\_le\_card\_mul\_finrank\_eigenspace\_le}.
This is the map input used in the cover-to-covolume argument of
Section~\ref{sec:cov-lattice}.

\subsection{One integral section with all component determinants nonzero}
\label{sec:sr-bounded-determinant-selection}

We next prove the simultaneous integral selection behind the preceding
maps.  Let
\[
                      R=\prod_{Z\in\mathcal Z}R_Z
\]
be a finite product of integral \(K\)-algebras, each embedded in a
field \(\Omega_Z\), and suppose \(S\) permutes the factors through
its action.  Write \(e_Z\) for the component idempotents.
Let \(W\subset R\) be a finite dimensional subspace decomposed as
\(\bigoplus_Z e_ZW\).
For each \(Z\), choose integral elements
\(w_{Z,1},\ldots,w_{Z,n_Z}\) spanning \(e_ZW\) over \(K\).
Suppose there is \(u_0\in W\) for which
\[
                       \Delta(u_0)_Z\ne0\quad(Z\in\mathcal Z).
\]

\begin{srproposition}\label{prop:sr-bounded-determinant-selection}
There are integers \(0\leq m_{Z,i}\leq s^2\) such that
\[
 u=\sum_{Z,i}m_{Z,i}w_{Z,i},\qquad
                       \Delta(u)_Z\ne0\quad(Z\in\mathcal Z).
\]
In particular \(\Delta(u)\) is a nonzero divisor in \(R\).
Suppose \(n_Z\leq N\), and at every archimedean evaluation point
only its one component contributes.
If the nonnegative pointwise norms satisfy the triangle inequality
and every \(w_{Z,i}\) has norm at most \(\beta\Lambda\) on that
component, then \(u\) has pointwise norm at most
\begin{equation}\label{eq:sr-short-determinant-bound}
                              Ns^2\beta\Lambda .
\end{equation}
The bounds are independent of \(|\mathcal Z|\).
\end{srproposition}
\begin{proof}
Consider changing a single coefficient, so the element has the form
\(v+t w_{Z,i}\).  For a component \(Z'\), its group determinant
is a polynomial in \(t\) of degree at most \(s\).
If it is nonzero at the current coefficient, this polynomial
is nonzero, and has at most \(s\) roots in the characteristic-zero
field \(\Omega_{Z'}\).

Only the components in the \(S\)-orbit of \(Z\) can be affected.
Indeed \(\rho_g(w_{Z,i})\) is supported on the component to which
\(\rho_g\) sends \(Z\); all entries in \(A(w_{Z,i})\) vanish at
every other component.  That orbit has at most \(s\) members.
Consequently there are at most \(s^2\) bad integers for this
coefficient.  Among \(0,\ldots,s^2\) choose a good one.
Begin with the coefficients of \(u_0\) in the spanning family and
replace them one at a time.  At every step all component determinants
remain nonzero, so the finite procedure gives the asserted \(u\).
It is integral because its coefficients are ordinary nonnegative
integers.  A product of domains has a nonzero divisor whenever all
of its components are nonzero, proving the next assertion.

At an evaluation point on \(Z\), all generators supported on other
components vanish.  The triangle inequality gives
\[
            \|u\|\leq\sum_{i=1}^{n_Z}m_{Z,i}\|w_{Z,i}\|
                          \leq Ns^2\beta\Lambda ,
\]
which proves the norm estimate.
\end{proof}

The initial nonzero-determinant condition can itself be verified
by orbit evaluations.  Suppose for each component one can find an
evaluation in its function field, together with compatible trivializations
of the section fibers, such that the functionals
\[
                    w\longmapsto \operatorname{ev}_Z(\rho_g w),
                         \qquad g\in S ,
\]
on \(W\) are linearly independent after extending scalars to that
field.  Their joint map is then onto.  Choose values equal to one
at the identity and zero at all other \(g\).
The evaluation of \(A(w)\) is the permutation matrix
\(({\bf1}_{hg=1})_{h,g}\), and has determinant \(1\) or \(-1\).
Thus the determinant polynomial on \(W\) is not identically zero
on that component.

These finitely many nonzero polynomials have a simultaneous
nonvanishing specialization with coefficients in \(K\).
One way to check this even when their coefficients lie in different
extension fields is the following.  For a polynomial with coefficients
in an extension of \(K\), apply a \(K\)-linear functional on the
finite span of its coefficients that is nonzero on one coefficient.
This gives a nonzero polynomial over \(K\); nonvanishing of this
polynomial implies nonvanishing of the original one.
The product of the finitely many resulting polynomials is nonzero.
An infinite field has a point at which a given nonzero polynomial
does not vanish, by induction on its number of variables.
This gives the required \(u_0\).  Thus orbit separation and the
bounded integral generators give one integral section satisfying
all the determinants simultaneously.

Here is the uniform numerical consequence used in the actual descent.
Put \(z\geq1\).  Suppose constants \(C_1,C_3\geq0\), chosen before
all covers, give
\[
 \log s\leq120z^2,\qquad
 N\leq e^{C_3z^3},\qquad \beta\leq e^{C_1z^3},\qquad
                        \Lambda\leq e^{2C_3z^9}.
\]
The same selected section therefore has pointwise norm at most
\[
                         e^{(C_1+3C_3+240)z^9}.
\]
No component count occurs in the constant.
For an invariant Petersson \(L^2\) metric, a pointwise bound \(B\)
for multiplication by \(u\) gives
\(\|f u\|_2\leq B\|f\|_2\), by integration of the squared
pointwise inequality.  Pullback isometry gives this same bound for
every \(\rho_hu\).  Hence
\[
                         C=C_1+3C_3+240
\]
is a uniform exponent in the preceding descent proposition.

At the selected type level the source has the sharper group estimate
\(\log s\leq45z^2\).
If \(b\geq z^{12}\) and \(z\geq2C+200\), the maps satisfy the
required sum-of-squares bound with multiplier \(e^{2b}\):
\[
 \log(s^3e^{2Cz^9})
       \leq135z^2+2Cz^9
       \leq(135+2C)z^9
       \leq z^{10}\leq z^{12}\leq b\leq2b .
\]
This fixes the constants before selecting the geometric cover.
It does not choose a new section separately for each descent map.

For the archive these bounds come from
\srcid{ExistsShortIntegralVectorOfDeterminantBound},
using \srcid{IntegralCuspSlopePowerBound},
\srcid{ShortGeneratorNormComparison},
\srcid{ExistsCuspRegularLevelDeterminant} and
\srcid{CoveringDomainLatticeRegularity}.
The integer coefficient argument is the mathematical content of
\srcid{CardLeOfDetEqZero.exists\_short}.
The proof above supplies the finite selection and descent calculations.
To apply it to every member of the required family one must still
supply the geometric orbit separation, the short spanning families,
and their stated uniform metric comparisons.

\subsection{The simultaneous geometric input still required}
\label{sec:sr-geometric-selection-residual}

The preceding constructions reduce, but do not close, the geometric
existence problem.  Their exact residual inputs are as follows.

First, on every chosen cover the scalar Hodge-weight cusp space must
have the orbit-evaluation separation used above, and its actual
integral lattice must have the short spanning family with
\(N,\beta,\Lambda\) as displayed.
The source constructs orbit separation at its weight
\(c_0\,\mathrm{hodgeWt}\), where the specified \(c_0\) is large enough
for the quotient point-separation theorem.
Its proof uses a descended linearized Hodge line on a finite quotient,
the canonical-bundle identity, nefness and bigness, boundary vanishing
and the applicable point-separation theorem.
Qualitative existence at some larger weight would not prove this
statement at the prescribed weight.

Second, the theta realization, its local factors and the arithmetic
model must be tied to that same cover and those same metrics.
The definite Siegel--Weil theorem is the value of a doubled
Eisenstein series on a specified box section.
In the transposed doubling convention its parameter is
\((d-n)/2\), and its value is the normalized theta average multiplied
by \(1\) when \(n<d\), and by \(2\) when \(d\leq n\).
Meromorphic continuation and this definite value alone do not give
the agreement of two indefinite theta integrals with their ramified
matching maps.
The latter also requires the local intertwining, nonvanishing,
tower and constant-term comparisons, and the extension from
polynomial test vectors to the actual Schwartz data.
The inspected source producer is
\srcid{TheoremKOfMod.theoremK\_head2}; its final assembly
\srcid{SphericalClosure.CloseSph.closeSph} uses precisely these
comparisons.  The new \srcid{ThkBaseNormZeta} theorem gives the
holomorphic value from anisotropic Siegel--Weil and continuation
in the transposed box coordinates; it is one part of this chain.

Third, the arithmetic intersection and cover construction must give
the top inequality \eqref{eq:sr-top} with the four-term error
\eqref{eq:sr-error} for the same selected data.
The descent maps just constructed supply the map and rank part of
the cover-to-covolume argument.  They do not by themselves construct
the deep covers, the degree and rank shares, the component
Hilbert--Samuel comparison, or the height lower bound.
Their source assembly is
\srcid{HodgeBundleWeights.TopCovolG.topCovolG\_alt}, followed by
\srcid{ThirdPlaceSeries.TopCovolSlim2Use.topCovolSlim2\_use\_alt\_alt}.

The selection procedure in the source is more precise than taking
unrelated local choices.  Its
\srcid{nonempty\_hyps\_definite\_side\_of\_closed\_bound}
first takes the maximum of the row-functional threshold \(C_{F,0}\)
and the evaluation-grid threshold \(C_{F,1}\).
It then uses the same extension and the same grid
\(\mathrm{GridA}(\max(C_{F,0},C_{F,1}))\)
for the row functions, local matching, covolume and linked norm bounds.
The later \srcid{nonempty\_hyps\_of\_allRows\_alt}
intersects the finitely many row conditions before selecting its one
family class and its one raw selection.
Its projections provide the spectral data, base rows, height data,
columns and indefinite realization from that raw selection.
This explains how the archive avoids separate incompatible choices.
A standalone construction must prove the component inputs of this
procedure, not just restate its nonempty-package conclusion.

For any finite list of estimates already proved uniformly on such
a selection, taking the maximum of their lower scale thresholds
makes them hold simultaneously.
For constants this argument has to be performed before the family
or host is chosen.  A statement
\(\forall E\,\exists C_E\) cannot supply the fixed \(C\) in
Definition~\ref{def:sr-supply}.
The principal coordinate bridge above does prove that the signed
correction and the complementary mean transfer simultaneously from
one analytically controlled host; the remaining top comparison and
its compatible uniform geometric realization are not established
by this section.

\section{Integral short generators and uniform metric constants}
\label{sec:sr-short-generators}

We prove the arithmetic lattice step in the construction of a short
descent section.  The proof keeps the dependence on the coefficient
field polynomial in its degree and root discriminant.  In particular
it does not require a bound for the index of a chosen short spanning
lattice.

\begin{srlemma}[Quadratic Euclidean transference]
\label{lem:sr-quadratic-transference}
Let \(L\subset\mathbb R^n\) be a full Euclidean lattice, and suppose
every nonzero element of \(L\) has norm at least \(\lambda>0\).
The dual lattice \(L^*\) contains \(n\) linearly independent vectors
of norm at most \(n^2/\lambda\).
\end{srlemma}

\begin{proof}
Minkowski's convex body theorem applied to a cube contained in the
Euclidean ball gives
\(\lambda_1(M)\le\sqrt n\,\det(M)^{1/n}\) for every full lattice \(M\).
Apply this first to \(L\) and then to \(L^*\).  Since
\(\det(L)\det(L^*)=1\), it gives
\[
 \lambda_1(L^*)\le n/\lambda.
\]
Choose a shortest nonzero dual vector and divide out its integral
content to get a primitive \(y_0\in L^*\), still with
\(\|y_0\|\le n/\lambda\).
Primitivity means that there exists \(x_0\in L\) with
\(\langle y_0,x_0\rangle=1\): in a \(\mathbb Z\)-basis of \(L\), the
coordinates of the integral functional \(y_0\) have greatest common
divisor one.

Put \(H=y_0^\perp\) and \(L_H=L\cap H\).
The homomorphism \(L\to\mathbb Z\), \(x\mapsto\langle y_0,x\rangle\),
is onto, so \(L=L_H\oplus\mathbb Zx_0\) as an abelian group.
Thus \(L_H\) is a full lattice in the \((n-1)\)-dimensional space \(H\),
and its shortest nonzero vector has norm at least \(\lambda\).
Inductively choose independent
\(w_1,\ldots,w_{n-1}\in L_H^*\subset H\) with
\(\|w_i\|\le(n-1)^2/\lambda\).
Let
\[
 c_i=\langle w_i,x_0\rangle-
       \operatorname{round}(\langle w_i,x_0\rangle),\qquad
 y_i=w_i-c_i y_0.
\]
Then \(|c_i|\le1/2\), and \(y_i\) is integral on both \(L_H\) and
\(x_0\), so \(y_i\in L^*\).
Its projection to \(H\) is \(w_i\), hence
\(y_0,y_1,\ldots,y_{n-1}\) are independent.  Orthogonality gives
\[
 \|y_i\|^2
 \le\frac{(n-1)^4+n^2/4}{\lambda^2}
 \le\frac{n^4}{\lambda^2}.
\]
Together with the bound for \(y_0\), this completes the induction.
The assertions for \(n=0\) and \(n=1\) follow directly; the displayed
induction is used for \(n\ge2\).
\end{proof}

\begin{srlemma}[Small independent algebraic integers]
\label{lem:sr-small-integers}
Let \(K\) be a number field of degree \(\delta\), discriminant \(D_K\),
and root discriminant \(\operatorname{rd}(K)=|D_K|^{1/\delta}\).
There exist \(\delta\) elements
\(\alpha_1,\ldots,\alpha_\delta\in\mathcal O_K\), independent over
\(\mathbb Q\), such that
\[
 |\sigma(\alpha_i)|\le
 \delta^3\operatorname{rd}(K)
 \quad\hbox{for every embedding }\sigma:K\hookrightarrow\mathbb C.
\]
\end{srlemma}

\begin{proof}
Embed \(K\) in its Minkowski Euclidean space, with norm squared
\(\sum_\sigma|\sigma(x)|^2\), counting both complex embeddings.
The Euclidean dual of the image of \(\mathcal O_K\) is, up to an
orthogonal reflection in the imaginary coordinates, the image of the
codifferent
\(\mathfrak D_K^{-1}=\{y:\operatorname{Tr}_{K/\mathbb Q}
(y\mathcal O_K)\subset\mathbb Z\}\).
This follows by choosing an integral basis and its trace-dual basis:
reflection changes the real Euclidean pairing
\(\operatorname{Re}\sum_\sigma\sigma(x)\overline{\sigma(y)}\)
into \(\operatorname{Tr}(xy)\).

The codifferent has ideal norm \(|D_K|^{-1}\).
For \(0\ne y\in\mathfrak D_K^{-1}\), the inclusion
\(y\mathcal O_K\subset\mathfrak D_K^{-1}\) implies
\(|N_{K/\mathbb Q}y|\ge |D_K|^{-1}\).
The arithmetic--geometric mean inequality therefore gives
\[
 \sum_\sigma|\sigma(y)|^2
 \ge \delta |N_{K/\mathbb Q}y|^{2/\delta}
 \ge \delta\,\operatorname{rd}(K)^{-2}.
\]
In particular every nonzero vector of this dual lattice has norm at
least \(\operatorname{rd}(K)^{-1}\).
Apply Lemma~\ref{lem:sr-quadratic-transference} to the dual lattice.
Its dual, the original integer lattice, has \(\delta\) independent
vectors of Euclidean norm at most
\(\delta^2\operatorname{rd}(K)\).
Each individual embedding is bounded by that norm.
Enlarging \(\delta^2\) to \(\delta^3\) gives the stated form.
\end{proof}

\begin{srproposition}[Short spanning vectors give integral generators]
\label{prop:sr-short-integral-generators}
Let \(V\) be an \(R\)-dimensional vector space over \(K\), and let
\(M\subset V\) be a finitely generated \(\mathcal O_K\)-submodule
spanning \(V\).  For every embedding \(\sigma\), let \(N_\sigma\) be
a seminorm satisfying
\[
 N_\sigma(av)=|\sigma(a)|N_\sigma(v),\qquad
 N_\sigma(v+w)\le N_\sigma(v)+N_\sigma(w).
\]
Suppose \(S\subset M\) spans \(V\) over \(K\) and
\(N_\sigma(s)\le A\) for every \(s\in S\) and every \(\sigma\).
Then \(M\), as an abelian group, is generated by at most \(2\delta R\)
vectors \(w_i\) such that
\begin{equation}\label{eq:sr-generator-bound}
 N_\sigma(w_i)\le
 \max\{1,\delta^4R\operatorname{rd}(K)A\}
 \quad\hbox{for every }i,\sigma.
\end{equation}
\end{srproposition}

\begin{proof}
The case \(R=0\) is immediate.  Choose a \(K\)-basis
\(v_1,\ldots,v_R\) from \(S\), and choose the integers \(\alpha_j\)
of Lemma~\ref{lem:sr-small-integers}.
The \(\delta R\) vectors \(e_{j,k}=\alpha_jv_k\) form a
\(\mathbb Q\)-basis of \(V\); they lie in \(M\), and every seminorm
is at most
\(A'=\delta^3\operatorname{rd}(K)A\).
Write \(D=\delta R\), enumerate this basis by \(e_1,\ldots,e_D\),
and put \(M_0=\sum_\ell\mathbb Ze_\ell\subset M\).

The group \(M\) is free of rank \(D\): it is finitely generated and
torsion free over \(\mathbb Z\), and its rational span is \(V\).
Choose any \(\mathbb Z\)-basis \(b_1,\ldots,b_D\) of \(M\).
No norm bound for this basis is needed.
If \(b_i=\sum_\ell q_{i\ell}e_\ell\), define
\[
 \widetilde b_i=\sum_\ell
       (q_{i\ell}-\lfloor q_{i\ell}\rfloor)e_\ell.
\]
Then \(b_i-\widetilde b_i\in M_0\), so \(\widetilde b_i\in M\).
The \(2D\) vectors \(e_\ell,\widetilde b_i\) generate \(M\), since
they generate every \(b_i\); and
\[
 N_\sigma(\widetilde b_i)
 \le\sum_\ell
       (q_{i\ell}-\lfloor q_{i\ell}\rfloor)N_\sigma(e_\ell)
 \le DA'.
\]
The same bound holds for \(e_\ell\), since \(D\ge1\).
Thus \(DA'=\delta^4R\operatorname{rd}(K)A\) proves
\eqref{eq:sr-generator-bound}.  The index \([M:M_0]\) never enters.
\end{proof}

Here is the precise uniformity consequence used for descent sections.
Let \(z\ge1\) be a size parameter, let \(S_{\rm wt}\ge0\) be the sum
of absolute values of the highest-weight entries, and suppose the
chosen coefficient fields and component cusp lattices satisfy
\begin{align}
 \log\delta&\le391z^2,&
 \log\operatorname{rd}(K)&\le391z^2,\label{eq:sr-field-metric-size}\\
 \log\max(1,R)&\le C_bz^3,&
 \log A&\le C_2(z^3S_{\rm wt}+z^5).\label{eq:sr-span-metric-size}
\end{align}
Here \(C_b,C_2\ge0\) are fixed before fields, components, levels, and
weights are chosen.  Proposition~\ref{prop:sr-short-integral-generators}
gives generators of norm at most \(\Lambda\ge1\), in number \(N\), with
\begin{equation}\label{eq:sr-uniform-short-generators}
 \log\Lambda\le C_3(z^3S_{\rm wt}+z^5),\qquad
 \log\max(1,N)\le C_3z^3,\qquad
 C_3=2000+C_b+C_2.
\end{equation}
Indeed the first logarithm is bounded by the maximum of zero and
\[
 4\log\delta+\log R+\log\operatorname{rd}(K)+\log A
 \le1955z^2+C_bz^3+C_2(z^3S_{\rm wt}+z^5),
\]
which is at most the first right side of
\eqref{eq:sr-uniform-short-generators}; and
\(\log N\le\log2+\log\delta+\log R\le(392+C_b)z^3\).
For a zero component space take the empty generating family and
\(\Lambda=1\).

The component sup norm and rank bounds have a common Hilbert-space
normalization.  Let \(\mathcal H\) be a finite-dimensional Hilbert
space of line-bundle sections on a component with a probability measure \(\mu\),
and assume its norm is the corresponding \(L^2\) norm.
If point evaluation has squared norm at most \(\beta\), then
\[
 |f(y)|^2\le\beta\|f\|_2^2,\qquad
 \dim\mathcal H\le\beta.
\]
To prove the second assertion, choose an orthonormal basis \(f_i\).
The evaluation estimate is equivalent to
\(\sum_i|f_i(y)|^2\le\beta\).
Integrating and using \(\mu(X)=1\) gives
\(\dim\mathcal H=\sum_i\int|f_i|^2\le\beta\).
Thus, if
\begin{equation}\label{eq:sr-bergman-input}
 1\le\beta\le e^{C_1z^3},
\end{equation}
both the pointwise factor \(\sqrt\beta\le\beta\) and the logarithm of
the component rank are controlled by this one fixed \(C_1\).
The measure must be the same normalized component measure used in
the integral lattice metric; an unnormalized quotient measure would
insert its volume into the rank inequality.

For completeness we now combine these bounds with the simultaneous
group determinant construction.  Suppose the prescribed scalar
Hodge weight satisfies \(S_{\rm wt}\le z^5\), the finite level group
has cardinal \(s\) with \(\log s\le120z^2\), and a cusp section with
nonzero group determinant exists on every component.
Apply the componentwise integer-coefficient selection proved in
Section~\ref{sec:sr-compatible-descent}.  It uses coefficients in
\(\{0,\ldots,s^2\}\) in a single simultaneous choice.
The resulting integral section has pointwise norm at most
\(N s^2\beta\Lambda\), hence multiplication by any of its level
translates has norm at most \(B\), where
\begin{equation}\label{eq:sr-uniform-multiplier}
 \log B\le (C_1+3C_3+240)z^9.
\end{equation}
This follows from
\(\log N\le C_3z^3\),
\(\log\Lambda\le2C_3z^9\),
\(\log\beta\le C_1z^3\), and \(2\log s\le240z^2\).
Every constant in \eqref{eq:sr-uniform-multiplier} is chosen before
the level group or component is selected.

The integral maps
\(\varphi_h(f)=\sum_g\xi(g)^{-1}\rho_g(f\rho_hu)\)
therefore satisfy
\[
 \sum_h\|\varphi_h(f)\|^2\le s^3 B^2\|f\|^2.
\]
This uses only the isometry of level translations and the multiplier
bound just proved.  There is no division by \(s\) in the maps, so
their integrality is retained at primes dividing \(s\).

These arguments close the passage from a uniformly short integral
\(K\)-spanning family to the simultaneous integral descent maps.
For an actual Shimura family the remaining input
\eqref{eq:sr-span-metric-size} must concern the prescribed algebraic
theta sections, their common integral lattice, and the same
Petersson metric.  Abstract finite generation does not prove it.

More precisely, the frozen construction supplies this input in the
following order.  Its theta spanning assertion uses the identical
restricted tensor model, the integral Fourier datum
\(\widehat{\mathbf1}_{x_0}\), unimodular positive definite auxiliary
lattices trivial at the prescribed bad places, and integral harmonic
polynomials.  One denominator \(N_a\), valid for every section in that
theta family at the chosen weight, has
\(\log N_a\le C_{\rm int}z^2\kappa\); its raw pointwise estimates have
exponent \(C_{\rm raw}z^2\kappa\).  Multiplication by a reflex-ideal
scalar balances the lattice index at every embedding without changing
the \(K\)-span.  These are the substantive inputs to
\srcid{ExistsGoodCoverAtMinimalLevel}; the theorem name alone is not
a proof of them.
To pass from these even theta weights to a general dominant weight of
gap at least \(8fr+10\), one takes auxiliary even rank
\(8(fr+1)\), a key-unit twisting section, and a flat-line twisting
section with nonnegative finite valuations.  Their weights must add
exactly to the desired weight.  Their archimedean factors have
logarithms bounded by the key-unit term and
\[
 (|c_1|+q\log\operatorname{rd}(E))\sum_\tau|n_{0,\tau}|
 +\sum_{p\ {\rm bad}}\log p
 +\tfrac12\log\operatorname{rd}(K)+1,
 \qquad |c_1|\le9rL.
\]
Together with the proved bad-prime bound and the weight-size
inequalities, the source gives
\(C_2=11C_{\rm even}+C_{\rm bad}+500\) in
\eqref{eq:sr-span-metric-size}.
Those theta spanning, integral-denominator, key-unit, flat-line, and
Bergman constructions must all be established with these
normalizations to discharge the geometric input; the lattice and
constant arguments above do not replace them.

\section{A simultaneous finite-place theta regularizer}
\label{sec:sr-finite-regularizer}

This section constructs the regularizer used to pass from matched
quotients to ordinary convergent theta integrals.  The construction uses
only the split Schr\"odinger model at one finite place.  Its finite support,
nonzero degree, and compatibility with a finite set of relations are
proved explicitly.

Let \(k\) be a nonarchimedean local field with valuation ring \(\mathcal O\),
uniformizer \(\pi\), and residue field \(\mathbb F_Q\).  Give \(k\) the
absolute value \(|\pi|=Q^{-1}\).  Let \(d,r\) be nonnegative integers with
\(d/2\ge r+1\), and let \(\nu:k^\times\to\mathbb C^\times\) be unitary.
On functions on \(M_{r,d}(k)\) use the split Weil action
\begin{equation}\label{eq:sr-regularizer-action}
 (T_g\varphi)(Y)=\nu(\det g)|\det g|^{d/2}\varphi(g^{\mathsf t}Y),
 \qquad g\in\operatorname{GL}_r(k).
\end{equation}
A row frame \(Y\in M_{r,d}(\mathcal O)\) is primitive if its reduction
modulo \(\pi\) has row rank \(r\).  Equivalently its rows are independent
over \(k\) and their \(\mathcal O\)-span is saturated in
\(\mathcal O^d\).  Indeed, independence of the reductions implies that
every \(c\in k^r\) with \(cY\in\mathcal O^d\) belongs to \(\mathcal O^r\):
otherwise multiply \(c\) by a power of \(\pi\) to obtain a vector with at
least one unit coordinate whose product with \(Y\) is zero modulo \(\pi\).
The converse follows by lifting a nonzero relation modulo \(\pi\) and
dividing it by \(\pi\).

\begin{srproposition}[An explicit primitive-frame operator]
\label{prop:sr-primitive-regularizer}
There is a finitely supported sum
\[
 Z=\sum_{g\in\operatorname{GL}_r(k)}z_g T_g
\]
such that
\[
 Z{\bf1}_{M_{r,d}(\mathcal O)}
   ={\bf1}_{\{\text{primitive row frames}\}}.
\]
Writing \(x=\nu(\pi)Q^{-d/2}\) and \(\deg Z=\sum_gz_g\), one has
\begin{equation}\label{eq:sr-regularizer-degree}
 \deg Z=\prod_{i=0}^{r-1}(1-Q^i x),\qquad
 |\deg Z|\ge\frac12,\qquad
 \sum_g|z_g|\le e^{1/2}.
\end{equation}
In particular \(Z_0=(\deg Z)^{-1}Z\) has degree one and coefficient
\(\ell^1\)-norm less than \(4\).  With additive Haar measure on
\(M_{r,d}(k)\), its operator norm on \(L^2\) is less than \(4\).
\end{srproposition}

\begin{proof}
For every subspace \(L\le\mathbb F_Q^r\) of dimension \(j\), choose a
lift of a basis of \(L\), extended to a matrix \(A_L\in
\operatorname{SL}_r(\mathcal O)\).  For \(L=\mathbb F_Q^r\) take
\(A_L=1\).  Such an extension exists for proper \(L\): extend the basis
to an invertible matrix and adjust one of the complementary rows to
make its determinant one, then lift an element of
\(\operatorname{SL}_r(\mathbb F_Q)\) to
\(\operatorname{SL}_r(\mathcal O)\).
Put
\[
 \gamma_L=\operatorname{diag}(\underbrace{\pi^{-1},\ldots,\pi^{-1}}_j,
                         1,\ldots,1)A_L,\qquad g_L=\gamma_L^{\mathsf t}.
\]
Then \(\det g_L=\pi^{-j}\), and
\[
 \gamma_LY\in M_{r,d}(\mathcal O)
 \quad\Longleftrightarrow\quad
 Y\in M_{r,d}(\mathcal O)\ \hbox{and}\ L\overline Y=0.
\]
For the reverse implication to integrality, use
\(\gamma_L^{-1}=A_L^{-1}\operatorname{diag}(\pi,\ldots,\pi,1,\ldots,1)\),
whose entries are integral.

The M\"obius function from zero to a \(j\)-dimensional subspace in the
subspace lattice is
\(\mu(L)=(-1)^jQ^{j(j-1)/2}\).  This follows inductively from
\(\sum_{L\le U}\mu(L)=0\) for every nonzero \(U\); equivalently it is the
coefficient identity in the finite \(Q\)-binomial theorem
\[
 \sum_{L\le\mathbb F_Q^r}
    (-1)^{\dim L}Q^{\dim L(\dim L-1)/2}X^{\dim L}
 =\prod_{i=0}^{r-1}(1-Q^iX).
\]
The identity itself follows by induction on \(r\), splitting the
subspaces according to their intersection with a fixed hyperplane; a
\(j\)-dimensional subspace not contained in that hyperplane is specified
by its \((j-1)\)-dimensional intersection and \(Q^{r-j}\) possible
lifts.

Take
\[
 Z=\sum_{L\le\mathbb F_Q^r}\mu(L)x^{\dim L}T_{g_L}.
\]
The scalar in \(T_{g_L}\) is \(x^{-j}\); hence, at an integral \(Y\),
the value of \(Z{\bf1}_{M_{r,d}(\mathcal O)}\) is
\(\sum_{L\le\ker(\overline Y)}\mu(L)\).  This is one precisely when
\(\overline Y\) has independent rows and zero otherwise.  At a
nonintegral \(Y\) every summand is zero, by the preceding equivalence.

The same \(Q\)-binomial identity gives the degree in
\eqref{eq:sr-regularizer-degree}.  Set \(t_i=Q^{i-d/2}\).  Since
\(d/2\ge r+1\),
\[
 \sum_{i=0}^{r-1}t_i
 \le \sum_{j=2}^{\infty}Q^{-j}
 =\frac{1}{Q(Q-1)}\le\frac12.
\]
Consequently
\[
 |\deg Z|\ge\prod_{i=0}^{r-1}(1-t_i)
 \ge1-\sum_{i=0}^{r-1}t_i\ge\frac12.
\]
Taking absolute values of the coefficients before combining repeated
\(g_L\)'s, the \(Q\)-binomial identity with \(-|x|\) gives
\[
 \sum_g|z_g|
 \le\sum_LQ^{j(j-1)/2}|x|^j
 =\prod_{i=0}^{r-1}(1+t_i)
 \le\exp\left(\sum_it_i\right)\le e^{1/2}.
\]
Each \(T_g\) is an \(L^2\) isometry: the change of variables
\(Y\mapsto g^{\mathsf t}Y\) has modulus \(|\det g|^d\), cancelling the
square of the scalar in \eqref{eq:sr-regularizer-action}.  The operator
bound follows from the triangle inequality.
\end{proof}

The construction is unchanged by adapted integral coordinates.  More
precisely, let \(\Lambda\) be an \(\mathcal O\)-lattice in \(k^d\), let
\(P\in\operatorname{GL}_r(k)\), and take the local vacuum
\[
 \varphi_0(Y)=c\,{\bf1}_{M_r(\Lambda)}(PY).
\]
Use the preceding proof in a basis of \(\Lambda\) and replace every
\(g\) by \(P^{\mathsf t}g(P^{\mathsf t})^{-1}\).  Conjugation preserves
determinants and coefficient sums.  The resulting operator sends
\(\varphi_0\) to \(c\) times the primitive-frame indicator at \(PY\),
and satisfies the same degree and coefficient bounds.  Primitive
frames form a compact open subset of the full-rank locus.  For example,
compactness follows because they are a closed subset of
\(M_{r,d}(\mathcal O)\), and openness follows from their description by
the nonvanishing of an \(r\)-minor modulo \(\pi\).

Now let \(E/F\) be a CM extension and consider a small theta pair
\[
 \dim_E V'=r=a',\qquad \dim_E W=d=a+a',\qquad a'<a,
\]
with \(d\) even.  Thus \(a\ge a'+2\) and \(d/2\ge r+1\), exactly the
range of Proposition~\ref{prop:sr-primitive-regularizer}.  At a split
finite place \(v\) of \(F\), the fixed integral Schr\"odinger model
identifies the action of \(U(V')(F_v)\) with
\eqref{eq:sr-regularizer-action}, possibly after the adapted coordinate
change just described.  This identification includes its scalar
\(\nu(\det g)|\det g|^{d/2}\); replacing it by a scalar-free action would
give the wrong degree.

\begin{srproposition}[One place for a finite collection]
\label{prop:sr-simultaneous-regularizer}
Let \(\Psi_1,\ldots,\Psi_N\) belong to the algebraic restricted tensor
product of the fixed small theta model, and let
\(u_1,\ldots,u_N\in U(V')(\mathbb A_{F,f})\).
One can choose a single split finite place \(v\), one operator \(Z\)
supported in \(U(V')(F_v)\), and finite adelic elements \(u_i^{\,v}\)
whose \(v\)-components are one, such that
\[
 \deg Z\ne0,\qquad
 Z\omega(u_i)\Psi_i=\omega(u_i^{\,v})Z\Psi_i\quad(1\le i\le N).
\]
Every \(Z\Psi_i\) has a finite sum of pure tensor representations whose
\(v\)-factors are supported in the full-rank locus.  Assume the fixed arithmetic pair is the nondegenerate unitary Weil
pair of Section~\ref{sec:theta-decay-proof}, with its rational
Schr\"odinger coordinates, standard additive character, and the
unitary reduction-theoretic cover, height comparisons, and finite
quotient volume stated there.  Then every such vector has an ordinary
absolutely convergent theta integral of moderate growth, by
\eqref{eq:theta-full-decay} and the coefficient proof in
Subsection~\ref{subsec:sr-theta-coefficient-proof}.
\end{srproposition}

\begin{proof}
Write each \(\Psi_i\) as a finite sum of pure tensors.  Outside a finite
set of places all these tensors have the same prescribed vacuum
factor.  Enlarge this finite set so that the local characters and
integral models are unramified and each \(u_i\) belongs to the
hyperspecial group fixing the vacuum at every remaining place.
There are infinitely many rational primes splitting completely in a
Galois closure of \(E/\mathbb Q\); choose one outside the finite set and
take a place \(v\) of \(F\) above it.  This is the only use here of the
standard splitting-prime existence theorem; no size estimate for \(v\)
is required.

Use Proposition~\ref{prop:sr-primitive-regularizer} in the common vacuum
coordinates at \(v\), and transport \(Z\) into \(U(V')(F_v)\).  Its
degree stays nonzero.  The tensor factorization shows simultaneously
that all \(Z\Psi_i\) have primitive-frame factors there.  Define
\(u_i^{\,v}\) by deleting the \(v\)-component of \(u_i\).
The deleted component acts trivially on every original vacuum factor,
so \(\omega(u_i)\Psi_i=\omega(u_i^{\,v})\Psi_i\).
The operator \(Z\) commutes with \(\omega(u_i^{\,v})\), since they are
supported at disjoint places.  This proves the asserted identity.
Full-rank local support is precisely the hypothesis of
\eqref{eq:theta-full-decay}.  Under the fixed-pair and reduction
hypotheses in the statement, Subsection~\ref{subsec:sr-theta-coefficient-proof}
proves the coefficient majorant used by that estimate; the explicit
packing and minor argument in Section~\ref{sec:theta-decay-proof}
then proves convergence and moderate growth.
\end{proof}

The preceding proposition makes the quotient argument precise.
Let \(G=U(V')\), fix any Haar measure \(dg\) and a measurable fundamental
domain \(\mathcal F\) for \(G(F)\backslash G(\mathbb A_F)\), and write
\[
 I(\Psi)(h)=\int_{\mathcal F}\Theta_\Psi(g,h)\,dg
\]
whenever this integral is absolutely convergent.  The choice of Haar
measure is fixed throughout.  The theta transformation law and
rational left invariance imply
\begin{equation}\label{eq:sr-ordinary-translate}
 I(\omega(u)\Psi)=I(\Psi)
\end{equation}
whenever \(I(\Psi)\) exists.  Indeed
\(\Theta_{\omega(u)\Psi}(g,h)=\Theta_\Psi(gu,h)\);
right translation preserves quotient measure because a unitary group
is unimodular.  The two fundamental domains give the same integral
by rational left invariance.  In particular, finite sums satisfy
\begin{equation}\label{eq:sr-ordinary-degree}
 I(Z\Psi)=(\deg Z)I(\Psi)
\end{equation}
for every integrable \(\Psi\), regardless of the support of \(Z\).

\begin{srcorollary}[Integrable finite relations vanish]
\label{cor:sr-integrable-relations}
Suppose
\[
 \Psi=\sum_{i=1}^{N}b_i\bigl(\omega(u_i)\Psi_i-\Psi_i\bigr)
\]
belongs to the algebraic model, and its theta integral is absolutely
convergent.  Then \(I(\Psi)=0\).
\end{srcorollary}

\begin{proof}
Apply Proposition~\ref{prop:sr-simultaneous-regularizer} to the entire
finite collection.  Each \(Z\Psi_i\) is integrable and
\[
 Z\Psi=\sum_i b_i
 \bigl(\omega(u_i^{\,v})Z\Psi_i-Z\Psi_i\bigr).
\]
Every difference on the right has zero integral by
\eqref{eq:sr-ordinary-translate}; finite additivity is justified by
absolute convergence.  Thus \(I(Z\Psi)=0\).
On the other hand the assumed integrability of \(\Psi\) gives
\eqref{eq:sr-ordinary-degree}.  Divide by the nonzero degree.
\end{proof}

Consequently, if two small vectors have the same image in the finite
coinvariant quotient, both have convergent theta integrals, and their
difference is a finite relation of the displayed form, their ordinary
integrals agree.  If a representative \(\Psi_0\) in that class has
already been matched to a big vector \(\Phi\) with
\(I_V(\Phi)=cI_{V'}(\Psi_0)\), the same equality holds for every
integrable representative \(\Psi\) in the class.  There is no additional
volume scalar in this passage: the argument uses one Haar measure
throughout, and scaling it scales every small integral by the same
amount.

This constructs the finite-place regularizer and proves the passage
from a coinvariant class to any convergent representative.  The scalar
\(c\), its nonvanishing, the placewise matching maps, and the universal
character-independent rank-recursion signs still require their
separate theta-comparison proofs; they are not consequences of the
regularizer.  In the frozen source the construction here occurs in
\srcid{ExistsTranslatesEqPrimFrameLat},
\srcid{ExistsGoodPlaceTranslateCombinationLatticeIndica},
\srcid{ExistsCombinationMemRegularSubspace},
\srcid{TwistRegulariserIsGood}, and
\srcid{ThetaIntegralsAgreeOfIntegrableKernel}.

\subsection{The coefficient sum and its Schwartz majorant}
\label{subsec:sr-theta-coefficient-proof}

We give the coefficient calculation used in
\eqref{eq:theta-lattice-majorant}.
Let \(A=\mathbb A_F\), choose the standard nontrivial additive
character \(\psi:A/F\to\mathbb C^\times\), and give the additive
spaces self-dual Haar measure.  Use a rational polarization of
the symplectic space of the specified Weil pair, so its
Schr\"odinger space is \(\mathcal S(A^N)\) and its rational
Heisenberg coordinates are \(F^N\oplus F^N\).  In these coordinates
\[
 [S(a,b)f](u)=
 \psi\bigl(b\cdot u+\tfrac12 a\cdot b\bigr)f(u+a),
 \qquad Pf(u)=f(-u).
\]
The rational polarization identifies these coordinates with the
rational matrix space used in Section~\ref{sec:theta-decay-proof};
the action of \((g,h)\) on that space is
\(x\mapsto gx h^{\mathsf t}\), with the fixed changes of bases
included in the model.

For \(\Phi\in\mathcal S(A^N)\) set
\[
 c_\Phi(a,b)=
 \int_{A^N}\psi\bigl(b\cdot u+\tfrac12a\cdot b\bigr)
          \Phi(u+a)\overline{\Phi(-u)}\,du
 =\langle S(a,b)\Phi,P\Phi\rangle .
\]
Then
\begin{equation}\label{eq:sr-theta-square-coefficients}
 \left|\sum_{x\in F^N}\Phi(x)\right|^2
       =\sum_{a,b\in F^N}c_\Phi(a,b).
\end{equation}
All series in this identity converge absolutely.

Here are the analytic details justifying that assertion and its
uniform use with the infinite Schwartz variable.
At the archimedean places, \(c_f\) is a Schwartz function of \((a,b)\).
For every desired weight exponent its Schwartz seminorm is at most
a constant times \(p(f)^2\), for a continuous Schwartz seminorm \(p\).
To verify this, differentiate its integral in \(a,b\).
A derivative in \(b\) inserts a polynomial in \(u+a/2\);
a derivative in \(a\) differentiates \(f(u+a)\) and the phase.
To bound a monomial in \(b\), move derivatives of the phase onto
the two Schwartz factors by integration by parts.
Write each occurrence of \(a\) as \((u+a)-u\) and expand.
Every resulting integrand is a product of weighted derivatives of
\(f(u+a)\) and \(f(-u)\).
Bound the weighted derivatives by their Schwartz seminorms times
\((1+|u+a|)^{-m}(1+|u|)^{-m}\), with \(m\) larger than the real
dimension.  Its integral is bounded independently of \(a\).
This bounds every polynomially weighted derivative of \(c_f\).
Only finitely many derivatives and weights are needed for a fixed
bound, so their maximum, enlarged by a fixed constant, is the
continuous seminorm \(p\).
The same proof applies to the product of the real and complex
archimedean coordinates after the fixed linear change making the
character the standard Fourier character.

At a finite place, let \(\phi\) be locally constant with compact
support \(K\), and choose a compact open additive subgroup \(U\)
under which \(\phi\) is invariant.
Nonvanishing of the integrand implies \(a\in K+K\).
Translating the integration variable by \(t\in U\) multiplies
the coefficient by \(\psi(b\cdot t)\); hence it is zero unless
\(b\in U^\perp\).  The annihilator \(U^\perp\) is compact open.
The coefficient is locally constant, and thus is again
Schwartz--Bruhat, with compact support in
\((K+K)\times U^\perp\).
For a fixed finite adelic factor \(\phi_f\) this proves a bound
\[
 |c_{f\otimes\phi_f}(w)|
 \le p_{\phi_f,L}(f)^2\,
       {\bf1}_{C_f}(w_f)(1+|w|_\infty)^{-L}
 \qquad(w\in A^{2N})
\]
for every \(L\), some fixed compact \(C_f\), and a continuous
seminorm \(p_{\phi_f,L}\).
Restricted tensor products are handled by multiplication of the
local formulas; their vacuum factors agree outside finitely many
places.  Finite linear combinations follow by polarization and the
triangle inequality.

It remains to check \eqref{eq:sr-theta-square-coefficients}.
For fixed rational \(a\), apply additive adelic Poisson summation
in \(u\) to the Schwartz function
\(\Phi(u+a)\overline{\Phi(-u)}\).
Because \(a,b\) are rational, \(\psi(a\cdot b/2)=1\).
Self-dual measure for the standard character has
\(\operatorname{vol}(A^N/F^N)=1\), so there is no Poisson volume
factor, and the sum over rational \(b\) of its Fourier transform is
\[
 \sum_{u\in F^N}\Phi(u+a)\overline{\Phi(-u)}.
\]
Summing over rational \(a\), the bijection
\((a,u)\mapsto(u+a,-u)\) of \(F^N\times F^N\) gives the product
of the two theta sums in
\eqref{eq:sr-theta-square-coefficients}.
The coefficient sum is absolutely convergent by its proved Schwartz
bound.  The double sum of the products is absolutely convergent
because each theta sum of the absolute value of a Schwartz function
converges.  These facts justify both Poisson summation and the
rearrangement.

The unitary Weil action of the pair preserves the Schwartz space,
has
\(\omega(g,h)S(w)\omega(g,h)^{-1}=S((g,h)w)\), and commutes with
\(P\).  These assertions follow directly on the Schr\"odinger
generators: a linear change has the square-root modulus making it
unitary, multiplication by a quadratic character has modulus one
and is even, and the self-dual Fourier transform is unitary and
commutes with reflection.  Thus
\[
 c_{\omega(g,h)\Phi}(w)=c_\Phi((g,h)^{-1}w).
\]
Apply \eqref{eq:sr-theta-square-coefficients} to
\(\omega(g,h)\Phi\), take absolute values, and apply the preceding
Schwartz bound.  In the matrix coordinates of the fixed pair,
this is precisely
\[
 |\Theta_{f\otimes\phi_f}(g,h)|^2
 \le p_{\phi_f,L}(f)^2
     \sum_{x}
       {\bf1}_{C_f}\bigl((g^{-1}x(h^{-1})^{\mathsf t})_f\bigr)
       (1+|g^{-1}x(h^{-1})^{\mathsf t}|_\infty)^{-L}.
\]
For full-rank local data, apply the explicit compact-coset
Schr\"odinger separation in
Section~\ref{sec:theta-decay-proof} to the local factor before
this calculation.  On each piece the coefficient vanishes at
every rank-deficient local matrix.  Rational rank-deficient
matrices remain rank deficient under the two invertible group
matrices.  The sum can therefore be restricted to full row rank.
For a finite sum of pieces, take square roots of the separate
bounds, use the union of their finite compact supports, and add
their seminorms with the absolute values of their coefficients.
This proves \eqref{eq:theta-lattice-majorant} in both of its forms.

Consequently the lattice packing, full-rank minor estimate, and
two-chamber calculation in
Section~\ref{sec:theta-decay-proof} prove
\eqref{eq:theta-full-decay} for this fixed arithmetic pair.
The remaining standard reduction-theoretic assumptions are exactly
the unitary Siegel-set cover, its height comparisons, and finite
Haar volume of \(U_X(F)\backslash U_X(\mathbb A_F)\), as stated
there.  They apply to the nondegenerate Hermitian space of the
fixed pair; its unitary group is reductive and its center is the
norm-one torus of the CM extension, which is anisotropic over
\(F\).  At \(N_0=0\), \eqref{eq:theta-full-decay} is uniform in
\(g\); integration over this finite-volume quotient gives ordinary
absolute convergence.  Its polynomial bound in \(h\) gives moderate
growth of the resulting integral.  None of these assertions
constructs the arithmetic pair, nor asserts uniformity of its
Schwartz seminorm constants across varying fields.

\section{The lattice Gram operator in the theta spanning construction}
\label{sec:sr-theta-gram}

We isolate the quantitative part of the theta spanning argument and
its precise global kernel input.  The strict diagonal dominance below
is a consequence of a lattice count and finite Heisenberg
orthogonality; it is uniform in the field, the level, and the genus.

Let \(V_\mathbb Q\) have dimension \(\rho\) and a nondegenerate
alternating rational form \(B\).  For a full lattice \(L\), put
\(L^\#=\{v:B(v,L)\subset\mathbb Z\}\), and call \(L\) selfdual if
\(L^\#=L\).  For two selfdual lattices set
\[
 m(L,L')=[L:L\cap L'].
\]
The same index is obtained with \(L,L'\) interchanged, since
selfdual lattices have the same covolume.

\begin{srlemma}[Counting selfdual lattices]
\label{lem:sr-selfdual-count}
Fix a selfdual \(L\).  For every \(m\ge1\), there are at most
\(m^{2\rho}\) selfdual lattices \(L'\) with \(m(L,L')=m\).
The index equals one precisely when \(L'=L\).
\end{srlemma}

\begin{proof}
A free abelian group of rank \(\rho\) has at most \(m^\rho\)
subgroups of index \(m\).  For clarity, the Hermite normal form
count \(a_\rho(m)\) satisfies
\[
 a_\rho(m)=\sum_{d\mid m}d^{\rho-1}a_{\rho-1}(m/d).
\]
The formula follows by fixing the last diagonal pivot in triangular
normal form and its residue entries.  Induction and the bound
\(\#\{d:d\mid m\}\le m\) give
\(a_\rho(m)\le m^{\rho-1}\#\{d:d\mid m\}\le m^\rho\).

For \(L'\) as in the statement let \(J=L\cap L'\).
There are at most \(m^\rho\) possibilities for \(J\).
Both \(L\) and \(L'\) lie in \(J^\#\).
Duality reverses inclusion and preserves the index of an inclusion,
as follows by taking a basis and its \(B\)-dual basis.
Therefore
\([J^\#:L]=[L:J]=m\) and
\([J^\#:L']=[L':J]=m\).
For a fixed \(J\), the lattice \(L'\) is among the subgroups of index
\(m\) of the free group \(J^\#\), of which there are at most
\(m^\rho\).  Multiplication proves the bound.
If \(m=1\), then \(L\subset L'\); duality gives \(L'\subset L\).
\end{proof}

Here is the vacuum calculation producing the exponent \(n/2\).
In the finite Schr\"odinger representation associated with
\(J^\#/J\), the two complementary maximal isotropic groups
\(L/J\) and \(L'/J\) each have order \(m=m(L,L')\).
Their pairing is perfect: an element orthogonal to \(L'/J\)
lifts to \((L')^\#=L'\), and its intersection with \(L/J\) is zero.
In the realization by functions on \(L'/J\), one group acts by
translations and the other by its \(m\) distinct characters.
A unit eigenvector for the character group is a delta function;
a unit eigenvector for the translation group has all coordinates
of modulus \(m^{-1/2}\).
Consequently normalized vacua \(v_L,v_{L'}\), with the compatible
semicharacters, satisfy
\[
 |\langle v_L,v_{L'}\rangle|=m^{-1/2}.
\]
If their semicharacters disagree on \(J\), the inner product is
zero instead: acting by an element of \(J\) gives two different
eigenvalues.  Thus the inequality with \(m^{-1/2}\) always holds.
For \(n\) tensor copies it gives \(m^{-n/2}\).
This argument includes the semicharacter phases; omitting them
would fail at the even part of the Heisenberg extension.

In the level construction a test vector is a sum of \(b=|J_{\rm lev}|\)
translated normalized vacua.  Its labels run through one prescribed
finite level orbit, with unit-modulus character coefficients.
A pairing between two different labels vanishes: a common
Heisenberg eigenoperator distinguishes their residue classes.
For matching labels the same translation acts on both vacua and
preserves their inner product.  Therefore the test vectors
\(\Phi_{L,\Lambda}\), for every positive definite auxiliary
unimodular lattice \(\Lambda\), obey
\begin{equation}\label{eq:sr-test-gram}
 \langle\Phi_{L,\Lambda},\Phi_{L,\Lambda}\rangle=b,\qquad
 |\langle\Phi_{L,\Lambda},\Phi_{L',\Lambda}\rangle|
       \le b\,m(L,L')^{-n/2}.
\end{equation}
The equality of label sets and the distinguishing eigenoperator
are necessary compatibility conditions on the finite level
construction.  In the prescribed construction they are obtained
by the same residue representatives in both test vectors; the
coprimality of the two CM level ideals identifies equality of
their character pairings with equality of their residue classes.
This is the finite pairing calculation, rather than a choice of
independent normalizations for each \(L,\Lambda\).

\begin{srproposition}[A uniform diagonal gap]
\label{prop:sr-gram-gap}
Let \(\mathcal L\) be any set of selfdual lattices in \(V_\mathbb Q\),
let \(\mathcal C\) be a finite nonempty set of auxiliary lattice
classes with positive mass weights \(w_\Lambda\), and suppose
\eqref{eq:sr-test-gram} holds for one common \(b>0\).
If \(n\ge4\rho+4\), define
\[
 C(L,L')=\sum_{\Lambda\in\mathcal C}
       w_\Lambda\langle\Phi_{L,\Lambda},\Phi_{L',\Lambda}\rangle,
 \qquad a=b\sum_\Lambda w_\Lambda.
\]
Then \(a>0\), \(C(L,L)=a\), and
\begin{equation}\label{eq:sr-gram-row-gap}
 \sum_{\substack{L'\in\mathcal L\\L'\ne L}}|C(L,L')|
       \le0.65a.
\end{equation}
\end{srproposition}

\begin{proof}
The diagonal assertion follows from \eqref{eq:sr-test-gram}.
For the remaining sum, group the lattices by \(m(L,L')\).
Lemma~\ref{lem:sr-selfdual-count} gives
\[
 \sum_{L'\ne L}m(L,L')^{-n/2}
 \le\sum_{m=2}^\infty m^{2\rho-n/2}
 \le\sum_{m=2}^\infty m^{-2}<0.65.
\]
The last numerical inequality needs no special value of the zeta
function: the sum from \(2\) through \(10\), plus
\(\int_{10}^\infty t^{-2}\,dt\), is less than \(13/20\).
Apply the triangle inequality to the finite positive mass average.
These estimates also prove absolute convergence of every row.
\end{proof}

\begin{srproposition}[The actual inverse operator]
\label{prop:sr-gram-neumann}
Let \(\mathcal B\) be the Banach space of bounded families
\(F(L,x)\) with the fixed holomorphic weight on the common frame
domain.  Equip it with the supremum of the invariant pointwise
norms.  Let \(C\) satisfy
\eqref{eq:sr-gram-row-gap} and act in the lattice variable:
\[
 (\mathcal CF)(L,x)=\sum_{L'}C(L,L')F(L',x).
\]
Then \(\mathcal C\) is invertible and
\begin{equation}\label{eq:sr-gram-inverse}
 F'=\mathcal C^{-1}F
 =a^{-1}\sum_{k=0}^\infty(-R)^kF,\qquad
 R=a^{-1}\mathcal C-1,\qquad
 \|F'\|\le\frac{\|F\|}{0.35a}.
\end{equation}
The inverse preserves holomorphy, weight, and the simultaneous
automorphy law whenever the original test vectors have that law.
\end{srproposition}

\begin{proof}
The absolutely convergent row estimate gives \(\|R\|\le0.65\).
The geometric series converges in operator norm and multiplication
by \(1+R\) telescopes to the identity, giving
\eqref{eq:sr-gram-inverse}.
Normal convergence on compact subsets preserves holomorphy;
the weight relation is linear and passes through each row sum
and the uniform limit.

More explicitly, let a group \(\Delta\) act on the lattices and
frames and let \(\widetilde\mu(\gamma,L)\) be the common unitary multiplier.
The test-vector identity
\(\omega(\gamma)\Phi_{L,\Lambda}
 =\widetilde\mu(\gamma,L)\Phi_{\gamma L,\Lambda}\)
gives
\[
 C(\gamma L,\gamma L')\widetilde\mu(\gamma,L)
        \overline{\widetilde\mu(\gamma,L')}=C(L,L').
\]
If
\(F(\gamma L,\gamma x)=\widetilde\mu(\gamma,L)^{-1}F(L,x)\),
reindexing the absolutely convergent row shows that
\(\mathcal CF\) has the same law.  Thus every partial sum and
its limit in \eqref{eq:sr-gram-inverse} has that law.
\end{proof}

We state the exact analytic input which turns this inverse into
theta spanning.  This identifies the required version of the
doubled Siegel--Weil identity and the frame slab, without replacing
either by a theorem name.
Choose an orthonormal basis \(P_1,\ldots,P_t\) of the nonzero
harmonic highest-weight space with its Fock inner product.
Columnwise unitary action preserves that inner product, so
\[
 \sum_jP_j(uz)\overline{P_j(uw)}
       =\sum_jP_j(z)\overline{P_j(w)}.
\]
Let \(K(x,y)\) be the result of applying the associated paired
polynomial derivatives to the doubled Gaussian kernel.
Let
\(T_{\Lambda,j}(L,x)=\Theta_{\Phi_{L,\Lambda},P_j}(x)\).
For the common finite data the required kernel identity is
\begin{equation}\label{eq:sr-doubled-kernel-input}
 \kappa\sum_{\gamma\in\Delta}
    C(L,\gamma L')\overline{\widetilde\mu(\gamma,L')}K(x,\gamma y)
  =\sum_{\Lambda\in\mathcal C}w_\Lambda
       \sum_{j=1}^t
       T_{\Lambda,j}(L,x)
       \overline{T_{\Lambda,j}(L',y)},\qquad \kappa\ne0.
\end{equation}
Both the order of \(x,\gamma y\) and the conjugate multiplier are
part of the normalization.

Here the slab is the open set of frames
\[
 S_\delta=\{y:\ |\operatorname{Gram}_\tau(y)_{ab}
                  -\operatorname{diag}(1_p,-1_q)_{ab}|<\delta
                       \text{ for all }\tau,a,b\}.
\]
For sufficiently small fixed \(\delta>0\), the additional
archimedean assertions needed are
\begin{equation}\label{eq:sr-slab-reproduction-input}
 \int_{S_\delta}K(x,y)H(y)\,dy=\rho_0H(x),
 \qquad \rho_0>0,
\end{equation}
for every bounded holomorphic section \(H\) of the fixed weight,
and absolute integrability of \(K(x,\cdot)\) on this slab.
The invariant pointwise norm controls the ordinary absolute value
on the slab by one constant for the fixed weight.  In particular
\[
 \sum_{L'}\int_{S_\delta}
       |C(L,L')K(x,y)F'(L',y)|\,dy<\infty
\]
follows from the row bound and that slab bound.

\begin{srproposition}[Finite theta spanning from the normalized kernel]
\label{prop:sr-gram-theta-span}
Suppose the fixed level data satisfy
\eqref{eq:sr-test-gram},
\eqref{eq:sr-doubled-kernel-input} and
\eqref{eq:sr-slab-reproduction-input}, with the indicated absolute
convergence.  Suppose diagonal action of \(\Delta\) on
\(\mathcal L\times S_\delta\) has a measurable fundamental domain,
up to null sets, and preserves product counting and frame
Lebesgue measure.  Then every bounded holomorphic automorphic
family \(F\) of this weight is a finite linear combination, with
coefficients independent of \(L,x\), of the functions
\(T_{\Lambda,j}(L,x)\).
\end{srproposition}

\begin{proof}
Take \(F'=\mathcal C^{-1}F\).
Reproduction and absolute convergence give
\[
 \sum_{L'}\int_{S_\delta}
       C(L,L')K(x,y)F'(L',y)\,dy=\rho_0F(L,x).
\]
Unfold this integral over the measurable fundamental domain for
the diagonal action.
The automorphy law for \(F'\) contributes
\(\overline{\widetilde\mu(\gamma,L')}\), exactly the factor in
\eqref{eq:sr-doubled-kernel-input}.
Using that identity, the result is an integral of a finite sum
\[
 \kappa^{-1}\sum_{\Lambda,j}
   w_\Lambda\,T_{\Lambda,j}(L,x)
       \overline{T_{\Lambda,j}(L',y)}F'(L',y)
\]
over that fundamental domain.  Its value therefore lies in the
finite-dimensional span of the functions
\(T_{\Lambda,j}(L,x)\).

One need not assume in this last step that every displayed
coefficient integral is individually absolutely convergent.
Here is the requisite finite-dimensional argument.
If an integrable function of \(z\), for every parameter \(a\), is
of the form \(\sum_{j=1}^N e_j(a)t_j(z)\), choose a basis
\(e_{j_1},\ldots,e_{j_s}\) of the finite span of the \(e_j\)'s and
rewrite the sum in that basis.  There are \(s\) sample parameters
\(a_1,\ldots,a_s\) for which the evaluation matrix of this basis
is invertible.  Otherwise all evaluation functionals would span a
proper subspace of the dual and some nonzero basis combination
would vanish at every parameter.
The rewritten coefficient functions are consequently fixed
linear combinations of the \(s\) integrable functions at those
sample parameters.  They are integrable, and integration gives
a fixed linear combination of the basis \(e_{j_i}\)'s.
Apply this to \(a=(L,x)\) and the unfolded integrand.
Since \(\rho_0\kappa\ne0\), division proves the result.
\end{proof}

In applying this proposition, the normalized definite
Siegel--Weil formula of
Section~\ref{sec:sw-boundary} is used in the range \(n>4r\).
Its auxiliary Hermitian space is positive definite, so the
unitary theta quotient is compact at infinity and the genus mass
weights are positive.
The comparison must use the same auxiliary lattice class,
finite test function, doubled rational frame, and self-dual
Fourier measure as in \(T_{\Lambda,j}\).
In the Eisenstein expansion, maximal isotropic subspaces that
are graphs of unitary elements contribute the left side of
\eqref{eq:sr-doubled-kernel-input}.
The level support removes the graphs outside \(\Delta\).
The other maximal isotropic subspaces contain a nonzero rational
isotropic direction; their contributions vanish when the specified
finite test vectors have a nonzero-degree Heisenberg relation in
that direction.
Summing paired derivatives over the orthonormal \(P_j\)'s commutes
with the compact auxiliary unitary average by the displayed
Fock invariance, and the definite theta side becomes the finite
right side of \eqref{eq:sr-doubled-kernel-input}.
This explains both the Gram coefficient and the multiplier and
shows why independent local normalizations would be insufficient.

Sections~\ref{sec:sr-level-phase-proof} and~\ref{sec:sr-frame-slab-kernel}
supply the common finite test-vector norms and phases, the outside-level
support, the nonzero-degree isotropic relations, and the explicit
archimedean slab integrability and reproduction.
The multiplier in this section is their actual test-vector multiplier
\(\widetilde\mu\), which includes the fixed level character.
Subsection~\ref{subsec:sr-invariant-heat-jets} supplies the invariant
jet implication from the scalar identity.
The graph-term computation, passage of the isotropic annihilator
through the nongraph terms, the admissibility of the doubled test
family and identification of its Gaussian kernel remain to be
verified in this application.  In particular
the abstract proposition is not asserted to construct those inputs.
In the frozen source their producers are respectively
\srcid{EisensteinSeriesTermUnitaryGraph},
\srcid{ThetaTestFunctionNoCoinvariantsIsotropic},
\srcid{AdmissibleDoubledHermitianThetaFamily},
\srcid{GaussianKernelIntegrableOnThinSlab}, and
\srcid{SetIntegralCompOfStable}.
The finite lattice count, numerical gap, inverse operator, and
finite-dimensional unfolding implication have been proved above.

\section{Common level labels and their Heisenberg phases}
\label{sec:sr-level-phase-proof}

We verify the finite test-vector norm and overlap in
\eqref{eq:sr-test-gram} using the actual shifted level lattices.
The same calculation supplies one multiplier for all auxiliary
positive definite lattices.

Let \(L_0\) be the fixed polarized \(\mathcal O_E\)-lattice of
rank \(r\), selfdual for its rational alternating trace form
\(B\).  Thus \(B(av,z)=B(v,\overline a z)\).
Let \(\mathfrak a\) be the finite level ideal, with
\[
 (\mathfrak a,\overline{\mathfrak a})=1,\qquad
 (\mathfrak a,2)=1,\qquad M\in\mathfrak a,\quad M\ne0.
\]
Choose \(w\in\mathcal O_E\) with \(w+\overline w=1\).
All Heisenberg operators and all vacua below use this one \(w\).
Put \(\mathfrak j=\mathfrak a\overline{\mathfrak a}^{-1}\).
The lattices \(L\) under consideration are selfdual and agree
with \(L_0\) after localization at every prime dividing
\(\mathfrak a\overline{\mathfrak a}\).
This is equivalent to having a common denominator, prime to
those ideals, carrying each lattice into the other.

The auxiliary lattices all lie in one fixed Hermitian space,
and this is essential for the projections below.  More explicitly,
fix
\[
 V_+=E^n,\qquad
 h_+(v,v')=\sum_{i=1}^n v_i\overline{v'_i},\qquad
 \Lambda_0=\mathcal O_E^n.
\]
Assume \(n\) is even.  Every \(\Lambda\) below is an
\(\mathcal O_E\)-submodule of this \(V_+\), selfdual for this
one \(h_+\), and standard at the primes of \(M\).
The form is positive definite at each complex embedding.
The lattice \(\Lambda_0\) is selfdual: pairing with the coordinate
vectors shows that its Hermitian dual has all coordinates in
\(\mathcal O_E\), and the reverse inclusion follows by summing
products of integral coordinates.  No individual Hermitian form
is assigned to \(\Lambda\).

Put \(V=E^r\), with its fixed polarized trace form \(B\).
There is one rational symplectic space
\[
 \mathbb W=V_+\otimes_E V
        \simeq \operatorname{Mat}_{n,r}(E),\qquad
 B_n(Y,Z)=\sum_{i=1}^n B(Y_i,Z_i).
\]
On pure tensors its pairing is
\(B_n(v\otimes u,v'\otimes u')
  =B(h_+(v,v')u,u')\), by the adjointness of multiplication
by \(E\).  Thus the perfect tensor pairings used below are
restrictions of this same \(B_n\).
Choose a rational polarization of \(\mathbb W\) once, for example
by choosing a symplectic integral basis of \(L_0\), and take
\[
 {\cal H}=L^2(\mathbb A_f^{ng}),\qquad
 {\cal S}={\cal S}(\mathbb A_f^{ng}),\qquad
                   2g=\dim_{\mathbb Q}V.
\]
Use the standard rational additive character and its selfdual
Haar measure.  In the rational-function notation of the source,
\({\cal S}\) is the space of rational Schwartz--Bruhat functions
on \(\operatorname{Mat}_{n,g}(\mathbb Q)\): restriction gives the
identification, because \(\mathbb Q^{ng}\) is dense in
\(\mathbb A_f^{ng}\).  A common support denominator and an open
period subgroup make the inverse extension well defined; its
finite residue average is the same \(L^2\) inner product.
All vacua, projections and translated vectors in this section
are vectors in this one \({\cal H}\).

The group \(U(V,B)\) acts on \(\mathbb W\) on each of the \(n\)
rows, preserving \(B_n\).  We use one genuine unitary Weil
action \(W\) of this group on \({\cal S}\), with the covariance
convention
\[
 W_\gamma S(Y)W_\gamma^{-1}=S(Y\gamma^t).
\]
Here genuine means \(W_{\gamma'\gamma}=W_{\gamma'}W_\gamma\).
This action is fixed before the lattices are varied.
The standard metaplectic Schr\"odinger representation supplies
it: the action on \(n\) rows is the \(n\)-fold tensor power of
the action for \(V\); the two lifts of a symplectic map differ
by \(-1\), acting by \((-1)^n=1\) on that tensor power.
It therefore descends to the symplectic group when \(n\) is
even.  Fourier transform with selfdual measure, quadratic
characters of modulus one and normalized linear changes of
variables give its unitary operators and preserve \({\cal S}\).
This is the standard metaplectic background, with the additive
character, measure and parity now specified; see Weil,
\emph{Sur certains groupes d'op\'erateurs unitaires},
Acta Math.\ 111 (1964), 143--211,
\url{https://doi.org/10.1007/BF02391012}.

The theta-normalized action in the source is characterized by
\[
 \operatorname{thA}(W_\gamma\phi;\gamma x,z)
                         =\operatorname{thA}(\phi;x,z).
\]
The usual Poisson theta transformation fixes this scalar
normalization.  Its uniqueness, and hence its genuine group law,
can also be checked directly.  At one frame \(x\), the invertible
period-coordinate change and the nonvanishing Gaussian factor
reduce equality of these entire functions of \(z\) to equality of
\[
 \sum_{a\in\mathbb Q^{ng}}
       \phi(a)\exp\{\pi i\,\operatorname{tr}(a\tau a^t)
                              +2\pi i\,\operatorname{tr}(az^t)\},
 \qquad \operatorname{Im}\tau>0 .
\]
Put the two test functions over one support denominator.
Absolute Gaussian convergence permits Fourier coefficient
extraction on a real period torus; the coefficient at \(a\) is
its displayed nonzero Gaussian times \(\phi(a)\).
Thus equality gives equality of the test functions.
Applying the theta transformation twice proves the group law
by this uniqueness, without choosing phases separately for
different lattices.  In the frozen archive the precise producer
is \srcid{WeilActionPreservesThetaAndInnerProduct.wG\_laws};
its inputs are the fixed symplectic basis and evenness of \(n\),
and its conclusions include unitarity, linearity, the group law
and the above Heisenberg covariance.

With this fixed ambient model, set
\[
 A_{\Lambda,L}=\Lambda\otimes_{\mathcal O_E}\mathfrak jL,\quad
 H_{\Lambda,L}=\Lambda\otimes_{\mathcal O_E}
                         2\overline{\mathfrak a}L,\quad
 K_{\Lambda,L}=\Lambda\otimes_{\mathcal O_E}
                         2\overline{\mathfrak a}\mathfrak aL.
\]
The first lattice is selfdual.
Indeed tensoring the perfect Hermitian pairing on \(\Lambda\)
with the perfect alternating pairing on \(L\) is perfect;
and the twist by \(\mathfrak j\) preserves selfduality because
\(\mathfrak j\overline{\mathfrak j}=1\).
These statements can be checked over every local DVR and then
intersected inside the rational space.

Fix the Schr\"odinger operators with convention
\[
 S(u)S(v)=e^{-\pi iB_n(u,v)}S(u+v),\qquad
 S(u)S(v)=e^{-2\pi iB_n(u,v)}S(v)S(u).
\]
For a rational vector \(\delta\), set
\begin{equation}\label{eq:sr-common-eigenoperator}
 \epsilon_w(\delta)=e^{\pi iB_n(\delta,w\delta)},\qquad
 U_\delta=\epsilon_w(\delta)S(\delta).
\end{equation}
On any integral \(\mathcal O_E\)-stable lattice \(A\), these
operators form an additive representation.
In fact the quotient of \(U_uU_v\) by \(U_{u+v}\) is
\[
 \exp\{-\pi i[
 B_n(u,wv)+B_n(v,wu)+B_n(u,v)]\}
       =\exp\{-2\pi iB_n(u,wv)\}=1.
\]
Here adjointness and \(w+\overline w=1\) give the equality,
and integrality gives the final value.

There is a one-dimensional vacuum line fixed by all
\(U_\delta\), \(\delta\in A\), for every selfdual \(A\).
For completeness, choose a symplectic integral basis of one
selfdual lattice and its Schr\"odinger realization.
The standard compact-open indicator, translated and multiplied
by half-integral characters to account for the signs
\(\epsilon_w\) on this basis, gives its vacuum.
For a second commensurable selfdual lattice, reduce to the
finite Heisenberg quotient of their intersection.
In the model of functions on a finite maximal isotropic group,
translations and the distinct dual characters have each joint
eigenline of dimension one.
They give the corresponding vacuum and show uniqueness.
The construction uses finite sums of rational Heisenberg
translations, so the vacuum remains Schwartz--Bruhat.

The phases can be chosen compatibly with the prescribed genuine
unitary Weil action.  First choose unit vacua \(v_{0,L}\) for
\(A_{0,L}=\mathcal O_E^n\otimes\mathfrak jL\), once for each \(L\).
For a level-group element \(\gamma\), write
\[
 W_\gamma v_{0,L}=\mu(\gamma,L)v_{0,\gamma L}.
\]
The lines have just been proved one-dimensional, so the scalar
exists, has modulus one, and satisfies the cocycle relation
forced by the same genuine action \(W\).
For an arbitrary auxiliary \(\Lambda\), let \(P_{\Lambda,L}\)
be the orthogonal projector onto the
\(A_{\Lambda,L}\)-vacuum line, obtained by averaging
\(U_\delta\) over its finite orbit on \(v_{0,L}\), and set
\[
 v_{\Lambda,L}=
  \frac{P_{\Lambda,L}v_{0,L}}
       {\|P_{\Lambda,L}v_{0,L}\|}.
\]
The projection is nonzero: the finite Heisenberg calculation
of Lemma~\ref{lem:sr-selfdual-count}'s subsequent vacuum argument
gives overlap of modulus the inverse square root of the
relative index, since the semicharacters agree on the
intersection.
Weil covariance intertwines the two projectors.
Unitarity preserves their positive normalization denominators,
and hence
\begin{equation}\label{eq:sr-common-vacuum-multiplier}
 W_\gamma v_{\Lambda,L}
       =\mu(\gamma,L)v_{\Lambda,\gamma L}.
\end{equation}
In particular the multiplier is the same for every \(\Lambda\);
it is not chosen independently for each auxiliary theta vector.

Choose \(r\) vectors \(\beta_1,\ldots,\beta_r\in L_0\)
spanning \(L_0/\mathfrak aL_0\).
They form a basis over \(\mathcal O_E/\mathfrak a\):
local freeness gives cardinality
\(\operatorname N(\mathfrak a)^r\), so the surjection from the
free module of that same cardinality is bijective.
Let
\(J_{\rm lev}\subset\operatorname{GL}_r(\mathcal O_E/\mathfrak a)\)
be the prescribed finite subgroup, with character
\(\lambda:J_{\rm lev}\to\mathbb C^\times\).
Its values have modulus one, since \(J_{\rm lev}\) is finite.
For \(k\in J_{\rm lev}\), prescribe the residue of a row tuple by
\[
 \bar y_i(k)=\sum_j k_{ji}\bar\beta_j\quad(1\le i\le r),
 \qquad \bar y_i(k)=0\quad(r<i\le n).
\]
These are distinct labels, indexed exactly by \(J_{\rm lev}\).

There is a canonical residue isomorphism
\begin{equation}\label{eq:sr-level-quotient}
 H_{\Lambda,L}/K_{\Lambda,L}
              \simeq (L_0/\mathfrak aL_0)^n.
\end{equation}
To prove it, localize at a prime \(\mathfrak p\mid\mathfrak a\).
The elements \(2\) and
\(\overline{\mathfrak a}\) are units there,
\(\Lambda\) is standard there because \(M\in\mathfrak a\),
and \(L_{\mathfrak p}=(L_0)_{\mathfrak p}\).
Thus the localized quotient is precisely the usual reduction
of \(L_0^n\) modulo \(\mathfrak a\).
At primes outside \(\mathfrak a\) the quotient is zero.
The local identifications combine by the Chinese remainder
theorem, proving \eqref{eq:sr-level-quotient}.

Choose any representatives \(y_k\in H_{\Lambda,L}\) of the
prescribed labels and define the actual test vector
\begin{equation}\label{eq:sr-actual-level-test}
 \Phi_{\Lambda,L}
    =\sum_{k\in J_{\rm lev}}\lambda(k)S(y_k)v_{\Lambda,L}.
\end{equation}
This definition is independent of the representatives, with no
missing root-of-unity factor.
Indeed for \(h\in K_{\Lambda,L}\), one has
\(h\in A_{\Lambda,L}\), because the ideal ratio is
\(2\overline{\mathfrak a}^{\,2}\subset\mathcal O_E\).
Also \(y_k,h\in2(\Lambda\otimes L)\).
This untwisted tensor lattice is integral and stable under \(w\),
so
\[
 B_n(y_k,h)\in4\mathbb Z,\qquad
 B_n(h,wh)\in4\mathbb Z .
\]
Thus \(\epsilon_w(h)=1\), \(S(h)v_{\Lambda,L}=v_{\Lambda,L}\),
and the Heisenberg product law gives
\(S(y_k+h)v_{\Lambda,L}=S(y_k)v_{\Lambda,L}\).
The explicit factor \(2\) in the two quotient lattices is what
makes this representative change exact, including at two.

We can now use common representatives for two lattices \(L,L'\).
Their modules \(H_{\Lambda,L}\) and \(H_{\Lambda,L'}\) have the
same localization at primes of \(\mathfrak a\).
Their intersection therefore still surjects onto
\((L_0/\mathfrak aL_0)^n\), by the same local argument and CRT.
Choose \(y_k\) in this intersection.
The just-proved independence means that these simultaneous
representatives compute both test vectors in
\eqref{eq:sr-actual-level-test}.

\begin{srproposition}[The common operator separates distinct labels]
\label{prop:sr-level-label-orthogonality}
With those representatives,
\[
 \langle S(y_k)v_{\Lambda,L},
              S(y_{k'})v_{\Lambda,L'}\rangle=0
                      \qquad(k\ne k').
\]
\end{srproposition}

\begin{proof}
Put \(A=A_{\Lambda,L}\), \(A'=A_{\Lambda,L'}\).
For every \(\delta\in A\cap A'\), the very same operator
\(U_\delta\) from \eqref{eq:sr-common-eigenoperator} acts on the
two vectors with eigenvalues
\[
 e^{-2\pi iB_n(\delta,y_k)},\qquad
 e^{-2\pi iB_n(\delta,y_{k'})},
\]
respectively.  This follows directly by commuting \(S(\delta)\)
past \(S(y)\); the common factor \(\epsilon_w(\delta)\)
cancels against its eigenvalue on each vacuum.
If their inner product were nonzero, unitarity would force
these eigenvalues to agree.  Thus
\[
 B_n(\delta,y_k-y_{k'})\in\mathbb Z
                           \quad(\delta\in A\cap A').
\]
Consequently
\[
 y_k-y_{k'}\in(A\cap A')^\#=A^\#+(A')^\#=A+A'.
\]
The dual-intersection identity can be verified without a
duality convention: the exact sequence
\(0\to A\cap A'\to A\oplus A'\to A+A'\to0\)
splits as abelian groups since \(A+A'\) is free.
An integral functional on \(A\cap A'\) therefore extends to
one on \(A\oplus A'\), giving the asserted sum of the duals.

At \(\mathfrak p\mid\mathfrak a\), both \(A\) and \(A'\)
localize to
\(\mathfrak a_{\mathfrak p}(L_0)_{\mathfrak p}^{\,n}\):
the conjugate ideal is a unit there, and the auxiliary lattice
is standard there.
Thus \(y_k-y_{k'}\) has zero residue in
\((L_0/\mathfrak aL_0)^n\).
By \eqref{eq:sr-level-quotient} the two labels coincide.
The \(\beta_j\)'s are a residue basis, so \(k=k'\), a
contradiction.
\end{proof}

\begin{srproposition}[Exact level Gram normalization]
\label{prop:sr-exact-level-gram}
For the single construction above, with
\(b=|J_{\rm lev}|\),
\[
 \langle\Phi_{\Lambda,L},\Phi_{\Lambda,L'}\rangle
         =b\,\langle v_{\Lambda,L},v_{\Lambda,L'}\rangle.
\]
In particular
\[
 \|\Phi_{\Lambda,L}\|^2=b,\qquad
 |\langle\Phi_{\Lambda,L},\Phi_{\Lambda,L'}\rangle|
      \le b\,m(L,L')^{-n/2}.
\]
These are \eqref{eq:sr-test-gram}, now with the level phases fixed.
\end{srproposition}

\begin{proof}
Expand the two finite sums using their common representatives.
Proposition~\ref{prop:sr-level-label-orthogonality} removes all
terms with different labels.
For equal labels, the common \(S(y_k)\) is unitary and
\(|\lambda(k)|=1\), so the summand is exactly
\(\langle v_{\Lambda,L},v_{\Lambda,L'}\rangle\).

To compute the index in the vacuum overlap, note that
\(\Lambda\) is projective of rank \(n\), and the twist
\(\mathfrak j\) is an invertible module.
Local bases therefore give
\[
 A_{\Lambda,L}\cap A_{\Lambda,L'}
       =\Lambda\otimes\mathfrak j(L\cap L'),\qquad
 [A_{\Lambda,L}:A_{\Lambda,L}\cap A_{\Lambda,L'}]
       =[L:L\cap L']^n.
\]
For the index assertion, decompose each finite local quotient
into cyclic modules over its local DVR; tensoring a rank-\(n\)
free module makes \(n\) copies, and tensoring an invertible
ideal preserves its length.  Multiplication of the local
cardinalities proves the formula.
The normalized finite Heisenberg overlap is the inverse
square root of this index, as proved in
Section~\ref{sec:sr-theta-gram}.
This gives the stated bound and the diagonal norm.
\end{proof}

The multiplier on the actual level test vectors includes the fixed
level character.  If the residue action of \(\gamma\) on the
\(\beta\)-basis is \(k_\gamma\in J_{\rm lev}\), applying
\(\gamma\) sends the label \(k\) to \(k_\gamma k\).
Reindexing the finite sum, with the representative-independence
already proved, gives the exact formula
\[
 W_\gamma\Phi_{\Lambda,L}
   =\widetilde\mu(\gamma,L)\Phi_{\Lambda,\gamma L},
 \qquad
 \widetilde\mu(\gamma,L)
       =\mu(\gamma,L)\lambda(k_\gamma)^{-1}.
\]
Thus it is \(\widetilde\mu\) that enters the lattice Gram
automorphy identity and the inverse law for holomorphic families.
It remains independent of \(\Lambda\), and has modulus one.
This fixes the possible level-character normalization shift.

All choices in this section are made in the fixed standard
Hermitian space \(V_+\), the fixed rational symplectic space
\(\mathbb W\) and its single Schr\"odinger Hilbert space, from
one trace form, one semicharacter \(w\), one genuine Weil action,
the fixed level character, and common residue representatives.
The projection construction proves
\eqref{eq:sr-common-vacuum-multiplier} with a multiplier
independent of the auxiliary lattice.
It does not construct the arithmetic polarized lattice or its
integral model; those remain part of the family construction.
In the frozen source the displayed phase calculation is
\srcid{ThetaTestFunctionNormAndOverlapBound.pair_int};
the local residue implication is \srcid{cong_of_pair};
and the final diagonal sum is \srcid{phi_norm_overlap}.


The fixed-model assertion here agrees literally with the source
types.  In \srcid{FirstCoordinateAdmissible}, \srcid{Lt E n}
is a submodule of \(E^n\), and \srcid{Vset E n M} uses
\srcid{herm 1} for every auxiliary lattice.
In \srcid{ThShort}, \srcid{GenLat} has only the submodule
field \srcid{La}, and \srcid{GenLat.Unimod} uses the same
coordinate Hermitian pairing.  In \srcid{IsSystem},
\srcid{TF n g} is fixed independently of \(\Lambda\).
These are lattice variations inside a fixed space, so no
unproved isometric identification between varying rational
Hermitian forms is being used in \(P_{\Lambda,L}v_{0,L}\).
The arithmetic polarized model and the compatibility of the
global doubled theta kernel with its integral geometric
realization remain separate requirements.

\subsection{Zero overlap outside the prescribed level group}
\label{subsec:sr-level-support}

Continue in the single fixed Hermitian and Schr\"odinger model
of Section~\ref{sec:sr-level-phase-proof}.
Define the level group \(\Delta\) to consist of those rational
unitary maps \(\gamma\) which preserve \(L_0\) locally at the
primes of \(\mathfrak a\) and act on
\(L_0/\mathfrak aL_0\), in the \(\beta\)-basis, by a matrix
\(k_\gamma\in J_{\rm lev}\).
This is a subgroup: preservation and the residue matrices
are stable under multiplication and inverse.
Equivalently it is the subgroup generated by the rational
unitary maps satisfying this condition.
The latter is the source's definition \srcid{Level.grp}:
its \srcid{Cong} relation is equality in these local residue
modules, since a denominator prime to \(\mathfrak a\) is
invertible there.  Conversely, local residue equality can be
multiplied by a common denominator prime to
\(\mathfrak a\), giving exactly that \srcid{Cong} relation.

\begin{srproposition}\label{prop:sr-level-outside-zero}
For every rational unitary \(\gamma\notin\Delta\), every two
allowed \(L,L'\), and every allowed \(\Lambda\),
\[
       \langle\Phi_{\Lambda,L},
                        W_\gamma\Phi_{\Lambda,L'}\rangle=0 .
\]
\end{srproposition}

\begin{proof}
Suppose the inner product is nonzero.
Expand the two finite labelled sums and use Weil covariance.
Some pair of labels \(k,k'\) therefore satisfies
\[
 \langle S(y_k)v_{\Lambda,L},
              S(\gamma y'_{k'})W_\gamma v_{\Lambda,L'}\rangle
                                                        \ne0 .
\]
We regard rows of a tuple as column vectors when applying
\(\gamma\).
Put \(A=A_{\Lambda,L}\) and \(A'=\gamma A_{\Lambda,L'}\).
Both are selfdual lattices for the same \(B_n\).
The vector \(W_\gamma v_{\Lambda,L'}\) is fixed by
\(U_\delta\), \(\delta\in A'\): covariance gives this, and
\[
 B_n(\gamma z,w\gamma z)=B_n(z,wz)
\]
since \(\gamma\) is \(E\)-linear and unitary.
Consequently the common operators \(U_\delta\),
\(\delta\in A\cap A'\), have equal eigenvalues on the two
vectors with nonzero inner product.
Exactly as in
Proposition~\ref{prop:sr-level-label-orthogonality}, this yields
\begin{equation}\label{eq:sr-level-overlap-difference}
              y_k-\gamma y'_{k'}\in A+A'.
\end{equation}

Fix a prime \(\mathfrak q\mid\mathfrak a\), and put
\(e=v_{\mathfrak q}(\mathfrak a)>0\).
At this prime, \(\Lambda\) is standard,
\(L=L'=L_0\), and
\(\mathfrak j=\mathfrak a\overline{\mathfrak a}^{-1}\)
is \(\mathfrak q^e\) times a unit.
Choose a local basis of \(L_0\), and let \(U,U'\) be the
\(r\)-by-\(r\) matrices whose columns are the first \(r\)
rows of \(y_k,y'_{k'}\), in this basis.
These columns are integral and give bases modulo
\(\mathfrak q^e\), because their labels come from invertible
\(k,k'\) and the residue basis \(\beta\).
Their determinants are therefore units, so
\(U,U'\in\operatorname{GL}_r(\mathcal O_{E,\mathfrak q})\).

Equation~\eqref{eq:sr-level-overlap-difference} gives integral
matrices \(R,S\) with
\[
                 U-\gamma U'=\pi^e R+\pi^e\gamma S ,
\]
where \(\pi\) is any local uniformizer; the unit in the
fractional ideal has been absorbed in \(R,S\).
Thus
\[
       \gamma=(U-\pi^eR)(U'+\pi^eS)^{-1}.
\]
Both factors on the right are invertible integral matrices:
they reduce to \(U,U'\), respectively, modulo the maximal
ideal.  Hence \(\gamma\) preserves \(L_0\) at this prime and
has an integral inverse there.
We did not assume local integrality of \(\gamma\).

We may now reduce the same difference equation modulo
\(\mathfrak q^e\).
It gives \(U=\gamma U'\) in that quotient.
Using the \(\beta\)-basis and the definition of the two
labels, this says
\[
                  \gamma\beta k'=\beta k .
\]
This holds at every prime dividing \(\mathfrak a\).
The Chinese remainder theorem therefore gives one matrix
\[
                   k_\gamma=k(k')^{-1}\in J_{\rm lev}
\]
for the action of \(\gamma\) on
\(L_0/\mathfrak aL_0\).
Clearing denominators prime to \(\mathfrak a\) expresses the
same assertion as the source's \srcid{Level.Red}.
Thus \(\gamma\) is one of the generators defining \(\Delta\),
a contradiction.
\end{proof}

This proves the support condition \srcid{IsSystem.supp}.
Together with
Propositions~\ref{prop:sr-exact-level-gram}
and~\ref{prop:sr-level-isotropic-relation}, the common vacuum
multiplier, lattice-counting estimate and Neumann inverse,
the finite test-vector system now has explicit norm,
off-diagonal, automorphy, outside-level and isotropic
translation proofs in the same model.
The global doubled theta identity and the arithmetic
realization of that system still have to be matched to it.

\subsection{A nonzero-degree relation along every isotropic line}
\label{subsec:sr-level-isotropic-relation}

Use the fixed ambient model and the level test functions of
Section~\ref{sec:sr-level-phase-proof}.  Assume the level ideal is
proper.  Its conjugate is then proper as well.

\begin{srlemma}\label{lem:sr-level-line-detection}
Let \(0\ne\xi\in E^r\), and let \(y\) be any row tuple whose
first \(r\) rows span \(L_0/\mathfrak aL_0\).
For every allowed \(L,\Lambda\), there is a vector \(v\in E^n\)
such that
\[
 \delta=v\otimes\xi\in A_{\Lambda,L},\qquad
                         B_n(\delta,y)\notin\mathbb Z .
\]
Here the row tuple \(y\) can be any representative of a
nonzero level label.
\end{srlemma}

\begin{proof}
Put \(N=(E\xi)\cap\mathfrak jL\), an
\(\mathcal O_E\)-lattice of rank one on the line \(E\xi\).
It suffices to find \(\delta\in\Lambda\otimes N\) with the
stated nonintegral pairing, since every tensor in this module
has all its rows on \(E\xi\).
Localizing the intersection is legitimate: after bounding
denominators it is the kernel of a map of finitely generated
modules, and localization is exact.

Choose a prime \(\mathfrak q\mid\mathfrak a\), write
\(e=v_{\mathfrak q}(\mathfrak a)>0\), and put
\(\mathfrak p=\overline{\mathfrak q}\).
Coprimality of \(\mathfrak a,\overline{\mathfrak a}\)
gives \(\mathfrak p\ne\mathfrak q\), with
\[
 \mathfrak j_{\mathfrak p}
       =\mathfrak p^{-e}\mathcal O_{E,\mathfrak p},
 \qquad L_{\mathfrak p}=(L_0)_{\mathfrak p},\qquad
 \Lambda_{\mathfrak p}=\mathcal O_{E,\mathfrak p}^n .
\]
The last assertion holds at every prime over the underlying
rational prime \(\ell\): this prime divides \(M\), and
\(\Lambda\) is standard at \(M\).
Use completions in the rest of this local argument.

The saturated rank-one module
\(E_{\mathfrak p}\xi\cap(L_0)_{\mathfrak p}\)
has a primitive generator \(u\).
Indeed a generator lying in
\(\mathfrak p(L_0)_{\mathfrak p}\) could be divided by a
uniformizer and would not generate the intersection.
Consequently
\[
       N_{\mathfrak p}
            =\mathfrak p^{-e}\mathcal O_{E,\mathfrak p}u
                      \not\subset (L_0)_{\mathfrak p}.
\]
The trace pairing at \(\ell\) pairs the
\(\mathfrak p\)-component only with the
\(\mathfrak q\)-component.  This follows from
\(B(az,z')=B(z,\overline a z')\), applied to the idempotents
of \(E\otimes\mathbb Q_\ell\).
Selfduality of \(L_0\) says exactly that these two local
components are integral duals.  This formulation includes the
local different already present in the trace form.

Suppose every \(a\in N_{\mathfrak p}\) paired integrally with
each row \(y_i\), \(1\le i\le r\), at \(\mathfrak q\).
The same holds after multiplying each row by any local
integer, since the adjoint scalar preserves \(N_{\mathfrak p}\).
The rows span
\((L_0)_{\mathfrak q}/\mathfrak q^e(L_0)_{\mathfrak q}\).
They therefore express every \(z\in(L_0)_{\mathfrak q}\) as
a sum of such integer multiples plus an element of
\(\mathfrak q^e(L_0)_{\mathfrak q}\).
The remainder also pairs integrally with \(N_{\mathfrak p}\),
because
\(\mathfrak p^eN_{\mathfrak p}
 \subset(L_0)_{\mathfrak p}\).
It would follow that all of \(N_{\mathfrak p}\) pairs
integrally with \((L_0)_{\mathfrak q}\).
The local duality just established would give
\(N_{\mathfrak p}\subset(L_0)_{\mathfrak p}\), a contradiction.
Thus some row and some \(a\in N_{\mathfrak p}\) have
nonintegral \(\mathbb Q_\ell\)-pairing.

Place \(a\) in that row and zero in the other rows and prime
components of the \(\ell\)-adic completion of
\(\Lambda\otimes N\).
Global lattice points are dense in this completion: choose a
\(\mathbb Z\)-basis of the lattice and use the density of
\(\mathbb Z\) in \(\mathbb Z_\ell\) in each coordinate.
Approximate this local tensor closely enough that its pairing
with the fixed rational tuple \(y\) changes by an element of
\(\mathbb Z_\ell\).
The resulting global tensor \(\delta\) still has nonintegral
\(\ell\)-adic pairing, so \(B_n(\delta,y)\notin\mathbb Z\).
\end{proof}

\begin{srproposition}\label{prop:sr-level-isotropic-relation}
If in addition \(\xi\) is isotropic for the polarized
Hermitian form on \(E^r\), then there are a finite set
\(S\subset E^n\) and complex numbers \(a_v\) with
\[
 \sum_{v\in S}a_v\ne0,\qquad
 \sum_{v\in S}a_v S(v\otimes\xi)\Phi_{\Lambda,L}=0 .
\]
All these are translations in the same Schr\"odinger model.
\end{srproposition}

\begin{proof}
Let
\[
 {\cal B}_\xi=\{v\in E^n:v\otimes\xi\in A_{\Lambda,L}\}.
\]
The Hermitian isotropy gives
\(B(a\xi,b\xi)=0\) for every \(a,b\in E\).
It follows that \(S(v\otimes\xi)\), \(v\in{\cal B}_\xi\),
form an ordinary additive representation, with no Heisenberg
cocycle.  Also
\(\epsilon_w(v\otimes\xi)=1\), so these operators fix the
vacuum \(v_{\Lambda,L}\).

For a level representative \(y_k\), commuting a translation
past \(S(y_k)\) gives
\[
 S(v\otimes\xi)S(y_k)v_{\Lambda,L}
   =\chi_k(v)S(y_k)v_{\Lambda,L},\qquad
 \chi_k(v)=\exp\{-2\pi iB_n(v\otimes\xi,y_k)\}.
\]
This is a character of the additive group
\({\cal B}_\xi\).
Lemma~\ref{lem:sr-level-line-detection} supplies
\(v_k\in{\cal B}_\xi\) with \(\chi_k(v_k)\ne1\)
for each of the finitely many labels.
In the complex group algebra of \({\cal B}_\xi\) take
\[
 Z_\xi=\prod_{k\in J_{\rm lev}}
                   ([v_k]-\chi_k(v_k)[0]).
\]
Evaluation under \(\chi_k\) kills this product, by its
\(k\)-th factor.  Thus its translation action kills every
\(S(y_k)v_{\Lambda,L}\), and hence their weighted sum
\(\Phi_{\Lambda,L}\).
Its augmentation is
\[
       \prod_{k\in J_{\rm lev}}(1-\chi_k(v_k))\ne0 .
\]
Writing out the finite product as
\(Z_\xi=\sum_{v\in S}a_v[v]\) proves the assertion.
\end{proof}

This proves the finite test-function relation called
\srcid{IsSystem.cusp}, including the required nonzero
coefficient sum.  Passing this relation through an absolutely
convergent invariant integral makes that integral zero.
Using it to eliminate every nongraph term in a doubled
Eisenstein expansion still requires that expansion, its
identification of the corresponding isotropic stabilizer,
and the convergence or regularization justified for that
term.  Those analytic identifications are not consequences
of the finite group-algebra calculation alone.

\subsection{Invariant heat jets from the scalar Siegel--Weil identity}
\label{subsec:sr-invariant-heat-jets}

The scalar identity at zero elliptic variable does not, by
itself, give a Jacobi identity at arbitrary elliptic variable.
The following algebraic and analytic argument explains the
invariant differentiations actually needed for the theta
kernel.

Let \(T\) be a finite set, \(n\ge r\), and for each
\(\tau\in T\) write
\[
 U_\tau\in\operatorname{Mat}_{n,r}(\mathbb C),\qquad
 V_\tau\in\operatorname{Mat}_{r,n}(\mathbb C).
\]
Use the action
\(U_\tau\mapsto g_\tau U_\tau\),
\(V_\tau\mapsto V_\tau g_\tau^{-1}\), and put
\(s_{\tau,ab}=\sum_k(U_\tau)_{ka}(V_\tau)_{bk}\).

\begin{srlemma}\label{lem:sr-invariants-pairings}
Every complex polynomial invariant under all
\(g_\tau\in U(n)\) is a polynomial in the
\(s_{\tau,ab}\).
\end{srlemma}

\begin{proof}
Invariance under the unitary group implies invariance under
\(\operatorname{GL}_n(\mathbb C)\).
Indeed differentiating the invariant polynomial identity
along \(\exp(tK)\), for every skew-Hermitian \(K\), gives
zero infinitesimal action.
The complex span of these \(K\)'s is the full matrix algebra:
every matrix is the sum of a skew-Hermitian matrix and
\(i\) times a skew-Hermitian matrix.
The same infinitesimal identity therefore holds for every
complex matrix.  It gives invariance along its exponential.
Every invertible complex matrix has a logarithm, by putting
it in Jordan form and using the finite logarithm series on
each nilpotent part, so the full general linear group acts
invariantly.

For one \(\tau\), put
\(U_0=\binom{I_r}{0}\).
On the open set where \(U\) has rank \(r\), a general linear
change of rows carries \(U\) to \(U_0\).
Write the corresponding \(V\) as \([C,D]\), so \(C=VU\).
On the further open set where \(C\) is invertible, the
stabilizer of \(U_0\) contains
\[
       g=\begin{pmatrix}I_r&A\\0&I_{n-r}\end{pmatrix},
                   \qquad A=C^{-1}D.
\]
Right multiplication by \(g^{-1}\) sends \([C,D]\) to
\([C,0]\).  Thus an invariant polynomial \(F\) equals
\(F(U_0,[VU,0])\) on this open set.
The last expression is a polynomial in \(VU\).
The open set is nonempty and Zariski dense, so the polynomial
identity holds everywhere.
For \(n=r\) the same proof omits the empty lower block.
Apply this independently at each \(\tau\); the resulting
polynomial is obtained by fixing all \(U_\tau=U_0\) and
putting \(V_\tau=[C_\tau,0]\).
\end{proof}

For elliptic coordinates \(u_{\tau,ka},v_{\tau,bk}\), define
the constant-coefficient contraction operators
\[
 {\cal D}_{\tau,ab}
       =\sum_{k=1}^n
          \frac{\partial^2}
           {\partial u_{\tau,ka}\,\partial v_{\tau,bk}} .
\]
Let \(Z\) vary in an open complex matrix domain \(\Omega\).
Suppose \(D(Z,u,v)\) is holomorphic on
\(\Omega\times\mathbb C^{2|T|nr}\) and satisfies
\begin{equation}\label{eq:sr-exact-heat}
 {\cal D}_{\tau,ab}D
         =2\pi i\,\frac{\partial D}{\partial Z_{\tau,ab}},
                 \qquad D(Z,0,0)=0 .
\end{equation}
The coordinate order in \(Z_{\tau,ab}\) is fixed by this
display.  For a theta exponential
\[
 \exp\!\left(2\pi i\sum_k
       [\,x_kZ\overline{x_k}^{\,t}+x_ku_k+
                                      v_k\overline{x_k}^{\,t}\,]\right),
\]
direct differentiation gives precisely this heat equation,
with the corresponding row and column convention.

\begin{srproposition}\label{prop:sr-heat-invariant-zero}
For every invariant polynomial \(Q\) in the differential
symbols,
\[
                 Q(\partial_u,\partial_v)
                                    D(Z,0,0)=0 .
\]
\end{srproposition}

\begin{proof}
Lemma~\ref{lem:sr-invariants-pairings} writes
\(Q=p(s_{\tau,ab})\).
Substituting the commuting derivatives gives
\(Q(\partial_u,\partial_v)=p({\cal D}_{\tau,ab})\).
All the derivatives commute on holomorphic functions.
Repeated use of \eqref{eq:sr-exact-heat} therefore yields
\[
 p({\cal D}_{\tau,ab})D(Z,0,0)
   =p(2\pi i\,\partial_{Z_{\tau,ab}})D(Z,0,0)=0.
\]
The last equality follows by differentiating the identically
zero holomorphic function \(Z\mapsto D(Z,0,0)\).
\end{proof}

A Gaussian multiplier can be absorbed without changing this
conclusion.  Let
\[
 G(u,v)=c\exp\!\left(\sum_{\tau,k,a,b}
                     A_{\tau,ab}u_{\tau,ka}v_{\tau,bk}\right).
\]
It is invariant under the row action above.
For a polynomial differential operator \(Q\), Leibniz's rule
gives a unique polynomial \(p_G\) such that
\begin{equation}\label{eq:sr-absorb-gaussian}
   Q(\partial)(G F)(0)=p_G(\partial)F(0)
\end{equation}
for all holomorphic \(F\).
For example it is specified explicitly by
\[
              p_G(\xi)
                  =\left.Q(\partial_x)
                       \bigl(G(x)e^{\langle x,\xi\rangle}\bigr)
                                                     \right|_{x=0}.
\]
The right side is polynomial in \(\xi\).
Coefficient comparison on monomials, or direct Leibniz
expansion, proves \eqref{eq:sr-absorb-gaussian} for every
holomorphic \(F\); only finitely many derivatives occur.
If \(Q\) and \(G\) are invariant, the chain rule in this
exponential formula shows that \(p_G\) is invariant under
the dual row action.
Applying Proposition~\ref{prop:sr-heat-invariant-zero} to
\(p_G\) proves
\[
                   Q(\partial)(G D)(Z,0,0)=0 .
\]

The same argument permits fixed linear column substitutions.
More precisely, if the two sets of differentiating
directions intertwine the row action, their pullback of an
invariant bidifferential is invariant.
Right multiplication by the fixed matrices in the doubled
coordinates commutes with the left row action; the dual
variables carry the inverse action.
For a family \(P_j\) satisfying the common invariant-sum
identity, the pulled-back operator
\(\sum_jP_j(\partial_z)\overline{P_j}(\partial_w)\)
therefore has the required invariance.
The Gaussian factor is of the displayed row-contraction
form, so it can be absorbed as above.
This proves the invariant-jet implication used in the
doubled frame-kernel construction, once its stated column
coordinates are used in the same order.

In particular, to compare the class-averaged differentiated
theta series with differentiated Eisenstein series, it is
enough to prove the scalar identity at \((u,v)=(0,0)\),
the exact heat equation, local holomorphy, and normal
convergence permitting the stated differentiations.
One need not assert a stronger Jacobi identity.
This proposition proves that implication; it does not
supply the scalar Siegel--Weil identity or its
class-averaging normalization.

\section{An explicit frame-slab kernel}
\label{sec:sr-frame-slab-kernel}

We prove the archimedean integrability and reproduction used in
\eqref{eq:sr-slab-reproduction-input} for an explicit kernel.
This removes a separate analytic reproduction assumption once the
doubled Gaussian has been identified with this kernel in the chosen
highest-weight coordinates.

Put \(r=p+q\), \(J=\operatorname{diag}(1_p,-1_q)\), and
\[
 \mathcal D=\{Y\in\operatorname{GL}_r(\mathbb C):
             (Y^*JY)_{[1,p],[1,p]}>0\},\qquad U=U(p,q).
\]
Write \(\Delta_k(Y)\) for the leading \(k\)-minor,
\(\Delta_0=1\).  Let \(\lambda=(\lambda_1,\ldots,\lambda_r)\)
be integral, nonincreasing separately on the two blocks, with
\begin{equation}\label{eq:sr-classical-weight-gap}
 \lambda_i+n\le\lambda_j
           \quad(1\le i\le p<j\le r).
\end{equation}
These are the classical frame weights: the first block lies below
the second block.  They are not the algebraic weight array before
the dictionary to classical frames.

Define
\[
 m_\lambda(Y)=
    \det(Y)^{\lambda_r}
       \prod_{k=1}^{r-1}
          \Delta_k(Y)^{\lambda_k-\lambda_{k+1}},\qquad
 K_\lambda(Y,Z)=m_\lambda(JZ^*JY).
\]
Only the \(p\)-th minor can have a negative exponent in the
product; that minor is nonzero on \(\mathcal D\).
Indeed the first \(p\) columns of \(Y\) form a maximal positive
subspace, and projection to the first \(p\) coordinates is
injective on any such subspace.  Determinants are also nonzero.
Thus \(m_\lambda\) is holomorphic on \(\mathcal D\).

The same assertion for the cross matrix in \(K_\lambda\) follows
by writing the two positive \(p\)-planes as graphs
\((I,A)^t,(I,B)^t\) with \(\|A\|,\|B\|<1\).
Their cross pairing has the invertible factor
\(I-B^*A\), since \(\|B^*A\|<1\).
Consequently \(K_\lambda\) is holomorphic in \(Y\) and
antiholomorphic in \(Z\).
For an invertible upper triangular matrix \(b\), put
\(\chi_\lambda(b)=\prod_i b_{ii}^{\lambda_i}\).
Leading minors give, directly,
\begin{align}
 K_\lambda(Yb,Z)&=\chi_\lambda(b)K_\lambda(Y,Z),\label{eq:sr-kernel-weight}\\
 K_\lambda(Y,Zb)&=\overline{\chi_\lambda(b)}K_\lambda(Y,Z),\nonumber\\
 K_\lambda(uY,uZ)&=K_\lambda(Y,Z)\quad(u\in U),\nonumber\\
 K_\lambda(Z,Y)&=\overline{K_\lambda(Y,Z)},\qquad
 K_\lambda(I,I)=1.\nonumber
\end{align}
For the second weight identity, \(Jb^*J\) is lower triangular,
so its leading \(k\)-minor multiplies that of the cross matrix
by \(\prod_{i\le k}\overline{b_{ii}}\).
For symmetry use
\(JY^*JZ=J(JZ^*JY)^*J\); the two diagonal \(J\)'s cancel in
each leading minor.

\begin{srlemma}[Lowest weight uniqueness]
\label{lem:sr-frame-lowest-unique}
Suppose \(F\) is holomorphic on \(\mathcal D\), has the right
upper-triangular character \(\chi_\lambda\), and for every
diagonal unitary \(t\) satisfies
\(F(tY)=\chi_\lambda(t)F(Y)\).
Then \(F(Y)=F(I)m_\lambda(Y)\) on \(\mathcal D\).
\end{srlemma}

\begin{proof}
Both \(m_\lambda\) and \(F\) have these two transformation laws.
Their difference \(G=F-F(I)m_\lambda\) vanishes at \(I\).

First take an invertible lower triangular \(Y\in\mathcal D\) and
let \(D(z)=\operatorname{diag}(1,z,\ldots,z^{r-1})\).
The matrix
\(Y(z)=D(z)YD(z)^{-1}\), initially for \(z\ne0\), extends
polynomially to \(z=0\), since its entries below the diagonal
are \(z^{i-j}Y_{ij}\).  Its value at zero is the diagonal of \(Y\).
For \(0<|z|\le1\) it stays in \(\mathcal D\):
before the invertible right diagonal change, the positive
coordinates are scaled by at least \(|z|^{p-1}\), and the
negative coordinates by at most \(|z|^p\).
This preserves strict positivity of the first \(p\)-plane.
At \(z=0\) the diagonal matrix is also in \(\mathcal D\).
On \(|z|=1\), the two transformation laws cancel to give
\(G(Y(z))=G(Y)\).
The holomorphic function of the disk therefore has this constant
boundary value and is constant, by the maximum principle.
At zero its value is
\(\chi_\lambda(\operatorname{diag}Y)G(I)=0\).
Hence \(G(Y)=0\) on these lower triangular matrices.

On the open cell where every leading minor is nonzero,
Gaussian elimination writes a matrix as a lower triangular
matrix times an upper triangular matrix, so the right
transformation law proves \(G=0\) there.
That cell is dense in \(\mathcal D\):
replace \(Y\) by \(Y+sI\), avoiding the finite root set of each
of its monic leading-minor polynomials in \(s\).
Continuity proves \(G=0\) everywhere.
\end{proof}

\begin{srlemma}[The unitary minor bound]
\label{lem:sr-unitary-minors}
For \(u\in U(p,q)\),
\[
 1\le|\Delta_p(u)|,\qquad |\det u|=1,\qquad
 |\Delta_k(u)|\le|\Delta_p(u)|\quad(0\le k\le r).
\]
Consequently
\begin{equation}\label{eq:sr-frame-kernel-decay}
 |m_\lambda(u)|\le|\Delta_p(u)|^{-n}.
\end{equation}
If \(p=0\) or \(q=0\), interpret the last bound as \(1\).
\end{srlemma}

\begin{proof}
In block form \(u=\left(\begin{smallmatrix}A&B\\ C&D\end{smallmatrix}\right)\),
the relation \(u^*Ju=J\) gives \(A^*A-C^*C=I\).
Thus all singular values of \(A\) are at least one.
The determinant identity gives \(|\det u|=1\).

For the remaining minor inequality let \(R_p\) have the first
\(p\) rows of \(u\) and the last \(q\) rows of \(I\);
let \(R'_p\) have the first \(p\) rows of \(I\) and the last
\(q\) rows of \(u\).
The unitary relation gives
\((R'_p)^*R'_p=R_p^*R_p\), so \(V=R'_pR_p^{-1}\) is unitary.
Let \(R_k\) similarly use the first \(k\) rows of \(u\) and
the remaining rows of \(I\).
Every row of \(R_kR_p^{-1}\) is a row of either \(I\) or \(V\),
and therefore has norm one.
Hadamard's determinant inequality gives
\[
 |\Delta_k(u)/\Delta_p(u)|
       =|\det(R_kR_p^{-1})|\le1.
\]
In \(m_\lambda(u)\), every exponent
\(\lambda_k-\lambda_{k+1}\) except that at \(k=p\)
is nonnegative; bound its minor by \(|\Delta_p(u)|\).
The \(p\)-th factor is already a power of that minor, while
the full determinant has modulus one.
The resulting exponent is
\(\lambda_1-\lambda_r\le-n\), by
\eqref{eq:sr-classical-weight-gap}.
This proves \eqref{eq:sr-frame-kernel-decay}.
The definite cases use the same argument with all leading minors
bounded by one and the full determinant of modulus one.
\end{proof}

Let
\[
 S_\delta=\{Y:\ |(Y^*JY-J)_{ab}|<\delta\quad(1\le a,b\le r)\}.
\]
For sufficiently small \(\delta>0\), this open set lies in
\(\mathcal D\) and contains \(I\).
Indefinite Gram--Schmidt writes every member uniquely as
\(Y=ub\), where \(u\in U(p,q)\), \(b\) is upper triangular
with positive real diagonal, and \(b^*Jb=Y^*JY\).
All these \(b\)'s lie in a compact neighborhood of \(I\).
To see the uniform bound, the successive pivots in the
indefinite Cholesky recursion are near \(+1\) for the first block
and \(-1\) for the second; they remain bounded away from zero.
Every division and positive square root in the recursion is
therefore bounded on the closure of a sufficiently small
neighborhood of \(J\).

Ambient Lebesgue measure in these coordinates is
\(j(b)\,du\,db\), with \(j\) smooth and bounded on this compact
set of \(b\)'s.  This follows either from the change-of-variables
Jacobian or from its local principal-bundle coordinates.
There is no \(u\)-dependent factor: left multiplication by
\(u\) preserves ambient Lebesgue measure, since
\(|\det_\mathbb C u|^{2r}=1\), and preserves Haar measure \(du\).

\begin{srproposition}[Slab integrability]
\label{prop:sr-frame-slab-integrable}
If \(n\ge2r\), then for sufficiently small \(\delta>0\),
\(K_\lambda(X,\cdot)\) is bounded and absolutely integrable
on \(S_\delta\), for every \(X\in\mathcal D\).
\end{srproposition}

\begin{proof}
First take \(X=I\).  The weight relation, the compact bound on
\(b\), and Lemma~\ref{lem:sr-unitary-minors} reduce the integral to
\(\int_U|\Delta_p(u)|^{-n}\,du\).
Use the Cartan integration formula with
\(\ell=\min(p,q)\), \(t_1\ge\cdots\ge t_\ell\ge0\).
The positive-block determinant has modulus
\(\prod_i\cosh t_i\), and, up to a fixed constant, the radial
Jacobian is
\[
 \prod_{i<j}\sinh^2(t_i-t_j)\sinh^2(t_i+t_j)
 \prod_i\sinh^{2|p-q|}(t_i)\sinh(2t_i).
\]
Here is the Jacobian calculation.  Ordinary singular-value
decomposition of the off-diagonal block, followed by
\(A^*A-C^*C=I\) and its companion identities, puts \(u\), up
to left and right multiplication by \(U(p)\times U(q)\), into
\(\ell\) hyperbolic planes
\(\left(\begin{smallmatrix}\cosh t_i&\sinh t_i\\
\sinh t_i&\cosh t_i\end{smallmatrix}\right)\), with identity
on the unused coordinates.
On the regular part, differentiate this compact--diagonal--compact
map and divide out the compact centralizer.
For each pair of opposite restricted root spaces the two
compact directions are sent to their difference under
\(\exp(\operatorname{ad}H)\); its nonconstant determinant factor
is \(\sinh(\alpha(H))\) in each real root direction.
Matrix units coupling two hyperbolic planes give
\(e_i-e_j,e_i+e_j\), each with real multiplicity \(2\);
coupling a plane to an unused coordinate gives \(e_i\)
with multiplicity \(2|p-q|\); and the remaining imaginary
direction within a plane gives \(2e_i\) with multiplicity \(1\).
Their product is exactly the displayed Jacobian, up to a
positive constant.  The compact factors have finite volume.
The singular walls have lower dimension and contribute measure
zero; restricting to the ordered chamber only changes the
constant by the finite Weyl multiplicity.
Bounding \(\sinh t\le e^t\) and
\(\cosh t\ge e^t/2\), its integrand is at most a constant times
\[
 \exp\left[-\sum_{i=1}^{\ell}
                   (n-2r+4i-2)t_i\right]
 \le \exp\left[-2\sum_i t_i\right].
\]
The integral converges.  The same bounds and
\(|\Delta_p(u)|\ge1\) prove boundedness.
When \(\ell=0\), \(U\) is compact and the assertion is immediate.

For arbitrary \(X\), indefinite Gram--Schmidt gives
\(X=u_0b_0\).  Kernel invariance and its first weight relation
reduce \(K_\lambda(X,Y)\) to a constant times
\(K_\lambda(I,u_0^{-1}Y)\).
The slab and Lebesgue measure are invariant under \(u_0\).
This proves both conclusions.
\end{proof}

\begin{srproposition}[Slab reproduction]
\label{prop:sr-frame-slab-reproducing}
Under the same conditions there is \(\rho_\lambda>0\) such that
for every holomorphic \(F\) on \(\mathcal D\) with right
character \(\chi_\lambda\) and
\(\sup_{u\in U}|F(u)|<\infty\),
\begin{equation}\label{eq:sr-explicit-slab-reproduction}
 \int_{S_\delta}K_\lambda(X,Y)F(Y)\,dY
                 =\rho_\lambda F(X)\qquad(X\in\mathcal D).
\end{equation}
\end{srproposition}

\begin{proof}
The factorization \(Y=ub\) and compactness of the possible \(b\)'s
make \(F\) bounded on the slab.
All following integrals are therefore absolutely convergent.

Let \(T\subset U\) be the diagonal compact torus with probability
Haar measure and set
\[
 G(Y)=\int_T\chi_\lambda(t)^{-1}F(tY)\,dt .
\]
Compact averaging preserves holomorphy and the right character.
It gives the left torus character in
Lemma~\ref{lem:sr-frame-lowest-unique}, and \(G(I)=F(I)\).
Thus \(G(Y)=F(I)m_\lambda(Y)=F(I)K_\lambda(Y,I)\).

Kernel invariance and its weight at the base point show
\(K_\lambda(I,tY)=\chi_\lambda(t)^{-1}K_\lambda(I,Y)\).
Changing variables in the slab, which \(T\) preserves, consequently
gives
\[
 \int_{S_\delta}K_\lambda(I,Y)G(Y)\,dY
       =\int_{S_\delta}K_\lambda(I,Y)F(Y)\,dY .
\]
Using the formula for \(G\) proves reproduction at \(I\), with
\[
 \rho_\lambda=\int_{S_\delta}|K_\lambda(I,Y)|^2\,dY .
\]
This is finite by boundedness and integrability of the kernel.
It is positive because the slab is open, contains \(I\), and the
continuous kernel equals one at \((I,I)\).
For \(X=ub\), change variables by \(u\), then use the first
weight identity to reduce to the identity at \(I\) for
\(Y\mapsto F(uY)\).  This proves
\eqref{eq:sr-explicit-slab-reproduction}.
\end{proof}

Products over the archimedean embeddings have the same proof:
take the product of the domains, tori, slabs and kernels.
Their reproduction constant is the positive product of the
individual \(\rho_\lambda\)'s.

This also gives a short and exact test for the doubled Gaussian
normalization.  If the paired-derivative kernel \(K_{\rm Gau}\)
is holomorphic/antiholomorphic, unitary invariant, and has the
two characters in \eqref{eq:sr-kernel-weight}, then
Lemma~\ref{lem:sr-frame-lowest-unique} gives
\[
 K_{\rm Gau}(Y,I)=K_{\rm Gau}(I,I)m_\lambda(Y).
\]
Writing the second frame as \(ub\) proves
\(K_{\rm Gau}=K_{\rm Gau}(I,I)K_\lambda\) everywhere.
In orthonormal Gaussian coordinates its diagonal is computed
by
\[
 P(\partial_z)\overline P(\partial_w)
          \exp\!\left(\sum_j c_jz_jw_j\right)\big|_{z=w=0}
      =\sum_\alpha\alpha!c^\alpha|P_\alpha|^2,
      \qquad c_j>0 .
\]
For a nonzero orthonormal harmonic family this is strictly
positive.  Hence the normalized paired Gaussian has a positive
scalar multiple of the explicit kernel, and reproduction has
that same positive scalar multiplier.
For the actual theta model, the three stated structural
identities and the sign of the Gaussian coordinates must be
verified in its fixed frame and differential dictionary.
The argument above proves all the ensuing slab convergence,
positivity and reproduction, with no further global
automorphic input.

\section{A common theta denominator at all finite places}
\label{sec:sr-theta-common-denominator}

The denominator in the short theta spanning construction can be chosen
before its auxiliary lattice, Fourier label, harmonic polynomial, or
section.  We give its exact formula and the three valuation arguments.
The argument explicitly records the local geometric calculus it uses;
establishing that calculus for the chosen integral models remains part
of their construction.

Let \(E\) be the fixed Galois CM host, put
\[
 f=[E:\mathbb Q]/2,\quad r=p+q,\quad
 \Delta=|\operatorname{disc}E|,\quad
 d_P=\sum_{\tau,i}a_{\tau,i},\quad
 \kappa=d_P+frn/2,
\]
and let \(M\ge3\) be the fixed odd level.
Here \(n\) is even, and \(P\) ranges over the prescribed integral
harmonic polynomials of row degrees \(a_{\tau,i}\); every nonzero
such \(P\) is homogeneous of total degree \(d_P\).
The finite Fourier-label group has order
\[
 Q_M=M^{frn}.
\]
Take the actual Fourier datum of a level class \(x_0\):
its coefficients are
\(w_{x_0}(q)=\zeta_M^{\langle x_0,q\rangle}\).
They satisfy \(w_{x_0}(q)^M=1\), so they are algebraic integers
at every finite place, including places dividing \(M\).

The single integer used below is
\begin{equation}\label{eq:sr-common-theta-denominator}
 Z=Q_M\,\Delta^{d_P}M^{d_P}
   =M^{frn+d_P}\Delta^{d_P},\qquad N_a=Z^2.
\end{equation}
Its independence from the theta vector is essential.

Let \(\Lambda_a\) be an auxiliary positive definite unimodular
Hermitian lattice and
\(\Lambda_n=\Lambda_a\cap\mathcal O_E^n\).
For each embedding \(\tau:E\hookrightarrow K\), let
\(\det_\tau\Lambda_a,\det_\tau\Lambda_n\) be the determinant
fractional ideals in the common coefficient field \(K\).
For the fixed signature multiplicities \(r_\tau\), write
\(\mathcal T(J)=\prod_\tau\tau(J)^{r_\tau}\).
The scalar \(s\in K\) is chosen with the exact balancing condition
\begin{equation}\label{eq:sr-reflex-scalar-balance}
 (s)\mathcal T(\det\Lambda_n)
          \subset\mathcal T(\det\Lambda_a).
\end{equation}
This is an ideal condition; it is not replaced by an archimedean
norm estimate.

We use the following local reading of the algebraic theta section
\(Y_{\Lambda,P,x_0}\).
At a divisorial integral chart, after a finite extension of complete
DVRs, it is
\begin{equation}\label{eq:sr-theta-local-reading}
 Y_{\Lambda,P,x_0}
    =Q_M^{-1}\sum_{q}w_{x_0}(q)H_q,\qquad
 H_q=I_\Lambda^{-1}N^{-d_P}
       \sum_{z\in\mathcal Z_\Lambda}
                {\cal D}_P(F_{z,q}).
\end{equation}
Here \(N\) is an odd auxiliary denominator with \(M\mid N\),
\(I_\Lambda\) is the prescribed lattice index,
\(\mathcal Z_\Lambda\) is its isotropic torsion group, and the
\(F_{z,q}\) are the normalized symmetric translates of the
normalized multiplication pullback of the same primitive theta
section.  The differential \({\cal D}_P\) is that of the fixed
canonical connection and frame.
The theta scalar multiplying the source's short rule is included
in \(Y\); it is not an extra unit inserted after taking valuations.
Formula \eqref{eq:sr-theta-local-reading} is an equality in the
fraction field, valid with the same field embedding as \(s,P\).

Normalize a chart valuation as \(v=v_T/e\).
Under a finite extension of ramification index \(t\), use
\(v_{T'}/(te)\).  Its restriction is exactly \(v_T/e\).
Thus an inequality established after the extension descends
without a new loss.
The required connection bound, for every integral local germ of
the indicated weight and every line power, has loss
\begin{equation}\label{eq:sr-chart-ramification-loss}
 \rho(v)=\frac{d_P}{2}
   \sum_{\ell\mid\operatorname{cond}(E)}
       v(\ell)\frac{\operatorname{ord}_\ell\Delta}{[E:\mathbb Q]}
 \le d_P\,v(\Delta).
\end{equation}
All terms are nonnegative, and
\(v(\Delta)=\sum_\ell\operatorname{ord}_\ell(\Delta)v(\ell)\);
this proves the last inequality.
The more precise half-degree loss on the left is the one in the
ramified CM connection estimate.

We first explain two elementary pieces of the local argument.
For a homogeneous harmonic \(P\), normalized pullback along a CM
endomorphism \(\beta\) obeys
\begin{equation}\label{eq:sr-normalized-harmonic-pullback}
 {\cal D}_P([\beta]^*_{\rm norm}h)
        =\operatorname{Nm}_a(\beta)\,{\cal D}_P(h),
 \quad
 \operatorname{Nm}_a(\beta)
   =\prod_{\tau,i}\tau(\beta)^{a_{\tau,i}},
 \quad
 \operatorname{Nm}_a(m)=m^{d_P}.
\end{equation}
The formula uses the actual normalized pullback, its canonical
connection, and its induced integral tangent directions.
It follows by applying the differential chain rule in each
eigendirection: every letter in row \((\tau,i)\) gains
\(\tau(\beta)\), and every monomial has the same row degrees.
This also shows multiplicativity of \(\operatorname{Nm}_a\).

Suppose \(N=\beta\alpha n'\) in \(\mathcal O_E\), with \(n'\)
a unit at the chart.  If the translated section at level \(N^2\)
is the normalized \(\beta\)-pullback of a regular reached section
\(h\), then \eqref{eq:sr-normalized-harmonic-pullback} gives
\[
 \frac{\operatorname{Nm}_a(\alpha)}{N^{d_P}}
             {\cal D}_P(F)
     =(n')^{-d_P}{\cal D}_P(h).
\]
Regularity puts the germ of \(h\) in the integral power-series
ring.  The canonical connection estimate
\eqref{eq:sr-chart-ramification-loss} therefore gives valuation
at least \(-\rho(v)\); the factor \(n'\) has valuation zero.
The equality expressing the mixed symmetric translate as that
normalized pullback, and regularity of the reached section,
are the geometric assertions needed at this point.

There is also an exact cancellation of primary torsion.
Write \(\mathcal Z_\Lambda=A_\ell\times A^{(\ell)}\).
When \(\Lambda_a\) is standard at \(\ell\), every element of
\(A_\ell\) is killed by \(N\), and
\[
 |A_\ell|=\ell^{\operatorname{ord}_\ell I_\Lambda},
 \qquad I_\Lambda=|A_\ell|I^{(\ell)},\qquad
 \ell\nmid I^{(\ell)}.
\]
The normalized multiplication pullback \([N]^*_{\rm norm}\theta\)
is fixed by its normalized symmetric \(N\)-torsion translations.
Translations of coprime orders commute: their commutator is a
Weil-pairing value of order dividing both orders, hence equals
one.  Their product is the normalized translate of their product
point, by uniqueness of the symmetric normalization at odd level.
It follows that \(F_{a,a',q}\) is independent of \(a\in A_\ell\)
once \(a'\) is fixed.
This is exact equality of sections, including their phases.
Thus
\begin{equation}\label{eq:sr-primary-index-cancellation}
 I_\Lambda^{-1}\sum_{a,a'}{\cal D}_P(F_{a,a',q})
       =(I^{(\ell)})^{-1}\sum_{a'}{\cal D}_P(F_{0,a',q}).
\end{equation}
The notation \srcid{sqN} in the frozen source only identifies
the equal exponents \((N:\mathbb Z)^2\) and \(N^2:\mathbb N\).
No squaring of a section is used to obtain this cancellation.

Here are the three cases for the local reading.
They use the same connection, the same normalized translations,
and the same scalar condition \eqref{eq:sr-reflex-scalar-balance}.

\paragraph{Case \(\ell\nmid M\).}
The level translates have prime-to-\(\ell\) order.
The adapted lattice isogeny separates the \(\ell\)-primary
isotropic sum and its prime-to-\(\ell\) part.
The normalized isogeny trace cancels the
\(\ell\)-primary part of \(I_\Lambda\).
For the remaining integral harmonic germ its determinant factor
is
\(\prod_\tau\tau(\det A_\ell)^{r_\tau}\), where \(A_\ell\)
is the integral adapted lattice matrix.
The balancing condition makes
\begin{equation}\label{eq:sr-adapted-determinant-integral}
 v\left(s\prod_\tau
                  \tau(\det A_\ell)^{r_\tau}\right)\ge0.
\end{equation}
To verify the ideal step explicitly, let \(G\) be a local column
basis of \(\Lambda_n\), and let \(g\) carry
\(\Lambda_a\) to the standard lattice at \(\ell\).
The adapted matrix is a rational-integer unit times \(gG\);
the multiplying integer is prime to \(\ell\).
Hence
\[
 v(\mathcal T(\det\Lambda_n))
       =\sum_\tau r_\tau v(\tau(\det G)),\quad
 v(\mathcal T(\det\Lambda_a))
       =-\sum_\tau r_\tau v(\tau(\det g)).
\]
Applying \(v\) to \eqref{eq:sr-reflex-scalar-balance} proves
\eqref{eq:sr-adapted-determinant-integral}.
Using it in the normalized isogeny trace and applying
\eqref{eq:sr-normalized-harmonic-pullback} gives
\[
                         v(sH_q)\ge-\rho(v).
\]
The existence of the adapted isogeny and the normalized trace
identity, with this determinant factor, are required here.

\paragraph{Case \(\ell\mid M\), \(\ell^2\nmid M\).}
Take the rational integer \(\alpha=\ell\), so
\(\operatorname{Nm}_a(\alpha)=\ell^{d_P}\).
The odd-level torsion factorization expresses the mixed
translate as a normalized pullback of a reached section with
the residual auxiliary denominator prime to \(\ell\).
The preceding chain-rule argument and primary decomposition
therefore give
\[
                         v(\ell^{d_P}H_q)\ge-\rho(v).
\]
Write \(M=\ell M'\); \(M'\) is integral, so multiplication by
\((M')^{d_P}\) preserves the inequality.
Since the auxiliary lattice is standard at every prime of \(M\),
its two determinant ideals in
\eqref{eq:sr-reflex-scalar-balance} are units at the chart.
Thus \(v(s)\ge0\), and
\[
                         v(M^{d_P}sH_q)\ge-\rho(v).
\]
This uses the rational \(\ell\) as the CM endomorphism.
There is no need to choose an element of norm \(\ell\).

\paragraph{Case \(\ell^2\mid M\).}
Write
\[
 N=\ell^e n',\qquad M=\ell^k m,\qquad
 k\le e,\quad m\mid n',\quad \ell\nmid n',
\]
and take \(\alpha=\ell^k\), \(\beta=\ell^{e-k}\).
Then
\(\alpha\overline\alpha=\ell^{2k}\),
\(\beta\overline\beta=\ell^{2(e-k)}\).
A translated point \(t\) satisfies \(t^{NM}=1\), because
\(t^N\) is killed by the level ideal and hence by \(M\).
Consequently
\[
 t^{N^2}=1,\qquad
 (t^\beta)^{(n')^2\ell^{2k}}=1:
\]
the first exponent is a multiple of \(NM\), and the second is
\(\ell^{e+k}(n')^2=(NM)(n'/m)\).
These are precisely the torsion exponents needed to express the
mixed translate as a normalized \(\beta\)-pullback of a regular
reached section at line power \((n')^2\ell^{2k}\).

Use \eqref{eq:sr-primary-index-cancellation}; its remaining index
\(I^{(\ell)}\) is a unit.
Apply the pullback calculation with
\(\operatorname{Nm}_a(\alpha)=\ell^{kd_P}\).
Every remaining summand has valuation at least \(-\rho(v)\),
so the nonarchimedean triangle inequality gives
\[
                   v(\ell^{kd_P}H_q)\ge-\rho(v).
\]
The same determinant-ideal argument gives \(v(s)\ge0\), and
\(M/\ell^k=m\) is integral.  Thus again
\[
                         v(M^{d_P}sH_q)\ge-\rho(v).
\]

In each of the three cases, the integral Fourier coefficients and
\eqref{eq:sr-theta-local-reading} imply
\[
 v(sY_{\Lambda,P,x_0})
      \ge-v(Q_M)-d_Pv(M)-\rho(v).
\]
Multiplication by \(Z^2\), followed by
\eqref{eq:sr-chart-ramification-loss}, gives a nonnegative valuation.
The normalization of valuations under finite DVR extension noted
above descends it to the original divisorial chart.
If the original section already belongs to the horizontal cusp
space, this proves membership in the prescribed integral cusp
lattice, whose additional conditions are exactly these
nonnegative vertical valuations.
The same \(N_a=Z^2\) works for every auxiliary lattice, every label,
and every polynomial of the fixed row degrees.

Finally the size estimate is elementary and has no local choice
in it.  Suppose
\[
 \log M\le30L,\qquad
 \log\Delta\le7fL,\qquad f,r,L\ge1,\qquad z=frL .
\]
Since \(frn+d_P\le2\kappa\) and \(d_P\le\kappa\),
\[
 \log Z
   =(frn+d_P)\log M+d_P\log\Delta
   \le60L\kappa+7fL\kappa
   \le67z^2\kappa ,
 \qquad
 \log N_a\le134z^2\kappa .
\]
Thus the integer clearing factor contributes one fixed exponent
before the theta spanning family is chosen.

The denominator size, determinant-ideal implications, primary
factorization, and valuation bookkeeping above are complete.
For their application to the actual arithmetic models the local
reading \eqref{eq:sr-theta-local-reading}, the half-degree
connection bound \eqref{eq:sr-chart-ramification-loss}, regularity
and exact normalized pullback of reached sections, and the adapted
isogeny trace with the displayed determinant factor must be
established in those models.
These are stronger than integrality of the Fourier coefficients.
Their producers in the frozen source include
\srcid{HarmonicValuationLowerBoundHostFieldAllPlaces},
\srcid{HarmonicValueTranslatedProductValuationBound},
\srcid{ReachedSectionRegularOnIntegralModel},
\srcid{MixedTranslatePullbackEqReachedSection},
\srcid{ExistsAdaptedSimilitudeDataAtPrime},
\srcid{SumHvPrimaryKernel}, and
\srcid{ShortSeriesBoundCoprimeTorsionMul}.
The global assembly is
\srcid{ExistsBoundedIntegerGoodCoverThetaSeriesIntegral}.

\section{Classical global inputs with fixed doubling normalization}
\label{sec:sr-classical-doubling-inputs}

We match the two classical global inputs used by the theta
comparison to their precise published statements.
This establishes the anisotropic Siegel--Weil and Eisenstein
continuation inputs themselves.
It does not replace the separate local quotient comparison,
rank-recursion and covolume arguments.

Let \(E/F\) be a quadratic CM extension, with involution
\(c\), and fix a nontrivial character
\(\psi:\mathbb A_F/F\to\mathbb C^\times\).
Use the source's choice
\(\psi_\infty(t)=\exp(-2\pi i\operatorname{Tr}_{F_\infty/\mathbb R}t)\)
and the selfdual measures for it.
Let \(Y\) be a nondegenerate Hermitian or skew-Hermitian
space of dimension \(n\), and put
\[
             Y^\square=Y\oplus Y^{-},\qquad
             G^\square=U(Y^\square).
\]
The diagonal \(Y^\triangle=\{(y,y):y\in Y\}\) is a rational
maximal isotropic subspace.
Its stabilizer \(P\) is the Siegel parabolic, with Levi
\(\operatorname{Res}_{E/F}\operatorname{GL}_n\).
Thus the doubled group is quasi-split even when \(U(Y)\)
is an inner form of a different signature.

Write \(|\cdot|_E\) for the normalized adelic absolute value,
and let \(K\) be a standard maximal compact subgroup in
a rational hyperbolic frame carrying \(Y^\triangle\) to the
first isotropic coordinate space.
For the Iwasawa height
\(a(m(t)uk)=|\det t|_E\), use normalized induction
\begin{equation}\label{eq:sr-exact-induced-section}
 f_s(m(t)ug)
       =\chi(\det t)|\det t|_E^{\,s+n/2}f_s(g).
\end{equation}
The shift \(n/2\) is forced by the parabolic modulus.
Over a splitting field the unipotent radical is the space of
Hermitian \(n\)-by-\(n\) matrices.
For a diagonal \(t=\operatorname{diag}(t_i)\), its diagonal
coordinates contribute \(\prod_i|t_i|_E\), and each
off-diagonal \(E\)-coordinate contributes
\(|t_i|_E|t_j|_E\).
Their product is \(|\det t|_E^n\).
Conjugation invariance and triangular decomposition give the
same formula for every \(t\).
Hence \(\delta_P^{1/2}=|\det t|_E^{n/2}\).

\subsection{The definite theta integral and its constant}

Both ranks satisfy \(m,n\ge1\): these are explicit hypotheses
in the source's \srcid{SiegelWeilStd}, not restrictions added
for the application.  No zero-rank case is quantified there.
Let \(A\) have dimension \(m\), be of the opposite
Hermitian parity to \(Y\), and be definite at every infinite
place of \(F\).
Choose the standard unitary splitting character
\[
                    \chi|_{\mathbb A_F^\times}
                                      =\eta_{E/F}^{\,m}.
\]
Let \(\omega^\square\) be the corresponding Weil representation
for \(U(A)\times G^\square\) on its adelic Schwartz space.
For a Schwartz vector \(\Phi\) that is \(K\)-finite, the
Siegel--Weil section is
\[
 s_0=(m-n)/2,\qquad
 f_{\Phi,s}(g)
     =a(g)^{\,s-s_0}
                  (\omega^\square(g)\Phi)(0).
\]
The Weil action on the Siegel Levi has exponent \(m/2\).
Since \(s_0+n/2=m/2\), this is exactly
\eqref{eq:sr-exact-induced-section}, rather than an induction
shift chosen after the identity.

We use the following published theorem: the anisotropic
Siegel--Weil formula is holomorphic at \(s_0\) and equals
\(c_{m,n}\) times the probability-normalized theta integral,
where \(c_{m,n}=1\) for \(m>n\) and \(c_{m,n}=2\) for
\(m\le n\).  This is Theorem~7.1 of Gan--Qiu--Takeda,
\emph{The regularized Siegel--Weil formula (the second term
identity) and the Rallis inner product formula},
\url{https://arxiv.org/abs/1207.4709}.

Here is the full normalization check for our application.
A nonzero \(F\)-isotropic vector in \(A\) would remain
isotropic at every infinite place, contradicting definiteness.
Thus \(U(A)\) is anisotropic over \(F\) and its adelic
quotient is compact.
For any Haar measure \(\mu\) and any measurable fundamental
domain \(D_A\), put
\[
       {\rm Av}_{A,\mu}(H)
           =\mu(D_A)^{-1}\int_{D_A}H(h)\,d\mu(h).
\]
Its volume is finite and positive.
Replacing \(\mu\) by \(t\mu\), \(t>0\), multiplies numerator
and denominator by the same \(t\); changing fundamental
domain preserves the integral of an automorphic function.
The normalized integral therefore agrees with the
Tamagawa-normalized integral of the theorem.
For unitary groups the extra displayed factor
\(\tau(U(A))/[E:F]\) in that theorem is \(1\).
Consequently its conclusion is precisely
\begin{equation}\label{eq:sr-exact-anisotropic-sw}
 E(g,s_0,f_{\Phi})
       =\begin{cases}1,&n<m,\\2,&m\le n\end{cases}
          {\rm Av}_{A,\mu}
            \left(\sum_{x}
                  \omega^\square(g,h)\Phi(x)\right).
\end{equation}
All finite components of \(\Phi\), including ramified ones,
are arbitrary Schwartz vectors in this theorem.
No spherical-vector restriction or later adjustment of the
ramified Haar measures occurs in
\eqref{eq:sr-exact-anisotropic-sw}.
For the other Hermitian parity, choose
\(b\in E^\times\) with \(\overline b=-b\).
Replace the form on \(A\) by \(b\) times that form and the
form on \(Y\) by \(b^{-1}\) times that form.
Their Hermitian parities are reversed, their unitary groups
are unchanged, and their tensor symplectic pairing is
unchanged because the two scalars multiply to one.
This gives the skew-definite version with exactly the same
\(\psi\), measures and factor.

The source permits a rational Darboux realization \(d\), a
Siegel evaluation functional \(l\), and an action \(S\) of
\(U(A)\), rather than prescribing coordinate evaluation
in advance.
Its \srcid{IsDbl} conditions provide exactly their
normalization: rational vectors have rational Darboux
coordinates; \(l\) is invariant under the diagonal
modulations and under \(S\); and
\[
 \sum_{y\in A\otimes Y}
      l\bigl(\rho_d(y,-y)\Phi\bigr)
                               =\sum_x\Phi(x).
\]
There is a direct scalar check, requiring no unproved
uniqueness assertion for distributions.
Choose a rational polarization carrying the diagonal
\(A\otimes Y^\triangle\) to the modulation space.
Its opposite Lagrangian is \(A\otimes\{(y,-y)\}\).
The source's rationality condition makes its rational
translations the full \(F\)-lattice of that opposite
Lagrangian.  Rational changes of Schr\"odinger polarization
preserve the theta functional: linear changes reindex the
rational sum and Fourier changes use Poisson summation
with selfdual measures and the product formula.
Thus in these coordinates, writing
\[
 T_r\phi(t)=\phi(t+r),\qquad
 M_b\phi(t)=\psi(b\cdot t)\phi(t),
\]
the two conditions on \(l\) become the following statements
for \emph{all adelic} \(b\) and \emph{every} Schwartz \(\phi\):
\[
 l(M_b\phi)=l(\phi),\qquad
       \sum_{r\in F^{2mn}}l(T_r\phi)
                         =\sum_{r\in F^{2mn}}\phi(r).
\]
The rational Darboux change identifies the opposite
translation lattice with \(F^{2mn}\).
Apply the second equality also to \(M_b\phi\).
Since \(T_rM_b=\psi(b\cdot r)M_bT_r\), the first equality
gives
\[
 \sum_{r\in F^{2mn}}\psi(b\cdot r)
                 \bigl(l(T_r\phi)-\phi(r)\bigr)=0
                                     \quad(b\in\mathbb A_F^{2mn}).
\]
The series is absolutely summable: the source states an
unconditional complex \srcid{HasSum} for the left side at
\(b=0\), and the Schwartz theta sum on the right is
absolutely convergent.
Multiplication by the unitary characters gives one uniform
absolute majorant for every \(b\).
Integrating over the compact group
\((\mathbb A_F/F)^{2mn}\), of mass one, and using character
orthogonality selects \(r=0\).
It gives \(l(\phi)=\phi(0)\) for every Schwartz \(\phi\).
Thus both the evaluation functional and its scalar are
fixed by the exact \srcid{IsDbl} conditions.
In particular no freedom to multiply \(l\) by a nontrivial
constant remains.

The unspecified \(U(A)\)-action is fixed as well.
In the diagonal polarization \(U(A)\) preserves both
Lagrangians and acts on the translation coordinate space
\(X=(A\otimes_EY)(\mathbb A_F)\) by a linear map \(a_h\).
Its \(F\)-determinant has absolute value one:
the unitary determinant identity gives
\(N_{E/F}(\det_E h)=1\), and the determinant on \(X\)
is its \(n\)-th power.
Consequently the ordinary Weil action on this factor is
\[
                    (R_h\phi)(x)=\phi(a_h^{-1}x);
\]
there is no Jacobian factor.  It preserves Schwartz
functions, fixes evaluation at zero and satisfies the
same Heisenberg intertwining relation as \(S(h)\).
We now use those facts, rather than presume scalar
uniqueness of a representation.
For every Schwartz \(\phi\) and every Heisenberg
parameter \(w\), the intertwining relation in
\srcid{IsDbl.lin} and the equality \srcid{IsDbl.fix}
give
\[
 \begin{split}
 l\!\left(\rho(w)S(h)\phi\right)
   &=l\!\left(S(h)\rho(h^{-1}w)\phi\right)\\
   &=l\!\left(\rho(h^{-1}w)\phi\right)
    =l\!\left(\rho(w)R_h\phi\right).
 \end{split}
\]
The same relation holds in the transported rational
Darboux realization.
Since \(l=\operatorname{ev}_0\), taking \(w=(x,0)\)
among the pure translations says
\((S(h)\phi)(x)=(R_h\phi)(x)\) for every \(x\).
Thus \(S(h)=R_h\) on the entire Schwartz space.
This proves the compatibility with the ordinary
\(U(A)\)-Weil action, including its scalar, using precisely
the preservation, Egorov and fixed-functional conditions
present in \srcid{IsDbl}.  Behavior on functions outside
the Schwartz space is irrelevant to either side of the
identity.

We also check that the unspecified doubled action
\(\mathrm{om}\) has no remaining scalar character.
The exact definition of \srcid{IsUWeilSplit} is
\[
  \srcid{IsLagr}\ \wedge\
  \srcid{IsUSplitting}\ \wedge\ \srcid{SiegelRule}.
\]
Here \srcid{IsUSplitting} includes preservation of the
Schwartz space, Egorov intertwining, preservation of
the \(L^2\)-norm for Haar measure, and continuity after
every admissible linear functional.
Norm preservation and the inverse group element extend
each operator to a unitary operator on \(L^2(X)\).
The stated continuity gives continuous matrix
coefficients against Schwartz vectors; these are
admissible functionals, since pairing with a fixed
Schwartz vector is continuous on each fixed finite
Schwartz factor.
Density and unitarity give a continuous unitary action.

Compare this action with the ordinary Weil action for
the specified \(\chi\).
For each \(g\), their ratio commutes with all adelic
translations and modulations, by Egorov.
Such a bounded operator is scalar: commuting with
modulations makes it multiplication by a bounded
measurable function, and commuting also with
translations makes that function constant almost
everywhere.
The first assertion follows by Fourier inversion for
compactly supported test functions and then the
bounded multiplication spectral algebra.
Thus the ratio is \(c(g)\in U(1)\), and the group law
and continuity make \(c\) a continuous character.
Apply \srcid{SiegelRule} to
\(l=\operatorname{ev}_0\), which is admissible and
invariant under the diagonal modulation space.
Since the ordinary action obeys the same rule, a
Schwartz vector with \(\phi(0)\ne0\) gives \(c(p)=1\)
for every element of the adelic Siegel parabolic.

At each local place, a character trivial on this
parabolic is trivial on the whole doubled group.
Indeed the parabolic contains a Borel torus and its
positive unipotent root groups; characters are
conjugation invariant, so they are trivial on the
opposite root groups as well.
The root-group expressions for the Weyl
representatives and Bruhat decomposition then
generate the local group.
This applies also at split places, where the group
is general linear.
Finite products of local factors are dense in the
adelic restricted product, so continuity gives
\(c(g)=1\) everywhere.
The two actions agree on \(L^2\), and on Schwartz
functions agree pointwise by continuity and full
support of Haar measure.

The splitting character and doubled action in
\srcid{IsDbl.weil} are therefore exactly those in the
classical formula, not a representation identified
only up to an unspecified scalar.
Finally \srcid{IsHgt.val} gives the height just used.
The finite-dimensional \(K\)-orbit condition is precisely
\(K\)-finiteness.
Thus \eqref{eq:sr-exact-anisotropic-sw}, including its
holomorphy assertion, establishes the classical input
\srcid{SiegelWeilStd} with the source's quantifiers.
The later change from a standard hyperbolic frame to a
lattice-adapted frame transports this identity; it does not
change the factor \(1\) or \(2\).

\subsection{Meromorphic continuation in the required growth space}

For a continuous \(K\)-finite section \(f=f_{s_1}\) obeying
\eqref{eq:sr-exact-induced-section}, set
\(f_s=a^{s-s_1}f\) and initially
\[
                  E(g,s,f)
                     =\sum_{\gamma\in P(F)\backslash
                                  G^\square(F)}f_s(\gamma g).
\]
This quotient agrees with the rational parabolic quotient
in the source, since every rational maximal isotropic
frame is obtained from the fixed diagonal one.

We use Bernstein--Lapid, \emph{On the meromorphic
continuation of Eisenstein series}, JAMS~37 (2024), 187--234,
Theorem~2.3 and the corollary at the end of Section~8,
\url{https://doi.org/10.1090/jams/1020}.
They give convergence in a right half-plane and
meromorphic continuation; locally, a scalar polynomial
clears the poles into one Fr\'echet space of functions of
uniform moderate growth.

We check the section and topology hypotheses.
Restriction of \(f\) to \(K\) belongs to a finite-dimensional
continuous representation.
At the infinite places its matrix coefficients are smooth.
At finite places its compact image in a complex general
linear group is a totally disconnected compact Lie group,
hence finite; the section therefore has an open compact
stabilizer.
The parabolic transformation law extends this smoothness
from \(K\) to \(G^\square\).
The inducing representation of the Levi is the automorphic
character \(\chi\circ\det\), with the displayed spectral
twist.
It is an automorphic form, so the theorem applies to this
inducing datum without a cuspidality assumption.

At a fixed \(s_*\), restrict the theorem's several-variable
parameter to the one Siegel line \(s\).
Shrink a disk around \(s_*\) so that no other singular
hyperplane meets it.
A product of linear factors through \(s_*\) restricts to a
nonzero constant times \((s-s_*)^k\).
The right half-plane of convergence ensures that the
Siegel line is not contained in a singular hyperplane.
Thus, for some \(k\ge0\),
\[
                 H(s)=(s-s_*)^kE(\,\cdot\,,s,f)
\]
extends holomorphically across \(s_*\) in one fixed
uniform-moderate-growth Fr\'echet space.
Its zeroth seminorm supplies a fixed integer \(N\) and,
on a smaller closed disk, a fixed \(C\) with
\[
                     |H(s)(g)|\le C\,{\rm ht}(g)^N .
\]
Evaluation at \(g\), and all continuous linear growth
functionals, are holomorphic on this disk.

For clarity the latter statement gives exactly the weak
functional condition used by the source.
If a linear functional \(\lambda\) obeys
\[
 |\lambda(F)|\le C_N
       \quad\hbox{whenever}\quad |F(g)|\le{\rm ht}(g)^N ,
\]
homogeneity gives
\(|\lambda(F)|\le C_N\sup_g|F(g)|{\rm ht}(g)^{-N}\).
It is therefore continuous on the fixed growth space in
which \(H\) is holomorphic.
The composition \(s\mapsto\lambda(H(s))\) is holomorphic.
The height in the source is
\({\rm ht}_2(g)={\rm heightGL}_F(g)\) for the displayed
faithful matrix realization.
It is polynomially comparable to the algebraic matrix
height in the continuation theorem.
Indeed coordinates of either fixed faithful realization
are regular functions in the coordinates and inverse
determinant of the other; bounding the finitely many
coefficients and applying the local maximum norm bounds
these functions by a constant times a fixed power.
At all but finitely many finite places the coefficients
are integral, so the product of these constants is finite.
Applying this argument in both directions gives the
comparison.  The field, rank and matrix realizations are
fixed in this application; its constants may depend on
them, as the quantifiers of \srcid{EisBL} permit.
Multiplying the height by a constant or taking an
equivalent matrix height only alters \(C,N\).
This proves the \srcid{EisBL} input, including its
\srcid{IsGrowForm} quantifier, rather than only pointwise
meromorphic continuation.

These two propositions close the named classical inputs
to \srcid{theoremK\_head2}.
They do not by themselves prove its complete output:
\srcid{ThetaIntegralsAgree} additionally contains the
normalized real and finite quotient maps, all matching
identities, nonzero ordinary integrals and both
codimension-two recursions with character-independent
constants.
Those consequences require their displayed local and
global comparison arguments.
Nor do the two classical inputs imply
\srcid{CovolumeBoundOfX}, which belongs to the separate
integral lattice and covering-tower branch.

\section{Construction of the ramified split matching maps}
\label{sec:sr-ramified-matching-construction}

We now construct the finite matching maps at every place of \(T\).
This proves the ramified part of local matching, including its
normalization on arbitrary Schwartz--Bruhat functions.  It does not
supply the archimedean matching quotient, the inert quotient, or the
global comparison of integrals.

Let \(k=F_v\), let \(\mathcal O\) be its valuation ring, and put
\(q=\#(\mathcal O/\varpi)\), \(K_r=\mathrm{GL}_r(\mathcal O)\).
Suppose \(a'<a\), put \(d=a+a'\), and let \(\alpha,\beta\) be
smooth unitary characters of \(k^\times\).  Write
\(\nu=\alpha/\beta\) and assume \(\nu|_{\mathcal O^\times}\ne1\).
The Haar measure on each \(\mathrm{GL}_r(k)\) gives \(K_r\)
mass one.  Set \(X_0=[1_a\ 0]\), \(Y_0=[0\ 1_{a'}]\).
The two split Weil actions, in the designated linear and opposite
linear coordinates, are
\begin{align}
 A(g,h)\phi(X)
 &=\nu(\det g)^{-1}\beta(\det h)
       |\det g|^{-d/2}|\det h|^{-a/2}
                    \phi(g^{-1}Xh^{-t}),\label{eq:sr-ram-action-A}\\
 B(g',h)\phi'(Y)
 &=\nu(\det g')\alpha(\det h)
       |\det g'|^{d/2}|\det h|^{a'/2}
                    \phi'({}^tg'Yh).\label{eq:sr-ram-action-B}
\end{align}
The transpose in the second formula is part of the convention.

\paragraph{Finite support of the ramified Cartan sum.}
For \(r\ge1\), a smooth compactly supported function
\(f:M_r(k)\to\mathbb C\), and \(z\in\mathbb C\), define
\[
 Z_r(f,\nu,z)=
 \sum_{n_1\ge\cdots\ge n_r}
 \int_{K_r\operatorname{diag}(\varpi^{n_i})K_r}
              f(g)\nu(\det g)|\det g|^z\,dg .
\]
This sum has only finitely many nonzero summands.  To prove it,
choose integers \(c,e\) such that \(f=0\) outside
\(\varpi^{-c}M_r(\mathcal O)\), and \(f\) is invariant under
translation by \(\varpi^eM_r(\mathcal O)\).  If \(n_r<-c\),
the corresponding shell does not meet the support: its operator
norm is \(q^{-n_r}\), since multiplication on either side by
\(K_r\) preserves that norm.  Choose
\(u\in\mathcal O^\times\) with \(\nu(u)\ne1\).  On the shell
with diagonal representative \(t_n\), integrate first over
\(k_1,k_2\in K_r\).  The pushforward of their product probability
measure to \(K_rt_nK_r\) is its invariant probability measure.
Replace \(k_1\) by
\(k_1\operatorname{diag}(u,1,\ldots,1)\).  If \(n_1\ge e\),
the argument \(k_1t_nk_2\) changes by an element of
\(\varpi^eM_r(\mathcal O)\), whereas its character factor is
multiplied by \(\nu(u)\).  Thus this shell integral is zero.
Consequently all nonzero shells satisfy
\[
                       -c\le n_r\le n_1<e.
\]
This argument is independent of \(z\).  Each remaining term is
a constant times \(q^{-z\sum_i n_i}\), so \(Z_r\) is an entire
function of \(z\), in fact a Laurent polynomial in \(q^{-z}\).
It also proves linearity directly: take common \(c,e\) for a
finite collection of test functions.

For \(\Re z>r-1\), this sum is the ordinary absolutely convergent
integral over \(\mathrm{GL}_r(k)\).  Indeed additive matrix measure
and multiplicative Haar measure differ by a positive constant
times \(|\det g|^{-r}\).  Successive integral column elimination
on \(M_r(\mathcal O)\) reduces convergence to geometric series
with ratios \(q^{-(\Re z-i)}\), \(0\le i<r\).
A compact support is contained in a scalar dilate of this set.
This proves the assertion without a continuation theorem.
Changes of variables in that half-plane therefore continue to
every \(z\), because both sides are finite Laurent sums.

The same construction applies to \(f(g)=\phi(gX)\) when \(X\)
is any full-row-rank rectangular matrix: choose a right inverse
to \(X\).  It bounds \(g\) on the support, and a sufficiently
small matrix translation of \(g\) changes \(gX\) inside a
translation subgroup for \(\phi\).  In particular every occurrence
of \(Z_a\) below has the finite-support property just proved.

\paragraph{The two maps and their exact common target.}
Define
\begin{align}
 P_A\phi(h)
 &=\beta(\det h)|\det h|^{-a/2}
       Z_a\bigl(g\mapsto\phi(gX_0h^{-t}),\nu,d/2\bigr),
       \label{eq:sr-ram-PA}\\
 P_B\phi'(h)
 &=\alpha(\det h)|\det h|^{a'/2}
       \int_{\mathrm{GL}_{a'}(k)}
          \phi'(gY_0h)\nu(\det g)|\det g|^{d/2}\,dg .
       \label{eq:sr-ram-PB}
\end{align}
The second integral is absolutely convergent since
\(d/2>a'-1\).  Both maps are linear.

Let \(P\) be the upper block parabolic of type \((a,a')\).
The common target is the smooth normalized induced representation
\[
 \Pi_v=\operatorname{Ind}_{P}^{\mathrm{GL}_d(k)}
           \bigl((\alpha\circ\det)\boxtimes(\beta\circ\det)\bigr).
\]
Thus its functions obey, for
\(p=\left(\begin{smallmatrix}A&C\\0&D\end{smallmatrix}\right)\),
\begin{equation}
 F(ph)=\alpha(\det A)\beta(\det D)
       |\det A|^{a'/2}|\det D|^{-a/2}F(h).
 \label{eq:sr-ram-induction-law}
\end{equation}
The exponents are exactly \(\delta_P^{1/2}\).

Here is a verification for both transforms.  One has
\(X_0p^{-t}=A^{-t}X_0\) and \(Y_0p=DY_0\).
In the first integral change \(g\) to \(gA^t\);
its character and determinant factor contribute
\(\nu(\det A)|\det A|^{d/2}\).
Together with the outside factor this is
\[
 \beta(\det A\det D)|\det A\det D|^{-a/2}
        \nu(\det A)|\det A|^{d/2},
\]
which is the factor in \eqref{eq:sr-ram-induction-law}.
For \(P_B\), change \(g\) to \(gD^{-1}\); the resulting factor is
\[
 \alpha(\det A\det D)|\det A\det D|^{a'/2}
        \nu(\det D)^{-1}|\det D|^{-d/2},
\]
the same expression.  The first change of variables is justified
initially for large \(\Re z\), with \(d/2\) replaced by \(z\),
and then by the finite Laurent identity above.

A Schwartz function is fixed by some compact open subgroup
under right matrix multiplication.  Shrink the subgroup to kill
the smooth outside character.  The transforms are consequently
fixed by an open subgroup under right translation, as required
in the smooth induction.  Direct substitution in
\eqref{eq:sr-ram-action-A}--\eqref{eq:sr-ram-action-B} gives
\begin{align}
 P_A(A(g,1)\phi)&=P_A\phi,&
 P_B(B(g',1)\phi')&=P_B\phi',\label{eq:sr-ram-first-invariance}\\
 P_A(A(1,h_1)\phi)(h)&=P_A\phi(hh_1),&
 P_B(B(1,h_1)\phi')(h)&=P_B\phi'(hh_1).
 \label{eq:sr-ram-second-equivariance}
\end{align}
For the first two equalities use respectively the substitutions
\(g_{\rm int}\mapsto gg_{\rm int}\) and
\(g_{\rm int}\mapsto{}^tg'^{-1}g_{\rm int}\).
Their determinant factors cancel the displayed first-group
factors in the actions.  For the latter two equalities use the
representation laws; the first and second actions commute.

Each map is nonzero.  For \(P_A\), choose a principal congruence
subgroup \(K_a(e)\) on which \(\nu\circ\det\) is trivial, and
take \(\phi\) to be the characteristic function of
\[
       [\,K_a(e)\quad \varpi^eM_{a,a'}(\mathcal O)\,].
\]
At \(h=1\) its pullback under \(g\mapsto gX_0\) is precisely
\({\bf1}_{K_a(e)}\), so its transform is
\(\operatorname{vol}(K_a(e))>0\).  This involves only the unit
shell.  The identical construction around \(Y_0\), with the
last square block in \(K_{a'}(e)\), proves \(P_B\ne0\).
If \(a'=0\), replace this last argument by the function \(1\)
on the zero-dimensional matrix space; its transform is the
one-dimensional character \(\alpha\circ\det\).
The applications have \(a'\ge1\).

The one published representation theorem needed at this point
is Bernstein's irreducibility of normalized induction from a
unitary irreducible representation of a Levi subgroup of
\(\mathrm{GL}_d(k)\); see
\href{https://www.math.tau.ac.il/~bernstei/Publication_list/publication_texts/Bernstein-P-invar-SLN.pdf}
{Bernstein, \emph{\(P\)-invariant distributions}, Corollary in
\(\S0.2\), proved in \(\S8.2\)}.
It applies because both inducing representations are the
one-dimensional unitary representations \(\alpha\circ\det\)
and \(\beta\circ\det\).  The induction here is precisely the
normalized one, by \eqref{eq:sr-ram-induction-law}.
No restriction to cuspidal inducing data is being imposed.
Each nonzero invariant image in \(\Pi_v\) is therefore all of
\(\Pi_v\).  Thus \(P_A\) and \(P_B\) are surjective.

\paragraph{Transport to the specified Weil coordinates.}
We record the coordinate argument to specify the model involved.
At a split place,
\(E\otimes_F k=k\oplus k\), with conjugation interchanging the
summands.  Use the designated split lattice bases \(b,c\).
The two complementary tensor summands are Lagrangian; their
symplectic pairing is the matrix trace pairing, with the minus
sign dictated by the skew-Hermitian form.  In linear coordinates
the translation variable is
\[
 b^{-1}u_1c^{-t};
\]
in opposite coordinates it is
\[
                 -{}^tb\,J_Xu_0\,{}^tJ_Wc .
\]
The second formula is also the dual of the first for the
specified symplectic form.  Hence their translations and
modulations give the usual Schr\"odinger representation.
A linear change of coordinates followed, where necessary, by
Fourier transform in one Lagrangian gives isomorphisms
\[
 \iota:\mathcal S_v(V\otimes W)\longrightarrow
                         \mathcal S(M_{a,d}(k)),\qquad
 \iota':\mathcal S_v(V'\otimes W)\longrightarrow
                         \mathcal S(M_{a',d}(k)).
\]
These operations preserve Schwartz--Bruhat functions.

There is a unique lattice-vacuum line in each of these
realizations.  In the matrix model, invariance under integral
translations makes a vector constant on integral cosets;
invariance under integral modulations forces its support to
the annihilator lattice.  When \(\psi_v\) is unramified this is
the integral matrix lattice.  Thus the vacuum is a nonzero
multiple of its indicator.  Scale \(\iota,\iota'\) so that
both multiples equal \(1\).  This is compatible with every
Heisenberg operator.  In a non-self-dual additive convention,
the annihilator lattice must first be scaled by the inverse
different; the places of \(T\) in the present application are
unramified for \(\psi_v\), so that correction is absent.

The restrictions of the splitting characters to the diagonal
\(k^\times\) are trivial at a split place.  Thus their two
components are inverse characters.  For the large pair,
\((\beta/\alpha)(t\circ\sigma)=\nu(t)\) and
\(\beta(t\circ\sigma)=\beta(t)^{-1}\); for the small pair,
\((\alpha/\beta)(t)=\nu(t)\).  The dual-pair actions in these
coordinates preserve the respective Lagrangians.  Egorov's relation makes their operators scalar
multiples of the indicated linear substitutions.  The scalar is
fixed by their action at the origin, namely the Siegel rule with
the specified splitting characters.  For the large pair
\((\beta,\beta/\alpha)\) it gives precisely
\eqref{eq:sr-ram-action-A}; for the small pair
\((\alpha,\alpha/\beta)\) it gives
\eqref{eq:sr-ram-action-B}.  One can check the scalar at any
translated origin as well: translating an arbitrary test vector
reduces the whole equality to its value at zero.  Thus no
additional unit-modulus character can be inserted in these formulas.
The use of the same basis \(c\) for \(W\) ensures that the
second-group variable \(h\) in both formulas is the same
element of \(\mathrm{GL}_d(k)\).  The map from \(U(W)(k)\)
to this general linear group is onto, since its other split
component is uniquely determined by the first and the
skew-Hermitian form.

Define the actual quotient maps by
\[
                         p_v=P_A\iota,\qquad
                         p'_v=P_B\iota'.
\]
They are surjective onto the same nonzero irreducible
\(\Pi_v\).  Equation \eqref{eq:sr-ram-first-invariance} shows
that they annihilate the respective first-group difference
spaces.  Equation \eqref{eq:sr-ram-second-equivariance} gives
exact compatibility under the common second group.
Their matching criterion is pointwise equality
\[
              p_v\phi=p'_v\phi'
        \quad\Longleftrightarrow\quad
              P_A(\iota\phi)=P_B(\iota'\phi')
                    \quad\hbox{on }\mathrm{GL}_d(k).
\]
This is the exact criterion required at a place of \(T\).
Vacuum matching is required only outside \(T\).
The construction is placewise and applies to every allowed
character pair; it makes no choice of a character-dependent
global comparison constant.

\paragraph{Scope for the two rank recursions.}
The preceding construction supplies the ramified finite maps
themselves.  The standard matched vectors and their local cone
evaluations in Section~\ref{sec:ramified-whittaker} consequently
use these very quotient maps.  To infer the cone comparison for
every matched vector one still needs the codimension-two
multiplicity-one statement, together with existence and
equivariance of the descended cone and evaluation operators.
Bernstein irreducibility of \(\Pi_v\) alone does not prove that
branching statement.  Likewise this construction by itself does
not compare the two global theta integrals.  Those are separate
remaining steps in \srcid{ThetaIntegralsAgree}.


\section{Construction at unramified split places}
\label{sec:sr-unramified-split-matching}

We construct the matching maps at every split place outside \(T\).
This includes the possible poles of the unnormalized large-side
matrix integral; evaluation there means the finite normalized
Laurent polynomial constructed below.

Let \(k,\mathcal O,q,\varpi,a,a',d=a+a'\) be as in
Section~\ref{sec:sr-ramified-matching-construction}, with
\(a'<a\).  This time \(\alpha,\beta\), and hence
\(\nu=\alpha/\beta\), are unramified unitary characters.
All multiplicative Haar measures give
\(K_r=\mathrm{GL}_r(\mathcal O)\) mass one.
For a Schwartz--Bruhat function \(f\) on \(M_r(k)\), initially put
\[
 Z_r(f,s)=\int_{\mathrm{GL}_r(k)}
                 f(g)\nu(\det g)|\det g|^s\,dg,
                    \qquad \Re s>r-1 .
\]
Write
\[
 T=\nu(\varpi)q^{-s},\qquad
 D_r(T)=\prod_{i=0}^{r-1}(1-q^iT),\qquad
 c_r=\prod_{j=1}^r(1-q^{-j}).
\]
We prove that
\begin{equation}
                       D_r(T)Z_r(f,s)
                 \quad\hbox{is in }\mathbb C[T,T^{-1}],
 \label{eq:sr-unr-Laurent}
\end{equation}
and that for \(f={\bf1}_{M_r(\mathcal O)}\) this Laurent
polynomial is \(1\).

\paragraph{The integral of the standard lattice.}
Give additive matrix space measure with integral matrices of
mass one.  Then \(dg=c_r^{-1}|\det g|^{-r}\,dX\).
Let
\[
 I_r(t)=\int_{M_r(\mathcal O)}
                         \nu(\det X)|\det X|^t\,dX
\]
on its convergence half-plane; the rank-deficient matrices
have additive measure zero.  In the first column write
\(v=\varpi^j u\), \(u\) primitive, \(j\ge0\).
The corresponding column measure is
\((1-q^{-r})q^{-rj}\).  An element of \(K_r\) sends each
primitive \(u\) to the first basis vector; its determinant is
a unit, so neither character nor absolute-value weights
change.  The remaining determinant is that of an
\((r-1)\)-by-\((r-1)\) integral matrix.  Integration in the
unused first-row coordinates contributes \(1\).
Consequently
\[
 I_r(t)=\frac{1-q^{-r}}
               {1-\nu(\varpi)q^{-(t+r)}}I_{r-1}(t),
                          \qquad I_0(t)=1.
\]
Thus
\[
 I_r(t)=\prod_{j=1}^r
            \frac{1-q^{-j}}{1-\nu(\varpi)q^{-(t+j)}}.
\]
Putting \(t=s-r\) proves
\[
         Z_r({\bf1}_{M_r(\mathcal O)},s)=D_r(T)^{-1}.
\]
It proves absolute convergence for \(\Re s>r-1\) at the same
time.

\paragraph{The polynomial numerator for every test function.}
Here are the details needed to extend the preceding formula
to arbitrary \(f\), rather than merely spherical functions.
After a scalar dilation suppose its support is in
\(M_r(\mathcal O)\), and choose \(e\) such that \(f\) is
constant modulo \(\varpi^eM_r(\mathcal O)\).
We induct on \(r\).
Partition matrices by their residue modulo \(\varpi\).
On a residue class of positive rank \(k\), choose a unit
\(k\)-by-\(k\) pivot block \(A\), after integral row and column
permutations.  The block substitution
\[
 \begin{pmatrix}A&B\\ C&D\end{pmatrix}
       \longmapsto (A,B,C,S),\qquad S=D-CA^{-1}B,
\]
has additive Jacobian \(1\) and determinant
\(\det A\det S\).  The matrix \(A\) is integral invertible,
so its determinant contributes neither character nor norm.
The blocks \(A,B,C\) range over compact sets.
A finite compact-open partition, using the translation level
of \(f\), reduces their integrals to finite linear combinations
of determinant integrals in size \(r-k\), with the same
additive exponent \(t=s-r\).
Their multiplicative parameter is \(s-k\).
By induction their denominator divides
\[
          \prod_{i=k}^{r-1}(1-q^iT),
\]
which divides \(D_r(T)\).
The constants from the additive measures are independent of
\(s\).

The sole zero-residue class is \(X\in\varpi M_r(\mathcal O)\).
Substitute \(X=\varpi X_1\); its scaling factor is \(T^r\).
Repeat the preceding residue partition on this branch.
After at most \(e\) such zero branches, \(f\) is constant
on the remaining central box, and its integral is a scalar
multiple of a dilate of the standard-lattice integral.
This proves \eqref{eq:sr-unr-Laurent}.  The initial scalar
dilation contributes only a Laurent monomial in \(T\).
All operations preserve linearity in \(f\).

It follows that \(Z_r(f,s)\) has the specified rational
continuation and that its normalized value is defined at
every \(s\), even where \(D_r(T)=0\).  No division by a
vanishing value is used in evaluating the normalized map.

\paragraph{The source's finite convolution.}
Expand the convergent Cartan integral as a Laurent series
in an auxiliary variable \(X\), keeping \(s_0=d/2\) fixed:
\[
 \sum_{m\in\mathbb Z}z_mX^m,\qquad
 z_m=\!\!\sum_{\substack{n_1\ge\cdots\ge n_r\\
                         \sum_i n_i=m}}\!
       \int_{K_r\operatorname{diag}(\varpi^{n_i})K_r}
           f(g)\nu(\det g)|\det g|^{s_0}\,dg.
\]
Each inner sum is finite: compact support bounds every
\(n_i\) below, and the sum of the \(n_i\)'s is fixed.
There are only finitely many negative powers in the outer
series.  Multiplying by
\[
 D_{r,s_0}(X)=
       \prod_{i=0}^{r-1}
          (1-\nu(\varpi)q^{\,i-s_0}X)
\]
gives a Laurent polynomial by
\eqref{eq:sr-unr-Laurent}.  Its coefficient of \(X^m\) is
\[
        \sum_{j=0}^r [X^j]D_{r,s_0}(X)\,z_{m-j}.
\]
Only finitely many of these coefficients are nonzero, and
their sum is its value at \(X=1\).  This is exactly the
finite-support convolution called \(\mathrm{zetaUnr}\)
in the source; its polynomial is \(\mathrm{invL}\).
This calculation identifies that definition with the
normalized rational continuation on every Schwartz function.

\paragraph{The common quotient and the spherical vectors.}
Use the actions \(A,B\) and the designated split coordinates
of the preceding section.  At \(s_0=d/2\) put
\begin{align}
 P_A^{\rm unr}\phi(h)
 &=\beta(\det h)|\det h|^{-a/2}
   \left[
     D_a(\nu(\varpi)q^{-s})\,
     Z_a(g\mapsto\phi(gX_0h^{-t}),s)
   \right]_{s=s_0},\label{eq:sr-unr-PA}\\
 P_B^{\rm unr}\phi'(h)
 &=D_{a'}(\nu(\varpi)q^{-s_0})\,
     \alpha(\det h)|\det h|^{a'/2}
       \int_{\mathrm{GL}_{a'}(k)}
       \phi'(gY_0h)\nu(\det g)|\det g|^{s_0}\,dg .
       \label{eq:sr-unr-PB}
\end{align}
The small integral converges ordinarily because
\(s_0>a'-1\).
Moreover its normalizing factor is nonzero:
for \(0\le i<a'\),
\[
             |\nu(\varpi)q^{\,i-s_0}|<1.
\]
The large normalized expression is evaluated as a Laurent
polynomial, as proved above.

Both maps take values in
\[
 \Pi_v=\operatorname{Ind}_{P_{a,a'}}^{\mathrm{GL}_d(k)}
          \bigl((\alpha\circ\det)\boxtimes(\beta\circ\det)\bigr)
\]
with exactly the normalized parabolic law
\eqref{eq:sr-ram-induction-law}.  To verify this, do the same
changes of variables as in the preceding section on a
convergence half-plane, multiply by \(D_r\), and use equality
of rational functions before specializing to \(s_0\).
The identical argument proves first-group invariance and
second-group equivariance.  Smoothness follows from an open
right stabilizer of the input function, as before.

The common target at \(s_0\) must not be confused with a
fixed induction for all parameters.  Before specialization,
the unnormalized large map has parabolic multiplier
\[
 \alpha(\det A)\beta(\det D)
 |\det A|^{a'/2}|\det D|^{-a/2}
                         |\det A|^{\,s-s_0}.
\]
For the first-group action it satisfies
\[
 P_{A,s}(A(g,1)\phi)
                    =|\det g|^{\,s-s_0}P_{A,s}\phi .
\]
These identities follow on the convergence half-plane by the
changes of variables already given; multiply by \(D_a\)
and continue the rational identities.  The normalized
expressions are Laurent polynomials, and the extra scalar
factors are holomorphic and equal to \(1\) at \(s_0\).
Specialization therefore gives the asserted fixed induction
law and first-group invariance, including at a zero of
\(D_a\).  The second-group equivariance has no extra
parameter factor.

The two integral-matrix indicators have value \(1\) at
\(h=1\) under the respective maps, by the standard-lattice
calculation.  They are right \(K_d\)-invariant.  Their images
are equal as functions on all of \(\mathrm{GL}_d(k)\):
Iwasawa decomposition \(G=P_{a,a'}K_d\) and the fixed
parabolic transformation law determine a right-spherical
function from its value at \(1\).  One elementary proof of
this Iwasawa decomposition chooses an integral basis of the
saturated submodule obtained by intersecting the relevant
\(k\)-subspace with \(\mathcal O^d\), then extends it to an
integral basis of \(\mathcal O^d\); the resulting change of
basis lies in \(K_d\) and makes the matrix upper block
triangular.  The common vacuum image is nonzero.

Bernstein's unitary induction theorem applies to this same
\(\Pi_v\), so its two nonzero invariant images are all of
\(\Pi_v\).  Transporting the maps by \(\iota,\iota'\) therefore
gives a common nonzero irreducible quotient of the two
specified local Schwartz spaces, with first-group differences
annihilated, matching preserved by the common second group,
and equal nonzero images of the vacuum vectors.
These are all the finite matching requirements at a split
place outside \(T\).

Together with
Section~\ref{sec:sr-ramified-matching-construction}, this
constructs the finite maps at every split place.  The
construction makes no assertion about the inert places or
the archimedean quotient.  Exact kernel-equals-differences
statements at the additional good places, needed by the
global coinvariant argument, are also separate from the
surjectivity and compatibility proved here.


\paragraph{Zero-rank convention and actual range.}
For \(r=0\), set \(\mathrm{GL}_0(k)=\{1\}\), give it mass
one, set \(D_0=1\), and take the zero-dimensional matrix
Schwartz space to be \(\mathbb C\).
If \(a'=0<a\), the target is the one-dimensional
representation \(\alpha\circ\det\).
The small map is \(c\mapsto(h\mapsto\alpha(\det h)c)\).
The large normalized map has the same covariance under
the whole group \(P_{a,0}=\mathrm{GL}_a(k)\), and sends its
standard lattice vector to this function with \(c=1\).
All assertions follow with these conventions.
In the actual source's admissible triples,
\(2(p+q)+1<a'\), so \(a'\ge2\); the positive-rank proof
above covers every application.


\section{Ramified multiplicity one in codimension two}
\label{sec:sr-ramified-branching}

The irreducibility used to construct the maps in
Section~\ref{sec:sr-ramified-matching-construction} also supplies
the branching statement needed in the local rank recursions.
We prove that statement in matrix coordinates, retaining the
unit diagonal conditions that are essential to its application.

Let \(k\) be a finite extension of \(\mathbb Q_p\),
let \(d=m+2\), and suppose \(n\le m\).
Write a matrix of size \(n\)-by-\(d\) as
\[
                         Y=[\,M\mid x\mid y\,].
\]
Let \(\nu:k^\times\to\mathbb C^\times\) be smooth unitary and
ramified.  Fix \(l\in\mathcal O^\times\) with \(\nu(l)\ne1\).
Let \(P_0\) be the group of matrices
\[
 h=\begin{pmatrix}
       1_m&0&b\\ c&u&e\\0&0&v
    \end{pmatrix},
 \qquad u,v\in\mathcal O^\times .
 \tag{\(*\)}
\]
Here \(c\) is a row of length \(m\), \(b\) a column, and
\(e\in k\).  Its action on the three blocks is
\[
  [M,x,y]h=[\,M+xc,\ ux,\ Mb+ex+vy\,].
\]
For a complex vector space \(W\), call a linear map
\(L:\mathcal S(M_{n,d}(k))\to W\) admissible if
\begin{align}
 L\!\left(\nu(\det a)|\det a|^{d/2}
               \Phi({}^taY)\right)&=L(\Phi)
                       &&(a\in\mathrm{GL}_n(k)),\label{eq:sr-br-left}\\
                         L(\Phi(Yh))&=L(\Phi)
                       &&(h\in P_0).\label{eq:sr-br-right}
\end{align}
These are exactly the two invariances of the source's
\(\mathrm{Inv}\) condition.

Suppose \(W\) is a nonzero countable-dimensional irreducible
module for \(\mathrm{GL}_m(k)\), with action \(\rho\).
Suppose two admissible maps \(L,L'\) satisfy
\begin{equation}
 L(\Phi([Mg,x,y]))=\rho(g)L(\Phi),\qquad
 L'(\Phi([Mg,x,y]))=\rho(g)L'(\Phi)
                      \quad(g\in\mathrm{GL}_m(k)).
 \label{eq:sr-br-middle}
\end{equation}
If \(L\ne0\), then there is \(c\in\mathbb C\) with \(L'=cL\).

If \(n=0\), the matrix Schwartz space is \(\mathbb C\).
A nonzero equivariant \(L\) makes its image a nonzero
\(\rho\)-invariant trivial line in \(W\), so irreducibility
forces \(W\) to be that line.  Then every other equivariant
map is proportional to \(L\).
For the remainder of the proof suppose \(1\le n\le m\).

\paragraph{The local distribution tools.}
We use the localization principle and Frobenius descent for
Schwartz distributions on locally compact totally disconnected
spaces, in
\href{https://www.math.tau.ac.il/~bernstei/Publication_list/publication_texts/Bernstein-P-invar-SLN.pdf}
{Bernstein, \emph{\(P\)-invariant distributions}, \(\S1.4\)--\(\S1.5\)}.
The precise consequence needed is the following.
In a constant-rank matrix chart, let an invariant map to a
Grassmannian record a kernel, and suppose its fibers are
homogeneous spaces for the acting matrix group.
An equivariant distribution on a fiber vanishes if its
stabilizer contains an element on which the equivariance
character differs from the required modular character.
If every fiber has such an element, localization gives
vanishing on the whole chart.  Rank filtrations and the exact
sequence for an open subset and its closed complement then
give the same conclusion for a finite union of these strata.

We check the modular issue explicitly below.  All stabilizer
elements used there are semisimple with eigenvalues powers of
the unit \(l\).  Their adjoint eigenvalues on the corresponding
stabilizer Lie algebra have absolute value \(1\); hence their
modular character is \(1\).  The ambient general linear
groups are unimodular.  The group \(P_0\) is a unipotent group
extended by \((\mathcal O^\times)^2\), whose determinant action
on its unipotent radical has absolute value \(1\), and is
also unimodular.  Thus no modular factor cancels \(\nu(l)\).

For completeness these principles apply to algebraic linear
maps into \(W\).  Compose a map with any linear functional
on \(W\).  On \(\mathcal S(k^N)\), every algebraic functional
is a Schwartz distribution: this space is the inductive limit
of the finite-dimensional spaces of functions supported in a
fixed compact open box and constant modulo a fixed smaller
box.  The same holds for the locally closed matrix strata.
After proving a scalar vanishing statement, separation by the
algebraic dual of \(W\) proves its vector-valued version.

\paragraph{Removal of deficient total rank.}
On matrices of rank \(r<n\), the map \(Y\mapsto\ker Y\) takes
values in \(\operatorname{Gr}_{d-r}(k^d)\); its fibers are the
orbits of left \(\mathrm{GL}_n(k)\).  A nonzero \(n\)-by-\(r\)
minor gives local sections by Gaussian elimination, so this
is precisely a homogeneous-space family of the preceding type.
At a matrix \(Y_0\) of rank \(r<n\), choose a nonzero row
\(c\) with \(cY_0=0\), and a column \(z\) with \(cz=1\).
The projection \(P=zc\) is idempotent and
\[
                a_0=1+(l^{-1}-1)P
\]
fixes \(Y_0\) by left multiplication, has determinant \(l^{-1}\),
and has unit eigenvalues.  Equation \eqref{eq:sr-br-left}, with
\({}^ta=a_0\), supplies the nontrivial multiplier \(\nu(l)\).
Frobenius descent and localization therefore show that the
twisted coinvariants on each deficient-rank stratum are zero.
Induct over the finite rank filtration.  Consequently every
Schwartz function is congruent, modulo the first-group
differences in \eqref{eq:sr-br-left}, to a function compactly
supported on the open full-row-rank locus \(U\).
An admissible \(L\) is determined by its restriction to
\(\mathcal S(U)\).

\paragraph{The first exceptional column.}
On the open set \(x\ne0\), the \(\mathrm{GL}_n(k)\times P_0\)
orbits can be described explicitly.
Reduce modulo the line \(kx\) and put
\[
 r=\operatorname{rank}(\bar M),\qquad
 K=\ker(\bar M:k^m\to k^n/kx).
\]
One further invariant records whether \(y\) lies in
\(\operatorname{im}M+kx\).
For fixed \(r\) and this last binary value, \(K\) ranges over
a Grassmannian and is a complete orbit invariant.
To verify this, first send \(x\) to the first basis vector.
The row \(c\) in \((*)\) removes the first row of \(M\);
left changes of basis on the remaining rows identify matrices
with the same kernel \(K\).  The column \(b\) and entry \(e\)
remove the component of \(y\) in \(\operatorname{im}M+kx\).
A nonzero remaining component can be sent to a fixed last
basis vector by a change of basis on the complement.
On each pivot chart these choices are rational continuous
functions of the nonzero minors.  Thus the invariant map
has homogeneous fibers and the required local sections.

Here is a stabilizer element with nontrivial multiplier at
any such \(Y_0\).  Choose a row \(c_0\) with \(c_0x_0=1\).
Let
\[
 P=x_0c_0,\quad
 Q=e_x(c_0Y_0),\quad
 a_0=1+(l^{-1}-1)P,\quad h_0=1+(l-1)Q,
\]
where \(e_x\) is the column of the exceptional \(x\)-position.
Both \(P,Q\) are idempotent and \(PY_0=Y_0Q\).
It follows that \(a_0Y_0h_0=Y_0\).
The matrix \(h_0\) belongs to \(P_0\): only its \(x\)-row
differs from the identity, its \(x\)-diagonal entry is \(l\),
and its \(y\)-diagonal entry is \(1\).
Its right multiplier is \(1\), while the left multiplier is
\(\nu(l)\), as before.  The modular characters are \(1\).
Localization proves that \(L\) annihilates all functions
supported on \(x\ne0\), and the same is true of \(L'\).

In particular, on \(U\) these maps descend through restriction
to the closed subset \(Z=\{x=0\}\cap U\).  This uses the exact
Schwartz sequence
\[
 0\longrightarrow\mathcal S(U\setminus Z)
 \longrightarrow\mathcal S(U)\longrightarrow\mathcal S(Z)
 \longrightarrow0 .
\]
In this nonarchimedean situation there are no transverse
derivatives of a distribution supported on \(Z\).

\paragraph{The remaining exceptional rank.}
Inside \(Z\), full rank of \(Y=[M,0,y]\) forces either
\(\operatorname{rank}M=n\), or
\(\operatorname{rank}M=n-1\) and \(y\notin\operatorname{im}M\).
The latter is the only exceptional stratum.
Its orbits are again the fibers of \(M\mapsto\ker M\).
At one of its points choose \(c_0M_0=0,\ c_0y_0=1\).
Now put
\[
       P=y_0c_0,\quad
       a_0=1+(l^{-1}-1)P,\quad
       h_0=\operatorname{diag}(1_m,1,l).
\]
Again \(a_0Y_0h_0=Y_0\), \(h_0\in P_0\), and the multiplier
is \(\nu(l)\ne1\), with trivial modular factors.
Gaussian pivot charts and localization prove that the
coinvariants of this exceptional stratum are zero.
It follows that an admissible map is already determined by
its restriction to
\[
                 Z_0=\{\,x=0,\ \operatorname{rank}M=n\,\}.
\]
Precisely, if the induced map on \(\mathcal S(Z_0)\) is zero,
it factors through the exceptional stratum and hence is zero
on \(\mathcal S(Z)\), then on \(\mathcal S(U)\), and finally
on the whole matrix Schwartz space.

\paragraph{Integration in the last column.}
On \(Z_0\), the shear with column \(b\) sends \(y\) to \(y+Mb\).
Since \(M\) is onto, these are all translations in \(y\).
Their coinvariant map is integration in \(y\):
\[
                f(M,y)\longmapsto\int_{k^n}f(M,y)\,dy .
\]
Here additive Haar measure gives \(\mathcal O^n\) mass one.
To check this with varying \(M\), work on a chart where an
\(n\)-by-\(n\) minor \(M_S\) is invertible and use the fiber
coordinate \(z=M_S^{-1}y\).  The shears supported in columns
\(S\) are ordinary translations \(z\mapsto z+t\).
For compactly supported locally constant functions their
coinvariants are Haar integration: subdivide into cosets of
one common compact open subgroup; the zero-sum coefficients
are finite sums of differences of translates.
The Jacobian \(|\det M_S|\) converts integration in \(z\) to
integration in \(y\).  A finite compact-open partition of the
base support joins these chart arguments.

Let
\(\mathcal S_{\rm full}(M_{n,m}(k))\) denote Schwartz functions
compactly supported on the full-row-rank locus.
Choose indicators \(\eta_x,\eta_y\) of \(\mathcal O^n\).
The restriction of \(L\) is therefore encoded by the linear map
\[
 T_L(\psi)
       =L\bigl(\psi(M)\eta_x(x)\eta_y(y)\bigr),
       \qquad \psi\in\mathcal S_{\rm full}(M_{n,m}(k)).
\]
It is independent of the chosen test functions with
\(\eta_x(0)=1\) and \(\int\eta_y=1\).  The previous paragraphs
prove that \(L\mapsto T_L\) is injective.

Equation \eqref{eq:sr-br-left} induces on this lower space
the first-group invariance
\begin{equation}
 T_L\!\left(\nu(\det a)|\det a|^{m/2}
                           \psi({}^taM)\right)=T_L(\psi).
 \label{eq:sr-br-lower-action}
\end{equation}
Indeed \(d/2=(m/2)+1\); changing \(y\) to \({}^ta y\)
in its Haar integral contributes \(|\det a|^{-1}\).
The value at \(x=0\) contributes no determinant.
Equation \eqref{eq:sr-br-middle} gives
\(T_L(\psi(Mg))=\rho(g)T_L(\psi)\), and likewise for \(L'\).

\paragraph{The lower irreducible quotient and scalarity.}
The full-rank matrix space is a principal
\(\mathrm{GL}_n(k)\)-bundle over its Grassmannian of kernels.
Integration along its fibers with weight
\(\nu(\det a)|\det a|^{m/2}\) identifies the quotient in
\eqref{eq:sr-br-lower-action} with
\[
 Q=|\det|^{-n/2}\otimes
       \operatorname{Ind}_{P_{m-n,n}}^{\mathrm{GL}_m(k)}
                    (1\boxtimes(\nu^{-1}\circ\det)).
 \label{eq:sr-br-Q}
\]
The representation on \(Q\) here is pure right pullback.
To check the convention, the integral
\[
 |\det h|^{n/2}\!
 \int_{\mathrm{GL}_n(k)}
       \psi\bigl(a[0\ 1_n]h\bigr)\nu(\det a)|\det a|^{m/2}\,da
\]
has upper-parabolic covariance
\(|\det A|^{n/2}\nu(\det D)^{-1}|\det D|^{-(m-n)/2}\).
Its right action differs from pure pullback by
\(|\det h|^{n/2}\), giving the twist in \eqref{eq:sr-br-Q}.
The integral is compactly supported in its fiber for each
full-rank Schwartz function, so no convergence continuation is
needed here.  On a local trivialization of the frame bundle,
fiber integration has kernel generated by the first-group
differences and is onto: take a compact open fiber bump of
weighted integral \(1\).  A finite partition over the compact
Grassmannian proves the quotient identification.

Both inducing representations in \eqref{eq:sr-br-Q} are
one-dimensional unitary representations.  Bernstein's
irreducibility theorem, with the same normalization as in the
preceding section, makes \(Q\) irreducible; its scalar
\(|\det|^{-n/2}\) twist preserves irreducibility.
When \(m=n\), the parabolic has a zero block and the same
statement is the one-dimensional character directly.

The nonzero map \(T_L:Q\to W\) is an isomorphism: its kernel
and image are invariant, and both modules are irreducible.
Thus \(T_{L'}T_L^{-1}\) is an endomorphism of \(W\) commuting
with \(\rho\).  Such an endomorphism is scalar.
Here is the countable-dimensional form of Schur's argument.
If \(D\) is a nonscalar commuting endomorphism, every
\(D-\lambda\), \(\lambda\in\mathbb C\), is invertible by
irreducibility.  No nonzero polynomial in \(D\) annihilates
a nonzero vector; factoring that polynomial into linear
factors proves this assertion.  Hence for a fixed \(w\ne0\)
the vectors \((D-\lambda)^{-1}w\), \(\lambda\in\mathbb C\),
are linearly independent: a finite relation, multiplied by
the product of its distinct \(D-\lambda\)'s, would give a
nonzero polynomial annihilating \(w\).
This contradicts countable dimension.
Therefore \(T_{L'}=cT_L\), and injectivity of
\(L\mapsto T_L\) proves \(L'=cL\).

\paragraph{Application to the specified local rank steps.}
At a ramified split place the characters on the relevant
opposite model are \(\nu_v\) or \(\nu_v^{-1}\); both are
ramified.  The adapted frame
\(C=[\,\iota\mid H.x\mid H.y\,]\) is invertible and has Gram
matrix \(J_{W_1}\oplus
\left(\begin{smallmatrix}0&1\\-1&0\end{smallmatrix}\right)\).
Use it to order the complementary two columns as \(x,y\).
For \(h\in P_0\) with diagonal entries \(u,v\) on these
columns, its Levi part is represented by a local split
idele \(z\) with components \((u,v^{-1})\).
Its norm is \(1\), and the base-character restriction gives
\[
                \chi_X(z)=\chi_{X,t}(uv)
                          =\chi_{X,t}(\det h).
\]
The source's rank-step equivariance has factor
\(\chi_X(z)|z|^r\), whereas the opposite linear Weil
operator has factor
\(\chi_{X,t}(\det h)|\det h|^{n/2}\).
They agree, so cancellation leaves exactly the pure
right invariance \eqref{eq:sr-br-right}.
The first-group difference condition gives
\eqref{eq:sr-br-left}.

The common lower quotient is \(W=\Pi_{1,v}\).
Let \(\tau\) be its second-group action and identify
\(U(W_1)(k)\) with \(\mathrm{GL}_m(k)\) by the same adapted
frame.  Put
\[
 \rho(g)=
   \bigl(\chi_{X,t}(\det g)|\det g|^{n/2}\bigr)^{-1}
                      \tau(g).
\]
This scalar determinant twist is a representation and
preserves irreducibility.  The same cancellation, with the
middle block \(g\), is precisely
\eqref{eq:sr-br-middle}.  It is irreducible by the finite
matching hypotheses and countable-dimensional because it is
a quotient of the local Schwartz space.  That Schwartz space
has a countable basis obtained from compact open boxes and
their cosets.  The rank inequalities in the two steps give
\(n\le d_1=m\), as required above.

Thus two nonzero descended cone and evaluation maps of this
specified type are proportional.  Evaluation on the standard
matched vectors computed in
Section~\ref{sec:ramified-whittaker} fixes the proportionality
constant to \(\epsilon_v(-\delta)\) in the beta step, and to
\(\epsilon_v(1-\delta)\) in the alpha step.
This proves the ramified multiplicity-one input used for that
extension from the standard vectors.  Existence of the
descended operators and the global constant-term argument
are not consequences of multiplicity one and remain separate
requirements in the full rank recursion.


\section{Good split places and the exact coinvariant condition}
\label{sec:sr-good-split-coinvariants}

The finite matching constructions imply the stronger kernel
statement needed at the source's good places.  We give both
ingredients: nontriviality of the local character, and the
identification of the full Schwartz coinvariant quotient.
This proves the \(\mathrm{Coinv}\) component of
\srcid{ThetaIntegralsAgree}, with its original quantifiers
over any compatible matching maps.

\paragraph{Why the local character is nontrivial.}
Let \(E/\mathbb Q\) be Galois and CM, with CM type \(\Sigma\).
Let \(\nu\) be a unitary Hecke character whose infinity type is
\[
        \nu_\infty(z)=
          \prod_{\tau\in\Sigma}
             \left(\frac{\tau z}{|\tau z|}\right)^{\kappa_\tau},
                       \qquad \kappa_\tau>0.
\]
Suppose \(\nu\) is unramified away from a finite set \(T\).
Let a rational prime \(p\) split completely in \(E\), and
let \(P\mid p\) be outside \(T\).  Then the local
character at \(P\) is nontrivial.

Indeed assume that the character at \(P\) is trivial.
Choose \(h>0\) with \(P^h=(\pi)\), \(\pi\in\mathcal O_E\);
the ideal class group is finite.  At every other finite
prime \(\pi\) is a unit.  Its character value is \(1\) outside
\(T\); at the finitely many primes of \(T\), its values are
roots of unity.  For the latter assertion, the continuous
character of the compact local unit group has finite image:
it is a compact totally disconnected subgroup of
\(\mathbb C^\times\), hence finite.  Choose \(M>0\) killing
all these values.  The finite-idele value at \(\pi^M\) is
therefore \(1\), and the Hecke relation implies
\[
        \prod_{\tau\in\Sigma}
          \left(\frac{\tau\pi}{|\tau\pi|}\right)^{M\kappa_\tau}
                          =1.
\]
Squaring and using
\(|\tau\pi|^2=\tau\pi\,\overline{\tau\pi}\) gives
\begin{equation}
 \prod_{\tau\in\Sigma}(\tau\pi)^{M\kappa_\tau}
     =\prod_{\tau\in\Sigma}(\overline{\tau\pi})^{M\kappa_\tau}.
 \label{eq:sr-good-character-contradiction}
\end{equation}

Fix \(\tau_0\in\Sigma\).  Since \(E/\mathbb Q\) is Galois,
write \(\tau=\tau_0g_\tau\), \(g_\tau\in\operatorname{Gal}(E/\mathbb Q)\),
with \(g_{\tau_0}=1\), and let \(c\) be complex conjugation.
Pull \eqref{eq:sr-good-character-contradiction} back to \(E\)
and take its \(P\)-adic valuation.  The left side has a
strictly positive contribution from \(g_{\tau_0}=1\).
The right side has none: a factor could have positive
valuation only if \(g_\tau c\) fixed \(P\).  A completely split
prime has trivial decomposition group, so this would force
\(g_\tau=c\), that is \(\tau=\bar\tau_0\), which is not in the
CM type \(\Sigma\).  This contradiction proves the assertion.
The same argument applies to every \(P\mid p\) outside \(T\).

In the present packet \(\nu=\alpha/\beta\) and
\(\kappa_\tau=m_\alpha(\tau)-m_\beta(\tau)>d>0\).
The source's \(\mathrm{Good}\) condition says that \(v\notin T\)
and its underlying rational prime splits completely in \(E\).
Thus every designated local component \(\nu_t\), for either
embedding \(t:E\to F_v\), is nontrivial.  Such a place splits
in \(E/F\), since complete splitting gives local degree one
in every intermediate field.

\paragraph{The full-rank quotient.}
Let \(k\) be nonarchimedean of characteristic zero, let
\(1\le n\le d\), and let \(\nu\ne1\) be smooth unitary.
Consider either of the two matrix actions
\[
 \nu(\det g)^{-1}|\det g|^{-d/2}\phi(g^{-1}X),
 \qquad
 \nu(\det g)|\det g|^{d/2}\phi({}^tgX).
\]
The additional character and determinant factor for the
second group does not change the first-group differences.
Let \(\mathcal R\) be the span of those differences, and
\(\mathcal S_{\rm full}\) the functions compactly supported
on full-row-rank matrices.

Every class in \(\mathcal S(M_{n,d}(k))/\mathcal R\)
has a representative in \(\mathcal S_{\rm full}\).
To see this, filter the closed deficient-rank locus by rank.
On its rank-\(r\) stratum, the kernel map to
\(\operatorname{Gr}_{d-r}(k^d)\) has homogeneous fibers under
\(\mathrm{GL}_n(k)\); Gaussian pivot charts give local
sections.  The stabilizer of
\(\left(\begin{smallmatrix}1_r&0\\0&0\end{smallmatrix}\right)\)
contains
\[
                \begin{pmatrix}1_r&B\\0&D\end{pmatrix},
              \qquad D\in\mathrm{GL}_{n-r}(k).
\]
Its determinant ranges over all of \(k^\times\).
The modular character is a positive real power of
\(|\det D|\), as is the modulus in either displayed action.
For equivariant distributions on this homogeneous space to
exist, their character restricted to this stabilizer must
equal that positive modular character.
This is impossible: choose \(D\) with
\(\nu(\det D)\ne1\).  A nontrivial unit-modulus number cannot
be a positive real number.  Frobenius descent and Bernstein
localization, with the kernel map just specified, give zero
coinvariants on each deficient-rank stratum.
The exact Schwartz sequences for the rank filtration show
that the quotient of the whole Schwartz space by
\(\mathcal S_{\rm full}\) has zero coinvariants.
Right exactness of coinvariants proves the claimed
surjectivity from \(\mathcal S_{\rm full}\).

The full-row-rank frame bundle over the Grassmannian has
free left \(\mathrm{GL}_n(k)\)-fibers.
Fiber integration with the displayed character and norm
weight identifies its coinvariants with a normalized
parabolic induction of two unitary determinant characters,
up to a scalar determinant twist for the right action.
The verification is the same fiber-bump and finite
partition argument as for
\eqref{eq:sr-br-Q}.  For the two complementary matrix models
with \(n=a\) or \(a'\), it is the common induction
\[
 \operatorname{Ind}_{P_{a,a'}}^{\mathrm{GL}_d(k)}
          ((\alpha\circ\det)\boxtimes(\beta\circ\det))
\]
in the action convention of their matching maps.
It is irreducible by Bernstein's unitary induction theorem.

Consequently the natural surjection from the full-rank
coinvariant quotient onto the whole Schwartz coinvariant
quotient is either zero or an isomorphism.  It is nonzero.
At a ramified split place the explicit full-rank bump in
Section~\ref{sec:sr-ramified-matching-construction} has
nonzero transform.  At an unramified split place with
\(\nu\ne1\), the same bump works for the normalized maps:
the large-side normalizing factor does not vanish, because
a zero of
\(\prod_{i<a}(1-\nu(\varpi)q^{\,i-d/2})\)
would require \(i=d/2\) and \(\nu(\varpi)=1\).
An unramified nontrivial character has
\(\nu(\varpi)\ne1\).  The small-side factor is always
nonzero by \(a'<d/2\).
Thus the whole coinvariant quotient is itself this nonzero
irreducible representation.

This argument uses right exactness only.  It does not
assume that coinvariants preserve injections.

\paragraph{Any matching map has the exact kernel.}
Now let \(p\) be any of the source's finite matching maps:
it is onto a nonzero irreducible \(\Pi\) and annihilates
\(\mathcal R\).
The induced map
\[
       \mathcal S(M_{n,d}(k))/\mathcal R\longrightarrow\Pi
\]
is a nonzero surjection between irreducible modules, hence
an isomorphism.  Therefore, on Schwartz functions,
\[
                       p\phi=0\quad\Longleftrightarrow\quad
                                      \phi\in\mathcal R.
\]
Moreover \(\mathcal S_{\rm full}\) surjects under \(p\),
because it already surjects onto the whole coinvariant
quotient.  Both assertions survive the linear and opposite
Weil coordinate isomorphisms.

If \(\ell\) is any linear functional on the ambient
function space that annihilates \(\mathcal R\), its
restriction to Schwartz functions factors uniquely through
\(p\).  Define \(\ell_0(p\phi)=\ell(\phi)\); the kernel
identity proves well-definedness and surjectivity of \(p\)
gives its entire domain.  This is exactly the functional
factorization demanded in \(\mathrm{Coinv}\), rather than
a statement about merely continuous functionals.

Apply this argument at a good \(v\) using either embedding
\(t:E\to F_v\).  The corresponding opposite linear
realization sends the source's
\(\mathcal S^\circ_{\mathrm{loc}}(t)\) to the full-row-rank
Schwartz space.  Since the local character is nontrivial
for each \(t\), its image under \(p'_v\) is all of \(\Pi_v\).
Changing the split embedding may change that full-rank
subspace, but the surjectivity argument applies separately
to each one and uses the same already chosen quotient map.

\paragraph{Stability of the global small subspace.}
At the distinguished split place \(v_1\), the opposite linear
Weil action consists of invertible left and right matrix
substitutions times a scalar, as in
\eqref{eq:sr-ram-action-B}.  Such substitutions preserve
matrix rank.  The local full-rank subspace is therefore
stable under both members of the dual pair.
At all other finite places the local actions preserve
Schwartz--Bruhat functions; at almost all places the
adelic group element is integral and fixes the vacuum.
At infinity the Weil action preserves Schwartz functions.
The restricted tensor product action consequently carries
each defining finite sum of pure tensors for
\(\mathcal S^\circ\) to another such finite sum, with its
\(v_1\)-factor still in the same full-rank subspace.
Thus \(\omega'(\gamma)\mathcal S^\circ
\subseteq\mathcal S^\circ\) for every adelic pair \(\gamma\).

Combining the exact kernels, the full-rank surjections for
every split embedding, the factorization of every
annihilating functional, and this global stability proves
all clauses of \(\mathrm{Coinv}\).
It does not supply the nonzero global theta integral or the
two character-uniform rank-step constants in
\srcid{ThetaIntegralsAgree}.

\section{Constructing the inert comparison in every parity}
\label{sec:sr-inert-matching-construction}

We construct the finite matching maps at an inert place outside
the prescribed ramification set. The ranks of the two Hermitian
spaces need not be even. This distinction matters: results stated
only for even Hermitian rank do not supply this clause of the packet.

Write \(k=F_v\), \(L=E_w\), and let \(\overline{\phantom{x}}\)
denote the nontrivial automorphism of the unramified quadratic
extension \(L/k\). Put
\[
 d=a+a'=2n,\qquad a>a'\ge1,
 \qquad K_H=U(\Lambda_H),\quad K=U(\Lambda_W),
\]
where \(H\) is either \(V\) or \(V'\), and \(\Lambda_H,\Lambda_W\)
are the actual localizations of the prescribed lattices.
The additive character is the prescribed \(\psi_v\).
Haar measure on each compact group in this section has total mass one.

\subsection{Why the actual local data are unramified}
The base hypotheses at an inert place give odd residue characteristic,
residue cardinality greater than three, and
\[
 \{x\in k:\psi_v(x\mathcal O_k)=1\}=\mathcal O_k.
\]
For the standard global character the group on the left is
\(\mathfrak D_{k/\mathbb Q_p}^{-1}\). This follows directly by
localizing the trace-dual definition of the codifferent, or by
writing the character as the standard character of \(\mathbb Q_p\)
composed with the trace. Consequently the different exponent of
\(F/\mathbb Q\) at \(v\) is zero. The tower formula for differents
and the unramifiedness of \(E/F\) give
\(\mathfrak D_{L/\mathbb Q_p}=\mathcal O_L\).
This checks a potentially relevant normalization: the source's
skew lattice is initially defined using the absolute trace dual.
Here that dual is the ordinary \(\mathcal O_L\)-dual.
It is therefore a genuinely unimodular skew-Hermitian lattice,
just as the two Hermitian lattices are unimodular.

Choose a trace-zero unit \(\delta\in\mathcal O_L^\times\).
Multiplying a skew-Hermitian form by \(\delta^{-1}\) makes it
Hermitian and changes none of these duality assertions.
Such a unit exists by lifting an element of the residue field
whose conjugate differs from itself; subtract its conjugate.
The norm map on units is surjective. One way to verify this is
to use the surjectivity of the norm between the two finite
residue fields and then lift successively: the first-order norm
of \(1+\varpi^r x\) is \(1+\varpi^r\operatorname{Tr}(x)\).
The residue trace is onto, so each correction can be made.

Reduction modulo \(\varpi\), followed by Hensel lifting, now
splits a hyperbolic plane off every even unimodular Hermitian
lattice. To give the lifting step explicitly, start with a primitive
isotropic residue vector, lift it, choose a partner whose pairing
is a unit, and adjust along that partner to make its norm zero.
The equation has a unit derivative because two is a unit.
Normalize the pairing and replace the partner by an isotropic
partner by subtracting half its norm times the first vector.
The orthogonal complement is again unimodular.
Induction gives a hyperbolic basis of \(\Lambda_W\), since its
rank \(2n\) is even. In particular \(U(W)\) is the split-rank
unitary group \(U(n,n)(k)\), with its specified hyperspecial subgroup
\(K\). For either parity of \(a\) and \(a'\), the same argument
gives the usual unimodular local Hermitian models of \(V,V'\).

Let \(\alpha_v,\beta_v\) be the two local splitting characters.
They are unramified and their restrictions to \(k^\times\)
are respectively \(\eta^{a'}\) and \(\eta^a\), where \(\eta\)
is the unramified quadratic character of \(k^\times\).
The common uniformizer of \(k\) and \(L\) therefore satisfies
\[
 \alpha_v(\varpi)=(-1)^{a'}=(-1)^a=\beta_v(\varpi).
\]
Their values on units are one, so
\begin{equation}\label{eq:sr-inert-character-equality}
 \alpha_v=\beta_v=\chi,\qquad \chi^2=1,\qquad
                 \alpha_v/\beta_v=1.
\end{equation}
Thus the first member of each dual pair has trivial auxiliary
character, and the second member has the same unramified splitting
character \(\chi\). This verifies the splitting data, including
the case of odd \(a,a'\).

\subsection{Transport to the ordinary local model}
Use a hyperbolic polarization \(W=Y\oplus Y^\vee\) with
\(\Lambda_W=\Lambda_Y\oplus\Lambda_{Y^\vee}\).
For \(H=V,V'\), put
\[
 \mathcal S_H=\mathcal S(H\otimes_LY^\vee)
              \simeq\mathcal S(H^n),\qquad
                     e_H=1_{\Lambda_H^n}.
\]
All additive measures are self-dual for \(\psi_v\).
Because the tensor lattice is selfdual, its volume is one
in dual coordinates; the Fourier transform carries \(e_H\)
to the corresponding dual indicator.
Integral unipotents fix \(e_H\), since their quadratic phases
are integral and two is a unit. Integral Levi matrices fix it
by the determinant formula; their splitting-character factor is one.
The Fourier generators fix it as well, by the preceding selfdual
Fourier identity and the conductor-zero, odd-residue normalization.
These are the lattice-model generators of both hyperspecial groups,
so \(e_H\) is fixed by \(K_H\times K\).

This model is the actual model in different coordinates, not an
additional choice of an unrelated Weil representation.
Indeed a lattice frame of \(\Lambda_H,\Lambda_W\) identifies their
tensor lattice with integral matrices and identifies the symplectic
pairing with the trace pairing just used. The Stone--von Neumann
intertwiner transports the prescribed Heisenberg operators to
ordinary Schr\"odinger translations and modulations.
The phase in the prescribed vacuum condition is one here:
the tensor-lattice pairings and the integral element \(\theta\)
occurring in that condition are integral, and two is a unit.
Its lattice-fixed line is one-dimensional: modulation by the
integral lattice forces support in its annihilator, and integral
translations then force a constant on that lattice.
Rescale the intertwiner to send the prescribed nonzero vacuum
exactly to \(e_H\). This also handles any original normalization
of that vacuum.

Egorov covariance implies that the transported Weil operators
differ from the ordinary ones by a character of
\(U(H)\times U(W)\). The first member preserves the whole
Schr\"odinger polarization \(H\otimes_LY^\vee\).
The specified Siegel rule on this polarization fixes its action
for every \(g\in U(H)\): it is pullback, since its determinant
has norm one and its auxiliary character is one.
On the second member the same rule fixes the scalar discrepancy
on the Siegel parabolic \(P\), including its Levi.
A character trivial on \(P\) is trivial on its opposite as well,
because characters are conjugation invariant.
The two opposite parabolics generate \(U(W)\):
their root groups generate the special unitary subgroup, and
the determinant subgroup is already represented by the Levi
(the latter assertion is Hilbert's theorem 90).
Thus the scalar discrepancy is one. This argument includes
\(U(1,1)\); for a compact rank-one first member its entire action
was already fixed by the polarization rule.
The source and ordinary actions, their compact stabilizers
and their vacua have all been identified simultaneously.

Evaluation at zero now supplies a nonzero first-member-invariant
linear functional
\[
       e_H^0(\phi)=\phi(0),\qquad e_H^0(e_H)=1.
\]
The invariance uses \eqref{eq:sr-inert-character-equality}:
the first member acts by pullback with no remaining character.
Let \(P\) be the Siegel parabolic of \(U(W)\), stabilizing \(Y\).
For a Levi element with action \(t\in\mathrm{GL}_n(L)\) on \(Y\),
the prescribed Siegel rule gives
\begin{equation}\label{eq:sr-inert-parabolic-functional}
 e_H^0\bigl(\omega_H(p)\phi\bigr)
   =\chi(\det t)|\det t|_L^{\dim_L H/2}\,e_H^0(\phi).
\end{equation}
The unipotent radical has scalar one in this formula.

\subsection{The two complementary spherical functions coincide}
Set \(r_a(p)=\chi(\det t)|\det t|_L^{a/2}\) and
\(r_{a'}(p)=\chi(\det t)|\det t|_L^{a'/2}\).
The modular character of \(P\) is
\[
                   \delta_P(p)=|\det t|_L^n.
\]
For example the unipotent radical is the additive space of
Hermitian \(n\)-by-\(n\) matrices; conjugation by \(t\)
has \(k\)-determinant of absolute value
\(|\det t|_L^n\), as is seen first for diagonal \(t\)
and then by multiplicativity. Hence
\begin{equation}\label{eq:sr-inert-complementary-character}
                          r_a r_{a'}=\delta_P.
\end{equation}

Define
\[
 f_a(h)=e_V^0(\omega_V(h)e_V),\qquad
 f_{a'}(h)=e_{V'}^0(\omega_{V'}(h)e_{V'}).
\]
They are right \(K\)-invariant, equal to one on \(K\), and
satisfy \(f_a(ph)=r_a(p)f_a(h)\), with the analogous rule
for \(a'\). The Iwasawa decomposition \(U(W)=PK\)
therefore determines each function uniquely.
Let \(I_a,I_{a'}\) be the smooth induced spaces with these
transformation rules. The compact-picture bilinear pairing
\[
                    B(F,F')=\int_K F(k)F'(k)\,dk
\]
is invariant under simultaneous right translation.
To check the assertion, decompose \(kh=p(k,h)k_h\).
The change of variable on \((P\cap K)\backslash K\)
has density \(\delta_P(p(k,h))^{-1}\).
The product of the two transformation factors is
\(\delta_P(p(k,h))\) by
\eqref{eq:sr-inert-complementary-character}, and cancels it.
The inner integration over \(P\cap K\) has probability mass one.
This proves the assertion directly, also at reducible induction
parameters.

Normalize the compact averages of the evaluation functionals by
\[
 \ell_a(\phi)=\int_K e_V^0(\omega_V(k)\phi)\,dk,\qquad
 \ell_{a'}(\phi')=\int_K e_{V'}^0(\omega_{V'}(k)\phi')\,dk.
\]
Each integral is a finite average for every smooth vector.
Both functionals annihilate the first-member augmentation,
are \(K\)-invariant, and take the value one on their vacua.
Their spherical coefficients are
\[
 \Omega_a(h)=\ell_a(\omega_V(h)e_V)
              =B(R(h)f_a,f_{a'}),\qquad
 \Omega_{a'}(h)=B(f_a,R(h)f_{a'}).
\]
The invariant pairing proves
\(\Omega_a(h)=\Omega_{a'}(h^{-1})\).

For completeness the last inverse may be removed using the
specified hyperspecial group, without changing that group.
Integral hyperbolic elimination gives
\[
 h=k_1\operatorname{diag}
   (\varpi^{\lambda_1},\ldots,\varpi^{\lambda_n},
    \varpi^{-\lambda_1},\ldots,\varpi^{-\lambda_n})k_2,
 \quad k_1,k_2\in K.
\]
One obtains this by diagonalizing the relative position of
the two selfdual lattices \(\Lambda_W,h\Lambda_W\):
their elementary divisors occur in opposite pairs.
At each step a primitive extremal vector and a dual partner
are carried to a hyperbolic pair by integral isometries,
and the orthogonal complement gives the next step.
The integral Weyl element exchanging the two vectors of
every hyperbolic pair conjugates the displayed diagonal
to its inverse. Thus \(h^{-1}\in KhK\).
The two coefficients are bi-\(K\)-invariant, and consequently
\begin{equation}\label{eq:sr-inert-common-spherical-function}
                    \Omega_a(h)=\Omega_{a'}(h)=:\Omega(h),
                       \qquad\Omega(1)=1.
\end{equation}

\subsection{Constructing an actual irreducible common image}
We use two precise published local facts. For unramified quadratic
unitary dual pairs with selfdual lattices in odd residue
characteristic, the joint hyperspecial fixed space is generated
from the lattice vacuum by either spherical Hecke algebra;
an irreducible spherical quotient receives a nonzero vacuum
and has a one-dimensional hyperspecial fixed line.
These are Howe's theorems as stated in Harris,
\emph{Cohomological automorphic forms on unitary groups, II},
Appendix~2.A, Theorems~2.A.1--2.A.2.
\footnote{\url{https://webusers.imj-prg.fr/~michael.harris/periodrelations.pdf}}
For a nonzero big theta lift of an irreducible admissible
representation, finite length and Howe duality give a nonzero
irreducible quotient. We use Gan--Takeda,
\emph{A proof of the Howe duality conjecture},
Theorem~1.2 and its introductory finite-length assertion.
\footnote{\url{https://arxiv.org/abs/1407.1995}}
Their field, lattice, residue, character and splitting hypotheses
were checked above; no even-rank restriction is being added.

Here is the exact quotient argument needed from those facts.
Let
\[
 Q_a=\mathcal S_V/
    \langle\omega_V(g,1)\phi-\phi:g\in U(V),\ \phi\in\mathcal S_V\rangle .
\]
Evaluation at zero proves \(Q_a\ne0\).
Moreover
\begin{equation}\label{eq:sr-inert-coinvariant-fixed-line}
                            Q_a^K=\mathbb C[e_V].
\end{equation}
Indeed compact averaging is exact on smooth complex
representations. A \(K\)-fixed vector of \(Q_a\) therefore
has a \(K\)-fixed lift; average this lift over \(K_V\) as well.
The second averaging does not change its image, because the
whole first member acts trivially on \(Q_a\).
The resulting joint fixed vector is a Hecke translate of
\(e_V\) under the first spherical Hecke algebra.
On first-member coinvariants a Hecke function acts by
its integral, so its image is a scalar multiple of \([e_V]\).
This proves \eqref{eq:sr-inert-coinvariant-fixed-line}.

Let \(q_a:Q_a\twoheadrightarrow\pi_a\) be the nonzero
irreducible theta quotient. Its vacuum is nonzero by the cited
unramified theorem. Equation
\eqref{eq:sr-inert-coinvariant-fixed-line} and exactness of
\(K\)-averaging give \((\ker q_a)^K=0\).
The \(K\)-invariant functional \(\ell_a\) annihilates this
kernel: on a kernel vector its value is unchanged by
\(K\)-averaging, and that average is zero.
Thus \(\ell_a\) descends to a nonzero functional on \(\pi_a\).
The coefficient map
\[
 C_a(\phi)(h)=\ell_a(\omega_V(h)\phi)
\]
therefore factors through \(\pi_a\). Its induced map is
nonzero, since \(C_a(e_V)(1)=1\), so is injective by
irreducibility. Its image is an irreducible representation
and is generated by \(C_a(e_V)=\Omega\), since every nonzero
vector in an irreducible representation is cyclic.
The same reasoning applies to \(a'\).

It follows from \eqref{eq:sr-inert-common-spherical-function}
that both coefficient maps have the same actual image
\[
            \Pi_v=
              \operatorname{span}_{\mathbb C}
                 \{h\longmapsto\Omega(hh_0):h_0\in U(W)\}.
\]
Set \(p_v=C_a\) and \(p'_v=C_{a'}\), with codomain \(\Pi_v\).
They are surjective and annihilate their respective
first-member augmentations. For every \(h_0\in U(W)\),
\[
 C_a(\omega_V(h_0)\phi)(h)=C_a(\phi)(hh_0),
\]
and the same equality holds for \(C_{a'}\).
Thus they induce the same second-member action, not merely
isomorphic abstract spaces. Their vacuum images are the
same nonzero function \(\Omega\). Finally every equivariant
endomorphism of \(\Pi_v\) is scalar: it preserves the
one-dimensional \(K\)-fixed line, so acts by a scalar on
\(\Omega\), and hence on all its translates.

Transport \(p_v,p'_v\) back by the vacuum-normalized local
intertwiners above. This constructs exactly the inert
finite matching datum for the prescribed models, including
surjectivity, augmentation, common action, irreducibility,
scalar endomorphisms, and equality of the actual vacua.

The positive-rank statement covers every application here:
the packet inequalities give \(a'\ge2\) and \(n\ge1\).
If a zero Hermitian rank is included in an auxiliary formulation,
its Schwartz space is \(\mathbb C\), its Weil action and
normalized vacuum coefficient are trivial, and its common
image is the constant functions. The same compact pairing
with \(r_0=1\) gives the constant coefficient on its
complementary positive-rank side. If all ranks are zero
both maps are the identity of \(\mathbb C\).
No \(\mathrm{GL}_0\) stabilizer or positive-rank Howe assertion
is used for those trivial spaces.

\subsection{Scope and simultaneous choices}
Together with the split-unramified and ramified constructions,
this proves the finite matching clause at every finite place
of a fixed actual triple. At a place outside the ramification
set, normalize by the prescribed vacuum, as above.
At a ramified place use the fixed full-rank local normalization
already constructed. The resulting local maps may all be
chosen at once; each codomain is the specified common
function image, so no incompatible choice of quotient
isomorphisms is needed. Their tensor normalization is fixed
before choosing polynomial vectors or integral theta tests.
This assertion does not impose a uniform norm estimate on
these maps as the global triple varies.

The archimedean common Fock quotient, the global nonzero
comparison constant and the two codimension-two recursion
identities still have to be supplied to conclude the full
global matching and theta-integral assertions. The construction
above closes the inert finite clause itself; it does not
replace any of those remaining global assertions.

\section{Ordinary convergence on the large theta side}
\label{sec:sr-large-theta-convergence}

We prove the ordinary-convergence clause of
\(\srcid{ThetaIntegralsAgree.Decay}\).
Let \(E/F\), \(V\), \(W\), the additive character, splitting
characters and rational Schr\"odinger realization be those
of the theta packet.
Write
\[
       [F:\mathbb Q]=f,\qquad
       \dim_E V=a,\qquad \dim_E W=d=a+a'.
\]
The Hermitian signatures of \(V\) are
\((a-p,p)\), with the same \(p\) at every real place.
Let \(r\) be its Witt index over \(F\).
Then \(r\le p\).
The following argument proves convergence whenever
\begin{equation}\label{eq:sr-large-weil-range}
                         d<2(a-r).
\end{equation}
In the source,
\[
     (a-a')/2>2(p+q)+1
\]
implies \(a-a'>2p\ge2r\), hence
\eqref{eq:sr-large-weil-range}.

Fix an adelic Schwartz function \(\Phi\) in the actual
Weil model.  For every \(h\in U(W)(\mathbb A_F)\) we claim
\[
 \int_{U(V)(F)\backslash U(V)(\mathbb A_F)}
                 |\Theta_V(\Phi;g,h)|\,dg<\infty .
\]
The claim uses ordinary integration; no finite-place
regularizer is inserted.

\subsection{A mixed model and its lattice bound}

Choose a rational Witt decomposition
\[
                       V=U\oplus V_0\oplus U^*,
                      \qquad\dim_EU=r .
\]
In the mixed Schr\"odinger model for \(V\otimes_EW\),
the variables belonging to \(U\oplus U^*\) form
\(W(\mathbb A_E)^r\); the remaining variables are
those of a rational Schr\"odinger model for \(V_0\otimes_EW\).
Write their infinite coordinates as
\[
                  (u_1,\dots,u_r,\xi).
\]
Each \(u_i\) has real dimension \(2fd\).
The change from the initially specified rational model
uses rational linear transformations and partial Fourier
transforms.
It preserves the theta functional by rational reindexing
and Poisson summation with selfdual measures.
It preserves the Schwartz class as well.
The normalization of the Weil action fixes its scalar,
so that this is a statement about the specified theta
kernel.

Let \(A\) be the split torus acting by \(e^{t_i}\) on
the \(i\)-th line of \(U\) and by \(e^{-t_i}\) on the
paired line of \(U^*\), with the same real scalar at
each real place of \(F\).
The unitary Schr\"odinger action is
\begin{equation}\label{eq:sr-large-cartan-action}
 \omega(a_{\mathbf t})\phi(u_1,\ldots,u_r,\xi)
   =e^{-fd\sum_i t_i}
       \phi(e^{-t_1}u_1,\ldots,e^{-t_r}u_r,\xi).
\end{equation}
Indeed the real determinant on one \(u_i\)-block is
\(e^{2fd t_i}\), and unitarity supplies its inverse
square root.
The determinant character of the hyperbolic torus is
one.  Thus no splitting-character modulus changes
this exponent.

We record the elementary estimate used to sum these
functions over rational points.
If \(L\) is a fixed full lattice in a real space of
dimension \(D\), and \(M>D\), then
\begin{equation}\label{eq:sr-large-shifted-lattice}
       \sum_{z\in L}(1+\|e^{-t}(z+b)\|)^{-M}
                  \le C_{L,M}e^{Dt}
          \quad(t\ge0,\ b\in\mathbb R^D).
\end{equation}
The number of points in a translated ball of radius
\(R\) is bounded by \(C_L(1+R)^D\), uniformly in its
center: place disjoint balls of a fixed radius around
them and compare volumes.
Divide the sum into dyadic annuli of radii
\(e^t2^j\).  Their contributions are bounded by
\(C_Le^{Dt}2^{j(D-M)}\), a convergent geometric series.
The same estimate with \(t=0\) is uniform in \(b\).

Let \(N\) be the unipotent radical of a rational minimal
parabolic compatible with the chosen Witt flag.
Its mixed-model action has the form
\[
 (\omega(n)\phi)(u,\xi)
     =\psi(Q_n(u,\xi))\,
                   \phi(B_nu,\xi+L_nu),
\]
with unit-modulus character factor and a block unitriangular
\(B_n\).
The block order may be chosen so that its \(i\)-th coordinate
is \(u_i\) plus a linear expression in later coordinates.
This formula follows directly from the three types of
unipotent matrices: changes within \(U\) give \(B_n\),
maps between \(V_0\) and the hyperbolic pairs give
Heisenberg translations and modulations in \(\xi\),
and the central Hermitian block gives the quadratic
character.
Only the linear translation part of each Heisenberg
operator affects absolute values.

If \(n\) ranges over a compact subset of \(N_\infty\),
and the remaining right-hand factors range over compact
sets, all Schwartz seminorms needed below have common
bounds.
For the ordinary Weil representation this follows from
the continuity of linear substitutions and partial
Fourier transforms on the Schwartz space; compact sets
are covered by finitely many such coordinate charts.
At finite places a fixed compactly supported
Schwartz--Bruhat factor restricts rational points to
one fixed fractional lattice and has bounded absolute
value.
Enlarge that lattice to a product lattice in the
displayed rational coordinates.

Apply a product Schwartz bound, then sum first over
\(\xi\).
The \(t=0\) case of
\eqref{eq:sr-large-shifted-lattice} is uniform in its
shift \(L_nu\).
Sum next over \(u_1\), then \(u_2\), and so forth.
The triangular form makes each translation a fixed
shift while that block is being summed.
Equation \eqref{eq:sr-large-shifted-lattice} contributes
\(e^{2fd t_i}\) from the \(i\)-th block.
Together with the prefactor in
\eqref{eq:sr-large-cartan-action}, this proves
\begin{equation}\label{eq:sr-large-theta-bound}
             |\Theta_V(\Phi;n a_{\mathbf t}m k,1)|
                 \le C_\Phi e^{fd\sum_i t_i}
\end{equation}
on each such set, for \(t_i\ge0\).
This summation argument includes the zero rational
vector and all singular orbits.

\subsection{Integration on reduction sets}

We use the usual arithmetic reduction theorem in the
form of finitely many adelic classes and finitely many
Siegel sets for their arithmetic groups.
For a compatible minimal parabolic these sets have
compact unipotent and right-hand factors and split
coordinates
\[
                    t_1\ge t_2\ge\cdots\ge t_r\ge0 ,
\]
after harmless fixed shifts of the origin.
The anisotropic Levi factor has compact arithmetic
quotient.  The norm-one part of its restriction-of-scalars
torus is compact modulo units.
Thus the split coordinates may be taken as the common
\(F\)-scalings used in
\eqref{eq:sr-large-cartan-action}.
No estimate uniform in the field or the reduction
sets is asserted here; ordinary convergence requires
none.

For precision, the positive relative root spaces of
\(U(V)\) are
\[
 \begin{array}{c|c}
        e_i-e_j,\ e_i+e_j\ (i<j)&2\\
        2e_i&1\\
        e_i&2(a-2r).
 \end{array}
\]
The second column gives their dimensions over \(F\).
They can be read off from the off-diagonal
\(E\)-entries, the one-dimensional skew-Hermitian diagonal
entry, and the \(V_0\)-row entries in a Witt matrix.
Summing their weights gives
\[
                  2\rho(\mathbf t)
                      =2\sum_{i=1}^r(a-2i+1)t_i .
\]
Restriction of scalars multiplies this by \(f\).
Consequently Haar measure on the reduction set is
bounded by a constant times
\[
             e^{-2f\sum_i(a-2i+1)t_i}
                           \,dt_1\cdots dt_r
\]
times the finite measures of its compact factors.

Multiply this density by
\eqref{eq:sr-large-theta-bound}.
Put \(t_{r+1}=0\) and \(s_j=t_j-t_{j+1}\ge0\).
The exponent becomes
\[
 \begin{split}
 f\sum_{i=1}^r\{d-2(a-2i+1)\}t_i
        &=f\sum_{j=1}^r
             j(d-2a+2j)s_j .
 \end{split}
\]
Every coefficient is strictly negative by
\eqref{eq:sr-large-weil-range}.
The integral is bounded by the product
\[
                 C_\Phi\prod_{j=1}^r
                    \frac1{fj(2a-d-2j)}<\infty .
\]
When \(r=0\), the arithmetic quotient is compact and
the continuous theta function is integrable directly.
This proves convergence on all reduction sets and hence
on the quotient.

Finally fix any \(h\).
The vector \(\omega(1,h)\Phi\) is Schwartz, so the same
argument applies to it.
The two unitary-group actions commute, and its theta
function at \(g\) is the original kernel at \((g,h)\).
Automorphy under \(U(V)(F)\) follows from rational
reindexing and Poisson summation; the splitting
characters are Hecke characters, hence are trivial on
rational determinants.
Integrability is therefore independent of the
fundamental domain.
Any Haar measure differs by a positive scalar, which
does not affect finiteness.
We have proved exactly
\[
          \forall\Phi\in\mathcal S,\qquad
               \srcid{Model.Conv}(M,\omega,\Phi).
\]

This proves the large-side ordinary-convergence
assertion of the packet. The next subsections identify
the prescribed split bar test with the global rank cutoff
and establish the small-side clauses separately.



\subsection{Matching the specified bar test and the global coefficient cutoff}
\label{subsec:sr-exact-bar-cutoff}

We verify that the local test used above is precisely the
source's \(\mathrm{SoLoc}\), including its existential
choice of split intertwiner. Let \(v_1\) be the designated
split finite place, \(k=F_{v_1}\), and \(t:E\to k\) its
chosen embedding. Write
\[
 P=J_V(t),\qquad Q=J_W(t),\qquad
             \operatorname{coB}(z)=-Pz_{t\sigma}Q^{\mathsf t}.
\]
Both \(P,Q\) are invertible. Thus
\(\operatorname{rank}\operatorname{coB}(z)
 =\operatorname{rank}z_{t\sigma}\).
The local condition supplies a Schwartz bijection
\(\iota\) satisfying the split bar-model intertwining
laws, with both auxiliary scalar normalizations equal
to one, and such that \(\iota\phi\) vanishes off the
full-row-rank locus. No special choice of this
\(\iota\) will be made.

Put \(T=\operatorname{supp}(\iota\phi)\).
It is compact and full rank. Since \(2\ne0\) in \(k\),
\(2T\) is full rank as well. Choose a compact open
additive subgroup \(\Lambda\) of \(M_{n,d}(k)\) with
\(2T+\Lambda\) full rank, and let \(K\) be its annihilator
under \((Y,U)\mapsto\psi_v(\operatorname{tr}(YU^{\mathsf t}))\).
One may shrink \(\Lambda\) so that \(T\) meets only
finitely many of its cosets; compactness in the
discrete quotient gives this finiteness for every
compact open \(\Lambda\).
Nondegeneracy of the trace pairing and local Fourier
duality give \(K^\perp=\Lambda\).
The explicit scaled-integral-lattice construction in
Section~\ref{sec:theta-decay-proof} gives the same
annihilator if the additive character has nontrivial
conductor.

Partition \(\iota\phi\) by these cosets and pull the
pieces back by \(\iota^{-1}\). They are Schwartz and
sum to \(\phi\), since \(\iota\) is a linear bijection
on Schwartz functions. Fix a piece with image supported
in \(Y_0+\Lambda\). For \(U\in K\), the defining
bar-model multiplication law says that its local
Schr\"odinger translate by the split vector
\(w_U=\operatorname{spM}(-U,0)\) is
\[
       S(M.\operatorname{co}(w_U))\phi
          =\psi_v(\operatorname{tr}(Y_0U^{\mathsf t}))\phi.
\]
This is the \(\mathrm{IsLinBar.mult}\) law itself;
it does not require \(\iota\) to be an isometry.

The symplectic pairing in the specified rational model is
\[
 \langle z,w_U\rangle
                 =\operatorname{tr}(\operatorname{coB}(z)U^{\mathsf t}).
\]
For the reflected coefficient \(c_\Phi\), the elementary
Schr\"odinger commutator calculation therefore gives
\[
 c_\Phi(M.\operatorname{co}(z))
   =\psi_v\!\left(
        \operatorname{tr}((2Y_0-\operatorname{coB}(z))U^{\mathsf t})
             \right)c_\Phi(M.\operatorname{co}(z)).
\]
If \(z_{t\sigma}\) has rank less than \(n\), then
\(\operatorname{coB}(z)\notin2Y_0+\Lambda\).
An element \(U\in K\) makes the displayed scalar
different from one. The coefficient is consequently
zero. This proves the cutoff for each piece and every
choice of the other tensor factors, including the
archimedean Schwartz factor.

Here is why the cutoff transfers to the global rational
sum. The model's rational symplectic-coordinate map is
an injective \(F\)-linear map from \(M_{n,d}(E)\) into
\(F^{nd}\oplus F^{nd}\). Both spaces have dimension
\(2nd\) over \(F\), so it is a bijection.
The coefficient sum can therefore be indexed by
\(x\in M_{n,d}(E)\). For every \(g,h\), a matrix
\(x\) of deficient \(E\)-rank stays deficient after
extension by \(t\sigma\) and the invertible left and
right actions. The preceding local coefficient is
zero at \(g^{-1}x(h^{-1})^{\mathsf t}\).
The reflected theta sum equals the original theta sum:
reflection is the central \((-1,1)\) Weil operator
up to a unit scalar, and rational reindexing sends
\(x\) to \(-x\). The coefficient identity thus bounds
an absolute square. These are exactly the parity and
square identities proved in
Subsection~\ref{subsec:sr-theta-coefficient-proof}.
Taking square roots, adding the finitely many piece
bounds, and enlarging their compact supports proves
\eqref{eq:theta-lattice-majorant} with rational full
row rank for every element of the specified
\(\mathcal S^\circ\).
This matches the source's \(\mathrm{soCut}\) and
\(\mathrm{sqSo}\) constructions without imposing
a more restrictive local test.

\subsection{The rational-translate height bound}
\label{subsec:sr-rational-torus-height}
We also spell out the height assertion used when the
second variable is allowed to vary. Its rational
translate is essential. It cannot be inferred from
submultiplicativity after simply discarding the
rational translating matrix.

In a fixed rational Witt basis put
\(a_{\mathbf v}\) diagonal with weights
\(v_1,\ldots,v_y,-v_1,\ldots,-v_y,0,\ldots,0\)
at all infinite places and identity at finite places.
Write \(V=\sum_jv_j\), with
\(v_1\ge\cdots\ge v_y\ge0\).
For the multiplicative adelic matrix height
\[
 H_0(b)=\prod_w
        \max(1,\|b_w\|_w,\|b_w^{-1}\|_w)^{n_w},
\]
use the sup matrix norm at finite places and the
Euclidean operator norm at infinity; \(n_w\) is the
local degree in the product formula. At all but
finitely many places the local factor is one.
For every \(\delta\in\mathrm{GL}_d(E)\), if \(y>0\),
\begin{equation}\label{eq:sr-height-rational-diagonal}
                   H_0(\delta a_{\mathbf v})
                                      \ge e^{D v_1},
                      \qquad D=[E:\mathbb Q].
\end{equation}
Indeed choose a nonzero entry \(\alpha\) in the row
of \(\delta^{-1}\) corresponding to weight \(-v_1\).
In \((\delta a_{\mathbf v})^{-1}
 =a_{\mathbf v}^{-1}\delta^{-1}\), that entry is
\(e^{v_1}\alpha\) at every infinite place and \(\alpha\)
at every finite place. Each matrix-norm factor is at
least the absolute value of this entry. Multiplication
over all places and the product formula for
\(\alpha\ne0\) give \eqref{eq:sr-height-rational-diagonal}.
A fixed rational basis change contributes only a
fixed multiplicative constant, by submultiplicativity
applied on both sides.

The matrix height \(H\) in the Weil model and \(H_0\)
are polynomially comparable, with constants depending
on the fixed arithmetic model. To verify this directly,
expand the \(E\)-matrix entries in a fixed \(F\)-basis
and compare their component norms with the finitely
many embedding norms at each place.
Finite-dimensional norm comparison gives local
constants in each direction; outside the finitely
many places where this basis, its inverse or the
different is not integral, those constants are one.
Taking their product gives a finite global constant.
Replacing component maxima by their local-degree
products changes the comparison by at most a fixed
integer power. Apply the same calculation to inverse
matrices. Thus \(H_0\le C H^b\) for fixed \(C,b\),
and conversely after changing them.
Since \(V\le yv_1\), \eqref{eq:sr-height-rational-diagonal}
gives fixed \(C_1,b_1\) such that
\[
      e^V\le C_1 H(\delta a_{\mathbf v})^{b_1}
                                  \quad\hbox{for every rational }\delta.
\]
For \(y=0\) this holds with \(e^V=1\).
The upper bound \(H(a_{\mathbf v})\le C_2e^{b_2V}\)
follows from its finitely many eigenweights and the
same fixed-basis comparisons.

Now write a reduced second variable as
\(\gamma h=a_{\mathbf v}\kappa\), with \(\kappa\)
in the fixed compact reduction set.
Use the preceding lower bound with
\(\delta=\gamma^{-1}\).
Since \(\gamma^{-1}a_{\mathbf v}=h\kappa^{-1}\),
compact boundedness, inversion invariance and
submultiplicativity give
\begin{equation}\label{eq:sr-height-second-variable}
             e^V\le C_3H(h)^{b_1}.
\end{equation}
This supplies the exact height step in the
two-chamber proof of
\eqref{eq:theta-full-decay}, including all its
rational translates.



\subsection{Full-rank tests on the small theta side}
\label{subsec:sr-small-theta-ordinary-convergence}

The full-rank test condition also gives ordinary absolute
convergence when the first space is the smaller member.
For this assertion the second adelic variable is fixed;
a height estimate uniform in that variable is unnecessary.

Let \(V\) have \(E\)-dimension \(n\), let \(W\) have
\(E\)-dimension \(d\ge n\), and use the actual rational
Weil model and the full-rank subspace \(\mathcal S^\circ\)
at its distinguished split finite place. For every
\(\Phi\in\mathcal S^\circ\) and every fixed
\(h\in U(W)(\mathbb A_F)\),
\begin{equation}\label{eq:sr-small-theta-absolute}
 \int_{U(V)(F)\backslash U(V)(\mathbb A_F)}
             |\Theta(\Phi;g,h)|\,dg<\infty .
\end{equation}
No dimension inequality beyond \(n\le d\) is needed.

Here are the details of the cusp estimate. The coefficient
identity and compact-coset separation in
Subsection~\ref{subsec:sr-theta-coefficient-proof} give, for
each exponent \(L\), the full-row-rank majorant
\[
 |\Theta_{f\otimes\phi}(g,1)|^2
 \le p_{\phi,L}(f)^2
       \sum_{\substack{x\in M_{n,d}(E)\\ \operatorname{rank}x=n}}
       {\bf1}_{C_f}((g^{-1}x)_f)
                 (1+|g^{-1}x|_\infty)^{-L}.
\]
Its continuous seminorm and finite compact support are
chosen before \(f\) and \(g\). In particular the rank
condition concerns the coefficient matrix, after taking
the absolute square; it is not an assertion that the
original theta summation variable has full row rank.

Put \(D=[E:\mathbb Q]\), and let \(r\) be the rational
Witt index of \(V\). Reduction theory gives finitely
many rational Witt frames and representatives
\[
       \gamma g=a_{\mathbf u}\kappa,\qquad
       u_1\ge\cdots\ge u_r\ge0,\qquad \kappa\in\Omega ,
\]
with \(\Omega\) compact. The archimedean eigenweights
in such a frame are \(u_1,\ldots,u_r\),
\(-u_1,\ldots,-u_r\), and \(n-2r\) zeros;
the torus is the identity at every finite place.
Rational invariance reindexes the matrix sum.
The inverse matrices of \(\kappa\) have uniformly
bounded archimedean operator norms, and their finite
components enlarge \(C_f\) to one fixed compact set.
A fixed rational change to the Witt basis therefore
leaves only matrices
\(\xi\in M_{n,d}(m^{-1}\mathcal O_E)\)
for one integer \(m\ge1\). Up to fixed positive
constants their infinite norm is
\[
 s_{\mathbf u}(\xi)=
       \sum_{v,i,j}e^{-\lambda_i(\mathbf u)}
                                    |\tau_v(\xi_{ij})|.
\]
All constants in this paragraph can be taken
simultaneously over the finite frame list.

For \(L_0=(D+1)nd\), the one-coordinate lattice packing
argument gives
\begin{equation}\label{eq:sr-small-all-lattice}
 \sum_{\xi\in M_{n,d}(m^{-1}\mathcal O_E)}
                   (1+s_{\mathbf u}(\xi))^{-L_0}
                    \le C\exp\!\left(Dd\sum_i u_i\right).
\end{equation}
To see the exponent directly, use
\[
 (1+\sum_{i,j}A_{ij})^{-L_0}
            \le\prod_{i,j}(1+A_{ij})^{-(D+1)}
\]
with \(A_{ij}=e^{-\lambda_i}\sum_v|\tau_v(\xi_{ij})|\),
and enlarge to independent integral-ideal lattice entries.
Packing in real dimension \(D\) bounds each factor sum
by \(C_m e^{D\max(0,\lambda_i)}\).
There are \(d\) entries per row and the positive
eigenweights sum to \(\sum_i u_i\), giving
\eqref{eq:sr-small-all-lattice}.

If \(r>0\), every full-row-rank \(\xi\) has a nonzero
entry \(\alpha\) in its row with weight \(-u_1\).
The algebraic integer \(m\alpha\ne0\) has absolute norm
at least one. Consequently
\(\max_v|\tau_v(\alpha)|\ge m^{-1}\), and
\[
              1+s_{\mathbf u}(\xi)\ge m^{-1}e^{u_1}.
\]
Write \(U=\sum_i u_i\le r u_1\).
For every integer \(q\ge0\), take
\(M=Ddr+2qr\). Multiplying the preceding estimate to
the power \(M\), and then using
\eqref{eq:sr-small-all-lattice}, gives
\[
 \begin{split}
 \sum_{\operatorname{rank}\xi=n}
       (1+s_{\mathbf u}(\xi))^{-(L_0+M)}
 &\le m^M C e^{DdU-Mu_1}\\
 &\le m^M C e^{-2qU}.
 \end{split}
\]
The square majorant therefore proves
\[
           |\Theta_{f\otimes\phi}(g,1)|
                        \le C_{\phi,q}p_{\phi,q}(f)e^{-qU}
\]
on each reduction set. The finite constants can be
absorbed into a continuous seminorm. When \(r=0\),
the same conclusion with \(U=0\) follows from the
unrestricted lattice bound on the compact reduction set.

In particular \(q=0\) bounds the theta function uniformly
in \(g\). The rational unitary quotient has finite
Haar volume, so this bound proves its absolute
integrability. More generally, the upper torus height
bound gives \(H_Q(g)\le A e^{bU}\). Taking \(q\ge bR\)
gives rapid decay against \(H_Q(g)^R\), for every
integer \(R\ge0\), with the second variable still fixed.

Finally the two Weil actions commute. The action of
\((1,h)\) preserves \(\mathcal S^\circ\), by the
split-matrix rank stability proved in
Section~\ref{sec:sr-good-split-coinvariants}.
Replace \(\Phi\) by \(\omega(1,h)\Phi\) in the preceding
argument. Its archimedean action is continuous on
Schwartz space, and its finite action produces a finite
sum of the allowed pure tensors. This proves
\eqref{eq:sr-small-theta-absolute} for every fixed \(h\)
and every \(\Phi\in\mathcal S^\circ\).
The same bound is uniform for \(h\) in a compact set:
the Schwartz seminorms have bounded operator constants
there and the finite factors have common compact support
and a common local constancy level.

This establishes ordinary convergence on the small side
with its actual full-rank test condition.
The bounds here allow constants depending on the fixed
arithmetic model and on \(h\); they do not assert
a uniform metric or covolume estimate over the varying
family, nor the global equality of the two integrals.

\subsection{The complete ordinary-convergence and decay clause}
\label{subsec:sr-theta-decay-closed}
For every actual triple in the packet,
the large-side ordinary convergence proved above,
the just-matched full-rank coefficient cutoff, and
the two-chamber argument of
Section~\ref{sec:theta-decay-proof} give the following.
For every small-side \(\Phi'\in\mathcal S^\circ\)
and every integer \(R\ge0\), there are \(C\ge0,k\ge0\)
such that
\[
       |\Theta'(\Phi';g,h)|H_Q(g)^R\le C H(h)^k
                                      \quad\hbox{for all }g,h .
\]
The exponent in that chamber proof is chosen before
the finite test factor; all seminorms remain continuous.
Finite sums of pure tensors are handled by the
triangle inequality and by replacing their exponents
with the maximum, since \(H(h)\ge1\).
At \(R=0\), integration over the finite-volume first
quotient gives
\[
 \int |\Theta'(\Phi';g,h)|\,dg\le C' H(h)^k,\qquad
                  |I'(\Phi';h)|\le C'H(h)^k .
\]
The integrand is continuous, hence measurable.
These inequalities prove ordinary absolute convergence
and exactly the moderate-growth assertion for \(I'\).
They use its specified Haar measure; a fixed positive
normalizing scalar only changes \(C'\).

Together with the unrestricted large-side convergence,
these are all clauses of the source's \(\mathrm{Decay}\)
predicate. They hold for the actual bar-test subspace,
rather than a substitute subspace. They do not
establish \(\mathrm{OneConst}\), the two rank recursions,
or the covolume/metric bound; those assertions have
different quantifiers and content.

\section{The uniform cusp-weight mean on unitary period quotients}
\label{sec:sr-uniform-cusp-mean}

We prove the numerical mean estimate used by the covolume
construction.  The constants in this section are chosen before the
field, Hermitian space, lattice and arithmetic level.
The construction of an integral model with its tied period map is a
separate issue; no such model is inferred from the mean estimate.

Let \(E/F\) be a quadratic CM extension, put
\[
 [F:\mathbb Q]=f,\qquad [E:\mathbb Q]=d=2f,\qquad
 R=|\operatorname{disc}E|^{1/d},
\]
and fix an \(E\)-Hermitian space \(V\) of dimension \(N=p+q\),
of signature \((p,q)\) at every real place, where
\[
                   q\ge1,\qquad 3q+1\le N .
\]
Let \(\Lambda\subset V\) be a full \(\mathcal O_E\)-lattice.
We use the canonical homogeneous Hodge metric of a compatible
polarized Riemann-lattice realization.  Its alternating form \(B\) is integral on \(\Lambda\), its positive
Hermitian majorants have imaginary part \(B\), and
\[
 B(ax,y)=B(x,\overline a y)\quad(a\in\mathcal O_E).
\]
The \(\mathcal O_E\)-action is the one induced by the fixed
\(E\)-vector space.  In particular its operator norm, for every
compatible majorant, is at most
\(\max_\sigma|\sigma(a)|\).
Let \(\Gamma\) be any arithmetic subgroup of the integral
unitary group of this same lattice and form.
Write \(X_\Gamma=\Gamma\backslash G_\infty\), with invariant
probability measure, where
\(G_\infty=\prod_{\tau:F\hookrightarrow\mathbb R}U(p,q)\).
Passing to the quotient by a maximal compact subgroup makes no
difference to any function below.

For a majorant indexed by \(g\), let \(\|\cdot\|_g\) be its
Euclidean norm and
\[
                  \lambda(g)=
                     \min_{0\ne x\in\Lambda}\|x\|_g .
\]
Every minimum is positive and attained.  This follows from
discreteness of a full lattice and properness of a positive
Euclidean norm.  The minimum is a continuous function of \(g\):
on a relatively compact neighborhood only finitely many lattice
vectors can have norm below any fixed bound, and the minimum is
locally the minimum of those finitely many continuous functions.

For a nonzero saturated \(\mathcal O_E\)-submodule \(I\subset\Lambda\)
that is \(B\)-isotropic, let \(r=\operatorname{rank}_{\mathcal O_E}I\).
Its real rank is \(2fr\), and \(r\le q\).
Indeed isotropy on the \(\mathcal O_E\)-module implies
\(B(ax,y)=0\) for every \(a\in E\); nondegeneracy of the trace
pairing then gives isotropy for the underlying \(E\)-Hermitian
form.  Its dimension is at most the Witt index \(q\).

Use the normalized cusp covolume
\begin{equation}\label{eq:sr-cusp-covol-normalization}
                  n_g(I)=
                   \frac{\operatorname{covol}_g(I)}{R^{fr}},
 \qquad
 U(g)=\max\left(0,\sup_I-\log n_g(I)\right).
\end{equation}
This is the source's normalization, including its discriminant
factor.  Here is a direct check against its determinant definition.
Choose independent \(v_1,\dots,v_r\in I\) with
\(J=\sum_i\mathcal O_Ev_i\) of finite index in \(I\), and let
\(H_\sigma\) be the Gram matrix of their \(\sigma\)-eigenprojections.
If \(\sigma\) runs through one CM type, that definition is
\[
                     n_g(I)=
                   \frac{\prod_\sigma\det H_\sigma}{[I:J]} .
\]
For \(x,y\in I\), the imaginary part of their polarized inner
product is zero.  Applying this also to \(ax,by\) and using
nondegeneracy of the embedding matrix shows that the Gram matrices
for conjugate eigenprojections are complex conjugates.
They have the same positive determinant.
In a tower basis \(\omega_jv_i\), where \((\omega_j)\) is an
integral basis of \(E\), the real Gram matrix is obtained from the
block diagonal matrices \(H_\rho\), for all embeddings \(\rho\),
by the embedding matrix \((\rho(\omega_j))\).
Taking determinants therefore gives
\[
 \det\operatorname{Gram}_{\mathbb R}(J)
       =|\operatorname{disc}E|^r
                        \left(\prod_\sigma\det H_\sigma\right)^2 .
\]
Changing from \(J\) to \(I\) divides this determinant by
\([I:J]^2\).  Since
\(R^{2fr}=|\operatorname{disc}E|^r\), taking positive square roots
gives exactly \eqref{eq:sr-cusp-covol-normalization}.
There is no suppressed factor of two in this computation.
For example in \(E=\mathbb Q(i)\), take an isotropic real
unit vector whose two conjugate eigenprojections have squared
norm \(1/2\).  The real Gram determinant for its \(1,i\)
translates is \(1\), whereas the displayed expression is
\(4(1/2)^2=1\).  More generally the embedding determinant is
the determinant over \emph{all} complex embeddings; the two
conjugate projection blocks already contain these factors.
Thus the proof below concerns the actual determinant cusp weight
for the canonical tied Hodge metric.

\subsection{Short integral multipliers and isotropic vectors}

There are \(d\) independent algebraic integers
\(a_1,\dots,a_d\) satisfying
\begin{equation}\label{eq:sr-cusp-short-integers}
                 \max_\rho|\rho(a_i)|\le dR .
\end{equation}
For completeness use the Euclidean embedding whose squared norm is
\(\sum_\rho|\rho(x)|^2\).
The dual of the integer lattice is the conjugate embedding of the
codifferent.  If \(0\ne b\) belongs to the codifferent, ideal
containment gives
\(|N_{E/\mathbb Q}b|\ge|\operatorname{disc}E|^{-1}\).
The arithmetic--geometric mean inequality yields
\[
                  \sum_\rho|\rho(b)|^2\ge d/R^2 .
\]
The Euclidean transference inequality
\(\lambda_d(L)\lambda_1(L^*)\le d\)
therefore supplies \(d\) independent integral vectors of norm
at most \(\sqrt d\,R\), hence at most \(dR\).
We use only this standard transference inequality from
Banaszczyk, \emph{New bounds in some transference theorems in
the geometry of numbers}, Theorem~2.1,
\url{https://doi.org/10.1007/BF01445125}.
Coordinate norms are bounded by Euclidean norm, proving
\eqref{eq:sr-cusp-short-integers}.

Set
\[
                 v_0=2fR=dR,\qquad r_0=v_0^{-1/2},
 \qquad a=2fq,\qquad k=2f(N-1).
\]
Then \(k\ge3a\).
Any two lattice vectors of norm less than \(r_0\) are
\(E\)-orthogonal.  In fact
\[
 |B(a_i x,y)|\le
              \|a_i x\|_g\|y\|_g
          \le v_0\|x\|_g\|y\|_g<1 .
\]
The left side is the absolute value of an integer, so it is zero.
The \(a_i\) span \(E\) over \(\mathbb Q\); the same vanishing holds
for every element of \(E\), and the nondegenerate trace pairing
identifies this with \(E\)-Hermitian orthogonality.
Consequently all vectors of norm below \(t\le r_0\) lie in one
real space of dimension at most \(a\).
Disjoint balls of radius \(\lambda(g)/2\) in that space give
\begin{equation}\label{eq:sr-cusp-packing}
 \#\{0\ne x\in\Lambda:\|x\|_g<t\}
       \le(1+2t/\lambda(g))^a
       \le(3t/\lambda(g))^a
                   \quad(\lambda(g)<t\le r_0).
\end{equation}

A bounded divisibility condition will make the counting function
unfold into a finite sum of orbital integrals.
For \(0\ne x\in\Lambda\), define
\[
                \iota(x)=
          [\,\Lambda\cap Ex:\mathcal O_Ex\,].
\]
For a shortest vector \(x\), this index is at most
\[
                    C_0=(v_0\sqrt d)^d .
\]
To see this, the real \(d\)-dimensional lattice
\(L_x=\Lambda\cap Ex\) has minimum \(\lambda(g)\).
Minkowski's convex-body theorem gives
\(\operatorname{covol}(L_x)\ge
(\lambda(g)/\sqrt d)^d\); the weaker constant used here follows
by putting a cube of volume \(2^d\) inside the Euclidean ball
of radius \(\sqrt d\).
The subgroup generated by \(a_1x,\dots,a_dx\) has covolume
at most \((v_0\lambda(g))^d\), by Hadamard's determinant bound.
It is contained in \(\mathcal O_Ex\), so division of the two
covolumes proves the stated index bound.

Let \(\mathcal V\) consist of the nonzero \(E\)-isotropic
\(x\in\Lambda\) with \(\iota(x)\le C_0\), and define
\[
                    c_t(g)=
                  \#\{x\in\mathcal V:\|x\|_g<t\}.
\]
This is a nonnegative measurable function.
It is zero when \(t\le\lambda(g)\), is bounded by
\eqref{eq:sr-cusp-packing} when \(t\le r_0\), and is at least
one when \(\lambda(g)<t\le r_0\): a shortest vector then
belongs to \(\mathcal V\).

\subsection{The exact homogeneous mean}

We now prove
\begin{equation}\label{eq:sr-cusp-count-homogeneous}
                 \int_{X_\Gamma}c_t\,d\mu=A t^k
                       \quad(t>0)
\end{equation}
with some finite \(A\ge0\); no uniform bound on \(A\) is assumed.

The two reduction-theoretic background facts needed here are
finite covolume for arithmetic groups with no rational
characters, and finiteness of rational parabolic double cosets.
They are Borel--Harish-Chandra's finite-volume theorem and
Borel, \emph{Introduction aux groupes arithm\'etiques},
Theorem~15.5 and Proposition~15.6, applied by restriction of
scalars to \(\mathbb Q\).
We spell out their applications, since neither a uniform
covolume nor a quantitative cusp count is needed.

There are finitely many \(\Gamma\)-orbits of rational isotropic
\(E\)-lines: their stabilizers are the line parabolics.
On a fixed such line, \(I_\ell=\Lambda\cap\ell\) is a
fractional ideal times a fixed vector.
The condition \([I_\ell:\mathcal O_Ex]\le C_0\) leaves finitely
many possibilities for the principal submodule
\(\mathcal O_Ex\): its quotient is an integral ideal of bounded
norm after identifying the fixed ideal class.
Generators of each principal submodule form one
\(\mathcal O_E^\times\)-orbit.
A subgroup of finite index in these units occurs in the line
stabilizer inside \(\Gamma\).
Indeed choose a rational hyperbolic complement and write its
Levi action as
\(\operatorname{diag}(u,1,\overline u^{-1})\).
Units in a suitable congruence subgroup preserve the fixed
lattice.  Their intersection with \(\Gamma\) has finite index
by arithmetic commensurability with the integral stabilizer.
This argument does not require \(\Gamma\) itself to be a
congruence subgroup.
There are therefore only finitely many
\(\Gamma\)-orbits in \(\mathcal V\).

Fix an orbit representative \(x\).
Its real stabilizer has unipotent radical and reductive
quotient \(U(x^\perp/Ex)\).
The latter has anisotropic center over \(F\), and hence no
nontrivial \(F\)-rational characters; its arithmetic
quotient has finite volume.
The unipotent arithmetic quotient is compact.
Thus the stabilizer quotient occurring in the orbit
unfolding has finite volume.

At one infinite place the orbit of a nonzero isotropic
vector is the cone \(Q(z)=h(z,z)=0\), with its vertex deleted.
The unitary group is transitive on it, by extending the
vector to a hyperbolic pair.
Its invariant measure is a positive constant times
\(\delta(Q(z))\,dz\), equivalently surface measure divided
by \(\|\nabla Q\|\).
The group preserves Lebesgue measure and \(Q\), so this
measure is invariant.
Under \(z\mapsto tz\), \(dz\) contributes \(t^{2N}\)
and the delta factor contributes \(t^{-2}\).
It therefore scales by \(t^{2N-2}\).
The measure of its intersection with a bounded majorant
ball is finite: in polar coordinates its radial density
is \(r^{2N-3}dr\), integrable at zero, and the angular
section is compact and smooth.
Taking products over the \(f\) real places gives degree
\(2f(N-1)=k\).

Tonelli's theorem unfolds the nonnegative count of this
orbit into the stabilizer volume times its invariant
cone integral.
The preceding paragraph shows that the result is a finite
constant times \(t^k\).
Summing over the finitely many orbits, and dividing by
the finite positive volume of \(X_\Gamma\), proves
\eqref{eq:sr-cusp-count-homogeneous}.
If there are no rational isotropic lines, the count is
zero and the same formula holds with \(A=0\).

\subsection{A uniform logarithmic tail}

The unknown constant \(A\) disappears in a short argument.
Put \(Z=r_0/\lambda\) and \(X=A r_0^k\).
From the lower bound on \(c_t\) we have
\[
             \Pr(Z>u)\le\min(1,Xu^{-k})\quad(u\ge1).
\]
If \(X=0\), then \(Z\le1\) almost everywhere.
If \(0<X\le1\), integration of this tail gives
\(\mathbb E\log^+ Z\le X/k\le1/k\).
If \(X>1\), the packing upper bound gives
\[
 X=\mathbb E c_{r_0}
       \le3^a\mathbb E(Z^a{\bf1}_{Z>1})
       \le3^a\frac{k}{k-a}X^{a/k}.
\]
For the last inequality split the layer-cake integral
at \(u=X^{1/k}\).
Taking logarithms and integrating the same tail once
more gives
\[
 \mathbb E\log^+ Z
 \le\frac{\log X}{k}+\frac1k
 \le \frac{a\log3+\log(k/(k-a))}{k-a}+\frac1k
 \le3 .
\]
The last bound uses \(k\ge3a\), \(a\ge1\),
\(\log3<2\), and \(\log(3/2)<1\).
In particular all logarithmic quantities just integrated
are integrable.  We have proved, uniformly in the field,
lattice and level,
\begin{equation}\label{eq:sr-cusp-minimum-mean}
             \mathbb E\log^+(r_0/\lambda)\le3 .
\end{equation}

Minkowski applied to a cusp sublattice \(I\) of rank
\(s=2fr\) yields
\(\operatorname{covol}_g(I)\ge
(\lambda(g)/\sqrt s)^s\).
Using \eqref{eq:sr-cusp-covol-normalization} and \(r\le q\),
\[
 \begin{split}
 -\log n_g(I)
   &\le fr\{2\log(1/\lambda(g))+\log(2frR)\}\\
   &\le fq\{2\log^+(1/\lambda(g))+\log(2fqR)\}.
 \end{split}
\]
The latter expression is nonnegative and independent of \(I\).
It bounds the supremum defining \(U\), and proves that
this supremum is finite.
It is measurable, being a supremum over a countable
collection of submodules of a countable lattice.
Since \(r_0=v_0^{-1/2}\) and \(v_0\ge1\),
\[
          \log^+(1/\lambda)
             \le\log^+(r_0/\lambda)+\tfrac12\log v_0 .
\]
Consequently
\begin{equation}\label{eq:sr-cusp-mean-explicit}
         0\le\mathbb E U
           \le fq\{6+\log(2fR)+\log(2fqR)\}.
\end{equation}

For the principal hosts \(R\ge2\).
If \(B_0\ge4fN^2\), then \(B_0\ge64\) and \(2fq\le B_0\).
Thus both logarithms in
\eqref{eq:sr-cusp-mean-explicit} are bounded by
\(\log(B_0R)\), while \(\log(B_0R)\ge\log128>3\).
For any real \(L\) this proves the precise estimate
\begin{equation}\label{eq:sr-cusp-mean-four}
       \mathbb E U
           \le4fq\{\log(B_0R)+|L|\}.
\end{equation}
One can in particular take the source's \(L=\operatorname{hostLval}\).
The constant is exactly \(C_u=4\), chosen before the host.

\subsection{Passage to the tied Hodge integral}

We identify the measure used above with the geometric
integral in the covolume construction.
In a matrix ball chart \(Z\in M_{p,q}(\mathbb C)\), the
Hodge determinant metric has potential, up to a fixed
positive scale and an additive constant,
\[
                         -\log\det(1-Z^*Z).
\]
Its Hessian sends a tangent matrix \(T\) to the positive
form with left and right factors
\((1-ZZ^*)^{-1}\) and \((1-Z^*Z)^{-1}\).
The determinant of this Hessian is
\[
 \det(1-ZZ^*)^{-q}\det(1-Z^*Z)^{-p}
                   =\det(1-Z^*Z)^{-N}.
\]
Taking products over the CM type shows that the top
Hodge curvature measure is a fixed positive scalar
times the invariant unitary measure.
The scalar cancels in a normalized mean.

This also handles the source's finite collections of
period sheets and generic multiplicities.
The integral of \(U\) is a sum of invariant integrals
with positive sheet multiplicities; its denominator
is the same sum with \(U\) replaced by \(1\).
Dividing gives a positive weighted average of the
means already bounded in
\eqref{eq:sr-cusp-mean-four}.
Branch loci and toroidal boundary divisors have zero
top-dimensional measure.  The established integrability
of \(U\) permits their removal without changing the
integral.
A compatible isometry of Riemann lattices preserves
the submodules, Gram determinants, indices and
minimum appearing in
\eqref{eq:sr-cusp-covol-normalization}; thus the cusp
weight on every sheet is the same actual function
pulled back from its lattice fiber.

For any tied canonical Hodge realization of the
principal packet, with \(q=j'\), \(N=j+j'\) and
\(f=|\mathrm{CMType}|\), we therefore obtain
\[
 \frac{\displaystyle
          \int U\,c_1(\overline{\mathcal M})^{\dim X}}
      {\displaystyle
          \int c_1(\overline{\mathcal M})^{\dim X}}
       \le4f j'
          \{\log(B_0\operatorname{hostRd})
                    +|\operatorname{hostLval}|\}.
\]
This proves the analytic numerical assertion
\srcid{MeanU} for the canonical determinant cusp weight.
It supplies the uniform mean part of the covolume
producer.  It does not supply the remaining
\srcid{HtU} lower bound at every algebraic point, nor
the simultaneous integral root-tower and Hodge
comparison witnesses needed to build
\srcid{CovolumeBoundOfX}.

\section{A quadratic CM period witness with the required metric}
\label{sec:sr-quadratic-period-witness}

The flat determinant comparison uses the quadratic subfield
\[
                  K=\mathbb Q(\sqrt{-D}),
                  \qquad D=\operatorname{Host.d}.
\]
This is the coordinate selected by
\(\operatorname{Host.rt}^2=-D\).
It is not the other quadratic coordinate
\(\mathbb Q(\sqrt{-\operatorname{Host.q}})\).
Here \(D\ge7\) is squarefree and \(D\equiv3\pmod4\);
thus \(\operatorname{disc}K=-D\).
Write \(\chi=\chi_{-D}\) for the primitive odd quadratic
Dirichlet character.

We construct the elliptic period witness and compute its
normalization.  We use the first main theorem of elliptic
complex multiplication, in the form that CM \(j\)-invariants
are integral and that the Hilbert class field acts
transitively through the ideal classes.
A reference with these precise assertions is Milne,
\emph{Modular Functions and Modular Forms},
Corollary~12.16 and Theorem~12.18,
\url{https://www.jmilne.org/math/CourseNotes/MF110.pdf}.
We also use potential good reduction for an elliptic curve
with integral \(j\), as in Silverman,
\emph{The Arithmetic of Elliptic Curves}, VII.5.5,
\url{https://www.math.ens.psl.eu/~obenoist/refs/Silverman.pdf}.
The analytic mean formula is derived below.

\subsection{The integral curve and its tangent generator}

Begin with the complex elliptic curve
\(\mathbb C/\mathcal O_K\), with its canonical principal
polarization and its \(\mathcal O_K\)-action.
Complex multiplication supplies a number field of definition.
Enlarge it to contain \(K\) and to define the endomorphisms.
Its integral \(j\)-invariant implies potential good reduction
at every finite prime.  A single finite extension gives good
reduction at all primes: there are only finitely many bad
primes, one may realize the required local extensions by
approximation, and good reduction persists under further
base change.
We obtain an elliptic scheme \(\mathcal C/\mathcal O_M\)
with its canonical principal polarization.
Its generic-fiber endomorphisms extend to the scheme by the
N\'eron property of a good elliptic model.

We may further arrange that the Hodge line
\(e^*\Omega^1_{\mathcal C/\mathcal O_M}\) is free.
Indeed represent it by an invertible fractional ideal \(I\).
Choose \(h\ge1\) with \(I^h=(\alpha)\), and adjoin an
\(h\)-th root \(b\) of \(\alpha\).
Then \((I\mathcal O_{M'})^h=(b)^h\).
The fractional-ideal group of the Dedekind ring
\(\mathcal O_{M'}\) has no torsion, so
\(I\mathcal O_{M'}=(b)\).
Taking this base change and, if desired, its normal closure
makes the Hodge line free without losing good reduction.

Choose a basis \(\omega\) of the Hodge line, and let \(v\)
be its dual integral tangent generator.
Use the \(\mathcal O_K\)-action for which
\[
               w=(1+\sqrt{-D})/2
\]
acts at the distinguished complex embedding by that scalar.
The same scalar identity holds on the one-dimensional
generic Lie algebra and on its integral lattice, since
both maps are integral and agree generically.
In particular
\[
                     w^2-w+(1+D)/4=0
\]
is the derivative relation required in the period witness.
The Rosati involution is conjugation on \(\mathcal O_K\).
We also choose the representing line bundle to satisfy
the source's primitive-pairing identity, not only equality
of polarizations.  On rigidified line bundles let
\(T(L)\) be the cubical pairing pulled back by \(1,w\).
The theorem of the cube makes \(T\) additive and makes
it kill \(\operatorname{Pic}^0\).
On the N\'eron--Severi group its action is multiplication
by \(w+\overline w=1\).
Therefore \(T^2=T\): the quotient \(T(L)/L\) is in
\(\operatorname{Pic}^0\), where \(T\) is zero.
Starting with the symmetric degree-one bundle of the
zero section, choose its image under \(T\).
It is again symmetric of degree one and represents the
same principal polarization, and now
\(T(L)\simeq L\) as rigidified line bundles.
The identity extends over the good model; the relative
Picard scheme is separated and the rigidification
removes a line pulled back from the base.
This is precisely \srcid{Ell.prim}.
The metric on the tangent space depends on the
polarization, so this choice does not change its
normalization.
At the distinguished embedding the period lattice is
homothetic to \(\mathcal O_K=\mathbb Z+\mathbb Zw\),
because the chosen complex curve was
\(\mathbb C/\mathcal O_K\).
No assertion that every Galois conjugate is in the principal
ideal class is needed.

For an embedding \(\sigma:M\hookrightarrow\mathbb C\), write
the period lattice for \(\omega\) as
\[
                   \Lambda_\sigma=
                     \Omega_\sigma(\mathbb Z+\tau_\sigma\mathbb Z),
       \qquad y_\sigma=\operatorname{Im}\tau_\sigma>0 .
\]
The principal polarization has Hermitian form
\(|z|^2/\operatorname{area}(\Lambda_\sigma)\).
The source defines the tangent norm by multiplying this
form by \(\pi\).
Since \(\omega(v)=1\), its squared tangent norm is therefore
\begin{equation}\label{eq:sr-cm-tangent-norm}
                  n_\sigma=
                   \frac{\pi}{|\Omega_\sigma|^2y_\sigma}.
\end{equation}
This fixes the normalization directly, without choosing
among conventions for Faltings metrics.

The canonical discriminant section is a nowhere-zero
section of the twelfth Hodge power on a good elliptic scheme.
Relative to \(\omega^{12}\) its coefficient
\(\Delta_\omega\) is an algebraic unit.
The analytic discriminant identity is
\[
             \sigma(\Delta_\omega)
               =(2\pi)^{12}\Omega_\sigma^{-12}
                                      \eta(\tau_\sigma)^{24}.
\]
It follows either from the usual Weierstrass
\(g_2^3-27g_3^2\) identity, with invariant differential
\(dz\), or by comparing its weight-twelve Fourier
expansion at the cusp.
The product formula for the unit gives
\[
 \frac1{[M:\mathbb Q]}\sum_\sigma\log|\Omega_\sigma|
       =\log(2\pi)+
          \frac2{[M:\mathbb Q]}\sum_\sigma
                                      \log|\eta(\tau_\sigma)|.
\]
Combining this with
\eqref{eq:sr-cm-tangent-norm} yields
\begin{equation}\label{eq:sr-cm-eta-mean}
 -\frac1{[M:\mathbb Q]}\sum_\sigma\log n_\sigma
       =\frac1{[M:\mathbb Q]}\sum_\sigma
                  \log\!\left(4\pi y_\sigma
                                      |\eta(\tau_\sigma)|^4\right).
\end{equation}
CM reciprocity distributes the ideal classes uniformly
among these embeddings.
Passing to a finite extension repeats each embedding the
same number of times.  Complex conjugation preserves
\(y|\eta(\tau)|^4\).
Consequently the right side is the uniform average over
the \(h_K\) ideal classes, with \(\tau_\mathfrak a\)
any corresponding upper-half-plane class point.

\subsection{The mean from Poisson summation}

Put
\[
 \mathcal E(\tau,s)
    =\sum_{(m,n)\ne(0,0)}
                  \frac{y^s}{|m+n\tau|^{2s}}
                  \qquad(\operatorname{Re}s>1).
\]
Mellin inversion of the Gaussian and one-dimensional
Poisson summation give
\[
 \begin{split}
 \mathcal E(\tau,s)
   ={}&2\zeta(2s)y^s
       +2\sqrt\pi\,
          \frac{\Gamma(s-\tfrac12)}{\Gamma(s)}
                         \zeta(2s-1)y^{1-s}\\
     &+\frac{8\pi^s\sqrt y}{\Gamma(s)}
        \sum_{m,n\ge1}(m/n)^{s-1/2}
           K_{s-1/2}(2\pi mn y)\cos(2\pi mn\,\operatorname{Re}\tau).
 \end{split}
\]
To obtain it, separate \(n=0\), write each remaining
power as its Gaussian Mellin integral, and Poisson-sum
the \(m\)-variable.
The zero Fourier mode is the gamma quotient in the first
line, and the two opposite nonzero modes combine into
the cosine in the second.
For fixed \(y>0\), the second line converges normally
near \(s=1\), by exponential decay in \(mn y\).

At \(s=1\) use
\(K_{1/2}(t)=\sqrt{\pi/(2t)}e^{-t}\),
\(\Gamma'(\tfrac12)/\Gamma(\tfrac12)
-\Gamma'(1)/\Gamma(1)=-2\log2\),
and
\[
 \log|\eta(\tau)|
     =-\pi y/12-
         \sum_{m,n\ge1}\frac{e^{-2\pi mn y}}n
                       \cos(2\pi mn\,\operatorname{Re}\tau).
\]
The term \(2\zeta(2)y=\pi^2y/3\) cancels the
linear \(y\)-term from the eta product.
The resulting Laurent expansion is
\begin{equation}\label{eq:sr-cm-limit-formula}
 \mathcal E(\tau,s)
       =\frac{\pi}{s-1}
           +2\pi\{\gamma-\log2-
                         \log(\sqrt y|\eta(\tau)|^2)\}
           +O(s-1).
\end{equation}
All identities here use ordinary Lebesgue measure and
the Gaussian whose Fourier transform is itself.

Let \(u_K=|\mathcal O_K^\times|\).
Writing a fractional ideal as
\(\Omega(\mathbb Z+\tau\mathbb Z)\) and using
\(\operatorname{area}(\mathfrak a)
   =N(\mathfrak a)\sqrt D/2\)
shows, by grouping nonzero elements modulo units, that
\[
 \zeta_K(s)=\frac1{u_K}
           \left(\frac2{\sqrt D}\right)^s
                         \sum_{[\mathfrak a]}
                                     \mathcal E(\tau_\mathfrak a,s).
\]
This also fixes the factor \(2/\sqrt D\).
On the other hand \(\zeta_K(s)=\zeta(s)L(s,\chi)\).
Comparing residues in
\eqref{eq:sr-cm-limit-formula} gives
\[
                     L(1,\chi)=
                        \frac{2\pi h_K}{u_K\sqrt D}>0 .
\]
Comparing constant terms divided by the residue gives
\[
 \gamma+\frac{L'(1,\chi)}{L(1,\chi)}
    =2\gamma-\log2-\tfrac12\log D
       -\frac1{h_K}\sum_{[\mathfrak a]}
                   \log\!\left(y_\mathfrak a
                                  |\eta(\tau_\mathfrak a)|^4\right).
\]
Adding \(\log(4\pi)\) and using
\eqref{eq:sr-cm-eta-mean} proves the exact period value
\begin{equation}\label{eq:sr-cm-period-exact}
 -\frac1{[M:\mathbb Q]}\sum_\sigma\log n_\sigma
       =x_D
       :=\gamma+\log(2\pi/\sqrt D)
                            -\frac{L'(1,\chi)}{L(1,\chi)} .
\end{equation}
The value is real, since \(\chi\) is real and
\(L(1,\chi)>0\).

Thus the constructed integral elliptic curve with
tangent generator satisfies precisely
\(\srcid{Ell.HasPeriod}(x_D)\).
Conjugating its \(\mathcal O_K\)-action gives the witness
for the other root of \(-D\), with the same underlying
polarized curve and the same mean norm.
The construction can be based-changed to contain any
fixed number field needed by the host; the normalized
mean in \eqref{eq:sr-cm-period-exact} is preserved.

\subsection{The host discriminant correction}

For the larger CM host \(E\), the source defines
\[
 \operatorname{logOm}(E,\chi)
    =\gamma+\log(2\pi/\operatorname{rd}E)
                            -\frac{L'(1,\chi)}{L(1,\chi)} .
\]
Consequently the precise relation between its parameter
and the quadratic elliptic period is
\begin{equation}\label{eq:sr-cm-period-host}
 \operatorname{Re}\operatorname{logOm}(E,\chi)
       =x_D-\log\operatorname{rd}E+\log\sqrt D .
\end{equation}
Replacing \(\sqrt D\) by \(\operatorname{rd}E\) inside
the elliptic formula would lose this correction.

There is an equally concrete finite-place correction
when passing from the raw flat-frame norm to the
normalized determinant ratio.
For an embedding \(E\hookrightarrow M\), let
\[
 D_v=\log\|\mathfrak D_{E/\mathbb Q}\|_v
\]
use the source's ideal-log convention, in which the
normalized sum over finite places and embeddings is
\(-\log\operatorname{rd}E\).
The product formula gives the normalized finite sum of
\(\log|D|_v\) as \(-\log D\).
Thus the correction
\[
                       q(D_v-\log|D|_v)
\]
has normalized finite sum
\(q(\log D-\log\operatorname{rd}E)\).
Adding it to the finite log norm subtracts this quantity
from its Arakelov degree.
In particular the raw degree
\[
 \frac{p-q}{2}\operatorname{Re}\operatorname{logOm}(E,\chi)
          +q(\log D-\log\operatorname{rd}E)
\]
becomes exactly
\[
                 \frac{p-q}{2}
                       \operatorname{Re}\operatorname{logOm}(E,\chi).
\]
This is the correction in
\srcid{KroneckerCharacterPeriodFormula.degree\_frames}.
The calculation checks its sign and normalizes the
period witnesses just constructed.
It does not replace the separate assertion that a
given global flat frame has these integral local norms,
nor does it construct the root tower required for
\srcid{CovolumeBoundOfX}.

\section{Consolidation of the current project sources}
\label{sec:sr-source-review}
This section records the source review behind the present consolidation.
Its purpose is to distinguish a proved internal deduction, an imported
formal theorem and an assumption still awaiting a prose proof.
The review is of the files available on 21 November 2026.
Version identifiers and verified byte hashes are recorded in the
accompanying \srcid{short-route-sources.json}.

\subsection{The paper consolidation and historical review}
The
\href{https://principia-math.com/projects/d16ab288ed353bc5b10e292ace0ddbc1/paper}{project papers page}
has one current RH working manuscript from this account.
The earlier entry-point paper and the two earlier versions of the
translation were consolidated into this manuscript and their separate
attachments removed. Their source snapshots and metadata are preserved in
\href{https://github.com/ChristianKudeba/alpoge-rh-verification/blob/main/translation/archive/principia-before-194-page-consolidation.json}{the historical archive}.
Our separate Masullo audit remains a separate paper because it addresses
a different proof attempt. Other contributors' papers were left unchanged.
Before this consolidation the following five PDFs were identified and
reviewed; the older artifact links below are historical references.
\begin{description}
\item[Formal entry points.]
\href{https://principia-math.com/api/ws/files/402cfc03e76bc4f9eafd0eb930a80a20}{\srcid{alpoge-rh-formal-entry-points.pdf}}
records endpoint declarations and their formal dependencies.
Its six pages are not a six-page mathematical proof.
\item[Earlier working translation.]
\href{https://principia-math.com/api/ws/files/259cf87468cc745b1e81e358086a251e}{\srcid{alpoge-rh-working-translation.pdf}}
is an earlier version of our exposition, rather than an independent
argument that can discharge its successor's assumptions.
\item[The 163-page working translation.]
\href{https://principia-math.com/api/ws/files/8ca901b46a918a9293164d6a63d1ea7e}{\srcid{alpoge-rh-working-translation-163pp.pdf}}
is the supporting text incorporated in Part II.
Proposition~\ref{prop:sr-PB} now connects its printed minus descent to the
short route; no proof term is needed to make that connection.
\item[The separate analytic audit.]
\href{https://principia-math.com/api/ws/files/b561b940f8dd3d3b8f8a2d1d83b359dc}{\srcid{masullo-zero-count-audit.pdf}}
contains counterexamples and checks on the second proposed approach.
Section~\ref{sec:sr-shared-distribution} includes the checks relevant here.
\item[Masullo's draft.]
\href{https://principia-math.com/api/ws/files/e8e3c2bb97f9bf324d74bff487bfb435}{\emph{The Music of the Primes}}
is the 940-page version already examined in that audit.
Its prime-side normalization is comparable with
\eqref{eq:sr-prime-distribution}; its global spectral identification
cannot be used here without replacing the affected steps.
\end{description}
This review found no additional standalone proof closing the stated
inputs among these five files. The bytes of every listed PDF were hashed for this review.
In particular the Masullo version has not silently been replaced by
a revision resolving the identified issues.

\subsection{Anton's conductor-one reductions}
The first new specialization examined here was in the collaborative editor:
\href{https://principia-math.com/projects/d16ab288ed353bc5b10e292ace0ddbc1/editor?file=b64c07f1e23480429eacd6b354f04f26}{\srcid{ConductorOneRHCombined.lean}}.
The earlier 552-line snapshot has SHA-256
\[
 \begin{gathered}
 \texttt{e2bab4de455929689f2977514ae8a84b8}\\
 \texttt{004fbbbb19619cf7b6866d77c93d33e}.
 \end{gathered}
\]
The two displayed strings concatenate to one 64-digit hash.
For that snapshot, the author reports a successful Lean 4.33.1
compilation in 664 seconds.
The present consolidation independently checked the source bytes and
read the theorem signatures; it did not independently repeat that
compilation. The reported 33 direct imports reach 3,857 project
modules, besides Mathlib. The preceding 1,873-line version reached
3,856. Thus the shorter source is a useful specialization and
reorganization; its imported mathematical development remains extensive.

The top theorem is
\srcid{Principia.Ordinary.ConductorOneRH.rh}.
It is an implication with three explicit inputs.
The first is the oddness of the class number of every cyclotomic field
of \(2\)-power conductor; this is proved in
Lemma~\ref{lem:sr-oneprime} and its application above.
The other two inputs are as follows.

\paragraph{A uniform plus-length bound.}
In the source this is the predicate \srcid{Height}.
It requires a single natural exponent \(C\) and two fixed auxiliary
primes, before the host and the twist range are chosen.
For every sufficiently large principal host \(h\), every
\(t\geq2\) with \(h.\mathrm{Big}(t)\), and its CM field \(E\), it gives
\begin{equation}\label{eq:sr-principal-height}
 \sum_{\substack{\ell\in S\\\ell\ \mathrm{odd}}}\log\ell
       \sum_{a=1}^{t-1}s^+_{\ell,E}(a)
                       \leq(Nt)^C
       \quad\hbox{for every finite prime set }S .
\end{equation}
The host is for the trivial character. Its field and \(N\) are those
in Section~\ref{sec:sr-arithmetic}.
The size condition includes \(2t+1<q,d\) and excludes finitely many
fixed character-dependent constants. The two auxiliary-prime
inequalities and these fixed exclusions must hold for the same host.
Equation~\eqref{eq:sr-principal-height} is not a consequence of the minus
bound, since it controls the complementary conjugation eigenspaces.
Allowing \(C\) to depend on the individual host would change the assertion.

\paragraph{A construction of connected families.}
The final theorem also assumes the implication
\begin{quote}\raggedright
\srcid{HostExW (3 / 2) (GChi PChar.one).req}
\(\Longrightarrow\)
\srcid{FamilyClassConnected ((GChi PChar.one).cls (3 / 2))}.
\end{quote}
These are the archive's host-window and connected-family predicates.
They include geometric witnesses and their uniform estimates.
A prose deduction using this route must define those witnesses and prove
their properties. Naming the predicates does not do that.
Definition~\ref{def:sr-supply} instead states the exact signed-data
existence assertion used by the fully printed comparison in Part I.
The two interfaces should not be identified without constructing the
map between their witnesses.

The 552-line snapshot proves the following assembly within the imported
development. Its degree-one specialization removes the exceptional
odd-prime term. The minus estimate, the supplied dyadic estimate and
\eqref{eq:sr-principal-height} give the full strict-length budget.
The high-twist size theorem and the rank bound provide an exceptional
set with at most 24 elements. Finite-exception avoidance supplies a
common interval. This produces a principal cohomological window.
The assumed connected-family construction then gives the prime-sum
inequality, which excludes right-half-plane zeros. Modulus-one
Dirichlet \(L\) is \(\zeta\), so reflection gives ordinary RH.
The degree-one cancellation and the final identification of \(L\)
remove bookkeeping required for nontrivial characters.

Both combined revisions use
\[
 J_{\mathrm{archive}}=2^{26}(2C_KC_t/\epsilon)^{24}.
\]
The direct interval proof in \eqref{eq:sr-J} gives
\(J_{\mathrm{direct}}=2^{25}(2C_KC_t/\epsilon)^{24}\) under its
stated exceptional-set hypotheses.
These are different interval estimates. The sharper numerical
calculation does not itself establish those hypotheses.

\paragraph{Subsequent additions.}
The
\href{https://github.com/antoshashakov/riemann-hypothesis/blob/d85b449fa08d46836c934a7809fe3c9cae2e4a6a/verification/LogCubeRH.lean}{630-line cubic-logarithmic revision}
has the same direct import list and adds a weaker height option:
\eqref{eq:sr-logcube-height}.
Its pinned GitHub bytes match the updated Principia editor body,
with SHA-256
\[
 \begin{gathered}
 \texttt{d183191cac31584b7fa051847a87fbce3}\\
 \texttt{dac7ec4a44a9f5701fe4ad2f9904321}.
 \end{gathered}
\]
The author's report states that Lean checked both RH endpoints and
ten axiom reports. We checked the source bytes and its implications,
without repeating that Lean build.
Section~\ref{sec:sr-logcube} proves the new numerical absorption.

The further
\href{https://github.com/antoshashakov/riemann-hypothesis/blob/442bb98b2aa223ac4f28ec86e754e1a54b3ae081/verification/PrincipalHeightCounts.lean}{principal height counts},
\href{https://github.com/antoshashakov/riemann-hypothesis/blob/b8afa4ee2aca3cb3478b1f706cac6b03b258b9fe/verification/RegularCharacterHeight.lean}{regular-character vanishing}
and
\href{https://github.com/antoshashakov/riemann-hypothesis/blob/76ffabc62ab67b126afc128013ab6efa6f7d2a7c/verification/TwoCharacterCounts.lean}{two-character reduction}
have 342, 282 and 356 lines respectively.
Their uniform generic sum and nontrivial pointwise estimate remain
explicit inputs.
Section~\ref{sec:sr-two-counts} writes these inputs in mathematical
notation, proves the needed trivial-character vanishing and derives
the height bound from them with all prime weights retained.
This removes a separate CM-host assumption at 3 and 7.
It does not remove the need to prove the two rational estimates.

\subsection{New certificate constructions}
The newer repository work at
\href{https://github.com/ChristianKudeba/alpoge-rh-verification/tree/6186b8d5c272c08c90f6f5cce1a25d22232ed78d/verification/research/certificate-construction}{the pinned certificate-construction snapshot}
supplies explicit integer witnesses before the field realizations and
coefficient primes are chosen.
Its odd-minus certificates, Weber proof and dyadic-window adapters
assemble proved archive inputs. Its nonresonant plus certificate
uses \(\mathbb Q(\sqrt{-3})\), whose elementary arithmetic packet
is proved in Section~\ref{sec:sr-three-auxiliary}.
A pair of distinct auxiliary primes supplies two nonzero integers;
their squared product covers the one excluded prime of either
individual estimate. Products over a fixed set of rational character
orbits then give one integer independent of the later prime set.

The snapshot's notes explicitly retain the substantial
\srcid{PropGK} integer producer and the nontrivial resonant estimate
as inputs. They also identify \(\ell=3,a=1\) as a case that a
rank-four estimate stated only for \(a\geq2\) cannot cover.
Thus these constructions improve the arithmetic assembly without
providing the missing standalone counting proof.
The recorded kernel replay reports acceptance of the closed
odd-minus certificate, Weber theorem and the implication from a joint
window to the final cohomological window, using only the three standard
axioms. We inspected those records and the relevant sources; we did not
repeat their replay in this document build.
The joint-window producer still has mathematical content to translate.

\subsection{What imported results still need an internal proof}
The short file supplies the high-twist theorem by assembling the
class-counting ingredients. Its source lines 182--189 invoke
\srcid{lemmaK_of_steps} together with chain, prism, intersection,
theta, class, torsion and quaternionic bounds.
The explicit Poitou--Tate and local duality parameters are therefore
not its only mathematical content: the other ingredients are
imported proved results. The standalone manuscript still needs their
applicable statements and proofs, or a valid replacement for them.

The connected-family branch we inspected is
\href{https://github.com/ChristianKudeba/alpoge-rh-verification/blob/alpoge-source-7cbf7a0c/src/Development/AutomorphicArithmetic/ThetaCorrespondence/ThetaPairLifts/ThirdPlaceSeries/TopShape.lean}{\srcid{ThirdPlaceSeries.TopShape}}.
Its concluding comparison sets \(e_D=2,e_L=0\), turning the
exponent \(1/2+(e_D+e_L+7)\epsilon\) into \(1/2+9\epsilon\).
That substitution is elementary.
The hypotheses supplying the integral theta realizations, rank and
determinant bounds, genus mass, metric control, admissible twists and
connected geometric model are not elementary consequences of that
substitution. The longer route's supporting chapters establish several
comparisons from specified local data, while retaining the construction
of those data as inputs.

The community's generated
\href{https://principia-math.com/projects/d16ab288ed353bc5b10e292ace0ddbc1/editor?file=fd6fc09854184b61bca8dd800a4859b5}{\srcid{rh-unpacked.tex}}
earlier snapshot reports 1,067 translated modules and 1,001 reviewed modules
out of 175,187. We inspected relevant geometric and cohomological
entries in that snapshot. Subsequent project notes at 3:15 AM report
1,155 translated modules and 1,071 reviewed modules.
For example
\srcid{CuspidalModerateGrowthRapidDecay} separates the height used in
derivative growth from the quotient height used in its conclusion.
Its proof uses a cusp-scale decomposition, a unipotent-variation bound
and an iterated-gain theorem.
The following \srcid{CuspFormsRapidDecay} entry is a wrapper converting
that finite-derivative estimate into rapid decay.
These entries help track the proof, but appending their statements
would not prove the imported cusp-scale results or the uniform metric
estimates needed for the present families.
Likewise a published-hypothesis entry states an input; it is not an
independent proof of the input.

\subsection{The resulting proof ledger}
The odd prime-to-degree minus bound is now an internal consequence:
the strict characterwise descent, its derivative system and local
identities, the common Bernoulli integer and
Proposition~\ref{prop:sr-PB} form one chain.
The dyadic class-number base, the uniform 24-generator bound of
Proposition~\ref{prop:sr-dyadic-generators}, the elementary field degree,
the finite interval selection and numerical absorption are also proved here.
Section~\ref{sec:sr-dyadic-bounds} now proves a uniform free-lift
congruence, cokernel growth from depth seven, the matrix-logarithm
center bound, the 24-element exceptional set fixed on a whole horizon,
and the fixed-depth discriminant estimate. It also computes all local
cohomology corrections exactly and proves the resulting growth bound
of 25 per layer. Section~\ref{sec:sr-ideal-lattices} proves the
\srcid{SStein} fractional-ideal comparison and its covolume identity.
These are internal ingredients of the class-counting chain; its
remaining global estimates still require proofs.
The dyadic common-window
theorem is an internal deduction from \(K_2(C_K)\), not a proof of
that remaining high-twist input.
The analytic endpoint and signed finite comparison have complete
printed proofs. The new rational vanishing arguments and the assembly
from the two character estimates are printed in
Section~\ref{sec:sr-two-counts}.

To obtain an unconditional standalone proof by this route, the
following work remains.
\begin{enumerate}
\item Prove the uniform plus estimate
\eqref{eq:sr-principal-height}, or the weaker bound
\eqref{eq:sr-logcube-height}, with one constant before all hosts and
twist ranges. The two estimates in
Definition~\ref{def:sr-character-inputs} suffice for the polynomial
version, by Proposition~\ref{prop:sr-two-counts}. The height/deformation and automorphic counting
comparisons in Part II must be supplied with the construction and
compatibility inputs they currently assume.
\item Prove \(K_2(C_K)\) of Definition~\ref{def:sr-K2} with one
constant uniform in the principal pair and the twist.
The growth, center, exceptional-set and discriminant arguments of
Section~\ref{sec:sr-dyadic-bounds} then give the full dyadic
common-window bound by Theorem~\ref{thm:sr-dyadic-window}.
The remaining content is the quantitative high-twist
class-counting theorem, rather than the construction of the exceptional
set or its interval selection.
\item Construct the connected geometric families, or directly the
signed data of Definition~\ref{def:sr-supply}, on the required common
windows. This includes the relevant global theta identity, ramified
matching, cusp/metric and integral-lattice estimates with compatible
normalizations and constants.
\end{enumerate}

\clearpage
\part{Supporting translations from the longer route}
\paragraph{The source.}
This translation follows the Lean development released on 20 November 2026.
The classical endpoint in \srcid{Main.lean} is the theorem
\srcid{riemannHypothesis}, of type \srcid{RiemannHypothesis} in Mathlib.
Alpöge's announcement describes the origin of the argument. Andrew Sutherland
reports that the independent kernels nanoda and con-leche accepted the 
full export. The five endpoints use propositional extensionality, classical
choice and quotient soundness. Section~\ref{sec:kernel-reading} spells out
the named endpoint and the meaning of a proof-term reference.
These results concern the formal development. The proofs printed below can be
assessed separately by tracking the inputs used in each section.
\footnote{\href{https://alpo.ge/remark.pdf}{Alpöge's announcement};
\href{https://math.mit.edu/~drew/grh/}{Sutherland's verification report and
export}; \href{https://github.com/Deicyde/riemann-hypothesis/blob/40e56bab026552e089d7afcfeff08d2defb9c005/src/Main.lean}{the final formal declarations}.}
The arithmetic family must control characters of varying conductors. The
modulus-one specialization occurs in the final deduction for the zeta
function. We begin with that deduction so that the conclusion supplied by the
later estimates is explicit.

% A manuscript fragment. Requires amsmath, amssymb and hyperref.
\section{From primitive characters to the Riemann hypothesis}
\label{sec:character-to-rh}

Let \(N\) be a positive integer. A complex Dirichlet character of level
\(N\) is a character of \((\mathbb Z/N\mathbb Z)^\times\), extended by
zero to residue classes that are not units. For \(\Re s>1\) write
\[
       L(s,\chi)=\sum_{n=1}^{\infty}\frac{\chi(n)}{n^s}.
\]
Elsewhere \(L(s,\chi)\) denotes its analytic continuation at the points
where that continuation is regular. Write \(\mathbf 1_N\) for the
principal character of level \(N\). The all-character assertion used
in the last step of the proof is
\begin{equation}
\label{eq:all-character-nv}
   \Re s>\tfrac12,\qquad
   (\chi\ne\mathbf 1_N\ \text{or}\ s\ne 1)
       \quad\Longrightarrow\quad L(s,\chi)\ne0
\end{equation}
for every positive \(N\) and every character \(\chi\) of level \(N\).
The exception in the principal case removes the pole. Keeping this
disjunction also permits evaluation of a nonprincipal character at
\(s=1\). The primitive nonvanishing assertion in the source is phrased
using a total complex-valued function, including at exceptional points.
Its restriction to the regular points above suffices for the
classical argument.\footnote{\href{https://github.com/Deicyde/riemann-hypothesis/blob/40e56bab/src/Development/AutomorphicArithmetic/CharacterBoxes/EveryCharacter.lean\#L13-L18}{Definition of the all-character assertion and its deduction from the primitive assertion},
\href{https://github.com/Deicyde/riemann-hypothesis/blob/40e56bab/src/Formalism/SharpRestriction/RamificationFrobenius/TotientGaloisClasses/ExplicitFormulaErrorOfHostConditions.lean\#L11-L18}{the primitive assertion}.}

\subsection{Restoring the Euler factors at the level}
\label{sec:restore-euler}

Assume the right-half-plane nonvanishing produced by the primitive-character
argument, and fix \(N,\chi,s\) satisfying the hypotheses in
\eqref{eq:all-character-nv}. Put \(f=\operatorname{cond}(\chi)\).
Then \(f\) divides \(N\), and there is a primitive character \(\chi^*\)
of level \(f\) inducing \(\chi\). If \(s=1\) we have
\(\chi\ne\mathbf 1_N\), hence \(\chi^*\ne\mathbf 1_f\), since raising
the level of a principal character yields the principal character.
Thus the primitive character is evaluated at a regular point in all
cases under consideration. The primitive nonvanishing statement gives
\(L(s,\chi^*)\ne0\).

The change-of-level identity in the normalization used here is
\begin{equation}
\label{eq:dirichlet-change-level}
     L(s,\chi)=L(s,\chi^*)\prod_{p\mid N}
                     \bigl(1-\chi^*(p)\,p^{-s}\bigr),
\end{equation}
where \(p\) ranges over distinct rational primes. For \(p\mid f\) the
factor is \(1\), since \(\chi^*(p)=0\). For \(\Re s>1\) the identity
follows by deleting the Euler factors at primes dividing \(N\) from
the series for \(\chi^*\). Holomorphy gives the identity on
\(\mathbb C\setminus\{1\}\) by analytic continuation. If \(s=1\) is
allowed, both characters are nonprincipal and continuity gives the same
identity there.\footnote{\href{https://github.com/leanprover-community/mathlib4/blob/v4.33.1/Mathlib/NumberTheory/LSeries/DirichletContinuation.lean\#L125-L162}{Change of level and extension to \(s=1\).}}

For every \(p\mid N\), the value of \(\chi^*(p)\) is zero or a root of
unity, so \(\lvert\chi^*(p)\rvert\le1\). Also
\[
       \lvert\chi^*(p)\,p^{-s}\rvert
       \le p^{-\Re s}<1.
\]
The strict inequality follows from \(p>1\) and \(\Re s>\tfrac12\).
Consequently no factor in \eqref{eq:dirichlet-change-level} vanishes.
Multiplication by the nonzero value \(L(s,\chi^*)\) proves
\eqref{eq:all-character-nv}.\footnote{\href{https://github.com/Deicyde/riemann-hypothesis/blob/40e56bab/src/Development/AutomorphicArithmetic/CharacterBoxes/Imprimitive.lean\#L5-L35}{Reduction from arbitrary characters to primitive characters.}}

\subsection{The character of level one}
\label{sec:level-one-specialization}

The group \((\mathbb Z/\mathbb Z)^\times\) is trivial. Its unique
character takes the value \(1\) on every positive integer. Its
\(L\)-function therefore agrees with \(\zeta(s)\) for \(\Re s>1\),
and their continuations agree away from the pole. Specializing
\eqref{eq:all-character-nv} to \(N=1\) and \(\chi=\mathbf 1_1\)
gives
\begin{equation}
\label{eq:zeta-right-half}
             \zeta(s)\ne0
             \qquad(\Re s>\tfrac12,\ s\ne1).
\end{equation}
Both restrictions in \eqref{eq:zeta-right-half} are exactly those used in
the reflection argument.\footnote{\href{https://github.com/Deicyde/riemann-hypothesis/blob/40e56bab/src/Formalism/GammaFactorNonvanishing/CritLine.lean\#L29-L32}{Specialization to the character modulo one.}}

\subsection{Reflection, including the possible zeros of the cosine}
\label{sec:zeta-reflection}

\noindent\textbf{Proposition.}
Suppose \eqref{eq:zeta-right-half} holds. For every \(s\in\mathbb C\),
if \(\zeta(s)=0\), \(s\ne1\), and \(s\ne-2(n+1)\) for each
\(n\in\mathbb Z_{\ge0}\), then \(\Re s=\tfrac12\).

\medskip
\noindent\textit{Proof.}
Fix such a zero \(s\). The possibility \(\Re s>\tfrac12\) is excluded
by \eqref{eq:zeta-right-half}. Suppose then that \(\Re s<\tfrac12\),
and set \(w=1-s\). Since \(\zeta(0)=-\tfrac12\), we have \(s\ne0\),
and hence \(w\ne1\). We also have \(\Re w>\tfrac12\), which puts
\(w\) outside \(\{0,-1,-2,\dots\}\). These statements justify using
the zeta functional equation at \(w\), which gives
\begin{equation}
\label{eq:zeta-cos-reflection}
    \zeta(1-w)=2(2\pi)^{-w}\Gamma(w)
                       \cos\!\left(\frac{\pi w}{2}\right)\zeta(w).
\end{equation}
The complex power \((2\pi)^{-w}\) is nonzero because \(2\pi>0\).
The gamma factor is nonzero because \(\Re w>0\). Equation
\eqref{eq:zeta-right-half} at \(w\) gives \(\zeta(w)\ne0\).

It remains to test the cosine. If it vanished, the complex zeros of
cosine would give \(w=2k+1\) for some \(k\in\mathbb Z\), and
therefore \(s=-2k\). If \(k<0\), then \(s\ge2\), contradicting
\(\Re s<\tfrac12\). If \(k=0\), then \(s=0\), already excluded.
If \(k>0\), take \(n=k-1\) to contradict \(s\ne-2(n+1)\).
Thus every factor on the right of \eqref{eq:zeta-cos-reflection} is
nonzero, whereas the left side is \(\zeta(s)=0\). This excludes
\(\Re s<\tfrac12\) and proves the proposition.
\hfill\(\square\)

This proposition has precisely the quantifiers used for the classical
Riemann hypothesis in the released theorem. Its nonvanishing premise
follows from the all-character result through
\eqref{eq:zeta-right-half}.\footnote{\href{https://github.com/Deicyde/riemann-hypothesis/blob/40e56bab/src/Formalism/GammaFactorNonvanishing/RhOfRight.lean\#L5-L62}{The reflection argument},
\href{https://github.com/leanprover-community/mathlib4/blob/v4.33.1/Mathlib/NumberTheory/LSeries/RiemannZeta.lean\#L175-L186}{the zeta functional equation and the formal statement of RH},
\href{https://github.com/Deicyde/riemann-hypothesis/blob/40e56bab/src/Supplements/Kernels/MainTheorems.lean\#L7-L9}{application of the all-character result},
\href{https://github.com/Deicyde/riemann-hypothesis/blob/40e56bab/src/Main.lean\#L4-L6}{the two classical RH declarations}.}

\section{How the arithmetic input supplies the character theorem}
\label{sec:entry-assembly}

The all-character assertion \eqref{eq:all-character-nv} is obtained by
applying a theorem with two inputs. One input includes the existence of
local Langlands parameters in the normalization described below, together
with two uses of a local--global compatibility statement. The other
input supplies arithmetic windows for the auxiliary fields. The two
uses of local--global compatibility are instantiated by exactly the same
theorem. Although its first input lists three components, the proof
has to establish two different kinds of assertion to fill them.
\footnote{\href{https://github.com/Deicyde/riemann-hypothesis/blob/40e56bab/src/Development/AutomorphicArithmetic/ThetaCorrespondence/ThetaPairLifts/ThirdPlaceSeries/EveryCharacterOfLiteratureOfVWin.lean\#L210-L218}{The three components of the input},
\href{https://github.com/Deicyde/riemann-hypothesis/blob/40e56bab/src/Development/AutomorphicArithmetic/CharacterBoxes/MainConclusion.lean\#L6-L12}{their instantiation in the final argument}.}

\subsection{The local parameter used in the comparison}

Fix a number field \(E\), an algebraic closure \(\Omega\), a finite
place \(w\) of \(E\), and a place of \(\Omega\) above \(w\). Put
\(K=E_w\). For each \(n\ge1\), consider irreducible smooth complex
representations of \(\mathrm{GL}_n(K)\) up to isomorphism. The local
construction supplies a parameter \(r(\pi)\) of Weil--Deligne type,
together with a local Artin map \(a\) from the Weil group to \(K^\times\).
On characters it gives \(r(\chi)=\chi\circ a\) with zero monodromy. It
sends cuspidal representations to irreducible parameters and sends a
Langlands quotient for a multisegment to the direct sum of the parameters
of its segments. It preserves dimension and intertwines real powers of
the norm. For every \(\pi\) with central character \(\omega_\pi\), it
also satisfies
\[
   |\det r(\pi)(\sigma)|=|\omega_\pi(a(\sigma))|
\]
on the Weil group, with \(\det\) taken on the Weil action. Several
stronger assertions in this construction have a restriction on the
cuspidal support. An allowed cuspidal constituent has rank one, or is a
smooth-character twist of a representation with first principal
congruence invariants over a field topologically isomorphic to
\(\mathbb Q_p\). If \(\pi\) has a cuspidal support all of whose members
are allowed, \(r\) respects smooth character twists, the central
character equality holds without absolute values, and duals correspond.
For such a \(\pi\), equality \(r(\pi')=r(\pi)\) implies \(\pi'=\pi\).
If in addition \(\pi\) has first principal congruence invariants, its
parameter is tame.

More explicitly, suppose \(K\simeq\mathbb Q_g\) and choose an
extension place \(w'\) unramified over \(g\) with residue field
\(\mathbb F_{g^n}\). Suppose the restriction of the local representation
to the integral subgroup contains the inflation of a cuspidal type
of \(\mathrm{GL}_n(\mathbb F_g)\) with regular parameter
\(\theta:\mathbb F_{g^n}^{\times}\to\mathbb C^\times\).
The character of this type at a matrix with irreducible characteristic
polynomial and an eigenvalue \(z\in\mathbb F_{g^n}^{\times}\) is
\((-1)^{n-1}\sum_{j=0}^{n-1}\theta(z^{g^j})\). The monodromy of
\(r(\pi)\) is zero. If \(t\) belongs to inertia at \(w'\) and
\(z\) is the residue of its local Artin image there, regard \(t\)
as an element of the inertia group over \(E\). Its characteristic
polynomial under \(r(\pi)\) is
\[
       \prod_{j=0}^{n-1}\bigl(X-\theta(z^{g^j})\bigr).
\]
\footnote{\href{https://github.com/Deicyde/riemann-hypothesis/blob/40e56bab/src/Development/AutomorphicArithmetic/Foundations/LocalLanglands/IsRecTOfIsRec.lean\#L14-L80}{The restrictions on cuspidal support, the inertia polynomial and the defining properties of the parameter.}}

Here is the extension step used in constructing \(r\).
Choose the Artin map, and first construct parameters \(r_{\rm c}\)
on cuspidal representations. Multisegment exhaustion and the extension
lemma give a map \(r\) on all irreducible representations agreeing with
\(r_{\rm c}\) on cuspidal ones. Suppose \(\pi\) is represented by a
multisegment and let \(S\) be the multiset of its cuspidal terms,
including their norm twists. For \(\sigma\) in the Weil group let
\(P_{D,\sigma}(X)\) denote the characteristic polynomial of the Weil
action on \(D\). The segment calculation gives
\begin{equation}
\label{eq:entry-segment-polynomial}
      P_{r(\pi),\sigma}(X)
      =\prod_{\tau\in S}P_{r_{\rm c}(\tau),\sigma}(X).
\end{equation}
Indeed for a segment of length \(d\) with first term \(\rho\), the Weil
eigenvalues of the segment parameter collect the eigenvalues of the
successive norm twists of \(r_{\rm c}(\rho)\) at exponents
\(0,1,\ldots,d-1\). Norm-twist compatibility on cuspidals identifies
these with the parameters of the terms of the segment. Direct summation
multiplies characteristic polynomials, giving
\eqref{eq:entry-segment-polynomial} for the union of the segments.
At \(\sigma=1\), comparison of degrees in this identity gives
\(\dim r(\pi)=\sum_{\tau\in S}\dim r_{\rm c}(\tau)=n\). For general
\(\sigma\), comparison of products of eigenvalues gives
\(\det r(\pi)(\sigma)=\prod_{\tau\in S}\det r_{\rm c}(\tau)(\sigma)\).
The central character of \(\pi\) is the product of the central
characters of its cuspidal support. Taking absolute values therefore
gives the stated central character identity for every \(\pi\). If
all its constituents are allowed, use the exact central character
identity on each factor to obtain the exact identity for \(\pi\).
\footnote{\href{https://github.com/Deicyde/riemann-hypothesis/blob/40e56bab/src/Development/AutomorphicArithmetic/Foundations/LocalLanglands/LocalCentralUnits/ExtDim.lean\#L19-L78}{The segment identity, its degree and the two determinant conclusions.}}

For a one-dimensional representation, the cuspidal construction gives
the asserted formula from reciprocity directly. Each segment starts at
a cuspidal representation, so the multisegment extension rule with
\(r_{\rm c}\) also gives the required rule with \(r\) in its place.
The extension lemmas for twists, duals, injectivity and depth give
the corresponding properties above. For the inertia assertion the
representation containing the inflated finite cuspidal type is itself
cuspidal. Hence \(r\) agrees with \(r_{\rm c}\) on that representation,
and the inertia calculation for \(r_{\rm c}\) applies. This supplies
the local existence assertion at the chosen place, and the choices of
\(E,\Omega,w\) and a place over \(w\) were arbitrary.\footnote{\href{https://github.com/Deicyde/riemann-hypothesis/blob/40e56bab/src/Development/AutomorphicArithmetic/Foundations/LocalLanglands/LocalCentralUnits/RecTExists.lean\#L20-L66}{Construction of the local parameter and verification of its required laws.}}

\subsection{The global compatibility statement}

Let \(E\) be a CM field with complex conjugation \(c\), and let \(\ell\)
be a rational prime. Fix an algebraic closure \(\Omega\) of \(E\) and
a field isomorphism \(\iota:\overline{\mathbb Q}_{\ell}\to\mathbb C\).
For \(a\ge1\), take a regular algebraic cuspidal representation \(\pi\)
of \(\mathrm{GL}_a\) over \(E\), together with a Hecke character of
its maximal totally real subfield polarizing \(\pi\). The global
statement provides a continuous semisimple representation
\[
  R:\operatorname{Gal}(\Omega/E)\longrightarrow
       \mathrm{GL}_a(\overline{\mathbb Q}_\ell),
\]
an integer \(b\), and a multiset \(H_\tau\) of integers for every
embedding \(\tau:E\to\overline{\mathbb Q}_{\ell}\). If \(I_{\iota\tau}\)
denotes the multiset of infinitesimal-character parameters for \(\pi\)
at \(\iota\tau\), the equalities of multisets are
\[
  H_\tau=\{(a-1)/2-z:z\in I_{\iota\tau}\},
  \qquad H_{\tau c}=\{b-h:h\in H_\tau\}.
\]
The first equality uses the natural inclusion \(\mathbb Z\subset\mathbb C\).

At a finite place \(w\nmid\ell\), let \(x\) be the Weil--Deligne
parameter of \(R\) and \(y\) the local parameter of the algebraically
normalized component of \(\pi\), using the chosen \(\iota\) to compare
coefficients. Their dimensions agree. If all Weil commutators act
trivially in \(y\) and its monodromy \(N_y\) vanishes, then \(x=y\).
There is a second comparison when \(E_w\simeq\mathbb Q_g\) for a
prime \(g\) and either \(a=1\) or \(16a^2<g\). Suppose \(y\) is tame,
\(N_y=0\), and for each inertia element generating tame inertia
topologically modulo wild inertia every repeated eigenvalue \(z\) of
that element in \(y\) satisfies \(z^{\operatorname{N}w-1}=1\). Then
\(N_x=0\) and the characteristic polynomials of \(x\) and \(y\)
agree on every Weil element of degree zero.

At \(w\mid\ell\), with the compatible local period data of the
comparison theorem, a spherical local component gives a crystalline,
de Rham representation with jumps \(H_\tau\) at the embeddings it
induces. If \(N_y=0\) and either the commutator condition or the tame,
repeated-eigenvalue and prime-size conditions above hold, the representation
is de Rham and potentially semistable, with the same prescribed jumps.
Its potentially semistable Weil--Deligne parameter then satisfies the
two comparisons just stated for \(x\) and \(y\). These clauses specify
the local conclusions used in both occurrences of the global input.
\footnote{\href{https://github.com/Deicyde/riemann-hypothesis/blob/40e56bab/src/Development/AutomorphicArithmetic/Foundations/GaloisRepresentations/TypeConditions/WeightsConjBis.lean\#L15-L86}{The global comparison, local hypotheses and weights},
\href{https://github.com/Deicyde/riemann-hypothesis/blob/40e56bab/src/Development/AutomorphicArithmetic/Foundations/GaloisRepresentations/TypeConditions/WeightsConjAlt.lean\#L22-L33}{the commutator and inertia conditions}.}

\subsection{The arithmetic window and its quantifiers}

For a primitive Dirichlet character \(\psi\), \(0<\epsilon<1/52\) and
\(U>1\), the pair of auxiliary primes \(q,d\) satisfies
\begin{align*}
 q&\equiv3\pmod8,& \tfrac12 U^{\epsilon/5}
                   &<\log q\le U^{\epsilon/5},\\
 d&\equiv7\pmod{16},&
 \tfrac12 U^{\epsilon/5}(\log U)^2
                   &<\log d\le U^{\epsilon/5}(\log U)^2 .
\end{align*}
In addition \(q\mid d+1\) and every odd prime \(p\le U^{\epsilon/5}\)
divides \(d+1\). Choose any function \(e:\mathbb R\to\mathbb Z_{\ge0}\)
such that for all \(x\ge1\),
\[
       x^{2+8\epsilon}\le 2^{e(x)}\le 2x^{2+8\epsilon}.
\]
Such a function exists. For \(x\ge1\), take \(n\ge0\) with
\(2^n\le x^{2+8\epsilon}<2^{n+1}\), and put \(e(x)=n+1\).
For \(x<1\) put \(e(x)=0\).

For fixed \(e,q,d\), the number fields \(E\) in the assertion can be
described by a subgroup of a cyclotomic Galois group. There must be
a positive integer \(M\) divisible by \(2^{e+2}\), the conductor of
\(\psi\), \(q\), and \(d\), and a \(\mathbb Q\)-embedding
\(E\hookrightarrow\mathbb Q(\zeta_M)\). Its image is the fixed field
of the automorphisms \(\zeta_M\mapsto\zeta_M^u\) satisfying
\[
  u\equiv\pm1\bmod 2^{e+2},\qquad \psi(u)=1,\qquad
  \Bigl(\frac uq\Bigr)=\Bigl(\frac ud\Bigr)=1.
\]

Let \(S_\ell\) consist of all places of \(E\) dividing \(\ell\),
with a choice of a place of an algebraic closure above each one.
Use the Galois group \(G_{E,S_\ell}\) unramified outside \(S_\ell\)
at finite places, and the discrete module
\(A_{\ell,k}=(\mathbb Q_\ell/\mathbb Z_\ell)(-k)\), on which Galois
acts by the inverse \(k\)-th power of the cyclotomic character.
Let \(D_v\) be the decomposition group at the chosen place over \(v\).
Put
\[
   T_{\ell,k}(E)=\ker\!\left[
      H^1(G_{E,S_\ell},A_{\ell,k})\longrightarrow
           \prod_{v\in S_\ell}H^1(D_v,A_{\ell,k})\right],
   \qquad b_{\ell,k}=\operatorname{length}_{\mathbb Z_\ell} T_{\ell,k}(E).
\]
All cohomology here uses continuous cochains.
\footnote{\href{https://github.com/Deicyde/riemann-hypothesis/blob/40e56bab/src/Formalism/GaloisCohomology/Statement.lean\#L48-L78}{The places, cyclotomic character, sign of the twist and strict cohomology lengths}.}

For this \(\psi,\epsilon,e(\cdot)\), the window assertion supplies
\(J\ge1\) and \(X_0\) such that for every \(X\ge X_0\) and every
pair \(q,d\) satisfying the conditions above at \(U=X^{1/J}\),
one can choose \(Y\ge X^{1/J}\) with \(Y^2\le X\). For \emph{every}
\(Y\le x\le Y^2\), take \(t=\lfloor x^{401}\rfloor\). For every field
\(E\) defined above at \(e=e(x)\), and every choice of algebraic closure
and local data, all \(b_{\ell,k}\) with \(1\le k<t\) are finite.
Simultaneously for every finite set \(P\) of rational primes,
\begin{equation}\label{eq:entry-strict-length}
      \sum_{\ell\in P}(\log\ell)\sum_{k=1}^{t-1}b_{\ell,k}
                 \ \le\ \sqrt{\exp(x^\epsilon)} .
\end{equation}
In particular \(Y\) is chosen before \(x,E,P\). The scale for
choosing the auxiliary primes is \(X^{1/J}\), which may be smaller
than \(Y\). The source formulates the bound over all finite \(P\),
so its use never requires a prior finiteness assertion about the set
of contributing primes.
\footnote{\href{https://github.com/Deicyde/riemann-hypothesis/blob/40e56bab/src/Development/AutomorphicArithmetic/CharacterBoxes/ExistsPageN.lean\#L12-L49}{Definitions of the field, the prime pair and the window inequality},
\href{https://github.com/Deicyde/riemann-hypothesis/blob/40e56bab/src/Development/AutomorphicArithmetic/CharacterBoxes/ExistsPageN.lean\#L67-L77}{existence of the scale function}.}

\subsection{Where the two inputs are used}

On the global side, the geometric realization and the theta-lift
theorem first give the local--global comparison for this normalization
of the parameters. More precisely the geometric input supplies the
initial representations at the odd levels in the construction, while
weight existence supplies the congruences. The passage to the limit,
descent, and the comparison of traces give the simultaneous core
statement. From it one obtains spherical compatibility on the unlinked
support and the Iwahori-fixed comparison. Applying these to the local
and global parameters produces the form of the comparison required
above.\footnote{\href{https://github.com/Deicyde/riemann-hypothesis/blob/40e56bab/src/Development/AutomorphicArithmetic/Foundations/GaloisRepresentations/IntegralStructures/Development.lean\#L18-L31}{Local--global comparison from geometric and theta-lift inputs},
\href{https://github.com/Deicyde/riemann-hypothesis/blob/40e56bab/src/Development/AutomorphicArithmetic/Foundations/GaloisRepresentations/IntegralStructures/RowSrc.lean\#L4-L8}{the resulting Galois theorem}.}

Use this compatibility in the first and third components of the first
input. Between them use the local parameter constructed above. To
obtain the window input, combine the same compatibility with isobaric
base change. The four sets of data used in the window theorem are local
representation theory, automorphic results with base change and Galois
compatibility, Galois cohomology, and the CM and abelian-variety
results. The local data include multisegment classification and the
finite-field cuspidal characters. The CM data include deformation,
comparison of periods, good reduction, the elliptic CM theorem and the
Hecke--Tate functional equation. With the four sets supplied the
window theorem applies. Both inputs to the character theorem have now
been specified on the same Galois comparison.\footnote{\href{https://github.com/Deicyde/riemann-hypothesis/blob/40e56bab/src/Development/AutomorphicArithmetic/AutomorphicVarieties/ModularityArgument/ThresholdArguments/DenseThresholdCases/TopMainZj.lean\#L18-L74}{Local, automorphic, cohomological and CM inputs to the window},
\href{https://github.com/Deicyde/riemann-hypothesis/blob/40e56bab/src/Development/AutomorphicArithmetic/AutomorphicVarieties/ModularityArgument/ThresholdArguments/DenseThresholdCases/VwinOfRow.lean\#L4-L5}{application to the specified compatibility theorem}.}

The final arithmetic implication uses the theta integral equality and
local matching as two separate inputs. It also uses the covolume
estimate for the Hodge bundles, local Langlands existence and the
arithmetic windows just constructed. In the passage from the theta
inputs to the family of comparisons, the proof obtains two lower
thresholds \(C_0\) and \(C_1\). The first ensures the existence of
the required row functionals and the second ensures the evaluation
property on the grid of parameters. Choosing \(C=\max(C_0,C_1)\)
permits using both conclusions for one grid. The linked-bundle
covolume estimate applies with this same \(C\), as does the bound
for the contributions at two and three on the chosen family. These
give the required family of hypotheses. The primitive-character
argument and the Euler-factor calculation in
Section~\ref{sec:restore-euler} now yield
\eqref{eq:all-character-nv}.\footnote{\href{https://github.com/Deicyde/riemann-hypothesis/blob/40e56bab/src/Development/AutomorphicArithmetic/ThetaCorrespondence/ThetaPairLifts/ThirdPlaceSeries/EveryCharacterOfLiteratureOfVWin.lean\#L810-L813}{Inputs to the last arithmetic implication},
\href{https://github.com/Deicyde/riemann-hypothesis/blob/40e56bab/src/Development/AutomorphicArithmetic/ThetaCorrespondence/ThetaPairLifts/ThirdPlaceSeries/SiegelEisensteinSecondTermHypotheses.lean\#L44-L73}{choice of the common threshold and application to characters}.}

\section{The arithmetic window and the polynomial bound at level two}
\label{sec:arithmetic-translation}

Two statements involving arithmetic lengths occur close to the conclusion in
the source. The proof of nonvanishing for all characters uses the window
assertion denoted by \(\mathrm{VWin}\). A separate theorem in the module
\texttt{Small\allowbreak{}Class\allowbreak{}Numbers.\allowbreak{}Cor\allowbreak{}GKOf} applies a general estimate to the trivial
character modulo two. We give the closing deductions of both statements and
specify the results imported at each deduction. Merely importing a module does
not make every theorem in that module a premise of the theorem being proved.

\subsection{The four counts in the window argument}

Here is the window with its quantifiers, from \texttt{Exists\allowbreak{}Page\allowbreak{}N}. For a
number field \(E\), a prime \(\ell\), and a positive twist \(k\), let
\(s_{\ell,E}(k)\) be the \(\mathbb Z_\ell\)-length of the kernel of the
localization map in degree one for
\((\mathbb Q_\ell/\mathbb Z_\ell)(-k)\). The global group allows
ramification at the primes above \(\ell\); local restriction is taken at
each of those primes. Formally the length is defined for a choice of
algebraic closure of \(E\) and local data \(D_\ell\) at each prime.
The assertion \(\mathrm{VVle}(E,t,B)\) quantifies over every such closure
and every family \(D_\ell\). It requires finiteness for every prime and
every \(1\le k<t\), and, for every finite set \(S\) of primes, the inequality
\[
  \sum_{\ell\in S}\log\ell\sum_{1\le k<t}s_{\ell,E}(k)\ \le B.
\]
Shrinking \(t\) and increasing \(B\) preserves the assertion: the
weights and omitted lengths are nonnegative, and finiteness is inherited.

Fix a primitive character \(\chi\) of level \(m\), a real number
\(0<\varepsilon<1/52\), and any function \(a:\mathbb R\to\mathbb N\)
such that
\[
   x^{2+8\varepsilon}\le 2^{a(x)}\le 2x^{2+8\varepsilon}
       \qquad(x\ge 1).
\]
There are \(J\ge1\) and \(X_0\) such that the following statement holds
for every \(X\ge X_0\). Write \(U=X^{1/J}\) and take any primes
\(q\equiv3\bmod 8,\ d\equiv7\bmod16\), with \(q\mid d+1\), such that
\[
 \begin{split}
  \tfrac12U^{\varepsilon/5}&<\log q\le U^{\varepsilon/5},\\
  \tfrac12U^{\varepsilon/5}(\log U)^2&<\log d
         \le U^{\varepsilon/5}(\log U)^2,
 \end{split}
\]
and such that \(p\mid d+1\) for every odd prime \(p\le U^{\varepsilon/5}\).
There exists \(Y\ge U\) with \(Y^2\le X\) such that, for every
\(Y\le x\le Y^2\) and every realization of the following field \(E\),
\[
   \mathrm{VVle}\!\left(E,\lfloor x^{401}\rfloor,\,
                      \sqrt{\exp(x^\varepsilon)}\right).
\]
The field condition means that \(E\) embeds in a cyclotomic field at a
level divisible by \(2^{a(x)+2},m,q,d\), and, for a fixed primitive
root at that level, the automorphism given by exponent \(b\) fixes
the image of \(E\) pointwise precisely when
\[
 b\equiv\pm1\bmod 2^{a(x)+2},\qquad
 \chi(b)=1,\qquad
 \left(\frac bq\right)=\left(\frac bd\right)=1.
\]
In particular \(q,d\) are fixed throughout the interval for \(x\), while
the two-power exponent \(a(x)\) can vary there.

Let \(H_N\) denote the collection of local representation-theoretic,
automorphic, Galois-cohomological and complex-multiplication results called
\texttt{Literature\allowbreak{}Rank\allowbreak{}NHypothesis}. In the proof of
\texttt{vwin\_\allowbreak{}of\_\allowbreak{}row}, the input \(h\) is the Galois compatibility assertion
for the specified local parameter relation. The open-image record is formed
from \(h\). The local inputs are obtained from this record, the automorphic
inputs from this record together with base change and \(h\), and the Galois
cohomology and complex-multiplication inputs from the same record. These four
constructions give \(H_N\). No auxiliary field has yet been specialized to
discriminant \(-3\).

Apply \texttt{published\_\allowbreak{}hypotheses\_\allowbreak{}from\_\allowbreak{}patching\_\allowbreak{}alt} to \(H_N\) and the
weak descent theorem \texttt{desc\allowbreak{}Weak}, and call the resulting data \(H'\).
The size lemma called \(K\) supplies a real constant \(C_K\) from the
rank-\(N\) part of \(H'\). The endpoint theorem
\texttt{ends0\_\allowbreak{}alt} supplies two assertions. The first, together with
the rank-\(N\) inputs, produces unramified-coefficient automorphy. With
the two-dimensional Kato input, \(H'\), and the size lemma \(K\), this
produces the count in case I. The two counts in cases II and III are
obtained together by applying
\texttt{count\allowbreak{}II\_\allowbreak{}and\_\allowbreak{}count\allowbreak{}III\_\allowbreak{}of\_\allowbreak{}prints\allowbreak{}H\_\allowbreak{}of\_\allowbreak{}sp4\allowbreak{}In}
to \(H'\), \(K\), and its four-dimensional symplectic input. The second
endpoint, with \(H'\) and \(K\), gives the count in case V through the
reduced unitary lift. These are four distinct imported count theorems.

The last step uses \texttt{of\_\allowbreak{}four\_\allowbreak{}alt}. Its hypotheses consist of
\(H_N\), the existence of \(C_K\), and counts I, II, III, V. Apply
\texttt{thm\allowbreak{}Bab\_\allowbreak{}of\_\allowbreak{}lemma\allowbreak{}K} to the global duality input, the two local
cohomology inputs, the Ferrero--Washington input, and the size lemma.
Specialize the resulting auxiliary estimate to \(401\). The minus-eigenspace bound is supplied by \texttt{thm\allowbreak{}Bc\allowbreak{}Free\_\allowbreak{}of} applied to \texttt{ML.\allowbreak{}iv}. On the positive side the exceptional-prime sum uses \texttt{count\allowbreak{}IV} and \texttt{count\allowbreak{}Vs} alongside II and V. Case Z is supplied by \texttt{count\allowbreak{}Z} applied to \texttt{prop\allowbreak{}Z}. The theorem \texttt{vwin\_\allowbreak{}of\_\allowbreak{}counts} now applies to these
inputs and the four counts, yielding \(\mathrm{VWin}\). This describes the
use of the four estimates; their arithmetic proofs are dependencies of this
deduction. The source loci are \texttt{Vwin\allowbreak{}Of\allowbreak{}Row},
\texttt{Top\allowbreak{}Main\allowbreak{}Zj}, \texttt{Window\allowbreak{}Bound\allowbreak{}Of\allowbreak{}Literature\allowbreak{}Level}, and
\texttt{HVCore}, all in the \texttt{Dense\allowbreak{}Threshold\allowbreak{}Cases} directory.

\subsection{Character sums in the four counts}

The sets of characters and the range of summation are part of the
estimates. Fix the target \(\chi\). A host \(h\) determines its character
modulus \(N\), the two primes \(q_1,d\), and the subgroup \(H_h\) of
\((\mathbb Z/N\mathbb Z)^\times\) fixing its abelian field.
For coefficient ring \(\mathcal O_\ell\) and twist \(k\), let
\(\mathcal X_{\ell,k}\) be the characters modulo \(N\), valued in
\(\mathcal O_\ell\), trivial on \(H_h\), and with
\(\psi(-1)=(-1)^k\). Let \(b_{\ell,k}(\psi)\) be the strict length for
the character \(\psi^{-1}\chi_\ell^{-k}\), with strict local conditions
at the primes dividing \(\ell N\); put
\(v_{\ell,k}(\psi)=b_{\ell,k}(\psi)/r_\ell\) when the length is finite.
Write \(B_\chi\) for the exceptional primes of the target. Define
\(\psi^{(0)}_\ell\) by removing the \(\ell\)-power-order part of the
character and then removing its component at the prime \(\ell\) in
the level. The cases which occur in the count statements are
\[
\begin{array}{cl}
\mathrm I&
 \ell\notin B_\chi,\quad \ell-1\nmid 2k;\\
\mathrm Z&
 \ell\notin B_\chi,\quad\ell-1\mid 2k,\quad\psi=1;\\
\mathrm {II}&
 k\ge2,\quad 2\mid\operatorname{ord}(\psi_\ell^{(0)}),\quad
    (\ell\in B_\chi\ \text{or}\ \ell-1\mid 2k);\\
\mathrm {III}&
 \ell\notin B_\chi,\quad\ell-1\mid 2k,\quad
   \operatorname{ord}(\psi)\ \text{odd},\quad\psi\ne1;\\
\mathrm {IV}&
 \ell\in B_\chi,\quad k\ge2,\quad
   \operatorname{ord}(\psi_\ell^{(0)})\ \text{odd};\\
\mathrm {V}&\ell\in B_\chi,\quad k=1.
\end{array}
\]
The host-size condition at cutoff \(t\) requires that both \(q_1,d\)
exceed \(2t+1\) and all five host parameters \(q,o,p_K,h_*,p_1\)
in \texttt{Host.\allowbreak{}Big}. In case I there exist an integer \(C\) and test
primes \(r_1,r_2\), depending on the target, such that whenever the
host meets this condition and both auxiliary primes exceed both test
primes, for each \(1\le k<t\), every choice of coefficient rings,
and every finite set \(S\) of odd coefficient primes in case I,
all the relevant lengths are finite and
\[
     \sum_{\ell\in S}\log\ell
             \sum_{\psi\in\mathcal X_{\ell,k}}v_{\ell,k}(\psi)
            \le (Nt)^C .
\]
In II and III the assertion is pointwise. For each target there is an
integer \(C\) such that at all these hosts, all \(1\le k<t\),
all odd primes, all coefficient choices and all
\(\psi\in\mathcal X_{\ell,k}\) in the specified case,
\[
      b_{\ell,k}(\psi)<\infty,\qquad
      (\log\ell)\,v_{\ell,k}(\psi)\le(\ell Nk)^C .
\]
No additional condition involving \(r_1,r_2\) occurs in the pointwise
statement. Case V has this shape on the branch
\(R_1(\ell,\psi)\le 2^{c_2(\psi)}\) of the host cutoff; the source
assigns the other branch a separate name \texttt{Count\allowbreak{}Vs}. Case Z
uses the vanishing for the trivial character. Case IV uses a
cyclotomic-layer estimate summed over twists. In particular the
pointwise estimate does not authorize summing an infinite set of primes.
The finite-prime-set quantifier in \(\mathrm{VVle}\) is supplied after
the separate estimates have been assembled.

\subsection{The exceptional primes and the sum over the other primes}

The exceptional primes can be computed before any window is chosen.
Write \(o=\operatorname{ord}(\chi)\). Let \(p_K\) be the least prime
\(p\ge7\) congruent to \(3\bmod4\) and prime to \(m\), let
\(h_*\) be the class number of \(\mathbb Q(\sqrt{-p_K})\), and
let \(p_1\) be the least prime coprime to \(2mp_K\) and congruent
to one modulo the odd part of \(o\). The definition in \texttt{Stmt\allowbreak{}Fields}
is
\[
 B_\chi=\{3\}\ \cup\
 \{p:\ p\text{ an odd prime},\ p\mid o\,m\,p_K\,h_*\,p_1\}.
\]
For the primitive trivial target, \(o=m=1\), \(p_K=7\) and \(p_1=3\).
The field \(\mathbb Q(\sqrt{-7})\) has discriminant \(-7\).
Every ideal class has an integral representative of norm at most
\(2\sqrt7/\pi<2\) by the imaginary quadratic Minkowski bound.
Its norm is one. Thus \(h_*=1\) and \(B_{\mathbf1}=\{3,7\}\).
For a general target the product in the display can have additional
odd prime factors.

For a CM realization \(E\) of the host and local data for its
algebraic closure, write \(\lambda^+_{\ell,e}\) for the
\(\mathbb Z_\ell\)-length of the submodule of strict classes
where complex conjugation acts by \((-1)^e\), called
\texttt{plus\allowbreak{}Len}. We derive the assertion
\texttt{Sum\allowbreak{}Out} of \texttt{Height\allowbreak{}And\allowbreak{}Vanishing.\allowbreak{}Sum\allowbreak{}Out}. Take the
constants \(C_1,C_Z,C_2,C_3\) supplied by the four indicated
counts for \(\chi\), and take the two test primes supplied by count I.
For every host satisfying the size and test-prime conditions above,
every integer \(t\ge2\), and every finite set \(S\) of odd primes
disjoint from \(B_\chi\), put
\[
        M=Nt,\qquad K=\max\{C_Z,C_2,C_3\}.
\]
Here \(M\ge2\). All \(\lambda^+_{\ell,e}\) for \(\ell\in S\) and
\(1\le e<t\) are finite, and
\begin{equation}
 \sum_{\ell\in S}\log\ell\sum_{1\le e<t}\lambda^+_{\ell,e}
       \le M^{C_1+3K+4}.
 \label{eq:rh-arith-sumout}
\end{equation}
The level and \(t\) vary while the exponent and test primes remain fixed.

To prove this assertion, choose at each \(\ell\) an admissible coefficient
ring with the roots of unity required by the character comparison.
Fix \(e\) in the indicated range. The contribution from all \(\ell\in S\)
in case I to the weighted character sum is at most \(M^{C_1}\).
For the other primes, \(\ell-1\mid2e\), hence
\(\ell\le 2e+1\le2t\). There are at most \(2t\) such primes, by
injecting them into the integers from \(1\) to \(2t\). The case
decomposition in \texttt{Height\allowbreak{}And\allowbreak{}Vanishing.\allowbreak{}Cases} places each of their
characters in Z, II, or III. Since \(\ell N e\le2Mt\le M^3\), the
pointwise estimate gives
\[
    (\log\ell)v_{\ell,e}(\psi)\le M^{3K}.
\]
There are at most \(N\) characters at a given prime, by
\texttt{Card\allowbreak{}Le\allowbreak{}Level.\allowbreak{}x\allowbreak{}Card}. Consequently the entire weighted character
sum at this \(e\) is at most \(M^{C_1}+2M^{3K+1}\).
Its individual lengths are finite in Z, II and III by the pointwise
assertions, and in I by applying that count to a singleton prime.

The imported field-to-character comparison \texttt{red\allowbreak{}Prime} applies at
each prime of \(S\). Here the prime is prime to the index of the
subgroup defining the host field, by \texttt{bchi\allowbreak{}Facts}, and
\texttt{host\allowbreak{}Level} identifies the field at modulus \(N\).
With zero error term the summation lemma \texttt{Chr.\allowbreak{}sum\allowbreak{}Back} gives
finiteness on the field side and
\[
     \sum_{1\le e<t}\lambda^+_{\ell,e}
        \le \sum_{1\le e<t}\sum_{\psi\in\mathcal X_{\ell,e}}
                                      v_{\ell,e}(\psi).
\]
Multiply by the nonnegative \(\log\ell\), sum over \(S\), interchange
the finite sums, and apply the preceding estimate for each \(e\).
There are at most \(t\) twists. Since \(t\le M\), the total is at most
\[
   t(M^{C_1}+2M^{3K+1})
   \le M^{C_1+1}+2M^{3K+2}
   \le M^{C_1+3K+4}.
\]
For the last inequality, each term without its numerical coefficient
is at most \(M^{C_1+3K+2}\) and \(3\le M^2\).
This proves \eqref{eq:rh-arith-sumout}.

\subsection{Restoring the exceptional primes}

Fix a function \(j:\mathbb N\to\mathbb N\). For each
\(\ell\in B_\chi\) the source uses \texttt{Layer\allowbreak{}Le} at \(j(\ell)\):
for every number field \(L_0\) realizing the abelian-field datum
\(L_\chi(\ell)\) of the full cyclotomic field of level \(\ell m\),
and every cyclotomic layer \(L_1/L_0\) obtained by
adjoining \(\ell\)-power roots of unity and having degree
\(\ell^{j(\ell)}\), the single-prime strict-length bound
\(\mathrm{VVat}(\ell,L_1,t,B)\) holds. Explicitly, for every
algebraic closure of \(L_1\) and each choice of local data, each
\(s_{\ell,L_1}(e)\) for \(1\le e<t\) is finite and
\(\sum_{1\le e<t}s_{\ell,L_1}(e)\le B\). Assume a common \(B\ge0\)
for all the specified layers. The input \texttt{Sum\allowbreak{}In} then gives,
for each exceptional prime, a natural number
\(C_\ell=C_\ell(\chi,\ell,j(\ell))\) and
\[
       \log\ell\sum_{1\le e<t}\lambda^+_{\ell,e}
                   \le M^{C_\ell}+C_\ell B,
\]
with each length finite. The same host-size condition is used above and
below. For primes outside \(B_\chi\), write \(C_0\) for the exponent
in \eqref{eq:rh-arith-sumout}. Put
\[
       b=|B_\chi|,\quad D=\sum_{\ell\in B_\chi}C_\ell,\quad
       C=C_0+D+b+bD .
\]
We claim that for every finite set \(S\) of odd primes the left side
of \eqref{eq:rh-arith-sumout} is at most \(M^C+CB\).
Indeed, \(C_\ell\le D\). Summing the exceptional estimates and then
applying \eqref{eq:rh-arith-sumout} to \(S\setminus B_\chi\) gives
\[
     |S\cap B_\chi|(M^D+DB)+M^{C_0}.
\]
The terms independent of \(B\) total at most
\((b+1)M^{D+C_0}\), since \(M\ge2\). Now \(b+1\le2^b\le M^b\)
and \(D+C_0+b\le C\). The coefficient of \(B\) is at most \(bD\le C\).
This proves the claim. Finiteness for any given odd prime follows
by its singleton estimate on the corresponding side of the partition.
This is the proof of \texttt{Height\allowbreak{}And\allowbreak{}Vanishing.\allowbreak{}Point\allowbreak{}Arith.\allowbreak{}point}.
It makes explicit the dependence on the target and its layer indices,
and the fact that the positive-length estimate still has a layer
hypothesis at every exceptional prime.

\subsection{Recombining the two eigenspaces and the prime two}

Let \(c\) denote the action of complex conjugation on the twisted
cohomology over \(E\). Functoriality of the Galois action gives
\(c^2=1\) and its compatibility with localization preserves the strict
submodule. Write \(\lambda^-_{\ell,e}\) for the length in the
\((-1)^{e+1}\)-eigenspace. For an odd \(\ell\), the element \(2\)
is invertible in \(\mathbb Z_\ell\), because its reduction modulo
\(\ell\) is nonzero. Put \(s=(-1)^e\). The strict class \(x\) has the
decomposition
\[
        x=\frac{x+scx}{2}+\frac{x-scx}{2}.
\]
The two summands lie in the strict submodule. Applying \(c\) multiplies
the first by \(s\) and the second by \(-s\). Their eigenspaces intersect
trivially, since a member of the intersection is killed by \(2s\), a
unit. Thus addition is an isomorphism from their direct sum to the
strict submodule. Additivity of module length proves
\[
       s_{\ell,E}(e)=\lambda^+_{\ell,e}+\lambda^-_{\ell,e}.
\]
This equality also holds for extended lengths. If each signed length
is finite, their sum is finite.

Suppose that the positive and negative estimates over every finite set
of odd primes are \(B_+\) and \(B_-\), respectively, with the displayed
twist range, and that the full lengths at two are finite with
\(\sum_{1\le e<t}s_{2,E}(e)\le B_2\). For any finite set of primes
\(S\), split off \(S\cap\{2\}\). Its weighted contribution is at most
\((\log 2)B_2\), because all lengths are nonnegative. Apply the
eigenspace equality at every prime in \(S\setminus\{2\}\) and sum the
two estimates there. We obtain
\[
       \mathrm{VVle}(E,t,\,(\log 2)B_2+B_-+B_+).
\]
These are the proofs of \texttt{sha\allowbreak{}Len\_\allowbreak{}eq\_\allowbreak{}add} and
\texttt{vvle\_\allowbreak{}of\_\allowbreak{}parts} in \texttt{Length\allowbreak{}Eq\allowbreak{}Add\allowbreak{}Eigen}.

\subsection{From two congruence integers to an orbit sum}

This is the implication from \(\mathrm{PropGK}\) to \(\mathrm{CorGK}\)
proved in \texttt{Unitary\allowbreak{}Three\allowbreak{}Groups.\allowbreak{}Cor\allowbreak{}GSum\allowbreak{}Log}. Fix a number field
\(K\). The antecedent gives integers \(c_5,A\ge0\) depending on
\(K\). Given \(k\ge1,\ N\ge1\), a subgroup
\(H\le(\mathbb Z/N\mathbb Z)^\times\) of index \(u\), and a prime
\(p\nmid 30N|D_K|\), there is a nonzero integer \(n_p\) with
\[
   \log|n_p|\le c_5\varphi(u)\{N(k+1)\}^A\log p.
\]
For an algebraic closure \(\Omega\) of \(\mathbb Q\) equipped with an
embedding of \(K\), and an admissible ring \(\mathcal O_\ell\), write
\(\mathcal X_\ell(H)\) for the finite set of characters valued in this
ring whose kernel on the units modulo \(N\) is exactly \(H\).
For \(\ell\ne p\), provided at least one character in
\(\mathcal X_\ell(H)\) satisfies the good-data predicate, the
antecedent gives finiteness of all lengths in that orbit and
\[
       \sum_{\psi\in\mathcal X_\ell(H)}v_{\ell,k}(\psi)
               \le 2\operatorname{ord}_\ell(n_p).
\]
The integer \(n_p\) has already been chosen before \(\Omega,\ell\)
and its coefficient ring.

Choose distinct allowable primes \(p,p'\), a finite set \(S\) of
coefficient primes and, for each \(\ell\in S\), a witness for the
good-data predicate in its orbit. When \(S\) is empty the required sum
is zero. Otherwise such a witness gives \(k\ge1\), so we have both
\(n_p\) and \(n_{p'}\). For each \(\ell\in S\), at least one is
applicable, even when \(p\) or \(p'\) belongs to \(S\). The valuation
of the other nonzero integer is nonnegative. Thus the lengths are
finite and their sum at \(\ell\) is at most
\(2\operatorname{ord}_\ell(n_p)+2\operatorname{ord}_\ell(n_{p'})\).
For every nonzero integer \(n\), the distinct prime powers
\(\ell^{\operatorname{ord}_\ell(n)}\), \(\ell\in S\), are pairwise
coprime and each divides \(n\); their product divides \(|n|\).
Taking the logarithm gives
\(\sum_{\ell\in S}\log\ell\,\operatorname{ord}_\ell(n)\le\log|n|\).
Apply this inequality twice to obtain exactly
\[
   \sum_{\ell\in S}\log\ell
         \sum_{\psi\in\mathcal X_\ell(H)}v_{\ell,k}(\psi)
     \le 2c_5\varphi(u)\{N(k+1)\}^A(\log p+\log p').
\]
This also explains why assigning a distinct auxiliary prime to every
coefficient prime would change the content of the estimate.

\subsection{A consequence for the trivial character at level two}

Fix an algebraic closure \(\Omega\) of \(\mathbb Q\). For a rational prime
\(\ell\), let \(\mathcal O_\ell\) be a permitted coefficient ring, write
\(\chi_\ell\) for its cyclotomic character, and write
\(r_\ell\) for the ramification index of the coefficient ring.
In the notation of \texttt{Small\allowbreak{}Class\allowbreak{}Numbers.\allowbreak{}Class\allowbreak{}Number\allowbreak{}Three}, put
\[
  b_\ell(e)=
  \operatorname{length}_{\mathcal O_\ell}
  \operatorname{ShaU}\!\left(
    \operatorname{twist}_2(\ell,\chi_\ell^{-e}),\,2\ell,\,\top
  \right).
\]
The underlying discrete coefficient module is
\(\mathcal O_\ell\otimes_{\mathbb Z_\ell}(\mathbb Q_\ell/\mathbb Z_\ell)\),
acted on by the indicated character; \(\mathcal O_\ell\) is a DVR
subalgebra of \(\overline{\mathbb Q}_\ell\), finite over
\(\mathbb Z_\ell\). Here \(r_\ell=\operatorname{length}_{\mathcal O_\ell}
(\mathcal O_\ell/\ell\mathcal O_\ell)\).
The symbol \(\top\) denotes the full Galois group, and the subgroup
\(\operatorname{ShaU}\) consists of degree-one classes restricting
trivially to all inertia groups and to the decomposition groups at
places above \(2\ell\). The length is initially in
\(\mathbb N\cup\{\infty\}\). The
polynomial estimate asserts that there exist natural numbers \(C,A\),
independent of all the following choices, such that for every
\(\Omega\), every positive even integer \(e\), every choice of coefficient
rings for all primes, and every finite set \(S\) of primes satisfying
\[
   \ell\geq 5,\qquad \ell-1\nmid 2e
       \qquad(\ell\in S),
\]
each \(b_\ell(e)\) for \(\ell\in S\) is finite and
\[
     \sum_{\ell\in S}(\log\ell)\,\frac{b_\ell(e)}{r_\ell}
       \ \leq\ C(e+1)^A .
\]
The division by \(r_\ell\) occurs separately in each summand.

We give the specialization proof. Every Dirichlet character modulo two
is trivial because \((\mathbb Z/2\mathbb Z)^\times\) has one element.
The Galois character attached to this trivial character is also trivial.
Consequently the character that defines the general length
\(V_{\ell,e}(1)\) is \(1^{-1}\chi_\ell^{-e}=\chi_\ell^{-e}\), and
\(V_{\ell,e}(1)=b_\ell(e)\), including the possibility of infinite
length. This is the equality \texttt{vlen\_\allowbreak{}one}.

Take \(K=K_*(3)\), the imaginary quadratic field of discriminant
\(-3\). Its complex-place count is one and its degree is two. The
Minkowski bound for an integral representative in each ideal class is
\[
   \frac{4}{\pi}\frac{2!}{2^2}\sqrt{3}
   =\frac{2\sqrt{3}}{\pi}<2 .
\]
The representative has positive integral norm and therefore norm one.
It is the unit ideal. This proves \(h(K)=1\).

We verify the local input to the general estimate. Let a valuation on
\(\Omega\) lie over \(3\). Every integer has valuation of absolute value
at most one. If \(\ell\ne3\) and both \(3\) and \(\ell\) had absolute
value less than one, write \(1=3a+\ell b\) by B\'ezout's identity.
The ultrametric inequality would make \(1\) have absolute value less
than one. Hence \(\ell\) is a unit at this valuation. The
\(\ell\)-adic cyclotomic character is trivial on inertia there, as is
the character modulo two. But \(K\) is ramified at \(3\). Thus inertia
at this place cannot act trivially on \(K\) while the whole joint
kernel of the two preceding characters fixes \(K\). This proves the
required noncontainment. The remaining hypotheses of the good-data
predicate follow from \(h(K)=1\), \(\lvert D_K\rvert=3\),
\(\ell\geq5\), \(\ell-1\nmid2e\), and the evenness of \(e\), which gives
\(1=(-1)^e\). The proof uses ramification at \(3\) to test \(K\); \(3\)
is outside the summation range of this polynomial estimate.

Apply the general orbit estimate \(\mathrm{CorGK}\) at this \(K\). Write
\(c_5,a\) for its natural-number constants. Embed \(K\) in \(\Omega\),
take level \(2\), the full subgroup of \((\mathbb Z/2\mathbb Z)^\times\),
and the distinct test primes \(7,11\). Neither divides
\(30\cdot2\lvert D_K\rvert=180\). At each \(\ell\in S\) supply the trivial
character and the good data just verified. Each orbit in the estimate
is the singleton consisting of that character. Its contribution thus
has no orbit multiplicity; the index of the full subgroup is one and
\(\varphi(1)=1\). The finiteness clause of \(\mathrm{CorGK}\) becomes
the claimed finiteness. In its real inequality the sum over characters
is the single term \(b_\ell(e)/r_\ell\). Using
\(\log 7\leq 7\), \(\log 11\leq11\) and expanding the \(a\)-th power of
\(2(e+1)\) gives the claimed bound with
\[
     A=a,\qquad C=2c_5\,2^a(7+11)=36\,c_5\,2^a .
\]
In particular this use of \(7\) is as a test prime, and the assertion
places no requirement that \(7\) belong to \(S\).

The proof of the imported \(\mathrm{CorGK}\) is assembled in
\texttt{Small\allowbreak{}Class\allowbreak{}Numbers.\allowbreak{}Cor\allowbreak{}GKOf}. Starting from \(H'\), the
construction first obtains \(K\)'s size lemma, unramified-coefficient
automorphy from the first endpoint, and cohomological input from the
rank-\(N\) data and the size lemma. It derives a further endpoint from
\texttt{linked\allowbreak{}Unl\allowbreak{}E\_\allowbreak{}alt}. The two algebraic bounds, the auxiliary
algebraic lemma, the Kato input, the two rigidity arguments, the
even-level admissibility result and unramified-coefficient automorphy
give the single-character bound. Combine it with the count of
rank-three unitary representations to obtain \(\mathrm{PropGK}\);
the general theorem \texttt{U3.\allowbreak{}cor\allowbreak{}GK} then yields \(\mathrm{CorGK}\).
Instantiating \(H'\) with the rank-\(N\) patching-and-descent data gives
the unconditional input used in the specialization. This last
paragraph identifies imported results; it does not supply their
individual proofs.

% Requires amsmath, amssymb, hyperref.
\section{The odd-prime estimate for the minus part}
\label{sec:odd-prime-minus}

The arithmetic window of \eqref{eq:entry-strict-length} bounds the full
strict cohomology group. One of the estimates contributing to that
bound is an inequality for its minus part at all odd primes. Two
different calculations enter the inequality. Character valuations
control the primes away from the degree of the auxiliary field. An
annihilator controls the primes dividing the degree. The sums below
retain their factors of \(\log\ell\) throughout.

\subsection{Statement for a fixed field and range of twists}

Write \(E\) for the auxiliary CM field, \(D=[E:\mathbb Q]\), and
\(F\) for its conductor, the least positive integer such that \(E\)
embeds in \(\mathbb Q(\zeta_F)\). For an odd prime \(\ell\) and
\(k\ge1\), let \(m_{\ell,k}\) be the \(\mathbb Z_\ell\)-length of the
minus part of the strict group for the twist \(-k\). More precisely,
if \(c\) acts on the twisted cohomology by complex conjugation, the
minus part is \(T_{\ell,k}(E)\cap\ker(c-(-1)^{k+1})\). The plus part
has eigenvalue \((-1)^k\). Conjugation preserves the strict kernel and
squares to the identity. Since \(2\) is a unit in \(\mathbb Z_\ell\),
the two operators \((1+(-1)^kc)/2\) and \((1-(-1)^kc)/2\) restrict to
complementary idempotents of the strict kernel. They give a direct sum
of these subgroups and hence additivity of the two lengths, even
before finiteness has been established.
\footnote{\href{https://github.com/Deicyde/riemann-hypothesis/blob/40e56bab/src/Development/AutomorphicArithmetic/AutomorphicVarieties/LengthEqAddEigen.lean\#L68-L103}{Parity and length additivity for odd primes}.}

Fix a positive integer \(m\), a subgroup
\(H\subseteq(\mathbb Z/m\mathbb Z)^\times\), and put
\(d_0=[(\mathbb Z/m\mathbb Z)^\times:H]\). Choose \(e\ge2\) and
primes \(q\equiv3\bmod8\), \(d\equiv7\bmod16\) with
\(\gcd(qd,m)=1\). In the field construction of
Section~\ref{sec:entry-assembly}, replace the test \(\psi(u)=1\)
by \(u\bmod m\in H\). Then
\[
       \ell\mid D,\quad \ell\ne2\quad\Longrightarrow\quad
       \ell\mid d_0
       \qquad(\ell\text{ prime}).
\]
To prove this divisibility, choose the cyclotomic field in the
definition of \(E\), and let \(M\) be the order of its chosen
primitive root. Each automorphism raises the root to a power \(u\)
prime to \(M\). Let \(h=2^{e+1}d_0\). Since
\(\varphi(2^{e+2})=2^{e+1}\), we have \(u^h\equiv1\bmod 2^{e+2}\).
The class of \(u\) modulo \(H\) has order dividing \(d_0\), so
\(u^h\bmod m\in H\). The exponent \(h\) is even and hence both
quadratic symbols of \(u^h\) are \(1\). It follows from the
fixed-field definition that the \(h\)-th power of every cyclotomic
automorphism fixes \(E\). Its fixing subgroup is normal since the
cyclotomic Galois group is abelian, and the quotient has order
\([E:\mathbb Q]=D\). For any prime \(\ell\mid D\), Cauchy's theorem
now gives a quotient element of order \(\ell\). Its \(h\)-th power is
the identity, so \(\ell\mid h\). When \(\ell\) is odd this implies
\(\ell\mid d_0\), as asserted.
\footnote{\href{https://github.com/Deicyde/riemann-hypothesis/blob/40e56bab/src/Development/AutomorphicArithmetic/AutomorphicVarieties/StmtFields.lean\#L11-L45}{The cyclotomic field and subgroup data},
\href{https://github.com/Deicyde/riemann-hypothesis/blob/40e56bab/src/Development/AutomorphicArithmetic/AutomorphicVarieties/OddResidueBounds/DegExistsLift.lean\#L33-L99}{Powering in the cyclotomic quotient and the degree divisibility}.}

There is an absolute \(C\ge0\) from the annihilator estimate
such that, for every integer \(t\ge2\), every finite set \(S\) of
odd primes, and every allowed choice of local data, the lengths
\(m_{\ell,k}\) with \(1\le k<t\) are finite and
\begin{align}
 \sum_{\ell\in S}(\log\ell)\sum_{k=1}^{t-1}m_{\ell,k}
 &\le 3Dt^2\log(tF) \notag\\
 &\quad{}+
 C\sum_{\substack{\ell\mid d_0\\\ell\ {\rm odd\ prime}}}
       \ell D^2t^2\bigl(\log(t\ell F)\bigr)^2 .
 \label{eq:galois-odd-minus-bound}
\end{align}
The same \(C\) works for all the fields and all the finite sets of primes.
\footnote{\href{https://github.com/Deicyde/riemann-hypothesis/blob/40e56bab/src/Development/AutomorphicArithmetic/AutomorphicVarieties/OddResidueBounds/ThmBcOfPsiBound.lean\#L147-L177}{The field conductor and the split according to its degree},
\href{https://github.com/Deicyde/riemann-hypothesis/blob/40e56bab/src/Development/AutomorphicArithmetic/AutomorphicVarieties/OddResidueBounds/ThmBcOfPsiBound.lean\#L266-L295}{the uniform odd-prime conclusion and its inputs}.}

\subsection{Primes dividing the degree}

For an odd \(\ell\) the annihilator input gives a nonnegative integer
\(r\) killing the minus subgroup by multiplication by \(\ell^r\), with
\[
       r\log\ell\le C_0D(k+1)\log(k\ell F)
\]
for one absolute \(C_0\ge0\). A second input gives an integer \(g\)
bounding the length of the subgroup killed by \(\ell\) in the
ambient global \(H^1\), with \(g\le 8\ell D\log(\ell F)\).

To see exactly how these estimates combine, let \(P\) be the minus
subgroup and \(P_i=P[\ell^i]\). Its filtration starts with \(P_0=0\)
and ends with \(P_r=P\). Multiplication by \(\ell^{i-1}\) embeds
\(P_i/P_{i-1}\) into the \(\ell\)-torsion of the ambient global group.
Indeed its image is killed by \(\ell\), and the kernel on \(P_i\) is
exactly \(P_{i-1}\). Thus \(\operatorname{length}P\le rg\), including
the case \(r=0\). This proves finiteness. Since all factors are
nonnegative and \(\log(\ell F)\le\log(k\ell F)\), it also gives
\begin{align}
 m_{\ell,k}\log\ell
 &\le (r\log\ell)g \notag\\
 &\le 8C_0\ell D^2(k+1)\bigl(\log(k\ell F)\bigr)^2 .
 \label{eq:galois-annihilator-length}
\end{align}
Set \(C=8C_0\). This estimate actually applies at every odd \(\ell\).
Using it only for \(\ell\mid D\) keeps its factor of \(\ell\) out
of the unbounded prime sum.\footnote{\href{https://github.com/Deicyde/riemann-hypothesis/blob/40e56bab/src/Development/AutomorphicArithmetic/AutomorphicVarieties/OddResidueBounds/ThmBcOfPsiBound.lean\#L15-L39}{The annihilator, generator estimate and resulting finite length with constant \(8C_0\)}.}

\subsection{The character sum away from the degree}

Let \(\rho:(\mathbb Z/F\mathbb Z)^\times\to
\operatorname{Gal}(E/\mathbb Q)\) be the cyclotomic action. It is
surjective, and \([\left(\mathbb Z/F\mathbb Z\right)^\times:\ker\rho]=D\).
For a fixed \(k\), let \(\mathcal C\) be the set of complex characters
trivial on \(\ker\rho\) such that \(\chi(-1)=(-1)^{k+1}\).
Passing to the quotient embeds
\(\mathcal C\) in the character group of
\((\mathbb Z/F\mathbb Z)^\times/\ker\rho\). The complex character
group of a finite abelian group has its order, so
\(\lvert\mathcal C\rvert\le D\).

Here are the valuation estimates used for this step. For
\(\ell\nmid 2D\) the characterwise input gives finiteness of
\(m_{\ell,k}\) and a finite set \(T_\ell\) of allowed \(\ell\)-adic
characters with
\[
     m_{\ell,k}\log\ell
       \le\sum_{\chi\in T_\ell}-\log
              \lvert\alpha_F(k+1,\chi)\rvert_\ell.
\]
To specify the normalization, write \(f\) for the conductor of \(\chi\)
and \(\chi^*\) for the primitive character inducing it. With \(B_j(X)\)
the \(j\)-th Bernoulli polynomial, the scalar is
\[
    \alpha_F(j,\chi)=F\!\!\prod_{\substack{p\ {\rm prime}\\p\le j+1}}\!p
          \,\cdot f^{j-1}\sum_{a=1}^f\chi^*(a)\,B_j(a/f).
\]
The product-valuation estimate gives nonnegative integers \(w_\ell\),
the same ones for
all these choices, such that
\begin{align*}
  \sum_{\chi\in T_\ell}-\log
              \lvert\alpha_F(k+1,\chi)\rvert_\ell
       &\le w_\ell\log\ell,\\
  \sum_{\ell\in S'}w_\ell\log\ell
       &\le\lvert\mathcal C\rvert\log B(F,k+1),&
 B(F,j)&=16F^{j+1}j!
\end{align*}
for each finite prime set \(S'\). Applying the first inequality to
each member of an arbitrary \(S'\) disjoint from the divisors of
\(2D\), and then using the second, yields
\begin{equation}\label{eq:galois-good-prime-sum}
      \sum_{\ell\in S'}m_{\ell,k}\log\ell
          \le D\log B(F,k+1).
\end{equation}
The last inequality uses \(B(F,k+1)\ge1\). The characterwise input and
the product-valuation estimate are the two arithmetic premises of
this deduction.\footnote{\href{https://github.com/Deicyde/riemann-hypothesis/blob/40e56bab/src/Development/AutomorphicArithmetic/AutomorphicVarieties/OddResidueBounds/ThmBcOfPsiBound.lean\#L41-L108}{The quotient-character injection, valuation estimates and sum over primes},
\href{https://github.com/Deicyde/riemann-hypothesis/blob/40e56bab/src/Development/AutomorphicArithmetic/AutomorphicVarieties/OddResidueBounds/ZeroMemUnrHom.lean\#L79-L110}{the characterwise arithmetic estimates},
\href{https://github.com/Deicyde/riemann-hypothesis/blob/40e56bab/src/Development/AutomorphicArithmetic/AutomorphicVarieties/GaloisLibrary/BernoulliNumbers/MapGenBChar.lean\#L9-L31}{Bernoulli and denominator conventions}.}

\subsection{Summing over the twists}

If \(1\le k<t\), then \(tF\ge2\) and \(k+1\le t\). Consequently
\[
  16\le(tF)^4,\qquad F^{k+2}\le(tF)^{t+1},\qquad
  (k+1)!\le(k+1)^{k+1}\le(tF)^t.
\]
Multiplying and taking logarithms gives
\(\log B(F,k+1)\le(2t+5)\log(tF)\). For the primes in \(S\) away
from \(D\), interchange the two finite sums in
\eqref{eq:galois-odd-minus-bound} and apply
\eqref{eq:galois-good-prime-sum} at each \(k\). The result is at most
\[
     D(t-1)(2t+5)\log(tF)\le 3Dt^2\log(tF),
\]
since \(3t^2-(t-1)(2t+5)=t^2-3t+5>0\).

For a remaining prime use \(k+1\le t\) and
\(0\le\log(k\ell F)\le\log(t\ell F)\) in
\eqref{eq:galois-annihilator-length}. Its contribution for all twists is
at most \(C\ell D^2(t-1)t(\log(t\ell F))^2\), hence at most
\(C\ell D^2t^2(\log(t\ell F))^2\). Every such \(\ell\) is an odd
prime divisor of \(d_0\). Enlarging the sum to all these divisors
only adds nonnegative terms. Adding the two bounds proves
\eqref{eq:galois-odd-minus-bound}. The proof applies at each fixed
field and range of twists, so it can be applied at \(e(x),q,d\) and
\(t=\lfloor x^{401}\rfloor\) as soon as those data meet the field
construction's hypotheses.\footnote{\href{https://github.com/Deicyde/riemann-hypothesis/blob/40e56bab/src/Development/AutomorphicArithmetic/AutomorphicVarieties/OddResidueBounds/ThmBcOfPsiBound.lean\#L110-L145}{The factorial and conductor estimate},
\href{https://github.com/Deicyde/riemann-hypothesis/blob/40e56bab/src/Development/AutomorphicArithmetic/AutomorphicVarieties/OddResidueBounds/ThmBcOfPsiBound.lean\#L178-L289}{summing the two ranges}.}

In the passage to \eqref{eq:entry-strict-length} this odd-prime
estimate is combined with the dyadic and fixed-prime estimates, the
other cohomological counts and the global window construction. The
source applies the odd-prime theorem to its characterwise comparison
obtained from Poitou--Tate duality and the supplied odd-prime
input.\footnote{\href{https://github.com/Deicyde/riemann-hypothesis/blob/40e56bab/src/Development/AutomorphicArithmetic/AutomorphicVarieties/ModularityArgument/ThresholdArguments/DenseThresholdCases/HVCore.lean\#L18-L22}{Application inside the proof of the full window},
\href{https://github.com/Deicyde/riemann-hypothesis/blob/40e56bab/src/Development/AutomorphicArithmetic/AutomorphicVarieties/ModularityArgument/ThresholdArguments/DenseThresholdCases/WindowBoundOfLiteratureLevel.lean\#L25-L34}{the automorphic hypotheses furnishing the other counts}.}

\section{The conjugate weights of the global Galois representation}
\label{sec:galois-parallel-weight}

Let \(E/E^+\) be a CM field with maximal totally real subfield \(E^+\)
and complex conjugation \(c\). Fix \(a\ge1\) and let \(\pi\) be regular
algebraic and cuspidal on \(\mathrm{GL}_a\) over \(E\), polarized by a
Hecke character \(\chi\) of \(E^+\). For a complex embedding \(\sigma\)
let \(I_\sigma\) be the multiset of infinitesimal-character parameters
of \(\pi\). We derive the integer \(b\) and the multisets \(H_\tau\)
required by Section~\ref{sec:entry-assembly}. The relation with the
actual jumps of the Galois representation uses the local de Rham
comparison at the end of the argument.

\subsection{The exponent at infinity of the polarizing character}

Let \(\psi=\chi\circ\operatorname{N}_{E/E^+}\), and write \(s_\sigma\)
for the exponent of its archimedean character at \(\sigma\). The
infinitesimal-character rules for twisting and contragredients, applied
to the invariant pairing expressing the polarization, give
\begin{equation}\label{eq:galois-inf-char}
    I_{\sigma c}=\{-(z-s_\sigma):z\in I_\sigma\}.
\end{equation}
Here twisting by \(\psi^{-1}\) adds \(-s_\sigma\), and passage across
the pairing negates the resulting parameter. The multiset \(I_\sigma\)
has cardinality \(a\ge1\). Take one member \(z\). Regular algebraicity
gives \(z=(a-1)/2+n\), and likewise
\(s_\sigma-z=(a-1)/2+n'\), for integers \(n,n'\). Adding shows
\[
                  s_\sigma=a-1+n+n'\in\mathbb Z.
\]
The pullback by the norm gives \(s_{\sigma c}=s_\sigma\). Hence these
integers give a well-defined \(s_\phi\) indexed by the real embeddings
\(\phi:E^+\to\mathbb R\).

For \(x\in(E^+)^\times\) totally positive, put
\(u_\phi=\exp(\tfrac12\log\phi(x))\) and regard \(u\) as a positive
idele at the complex places of \(E\), equal to \(1\) at finite places.
Its norm to \(E^+\) has archimedean components \(\phi(x)\). Evaluating
the archimedean character of \(\psi\) at \(u\) multiplies the two
equal contributions above each real embedding. It follows that the
archimedean component of \(\chi\) satisfies
\begin{equation}\label{eq:galois-inf-weight}
              \chi_\infty(x)=\prod_\phi\phi(x)^{s_\phi}.
\end{equation}
\footnote{\href{https://github.com/Deicyde/riemann-hypothesis/blob/40e56bab/src/Development/AbelianVarietyTheory/Proofs/SatakeParameters/ExistsIsWeight.lean\#L25-L67}{The invariant pairing, integrality and conjugation of the exponents},
\href{https://github.com/Deicyde/riemann-hypothesis/blob/40e56bab/src/Development/AbelianVarietyTheory/Proofs/SatakeParameters/ExistsIsWeight.lean\#L67-L138}{the norm at the archimedean places}.}

In fact \(s_\phi\) is independent of \(\phi\). To see this in the form
used in the construction, take a unit \(v\) of \(E^+\). The finite
idele of \(v\), raised to the powers \(j!\), converges to \(1\).
Indeed the image of \(v\) in the unit group of every finite quotient
by a nonzero ideal has finite order, which divides \(j!\) eventually.
These congruences prove convergence in the product of the local unit
groups. Continuity of \(\chi\) shows that the modulus of its value
on this finite idele is \(1\), since a positive real \(r\) with
\(r^{j!}\to1\) must be \(1\). The value on the full principal idele
is \(1\), so \(\lvert\chi_\infty(v)\rvert=1\). Apply
\eqref{eq:galois-inf-weight} to the totally positive element \(v^2\)
and take logarithms to obtain
\[
               \sum_\phi s_\phi\log|\phi(v)|=0.
\]
By the product formula also \(\sum_\phi\log|\phi(v)|=0\). Choose
\(\phi_0\) minimizing \(s_\phi\). In degree one the exponents are
already constant. In higher degree, Dirichlet's unit theorem gives a
unit \(v\) with \(\log|\phi(v)|<0\) for all \(\phi\ne\phi_0\), since
the full log-unit lattice in the sum-zero hyperplane meets the open
cone given by those strict inequalities. Subtracting \(s_{\phi_0}\) times the
product formula from the preceding equality for this \(v\) gives
\[
         \sum_\phi (s_\phi-s_{\phi_0})\log|\phi(v)|=0 .
\]
Every summand is nonpositive, and a strict increase in any of the
integers \(s_\phi\) would make its summand negative. Thus there is
one integer \(s_0\) common to all the embeddings.
\footnote{\href{https://github.com/Deicyde/riemann-hypothesis/blob/40e56bab/src/Development/AutomorphicArithmetic/AutomorphicVarieties/BuildingsResidueChars/HeckeCharParallelWeight.lean\#L39-L87}{Convergence of factorial powers of the finite unit idele},
\href{https://github.com/Deicyde/riemann-hypothesis/blob/40e56bab/src/Development/AutomorphicArithmetic/AutomorphicVarieties/BuildingsResidueChars/HeckeCharParallelWeight.lean\#L178-L287}{the minimum exponent and the choice of unit}.}

\subsection{Computing the jumps at conjugate embeddings}

Set \(b=a-1-s_0\). Equation~\eqref{eq:galois-inf-char} gives
\(I_{\sigma c}=\{a-1-b-z:z\in I_\sigma\}\). If
\(\iota:\overline{\mathbb Q}_\ell\to\mathbb C\) is the fixed
coefficient field isomorphism, define
\[
       H_\tau=\{(a-1)/2-z:z\in I_{\iota\tau}\}
\]
for each \(\tau:E\hookrightarrow\overline{\mathbb Q}_\ell\).
Each member is an integer by regular algebraicity. In the multiset
for \(\tau c\), the term belonging to \(a-1-b-z\) equals
\[
       (a-1)/2-(a-1-b-z)=b-\bigl((a-1)/2-z\bigr).
\]
This proves \(H_{\tau c}=\{b-h:h\in H_\tau\}\) with multiplicity.
For each local embedding in the scope of the de Rham comparison,
its jumps clause identifies \(H_\tau\) with the jumps at \(\tau\),
using exactly the displayed infinitesimal-character formula as its
hypothesis. Thus the same choice of \(b\) and \(H_\tau\) works at
each such place of the global representation.
\footnote{\href{https://github.com/Deicyde/riemann-hypothesis/blob/40e56bab/src/Development/AbelianVarietyTheory/Proofs/SatakeParameters/ExistsIsWeight.lean\#L138-L150}{The formula for the global weight},
\href{https://github.com/Deicyde/riemann-hypothesis/blob/40e56bab/src/Development/AutomorphicArithmetic/Foundations/GaloisRepresentations/IntegralStructures/InertCharpoly/GaloisLocalGlobalOfLocalGlobalCompatibility.lean\#L127-L171}{integral weights, conjugate weights and their occurrence in the local comparison}.}

\section{The generator bound at an odd prime}
\label{sec:selmer-generators}

Let \(E\) be a number field of degree \(D\) admitting an embedding in
\(\mathbb Q(\mu_F)\), where \(F\ge1\). Fix a prime \(\ell\) and let
\(S_E\) be the finite places of \(E\) above \(\ell\). Write
\(G_{E,S_E}\) for the Galois group unramified at all finite places
outside \(S_E\). Take the discrete module
\(A=(\mathbb Q_\ell/\mathbb Z_\ell)(-k)\) for an integer \(k\).
We prove that the subgroup
\[
        V=H^1(G_{E,S_E},A)[\ell]
\]
is finite of \(\mathbb Z_\ell\)-length \(g\), with
\[
              g\le 8\ell D\log(\ell F).
\]
This is the ambient cohomology group in
\eqref{eq:galois-annihilator-length}. In particular the estimate
allows ramification at places above \(\ell\) and imposes no strict
local condition at those places. The same calculation works at
\(\ell=2\) for this generator bound.

\subsection{Restriction to the first cyclotomic layer}

Put \(K=E(\mu_\ell)\), \(n=[K:\mathbb Q]\), and
\(S=\{\mathfrak p\mid\ell\text{ in }K\}\). The action on the
\(\ell\)-th roots embeds \(\operatorname{Gal}(K/E)\) in
\(\mathbb F_\ell^\times\), so \([K:E]\mid\ell-1\) and
\(n\le\ell D\). Both \(E\) and the required roots embed compatibly
in \(\mathbb Q(\mu_{\ell F})\), by extending the given embedding of
\(E\) to an algebraic closure and taking the compositum there.

For each \(v\in V\), choose a continuous representative \(z\).
The equation \(\ell v=0\) gives \(\ell z=db\), with coboundary
\(db(\sigma)=\sigma b-b\). Multiplication by \(\ell\) on \(A\)
is onto. Choosing \(b'\) with \(\ell b'=b\), we may replace \(z\)
by \(z-db'\), so its values belong to \(A[\ell]\).
Fix an additive identification of this group with \(\mathbb F_\ell\).
On \(G_K=\operatorname{Gal}(\overline E/K)\) the residual
cyclotomic character is trivial. The restriction of the inflated
cocycle is therefore an additive continuous homomorphism \(t_v\).
It vanishes on inertia outside \(S\), since such inertia maps into inertia
at a place outside \(S_E\), which the original quotient kills.

The choices of representatives give an injection of sets from \(V\)
to these homomorphisms. Indeed, if the restrictions for two classes
agree, their difference \(f\) vanishes on \(G_K\) and has values in
\(A[\ell]\). Its cocycle law makes it constant on right cosets of
\(G_K\), and the action on its values factors through
\(\Delta=\operatorname{Gal}(K/E)\). Write \(d=|\Delta|\) and
\(b=-d^{-1}\sum_{\tau\in\Delta} f(\tau)\) in \(A[\ell]\), where
\(d\) is prime to \(\ell\). Summing
\(f(\sigma\tau)=f(\sigma)+\sigma f(\tau)\) over \(\tau\) gives
\(f(\sigma)=\sigma b-b\). The two original classes coincide.
This proof uses the homomorphisms only as an upper bound on the
number of classes and is independent of the choices made for
representatives.

\subsection{Counting the permitted Kummer classes}

Choose a primitive root \(\zeta_\ell\in K\). For each homomorphism
\(t:G_K\to\mathbb F_\ell\) just described, Hilbert's theorem 90
on its finite cyclic quotient gives \(\beta\ne0\) in the field
fixed by \(\ker t\) such that
\(\sigma\beta/\beta=\zeta_\ell^{t(\sigma)}\).
Indeed these roots form a multiplicative cocycle since
\(G_K\) fixes \(\zeta_\ell\). Then \(a=\beta^\ell\in K^\times\).
Changing \(\beta\) by a scalar in \(K^\times\) changes \(a\) by
an \(\ell\)-th power. Conversely, if \(a'/a=c^\ell\) with
\(c\in K^\times\), then \(\beta'/(\beta c)\in\mu_\ell\subset K\),
and the two homomorphisms coincide. For \(\mathfrak p\notin S\)
the field fixed by \(\ker t\) is unramified at \(\mathfrak p\).
At any place above it the normalized integer valuation therefore
restricts to \(v_{\mathfrak p}\). It follows that
\(v_{\mathfrak p}(a)=\ell v(\beta)\) is divisible by \(\ell\).
Thus we have injected the homomorphisms into the finite group
\[
 \mathcal K_S=\{a\in K^\times:
          \ell\mid v_{\mathfrak p}(a)\text{ for }\mathfrak p\notin S\}
                      /K^{\times\ell}.
\]

Let \(U_S\) be the \(S\)-unit group and \(\mathrm{Cl}_S\) the
ideal class group of the ring of \(S\)-integers. Factoring
\((a)=\mathfrak a^\ell\) in that ring gives an exact sequence
\[
 1\longrightarrow U_S/U_S^\ell\longrightarrow \mathcal K_S
       \longrightarrow\mathrm{Cl}_S[\ell]\longrightarrow 1 .
\]
Indeed unique ideal factorization determines \(\mathfrak a\).
Multiplication of \(a\) by \(c^\ell\) changes \(\mathfrak a\) by
\((c)\). A class mapping to the identity can thus be represented by
an \(S\)-unit. An \(\ell\)-th root in \(K^\times\) of an \(S\)-unit
is itself an \(S\)-unit by its valuations. Finally any ideal class
of order dividing \(\ell\) has a representative whose \(\ell\)-th
power is principal, proving surjectivity.

If \(K\) has signature \((r_1,r_2)\), Dirichlet's \(S\)-unit theorem
gives free rank \(r_1+r_2-1+|S|\). The torsion is cyclic and contributes
at most \(\ell\) modulo \(\ell\)-th powers. We have \(r_1+r_2\le n\)
and \(|S|\le n\), the latter by the prime decomposition of \(\ell\).
Also \(\mathrm{Cl}_S\) is a quotient of the ordinary ideal class
group, of order \(h_K\). Hence
\[
      |V|\le|\mathcal K_S|\le\ell^{2n} h_K .
\]

\subsection{The class number and the uniform constant}

We give the estimate on \(h_K\) used here. For any degree-\(n\)
number field of discriminant of absolute value \(\Delta_K\), every
ideal class has an integral representative of norm
\(\le X=\sqrt{\Delta_K}\). Indeed Minkowski's bound supplies the
factor \((4/\pi)^{r_2}n!/n^n\le1\). To verify the last inequality,
\(r_2\le n-1\), and \(2^{n-1}n!\le n^n\) by induction from \(n=1\),
using \((1+1/n)^n\ge2\).

For each rational prime \(p\), fix an ordering of the primes above
it with residue degrees \(f_{p,i}\), and pad the list to length \(n\)
with empty positions. Encode an integral ideal with exponents
\(a_{p,i}\) as an ordered \(n\)-tuple of positive integers. Its
\(i\)-th entry has \(p\)-adic valuation \(f_{p,i}a_{p,i}\) in an
occupied position and zero in an empty one. Prime factorization
recovers the ideal, and the tuple product equals its norm. The number
of tuples with product at most \(X\) is bounded by
\[
 X\left(\sum_{1\le j\le X}\frac1j\right)^{n-1}
       \le X(1+\log X)^{n-1}.
\]
For the first inequality, fix the first \(n-1\) entries and bound
the last by \(X\) divided by their product. For the second, compare
the harmonic sum after its first term with an integral.
Since \(X\ge1\), taking logarithms and increasing \(n-1\) to \(n\)
gives
\[
  \log h_K\le \tfrac12\log\Delta_K+
                       n\log(1+\tfrac12\log\Delta_K).
\]

Now \(K\subseteq\mathbb Q(\mu_M)\) for \(M=\ell F\ge2\).
The conductor--discriminant formula expresses \(\Delta_K\) as the
product of the conductors of the \(n\) characters of
\(\operatorname{Gal}(K/\mathbb Q)\). Each divides \(M\), so
\(\Delta_K\le M^n\). Also \(n\le\varphi(M)\le M\). Since
\(\log M\le M-1\),
\[
   1+\tfrac12\log\Delta_K\le 1+\tfrac12 M(M-1)\le M^2.
\]
The preceding class-number inequality yields
\(\log h_K\le(5/2)n\log M\le3n\log M\).
In particular the bound is uniform in the subfield and the local data.

The group \(V\) is a finite vector space over \(\mathbb F_\ell\).
Its dimension is its \(\mathbb Z_\ell\)-length \(g\), so
\[
 g\log\ell=\log|V|\le2n\log\ell+\log h_K
          \le 5n\log(\ell F).
\]
Convexity of \(u\mapsto1/u\) on \([1,2]\) and its tangent at \(3/2\)
give \(\log2\ge2/3>5/8\). Division by \(\log\ell\ge\log2\) and
\(n\le\ell D\) prove the stated estimate. Its proof applies before
restricting to the eigenspace or imposing the decomposition-group
conditions defining the strict subgroup.
\footnote{\href{https://github.com/Deicyde/riemann-hypothesis/blob/40e56bab026552e089d7afcfeff08d2defb9c005/src/Development/AutomorphicArithmetic/AutomorphicVarieties/OddResidueBounds/Imports.lean}{Assembly of the generator estimate},
\href{https://github.com/Deicyde/riemann-hypothesis/blob/40e56bab026552e089d7afcfeff08d2defb9c005/src/Development/AutomorphicArithmetic/AutomorphicVarieties/OddResidueBounds/UnrInj.lean}{restriction to the cyclotomic layer},
\href{https://github.com/Deicyde/riemann-hypothesis/blob/40e56bab026552e089d7afcfeff08d2defb9c005/src/Development/AutomorphicArithmetic/AutomorphicVarieties/OddResidueBounds/UnrSel.lean}{Kummer injection},
\href{https://github.com/Deicyde/riemann-hypothesis/blob/40e56bab026552e089d7afcfeff08d2defb9c005/src/Development/AutomorphicArithmetic/AutomorphicVarieties/OddResidueBounds/SelCard.lean}{the local conditions and the unit bound},
\href{https://github.com/Deicyde/riemann-hypothesis/blob/40e56bab026552e089d7afcfeff08d2defb9c005/src/Development/AutomorphicArithmetic/CyclotomicBounds/ClassNumberBounds/IdealCounting/ClassC.lean}{Minkowski and ideal counting},
\href{https://github.com/Deicyde/riemann-hypothesis/blob/40e56bab026552e089d7afcfeff08d2defb9c005/src/Development/AutomorphicArithmetic/AutomorphicVarieties/OddResidueBounds/LogH.lean}{the cyclotomic class-number estimate}.}

\section{A cyclotomic annihilator for the signed Selmer group}
\label{sec:selmer-annihilator}

Let \(N\ge1\), let \(E\subset\mathbb Q(\mu_N)\) be a CM field, and put
\(d=[E:\mathbb Q]\), \(\Gamma=\operatorname{Gal}(E/\mathbb Q)\).
Write \(c\) for complex conjugation in \(\Gamma\). Fix a prime
\(\ell\) and an integer \(e\ge1\). We assume either
\(\ell\ne2\), or \(\ell=2\) and \(\mu_4\subset E\).
Put \(k=e+1\), \(\epsilon=(-1)^k\), and
\(A=(\mathbb Q_\ell/\mathbb Z_\ell)(-e)\). The action on \(A\)
is multiplication by \(\chi_\ell^{-e}\), where \(\chi_\ell\) is
the cyclotomic character. Let \(S\) be the set of places of \(E\)
over \(\ell\). With \(G_{E,S}\) the quotient of the absolute
Galois group unramified outside \(S\) at finite places, set
\[
 H=H^1_{\rm cts}(G_{E,S},A),\qquad
 P=\ker\!\left(H\longrightarrow\prod_{v\in S}H^1(E_v,A)\right),
 \qquad P_\epsilon=\{x\in P:cx=\epsilon x\}.
\]
The action of \(\Gamma\) on \(H\) uses both conjugation on the
Galois group and the given cyclotomic action on the coefficients.
These groups are the \(H^1\), \(\mathrm{Sha}\) and
\(\mathrm{ShaMinus}\) of the source for these data.

\paragraph{Annihilator and length bounds.}
There is an integer \(r\ge0\), independent of \(x\in P_\epsilon\),
such that
\begin{equation}
 \ell^r P_\epsilon=0,\qquad
 r\log\ell\le 30d(e+1)\log(e\ell N).
 \label{eq:cyclotomic-kill}
\end{equation}
Consequently \(P_\epsilon\) has finite \(\mathbb Z_\ell\)-length and
\begin{equation}
 \operatorname{length}_{\mathbb Z_\ell}(P_\epsilon)\log\ell
 \le 240\ell d^2(e+1)\log^2(e\ell N).
 \label{eq:cyclotomic-length}
\end{equation}
In particular these inequalities apply at odd primes dividing \(d\).
The separate characterwise estimate in the prime-to-\(d\) case is a
stronger assertion.

\subsection{The weight-one class-group relation}

For \(L>2\) put \(K_L=\mathbb Q(\mu_L)\),
\(G_L=(\mathbb Z/L\mathbb Z)^\times\), and let \(\sigma_b\) send
\(\zeta_L\) to \(\zeta_L^b\). In formulas at level \(L\), the
symbol \([b]_L\) denotes the representative in \([0,L-1]\).
Fix an odd positive integer \(a\) prime to \(L\), and write
\[
 j_L(b)=\frac{a[a^{-1}b]_L-[b]_L}{L},\qquad
 h_L(b)=\frac{a-1}{2}-j_L(b).
\]
Both numbers are integers. For the usual Stickelberger element
\(\vartheta_L=L^{-1}\sum_{b\in G_L}[b]_L\sigma_b^{-1}\), coefficient
comparison gives the identity in \(\mathbb Q[G_L]\)
\[
 \sum_{b\in G_L}j_L(b)\sigma_b^{-1}
   =(a\sigma_a^{-1}-1)\vartheta_L.
\]
The right side belongs to \(\mathbb Z[G_L]\vartheta_L\cap
\mathbb Z[G_L]\). Stickelberger's theorem says that this ideal
annihilates the ideal class group of \(K_L\), for every cyclotomic
modulus \(L\). Also, \(\sum_b\sigma_b^{-1}\) annihilates that class
group, because the product of the conjugates of a fractional ideal
is the extension of its norm to \(\mathbb Q\), a principal ideal.
Subtracting the two relations gives
\begin{equation}
 \sum_{b\in G_L}h_L(b)\sigma_b^{-1}C=0
 \quad\text{for every } C\in\operatorname{Cl}(K_L)
 \label{eq:cyclotomic-class}
\end{equation}
in additive notation. The precise classical input, with these group
ring conventions, is R.~Sharifi, \emph{Iwasawa Theory}, Definition
4.5.1 and Theorem~4.5.6. The same theorem is proved in
L.~Washington, \emph{Introduction to Cyclotomic Fields}, second
edition, Theorem~6.10.

Suppose \(\eta:G_{K_L}\to A\) is a continuous homomorphism with
finite image, written with trivial coefficient action, unramified
at every finite place. The field \(K_L\) is totally imaginary.
The Hilbert class field theorem therefore identifies the finite
unramified abelian quotient through which \(\eta\) factors with a
quotient of \(\operatorname{Cl}(K_L)\). Artin reciprocity is equivariant
for the action of \(G_L\). For \(u\in G_{K_L}\) and any absolute
lifts \(\widetilde\sigma_b\) of \(\sigma_b\), applying \(\eta\) to
the Artin image of \eqref{eq:cyclotomic-class} gives
\begin{equation}
 \sum_b h_L(b)\eta(\widetilde\sigma_b^{-1}u\widetilde\sigma_b)=0.
 \label{eq:cyclotomic-hom}
\end{equation}
Conjugation with the inverse lift on the left gives the inverse
automorphism on the ideal class, as required by \eqref{eq:cyclotomic-class}.

\subsection{Lifting the weight and passing to a fixed group ring}

Set \(M=\ell N\), \(a=2M+1\), and let \(L=M\ell^m\) with \(m\ge0\).
Write \(\pi_L:G_L\to\Gamma\) for restriction. The ordinary Bernoulli
polynomials have \(B_1(X)=X-\tfrac12\). Define
\begin{align*}
 E_k(L,x)&=\frac{L^{k-1}}{k} B_k(x/L),\\
 \Theta(L)&=-\sum_{b\in G_L} E_k(L,[b]_L)\pi_L(b),\\
 Y(L)&=(1-a^k\pi_L(a))\Theta(L),\\
 W(L)&=-\sum_{b\in G_L}[b]_L^{\,e}h_L(b)\pi_L(b).
\end{align*}
Here \(\Theta(L),Y(L)\in\mathbb Q[\Gamma]\) and
\(W(L)\in\mathbb Z[\Gamma]\).
The source evaluates its inverse-indexed group-ring sum on the
inverse of the restriction map. This accounts for \(\pi_L(b)\)
as written in these three expressions.

The two needed identities are
\[
 Y(L)=Y(M),\qquad
 Y(M)-W(L)\in\ell^t\mathbb Z_\ell[\Gamma]
 \quad\bigl(m\ge t+v_\ell(D_k)+1,\ t\ge0\bigr),
 \qquad
 D_k=2k\prod_{p\le k+1}p ,
\]
where the product is over primes. For the equality, every lift of
a unit modulo \(M\) to a residue modulo \(M\ell^m\) is a unit,
since \(\ell\mid M\). Summing \(E_k(L,x)\) over its \(\ell^m\)
representatives gives \(E_k(M,[x]_M)\) by
\[
 \sum_{i=0}^{q-1}B_k\!\left(\frac{u+i}{q}\right)
    =q^{1-k}B_k(u),\qquad q=\ell^m .
\]
This distribution identity follows by summing the generating series
\(te^{ut}/(e^t-1)\). Restriction to \(\Gamma\) is constant on each
fiber, proving the assertion for \(\Theta\) and hence for \(Y\).

For the congruence, put \(P_k=\prod_{p\le k+1}p\).
The binomial expansion of \(B_k(X)\) and the von Staudt--Clausen
theorem imply
\[
 E_k(L,x)=\frac{x^k}{kL}-\frac{x^{k-1}}2+
               \frac{L}{kP_k}z_x,\qquad z_x\in\mathbb Z .
\]
Indeed, the term indexed by \(i\ge2\) in the Bernoulli expansion has
a factor \(L^{i-1}\), and \(P_k B_i\) is integral for \(2\le i\le k\).
Set \(x=[b]_L\), \(y=[a^{-1}b]_L\), \(j=j_L(b)\), so that
\(ay=x+Lj\). For some integers \(u,v\),
\[
 (ay)^k=x^k+kx^{k-1}Lj+L^2u,\qquad
 (ay)^{k-1}=x^{k-1}+Lv.
\]
Substitution yields
\[
 E_k(L,x)-a^kE_k(L,y)-x^{k-1}h_L(b)
 =\frac{L}{D_k}
    \bigl(2z_x-2a^k z_y-2P_k u+akP_k v\bigr).
\]
Multiplying by \(-\pi_L(b)\), summing and using \(e=k-1\) proves
the congruence. Taking \(t=0\) also gives \(Y:=Y(M)\in
\mathbb Z_\ell[\Gamma]\). Thus a single rational group-ring element
is integral at \(\ell\) and is approximated by the explicit \(W(L)\)
to every required precision.

\subsection{The Selmer restriction and cancellation of the twist}

Let \(U=G_{E(\mu_{\ell^\infty})}\). If \(x\in P\), choose a
continuous cocycle \(f:G_E\to A\) inflated from \(G_{E,S}\)
representing \(x\). It has finite image by compactness of \(G_E\)
and discreteness of \(A\). For each \(v\mid\ell\), choose a
decomposition group \(D_v\) in \(G_E\). Local triviality provides
\(t_v\in A\) such that
\(f(g)=(\chi_\ell(g)^{-e}-1)t_v\) on \(D_v\). Choose \(s\ge1\)
so that \(\ell^s\) kills the image of \(f\) and all the finitely
many \(t_v\). On \(U\), \(f\) is a homomorphism \(\kappa\);
coboundaries vanish there. In particular \(\kappa\) depends only
on \(x\).

For every \(m\ge s\), \(K_L=\mathbb Q(\mu_{M\ell^m})\) contains
\(E\) and \(\mu_{\ell^s}\). Put \(V_L=G_{K_L}\).
The restriction \(\eta=f|_{V_L}\) is a homomorphism, since \(V_L\)
acts trivially on the image of \(f\). We show it is unramified.
Away from \(\ell\), the relevant inertia is contained in inertia of
\(E\) killed by \(G_E\to G_{E,S}\), so \(f=0\) there. At \(v\mid\ell\)
the character of \(V_L\cap D_v\) acts trivially on \(t_v\), giving
\(f|_{V_L\cap D_v}=0\). For any \(h\in G_E\) and \(w\in V_L\), the
cocycle rule gives
\[
 f(hwh^{-1})=\chi_\ell(h)^{-e} f(w)
                    +(1-\chi_\ell(w)^{-e})f(h).
\]
The last term vanishes since \(\ell^s f(h)=0\). Conjugating the
chosen decomposition groups therefore covers every place above
\(\ell\) in \(K_L\). This proves the claimed ramification property
of \(\eta\) on the entire finite layer.

If \(u\in U\cap G_{K_M}\), it fixes all \(M\ell^m\)-th roots of
unity. Indeed \(u\) fixes \(\mu_M\) and every \(\ell\)-power root
of unity, and these roots generate \(\mu_{M\ell^m}\).
For \(b\in G_L\) put \(w_b=\widetilde\sigma_b^{-1}u\widetilde\sigma_b\).
Both \(V_L\) and \(U\) are normal under these lifts, so \(w_b\) lies
in their intersection. Thus \(\kappa(w_b)=\eta(w_b)\).
The action on cohomology, evaluated on \(u\), is
\[
 (\pi_L(b)x)|_U(u)=
             \chi_\ell(\widetilde\sigma_b)^{-e}\,\kappa(w_b).
\]
Changing the lift changes the representing cocycle by a coboundary,
which vanishes on \(U\). Also,
\(\chi_\ell(\widetilde\sigma_b)\equiv[b]_L\bmod\ell^s\). Hence
\([b]_L^e\chi_\ell(\widetilde\sigma_b)^{-e}\) acts as the identity on
\(\kappa(w_b)\). Equation~\eqref{eq:cyclotomic-hom} now gives
\[
 (W(L)x)|_U(u)=-\sum_b h_L(b)\eta(w_b)=0 .
\]
Taking \(m\ge s+v_\ell(D_k)+1\) gives \(W(L)x=Yx\), since
\(\ell^s x=0\). We have proved that the homomorphism \((Yx)|_U\)
vanishes on \(U\cap G_{K_M}\). For every \(u\in U\), the power
\(u^{\varphi(M)}\) fixes \(\mu_M\). Thus
\((\varphi(M)Yx)|_U=0\).

We justify the final passage back to global cohomology. Inflation
from \(G_{E,S}\) to \(G_E\) is injective on \(H^1\) because the
coefficient action factors through \(G_{E,S}\). A cocycle on \(G_E\)
whose restriction vanishes on \(U\) descends to
\(J=G_E/U=\chi_\ell(G_E)\) by the cocycle rule. This is an open
subgroup of \(\mathbb Z_\ell^\times\). Indeed its index equals
\([E\cap\mathbb Q(\mu_{\ell^\infty}):\mathbb Q]\), which is finite.
For odd \(\ell\), the group of units is the product of a finite
cyclic group of order prime to \(\ell\) and \(1+\ell\mathbb Z_\ell
\simeq\mathbb Z_\ell\), hence is procyclic. Its closed subgroups
are procyclic. For \(\ell=2\) our hypothesis puts \(J\) in
\(1+4\mathbb Z_2\simeq\mathbb Z_2\). In both cases an element
\(\gamma\) topologically generating \(J\) has infinite order.
The nonzero scalar \(\gamma^{-e}-1\) acts surjectively on the
divisible group \(A\). Given a continuous cocycle \(z\) on \(J\),
choose \(t\) with \(z(\gamma)=(\gamma^{-e}-1)t\). Subtract this
coboundary. The resulting cocycle vanishes on all integral powers
of \(\gamma\), and then on \(J\) by continuity. This proves
\[
             \varphi(M)Y P=0 .
\]

\subsection{The determinant on the signed quotient}

Let \(R=\mathbb Z_\ell[\Gamma]/(c-\epsilon)\). This is a free
\(\mathbb Z_\ell\)-module of rank \(d/2\). In fact \(c\) has order
two since \(E\) is CM and each pair \(g,cg\) contributes the
single basis element represented by \(g\), with \(cg=\epsilon g\).
The same integral basis works over \(\mathbb Z,\mathbb Q,\mathbb C\).
Write \(\Delta\) for the determinant of multiplication by \(Y\) in
this basis. For \(x\in P_\epsilon\), the map \(R\to H\) sending
\(b\mapsto b\varphi(M)x\) is well-defined. The image of \(YR\)
under this map vanishes. By the adjugate identity,
\(\Delta\cdot1\) belongs to \(YR\), whence
\begin{equation}
                  \Delta\varphi(M)x=0 .
 \label{eq:cyclotomic-detkill}
\end{equation}
This argument works integrally at primes dividing \(d\).

Define the integers
\[
 D=Mk(k+1)!,\qquad S_0=2a^kM^k k!,\qquad B=D S_0,\qquad q=d/2 .
\]
The denominators of all \(B_i\) for \(0\le i\le k\) divide
\((k+1)!\) by von Staudt--Clausen, including \(B_1=-1/2\).
Expanding \(E_k\) thus shows \(D E_k(M,x)\in\mathbb Z\) for each
integer \(x\), since the smallest power of \(M\) in the expansion
is \(M^{-1}\). It follows that \(DY\in\mathbb Z[\Gamma]\), and
\(X=D^q\Delta\) is an integer.

We prove \(0<|X|\le B^q\).
Over \(\mathbb C\), the eigenvalues of multiplication by \(Y\) on
the quotient are its evaluations at the characters \(\lambda\) of
\(\Gamma\) with \(\lambda(c)=\epsilon\). This follows by the
finite Fourier transform on the abelian group \(\Gamma\).
Set \(\xi(b)=\lambda(\pi_M(b))\), extended by zero on nonunits.
Then
\[
 \lambda(Y)=(1-a^k\xi(a))\,L(1-k,\xi).
\]
For completeness, the identity for the second factor follows from
\(L(s,\xi)=M^{-s}\sum_{b\in G_M}\xi(b)\zeta(s,[b]_M/M)\) and
\(\zeta(1-k,x)=-B_k(x)/k\). Here \(\zeta(s,x)\) is the Hurwitz
zeta function. If \(\xi^*\) induces \(\xi\) and has conductor \(f\),
then
\[
 L(1-k,\xi)=L(1-k,\xi^*)
   \prod_{\substack{p\mid M\\p\nmid f}}(1-\xi^*(p)p^{k-1}).
\]
Each omitted Euler factor is nonzero as \(k\ge2\). The gamma factor
of the primitive completed \(L\)-function is finite and nonzero at
\(1-k\). Indeed, writing \(\xi^*(-1)=(-1)^\delta\), its argument is
\((1-k+\delta)/2\), an integer plus one half, since
\(\delta\equiv k\bmod 2\). The primitive functional equation
relates the completed value there to its value for
\(\overline{\xi^*}\) at \(k>1\). The latter is finite and nonzero
by the absolutely convergent Euler product. This includes the
primitive trivial character and the zeta functional equation.
Finally \(|a^k\xi(a)|=a^k>1\). Every eigenvalue is therefore nonzero.

The Fourier series of \(B_k\) gives, for \(0\le u\le1\) and \(k\ge2\),
\[
 |B_k(u)|\le \frac{2k!}{(2\pi)^k}\sum_{n\ge1}n^{-k}\le k!.
\]
The sum for \(\Theta(M)\) has at most \(M\) terms, while
\(|1-a^k\xi(a)|\le2a^k\). Hence \(|\lambda(Y)|\le S_0\).
The product of these \(q\) eigenvalues gives
\(0<|X|=D^q|\Delta|\le B^q\).

Now set \(r=v_\ell(|X|)+v_\ell(\varphi(M))\).
Since \(\Delta\in\mathbb Z_\ell\) and \(X=D^q\Delta\),
\(v_\ell(\Delta)\le v_\ell(|X|)\). By
\eqref{eq:cyclotomic-detkill}, \(\ell^r x=0\) for all
\(x\in P_\epsilon\). Furthermore,
\[
 r\log\ell \le q\log B+\log M
                  \le d\log(BM).
\]
For \(T=3e\ell N\), we have \(M,k,2\le T\), \(a\le T^3\), and
\(k+1\le T^2\). Bounding \(k!\) by \(T^k\) and \((k+1)!\) by
\(T^{2(k+1)}\) gives \(BM\le T^{7k+6}\le T^{10k}\) because
\(k\ge2\). Since \(e\ell N\ge2\), \(T\le(e\ell N)^3\).
This proves \eqref{eq:cyclotomic-kill}.

For the length assertion, multiplication by \(\ell^j\) embeds
\(H[\ell^{j+1}]/H[\ell^j]\) into \(H[\ell]\) for every \(j\ge0\).
Section~\ref{sec:selmer-generators} proves
\(\operatorname{length}_{\mathbb Z_\ell}H[\ell]\le
8\ell d\log(\ell N)\). Filtering \(P_\epsilon\subset H[\ell^r]\)
by the successive kernels and applying \eqref{eq:cyclotomic-kill}
gives \eqref{eq:cyclotomic-length}.

\paragraph{Correspondence with the release.}
In the \texttt{OddResidueBounds} development, the class relation is
\texttt{clKill} and its reciprocity consequence is \texttt{homKill}.
The Bernoulli approximation is \texttt{congr} in \texttt{NearLevel}.
The restriction, power and integral annihilators are
\texttt{resKill}, \texttt{totKill} and \texttt{zKill}.
The finite-layer argument above uses the original Selmer cocycle
and its finitely many local trivializations to construct the
unramified character directly. The release obtains the same character
on a sufficiently deep layer by descent of its restriction.
The determinant, logarithmic estimate, exponent and filtration
correspond to \texttt{detVal}, \texttt{bdxLog}, \texttt{killPow}
and \texttt{lenPow}, respectively. The conclusion is the source's
\texttt{StickBound} with \(C=240\).

\section{The characterwise estimate on the minus side}
\label{sec:minus-characterwise}

Let \(E\subseteq\mathbb Q(\mu_N)\) be a CM field. Put
\(D=[E:\mathbb Q]\), \(G=\operatorname{Gal}(E/\mathbb Q)\), and
let \(\ell\nmid 2D\) be prime. Let \(\mathcal O\) be the ring of
integers of a finite extension of \(\mathbb Q_\ell\) in the chosen
algebraic closure, with uniformizer \(\varpi\), valuation
\(v(\varpi)=1\), and \(r_{\mathcal O}=v(\ell)\).
Write \(W=\operatorname{Frac}(\mathcal O)/\mathcal O\). A character
\(\psi:G\to\mathcal O^\times\) pulls back to a character \(\chi\) of
level \(N\). Its order divides \(D\) and is prime to \(\ell\).
Fix \(e\ge1\), put \(j=e+1\), and assume \(\chi(-1)=(-1)^j\).
Write \(\kappa\) for the \(\ell\)-adic cyclotomic character and put
\(\alpha=(\chi\kappa^e)^{-1}\), using the Galois character attached
to \(\chi\). Let \(V(\alpha)\) be the \(\mathcal O\)-length of the
strict subgroup of \(H^1(\mathbb Q,W(\alpha))\), with zero restriction
to all inertia groups and to the decomposition groups at \(\ell\).
The source proves
\begin{equation}
    V(\alpha)\le v\bigl(\alpha_N(j,\chi)\bigr).
    \label{eq:characterwise-bern}
\end{equation}
On the right, \(\alpha_N\) is the integral generalized Bernoulli
number defined in Section~\ref{sec:odd-prime-minus}. We give the
conductor and local conditions, the finite descent, its exceptional
case, and the passage to a bound on the \(E\)-group. We specify below
the cyclotomic derivative-system construction used by the descent.

\subsection{Conductor and local conditions}

Write \(f=\operatorname{cond}\chi\) and \(\xi=\chi^*\). Raising the
level from \(f\) to \(N\) preserves the Galois character, its order
and its parity. By the primitive-value convention,
\[
  \alpha_N(j,\chi)=(N/f)\alpha_f(j,\xi).
\]
The factor \(N/f\) is integral, so it suffices to prove
\eqref{eq:characterwise-bern} at level \(f\).

If \(p\mid f\), \(p\ne\ell\), inertia at \(p\) contains an element
on which \(\xi\) is nontrivial. Indeed the inertia image in the units
modulo \(f\) contains the full \(p\)-component, and triviality on
this component would lower the conductor. A nontrivial root of unity
of order prime to \(\ell\) differs from \(1\) by a unit in
\(\mathcal O\), since reduction is injective on these roots. The
cyclotomic character is unramified away from \(\ell\). Thus at \(p\)
some \(\tau\in I_p\) has \(\alpha(\tau)-1\) a unit. An
inertia coboundary is then a decomposition-group coboundary.
Indeed subtract a local coboundary so the cocycle \(c\) vanishes on
inertia. For \(g\) in the decomposition group, compare \(c(\tau g)\)
with \(c(g(g^{-1}\tau g))\). This gives
\((\alpha(\tau)-1)c(g)=0\), hence \(c(g)=0\).

Choose a lift \(\delta\) of complex conjugation. The parity gives
\(\alpha(\delta)=-1\), so \(\alpha(\delta)-1\) is a unit. Each
class has a unique representative \(c\) satisfying \(c(\delta)=0\):
subtract the coboundary of
\((\alpha(\delta)-1)^{-1}c(\delta)\).
For \(a_0\ne0\), denote by \(S(a_0)\) the module of these normalized
continuous cocycles killed pointwise by \(a_0\), vanishing on inertia
outside \(f\ell\), and coboundaries at the places dividing \(f\ell\)
with local primitives killed by \(a_0\). Every strict class belongs
to the image of one of these modules. The finite image of \(c\) and
one chosen local primitive above \(\ell\) have a common annihilator.
At conjugate places above \(\ell\) the new primitive has the form
\(\alpha(g)m-c(g)\) and is killed by the same annihilator. At a
ramified prime \(p\ne\ell\) the equation
\((\alpha(\tau)-1)m=c(\tau)\) gives the same property by the unit
just found. Outside \(f\ell\), inertia acts trivially, so a local
coboundary is zero. The modules with \(a_0=\varpi^h\) increase with
\(h\). A uniform finite upper bound for their lengths bounds the union:
a longer finite strict chain in the union would have a finite set of
witnesses already contained in one term.

\subsection{The residual cyclotomic exception}

Let \(\overline{\alpha}\) be reduction modulo \(\varpi\). If
\(\overline\alpha=\overline\kappa\), the preceding inertia observation
forces \(\alpha\) to be unramified outside \(\ell\): away from \(\ell\)
its inertia values have order prime to \(\ell\) and reduce to one.
Take \(t\) in inertia at a place over \(\ell\) which acts by inversion
on \(\mu_\ell\). Such \(t\) exists by total ramification of the local
cyclotomic extension. Normalize at \(t\) in this case.
Since \(\alpha(t)-1\) is a unit, vanishing of the cocycle at \(t\)
makes its local primitive at the chosen place zero.

Suppose a normalized cocycle is killed by \(\varpi\). Every
\(\mathbb F_\ell\)-linear coordinate on \(W[\varpi]\) takes it to a
cocycle with values in \(\mathbb F_\ell(1)\): the residual action is
multiplication by \(\overline\kappa\), which lies in the prime field.
Identify this target with \(\mu_\ell\). Kummer theory writes the
coordinate cocycle as \(\sigma\mapsto\sigma\beta/\beta\) with
\(a=\beta^\ell\in\mathbb Q^\times\). Its vanishing on inertia at
a place above each rational prime shows \(\ell\mid v_p(a)\):
\(\beta\) is in the local unramified extension there, whose value
group is \(\mathbb Z\). The sign of \(a\) is an \(\ell\)-th power
as \(\ell\) is odd. Rational prime factorization gives
\(a=b^\ell\) for \(b\in\mathbb Q^\times\).
Now \(\beta/b\in\mu_\ell\). Normalization at \(t\), which acts by
inversion, forces \(\beta/b=1\). Every coordinate of the cocycle
therefore vanishes. They separate points, so the cocycle vanishes.
Applying this argument to successive powers of \(\varpi\) proves
\(S(a_0)=0\) for all \(a_0\). This proves
\eqref{eq:characterwise-bern} in the exceptional branch.

\subsection{The normalized finite Stickelberger value}

Assume \(\overline\alpha\ne\overline\kappa\). For a fixed \(a_0\ne0\)
take \(m\ge\max(1,v_\ell(f))\), and set
\(P=f\ell^{m-v_\ell(f)}\). The sets of primes dividing \(P\) and
\(f\ell\) agree. Choose a lift \(\sigma_t\) of each
\(t\in(\mathbb Z/P\mathbb Z)^\times\). Write \([t]\) for the least
nonnegative representative and choose \(b\) a unit representative
modulo \(P\) coming from a \(g\) with
\(\overline\alpha(g)\ne\overline\kappa(g)\). Then
\(\xi(b)b^{e+1}-1\) is a unit, since multiplication of
\(\alpha(g)-\kappa(g)\) by \(-\alpha(g)^{-1}\) gives the same
quantity modulo \(\varpi\). Define
\[
 h_b(t)=\left\lfloor\frac{b[t]}P\right\rfloor-
        \left\lfloor\frac{b[-t]}P\right\rfloor,\qquad
 z=\sum_t h_b(t)\alpha(\sigma_t)^{-1}.
\]
The level-\(P\) values of \(\alpha(\sigma_t)^{-1}\) are congruent
modulo \(\ell^m\) to \([t]^e\xi([t])\). Write \(S\) for the sum
obtained by this replacement. Then \(|z-S|_\ell\le|\ell^m|_\ell\).

Here is the source's normalization of the other congruence. Put
\(d_j=2j\prod_{p\le j+1}p\) over primes and
\(B=B_{j,\xi_P}\), the generalized Bernoulli number of the
level-\(P\) character induced by \(\xi\). Then
\[
 \left|d_j\left(b^e\xi(b)S
        -2\{\xi(b)b^j-1\}\frac{B}{j}\right)\right|_\ell
        \le|\ell^m|_\ell .
\]
For completeness take \(u=[t]\), \(h=\lfloor bu/P\rfloor\),
\(y=bu-hP\). Both \(u,y\) are units modulo \(P\) and
\(h_b(t)=2h+1-b\). The polynomial
\(C(x)=P^{j-1}B_j(x/P)/j\) satisfies
\[
       C(x)=\frac{x^j}{jP}-\frac{x^{j-1}}2+
                        \frac{2P}{d_j}Z_x,\qquad Z_x\in\mathbb Z.
\]
This follows by the denominator calculation of
Section~\ref{sec:odd-prime-minus} applied to the terms of Bernoulli
index at least two. Expand \(y^j=(bu-hP)^j\) through its linear
term in \(P\), and \(y^{j-1}\) through its constant term.
The leading terms then cancel in
\[
 d_j\{b^e h_b(t)u^e+2(C(y)-b^j C(u))\}\in P\mathbb Z.
\]
Indeed their coefficient of \(b^e u^e\) is
\(h_b(t)-2h-1+b=0\), and \(j\mid d_j\) clears the remaining
binomial denominators. Multiply this relation by \(\xi_P(bt)\),
sum over units, and permute them by \(t\mapsto bt\). This gives
the asserted congruence.

Raising the level from \(f\) to \(P\) adds at most the Euler factor
\(1-\xi(\ell)\ell^e\), of \(\ell\)-adic norm one because \(e\ge1\).
Factors at the primes already dividing \(f\) are one. Consequently
\(|B/j|_\ell=|B_{j,\xi}/j|_\ell>0\). Choose \(m\) large enough that
\[
 |\ell^m|_\ell
     < |a_0|_\ell\,|d_j|_\ell\,|B_{j,\xi}/j|_\ell .
\]
This is possible before choosing \(b\) since the right side is fixed.
The elements \(b^e\xi(b)\), \(2\), and \(\xi(b)b^j-1\) are units.
The strong triangle inequality in the two congruences gives
\(|z|_\ell=|B_{j,\xi}/j|_\ell\). It follows that \(z\ne0\),
\(a_0z\mid\ell^m\), and \(z\mid\alpha_f(j,\xi)\) in \(\mathcal O\):
the last ratio has norm at most one because
\(\alpha_f(j,\xi)=f(\prod_{p\le j+1}p)B_{j,\xi}\).

\subsection{Local reciprocity and the finite descent}

The reciprocity used in the source is a consequence of the
finite-module Poitou--Tate sequence and the unramified-subgroup
theorem in local Tate duality. In the conventions of Milne's
\emph{Arithmetic Duality Theorems}, second edition, these are
I.4.10 and I.2.6 respectively. The local assertion applies at an
auxiliary prime \(q\ne\ell\) to unramified finite coefficients of
\(\ell\)-power order. The global assertion uses the sum of the
local invariant maps on \(H^2(\mu_{\ell^h})\). The archimedean
contribution vanishes for odd \(\ell\).\footnote{\href{https://www.jmilne.org/math/Books/ADTnot.pdf}{Milne, second edition, Chapter I, Theorems 2.6 and 4.10}.}

More explicitly, take a finite coefficient quotient \(R=\mathcal
O/(a)\), killed by \(\ell^h\), and cocycles \(c,n\) with characters
\(\alpha,\kappa\alpha^{-1}\) modulo \(a\). Assume their inertia
restrictions vanish outside \(N'q\), where \(\ell\mid N'\), the
first cocycle is a coboundary at the places dividing \(N'\) away
from \(q\), both characters are trivial modulo \(a\) at \(q\), and
\(c\) vanishes on \(I_q\). For each \(\mathbb Z/\ell^h\)-linear
functional \(\lambda:R\to\mathbb Z/\ell^h\), the cyclotomic
two-cocycle
\[
        (g,h')\longmapsto\lambda(n(g)\alpha(g)c(h'))
\]
has zero local class away from \(q\). At the primes outside the
ramification set it is an unramified cup product. At the primes
in \(N'\) away from \(q\) one factor is a coboundary. The global
invariant sum and injectivity of the local invariant map therefore
give zero local class at \(q\). Let \(\phi\) be a Frobenius at \(q\).
For an unramified \(\mathbb Z/\ell^h\)-valued homomorphism \(x\),
one has \(x(h')c(\phi)=x(\phi)c(h')\) on the decomposition group:
both sides are additive, coincide at \(\phi\), and vanish on
inertia. Consequently the cup product of
\(\lambda(n(\cdot)c(\phi))\) with \(x\) is the preceding zero local
class multiplied by \(x(\phi)\). Unramified orthogonality puts
\(\lambda(n(\cdot)c(\phi))\) in the unramified subgroup. Varying
\(\lambda\) now gives \(n(\tau)c(\phi)=0\) for all \(\tau\in I_q\).
These functionals separate points by decomposing the finite
\(\ell\)-group into cyclic summands. This is the local
annihilation relation used in the descent.

Choose an auxiliary prime \(r\equiv1\bmod P\), and let \(L\) range
over lists of distinct primes \(q\ne r\) with \(q\equiv1\bmod P\)
and \(\ell\nmid\operatorname{ord}_{q}(r)\).
Write \(H=G_{\mathbb Q(\mu_P)}\). The source constructs
\(\nu_L|_H:H\to W[\ell^m]\), continuous additive and equivariant
under conjugation by the dual character \(\kappa\alpha^{-1}\).
Its inertia vanishes away from \(Pr\prod_{q\in L}q\).
For a fixed generator \(w\) of \(W[\ell^m]\), the initial value on a
suitable inertia element at \(r\) is \(zw\) times a unit. At each
admissible new prime \(q\), a suitable inertia value of \(\nu_{qL}\)
is the Frobenius value of \(\nu_L\) times a unit. Here \(qL\) means
the list with \(q\) added. The construction starts with cyclotomic
elements with prescribed Stickelberger ideal factorizations, applies
the multiplicative derivatives
\(\prod_{i=0}^{q-2}(\tau_q^i x)^i\) for generators \(\tau_q\) of the
new cyclotomic factors, and takes Kummer classes. The relations on
these derivatives give the stated initial and transition values.

Write \((A_L)=\operatorname{Ann}_{\mathcal O}(\nu_L(H))\) and
\(J=v(\ell^m)\). The coefficient ideal contains \(\ell^m\), so
\(j_L=v(A_L)\le J\). For \(a_0\mid A_L\) let \(S_L\subset S(a_0)\)
consist of normalized cocycles vanishing on the decomposition groups
at the primes dividing \(r\prod_{q\in L}q\). For \(S_L\ne0\),
choose \(c\ne0\) in it. The source's Chebotarev selection supplies
an admissible \(q\) outside this finite list and a Frobenius
\(\phi\in H\) such that \(c(\phi)\ne0\) and \(\nu_L(\phi)\)
generates the \(\mathcal O\)-span of \(\nu_L(H)\). Thus
the new inertia value has annihilator \((A_L)\).
Since inertia at \(q\) lies in \(H\), we have
\(A_{qL}=A_L t\) for some nonzero \(t\in\mathcal O\), up to a unit.

The equivariant derivative has a global extension with controlled
image. Write \(\beta=\kappa\alpha^{-1}\). The assumption
\(\overline\alpha\ne\overline\kappa\) supplies \(d\in G_{\mathbb Q}\)
with \(\beta(d)-1\) a unit. Given any of the equivariant derivative
homomorphisms \(\nu\) on \(H\), define
\[
 n(g)=-(\beta(d)-1)^{-1}\nu([g,d]),\qquad
 [g,d]=gdg^{-1}d^{-1}.
\]
Every such commutator belongs to \(H\), since \(G_{\mathbb Q}/H\)
is abelian. This function is locally constant by continuity of the
commutator and local constancy of \(\nu\) on \(H\). The identity
\([gh,d]=g[h,d]g^{-1}[g,d]\) and equivariance give
\(n(gh)=n(g)+\beta(g)n(h)\). For \(u\in H\),
\(\nu([u,d])=(1-\beta(d))\nu(u)\), whence \(n(u)=\nu(u)\).
All values of \(n\) are multiples of values in \(\nu(H)\).
In particular the annihilator of \(\nu(H)\) kills the whole cocycle.
Away from \(P\), inertia lies in \(H\), so extension preserves the
required inertia conditions. The same observation identifies the
inertia values in the local comparison at the selected \(q\). This
proves the extension used with the duality relation above.
\footnote{\href{https://github.com/Deicyde/riemann-hypothesis/blob/40e56bab026552e089d7afcfeff08d2defb9c005/src/Development/AutomorphicArithmetic/AutomorphicVarieties/PrimitiveFrobeniusFixed/ImportsAlt.lean}{The commutator formula and its restriction and image properties}.}

To see the resulting cost, choose a generator of \(W[A_{qL}]\)
and write the new inertia value as \(y\) times this generator. If
the value is nonzero, \(v(y)=v(t)\) by its annihilator; if it is zero,
then \(A_L\) is a unit, and \(a_0\mid A_L\) forces \(S_L=0\).
Using the extension of the equivariant derivative to a dual cocycle
in the source's reciprocity step gives \(y\,c'(\phi)=0\) for every
\(c'\in S_L\). Thus \(t\) kills the image of evaluation on \(S_L\)
at \(\phi\), which has length at most \(v(t)\) in \(W\). Its kernel
is contained in \(S_{qL}\): the cocycles vanish on inertia at \(q\)
and the unramified quotient is topologically generated by Frobenius.
Conjugacy propagates the vanishing to the other places above \(q\);
the coefficient character acts trivially modulo \(a_0\) throughout
these decomposition groups. We obtain
\[
   \operatorname{length}S_L+j_L
       \le \operatorname{length}S_{qL}+j_{qL}.
\]
Since some \(c(\phi)\ne0\), \(t\) is a nonunit and \(j_{qL}>j_L\).
Since \(j_L\le J\), iteration ends with a zero submodule. Hence
\(\operatorname{length}S_L+j_L\le J\) on every such list.

The first step at \(r\) begins with the whole \(S(a_0)\).
Put \(k_0=\ell^m/z\). The initial inertia value \(zw\) has
annihilator \((k_0)\), and \(a_0\mid k_0\). The same evaluation gives
\(A_{\varnothing}=k_0t_0\) up to a unit and
\(\operatorname{length}S(a_0)\le
  \operatorname{length}S_{\varnothing}+v(t_0)\).
Combining this with the bound for the empty list gives
\[
  \operatorname{length}S(a_0)
        \le J-v(k_0)=v(z)\le v(\alpha_f(j,\xi)).
\]
This bound is independent of \(a_0\). Taking the increasing union
described above gives \eqref{eq:characterwise-bern}.

\subsection{Return to the minus eigenspace}

Take \(\mathcal O\) large enough to contain all character values
of \(G\). Since \(\ell\nmid D\), restriction to \(G_E\) identifies
the twisted strict group over \(\mathbb Q\) with the
\(\psi\)-eigenspace over \(E\). Indeed positive cohomology of \(G\)
on the \(\ell\)-primary invariants vanishes by averaging over \(G\);
the same restriction and averaging applies to the local subgroups,
whose indices divide \(D\). Untwisting by \(\psi\) produces the
indicated eigenspace. In the sign of complex conjugation under
consideration, the characters are those with
\(\psi(c)=(-1)^{e+1}\). Base change multiplies a finite
\(\mathbb Z_\ell\)-length by \(r_{\mathcal O}\), as is seen on the
simple factors. Equivalently the signed length comparison is
\[
 r_{\mathcal O}m_{\ell,e}
       \le\sum_{\psi(c)=(-1)^{e+1}} V((\psi\kappa^e)^{-1})
\]
in extended lengths; the character bound makes the right side finite.
For any nonzero \(a\in\mathcal O\), the normalized absolute value
satisfies \(v(a)\log\ell=-r_{\mathcal O}\log|a|_\ell\).
Multiplying the displayed inequality by \(\log\ell\), applying
\eqref{eq:characterwise-bern}, and cancelling \(r_{\mathcal O}>0\)
give the inequality used before
\eqref{eq:galois-good-prime-sum}. Distinct characters of \(G\) pull
back to distinct Dirichlet characters because the cyclotomic action
map onto \(G\) is surjective.

Sections~\ref{sec:minus-derivatives} and
\ref{sec:minus-cyclotomic-elements} prove the cyclotomic
construction and the two local derivative identities.
The simultaneous Chebotarev choice is proved in
Section~\ref{sec:minus-frobenius}. Together with the two duality
theorems and the preceding descent these yield
\eqref{eq:characterwise-bern}.
\footnote{\href{https://github.com/Deicyde/riemann-hypothesis/blob/40e56bab026552e089d7afcfeff08d2defb9c005/src/Development/AutomorphicArithmetic/AutomorphicVarieties/PrimitiveFrobeniusFixed/Iv.lean}{Precise input to the characterwise theorem},
\href{https://github.com/Deicyde/riemann-hypothesis/blob/40e56bab026552e089d7afcfeff08d2defb9c005/src/Development/AutomorphicArithmetic/AutomorphicVarieties/PrimitiveFrobeniusFixed/OverIff.lean}{global and local reciprocity},
\href{https://github.com/Deicyde/riemann-hypothesis/blob/40e56bab026552e089d7afcfeff08d2defb9c005/src/Development/AutomorphicArithmetic/AutomorphicVarieties/PrimitiveFrobeniusFixed/ZeNorm.lean}{finite Bernoulli value and Euler factor},
\href{https://github.com/Deicyde/riemann-hypothesis/blob/40e56bab026552e089d7afcfeff08d2defb9c005/src/Development/AutomorphicArithmetic/AutomorphicVarieties/PrimitiveFrobeniusFixed/RatPlace.lean}{the exceptional case},
\href{https://github.com/Deicyde/riemann-hypothesis/blob/40e56bab026552e089d7afcfeff08d2defb9c005/src/Development/AutomorphicArithmetic/AutomorphicVarieties/PrimitiveFrobeniusFixed/NuSysEx.lean}{cyclotomic derivatives},
\href{https://github.com/Deicyde/riemann-hypothesis/blob/40e56bab026552e089d7afcfeff08d2defb9c005/src/Development/AutomorphicArithmetic/AutomorphicVarieties/PrimitiveFrobeniusFixed/AddPowNsmul.lean}{Frobenius selection},
\href{https://github.com/Deicyde/riemann-hypothesis/blob/40e56bab026552e089d7afcfeff08d2defb9c005/src/Development/AutomorphicArithmetic/AutomorphicVarieties/PrimitiveFrobeniusFixed/ChiDvd.lean}{extension to the dual cocycle},
\href{https://github.com/Deicyde/riemann-hypothesis/blob/40e56bab026552e089d7afcfeff08d2defb9c005/src/Development/AutomorphicArithmetic/AutomorphicVarieties/PrimitiveFrobeniusFixed/FrobFix.lean}{length induction},
\href{https://github.com/Deicyde/riemann-hypothesis/blob/40e56bab026552e089d7afcfeff08d2defb9c005/src/Development/AutomorphicArithmetic/AutomorphicVarieties/PrimitiveFrobeniusFixed/PrimCase.lean}{primitive conductor comparisons},
\href{https://github.com/Deicyde/riemann-hypothesis/blob/40e56bab026552e089d7afcfeff08d2defb9c005/src/Development/AutomorphicArithmetic/AutomorphicVarieties/OddResidueBounds/AtPrime.lean}{the signed character sum}.}

\section{Cyclotomic elements with compatible divisors and residues}
\label{sec:minus-cyclotomic-elements}

We construct the elements \(X_n\) used in
Section~\ref{sec:minus-derivatives}. Keep its integers \(P,r,b\),
where \(0\le b<P\), and fix a place above \(r\) in
\(\overline{\mathbb Q}\). For a positive integer \(a\) put
\(F_a=\mathbb Q(\mu_a)\), \(G_a=(\mathbb Z/a\mathbb Z)^\times\),
and let \(\sigma_t\) act on roots of unity by \(t\). Write
\([t]_a\) for the least nonnegative residue, and put
\[
   h_a(t)=\left\lfloor b[t]_a/a\right\rfloor-
            \left\lfloor b[-t]_a/a\right\rfloor,\qquad
   \Theta_a=\sum_{t\in G_a} h_a(t)\sigma_t^{-1}.
\]
Use multiplicative group-ring notation \(I^{\Theta_a}\) for ideals
and \(y^{\Theta_a}\) for elements. For a nonzero integral ideal
\(I\) in \(F_a\), seek an element \(X_a(I)\) satisfying
\[
       X_a(I)\overline{X_a(I)}=1,\qquad
       (X_a(I))=I^{\Theta_a}.
\]
Here the bar denotes complex conjugation. Any two choices differ by
a root of unity of order dividing \(2a\). Indeed, their quotient
is an algebraic unit of absolute value one in every embedding,
since complex conjugation commutes with those embeddings.
Kronecker's theorem makes it a root of unity, and the roots of
unity of \(F_a\) have exponent dividing \(2a\).
For \(I=(y)\) we can take \(y^{\Theta_a}\), by
\(h_a(-t)=-h_a(t)\). The two identities are multiplicative in \(I\)
and \(X\), and allow division in the group of fractional ideals.

\subsection{Existence and the Jacobi sum divisor}

Every ideal class of \(F_a\) contains a prime ideal of residue
degree one over \(\mathbb Q\), avoiding any prescribed finite set.
Indeed, the Hilbert class field of \(F_a\) is Galois over
\(\mathbb Q\), since \(F_a/\mathbb Q\) is Galois. Apply Chebotarev
over \(\mathbb Q\) to the Artin element of the class, an element
of the Galois group over \(F_a\). Its restriction to \(F_a\) is
trivial. Choosing the prime above the resulting rational prime to
give that exact Artin element proves the assertion by Hilbert class
field reciprocity. Choose in this manner a degree-one prime
\(\mathfrak p\) in the inverse class of \(I\) over \(p\nmid a\).
If \(I\mathfrak p=(y)\), dividing \(y^{\Theta_a}\) by a choice
at \(\mathfrak p\) gives the desired choice at \(I\).

The reduction map \(\pi:\mathcal O_{F_a}\to\mathbb F_p\) at
\(\mathfrak p\) sends \(\mu_a\) bijectively to the \(a\)-th roots
of unity in \(\mathbb F_p\), so \(a\mid p-1\). Write
\(d=(p-1)/a\), and lift \(x\mapsto x^d\) uniquely to a character
\(\chi:\mathbb F_p^\times\to\mu_a\). Extend every multiplicative
character used in sums by zero at zero. Set
\[
 J_j=\sum_{x\in\mathbb F_p}\chi(x)\chi^j(1-x),\qquad
 J=\prod_{1\le j<b}J_j,\qquad X_a(\mathfrak p)=\overline J/J .
\]
For \(1\le j<b<a\), all of \(\chi,\chi^j,\chi^{j+1}\) are
nontrivial. In any complex embedding the Gauss sums
\(G(\lambda)=\sum_x\lambda(x)e^{2\pi i x/p}\) have
\(\lvert G(\lambda)\rvert^2=p\) for nontrivial \(\lambda\):
put \(x=ty\) in the squared absolute value and sum
\(e^{2\pi i(t-1)y/p}\) over \(y\ne0\), obtaining \(p-1\) at
\(t=1\) and \(-1\) elsewhere. Moreover
\(G(\chi)G(\chi^j)=J_jG(\chi^{j+1})\), by grouping according
to \(x+y\); the terms with \(x+y=0\) sum to zero.
Thus \(J_j\overline J_j=p\). Each \(J_j\) is integral and its
ideal divides the squarefree ideal \((p)\).

For \(t\in G_a\) let \(A=d[t]_a,\ B=d[jt]_a\), both between
\(1\) and \(p-2\). The reduction of \(\sigma_t J_j\) at
\(\mathfrak p\) is \(\sum_x x^A(1-x)^B\).
Expanding the second power and using the power sums of
\(\mathbb F_p^\times\) gives zero if \(A+B<p-1\). Otherwise the
only contributing positive multiple of \(p-1\) is \(p-1\), and
the coefficient is
\(-(-1)^{p-1-A}\binom{B}{p-1-A}\ne0\bmod p\).
The primes \(\sigma_t^{-1}\mathfrak p\) are distinct and exhaust
the primes above \(p\). Hence the exponent of
\(\sigma_t^{-1}\mathfrak p\) in \(J_j\) is
\[
       1-\left\lfloor\big([t]_a+[jt]_a\big)/a\right\rfloor .
\]
The floor is the increment in \(\lfloor j[t]_a/a\rfloor\)
when \(j\) increases to \(j+1\). For \(b\ge1\) telescoping
gives the exponent of \(J\) as \(b-1-\lfloor b[t]_a/a\rfloor\).
Consequently \(\overline J/J\) has the required exponent
\(h_a(t)\) at this prime. At \(b=0\) both the required
exponents and the empty-product exponents are zero. This proves
existence, including the divisor and the archimedean normalization.

\subsection{Norm from a new prime layer}

Fix \(a=Pn\) and an admissible \(q\nmid a\), with \(0\le b<P<q\).
Write \(E=F_a,\ E'=F_{qa},\ u=q\bmod a\). Restriction of the group
ring satisfies
\begin{equation}
       \left.\Theta_{qa}\right|_{G_a}
                  =(1-\sigma_u^{-1})\Theta_a .
 \label{eq:cyclotomic-distribution}
\end{equation}
To see this coefficientwise, the \(q\) representatives lifting
\(t\bmod a\) are \([t]_a+ia\) for \(0\le i<q\). Exactly one,
\(q[u^{-1}t]_a\), is excluded because it is divisible by \(q\).
For \(b>0\), \(ib\bmod q\) permutes the residues modulo \(q\);
summing the floors of
\(b([t]_a+ia)/(qa)\) over all \(i\) therefore gives
\(\lfloor b[t]_a/a\rfloor+\sum_{i<q}\lfloor ib/q\rfloor\).
Indeed the fractional part of \(b[t]_a/a\), divided by \(q\),
does not affect any of these floors, and adding its integral
part permutes the residues in the numerator. Subtract the missing
term \(\lfloor b[u^{-1}t]_a/a\rfloor\), then subtract the same
calculation with \(-t\). The common constant cancels, giving
\eqref{eq:cyclotomic-distribution}. For \(b=0\) it holds identically.

For a nonzero ideal \(I'\) of \(E'\), put \(I=N_{E'/E}I'\).
Norm commutes with conjugation and products of fractional ideals.
Thus the divisor of \(N X_{qa}(I')\) is
\(I^{(1-\sigma_u^{-1})\Theta_a}\). The latter is also the divisor
of \(X_a(I)/\sigma_u^{-1}X_a(I)\). Both elements have product
one with their complex conjugates. Their ratio is therefore an
element of \(\mu_{2a}\), by the uniqueness argument.

\subsection{Simultaneous reduction at every prime over \(q\)}

Suppose \(I'\) is prime to \(q\). We prove, with a single root
\(\eta\in\mu_{2qa}\), the congruence
\[
          X_{qa}(I')\equiv\eta\,\sigma_u^{-1}X_a(I)
                       \pmod{\mathfrak q}
          \qquad(\mathfrak q\mid q\ \text{in }E').
\]
Both sides are units at all such places by their divisors.

First take a degree-one prime \(\mathfrak p'\) above \(p\nmid qa\).
Its relative ideal norm is the contracted prime \(\mathfrak p\).
Use the Jacobi choices already constructed, with characters
\(\chi'\) at level \(qa\) and \(\chi\) at level \(a\), defined by
reduction at \(\mathfrak p'\) and \(\mathfrak p\) respectively.
For \(x\ne0\), \(\chi'(x)^q=\chi(x)\): these are \(a\)-th roots
with equal reductions modulo \(\mathfrak p'\). Also
\((\sigma_u^{-1}\chi(x))^q=\chi(x)\). In the residue field of
any \(\mathfrak q\mid q\), the \(q\)-power map is injective, so
\(\chi'(x)\equiv\sigma_u^{-1}\chi(x)\pmod{\mathfrak q}\).
The same argument applies after complex conjugation. The two pairs
of Jacobi products therefore have equal reductions, term by term.
Each product and its conjugate is a unit at \(q\), since its
factors divide \(p\). Dividing proves the congruence with
\(\eta=1\) for the Jacobi choices. Changing either choice multiplies
the ratio by a single global root of order dividing \(2qa\).

For a principal ideal \((y)\), with \(y\) integral and prime to
\(q\), use the explicit principal generators. Fix
\(\mathfrak q\mid q\) and write \(f(t)\) for the reduction of
\(\sigma_{\tilde t}^{-1}y\) at this place, for any lift
\(\tilde t\in G_{qa}\) of \(t\). This is independent of the lift
because \(\operatorname{Gal}(E'/E)\) fixes residues at \(q\):
\(\mu_q\) reduces to \(1\) and
\(\mathcal O_{E'}=\mathbb Z[\mu_{qa}]\). A lift of \(\sigma_u\)
acts on these residues by \(q\)-power, as this is its action on
the reductions of \(\mu_a\) and \(\mu_q\). Consequently
\(f(ut)^q=f(t)\). All \(f(t)\) are nonzero. Using
\eqref{eq:cyclotomic-distribution}, the reduction of
\(y^{\Theta_{qa}}\) is
\[
 \prod_{t\in G_a} f(t)^{h_a(t)-h_a(u^{-1}t)}
    =\prod_{t\in G_a}\big(f(ut)^{q-1}\big)^{h_a(t)}.
\]
Each of the \(q-1\) conjugates entering \(N_{E'/E}y\) has the
same reduction after applying \(\sigma_{ut}^{-1}\). Hence the
last product is exactly the reduction of
\(\sigma_u^{-1}((Ny)^{\Theta_a})\).
The equality holds at every place above \(q\).

For general \(I'\) prime to \(q\), choose a degree-one
\(\mathfrak p'\) in its inverse class with \(p\nmid qa\), and write
\(I'\mathfrak p'=(y)\). Apply the principal calculation to
\(y\) and the prime calculation to \(\mathfrak p'\). Divide
the two congruences, using units at each place, and apply uniqueness
of the choices at both levels. The resulting root \(\eta\) is a
product of the same globally chosen roots at every place, as required.

\subsection{The family at the fixed prime}

Let \(\mathfrak r_a\) be the prime cut out on \(F_a\) by our
fixed place at \(r\). For every admissible product \(n\) take
\(X_n=X_{Pn}(\mathfrak r_{Pn})\); arbitrary choices suffice.
Since \(\mathfrak r_{a}\) lies above \(r\), their divisors show
that they are units away from \(r\). For \(a=P\), \(r\equiv1\bmod P\)
implies complete splitting and unramifiedness. Thus for every
\(t\in G_P\), the valuation of \(\sigma_t X_1\) at the fixed
place, normalized by \(v(r)=1\), is precisely \(h_P(t)\).

At the layer \(a\mapsto qa\) put
\(s=[k(\mathfrak r_{qa}):k(\mathfrak r_a)]\).
The norm of \(\mathfrak r_{qa}\) is \(\mathfrak r_a^s\), for
which \(X_n^s\) is a normalized generator. The integer \(s\)
divides the absolute residue degree
\(\operatorname{ord}_{qa}(r)\), since \(r\nmid qa\).
By the Chinese remainder theorem, the latter divides the product
of the orders of \(r\bmod p\) over primes \(p\mid qn\), as
\(r\equiv1\bmod P\) and those primes are distinct and coprime
to \(P\). Every factor is prime to \(\ell\) by admissibility,
so \(\ell\nmid s\). Applying the norm and simultaneous reduction
just proved to \(X_{qn}\) and \(X_n^s\) gives exactly the two input
relations in Section~\ref{sec:minus-derivatives}.
\footnote{\href{https://github.com/ChristianKudeba/alpoge-rh-verification/blob/alpoge-source-7cbf7a0c/src/Development/AutomorphicArithmetic/AutomorphicVarieties/PrimitiveFrobeniusFixed/XJacPowNe.lean}{Jacobi sums and their prime divisors},
\href{https://github.com/ChristianKudeba/alpoge-rh-verification/blob/alpoge-source-7cbf7a0c/src/Development/AutomorphicArithmetic/AutomorphicVarieties/PrimitiveFrobeniusFixed/XEx.lean}{passage to arbitrary ideals},
\href{https://github.com/ChristianKudeba/alpoge-rh-verification/blob/alpoge-source-7cbf7a0c/src/Development/AutomorphicArithmetic/AutomorphicVarieties/PrimitiveFrobeniusFixed/ThDist.lean}{distribution of the signed floor coefficients},
\href{https://github.com/ChristianKudeba/alpoge-rh-verification/blob/alpoge-source-7cbf7a0c/src/Development/AutomorphicArithmetic/AutomorphicVarieties/PrimitiveFrobeniusFixed/XNormDiv.lean}{norm relation},
\href{https://github.com/ChristianKudeba/alpoge-rh-verification/blob/alpoge-source-7cbf7a0c/src/Development/AutomorphicArithmetic/AutomorphicVarieties/PrimitiveFrobeniusFixed/XCongJac.lean}{Jacobi congruence},
\href{https://github.com/ChristianKudeba/alpoge-rh-verification/blob/alpoge-source-7cbf7a0c/src/Development/AutomorphicArithmetic/AutomorphicVarieties/PrimitiveFrobeniusFixed/XCongPrin.lean}{principal congruence},
\href{https://github.com/ChristianKudeba/alpoge-rh-verification/blob/alpoge-source-7cbf7a0c/src/Development/AutomorphicArithmetic/AutomorphicVarieties/PrimitiveFrobeniusFixed/AlgStep.lean}{division to obtain simultaneous congruences},
\href{https://github.com/ChristianKudeba/alpoge-rh-verification/blob/alpoge-source-7cbf7a0c/src/Development/AutomorphicArithmetic/AutomorphicVarieties/PrimitiveFrobeniusFixed/PlaceEmb.lean}{assembly at the chosen place}.}

\section{Cyclotomic derivatives at the initial and auxiliary primes}
\label{sec:minus-derivatives}

We establish the two local relations for \(\nu_L\) used in
Section~\ref{sec:minus-characterwise}. Fix \(M=\ell^m\), \(m\ge1\),
\(\ell\) odd, \(P=N_0M\) with \(\ell\nmid N_0\), and a prime
\(r\equiv1\bmod P\). Put \(K_n=\mathbb Q(\mu_{Pn})\) and
\(H=G_{K_1}\). Each index \(n\) is the product of a list \(L\) of
distinct admissible primes, in the sense of the preceding section.
Let \(\alpha\) be the character there and \(\beta=\kappa\alpha^{-1}\).
On \(H\), \(\alpha\equiv1\bmod M\). Choose
\(\delta\) with \(\alpha(\delta)-\kappa(\delta)\) a unit, and an
integer \(a\) giving the action of \(\delta\) on \(\mu_M\).
Thus \(\alpha(\delta)-a\) is a unit. Write \(T(v)=\delta(v)/v^a\).
The map \(T\) is multiplicative and kills \(\mu_M\).

\subsection{The input cyclotomic elements and the derivatives}

The cyclotomic input supplies \(X_n\in K_n^\times\), units away
from \(r\), and integers \(s(n,q)\ge1\) prime to
\(\ell\) for each admissible extension \(n\mapsto qn\).
For \(\tau_q\) acting as a generator on \(\mu_q\) and fixing the
prime-to-\(q\) cyclotomic factors, and \(\phi\) acting on \(K_n\)
by the exponent \(q\), the input says
\[
  \operatorname N_q X_{qn}
    =\eta\,(X_n/\phi^{-1}X_n)^{s(n,q)},\qquad
  X_{qn}\equiv\eta'\,(\phi^{-1}X_n)^{s(n,q)}\quad\bmod\mathfrak q
\]
at all places \(\mathfrak q\mid q\). Here \(\operatorname N_q\)
denotes the norm \(K_{qn}/K_n\). Both error factors are roots
of unity of order dividing \(2Pqn\), and the same \(\eta'\) works
for all these places. At a fixed place \(A_0\mid r\), normalized by
\(v_{A_0}(r)=1\), the prescribed initial divisor is
\(v_{A_0}(\sigma_t X_1)=h_b(t)\) for every lift \(\sigma_t\) of
\(t\in(\mathbb Z/P\mathbb Z)^\times\), with \(h_b\) as in the
preceding section. This is the input called \texttt{laxEx}. Its
construction uses cyclotomic ideal factorizations, relative norms
and congruences of chosen generators.

Put \(x_n=T(X_n)^{2N_0n}\) and \(c(n,q)=q\,s(n,q)\).
If \(\eta^{2Pqn}=1\), then \(\eta^{2N_0qn}\in\mu_M\), so
\(T(\eta)^{2N_0qn}=1\). Norm and the automorphism \(\delta\) commute
on these abelian cyclotomic fields. Conjugating a place above \(q\)
gives another place above \(q\), so \(T\) preserves the congruences
with the indicated transformed error. Applying \(T\) and raising
to \(2N_0qn\) gives
\begin{equation}
 \operatorname N_q x_{qn}=(x_n/\phi^{-1}x_n)^{c(n,q)},\qquad
 x_{qn}\equiv(\phi^{-1}x_n)^{c(n,q)}\pmod{\mathfrak q}.
 \label{eq:deriv-norm}
\end{equation}
The exponents here explain the factor \(q\) in the final transition
unit. Since \(q\equiv1\bmod P\), it is prime to \(\ell\).

For \(v\in K_{qn}^\times\) set
\(D_q(v)=\prod_{i=0}^{q-2}(\tau_q^iv)^i\). Let \(D_L\) be the
composition over \(L\) and set \(y_L=D_L(x_n)\). All these
operators commute with automorphisms on the cyclotomic field
and with each other. Telescoping the consecutive powers and using
\(\tau_q^{q-1}=1\) there gives
\begin{equation}
 \frac{\tau_q(D_q v)}{D_q v}
       =\frac{v^{q-1}}{\operatorname N_q v}.
 \label{eq:deriv-telescoping}
\end{equation}
For every \(g\in H\), \(g(y_L)/y_L\) is an \(M\)-th power in
\(K_n^\times\). To verify this, use induction on \(|L|\).
The empty-list element belongs to \(K_1\). For \(q\in L\) write
\(n=qn'\) and \(L'=L\setminus\{q\}\). Equations \eqref{eq:deriv-norm} and \eqref{eq:deriv-telescoping} give
\[
 \frac{\tau_q(y_L)}{y_L}
  =D_{L'}(x_n)^{q-1}
       \left(\frac{\phi^{-1}y_{L'}}{y_{L'}}\right)^{c(n',q)} .
\]
The first factor is an \(M\)-th power because \(M\mid q-1\).
For the second use the induction hypothesis since \(\phi^{-1}\)
fixes \(K_1\). The restrictions of all these \(\tau_q\) generate
\(\operatorname{Gal}(K_n/K_1)\) by the Chinese remainder theorem.
The elements \(g\) satisfying the assertion form a subgroup, since
\(K_n\) is Galois and the \(M\)-th powers are preserved by every
automorphism. This proves the assertion.

\subsection{The Kummer classes and their projection}

Choose a primitive \(M\)-th root \(\zeta\), and \(Y\) with \(Y^M=y_L\).
For \(g\in H\) take \(v_g\in K_n\) with \(v_g^M=g(y_L)/y_L\).
Define \(\kappa_L(g)\in\mathbb Z/M\) by
\begin{equation}
 \zeta^{\kappa_L(g)}
       =\frac{g(TY)}{TY\,T(v_g)} .
 \label{eq:deriv-kummer}
\end{equation}
The right side has \(M\)-th power one since \(g\) and \(\delta\)
commute on \(K_n\). Changing either root changes it by an element of
\(\mu_M\), which \(T\) kills. Taking
\(v_{gh}=g(v_h)v_g\) proves additivity, since \(T(gv_h)=g(Tv_h)\)
and \(g\) fixes \(\mu_M\). Local constancy follows by restricting to
a finite extension containing \(Y,\delta Y\) and the finitely many
roots \(v_g\) required for the finite possible actions on \(K_n\).
At a prime \(p\nmid Prn\), inertia fixes \(K_n\). The element \(y_L\)
is a \(p\)-adic unit and \(p\nmid M\), so it has an \(M\)-th root
in the local unramified extension, by reduction and Hensel's lemma.
The same holds for \(\delta y_L\) at this place. Thus the chosen
roots are unramified there, and \eqref{eq:deriv-kummer} gives zero on inertia.

Let \(w\) be a generator of the cyclic module \(W[M]\). Using the
lifts fixed in the preceding chapter, set
\begin{equation}
    \nu_L(h)=\sum_t \kappa_L(\sigma_t^{-1}h\sigma_t)
                                 \,\beta(\sigma_t)w,\qquad h\in H .
 \label{eq:deriv-projection}
\end{equation}
This is independent of representatives for the residues of
\(\kappa_L\), since \(Mw=0\). It is locally constant, additive on
\(H\), killed by \(M\), and has the inertia vanishing just established.
Conjugation by \(g\in G_{\mathbb Q}\) multiplies \eqref{eq:deriv-projection} by \(\beta(g)\):
reindex the sum through the action of \(g\) modulo \(H\).
Conjugation by elements of \(H\) preserves the additive homomorphism
\(\kappa_L|_H\), and \(\beta|_H\equiv1\bmod M\).

\subsection{The initial value at \(r\)}

Take \(\rho^M=r\) and choose \(t_0\in I_{A_0}\) with
\(t_0(\rho)=\zeta\rho\). Such an element exists since \(r\nmid M\)
and the \(M\)-th root of the uniformizer generates a totally tamely
ramified extension over the unramified field containing \(\mu_M\).
In the empty-list case \(y=x_1\) belongs to \(K_1\). Write
\(K(t)=v_{A_0}(\sigma_t T(y))\); these are integers. Let \(Y^M=y\)
and \(W_t=\sigma_t T(Y)\). Then \(v_{A_0}(W_t)=K(t)/M\), so
\(W_t/\rho^{K(t)}\) is a unit. The \(M\)-th power of this unit is
fixed by \(t_0\), since \(T(y)\) and its conjugates belong to the
unramified cyclotomic field at \(r\). Its quotient under \(t_0\) is
therefore an \(M\)-th root of unity reducing to one. It equals one,
as \(r\nmid M\). Thus \(t_0(W_t)/W_t=\zeta^{K(t)}\). In \eqref{eq:deriv-kummer} the
element \(v_g\) for \(g=\sigma_t^{-1}t_0\sigma_t\) has \(M\)-th
power one. Applying \(\sigma_t\) to \eqref{eq:deriv-kummer} gives
\[
       [t]\kappa_{\varnothing}(\sigma_t^{-1}t_0\sigma_t)
                         =K(t)\pmod M .
\]

Define \(\mathcal Z(k)=\sum_t k(t)\alpha(\sigma_t)^{-1}\).
Because \(\kappa(\sigma_t)\equiv[t]\bmod M\) for the cyclotomic
character \(\kappa\), equation \eqref{eq:deriv-projection} gives
\(\nu_{\varnothing}(t_0)=\mathcal Z(K)w\).
If \(\delta\) corresponds to \(d\in(\mathbb Z/P\mathbb Z)^\times\),
the operator \(T\) transforms the divisor vector \(k(t)\) into
\(k(td)-a k(t)\). Permuting \(t\mapsto td\) shows that \(\mathcal Z\)
of this new vector equals
\((\alpha(\delta)-a)\mathcal Z(k)\) modulo \(M\), since
\(\alpha|_H\equiv1\bmod M\). Now \(T(y)=T(T(X_1)^{2N_0})\).
Apply this transformation twice and the multiplicity \(2N_0\) once,
starting from \(h_b(t)\). As \(z=\mathcal Z(h_b)\),
\begin{equation}
      \nu_{\varnothing}(t_0)
             =2N_0(\alpha(\delta)-a)^2 zw .
 \label{eq:deriv-initial}
\end{equation}
The multiplier is a unit because \(\ell\nmid2N_0\) and
\(\alpha(\delta)-a\) is a unit. This proves the initial relation with
its actual coefficient.

\subsection{The transition from Frobenius to inertia}

Let \(q\notin L\) be admissible, \(n=\prod_{p\in L}p\),
\(c=c(n,q)\) and \(k=(q-1)/M\). At a fixed place over \(q\) take
an arithmetic Frobenius \(\phi\) and \(t_q\) in inertia whose action
on \(\mu_q\) agrees with \(\tau_q\). Cyclotomic total ramification
at \(q\) supplies \(t_q\). The actions of \(t_q\) and \(\tau_q\)
agree on \(K_{qn}\) since both fix \(K_n\). Put
\[
       y=D_L(x_n),\qquad u=D_L(x_{qn}),\qquad y'=D_q u .
\]
At \(q\) all these are units. Applying \(D_L\) to \eqref{eq:deriv-norm} gives
\(\operatorname N_q u=(y/\phi^{-1}y)^c\) and
\(u\equiv(\phi^{-1}y)^c\) at every place over \(q\).
Take \(Y^M=y'\), \(V^M=y\), and
\(b_0\in K_n\) with \(b_0^M=(\phi^{-1}y)/y\), using the
\(M\)-th power assertion above and \(\phi\in H\).
By \eqref{eq:deriv-telescoping},
\[
       (u^k b_0^c)^M=t_q(y')/y'.
\]
Thus \eqref{eq:deriv-kummer} computes \(\kappa_{qL}(t_q)\) with \(u^k b_0^c\) as its
adjusting root. At \(\phi^{-1}\), the corresponding root for
\(\kappa_L\) is \(b_0\). Dividing the two formulas gives the identity
\begin{equation}
 \frac{\zeta^{\kappa_{qL}(t_q)}}{(\zeta^{\kappa_L(\phi^{-1})})^c}
  =\frac{t_q(TY)}{TY}
     \left(\frac{TV}{(\phi^{-1}TV)^q}\right)^c
      \bigg/ T\!\left(\frac{u}{(\phi^{-1}y)^c}\right)^k .
 \label{eq:deriv-transition-id}
\end{equation}
For the algebra here, \((\phi^{-1}TV)^M=T(\phi^{-1}y)\) since
\(\delta\) and \(\phi\) commute on \(K_n\), and \(q=1+kM\).
The factors \(T(b_0)^c\) in the two expressions cancel.

Every factor on the right of \eqref{eq:deriv-transition-id} reduces to one. For the first use
inertia on the unit \(TY\). For the second use that Frobenius acts
as the \(q\)-power on reductions, applied to the unit
\(\phi^{-1}TV\). For the denominator use the congruence for \(u\)
at every place over \(q\), which also applies at the place conjugate
by \(\delta\). The left side is an \(M\)-th root of unity and \(q\nmid
M\), so injectivity of reduction on \(\mu_M\) makes it one.
Additivity on \(H\) gives
\[
            \kappa_{qL}(t_q)=-c\,\kappa_L(\phi)\pmod M .
\]
This computation applies to
\((\sigma_t^{-1}t_q\sigma_t,\sigma_t^{-1}\phi\sigma_t)\) at the
conjugate place: the new inertia element still agrees with
\(\tau_q\) on \(\mu_q\), because the cyclotomic Galois group is
abelian. Sum the congruences with the weights in \eqref{eq:deriv-projection}. The result is
\begin{equation}
                 \nu_{qL}(t_q)=-q\,s(n,q)\,\nu_L(\phi).
 \label{eq:deriv-transition}
\end{equation}
The factor is a unit because \(\ell\nmid q s(n,q)\). This proves
the transition used in the length induction.

Section~\ref{sec:minus-cyclotomic-elements} constructs the input
\(X_n\), with the prescribed divisor and the same root of unity in
the congruence at every prime above \(q\).
\footnote{\href{https://github.com/ChristianKudeba/alpoge-rh-verification/blob/alpoge-source-7cbf7a0c/src/Development/AutomorphicArithmetic/AutomorphicVarieties/PrimitiveFrobeniusFixed/PlaceEmb.lean}{Cyclotomic elements and divisor data},
\href{https://github.com/ChristianKudeba/alpoge-rh-verification/blob/alpoge-source-7cbf7a0c/src/Development/AutomorphicArithmetic/AutomorphicVarieties/PrimitiveFrobeniusFixed/SysOfLax.lean}{removal of root-of-unity factors},
\href{https://github.com/ChristianKudeba/alpoge-rh-verification/blob/alpoge-source-7cbf7a0c/src/Development/AutomorphicArithmetic/AutomorphicVarieties/PrimitiveFrobeniusFixed/DerInv.lean}{invariance modulo powers},
\href{https://github.com/ChristianKudeba/alpoge-rh-verification/blob/alpoge-source-7cbf7a0c/src/Development/AutomorphicArithmetic/AutomorphicVarieties/PrimitiveFrobeniusFixed/KolHom.lean}{Kummer classes},
\href{https://github.com/ChristianKudeba/alpoge-rh-verification/blob/alpoge-source-7cbf7a0c/src/Development/AutomorphicArithmetic/AutomorphicVarieties/PrimitiveFrobeniusFixed/KolInit.lean}{initial local computation},
\href{https://github.com/ChristianKudeba/alpoge-rh-verification/blob/alpoge-source-7cbf7a0c/src/Development/AutomorphicArithmetic/AutomorphicVarieties/PrimitiveFrobeniusFixed/KolStep.lean}{transition local computation},
\href{https://github.com/ChristianKudeba/alpoge-rh-verification/blob/alpoge-source-7cbf7a0c/src/Development/AutomorphicArithmetic/AutomorphicVarieties/PrimitiveFrobeniusFixed/NuSysEx.lean}{projection to the dual character and the two units}.}

\subsection{One prime for the two Frobenius conditions}
\label{sec:minus-frobenius}

Write \(U=G_{\mathbb Q(\mu_P)}\), \(W=\operatorname{Frac}(\mathcal O)/\mathcal O\),
and let \(\varpi\) be a uniformizer of \(\mathcal O\).
Let \(c,\nu:U\to W[M]\) be continuous additive homomorphisms such
that \(c(gug^{-1})=\alpha(g)c(u)\) and
\(\nu(gug^{-1})=\beta(g)\nu(u)\) for \(u\in U\) and
\(g\in G_{\mathbb Q}\). Assume \(c\) is nonzero on \(U\).
Given any finite set of excluded primes, we shall select an
admissible \(q\) outside the set and a Frobenius \(\phi\in U\) at
a place above \(q\) for which
\[
    c(\phi)\ne0,\qquad \mathcal O\,\nu(\phi)=
          \mathcal O\,\nu(U).
\]
Here \(\mathcal O\,\nu(U)\) means the submodule generated by all
the values. In the application \(c\) is a nonzero normalized strict
cocycle restricted to \(U\). Indeed, if the restriction were zero,
the cocycle, with \(\alpha|_U=1\) on \(W[M]\), would descend to
the abelian cyclotomic quotient. Commuting \(g\) with the
chosen lift of inversion \(j\) there gives
\((1-\alpha(j))c(g)=(1-\alpha(g))c(j)=0\). As
\(\alpha(j)\equiv-1\bmod M\) and \(\ell\) is odd, this would
force \(c=0\).

Choose \(\rho^M=r\) and let \(V=\{u\in U:u\rho=\rho\}\).
The homomorphism \(u\mapsto u\rho/\rho\) embeds
\(U/V\) in \(\mu_M\). For \(g\in G_{\mathbb Q}\) and \(u\in U\),
the ratio \(g^{-1}\rho/\rho\) lies in \(\mu_M\) and is fixed by
\(u\). Therefore
\[
       (gug^{-1})\rho/\rho=g(u\rho/\rho).
\]
Consequently any \(\gamma\)-equivariant additive map
\(f:U\to W'\) into an \(M\)-torsion \(\mathcal O\)-module that
vanishes on \(V\) satisfies
\((\gamma(g)-\kappa(g))f(u)=0\) for every \(g,u\). This follows
by computing \(f(gug^{-1})\) on \(U/V\), where conjugation acts
by the exponent \(\kappa(g)\bmod M\).

Put \(B=\ker c\). The unit \(\alpha(\delta)-\kappa(\delta)\) and
\(c|_U\ne0\) show that \(V\) is not contained in \(B\).
For \(C=\mathcal O\,\nu(U)\ne0\), put
\(A=\{u\in U:\nu(u)\in\varpi C\}\). The values generate a finite
cyclic module: \(W[M]\) is a quotient of the discrete valuation
ring \(\mathcal O\) by a nonzero principal ideal, and all its
submodules are cyclic. The image modulo \(\varpi C\) is nonzero.
At a lift \(j\) of inversion we have
\(\beta(j)-\kappa(j)=\beta(j)(1-\alpha(j))\), a unit. Applying
the preceding calculation to \(f=\nu\bmod\varpi C\) proves that
\(V\) is not contained in \(A\). Each \(\phi\in U\setminus A\)
generates \(C\) under evaluation by \(\nu\). For \(C=0\) take
\(A=B\); the generator condition then holds automatically.

For the third subgroup put \(S=G_{\mathbb Q(\mu_{P\ell})}\).
Choose \(u\) acting on \(\mu_{P\ell}\) by \(1+P\), a unit
modulo \(P\ell\), and \(g\) acting by inversion. Then \(u\in U\),
and \(v=gug^{-1}u\) lies in \(V\) by the displayed radical formula.
Its exponent on \(\mu_{P\ell}\) is \((1+P)^2\); this differs
from \(1\bmod P\ell\) since \(\ell\nmid 2+P\).
Thus \(V\) is not contained in \(S\).

Here is the elementary simultaneous-avoidance argument.
A group cannot be covered by two proper subgroups: if neither
contains the other, the product of an element exclusive to each lies
in neither. If \(V\) were covered by \(A,B,S\), choose
\(x\in V\setminus(B\cup S)\) and \(y\in V\setminus(A\cup S)\).
Then \(x\in A,\ y\in B,\ xy\in S\). The element \(x^2y\) cannot
lie in \(A\), since \(y\notin A\), or in \(S\), since \(xy\in S\)
and \(x\notin S\). If \(x^2y\in B\), then \(x^2\in B\).
Also \(x^M\in B\) since \(M c(x)=0\). As \(M\) is odd, these
two powers imply \(x\in B\), a contradiction. Hence
\(V\setminus(A\cup B\cup S)\) is nonempty. This set is open:
the stabilizers are open and the kernels defining \(A,B\) are
open and closed by continuity to finite discrete groups.

Choose any coset of an open normal subgroup of \(G_{\mathbb Q}\)
contained in this set. Chebotarev's theorem in the corresponding
finite Galois extension gives an unramified \(q\) outside the
excluded set and the divisors of \(P\ell r\) with the conjugacy
class of that coset as Frobenius. A choice of place above \(q\)
realizes the prescribed representative, and an arithmetic Frobenius
lift \(\phi\) then belongs to our open set. The cyclotomic action
gives \(q\equiv1\bmod P\) and \(q\not\equiv1\bmod P\ell\).
The element \(\rho\) is a unit at this place since \(q\ne r\);
as \(\phi\rho=\rho\), its reduction is fixed by \(q\)-power
Frobenius and hence lies in \(\mathbb F_q^\times\). It follows
that \(r^{(q-1)/M}\equiv1\bmod q\). Because
\(v_\ell(q-1)=m\), the order of \(r\bmod q\) is prime to
\(\ell\). These are precisely the admissibility conditions used
in the derivative relations. The exclusions from \(A\) and \(B\)
give the two simultaneous evaluations.
\footnote{\href{https://github.com/ChristianKudeba/alpoge-rh-verification/blob/alpoge-source-7cbf7a0c/src/Development/AutomorphicArithmetic/AutomorphicVarieties/PrimitiveFrobeniusFixed/AddPowNsmul.lean}{Selection of the auxiliary prime and simultaneous avoidance},
\href{https://github.com/ChristianKudeba/alpoge-rh-verification/blob/alpoge-source-7cbf7a0c/src/Development/AutomorphicArithmetic/AutomorphicVarieties/PrimitiveFrobeniusFixed/RadEqv.lean}{action on the radical quotient},
\href{https://github.com/ChristianKudeba/alpoge-rh-verification/blob/alpoge-source-7cbf7a0c/src/Development/AutomorphicArithmetic/AutomorphicVarieties/PrimitiveFrobeniusFixed/ChebQ.lean}{Chebotarev with a finite set of exclusions}.}

\section{An individual Selmer estimate: the vanishing in case Z}
\label{sec:case-z}

Case Z enters \texttt{vwin\_\allowbreak{}of\_\allowbreak{}counts} through \texttt{count\allowbreak{}Z} and
\texttt{prop\allowbreak{}Z}. Here the source proves an exact zero. We derive the
estimate for both residue classes of the twist and then pass to the
character and level used in the window.

\subsection{The strict group and its torsion}

Fix an algebraic closure \(\Omega\) of \(\mathbb Q\), let
\(G=\operatorname{Gal}(\Omega/\mathbb Q)\), and write
\(\kappa:G\to\mathbb Z_\ell^\times\) for the cyclotomic character.
Let \(\mathcal O\) be a discrete valuation ring and a finite domain
over \(\mathbb Z_\ell\) as in \texttt{Coef}. Put
\[
  W=\mathcal O\otimes_{\mathbb Z_\ell}(\mathbb Q_\ell/\mathbb Z_\ell),
     \qquad \alpha=\kappa^{-e}.
\]
The action on \(W\) is scalar multiplication by \(\alpha(g)\).
All cohomology uses continuous cocycles and the discrete topology on
\(W\). Denote by \(T_{\ell,e}\) the submodule of \(H^1(G,W)\)
whose classes restrict to zero on every inertia group and on every
decomposition group above \(\ell\). We prove
\begin{equation}
 \ell\ge5,\quad e\ge2,\quad 2\mid e,\quad \ell-1\mid 2e
           \quad\Longrightarrow\quad T_{\ell,e}=0 .
 \label{eq:rh-z-vanish}
\end{equation}
In the source, \(\operatorname{length}_{\mathcal O}T_{\ell,e}\) is
\(\mathrm{Chr.V}(\ell,e,1)\).

Every element of \(W\) is killed by some power of \(\ell\).
Multiplication by a nonzero element of \(\mathcal O\) on \(W\) is
surjective: \(\mathcal O\) is finite and torsion-free over
\(\mathbb Z_\ell\), so \(W\) identifies with
\(\operatorname{Frac}(\mathcal O)/\mathcal O\).
A continuous cocycle on the compact group \(G\) has finite image in
\(W\). A common power of \(\ell\) therefore kills all its values.
Every cohomology class is consequently killed by some power of \(\ell\).
It suffices to show \(T_{\ell,e}[\ell]=0\), since a nonzero element
of order \(\ell^n\) gives one of order \(\ell\) after multiplication by
\(\ell^{n-1}\). We use the coboundary convention
\(d m(g)=(\alpha(g)-1)m\).

\subsection{The trivial residual character}

Suppose first that \(\ell-1\mid e\). For any
\(u\in\mathbb Z_\ell^\times\), reduction modulo \(\ell\) gives
\(u^{-e}\equiv1\bmod\ell\), because
\(\mathbb F_\ell^\times\) has order \(\ell-1\). Thus the action on
\(W[\ell]\) is trivial. For \(x\in T_{\ell,e}[\ell]\), choose a
representative \(z\). Write \(\ell z=dm\) and take \(m'\in W\)
with \(\ell m'=m\). The representative \(f=z-dm'\) satisfies
\(\ell f=0\) pointwise. Its cocycle identity is now the homomorphism
identity \(f(gh)=f(g)+f(h)\). At a prime other than \(\ell\),
the cyclotomic character is unramified. The restriction of \(f\)
to inertia is a coboundary there, hence is zero.

We need the following form of the rational cyclotomic calculation
\texttt{Kwl\allowbreak{}Onto.\allowbreak{}kwl}. If \(b:G\to\mathbb F_\ell\) is a continuous
additive homomorphism unramified away from an odd prime \(\ell\), then
\(b\) vanishes on the subgroup fixing \(\mu_{\ell^2}\).
Indeed, the Kronecker--Weber theorem with its ramification condition
places the field fixed by \(\ker b\) in \(\mathbb Q(\mu_{\ell^a})\)
for some \(a\ge2\). The unit group modulo \(\ell^a\) is cyclic of
order \(\ell^{a-1}(\ell-1)\). The kernel of reduction modulo
\(\ell^2\) has order \(\ell^{a-2}\). In the cyclic group it consists
of powers of a generator whose exponent is divisible by
\(\ell(\ell-1)\), so each is an \(\ell\)-th power. Every additive
map into \(\mathbb F_\ell\) kills these powers. This proves the claim,
including \(a=2\).

Apply the claim to each composite of \(f\) with a linear functional
\(W[\ell]\to\mathbb F_\ell\). Such functionals separate points, so
\(f\) vanishes on the subgroup fixing \(\mu_{\ell^2}\), hence on
\(\ker\kappa\). The restriction-to-the-cyclotomic-kernel step of
\texttt{Tors\allowbreak{}Class.\allowbreak{}z\allowbreak{}One} now gives \(x=0\). For clarity its injectivity
in this specialization has the following direct verification. Both the
action and the cocycle factor through \(\Gamma=\mathbb Z_\ell^\times\).
For odd \(\ell\), this group has a topological generator \(\gamma\):
its finite factor has order \(\ell-1\), coprime to the orders of all
finite quotients of its procyclic factor \(1+\ell\mathbb Z_\ell\).
Since \(e>0\), the scalar \(\gamma^{-e}-1\) is nonzero.
Choose \(v\in W\) with \((\gamma^{-e}-1)v=f(\gamma)\).
The continuous cocycle \(f-dv\) vanishes on \(\gamma\), hence on all
its integer powers by the cocycle identity, and then everywhere by
density. This proves the vanishing in the first case.

\subsection{The quadratic residual character}

Suppose now that \(\ell-1\mid 2e\) and \(\ell-1\nmid e\).
Write \(n=(\ell-1)/2\). Then \(e=nk\) for an odd integer \(k\).
Since \(e\) is even, \(n\) is even and \(\ell\equiv1\bmod4\).
The residual character \(\overline\kappa^{-e}\) is the quadratic
character of \(\mathbb F_\ell^\times\).

Choose a primitive \(\ell\)-th root \(\zeta\) in \(\Omega\) and let
\(\eta\) be the Legendre symbol, extended by \(\eta(0)=0\).
Set \(s=\sum_{a\bmod\ell}\eta(a)\zeta^a\). Acting on \(\zeta\) by
the exponent \(u\) multiplies \(s\) by \(\eta(u)\), by the change of
variable \(a\mapsto ua\). Replacing \(\zeta\) by \(\zeta^{-1}\)
fixes \(s\) since \(\eta(-1)=1\). Multiplying these two sums and
writing \(a=tb\) in the terms with \(b\ne0\) gives
\[
 s^2
   =\sum_{t\bmod\ell}\eta(t)\sum_{b\ne0}\zeta^{b(t-1)}
   =(\ell-1)-\sum_{t\ne1}\eta(t)=\ell .
\]
Here the sum of the nontrivial character over all residues is zero.
Thus \(K=\mathbb Q(s)=\mathbb Q(\sqrt\ell)\) is quadratic with
discriminant \(\ell\), and its absolute Galois group \(G_K\) is the
kernel of the residual quadratic character. For \(g\in G_K\),
\(\alpha(g)-1\in\ell\mathcal O\). In particular \(G_K\) acts
trivially on \(W[\ell]\).

Choose \(u\in\mathbb Z_\ell^\times\) reducing to a generator of
\(\mathbb F_\ell^\times\). Then \(\overline u^e\ne1\).
At every place \(A\) above \(\ell\) the cyclotomic character maps
inertia onto \(\mathbb Z_\ell^\times\). Indeed the local cyclotomic
extensions of level \(\ell^a\) are totally ramified with full unit
Galois group, as follows by shifting the cyclotomic polynomial by
one to obtain an Eisenstein polynomial. Compactness passes surjectivity
to the inverse limit. Choose \(\delta\) at one such place with
\(\kappa(\delta)=u\). Then \(\alpha(\delta)-1\) is a unit and
\(\delta s=-s\).

Given \(x\in T_{\ell,e}[\ell]\), normalize a representative \(c\)
to have \(c(\delta)=0\), by subtracting
\(d((\alpha(\delta)-1)^{-1}c(\delta))\). Since \(\ell x=0\),
we have \(\ell c=dw\). Evaluation at \(\delta\) gives
\((\alpha(\delta)-1)w=0\), hence \(w=0\). Thus \(c(g)\in W[\ell]\)
for every \(g\). The restriction \(c|_{G_K}\) is an additive
homomorphism. Away from \(\ell\) it vanishes on inertia, since the
character is unramified there and the strict class is locally zero.

At a place \(A\mid\ell\) the strict condition gives \(m_A\in W\)
such that \(c(g)=(\alpha(g)-1)m_A\) for all \(g\) in its
decomposition group. Take \(\tau\) in inertia at this same place
with \(\kappa(\tau)=u\). Then
\((\alpha(\tau)-1)\ell m_A=\ell c(\tau)=0\). The factor is a unit,
so \(m_A\in W[\ell]\). Therefore \(c(g)=0\) for every
\(g\) in the decomposition group that also belongs to \(G_K\).
The restriction to \(G_K\) is unramified at all finite places.
Its target has odd exponent, so it also vanishes on all real inertia.
By the unramified class-field correspondence it factors through the
ideal class group of \(K\).

We give the class-number inequality used here in \texttt{HQuad}.
Put \(X=\sqrt{\ell}\). The quadratic Minkowski bound supplies every
class with an integral ideal of positive norm at most \(X\).
For each integer \(a\ge1\), the number of ideals of norm \(a\) is
at most the number \(d(a)\) of ordered factorizations \(a=ij\).
On a prime power \(p^v\) the split case admits \(v+1\) possible
pairs of exponents above \(p\), and the inert and ramified cases
each admit at most one possibility. Unique factorization of ideals
multiplies these bounds over rational primes. Thus
\[
    h(K)\le\sum_{a\le X}d(a)
      =\sum_{i\le X}\lfloor X/i\rfloor
      \le X(1+\log X)<X^2=\ell .
\]
The harmonic-sum bound follows by integration of \(1/t\), and the
strict inequality from \(\log X<X-1\) for \(X>1\).
Hence \(h(K)\) is coprime to \(\ell\). Every value of \(c|_{G_K}\)
is annihilated both by \(h(K)\) and by \(\ell\), so this restriction
is zero. For \(g\notin G_K\), write \(g=\delta h\), \(h\in G_K\).
Then \(c(g)=c(\delta)+\alpha(\delta)c(h)=0\).
This proves the second case of \eqref{eq:rh-z-vanish}.

The normalization and local restriction just used are the zero-length
specialization of \texttt{Primitive\allowbreak{}Frobenius\allowbreak{}Fixed.\allowbreak{}Sha\allowbreak{}To\allowbreak{}Sh\allowbreak{}T} and
\texttt{Level\allowbreak{}Cohomology.\allowbreak{}Zq\allowbreak{}Field.\allowbreak{}z\allowbreak{}Quad}. In the general-length
formulation the source takes the union over annihilators \(a\ne0\) of
normalized cocycles. Every such annihilator divides a power of \(\ell\)
in this discrete valuation ring. The preceding torsion argument gives
the same conclusion at length zero. The class-field inputs in the
two cases are respectively ramified Kronecker--Weber over \(\mathbb Q\)
and the unramified reciprocity map over the real quadratic field.

\subsection{The estimate at the host level}

In case Z the coefficient prime is odd and outside \(B_\chi\).
Since \(3\in B_\chi\), we have \(\ell\ge5\). The character at level
\(N\) is trivial, and the host also requires
\(\psi(-1)=(-1)^e\). In the characteristic-zero coefficient ring this
forces \(e\) even. The window has \(0<e<t\), so \(e\ge2\).
The last case-Z condition is \(\ell-1\mid2e\); all hypotheses of
\eqref{eq:rh-z-vanish} hold. Passing from the strict condition at
\(\ell\) to the one at \(\ell N\) adds decomposition-group conditions.
It takes a submodule of the group just shown to vanish. The arithmetic
length \(b_{\ell,e}(1)\) of Section~\ref{sec:arithmetic-translation}
is consequently zero at every such host. Its normalized logarithmically
weighted length is zero as well. This proves \texttt{Count\allowbreak{}Z} with the
single exponent \(C_Z=0\), uniformly in all the windows and choices of
coefficients.

\section{Polynomial estimates for the odd-prime positive eigenspace}
\label{sec:rh-plus-nonzero}

We derive the individual estimate \texttt{Height\allowbreak{}And\allowbreak{}Vanishing.\allowbreak{}Count\allowbreak{}Vs}
on the positive-eigenspace side of the strict window. In
\texttt{Joint\allowbreak{}Win\allowbreak{}Of\allowbreak{}Counts.\allowbreak{}point\allowbreak{}HV\_\allowbreak{}of\_\allowbreak{}counts}, this estimate is combined
with counts II, IV and the other branch of V to give \texttt{Sum\allowbreak{}In}.
The result concerns the twist \(e=1\), where the source gives an upper
bound for the length. We also derive the estimate used outside the
exceptional set for arbitrary positive twists in case I.

\subsection{The conductor condition bounds the degree}

Fix a primitive target \(\chi\) of conductor \(m\) and order \(o\),
and write \(M=\max B_\chi\), which exists since \(3\in B_\chi\).
Take any host of character modulus \(N\) satisfying the size conditions
of Section~\ref{sec:arithmetic-translation}, an odd prime
\(\ell\in B_\chi\), any admissible coefficient ring \(\mathcal O\), and
\(\psi\in\mathcal X_{\ell,1}\). The cutoff is \(t>1\), so \(e=1\) is
in the range. Write \(\psi_2\) for the \(2\)-component of \(\psi\)
at the level and \(c=v_2(\operatorname{cond}\psi_2)\).
The conductor of \(\psi_2\) is \(2^c\).

Remove the \(\ell\)-power-order part of \(\psi\), then its component
at two, to obtain \(\theta\); let \(a\) denote its order. Put
\[
  f=\operatorname{ord}_{\operatorname{lcm}(4,a)}(\ell),\qquad
  w=v_2(\ell^f-1),\qquad R=16o(\ell m)^2\,4^w.
\]
The hypothesis for this branch is \(2^c<R\). The host character
comparison \texttt{host\allowbreak{}Th1} gives \(a\mid2o\) and \(\ell\nmid a\).
The modulus \(d=\operatorname{lcm}(4,a)\) divides \(4o\) and is
prime to \(\ell\). Thus \(f\ge1\), \(f\mid\varphi(d)\), and
\(\varphi(d)\mid\varphi(4o)\). The last divisibility follows directly
on comparing the factors \(p^{v-1}(p-1)\) in Euler's product.
It follows that \(2^w\le\ell^f-1<\ell^{\varphi(4o)}\), giving
\[
        R\le 16o(\ell m)^2\ell^{2\varphi(4o)}.
\]
The source's component calculation \texttt{host\allowbreak{}Comp} bounds the order
of every component of \(\psi\) at an odd prime dividing \(N\) by
divisibility into \(2o\). The order of \(\psi_2\) divides
\(\varphi(2^c)\), since it factors through the units of its conductor,
and \(\varphi(2^c)\mid 2^c\), including \(c=0\).
The Chinese remainder decomposition now gives \(\psi^{2o\,2^c}=1\).
Thus \texttt{Pow\allowbreak{}Totient\allowbreak{}Conductor.\allowbreak{}vs\allowbreak{}Ord} gives
\[
       \operatorname{ord}(\psi)\,\ell\le
       D_0,\qquad
       D_0=2o\bigl(16o(Mm)^2 M^{2\varphi(4o)}\bigr)M .
\]
This number depends only on the target, even when \(N\) varies.

Inside \(\mathbb Q(\mu_{\ell N})\) take the fixed field \(L\) of
the kernel of the map from the units modulo \(\ell N\) given by
\[
            u\longmapsto \bigl(\psi(u\bmod N),u\bmod\ell\bigr).
\]
Then \(\mu_\ell\subset L\), \(L/\mathbb Q\) is abelian, and \(\psi\)
is trivial on its absolute Galois group. Its finite image in the first
factor is cyclic of order \(\operatorname{ord}(\psi)\), since it is a
finite subgroup of a field's multiplicative group. The second factor
has order \(\ell-1\). Consequently \(n=[L:\mathbb Q]\le
\operatorname{ord}(\psi)(\ell-1)\le D_0\). This is the construction in
\texttt{Vs\allowbreak{}Field.\allowbreak{}vs\allowbreak{}Field}. Put \(x=\ell N\) and
\(\Delta=|\operatorname{disc}L|\). The discriminant tower identity
gives \(\Delta^{1/n}\le
|\operatorname{disc}\mathbb Q(\mu_x)|^{1/\varphi(x)}\).
The cyclotomic discriminant formula bounds the right side by \(x\);
explicitly it is \(x\prod_{p\mid x}p^{-1/(p-1)}\).
Thus \(\log\Delta\le n\log x\), as in \texttt{Win.\allowbreak{}disc}.

\subsection{Descended length and the numerical comparison}

The arithmetic input \texttt{Lemma\allowbreak{}K} is a size assertion for \(H^2\).
At twist \(e\) over a totally complex field \(L\) containing the
specified roots of unity, and for every choice of its local data,
it gives a finite \(H^2\) for the compact integral cyclotomic twist
\(\mathbb Z_\ell(e+1)\) in the source's Galois problem, and
\[
 \log\# H^2 \le
 \exp\big((2en)^{C_K}\log(2\Delta)\big).
\]
In the present case \(L\) is totally complex since it contains a
primitive odd-prime root of unity. Write \(K=\max(0,\lceil C_K\rceil)\)
and \(A=(2D_0)^K(D_0+1)\). Replacing \(C_K\) by \(K\) increases
the right side. At \(e=1\), write \(k\) for the \(\mathbb Z_\ell\)-length
of this finite group, so its cardinality is \(\ell^k\). Because
\(\log\ell\ge\log2>1/2\),
\[
      k\le2\exp\big((2n)^K\log(2\Delta)\big)\le2x^A.
\]
For the second inequality use \(x\ge2\) and \(1\le n\le D_0\) to get
\(\log(2\Delta)\le\log x+n\log x\le(D_0+1)\log x\).
The exponent of \(x\) is therefore independent of the host and
coefficient choices.

We record the arithmetic comparisons used to descend this estimate.
The global-to-local inequality \texttt{Prop\allowbreak{}A.\allowbreak{}sha\allowbreak{}Le\allowbreak{}V} uses
Poitou--Tate duality, local finiteness and local Tate duality to
bound the strict \(\mathbb Z_\ell\)-length over \(L\) by \(k\).
The character restriction bound \texttt{Chr.\allowbreak{}up\allowbreak{}V} allows one error
from local degree-zero invariants for each generator of
\(\operatorname{Gal}(L/\mathbb Q)\). Take the full group as generating
set; it has \(n\) elements. The Galois character of \(\psi\) factors
through the units modulo \(N\), so
\texttt{H0\_\allowbreak{}local\_\allowbreak{}card\_\allowbreak{}bounded\_\allowbreak{}of\_\allowbreak{}cls\_\allowbreak{}ker\_\allowbreak{}le\_\allowbreak{}ker} applies.
With \(r_\ell=\operatorname{length}_{\mathcal O}(\mathcal O/\ell)\)
the resulting inequalities are
\[
    b_{\ell,1}(\psi)\ \le\ (k+n\nu)\,r_\ell,\qquad
    \nu=
    \begin{cases}
       \max(v_\ell(N),1),&\ell\mid N\ \text{or}\ \ell-1\mid 1,\\
       0,&\text{otherwise}.
    \end{cases}
\]
The displayed local constant specializes the definition
\texttt{LData\allowbreak{}Card\allowbreak{}Le\allowbreak{}Finrank.\allowbreak{}nu\allowbreak{}Bound} to odd \(\ell\) and \(e=1\).
It is at most \(2N\), since \(v_\ell(N)\le N\) and \(N\ge2\).
The displayed inequality also uses monotonicity when the strict
decomposition conditions at the primes dividing \(N\) are imposed.
All three comparisons here are separate arithmetic inputs.

The ramification index \(r_\ell\) is a positive finite integer,
so \(b_{\ell,1}(\psi)\) is finite and division gives
\(v_{\ell,1}(\psi)\le2x^A+2D_0 N\).
Since \(0\le\log\ell\le\ell\le x\) and \(N\ell=x\),
\[
\begin{aligned}
  (\log\ell)v_{\ell,1}(\psi)
   &\le 2x^A\ell+2D_0 x\\
   &\le 2(D_0+1)x^{A+1}\\
   &\le x^{A+D_0+3}.
\end{aligned}
\]
Indeed \(D_0+1\le2^{D_0}\) and \(x\ge2\), so
\(2(D_0+1)\le x^{D_0+2}\). Hence the asserted count holds with the
single exponent \((2D_0)^K(D_0+1)+D_0+3\). The two alternatives for
the conductor threshold are what permit the source to use this
degree bound for every occurrence of the strict inequality \(2^c<R\).

\subsection{The uniform case I estimate over a finite set of primes}

We now follow \texttt{Unitary\allowbreak{}Three\allowbreak{}Groups.\allowbreak{}Count\allowbreak{}I}. In the actual
window argument, \texttt{Linked\allowbreak{}Unl\allowbreak{}E.\allowbreak{}count\_\allowbreak{}i\_\allowbreak{}of\_\allowbreak{}printed\_\allowbreak{}bound\_\allowbreak{}k\_\allowbreak{}kato\_\allowbreak{}automorphy}
obtains its orbit estimate by applying \texttt{U3.\allowbreak{}cor\allowbreak{}GK} to
\texttt{prop\allowbreak{}GK\_\allowbreak{}of\_\allowbreak{}number\allowbreak{}U3\allowbreak{}U\_\allowbreak{}of\_\allowbreak{}one\allowbreak{}Char\allowbreak{}U}. Apply the orbit estimate
of Section~\ref{sec:arithmetic-translation} with
\(F=\mathbb Q(\sqrt{-p_K})\), and let \(c_5,A\) be its constants for
this field. Choose primes \(r_1>30mp_K\) and \(r_2>r_1\), and choose
\(k\ge0\) with \(2c_5(\log r_1+\log r_2)<16^k\). The source
claims the exponent \(k+1+A\) in count I for this target.

Indeed take a host for which each \(r_i\) is smaller than both
host primes \(q_1,d\), a twist \(1\le e<t\), an arbitrary family of
coefficient rings, and a finite set \(S\) of odd primes in case I.
The host-level facts give \(16\mid N\) and every prime divisor of
\(N\) equals two, divides \(m\), equals \(q_1\), or equals \(d\).
Furthermore \(|\operatorname{disc}F|=p_K\). Thus neither \(r_i\)
divides \(30N|\operatorname{disc}F|\). The test primes are eligible
simultaneously for every host satisfying these inequalities.

Let \(G_N=(\mathbb Z/N\mathbb Z)^\times\) and form the finite set
\(\mathcal H\) of kernels of characters in
\(\bigcup_{\ell\in S}\mathcal X_{\ell,e}\), taking kernels in \(G_N\).
For \(H\in\mathcal H\), let
\[
  S_H=\{\ell\in S:\ \ker\psi=H
           \text{ for some }\psi\in\mathcal X_{\ell,e}\}.
\]
For every \(\ell\in S_H\), the host lemma \texttt{gd\allowbreak{}Of\allowbreak{}Host} makes that
character a witness for the good-data predicate in the orbit estimate.
In particular every length occurring in the desired sum is finite.
Partition the character sum at each prime by its kernel, and enlarge
each nonempty part to the entire orbit of characters with that kernel.
The summands are nonnegative, and \(\log\ell\) is positive, so this
operation gives an upper bound. Applying the orbit estimate to each
\(H\) with its own \(S_H\) now gives
\[
 \sum_{\ell\in S}(\log\ell)\sum_{\psi\in\mathcal X_{\ell,e}}
       v_{\ell,e}(\psi)
 \le
 2c_5(\log r_1+\log r_2)\{N(e+1)\}^A
        \sum_{H\in\mathcal H}\varphi([G_N:H]).
\]
For each such \(H\), the quotient \(G_N/H\) is finite cyclic, being
isomorphic to the image of a multiplicative character in a field.
If its order is \(h\), the faithful complex characters on this
quotient number \(\varphi(h)\), by sending a generator to a primitive
\(h\)-th root of unity. Inflating them to \(G_N\) gives characters
with kernel exactly \(H\). Distinct \(H\)'s give disjoint sets.
The finite abelian group \(G_N\) has \(|G_N|\) complex characters:
decompose it as a direct product of cyclic groups and count the
choices on each factor. Therefore
\(\sum_{H\in\mathcal H}\varphi([G_N:H])\le |G_N|=\varphi(N)\le N\).
This proves the counting step \texttt{orb\allowbreak{}Sum} without dependence on
the number of coefficient primes.

Finally \(16\le N\), \(e+1\le t\) and \(t\ge1\).
The right side is at most
\[
            16^k\,N\,\{N(e+1)\}^A\le (Nt)^{k+1+A}.
\]
This proves the numerical and orbit deductions in the case I estimate.
Their arithmetic inputs are the good-data assertion for eligible
characters and the rank-three estimate that supplies the orbit bound.

\subsection{Why the sum of valuations is bounded by one integer}

The step \texttt{FProp\allowbreak{}GKU.\allowbreak{}prop\allowbreak{}GK\_\allowbreak{}of\_\allowbreak{}number\allowbreak{}U3\allowbreak{}U\_\allowbreak{}of\_\allowbreak{}one\allowbreak{}Char\allowbreak{}U}
combines two substantial inputs. The first is an integer estimate for
an entire collection of rank-three trace comparisons, denoted by
\texttt{Number\allowbreak{}U3\allowbreak{}U\_\allowbreak{}s}. The second, \texttt{One\allowbreak{}Char\allowbreak{}U\_\allowbreak{}s}, bounds a
single normalized Selmer length by trace differences. Here is their
precise use for a fixed imaginary quadratic \(F\), \(e\ge1\), \(N\ge1\),
and subgroup \(H\le G_N\).

Choose a prime \(p\nmid30N|\operatorname{disc}F|\), a prime \(w\)
of \(F\) over \(p\), and write \(q=|\mathcal O_F/w|\). In particular
\(p\ge7\) and \(q\) is coprime to \(N\). Take as index set all complex
characters \(\xi\) of modulus \(N\) with kernel \(H\). For each index
use the root of unity \(z_\xi=\xi(q)^{-1}\). Applying a field
automorphism of \(\mathbb C\) to character values permutes this set and
sends each \(z_\xi\) to the value indexed by the conjugate character.
Thus the collection satisfies the Galois stability required of the
integer estimate. The index set has at most \(\varphi([G_N:H])\)
elements, by the faithful cyclic-character calculation above (and is
empty if no character has kernel \(H\)).

The integer estimate supplies a nonzero rational integer \(n_p\),
chosen before the coefficient prime, with
\[
     \log |n_p|\le c_5\varphi([G_N:H])\{N(e+1)\}^{A}\log p .
\]
Fix an odd prime \(\ell\ne p\), an embedding of the coefficient ring into
\(\overline{\mathbb Q}_\ell\), and a field isomorphism
\(\iota:\overline{\mathbb Q}_\ell\longrightarrow\mathbb C\).
Suppose the \(\ell\)-adic orbit with kernel \(H\) has a good-data
witness. Composition with \(\iota\) identifies that orbit with the
complex index set. Indeed composition preserves kernels and is injective
on characters. Starting with the witness, any other complex character
with the same kernel is a power of its image by an exponent coprime to
the cyclic quotient order. The same power of the witness gives a
preimage in the coefficient ring.

Choose a Frobenius element \(\sigma\) at a place above \(w\) with
the convention that it acts by the \(q\)-power on unramified residue
data. Its cyclotomic images are \(q\bmod N\) and \(q\in\mathbb Z_\ell\).
For the character \(\psi\) corresponding to \(\xi\), the
reducible comparison trace evaluated at \(\sigma^{-1}\) is therefore
\[
 T_\psi=\iota^{-1}\big(q^{-e}\xi(q)^{-1}+1+q^e\xi(q)\big).
\]
For an eligible three-dimensional representation \(\rho\), put
\(d_{\psi,\rho}=\operatorname{tr}(\rho(\sigma^{-1}))-T_\psi\).
The two eligible families in the printed input are
\(\mathrm{Adm3U}_{s}\) at level \(N\) and \(\mathrm{Adm2A}\) at
level \(N|\operatorname{disc}F|\). The uniform integer theorem,
applied to singleton sets of representations, gives
\(d_{\psi,\rho}\ne0\) for every eligible \(\rho\). The good-data
conditions transfer from the witness to each character with the same
kernel, by \texttt{gd\allowbreak{}Facts}. The single-character input now provides
finiteness and a finite set \(X_\psi\) of pairwise nonisomorphic
eligible representatives, with
\[
       v_{\ell,e}(\psi)\le
              2\sum_{\rho\in X_\psi}
                 \bigl(-\log_\ell|d_{\psi,\rho}|_\ell\bigr).
\]
Apply the uniform integer theorem once more with this simultaneous
choice of the \(X_\psi\) indexed by all the complex characters. It gives
\[
       |n_p|_\ell\le\prod_{\psi}\prod_{\rho\in X_\psi}
                                      |d_{\psi,\rho}|_\ell .
\]
Every factor is nonzero. Taking negative logarithms to base \(\ell\)
reverses the inequality, and
\(-\log_\ell|n_p|_\ell=\operatorname{ord}_\ell(n_p)\).
Summing the individual estimates thus gives
\[
              \sum_{\ker\psi=H}v_{\ell,e}(\psi)
                         \le2\operatorname{ord}_\ell(n_p).
\]
The simultaneous product and its uniform integer are essential inputs
at this point: individually bounded trace differences would not imply
this bound for their sum. The proof of the integer theorem and that of
the single-character theorem are upstream of the calculation just given.

\section{The integer controlling the Frobenius trace differences}
\label{sec:common-trace-integer}

Fix an imaginary quadratic field \(K\), a positive integer \(N\),
a twist \(e\ge1\) and a prime \(p\ge7\) prime to
\(N\operatorname{disc}(K)\). Let \(w\mid p\), let
\(h=[\mathcal O_K/w:\mathbb F_p]\in\{1,2\}\), and set
\(q=p^h\) and \(Q=q^e\). Let \(O\) be a finite index set with
roots of unity \(z_o\in\mathbb C\) such that every field automorphism
of \(\mathbb C\) permutes the indexed values. Repetitions are
allowed and their multiplicities matter. Write
\[
                T_o=Q^{-1}z_o+1+Qz_o^{-1}.
\]
For the orbit in Section~\ref{sec:rh-plus-nonzero}, \(O\) consists
of the characters \(\xi\) with the prescribed kernel and
\(z_\xi=\xi(q)^{-1}\).

\subsection{The product and its size}

A finite family of trace values consists of a finite set \(S\)
and numbers \(t_s\in\mathbb C\) such that every automorphism of
\(\mathbb C\) permutes the indexed values and
\[
           |t_s|\le\sqrt q+2,\qquad
           Qt_s\ \text{is an algebraic integer}.
\]
At a coefficient prime \(\ell\nmid q\) and for a field
isomorphism \(\iota:\overline{\mathbb Q}_\ell\to\mathbb C\), a
family of admissible representations is \emph{classified} by \(S\)
at \(g\) if each \(\rho\) has a label \(s(\rho)\in S\) with
\[
 \iota\bigl(\operatorname{tr}\rho(g)\bigr)=t_{s(\rho)},\qquad
 s(\rho)=s(\rho')\ \Longrightarrow\ \rho\simeq\rho'.
\]
Several labels can have equal numeric traces. The implication
about labels preserves the multiplicity of their contributions.

Define the product before choosing any coefficient prime,
representations or isomorphism by
\begin{equation}
   n_p=\prod_{o\in O}\prod_{s\in S} F_{o,s},
   \qquad
   F_{o,s}=Q(t_s-T_o)=Qt_s-z_o-Q-Q^2z_o^{-1}.
   \label{eq:common-trace-product}
\end{equation}
We claim that \(n_p\in\mathbb Z\setminus\{0\}\), that
\begin{equation}
  \log|n_p|\le |O|\,|S|\,(2e\log q+1),
  \label{eq:common-trace-log}
\end{equation}
and that for every classified family and every finite selection
\(X_o\) of pairwise nonisomorphic members for each \(o\),
\begin{equation}
 |n_p|_\ell\le
 \prod_{o\in O}\prod_{\rho\in X_o}
 \left|\operatorname{tr}\rho(g)-\iota^{-1}(T_o)\right|_\ell .
 \label{eq:common-trace-subproduct}
\end{equation}

To separate \(t_s\) and \(T_o\), note that \(Q\ge q\ge7\) and
\(\sqrt q\le(q+9)/6\) by \((\sqrt q-3)^2\ge0\).
It follows that \(Q-\sqrt q-3\ge4/3>Q^{-1}\). By the reverse
triangle inequality,
\[
              |T_o|\ge Q-1-Q^{-1}>\sqrt q+2 .
\]
Every factor of \eqref{eq:common-trace-product} is therefore
nonzero. Each factor is integral because \(Qt_s\), \(z_o\),
\(z_o^{-1}\) and \(Q\) are algebraic integers. In addition,
\[
 |F_{o,s}|\le Q^2+Q(\sqrt q+3)+1\le2Q^2 .
\]
The last inequality follows from \(Q(Q-\sqrt q-3)>1\).

All factors belong to a finite normal extension of \(\mathbb Q\)
in \(\mathbb C\). Each automorphism of this extension extends to an
automorphism of \(\mathbb C\) by extending first to an algebraic
closure and then using a transcendence basis. The resulting
permutations of \(O\) and \(S\) jointly permute \(O\times S\).
Their product \(n_p\) is consequently rational. A rational algebraic
integer is an integer. The empty product is \(1\). Taking logarithms
and using \(\log2\le1\) proves \eqref{eq:common-trace-log}.

Finally transport all the factors by \(\iota^{-1}\).
Every resulting algebraic integer has absolute value at most one
in \(\overline{\mathbb Q}_\ell\). For each \(o\), the map
\(\rho\mapsto s(\rho)\) is injective on \(X_o\). Discarding the
factors whose labels are outside its image can only increase
the \(\ell\)-adic absolute value of the product. Since
\(|Q|_\ell=1\), the remaining factors give
\eqref{eq:common-trace-subproduct}. In particular every trace
difference in its right side is nonzero. This proves all three
assertions for arbitrary trace data with the stated hypotheses.

\subsection{Trace data from the two kinds of eigensystems}

We describe exactly which local estimates give these hypotheses.
The rank-three CM eigensystems for the inverse cyclotomic twist
have three parameters \(b_i\) with the convention that the
eigenvalues of the twisted representation at \(\sigma^{-1}\) are
the \(b_i^{-1}\), where \(\sigma\) is arithmetic Frobenius at \(w\).
The weight is \(-2\) and the integral weight vector is
\((e-1,-1,-e-1)\). The local estimates on the parameters are
\[
     q^{-1}q^{-1/2}< |b_i| < q^{-1}q^{1/2},\qquad
     \prod_{i=1}^3 |b_i|=q^{-3},\qquad
     q^{e-1}b_i^{-1}\ \text{is integral}.
\]
Here integrality means algebraic integrality over \(\mathbb Z\).
Undoing the twist divides each eigenvalue by \(q\), since
\(\chi_\ell(\sigma)=q\). The trace label for the original
three-dimensional representation is
\[
                         t_s=\sum_{i=1}^3(qb_i)^{-1}.
\]
Thus \(Qt_s\) is integral. Put \(a=\sqrt q\) and
\(x_i=|(qb_i)^{-1}|\). We have \(1/a\le x_i\le a\) and
\(\prod_i x_i=1\). Reorder them so that \(x_1\ge x_2\ge x_3\).
If \(x_2\le1\), their sum is at most \(a+2\).
If \(x_2\ge1\), then \(x_1+x_2\le x_1x_2+1\) and
\(x_1x_2=x_3^{-1}\le a\), while \(x_3\le1\).
The same bound holds in both cases, so \(|t_s|\le\sqrt q+2\).

For the rank-two systems over \(\mathbb Q\), the prime has residue
size \(p\), the weight is zero and the integral weight vector is
\((e,-e)\). Their two parameters satisfy
\[
   p^{-1/2}<|b_i|<p^{1/2},\qquad |b_1b_2|=1,\qquad
             p^e b_i^{-1}\ \text{is integral}.
\]
The representation over \(K\) is isomorphic to
\(\mathbf1\oplus W|_{G_K}\). On the unramified quotient at the
chosen place, arithmetic Frobenius over \(K\) restricts to the
\(h\)-th power of arithmetic Frobenius over \(\mathbb Q\).
Consequently its inverse gives
\[
                   t_s=1+b_1^{-h}+b_2^{-h}.
\]
The identity \(Qb_i^{-h}=(p^eb_i^{-1})^h\) gives the claimed
integrality. Of the two numbers \(|b_1|^{-h}, |b_2|^{-h}\) with
product one, the larger is at most \(\sqrt q\) and the smaller
at most one. This again gives \(|t_s|\le\sqrt q+2\).
In both ranks the operations used to make \(t_s\) commute with
field automorphisms. A permutation of the parameter multisets
therefore permutes the trace labels.

The absolute value bounds here are precisely the Satake bounds
used in the released argument. For a unitary cuspidal
representation of \(\mathrm{GL}_n\) at a spherical place of residue
size \(u\), the Jacquet--Shalika estimate puts the absolute values
of its normalized Satake parameters strictly between \(u^{-1/2}\)
and \(u^{1/2}\). Their product has absolute value one by unitarity
of the central character. Twisting to weight zero gives the
displayed inequalities after reinstating the normalizing powers of
\(q\) or \(p\). See H.~Jacquet and J.~Shalika, \emph{On Euler
products and the classification of automorphic representations I},
Amer.~J.~Math.~103 (1981), 499--558. Integrality of the scaled
inverse parameters is the separate integral Satake estimate in the
source.

\subsection{Global labels and the quadratic restriction argument}

A label for an automorphic representation over a CM field is its
full system of normalized unramified Hecke multisets away from
the prescribed level. The test place \(w\) is only one entry.
Suppose the relevant cuspidal representations at weight vector
\(m\) and level \(M\) have at most \(C(M(H+1))^A\) pairwise
nonisomorphic members, where \(H\ge\max_i|m_i|\). Suppose also
that the system depends only on the isomorphism class, its Galois
conjugates occur in the same collection, and the attached
semisimple Galois representations agree with it at places
away from \(M\ell\). The set of distinct global systems then has
the same upper bound. Indeed any finite subset of the image lifts
to pairwise nonisomorphic preimages by choosing one form for each
system. Thus it is a finite set of labels independent of \(\ell\).
Conjugation permutes these labels. Equality of labels gives equality
of the characteristic polynomials of the attached representations
at all but finitely many Frobenius classes. By Chebotarev these
classes are dense up to conjugacy. Continuity gives equality of
traces everywhere, and the characteristic-zero Brauer--Nesbitt
theorem gives isomorphism for semisimple representations.
These statements construct exactly the classification used above
when the indicated automorphic estimates and compatibility hold.

For the rational rank-two systems the source first constructs
finite collections over auxiliary imaginary quadratic fields. Here
is the finite counting argument allowing those collections to
classify the rational systems. Choose six imaginary quadratic
fields \(K_i\) of pairwise coprime discriminants. For example
one may take \(K_i=\mathbb Q(\sqrt{-d_i})\) with
\(d_i=3,7,11,19,23,31\). At most one of these discriminants is
divisible by the test prime \(p\), leaving five choices. Write
\(\delta_i\) for the distinct quadratic characters of
\(G_{\mathbb Q}\) with kernels \(G_{K_i}\).

For an irreducible two-dimensional representation \(W\) over
\(\overline{\mathbb Q}_\ell\), at most four distinct characters
\(\delta\) satisfy \(W\simeq W\otimes\delta\). Each such
intertwiner is a nonzero vector in the four-dimensional
\(\operatorname{End}(W)\) on which conjugation acts by
\(\delta\) (or by its inverse, equal to it here). Nonzero
eigenvectors belonging to distinct characters of a group are
linearly independent. Indeed a minimal relation with at least
two terms yields a shorter nonempty relation by subtracting a
scalar multiple of its translate by an element distinguishing
two of the characters. Consequently one of the five remaining
\(i\) satisfies \(W\not\simeq W\otimes\delta_i\).

For \(H_i=G_{K_i}\), restriction to \(H_i\) is semisimple on
this irreducible \(W\). In fact choose a simple \(H_i\)-submodule
and sum its two conjugates under coset representatives. The span
is nonzero and \(G_{\mathbb Q}\)-stable, hence all of \(W\), and
a sum of simple submodules is semisimple. Frobenius reciprocity
and \(\operatorname{Ind}_{H_i}^{G_{\mathbb Q}}\mathbf1
\simeq\mathbf1\oplus\delta_i\) give
\[
 \operatorname{End}_{H_i}(W)
    \simeq\operatorname{End}_{G_{\mathbb Q}}(W)
       \oplus\operatorname{Hom}_{G_{\mathbb Q}}
                             (W,W\otimes\delta_i)
\]
as vector spaces. Its dimension is one for our chosen \(i\), so
\(W|_{H_i}\) is irreducible. The same adjunction for two irreducible
representations \(W,W'\) shows that isomorphic restrictions force
\(W'\simeq W\) or \(W'\simeq W\otimes\delta_i\). Hence at most two
isomorphism classes of rational representations occur over a fixed
restricted class. The arithmetic hypothesis presently needed is
that an admissible rational system with irreducible restriction to
\(H_i\) belongs there to the indicated class of polarized CM systems.

Choose a place above \(p\) in each of these five fields and let its
residue degree be \(f_i\), equal to one or two. From a restricted
label \(s\in S_i\) and its parameter multiset \(B_i(s)\), form all
two-element multisets \(B\) such that
\(\{b^{f_i}:b\in B\}=B_i(s)\), with multiplicity. For each \(s\)
there are at most \(f_i^2\le4\) possibilities. Enumerate the two
target entries, choose one of their \(f_i\)-th roots at each
position, and forget the order. Append two copies of each label
\((i,s,B)\). The number of labels is at most
\[
                       8\sum_i |S_i|.
\]
Galois conjugation permutes these labels, acting on \(s\) and \(B\)
and leaving the copy index fixed. The bounds on the absolute values
and their product pass to \(f_i\)-th roots. The needed integrality
passes as well, since \((p^e b^{-1})^{f_i}\) integral implies
\(p^eb^{-1}\) integral by transitivity of integrality.

To classify a rational representation, choose one of its irreducible
restrictions as above, take the corresponding restricted label, and
choose the actual parameter multiset at \(p\). The power relation at
Frobenius ensures that it is among the listed lifts. At each fixed
triple \((i,s,B)\) assign distinct copy indices to its at most two
underlying rational isomorphism classes. This gives classification
for each \(\ell\) and \(\iota\). The two-copy label set used by all
\(\ell\) and \(\iota\) was fixed before making these assignments.
The automorphic restriction assertion, the estimates for the
polarized CM class, and the level and newform criteria which supply
the requisite rank-two representation are inputs to this reduction.

\subsection{The polynomial bound uniform in the prime of coefficients}

Write \(D_K=|\operatorname{disc}K|\) and \(Y=N(e+1)\).
Suppose the two global classifications above yield label sets
\(S_1,S_2\), with
\[
             |S_1|\le c_1Y^{A_1},\qquad
             |S_2|\le c_2(D_KY)^{A_2}.
\]
For \(S_1\) one uses level \(N\). The weight bound \(H=e+1\)
would give \(N(e+2)\), and \(e+2\le2(e+1)\) allows the extra
factor \(2^{A_1}\) to be absorbed into \(c_1\).
For \(S_2\) the rational level is \(ND_K\) and \(H=e\).
The sets fit together by disjoint union \(S=S_1\sqcup S_2\).
If a representation satisfies both descriptions, give it its
first label. Equality of disjoint-union labels then still implies
isomorphism. We obtain
\[
 |S|\le(c_1+c_2D_K^{A_2})Y^{A_1+A_2},\quad
 2e\log q+1\le 5Y\log p ,
\]
using \(q=p^h,\ h\le2,\ e+1\le Y\) and \(1<\log p\).
Equations~\eqref{eq:common-trace-log} and
\eqref{eq:common-trace-subproduct}, with \(g=\sigma^{-1}\),
give the required common integer with
\[
 \log|n_p|\le 5(c_1+c_2D_K^{A_2})\,
                 |O|\,Y^{A_1+A_2+1}\log p .
\]
Both the size bound and the simultaneous subproduct inequality
apply on the same \(e,N,p,w,O\). Only the subcollections of
representations vary with the coefficient prime.

\paragraph{Dependencies of the assertion for the actual families.}
The release's \texttt{number\allowbreak{}U3\allowbreak{}U\_\allowbreak{}of\_\allowbreak{}prints} uses the CM
and rational counting arguments in the dense-threshold file
\texttt{Frobenius\allowbreak{}Trace\allowbreak{}Count\allowbreak{}Bound\allowbreak{}Of\allowbreak{}Cm\allowbreak{}Eigensystem\allowbreak{}Count}.
Its inputs include the polynomial count of regular algebraic
cuspidal systems of the given level and weights, stability under
field automorphisms, the integral Satake bounds, and compatibility
of the admitted Galois representations with the global Hecke
systems. On the rational side it also uses the polarized CM
class counts, passage to irreducible quadratic restrictions, and
the newform-admissibility comparison. The local parameter estimates
for these two families and the common-product theorem occur in
\texttt{Admissible\allowbreak{}Residues.\allowbreak{}Frobenius\allowbreak{}Trace\allowbreak{}Count\allowbreak{}Bound\allowbreak{}Of\allowbreak{}Rat\allowbreak{}Eigensystem\allowbreak{}Count}.
The archimedean bound in the associated trace data is
\(\sqrt q+2\). The preceding construction proves the common integer
from the stated automorphic inputs. The separate deformation
inequality for one Selmer character, denoted by
\texttt{One\allowbreak{}Char\allowbreak{}U\_\allowbreak{}s}, is used in Section~\ref{sec:rh-plus-nonzero}.

\section{Counting the cuspidal spectrum that receives the deformation points}
\label{sec:spectrum-deformation}

For a CM number field $K$, set $d_K=[K:\mathbb Q]$. Let $n\ge2$,
$M\ge1$, $H\ge0$ be integers, and let $m_0>\cdots>m_{n-1}$ be
integers satisfying $|m_i|\le H$. Write $K_n(M)$ for the principal
congruence subgroup of $\mathrm{GL}_n(\widehat{\mathcal O}_K)$ of level
$M$, viewed inside the finite adelic group. The collection
$\mathrm{Sig}_K(m,M)$ consists of cuspidal representations with
contragredient equal to the complex-conjugation pullback in the sense of
the source, a nonzero $K_n(M)$-fixed vector, and infinitesimal character
\[
                    I_m=\{(n-1)/2+m_i:0\le i<n\}
\]
at every embedding $K\hookrightarrow\mathbb C$. For fixed $K,n$
and any finite set $S$ of pairwise nonisomorphic members, we will obtain
\begin{equation}
 |S|\le C(K,n) M^{s}(H+1)^s,\qquad s=d_K n^2 .
 \label{eq:spectrum-count}
\end{equation}
Here the bounds in the source construction are kept separate until
the final multiplication. The quantified source theorem
\texttt{card\_k\_of\_printed\_results} permits unnamed constants and
uses the coarser exponent $s+s$ on $M(H+1)$. Both estimates follow from
the same two inequalities. The constants depend on $K,n$.

\subsection{One coefficient system for all the representations in $S$}

For $0\le k<n$ put $\nu_k=k+1+H-m_{n-1-k}$ and $H'=n+2H$.
These are nonnegative integers bounded by $H'$. They are
nonincreasing, since
\[
       \nu_{k+1}-\nu_k=1-(m_{n-2-k}-m_{n-1-k})\le0 .
\]
Let $X$ be a generic $n$ by $n$ matrix. For $0\le h<H'$ the set
$\{k:\nu_k>h\}$ is an initial segment of size $b_h$.
Take the product $f(X)$ of the top-left principal minors of orders
$b_h$, with the empty minor equal to one. This polynomial is nonzero
by evaluation at the identity, and has degree at most $nH'$.
For an upper triangular $B$, the top-left $b_h$ by $b_h$ block of
$XB$ is the corresponding block of $X$ multiplied by that of $B$.
Consequently
\[
             f(XB)=f(X)\prod_{i=0}^{n-1} B_{ii}^{\nu_i}.
\]
The space of polynomials of degree at most $nH'$ in the $n^2$
entries of $X$ is stable under the right substitution $X\mapsto XB$.
Each monomial has each exponent between zero and $nH'$, giving
dimension at most $(nH'+1)^{n^2}$. The highest-weight scalar
calculation in \texttt{Weyl.existsScalar} extracts a nonzero module $P$
of dimension at most this bound, with the central quantum determinant
acting as $\prod_k(u-k+\nu_k)$.

Use one copy at every embedding of $K$, with the determinant shift
$N_0=n+H$ in \texttt{exists\_tens}. Their external tensor product
has dimension at most $(nH'+1)^{n^2d_K}$. The logarithmic differential
at a given embedding acts through that factor. The central polynomial
therefore acts by the scalar
\[
 \prod_{k=0}^{n-1}(u-N_0-k+\nu_k)
       =\prod_{k=0}^{n-1}(u-(n-1)-m_{n-1-k}).
\]
An irreducible submodule $\Xi$ inherits this scalar, and the dimension
can only decrease. In the infinitesimal-character convention the
polynomial is evaluated at $u+(n-1)/2$. Its value is
$\prod_i(u-((n-1)/2+m_i))$, giving $\Xi$ the character $-I_m$
at every embedding. In particular the sign of the dual coefficient
system is fixed throughout $S$. Since
$n(n+2H)+1\le(n+1)^2(H+1)$, we have
\[
                 \dim\Xi\le((n+1)^2)^s (H+1)^s .
\]
This uses the highest-weight scalar calculation and tensor realization
in the source. The transformation law of $f$, dimension estimate,
reversal of indices and infinitesimal-character shift were computed
above.

\subsection{The level factor in cohomology}

Put $G=\mathrm{GL}_n(\mathbb A_{K,f})$ and
$\Gamma=\mathrm{GL}_n(K)$, acting on the left. For the fixed level
$U_0=K_n(3)$ there are finitely many double cosets
$\Gamma g_i U_0$. This is the finite-double-coset input, together
with the finite index of $U_0$ at integral level one. Write
$\Gamma_i=\Gamma\cap g_iU_0g_i^{-1}$. The arithmetic subgroup
cohomology theorem \texttt{bddArith} gives a constant $c_i$
such that for each finite-dimensional complex $\Gamma_i$-module $V$,
\[
          \sum_{j\ge0}\dim H^j(\Gamma_i,V)\le c_i\dim V .
\]
Its hypotheses apply because the $\Gamma_i$ are commensurable with
$\mathrm{GL}_n(\mathcal O_K)$ in $\Gamma$. The underlying geometric
input is a bounded arithmetic-group complex; the commensurability
comparison is part of that theorem.

Set $U'=K_n(3M)$. Every fiber of $G/U'\to G/U_0$ has
$D=[U_0:U']$ elements: the fiber above $gU_0$ is $gU_0/U'$.
For a coefficient representation of dimension $a$ the function
space on such a fiber has dimension $Da$. The inverse image of a
$\Gamma$-orbit in $G/U_0$ contributes a coinduced representation
from $\Gamma_i$, with this fiber function space as its coefficients.
Concretely a section is determined on each translate by its values
on the fiber at $g_iU_0$, subject to the $\Gamma_i$-action there.
Shapiro's lemma identifies its cohomology with that of $\Gamma_i$
on the fiber. Summing over the finitely many orbits gives
\[
 \sum_{j\ge0}\dim H^j\!\left(\Gamma,
        \operatorname{Map}(G/U',\mathbb C^a)\right)
                  \le \left(\sum_i c_i\right) Da .
\]
The action on functions is
$(\gamma f)(x)=\xi(\gamma)f(\gamma^{-1}x)$ for the coefficient
representation $\xi$. This is the cohomology quantity
\texttt{cohDim} used by the spectrum count.

Reduction modulo $3M$ embeds the quotient by its principal-congruence
kernel into $\mathrm{GL}_n(\widehat{\mathcal O}_K/
3M\widehat{\mathcal O}_K)$. The coefficient ring has
$(3M)^{d_K}$ elements by the rank of the integral lattice. The
number of invertible matrices is at most the number of all matrices,
so $D\le(3M)^s$. This proves
$\mathrm{cohDim}\le c_2 M^s a$ with
$c_2=3^s\sum_i c_i$, independent of the level and the coefficients.

Every $\pi\in S$ has the opposite infinitesimal character to $\Xi$.
The cohomological criterion \texttt{ClozelCohC} gives nonzero relative
Lie-algebra cohomology for $\pi\otimes\Xi$. Every fixed vector at
level $M$ remains fixed at level $3M$ because reduction modulo $M$
factors through reduction modulo $3M$. The cuspidal cohomology
injection \texttt{BorelInjC}, applied to the finite pairwise
nonisomorphic set $S$ and the open compact group $U'$, gives
$|S|\le\mathrm{cohDim}$ for $a=\dim\Xi$. Combining the inequalities
proves \eqref{eq:spectrum-count} with
$C(K,n)=c_2((n+1)^2)^s$. The argument applies to every finite set of
representatives and therefore also bounds the number of isomorphism
classes.

For rank three over the imaginary quadratic field used in the
single-character estimate, $s=18$. After the inverse cyclotomic
twist the weight vector is $(e-1,-1,-e-1)$, with $e\ge1$.
It decreases strictly, its opposite pairs sum to $-2$, and one
can take $H=e+1$. Thus the bound for this CM spectrum is at most
$C(K,3)2^{18}[M(e+1)]^{18}$, since $e+2\le2(e+1)$.
The count concerns global systems. Section~\ref{sec:common-trace-integer}
explains how global labels give traces at a fixed unramified prime.
The uniform passage over all residue characteristics also uses local
compatibility, Satake integrality and conjugation of the systems.

\subsection{Which mixed points contribute to the two spectra}

Write $\Omega=\overline{\mathbb Q}$, $G_F=\operatorname{Gal}(\Omega/F)$,
and use the deformation data of Section~\ref{sec:single-character}.
In this argument $K$ is the imaginary quadratic field cut out by
the quadratic multiplier $\eta$, so $\ker\eta=G_K$. The coefficient
ring $\mathcal O$ has absolute ramification index one here, and
$\iota:\overline{\mathbb Q}_\ell\simeq\mathbb C$ is fixed.
Let $r$ be a geometric point of the marked three-dimensional
representation and $X_g$ its matrix on $\ker\eta$. Index the rows
and columns by $0,1,2$. A mixed point is neither upper nor lower
triangular in the normalization. The theorem \texttt{twoPt\_shape}
gives two cases. In type I the $X_g$ act irreducibly on three
coordinates. In type II the middle coordinate and the outer pair
form blocks, and the two by two outer corner acts irreducibly.
The two cases cannot overlap because the middle coordinate is an
invariant line in type II.

Let $B$ be the chosen valuation ring of $\Omega$ above $\ell$.
The endomorphism and automorphy input \texttt{GdEndAlt} supplies one
marking with an integral finiteness conclusion and integral ordinary
endpoint statements for ranks three and two. The integral-point
theorem \texttt{exists\_ordPt\_of\_point} says precisely what is
needed to use these statements. For every permitted local deformation
condition, its marked ring, and an evaluation of the algebra of
matrix entries giving $r$, there exist a finite DVR
$\mathcal O\subset A\subset\overline{\mathbb Q}_\ell$, a unit $t$
of $\overline{\mathbb Q}_\ell$, and an ordinary $A$-point $r'$
whose scalar extension is the torus conjugate of $r$ by
$T=\operatorname{diag}(t,1,t^{-1})$. Its construction uses finiteness
of the weight-zero algebra, a graded finite-point theorem and the
local deformation conditions. The same marking works for every such
point. We spell out the use of conjugation because it changes the
coordinates of the outer corner.

On $\ker\eta$ one has $X'_g=TX_gT^{-1}$. In type I this carries
invariant subspaces bijectively to invariant subspaces. The ordinary
irreducibility theorem \texttt{matIrr\_open\_gd} upgrades the resulting
irreducibility to each open subgroup of $G_{\mathbb Q}$ contained in
$\ker\eta$. This stronger conclusion is an explicit hypothesis of
the rank-three ordinary endpoint. Applying it gives conjugate
self-dual automorphy for the restriction $\rho'=r'|_{G_K}$ after
the twist $\rho'\otimes\chi_\ell^{-1}$, including the source's
local witness for the continuous representation. Conjugating by
$T^{-1}$ intertwines this twist with
$\rho\otimes\chi_\ell^{-1}$. The same automorphy conclusion therefore
holds at the original point. Admission to the counted spectrum also
requires the level conditions and the local assertion denoted
$\mathrm{Unl}(\ell,M,(e,0,-e),\rho)$. The separate result
\texttt{admissible\_typeI\_even\_level\_alt} supplies these for the
admissible type-I data.

In type II, the point theorem gives $(X_g)_{11}=1$ for $g\in G_K$.
For an extension $W:G_{\mathbb Q}\to\mathrm{GL}_2(
\overline{\mathbb Q}_\ell)$ agreeing with the outer corner on $G_K$,
the marking requires $\operatorname{tr}W(c)=0$ at the designated
complex conjugation $c$. Let $D_t=\operatorname{diag}(t,t^{-1})$
and $W'=D_t W D_t^{-1}$. The odd-weight matrix entries of $X'_g$
still vanish, its middle entry is one, and its corner on $G_K$ is
$D_t\operatorname{corner}(X_g)D_t^{-1}=W'(g)$. Its corner remains
irreducible. Also $\operatorname{tr}W'(c)=0$. The rank-two ordinary
endpoint gives a classical newform $f$ of weight $2e+1$ attached to
$W'\otimes\chi_\ell^{-e}$, together with the witness
$\mathrm{WitAQ}(\iota,W')$ for weak admissibility. Conjugacy
transports both assertions to $W$. Thus the untwisted rank-three
restriction on $G_K$ has the form $\mathbf1\oplus W|_{G_K}$ and
the rational spectrum supplies this branch's labels. In the assembled
integer bound, the rank-two level is $M|\operatorname{disc}K|$ and
the integral weight vector is $(e,-e)$. The rank-three
CM count above applies to the type-I branch. The rational count
and its quadratic restriction argument appear in
Section~\ref{sec:common-trace-integer}.

\paragraph{Scope of the arithmetic inputs.}
The cohomological criterion, cuspidal cohomology injection, finite
double coset theorem, arithmetic-stabilizer bound, and scalar
highest-weight calculation enter the spectrum count. The passage from
integral ordinary points to automorphy uses the separate endpoint
theorem, proven in the release through linked ordinary systems,
descent, and unitary automorphy. The conjugacy argument above proves
the transfer to the original point. It uses the marked shape and
open-subgroup irreducibility theorems as inputs. The subsequent
rigidity argument applies adjoint Selmer vanishing to type I and
rank-two adjoint and Kato inputs to type II.

\paragraph{Source map.}
The definitions of the spectrum and weights are in
\texttt{AutomorphicForms.IsoOfFrob}. The coefficient construction is
\texttt{TWeylC}, with the minor polynomial in \texttt{TopPol} and the
tensor and scalar steps in \texttt{ExistsTens} and \texttt{ExistsScalar}.
The orbit argument is \texttt{TemperedLocalComponents.CohBoundFree}
and the final cohomological comparison is \texttt{TCardK}. The
rank-three specialization is in
\texttt{FrobeniusTraceCountBoundOfCmEigensystemCount}.
The point statements are \texttt{UnitaryThreeGroups.PtShape},
\texttt{CurveDVRSpecialization.WeightZeroAlgebraLeEntryAlgebra},
\texttt{MatIrrOpenGd}, \texttt{AutEnd3AltAlt}, and
\texttt{AutomorphyU3OfEndoscopy}. The precise modular-form
normalization is in \texttt{AutomorphyDescent.AutNewformW}.

\footnote{\href{https://github.com/Deicyde/riemann-hypothesis/blob/40e56bab026552e089d7afcfeff08d2defb9c005/src/Development/AutomorphicArithmetic/AutomorphicVarieties/ModularityArgument/ThresholdArguments/DenseThresholdCases/TCardK.lean}{The counting deduction},
\href{https://github.com/Deicyde/riemann-hypothesis/blob/40e56bab026552e089d7afcfeff08d2defb9c005/src/Development/AutomorphicArithmetic/AutomorphicVarieties/AutomorphicForms/TemperedLocalComponents/CohBoundFree.lean}{the level and cohomology exponents}, and
\href{https://github.com/Deicyde/riemann-hypothesis/blob/40e56bab026552e089d7afcfeff08d2defb9c005/src/Development/AutomorphicArithmetic/AutomorphicVarieties/ModularityArgument/ThresholdArguments/DenseThresholdCases/AutomorphyU3OfEndoscopy.lean}{the integral endpoint comparison}.}

\section{Irreducibility of the ordinary point on every open subgroup}
\label{sec:spectrum-open-irreducibility}

The rank-three automorphy endpoint in
Section~\ref{sec:spectrum-deformation} needs irreducibility on every
open subgroup contained in $G_K$. We prove this for the ordinary
integral type-I points in that section. The essential restriction on
the residual character is visible in the proof.

Write $G=G_{\mathbb Q}$, $H=\ker\eta=G_K$, and $I_\ell$ for inertia at
the chosen valuation above $\ell$. Let $\mathcal O\subset A\subset
\overline{\mathbb Q}_\ell$ be the finite DVR of the ordinary point. Use
$k_A$ for its residue field and write
$\Psi=\chi_\ell^e\psi_{\mathrm{Gal}}$ in the coefficient units, with
$\psi_{\mathrm{Gal}}$ of Dirichlet level $N$ and $e\ge1$. On $H$ the
matrix part of the point is an ordinary representation, denoted by
$R$. In the marked integral basis its reduction has diagonal entries
$\overline\Psi,1,\overline\Psi^{-1}$ on $H$ (equivalently one may
use just this residual trace identity below). At the decomposition
group at $\ell$ intersected with $H$, the ordinary flag gives a
basis in which the matrices are upper triangular. On inertia the
diagonal characters in that basis are $\Psi,1,\Psi^{-1}$.
The arithmetic input \texttt{dat\_facts} gives
\begin{equation}
 I_\ell\subset H,\qquad
 \exists x\in I_\ell:\quad \Psi(x)^2-1\in\mathcal O^\times .
 \label{eq:open-residue-separation}
\end{equation}
In particular this difference also has nonzero image in $k_A$.
The hypothesis that $R$ acts irreducibly on $H$ is precisely type I.

\subsection{A three-dimensional permutation argument}

Let $H$ act irreducibly on $V=L^3$ over a field $L$, let $U$ be a
normal subgroup of $H$, and suppose some $u\in U$ has three distinct
eigenvalues in $L$. Suppose $R(u)$ and $R(x)$ are upper triangular
in a common basis for some $x\in H$. If $U$ acts reducibly then there
is a $g\in H$ for which
\begin{equation}
               \operatorname{tr} R(gx^j)=0\qquad(j\ge0).
 \label{eq:open-vanishing-coset}
\end{equation}
To see this, diagonalize $R(u)$ by an upper triangular change of
basis. For each diagonal position its eigenvector is obtained by
choosing entry one at that position, zero below it, and solving
successively in the rows above it. The differences of diagonal
entries that must be divided by are all nonzero. Thus $R(x)$
remains upper triangular. Denote the three resulting eigenlines
by $L_0,L_1,L_2$.

Any $U$-invariant subspace is a sum of some of these lines: the
projection to the $i$th line is the polynomial in $R(u)$ given by
$\prod_{a\ne i}(R(u)-\lambda_a)/(\lambda_i-\lambda_a)$.
Take a nonzero proper $U$-invariant subspace $W$. If it is a line,
then a $U$-invariant eigenline has been found. If it is a plane,
irreducibility on $H$ gives $h\in H$ with $hW\ne W$.
Normality makes $hW$ another $U$-invariant plane. Their intersection
is a $U$-invariant line. Every translate by $H$ of a $U$-invariant
line is again $U$-invariant, hence one of the three eigenspaces.
The span of its orbit is nonzero and $H$-invariant. Irreducibility
implies that the orbit consists of all three lines. Thus the action
on the three lines gives a transitive subgroup of $S_3$.
Its order is divisible by three, so it contains a three-cycle
by Cauchy's theorem. Choose $g$ realizing this cycle.

Every $R(h)$ is now a monomial matrix. The permutation of an upper
triangular invertible monomial matrix fixes each index: if it sends
column $i$ to row $p(i)$, then $p(i)\le i$ for every $i$, and
$\sum_i p(i)=\sum_i i$. In particular $R(x)$ is diagonal.
$R(gx^j)$ therefore has the same permutation as $R(g)$ and all its
diagonal entries vanish. This proves \eqref{eq:open-vanishing-coset}.

\subsection{Inertia supplies both matrices}

Let $V_0$ be an arbitrary open subgroup of $G$ contained in $H$,
and set $U=\bigcap_{g\in G}gV_0g^{-1}$. Openness in the profinite
group gives finite index, so the normal core $U$ also has finite
index. It suffices to prove irreducibility on $U$, as every
$V_0$-invariant subspace is $U$-invariant. Fix the element $x$ in
\eqref{eq:open-residue-separation}. The cyclotomic inertia theorem
over $\mathbb Q_\ell$ gives $\sigma\in I_\ell$ with
$\chi_\ell(\sigma)=1+\ell$. Put
\[
    m=[G:U]\, |(\mathbb Z/N\mathbb Z)^\times|,\qquad u=\sigma^m .
\]
Here $m>0$, $u\in U\cap I_\ell$, and $\psi_{\mathrm{Gal}}(u)=1$,
since its order divides the order of the Dirichlet unit group.
It follows that $\Psi(u)^2=(1+\ell)^{2em}\ne1$ in characteristic
zero. In the ordinary basis the three eigenvalues of $R(u)$ are
therefore distinct. Both $u$ and $x$ belong to inertia, so the
ordinary flag triangularizes both matrices.

If $U$ acted reducibly, \eqref{eq:open-vanishing-coset} would produce
$g\in H$. Compute its vanishing traces in the integral basis, and
reduce modulo the maximal ideal of $A$. Put $a=\overline\Psi(g)$,
$b=\overline\Psi(x)$ in $k_A^\times$. The residual trace identity
gives
\[
                 ab^j+1+(ab^j)^{-1}=0\quad(j\ge0).
\]
Then each of $a,ab,ab^2$ is a root of $X^2+X+1$.
Since $b^2\ne1$ by \eqref{eq:open-residue-separation}, $b\ne1$
and the three units are pairwise distinct. A quadratic polynomial
over a field has at most two distinct roots. This contradiction proves
irreducibility on $U$ and hence on $V_0$. The proof works for each
open $V_0\subseteq H$ and supplies the required quantifier.

\paragraph{Source and inputs.}
This is \texttt{matIrr\_open\_gd}. The argument above derives the open
irreducibility conclusion from the residual trace, the ordinary flag,
\eqref{eq:open-residue-separation}, and cyclotomic surjectivity on
inertia. For the given good deformation data the residual and ordinary
conditions belong to \texttt{OrdPt}; \texttt{dat\_facts} supplies the
displayed separation by the character and unramifiedness conditions
of the datum. The characteristic-zero and residual eigenvalue
conditions have different uses. The power $u$ distinguishes the
three eigenlines in characteristic zero, and $x$ distinguishes three
residual values after applying the Frobenius-independent trace test.
\footnote{\href{https://github.com/Deicyde/riemann-hypothesis/blob/40e56bab026552e089d7afcfeff08d2defb9c005/src/Development/AutomorphicArithmetic/AutomorphicVarieties/CurveDVRSpecialization/MatIrrOpenGd.lean}{The invariant subspaces, trace obstruction and arithmetic specialization},
\href{https://github.com/Deicyde/riemann-hypothesis/blob/40e56bab026552e089d7afcfeff08d2defb9c005/src/Development/AutomorphicArithmetic/AutomorphicVarieties/UnitaryThreeGroups/DatFacts.lean}{inertia and residual separation}.}

\section{Ordinary rank-three points and automorphic descent}
\label{sec:ordinary-global-descent}

Let $\Omega$ be an algebraic closure of $\mathbf Q$, let $K$ be the marked imaginary quadratic field and write $G_F=\operatorname{Gal}(\Omega/F)$ for an embedded subfield. Fix the coefficient ring $\mathcal O$ at $\ell$, $\iota:\overline{\mathbf Q}_\ell\simeq\mathbf C$, the positive weight parameter $e$, the even positive modulus $N$, and the character $\psi$ of the good datum. Write $\epsilon$ for the $\ell$-adic cyclotomic character, $d$ for the character obtained from $\psi$ on the Galois group, and $\Psi=\epsilon^e d$, after embedding the coefficients. The character $\eta$ cuts out $K$. Fix the marked conjugation $c$, the place $B$ over $\ell$, and the element $g_1$ supplied by the good frame. For every finite $\mathcal O$-subalgebra $A\subset\overline{\mathbf Q}_\ell$ which is a discrete valuation ring and every marked ordinary point
\[
                R:G_{\mathbf Q}\longrightarrow
                 \mathrm{GL}_3(A)\rtimes\{1,-1\},
\]
with the sign acting by the framed inverse-transpose involution, let $\rho$ be its matrix representation on $G_K$ after extending the coefficients. Suppose that these matrices act irreducibly on every open subgroup of $G_{\mathbf Q}$ contained in $\ker\eta$. Section~\ref{sec:spectrum-open-irreducibility} proves this condition for the indicated type I points. The conclusion of the descent is that
\[
                       r=\rho\otimes\epsilon^{-1}
\]
is attached to a regular algebraic conjugate self-dual cuspidal representation over $K$ and satisfies the additional local assertion below. We give the reduction to the linked automorphy input over the auxiliary field and the cyclic base-change and attachment inputs in [1--6].

\subsection{Field, marked point and polarization}

Put $M=\ell N|\operatorname{disc}K|$. Any inertia group with nontrivial $\eta$ lies at a prime dividing $|\operatorname{disc}K|$. Indeed, at a prime away from the discriminant two embeddings of $\mathcal O_K$ which agree after reduction at a chosen prime of $\Omega$ agree in $\Omega$. This follows from the invertibility there of the discriminant pairing, or from the unramified local decomposition of $\mathcal O_K$. An inertia element acts identically on the residue field, so the original embedding and its translate agree. This gives the ramification condition at $M$, and $\ell N\mid M$. The parity assumption $\psi(-1)=(-1)^e$ and $\epsilon(c)=-1$ give $\Psi(c)=1$. Complex conjugation squares to the identity. Thus the condition at $B$, the ramification condition, and the divisibility and positivity conditions used in the link hold with the original $e,N$.

The imported frame theorem supplies an abelian CM field $L/\mathbf Q$, a frame $\Phi$, and a Galois link with $\Phi$ satisfying the two size and admissibility predicates and the seed condition of [1]. The link gives $G_L\subset\ker\eta$. For any allowed embedding of $K$, its image is the fixed field of $\ker\eta$ because the two embeddings of a quadratic Galois extension have the same image. Hence every element of that image is fixed by $G_L$ and belongs to $L$. This constructs $K\hookrightarrow L$ compatibly with the given embeddings in $\Omega$. It gives a finite abelian Galois extension $L/K$. Since $G_L$ is open in $G_{\mathbf Q}$ and lies in $\ker\eta$, $\rho|_{G_L}$ is irreducible. A nonzero scalar twist preserves the invariant subspaces. In particular $r|_{G_L}$ is irreducible.

Here is the precise use of automorphy over $L$. For every finite discrete valuation coefficient extension as above, the rank-three conclusion of the link applies to a normalized frame lift whose local inertia characters are $\Psi^{a_i}$, $(a_1,a_2,a_3)=(1,0,-1)$. It applies to any irreducible representation over $L$ whose matrices are the image of that lift evaluated on the map from $G_L$ to the frame group. Construction of the normalized lift from $R$ gives those matrices as $\rho|_{G_L}$, and its three inertia formulas give exactly these powers. The ring map $\mathcal O\to A$ is local by integrality. The rank-three twist built into the link is precisely $\epsilon_L^{-1}$. Consequently the conclusion over $L$ applies to $r|_{G_L}$. In [1] this link uses an ordinary rank-two seed, a rank-three density theorem for the unitary family, and propagation of the adjoint lift of the seed. These are global inputs even when one starts with an already irreducible marked point.

Let $J$ be the rank-three reversal matrix, so $J^2=1$. The normalized matrix component at $c$ is the identity, and its sign is $-1$. For $s\in G_K$ and $s'=csc^{-1}$ the equality $s'c=cs$ in the semidirect product therefore gives
\[
             \rho(s')=J\,{}^t\!\rho(s)^{-1}J,\qquad
             {}^t\!\rho(s)J\rho(s')=J .
\]
As $\epsilon(s')=\epsilon(s)$, twisting both factors gives the exact multiplier
\begin{equation}
 {}^t\!r(s)Jr(s')=\epsilon(s)^{-2}J
                 =\epsilon(s)^{1-3}J.                  \label{eq:ord-descent-pair}
\end{equation}
The multiplier on the real subfield $K^+$ is the same cyclotomic power. It is algebraic by the cyclotomic character construction, with local exponent $-2$ and trivial finite factor, and its value at each complex conjugation is $1$. The conjugation $c$ fixes $K^+$ because its image under every relevant real embedding is real.

Continuity after extending the coefficients can also be checked explicitly. Every element of $A$ is integral over $\mathbf Z_\ell$ and has norm at most one. Write $\ell=u\varpi^h$ for a unit $u\in A^\times$ and uniformizer $\varpi$. Then every element of $\mathfrak m_A^{hj}$ is divisible by $\ell^j$ and has norm at most $|\ell|_\ell^j$. Each entry of $R$ is locally constant modulo $\mathfrak m_A^{hj}$ by the adic continuity condition of the point. Taking $j\to\infty$ gives continuity of the entries of $\rho$. The inverse entries are continuous by applying the same argument at $s^{-1}$. Multiplication by the continuous scalar $\epsilon^{-1}$ gives continuity of $r$.

\subsection{Inertia away from \texorpdfstring{$\ell$}{l} and flags above it}

Fix a finite place $v\nmid\ell$ of $K$, a prolongation, and its decomposition and inertia groups $D_v,I_v$. The inertia group acts trivially on every $\ell^j$-th root of unity. Indeed reduction is injective on roots of unity of order prime to the residue characteristic, and inertia acts trivially on the reductions. Thus $\epsilon(I_v)=1$. The character $\Psi|_{G_K}$ is continuous and its restriction to $I_v$ has image in the finite image of $d$. If $s\in I_v$ and $\Psi(s)=1$, then $\eta(s)=1$ and the away-from-$\ell$ condition in the marked ordinary point gives $R(s)=1$. The same holds for $r(s)$ since $\epsilon(s)=1$.

The image $\Psi(I_v)$ is a cyclic group of roots of unity, of order $m$, say. The representation of $I_v$ factors through this group. Choose $u\in I_v$ mapping to a generator $\xi$. The operator $r(u)$ satisfies $X^m-1$, which is separable in characteristic zero. In an eigenbasis its eigenvalues are $\xi^{b_1},\xi^{b_2},\xi^{b_3}$ with $b_i\ge0$. For every inertia element $s$ it follows that
\[
       \det(X-r(s))=\prod_{i=1}^3(X-\Psi(s)^{b_i}).
\]
Each $\Psi^{b_i}$ extends continuously to $D_v$ by restriction of the global character. This is the full local characteristic-polynomial condition called \texttt{IneAb} in [2].

For $v\mid\ell$, choose a place of $\Omega$ over $v$. The frame conjugacy condition transports it to the marked place $B$. Conjugating the ordinary flag back supplies the upper triangular form of $\rho$ there. For a transporter with positive sign this uses matrix conjugation. For a transporter with negative sign it also uses $b\mapsto J\,{}^tb^{-1}J$. This involution takes an upper triangular matrix to an upper triangular matrix and sends its diagonal $(x_1,x_2,x_3)$ to $(x_3^{-1},x_2^{-1},x_1^{-1})$. Since $a_{4-i}=-a_i$, it preserves the prescribed three powers of $\Psi$ in their order. The ordinary characters of the successive one-dimensional quotients of $\rho$ therefore restrict on inertia to $\epsilon^{ea_i}d^{a_i}$. Twisting them all by $\epsilon^{-1}$ gives
\[
        \epsilon^{\,e-1}d,\qquad
        \epsilon^{-1},\qquad
        \epsilon^{-e-1}d^{-1}.
\]
Transport to any other prolongation over $v$ by a Galois element conjugates the flag and the decomposition group together. This proves the flag hypothesis for all prolongations in [2].

The descent of the local window uses a further property of the selected frame. Write $U_1(r,v)$ if at each prolongation there is a decomposition-group flag with diagonal characters $\kappa_i$ such that $\kappa_i\ne\kappa_j\epsilon$ for every $i,j$. Write $U_\infty(r,v)$ for the analogous condition with $\kappa_i^a\ne(\kappa_j\epsilon)^a$ for every integer $a\ge1$. The release proves for the even level frame that for each $w\mid\ell$ of $L$, with $v=w|_K$,
\[
             U_1(\rho,v)\ \Longrightarrow\
                     U_1(\rho|_{G_L},w).
\]
This preservation uses the linked noncoincidence theorem in [2]. A common scalar twist cancels from the two sides of every character equality, including their positive powers. Thus the same implication holds for $r$.

\subsection{Correcting the descent at a prime cyclic layer}

We detail the part of the global argument that identifies the specified representation downstairs. For a cyclic extension $E/F$ of prime degree in the solvable CM tower, let $r_F$ be the restriction of $r$ to $G_F$ and $r_E=r_F|_{G_E}$. Both are irreducible, since their groups contain $G_L$. Suppose that a regular algebraic cusp form $\Pi$ over $E$ is attached to $r_E$. At all but finitely many places, Frobenius classes at conjugate places over $F$ have conjugate matrices under $r_F$. Thus the attached Satake multisets of $\Pi$ at those places agree with their conjugates. The prime-degree cyclic descent assertion supplies a cusp form $\pi_0$ over $F$, weakly base changing to $\Pi$ and carrying the descended archimedean parameters [3,4]. Every complex embedding of $F$ extends to $E$, so $\pi_0$ inherits regular algebraicity. The attachment of a Galois representation downstairs requires an additional argument. The dual and the conjugate predicted by the polarization have weak base changes related by the predicted character, using the upstairs Frobenius data. Two descents of the same weak base change differ by a finite norm character, by the fiber assertion of the prime-degree theorem. In the release this norm character is converted into a Galois character on $G_F$ and extended to $G_{F^+}$. Dividing the previous polarizing character by this finite-order correction leaves its value constant on complex conjugations. Comparison of unramified Satake multisets and strong multiplicity one then polarize $\pi_0$ with this corrected character. The polarized Galois-attachment theorem supplies a continuous semisimple $r_0$ attached to $\pi_0$ [5].

The restriction of $r_0$ to $G_E$ is semisimple by normality and finite index. On restriction, its unramified parameters are those of $\Pi$, by the weak base-change identity. Uniqueness of the Satake multiset gives equality of the characteristic polynomials of $r_0|_{G_E}$ and $r_E$ on almost all Frobenius classes. The Frobenius isomorphism criterion (Chebotarev density and the characteristic-zero semisimple character criterion) gives
\[
                     T\,r_0(h)T^{-1}=r_F(h)
                         \qquad(h\in G_E).
\]
To see exactly what remains on $G_F$, put $r'_0=T r_0 T^{-1}$. For $g\in G_F$ both $r_F(g)$ and $r'_0(g)$ act by the same conjugation on $r_E(G_E)$. For $h\in G_E$ this follows by applying the equality to $ghg^{-1}\in G_E$. Hence $r'_0(g)^{-1}r_F(g)$ commutes with $r_E(G_E)$. Schur's lemma makes it a scalar $\delta(g)$ since the coefficients are algebraically closed and $r_E$ is irreducible. Multiplicativity of both representations gives $\delta(gg')=\delta(g)\delta(g')$. The scalar is one on $G_E$, so $\delta$ is a finite character of $\operatorname{Gal}(E/F)$. Use its corresponding finite-order Hecke character in the attachment convention of the release to twist $\pi_0$. Twisting preserves cuspidality and regular algebraicity and multiplies the attached representation by $\delta$. This yields an automorphic representation attached to $r_F$ itself.

Iteration is legitimate on this tower. The finite abelian group $\operatorname{Gal}(L/K)$ has a composition series with cyclic prime quotients. The intermediate fields contain $K$, and are CM. Since $L/\mathbf Q$ is abelian, complex conjugation induces a central automorphism of each intermediate field; its fixed field there is totally real, while its restriction to $K$ is nontrivial. Thus each intermediate field is quadratic over its totally real subfield. At every step the restriction of the original $r$ to the upper group is irreducible, and the original polarization restricts. The release proves this induction more generally for solvable CM extensions [3]. Applying it down to $K$ produces a regular algebraic cusp form attached to $r$. Equation~\eqref{eq:ord-descent-pair} then identifies its conjugate with its contragredient. In the attachment normalization the factor $\epsilon^{1-3}$ is precisely the prescribed duality twist. The release applies its comparison of dual Satake parameters and strong multiplicity one, using the trivial auxiliary Hecke character [5].

\subsection{The additional local assertion}

For a continuous representation, the assertion called \texttt{WitA} consists of two clauses [6]. For each $v\mid\ell$, $U_1(r,v)$ implies $U_\infty(r,v)$. For every regular algebraic cusp form $\pi$ attached to $r$, there exists a real $t$ such that, simultaneously for every finite place $v$ with $v\nmid\ell$ or $U_\infty(r,v)$, its local component satisfies the source condition $\operatorname{WinAt}(\pi,t,v)$. In this condition the support satisfies the toral condition $\operatorname{TorAt}(\pi,v)$, and for every character constituent $\xi$ of any such local cuspidal support one has, with $q_v$ the residue cardinality and $\varpi_v$ the chosen uniformizer,
\[
         q_v^{\,t-1/2}<|\xi(\varpi_v)|<q_v^{\,t+1/2}.
\]
The parameter $t$ may depend on $\pi$ but does not depend on $v$.

For the first clause take a place $w$ above $v$. The even-level preservation established above gives $U_1(r|_{G_L},w)$ from $U_1(r,v)$. The first clause over $L$ gives $U_\infty(r|_{G_L},w)$. The flag over $K$ descends this power inequality. Restricting a hypothetical equality of two powers at $v$ would give the equality at $w$. The multiset of diagonal characters of any two flags on the same restriction is equal, by equality of traces and linear independence of distinct characters in characteristic zero, so a change of the flag above $w$ cannot remove the equality. For the second clause use the upstairs form with the global parameter $t$ furnished by its local assertion. At a prime-degree step weak base change relates the two attached forms. The window comparison of [6] applies at every $v\nmid\ell$ using the inertia characteristic-polynomial condition, and at $v\mid\ell$ satisfying $U_\infty$ using the ordinary comparison at $\ell$. The same $t$ works at each such $v$. This is the second input beyond the descent of the cusp form. The linked density theorem over $L$, the polarized attachment used in the cyclic step, and this local window comparison are used as theorem inputs here.

\paragraph{Numbered sources.}
\begin{enumerate}
\item The release's \href{https://github.com/Deicyde/riemann-hypothesis/blob/40e56bab026552e089d7afcfeff08d2defb9c005/src/Development/AutomorphicArithmetic/AutomorphicVarieties/ModularityArgument/ThresholdArguments/DenseThresholdCases/GdEndOf.lean}{good frame and rank-three application} and \href{https://github.com/Deicyde/riemann-hypothesis/blob/40e56bab026552e089d7afcfeff08d2defb9c005/src/Development/AutomorphicArithmetic/AutomorphicVarieties/ModularityArgument/ThresholdArguments/DenseThresholdCases/AutEnd3Alt.lean}{rank-three endpoint proof} give the marked quantifiers and the choice of twist. The \href{https://github.com/Deicyde/riemann-hypothesis/blob/40e56bab026552e089d7afcfeff08d2defb9c005/src/Development/AutomorphicArithmetic/AutomorphicVarieties/ModularityArgument/ThresholdArguments/DenseThresholdCases/SymTop.lean}{linked density theorem} identifies the rank-two seed and unitary rank-three inputs.
\item The proofs of the \href{https://github.com/Deicyde/riemann-hypothesis/blob/40e56bab026552e089d7afcfeff08d2defb9c005/src/Development/AutomorphicArithmetic/AutomorphicVarieties/ModularityArgument/UpperTriangularPlaces/OrdinaryFlagThree.lean}{ordinary flag at each place}, the \href{https://github.com/Deicyde/riemann-hypothesis/blob/40e56bab026552e089d7afcfeff08d2defb9c005/src/Development/AutomorphicArithmetic/AutomorphicVarieties/ModularityArgument/AutomorphyLifting/AutomorphyDescent/InertiaRankTwo/LocalConditionRankThreeLinkOfInertiaOrdinary.lean}{away-from-$\ell$ inertia condition}, and the \href{https://github.com/Deicyde/riemann-hypothesis/blob/40e56bab026552e089d7afcfeff08d2defb9c005/src/Development/AutomorphicArithmetic/AutomorphicVarieties/ModularityArgument/ThresholdArguments/DenseThresholdCases/DnEnds.lean}{even-level preservation for the linked field} give the separate local hypotheses. The \href{https://github.com/Deicyde/riemann-hypothesis/blob/40e56bab026552e089d7afcfeff08d2defb9c005/src/Development/AutomorphicArithmetic/AutomorphicVarieties/CurveDVRSpecialization/OrthogonalPairingFromIntertwiner.lean}{matrix pairing and coefficient continuity lemmas} prove the corresponding assertions before descent.
\item The release's \href{https://github.com/Deicyde/riemann-hypothesis/blob/40e56bab026552e089d7afcfeff08d2defb9c005/src/Development/AutomorphicArithmetic/AutomorphicVarieties/ModularityArgument/ThresholdArguments/DenseThresholdCases/CuspUnique.lean}{prime cyclic step} and \href{https://github.com/Deicyde/riemann-hypothesis/blob/40e56bab026552e089d7afcfeff08d2defb9c005/src/Development/AutomorphicArithmetic/AutomorphicVarieties/ModularityArgument/BaseChangeUnits/SolvableEquivalenceOfBaseChange.lean}{solvable degree induction} give the arithmetic formulation in terms of the specified $r$ and its restriction.
\item J.~Arthur and L.~Clozel, \href{https://www.math.toronto.edu/arthur/pdf/30.pdf}{\emph{Simple Algebras, Base Change, and the Advanced Theory of the Trace Formula}}, Annals of Mathematics Studies 120 (1989), Chapter 3, for prime-degree cyclic base change. L.~Clozel and C.~S.~Rajan, \href{https://doi.org/10.1515/crelle-2020-0023}{\emph{Solvable base change}}, Theorem 3, describes the image and fibers for general solvable base change. The precise prime-degree Satake and infinity-type statements used here are in the release's \href{https://github.com/Deicyde/riemann-hypothesis/blob/40e56bab026552e089d7afcfeff08d2defb9c005/src/Development/AutomorphicArithmetic/AutomorphicVarieties/ModularityArgument/BaseChangeUnits/CycCOfCyclic.lean}{cyclic base-change hypotheses}.
\item The release's \href{https://github.com/Deicyde/riemann-hypothesis/blob/40e56bab026552e089d7afcfeff08d2defb9c005/src/Development/AutomorphicArithmetic/AutomorphicVarieties/ModularityArgument/ThresholdArguments/DenseThresholdCases/SVRSol.lean}{polarized attachment at the cyclic descent step} and \href{https://github.com/Deicyde/riemann-hypothesis/blob/40e56bab026552e089d7afcfeff08d2defb9c005/src/Development/AutomorphicArithmetic/AutomorphicVarieties/ModularityArgument/ThresholdArguments/DenseThresholdCases/VWCsdOf.lean}{passage from a polarized Galois representation to a conjugate self-dual cusp form} supply the two different attachment assertions.
\item The \href{https://github.com/Deicyde/riemann-hypothesis/blob/40e56bab026552e089d7afcfeff08d2defb9c005/src/Development/AutomorphicArithmetic/AutomorphicVarieties/ModularityArgument/AutomorphyLifting/WeakOfAtt.lean}{definitions of the flag, power condition and automorphic window}, the \href{https://github.com/Deicyde/riemann-hypothesis/blob/40e56bab026552e089d7afcfeff08d2defb9c005/src/Development/AutomorphicArithmetic/AutomorphicVarieties/ModularityArgument/ThresholdArguments/DenseThresholdCases/TAUpDownOf.lean}{descent of the local window}, the \href{https://github.com/Deicyde/riemann-hypothesis/blob/40e56bab026552e089d7afcfeff08d2defb9c005/src/Development/AutomorphicArithmetic/AutomorphicVarieties/ModularityArgument/LeviSmoothLocalization/WinUpDown.lean}{strict uniformizer bounds}, and the \href{https://github.com/Deicyde/riemann-hypothesis/blob/40e56bab026552e089d7afcfeff08d2defb9c005/src/Development/AutomorphicArithmetic/AutomorphicVarieties/ModularityArgument/ThresholdArguments/DenseThresholdCases/DnDown.lean}{final descent of both witnesses} specify the local part of the conclusion.
\end{enumerate}

\section{The deformation inequality for a single character}
\label{sec:single-character}

The estimate used in Section~\ref{sec:rh-plus-nonzero} first bounds one
strict Selmer group by distances between traces. We give the lower
comparison with the saturated cotangent module, the multiplicities in
the upper comparison, and the passage to isomorphism classes. We state
which deformation and automorphy results enter these comparisons.

\subsection{The arithmetic hypotheses and the estimate}

Take an algebraic closure \(\Omega/\mathbb Q\), an imaginary quadratic
subfield \(K\), a prime \(\ell\), and the integers \(\mathcal O\) in a
finite extension \(F/\mathbb Q_\ell\) inside \(\overline{\mathbb Q}_\ell\).
Write \(e_\mathcal O=v_{\mathcal O}(\ell)\), with
\(v_{\mathcal O}(\varpi)=1\). Fix an even positive integer \(N\), an
integer \(e\ge1\) and an \(\mathcal O\)-valued Dirichlet character
\(\psi\bmod N\). Write \(\psi_\ell\) for its component at the
\(\ell\)-part of the level and \(\psi'=\psi/\psi_\ell\) on the units.
The good-data assumptions of the source are
\[
\begin{gathered}
 \ell\ne2,\qquad \ell\nmid\operatorname{ord}(\psi),\qquad
 \psi_\ell^2=1,\qquad \psi(-1)=(-1)^e,\qquad
 \ell-1\nmid 2e,\\
 \ell\nmid |\operatorname{disc}K|\,h_K,\qquad
 \gcd(|\operatorname{disc}K|,N)=1,\\
 \ker\overline\chi_\ell\cap\ker\psi'\ \not\subseteq G_K .
\end{gathered}
\]
Here \(\overline\chi_\ell\) is the mod-\(\ell\) cyclotomic character,
\(h_K\) is the class number, and the last line uses Galois
characters on \(G_{\mathbb Q}\). The prime-to-\(\ell\) order assumption
makes \(\psi'\) precisely the character called \(\mathrm{chiP}\)
in the source. In particular \(\ell\ge5\).

Put \(\Gamma=G_{\mathbb Q}\), let \(\eta\) be the quadratic character
with kernel \(G_K\), and put \(\Psi=\chi_\ell^e\psi\),
\(\alpha=\Psi^{-1}\), and \(W=(F/\mathcal O)(\alpha)\).
Let \(\mathcal D\) consist of the decomposition groups at places above
primes dividing \(\ell N\) and the inertia groups at all finite primes.
Define
\[
 V=\ker\left(H^1_{\rm cont}(\Gamma,W)\longrightarrow
                      \prod_{D\in\mathcal D}H^1_{\rm cont}(D,W)\right).
\]
Fix a field isomorphism \(\iota:\overline{\mathbb Q}_\ell\to\mathbb C\)
and an element \(\sigma\in G_K\). Set
\[
 t=\Psi(\sigma)+1+\Psi(\sigma)^{-1},\qquad
 v_\ell(z)=-\log_\ell |z|_\ell\quad(z\ne0).
\]
The admissible \(\rho:G_K\to{\rm GL}_3(\overline{\mathbb Q}_\ell)\)
are the two classes used in the integer estimate. The first consists
of irreducible representations of level \(N\), with the source's
weight tuple \((e,0,-e)\) and its unlinked local condition. Their
twists \(\rho\otimes\chi_\ell^{-1}\) satisfy the conjugate-self-dual
automorphy condition. In the second class,
\(\rho\simeq{\bf1}\oplus W_2|_{G_K}\), where \(W_2\) is an irreducible
two-dimensional \(\overline{\mathbb Q}_\ell\)-representation of
\(G_{\mathbb Q}\). The source requires its level condition
at \(N|\operatorname{disc}K|\), weights \((e,-e)\), determinant of
cyclotomic weight zero, a newform of weight \(2e+1\) attached to
\(W_2\otimes\chi_\ell^{-e}\), and the rational automorphic witness
called \(\mathrm{WitAQ}\). These are the predicates
\(\mathrm{Adm3U}_s\) and \(\mathrm{Adm2A}\), respectively. In
particular the second class still supplies three-dimensional traces.

Suppose \(\operatorname{tr}\rho(\sigma)\ne t\) for every admissible
\(\rho\). Under the deformation, rigidity and automorphy inputs specified
below, the source gives a finite set \(X\) of pairwise nonisomorphic
admissible representations such that \(V\) has finite length and
\begin{equation}
 \frac{\operatorname{length}_{\mathcal O}V}{e_\mathcal O}
 \ \le\ 2\sum_{\rho\in X}v_\ell(\operatorname{tr}\rho(\sigma)-t).
 \label{eq:single-def-ineq}
\end{equation}
In the prime sum, \(\sigma\) is specialized to the inverse of the
chosen Frobenius element. It is arbitrary in this assertion.

\subsection{The strict cocycle in a marked deformation}

Write \(c\) for a complex conjugation. Then \(c^2=1\) and
\(\Psi(c)=1\). For now take the marked datum supplied over an
unramified coefficient ring. It gives \(g_1\in\Gamma\) with
\(\Psi(g_1)^2-1\) a unit. At the selected group \(D_\ell\) over
\(\ell\) it gives \(s\in D_\ell\) with \(\Psi(s)^2-1\) a unit.
If \(u^2-1\) is a unit, then so are \(u-1\) and
\(u^{-1}-1=-u^{-1}(u-1)\). These two markings will be used for
different purposes.

Every class \(x\in H^1(\Gamma,W)\) has a unique cocycle \(f_x\)
representing it with \(f_x(g_1)=0\). In fact if \(z\) represents \(x\),
put
\[
 m=(\alpha(g_1)-1)^{-1}z(g_1),\qquad
 f_x(g)=z(g)-(\alpha(g)-1)m .
\]
Changing \(z\) by the coboundary of \(u\) changes \(m\) by \(u\).
A coboundary vanishing at \(g_1\) has zero parameter. Hence this
normalization is \(\mathcal O\)-linear and injects cohomology into
the locally constant functions on \(\Gamma\).

Let \(b\in\mathcal O\setminus\{0\}\), \(Q_b=\mathcal O/(b)\).
For \(x\in V[b]\) the function \(bf_x=f_{bx}\) vanishes identically.
The map \(y\bmod b\mapsto y/b\bmod\mathcal O\) identifies \(Q_b\)
with \((F/\mathcal O)[b]\). In these coordinates write the cocycle
as \(u_x:\Gamma\to Q_b\). On any \(D\in\mathcal D\) there exists
\(m_D\in F/\mathcal O\) with
\(f_x(g)=(\alpha(g)-1)m_D\) for \(g\in D\). In particular
\((\alpha(s)-1)bm_{D_\ell}=0\) and the factor in parentheses is a
unit. Thus \(bm_{D_\ell}=0\) and the local coboundary at \(\ell\) has a
parameter in \(Q_b\) itself. At elements \(g\) of the imposed
inertia kernel, the marked datum has \(\Psi(g)=1\) and \(g\) belongs
to a group in \(\mathcal D\). The same formula gives \(f_x(g)=0\).

Also \(0=f_x(c^2)=2 f_x(c)\), hence \(f_x(c)=0\).
The identities \(f_x(1)=0\) and
\(f_x(gh)=f_x(g)+\alpha(g)f_x(h)\) show first
\(f_x(g_1^{-1})=0\) and then \(f_x(cg_1^{-1}c)=0\).
These vanishings impose the two marked off-diagonal conditions in
the normalized representation.

Let \(J_3\) be the antidiagonal identity and use on \({\rm GL}_3\)
the involution \(M\mapsto J_3 M^{-{\rm t}} J_3\).
The polarized reducible representation of \(\Gamma\) has matrix part
\(\operatorname{diag}(\Psi,1,\Psi^{-1})\) and component \(\eta\).
Put
\[
 B(u)=\begin{pmatrix}0&0&0\\u&0&0\\0&-u&0\end{pmatrix}.
\]
The differential of the involution fixes \(B(u)\), since
\(-J_3 B(u)^{\rm t}J_3=B(u)\). Diagonal conjugation sends it to
\(B(\alpha(g)u)\). Thus \(\gamma_x(g)=B(u_x(g))\) is an adjoint
cocycle of weight \(-1\). Its \((2,1)\) entry recovers \(u_x\).
It vanishes on the prescribed kernel, satisfies the marked
conditions, and is a coboundary at \(D_\ell\). This gives an
injective additive map from \(V[b]\) to the normalized matrix
cocycles with coefficient module \(Q_b\).

Let \(R\) be the ring representing the marked global deformation
condition, \(r\) its universal representation and \(x_0:R\to\mathcal O\)
the reducible point. The ring \(\mathcal O\oplus Q_b\) has product
\((a,u)(a',u')=(aa',au'+a'u)\). Multiplying the base representation
by \(1+\gamma_x(g)\) gives a representation over this ring because
the cross-products of the infinitesimals vanish and the remaining
identity is exactly the adjoint cocycle equation. The local coboundary
can be lifted to a local change of basis by lifting its scalar
parameter to \(\mathcal O\). The local deformation condition includes
these conjugates of the base lift. Consequently the representing
property gives a point of the form \((x_0,D_x)\). Multiplicativity
gives the Leibniz rule
\(D_x(uv)=x_0(u)D_x(v)+x_0(v)D_x(u)\).
On each universal matrix entry the value is the corresponding
entry of \(\gamma_x(g)\operatorname{diag}(\Psi(g),1,\Psi(g)^{-1})\).
In particular these values determine \(u_x\).

\subsection{The saturated cotangent lattice}

Write \(\pi:A\to R\) for the map from the graded model, \(x_A=x_0\pi\),
\(I=\ker x_A\), and
\[
 I^{(2)}_{\rm sat}
   =\{a\in A:\ \exists\,0\ne d\in\mathcal O,\ da\in I^2\}.
\]
The weight \(-1\) part of the kernel maps to a lattice \(L_-\) in
\(A/I^{(2)}_{\rm sat}\). In the presentation let
\(F_1,\ldots,F_{r+1}\) be the equations of weight \(-1\), and put
\(C_-=L_-/\sum_i\mathcal O[F_i]\). Here \(r\) is the rank of
\(L_-\), and brackets indicate images in the saturated quotient.

We use the following particular lifting property of the graded
presentation. For the \(D_x\) above, their pullbacks to \(A\) are
reductions modulo \(b\) of \(\mathcal O\)-valued derivations at
\(x_A\) supported in weight \(-1\). This is the assertion constructed
in the source's \(\mathrm{ModelLower.main\_der}\). Its local argument
lifts the weight-\(-1\) coboundary parameter at \(\ell\), hence gives
a compatible integral local derivation. The presentation theorem then
extends that derivation across the weighted ambient variables. It
only prescribes their reductions on the global quotient. The
extension is a derivation on \(A\), whereas \(D_x\) on \(R\) can
have intrinsically torsion coefficients.

For an integral lift \(\widetilde D_x\), Leibniz gives
\(\widetilde D_x(I^2)=0\). If \(df\in I^2\) with \(d\ne0\),
then \(d\widetilde D_x(f)=0\) in the domain \(\mathcal O\), so
\(\widetilde D_x(f)=0\). The reduction kills \(I^{(2)}_{\rm sat}\).
Each \(F_i\) maps to zero in \(R\), so \(D_x\pi(F_i)=0\).
The resulting map
\[
                 V[b]\longrightarrow
            \operatorname{Hom}_{\mathcal O}(C_-,Q_b)
\]
is injective. Indeed agreement on \(C_-\) gives agreement on all
of \(A_{-1}\). The pullbacks vanish on the other homogeneous pieces
by the support of their integral lifts, so they agree on \(A\).
Every universal matrix coordinate has a homogeneous lift to \(A\).
The formula for the \((2,1)\) coordinate now recovers \(u_x\), hence
\(x\). In particular saturation here uses torsion-free target
coefficients before reduction.

The presentation's spanning condition says that the classes of the
\(F_i\) span \(L_-\otimes_{\mathcal O}F\), and \(L_-\) is finite free.
Over the DVR write \(C_-\simeq\bigoplus_i\mathcal O/\varpi^{c_i}\).
If \(v_{\mathcal O}(b)=j\), its Hom group above has length
\(\sum_i\min(c_i,j)\), since the image of each cyclic generator can
be any element of \(Q_b\) killed by \(\varpi^{c_i}\).
An injection of finite sets of modules of finite length over
\(\mathcal O\) compares their lengths, since each module of length \(k\)
has \((\#\mathcal O/\varpi)^k\) elements by a composition series.
Thus \(\operatorname{length}V[b]\le\sum_i c_i\).
Every continuous cocycle into the discrete group \(W\) has finite
image because \(\Gamma\) is compact. Some power of \(\varpi\) kills
that image. Thus \(V=\bigcup_j V[\varpi^j]\). The lengths in this
increasing chain are bounded integers, and equality of two successive
finite lengths implies equality of the corresponding nested modules.
After its last strict increase the chain is constant. We conclude
\begin{equation}
                   \operatorname{length}_{\mathcal O} V
                        \le \sum_i c_i .
 \label{eq:single-cot-lower}
\end{equation}

\subsection{The two multiplicities in the intersection estimate}

We spell out the source of the coefficient \(2\) in
\eqref{eq:single-def-ineq}. Lift the three diagonal coordinates of
\(r(\sigma)\) to elements of \(A_0\) and let \(a\in A_0\) be their
sum minus \(t\), which belongs to \(\mathcal O\) in our normalization.
Then \(x_A(a)=0\). At each point \(f\) of the deformation ring,
\[
                 f(\pi(a))=\operatorname{tr}\rho_f(\sigma)-t
\]
because \(\eta(\sigma)=1\). Scaling a homogeneous coordinate of
weight \(k\) by \(z^k\), \(z\in\overline{\mathbb Q}_\ell^\times\),
fixes this trace. Call a geometric point mixed if both a positive
and a negative weight have some coordinate with nonzero value.
Take \(T\) meeting every mixed orbit on the equations. The
point and rigidity results mentioned below yield a finite such set.

One may change generators among the \(r+1\) weight \(-1\) equations.
Smith normal form gives generators with unchanged span and elements
\(y_i\in A_{-1}\) such that
\[
  F_i\equiv\varpi^{c_i}y_i\pmod{I^{(2)}_{\rm sat}}
     \quad(1\le i\le r),\qquad
  F_{r+1}\equiv0\pmod{I^{(2)}_{\rm sat}} .
\]
The other equations come in \(s\) of weight \(-2\), \(r'\) of weight
\(1\), \(s'\) of weight \(2\), and \(m\) of weight zero. Denote them
by \(G_i,E_i,E'_i,Z_i\), respectively. The chart used for the
intersection kills the \(Z_i\). One can replace \(a\) by a
weight-zero representative modulo the equations whose image in the
chart is nonzero. Its value and hence its valuation at every point
of \(T\) is unchanged.

For clarity the graded ring in the comparison is the three-variable
Rees algebra \(\mathcal R\) over \(A_0\) of the following
\(A_0\)-submodules of \(A\),
\[
    A_{-2},\quad A_2,\quad A_0F_{r+1}^2+A_{-2}^2 A_2.
\]
Its component at \(d=(d_0,d_1,d_2)\in\mathbb N^3\) has coefficients
in the product of the respective \(d_i\)-th powers. Let \(s_*\)
denote the following sequence of homogeneous sections. In degree
\((1,0,0)\) take \(F_1^2,\ldots,F_r^2\) and all \(G_i\). In
\((0,1,0)\) take all \(E_i^2,E'_i\). In \((0,0,1)\) take
\(F_{r+1}^2\). Let \(\mathcal J\) consist of polynomials whose
coefficients are killed by the chart map. It is a homogeneous
ideal. Write \(M=\mathcal R/\mathcal J\) and \(\chi(-;-)_d\)
for the Euler characteristic of the prescribed homogeneous sections
on the indicated module in multidegree \(d\), using lengths.

The lower chart estimate and upper point estimate in the source
give the following chain for all sufficiently large multidegrees
in the product order, with \(v_{\mathcal O}=e_{\mathcal O}v_\ell\):
\begin{align}
  2^{r+r'}\sum_i c_i
       &\le \chi(s_*;M/aM)_d
        \le \chi(a;M/(s_*)M)_d \notag\\
       &\le 2^{r+1+r'}\sum_{\xi\in T}v_{\mathcal O}(\xi(a)).
       \label{eq:single-rees}
\end{align}
The middle comparison uses regularity of \(a\) on \(M\). Indeed, in
each coefficient the chart sends it to a nonzero power series in a
domain, so \(a h\in\mathcal J\) implies \(h\in\mathcal J\). Applying
the nonzerodivisor rule for Euler characteristics, permuting the
sections, and applying the upper bound by the quotient by \(s_*\)
gives the stated comparison. The source proves the required
admissibility for this sequence before it uses these Euler rules.

Here is the coefficient calculation on the lower side. The leading
chart quotient modulo \(\varpi\) has as Hilbert function the number
of monomials with total negative weight \(2d_0+4d_2\) and total
positive weight \(2d_1+2d_2\). This identification and the lower
Euler bound on pulling out scalar factors from homogeneous sections
are inputs of the chart argument. In this quotient the first \(r\)
sections have leading parts \(\varpi^{2c_i}y_i^2\). The latter Euler
bound gives \(\sum_i2c_i\) times a finite difference of the Hilbert
function. The difference uses \(r-1+s\) steps in direction zero,
\(r'+s'\) in direction one, and one in direction two.

We compute this difference directly. The chart after cutting the
weight-zero equations has \(r,s,r'+1,s'+1\) variables of weights
\(-1,-2,1,2\), respectively. Thus its bivariate monomial generating
function is
\[
 H(X,Y)=\frac{1}{(1-X)^r(1-X^2)^s(1-Y)^{r'+1}(1-Y^2)^{s'+1}} .
\]
The coefficient is read at \(X^{2d_0+4d_2}Y^{2d_1+2d_2}\).
The differences therefore multiply \(H\) by
\((1-X^2)^{r-1+s}(1-Y^2)^{r'+s'}(1-X^4Y^2)\).
Writing \(P=(1+X)^{r-1}(1+Y)^{r'}\), the product equals
\[
 \frac{P}{(1-X)(1-Y)}
   +\frac{PY^2(1+X+X^2+X^3)}{(1-Y)(1-Y^2)} .
\]
This identity follows by splitting \(1-X^4Y^2=(1-Y^2)+Y^2(1-X^4)\).
The second summand has bounded \(X\)-degree. At large degrees in
both variables the first summand has coefficient
\(P(1,1)=2^{r-1+r'}\). Multiplying this by \(\sum_i2c_i\)
gives exactly the leftmost side of \eqref{eq:single-rees}.

On the upper side the generic local length of each weight piece at
the restriction of a mixed point to \(A_0\) is at most one for the
original equations. This uses the point rigidity and the graded
slice theorem. Replacing a homogeneous equation \(u\) of weight
\(k\) by its square at most doubles a uniform bound on the local
lengths of all the weight pieces. Indeed in a graded quotient \(P'\)
the map \(P'/u^2P'\to P'/uP'\) has kernel \(uP'/u^2P'\). The latter
is a quotient of \(P'/uP'\) via multiplication by \(u\), with a
shift by \(k\). This remains a right-exact sequence on each graded
piece and on localization. Thus a weight-\(j\) piece of the squared
quotient has length at most the sum for weights \(j\) and \(j-k\)
of the original quotient. Iterating over the \(r+1\) equations \(F_i\)
and the \(r'\) equations \(E_i\) gives \(2^{r+1+r'}\). The other
equations are unchanged. The upper Euler estimate uses a surjection
onto the relevant large Rees pieces, then the length-versus-valuation
formula on their support. Only the restrictions of mixed points can
contribute. Their values of \(a\) are nonzero and have nonnegative
valuation. Summing over \(T\) majorizes the sum over distinct
restrictions even when several elements restrict to the same point.

The algebraic results used for these two estimates occur in
\(\mathrm{AdmissibleMonoIdeal.low\_holds}\),
\(\mathrm{TwoSigned.up\_holds}\), and
\(\mathrm{AlgMemCoef.core\_holds}\). The calculations above prove
the two powers of two and the regularity check. They use, as separate
inputs, the chart identification, the Euler rules and support
comparison, and the generic local-length-one assertion. Since
\(r\ge1\), cancellation of the positive factor \(2^{r+r'}\) in
\eqref{eq:single-rees} yields
\[
               \sum_i c_i\le
                 2\sum_{\xi\in T}v_{\mathcal O}(\xi(a)).
\]

\subsection{Identification of the points and coefficient descent}

The cohomological hypothesis used for the construction consists of
the local deformation inputs, normalized Poitou--Tate duality, global
cohomological finiteness and the global Euler characteristic theorem,
and the size bound for \(H^2\) called Lemma~\(K\).
Together with the presentation theorem called Lemma~\(P\), these
give the marked ring and its graded model, including the integral
lifting property used above. Over unramified coefficients the
automorphy input supplies the marking, finiteness of the subalgebra
of weight-zero matrix monomials over \(\mathcal O\), automorphy for
points of the first type, and modularity for the two-dimensional
extensions occurring at points of the second type.

More precisely, the point condition is deduced in
\(\mathrm{FPtHypU}\). It identifies a point that is both upper and
lower with the base point, makes the prescribed tangent vectors at
the base coboundaries by nonpositive-weight Bloch--Kato vanishing,
and makes tangent vectors at mixed points coboundaries by adjoint
Selmer vanishing. For points of the first type the last vanishing
uses \(\mathrm{RigidityFirstCase}\) at the automorphic representation.
For points of the second type, the proof uses the rank-two Kato input
\(\mathrm{KatoTwoW}\), \(\mathrm{RigiditySecondCase}\) for the
two-dimensional extension of determinant \(\eta\), and its modularity.
These give the point, tangent and spanning hypotheses of the graded
intersection estimate. The first type is admissible by
\(\mathrm{AdmI}_s\). The second is admissible by the theorem on
modular extensions of type II. These are arithmetic inputs, additional
to the linear and intersection calculations above.

The local-point comparison in the graded presentation puts every
mixed point, up to the scaling action, over a point of \(R\). Conversely
it identifies mixed points of \(R\) with those graded points having
both signs. The matrix coordinates of weight \(j\) scale by \(z^j\). On representations this is conjugation by
\(\operatorname{diag}(z,1,z^{-1})\). Choose one point of \(R\) per
mixed orbit. The marked shape theorem says that for two mixed points
an isomorphism between their restrictions to \(G_K\) is induced by
such a diagonal conjugacy. Thus restriction to \(G_K\) gives a set
\(X\) with one representative per isomorphism class in its image,
and summing over these representatives preserves the valuation
sum. Each term is the trace difference above. Combining the upper
estimate with \eqref{eq:single-cot-lower} and dividing by
\(e_\mathcal O\) proves \eqref{eq:single-def-ineq} for the
unramified coefficient choice.

Finally the order \(m_\psi\) of \(\psi\) is prime to \(\ell\).
Choose \(f>0\) with \(m_\psi\mid\ell^f-1\). The required roots
of unity lift by Hensel's lemma from \(\mathbb F_{\ell^f}\) to
the unramified degree-\(f\) extension, so the same character values
occur in an unramified coefficient ring. To compare with the original
ring take a finite compositum. For an inclusion of coefficient rings
\(\mathcal O_1\subset\mathcal O_2\) the latter is finite free, and
\((F_1/\mathcal O_1)\otimes_{\mathcal O_1}\mathcal O_2
 \simeq F_2/\mathcal O_2\).
After choosing an \(\mathcal O_1\)-basis of \(\mathcal O_2\),
continuous cochains and restrictions for this scalar extension are
finite direct sums of their originals, so the kernels defining
the strict group commute with extension. For a finite length
module a composition series shows that its length is multiplied by
the relative ramification index, since the extension of a residue
field factor is \(\mathcal O_2/\varpi_1\). Finiteness also descends
by finite freeness. The denominator \(v_{\mathcal O_i}(\ell)\)
changes by the same index. Thus the normalized length, its finiteness,
and the target trace are the same in both coefficient realizations.
This is the coefficient comparison of \(\mathrm{CoefIndep}\) in the
proof of \(\mathrm{FOneCharU}\).

\section{The chart Hilbert function and the length of a mixed weight space}
\label{sec:deformation-chart}

We prove three assertions used in Section~\ref{sec:single-character}, namely
the choice of a nonzero trace representative on the chart, the
Hilbert function of the leading chart, and the local multiplicity
bound for the original equations. These assertions concern graded
algebras once the deformation presentation and its tangent
conditions have been constructed.

\subsection{Chart coordinates and leading monomials}

Keep the DVR \(\mathcal O\), its fraction field \(F\), and the
augmentation \(x_A:A\to\mathcal O\). The algebra \(A\) is
\(\mathbb Z\)-graded, finitely generated over \(A_0\) by homogeneous
elements, and \(x_A\) is zero on every nonzero weight. The cut chart
is an algebra homomorphism
\[
           \phi:A\longrightarrow F[[X_i:i\in J]]
\]
with \(J\) finite and weights \(w_i\in\{-2,-1,1,2\}\). Its constant
coefficient is \(x_A\), it carries homogeneous elements to series of
the same weight, and its linear terms separate the axes. The last
condition means that for every \(i\) some element maps to a series
whose coefficient of \(X_i\) is nonzero and whose other linear
coefficients vanish. Taking its homogeneous component of weight
\(w_i\) preserves these properties. There are \(r\ge1\) negative
unit-weight axes and \(r'+1\ge1\) positive unit-weight axes.

If \(\phi(a)=0\) for the weight-zero trace lift \(a\), it can be
changed by a weight-zero relation. Indeed, the coefficient of a chosen
weight-\(-1\) variable gives a linear functional on \(A_{-1}\)
which is nonzero by axis separation. It kills
\((\ker x_A)^{(2)}_{\rm sat}\) because products of two series with
zero constant term have zero linear coefficients, and coefficients
are taken in a field of fractions. The weight-\(-1\) equations
span the saturated cotangent lattice over \(F\). Consequently
\(\phi(F_i)\ne0\) for some \(i\). Similarly choose \(g\in A_1\)
with nonzero linear coefficient. Then \(a'=a+gF_i\) has the same
values as \(a\) on solutions of the equations and at the origin,
and \(\phi(a')=\phi(g)\phi(F_i)\ne0\). If \(\phi(a)\ne0\) already,
take \(a'=a\). This proves the assertion used in choosing
\(\mathcal J\) in the Rees argument.

For a monomial \(X^\nu\) put
\[
 n(\nu)=\sum_{w_i<0}|w_i|\nu_i,\qquad
 p(\nu)=\sum_{w_i>0}w_i\nu_i .
\]
Its weight is \(p(\nu)-n(\nu)\). The filtration given by the
inequality \(n(\nu)\ge k\) is multiplicative, since \(n\) is
additive on monomials. Put \(F_*=F_{r+1}\) in Smith coordinates, and
\[
 L_0=A_{-2},\qquad L_1=A_2,\qquad
 L_2=A_0F_*^2+A_{-2}^2A_2,\qquad
 L(d)=L_0^{d_0}L_1^{d_1}L_2^{d_2}.
\]
The first submodule has negative filtration at least two, the second
at least zero. Since a nonzero scalar multiple of \(F_*\) lies in
\((\ker x_A)^2\), its image in the chart has no constant or linear
term. A monomial of weight \(-1\) with negative weight at most one
would be a single weight-\(-1\) variable. Thus \(\phi(F_*)\) has
negative filtration at least two, and every element of \(L_2\)
has filtration at least four. It follows that every
\(h\in L(d)\) has weight \(2(d_1-d_0-d_2)\), and
\(\phi(h)\) has \(n\ge u\), with
\[
                 u=2d_0+4d_2,\qquad v=2d_1+2d_2 .
\]
On the layer \(n=u\) the weight identity forces \(p=v\).

Let \(S_d\) be the finite set of monomials with these two weights
\((u,v)\). The coefficient map from \(L(d)\) to \(F^{S_d}\)
records that first layer. Its image \(H_d\) is an \(\mathcal O\)-lattice
of rank \(\#S_d\), as we now prove.

First, multiplying by \(b\in A_0\) acts on this layer by \(x_A(b)\).
Every nonconstant monomial in \(\phi(b)\) has \(n=p>0\), since
there are no variables of weight zero. Thus \(H_d\) is finitely
generated over \(\mathcal O\) once \(L(d)\) is finitely generated
over \(A_0\). Every weight component is finitely generated over
\(A_0\). To see this, use fixed homogeneous algebra generators and,
for a fixed weight, choose the finitely many minimal exponent
vectors of that weight in coordinatewise order. Their finiteness
is Dickson's lemma. Every other monomial of that weight is a multiple
of a minimal one by a weight-zero monomial. This proves the assertion
for the components \(A_{-2},A_2\), and products and finite sums of
finitely generated submodules are finitely generated. The claimed
finiteness of \(H_d\) follows.

Next we isolate every individual coordinate of \(F^{S_d}\). For a
weight-\(\pm1\) variable, take the homogeneous element separating
its linear axis. Its lowest possible pair \((n,p)\) is \((1,0)\)
or \((0,1)\). All the monomials on that layer are linear and only
the desired coefficient is nonzero. For a weight-\(\pm2\) variable,
axis separation disposes of other linear terms. On its lowest
layer the remaining monomials are products of two unit-weight
variables of the same sign. We have already constructed elements
whose products isolate each such quadratic monomial on this layer.
Subtract suitable \(F\)-multiples of these products to cancel their
coefficients. Clearing all denominators yields a homogeneous element
of \(A_{\pm2}\) with a nonzero coefficient on the desired axis and
no other term on the lowest layer. The argument also supplies such
an element for each product of two same-sign unit-weight variables
by direct multiplication. When two such isolating elements are
multiplied, terms at a larger negative weight cannot contribute to
the least negative layer of the product. On that layer, the product
therefore has just the chosen product monomial with nonzero coefficient.

For \(\nu\in S_d\) partition its negative variables into
\(u/2\) groups of weight \(-2\) and its positive variables into
\(v/2\) groups of weight \(2\). This is possible because the
number of weight-one factors on each side is even after the
weight-two factors are removed. The preceding construction gives
an element isolating \(X^\nu\) in
\[
 A_{-2}^{\,d_0+2d_2}A_2^{\,d_1+d_2}
    = L_0^{d_0}L_1^{d_1}(A_{-2}^2A_2)^{d_2}
    \ \subseteq L(d).
\]
Hence \(F H_d=F^{S_d}\). This finite torsion-free module over a DVR
is free, with the asserted rank.

Define the homogeneous ideal \(\mathfrak q\) of the Rees algebra
\(\mathcal R=\bigoplus_d L(d)T^d\) by requiring each coefficient in
degree \(d\) to have zero layer at \(n=u\). The filtration bound
and its multiplicativity show it is an ideal. At each \(d\), its
component is exactly the kernel of the coefficient map. It follows
that the \(d\)-component of \((\mathcal R/\mathfrak q)/\varpi\)
is \(H_d/\varpi H_d\). Its \(A_0\)-action factors through the
surjection \(x_A:A_0\to\mathcal O\), so its length is
the \(\mathcal O\)-length. This equals \(\#S_d\), as asserted for
every \(d\in\mathbb N^3\), including the boundary of the cone.
This proves the leading-chart Hilbert function translated from
\(\mathrm{ReesGradedPieceSurjection}\).

\subsection{Removing the torus direction from a mixed point}

We now prove the local-length bound used before squaring the equations.
Let \(B=A/I_0\) where \(I_0\) is their homogeneous ideal. Reserve \(I_0\) for
this ideal and \(\ker x_A\) for the augmentation ideal. Write \(D=B_0\).
The two model hypotheses we use are that \(D\) is a finite
\(\mathcal O\)-module and that at every mixed geometric point in
characteristic zero the tangent space to the equations has dimension
at most one. In the source the latter hypothesis states that any
two point derivations annihilating \(I_0\) are linearly dependent.
Since \(A\) is finitely generated over \(A_0\), \(B\) is finitely
generated over \(D\). Let \(k\) be an algebraic closure of \(F\).

Fix a mixed geometric point. There are \(f\in B_{-2}\),
\(g\in B_2\) with \(fg\) nonzero there. For the model in this
application the generators over degree zero have weights
\(\pm1,\pm2,0\). If some element of negative weight has nonzero
value then a monomial among its terms does also, and at least one
negative-weight generator occurring in that monomial has nonzero
value. If it has weight \(-1\) take its square. This supplies \(f\). Positive weights give \(g\) in the same way. Let \(\mathfrak p\) be the kernel
of the geometric point restricted to \(D\). Define the graded ring
\[
     H=(B\otimes_{\mathcal O}F)[(fg)^{-1}] .
\]
Both \(f\) and \(g\) are units in \(H\) and have their original
degrees. We claim that \(H\) is reduced.

Set \(C=H/(f-1)\). It is of finite type over \(F\). At any geometric
point \(\xi\) of this slice, \(\xi(f)=1\) and \(\xi(g)\ne0\), so
\(\xi\) is mixed. The Euler tangent vector on the original graded
algebra sends a homogeneous \(h\) of weight \(w\) to
\(w\,\xi(h)\). The Leibniz rule follows by adding weights. The
homogeneity of the equations makes it a tangent vector on \(B\).
Its value on \(f\) is \(-2\). A tangent vector on the slice pulls
back to a tangent vector on \(B\) vanishing on \(f\). These
two vectors are linearly dependent by the mixed-point hypothesis.
Evaluation on \(f\) shows the slice vector vanishes on \(B\).
It also vanishes on every inverted element by the quotient rule,
hence on \(C\). This holds at every geometric point of the slice.

For completeness base change to \(k\). At a maximal ideal
\(\mathfrak m\) of \(C\otimes_F k\) the residue field is \(k\) by
the Nullstellensatz. The dual of \(\mathfrak m/\mathfrak m^2\) is
the geometric tangent space just shown to vanish. Nakayama's lemma
in the Noetherian local ring now gives \(\mathfrak m_{\mathfrak m}=0\).
Each such local ring is a field. A nilpotent which survives in the
algebra survives at some maximal ideal containing its annihilator,
which would contradict the field property there. Thus \(C\otimes_F k\)
and, by faithful scalar extension, \(C\) are reduced.

The assertion lifts from the slice to the open graded ring without
any descent assumption on weight spaces. Introduce an invertible
variable \(U\) with \(U^2=f^{-1}\). There is an isomorphism
\[
 H[U,U^{-1}]/(U^2-f^{-1})\ \simeq\ C[V,V^{-1}]
\]
sending homogeneous \(h\in H_w\) to \(\bar h V^w\) and \(U\) to \(V\).
Its inverse sends \(\bar h\) for a homogeneous representative to
\(hU^{-w}\), and sends \(V\) to \(U\). Indeed \(fU^2=1\) makes
the map factor through the slice, and checking on homogeneous
elements and the invertible variable proves the two compositions
are the identities. The left-hand ring is free of rank two over
\(H\) by the monic equation and invertibility of its constant
coefficient. The right-hand ring is reduced since \(C\) is reduced.
The inclusion of \(H\) is injective, giving the claim.

Now localize \(B\) at \(D\setminus\mathfrak p\), obtaining a graded
ring \(B'\). This factors through \(H\), since all nonzero scalars
of \(\mathcal O\) and \(fg\) lie outside \(\mathfrak p\). Thus \(B'\) is a further localization of \(H\).
Thus \(B'\) is reduced. Its degree-zero ring \(D_{\mathfrak p}\)
is local, reduced, and a localization of the finite-dimensional
\(F\)-algebra \(D\otimes_{\mathcal O}F\). It is therefore a field.

For an arbitrary integer \(w\), the element \(h_*=f^w\) is a
unit of degree \(-2w\) in \(B'\). If \(B'_w\ne0\), choose
\(0\ne y\in B'_w\). Reducedness makes \(y^2 h_*\) nonzero in
the field \(D_{\mathfrak p}\). For each \(z\in B'_w\) the identity
\[
              z=\frac{zyh_*}{y^2 h_*}\,y
\]
then expresses \(z\) as a scalar multiple of \(y\). Hence every
weight piece has dimension at most one over \(D_{\mathfrak p}\).
Pullback along the surjection from the localization of \(A_0\)
does not change its module length. This proves the generic
local-length-one bound at every restriction of a mixed geometric
point. It supplies the input to the squaring argument in
\eqref{eq:single-rees}. The corresponding source proofs are
\(\mathrm{SliceHolds}\) and
\(\mathrm{LocLenQuotientWeightPieceLeOne}\).

The arithmetic construction of the graded model, its chart and the
mixed-point tangent bound are inputs to this section. The lower
Euler estimate for extracting scalar factors, the support comparison
for the large Rees pieces and the global automorphy statements in
Section~\ref{sec:single-character} are further steps of the original
argument.

\section{Support of the Rees pieces and the upper valuation estimate}
\label{sec:deformation-support}

We prove the last inequality in \eqref{eq:single-rees} from the
graded presentation and its point conditions. Write \(D=A_0\) and
\(I_0\) for the homogeneous ideal of the original equations. Use
\(L_0,L_1,L_2,\mathcal R,s_*\) from the preceding two sections and
write \(Q=\mathcal R/J\), where \(J\) is any homogeneous ideal
containing all the sections \(s_*\) and all the degree-zero \(Z_i\).
The module to which the upper estimate applies is \(P_d=Q_d\).
In the application \(J\) also contains the coefficient kernel of
the chart map. We need only the inclusions just stated.

We recall the assumptions about points of the presentation used here.
The ring \(D\) is Noetherian local, with the same residue field as
\(\mathcal O\), and \(D/(D\cap I_0)\) is finite as an
\(\mathcal O\)-module. The weights of a finite set of generators
for \(A\) over \(D\) lie between \(-2\) and \(2\). The weight
\(-1,-2\) equations span their respective saturated cotangent
lattices after extending scalars to \(F\). At a fixed geometric
point on the equations other than the augmentation, the model
excludes simultaneous existence of two solutions with the same
values on \(D\), one supported on nonpositive weights but nonzero
on some negative weight, and one supported on nonnegative weights
but nonzero on some positive weight. This is the specific geometric
assertion called \(M3\). The mixed-point tangent condition used in
Section~\ref{sec:deformation-chart} is also assumed.

\subsection{Equations seen by a relevant prime}

Let \(U_i\) be the ideal of \(\mathcal R\) generated by \(L_i T_i\),
and put \(\mathfrak b=U_0U_1U_2\). Call a prime \(\mathfrak P\)
of \(\mathcal R\) relevant if it contains \(J\) and does not contain
\(\mathfrak b\). For such a prime choose
\[
 f\in A_{-2},\quad g\in A_2,\quad h\in L_2,\qquad
 fT_0,\ gT_1,\ hT_2\notin\mathfrak P .
\]
The weights of \(f,g,h\) are \(-2,2,-2\). Write
\(\mathfrak p=\mathfrak P\cap D\). We show
\begin{align}
 D\cap I_0&\subseteq\mathfrak p, \label{eq:rees-contract}\\
 \Delta\subseteq\mathfrak p&\ \Longrightarrow\
         \mathfrak p=\mathfrak m_D,\qquad
 \Delta=(A_{-2}A_2)D\ \subseteq D . \label{eq:rees-boundary}
\end{align}
In the formula for \(\Delta\) the ideal is generated by the indicated
products.

First suppose \(\tau=uv\notin\mathfrak p\) for \(u\in A_{-2}\),
\(v\in A_2\). Localize at \(\tau\), so \(u,v\) are units. The
weight \(2j\) component of the localized algebra is \(u^{-j}D_\tau\)
by multiplication by the unit \(u^j\). In particular the three
modules \(L_0,L_1,L_2\) become \(uD_\tau,u^{-1}D_\tau,uD_\tau\);
the last equality uses \(A_{-2}^2 A_2\subseteq L_2\subseteq A_{-2}\).
Thus the localized Rees algebra is the polynomial algebra over
\(D_\tau\) on \(uT_0,u^{-1}T_1,uT_2\). Relevance says none of
these variables belongs to the localized prime. The sections force
the coefficient of every corresponding weight-normalized generator
\(F_i^2,G_i,E_i^2,E'_i,Z_i\) into its contraction. For a homogeneous
\(a\) of weight opposite to an unsquared equation, \(a\) times that
equation has degree zero. Its square is a scalar multiple of the
corresponding normalized squared equation in the odd-weight cases,
so the product belongs to the same prime. In even weights one uses
the normalized unsquared equation directly. Every element of
\(D\cap I_0\) is a sum of such degree-zero products by homogeneity.
This proves \eqref{eq:rees-contract} on this open set.

Suppose next that \(\Delta\subseteq\mathfrak p\). We construct
one-sided points directly. For either sign \(\epsilon=\pm1\), put
\(t=g,i=1\) if \(\epsilon=1\), and \(t=f,i=0\) if \(\epsilon=-1\).
The even components satisfy
\(A_{2\epsilon n}=(A_{2\epsilon})^n\) for \(n\ge0\), with products
over \(D\). Indeed in any generator monomial of weight \(2\epsilon n\)
the factors of sign \(\epsilon\) have total absolute weight at least
\(2n\). One can choose \(n\) disjoint blocks of weight \(2\epsilon\),
each a single weight-two factor or two weight-one factors. All
remaining factors have net weight zero.

In the fraction field of \(\mathcal R/\mathfrak P\), define on
homogeneous \(a\in A_{2\epsilon m}\), \(m\in\mathbb Z\),
\[
 \theta_\epsilon(a)=
       \frac{[(at^k)T_i^{\,m+k}]}{[tT_i]^k},
             \qquad k\ge\max(0,-m).
\]
The even-component identity makes the numerator a Rees element.
Increasing \(k\) multiplies numerator and denominator by the same
nonzero element. Choosing common \(k\)'s proves additivity, and
using the sum of the two chosen \(k\)'s proves multiplicativity.
Thus \(\theta_\epsilon\) is a ring map on the even Veronese subring.
It induces the original map on \(D\), and its value on \(t\) is
nonzero. For \(a\in A_{-2\epsilon}\), the numerator with \(k=1\)
is \(at\), which belongs to \(\Delta\), so its value is zero.
The even-component identity kills every weight of opposite even sign.

All squared odd-weight equations on the chosen side, and the
unsquared equations of weight \(2\epsilon\), map to zero since their
sections belong to \(\mathfrak P\). For the negative side this also
holds for \(F_{r+1}^2T_0\), which was imposed in slot two. Indeed,
\[
        (F_{r+1}^2T_0)(hT_2)=(hT_0)(F_{r+1}^2T_2)
\]
and \(hT_2\notin\mathfrak P\). The degree-zero equations map to
zero as well. The algebra \(A\) is integral over its even subring:
an odd element squares into that subring, and every element is a
sum of an even and an odd element. By lying over, extend the map
to \(A\) with values in a field extension. Odd equations on the
chosen side vanish because their squares do. Every element of odd
opposite sign vanishes for the same reason. We have obtained a
solution of the equations on which only the indicated sign and
weight zero can be nonzero, with \(t\) surviving. Its restriction
to \(D\) has kernel \(\mathfrak p\). In particular this proves
\eqref{eq:rees-contract} also on the boundary.

It remains to exclude a boundary prime in characteristic zero.
By \eqref{eq:rees-contract}, \(D/\mathfrak p\) is finite over
\(\mathcal O\). If \(\mathfrak p\) meets \(\mathcal O\) nontrivially,
it contains \(\varpi\), so this quotient is a finite-dimensional
domain over the residue field and is a field. Thus
\(\mathfrak p=\mathfrak m_D\). Otherwise choose an embedding
of \(\operatorname{Frac}(D/\mathfrak p)\) in \(k=\overline F\).
Both one-sided solutions can be realized in \(k\) over this fixed
embedding. For each sign the equations, the requirement that
opposite generators vanish, and \(t\ne0\) define a finite type
algebra over \(\operatorname{Frac}(D/\mathfrak p)\). The constructed
field-valued solution shows it is nonzero. Extending scalars to
\(k\) preserves nonzeroness, and a maximal ideal gives the desired
\(k\)-point.

Sending every nonzero-weight element of the solution of negative
sign to zero and preserving its values on \(D\) produces a fixed
point. One way to see multiplicativity is to send a homogeneous
element of weight \(j\le0\) to its value times \(T^{-j}\) and send
all \(j>0\) to zero. The original solution vanishes on positive
weights, so this gives a map to \(k[T]\) whose specialization at
zero is the fixed point. The fixed point still satisfies the
homogeneous equations. If it differs from the augmentation, the
two one-sided solutions contradict \(M3\).

If it is the augmentation, the one-sided solution \(\xi_-\) agreeing
with it on \(D\) must vanish also on negative weights. Here are the
details. A homogeneous element of \((\ker x_A)^2\) of weight \(-1\)
is a sum of products of two homogeneous augmentation elements whose
weights sum to \(-1\). At least one factor has nonnegative weight,
and its value at \(\xi_-\) is zero. The same holds for the saturation
since nonzero scalars remain nonzero. Spanning by the \(F_i\) over
\(F\) now makes \(\xi_-\) vanish on all of \(A_{-1}\). In weight
\(-2\), products of two augmentation elements either have a
nonnegative-weight factor or have two factors of weight \(-1\).
So the saturated square vanishes there as well, and spanning by the
\(G_i\) kills \(A_{-2}\). This contradicts \(\xi_-(f)\ne0\).
This proves \eqref{eq:rees-boundary}.

\subsection{Passing to large multidegrees}

Put \(I_D=D\cap I_0\). For every prime of \(\mathcal R\) containing
\(J\), either \(\mathfrak b\) or \(I_D\) lies in it. Similarly every
prime containing \(J+\Delta\mathcal R\) contains \(\mathfrak b\)
or \(\mathfrak m_D\). The Noetherian property now gives a single
integer \(n\) satisfying
\[
  (I_D\mathfrak b)^n\subseteq J,\qquad
  (\mathfrak m_D\mathfrak b)^n\subseteq J+\Delta\mathcal R.
\]
Indeed \(\mathcal R\) has finitely many generators over \(D\) by
the module finiteness proved in Section~\ref{sec:deformation-chart},
and an ideal with finitely many generators in a radical has a
power in the ideal. If all three \(d_i\ge n\), each element of
\(\mathcal R_d\) belongs to \(\mathfrak b^n\), by expanding it
into products of generating sections. Comparing homogeneous
components therefore gives
\[
                     I_D^nP_d=0,\qquad
                  \mathfrak m_D^nP_d\subseteq\Delta P_d .
\]
The ring \(D/I_D^n\) is finite over \(\mathcal O\) because its successive
pieces \(I_D^j/I_D^{j+1}\) are finite over \(D/I_D\). Thus \(P_d\)
is finite over \(\mathcal O\) and so is \(D/\operatorname{ann}P_d\).

For any \(\sigma:D\to k\) in characteristic zero whose localization
\((P_d)_{\ker\sigma}\) is nonzero, the first equality implies that
\(\sigma(I_D)=0\). If \(\sigma(\Delta)=0\), then \(\varpi\in\mathfrak m_D\)
is invertible at \(\ker\sigma\), and the second inclusion gives
\((P_d)_{\ker\sigma}=\Delta(P_d)_{\ker\sigma}\). Nakayama's lemma
contradicts its nonzeroness. Consequently every such \(\sigma\)
has \(\sigma(uv)\ne0\) for some \(u\in A_{-2},v\in A_2\).

This \(\sigma\) extends to a mixed geometric solution of the original
equations. In fact \((A/I_0)_0=D/I_D\) is a direct summand of
\(A/I_0\) as a \(D/I_D\)-module by degree projection.
Consequently \((A/I_0)\otimes_{D,\sigma} k\) contains \(k\) as a
direct summand. This is a nonzero algebra of finite type over \(k\),
hence has a \(k\)-point. The product \(uv\) is nonzero there,
forcing both signs to occur. Its restriction to \(D\) is one of
the restrictions of the finite set \(T\) meeting the mixed orbits,
since torus scaling fixes \(D\).

For such a \(\sigma\) the open-set polynomial description above
also bounds local multiplicities. Invert \(\tau=uv\) and write
the weighted sections as scalar coefficients times the three
polynomial variables. Put \(I_2=(F_i^2,G_i,E_i^2,E'_i,Z_i)\) in
\(A\). In any multidegree with all \(d_i\ge1\), the image of the
relations in that degree contains the ideal generated by those
scalar coefficients and the \(Z_i\) times the degree-\(d\) monomial.
These scalar coefficients generate \((I_2)_\tau\cap D_\tau\)
because every generator of \(I_2\) has even weight and, after
inversion, each even-weight piece is a unit multiple of \(D_\tau\).
It follows that \((P_d)_\tau\) is a quotient of the indicated
even-weight piece of \((A/I_2)_\tau\), after identifying that
piece with \(D_\tau/((I_2)_\tau\cap D_\tau)\) by the weight unit.
The same holds after localization at \(\ker\sigma\). Each weight
piece for the original equations has localized length at most one
by Section~\ref{sec:deformation-chart}. The repeated squaring
argument of Section~\ref{sec:single-character} therefore proves
\[
             \operatorname{length}_{D_{\ker\sigma}}
                 (P_d)_{\ker\sigma}\le 2^{r+1+r'} .
\]

\subsection{Euler characteristic as a sum of valuations}

Let \(P=P_d\) and denote \(m_\sigma=
\operatorname{length}_{D_{\ker\sigma}}P_{\ker\sigma}\).
For \(a\in D\) with \(\sigma(a)\ne0\) on the generic support,
both the kernel and cokernel of \(a:P\to P\) have finite length.
Moreover
\[
  \operatorname{length}_D(P/aP)-\operatorname{length}_D P[a]
       =\sum_{\sigma:D\to k,\;m_\sigma>0}
                    m_\sigma v_{\mathcal O}(\sigma(a)).
\]
To prove the identity remove the \(\mathcal O\)-torsion \(P_{\rm tor}\).
It has finite length, so its kernel and cokernel under an endomorphism
have the same length. The snake sequence shows that their difference
is additive in short exact sequences. On \(P/P_{\rm tor}\), which
is finite free over \(\mathcal O\), multiplication by \(a\) has
nonzero determinant. On the generic fiber all simple factors are
residue fields of the finite-dimensional algebra
\((D/\operatorname{ann}P)\otimes_{\mathcal O}F\), and \(a\) is
nonzero in each occurring field. Smith normal form gives
\(v_{\mathcal O}(\det a)\) for the length of its cokernel. The
\(D\)- and \(\mathcal O\)-lengths of finite-length factors agree
because their local residue fields agree.

In a composition series over the generic algebra each residue
field \(E/F\) occurs as many times as the length of the component
localized at its corresponding maximal ideal. These lengths are
precisely the \(m_\sigma\), since \(\varpi\) is already invertible
at \(\ker\sigma\). The determinant on such a simple factor is
\(\operatorname{Norm}_{E/F}(a)\). The extension is separable in
characteristic zero. Its norm factors as the product over embeddings
\(E\hookrightarrow k\), which are exactly the \(\sigma\)'s with this
kernel. Taking valuations proves the formula, including its
multiplicities.

Each contributing \(\sigma\) restricts from \(T\) and factors
through \(D/I_D\), a finite \(\mathcal O\)-algebra. Thus
\(\sigma(a)\) is integral over \(\mathcal O\) and its valuation
is nonnegative. The one-element Euler characteristic in degree
\(d\) is exactly the kernel-cokernel difference just computed.
We have proved, for all sufficiently large \(d\),
\[
        \chi(a;Q)_d
            \le 2^{r+1+r'} \sum_{\xi\in T}v_{\mathcal O}(\xi(a)).
\]
The argument covers each restriction once before enlarging the
sum to \(T\). The source loci are \(\mathrm{ExistsSecNotMem}\),
\(\mathrm{SameDegreeZeroOfRestriction}\),
\(\mathrm{EventuallyMQGrFiniteAnnihilator}\),
\(\mathrm{LocalizedLengthLeOfSurjective}\) and
\(\mathrm{HerbrandFinLen}\). The proof here uses the stated finite
presentation, point condition, and cotangent spanning hypotheses
from the arithmetic model. The lower bound and middle Euler
comparison in \eqref{eq:single-rees} are separate assertions.

\section{Euler inequalities for the deformation presentation}
\label{sec:deformation-euler}

We finish the two remaining inequalities in \eqref{eq:single-rees}.
The proof uses a positivity argument for graded Koszul complexes and a
short exact sequence which extracts the Smith exponents separately.

\subsection{Positivity for homogeneous sections}

Let \(D\) be a Noetherian local ring with maximal ideal \(\mathfrak m\).
Consider a standard \(\mathbb N^k\)-graded algebra \(R\) over \(D\),
generated by finitely many homogeneous elements in each unit-vector
degree \(e_i\). Let \(\mathfrak b=\prod_i(R_{e_i})R\), where the
product is of the corresponding generated ideals. For a finite graded
\(R\)-module \(P\), consider a list \(S\) of homogeneous sections
of degrees in \(\mathbb N^k\). Say the list satisfies the finiteness
condition on \(P\) if some power of \(\mathfrak m\mathfrak b\) is
contained in \(\operatorname{ann}P+(S)\).
For large multidegree \(d\), every finite graded module annihilated
by this latter ideal has finite \(D\)-length in degree \(d\).
Indeed monomials generating a given component, times one of finitely
many module generators, have at least \(n\) section factors of every
unit degree whenever all coordinates of \(d\) are sufficiently large.
Thus \(\mathfrak m^n\) kills that component if
\((\mathfrak m\mathfrak b)^n\) kills the module. The component is
finitely generated over \(D\), and \(D/\mathfrak m^n\) is Artinian.

Write \(\chi(S;P)_d\) for the alternating sum of the lengths, in degree
\(d\), of the homology of the graded Koszul complex on \(S\). Its
homology is finite as a module, and every section in \(S\) annihilates
it by the Koszul homotopy. The preceding observation therefore makes
the definition valid in all sufficiently large multidegrees. A short
exact sequence of coefficient modules yields additivity of \(\chi\)
by the homology sequence. A permutation of the list gives the signed
permutation isomorphism of Koszul complexes, preserving \(\chi\).
If \(s\) has degree \(p\), and \(S'\) is the rest of the list, then
\[
 \chi(s,S';P)_d=\chi(S';P/sP)_d
                     -\chi(S';P[s](-p))_d .
\]
Here the degree-\(d\) part of \(P(-p)\) is \(P_{d-p}\).
To see the equality, take first the homology of
\(P(-p)\xrightarrow{s}P\) in the Koszul double complex; the two rows
left are exactly the indicated kernel and cokernel with the
remaining sections. The Euler sums can be taken on the page after
their remaining Koszul homologies have been computed, where the
terms have finite length. This also connects the Koszul definition
with the recursive Euler characteristic in the source.

We prove \(\chi(S;P)_d\ge0\) for all sufficiently large \(d\).
Induct on the number of sections. The case of the empty list follows
from the definition of length. By additivity it suffices to handle
a shift of \(R/\mathfrak p\) for a homogeneous prime
\(\mathfrak p\supseteq\operatorname{ann}P\). Indeed there is a
finite homogeneous prime filtration of \(P\). Choose a nonzero
homogeneous element with maximal annihilator among such elements.
Its homogeneous annihilator is prime since, for homogeneous
\(a,b\), \(abx=0\) and \(bx\ne0\) imply that \(x\) and \(bx\) have
the same annihilator. Homogeneous primality is primality here by
using any translation-invariant total order on \(\mathbb Z^k\) and
the extremal nonzero homogeneous components of products in the
quotient. Repeat on the quotient by the generated cyclic module. The ascending chain condition terminates the filtration. The
finiteness condition on the original module holds for each factor
and is unaffected in large degrees by its grading shift.

If \(\mathfrak p\) contains \(\mathfrak b\), the large components of
the factor and of the finitely many shifts in its complex vanish.
Otherwise there is a section outside \(\mathfrak p\) in each unit
degree. Write \(P_0=R/\mathfrak p\) and take the first section \(s\).
If \(s\notin\mathfrak p\) it is regular on \(P_0\), so the recursion
gives \(\chi(S;P_0)_d=\chi(S';P_0/sP_0)_d\), nonnegative by
induction. If \(s\in\mathfrak p\), take \(z\notin\mathfrak p\) of
the same degree \(p\), by multiplying sections of the unit degrees
(the empty product is \(1\)). Now \(S'\) already satisfies the
finiteness condition on \(P_0\). The recursion for the zero action
of \(s\) gives
\(\chi(S';P_0)_d-\chi(S';P_0(-p))_d\). Additivity for the injection
\(P_0(-p)\xrightarrow{z}P_0\) identifies this difference with
\(\chi(S';P_0/zP_0)_d\), again nonnegative by induction.
The finite number of factors permits one common lower threshold for
the multidegree. This proves positivity. In particular passage to a
quotient can only decrease \(\chi(S;-)_d\) under the finiteness
condition on the original module, by additivity for its kernel.

The same recursion bounds the Euler characteristic by successively
discarding annihilator terms:
\[
      \chi(S,U;P)_d\le\chi(U;P/(S)P)_d
\]
when the combined list satisfies the finiteness condition on \(P\).
At each step \(P[s]\) and \(P/sP\) are killed by \(s\), so the
remaining list satisfies the condition on them. The subtracted
term is eventually nonnegative by the result just proved.

\subsection{Finiteness on the chart and the middle comparison}

Apply this with \(D=A_0\), \(R=\mathcal R\), \(S=s_*\), and the
trace lift \(a\) chosen in Section~\ref{sec:deformation-chart}.
Let \(\mathcal J\) be the coefficient ideal of series killed by
the chart and \(M=\mathcal R/\mathcal J\). The list \(S,a\)
satisfies the finiteness condition on \(M\). To see it, take a
prime containing \(\mathcal J+(S,a)\) which avoids \(\mathfrak b\).
It contains all \(Z_i\) because the chart kills them. By
\eqref{eq:rees-contract} its contraction to \(D\) contains \(I_D\).
A vertical contraction equals \(\mathfrak m_D\) by the finiteness
of \(D/I_D\). A horizontal contraction admits an embedding in
\(\overline F\). It cannot contain \(\Delta\) by
\eqref{eq:rees-boundary}. As proved in Section~\ref{sec:deformation-support},
this embedding extends to a mixed geometric solution of the
equations. Its value on \(a\) would be zero, contradicting the
nonvanishing assumption. Thus each prime in question contains
\(\mathfrak m_D\). Taking radicals and then a power gives the stated
finiteness condition.

The chosen \(\phi(a)\) is a nonzero series and \(a\) is regular on
\(M\). Indeed multiplication by \(a\) acts on each coefficient via
multiplication by \(\phi(a)\) in a domain. Taking the Euler sum first
in \(a\), then moving \(a\) to the end, and discarding the
annihilator terms for \(S\) gives
\[
 \chi(S;M/aM)_d=\chi(a,S;M)_d=\chi(S,a;M)_d
              \le\chi(a;M/(S)M)_d .
\]
This proves the middle inequality of \eqref{eq:single-rees}.

\subsection{Extracting the uniformizer powers}

Use the leading ideal \(\mathfrak q\) of
Section~\ref{sec:deformation-chart} and put \(G=\mathcal R/\mathfrak q\).
The kernel of the chart lies in the leading kernel in every degree.
Also \(a\in\mathfrak q\), since \(x_A(a)=0\) makes every monomial
in its degree-zero image have positive negative weight. Consequently
\(M/aM\) surjects onto \(G\). The quotient comparison gives
\[
                         \chi(S;G)_d\le\chi(S;M/aM)_d .
\]
The leading coefficient lattice in every degree is torsion-free,
so \(\varpi\) is regular on \(G\). The special fiber
\(\overline G=G/\varpi G\) has finite-length components, with the
explicit Hilbert function \(H(d)=\operatorname{length}_D\overline G_d\)
proved in Section~\ref{sec:deformation-chart}. For any list of
sections \(U\), the chain modules in its degree-\(d\) Koszul complex
on \(\overline G\) all have finite length. Taking the alternating
sum on the chains or on their homologies gives
\[
        \chi(U;\overline G)_d=
              \prod_{u\in U}(1-T_{\deg u})H(d),
       \qquad T_p H(d)=H(d-p).
\]

The classes of the first \(r\) sections of \(S\) on \(G\) are
\(\varpi^{2c_i} y_i^2 T_0\). In fact the difference
\(F_i-\varpi^{c_i}y_i\) has weight \(-1\) and belongs to the
saturated square, hence has negative filtration at least two.
The sum \(F_i+\varpi^{c_i}y_i\) has filtration at least one.
Their product has filtration at least three, so it vanishes on the
layer \(n=2\) of the first Rees component. This proves the claim.

Let \(y\) be any section of degree \(p\), \(c\ge0\), and let \(U\)
be a remaining list for which \(\varpi^c y,U\) satisfies the
finiteness condition on \(G\). The following extraction is exact.
\[
   \chi(\varpi^c y,U;G)_d
       =\chi(y,U;G)_d+c\,\chi(U;\overline G)_d .
\]
Indeed \(G[\varpi^c y]=G[y]\) by regularity, and there is a short
exact sequence
\[
   0\longrightarrow G/yG
       \xrightarrow{\ \varpi^c\ }G/(\varpi^c y)G
       \longrightarrow G/\varpi^cG\longrightarrow0 .
\]
Injectivity follows by cancelling \(\varpi^c\) in \(G\).
The final module has \(c\) successive factors isomorphic to
\(\overline G\), by the same cancellation. Additivity for \(U\)
and the Euler recursion prove the identity. Stripping a scalar
factor enlarges the ideal of sections, so the required finiteness
condition is preserved at subsequent steps.

Apply the identity in turn to the first \(r\) sections, moving each
into first position by Koszul permutation. In the term measured on
\(\overline G\), all the other \(r-1\) sections have the same degree
\(e_0\), regardless of the scalar powers already stripped. Thus each
extraction contributes \(2c_i\) times the same difference
\[
        \mathcal D H=(1-T_{e_0})^{r-1+s}
                        (1-T_{e_1})^{r'+s'}(1-T_{e_2})H.
\]
The Euler characteristic on \(G\) with all \(r\) scalar powers
stripped is eventually nonnegative. We conclude that
\[
 \chi(S;M/aM)_d\ge \chi(S;G)_d
                \ge\left(\sum_{i=1}^r2c_i\right)\mathcal D H(d).
\]
The generating-function computation in
Section~\ref{sec:single-character} gives
\(\mathcal DH(d)=2^{r-1+r'}\) for sufficiently large \(d\).
This proves the leftmost inequality in \eqref{eq:single-rees}.
Combining it with the middle comparison above and the bound of
Section~\ref{sec:deformation-support} proves the full graded
intersection estimate for the stated model.

\subsection{The cut chart and finiteness of the set of orbits}

We give the last two geometric constructions used above from the
model hypotheses. Before cutting, smoothness in the source gives
a graded chart \(\phi_0\) over \(F\) with \(m\) variables of weight
zero, in addition to the \(r,s,r'+1,s'+1\) variables of nonzero
weights already specified. Its linear terms separate the axes, and
an augmentation element whose image has zero linear term belongs
to the saturated square. Thus the map from the saturated cotangent
lattice of weight \(-1\) to the vector space of linear coefficients
on the \(r\) weight-\(-1\) variables is injective. Every coordinate
axis is met by the image, up to a nonzero scalar, by projecting an
axis-separating element to weight \(-1\). This proves that the
lattice has rank \(r\). Its finite generation over \(\mathcal O\)
follows from finite generation of \(A_{-1}\) over \(D\): on the
cotangent classes \(D\) acts through the augmentation, since
\((b-x_A(b))h\in(\ker x_A)^2\) for every \(h\in A_{-1}, b\in D\).
The weight-zero analogue, applied after subtracting constant terms,
shows that the linear coefficient vectors of the \(m\) equations
\(Z_i\) span all \(m\) zero-weight directions because those equations
span their saturated cotangent lattice over \(F\). They therefore
give an invertible square matrix of linear terms in the zero-weight
variables. There are no linear terms of other weights.

In the formal series ring over \(F\), replace the \(m\) zero-weight
variables by \(\phi_0(Z_i)\), keeping the nonzero-weight variables.
This substitution has zero constant part and invertible linear part.
It has a unique inverse by formal degree recursion: solve the
linear part first, and at each subsequent degree apply its inverse
matrix to the error in that homogeneous degree. Each coefficient
stabilizes at its own degree. The substitution preserves weights, and
so does its inverse by the same recursion. Compose the original
chart with this inverse and then set all zero-weight variables to
zero. This kills the \(Z_i\), leaves the constant coefficient equal
to \(x_A\), and preserves the nonzero-weight linear terms, because
the linear substitution was block diagonal and the nonzero block
was the identity. This is exactly the cut chart used in the leading
coefficient calculation.

For orbit finiteness take finite \(D\)-module generating sets of
\(A_{-2}\) and \(A_2\), whose existence was shown in
Section~\ref{sec:deformation-chart}. Every mixed point lies where
some product \(fg\) of such generators is nonzero. On this locus
each orbit has a point with \(f=1\), since under torus scaling \(f\) is
multiplied by \(z^{-2}\), and the base field is algebraically closed.
The slice algebra on \(f=1\), \(g\ne0\), after base change to
\(\overline F\), was shown there to have a field as every local ring
at a maximal ideal. It is Noetherian of dimension zero. To see the
dimension assertion, any prime is contained in a maximal ideal and
the localization there is a field, so the two primes coincide.
A zero-dimensional Noetherian ring is Artinian and has finitely many
maximal ideals. Thus this slice has finitely many geometric points.
The finitely many pairs \((f,g)\) cover the mixed locus, giving
finitely many mixed orbits and the desired finite set of
representatives \(T\).

The source comparisons occur in \(\mathrm{CutHolds}\),
\(\mathrm{LaDifPerm}\), \(\mathrm{AdmissibleMonoIdeal}\), and
\(\mathrm{AlgMemCoef}\). This completes the graded intersection
calculation from the smooth model, the cotangent spanning conditions,
finiteness over the coefficient DVR in weight zero, and the two
geometric point conditions. Constructing the arithmetic model and
establishing those conditions for its Galois points use the
deformation and automorphy theorems specified in
Section~\ref{sec:single-character}.

\section{Selmer rigidity of a modular type-II point}
\label{sec:kato-rigidity}

Let \(K/\mathbb Q\) be the imaginary quadratic field of the window,
\(\theta\) its character, \(G=G_{\mathbb Q}\), and
\(\epsilon=\epsilon_\ell\) the cyclotomic character.
We prove the characteristic-zero Selmer vanishing used at a modular
type-II point. The arithmetic inputs are the finiteness theorem of
Kato and the adjoint vanishing theorem of Newton and Thorne. The
identification of the point as modular is a separate hypothesis.

\subsection{The local conditions and the two-dimensional summands}

For a continuous representation \(V\), write \(H^1_f(F,V)\) for
the classes of continuous cocycles whose restrictions to inertia
are coboundaries at every finite place away from \(\ell\), and whose
restrictions at \(v\mid\ell\) become coboundaries in
\(V\otimes_{\mathbb Q_\ell} B_{\rm cris}\). The latter condition says
that \(c(s)\otimes1=sx-x\) on the local absolute Galois group for
one element \(x\) of this tensor product with diagonal action.
Restrictions, finite direct sums and continuous equivariant linear
maps preserve these conditions. They agree with the conditions called
\(\mathrm{IsBK}\) and \(\mathrm{finiteCocycles}\) in the original
release. Choosing an embedding into a local algebraic closure selects a
decomposition group. Conjugating the embedding gives the same condition:
transport the local witness by the conjugating element and adjust by
the value of the cocycle on that element.

Let \(e\ge1\), \(k=2e+1\), and let
\(W:G\longrightarrow {\rm GL}_2(\overline{\mathbb Q}_\ell)\)
be continuous. Fix
\(\iota:\overline{\mathbb Q}_\ell\simeq\mathbb C\).
Suppose a normalized cuspidal newform \(f\) of weight \(k\), level
\(M_f\) and nebentypus \(\omega_f\) is attached to \(W(-e)\).
Here attachment requires semisimplicity, and for all but finitely
many primes \(p\), unramifiedness and the characteristic polynomial
\[
 X^2-a_p(f)X+\omega_f(p)p^{k-1}
 \quad\hbox{at geometric Frobenius, after applying }\iota.
\]
The original newform predicate imposes the good-prime eigenrelations,
normalization and exclusion of lower-level cusp forms with the same
good eigenvalues. Classical newform decomposition yields a primitive
eigenform with these good eigenvalues: decompose into old and new
eigensystems, choose a primitive constituent, and observe that a
constituent of smaller conductor would itself be a forbidden
lower-level form. At its primitive conductor the newvector line has
dimension one and the normalization fixes the scalar.

\paragraph{Proposition.}
Under these hypotheses
\[
 H^1_f(\mathbb Q,W)=H^1_f(\mathbb Q,W\otimes\theta)=0 .
 \tag{K1}
\]

\paragraph{Proof.}
Put \(E=F_\lambda\), where \(F\) is the coefficient field of \(f\)
and \(\lambda\) is the place induced by \(\iota^{-1}\).
Kato's representation \(V_E(f)\) has the displayed polynomial
at the inverse of arithmetic Frobenius (Example 13.3 of his paper).
Its semisimplicity and the Chebotarev density theorem give equality of
the traces of \(V_E(f)\) and \(W(-e)\) on all of \(G\) after scalar
extension. The characteristic-zero Brauer--Nesbitt theorem therefore
gives
\[
 W\simeq U\otimes_E\overline{\mathbb Q}_\ell,\qquad U=V_E(f)(e).
\]
In particular this is the positive Tate twist of Kato's cohomological
representation. The characteristic polynomial at geometric Frobenius
fixes the sign.

Take a \(G\)-stable \(\mathcal O_E\)-lattice \(T\) of \(U\).
Kato, \emph{\(p\)-adic Hodge theory and values of zeta functions of
modular forms}, Ast\'erisque 295 (2004), \S\S13.3, 14.1 and
Theorem 14.2(1), proves that for a finite abelian extension
\(L/\mathbb Q\) the discrete Selmer group of a stable lattice of
\(V_E(f)(r)\) is finite whenever \(1\le r\le k-1\) and \(2r\ne k\).
The discrete group is contained in \(H^1(L,U/T)\) and its local
conditions are the images of \(H^1_f(L_v,U)\).
Here \(r=e\) satisfies these inequalities and \(L=K\) is abelian.
The theorem applies to CM newforms as well.

We give the passage from discrete finiteness to vector-space
vanishing. The reduction map sends \(H^1_f(K,U)\) into this finite
group. Its image is divisible because its domain is an \(E\)-vector
space, so its image is zero. Let \(c\) represent any of its classes
and let \(\varpi\) be a uniformizer. For each \(j\ge1\), the reduction
of \(\varpi^{-j}c\) is therefore a coboundary. Lift the element of
\(U/T\) giving this coboundary to \(b_j\in U\). Then
\[
 c(s)-(s-1)\varpi^j b_j\in\varpi^jT\qquad(s\in G_K).
\]
The coboundaries form the image of the finite-dimensional space \(U\)
in the normed space of continuous functions \(G_K\to U\), with the
uniform norm from \(T\). This image is finite dimensional over the
complete local field \(E\), hence closed. The display gives uniform
convergence to \(c\), proving \(H^1_f(K,U)=0\).

We record the scalar-extension argument because the cocycle in
the statement takes values in an algebraic closure. Set
\(B=\overline{\mathbb Q}_\ell\). Continuous \(E\)-linear maps
\(\lambda:B\to E\) separate points. Indeed, on the line generated
by \(a\ne0\) prescribe \(\lambda(a)=1\), with norm bound
\(C=|a|^{-1}\). To extend a functional from a subspace \(D\)
of \(B\) to \(D+Ey\), choose its value at \(y\) in the balls
\[
 \{b\in E: |b-\lambda(x)|\le C|y-x|\},\qquad x\in D.
\]
The ultrametric inequality makes these balls nested. Their
intersection is nonempty: choose one of them, use its compactness
in the local field \(E\), and use the finite intersection property
for the smaller closed balls. The chosen value preserves the
bound on \(D+Ey\). The union of a chain of extensions preserves
the bound, so Zorn's lemma gives the required map on \(B\).

Apply every such \(\lambda\) coordinatewise to a finite cocycle
\(c:G_K\to U\otimes_E B\). This gives a continuous \(U\)-valued
cocycle satisfying the same local conditions, since the coefficient
map carries each local witness to a witness and extends by tensoring
with the identity on \(B_{\rm cris}\). Thus every projected cocycle
is a coboundary. To see that \(c\) is one, write the equations
\((U(s)-1)x=c(s)\) in an \(E\)-basis. Their left sides are rows
in the finite-dimensional \(E\)-dual of \(U\). Choose among all
rows a basis for their span, of dimension \(m\). Every other row is an \(E\)-linear
combination of the chosen ones. The same linear relation on the
right sides holds after every \(\lambda\), since the projected
equations have a solution. Separation of points gives each relation
in \(B\). The independent rows define a surjection to \(B^m\), so we can solve the basis equations.
This solution satisfies all equations. We have proved
\(H^1_f(K,W)=0\) without a restriction on the field generated by
the values of a cocycle.

Finally \(W\) and \(W\otimes\theta\) restrict to the same
representation on \(G_K\). Restriction preserves the local conditions.
Inflation--restriction identifies the kernel of restriction in
degree one with \(H^1({\rm Gal}(K/\mathbb Q),V^{G_K})\), for either
\(V\). Averaging over this group of order two kills that cohomology
in characteristic zero. Both restriction maps are injective,
proving (K1). \(\square\)

\subsection{The polarized tangent module}

We spell out the input supplied by the type-II classification.
Relative to the order \(0,1,2\), the matrix part of each \(r(s)\) is
\[
 P_s=\begin{pmatrix}a&0&b\\0&1&0\\c&0&d\end{pmatrix},\qquad
 M_s=\begin{pmatrix}a&b\\c&d\end{pmatrix}\in{\rm SL}_2 .
\]
Its component is \(\theta(s)\). At the chosen complex conjugation
\(c_\infty\), \(P_{c_\infty}=1\) and \(\theta(c_\infty)=-1\).
On \(G_K\) the attached two-dimensional representation is \(M_s\);
on the other coset it is \(M_sD\), with \(D={\rm diag}(-1,1)\).
The irreducibility supplied by the window is that of
\(W|_{G_{K(\mu_{\ell^\infty})}}\) over \(\overline{\mathbb Q}_\ell\).

The tangent module is the full space of three-by-three matrices.
On \(G_K\) the action is ordinary conjugation by \(P_s\). The chosen
complex conjugation acts by \(X\mapsto-J_3X^tJ_3\), where \(J_d\) is
the antidiagonal identity matrix of size \(d\). For \(t=1,-1\) set
\[
 p_t(X)=\binom{X_{01}+tX_{12}}{X_{21}-tX_{10}},\qquad
 s_t(u,v)=\begin{pmatrix}0&u&0\\-tv&0&tu\\0&v&0\end{pmatrix}.
\]
Let \(A(X)\) extract the outer corner on \(0,2\) and let \(i_A\)
insert it there. Write \(i_C(X_{11})\) for the central entry alone.
Direct substitution gives
\[
 X=i_A A(X)+\tfrac12s_+(p_+(X))
                +\tfrac12 s_-(p_-(X))+i_C(X_{11}).
 \tag{K2}
\]
Here \(+\) and \(-\) denote \(1\) and \(-1\). To verify the actions,
the two columns \((X_{01},X_{21})^t\) and
\((X_{12},-X_{10})^t\) both transform by \(M_s\) on \(G_K\).
For the second column use
\((X_{10},X_{12})\mapsto(X_{10},X_{12})M_s^{-1}\) and
\(\det M_s=1\). Complex conjugation exchanges these two columns with
the factor \(D\): the first becomes \(D\) times the second and the
second becomes \(D\) times the first. Thus their sum realizes \(W\),
their difference \(W\theta\). The center realizes \(\theta\).

On the outer corner \(G_K\) acts by conjugation. At \(c_\infty\)
the formula on two-by-two matrices is \(Y\mapsto-J_2Y^tJ_2\).
This is multiplication by \(-1\) on scalars. For
\(Y=\left(\begin{smallmatrix}u&v\\w&-u\end{smallmatrix}\right)\)
it equals \(DYD^{-1}\). Since \(2\) is invertible, the corner
decomposes as the trace-zero part with action \({\rm ad}^0 W\)
and a second copy of \(\theta\). Hence
\[
 {\rm ad}\,r \simeq
 {\rm ad}^0 W\ \oplus\ W\ \oplus\ (W\theta)\
                  \oplus\ \theta\ \oplus\ \theta .
 \tag{K3}
\]

\paragraph{Proposition.}
A modular type-II point with the block and irreducibility properties
above satisfies \(H^1_f(\mathbb Q,{\rm ad}\,r)=0\).

\paragraph{Proof.}
The newform \(f\) supplies a regular algebraic cuspidal
\({\rm GL}_2/\mathbb Q\) representation since \(k\ge2\).
In J.\ Newton and J.\ A.\ Thorne, \emph{Adjoint Selmer groups of
automorphic Galois representations of unitary type},
J.\ Eur.\ Math.\ Soc.\ 25 (2023), Theorem 5.4, the full adjoint
Selmer group of this representation vanishes if the form is non-CM,
or if its CM field is not contained in \(\mathbb Q(\mu_{\ell^\infty})\).
In the remaining possibility, with CM field \(L\) contained in that
cyclotomic field, the two-dimensional Galois representation is
induced from a character of \(G_L\). It splits into two characters
over \(G_L\) and hence is reducible on
\(G_{K(\mu_{\ell^\infty})}\subset G_L\). Twisting by \(\epsilon^e\)
preserves this reducibility, contrary to the specified hypothesis on
\(W\). The theorem therefore applies in every allowed case. Scalar
twists preserve the adjoint. The dual representation gives an
isomorphic adjoint via transposition, so either standard normalization
of an attached modular representation gives the same conclusion.
The trace-zero piece is a direct summand and
\(H^1_f(\mathbb Q,{\rm ad}^0W)=0\).

For the remaining character, restriction to \(K\) again injects,
and \(\theta|_{G_K}=1\). Over \(K\), the finite local condition
for the trivial representation \(E\) is unramified everywhere.
At places \(v\mid\ell\) this uses the standard local Bloch--Kato
fundamental exact sequence: since
\(D_{\rm dR}(E)/{\rm Fil}^0=0\), it identifies the finite local condition
with the cokernel of \(1-\varphi\) on \(D_{\rm cris}(E)\).
This group has dimension one over \(E\), since its kernel has
dimension one by unramified Frobenius descent. The unramified
line lies in it. Indeed an extension of the trivial representation
by itself with unramified action is crystalline (equivalently use
an unramified period solving \(\varphi^{f_v}(b)-b=1\), where
\(f_v\) is the residue degree).
A continuous additive character of \(G_K\) unramified at every
finite prime factors through the finite ideal class group by global
class field theory. Its image in \(E\) is zero. The separation
argument above also applies to \(B\)-valued classes. Thus
\(H^1_f(\mathbb Q,\theta)=0\). Apply (K1) and (K3). \(\square\)

\subsection{The Euler-system argument in the original release}

The archive reaches (K1) by constructing zeta classes. We describe
this separate derivation and the intermediate results it uses.
The original \(\mathrm{KatoTwoW}\) statement assumes an odd prime
\(\ell\) with \(\ell-1\nmid 2e\) and
\(\ell\nmid |\operatorname{disc}K|h_K\), an even
\(N>0\) prime to \(\operatorname{disc}K\), and a finite character
\(\psi\) factoring through \(\mathbb Q(\mu_N)\), of order prime
to \(\ell\) and quadratic on inertia above \(\ell\).
It imposes \(\Psi(c_\infty)=1\) for \(\Psi=\epsilon^e\psi\).
The level is \(N|\operatorname{disc}K|\) with ordered weights
\((e,-e)\), and the ordinary inertia polynomial after a finite
cyclotomic extension has roots \(\epsilon^e,\epsilon^{-e}\).
One assumes \(\det W=\theta\), inertia trace
\(\Psi+\Psi^{-1}\), and everywhere the residual trace is
\(\overline{\theta\Psi+\Psi^{-1}}\). At one place above \(\ell\)
there is a decomposition-group flag with diagonal inertia characters
\(\epsilon^e\chi,\epsilon^{-e}\chi\), with \(\chi\) factoring through
\(\mathbb Q(\mu_\ell)\), and \(\theta\) is trivial on this inertia
group. The irreducibility condition is the one in (K3).
The type-II point supplies these arithmetic properties through
the result \(\mathrm{used\_of\_alt}\). The archived version also
takes a rational cuspidal automorphic witness and the newform
attached to \(W(-e)\).

The comparison of the Galois row with quadratic base change and
its unramified local parameters provides a spherical compatibility
statement for the representation and its finite cyclotomic twists.
Comparison with an ordinary newform gives a newform \(g\) of
weight \(2e+1\), prime-to-\(\ell\) level, and
\(\lvert\iota^{-1}a_\ell(g)\rvert_\ell=1\). It is attached to
\[
 W\otimes\epsilon^{-e}(\det W)^{-1}\chi'
\]
for a character \(\chi'\) factoring through \(\mathbb Q(\mu_\ell)\).
For the selected flag with lower diagonal character \(\kappa_1\)
one has \(\kappa_1\epsilon^e(\det W)^{-1}\chi'=1\) on inertia.
The comparison also supplies nonexceptionality: for each admissible
decomposition-group flag, its first character differs from \(\epsilon\)
on some element of that decomposition group. These are the
unramified-source conditions in \(\mathrm{SrcL}\). After conjugation
to an ordinary integral basis they feed the construction
\(\mathrm{Final.zetaUr}\).

In this basis let \(A\) be the chosen place above \(\ell\) and \(D_A\)
its decomposition group. Write \({\rm Gr}\) for unramifiedness away
from \(\ell\) together with the condition that the lower coordinate
of the cocycle is a coboundary on \(D_A\). Write \({\rm Str}\) for
the same away condition together with being a coboundary in both
coordinates on \(D_A\). The Tate dual is
\(W'=W\otimes\epsilon(\det W)^{-1}\), via the alternating
form \(\omega(v,w)=v_0w_1-v_1w_0\). The zeta construction gives
an integral cocycle \(y\) unramified away from \(\ell\) whose lower
coordinate is not a coboundary on \(D_A\).

More precisely its norm relations supply the following data. There
is a finite exceptional set such that for \(n\ge1\), every auxiliary
prime \(q\) outside that set with \(\ell^n\mid q-1\), and every
surjective \(\psi_q:G\to\mathbb Z/\ell^n\) factoring through
\(\mathbb Q(\mu_q)\), one has integral vectors \(Y(s,x)\) for
\(x\in\mathbb Z/\ell^n\), a scalar \(P\), a unit \(\alpha\) of
absolute value one, and an integral vector \(b\), with
\begin{align*}
 Y(st,x)&=Y(s,x)+W(s)Y(t,x-\psi_q(s)),\\
 \sum_xY(s,x)&=Py(s)+(W(s)-1)b,\\
 P&\equiv \alpha\det(W({\rm Frob}_q)-1)\pmod{\ell^n}.
\end{align*}
The first family is continuous. The last congruence holds for each
admissible Frobenius lift, and the source imposes integral
inertia-coboundary witnesses for the family at all other primes
away from \(\ell\). Smoothed zeta classes, their integrality and
Hodge--Tate law, a nonvanishing Hecke projection, and a local
projection formula are the inputs to these assertions.

Here is the height contradiction that uses the norm relations.
Suppose that a strict dual cocycle \(z\) represents a nonzero class.
After enlarging the integral coefficient ring and multiplying \(y,z\)
by nonzero powers of \(\ell\), their values and their required local
witnesses are integral. Put
\[
 H=\ker\epsilon\cap\ker\det W,\qquad B_0=\ker\epsilon\cap\ker W.
\]
The cyclotomic-kernel lemmas in the release assert that \(y|_{B_0}=0\)
would give the ordinary condition on \(y\), and \(z|_{B_0}=0\)
would make \(z\) a global coboundary. The pairing and commutator
lemmas for the irreducible ordinary representation therefore supply
\(g,h\in B_0\) and \(\phi\in[H,H]\) such that
\[
 d'=\omega(y(h),z(g))-\omega(y(g),z(h))\ne0,\qquad
 d=\det(W(\phi)-1)\ne0.
\]
The elements are fixed for every precision \(n\).

For each \(n\ge1\), the cup product
\(\omega(y(s),W'(s)z(t))\) is locally a coboundary modulo
\(\ell^n\), by strictness at \(\ell\) and the unramified local cup
product away from \(\ell\). The global cochain gluing theorem
supplies an integral cochain \(u\) whose coboundary is this cup
product modulo \(\ell^n\), with the unramified-outside-\(S\)
condition. Define
\[
 T(s)=\det(W(s)-1)u(s)
  -\omega\bigl(\operatorname{adj}(W(s)-1)y(s),z(s)\bigr).
\]
The tame-lift theorem gives
\(\det(W(s)-1)T(s)\equiv0\bmod\ell^n\) on \([H,H]\).
Write \(c=[g,h]\). For clarity the commutator-height formula can
be seen directly. On \(B_0\) both actions are trivial, so \(y,z\)
restrict to homomorphisms and \(y(c)=z(c)=0\).
The cochain equation becomes
\[
 u(st)\equiv u(s)+u(t)-\omega(y(s),z(t))\pmod{\ell^n}
 \qquad(s,t\in B_0).
\]
Apply it to \(gh\) and \(hg\). Since \(gh=c\,hg\), it gives
\(u(c)\equiv d'\). The equation on \((\phi,c)\) now gives
\(u(\phi c)-u(\phi)\equiv d'\), because \(\epsilon(\phi)=1\).
The matrix and both cocycles at \(\phi c\) equal their values at
\(\phi\), hence \(T(\phi c)-T(\phi)\equiv dd'\). Both \(\phi\) and
\(\phi c\) belong to \([H,H]\), so the tame-lift theorem and
integrality of \(d\) give \(d^2d'\equiv0\bmod\ell^n\).
This holds for every \(n\) and contradicts \(d^2d'\ne0\).
Thus the strict dual group vanishes. The arithmetic input to
this calculation consists of the kernel and pairing lemmas,
local-to-global cochain gluing and the tame-lift height congruence.

The release upgrades strict dual vanishing to ordinary dual
vanishing by reciprocity for integral cup products and local
duality. After subtracting a coboundary, the lower coordinate
of an ordinary dual class is zero on \(D_A\). The local orthogonality
alternative makes either the lower coordinate of \(y\) a coboundary
or this adjusted dual class strict. The first option contradicts the
zeta class, and strict dual vanishing settles the second. Global
Selmer duality then gives ordinary vanishing for \(W\): a nonzero
ordinary class for \(W\) would produce dual finite-level classes
of unbounded order, and the compactness argument for these classes
produces a nonzero ordinary dual class.

Finally a finite class for \(W\) is ordinary. Its lower projection
is crystalline. For an upper triangular inertia matrix with diagonal
entries \(a,b\), the normal form gives \(a^m=\epsilon^{em}\)
for some \(m>0\) and \((ab)^2=1\). Thus
\((b\epsilon^e)^{2m}=1\) on inertia. The lower character has
weight \(-e\) in the cyclotomic-exponent convention. The
one-dimensional crystalline finite-class vanishing in negative weight
gives a coboundary on \(D_A\), as required. Applying this construction
to \(W\) and \(W\det W\) gives precisely (K1): the determinant
is quadratic and factors through a finite cyclotomic field, so the
spherical compatibility statement over all twists also applies to
the latter representation.

The original declarations for this route include
\(\mathrm{katoTwoW\_of\_printsH}\), \(\mathrm{twoS}\),
\(\mathrm{vanishOrdL\_of\_zetaUr}\), \(\mathrm{es\_vanish}\),
and \(\mathrm{StrictKill.main}\). Its type-II tangent proof projects
a cocycle by (K2), applies polarized-adjoint rigidity to the corner
and the nonpositive weight result to the center, then reassembles
the four coboundaries. The two propositions above supply a
classical derivation of the entire vanishing from the stated point
hypotheses. Establishing those hypotheses for all points of the
deformation ring, including production of \(f\), belongs to the
automorphy and classification arguments.

\section{Strong irreducibility and the automorphic type-I tangent space}
\label{sec:type-one-rigidity}

Let \(K/\mathbb Q\) be the imaginary quadratic field,
\(H=G_K\), and \(\epsilon=\epsilon_\ell\). In a normalized type-I
point \(r\), let \(\rho:H\to {\rm GL}_3(B)\) be its matrix part,
where \(B=\overline{\mathbb Q}_\ell\). The type-I condition is that
\(\rho\) has no proper nonzero invariant \(B\)-subspace. Write
\(\Psi=\epsilon^e\psi\), with \(e\ge1\), \(\psi\) of finite
cyclotomic conductor dividing \(N\). The ordinary and residual
conditions on the point supply an integral model with the following
two properties:
\begin{enumerate}
\item On decomposition at \(\ell\) inside \(H\) there is a common
invariant full flag; on inertia its successive eigencharacters
are \(\Psi,1,\Psi^{-1}\), in the order fixed by the window.
\item Reduction at the maximal ideal of the integral model satisfies
\[
 \overline{{\rm tr}\,\rho(s)}
   =\overline{\Psi(s)}+1+\overline{\Psi(s)}^{-1}
           \qquad(s\in H).
\]
\end{enumerate}
The residue hypothesis \(\ell-1\nmid 2e\) is essential to the
argument below. We first prove the open-subgroup irreducibility
assertion used in the original release.

\paragraph{Proposition.}
For every open subgroup \(V\) of \(H\), \(\rho|_V\) is irreducible.

\paragraph{Proof.}
Take an open normal subgroup \(U\) of \(G_\mathbb Q\) contained
in \(V\), obtained by intersecting its finitely many conjugates.
It is enough to prove irreducibility on \(U\). Choose a place over
\(\ell\). Its inertia subgroup lies in \(H\), since \(\ell\) does
not divide the discriminant of \(K\). The cyclotomic character on
this inertia subgroup surjects onto \(\mathbb Z_\ell^\times\):
the extensions \(\mathbb Q_\ell(\mu_{\ell^j})/\mathbb Q_\ell\)
are totally ramified with Galois groups
\((\mathbb Z/\ell^j\mathbb Z)^\times\), and passage to the inverse
limit gives the assertion. Choose \(\sigma\) on inertia with
\(\epsilon(\sigma)=1+\ell\). Put
\(m=[G_\mathbb Q:U]\cdot|(\mathbb Z/N)^\times|\). Then
\(u=\sigma^m\) lies in \(U\), and \(\psi(u)=1\). The three
eigenvalues of \(\rho(u)\) are distinct since
\(\Psi(u)^2=(1+\ell)^{2em}\ne1\).

Choose also \(x\) on inertia whose cyclotomic character generates
\(\mathbb F_\ell^\times\) after reduction. The window assumes
\(\psi^2=1\) on this inertia subgroup. Consequently
\(\overline{\Psi(x)}^2=\overline{\epsilon(x)}^{2e}\ne1\).
Both \(x\) and \(u\) preserve the ordinary flag. Each subspace
in the flag is spanned by some eigenlines of the regular semisimple
element \(\rho(u)\): its spectral projections are polynomials in
\(\rho(u)\).

Suppose \(\rho|_U\) is reducible. The restriction is semisimple:
take a nonzero simple \(U\)-submodule and all its \(H\)-translates.
Normality of \(U\) makes each translate simple, and their sum is
the entire space by \(H\)-irreducibility. A sum of simple
submodules is semisimple. The group \(H\) permutes its isotypic
components transitively, so they have equal dimension and equal
multiplicity of their respective simple constituents.
As the space has dimension three, reducibility on \(U\) now gives
two possibilities. One is three copies of one character, on which
\(u\) would be scalar. This contradicts its distinct eigenvalues.
The other is three pairwise distinct characters, with three lines
permuted transitively by \(H\). These lines are the eigenlines of
\(\rho(u)\). The permutation induced by \(x\) preserves each stage
of the complete ordinary flag and hence fixes the three lines
individually.

The transitive permutation image of \(H\) on three letters has
order divisible by three. Choose \(g\in H\) inducing a three-cycle.
Then \(\rho(gx^j)\) has zero diagonal in the eigenbasis of
\(\rho(u)\), for every \(j\ge0\). Its trace therefore vanishes.
In the residue field put \(a=\overline{\Psi(g)}\) and
\(b=\overline{\Psi(x)}\). The residual trace formula says
\[
 a b^j+1+(a b^j)^{-1}=0\qquad(j\ge0).
\]
Thus all the nonzero elements \(a b^j\) are roots of
\(T^2+T+1\). In particular among \(a,ab,ab^2\) two coincide,
forcing \(b=1\) or \(b^2=1\). This contradicts the choice of \(x\).
Irreducibility on \(U\) and on \(V\) follows. \(\square\)

\paragraph{Lemma.}
Let \(V\) be a finite-dimensional characteristic-zero continuous
representation of \(H\) over \(B\) which is irreducible on every
open subgroup of \(H\). For any continuous abelian quotient
\(H\to C\), the restriction of \(V\) to its kernel is irreducible.

\paragraph{Proof.}
Let \(M\) be the Zariski closure of the image in \({\rm GL}(V)\);
write \(M^0\) for its identity component. The preimage of \(M^0\)
is open in \(H\): there are finitely many components, each closed
in the \(\ell\)-adic topology, so its complement is closed as well.
Its image is Zariski dense in \(M^0\), since this component is also
open in \(M\) for the Zariski topology. In particular \(M^0\)
acts irreducibly. Its unipotent radical has a nonzero fixed vector
by triangularization. The whole fixed space is \(M^0\)-stable by
normality, so the unipotent radical fixes \(V\) pointwise and is
trivial in this faithful realization. It follows that \(M^0\) is
reductive. Its connected center acts by scalars by Schur's lemma.
A connected reductive group is the product of its connected center
and its derived subgroup. Hence \((M^0)^{\rm der}\) already acts
irreducibly on \(V\).

All group commutators of the preimage of \(M^0\) belong to the
kernel of \(H\to C\). The Zariski closure of its image there
contains the commutators of all pairs in \(M^0\), by density
and the algebraicity of the commutator map. It therefore
contains \((M^0)^{\rm der}\). A subspace invariant under the kernel
is invariant under its Zariski closure, and the preceding
irreducibility finishes the proof. \(\square\)

\paragraph{The polarized adjoint.}
In addition assume the type-I automorphy conclusion: \(\rho(-1)\)
is attached to a regular algebraic, conjugate self-dual cuspidal
automorphic representation \(\pi\) of \({\rm GL}_3(\mathbb A_K)\).
This is the automorphic witness in \(\mathrm{IsAutIAlt}\), whose
definition also carries local compatibility information. Attachment
compares the unramified normalized Satake roots with geometric
Frobenius; its factor for rank three is \(({\rm N}v)^{(3-1)/2}\).
With this normalization the semisimple representation
\(r_{\pi,\iota}\) of the automorphic Galois theorem identifies with
\(\rho(-1)\) by Chebotarev and Brauer--Nesbitt. The preceding lemma
for the cyclotomic quotient shows that \(\rho\), and hence its scalar
twist, is irreducible on \(G_{K(\mu_{\ell^\infty})}\).

Write \(J=J_3\) for the antidiagonal identity matrix, and choose
the normalized complex conjugation \(c\), whose matrix part is \(1\).
The semidirect-product law of the point gives
\[
 \rho(chc^{-1})=J\,{}^t\!\rho(h)^{-1}J^{-1}\qquad(h\in H).
\]
For \(V=\rho(-1)\), \(\epsilon(chc^{-1})=\epsilon(h)\) gives
\[
 V(chc^{-1})=\epsilon(h)^{-2}J\,{}^t\!V(h)^{-1}J^{-1}.
\]
Thus \(J\) supplies the symmetric polarization
\(V^c\simeq V^\vee\otimes\epsilon^{1-3}\).
Any two such polarizations on the irreducible representation
differ by a scalar, which cancels in the adjoint. On \(H\) the
polarized adjoint acts by conjugation by \(V(h)\), or equivalently
by \(\rho(h)\); on \(c\) it acts on all nine matrix entries by
\(X\mapsto-J\,{}^t\!XJ^{-1}\). These are precisely the actions
in the tangent module of the point.

In \emph{On the vanishing of adjoint Bloch--Kato Selmer groups of
irreducible automorphic Galois representations}, Theorem 6.1,
J.\ A.\ Thorne proves that for a CM field \(F\), a polarizable
regular algebraic cuspidal \(\pi\) over \(F\), and
\(r_{\pi,\iota}\) irreducible over \(F(\mu_{\ell^\infty})\), the
Selmer group of the polarized adjoint over \(F^+\) vanishes.
With \(F=K\) the base \(F^+\) is \(\mathbb Q\). The local
conditions are unramified away from \(\ell\) and the crystalline
kernel at \(\ell\), as in Section~\ref{sec:kato-rigidity}.
We have checked the polarization and the irreducibility, so
\[
 H^1_f(\mathbb Q,{\rm ad}\,r)=0.
\]
This applies after the type-I automorphic witness has been obtained.

The original route in \(\mathrm{rigIUE\_alt}\) factors the local
point through an integral valuation ring, uses the ordinary local
conditions, establishes the open-subgroup assertion
\(\mathrm{matIrr\_open\_gd}\), and obtains a free-ratio condition
from the automorphic twist. The subsequent result
\(\mathrm{rigEnd3\_alt}\) passes to the auxiliary CM field of the
frame, transports automorphy and irreducibility there, applies
adjoint lifting, and descends the vanishing. The propositions above
give the open-subgroup argument and a direct application of the
published theorem on the original quadratic field. The existence of
the required automorphic representation is an input to both versions.

\section{Local matching and the global theta integral}
\label{sec:theta-match}

This is the theorem about two unitary Weil representations used in
\texttt{theta\_\allowbreak{}integrals\_\allowbreak{}agree\_\allowbreak{}and\_\allowbreak{}local\_\allowbreak{}matching\_\allowbreak{}of\_\allowbreak{}literature}.
We give its local maps, rank restrictions, integral domains and
normalizations. The final subsection records which estimates and
representation-theoretic results enter the deductions below.

\subsection{Spaces, signatures and sections}

Write \(E/F\) for the CM extension with involution \(\sigma\).
The base of the theorem has \(E/\mathbb Q\) Galois,
\(\mathcal O_E/\mathcal O_F\) unramified and a standard additive
character \(\psi\). At every inert place the residue characteristic
is odd, the residue-field size exceeds three, and \(\psi_v\) is
unramified. The base also specifies split embeddings at a finite
set \(P\) of places, unramified for \(\psi\), and at a distinguished
split place \(v_1\). The two Hermitian spaces \(V,V'\) and the
skew-Hermitian space \(W\) have dimensions \(a,a',d=a+a'\) with
\(a'<a\) and \(d\) even. For the signature data in the source the
pairs are \((a-p,p),(q,a'-q),(d-p-q,p+q)\), respectively. Its
inequalities are
\[
 2(p+q)+1<(a-a')/2,\qquad 2(p+q)+1<a',\qquad
 d<m_\alpha(w)-m_\beta(w)\quad\text{at every real place }w .
\]
A pair of Hecke characters \(\alpha,\beta\), with the prescribed
infinite exponents and restrictions to the base ideles, determines
\(\nu=\alpha/\beta\). Write \(T=S_0\cup Q\subset P\), where \(S_0,Q\)
are disjoint. The characters \(\alpha,\beta\) are unramified outside
\(T\), and \(\nu\) is ramified at places of \(E\) over \(T\).
The distinguished place \(v_1\) lies outside \(T\) and splits
completely in the sense of the base. The Weil representations
\(\omega,\omega'\) use the pairs of splitting characters
\((\beta,\beta/\alpha)\), \((\alpha,\alpha/\beta)\), respectively.
Fix the models, polarizations, Haar measures and rational fundamental
domains specified with these data.

For \(Y=V\) or \(V'\) of dimension \(r\), and for a Schwartz
function \(\Phi\) in its adelic model, put
\[
 \Theta_Y(\Phi;g,h)=
   \sum_{x\in F^{r\times d}}[\omega_Y(g,h)\Phi](x),\qquad
 I_Y(\Phi)(h)=
   \int_{U(Y)(F)\backslash U(Y)(\mathbb A_F)}
         \Theta_Y(\Phi;g,h)\,dg .
\]
The second integral here uses the model's fundamental domain and measure.
The small-side subspace \(\mathcal S^\circ\) is spanned by pure tensors
that are vacuum vectors at almost all finite places. At \(v_1\), in
the specified split realization, their local factor vanishes on all
matrices of rank less than \(a'\). At infinity they have Schwartz
factors. Denote by \(\mathcal P\) the polynomial Fock subspace
in the corresponding archimedean model.

\subsection{The common quotients at finite places}

For every finite \(v\) there is a nonzero irreducible complex space
\(\Pi_v\) with maps \(p_v,p'_v\) onto it from the respective
Schwartz--Bruhat spaces. Both maps annihilate the linear span of
differences under the first member of the dual pair. Equality of
images is preserved by the action of the second member. At \(v\notin T\)
the two vacuum vectors have equal, nonzero images. For a good split
place \(v\notin T\) in the sense of the theorem (lying over a rational
prime split completely in \(E\)), the kernel on either Schwartz space
equals this span of differences. The small-side functions supported
on the full-rank locus in split coordinates surject onto \(\Pi_v\).
Every linear functional on the small Schwartz space arising from a
functional on the ambient functions and annihilating the differences
there factors through \(p'_v\) on Schwartz functions.

Here is the realization of the matching condition at \(v\in T\).
Let \(k=F_v\), choose the designated embedding \(E\longrightarrow k\),
and write \(\varpi\) for a uniformizer, \(\mathcal O\) for the
valuation ring, \(K_r=\operatorname{GL}_r(\mathcal O)\).
Use Haar measures of mass one on \(K_r\).
Use \(\alpha,\beta,\nu\) also for the local characters in these
coordinates. Put \(X_0=[1_a\ \ 0]\) and \(Y_0=[0\ \ 1_{a'}]\).
The argument of both transforms below is \(h\in\operatorname{GL}_d(k)\).
With \(h^{-t}=(h^{-1})^t\), they are
\begin{align*}
 P_A(\phi)(h)
   &=\beta(\det h)|\det h|^{-a/2}
     \mathop{\sum\nolimits^{\mathrm{fin}}}_{n_1\geq\cdots\geq n_a}
       \int_{K_a\operatorname{diag}(\varpi^{n_i})K_a}
       \phi(gX_0h^{-t})\nu(\det g)|\det g|^{d/2}\,dg,\\
 P_B(\phi')(h)
   &=\alpha(\det h)|\det h|^{a'/2}
       \int_{\operatorname{GL}_{a'}(k)}
       \phi'(gY_0h)\nu(\det g)|\det g|^{d/2}\,dg .
\end{align*}
The first expression uses the source's finite-support sum of shell
integrals indexed by decreasing integer tuples. The ramified
calculation establishes finite support on the Schwartz functions
to which the linearity argument is applied. The local Weil coordinate
changes \(\iota,\iota'\) send the two vacuum functions to indicators
of integral matrices of their respective sizes. The matching criterion is
\[
       p_v\Phi=p'_v\Phi'
        \quad\Longleftrightarrow\quad
       P_A(\iota\Phi)=P_B(\iota'\Phi')
           \quad\text{as functions on }\operatorname{GL}_d(k).
\]
There is no equality of vacuum images required at \(v\in T\).

To see how these transforms give the claimed quotient, use the
local change-of-variables and shell lemmas: both transforms are
linear on Schwartz--Bruhat functions, take values in the same
normalized parabolically induced representation, annihilate the
first group's differences, and intertwine the second group's action.
Each image contains a nonzero vector. For instance take a small
principal congruence subgroup on which \(\nu\circ\det\) is trivial.
At \(h=1\) localize the test function around the image of the
identity matrix under \(g\mapsto gX_0\) or \(g\mapsto gY_0\).
These inclusions have coordinate projections for left inverses,
so a compact open neighbourhood may be chosen with pullback
precisely the chosen subgroup. The support is then in the unit
shell, all factors in the integrand there are one if the
characteristic function is used, and its positive Haar volume is
the resulting value. This is the bump-function argument in
\texttt{p\allowbreak{}Ram\_\allowbreak{}ne\_\allowbreak{}zero} and its small-side counterpart.

The needed Bernstein irreducibility statement is the irreducibility
of the normalized product of two determinant representations of
\(\operatorname{GL}_a(k)\) and \(\operatorname{GL}_{a'}(k)\) for
smooth unitary characters when the dimensions or the characters
differ. Here \(a\ne a'\). An invariant subspace containing a
nonzero vector in this representation is the entire representation.
Both maps are therefore onto. Transport them through the split
Weil coordinates. Equality in their codomain, which is a subspace
of functions, is equivalent to pointwise equality of the transforms;
this proves the displayed ramified matching criterion.
In rank descent the source separately uses ramified multiplicity one:
for a nonzero equivariant map between the successive local models,
every other equivariant map of the same type is its scalar multiple.
It derives both local results using the geometric-lemma filtration.

For \(v\notin T\), use the split unramified maps when an embedding
\(E\longrightarrow F_v\) exists, and the inert spherical maps
otherwise. The inert input identifies the vacuum vectors in a common
spherical quotient. The source obtains it from its inert vacuum
comparison, spherical-function result and local representation
existence. These three cases exhaust the finite places. Choose the
pair of maps constructed at each \(v\); all the requirements on
them above are placewise. At infinity the two polynomial modules
surject onto a common nonzero irreducible \((\mathfrak g,K)\)-module
\(\Pi_\infty\). The maps kill first-group Lie derivatives and identify
first-group maximal-compact translates. They intertwine the
second group's Lie and maximal-compact actions on matching pairs.
The formal proof supplies this archimedean statement by its Fock model
and local lift comparison.

\subsection{The constant and the rank induction}

Two finite sums of pure tensors with polynomial archimedean factors
are polynomially matched when their images in
\(\Pi_\infty\otimes\bigotimes_{v\in U}\Pi_v\) are equal for a
finite set \(U\supset T\) outside which all their factors are
vacua. General matched data are finite sums with fixed finite
Schwartz factors whose archimedean factors are limits of sequences
of such polynomially matched sums. The small member and every one
of its approximants must lie in \(\mathcal S^\circ\).
For each choice of matching maps there is a unique \(c\ne0\) for which
\begin{equation}
       I_V(\Phi)(h)=c I_{V'}(\Phi')(h)
       \qquad\text{for all matched pairs and every }h.
       \label{eq:theta-identity}
\end{equation}

Here are the two rank changes controlling this constant. They also
specify which scalar factors may depend on local characters. Put
\(\delta=(a-a')/2\), and let \(\mathfrak b,\mathfrak c\) be the
fractional ideals of the two adapted planes for a step. For its
universal signs \(e_1,e_2,s\in\{1,-1\}\), set
\begin{align*}
 D(z)&=\nu(\mathfrak b)^{e_1}\nu(\mathfrak c)^{e_2}
       \prod_{v\in T}\nu_v(s)\epsilon_v(z),\\
 \epsilon_v(z)&=q_v^{-f_vz}\nu_v(\varpi_v)^{f_v}
      \sum_{y\in(\mathcal O_v/\varpi_v^{f_v})^\times}
        \nu_v(\widetilde y)^{-1}\psi_v(\widetilde y/\varpi_v^{f_v}).
\end{align*}
Here \(f_v\) is the conductor exponent, \(q_v\) the residue cardinality
and the sum uses lifts of the unit classes. The ideal character
values have the source's uniformizer conventions. At a beta step,
\(q>0\), the ranks change from \((a,a')\) to \((a,a'-2)\)
and \(q\) drops by one. With \(c_1\) the lower-rank constant,
\[
        c\,L^{S_0}(\nu,\delta+1)
          =[C D(-\delta)]^{-1}c_1 .
\]
At an alpha step, \(q=0,\ p>0\), the ranks change to
\((a-2,a')\), and \(p\) drops by one. In that case
\[
        c=C D(1-\delta)L^{S_0}(\nu,\delta)c_1 .
\]
Here \(L^{S_0}\) is the partial character \(L\)-function with the
specified \(S_0\) Euler factors omitted. Each step has its own
universal triple of signs. Its \(C\ne0\) is chosen after the planes,
geometry and archimedean maps, before the characters and finite
matching data.

To prove the deduction from the rank-step inputs, induct on \(p+q\).
The \(p=q=0\) case is the anisotropic polynomial identity. If \(q>0\)
use the beta step; if \(q=0\) and \(p>0\) use the alpha step.
The source's existence theorem for adapted planes supplies a
lower-rank admissible triple. Local matching supplies its models,
and the induction hypothesis supplies \(c_1\).
The comparisons of the constant terms for the zero orbit and
the relevant nonzero cone give the displayed step relations and
the polynomially matched identity for the larger pair. The
nonvanishing factor result gives \(D(z)\ne0\) at the specified
arguments and nonvanishing of the partial \(L\)-value.
Hence the new scalar is nonzero.

There is also a small-side test vector with nonzero integral.
When \(q=0\) this is the anisotropic nonvanishing conclusion.
For \(q>0\), the nonzero term in beta descent supplies a big-side
vector with nonzero integral. Choose its polynomially matched
small-side partner. Its integral must be nonzero by the polynomial
identity. Evaluating at an \(h\) where this integral is nonzero
forces any two admissible constants to be equal. To extend to
general matched data, pass to the archimedean approximants in
the definition above. The big-side integral is continuous on
each fixed finite Schwartz factor by its integrable majorant;
the small-side one is continuous along the indicated sequences
in \(\mathcal S^\circ\). There are only finitely many fixed
factors. Taking both limits in the polynomial identity proves
\eqref{eq:theta-identity}.

\subsection{Indefinite convergence and regularized representatives}

The large-side integral converges absolutely on every adelic
Schwartz function in the stated range. On the small side take
\(\Phi'\in\mathcal S^\circ\). Put \(H_Q(g)=\inf_{\gamma\in U(V')(F)}
H(\gamma g)\), where \(H\geq1\) is the matrix height of the model.
For each integer \(N\geq0\) there are \(C,k\) such that
\[
 |\Theta_{V'}(\Phi';g,h)|H_Q(g)^N\leq C H(h)^k .
\]
Moreover \(\int|\Theta_{V'}(\Phi';g,h)|\,dg\) is finite and bounded
by a fixed power of \(H(h)\) times a constant. Here the input from
unitary reduction theory consists of an integrable majorant for
each pure tensor in the full-rank subspace and a seminorm bound
with each power \(N\). The finite-sum deduction is explicit.
Write \(\Phi'=\sum_{i=1}^m f_i\otimes\phi_i\) in the defining span.
For its \(i\)-th summand use the integral estimate
\(\int|\Theta_i(g,h)|\,dg\leq C_i H(h)^{k_i}\) supplied by that
majorant. Form the finite sum, take \(k=\max_i k_i\), and use
\(\int|\sum_i\Theta_i|\leq\sum_i\int|\Theta_i|\).
The zero sum is immediate. The same argument applied pointwise to
the seminorm bounds gives the displayed rapid decrease. This also
shows directly why the integral is a moderate function of \(h\).

The subsequent twisted square comparison gives a regularized
representative for each small-side algebraic datum \(\Phi'\).
In its model let \(\mathcal R\) be the span of
\(\omega'(u,1)\Psi-\Psi\), with \(u\) in the allowed finite-adelic
subgroup and \(\Psi\) algebraic. There is \(\Phi_0\) in the same
algebraic space with \(\Phi'-\Phi_0\in\mathcal R\) and with a
convergent moderate theta integral. For any matched big-side
\(\Phi\) with \(p\Phi=p'\Phi_0\), the twisted analogue of
\eqref{eq:theta-identity} holds for \(\Phi,\Phi_0\).
The following argument from the archive shows precisely when the
integral of \(\Phi'\) itself can take its place.

Suppose \(\Psi\in\mathcal R\) and its kernel is integrable as a
function of \(g\) for every \(h\). Write
\(\Psi=\sum_i b_i(\omega'(u_i,1)\Psi_i-\Psi_i)\).
Choose a regularizer \(Z=\sum_k z_k\omega'(k,1)\), a finite sum in
the allowed subgroup with \(\deg Z=\sum_kz_k\ne0\), for this finite
list. The source constructs it at a completely split auxiliary place
outside the exceptional sets for both the kernels and the finite
translates. Its regularized vectors \(Z\Psi_i\) have convergent
moderate integrals, and for each \(i\) there is \(u'_i\) in the
same subgroup satisfying
\(Z\omega'(u_i,1)\Psi_i=\omega'(u'_i,1)Z\Psi_i\).
Right translation on the adelic quotient preserves the integral
of each integrable kernel. Thus
\[
 I_{V'}(Z\Psi)=
   \sum_i b_i\{I_{V'}(\omega'(u'_i,1)Z\Psi_i)-I_{V'}(Z\Psi_i)\}=0 .
\]
Since the kernel of \(\Psi\) itself is integrable, right translation
and the finite expansion of \(Z\) also give
\(I_{V'}(Z\Psi)=(\deg Z)I_{V'}(\Psi)\).
Division by the nonzero degree proves \(I_{V'}(\Psi)=0\).
If the original \(\Phi'\) has integrable kernel, apply this to
\(\Psi=\Phi'-\Phi_0\). Since the matching maps annihilate
\(\mathcal R\), \(p'\Phi'=p'\Phi_0\); since
\(I_{V'}(\Phi')=I_{V'}(\Phi_0)\), the scalar identity holds with
the ordinary integral of \(\Phi'\). For a general algebraic
\(\Phi'\) the comparison uses the regularized representative.

\paragraph{Inputs in the release.}
At the top edge \texttt{theorem\allowbreak{}K\_\allowbreak{}head2} takes two propositions.
\texttt{Eis\allowbreak{}BL} gives the half-plane of convergence of a doubled
Eisenstein family and weak holomorphic continuation after a finite
local pole-clearing factor. \texttt{Siegel\allowbreak{}Weil\allowbreak{}Std} gives, for a
definite first member of dimension \(m\) and second member of
dimension \(n\), the holomorphic value of the specified Eisenstein
section at \((m-n)/2\). It is the average of the theta sum over the
first group multiplied by \(1\) if \(n<m\) and \(2\) otherwise.
These two propositions are supplied by \texttt{Rg\allowbreak{}Eis.\allowbreak{}Hd.\allowbreak{}eis\allowbreak{}BL}
and \texttt{Siegel\allowbreak{}Weil\allowbreak{}Formula.\allowbreak{}Fin.\allowbreak{}siegel\allowbreak{}Weil\allowbreak{}Std}. Neither supplier
uses the nominal literature-bundle parameter in its proof.

The fourteen arguments in the source's next assembly
\texttt{Tr.\allowbreak{}close\allowbreak{}Sph} are the finite coinvariant dictionary,
the full-rank small-space coinvariant theorem, ramified
multiplicity one, the ramified finite matching maps, split
unramified finite matching maps, the spherical triple comparison,
inert finite matching maps, the second-group real Lie derivative
comparison, archimedean Lie equivariance, the common real quotient,
the unitary Siegel-set theorem, the Eisenstein continuation theorem,
the definite Siegel--Weil value, and existence of the requisite
Weil representations. The spherical input uses
\texttt{mod\allowbreak{}St2}; the ramified and split finite arguments use the
geometric-lemma filtration; the real quotient uses the local lift
and infinitesimal action. Inside \texttt{close\allowbreak{}Sph}, the
nontrivial analytic and algebraic steps cited above include
\texttt{pol\allowbreak{}Zero}, \texttt{non\allowbreak{}Zero\allowbreak{}An}, \texttt{end\allowbreak{}AB\_\allowbreak{}sph},
\texttt{loc\allowbreak{}A\_\allowbreak{}of}, \texttt{loc\allowbreak{}B\_\allowbreak{}of}, the plane towers,
\texttt{dom\allowbreak{}Weil\allowbreak{}U}, \texttt{so\allowbreak{}Semi\allowbreak{}U} and the unitary Siegel-set
estimates. Their underlying proofs are imported in this chapter.
The correspondence with the definitions is
\texttt{Weil\allowbreak{}Tensor\allowbreak{}Models.\allowbreak{}Is\allowbreak{}Matched}, \texttt{Decay},
\texttt{One\allowbreak{}Const}, \texttt{Coinv}, \texttt{Step\allowbreak{}Beta} and
\texttt{Step\allowbreak{}Alpha}.

The additional application to regularized representatives is
\texttt{twisted\_\allowbreak{}square\_\allowbreak{}theta\_\allowbreak{}integrals\_\allowbreak{}eq}; its inputs are
the two tensor-model theorems and
\texttt{Global\allowbreak{}Theory.\allowbreak{}Fun\allowbreak{}Pr.\allowbreak{}Theta\allowbreak{}Local\allowbreak{}Inputs}. The removal of the
regularizer for an already integrable kernel uses
\texttt{twist\_\allowbreak{}regulariser\_\allowbreak{}is\allowbreak{}Good},
\texttt{theta\allowbreak{}Integral\_\allowbreak{}eq\_\allowbreak{}zero\_\allowbreak{}of\_\allowbreak{}finite\allowbreak{}Relation} and
\texttt{theta\allowbreak{}Integrals\_\allowbreak{}agree\_\allowbreak{}of\_\allowbreak{}integrable\_\allowbreak{}kernel}.
The auxiliary-place construction uses placewise seesaw models,
the full-rank regularization theorem, commutation of separated
finite-place operators and existence of a completely split place
outside a prescribed finite set. These are source results used by
the argument in this subsection.

\section{The doubled Eisenstein series and its definite theta value}
\label{sec:sw-boundary}

Use the base \(B\), the CM extension \(E/F\), its involution \(\sigma\),
and its standard additive character \(\psi\) from
Section~\ref{sec:theta-match}. Fix \(m,n\geq1\), spaces \(A,Y\)
of dimensions \(m,n\), and their signature data \((r,s),(r',s')\).
The hypotheses allow either hermitian \(A\) and skew-hermitian \(Y\),
or skew-hermitian \(A\) and hermitian \(Y\). In both cases
\(\min(r,s)=0\). Write \(H=U(A)\), \(G=U(Y\oplus(-Y))\).
Its Siegel parabolic \(P\) fixes the diagonal maximal isotropic
subspace. Fix the prescribed unitary Hecke character \(\chi\) with
restriction \(\eta_{E/F}^{\,m}\) on the base ideles, a doubled
Darboux model with actions \(\omega\) of \(G(\mathbb A_F)\) and
\(S\) of \(H(\mathbb A_F)\), and the model's distinguished linear
functional \(\ell\). Choose a standard rational doubled frame \(b\).
In it the Gram matrix is
\(\left(\begin{smallmatrix}0&I\\ \epsilon I&0\end{smallmatrix}\right)\)
with \(\epsilon=\pm1\), and its first columns span the diagonal.
The associated subgroup \(K\) is unitary for the positive standard
inner products at infinity, with both the matrix and its inverse
integral in the finite frame. Choose the Iwasawa height \(a\) with
\(a(pk)=|\det t|_{\mathbb A_E}\) for the Levi action \(t\) of \(p\)
on the diagonal.

Write \(D\) for a fundamental domain of \(H(F)\) in
\(H(\mathbb A_F)\) with respect to any Haar measure \(\mu\).
Let \(\Phi\) be adelic Schwartz in the model such that
\(\operatorname{span}_{\mathbb C}\{\omega(k)\Phi:k\in K\}\)
is finite-dimensional. Put
\begin{align*}
 s_0&=(m-n)/2,\qquad f_0(g)=\ell(\omega(g)\Phi),\\
 f_z(g)&=a(g)^{z-s_0}f_0(g),\\
 I_\Phi(x)&=\mu(D)^{-1}\int_D
     \Theta_F\bigl(\omega(x)S(h)\Phi\bigr)\,d\mu(h).
\end{align*}
Here \(\Theta_F\) sums a test function over rational points in its
Darboux coordinates. For every such choice and every \(x\),
the Eisenstein series
\[
 E_\Phi(z,x)=\sum_{\gamma\in P(F)\backslash G(F)} f_z(\gamma x)
\]
has a meromorphic extension holomorphic at \(s_0\), with value
\begin{equation}
 E_\Phi(s_0,x)=
 \begin{cases} I_\Phi(x),& n<m,\\ 2I_\Phi(x),&m\leq n.\end{cases}
 \label{eq:sw-def-value}
\end{equation}
The measure in this identity is an average. Scaling \(\mu\) by a
positive constant cancels between its integral and its denominator.
The arguments below describe the continuation, an adelic Fourier
normalization, and the two terms in the boundary calculation leading
to the change in rank.

\subsection{Removing the global multiplier locally}

For any section \(f\) at parameter \(s_1\), write
\(f_z(g)=a(g)^{z-s_1}f(g)\). The assertion here allows any
nonnegative integer \(c\) with the prescribed restriction
\(\chi|_{\mathbb A_F^\times}=\eta_{E/F}^{\,c}\). Let \(\mathcal M\)
be the complex span of continuous functions of growth bounded by a
power of the matrix height \(h_2\) of \(G\). A linear functional
\(L\) on \(\mathcal M\) satisfies the growth condition if for every
\(N\geq0\) there is \(C_N\) bounding \(|L(v)|\) whenever
\(|v(g)|\leq h_2(g)^N\) for every \(g\).

Two analytic inputs enter the proof of the continuation.
First, in a right half-plane \(\Re z>t_0\), the coset series has
sum \(v(z)(x)\), with \(v(z)\in\mathcal M\) and \(L(v(z))\)
holomorphic for every such \(L\). Second, there are an entire
function \(A_0\) nonzero at some point, a function \(J(z)\) in
\(\mathcal M\) with \(L(J(z))\) entire, and \(t_1\in\mathbb R\)
such that on \(\Re z>t_1\) the series obtained by multiplying
each coset summand by \(A_0(z)\) sums to \(J(z)(x)\). These are
the half-plane and multiplier theorems in the source. Its
multiplier construction uses the rank-one case and induction on
the cardinality of the indexing set of a general linear block.

Let \(t=\max(t_0,t_1)\). Uniqueness of the sum on \(\Re z>t\)
gives \(J(z)(x)=A_0(z)v(z)(x)\) for all \(x\), including zeros
of \(A_0\). Use \(v(z)\) there and \(J(z)/A_0(z)\) elsewhere
away from the zeros. Values at remaining isolated zeros can be
assigned arbitrarily in taking a meromorphic representative.
Point evaluation at \(x\) satisfies the growth condition with
\(C_N=h_2(x)^N\), proving pointwise holomorphy in the right
half-plane. Now fix any \(z_*\). Since \(A_0\) is a nonzero
entire function, its power series at \(z_*\) has a first
nonzero coefficient. Factoring out the corresponding power gives
\(A_0(z)=(z-z_*)^k u(z)\) on an open neighbourhood \(U\) on
which \(u\) is holomorphic and nonvanishing. Set
\(J_*(z)=u(z)^{-1}J(z)\) there. For every \(x\) and every
\(z\in U\setminus\{z_*\}\),
\[
                  J_*(z)(x)=(z-z_*)^k E(z,x).
\]
For each allowed \(L\), the equality
\(L(J_*(z))=u(z)^{-1}L(J(z))\) proves holomorphy on \(U\).
The order \(k\) and the neighbourhood work simultaneously for
all \(x\) and all the functionals. This proves the weak
continuation claimed in \texttt{EisBL} from the two analytic
inputs.

\subsection{The rational Fourier pairing and its normalization}

Let \(V=\operatorname{Mat}_{m,n}(E)\) and choose the two elements
\(e_1,e_2\) of the \(F\)-basis of \(E\) in the datum. Use
\(Q=\{1,\ldots,m\}\times\{1,\ldots,n\}\times\{1,2\}\).
Write \(D_q\) for the elementary matrix with entry \(e_k\) in
position \((i,j)\), where \(q=(i,j,k)\). The coordinate isomorphism
\(C\colon V\to F^Q\) expands a matrix in the \(D_q\). For the
Gram matrix \(J_A\) define an \(F\)-bilinear pairing
\[
 B_A(Y,Z)=\operatorname{Tr}_{E/F}
              \operatorname{tr}(Y^t J_A\,\sigma Z).
\]
For any \(\varepsilon=\pm1\) the rational linear map
\(L_\varepsilon\colon F^Q\to F^Q\) given by
\((L_\varepsilon C Y)_q=\varepsilon B_A(Y,D_q)\) is bijective.
Indeed, vanishing of all its coordinates says
\[
 \operatorname{Tr}_{E/F}\bigl((Y^tJ_A)_{ji}\sigma(e_k)\bigr)=0
                         \qquad(i,j,k).
\]
The two \(\sigma(e_k)\) also span \(E\) over \(F\).
Nondegeneracy of the trace pairing for the separable extension
therefore gives \((Y^tJ_A)_{ji}=0\) for every \(j,i\).
Invertibility of \(J_A\) gives \(Y=0\). The source and target are
the same finite-dimensional vector space, proving bijectivity.

Extend \(C\) adelically. Its inverse sends \(c\) to
\(\sum_q c_q D_q\), so it identifies the diagonal rational
lattice \(F^Q\) with \(V\). Normalize additive Haar measure
\(\nu\) on \(V(\mathbb A_F)\) by rational covolume one.
The pullback of the product adelic measure on \(\mathbb A_F^Q\)
also has covolume one. Indeed, the inverse image under \(C\) of
a fundamental domain is a fundamental domain by the rational
lattice bijection, and its measure for the pullback is one.
Uniqueness of Haar measure up to scalar now identifies the two
measures. For adelic \(Z\) and rational \(Y\), the scalar
\(\varepsilon B_A(Y,Z)\) is \(\sum_q (CZ)_q(L_\varepsilon C Y)_q\)
under the diagonal embedding of the rational factors.

Consequently, for every adelic Schwartz function \(h\) in this
coordinate model,
\begin{equation}
 \sum_{Y\in V}\int_{V(\mathbb A_F)}
        \psi(\varepsilon B_A(Y,Z))h(Z)\,d\nu(Z)
       =\sum_{Y\in V}h(Y).
 \label{eq:sw-poisson}
\end{equation}
Make the change of variable \(c=CZ\) in each integral. With the
indicated sign convention the result is the adelic Fourier
transform of \(h\circ C^{-1}\) at \(L_\varepsilon C Y\).
The frequencies run over \(F^Q\) exactly once, by the preceding
bijection. Apply adelic Poisson summation on \(\mathbb A_F^Q\)
with standard character and rational covolume one, then reindex the
result by \(C^{-1}\). This proves \eqref{eq:sw-poisson} and explains
why this use of a Weyl element introduces no scalar of local origin.
The general coordinate Poisson formula is an input here.

\subsection{Splitting a boundary coefficient}

Take a lower-rank datum that omits index \(i_0\) in the \(n\)-dimensional
frame, with embedding \(j\colon G'\to G\) of the doubled group of
rank \(n-1\). Write \(N,N'\) for the additive Siegel unipotent
coordinates and \(n(\beta),n'(\beta')\) for their embeddings as
upper block unipotents. Let \(\pi\beta\) be the submatrix left
after removing the \(i_0\) row and column. Let \(T\in N(F)\) have
zero row and column \(i_0\), and set \(T'=\pi T\). For a function
\(B_0\) and a rational frequency \(T\) write
\[
 \mathcal F_T B_0(x)
 =\int_{N(F)\backslash N(\mathbb A_F)}
      B_0(n(\beta)x)\,
      \psi\bigl(-\operatorname{Tr}_{E/F}\operatorname{tr}(T\beta)\bigr)
      \,d\beta ,
\]
with normalized additive Haar measure. Here \(\operatorname{Tr}_{E/F}\)
also denotes the trace after adelic base change. In the equations
below the source defines this integral on a chosen fundamental
domain; its invariance under changes of domain follows from
rational periodicity. Let \(u_M(u)\) have upper Levi block
\(I+\sum_{j\ne i_0}u_j E_{i_0 j}\); the variable \(u\) runs
over an additive fundamental domain for \(E^{n-1}\) in
\(\mathbb A_E^{n-1}\), again of volume one.

For \(f_z=a^{z-s_1}f\) as above let \(E_n(z,x)\) be its coset
sum and, at the fixed \(x\), put
\[
 E'_{x}(z,g')=\sum_{q'\in P'(F)\backslash G'(F)}
                            f_z(j(q'g')x).
\]
The boundary identity in the release holds, for all \(x\), in a
common right half-plane depending on the model and section:
\begin{align}
 \int \mathcal F_T E_n(z,u_M(u)x)\,du
    &=\mathcal F_{T'} E'_{x}(z,1)+C_2(z,x),\label{eq:sw-cells}\\
 C_2(z,d_w(\tau)x)&=\exp\!\bigl(\tau(n-2z)\bigr)C_2(z,x).
 \label{eq:sw-dil}
\end{align}
The second equation holds for every real place \(w\), real
\(\tau\), and any model dilation there along the omitted
coordinate with that parameter.

To justify interchanging the coset sum with both integrals, take
fundamental domains with compact closures in the two additive
quotients. Such domains exist by cocompactness of the rational
lattices. The product \(C\) of the closures has compact image
under \((u,\beta)\mapsto n(\beta)u_M(u)x\). The right-half-plane
majorant for the section on that compact image supplies a summable
family \(b_q\) bounding the absolute value of its \(q\)-th summand,
uniformly there. The character factor is continuous on the compact
closure, with some bound \(M\). For the summand \(k_q\), including
its character, the product probability measure gives
\[
   \sum_q\iint |k_q(u,\beta)|\,d\beta\,du
                  \leq M\sum_q |b_q|<\infty .
\]
Continuity of each summand supplies measurability. Fubini and absolute
summability permit taking both integrals term by term. Rational
invariance of the Eisenstein series, the trace relation for \(T,T'\)
and rational invariance of the character imply independence of
the fundamental domains. Equivalently, these are integrals of
periodic functions for invariant measures on the additive quotients.

In the doubled frame partition the cosets by asking whether the
\(i_0\)-th column of the lower-left block of a representative
vanishes. For the vanishing class, the source's rational parabolic
orbit bijection gives a representative \(p_0j(\gamma')\) for each
lower-rank coset \(\gamma'\), with \(p_0\in P(F)\). The block
factorization used here asserts
\[
 j(h)n(\beta)u_M(u)x
  =p_1j(hn'(\pi\beta))x,\qquad
             \det(\operatorname{Levi}(p_1))=1 .
\]
The rational factor \(p_0\) and determinant-one parabolic \(p_1\)
leave \(f_z\) unchanged under left multiplication, by its
induction law. Expanding the trace as a double sum, every term
involving \(i_0\) vanishes, hence
\(\operatorname{tr}(T\beta)=\operatorname{tr}(T'\pi\beta)\).
Thus the corresponding summand is independent of \(u\) and
depends on \(\beta\) only through \(\pi\beta\). The \(u\)-integral
contributes one. Summing the lower-rank summands gives a continuous
rational-periodic function on \(N'(\mathbb A_F)\). Integration after
projection gives its normalized integral on the smaller quotient
(the fiber integration lemma in the source). This proves that the
first class gives precisely \(\mathcal F_{T'}E'_x(z,1)\).
The absolutely convergent sum over the complementary class defines
\(C_2(z,x)\), proving \eqref{eq:sw-cells}.

Here is the computation of the second class under dilation. Its
unfolding supplies a Weyl representative \(w_1\), coordinates
\(c\in\mathbb A_E^{n-1}\times\mathbb A_F\) with measure \(\kappa\),
and integrals of the form
\[
 \int f_z\bigl(w_1 c_M(c)j(h)x\bigr)\,d\kappa(c).
\]
The entire expression is a fixed scalar times a sum over lower
cosets of Fourier integrals of these inner integrals. The
coordinate identities used by the unfolding have three parts.
Writing \(d_\tau=d_w(\tau)\), the image \(j(h)\) commutes with
\(d_\tau\). There is \(p\in P(\mathbb A_F)\) such that
\(w_1d_\tau=pw_1\), with Levi determinant \(\Delta\) satisfying
\(|\Delta|_{\mathbb A_E}=e^{-2\tau}\) and \(\chi(\Delta)=1\).
Finally there is a measurable equivalence \(c\mapsto s_\tau(c)\)
with
\[
 c_M(c)d_\tau=d_\tau c_M(s_\tau(c)),\qquad
               (s_\tau)_*\kappa=e^{2n\tau}\kappa .
\]
These are the specific imported coordinate and Jacobian formulas
in the unfolded cell. The transformation law of a section at \(s_1\)
is \(f(pg)=\chi(\Delta)|\Delta|_{\mathbb A_E}^{s_1+n/2}f(g)\).
It follows, with positive real bases for all complex powers, that
\[
 f_z(pg)=
 e^{-2\tau(z-s_1)} e^{-2\tau(s_1+n/2)} f_z(g)
          =e^{-\tau(2z+n)}f_z(g).
\]
Insert \(d_\tau x\) in the inner integral. Move \(d_\tau\) to the
left past \(j(h)\) and \(c_M(c)\), apply the parabolic law to
\(w_1d_\tau=pw_1\), then push \(\kappa\) forward by \(s_\tau\).
The scalar is \(e^{-\tau(2z+n)}e^{2n\tau}=e^{\tau(n-2z)}\).
It is independent of the lower coset and its Fourier variable and
so passes unchanged through the remaining integral and sum. This
proves \eqref{eq:sw-dil}. The bounds, first orbit identification
and second orbit unfolding may have three different convergence
thresholds. Taking their maximum yields the single half-plane in
\eqref{eq:sw-cells}--\eqref{eq:sw-dil}.

\subsection{Recovering the value across the rank ranges}

We record how \eqref{eq:sw-cells} enters the induction. For
\(m=n\) the source proves a nondegenerate Whittaker integral
identity with scalar \(2\). For \(m>n\) the corresponding scalar
is \(1\). At a nondegenerate frequency the Eisenstein coefficient
unfolds to the integral over the full Siegel unipotent of the
section on the big Weyl cell against its Fourier character.
The theta coefficient is the orbital functional on the same
frequency. Their comparison in the source combines the finite
local integrals, unramified integrals, real-place Whittaker
integrals, global factorization and a representation-density
formula. The comparison fixes the scalar at nondegenerate
frequencies. To lift the equality to \emph{all} coefficients the
source separately uses the bounds for coefficient functionals,
a vanishing result for singular tests and the Fock-space
representation argument. These results are inputs to this chapter.

For clarity, consider passage from a single polynomial Fock lift
to a general \(K\)-finite datum. Fix a real place and write the
transformed \(\Phi\) as a finite sum of polynomial Gaussian factors
there times away-place factors. Lift the \(i\)-th summand to
\(\Phi_i\) in the original model. Injectivity of the model
transform on Schwartz functions gives
\(\Phi=\sum_{i=1}^k\Phi_i\), since the difference has transform
zero. Section formation and the theta integral are linear on
these finite sums. Suppose each \(E_{\Phi_i}\) has holomorphic
value \(cI_{\Phi_i}\) at \(s_0\), as supplied by the
polynomial-lift result. Choose a convergence threshold \(t_i\)
for the coset sum at the chosen \(x\). On
\(\Re z>\sum_i|t_i|\) every term converges absolutely, hence
\(E_\Phi=\sum_i E_{\Phi_i}\) there. The finite sum of the
continuations is holomorphic at \(s_0\) and gives
\(c\sum_i I_{\Phi_i}=cI_\Phi\). This is also valid for the empty
sum. Meromorphic uniqueness identifies this germ with the
continuation constructed above.

For \(m<n\) every Gram matrix represented by \(m\) rows has rank
at most \(m\). Full-rank theta coefficients therefore vanish.
The source's lower-rank step uses the boundary decomposition
\eqref{eq:sw-cells}, the covariance \eqref{eq:sw-dil}, the
restriction of the vacuum to the zero \(i_0\)-th column, and
the corresponding identity for singular orbital integrals.
The step also uses holomorphicity of the leading coefficient in
the real-group variable and the singular-test vanishing theorem.
Its conclusion transfers the value with scalar \(2\) from all
lower-rank vacuum vectors to all vacuum vectors of the given rank.
In the model a vacuum here is a finite sum of products of
archimedean Gaussians with finite Schwartz--Bruhat factors.
The induction is on \(k=n-m\). At \(k=0\) the equal-rank value
already proved applies to the vacua, which belong to the polynomial
Fock space and are \(K\)-finite. If \(k>0\), then \(n\geq2\); the
lower-rank existence theorem supplies the datum omitting \(i_0\).
Its \(m\), first space, base, character, Haar measure and
fundamental domain are unchanged. Its \(n'\) equals \(n-1\)
because its frame indices inject onto exactly the complement of
\(i_0\). Thus \(n'-m=k-1\). Apply induction to that datum and the
stated lower-rank step to the original datum. The vacuum generation
and differentiation theorem then gives the identity for all
\(K\)-finite vectors. This proves the scalar and rank bookkeeping
in \eqref{eq:sw-def-value}, conditional in this presentation on
the local integrals, singular-coefficient theorem and generation
results just specified.

The classical comparison is Shunsuke Yamana, \emph{The Siegel--Weil
formula for unitary groups}, Pacific Journal of Mathematics
\textbf{264} (2013), Theorem~2, p.~237,
\url{https://doi.org/10.2140/pjm.2013.264.235}.
Specialize his central division algebra to \(E\), so its degree is
one. His character on base ideles is the \(m\)-th power of the
quadratic character, the group on the doubled side is split in
the sense of admitting a complete polarization, and \(s_0=(m-n)/2\).
If \(r_0\) is the rational Witt index of the first space, his
convergence criterion includes \(r_0=0\), and his value theorem
treats \(m\leq n\) and \(m-r_0>n\). Definiteness at an
archimedean place forces \(r_0=0\), so these ranges cover the two
cases here. The source additionally transports between its Darboux
functional and the standard Schr\"odinger model and uses its own
normalized Haar average.

\subsection{Transport to the stabilizer used in the theta identity}

The rank identity above was stated for the standard compact subgroup
of a rational doubled frame. The theta integrals of
Section~\ref{sec:theta-match} are allowed an adapted stabilizer
\(K_1\). The frame-conjugacy lemma provides a parabolic element
\(p\) with Levi action \(t_p\), and standard subgroup
\(K_0=p^{-1}K_1p\). Suppose \(a_1\) is the height for \(K_1\).
Put \(d=|\det t_p|_{\mathbb A_E}>0\) and
\[
                 a_0(g)=d\,a_1(gp^{-1}).
\]
This is the height for \(K_0\). Indeed, if \(k\in K_0\), then
\(pkp^{-1}\in K_1\). For \(q\in P(\mathbb A_F)\) with Levi
action \(t_q\), write \(qkp^{-1}=(qp^{-1})(pkp^{-1})\).
The defining rule for \(a_1\) gives
\(a_1(qkp^{-1})=|\det(t_qt_p^{-1})|_{\mathbb A_E}\).
Multiplying by \(d\) gives \(|\det t_q|_{\mathbb A_E}\).
Every element admits the required decomposition by transporting
the Iwasawa decomposition for \(K_1\).

Take \(\Phi_0=\omega(p^{-1})\Phi\). If \(N\) is a
finite-dimensional \(K_1\)-invariant span containing \(\Phi\),
then \(\omega(p^{-1})N\) contains the \(K_0\)-span of
\(\Phi_0\), since for \(k\in K_0\),
\(\omega(k)\Phi_0=\omega(p^{-1})\omega(pkp^{-1})\Phi\).
Apply \eqref{eq:sw-def-value} for the standard data at \(xp\).
For every rational coset representative \(\gamma\) the section
there satisfies
\[
 a_0(\gamma xp)^{z-s_0}
  \ell\bigl(\omega(\gamma xp)\Phi_0\bigr)
 =d^{z-s_0}a_1(\gamma x)^{z-s_0}
                   \ell\bigl(\omega(\gamma x)\Phi\bigr).
\]
The factor \(d^{s_0-z}=\exp((s_0-z)\log d)\) is entire and
equals one at \(s_0\). Multiplying the continued standard
Eisenstein series by this factor thus gives the continued
adapted series with the same special value.
On the theta side the parabolic commutation lemma gives
\(\omega(p^{-1})S(h)\Phi=S(h)\omega(p^{-1})\Phi\) on Schwartz
functions. Hence
\(\omega(xp)S(h)\Phi_0=\omega(x)S(h)\Phi\), also giving exactly
the adapted theta average.

The input to that commutation lemma is injectivity of all the
translated theta functionals. It has a direct Fourier proof.
In Schr\"odinger coordinates, suppose
\(\Theta_F(\rho(x,y)\phi)=0\) for all \(x,y\in\mathbb A_F^\kappa\),
where \(\kappa\) is the finite coordinate set. For fixed \(x\)
divide out the nonzero scalar \(\psi(\langle x,y\rangle/2)\)
in the Schr\"odinger action. Then
\[
    \sum_{\xi\in F^\kappa}\phi(\xi+x)
                    \psi(\langle\xi,y\rangle)=0
                \qquad(y\in\mathbb A_F^\kappa).
\]
The coefficients are absolutely summable by the rational lattice
bound for the Schwartz translate. Every nonzero \(\xi\) gives a
nontrivial character on the additive quotient. Indeed choose
\(j\) with \(\xi_j\ne0\) and put in coordinate \(j\) an adele
\(a/\xi_j\) with \(\psi(a)\ne1\), using zero in the other
coordinates. Averaging over the quotient annihilates every such
character by translation invariance. Termwise averaging is
permitted by absolute summability. The zero coefficient is
\(\phi(x)\), so \(\phi=0\). In the doubled model the theta
expansion and the Heisenberg composition law transfer this
injectivity to translates of \(\ell\). The covariance of the two
actions then gives equal values of all these functionals on
\(\omega(p)S(h)\Phi\) and \(S(h)\omega(p)\Phi\), hence equality
of the vectors. This is the argument of \texttt{theta\_inj},
\texttt{dbl\_inj}, \texttt{dbl\_comm} and \texttt{sw\_transl}.

\paragraph{Source correspondence.}
The two top declarations are
\srcid{RgEis.Hd.eisBL} and
\srcid{SiegelWeilFormula.Fin.siegelWeilStd}. The local removal
of the multiplier is \texttt{BRoad.eisBL\_of}. The rational
pairing and Haar arguments are \texttt{rp\_nondeg},
\texttt{integral\_cov\_gen}, \texttt{int\_eq\_four} and
\texttt{TW.Twp.poisson}. The boundary calculation uses
\texttt{Act.absConv}, \texttt{C4.firstEq\_of},
\texttt{Act.secLaw\_of}, \texttt{Act.lhs\_eq\_tsum} and
\texttt{Act.ctSplit\_of\_parts}. The first-cell block
factorization, normalized projection of Haar measure, second-cell
parametrization and its Jacobian occur upstream in the proofs of
\texttt{key\_id}, \texttt{fiber}, \texttt{secUnf\_of} and
\texttt{cLawEx}. The rank assembly is
\texttt{Atop.mainC\_of}, \texttt{rank\_difference\_induction}
and \texttt{Atop.fieldOf}. The analytic and arithmetic inputs
identified above remain imported at this stage of the prose
translation.

\section{Singular Fourier tests for the archimedean vacuum}
\label{sec:sw-singular-vacuum}

Suppose \(0\le m<n\). In the definite Siegel--Weil calculation of
Section~\ref{sec:sw-boundary}, the local vacuum at the special
parameter has Fourier frequencies of rank at most \(m\). We prove
this assertion after arbitrary left and right translations and
after averaging by a compact group. The singularities here concern
the rank of the Fourier frequency.

Use \(J=\left(\begin{smallmatrix}0&I\\-I&0\end{smallmatrix}\right)\)
and \(G=\{g\in\operatorname{GL}_{2n}(\mathbb C):g^*Jg=J\}\).
Write \(A_g,B_g,C_g,D_g\) for its four \(n\)-by-\(n\) blocks, and set
\[
 U_g=iC_g+D_g,\quad V_g=iA_g+B_g,\quad
 n(S)=\begin{pmatrix}I&S\\0&I\end{pmatrix},\qquad
 f(g)=\det(g)^e\det(U_g)^{-m},\qquad e\in\mathbb Z.
\]
Here \(f\) is \texttt{fz} at parameter zero, the local parameter
measured relative to the distinguished global value. Identify
\(s\in\mathbb R^{n\times n}\) with the Hermitian matrix
\[
 S(s)_{jk}=\tfrac12(s_{jk}+s_{kj})+
                   \tfrac i2(s_{jk}-s_{kj}).
\]
This real-linear bijection fixes the coordinate measure \(ds\).
Call \(\psi\in C_c(\mathbb R^{n\times n},\mathbb C)\) a singular
test if for each Hermitian \(T\) with \(\det T=0\),
\[
       \int \psi(s)\exp(2\pi i\operatorname{tr}(T S(s)))\,ds=0 .
\]
Then for each \(y,b\in G\),
\begin{equation}
           \int f\bigl(y n(S(s)) b\bigr)\psi(s)\,ds=0.
 \label{eq:sw-singular-vacuum}
\end{equation}

\paragraph{A matrix Gaussian with singular frequencies.}
For a positive definite Hermitian \(P\), a Hermitian \(H\), and
complex Lebesgue measure in the \(nm\) entries of an
\(n\)-by-\(m\) matrix \(Z\), one has
\begin{equation}
 \det(P-iH)^{-m}
 =\pi^{-nm}\int_{\mathbb C^{n\times m}}
   e^{-\operatorname{tr}(Z^*PZ)}e^{i\operatorname{tr}(Z^*HZ)}\,dZ.
 \label{eq:sw-singular-gaussian}
\end{equation}
Indeed, put \(Z=P^{-1/2}W\). The real Jacobian for each column is
\(\det P^{-1}\). Diagonalize the Hermitian
\(P^{-1/2}HP^{-1/2}\) by a unitary change preserving both measure
and the norm. Each eigenvalue \(\lambda\) gives a factor
\[
 \int_{\mathbb C} e^{-(1-i\lambda)|z|^2}\,d\Re z\,d\Im z
       =2\pi\int_0^\infty r e^{-(1-i\lambda)r^2}\,dr
       =\frac{\pi}{1-i\lambda}
\]
in each column. Multiplication of these factors proves the
formula, including the normalization. All changes and factorizations
are justified by the absolute Gaussian majorant. For \(m=0\) the
integral and the determinant power are both one.

Fix a Hermitian \(H_0\). Insert \(H=H_0+S(s)\) in
\eqref{eq:sw-singular-gaussian} and integrate against \(\psi(s)\).
The absolute integral is bounded by
\(\|\psi\|_1\int e^{-\operatorname{tr}(Z^*PZ)}\,dZ<\infty\).
For fixed \(Z\), the \(S(s)\)-dependent factor becomes
\(\exp(2\pi i\operatorname{tr}(T S(s)))\) with
\(T=ZZ^*/(2\pi)\). It is Hermitian and has rank at most \(m<n\).
The definition of a singular test therefore annihilates the inner
\(s\)-integral, and Fubini gives
\begin{equation}
      \int\psi(s)\det(P-i(H_0+S(s)))^{-m}\,ds=0.
 \label{eq:sw-singular-hermitian}
\end{equation}
In particular the zero frequency is included. Transposing the
rectangular coordinate array gives the row convention of the
source, with the same trace pairing and determinant.

\paragraph{Factorization on the open matrix cell.}
Let \(y,b\in G\) and first assume \(C=C_y\) invertible. Write
\(D=D_y\), \(U=U_b\), \(V=V_b\). Since
\(\binom{V}{U}=b\binom{iI}{I}\), preservation of \(J\) gives
\[
          U^*V-V^*U=2iI .
\]
If \(Uv=0\), evaluating the identity against \(v\) gives
\(2i\|v\|^2=0\), so \(U\) is invertible. This argument holds for
every \(b\in G\). Also \(yJy^*=J\), since \(J^{-1}=-J\) and
\(y^*Jy=J\). Its bottom-right block gives \(CD^*=DC^*\).
It follows that \(C^{-1}D=(C^{-1}D)^*\).

Put \(L=VU^{-1}\). Multiplying the preceding identity for \(U,V\)
by \((U^{-1})^*\) and \(U^{-1}\) shows that
\[
 L-L^*=2iP,\qquad P=(U^{-1})^*U^{-1}>0,\qquad
 H_0=\tfrac12(L+L^*)+C^{-1}D=H_0^* .
\]
For any Hermitian \(S\), direct block multiplication gives
\[
 \begin{split}
 U_{y n(S)b}
     &=CV+(D+CS)U\\
     &=C(L+C^{-1}D+S)U\\
     &=C\bigl(i(P-i(H_0+S))\bigr)U .
 \end{split}
\]
Thus
\(\det U_{y n(S)b}=\kappa\det(P-i(H_0+S))\), where
\(\kappa=i^n\det C\det U\ne0\).
Since \(\det n(S)=1\),
\[
 f(y n(S)b)
   =(\det y\det b)^e\,\kappa^{-m}\det(P-i(H_0+S))^{-m}.
\]
Its integral vanishes by \eqref{eq:sw-singular-hermitian}.

\paragraph{The remaining matrices and the compact average.}
For arbitrary \(y\in G\), multiply on the right by
\(\bar n(t)=\left(\begin{smallmatrix}I&0\\tI&I\end{smallmatrix}\right)\)
with \(t\) real. This belongs to \(G\), and the new lower-left block
is \(C_y+tD_y\). The polynomial \(p(T)=\det(C_y+TD_y)\) is nonzero,
since \(p(-i)=(-i)^n\det(iC_y+D_y)\ne0\). Its real roots form a
finite set. Thus for all sufficiently small nonzero real \(t\)
the result on the open cell applies to \(y\bar n(t)\).
The determinants defining \(f\) are everywhere nonzero on \(G\),
so \((t,s)\mapsto f(y\bar n(t)n(S(s))b)\) is continuous. It is
bounded on \([-1,1]\times\operatorname{supp}\psi\).
Dominated convergence as \(t\to0\) proves
\eqref{eq:sw-singular-vacuum} for the original \(y\).

Let \(K\) be a compact group with a continuous homomorphism
\(\iota\) onto \(G\cap U(2n)\), let \(\mu\) be its Haar measure,
and let \(\rho\) be a continuous complex weight on \(K\).
For \(a,b,x\in G\), the associated averaged test has value
\[
 \int_K \rho(k)\left(\int
       f(x\iota(k)a n(S(s)) b)\,\psi(s)\,ds\right)d\mu(k)=0.
\]
Indeed the inner integral vanishes at each \(k\) by
\eqref{eq:sw-singular-vacuum}, with \(y=x\iota(k)a\).
The product of \(K\) and the support of \(\psi\) is compact,
so the double integral is absolutely convergent. This covers
the allowed weights of the local vacuum projection.
The argument evaluates the section at its distinguished parameter.
Derivatives in the parameter, and their use in controlling poles
of the continued Eisenstein series, require an additional argument.

\paragraph{Source correspondence.}
The tagged source defines \texttt{IsSingTest} using the displayed
coordinate pairing in
\srcid{SiegelWeilProof.HolomorphicContinuation.IsSingTest}.
Within its \texttt{CayleyCosetSums} directory, the rank bound and
Gaussian integral appear in \texttt{TraceTransposeMul} and
\texttt{InnerGen.core2}; the block identities and factorization
are \texttt{SecondCoordMulInner} and \texttt{Imports.factor}.
The extension from the open cell is
\texttt{VanishInv.vanish\_inv} and
\texttt{EventuallyUnit.vanish\_all}. Its compact-average conclusion
is \texttt{SingularTestsKillVacuum.lt3}.
\footnote{\href{https://github.com/ChristianKudeba/alpoge-rh-verification/tree/alpoge-source-7cbf7a0c/src/Development/WeilBound/RepresentationTheoryGL/SiegelWeilProof/HolomorphicContinuation/CayleyCosetSums}{Tagged local proof files}.
The global unitary identity in the anisotropic range is
S.~Yamana, \emph{The Siegel--Weil formula for unitary groups},
Pacific J. Math. 264 (2013),
\href{https://msp.org/pjm/2013/264-1/pjm-v264-n1-p11-p.pdf}{Theorem~2, p.~237}.}

\section{Whittaker values at exceptional finite places}
\label{sec:ramified-whittaker}

There are two local normalizations in the theta comparison. At a
nonsingular Hermitian frequency the Whittaker coefficient measures a
Gram fiber. In the split rank descent a one-dimensional Fourier
transform introduces a ramified Gauss sum. We compute both
normalizations. This distinguishes the scalar in a rank step from
the measure used to count solutions to a Gram equation.

\subsection{The weighted Gram fiber}

Let \(E/F\) be the CM extension with conjugation \(\sigma\), \(v\) any
finite place of \(F\), \(k=F_v\), and \(A=k\otimes_F E\). Use the
involution induced by \(\sigma\) and write \(b^*={}^{t}\bar b\).
Let \(J\) be a nonsingular Hermitian form in \(M\) variables over
\(E\). Let \(V=(A^M)^n\), \(H=\operatorname{Herm}_n(A)\), and let
\(Q(Y)\) be the Hermitian Gram matrix of the \(n\) vectors in \(Y\)
with respect to \(J\). Fix \(T\in H\) with \(\det T\in A^\times\).
This determinant condition also makes sense when \(A=k\times k\).
Write \(\psi_v\) for the local component of the specified nontrivial
unitary adelic additive character. On \(H\) the pairing used with
this character is
\[
 [b,D]=\tfrac12\operatorname{Tr}_{A/k}\operatorname{tr}(bD).
\]
Take the prescribed self-dual Haar measures \(\mu\) on \(V\) and
\(\eta\) on \(H\), with \(\eta\) self-dual for \(\psi_v([b,D])\).
Take any compact open additive subgroups \(L_1,L_2\subset A^n\) and
an additive Haar measure \(\lambda\) there. For \(b\in H\) set
\[
 h(b)=\frac{\lambda(L_1+bL_2)}{\lambda(L_1)},\qquad
 s_0=(M-n)/2 .
\]
Let \(\phi\) be locally constant and compactly supported on \(V\).
The density input is a complex-valued function \(g\), continuous on
the open set \(H^\times\) of nonsingular forms, such that
\[
 \int_{Q^{-1}(C)}\phi(Y)\,d\mu(Y)=\int_Cg(D)\,d\eta(D)
 \quad(C\subset H^\times\text{ compact}).
\]
The low-rank and equal-rank density theorems used by the global
formula supply such functions. The following local argument applies
to every \(g\) satisfying this condition.

Put
\[
 U_T(b)=\int_V\phi(Y)\psi_v([b,Q(Y)-T])\,d\mu(Y),
 \qquad F_z(b)=h(b)^{-(z-s_0)}U_T(b).
\]
There is an entire function \(W_T(z)\) equal to
\(\int_H F_z(b)\,d\eta(b)\) on a right half-plane, and
\begin{equation}
                       W_T(s_0)=g(T).
 \label{eq:exceptional-whit-density}
\end{equation}
We prove this also for nonstandard test functions and lattices.

\paragraph{Extraction of the density.}
Give matrices the maximum of their entry coordinate norms in a fixed
\(k\)-basis of \(A\). Write \(B_R=\{b\in H:\|b\|\leq R\}\) and
\(P_R=\{D\in H:\psi_v([b,D])=1\ \forall b\in B_R\}\). For \(R>0\),
\(B_R\) is a compact open additive subgroup. Orthogonality gives
\[
  \int_{B_R}\psi_v([b,D])\,d\eta(b)=\eta(B_R){\bf1}_{P_R}(D).
\]
Indeed translation by any \(b_0\in B_R\) multiplies this integral by
\(\psi_v([b_0,D])\). If \(D\notin P_R\), choose \(b_0\) with
nontrivial multiplier and subtract.

Here is an explicit bound on the annihilator. Choose \(y_0\in k\)
with \(\psi_v(y_0)\ne1\). Nondegeneracy of the trace form on the
separable algebra \(A\) gives \(c>0\) such that for each \(t\in A\)
some \(a\) satisfies \(\|a\|\leq1\) and
\(\lvert\operatorname{Tr}_{A/k}(at)\rvert\geq c\|t\|\).
For example, use a fixed trace-dual basis, scaled into the unit
ball by one scalar, and the equivalence of coordinate norms.
Choose \(C_\sigma\) with \(\|\bar a\|\leq C_\sigma\|a\|\).
For \(D\ne0\) choose a largest entry \(D_{ij}\) and apply the trace
bound to get \(a\) and \(w=\operatorname{Tr}(aD_{ij})\ne0\).
With \(X\) supported at \((j,i)\) with entry \(a\), the matrix
\(b=(y_0/w)(X+X^*)\) is Hermitian. Trace invariance under the
involution and cyclicity give \([b,D]=y_0\). Thus
\[
 D\in P_R\quad\Longrightarrow\quad
 \|D\|\leq\frac{|y_0|\max(1,C_\sigma)}{cR}.
\]
The factor \(1/2\) in the pairing cancels the two traces of \(X\)
and \(X^*\), including in residue characteristic two.
In particular \(P_R\) is closed and bounded, hence compact.
Fourier inversion at zero applied to \({\bf1}_{B_R}\), whose transform
is the displayed compactly supported function, now gives
\(\eta(B_R)\eta(P_R)=1\).

Fubini applies to \(B_R\times\operatorname{supp}\phi\). It gives
\[
 \int_{B_R}U_T(b)\,d\eta(b)=\eta(B_R)
                  \int_{Q^{-1}(T+P_R)}\phi(Y)\,d\mu(Y).
\]
For all sufficiently large \(R\), the compact set \(T+P_R\) lies in
\(H^\times\), and the last integral equals \(\int_{T+P_R}g\,d\eta\).
The absolute difference between the left side and \(g(T)\) is at
most \(\sup_{D\in T+P_R}|g(D)-g(T)|\), by the volume product above.
Continuity at \(T\) makes this tend to zero.

\paragraph{Cancellation on a large shell.}
Fix \(q_0>1\). We show that for every sufficiently large \(R\) and
every complex \(z\),
\[
          \int_{R<\|b\|\leq q_0R}F_z(b)\,d\eta(b)=0 .
\]
For \(a=1+x\) with \(\|x\|\leq r_*\) sufficiently small, the
transform \(b\mapsto a^*ba\) preserves each norm ball and its Haar
measure. Indeed choose \(r_*\) so the correction
\(x^*b+bx+x^*bx\) has norm strictly less than \(\|b\|\) for \(b\ne0\),
and \(1+x\) has an inverse. Preservation of a compact open ball of
positive volume fixes the Haar scaling at one. Also \(aL_2=L_2\)
and \(a^*L_1=L_1\), since both lattices are compact open and these
automorphisms and their inverses approach the identity. It follows
that \(h(a^*ba)=h(b)\). The integral
\(J_\phi(b)=\int_V\phi(Y)\psi_v([b,Q(Y)])\,d\mu(Y)\) is unchanged:
make the corresponding change of basis on the tuple \(Y\). A single
congruence neighborhood fixes \(\phi\) under this change. One can
choose it by expressing \(\phi\) as a finite linear combination of
indicators of cosets of one compact open subgroup and stabilizing
those cosets. Shrinking \(r_*\) ensures all of these properties.

The only varying factor of \(F_z\) is now the frequency character.
Expansion gives
\[
 [T,a^*ba]-[T,b]
 =\operatorname{Tr}_{A/k}\operatorname{tr}(Tbx)+[T,x^*bx].
\]
The two terms linear in \(x,x^*\) have the same field trace, which
cancels the \(1/2\) in the pairing. There exists \(C_Q\geq0\) such
that \(\lvert[T,x^*bx]\rvert\leq C_Q\|x\|^2\|b\|\), by continuity
and finite dimensionality of multiplication, trace and involution.
Choose \(\varepsilon>0\) with \(\psi_v\) trivial on
\(\{|t|\leq\varepsilon\}\).

The Fourier character in the linear term can be made nontrivial
uniformly in the direction of \(b\). Choose \(C_m\geq0\) with
\(\|AB'\|\leq C_m\|A\|\|B'\|\) for matrices and set
\[
 C_\nu=\frac{|y_0|(C_m\|T^{-1}\|+1)}{c}.
\]
For \(b\ne0\), let \(Y=Tb\), choose a largest entry \(Y_{ij}\) and
choose \(a'\) of norm at most one with
\(\lvert w\rvert=|\operatorname{Tr}(a'Y_{ij})|\geq c\|Y\|\).
Take \(x_0\) supported at \((j,i)\) with entry \(-y_0a'/w\).
Then
\[
 \psi_v(-\operatorname{Tr}\operatorname{tr}(Tbx_0))=\psi_v(y_0)\ne1,
 \qquad \|x_0\|\leq C_\nu/\|b\|.
\]
The bound follows by writing \(b=T^{-1}Y\).

Take
\[
 R\geq\max\{1,C_\nu/r_*, C_Q C_\nu^2 q_0/\varepsilon\},
 \qquad r=C_\nu/R .
\]
On the shell and \(\|x\|\leq r\) the quadratic term lies in the
kernel of \(\psi_v\). Consequently substitution by \(1+x\) multiplies
the integrand at \(b\) by precisely
\(\psi_v(-\operatorname{Tr}\operatorname{tr}(Tbx))\).
Average the substitutions over the additive matrix ball of radius
\(r\), with normalized Haar measure. Each substitution preserves the
shell integral. For each \(b\) in the shell the multiplier has mean
zero, by translation by the \(x_0\) just constructed. Fubini on the
compact sets proves the cancellation.

\paragraph{The half-plane and continuation.}
Choose \(0<|\pi|<1\) in \(k\), put \(q_1=|\pi|^{-1}\), and let
\(K>0\) be the Haar scaling \(\eta(B_{rq_1})=K\eta(B_r)\) of
multiplication by \(\pi^{-1}\). There is \(R_1>0\) such that
\[
       \|b\|\geq R_1q_1^j\quad\Longrightarrow\quad h(b)\geq 2^j
       \qquad(j\geq0).
\]
Indeed take \(\rho\geq0\) bounding the norm on \(L_1\) and take a
small nonzero scalar \(t_0\) with \(u t_0 e_i\in L_2\) for each
coordinate vector \(e_i\) and every \(u\in\mathcal O_k\). Put
\(R_1=(\rho+1)/|t_0|\). Choose a column \(b e_i\) containing an
entry of norm \(\|b\|\). For every binary word of length \(j\) take
\(u=\sum_{\ell=0}^{j-1}d_\ell\pi^\ell\), \(d_\ell\in\{0,1\}\).
Distinct words differ in norm by at least \(|\pi|^j\), by their
first differing digit. The \(2^j\) points \(u t_0 b e_i\) are
distinct modulo \(L_1\), since their differences have norm at
least \(|\pi|^j|t_0|\|b\|\geq \rho+1\). Their \(L_1\)-cosets
all lie in \(L_1+bL_2\), which proves the asserted lower bound.

Write \(s=\Re z-s_0\) and \(V_1=\eta(B_{R_1})\).
For \(s>\log_2(2K+2)\) we have \(K2^{-s}<1\). On the shell
\(R_1q_1^j\leq\|b\|\leq R_1q_1^{j+1}\), the integral
of \(|F_z|\) is at most
\(\|\phi\|_{L^1(\mu)}K V_1(K2^{-s})^j\). The compact ball adds a
finite contribution. This proves absolute convergence on the
stated half-plane.

Apply shell cancellation with \(q_0=2\). Fix \(R\) beyond its cutoff.
For every \(z\) the integrals over \(B_{2^jR}\) all equal the integral
over \(B_R\). On this compact ball \(1\leq h\leq C\), since \(L_2\)
and the ball are compact and the Minkowski sums lie in a bounded
set. The index \(h\) is measurable, indeed locally constant, by
local constancy of the sets \(L_1+bL_2\) in \(b\).
Differentiation under the compact integral is valid uniformly on
compact \(z\)-sets, since each derivative contributes a bounded
power of \(\log h\). The integral over \(B_R\) therefore defines
an entire function. On the half-plane of absolute convergence the
exhaustion gives the integral over \(H\). At \(s_0\) the weight is
one, so the exhaustion and the density limit give
\eqref{eq:exceptional-whit-density}. The local computation applies
at each member of the finite exceptional set in the global
factorization, independently of any formula for spherical vectors.

\subsection{The conductor and phase in split rank descent}

Now take \(v\in T=S_0\cup Q\) as in Section~\ref{sec:theta-match}.
This is a split place with designated embedding \(E\to k\). Write
\(\nu\) for the corresponding local component of the ratio of the
two characters. It is unitary and ramified, and \(\psi_v\) has
annihilator \(\mathcal O_k\) on \(\mathcal O_k\). Write \(\varpi\)
for a uniformizer, \(q=|\mathcal O_k/\varpi|\), \(U_0=\mathcal
O_k^\times\), and \(U_i=1+\varpi^i\mathcal O_k\) for \(i>0\).
Let \(f\) be the least \(i\) with \(\nu(U_i)=1\). It is positive.
Continuity gives a finite conductor: choose \(U_i\) on which all
values lie within \(1/2\) of one. Every power of each of its
character values then lies in the same disc. A subgroup of the
circle contained in this disc is trivial. Minimality gives a
\(u_0\in U_{f-1}\) with \(\nu(u_0)\ne1\), also when \(f=1\).

Put \(D_f=\mathcal O_k/\varpi^f\mathcal O_k\). For \(x\in D_f\) define
\[
 G(x)=\sum_{y\in D_f^\times}\nu(\widetilde y)^{-1}
                     \psi_v(\widetilde x\widetilde y/\varpi^f).
\]
The tilde denotes any integral lift. All lifts of a unit are units,
and replacing one by another multiplies it by an element of \(U_f\).
Changes of either lift change the additive argument by an integer
of \(k\). Thus both factors are well defined. In additive Haar
measure \(dt\) with \(\operatorname{vol}(\mathcal O_k)=1\), one has
the exact integral
\[
 \int_{\mathcal O_k^\times}\nu(t)^{-1}
                  \psi_v(\widetilde x t/\varpi^f)\,dt=q^{-f}G(x),
\]
since each additive coset in the support has mass \(q^{-f}\).

For a unit \(a\), the substitution \(y\mapsto ay\) gives
\(G(ax)=\nu(a)G(x)\). For a nonunit \(x\), multiplication of \(x\)
by \(u_0\) does not change \(x\bmod\varpi^f\), so
\(G(x)=\nu(u_0)G(x)\), hence \(G(x)=0\).
The additive pairing on \(D_f\) is perfect by the conductor
assumption on \(\psi_v\): if
\(\psi_v(xd/\varpi^f)=1\) for every integral \(x\), then
\(d/\varpi^f\in\mathcal O_k\). Translation on the finite group
therefore gives
\[
 \sum_{x\in D_f}\psi_v(xd/\varpi^f)
     =\begin{cases}q^f&d=0\text{ in }D_f,\\0&d\ne0\text{ in }D_f.
       \end{cases}
\]
Multiply \(G(x)\) by \(\psi_v(-x/\varpi^f)\) and sum in \(x\).
Only the term \(y=1\) survives, and gives \(q^f\). This proves the
source's nonvanishing \(\operatorname{gauss}(\psi_v,\nu)=G(1)\ne0\).

The same finite calculation gives the size when the characters are
unitary. Expand \(\sum_x|G(x)|^2\) and sum in \(x\) first.
Orthogonality leaves \(q^f|D_f^\times|\). Summing instead over
units \(x\) and using \(G(x)=\nu(x)G(1)\) gives
\(|D_f^\times|\,|G(1)|^2\). Thus
\[
 |G(1)|=q^{f/2},\qquad
 w_v=\nu(\varpi)^f\,q^{-f/2}G(1),\qquad |w_v|=1 .
\]
At the integer argument \(z\) of the source's rank step, its local
factor with the chosen sign \(e\in\{1,-1\}\) is exactly
\begin{equation}
 \nu(e)\epsilon_v(z)=\nu(e)\,q^{f(1/2-z)}w_v .
 \label{eq:rankstep-local-phase}
\end{equation}
Indeed the definition of \(\epsilon_v(z)\) is
\(q^{-fz}G(1)\nu(\varpi)^f\). Replacing \(\varpi\) by \(a\varpi\)
for a unit \(a\) changes the sum by \(\nu(a)^{-f}\) and the
uniformizer character factor by \(\nu(a)^f\). These cancel.
Replacing \(\psi_v(t)\) by \(\psi_v(at)\) multiplies \(w_v\) by
\(\nu(a)\). The character \(\nu(e)\) squares to one.
The hypotheses give the phase \(w_v\) by the displayed sum. A
numerical root of unity for this phase would require further
specialization. Even a finite-order character on units does not
specify the value of \(\nu(\varpi)\).

For the beta and alpha steps the argument is respectively
\(z=-\delta\) and \(z=1-\delta\), where
\(\delta=(a-a')/2\). Each exceptional place contributes
\eqref{eq:rankstep-local-phase}, with the sign of its rank-step
formula. The finite character sum computes these character-dependent
factors. The geometric factor in the rank step, the surjectivity
of the two maps onto the split induced representation and the
identity of local tensor models are separate inputs proved in the
respective descent and local matching arguments.

\subsection{A matched pair and the sign of the cone integral}

We can determine the local sign in \eqref{eq:rankstep-local-phase}
from the standard vectors used in the source. Put
\(\ell_\nu(t)={\bf1}_{\mathcal O_k^\times}(t)\nu(t)^{-1}\), extended
by zero at \(t=0\). If \(e\) selects \(r\) columns of an \(r\)-by-\(d\)
matrix \(Y\) in their indicated order, set
\[
  S_e(Y)=\ell_\nu(\det Y_e)
  \prod_{i,j}{\bf1}_{L_j}(Y_{ij}),\qquad
  L_j=\begin{cases}
       \mathcal O_k&j\in e,\\
       \varpi^f\mathcal O_k&j\notin e .
       \end{cases}
\]
Let \(e_A\) select the first \(a\) columns and \(e_B\) the last \(a'\).
These standard vectors are compactly supported, locally constant,
and have equal nonzero images in the two transforms \(P_A,P_B\)
of Section~\ref{sec:theta-match}.

For a matrix \(y\) of size \(a'\)-by-\(a\), write
\(u(y)=\left(\begin{smallmatrix}1&0\\y&1\end{smallmatrix}\right)\).
The matrices inserted in the two integrals have frames
\([1,-{}^ty]\) and \([y,1]\). In either transform the selected
minor is \(g\). On the support of the integrand its entries are
integral and its determinant a unit, so \(g\in\mathrm{GL}_r(\mathcal
O_k)\). The inverse is integral too. Thus the off-minor condition
is equivalent to \(y\in\varpi^f M_{a',a}(\mathcal O_k)\). The
weight \(\nu(\det g)\) cancels \(\ell_\nu(\det g)\), both absolute
determinant factors are one and the compact group has measure one.
Both transforms therefore give the indicator of this same set in
\(y\), with value one at zero.

To verify equality at arbitrary \(h\), write
\(h=\left(\begin{smallmatrix}A&B\\C&D\end{smallmatrix}\right)\).
When \(D\) is singular, its complementary minor in \(h^{-1}\)
is singular as well, by the complementary-minor formula for an
invertible matrix. Thus each selected determinant vanishes in the
corresponding frame and both integrands are zero. When \(D\) is
invertible put \(y=D^{-1}C\). Then \(h=p u(y)\) with \(p\) upper
block triangular. In the two transforms, left multiplication by \(p\)
has the same scalar. Indeed, writing \(A_0,D_0\) for its diagonal
blocks, changes of the integration variables by \(A_0^{-t}\) on
the large side and by \(D_0\) on the small side give
\[
 \alpha(\det A_0)\beta(\det D_0)
            |\det A_0|^{a'/2}|\det D_0|^{-a/2}.
\]
Here \(\nu=\alpha/\beta\). The covariance on the large side uses
the ramified shell transformation law of \(P_A\).
This proves equality as functions on all of \(\mathrm{GL}_d(k)\).
Transport the vectors to the two Weil models and write them
\(\phi,\phi'\). The construction at the lower rank gives matched
vectors \(\phi_1,\phi'_1\) with common nonzero image.

In the beta step \(a'=b+2\), put \(r=a\) and \(j=\delta+1\).
In the alpha step \(a=b+2\), put \(r=a'\) and \(j=\delta\). The
zero-orbit evaluation of the standard vector is \(q^{-fr}\) times
its lower standard vector: its \(r\) independent restricted
coordinates range over \(\varpi^f\mathcal O_k\), and the remaining
integrated unit balls have mass one. We compute the cone evaluation
in the same coordinates. Number its distinguished rows and selected
columns from \(0\) to \(b+1\), put \(l=b+1\), and let \(e\) select
the block undergoing rank descent. In the beta chart the free
coordinate positions \(P\) are those in row \(l\), together with
those in column \(e(l)\) outside row \(0\). In the alpha chart
they are those in column \(e(0)\), together with those in row \(0\)
outside column \(e(l)\). The corner \((0,e(l))\) is fixed to the
first split variable \(s\), and the opposite corner \((l,e(0))\)
has an integration variable \(u\) with character \(\psi_v(-tu)\),
where \(t\) is the second split variable.

In the beta chart, expansion along row \(0\) and then the column
\(e(0)\), and in the alpha chart expansion along column \(e(l)\)
and then row \(l\), give the same selected determinant
\[
                          -su\det(Y_{\rm mid}).
\]
The sign is the sign of interchanging the first and last of the
selected columns. On the integral domain of the coordinate
indicators all factors of this product are integral, so the
determinant is a unit exactly when they all are. Its character
factor splits as
\(\ell_\nu(s)\ell_\nu(-u)\ell_\nu(\det Y_{\rm mid})\).
Change \(u\) to \(-u\) in the corner integral. The two minus signs
cancel, giving
\[
       I(t)=\int_{\mathcal O_k^\times}\nu(u)^{-1}\psi_v(tu)\,du .
\]
There are exactly \(r\) free positions with column outside \(e\):
the beta chart has \((l,e_A(i))\) for the \(a\) unselected columns, and the alpha chart
has \((0,e_B(i))\) for the \(a'\) unselected columns.
Each contributes \(q^{-f}\). All other coordinate integrals carry
unit volume. The surviving factor at the split variables is
therefore
\[
                       q^{-fr}\ell_\nu(s) I(t)
\]
times the lower standard vector.

For completeness, \(I(t)\) is supported on
\(\varpi^{-f}\mathcal O_k^\times\). For
\(|t\varpi^f|<1\) use \(u_0\in U_{f-1}\) as above: substitution
by \(u_0\) preserves \(\psi_v(tu)\) and multiplies the integral
by a nontrivial scalar. For \(|t\varpi^f|>1\) choose
\(h\in\varpi^f\mathcal O_k\) with \(\psi_v(th)\ne1\), using the
annihilator of \(\mathcal O_k\). Translation by \(h\) preserves
\(\ell_\nu\) and forces the integral to vanish. For a unit \(v'\)
one has \(I(\varpi^{-f}v')=\nu(v')I(\varpi^{-f})\), and
\(I(\varpi^{-f})=q^{-f}G(1)\) by the coset integral.

The two components of the split idele have character
\(\nu(s)\nu(t)^{-1}\) and norm \(|s|\,|t|\). The local integral
uses multiplicative Haar measure giving the product of the unit
groups mass one. Its support is a single translate of that subgroup,
\(s\in\mathcal O_k^\times\), \(t=\varpi^{-f}v'\) with \(v'\) a
unit. On it \(\ell_\nu(s)\nu(s)=1\), the two values of \(\nu(v')\)
cancel and \((|s||t|)^j=q^{fj}\). The local cone operator therefore
takes the standard vector to its lower vector times
\[
 q^{-fr}\, q^{-f}G(1)\nu(\varpi)^f q^{fj}
                  =q^{-fr}\epsilon_v(1-j).
\]
The compact stabilizer average has mass one and fixes the vector.
Comparison with the zero-orbit evaluation cancels exactly \(q^{-fr}\).
Writing \(p_1,p'_1\) for the lower quotient maps, the beta formula is
\[
 p'_1(\mathrm{cone}_v\phi')=\epsilon_v(-\delta)\,
                                p_1(\mathrm{ev}_v\phi).
\]
The alpha formula is
\(p_1(\mathrm{cone}_v\phi)=\epsilon_v(1-\delta)\,
p'_1(\mathrm{ev}'_v\phi')\). The respective vector on the right
before multiplication is nonzero because the lower standard vectors
match and the value of their common transform at the identity is one.

Thus the source's witnesses for the local sign \(e\) are \(+1\) in
both steps. To extend from these vectors to all matched fibers,
the source uses ramified multiplicity one for the equivariant
maps on the common quotient. In the beta case, let \(D,D'\) be
the linear maps it obtains by descending the evaluation and the
cone, respectively. On the common image \(\xi=p\phi=p'\phi'\)
the computation gives \(D'\xi=\epsilon_v(-\delta)D\xi\) and
\(D\xi\ne0\). Since \(D'=cD\) for some scalar, evaluation at
\(\xi\) forces \(c=\epsilon_v(-\delta)\). The same reasoning
with the roles interchanged gives \(\epsilon_v(1-\delta)\) in
the alpha case. Existence and equivariance of the descended maps,
the local split coordinate changes and ramified multiplicity one
are the representation-theoretic inputs to this last argument.

\paragraph{Correspondence with the release.}
In \texttt{SiegelWeilFormula.LocalCharacterIntegrals} the statements
used above are \texttt{shellP\_core}, \texttt{perp\_bound},
\texttt{vol\_mul\_vol\_perp}, \texttt{dualP}, \texttt{convP} and
\texttt{locFinX\_of}. The global applications are
\texttt{whitGenLtOf} and \texttt{whitGenEqOf} and import the
representation-density hypotheses. The split calculation uses \texttt{std\_pair}, \texttt{int\_cone}, \texttt{idele\_int}, \texttt{eps\_eq}, \texttt{stdB}, \texttt{stdA\_gen}, \texttt{locB\_of} and \texttt{locA\_of}. In \texttt{WeilTensorModels}
the functions are \texttt{gauss}, \texttt{eps} and \texttt{stepFac}.
\texttt{ThK.GaussNe.gauss\_ne\_zero} proves the required
nonvanishing and \texttt{ThKG.Loc.stepFac\_eq} identifies the finite
product. The modulus and the change-of-uniformizer computation above
follow by the finite sums displayed here.

\section{Full-rank theta decay and the indefinite theta integral}
\label{sec:theta-decay-proof}

Keep the arithmetic base \(B\), its CM extension \(E/F\) and involution
\(\sigma\), its additive character \(\psi\), the embeddings
\(\tau_w:E\to\mathbb C\) above the infinite places \(w\) of \(F\), and
its distinguished split finite place \(v_1\) with a chosen embedding
\(\tau_1:E\to F_{v_1}\). Put \(D=[E:\mathbb Q]\).
Let \(X,W\) be the nondegenerate Hermitian and skew-Hermitian spaces
of dimensions \(n,d\) and signatures \((r,s),(r',s')\) in the pair
of Section~\ref{sec:theta-match}. Write \(U_X,U_W\) for their adelic
unitary groups, \(M\) for the specified Schr\"odinger model,
\(\omega\) for its Weil action, \(\Theta_\Phi(g,h)\) for its theta
sum, and \(\mathcal F_X\) for its rational fundamental domain with
the model's Haar measure \(\mu_X\). The pair hypothesis includes the
good base and model, the Hecke characters and their required restrictions,
and the Weil action on this model.

Let \(\mathcal S_\infty\) denote the Schwartz space on the infinite
part of \(\mathbb A_F^{nd}\), and \(\mathcal S_f\) the compactly
supported, locally constant functions on the finite part. In the
chosen split chart, \(M.\mathrm{SoLoc}\) consists of local
Schwartz--Bruhat functions whose image vanishes on matrices of row
rank less than \(n\). The image is taken under a choice of admissible local
intertwiner in the definition of the chart. In the global space
\(M.\mathrm{So}\) take the span of pure tensors with this condition at
\(v_1\), arbitrary Schwartz factors elsewhere and the specified
vacuum at almost every finite place. Set
\[
 \mathcal S_f^\circ=
 \{\,\phi\in\mathcal S_f:
          f\otimes\phi\in M.\mathrm{So}\quad
               \text{for every } f\in\mathcal S_\infty\,\}.
\]
Each member of \(M.\mathrm{So}\) is a finite sum of such tensors.
Let \(H\) be the matrix height of the model and
\(H_Q(g)=\inf_{\gamma\in U_X(F)}H(\gamma g)\).

For each \(N_0\in\mathbb N\) there is \(k\in\mathbb N\) such that
for each \(\phi\in\mathcal S_f^\circ\) there is a continuous seminorm
\(p_\phi\) on \(\mathcal S_\infty\) satisfying
\begin{equation}
 |\Theta_{f\otimes\phi}(g,h)|H_Q(g)^{N_0}
       \le p_\phi(f) H(h)^k
 \qquad(f\in\mathcal S_\infty,\ g\in U_X,\ h\in U_W).
 \label{eq:theta-full-decay}
\end{equation}
The exponent \(k\) is uniform in \(\phi\). For unrestricted finite
Schwartz data, the hypothesis
\begin{equation}
 d+2\min(r,s)<2n\quad\text{or}\quad\min(r,s)=0
 \label{eq:theta-weil-range}
\end{equation}
gives a possibly different uniform \(k\) and, for each \(\phi\),
a continuous seminorm \(p_\phi\) and a nonnegative measurable
\(G_\phi\in L^1(\mathcal F_X,\mu_X)\) such that
\[
 |\Theta_{f\otimes\phi}(g,h)|
       \le p_\phi(f)H(h)^kG_\phi(g).
\]
We derive these statements with the same local data and measure.

\subsection{The square of the kernel}

For \(z\in\operatorname{Mat}_{n,d}(\mathbb A_F\otimes_F E)\)
write \(|z|_\infty=\sum_{w,i,j}|\pi_{\tau_w}(z_{ij})|\). Denote its
finite projection by \(z_f\). For compact \(C\) in this matrix space put
\[
 b_{C,N}(z)=\mathbf 1_{\{z_f\in C_f\}}\,(1+|z|_\infty)^{-N},
 \qquad
 L_{C,N,Y}(g,h)=\sum_{x\in Y} b_{C,N}(g^{-1}x(h^{-1})^t).
\]
Here \(C_f\) is the projection of \(C\), rational matrices are
embedded diagonally, and \(Y\) is either all rational matrices or
\(Y_n=\{x:\operatorname{rank}_E x=n\}\).
For each \(N\) and \(\phi\in\mathcal S_f\), the coefficient estimate produces
a compact \(C\) and a continuous seminorm \(p\) with
\begin{equation}
 |\Theta_{f\otimes\phi}(g,h)|^2\le p(f)^2 L_{C,N,Y}(g,h)
 \label{eq:theta-lattice-majorant}
\end{equation}
for \(Y\) all matrices. For \(\phi\in\mathcal S_f^\circ\) it produces
the same conclusion with \(Y=Y_n\).

Here are the local cancellation and the role of the square in the
second assertion. Write \(P\Phi(x)=\Phi(-x)\), and let
\[
 c_\Phi(w)=\int(\operatorname{Schr}_\psi(w)\Phi)(u)
                    \overline{\Phi(-u)}\,du .
\]
If \(\operatorname{Schr}_\psi(v)\Phi=\psi(a)\Phi\), translation
invariance of additive Haar measure gives
\[
 c_\Phi(w)=\psi(v_1\cdot w_2-w_1\cdot v_2+2a)c_\Phi(w).
\]
The subscripts on \(v,w\) in this identity denote their two
polarization coordinates. Indeed translation of the integration
variable by \(v_1\) contributes the commutator of the two Schr\"odinger
operators, and reflection followed by conjugation contributes the
second copy of \(a\). The parity operator on the theta sum gives
\(\Theta_{P\Phi}(g,h)=\Theta_\Phi(g,h)\), by covariance and the
bijection \(x\mapsto-x\) on the rational summation set. The product
bound for this pair is thus a bound for an absolute square.

In the split chart take a finite subdivision of the support of each
local full-rank factor into compact open pieces. The source's local
separation lemma chooses the pieces so that for every rank-deficient
matrix in the complementary split coordinate there is a Schr\"odinger
operator acting on the piece by a scalar for which the displayed
multiplier differs from \(1\). The complementary coordinate is
multiplied on the left and right by the invertible matrices of the
two forms and by a sign. Its rank is preserved. The coefficient
therefore vanishes there. For a rational matrix of rank less than
\(n\), extending scalars to this component and multiplying by
\(g^{-1},(h^{-1})^t\) still gives rank less than \(n\).
After this cancellation, the rational coefficient-sum formula,
Weil covariance and the Schwartz coefficient bound give
\eqref{eq:theta-lattice-majorant} on each piece. If \(p_0\) is
the seminorm for the two coefficients, use
\(p_0(f)+p_0(Pf)\) to bound their geometric mean. Formula
\eqref{eq:theta-lattice-majorant} is stable under finite sums. Take square roots, replace the finitely many compact sets by their
union, then add the seminorms with the absolute values of the scalar
coefficients. This explains the use of the full-rank locus after
taking the square. The remaining analytic inputs at this point are the coefficient-sum
formula, covariance and the Fourier seminorm estimate.

For completeness the separating subgroup has an explicit construction.
In \(k=F_{v_1}\), with valuation ring \(\mathcal O\), choose \(b\)
with \(\psi_{v_1}(b)\ne1\). Write
\(\langle Y,U\rangle=\operatorname{tr}(YU^t)\).
Let \(T\) be the compact nonzero set of the local factor in the split
chart. Compactness here follows from local constancy and compact
support. The full-rank locus is open, since its complement is the
common zero locus of all \(n\)-row minors. Since \(2\) is nonzero
in \(k\), \(2T\) is contained in the full-rank locus. Compactness
gives a neighborhood \(V_0\) of zero such that \(2T+V_0\) remains
in that locus. Take \(a\ne0\) with
\(a\operatorname{Mat}_{n,d}(\mathcal O)\subset V_0\) and put
\[
 K=(b/a)\operatorname{Mat}_{n,d}(\mathcal O),\qquad
 \Lambda=\{Y:\psi_{v_1}(\langle Y,U\rangle)=1\ (U\in K)\}.
\]
This subgroup is open by local constancy of the character, continuity
of the pairing and compactness of \(K\). It lies in
\(a\operatorname{Mat}_{n,d}(\mathcal O)\). Indeed if some
\(Y_{ij}/a\notin\mathcal O\), its inverse belongs to \(\mathcal O\);
taking \(U\in K\) with the single entry
\((b/a)(Y_{ij}/a)^{-1}\) gives \(\langle Y,U\rangle=b\).

The image of \(T\) in the discrete quotient by \(\Lambda\) is
finite. Restrict the function to each coset meeting \(T\), choose
\(Y_0\) in its intersection with \(T\), and let \(Z\) be rank deficient.
Since \(2Y_0+\Lambda\) is full rank, \(Z-2Y_0\notin\Lambda\).
Some \(U\in K\) thus has
\(\psi_{v_1}(\langle Z-2Y_0,U\rangle)\ne1\). On this whole
piece \(\psi_{v_1}(\langle Y,U\rangle)
=\psi_{v_1}(\langle Y_0,U\rangle)\).
The split polarization realizes this multiplication operator as the
Schr\"odinger action of the local vector with split coordinates
\((-U,0)\). Its eigenvalue is the common value just computed, and
the symplectic pairing with the complementary coordinate \(Z\)
gives \(\langle Z,U\rangle\) in the sign convention of the model.
The coefficient identity gives multiplier
\(\psi_{v_1}(\langle 2Y_0-Z,U\rangle)\ne1\), proving the vanishing.
This also covers residue characteristic two, since the argument uses
\(2\ne0\) in the characteristic-zero field and never divides by
\(2\) in the residue field.


\subsection{The two logarithmic chambers}

Reduction theory for unitary groups supplies rational Witt frames
of lengths \(q\le\min(r,s)\) and \(y\le\min(r',s')\), tori
\(a_X,a_W\), and compact sets \(\Omega_X,\Omega_W\). There are
rational \(\gamma,\delta\) such that
\(\gamma g=a_X(u)\kappa\), \(\delta h=a_W(v)\kappa'\) with
\(\kappa,\kappa'\) in these compact sets and
\[
 u_1\ge\cdots\ge u_q\ge0,\qquad v_1\ge\cdots\ge v_y\ge0.
\]
Here \(u,v\) henceforth denote real tuples. Put \(U=\sum u_i\),
\(V=\sum v_j\), and
\[
 B(u,v)=\max_{0\le i\le q}
   \left(\sum_{a=1}^i u_a-\sum_{b=1}^{\min(i,y)}v_b\right).
\]
The choices \(i=0,q\) give \(B\ge0\) and \(U\le V+B\). The frame
weights \(\lambda_i\) are the \(+u_i,-u_i\) and \(n-2q\) zeros.
The weights \(\lambda'_j\) are the \(+v_j,-v_j\) and \(d-2y\) zeros.
The positive weight sum is
\begin{align*}
 P(u,v)&=\sum_{i,j}\max(0,\lambda_i+\lambda'_j)\\
 &=2\sum_{i=1}^{q}\sum_{j=1}^{y}\max(u_i,v_j)
                  +(d-2y)U+(n-2q)V\\
 &\le dU+nV .
\end{align*}
Indeed each block of four signed pair sums contributes
\(u_i+v_j+|u_i-v_j|\). The frame height bounds and boundedness on
compact sets give constants \(A,A',k_0,k_1\) independent of \(g,h\) with
\[
 H_Q(g)\le A'e^{k_0U},\qquad e^V\le A H(h)^{k_1}.
\]
For the second inequality apply the inverse torus height bound
to the rational translate \(\delta^{-1}a_W(v)=h(\kappa')^{-1}\)
and use submultiplicativity and compactness. For the first use the
rational representative \(\gamma g\). The lattice sums are invariant
under rational left translation in either group. Right translation
over the two compact sets can multiply a bound by a constant and
enlarge \(C\). These constants are independent of \(u,v,f\).

\subsection{Counting the rational matrices}

Take rational Witt bases \(b,b'\), write \(\xi=b^{-1}x(b'^{-1})^t\),
and put \(z=a_X(u)^{-1}x(a_W(v)^{-1})^t\).
There is an integer \(m\ge1\) and a constant \(C_0\ge1\) such that
\[
 z_f\in C_f\ \Longrightarrow\ m\xi_{ij}\in\mathcal O_E,\qquad
 \sum_{i,j,w} e^{-\lambda_i-\lambda'_j}|\tau_w(\xi_{ij})|
                    \le C_0|z|_\infty .
\]
The finite tori have identity finite components, so the finite
coordinates of \(\xi\) lie in the image of a fixed compact set
under a fixed change of basis. A compact finite adelic set has a
common denominator, by a finite subcover of the sets of bounded
denominators. At infinity the two fixed basis changes have bounded
operator norms and the tori act diagonally, giving the second bound.

If \(\Lambda\) is a discrete subgroup of a real normed vector
space of dimension \(D\), then for \(L>D\) an integer,
\begin{equation}
 \sum_{\ell\in\Lambda}(1+e^{-a}\|\ell\|)^{-L}
       \le K_\Lambda e^{D\max(0,a)} .
 \label{eq:theta-onevar}
\end{equation}
Choose \(\epsilon>0\) such that the open balls of radius \(\epsilon\)
about the lattice points are disjoint. Put \(c=e^{-\max(0,a)}\).
For \(y\) in the ball about \(\ell\),
\(1+c\|y\|\le(1+\epsilon)(1+c\|\ell\|)\).
The summand in \eqref{eq:theta-onevar} is bounded by
\((1+c\|\ell\|)^{-L}\). Integrate the latter bound on each ball,
sum and divide by the volume of a ball. The result is at most
\[
 \frac{(1+\epsilon)^L}{\operatorname{vol}(B_\epsilon)}
 c^{-D}\int (1+\|y\|)^{-L}\,dy .
\]
The integral converges since the volumes of shells of radii \(2^j\)
are \(O(2^{jD})\). This proves \eqref{eq:theta-onevar}.

Choose \(N_b=(D+1)nd\). Set
\(A_{ij}=e^{-\lambda_i-\lambda'_j}\sum_w|\tau_w(\xi_{ij})|\).
The coordinate estimate gives \(1+A_{ij}\le C_0(1+|z|_\infty)\).
Multiplying these inequalities and inverting gives
\[
 (1+|z|_\infty)^{-N_b}\le
       C_0^{N_b}\prod_{i,j}(1+A_{ij})^{-(D+1)} .
\]
The rational change of basis is injective. Enlarge the summation
to all matrices with independent entries in \(m^{-1}\mathcal O_E\)
and use Tonelli's theorem to factor the nonnegative sum. In the
Minkowski embedding the norm of an element is at most
\(\sum_w|\tau_w(\xi_{ij})|\) for the supremum norm on the archimedean
coordinates, since the two embeddings over any real place have the same
absolute value. Apply \eqref{eq:theta-onevar} in each coordinate,
with \(a=\lambda_i+\lambda'_j\). The result is
\begin{equation}
 L_{C,N_b,\mathrm{all}}(a_X(u),a_W(v))\le K_C e^{DP(u,v)}.
 \label{eq:theta-all-torus}
\end{equation}

Here is the gain from full row rank. For \(x\in Y_n\) with
\(z_f\in C_f\), choose \(i\) attaining the maximum \(B(u,v)\).
The \(i\) rows of \(\xi\) with weights \(-u_1,\ldots,-u_i\)
are linearly independent, so some \(i\) distinct columns
\(c_1,\ldots,c_i\) give a square minor \(T_0\) with nonzero
determinant. The determinant \(m^i\det T_0\) is a nonzero algebraic
integer. Its absolute field norm is an integer at least \(1\),
so at an embedding \(\tau_w\) we have
\(m^i|\tau_w(\det T_0)|\ge1\). Use that embedding in the matrix
\[
 T_{ab}=e^{u_a-\lambda'_{c_b}}\,\tau_w((T_0)_{ab}).
\]
The sum of the chosen \(i\) distinct column weights is at most
\(\sum_{b\le\min(i,y)}v_b\), since the only positive weights
available are \(v_1,\ldots,v_y\) in decreasing order. Consequently
\[
 e^{B(u,v)}
 \le m^i|\det T|
 \le i!\big[m C_0(1+|z|_\infty)\big]^i
 \le \big[n!m C_0(1+|z|_\infty)\big]^{n+1}.
\]
The middle inequality follows by bounding the \(i!\) terms of the
determinant individually using the coordinate estimate above.
The empty minor at \(i=0\) has determinant \(1\). With
\(c=(n!m C_0)^{-1}\), we have proved
\(c\,e^{B/(n+1)}\le1+|z|_\infty\).
For any \(M_1\in\mathbb N\), raise the last inequality to
\(M_1(n+1)\), invert it, and use \eqref{eq:theta-all-torus} on
the remaining lattice sum. Since all summands are nonnegative, its
restriction to \(Y_n\) can only decrease it. Thus
\begin{equation}
 L_{C,N_b+M_1(n+1),Y_n}(a_X(u),a_W(v))
       \le K_{C,M_1}\exp\{DP(u,v)-M_1 B(u,v)\}.
 \label{eq:theta-full-torus}
\end{equation}
The exponent needed on the left is chosen before the compact set
\(C\), as required when applying the Schwartz estimate.

\subsection{From the chamber to an integrable majorant}

Take \(M_1=Dd+2k_0N_0\). Applying
\eqref{eq:theta-lattice-majorant} and \eqref{eq:theta-full-torus}
and taking square roots, the exponent after multiplication by
\(H_Q(g)^{N_0}\) is at most
\begin{align*}
 \tfrac12(DP-M_1B)+k_0N_0U
 &\le \tfrac12\{M_1(U-B)+DnV\}\\
 &\le \tfrac12(M_1+Dn)V
 \le (M_1+Dn)V .
\end{align*}
The height bound on \(e^V\) proves \eqref{eq:theta-full-decay}
with \(k=k_1(M_1+Dn)\). Absorb all fixed compact-set and height
constants into the continuous seminorm. In particular \(N_0=0\)
permits \(G_\phi=1\). The rational quotient of this unitary group
has finite Haar volume. For fixed \(h\) this proves absolute
integrability and gives the bound
\[
 \int_{\mathcal F_X}|\Theta_{f\otimes\phi}(g,h)|\,d\mu_X(g)
       \le\mu_X(\mathcal F_X)p_\phi(f)H(h)^k .
\]

For unrestricted finite data use
\(G_\phi(g)=L_{C',N_b,\mathrm{all}}(g,1)^{1/2}\).
To see why this function suffices, reduce the second group into
\(a_W(v)\Omega_W\). Each diagonal archimedean stretch in
\(a_W(v)\), or its inverse, has operator norm at most \(e^V\)
in Witt coordinates. It is the identity at the finite places.
The fixed changes of basis therefore give
\[
 L_{C,N_b,\mathrm{all}}(g,h)\le K e^{N_bV}
                L_{C',N_b,\mathrm{all}}(g,1).
\]
Use \(e^V\le A H(h)^{k_1}\) and take a square root in
\eqref{eq:theta-lattice-majorant}. This gives the asserted majorant
with \(k=k_1N_b\), as in the source, since \(H(h)\ge1\).
The indicator of the finite projection of a compact set is Borel,
and the summands are measurable; \(G_\phi\) is measurable.

It remains to integrate over the first group. For a nonnegative
measurable rationally left-invariant function majorized on
\(a_X(u)\Omega_X\) by \(b(u)\), the Siegel integration inequality gives
\[
 \int_{\mathcal F_X}G(g)\,d\mu_X(g)
 \le K\int_{u_1\ge\cdots\ge u_q\ge0} b(u)e^{-D\rho_n(u)}\,du,
 \qquad
 \rho_n(u)=\sum_{i=1}^q(n-2i+1)u_i .
\]
The compact set and constant here may depend on \(\mu_X,\mathcal F_X\).
Equation \eqref{eq:theta-all-torus} at \(v=0\) gives
\(b(u)=K'\exp(DdU/2)\) for \(G_\phi\). Use the differences
\(z_j=u_j-u_{j+1}\) with \(u_{q+1}=0\). This takes the chamber to
the positive orthant and has Jacobian \(1\). The cumulative coefficient
of \(z_j\) in \(D(dU/2-\rho_n(u))\) is
\[
 D\sum_{i=1}^j(d/2-n+2i-1)=Dj(d/2-n+j).
\]
Equation \eqref{eq:theta-weil-range} implies \(q=0\) or
\(d+2q<2n\), because \(q\le\min(r,s)\). Thus every coefficient
for \(j\ge1\) is strictly negative, and with product Lebesgue measure
the chamber integral, omitting its fixed prefactor, equals
\begin{equation}
 \prod_{j=1}^q \frac{1}{Dj(n-d/2-j)}.
 \label{eq:theta-integral-product}
\end{equation}
In rank zero the empty product is \(1\). This completes the
integrability calculation for unrestricted finite data. Decomposing a
global Schwartz function into finitely many pure tensors proves the
same convergence for it.

In the triple of dimensions \(a,a',a+a'\), the large first space
has signature \((a-p,p)\), the smaller one has \((q',a'-q')\),
and their common second space has signature
\((a+a'-(p+q'),p+q')\). The triple hypotheses give
\(a+a'+2\min(a-p,p)<2a\). Apply the unrestricted bound to the
large first space and the full-rank bound to the small first space
with its prescribed split test data. The pointwise constants depend
on the arithmetic model, on the chosen compact sets and on the
Schwartz data. The estimates above supply the convergence wherever
the family supplies these hypotheses; they give no uniform bound
on those constants as the arithmetic model varies.

\subsection{Decay in the variable left after integration}

For the later cuspidal member of the theta span, decay is also
needed in the remaining variable \(h\in U_W\). The source applies
the cuspidal part of unitary reduction theory under the following
hypotheses on a function \(A\) on \(U_W\). The function is left
invariant under rational points and right invariant under an open
finite level. It is archimedean smooth and lies in a family closed
under differentiation along real one-parameter subgroups. All
members of this family have polynomial growth with a common
exponent, with individual constants allowed. For each nonzero
rational totally isotropic subspace the constant term over the
associated unipotent radical vanishes. In coordinates, one can
test these subspaces on matrices \(C\) of full column rank with
\(C^*JC=0\), where \(J\) is the Hermitian matrix representing the
group. The column span and its orthogonal complement specify the
unipotent radical. For the skew-Hermitian group the source first
identifies it with a Hermitian unitary group by scaling the
conjugate matrix of its form by a nonzero scalar of the appropriate
sign under \(\sigma\).

Under these hypotheses the cuspidal estimate in the imported
reduction theorem gives a constant \(C_N\) on a Siegel set for each
\(N\). For an arbitrary \(h\), choose a rational translate
\(\delta h\) in that set. Left invariance and the definition of the
quotient height give
\[
 |A(h)|H_Q(h)^N\le |A(\delta h)|H(\delta h)^N\le C_N .
\]
In \texttt{SigRapid.sigRapid\_of}, invariance, smoothness and
uniform moderate growth are supplied by \texttt{sig\_nice} for
members of the specified theta span. Vanishing of the constant
terms is a separate hypothesis of this implication.
\texttt{CuspRapidOf.rapidG} and \texttt{CuspRapidOf.cuspRapid\_of}
make the passage from the Siegel set to the rational quotient.
This use of the cuspidal reduction theorem and the explicit
full-rank estimate \eqref{eq:theta-full-decay} concern different
variables in the iterated integrals.

\paragraph{Source and inputs.}
In the pinned source tree, under
\srcid{Development/WeilBound/ThetaLifts/FrameEstimates},
\srcid{SquareIntegrableWeilRange} defines the lattices and chambers.
The local separation and the square estimate are
\texttt{HermitianFormDetNeZero.soCut\_of},
\texttt{ThetaSumNeg.qb\_piece} and \texttt{ThetaSumNeg.sqSo\_of'}.
The one-coordinate packing, unrestricted sum, minor and full-rank
sum are \texttt{OneVar.oneVar\_of}, \texttt{LatTor.latTor\_of},
\texttt{ExistsColsDetNeZero.minor\_of} and
\texttt{LatFull.latFull\_of}. The two global estimates are
\texttt{SoSemi.soSemiU\_of} and \texttt{DomWeil.domWeilU\_of},
assembled in \srcid{WeilConv}. These names combine the identifying
file with the declaration. The imported reduction theorem
\srcid{GlobalTheory.Published.SiegelSetUnitary} supplies the cover,
height comparisons and integration inequality.
\srcid{FundamentalDomainFiniteVolume} supplies finiteness of the
model's quotient measure. The local coset argument appears in \texttt{RankIffRight.exists\_cut}.
The Heisenberg product formula and its Schwartz coefficient estimate, and the construction
of each arithmetic pair with these hypotheses, remain upstream
inputs to this chapter.

\section{Construction and rank of the boundary period lattice}
\label{sec:boundary-periods}

Let \(T\) be the complete discrete valuation ring on a horizontal
boundary trait, \(K=\operatorname{Frac}T\), \(\pi\) a uniformizer, and
\(G/T\) the semiabelian model with principally polarized abelian
generic fiber. At a nonsplit point the source first splits the toric
part of \(G_{\kappa(T)}\). In the horizontal case its finite splitting
field is separable because \(\kappa(T)\) has characteristic zero.
Lifting this extension to the finite unramified extension of \(T\)
gives the split special-fiber torus. Write \(T\) for the resulting
ring. When a split presentation has already been supplied, the
construction below starts there. The torus need not have a
complement in the special fiber.

\subsection{The formal torus and the two lattices}

Put \(S_e=\operatorname{Spec}(T/\pi^{e+1})\). The torus with character
group \(X\) has over \(S_0\) a kernel embedding into \(G_0\), with
abelian quotient \(B_0\). On every thickening there is an extension
\[
  1\longrightarrow \operatorname{Hom}(X,\mathbb G_m)_{S_e}
  \longrightarrow G_e\longrightarrow B_e\longrightarrow 0.
\]
Here one uses the formal Raynaud extension theorem for a smooth
semiabelian scheme whose special torus is isotrivial. In the present
case that torus is split. The morphisms commute with restriction to
each preceding thickening. The construction used in the release starts
with an even trivialization of the cubical line on the special torus,
lifts that trivialization and the torus kernel across each nilpotent
thickening, and descends the line to \(M_e\) on \(B_e\). Restricting a
descended cubical line back to the preceding level gives its previous
descent. The finite-level quotient is a torus torsor, so the
weight-\(x\) functions on it form an invertible sheaf. Multiplication
identifies the tensor product of the weight-\(x\) and weight-\(x'\)
sheaves with the weight-\(x+x'\) sheaf, as one sees on any
trivialization of the torsor. This supplies the homomorphism into the
dual abelian scheme on every level.

The rings \(T/\pi^{e+1}\) exhaust the formal completion, each \(B_e\)
is proper, and the compatible \(M_e\) have ample special fiber.
Grothendieck algebraization therefore gives a projective \(B/T\)
with an invertible sheaf \(M\) inducing these data. Full faithfulness
on proper schemes applied to the group operations and their identities
gives the group law on \(B\). Smoothness and geometric connectedness
yield the abelian scheme. The Raynaud extension \(G^\natural\) is an
extension of this \(B\) by the split torus. Its formal completion is
identified with the formal completion of \(G\). This accounts for the
presence of a torus at every infinitesimal level even though \(G_K\)
is abelian.
\footnote{\href{https://www.kwlan.org/articles/cpt-PEL-type-thesis-revision-dense.pdf}{K.-W.~Lan,
\emph{Arithmetic Compactifications of PEL-Type Shimura Varieties},
\S3.3.3, equations (3.3.3.1)--(3.3.3.2) and Proposition 3.3.3.6};
\href{https://stacks.math.columbia.edu/tag/089A}{Grothendieck
algebraization, Stacks Project, Theorem 30.28.4}.}

In the usual degeneration datum \(X\) is the character group, \(Y\)
the lattice of periods, and \(\phi:Y\to X\) the map attached to the
polarization. Lan's degree formula is
\[
  \deg\lambda_K=[X:\phi Y]^2\,\deg\lambda_B.
\]
A principal polarization on \(G_K\) therefore identifies \(Y=X\)
through \(\phi\), and \(\lambda_B\) identifies \(B\) with its dual.
This explains why one and the same lattice and a map \(c:X\to B(T)\)
occur in the local jet formulas. More intrinsically, \(c\) corresponds
under this principal identification to the preceding homomorphism
from weights into \(B^\vee(T)\). Its value is determined by the
rigidified weight bundle, so tensor product gives
\(c(x+x')=c(x)+c(x')\).
\footnote{\href{https://www.kwlan.org/articles/cpt-PEL-type-thesis-revision-dense.pdf}{Lan,
Lemma 3.4.4.3 and Theorem 4.2.1.14}. The hypotheses on the base in
\S4.1 hold with the complete normal noetherian domain \(T\) and its
radical maximal ideal. The model is semiabelian, with split special
torus and abelian generic fiber. The rigidified symmetric cubical line
is ample by Lemma 4.2.1.6. Its formal descent is the compatible system
of descended lines above, as required by Definition 4.2.1.1.}

\subsection{The valuation of the period factors}

Write \(f_x\) for the Fourier coefficients of a nonzero generator of
the level-one theta space. In the principal case the space is
one-dimensional. Every coefficient map is nonzero, and translating
the \(x\)-coefficient through a period \(y\) gives
\[
 f_{x+y}=\psi(y)\tau(y,x)\,t_{c(y)}^* f_x .
\]
The equality uses the canonical isomorphisms of the corresponding
line fibers, with the scalar convention of the displayed equation.
Lan proves nonvanishing and this formula in
\S\S4.3.2--4.3.3. In the release the nonvanishing step uses the
degree-two addition formula and torus-torsion weights. Its support is
a coset of a subgroup of odd index; the exclusion of prime-index
subgroup obstructions makes that subgroup the full lattice.

Let \(Q(y)=v_T(\psi(y))\), evaluated relative to an integral generator
of its one-dimensional line. Any two such generators differ by a
unit, so \(Q(y)\in\mathbb Z\) is unambiguous. Write \(B(y,x)\) for
\(v_T(\tau(y,x))\) relative to the integral biextension. Comparison of
the two orders of translating \(f_0\) gives
\[
  Q(y+z)=Q(y)+Q(z)+B(y,z),\qquad B(y,z)=B(z,y).
\]
The biextension law gives additivity of \(B\) in each argument.
Rigidification gives \(Q(0)=0\). The even cubical normalization of
the chosen line gives \(Q(-y)=Q(y)\). Taking \(z=-y\) yields
\(B(y,y)=2Q(y)\); by induction \(Q(ny)=n^2Q(y)\) for \(n\in\mathbb Z\).

Here is the convergence input giving strict positivity. Since \(f_0\)
is a nonzero integral section on the proper abelian quotient, some
germ has order at most \(m\) for a fixed \(m\). Indeed a section
vanishing at every finite level is zero by formal separation.
Translation by the integral point \(c(y)\) moves this witness to a
germ where \(t_{c(y)}^*f_0\) has order at most the same \(m\). At any
fixed finite level a section on the torus torsor has only finitely
many weights. On a finite affine cover trivializing the torsor this
is finite Laurent support, and one takes the union of the supports.
Consequently, for fixed \(n\), away from finitely many \(y\) all germs
of \(f_y\) belong to \(\pi^{m+n}\). In integral frames the preceding
equation at \(x=0\) now reads \(b f_y=a h\) up to a unit, with
\(\psi(y)=a/b\) and \(h=t_{c(y)}^*f_0\) at its witness.
If \(v_T(a/b)<n\), write \(a/b=u\pi^j\) with \(u\in T^\times\)
and \(j<n\). The local ring is \(T\)-flat, so multiplication by \(\pi\)
is injective. Clearing denominators in \(f_y=u\pi^j h\) then gives
\(h\in\pi^{m+n-j}\), a contradiction. Thus \(\{y:Q(y)<n\}\) is
finite. If \(0\ne y\) had \(Q(y)\le0\), the distinct multiples
\(ky\), \(k\ge1\), would all belong to \(\{Q<1\}\).
This proves \(Q(y)>0\) for nonzero lattice vectors. The integral
symmetric matrix of \(B\) is positive definite on the real span:
a negative direction would contain rational approximants of negative
square, and a singular rational matrix would have a rational kernel
vector. A finite extension of complete discrete valuation rings of
ramification index \(e\) sends \(Q,B\) to \(eQ,eB\) for
integer-normalized orders. This statement precedes the subsequent
use of a divisible value group on \(\overline K\).
\footnote{\href{https://www.kwlan.org/articles/cpt-PEL-type-thesis-revision-dense.pdf}{Lan,
equations (4.3.1.2), (4.3.1.7), Proposition 4.3.1.9, and
\S4.2.4}. In the release see
\srcid{ExistsDegenerationDataFunctionalEquationDegreeOn},
\srcid{PsiTendstoZero}, \srcid{CubeRelScalarsMemMaximalIdeal}
and \srcid{EvenQuadraticPosOfFiniteNonpos} in
\srcid{UniformizationProofs}.}

The polarization and the integral action of \(\mathcal O_E\) also
give the balance identity required in Section~\ref{sec:period-basis}.
Write \(\mathcal P_B\) for the Poincaré biextension. On \(B\) the
Rosati condition is \(a^\vee\lambda_B=\lambda_B\bar a\). The
defining property of the dual morphism identifies
\((a\times1)^*\mathcal P_B\) with \((1\times a^\vee)^*\mathcal P_B\).
Under this identification the induced actions on the period
trivialization satisfy
\[
        (a\times1)^*\tau=(1\times\bar a)^*\tau
                \quad\text{on }Y\times X .
\]
This is the endomorphism part of the degeneration equivalence:
the integral action on the generic polarized abelian variety
corresponds to actions on its abelian part, its two lattices and its
period map. The left action on \(X\) is obtained from its natural
right action by conjugation. The displayed equality is an equality
in integral line fibers after extending scalars to \(K\); the
canonical identification of their integral lattices comes from the
Poincaré biextension over \(T\). In independent integral frames its
coefficient is a unit. Taking valuations thus gives
\(B(ay,x)=B(y,\bar a x)\), with no scalar error. This also gives the
conjugation in the torus action used below.
\footnote{\href{https://www.kwlan.org/articles/cpt-PEL-type-thesis-revision-dense.pdf}{Lan,
Definition 5.1.1.3(4) and Theorem 5.1.1.4}. The release takes the
valuation of an origin-normalized Rosati generator in
\href{https://github.com/ChristianKudeba/alpoge-rh-verification/blob/alpoge-source-7cbf7a0c/src/ReciprocityTheory/ReductiveGroups/ArithmeticGeometry/DiscreteValuationRings/ExistsRosatiGen.lean}{\srcid{bOf_balance}}.}

\subsection{The rank bound from the PEL determinant condition}

Let \(d=[E:\mathbb Q]\), \(r=\operatorname{rank}_{\mathbb Z}X\), and
\(s=\dim_E(X\otimes_{\mathbb Z}\mathbb Q)\). For a generic point of
the moduli datum of signature \((p_0,q_0)\) we prove
\[
   r=ds,\qquad s\le\min(p_0,q_0).
\]
The first equality follows by restriction of scalars.
Put \(F=\overline K\), and fix an embedding
\(\sigma:\mathcal O_E\to F\). It suffices to exhibit \(s\) independent
\(\sigma\)-eigenvectors in \(\operatorname{Lie}(G_F)\).
We use row vectors and let \(M(a)\) be the integral matrix of \(a\)
on the Lie algebra of \(G\), in a basis over \(T\). The toric Lie
algebra at level \(e\) consists of the additive forms on \(X\) with
values in \(T/\pi^{e+1}\). This description follows by evaluating
a character on dual numbers as \(1+\varepsilon t\). In a
\(\mathbb Z\)-basis of \(X\), define \(A(\bar a)\) to be the matrix
of \(\bar a\) on \(X\) with column coordinates. The conjugation is
the source's convention for the toric part of the endomorphism.

The kernel embedding on level \(e\) induces an injection on tangent
vectors. Lift its matrix to an \(r\times g\) matrix \(C\) over
\(T\), where \(g=\dim G_K\). Equivariance and injectivity give
\[
   CM(a)\equiv A(\bar a)C\pmod{\pi^{e+1}},
   \qquad CD=I_r
\]
for some \(g\times r\) matrix \(D\) over \(T\). For the second
assertion, a residue row in the kernel of \(C\bmod\pi\) lifts to
\(w\) with \(wC\in\pi T^g\). Then \(\pi^ew\) is in the level-\(e\)
kernel, so \(\pi^ew\in\pi^{e+1}T^r\) and the original residue row
is zero. Thus \(C\) has an invertible \(r\times r\) minor over \(T\),
which gives a right inverse. Neither the lifts \(C\) nor their right
inverses need to be chosen compatibly as \(e\) varies.

Let \(V\) be the integral closure of \(T\) in \(F\).
It is a domain with fraction field \(F\). Its element \(\pi\) is a
nonunit. Every \(0\ne u\in V\) divides some power of \(\pi\) in \(V\):
a monic integral equation for \(u\) can be chosen with nonzero
constant term, giving \(u\mid b\) for some \(0\ne b\in T\);
then \(b\) is a unit times a power of \(\pi\). Also
\(\sigma(\mathcal O_E)\subset V\). In an \(E\)-basis of \(X_{\mathbb Q}\),
applying \(\sigma\circ\text{conjugation}\) to each coordinate gives \(s\)
independent additive forms with
\(\varphi(\bar a x)=\sigma(a)\varphi(x)\). Clearing a common
denominator on the integral basis of \(X\) gives an \(s\times r\)
matrix \(\Phi\) over \(V\) of full row rank satisfying
\(\Phi A(\bar a)=\sigma(a)\Phi\). Form the matrix \(H\) by adjoining
horizontally the blocks \(M(a)-\sigma(a)I_g\) for integral additive
generators \(a\) of \(\mathcal O_E\). At every level the previous
congruence gives matrices over \(V\) such that
\[
    CD=I_r,\qquad \Phi C H=\pi^{e+1}Z .
\]

Over \(F\) choose a generalized inverse of \(H\) by inverting its
restriction to a complement of its kernel on its image. Choose a
right inverse to \(\Phi\). Clearing denominators gives integral
matrices \(P,\Psi\) and \(0\ne b,b'\in V\) with
\(HPH=bH\) and \(\Phi\Psi=b'I_s\). Take \(e\) with
\(\pi^e=bb' u\), and use matrices \(C,D,Z\) for this \(e\). Define
\[
   W=b\Phi C-\Phi CH P,\qquad N=D\Psi .
\]
Then
\[
  WH=0,\qquad WN=bb'\bigl(I_s-\pi u ZP D\Psi\bigr).
\]
The determinant in parentheses is \(1\bmod \pi\). Since
\(V/(\pi)\ne0\), it is nonzero. Thus the \(s\) rows of \(W\) are
independent over \(F\). The identity \(WH=0\) puts each row in the
\(\sigma\)-eigenspace for all additive generators, hence for every
element of \(\mathcal O_E\). Finally choose \(\sigma\) and its
conjugate. The distinguished imaginary generator of the host maps
under these embeddings to the two opposite nonzero roots in the
determinant condition. Their eigenspaces have dimensions \(p_0,q_0\)
in some order. Applying the bound to both gives the assertion.
\footnote{The source proof is
\href{https://github.com/ChristianKudeba/alpoge-rh-verification/blob/alpoge-source-7cbf7a0c/src/Development/AbelianVarietyTheory/Proofs/UnitaryExistence/ExistsMulEqSmulOne.lean}{\srcid{ExistsMulEqSmulOne}},
especially \srcid{exists_matrix_level}, \srcid{exists_solutions},
\srcid{finrank_le_of_lieMult}, and \srcid{rank_le_signature}.
The argument uses the actual \(\mathcal O_E\)-equivariance on the
formal trait and the generic moduli determinant condition. The
trace-normalized pullback of the theta line and its integral torsion
identification are additional structures used in the jet calculation
of Sections~\ref{sec:theta-lift-valuation}
and~\ref{sec:local-theta-jets}.}

\section{A short integral basis for the period lattice}
\label{sec:period-basis}

We prove the numerical lattice estimate used in the jet comparison.
Let \(E\) be the CM host field, \(d=[E:\mathbb Q]\),
\(\Delta=|\operatorname{disc}E|\), and \(R=\Delta^{2/d}/(2d)\).
Let \(Y\) be the period lattice, an \(\mathcal O_E\)-module free
of finite rank \(\rho\) over \(\mathbb Z\). Write
\(B:Y\times Y\to\mathbb R\) for the positive definite symmetric
form supplied by the valuation of the period pairing. It is balanced
in the sense that \(B(sy,z)=B(y,\bar s z)\).
With the source's normalization \(Q(y)=B(y,y)/2\). Let
\(\mathcal S\) be the real vector space of balanced symmetric forms.
Let \(\phi(B)\) be the infimum of \(\lambda(B)\) over real linear
functionals on \(\mathcal S\) integral on the integral forms and
strictly positive on every nonzero positive semidefinite form.
The chart identifies this infimum with \(\phi_T\).
Such functionals exist: for an integral basis \(v_i\) of \(Y\) and
the host element \(w_0\in\mathcal O_E\) with \(w_0+\bar w_0=1\),
\(\lambda(C)=\sum_i C(w_0v_i,v_i)=\tfrac12\sum_i C(v_i,v_i)\)
has both properties. If the last sum vanishes for a positive
semidefinite \(C\), the whole basis lies in its radical.

There is a \(\mathbb Z\)-basis \(b_i\) of \(Y\) satisfying
\begin{equation}
  Q(b_i)\le R\rho^3\phi(B),\qquad
  |B(b_i,z)|\le 2\sqrt{\rho}\,Q(z)\quad(z\in Y).
  \label{eq:period-short-basis}
\end{equation}
The rank-zero case is empty. Assume \(\rho>0\).

\subsection{The dual minimum and the factor from complex multiplication}

Fix an admissible \(\lambda\), a basis \(v_j\) of \(Y\), and its
real dual coordinates \(v_j^*\). There are \(y_s,z_s\in Y\) with
\(\lambda(C)=\sum_s C(y_s,z_s)\) for every \(C\in\mathcal S\).
Indeed quotient \(Y\otimes_{\mathbb Z}Y\) by the common integral
annihilator of all the forms in \(\mathcal S\). The quotient is
torsion-free and finite free. The defining symmetry and balancing
relations have integral coefficients, so the real dual of this quotient
is precisely \(\mathcal S\): the common annihilator is the saturation
of the integral span of those relations. Integral forms are precisely
the integral dual. A functional integral on that dual is thus in the
quotient lattice, where its representative is a finite integral sum of
pure tensors. Absorb integral coefficients in the first factor.

In \(Y_{\mathbb C}\) let \(p_\sigma\) denote the idempotent of
\(E\otimes_{\mathbb Q}\mathbb C=\prod_\sigma\mathbb C\) corresponding
to \(\sigma\); write \(\eta_\sigma(y)=\eta_{\mathbb C}(p_\sigma y)\)
for a real covector \(\eta\). Define
\[
 C_\eta(y,z)=\sum_\sigma
       \operatorname{Re}(\eta_\sigma(y)\overline{\eta_\sigma(z)}),
       \qquad L[\eta]=\lambda(C_\eta).
\]
The form \(C_\eta\) is balanced and positive semidefinite, and
is nonzero when \(\eta\ne0\), by \(\sum_\sigma p_\sigma=1\).
Thus \(L\) is a positive definite quadratic form on the real dual.
In the chosen coordinates let \(L,G\) denote the matrices of this
quadratic form and of \(B\), respectively. Then
\(\operatorname{tr}(LG)=\lambda(B)\). Indeed polarize the
definition of \(C_\eta\), multiply its value on \(v_i^*,v_j^*\)
by \(B(v_i,v_j)\), and sum over \(i,j\).
The result at \(y,z\) before applying \(\lambda\) is
\(\sum_\sigma \operatorname{Re} B_{\mathbb C}
(p_\sigma y,\overline{p_\sigma z})=B(y,z)\).
The last equality follows from balancing, which pairs only
the \(\sigma\) and \(\bar\sigma\) components.

For an integral covector \(0\ne\eta\in Y^*\), let \(t_y\in E\)
represent the rational linear functional
\(c\mapsto\eta(cy)\) under the trace pairing. Then
\(t_y\in\mathfrak D_E^{-1}\) because this pairing with
\(\mathcal O_E\) is integral. In the idempotent coordinates
\(\eta_\sigma(y)=\sigma(t_y)\), by
\(\eta(cy)=\sum_\sigma \sigma(c)\sigma(t_y)\).
Some \(t_y\ne0\), and all of its embeddings are nonzero.
Set
\[
 t=\sum_s(t_{y_s}\overline{t_{z_s}}+
                     \overline{t_{y_s}}t_{z_s})
            \in\mathfrak D_E^{-2}.
\]
For each \(\sigma\), \(\sigma(t)\) is twice the value of \(\lambda\)
on \((y,z)\mapsto\operatorname{Re}
(\eta_\sigma(y)\overline{\eta_\sigma(z)})\), which is a nonzero
positive semidefinite balanced form. Thus every \(\sigma(t)>0\).
As \((t)\mathfrak D_E^2\) is a nonzero integral ideal,
\(\prod_\sigma \sigma(t)=|N_{E/\mathbb Q}t|\ge\Delta^{-2}\).
The arithmetic--geometric mean inequality gives
\begin{equation}
       L[\eta]=\tfrac12\sum_\sigma\sigma(t)
          \ge \tfrac d2\Delta^{-2/d}=:\mu^2 .
      \label{eq:period-dual-min}
\end{equation}

Take nonzero \(\theta\in\mathcal O_E\) with \(\bar\theta=-\theta\),
by subtracting the conjugate of a nonreal integral element.
Its action matrix \(K\) on \(Y\) is invertible over \(\mathbb R\).
Indeed \(Y_{\mathbb Q}\) is an \(E\)-module. Balancing gives
\(K^tG=-GK\); the polarized expression for \(L\), with
\(\overline{\sigma(\theta)}=-\sigma(\theta)\), gives \(KL=-LK^t\).
Write \(L=S^2\) with \(S\) positive definite symmetric, and put
\(D=SGS,\ J=S^{-1}KS\). Then \(J^t=-J\) and \(JD=DJ\).
Every eigenspace of \(D\) is preserved by the invertible
skew-symmetric \(J\); in particular its top eigenspace has even
dimension. As \(D\) is positive definite, its largest eigenvalue
is at most \(\tfrac12\operatorname{tr}D\).
For each column \(x\in\mathbb R^\rho\) this gives
\[
       x^tGx\le\tfrac12\lambda(B)\,\|S^{-1}x\|^2 .
\]

The Euclidean successive-minima transference inequality
\(\lambda_1(\Lambda)\lambda_\rho(\Lambda^*)\le\rho\)
applied to the lattice \(\Lambda=S\mathbb Z^\rho\) and its
reciprocal \(S^{-1}\mathbb Z^\rho\) now gives \(\rho\) independent
\(c_i\in\mathbb Z^\rho\) with
\(\|S^{-1}c_i\|\le\rho/\mu\), by \eqref{eq:period-dual-min}.
Discrete full lattices attain their successive minima.
In the \(B\)-norm, their associated lattice vectors therefore have
\[
    Q(c_i)\le \frac14\lambda(B)\,\rho^2/\mu^2
            =R\rho^2\lambda(B).
\]
The least possible maximum of \(Q\) on \(\rho\) independent
lattice vectors is attained, since every bounded \(B\)-ball in
\(Y\) is finite. The same minimum is thus at most
\(R\rho^2\phi(B)\) after taking the infimum over \(\lambda\).
\footnote{\href{https://doi.org/10.1007/BF01445125}{W. Banaszczyk,
New bounds in some transference theorems in the geometry of numbers,
Mathematische Annalen 296 (1993), 625--635}.
The archive gives a Gaussian Poisson proof of the transference input
in \href{https://github.com/ChristianKudeba/alpoge-rh-verification/blob/alpoge-source-7cbf7a0c/src/Formalism/PeriodIntegrals/OrbitalBoundsConclusion/TransferenceMinima.lean}{its transference theorem}.}

\subsection{Converting independent vectors into the required basis}

Take successive shortest vectors \(u_i\) outside the subspaces
\(V_{i-1}=\langle u_1,\ldots,u_{i-1}\rangle_{\mathbb R}\), and put
\(V_i=V_{i-1}+\mathbb R u_i\). Such vectors exist by discreteness.
The preceding independent tuple shows that each
\(\|u_i\|_B^2/2\le R\rho^2\phi(B)\), and the norms are nondecreasing.
The quotient \((Y\cap V_i)/(Y\cap V_{i-1})\) is free cyclic:
its denominator is saturated, and the ranks of the numerator and
denominator are \(i,i-1\). Choose a lift \(b_i\) of the positive
generator with smallest norm. By induction these lifts form an
integral basis on every step of the flag.

Let \(u_j^*\) be the Gram--Schmidt vectors for the \(u_j\),
with the same \(B\)-metric. The perpendicular component of \(b_i\)
along \(u_i^*\) is a positive multiple at most one of \(u_i^*\),
since the nonzero class of \(u_i\) is an integral multiple of the
generator. Subtracting integer multiples of
\(u_{i-1},\ldots,u_1\) successively reduces the earlier
Gram--Schmidt coordinates in the same lift to absolute value
at most \(1/2\). The minimum norm choice and orthogonality imply
\[
 \|b_i\|_B^2\le\|u_i^*\|_B^2+
       \tfrac14\sum_{j<i}\|u_j^*\|_B^2
       \le \rho\|u_i\|_B^2 .
\]
This proves the first inequality of \eqref{eq:period-short-basis}.
For \(z\in Y\cap V_{i-1}\), the lifts \(b_i+z,b_i-z\) have
norm at least \(\|b_i\|_B\). Expanding the two squares gives
\(2|B(b_i,z)|\le\|z\|_B^2\). For \(z\notin V_{i-1}\), minimality
gives \(\|u_i\|_B\le\|z\|_B\), so Cauchy--Schwarz gives
\(|B(b_i,z)|\le\sqrt{\rho}\,\|z\|_B^2\).
These bounds imply the second inequality in both cases.

In applying \eqref{eq:period-short-basis} to
Section~\ref{sec:cusp-jets}, the geometric inputs are the period
pairing with diagonal \(2Q\), identification of its infimum with
the frame slope \(\phi_T\), and the toric rank inequality
\(\rho\le d\,q_-\), where \(q_-\) denotes the negative signature
rank. The two numerical basis parameters are therefore exactly
\(R\rho^3\phi_T\) and \(2\sqrt{\rho}\).
\footnote{\href{https://github.com/ChristianKudeba/alpoge-rh-verification/blob/alpoge-source-7cbf7a0c/src/Formalism/SharpRestriction/Divisibility/GeneratorPowers/PhiCTie.lean}{Integral functionals on balanced forms},
\href{https://github.com/ChristianKudeba/alpoge-rh-verification/blob/alpoge-source-7cbf7a0c/src/Formalism/SharpRestriction/Divisibility/GeneratorPowers/PhiCMinI.lean}{the codifferent lower bound},
\href{https://github.com/ChristianKudeba/alpoge-rh-verification/blob/alpoge-source-7cbf7a0c/src/Formalism/SharpRestriction/Divisibility/GeneratorPowers/PhiCMinK.lean}{the trace identity, CM symmetry and transfer},
\href{https://github.com/ChristianKudeba/alpoge-rh-verification/blob/alpoge-source-7cbf7a0c/src/Formalism/SharpRestriction/Divisibility/GeneratorPowers/PhiCMinima.lean}{minimization over the cone of forms},
\href{https://github.com/ChristianKudeba/alpoge-rh-verification/blob/alpoge-source-7cbf7a0c/src/ReciprocityTheory/ReductiveGroups/ArithmeticGeometry/FirstCohomology/ReducedBasis.lean}{basis reduction and the mixed estimate},
\href{https://github.com/ChristianKudeba/alpoge-rh-verification/blob/alpoge-source-7cbf7a0c/src/Development/AbelianVarietyTheory/Proofs/DivisorProofs/ExistsShortPeriodLatticeBasisBalancedSlope.lean}{specialization to the period lattice of the chart}.}

\section{Cusp jets on the covers used in the degree estimate}
\label{sec:cusp-jets}

The geometric estimate takes place on a fixed good cover in a family.
Fix a sufficiently small positive \(\epsilon\), take a family meeting
the largeness conditions for this \(\epsilon\), and let \(x_1\) be its
lower endpoint. There exists \(x_0\) such that the following statements
hold for every \(x>x_1\) with \(x\ge x_0\), every index \(\iota\)
above \(x\), every point satisfying the prescribed point and
reciprocity conditions, and every stage \(i\) with \(i+1<m\).
For each such stage they hold for every good cover \(Y\).
The bounds for \(x_0\) can depend on the family. In the conjunction
of two estimates the upper bound on \(\epsilon\) is the minimum of
their admissible upper bounds, the largeness bounds are their
componentwise maxima, and the eventual threshold is the maximum of
the two \(x_0\)'s. Thus one can apply all the estimates below at the
same \(x,\iota,i,Y\). In a later restriction to an arithmetic
window, the restriction must lie above this eventual threshold.
\footnote{\href{https://github.com/Deicyde/riemann-hypothesis/blob/40e56bab/src/Development/AutomorphicArithmetic/ThetaCorrespondence/ThetaPairLifts/ThirdPlaceSeries/LiftLambdaBound.lean\#L17-L32}{Quantifiers of the geometric cover condition},
\href{https://github.com/Deicyde/riemann-hypothesis/blob/40e56bab/src/Development/AutomorphicArithmetic/ThetaCorrespondence/ThetaPairLifts/ThirdPlaceSeries/ExistsCuspWithThetaJetValueLe.lean\#L24-L43}{simultaneous admissible range and eventual threshold}.}

\subsection{A theta rule of bounded valuation}

Let \(E\) be the CM host field of degree \(d=[E:\mathbb Q]\).
For a dominant array \(a\) of theta exponents on the datum, write
\(\delta=D_F(a)\) for its nonnegative degree as used in the jet
estimate. Let \(n\ge1\) be the number of theta columns, at least
the sum of twice the ranks of the indicated lattice summands. Put
\[
            R=\frac{|\operatorname{disc}E|^{2/d}}{2d}.
\]
On the characteristic-zero-residue integral chart of the proof,
let \(T\) be the complete discrete valuation ring, \(v_T\) its
integer-valued order with \(v_T(0)=+\infty\), and \(\phi_T\ge0\) its
given slope parameter. The chart has a truncated frame and a PEL
uniformization with the specified torsion isomorphism. The auxiliary
element \(\alpha\in\mathcal O_E\) satisfies
\(\alpha\overline\alpha=q\), with \(q\) dividing the level and the
level ideal contained in \((\overline\alpha)\).

Let \(N\) bound \(d\) times the negative signature rank of the
datum. Suppose the numerical data satisfy
\begin{equation}\label{eq:jets-chart-num}
    4\delta+4\le q,\qquad 2\delta\sqrt N<q,\qquad
    \frac{\delta^2 R N^3}{q}\le c .
\end{equation}
Then some theta rule \(f\) of type \((a,n)\) has nonzero evaluation
on this frame, with
\begin{equation}\label{eq:jets-value}
                   v_T(f)\le c\,\phi_T .
\end{equation}

The inputs to this inference are the characteristic-zero harmonic
polynomial with local invariant at most zero, the existence of the
power data and short basis on the chart, and the local comparison of
the infimum of the normalized theta valuations with the local
invariant. Fix a uniformizer and such a polynomial. The short
basis has parameter \(0\le\rho\le N\), so that
\(2\delta\sqrt{\rho}<q\). Apply the local comparison with valuation
normalization \(1\) and with these choices. The local invariant term
is nonpositive, leaving
\[
    \inf_f v_T(f)\le B,\qquad
    B=\frac{\delta^2 R\rho^3\phi_T}{q}
       \le \frac{\delta^2 RN^3}{q}\phi_T\le c\phi_T .
\]
The infimum runs over the rules obtained from this polynomial.
If every one of them had zero evaluation or order strictly larger
than \(c\phi_T\), each extended order would be at least
\(\lfloor B\rfloor+1\). The same would hold for the infimum,
contradicting \(\lfloor B\rfloor+1>B\). Thus a member with nonzero
evaluation satisfies \eqref{eq:jets-value}. The discreteness at this
step avoids any loss in passing from an infimum to a section.
\footnote{\href{https://github.com/Deicyde/riemann-hypothesis/blob/40e56bab/src/Development/AbelianVarietyTheory/Proofs/DivisorProofs/ReflectionSumJetNumber.lean\#L52-L98}{The short-basis, harmonic-polynomial and infimum inputs, followed by the argument with the integral valuation}.}

In the family application \(q=g_0g_1\), the product of the two host
primes, and \(N=2fm\) in the rank notation of the family. The
numerical jet condition at \(g_0\) supplies the first two inequalities
of \eqref{eq:jets-chart-num} with denominator bound \(g_0\), together
with
\[
       \frac{RN^3}{g_0}\left(\frac{rN}{4}+2\delta^2\right)\le c.
\]
Here \(r\ge0\) is the rank parameter in that condition. As
\(q\ge g_0>0\) and \(RN^3\ge0\),
\[
       \frac{\delta^2RN^3}{q}
       \le\frac{\delta^2 RN^3}{g_0}
       \le\frac{RN^3}{g_0}\left(\frac{rN}{4}+2\delta^2\right).
\]
The first two inequalities remain valid on replacing \(g_0\) by
\(q\). This proves the three bounds required on the chart.
\footnote{\href{https://github.com/Deicyde/riemann-hypothesis/blob/40e56bab/src/Development/AbelianVarietyTheory/Proofs/DivisorProofs/ReflectionSumJetNumber.lean\#L42-L50}{The degree comparison},
\href{https://github.com/Deicyde/riemann-hypothesis/blob/40e56bab/src/Development/AutomorphicArithmetic/ThetaCorrespondence/ThetaPairLifts/ThirdPlaceSeries/ExistsThetaRuleValueLeCoverCharZero.lean\#L47-L72}{specialization to the two host primes and the field degree}.}

\subsection{Descent to the cusp and the factor \(1/2500\)}

Take a component \(Z\), a toroidal flag chart, a boundary point
\(y\) outside its distinguished open locus, and a local map of the
stalk at \(y\) to a complete discrete valuation ring, compatible
with the given generic frame. The local map on this chart implies
that the residue field has characteristic zero. The chart-extension
theorem provides a valuation ring \(T'\), an extended frame admitting
the truncated uniformizing chart, and a positive integer \(e_T\).
The slope of the extended frame is \(e_T\phi_T\), and its extended
order on evaluations of a theta rule is \(e_T\) times the order on
the original frame. Use the canonical connection in this comparison.

On the extended frame choose \(f\) by \eqref{eq:jets-value}.
Its value there is nonzero and its order is at most
\(c e_T\phi_T\). If its value on the original frame were zero, the
comparison of extended orders would identify a finite number with
\(+\infty\). Thus its order there is finite as well. Divide the
order comparison by \(e_T>0\) to get the original bound \(c\phi_T\).
By dominance and positivity of the transformed theta weight, the
theta cuspidality theorem makes this a cusp form. Restriction to
the component gives a cusp form with the same evaluation because
the chosen point factors through \(Z\).
\footnote{\href{https://github.com/Deicyde/riemann-hypothesis/blob/40e56bab/src/Development/AutomorphicArithmetic/ThetaCorrespondence/ThetaPairLifts/ThirdPlaceSeries/ExistsCuspWithThetaJetValueLe.lean\#L44-L65}{Descent of nonvanishing and the valuation bound},
\href{https://github.com/Deicyde/riemann-hypothesis/blob/40e56bab/src/Development/AutomorphicArithmetic/ThetaCorrespondence/ThetaPairLifts/ThirdPlaceSeries/ThetaValueOnCuspNeighbourhood.lean\#L29-L57}{cuspidality and the equality after restriction to the component}.}

For the two weights in use, denote the weight of the quotient by
\(u_1\) and its companion upper weight by \(u_b\). The theta
construction first shifts each weight by its half-column weight.
There are \(n\) columns with \(8\mid n\), so \(4\mid n/2\).
The four-divisibility boundary comparison removes this shift and
preserves the nonzero value's order. The numerical estimate for
\(u_1\) follows from the family size inequalities. The upper-weight
degree comparison gives a theta degree for \(u_b\) no larger than
the one for \(u_1\), so the same numerical jet bound works for both.
Its coefficient is \(c=1/2500\). Extending the cusp form to the
compactified flag space under the extension isomorphism preserves its
evaluation. Thus for each of these two weights there is a section
\(w\) in the specified cusp subspace with
\[
            w(y')\ne0,\qquad v_T(w(y'))\le\phi_T/2500.
\]
Only the section is chosen after the component, boundary point and
valuative frame.
\footnote{\href{https://github.com/Deicyde/riemann-hypothesis/blob/40e56bab/src/Development/AutomorphicArithmetic/ThetaCorrespondence/ThetaPairLifts/ThirdPlaceSeries/JetsOfConstantSection.lean\#L26-L49}{The shift by the half-column weight and the boundary comparison},
\href{https://github.com/Deicyde/riemann-hypothesis/blob/40e56bab/src/Development/AutomorphicArithmetic/ThetaCorrespondence/ThetaPairLifts/ThirdPlaceSeries/JetNumUb.lean\#L23-L34}{comparison of the two degrees},
\href{https://github.com/Deicyde/riemann-hypothesis/blob/40e56bab/src/Development/AutomorphicArithmetic/ThetaCorrespondence/ThetaPairLifts/ThirdPlaceSeries/ExistsCuspValueLePhiOfHilbertSamuel.lean\#L24-L50}{the two weights, simultaneous threshold and final valuation}.}

\subsection{From boundary orders to the degree comparison}
Let \(\nu\) be a discrete valuation centered at a boundary point. Its
potential \(\phi\) is represented by a valuative frame and a real \(h>0\)
for which \(\phi=h\phi_T\). The additive order of a germ of nonzero
value is \(h\) times the order of its pullback to the valuation ring.
Choose the cusp section constructed above. A change of local
trivialization multiplies its representative by a unit, so
\[
  \operatorname{ord}_{\nu}(w)=h v_T(w(y'))\le h\phi_T/2500=\phi/2500 .
\]
An independent Hodge-section theorem gives a nonzero section \(\theta\)
of weight \(8(p+q)\) times the Hodge weight whose order is at least
\(8\phi/625\) wherever its germ has finite order. Here \(p,q\) are the
two signature ranks. Since \(\phi\ge0\), we obtain
\(3\operatorname{ord}_\nu(\theta)\ge8\operatorname{ord}_\nu(w)\).
In the multiplicative convention this gives
\(\nu(\theta)^3\le\nu(w)^8\). If \(\nu(\theta)=0\), the same inequality
is immediate. Combining this with base-point freeness on the open
locus and extension of cusp forms supplies the valuation-multiplicity
premise in the component-image degree comparison. That comparison
also uses compactification proportionality, the top Chern number
formula and birationality at the prescribed weight gaps.
\footnote{\href{https://github.com/Deicyde/riemann-hypothesis/blob/40e56bab/src/Development/AutomorphicArithmetic/ThetaCorrespondence/ThetaPairLifts/ThirdPlaceSeries/ExistsSectionOrderBoundHilbertSamuel.lean\#L22-L46}{Passing to germs};
\href{https://github.com/Deicyde/riemann-hypothesis/blob/40e56bab/src/Development/AutomorphicArithmetic/ThetaCorrespondence/ThetaPairLifts/ThirdPlaceSeries/ExistsNonzeroSectionWeightBound.lean\#L79-L82}{the Hodge section};
\href{https://github.com/Deicyde/riemann-hypothesis/blob/40e56bab/src/Development/AutomorphicArithmetic/ThetaCorrespondence/ThetaPairLifts/ThirdPlaceSeries/ValuationMultiplicityOfConsecutiveRun.lean\#L48-L97}{the powers three and eight};
\href{https://github.com/Deicyde/riemann-hypothesis/blob/40e56bab/src/Development/AutomorphicArithmetic/ThetaCorrespondence/ThetaPairLifts/ThirdPlaceSeries/CoveringDegreeMultiplicative.lean\#L25-L56}{application to the cover degree}.}

\subsection{The metric loss in the lift bound}
Write \(X,b\) for the size parameters with \(b=\exp L\). The Gram
condition identifies the squared norm on the integral cusp lattice with
the norm of the chosen frame. The Tamagawa measure comparison bounds
the maximal slope by
\[
 2\lambda_{\rm sup}+2\log C_{\rm sup}+
 [F:\mathbb Q]^{-1}\log|\operatorname{disc}F|,
\]
where \(F\) is the coefficient field. The three inputs are the mean
logarithmic sup-norm estimate, the logarithm of the supremum constant,
and the root-discriminant estimate. Let \(A\) be the size exponent.
The weight parameters obey \(0\le s\le8frb\) and \(0\le t\le b\).
Since \(fr,b\ge1\), this gives \(s+t+1\le10frb\). The logarithmic
estimates on these parameters and \(frL\le X^A\le b\) enter the
quantitative comparison for the three terms above, which bounds their
sum by \(X^4(X^A)^2b=X^{2A+4}b\). Thus the cusp-slope exponent is
\(A_1=2A+4\). The all-places measure comparison of constant \(C_m\)
and the host estimate of exponent \(A_h\) give the lift bound with
exponent \(A_1+(A_h+1)+\lceil2C_m\rceil\). Taking the maximum of
the cusp, measure, host, Gram and weight thresholds makes all five
inputs applicable at the same point.
\footnote{\href{https://github.com/Deicyde/riemann-hypothesis/blob/40e56bab/src/Development/AutomorphicArithmetic/ThetaCorrespondence/FockModels/TamagawaPeterssonLambdaSqNonneg.lean\#L596-L669}{Metric comparison};
\href{https://github.com/Deicyde/riemann-hypothesis/blob/40e56bab/src/Development/AutomorphicArithmetic/ThetaCorrespondence/ThetaPairLifts/ThirdPlaceSeries/CoveringLambdaInvariantBound.lean\#L21-L37}{all-places exponent};
\href{https://github.com/Deicyde/riemann-hypothesis/blob/40e56bab/src/Development/AutomorphicArithmetic/ThetaCorrespondence/ThetaPairLifts/ThirdPlaceSeries/TopCovolG.lean\#L50-L58}{degree and metric bounds used in the covolume argument}.}

\section{From a local harmonic value to a theta rule}
\label{sec:local-theta-jets}

Let \(T\) be the complete discrete valuation ring of an integral chart,
\(K\) its fraction field and \(\pi\) a uniformizer. Write
\(v_T(\pi)=1\), \(v_T(0)=+\infty\), and
\(v_e(u)=v_T(u)/e\) for a positive integer \(e\). The chart carries
a truncated frame \(y\), the canonical connection, and a polarized
PEL uniformization with the specified isomorphism on torsion.
Write \(M\) for the positive level and \(\mathfrak m_b\) for its
chosen level ideal. The level condition includes
\(\mathfrak m_b+\overline{\mathfrak m}_b=\mathcal O_E\), where \(E\)
is the CM host. On the period lattice \(X\) take the integral
quadratic form \(Q\) furnished by \(\pi\), with
\(B(u,z)=Q(u+z)-Q(u)-Q(z)\).
Assume a basis \(b_j\) satisfies
\[
  Q(b_j)\le\Phi,\qquad |B(b_j,z)|\le M_0 Q(z)
       \quad(z\in X),\qquad \Phi\ge0 .
\]
On a zero lattice one can replace \(M_0\) by \(\max(M_0,0)\).
Let \(P\) be a harmonic polynomial of row degrees \(a_{\tau,i}\)
in \(n\ge1\) columns. Write \(\delta=\sum_{\tau,i}a_{\tau,i}\).
Every monomial in its support has degree \(\delta\).

For an integral power series \(f\) on the \(n\) copies of the chart,
write \(\mathcal D_{P,k}f\in K\) for its constant-term covariant
differential with polynomial \(P\) and line-bundle power \(k\).
It is linear in \(f\) and factors through the quotient by the
\((\delta+1)\)-st power of the ideal generated by the chart
coordinates. The coefficient field of \(P\) maps to \(K\) by the
fixed structural embedding. The local quantity occurring in the
cusp-jet estimate is
\[
       \nu_{\rm loc}(C;e,M,P)=
          \inf_{f\in T[[\boldsymbol Z]]}v_e(\mathcal D_{P,M}f).
\]
For \(x=(x_1,\ldots,x_n)\) in the discrete level-label set
\(\mathrm{La}^n\), denote the unique corresponding theta rule
by \(F_{x,P}\). Its evaluation in the frame \(y\) is in \(K\).

\paragraph{Two local hypotheses.}
The comparison below applies in either of the following settings:
\begin{enumerate}
\item The frame is algebraic, \(q=M\), \(M\) is a unit of \(T\),
\(\alpha\overline\alpha=M\) for \(\alpha\in\mathcal O_E\), and
\(\mathfrak m_b=(\overline\alpha)\).
\item The residue field of \(T\) has characteristic zero,
\(q\) is a divisor of \(M\), \(\alpha\overline\alpha=q\), and
\(\mathfrak m_b\subseteq(\overline\alpha)\). This formulation admits
the integral frame on the horizontal valuation without an algebraicity
hypothesis on that frame. Every nonzero rational integer is a unit of
\(T\) in this case.
\end{enumerate}
In both cases impose \(4\delta+4\le q\) and \(\delta M_0<q\). Then
\begin{equation}\label{eq:local-theta-comp}
 \inf_{x\in\mathrm{La}^n}v_e\bigl(F_{x,P}(y)\bigr)
 \ \le\ \nu_{\rm loc}(C;e,M,P)+\frac{\delta^2\Phi}{qe}.
\end{equation}
Both orders here are nonarchimedean; the value \(+\infty\) records a
zero evaluation. Positive residue characteristic \(p\) is permitted in
the first case when \(p\nmid M\). The Hermitian norm and its integral at
a complex embedding enter the metric comparison of
Section~\ref{sec:cusp-jets}.
\footnote{\href{https://github.com/ChristianKudeba/alpoge-rh-verification/blob/alpoge-source-7cbf7a0c/src/Development/AbelianVarietyTheory/Proofs/DivisorProofs/IInfThRuleLeNuLocOf.lean}{Full-level theorem and algebraic frame};
\href{https://github.com/ChristianKudeba/alpoge-rh-verification/blob/alpoge-source-7cbf7a0c/src/Development/AbelianVarietyTheory/Proofs/DivisorProofs/IInfThetaRuleLeLocalValueCharZero.lean}{divisor-level theorem in residue characteristic zero}.}

\paragraph{Changing the power in residue characteristic zero.}
In an integral biextension gauge, write \(S\) for the sum, over the
\(n\) columns, of the scalar gauge series at the distinguished
endomorphism. Its constant coefficient is zero. The formula for
\(\mathcal D_{P,k}\) applies the fixed differential to
\(f\exp(kS)\). The exponential is defined over \(T\), since \(T\)
contains \(\mathbb Q\). Put \(f'=f\exp((M-q)S)\). Each coefficient
is a finite sum in \(T\), and \(f'\exp(qS)=f\exp(MS)\).
The gauge-independence identity for harmonic \(P\) carries this equality
of differential values back to the chosen integral chart. Thus
\(\mathcal D_{P,q}f'=\mathcal D_{P,M}f\). Existence of the integral
biextension and its normalized gauge is used here. At full level
one simply takes \(f'=f\).
\footnote{\href{https://github.com/ChristianKudeba/alpoge-rh-verification/blob/alpoge-source-7cbf7a0c/src/Development/AbelianVarietyTheory/Proofs/DivisorProofs/HarmonicValueIntegralGermFramePowerCharZero.lean}{Integral gauge, power change, and harmonic gauge independence}.}

\paragraph{The small characters and the quadratic cost.}
Use the product period lattice \(X^n\) and
\(Q^{(n)}(z)=\sum_i Q(z_i)\). Putting \(b_j\) in the \(i\)-th
position gives a basis with the same bound \(\Phi\). Its mixed term
with \(z\) is \(B(b_j,z_i)\), bounded in absolute value by
\(\max(M_0,0)Q^{(n)}(z)\). In the following calculation \(Q,B,b_j\)
refer to this product and \(M_0\) to the resulting nonnegative bound.
Take \(\chi=\sum_j m_j b_j\), \(m_j\ge0\), \(\sum_jm_j\le\delta\).
The discriminant inequality for the positive quadratic form gives
\(\lvert B(s,t)\rvert\le 2\sqrt{Q(s)Q(t)}\): first use positivity
on the rational line \(s+ut\) by clearing denominators, then use
continuity in \(u\). Consequently
\[
 \sqrt{Q(\chi)}\le\sum_j m_j\sqrt{Q(b_j)}\le\delta\sqrt{\Phi},
 \qquad |B(\chi,z)|\le\delta M_0 Q(z).
\]
For \(z\ne0\), \(Q(z)\ge1\), whence
\[
 Q(\chi+qz)-Q(\chi)=q^2Q(z)+qB(\chi,z)
       \ge q(q-\delta M_0)Q(z)>0.
\]
Thus \(\chi\) is the unique minimum in \(\chi+qX^n\). This
strictness selects a single coefficient after reduction modulo \(\pi\).

Here are the geometric steps that turn these characters into all
jets. Let \(J=T[\boldsymbol Z]/(\boldsymbol Z)^{\delta+1}\).
Start on the abelian quotient of the semiabelian fiber. Write
\(q=4\delta+r\), \(r\ge4\), and fix a character \(\chi\).
The source's twisted-section lemma supplies a section
\(\beta_0\) of the \(\chi\)-twisted power \(r\) whose value in the
closed identity fiber is a unit. In degree four the very ample
sections \(s_i\) give an affine chart \(s_0\ne0\) at that point;
the ratios \(t_i=s_i/s_0\) generate the coordinate algebra of the
chart and their images generate the abelian quotient's jet algebra.
The latter is generated as a \(T\)-module by products of at most
\(\delta\) ratios. Indeed, write each ratio as its augmentation plus
an element of the augmentation ideal. All products of \(\delta+1\)
elements of that ideal vanish. Moreover \(t_0=1\), so shorter
products can be padded to length \(\delta\).
The factors \(s_{i_1}\cdots s_{i_\delta}\beta_0\) have precisely
the required power \(q\) and twist \(\chi\); in the local frame
their jets are the corresponding products of ratios times the fixed
unit \(s_0^\delta\beta_0\). Every abelian quotient jet therefore
occurs as a coefficient section \(\beta\) at \(\chi\).

Extend \(\beta\) by the functional equation of the degeneration in
the coset \(\chi+qX^n\), setting the other cosets to zero.
The coefficient valuation change on translation is
\((Q(\chi+qz)-Q(\chi))/q\). Thus the initial coefficient is
integral and the others in its coset are divisible by \(\pi\):
the change is positive and the coefficients are defined over the
discretely valued base. The integral expansion theorem realizes
this family as a section. Modulo \(\pi\) its jet is the abelian
coefficient times the torus character for \(\chi\), up to a fixed
unit from the frame. To see why these products span, the special
fiber torus supplies augmentation parameters of the form
\(v(b_j)/v(b_j)(0)-1\), where \(v(b_j)\) denotes the character
function in a trivialization. Together with the jets pulled back
from the abelian quotient these generate the augmentation ideal.
This is the tangent-space assertion for the semiabelian extension:
a tangent vector killed by the quotient differential is a torus
tangent, and the basis characters detect it. Passing successively
through the powers of the nilpotent augmentation ideal proves
generation of the jet algebra. Multiplying character functions
corresponding to \(b_j\)'s adds their characters up to a unit
pulled back from the abelian quotient. Expanding monomials of
degree at most \(\delta\) now produces the products indexed by the
\(\chi\)'s above. The fixed common frame unit preserves their span.
Nakayama's lemma lifts spanning modulo \(\pi\) to spanning over \(T\),
since \(J\) is finite free. A finite collection of the resulting
sections \(\sigma_i\) suffices. Each has expansion \(\theta_i\)
supported in the coset of one \(\chi_i\), with
\(Q(\chi_i)\le\delta^2\Phi\).
\footnote{\href{https://github.com/ChristianKudeba/alpoge-rh-verification/blob/alpoge-source-7cbf7a0c/src/Development/AbelianVarietyTheory/Proofs/JetHarmonics/SmallCombinationStrictMinMod.lean}{Strict coset minimum};
\href{https://github.com/ChristianKudeba/alpoge-rh-verification/blob/alpoge-source-7cbf7a0c/src/Development/AbelianVarietyTheory/Proofs/JetHarmonics/MemBasicOpenOfIsUnit.lean}{twisted abelian jets from degree-four sections};
\href{https://github.com/ChristianKudeba/alpoge-rh-verification/blob/alpoge-source-7cbf7a0c/src/Development/AbelianVarietyTheory/Proofs/JetHarmonics/IntegralFamilyOfStrictMin.lean}{integral extension in a single coset};
\href{https://github.com/ChristianKudeba/alpoge-rh-verification/blob/alpoge-source-7cbf7a0c/src/Development/AbelianVarietyTheory/Proofs/JetHarmonics/SmallCharactersSpanJets.lean}{special-fiber character spanning and integral lifting}.}

\paragraph{A scalar paying for all the characters at once.}
Extend \(v_T\) to the algebraic closure \(\overline K\), normalized
by \(v(\pi)=1\). In a frame of the \(q\)-th power on the truncated
chart, write \(f'\bmod(\boldsymbol Z)^{\delta+1}=\sum_i c_i
\operatorname{jet}(\sigma_i)\) with \(c_i\in T\).
The valuation condition on an expansion at level \(q\) and lower
bound zero means that the coefficient at \(\gamma\) has section
valuation at least \(Q(\gamma)/q\). For \(\theta_i\) the same
inequality holds with the lower bound decreased by \(Q(\chi_i)/q\):
it holds at the integral coefficient \(\chi_i\), propagates to its
coset by the functional equation, and holds at every zero coefficient.
Put \(k=\lfloor\delta^2\Phi\rfloor\). As \(Q(\chi_i)\) is an integer,
\(Q(\chi_i)\le k\) for every \(i\). Choose \(\lambda\in\overline K\)
with \(\lambda^q=\pi^k\). Then
\[
     0\le v(\lambda)=k/q\le\delta^2\Phi/q .
\]
The section \(\sigma=\sum_i\lambda c_i\sigma_i\) has expansion
\(s=\sum_i\lambda c_i\theta_i\) with lower bound zero. The
coefficients \(c_i\) are integral, and multiplying by \(\lambda\)
compensates every loss \(Q(\chi_i)/q\).
In the indicated frame its germ \(g\) satisfies
\(g\equiv\lambda f'\pmod{(\boldsymbol Z)^{\delta+1}}\).
Only one scalar is needed, regardless of the number of sections.
\footnote{\href{https://github.com/ChristianKudeba/alpoge-rh-verification/blob/alpoge-source-7cbf7a0c/src/Development/AbelianVarietyTheory/Proofs/JetHarmonics/ExistsPowFrame.lean}{The uniformized jet construction and floor of the quadratic bound};
\href{https://github.com/ChristianKudeba/alpoge-rh-verification/blob/alpoge-source-7cbf7a0c/src/Development/AbelianVarietyTheory/Proofs/JetHarmonics/ExistsGenOfTop.lean}{valuation propagation from a coset representative}.}

\paragraph{Selecting a translate.}
The containment of the level ideal and its conjugate implies
\((\alpha)+(\overline\alpha)=\mathcal O_E\). The invertibility
condition used by the theta-group argument is \(q\in T^\times\).
Choose the product theta frame and its pullback by \(\alpha\).
The integral spanning theorem in Section~\ref{sec:theta-translate-span}
expresses \(s=\sum_j d_j t_j\) with \(v(d_j)\ge0\). Each \(t_j\)
is a normalized translate of the pullback, at a point whose
coordinates are killed by \(\overline\alpha\).
The expansion bijection gives the section represented by each
\(t_j\); let \(h_j\) be its germ. Uniqueness and linearity of
expansion and germs give \(g=\sum_j d_jh_j\).
Write \(a_0=\mathcal D_{P,M}f\) and suppose \(a_0\ne0\).
Since the differential uses only the \(\delta\)-jet, after base
change to \(\overline K\) we have
\[
     \lambda a_0=\mathcal D_{P,q}g
                 =\sum_j d_j\mathcal D_{P,q}h_j .
\]
Some \(j\) therefore satisfies
\(v(\mathcal D_{P,q}h_j)\le v(a_0)+v(\lambda)\): if all summands had
larger order, so would the sum. The selected evaluation is nonzero.
\footnote{\href{https://github.com/ChristianKudeba/alpoge-rh-verification/blob/alpoge-source-7cbf7a0c/src/Development/AbelianVarietyTheory/Proofs/JetHarmonics/ExistsTermLe.lean}{Integral span of product translates and valuation selection};
\href{https://github.com/ChristianKudeba/alpoge-rh-verification/blob/alpoge-source-7cbf7a0c/src/Development/AbelianVarietyTheory/Proofs/DivisorProofs/JetsOnCopiesAny.lean}{product version for a general integral frame}.}

\paragraph{Identifying the value of the selected translate.}
The product-translation formula factors the selected germ, to order
\(\delta\), into a scalar \(u_0\) of valuation zero times the product
of the one-column germs. Each normalized one-column identification
contributes a root of unity \(u_i\); the normalized lifts have
finite order. Thus \(u=u_0\prod_i u_i\) has valuation zero.
Write \(M=qk'\). For a translated point \(b\) killed by
\(\overline\alpha\), the pullback by \(k'\overline\alpha\) takes the
translated power-\(q\) section to the power-\(M^2\) section defining
the level evaluation, with label representing \(\alpha b\).
The powers agree since \(q\,(k')^2\alpha\overline\alpha=M^2\).
The normalized pullbacks respect the level torsion trivializations;
applying the construction in each column gives an \(x\in\mathrm{La}^n\).
In the covariant evaluation the normalized chain rule for \(\beta=k'\overline\alpha\)
introduces, in a monomial of the given row degrees, the factor
\(N_a(\beta)\). The level definition divides the power-\(M^2\)
evaluation by \(M^\delta\). Since \(\alpha\beta=M\),
\(N_a(\alpha)N_a(\beta)=M^\delta\). Solving for the original
translated value gives the common factor
\[
    N_a(\alpha)=\prod_{\tau}\iota_\tau(\alpha)^{\sum_i a_{\tau,i}} .
\]
Here \(\iota_\tau\) denotes the corresponding embedding into the
coefficient field, viewed in \(\overline K\). Consequently the
chart-independent evaluation of the rule satisfies
\[
     \mathcal D_{P,q}h_j=u\,N_a(\alpha)\,F_{x,P}(y).
\]
Each \(\iota_\tau(\alpha)\) is integral over \(\mathbb Z\); an
element integral over \(\mathbb Z\) has nonnegative \(v\), by
applying the ultrametric inequality to its monic equation.
It follows that
\[
 v_T(F_{x,P}(y))\le v(\mathcal D_{P,q}h_j)
             \le v_T(a_0)+\delta^2\Phi/q .
\]
When \(a_0=0\) the comparison with this individual \(f\) holds in
the extended order as its right hand side is \(+\infty\). Divide by
\(e>0\) and take the infimum over all integral \(f\), obtaining
\eqref{eq:local-theta-comp}. The use in
Section~\ref{sec:cusp-jets} has \(e=1\).
\footnote{\href{https://github.com/ChristianKudeba/alpoge-rh-verification/blob/alpoge-source-7cbf7a0c/src/Development/AbelianVarietyTheory/Proofs/DivisorProofs/HarmonicValueTranslateDivisorLevelIntegralChart.lean}{Normalized pullback, the character factor, and its integrality};
\href{https://github.com/ChristianKudeba/alpoge-rh-verification/blob/alpoge-source-7cbf7a0c/src/Development/AbelianVarietyTheory/Proofs/JetHarmonics/HarmonicValueTranslateEqUnitMulNormaliserThetaVa.lean}{full-level version of the translate identification}.}

The comparison uses the polarized uniformization and torsion
identification on the prescribed chart, the normalized theta-lift and unit-ball properties in
Section~\ref{sec:theta-translate-span}, and the chart-independence
and pullback identities for the canonical connection. The preceding
calculation gives the interpolation cost and its transfer to a rule;
a derivation of those geometric inputs from the definitions is still
required for a standalone account. A finite place at which \(p\mid M\)
requires its own local argument since the integral spanning step above
uses that \(q\) is a unit.

\section{The connection on a uniformizing chart}
\label{sec:jet-connection}

Fix a framed point on the characteristic-zero fiber. In one copy of
a formal group chart let \(Z\) be the coordinates, \(F_0\) the
formal group law, \(A_a(Z)\) the action of \(a\in\mathcal O_E\),
and \(D_u\) the invariant derivation in tangent direction \(u\).
Thus \(D_u f(Z)\) is the coefficient of \(t\) in
\(f(F_0(Z,tu))\). All calculations with exponentials below take
place over a characteristic-zero field. The lattice on which
values are eventually bounded is the one in
Section~\ref{sec:local-theta-jets}.

\subsection{A local formula and its independence of gauge}

The rigidified cubical line provides multiplicative factors
\(\Phi_1,\Phi_2\) for the two partial addition laws of the
biextension, a symmetry factor, and an adjoint factor \(J_a\).
A canonical gauge is a series \(h(Z,W)\), zero on each axis, whose
partial additive coboundaries are the logarithms of the first two
factors. For instance,
\[
 \log\Phi_2(X,Y,Z)=h(X,Y)+h(X,Z)-h(X,F_0(Y,Z)).
\]
Its difference under interchanging the arguments is the logarithm
of the symmetry factor. The adjoint condition reads
\(\log J_a(X,Y)=h(A_aX,Y)-h(X,A_{\bar a}Y)\). With the element
\(w_0\) defining the primitive theta line, put
\(H_h(Z)=h(Z,A_{w_0}Z)\). If the germ in the primitive frame of a
section of the \(k\)-th power is \(f\), the chart evaluates a
monomial in the connection by
\[
      \left.D_{u_1}\cdots D_{u_j}\bigl(f\,\exp(kH_h)\bigr)
      \right|_{Z=0}.
\]
For multiple columns apply the corresponding derivatives in the
specified copies, with the product factor
\(\exp(k\sum_b H_h(Z_b))\).

We give the gauge construction to indicate where denominators enter.
Pass to the normalized formal logarithm of the commutative smooth
formal group. It is obtained by termwise integration of the
closed invariant differentials dual to the tangent basis. They are
closed since the invariant tangent fields commute, and their
constant matrix is the identity. The resulting coordinates turn
the group law into addition, because the pullback of each invariant
differential by addition is the sum of its two pullbacks. For a
symmetric, normalized additive cocycle \(C(x,y)\), allow inactive
coordinates as parameters, and put
\(g_j(x)=\partial_{y_j} C(x,0)\). Differentiating the cocycle equation
gives \(\partial_{y_j}C(x,y)=g_j(x+y)-g_j(y)\).
Symmetry, followed by taking the mixed partial derivatives, gives
\(\partial_i g_j=\partial_j g_i\). Choose the formal primitive \(h_1\)
of \(-\sum_jg_j\,dx_j\) with zero constant. It has zero linear part
in \(x\), and
\(C(x,y)=h_1(x)+h_1(y)-h_1(x+y)\), by differentiation in \(y\)
and evaluation at \(y=0\). Apply this to the first multiplication
law. After this change of frame, the interchange law for a
biextension makes the logarithm of the second defect additive in
its passive variable. Applying the same primitive in its active
variable preserves that additivity, since integration is linear in
the coefficients. This kills both defects.

The remaining endomorphism defects are biadditive. In their
\(\mathbb Q\)-vector space let \(l(a)\) act on the first argument
and \(r(a)\) by \(\bar a\) on the second. The logarithms \(\delta(a)\)
of the adjoint factors satisfy
\[
  \delta(ab)=l(b)\delta(a)+r(a)\delta(b).
\]
Indeed this is the logarithm of the composition law for the adjoint
trivializations, and \(l,r\) commute. In \(E\otimes_\mathbb Q E\)
take the separability element \(\sum_i x_i\otimes y_i\), characterized
by \(\sum_i x_i y_i=1\) and
\((a\otimes1)e=(1\otimes a)e\); it is the projector onto the
multiplication factor of the finite étale algebra \(E\otimes E\).
Extend \(l,r,\delta\) by rational linearity. Apply
\(x\otimes y\mapsto l(x)\delta(y)\) to the last identity to obtain
\(\delta(a)=l(a)v-r(a)v\), where \(v=\sum_i l(x_i)\delta(y_i)\).
Adding this biadditive correction kills the adjoint defects without
changing the multiplication laws. Finally, the residual logarithm
of the symmetry factor is alternating, biadditive, and commutes
with the adjoint actions. Subtracting half of it supplies the
last correction since \(h(Y,X)-h(X,Y)\) is the required symmetry
logarithm. This constructs \(h\).
\footnote{See \href{https://github.com/ChristianKudeba/alpoge-rh-verification/blob/alpoge-source-7cbf7a0c/src/Development/AbelianVarietyTheory/Proofs/FormalGroupCharts/CanonGaugeOfLog.lean}{\srcid{canonGauge_of_log}} and
\href{https://github.com/ChristianKudeba/alpoge-rh-verification/blob/alpoge-source-7cbf7a0c/src/Development/AbelianVarietyTheory/Proofs/FormalGroupCharts/DerivationInnerOfSeparabilityElement.lean}{the endomorphism correction}.}

For two choices \(h,h'\), their difference in logarithmic coordinates
is a bilinear form \(L\). In fact its two additive coboundaries
vanish, so differentiation over the characteristic-zero field makes
it linear in each argument. The adjoint law gives
\(L(au,v)=L(u,\bar a v)\). If \(u,v\) have eigencharacters
\(\tau,\sigma\) with \(\sigma\ne\bar\tau\), choose \(a\) with
\(\tau(a)\ne\sigma(\bar a)\). Subtraction gives \(L(u,v)=0\).
The diagonal change of gauge is a homogeneous quadratic polynomial
\(q(Z)=L(Z,dA_{w_0}Z)\) in logarithmic coordinates.

For a polynomial \(P\) in the derivative symbols, let
\(\langle P,S\rangle=P(\partial)S(0)\) in these coordinates.
Repeated differentiation of a monomial gives
\(\langle P,Z_iS\rangle=\langle\partial_iP,S\rangle\).
The letters in the source are tangent directions \(u_j\). Equivalently,
substitute \(Z=\sum t_j u_j\) in each column and use ordinary
directional differentiation in the \(t_j\); the same identity
applies even if the substitution is not an isomorphism. With
\(q_{\mathrm{tot}}=\sum_b q(Z_b)\), it follows that
\[
  \langle P,e^{kq_{\mathrm{tot}}}S\rangle
   =\sum_{j=0}^{\lfloor\deg P/2\rfloor}
       \frac{k^j}{j!}\,
          \langle q_{\mathrm{tot}}(\partial)^j P,S\rangle .
\]
The quadratic coefficients have support only in conjugate pairs.
The same form occurs on every column. Thus
\(q_{\mathrm{tot}}(\partial)\) is a linear combination of the
operators
\(\sum_b\partial_{(\tau,i,b)}\partial_{(\bar\tau,j,b)}\).
The defining harmonic equations annihilate \(P\) under each of them.
Only \(j=0\) contributes. Individual mixed derivatives may change
with the gauge; the asserted invariant is the harmonic combination.

For a change of formal coordinates, compose the two charts with
their normalized formal exponential in the marked tangent
directions. The resulting group homomorphisms agree since they
have the same differential and additive coordinates determine a
homomorphism in characteristic zero. A change of primitive frame
may initially give different frames at adjacent truncation orders.
Their ratio on the lower truncation is a unit normalized at the
axes. The equal cubical multiplication, symmetry and adjoint
factors imply that its logarithm in two arguments is biadditive,
symmetric and balanced, through that order. In logarithmic
coordinates it is a bilinear form through the truncation. Successive
subtractions of these bilinear forms make the frames compatible
to the desired order. Their cumulative diagonal change has the form
\(\exp(kq_{\mathrm{tot}})\) just treated. The total degree
\(D=\deg P\) suffices. Writing \(I=(Z)\), an invariant derivative
sends \(I^j\) into \(I^{j-1}\), by the Leibniz rule. Consequently
terms in \(I^{D+1}\) make no contribution to the evaluated
derivative. These facts give chart-independence at the exact order
used for the jet.
\footnote{\href{https://github.com/ChristianKudeba/alpoge-rh-verification/blob/alpoge-source-7cbf7a0c/src/Development/AbelianVarietyTheory/Proofs/FormalGroupCharts/CanonGaugeIndependenceHarmonic.lean}{Gauge change and conjugate eigendirections};
\href{https://github.com/ChristianKudeba/alpoge-rh-verification/blob/alpoge-source-7cbf7a0c/src/Formalism/SharpRestriction/Divisibility/GeneratorPowers/FrRatio.lean}{ratios of truncated frames};
\href{https://github.com/ChristianKudeba/alpoge-rh-verification/blob/alpoge-source-7cbf7a0c/src/Formalism/SharpRestriction/Divisibility/GeneratorPowers/CIndVerb.lean}{reduction of chart-independence to the finite jet}.}

\subsection{The normalized pullback factor}

Let \(\beta\bar\beta=N>0\) and write \(X=Z\) for the universal
formal point. The integral ratio between the normalized pullback
frame and the old frame is the following unit:
\[
 W(X)=
 \frac{J_\beta(X,A_{w_0\beta}X)}
 {\prod_{i=0}^{N-1}\Phi_2(X,A_{iw_0}X,A_{w_0}X)} .
\]
The normalization of the cubical line on the zero face supplies
constant one in this formula. To see the frame identity, identify
the primitive line at \(X\), divided by its origin fiber, with the
biextension line at \((X,A_{w_0}X)\).
The denominator composes the second partial addition
\(N\) times in the \(w_0\)-direction. The numerator transports the
\(\beta\)-action from the first to the second factor by Rosati,
where \(\bar\beta w_0\beta=Nw_0\). Both identifications agree at the
origin. Therefore a section with germ \(f\) in power \(k\) has
normalized pullback germ \(f'=W^k(f\circ A_\beta)\) in power \(kN\).

The logarithmic formula makes the entire adjustment explicit.
The denominator telescopes to
\(NH_h(X)-h(X,A_{Nw_0}X)\), since \(h(X,0)=0\).
The logarithm of its numerator is
\(H_h(A_\beta X)-h(X,A_{Nw_0}X)\). Hence
\[
      W=\exp\{H_h\circ A_\beta-NH_h\},\qquad
      f'e^{kNH_h}=(f e^{kH_h})\circ A_\beta .
\]
Applying an invariant derivative in direction \(u_{\tau,i}\) to a
pullback by \(A_\beta\) multiplies it by \(\tau(\beta)\). This follows
by differentiating the group-homomorphism law on the dual numbers.
Iteration, multiplication over the columns, and the fixed row
degrees of \(P\) give the character
\[
      \prod_{\tau,i}\tau(\beta)^{a_{\tau,i}} .
\]
All the harmonic monomials have that same factor.
\footnote{\href{https://github.com/ChristianKudeba/alpoge-rh-verification/blob/alpoge-source-7cbf7a0c/src/Development/AbelianVarietyTheory/Proofs/JetHarmonics/ThreePointSeriesExponentialOfDefect.lean}{Normalized cubical and adjoint factors};
\href{https://github.com/ChristianKudeba/alpoge-rh-verification/blob/alpoge-source-7cbf7a0c/src/Development/AbelianVarietyTheory/Proofs/JetHarmonics/HarmonicValueNormalisedPullbackChainRule.lean}{pullback germs and the character on the row degrees}.}

\subsection{Comparison on the uniformized frame}

Take the integral chart over \(T\), the framed generic point \(y\),
its polarized uniformization, and a generating frame on the jets
compatible with the rigidified primitive frame. Let \(F\) be a
characteristic-zero extension field of \(\operatorname{Frac}T\).
Suppose the uniformized sections \(\sigma_b\) correspond to the
geometric sections \(s_b\) under the specified generic-fiber
identification. Write \(f_b\) for their uniformized germs. To compare
the two evaluations it suffices, on each truncated jet, to make the
same base change in the coordinates, the tangent directions, the
cubical line and its generating frame. Indeed the two sections
restrict to the same section of the identified line. Dividing by
the common power of its generating frame gives the same truncated
series; equality may be checked on a trivializing open cover.
The canonical gauge also pulls back, since its identities are
identities of series and the homomorphism is between
characteristic-zero fields. The chart-independence just proved
identifies the scalar on these germs with the scalar defining the
connection on the original framed point. Finally, derivatives in
different columns act on independent variables, so evaluating the
box product of the \(f_b\) gives precisely the sum of products of
one-column values used to define the theta rule. This proves the
equality for the tied sections at any power \(k\) and harmonic \(P\).
\footnote{\href{https://github.com/ChristianKudeba/alpoge-rh-verification/blob/alpoge-source-7cbf7a0c/src/Development/AbelianVarietyTheory/Proofs/JetHarmonics/HarmonicValueConnectionIntegralChart.lean}{\srcid{harmonic_value_of_tied_sections}}.}

This field-valued comparison can also be used at a finite residue
prime. An integral assertion about the connection requires additional
control. Here is the unit hypothesis in one such assertion. Let
\(u,u'\) be integral directions of character \(\sigma\). Let \(B_h\)
be the matrix of the bidegree-\((1,1)\) part of \(h\). Taking that
component of the adjoint equation shows that
\[
       (dA_a)^t B_h-B_h dA_{\bar a}
\]
has integral entries. The logarithm of the integral adjoint unit
has the same bidegree-\((1,1)\) coefficients as the unit itself:
its rigidification makes it one on both axes, so higher powers of
its augmentation have bidegree at least \((2,2)\). Contracting the
matrix between \(u^t\) and \(u'\) proves that
\((\sigma(a)-\sigma(\bar a))u^tB_hu'\) is integral.
Thus an \(a\) with \(\sigma(a)-\sigma(\bar a)\in T^\times\) gives
integrality of this contraction. The additional operator \(dA_{w_0}\)
on \(u'\) multiplies it by the integral eigenvalue \(\sigma(w_0)\).
This unit-separation hypothesis belongs to this coefficient argument.
The separate projector for integral theta translates in
Section~\ref{sec:theta-translate-span} uses \(q\in T^\times\).
On the horizontal boundary both arguments are in residue
characteristic zero, where a nonzero element of the coefficient
field has valuation zero under the structural map through \(T\).
\footnote{\href{https://github.com/ChristianKudeba/alpoge-rh-verification/blob/alpoge-source-7cbf7a0c/src/Development/AbelianVarietyTheory/Proofs/JetHarmonics/CanonicalGaugeEigenvectorPairingIntegral.lean}{The precise same-character integrality lemma}.}

\subsection{Why the boundary weight shift has zero cost}

Let \(e_Z\) be the idempotent of the component. If a cusp section
\(f=e_Z f\) has nonzero evaluation, then \(e_Z(y)=1\).
The source constructs a distinguished unit section \(K_4\) and a
unit rule \(\rho\) with \(e_Z\rho^4=a\,e_ZK_4\),
\(a\) a nonzero coefficient-field element depending on \(Z\).
The integral-model theorem for the theta scalar gives
\(v_T(K_4(y))=0\) here. The map from the stalk of the horizontal
boundary sends the coefficient field into \(T\). Applying this to
both \(a\) and \(a^{-1}\) gives \(v_T(a)=0\). Evaluating the fourth
power identity gives \(4v_T(\rho(y))=0\), hence the weight-one
twisting section \(t=e_Z\rho^{-1}\) has order zero. Its inverse on
the component is \(e_Z\rho\). The input on parallel weights gives
another invertible twisting section \(t_{-w}\) of weight
\(\operatorname{par}(-w)\), of order zero at the same point.
Its boundedness and that of its inverse preserve the cusp condition;
the same property is part of the construction of \(t\).
Thus \(f\,t_{-w}t^{-m}\) removes both \(\operatorname{par}(w)\) and
the scalar weight \(m\), preserving nonvanishing and the exact order.
The boundary theorem is stated for \(4\mid m\); in this application
\(m=n/2\) and \(8\mid n\). The final multiplication uses only the
integer power once the weight-one twist is in hand. The fourth
power relation enters in constructing and controlling that twist.

Section~\ref{sec:key-different} derives the global fourth root from
the complex theta normalization and gives the finite-place different
bound. The model unit theorem for \(K_4\) and the construction of the
parallel-weight twist also enter the boundary argument.
\footnote{\href{https://github.com/ChristianKudeba/alpoge-rh-verification/blob/alpoge-source-7cbf7a0c/src/Development/AbelianVarietyTheory/Proofs/DivisorProofs/IntegralFramePointStalk.lean}{Construction of the weight-one twist and its valuation};
\href{https://github.com/ChristianKudeba/alpoge-rh-verification/blob/alpoge-source-7cbf7a0c/src/Development/AutomorphicArithmetic/ThetaCorrespondence/ThetaPairLifts/ThirdPlaceSeries/BoundaryValuationTwistInvariant.lean}{the boundary twist at the prescribed point}.}

\section{The fourth root on each component and the different at a finite prime}
\label{sec:key-different}

The weight-one twist of Section~\ref{sec:jet-connection} comes from a
weight-four theta rule. Here we explain its descent from the complex
chart and compute the precise determinant loss in a vertical
valuation. Write \(k\) for the coefficient number field and
\(m_\tau\) for the multiplicity of the embedding \(\tau:E\hookrightarrow k\)
in the Lie algebra. Repetitions of a single embedding occur through
the integer \(m_\tau\). Embeddings of multiplicity zero may be retained.

\subsection{Descending the fourth root}

On each smooth framed chart with coordinate ring \(A\), the
functorial theta scalar in a primitive frame is a unit \(u\in A^\times\).
Its transformation law under a triangular change of frame is the
fourth power of the frame character. The chartwise values therefore
give the weight-four rule \(K_4\), and the reciprocal values glue
in the opposite weight. Their product has value one at each generic
component, hence is the identity. The complex uniformization gives
one more piece of data: after embedding \(k\) in \(\mathbb C\), the
complex-valued function \(C u\), \(C\in\mathbb C^\times\), has a
continuous fourth root on the complex points of the chart. The root
is built from the determinant of a matrix of holomorphic periods,
divided by the square of the scalar relating the primitive theta
frames, with the common normalization from the Hermitian datum.
The adapted-period charts and the transformation law of the normalized
frames make this value independent of the representative and
continuous. These are hypotheses on the actual normalized theta
uniformization. A fourth root at individual complex points would
give no descent.

Here is the algebraic descent. Put
\(B=\mathbb C\otimes_k A\), a smooth integral affine algebra, and
\(u'=Cu\) in \(B^\times\). The integral-domain assertion here is
the geometric integrality of the chart of the chosen component.
The finite cover \(t^4=u'\) is étale, since \(4t^3\) is invertible.
On any connected component of its complex points the ratio of \(t\)
to the supplied continuous fourth root belongs to \(\mu_4\), so it
is constant. An integral algebraic component has connected complex
points. Smoothness makes the algebraic components disjoint, and
each maps onto the connected base, as the image of a finite étale
component is both open and closed. For one such component the
constant-ratio assertion gives at most one point above each complex
point of the base. Its rank is consequently one, and the component
maps isomorphically to the base. This gives \(b'\in B^\times\)
with \(b'^4=u'\).

Choose a fourth root of \(C\) in \(\mathbb C\) and absorb it into
\(b'\), so that \(b'^4=u\). Each automorphism of \(\mathbb C/k\)
takes \(b'\) to a fourth root of the same unit. The ratio is a
constant fourth root of unity: the polynomial \(t^4-1\) splits into
linear factors over \(\mathbb C\), and \(B\) is a domain. Expand
\(b'=\sum_j c_j a_j\) in finitely many members of a \(k\)-basis of
\(A\), with \(c_1\ne0\). All ratios \(c_j/c_1\) are fixed by every
automorphism over \(k\), so belong to \(k\). For a number field,
the fixed field is \(k\): an algebraic element outside \(k\) has a
different conjugate, and a transcendental element may be moved in
a transcendence basis, in either case extending the isomorphism
to the algebraic closure. Write \(b'=c_1 b\) with \(b\in A\).
Then \(u=a b^4\) with \(a=c_1^4\in k^\times\), because
\(1\otimes u=c_1^4(1\otimes b^4)\) and \(u,b^4\in A\) are nonzero.

Fix a generic frame on each component \(Z\), and absorb the
frame character into \(b\). This gives a rational element \(b_Z\)
in its function field. In another frame with character \(j\),
\((j b_Z)^4\) is a nonzero coefficient-field multiple of the local
unit representing \(K_4\). Both \(j b_Z\) and its inverse satisfy
monic equations over the smooth chart ring. Normality therefore
makes them regular there. The rational rules of opposite weights
glue to inverse regular rules. Equality of their fourth powers with
\(K_4\) up to \(a_Z\in k^\times\) can be checked at the generic
frame. This proves \(e_Z\rho^4=a_Ze_ZK_4\) once the continuous
root and the chartwise unit above are given.
\footnote{\href{https://github.com/ChristianKudeba/alpoge-rh-verification/blob/alpoge-source-7cbf7a0c/src/Development/AbelianVarietyTheory/Proofs/IwasawaDecomposition/ContinuousOnKeySectionComplexValue.lean}{Continuity from the adapted period chart and the theta scalar};
\href{https://github.com/ChristianKudeba/alpoge-rh-verification/blob/alpoge-source-7cbf7a0c/src/Development/AbelianVarietyTheory/Proofs/IwasawaDecomposition/ExistsKeySectionUnif.lean}{descent of the power up to a base-field scalar};
\href{https://github.com/ChristianKudeba/alpoge-rh-verification/blob/alpoge-source-7cbf7a0c/src/Development/AbelianVarietyTheory/Proofs/DivisorProofs/ExistsKeySectionFourthRootPerComponent.lean}{the componentwise rational root and normality argument}.}

\subsection{The integral scalar in two different Lie frames}

Let \(T\) be a complete discrete valuation ring of characteristic-zero
fraction field \(F\). Fix an integral framed point and a prime
\(w\) of \(k\), with the frame valuation restricting to \(e v_w\),
\(e>0\). Take the split eigendirections \(u_{\tau i}\) furnished by
the point, \(1\le i\le m_\tau\). For each \(\tau\) they span the
entire integral \(\tau\)-eigenspace in the tangent module. They
form a basis collectively on the fraction field. Write \(U\) for
the square matrix whose rows are these vectors in an integral Lie
basis and set \(D=\det U\). The good-model theorem for the theta
scalar applies to an integral generating theta section in an
integral Lie basis. The assumptions in the source application
include \(p_0+q_0\ge1\) and the level predicate
\srcid{GoodAt} at every rational prime. The primitive scalar is
then a unit \(\eta\in T^\times\); its comparison with the eigenframe
gives
\[
                            K_4(y)=D^4\eta .
\]
In particular, in this equation unitness refers to \(\eta\).
The eigenframe can have positive colength at ramified primes.

For clarity on the scope of the good-model input, the release
uses the invertibility of the theta pushforward, its model
uniformization, the arc property, and the split-fiber theorem. On
the normal affine chart it starts with a unit on the dense
generic open. A rational function on a normal noetherian affine
domain extends if it belongs to all its height-one valuation rings.
The proof applies this test to both the function and its reciprocal.
At the boundary the arc comparison supplies bounded complex values
for both; at the good-reduction tests the good-reduction unit theorem
does the same algebraically. The local extensions multiply to one
by density. The resulting chart unit pulls back to the trait, where
a change between two integral generating frames multiplies it by an
eighth power of a unit and the inverse fourth power of a determinant
unit. The arc comparison and good-reduction assertions are separate
geometric inputs. An arbitrary chart centered at a bad level prime
does not acquire these hypotheses from the field-valued connection.

\subsection{The different controls the colength}

Put \(\mathfrak D=\mathfrak D_{E/\mathbb Q}\).
For each \(\tau\) choose a nonzero
\(\delta_\tau\in\mathfrak D\) such that
\[
    v_w(\tau(\delta_\tau))
      =\operatorname{ord}_w(\tau(\mathfrak D)\mathcal O_k).
\]
Such an element exists: the extended ideal is contained in the
power of \(w\) on the right and fails containment in the next power.
If all images of elements of \(\mathfrak D\) lay in the next power,
so would the generated ideal.

Let \(\omega_j\) be a \(\mathbb Z\)-basis of \(\mathcal O_E\) and
\(\omega_j^*\) its trace dual. By definition of the inverse different
\(\delta_\tau\omega_j^*\in\mathcal O_E\). Let \(A(a)\) be the matrix
for \(a\) in the integral Lie basis. Set
\[
   d_\tau=\tau(\delta_\tau),\qquad
   P_\tau=\sum_j \tau(\delta_\tau\omega_j^*)A(\omega_j).
\]
All the coefficients on the right are integral in \(T\), via the
given embedding. Over \(F\) the \(d_\tau^{-1}P_\tau\) are the
character projectors. Indeed
\(\sum_j\tau(\omega_j^*)\sigma(\omega_j)=\delta_{\tau\sigma}\).
To verify this, form the square embedding matrices of the two
bases. Their pairing under the trace is the identity, and inverting
in the opposite order gives exactly this relation. The embeddings
indexed here are all embeddings of \(E\), whose images belong to
\(F\). The relation acts on each eigenspace as claimed and yields
\(\sum_\tau d_\tau^{-1}P_\tau=I\).

Every column of \(P_\tau\) is an integral eigenvector of character
\(\tau\). Resolve it in the generating vectors \(u_{\tau i}\).
Stacking the integral coefficients gives a matrix \(G\) whose
\((j,(\tau,i))\)-entry is the coefficient of \(u_{\tau i}\) in
column \(j\) of \(P_\tau\). The transpose of the projector sum is
then
\[
  G\,\operatorname{diag}(d_\tau^{-1}
           \text{ repeated }m_\tau\text{ times})\,U=I .
\]
Taking determinants yields the integral identity
\(\det G\cdot D=\prod_\tau d_\tau^{m_\tau}\).
Thus
\[
  0\le \frac{v_T(K_4(y))}{e}
   =\frac{4v_T(D)}{e}
   \le 4\sum_\tau m_\tau
          \operatorname{ord}_w(\tau(\mathfrak D)\mathcal O_k).
\]
The estimate is zero wherever all the displayed different ideals
are units. If the fourth root on \(Z\) uses \(a_Z\), then the
weight-one section \(e_Z\rho^{-1}\), at any nonzero evaluation
of \(e_Z\), has normalized order in
\[
 \left[-\frac{v_w(a_Z)}4
       -\sum_\tau m_\tau\operatorname{ord}_w(\tau(\mathfrak D)\mathcal O_k),
       \ -\frac{v_w(a_Z)}4\right].
\]
This is the vertical contribution; on a horizontal trait the
coefficient field lies in \(T\) and the valuation comparison in
Section~\ref{sec:jet-connection} directly gives order zero.
\footnote{\href{https://github.com/ChristianKudeba/alpoge-rh-verification/blob/alpoge-source-7cbf7a0c/src/Development/AbelianVarietyTheory/Proofs/DivisorProofs/ColengthDvdOfNeZero.lean}{Trace-dual projectors and the determinant divisibility};
\href{https://github.com/ChristianKudeba/alpoge-rh-verification/blob/alpoge-source-7cbf7a0c/src/Development/AbelianVarietyTheory/Proofs/DivisorProofs/ThetaKeySectionValuationBoundsDifferent.lean}{the exact normalized vertical bounds};
\href{https://github.com/ChristianKudeba/alpoge-rh-verification/blob/alpoge-source-7cbf7a0c/src/Development/AbelianVarietyTheory/Proofs/DivisorProofs/ThetaKeyScalarIsUnitOnIntegralModels.lean}{the model hypotheses and height-one extension}.}

\section{The auxiliary powers in the theta-rule DVR test}
\label{sec:key-rule-dvr}

The good-reduction step in Section~\ref{sec:key-different} uses a
unit normalization for the theta scalar on an integral polarized
abelian family. We explain how the local construction can give
unit values at \emph{every} residue characteristic even though
the construction of an individual powered invariant uses an
auxiliary integer which must be invertible. The word level has
two different roles here. The auxiliary integer in this argument
is chosen for a determinant calculation on sections of the
polarizing line. The fixed level of the moduli problem satisfies
the separate condition in Section~\ref{sec:key-different}.

\subsection{Choosing transverse-intersection degrees at a given prime}

For the numerical construction let \(d\geq1\) denote the dimension
parameter for the intersection. Given a prime \(p\), set
\[
       n=4p,\qquad b=dn+1,\qquad b'=b+p,\qquad
       P=\prod_{j=0}^d (b-jn)(b'-jn).
\]
Then \(n\ge4\) and \(dn<b<b'\). Each factor in \(P\) is \(1\bmod p\),
so \(P>0\) and \(p\nmid P\), including at \(p=2\). The exponent in
the determinant construction is
\[
     c=2P^{3d+2}\,d!\,n^d(b'-b)>0 .
\]
It can be divisible by \(p\). Only \(P\) is required to be a unit
in this part of the construction.

For every \(m\geq0\) choose a four-squares representation
\(m=a^2+b_1^2+b_2^2+b_3^2\). The integer matrix
\[
 U_m=
 \begin{pmatrix}
 a&-b_1&-b_2&-b_3\\
 b_1&a&-b_3&b_2\\
 b_2&b_3&a&-b_1\\
 b_3&-b_2&b_1&a
 \end{pmatrix}
\]
has \(U_m^{\,t}U_m=mI_4\): the columns have squared norm \(m\)
and their pairwise dot products cancel termwise. These matrices
supply the square-determinant comparisons at degrees \(b-jn\)
and \(b'-jn\). This proves all the arithmetic admissibility
conditions.

Here is the geometric input used with these integers. A tuple of
\(d\) sufficiently general sections of \(L^n\) on the special
fiber cuts out \(d!n^d\) transverse points for the principal
polarization. The release picks a basis of \(H^0(L^n)\), gives
the tuple algebraically independent coefficients over the residue
field, and passes to a local injection into a complete DVR in
which the finitely many coefficients have lifts. Its transversality
theorem gives the required étale zero locus on the extended residue
field. After another injective local extension, the determinant
comparison theorem supplies a unit from the two section Koszul
complexes, their evaluation at those points, the Jacobian product,
and the comparisons on fourth powers given by the \(U_m\). This
step requires \(P\) to be invertible. The separate theorems on
independence of the sections and change of base make the resulting
powered invariant functorial. These intersection, determinant,
and functoriality assertions are geometric inputs to the following
valuation argument.
\footnote{\href{https://github.com/ChristianKudeba/alpoge-rh-verification/blob/alpoge-source-7cbf7a0c/src/Development/AbelianVarietyTheory/Proofs/ThetaLawProofs/ExistsAdmissibleKeyNumericalDataPrimeTo.lean}{Admissible prime-to-\(p\) parameters and the quaternions};
\href{https://github.com/ChristianKudeba/alpoge-rh-verification/blob/alpoge-source-7cbf7a0c/src/Development/AbelianVarietyTheory/Definitions/KeyFormula/DegProd.lean}{degrees and powers in the determinant invariant};
\href{https://github.com/ChristianKudeba/alpoge-rh-verification/blob/alpoge-source-7cbf7a0c/src/Development/AbelianVarietyTheory/Proofs/ThetaLawProofs/ExistsIntegralKeyInvariantRuleCharZero.lean}{residue-field extension, transverse intersections and comparison}.}

\subsection{Changing the auxiliary integer}

In abelian dimension \(g\ge1\), the normalization step uses the
previous numerical construction at \(d=2g\). For each prime \(p\)
it produces a positive exponent \(N_p=2c\), an integer \(P_p\)
prime to \(p\), and a scalar \(F_p(x)\) functorial on framed
data \(x\) over fields of characteristic zero. At a DVR where
\(P_p\) is a unit its value on the generic fiber of integral
data is a unit. On a complex datum in normalized half-characteristic
coordinates it has value
\[
          F_p(x)=\bigl(\det(M)^4/c_0^8\bigr)^{N_p},
\]
where \(c_0\) is the nonzero theta scaling. These assertions are
the output of the normalized determinant construction just invoked.

For any two auxiliary primes \(p,q\), and any field \(K\) of
characteristic zero, one has
\[
                         F_p(x)^{N_q}=F_q(x)^{N_p}.
\]
To see this, descend the framed datum to a countable subfield
\(K_0\subset K\), as permitted by finite data descent in the
source. Embed \(K_0\) into \(\mathbb C\) and extend the datum
there. The two sides are equal to the same power of
\(\det(M)^4/c_0^8\), by the complex uniformization. Functoriality
and injectivity of \(K_0\hookrightarrow\mathbb C\) give equality
over \(K_0\). Extending scalars back gives the formula over \(K\).

Fix one choice \(F_2\). In residue characteristic zero \(P_2\)
is invertible. If the residue characteristic is \(p>0\) and
\(P_2\) is a nonunit, use \(F_p\), for which \(P_p\) is a unit.
For a DVR valuation \(v\) on \(K\) and the generic fiber of the
integral frame,
\[
           N_p v(F_2(x))=N_2 v(F_p(x))=0 .
\]
Since \(N_p>0\), this gives \(v(F_2(x))=0\).
Thus the fixed powered rule is a unit at all residue
characteristics. This computation uses no \(p\)-adic division by
the integer \(N_p\) in the valuation ring.

The exponent-one scalar used in the manuscript also requires
a coherent root over arbitrary characteristic-zero fields.
The complex candidate \(t(x)=\det(M)^4/c_0^8\) is fixed by the
theta transformation law. The release establishes its descent
using connected families of polarized varieties, equivariance of
the complex value, and its specialization on smooth curves. The
elementary uniqueness step is as follows. Suppose a smooth
integral complex base carries a unit \(u\) and a continuous
normalized value \(t\) with \(t^N=u^N\) on the complex points.
Their ratio takes values in the finite discrete group \(\mu_N\);
the complex points of this base are connected, so the ratio is
constant. An equality \(u=t\) in one fiber identifies them on
the entire base. Establishing the indicated families and their
arithmetic comparison is additional to this observation.
Once descent gives the exact rule \(\kappa\) functorial in the
field and with the displayed normal value, apply the same
countable-field argument to obtain
\(\kappa(x)^{N_2}=F_2(x)\). Taking valuations gives
\(v(\kappa(x))=0\), proving the unit conclusion for an integral
proper theta datum at every residue characteristic under the
stated inputs.
\footnote{\href{https://github.com/ChristianKudeba/alpoge-rh-verification/blob/alpoge-source-7cbf7a0c/src/Development/AbelianVarietyTheory/Proofs/ThetaLawProofs/ThetaKeyRuleExistsOfAlgebraicInvariants.lean}{Cross-power equality and replacement of the auxiliary integer};
\href{https://github.com/ChristianKudeba/alpoge-rh-verification/blob/alpoge-source-7cbf7a0c/src/Development/AbelianVarietyTheory/Proofs/ThetaLawProofs/ThetaFrameValueEqUnitSmoothConnectedBase.lean}{constancy of the root on smooth connected bases};
\href{https://github.com/ChristianKudeba/alpoge-rh-verification/blob/alpoge-source-7cbf7a0c/src/Development/AbelianVarietyTheory/Proofs/UniformizationProofs/KeyFormulaOfConnFam.lean}{the geometric and specialization inputs for the exponent-one rule}.}

\section{The integral span of finite-order theta translates}
\label{sec:theta-translate-span}

We prove the spanning step in Section~\ref{sec:local-theta-jets}
from the valuation and theta-group properties of the uniformization.
Work over the valued algebraic closure \(F\) of the fraction field
of the chart. Put \(\mathcal V=\{u\in F:v(u)\ge0\}\), and let
\(\mathcal U\) be the level-\(q\) expansion module with normalized
coefficient valuation lower bound zero. A member has positive depth
if the same bound holds with some positive rational added at every
coefficient. In particular a positive-valuation multiple of a
member has positive depth. Here \(v(q)=0\).

Let \(A\) denote the polarized generic abelian variety in the product
uniformization and write \(g=\dim A\).
The CM ideals \((\alpha)\) and \((\overline\alpha)\) are coprime
and \(\alpha\overline\alpha=q\). Define
\(H_1=\ker[\overline\alpha]\) and \(H_2=\ker[\alpha]\) on geometric
points; both consist of \(q\)-torsion points. Choose
\(u\alpha+u'\overline\alpha=1\). On \(A[q]\) the two summands
\([u\alpha]\) and \([u'\overline\alpha]\) have image in \(H_1\)
and \(H_2\), respectively, and their sum is the identity.
Their kernels give \(H_1\cap H_2=0\). Hence
\(A[q]=H_1\oplus H_2\).

The Weil pairing separates the two groups. In the unimodular
alternating lattice \(L\) supplied by the specified Tate isometry,
represent distinct \(x,x'\in H_1\) by \(a,a'\) modulo \(qL\).
There is \(b_0\in L\) with \(q\nmid\langle b_0,a-a'\rangle\):
otherwise the integer-valued functional
\(b\mapsto\langle b,a-a'\rangle/q\) would imply
\(a-a'\in qL\) by unimodularity. The adjoint of the CM action is
conjugation and \(\overline\alpha(a-a')\in qL\), so
\[
   q\mid\langle u\alpha b_0,a-a'\rangle,\qquad
   q\nmid\langle u'\overline\alpha b_0,a-a'\rangle.
\]
The class \(z\) of \(u'\overline\alpha b_0\) belongs to \(H_2\);
the two pairing values \(e_q(z,x)\) and \(e_q(z,x')\) differ.
The same reasoning with conjugates interchanged gives separation
in the other direction. Thus each group injects into the
\(\mu_q\)-valued character group of the other. For a finite abelian
group of exponent dividing \(q\) this dual has the same cardinality,
by cyclic decomposition. Since \(\#A[q]=q^{2g}\), it follows that
\(\#H_1=\#H_2=q^g\).
\footnote{\href{https://github.com/ChristianKudeba/alpoge-rh-verification/blob/alpoge-source-7cbf7a0c/src/ReciprocityTheory/ReductiveGroups/ArithmeticGeometry/FirstCohomology/AlphaPair.lean}{Unimodular lattice pairing and separation};
\href{https://github.com/ChristianKudeba/alpoge-rh-verification/blob/alpoge-source-7cbf7a0c/src/ReciprocityTheory/ReductiveGroups/ArithmeticGeometry/FirstCohomology/AlphaBarEnum.lean}{enumeration of the two kernels}.}

Let \(\theta_\alpha\) denote the pullback of the primitive theta
frame. The forward pullback estimate puts it in \(\mathcal U\).
The primitive coefficient at zero and the level-frame inequality
show that \(\theta_\alpha\) cannot have positive depth: an improved
bound \(r>0\) would imply \(r\le v(1)-Q(0)=0\).
For \(x\in H_1\), choose a finite-order lift to the theta group
and let \(t_x\) be the associated translate of \(\theta_\alpha\),
including the scalar normalizing its order. The finite-order
valuation invariance transfers both assertions about depth to \(t_x\).
In particular each \(t_x\ne0\). For each \(z\in H_2\), the canonical
lift acts by a linear operator \(T_z\) fixing \(\theta_\alpha\);
it preserves \(\mathcal U\).
Using the commutator formula for the theta group with the polarized
Tate isometry, we get
\[
                    T_z t_x=e_q(z,x)t_x .
\]
Section~\ref{sec:theta-lift-valuation} proves the canonical-lift and
valuation identities in the normalized expansion.

Here is the resulting integral linear algebra. If \(\xi,\xi'\) are
distinct \(q\)-th roots of unity, \(v(\xi)=v(\xi')=0\) and
\(v(\xi-\xi')=0\). Indeed, were the two residues equal to \(c\), then
\(c\ne0\) and reducing
\[
          \sum_{i=0}^{q-1}\xi^i(\xi')^{q-1-i}=0
\]
would give \(qc^{q-1}=0\), contrary to \(v(q)=0\).
Given a fixed \(x\in H_1\), for each \(y\in H_1\setminus\{x\}\)
choose \(z_y\in H_2\) separating \(x,y\). On the span of these
eigenvectors form the product, in any fixed order,
\[
   D_x=\prod_{y\ne x}
   \frac{T_{z_y}-e_q(z_y,y)\operatorname{id}}
        {e_q(z_y,x)-e_q(z_y,y)} .
\]
All denominators are units of \(\mathcal V\). Each factor preserves
\(\mathcal U\), acts as the identity on \(t_x\), and annihilates
\(t_y\) for its particular \(y\). Thus \(D_x t_w\) is \(t_x\)
if \(w=x\) and zero otherwise. Applying \(D_x\) to any linear
relation shows the \(t_w\) are independent.

Section~\ref{sec:theta-unit-basis} gives \(\mathcal U\)
as the \(\mathcal V\)-span of a linearly independent family of order
\(q^{\operatorname{rank}X+\dim B}\), where \(X\) is the product
period lattice and \(B\) the abelian part. The semiabelian dimension
formula gives \(\operatorname{rank}X+\dim B=g\). Thus the \(q^g\)
translates span the same \(F\)-space. For \(s\in\mathcal U\),
write \(s=\sum_w d_wt_w\) there. Then \(d_xt_x=D_xs\in\mathcal U\).
If \(v(d_x)<0\), writing \(t_x=d_x^{-1}(d_xt_x)\) would give
positive depth for \(t_x\). Thus \(d_x\in\mathcal V\), for every \(x\).
This proves the exact integral span required for the valuation
selection on the original product and identifies where a prime
dividing the level would enter the separation argument.
\footnote{\href{https://github.com/ChristianKudeba/alpoge-rh-verification/blob/alpoge-source-7cbf7a0c/src/ReciprocityTheory/ReductiveGroups/ArithmeticGeometry/FirstCohomology/IntegralRingEmbedding.lean}{Integral eigenvector elimination};
\href{https://github.com/ChristianKudeba/alpoge-rh-verification/blob/alpoge-source-7cbf7a0c/src/ReciprocityTheory/ReductiveGroups/ArithmeticGeometry/FirstCohomology/OneCPf.lean}{pairing, primitivity, and the rank comparison};
\href{https://github.com/ChristianKudeba/alpoge-rh-verification/blob/alpoge-source-7cbf7a0c/src/ReciprocityTheory/ReductiveGroups/ArithmeticGeometry/FirstCohomology/CmLiftConditionOfLaws.lean}{assembly of geometric hypotheses};
\href{https://github.com/ChristianKudeba/alpoge-rh-verification/blob/alpoge-source-7cbf7a0c/src/Development/AbelianVarietyTheory/Proofs/JetHarmonics/ExistsTermLe.lean}{transport of the result to the product degeneration}.}

The polarized torsion isometry and the integral period and pullback
identifications connect the expansion model to the abelian variety.
Sections~\ref{sec:theta-lift-valuation} and \ref{sec:theta-unit-basis}
supply the required valuation and rank assertions in that model.

\section{The canonical theta lift and preservation of valuation}
\label{sec:theta-lift-valuation}

We keep the chart and uniformization of
Section~\ref{sec:theta-translate-span}, and write \(F=\overline K\)
and \(\mathcal V=\{c\in F:v(c)\ge0\}\). The two statements used in the
integral spanning argument have different hypotheses. A point killed
by \(\alpha\) has a lift fixing the pulled-back theta function.
An arbitrary finite-order lift preserves every normalized coefficient
valuation bound, even when the theta function being translated is
unrelated to that pullback.

\paragraph{The lift fixing the pullback.}
On the generic abelian variety let \(f=[\alpha]\). The equality
\(f[\overline\alpha]=[\overline\alpha]f=[q]\) implies that \(f\) is
surjective and that its kernel lies in \(A[q]\). Hence \(f\) is an
isogeny. We work over the algebraically closed field \(F\) of
characteristic zero, so this is a separable isogeny.
If \(\mathcal L\) is the ample theta bundle, set
\(\mathcal M=f^*\mathcal L\), with the fixed normalized identification
with the level-\(q\) bundle. For \(x\in\ker f\), \(fT_x=f\).
Pullback of this equality gives
\(T_x^*\mathcal M\simeq\mathcal M\). These identifications satisfy the
group law by functoriality, so they split the theta group above
\(\ker f\). Each fixes \(f^*s\) for every section \(s\), and its order
divides \(q\). This is precisely the pullback construction in Mumford's
level-subgroup correspondence: the isogeny, the descended ample
line bundle, and the identification of its pullback are specified here.
\footnote{\href{https://pazuki.perso.math.cnrs.fr/index_fichiers/Mumford66.pdf}{D.~Mumford,
\emph{On the equations defining abelian varieties I}, \S1,
Proposition~1, pp.~290--291}.}

We describe the compatibility required to use this lift on the
uniformized coefficients. Let \(p\) represent \(x\) in the
uniformization at level \(q\), with lattice coordinate \(y\in X\)
and abelian coordinate \(b\). Since \(f(x)=0\), the torsion
isomorphism gives a period \(z\in X\) with
\(\alpha y+qz=0\). Multiplication by \(\overline\alpha\) and
torsion freeness of \(X\) give \(y=-\overline\alpha z\).
On the abelian part \([\alpha]b=c(-z)\), where \(c\) is the period
map to the abelian part. The generator identifying the CM adjoint
with Rosati is rigidified to one at the origin. Write \(R_\chi\)
for its value on \((c(\chi),b)\). If \(G_{\overline\alpha\chi}\)
is the theta-torsor factor in the chosen representative, then the
multiplication rule for addition of the compensating period \(z\)
gives
\[
 G_{\overline\alpha\chi}R_\chi=\tau(-z,\chi).
\]
Indeed the class of \([\alpha]p\) plus the image of \(z\) is zero.
In the addition formula the cubical factor multiplying
\(\tau(-z,\chi)\tau(z,\chi)\) is the same one in the period
cocycle for \((-z,z)\); cancelling it and using \(\tau(0,\chi)=1\)
gives the displayed identity in common frames.

The coefficient of \(\theta_\alpha\) outside
\(\overline\alpha X\) vanishes. Translation by \(y\) preserves
that subset. At an index \(\overline\alpha\chi\), its translate
has index \(\overline\alpha(\chi-z)\). There are two ways to compare
the corresponding lines: pull back the \((-z)\)-period functional
equation of the original theta function, or translate its pullback
by \(b\). At the origin choose the lift generator so that these
comparisons agree, using the displayed identity and the specified
rigidification of the power-\(q\) bundle. The ratio of the two
cubical identifications at any point is a scalar: the line maps
differ by a global unit on the proper geometrically connected
abelian part. Its value at the origin is one, so they agree
everywhere. This verifies
\(\Phi_x\theta_\alpha=\theta_\alpha\) in the expansion action.
For a representative of a \(q\)-torsion point the \(q\)-th power
of its lift is central, hence acts by a scalar \(c\). As
\(\theta_\alpha\) is primitive and nonzero in our use, this
eigenvalue identity forces \(c=1\). In the version that allows the
zero theta function, any scalar \(q\)-th root of \(c^{-1}\) rescales
the operator to finite order and still fixes zero.
\footnote{\href{https://github.com/ChristianKudeba/alpoge-rh-verification/blob/alpoge-source-7cbf7a0c/src/ReciprocityTheory/ReductiveGroups/ArithmeticGeometry/FirstCohomology/CanonicalLiftUnpacking.lean}{Period representing the zero torsion point};
\href{https://github.com/ChristianKudeba/alpoge-rh-verification/blob/alpoge-source-7cbf7a0c/src/ReciprocityTheory/ReductiveGroups/ArithmeticGeometry/FirstCohomology/CanonicalLiftOrigin.lean}{Rosati normalization and cancellation at the origin};
\href{https://github.com/ChristianKudeba/alpoge-rh-verification/blob/alpoge-source-7cbf7a0c/src/ReciprocityTheory/ReductiveGroups/ArithmeticGeometry/FirstCohomology/CanonicalLiftFrobenius.lean}{comparison at each character};
\href{https://github.com/ChristianKudeba/alpoge-rh-verification/blob/alpoge-source-7cbf7a0c/src/ReciprocityTheory/ReductiveGroups/ArithmeticGeometry/FirstCohomology/CanonLiftPf.lean}{finite-order assertion}.}

\paragraph{The valuation shift for a general lift.}
For a level-\(q\) family \(s\) and \(r\in\mathbb Q\), denote by
\(\mathcal U_r(s)\) the following predicate. For every character
\(\chi\), its coefficient \(s_\chi\) is a finite combination of
base changes of integral sections at \(\chi\), with scalar
coefficients of valuation at least \(r+Q(\chi)/q\).
In particular the earlier \(\mathcal U\) consists of those \(s\)
satisfying \(\mathcal U_0(s)\).

Consider a lift represented by a torsion datum with lattice shift
\(y\), abelian coordinate \(b\), line generator \(w\), and scalar
\(\zeta\in F^\times\); write \(S\) for its operator. Properness
of the abelian part provides an extension \(m\) of \(b\) to
\(\mathcal V\). An integral trivialization of the line for \(w\)
there has ratio \(c_w\in F^\times\) to \(w\), with convention
\(w=c_w\) times the integral generator. We claim
\begin{equation}\label{eq:theta-liftshift}
 \mathcal U_r(s)\quad\Longleftrightarrow\quad
 \mathcal U_{r+C}(Ss),\qquad
       C=v(c_w)+v(\zeta)-Q(y)/q .
\end{equation}

For the index-dependent theta-torsor factor, choose an integral
generator \(G\) of the biextension line at \((m,c(\chi))\); write
the factor as \(t_\chi G\). We compute
\begin{equation}\label{eq:cubical-ratio}
                        qv(t_\chi)=B(y,\chi).
\end{equation}
The torsion relation is \(qm=c(y)\). The period multiplier at
level \(q\) has order \(qQ(y)+B(y,\chi)\): its factors have
orders \(qQ(y)\), \(Q(y+\chi)-Q(y)-Q(\chi)\), and the rigidification
order zero. Denote the scalar of that multiplier in an integral
frame by \(c_0\), the scalar of the \(q\)-th power of the period
generator at \(y\) by \(c_\psi\), and the scalar of the \(q\)-th
power of the rigidification by \(c_1\). Their valuations are
\(qQ(y)+B(y,\chi),qQ(y),0\), respectively.

To see the cancellation, trivialize all factors in a common chart.
Write \(P,R\) for the last two powered generators, \(\tau\) for the
period biextension factor, and \(D\) for the cubical identification
iterated from \(m\) to \(qm\). The family multiplier is
\(\tau P R\), and the theta-torsor chain at \(p\) gives
\((t_\chi G)^q=\tau D\). Write \(P=c_\psi E_\psi\),
\(R=c_1 E_1\), and \(\tau PR=c_0 E\) with the respective integral
generators \(E_\psi,E_1,E\). It follows that
\[
       c_0 A=t_\chi^q c_\psi c_1,\qquad
       A=E E_\psi^{-1} E_1^{-1} D G^{-q}.
\]
All the factors in \(A\) are integral generators of their lines.
Their product is a generator of the trivial line on
\(\operatorname{Spec}\mathcal V\): the tensor factors for \(m\)
and \(qm\) cancel by \(qm=c(y)\) and the cubical identity.
Therefore \(A\in\mathcal V^\times\). Taking valuations in the
previous equality proves \eqref{eq:cubical-ratio}, including its
sign and its factor \(q\).
\footnote{\href{https://github.com/ChristianKudeba/alpoge-rh-verification/blob/alpoge-source-7cbf7a0c/src/ReciprocityTheory/ReductiveGroups/ArithmeticGeometry/DiscreteValuationRings/ExistsValGenHasVal.lean}{Orders of the factors of the period multiplier};
\href{https://github.com/ChristianKudeba/alpoge-rh-verification/blob/alpoge-source-7cbf7a0c/src/Development/AbelianVarietyTheory/Proofs/RingsOfIntegers/ValuationThetaRatioEqBilinear.lean}{cubical trivialization and scalar identity}.}

A coefficient at \(\chi\) is carried to one at \(\chi+y\)
by translation through \(m\) and multiplication by the integral
line isomorphism just described and the scalar \(\zeta c_w t_\chi\).
Translation through \(m\), and its inverse, preserve integral
sections over \(\mathcal V\). They can equivalently be expressed
as combinations of sections over \(T\). Indeed \(T\) is a discrete
valuation ring and \(\mathcal V\subset F\) is torsion free over \(T\),
hence flat. The abelian scheme over \(T\) is proper, thus
quasi-compact and quasi-separated, and each coefficient sheaf is
invertible. Flat base change for \(H^0\) identifies the section
module after extension with the tensor product of the section
module over \(T\) and \(\mathcal V\). This proves the instance of
integral base change needed here, including for a nondiscrete
\(\mathcal V\).
\footnote{\href{https://stacks.math.columbia.edu/tag/0539}{Stacks Project,
Lemma~15.22.10, flatness over a valuation ring};
\href{https://stacks.math.columbia.edu/tag/02KH}{Lemma~30.5.2, flat base change for sections}.}

The gain in the unnormalized lower bound of the coefficient
is \(v(c_w)+v(\zeta)+B(y,\chi)/q\). Subtracting the new
\(Q(\chi+y)/q\) and restoring the old \(Q(\chi)/q\) gives exactly
\(C\). This proves the forward implication in
\eqref{eq:theta-liftshift}; translation by \(-m\) and the inverse
line generator gives the reverse implication. No nonzero
coefficient or attained minimum needs to be chosen. In the usual
analytic notation for a Raynaud extension, the underlying period
functional equation for coefficients is
\(a_{u+\lambda(z)}=t(z,u)c(z)T_{\Phi(z)}^*a_u\), interpreted in
the corresponding line bundles. Foster, Rabinoff, Shokrieh and Soto
give both this formula and the equality of its model norms over
a complete algebraically closed nonarchimedean field. It applies
after passing from \(F\) to its completion in the semiabelian
case here. The coefficient-span calculation above works directly
over \(F\) and supplies \eqref{eq:theta-liftshift} for the chosen
normalization.
\footnote{\href{https://arxiv.org/pdf/1705.08212}{T.~Foster,
J.~Rabinoff, F.~Shokrieh, A.~Soto, \emph{Non-Archimedean and
tropical theta functions}, Proposition~3.19 and Corollary~3.21};
\href{https://github.com/ChristianKudeba/alpoge-rh-verification/blob/alpoge-source-7cbf7a0c/src/ReciprocityTheory/ReductiveGroups/ArithmeticGeometry/DiscreteValuationRings/TWMain.lean}{the full scalar shift in the coefficient filtration}.}

\paragraph{Eliminating the shift for a finite-order lift.}
Suppose \(S^h=1\), \(h\ge1\). For fixed \(s\) the subset
\(D_s=\{r\in\mathbb Q:\mathcal U_r(s)\}\) is downward closed,
and \eqref{eq:theta-liftshift} gives \(D_s=D_s+hC\).
If \(C=0\) the assertion is immediate. If \(C\ne0\), the only
downward-closed subsets of \(\mathbb Q\) with that period are the
empty set and all of \(\mathbb Q\): from one element \(u\),
\(u+j|hC|\) belongs to the set for every \(j\ge0\). Thus in this
case \(D_s+C=D_s\) also. It follows that every finite-order
lift preserves \(\mathcal U_r\) for every \(r\), including positive
depth and zero depth. This proof covers the zero family, whose
valuation bounds are unbounded above.
\footnote{\href{https://github.com/ChristianKudeba/alpoge-rh-verification/blob/alpoge-source-7cbf7a0c/src/ReciprocityTheory/ReductiveGroups/ArithmeticGeometry/DiscreteValuationRings/PeriodicValuation.lean}{Periodicity and exact finite-order invariance}.}

The use of these identities assumes the polarized uniformization,
its torsion isometry and its rigidified cubical period data. Section~\ref{sec:theta-unit-basis} proves the integral rank of the
theta expansion module and the primitivity assertions used by
Section~\ref{sec:theta-translate-span}.

\section{The unit ball of theta expansions and its primitive frame}
\label{sec:theta-unit-basis}

Keep the valued algebraic closure \(F\), its valuation ring
\(\mathcal V\), and the PEL degeneration over the complete discrete
valuation ring \(T\). Let \(B/T\) be the abelian part of dimension
\(d\), with its principal theta bundle \(L\), and let \(X\) be the
period lattice of rank \(r\). At level \(q\ge1\) the coefficient for
\(\chi\in X\) is a section of a line bundle \(L_{q,\chi}\) on \(B\).
In a rigidification its class is \(L^{\otimes q}\) times the
degree-zero twist attached to \(\chi\). We compute the unit ball of
the expansion space in the coefficient valuation of
Section~\ref{sec:theta-lift-valuation}.

\paragraph{Integral coordinates at one character.}
The module \(V_{q,\chi}=H^0(B,L_{q,\chi})\) is finite over \(T\):
the base is noetherian, \(B\) proper, and the line bundle coherent.
It is torsion free, because \(B\) is flat over \(T\) and a line
bundle is locally free. Thus \(V_{q,\chi}\) is free over \(T\).
Flat base change identifies \(V_{q,\chi}\otimes_T F\) with the
sections on \(B_F\). On that fiber the bundle is ample, since its
degree-zero twist preserves numerical equivalence. Its polarization
homomorphism is \(q\lambda_L\), of degree \(q^{2d}\) by the principal
polarization. Riemann--Roch for abelian varieties gives
\(\chi(L_{q,\chi})^2=q^{2d}\); ampleness gives positive Euler
characteristic and vanishing of higher cohomology. Hence
\(\operatorname{rank}_T V_{q,\chi}=q^d\).
\footnote{\href{https://stacks.math.columbia.edu/tag/02O6}{Stacks Project,
Lemma~30.19.2, proper cohomology over an affine noetherian base};
\href{https://stacks.math.columbia.edu/tag/02KH}{Lemma~30.5.2, flat base change};
\href{https://jmilne.org/math/CourseNotes/AV.pdf}{J.~S.~Milne,
\emph{Abelian Varieties}, Chapter~I, Theorems~7.2 and 11.1};
\href{https://wiki.epfl.ch/shimvar/documents/IHESlectures-Genestier-Ngo.pdf}{A.~Genestier
and B.~C.~Ng\^o, \emph{Lectures on Shimura varieties},
Theorems~2.2.1--2.2.2, Riemann--Roch and the index of an ample class}.}

Choose a \(T\)-basis \(b_{\chi,1},\ldots,b_{\chi,q^d}\).
Writing a coefficient as \(\sum_i c_i b_{\chi,i}\) over \(F\),
its section valuation is at least \(h\in\mathbb Q\) if and only
if \(v(c_i)\ge h\) for each \(i\). The reverse direction follows
by taking this basis as the integral sections in the definition.
For the forward direction, expand the integral sections in any
given representation in this basis. Their coordinates lie in \(T\),
so each resulting coefficient is a sum of terms of valuation at
least \(h\). This interpretation gives all rational thresholds,
including when \(F\) is nondiscretely valued.

\paragraph{Coset coordinates for the whole family.}
Choose representatives \(e_j\) for \(X/qX\), of cardinality \(q^r\).
A theta expansion is determined by its coefficients at the \(e_j\).
Indeed, its functional equation determines the coefficient at
\(e_j+qz\) by translating that at \(e_j\) by \(c(z)\) on the
abelian part and tensoring with the period multiplier. Each
\(\chi\) has one such expression after fixing \(e_j\), because \(X\)
has no torsion. Conversely assign an arbitrary coefficient at each
representative and use this rule for all \(z\). Applying the rule
twice agrees with its application to the sum of the periods by
the period cocycle. This constructs a family and proves that
the joint coordinate map is a linear isomorphism
\[
     s\longmapsto(c_{j,i}(s))_{j,i},\qquad
     \mathcal F_q \ \simeq\
               \bigoplus_{j=1}^{q^r} F^{q^d}.
\]
Here \(\mathcal F_q\) denotes the space of expansion families
with the given cubical descent law.

In \(\chi=e_j+qz\) the shift in the lower bound of the coefficient
is \(qQ(z)+B(z,e_j)\). The identity
\[
 qQ(z)+B(z,e_j)
       =\frac{Q(e_j+qz)-Q(e_j)}{q}
\]
shows that the normalized bound at \(e_j\) propagates to its
whole coset. Thus the full filtration is given in these coordinates by
\[
 \mathcal U_h(s)\quad\Longleftrightarrow\quad
          v(c_{j,i}(s))\ge h+Q(e_j)/q\quad\hbox{for all }j,i .
\]
For each \(j\), choose \(a_j\in F^\times\) with
\(a_j^q=\pi^{Q(e_j)}\), possible as \(F\) is algebraically closed.
Lift the coordinate vector whose only nonzero entry is \(a_j\) in
position \((j,i)\) to its family. The resulting \(q^{r+d}\) families
are independent over \(F\) and freely generate \(\mathcal U_0\)
over \(\mathcal V\). Division by \(a_j\) gives the integral
coordinates of any element of the unit ball.
\footnote{\href{https://github.com/ChristianKudeba/alpoge-rh-verification/blob/alpoge-source-7cbf7a0c/src/Development/AbelianVarietyTheory/Proofs/RingsOfIntegers/ExistsBasisOfEq.lean}{Valuation in integral coefficient bases};
\href{https://github.com/ChristianKudeba/alpoge-rh-verification/blob/alpoge-source-7cbf7a0c/src/ReciprocityTheory/ReductiveGroups/ArithmeticGeometry/DiscreteValuationRings/ExistsFamRep.lean}{extension from one coset};
\href{https://github.com/ChristianKudeba/alpoge-rh-verification/blob/alpoge-source-7cbf7a0c/src/Development/AbelianVarietyTheory/Proofs/RingsOfIntegers/ExistsBasisOfCoord.lean}{scaled coordinate basis of the unit ball}.}

\paragraph{The level-one frame and its pullback.}
At \(q=1\) there is one coset with representative zero, and its
coefficient module is free of rank one. Choose a \(T\)-generator
\(\sigma\); its reduction is nonzero. Indeed the exact sequence
\(
 0\longrightarrow L_{1,0}\xrightarrow{\pi}L_{1,0}
 \longrightarrow L_{1,0}|_{B_{\kappa(T)}}\longrightarrow0
\)
identifies the kernel of reduction on sections with
\(\pi V_{1,0}\). Extend \(\sigma\) from zero by the period
functional equation to \(\theta\). As \(Q(0)=0\),
\(\theta\in\mathcal U_0\) at level one, and its zero coefficient
is the base change of a primitive integral section. Positivity of
\(Q\) also implies that in the base-field integral expansion all
coefficients outside zero are divisible by \(\pi\).

Let \(\alpha\overline\alpha=q\). In the level-\(q\) pullback the
coefficient at \(\overline\alpha\chi\) is the pullback of the
\(\chi\)-coefficient by \([\alpha]\), using the normalized
identification of its line; all other coefficients vanish.
Here the choice of theta bundle in the PEL uniformization matters.
Write \(L_x\) for its pullback to a point \(x\),
\(\ell(x)=L_x\otimes L_0^{-1}\), and
\[
 \mathcal P(x,y)=L_{x+y}\otimes L_0\otimes L_x^{-1}\otimes L_y^{-1}.
\]
The supplied integral isomorphism for the distinguished element \(w_0\)
identifies \(\ell(x)\) with \(\mathcal P(x,w_0x)\), respecting the
rigidification at zero. This is the extra line-bundle datum behind the
pullback formula. By the Rosati adjoint condition on the biextension,
bilinearity, and commutation of the CM endomorphisms, there are
rigidified isomorphisms over \(T\)
\[
 \ell(\alpha x)\simeq\mathcal P(\alpha x,w_0\alpha x)
  \simeq\mathcal P(x,\overline\alpha w_0\alpha x)
  =\mathcal P(x,qw_0x)\simeq\ell(x)^{\otimes q}.
\]
The coefficient line at \(x\) for \((k,\chi)\) is
\(L_x^{\otimes k}\otimes\mathcal P(x,c(\chi))\).
The same adjoint identity sends its biextension factor under
\(\alpha\) to \(\mathcal P(x,\overline\alpha c(\chi))
=\mathcal P(x,c(\overline\alpha\chi))\). Taking the preceding
isomorphism to the \(k\)-th tensor power and inserting the
rigidification of \(L_0^{\otimes k(q-1)}\) therefore identifies
the pulled-back coefficient line with \(L_{kq,\overline\alpha\chi}\).
The generic generator in the normalized pullback agrees with this
integral identification: their quotient is a global unit on the
proper geometrically connected abelian part over the generic field,
hence a scalar, and their common normalization at the identity sets
it equal to one. This proves the asserted absence of a scalar
valuation loss. The step uses the integral \(w_0\)-isomorphism
as well as Rosati compatibility.
\footnote{\href{https://github.com/ChristianKudeba/alpoge-rh-verification/blob/alpoge-source-7cbf7a0c/src/ReciprocityTheory/ReductiveGroups/ArithmeticGeometry/DiscreteValuationRings/PullGenPf.lean}{Integral square identity, character twist, and comparison of normalized generators}.}
The balanced CM action gives
\[
    2Q(\overline\alpha\chi)
       =B(\overline\alpha\chi,\overline\alpha\chi)
       =B(\chi,\alpha\overline\alpha\chi)
       =2qQ(\chi).
\]
Thus the required normalized bounds agree. Pullback of integral
sections preserves their integrality and gives
\(\theta_\alpha\in\mathcal U_0\) at level \(q\).
\footnote{\href{https://github.com/ChristianKudeba/alpoge-rh-verification/blob/alpoge-source-7cbf7a0c/src/Development/AbelianVarietyTheory/Proofs/JetHarmonics/ExistsThetaFrame.lean}{Primitive level-one frame};
\href{https://github.com/ChristianKudeba/alpoge-rh-verification/blob/alpoge-source-7cbf7a0c/src/ReciprocityTheory/ReductiveGroups/ArithmeticGeometry/DiscreteValuationRings/QuadraticAlgebraValuation.lean}{CM scaling of \(Q\)};
\href{https://github.com/ChristianKudeba/alpoge-rh-verification/blob/alpoge-source-7cbf7a0c/src/ReciprocityTheory/ReductiveGroups/ArithmeticGeometry/DiscreteValuationRings/PullbackValuationInvarianceOf.lean}{forward pullback on the coefficient bounds}.}

To exclude positive depth, suppose \(\mathcal U_h(\theta_\alpha)\)
holds for some \(h>0\). At the zero index, \(Q(0)=0\). Therefore
this coefficient is an integral combination with scalars in the
maximal ideal of \(\mathcal V\), and it vanishes upon reduction
to the residue field of \(\mathcal V\). The normalized pullback
identification at zero is an isomorphism of the integral lines,
so \([\alpha]^*\overline\sigma=0\) on this abelian special fiber.
Further pullback by \([\overline\alpha]\) gives
\([q]^*\overline\sigma=0\). Multiplication by \(q\) on an abelian
variety is finite faithfully flat, so pullback detects a zero
section. But \(\overline\sigma\ne0\): the original reduction was
nonzero and the extension of residue fields is faithfully flat.
This contradiction proves primitivity of the pullback for the
coefficient valuation. Together with
Section~\ref{sec:theta-lift-valuation}, it gives primitivity of
every normalized finite-order translate used in the span argument.
\footnote{\href{https://github.com/ChristianKudeba/alpoge-rh-verification/blob/alpoge-source-7cbf7a0c/src/ReciprocityTheory/ReductiveGroups/ArithmeticGeometry/FirstCohomology/LevelFrameOfLaws.lean}{Pullback of the zero coefficient and descent of divisibility through multiplication by \(q\)}.}

This calculation assumes the PEL degeneration on the given chart,
its rigidified integral \(w_0\) and Rosati comparisons, and the
cocycle of its period multipliers. The scalar and rank calculations
are valid in either residue characteristic once these data are
given. The use of an integral eigenbasis of torsion translates in
Section~\ref{sec:theta-translate-span} additionally requires
\(q\) to be a unit.

\section{Cusp growth, component volume and the degree of the image}
\label{sec:component-volume}
We give the two inequalities behind \eqref{eq:rankshare-vol}. Let
\(m=p+q\), let \(\Phi\) be the selected set of complex embeddings of the
CM field, and set \(D=|\Phi|\binom m2\). For a weight \(w\) set
\[
 V(w)=\prod_{\tau\in\Phi}\prod_{0\le i< j<m}(w_{\tau i}-w_{\tau j}),
 \qquad \eta_{\tau i}=m-1-i .
\]
Only differences of \(\eta\) matter below. On a good deep component \(Z\),
let \(S_Z(rw)\) be the cusp space, taken at scalar twist zero, and put
\(g_Z(w)=D!\limsup_{r\to\infty}\dim S_Z(rw)/r^D\). The source's
extended nonnegative growth agrees with this expression once finite
dimensionality holds. All weights used in the volume comparisons are
strictly decreasing at each embedding.

\paragraph{Nefness and restriction to the boundary.}
On the smooth projective flag compactification let \(L_w\) be the weight
bundle. For a weakly decreasing integer vector \(w_0,\ldots,w_{m-1}\)
at one embedding put \(g_i=w_i-w_{i+1}\) for \(0\le i<m-1\).
Write \(h_i\) for the vector with ones in positions \(0,\ldots,i\) and
zero elsewhere, \(t_i=h_i-\mathbf 1\), and
\(c=w_{m-1}+\sum_{i\ge p}g_i\), where the sum ends at \(m-2\). Then
\[
 w=\sum_{i<p}g_i h_i+\sum_{p\le i<m-1}g_i t_i+c\mathbf 1
\]
with indices in both sums restricted to \(0\le i<m-1\).
Indeed the last coordinate is \(w_{m-1}\) and taking successive differences
recovers every \(g_i\). The source's head and tail bundle estimates give
nefness of \(L_{h_i}\) for \(i<p\) and \(L_{t_i}\) for \(i\ge p\).
Both determinant weights \(\mathbf 1,-\mathbf 1\) give nef bundles.
The weight-to-Picard-class map is additive. On every integral curve all
degrees in this decomposition are nonnegative, using \(|c|\) copies
of the appropriate determinant bundle. Summing over embeddings proves
that \(L_w\) is nef. The head, tail and determinant estimates use the
semiabelian and torus bundles on the compactification.

Let \(B\) be the reduced boundary. The extension and boundary-vanishing
theorems identify the cusp conditions in \(H^0(L_{rw})\). The kernel of
restriction to \(B\) lies in the cusp space, and every cusp section lies
in \(H^0(L_{rw})\). Thus
\[
 h^0(X,L_{rw})-h^0(B,L_{rw}|_B)
 \ \le\ \dim S_Z(rw)\ \le\ h^0(X,L_{rw}) .
\]
Here \(L_{rw}\simeq L_w^{\otimes r}\). The interior is dense, so
\(\dim B<D\). Nef asymptotic Riemann--Roch gives
\(h^0(X,L_w^{\otimes r})/r^D\to L_w^D/D!\), and, on the reduced
boundary, \(h^0(B,L_w^{\otimes r}|_B)/r^D\to0\).
The squeeze thus proves the existence of the cusp growth limit and
\begin{equation}\label{eq:volume-nef-growth}g_Z(w)=L_w^D.\end{equation}

\paragraph{Normalization of the flag degree.}
Use the compactification with toric boundary charts and the source's
Mumford proportionality formula. It identifies \(L_w^D\) with the
covolume \(c_Z\) of the uniformizing group, in the chosen measure,
times the degree polynomial of the compact dual flag. In the source's
alternating coefficient formula, expand the \(D\)-th power of the weight
linear form. For a single embedding its factor, after extracting the common
multinomial numerator \(D!\) across all embeddings, is
\[
 \sum_{\pi\in\mathfrak S_m}\operatorname{sgn}(\pi)
 \prod_{i=0}^{m-1}\frac{w_{\tau i}^{\,m-1-\pi(i)}}{(m-1-\pi(i))!}
 =\frac{\prod_{i<j}(w_{\tau i}-w_{\tau j})}
             {\prod_{j=0}^{m-1}j!}.
\]
The sum is a determinant with descending powers. Reversing its columns
and applying the Vandermonde determinant produces the displayed sign.
Since \(\prod_jj!=\prod_{i<j}(j-i)\), multiplying the embedding factors
gives
\begin{equation}\label{eq:volume-flag}
 g_Z(w)=c_Z D!\,\frac{V(w)}{V(\eta)} .
\end{equation}
The constant \(c_ZD!/V(\eta)\) is strictly positive. To compare two
components one uses the equal-covolume result at the actual deep level.
The good local level at each rational prime is inherited from the
original two-prime construction. Complex uniformization supplies
full-level self-dual data for each component with the same covolume as
that component. Their complex types and signatures agree. Transport to
one Hermitian space preserves the covolume; the self-dual lattice
comparison there identifies the two volumes. In particular the
right side of \eqref{eq:volume-flag} is independent of \(Z\) for fixed
\(w\). This use of the self-dual lattice theorem is separate from the
calculation of the flag degree.

\paragraph{The two inequalities for the image degree.}
Set \(u'=u'(u)\) equal to the lower weight of the window. Write
\(\delta_Z\) for the degree of the birational image of the complete cusp
system at weight \(u\). Then
\begin{equation}\label{eq:volume-sandwich}
 g_{Z'}(u')\ \le\ \delta_Z\ \le\ g_{Z'}(u)
 \qquad\text{for every pair }Z,Z'.
\end{equation}
We prove this first with \(Z'=Z\). Use the function field \(\Omega\)
of the generic component with the source's evaluation trivialization.
Write \(W\) for the image of \(S_Z(u)\) there and form the standard
graded domain \(R=\bigoplus_{j\ge0}W^j X^j\subset \Omega[X]\),
where \(W^j\) denotes the span of products. Multiplication of cusp
forms at these even weights identifies \(W^j\) with a subspace of
the generic evaluations of \(S_Z(ju)\). The Hilbert polynomial of the
projective image has leading term \(\delta_Z j^D/D!\). Dividing the
inclusion of dimensions by \(j^D/D!\) gives \(\delta_Z\le g_Z(u)\).

For the reverse comparison we spell out where the valuative estimate
enters. It supplies \(a\ge1\) and a multiplier \(\theta\) with nonzero
generic value. If \(k\ge1\), \(f\in S_Z(ak u')\), put
\(z=\operatorname{ev}(\theta^k f)\). The valuation inequalities on the
cone imply that \(zX^{ak}\) is integral over \(R\). Birationality and
the ratio condition give \(zX^{ak}\in\operatorname{Frac}(R)\). The latter
condition fixes the field in which normalization is taken. The integral
closure \(\overline R\) is finite over \(R\), since \(R\) is a finitely
generated domain over a number field. Clearing denominators of a finite
set of module generators gives a nonzero \(c\in R\) with
\(c\overline R\subset R\).

Choose a nonzero homogeneous coefficient \(c_e\) of \(c\), so that
\(c_e\in W^e\). The degree \(ak+e\) coefficient of \(czX^{ak}\in R\)
gives \(c_ez\in W^{ak+e}\). Multiplying by \(c_e^{a-1}\) yields
\(c_e^az\in W^{a(k+e)}\). All the multipliers have nonzero generic
value. Evaluation of a cusp form on the generic component is injective
on the space defined on that component. Consequently
\[
 \dim S_Z(ak u')\le \dim W^{a(k+e)} .
\]
Divide by \(k^D/D!\) and pass to the limit. By
\eqref{eq:volume-nef-growth} the left side tends, after normalization,
to \(a^Dg_Z(u')\). The Hilbert polynomial gives the limit
\(a^D\delta_Z\) on the right, since \(e\) is fixed. Cancelling \(a^D\)
proves the other half of \eqref{eq:volume-sandwich}. Equality of
component covolumes now allows \(Z'\) to replace \(Z\) in the growth
factors. In the window \(1-\rho=V(u')/V(u)>0\), so
\(g_Z(u')=(1-\rho)g_Z(u)\) by \eqref{eq:volume-flag}. This gives
\eqref{eq:rankshare-vol}.

\paragraph{Source correspondences and geometric inputs.}
The weight decomposition is \srcid{EqLastAddGaps}; nefness is
\srcid{NefLineBundleWeightGap}. The limit on the boundary and the
normalization of the covolume appear in \srcid{ExistsEffCartierCompl}
and \srcid{FlagDegreeFormula}. The conductor calculation is in
\srcid{ConeDegreeGrowthChowFormIntegrality}, and its application to cusp
forms is \srcid{CuspFormGrowthLeImageDegreeLowerWeight}. The two-sided
comparison is \srcid{ComponentImageDegreeBetweenOfProportionality}.
These arguments use the toroidal model and comparison theorem, nef
asymptotic Riemann--Roch, the comparison of self-dual lattices at good
level, and the existence of the valuative multiplier at \(u,u'\).

\section{Why the uniformizing lattices have equal covolume}
\label{sec:selfdual-covol}
Here we give the index argument used in \eqref{eq:volume-flag}.
Let \(E/F\) be the CM extension with \(F=E^+\) and involution
\(x\mapsto\bar x\). Let \(r=p+q\), with \(p,q\ge1\).
Write \(G\) for the matrix of a nondegenerate Hermitian form.
At the deep level the source takes the positive integer \(M=l.M'\)
(the combined congruence modulus), with \(M\ge3\).
At the chosen embedding, the components under comparison have the
same CM type, signature and archimedean measure.

\paragraph{Reduction to two lattices on one Hermitian space.}
The full-level presentation of each component preserves its covolume
and supplies a trace-self-dual lattice. To pass from two presentations
to a single ambient space the source uses a rational similitude
\(P\) supplied by the classification of the polarized data. In its
normalization, for a move from \(X\) to \(Y\) the ratio of Hermitian
forms satisfies
\[
 h_Y(Pv,Pw)=\frac{\kappa_X}{\kappa_Y}\,h_X(v,w).
\]
Thus the scalar-weighted trace pairing is carried exactly into the
second one, the lattice becomes \(P\Lambda\), and reduction modulo
\(M(P\Lambda)\) is transported by \(P\). In particular
\(P\Gamma P^{-1}\) is the transported group. Its flag-domain
identification preserves the invariant volume. In what follows the
two full lattices, say \(\Lambda,\Lambda'\), have the same \(G\),
CM type, trace parameter and signature. The local self-dual isometry
theorem, assembled over the finite places, supplies
\(g\in U(G)(\mathbb A_{F,f})\) with
\(g\widehat\Lambda=\widehat{\Lambda'}\). The proof only needs a
finite-adelic isometry at this point.

\paragraph{Congruence groups and intersections.}
Write \(\iota\) for the diagonal embedding of rational matrices into
finite-adelic matrices. Let \(K\) be the subgroup of determinant-one
isometries preserving \(\widehat\Lambda\) which act trivially modulo
\(M\widehat\Lambda\). Put \(K'=gKg^{-1}\). The full-level congruence
description identifies \(\Gamma=\iota^{-1}K\) and
\(\Gamma'=\iota^{-1}K'\) with the two uniformizing groups. Here
isometries are taken in the same ambient rational unitary group.
In particular this description applies the determinant-one condition
to the level group; \(g\) can have arbitrary norm-one determinant.

We record how intersection indices pass to the rational points.
If \(J=K\cap K'\), reduction to cosets gives injective maps
\[
 \Gamma/(\Gamma\cap\Gamma')\longrightarrow K/J,\qquad
 \Gamma'/(\Gamma\cap\Gamma')\longrightarrow K'/J ,
\]
where each quotient denotes cosets and carries no normality assumption.
Both maps are surjective. Indeed \(J\) contains a deeper principal
congruence subgroup. Explicitly, denominators of \(g,g^{-1}\)
are bounded on the completed lattice by positive integers \(b,a\).
Then any isometry congruent to the identity modulo \(Mab\) on
\(\widehat\Lambda\) lies in \(K\) and its conjugate back by \(g\)
is congruent to the identity modulo \(M\). This follows by applying
\(g^{-1}(s-1)g\) to a lattice vector and cancelling the two denominators;
the condition for the inverse follows from the congruence subgroup
property. The source obtains rational isotropy of \(G\) from indefinite
signature and the trace-self-dual data and applies strong approximation
for its special unitary group. It gives, for every \(s\in K\cup K'\),
a rational isometry \(\gamma\) with
\(\iota(\gamma)^{-1}s\) in this deeper subgroup. Thus each of the
displayed cosets has a rational representative. The same argument works
with \(g\) replaced by a power, or by a rational isometry.

\paragraph{The ratio of two indices.}
For commensurable groups \(A,B\) write
\[
 r(A,B)=\frac{[A:A\cap B]}{[B:A\cap B]} .
\]
If a third group is commensurable with both, one can compute all ratios
using their triple intersection. This gives \(r(A,B)r(B,C)=r(A,C)\).
Conjugation carries a pair of indices into a pair of the same sizes.
Consequently \(r(K,g^jKg^{-j})=r(K,K')^j\) for positive \(j\).
The indices here are finite since the groups intersect in a deeper open
congruence subgroup of the compact level group.

Put \(z=\overline{\det g}\) and \(s=(zI)g^r\).
Unitarity gives \(z\det g=1\), so \(zI\) is unitary and
\(\det s=z^r(\det g)^r=1\). The central scalar drops out under
conjugation. Apply strong approximation at level \(M\) to \(s\), writing
\(s=\iota(\gamma)k\) for a rational unitary \(\gamma\) and \(k\in K\).
Then
\[
 g^rKg^{-r}=sKs^{-1}=\iota(\gamma)K\iota(\gamma)^{-1}.
\]
The corresponding rational groups are \(\Gamma\) and
\(\gamma\Gamma\gamma^{-1}\). Their quotient volumes in the archimedean
flag domain are equal by the isometry induced by \(\gamma\).
For two commensurable uniformizing groups \(A,B\) acting freely at
this level, the common covering for \(A\cap B\) has volume both
\([A:A\cap B]\operatorname{covol}(A)\) and
\([B:A\cap B]\operatorname{covol}(B)\), with identical measure
normalization. Applying this to the two rationally conjugate groups,
and using the coset bijections, gives \(r(K,g^rKg^{-r})=1\).
The ratio \(r(K,K')\) is positive, so its \(r\)-th power being one
implies \(r(K,K')=1\). Using the coset bijections for \(g\) and then
the common-cover formula yields
\(\operatorname{covol}(\Gamma)=\operatorname{covol}(\Gamma')\).

\paragraph{Source and scope.}
The two reductions to full-level data and a common Hermitian form are
\srcid{ExistsFullLevelGroupDataSameCovolume} and
\srcid{ExistsTransportSelfDualDataOneHermitianSpace}.
The index proof is \srcid{CovolumeEqSelfDualLatticesSameHermitianSpace};
the component application is \srcid{ComplexUniformisationComponentsCovolumeEq}.
The input theorems here are the full-level uniformization, rational
similitude and isotropy, finite-local equivalence of the relevant
trace-self-dual lattices, and special unitary strong approximation.

\section{Counting the sections carried by a component}
\label{sec:rank-share}
Let \(Z\) be a component at the chosen deep level, \(D\) the dimension of
its flag variety, and \(T\) the integer in the weight construction. Let
\(\delta>0\) be the degree of the image of the cusp linear system of
weight \(u\) on \(Z\). The source uses the birationality of this system.
The cusp dimension at weight \(\lambda\) is \(d_\lambda\). We prove
the estimate used in the height calculation:
\begin{equation}\label{eq:rank-share-count}
 \left(1-7(D^2/T+D\rho)\right)d_\lambda
 \ \le\ \binom{T-1}{D}\delta .
\end{equation}
For the terminal weight \(\lambda_T\) there is also
\(\binom{T-1}{D}\delta\le d_{\lambda_T}\).
The upper inequality applies to all the weights satisfying the dominance
conditions below. The lower one uses sections of a particular adjoint bundle.

\paragraph{Geometric data.}
Assume \(p,q\ge1\) and \([E:\mathbb Q]/2>2\), as in the selection of the
component. Here \(D\ge2\), \(T\ge D+1\), and \(0\le\rho\le1/2\).
Write \(v\) for the top-degree cusp growth of \(u\). The proportionality
and valuation-multiplicity comparison gives
\begin{equation}\label{eq:rankshare-vol}
 (1-\rho)v\le\delta .
\end{equation}
This comparison uses the strictly decreasing weights \(u,u'\) and the
equality of component volumes at deep level. Put \(n=T+D+3\). At each
embedding assume that both \(\lambda\) and \(nu-\lambda\) are weakly
decreasing, and that all consecutive gaps of \(u\) are at least
\[
 (D+3)\left(\binom{D+1}{2}+4(p+q)+D+3\right).
\]
The geometric inputs in the source give the following model.
On a projective modification \(X\) of the compactified flag component the
linear system of \(u\) has generated moving bundle \(H\) and effective
fixed divisor \(\Delta\). Write \(U=H+\Delta\) and \(B=L_\lambda\)
in additive line-bundle notation, using pullbacks on \(X\). Both
\(U,B\) and also \(nU-B\) are nef. The morphism of \(H\) is birational
onto the component image, so the projection formula gives \(H^D=\delta\).
The comparison between cusp growth and nef intersection gives \(U^D\le v\).
Cusp forms extend across the compactification and inject into \(H^0(X,B)\).
These assertions about compactification, growth and nefness are the
geometric hypotheses of the following calculation.

\paragraph{Upper bound by hyperplane sections.}
For \(a\ge0\), an integral projective \(D\)-fold, a generated \(H\) of
top intersection \(\delta>0\), and a line bundle \(B\), we claim
\[
 B H^{D-1} < (a+1)\delta
 \quad\Longrightarrow\quad h^0(B)\le\binom{a+D}{D}\delta .
\]
Take a generic member \(X_1\) of the linear system after extending scalars
to the rational function field of its coefficients. Cohomology dimensions
are preserved by this field extension. The generic divisor is integral
over that field when nonempty: on each open where one generating section
is a unit, the universal linear equation eliminates its coefficient.
The incidence scheme is locally an affine bundle over the integral base,
and its generic fiber is a localization. A nonempty section drops the
dimension by one. Restriction to it preserves \(\delta\) as
\((H|_{X_1})^{D-1}\). For dimension one the restriction has length
\(\delta\) over the coefficient field. The exact sequence for the section
bounds \(h^0(B)\) by the sum of \(h^0(B-H)\) and
\(h^0(X_1,B|_{X_1})\).

For clarity, in dimension \(r>0\) use the parameter \(b=a+1\). If
\(BH^{r-1}<0\), a nonzero section would have an effective zero divisor
whose intersection with \(r-1\) nef classes is nonnegative. Thus the
bound for \(b=0\) is zero. Passing to \(B-H\) subtracts \(\delta\)
from the mixed intersection and replaces \(b\) by \(b-1\). Restriction
to \(X_1\) leaves \(b\) fixed and replaces \(r\) by \(r-1\), since
\((B|_{X_1})(H|_{X_1})^{r-2}=BH^{r-1}\) for \(r\ge2\).
In dimension zero the section dimension of an invertible sheaf equals
the length, giving \(\delta\). For \(b\ge1\), induction in \(r,b\)
therefore gives
\[
 \left(\binom{b+r-2}{r}+\binom{b+r-2}{r-1}\right)\delta
   =\binom{b+r-1}{r}\delta
\]
whenever \(b+r\ge2\), using the zero initial value for \(b=0\).
This proves the claim.

On the modification, replacing \(H\) by \(U\) one factor at a time
in a product with other nef factors can only increase the intersection,
because \(U-H=\Delta\) is effective. Nefness of \(nU-B\) gives
\[
 BH^{D-1}\le nUH^{D-1}\le nU^D\le nv .
\]
Take \(a=\lfloor n/(1-\rho)\rfloor\). Inequality
\eqref{eq:rankshare-vol} and \(\delta>0\) imply
\(nv<(a+1)\delta\), including when \(n/(1-\rho)\) is an integer.
Thus \(d_\lambda\le\binom{a+D}{D}\delta\).

\paragraph{The loss in replacing \(a+D\) by \(T-1\).}
Write \(N=a+D\), \(M=T-1\), and \(x=4D/T+\rho\).
We have \(D\le M\le N\), \(a\ge n\), and
\(a(1-\rho)\le T+D+3\). Consequently
\[
 N-M\le \rho a+2D+4\le x(a+1)=x(N-D+1).
\]
Indeed \(2D+4\le4D\) and \(4D(a+1)/T\ge4D\).
If \(x\le1\), then for \(0\le j<D\)
\((M-j)/(N-j)\ge1-x\); multiplying gives
\(\binom{M}{D}\ge(1-Dx)\binom{N}{D}\) by Bernoulli's inequality.
The last inequality also holds when \(x>1\), since its right side is
nonpositive. Finally \(Dx=4D^2/T+D\rho\le7(D^2/T+D\rho)\).
If the coefficient on the left of \eqref{eq:rank-share-count} is
nonpositive the result is immediate. Otherwise multiply the upper bound
for \(d_\lambda\) by that coefficient and apply the binomial estimate.

\paragraph{Lower bound at the terminal weight.}
In the component data, the generated moving line bundle \(H\) has
\(H^D=\delta\). The adjoint sections of
\(K_X+A+TH\) inject into the terminal cusp space under the chosen frame.
The auxiliary divisor \(A\) is nef and big. We use the source's
Kawamata--Viehweg vanishing for
\(K_X+A+jH\), \(j\ge1\), in positive cohomological degrees on the
smooth projective model in characteristic zero. All identifications of
these ordinary sections with subspaces of cusp forms are made before
counting them.

Here is the counting lemma. Let a coherent sheaf \(F\) have support
dimension at most \(r\), and suppose that every positive cohomology group
of \(F(jH)\) vanishes for \(j\ge j_0\). For \(m\ge0\),
\[
 h^0(F((j_0+m)H))\ge \binom{m}{r}(H^r\cdot F).
\]
In dimension zero the dimension of sections equals the length over the
ground field and twisting preserves it. For \(r>0\), choose a section of
\(H\) outside the finitely many associated primes of \(F\). This is
possible over the infinite coefficient field because \(H\) is generated.
Multiplication by it is injective. The quotient \(F_1\) after restriction
has support dimension at most \(r-1\), with
\(H^{r-1}\cdot F_1=H^r\cdot F\).
The long exact cohomology sequence gives the same vanishing for \(F_1\)
starting at \(j_0+1\), and gives
\[
 h^0(F((j+1)H))-h^0(F(jH))=h^0(F_1((j+1)H)),\quad j\ge j_0 .
\]
Summing the dimension-\((r-1)\) bound for \(m\) successive differences,
and using \(\sum_{i=0}^{m-1}\binom{i}{r-1}=\binom{m}{r}\), proves
the claim. Apply it to \(F=K_X+A,\ j_0=1,\ m=T-1,\ r=D\). Twisting
the structure sheaf by an invertible sheaf preserves the top intersection.
The section injection therefore gives \(d_{\lambda_T}\ge
\binom{T-1}{D}\delta\), as required.

\paragraph{Source correspondences.}
In the original \srcid{DivisorProofs} directory the geometric reductions are
\srcid{RankComponentWeightInvariance} and
\srcid{ShareStatementOfKnownTheorems}. The binomial comparison is
\srcid{ChooseMulOneSubLeChoose}. For the lower bound see
\srcid{GlobalSectionRankLowerBoundOfHilbertSamuel},
\srcid{ChooseMulIntersectionNumberLeSectionDim}, and
\srcid{ChooseMulComponentDegreeLeDimension}.

\section{A Chow-form determinant and the arithmetic slope}
\label{sec:chow-determinant}
This section gives the determinant estimate used in the component
Hilbert--Samuel inequality. Let \(F\) be the coefficient number field.
We first state exactly the geometric and metric data used in the
calculation. Let \(V_j\), \(0\le j\le T\), be the graded section spaces
with full integral lattices and Hermitian norms at infinity. Put
\(D\ge1,\ T\ge D+1,\ d=\dim_F V_T\). Multiplication by a coordinate
linear form increases the index \(j\) by one. The system has
bases over \(F\) at indices \(0,D,T\) whose vectors lie in the respective lattices and have archimedean norm bounds
\(\Lambda_0,\Lambda_D,\Lambda_T\), all at least one. The complex
multiplication bound for a unit coordinate vector is \(\sqrt\beta\)
with \(\beta\ge1\).

\paragraph{Filtration at an intersection of hyperplanes.}
The projective image admits an irreducible Chow form \(f(s,l)\),
homogeneous of degree \(\delta\ge1\) in each of \(D+1\) hyperplane
blocks. Here \(s=(s_1,\ldots,s_D)\) and \(l\) is the last block.
There is a nonzero polynomial \(Q(s)\) such that at every \(F\)-point
with \(Q(s)\ne0\) the section splits into \(\delta\) distinct reduced
points \(P_1,\ldots,P_\delta\) over a finite extension. Thus
\[
 f(s,l)=c_s\prod_{\nu=1}^\delta l(P_\nu),\qquad c_s\ne0.
\]
\paragraph{The residue map and the subspace it filters.}
Write \(Y\) for the smooth projective \(D\)-fold over the splitting
field, \(L\) for the generated hyperplane bundle and \(A,L_2\) for
the two auxiliary nef and big bundles. Use additive notation for
tensor products. The data include a section
\(\epsilon\in H^0(Y,A-L_2-DL)\) nonvanishing on the prescribed
open \(U\). If \(Z\) is the common zero scheme of \(s_1,\ldots,s_D\),
the good-section conditions are that \(Z\) be finite, \'etale and
contained in \(U\). In this paragraph its \(\delta\) points are
rational after the indicated finite extension. Define
\[
 G_j=\operatorname{im}\{H^0(Y,K_Y+A+jL)\to H^0(U,K_Y+A+jL)\}.
\]
The coefficient space satisfies \(G_j\subset V_j\) for \(j\le T\)
and \(G_0=V_0\). The inclusion can matter for a positive \(j\).

On \(Z\) the ordered sections define the nowhere-zero Jacobian
\(J_s=ds_1\wedge\cdots\wedge ds_D\) in
\((K_Y+DL)|_Z\). Indeed the conormal map is an isomorphism at
each transverse point. In local frames the degree-\(j\) residue
of \(m\in V_j\) is \(m|_Z/(\epsilon J_s)\), with values in
\((L_2+jL)|_Z\). The choice of frames identifies its values with
one scalar at each of the \(\delta\) points.

We show that the degree-zero residue map is onto. Any prescribed
\(v\in H^0(Z,L_2|_Z)\) gives \(J_s v\) in
\(H^0(Z,(K_Y+L_2+DL)|_Z)\). This latter section lifts to
\(H^0(Y,K_Y+L_2+DL)\). To see the lifting, use the Koszul
resolution of \(\mathcal O_Z\) by the \(D\) equations, tensored by
\(K_Y+L_2+DL\). The term of index \(i\) is the direct sum of
\(\binom Di\) copies of \(K_Y+L_2+(D-i)L\), \(0\le i\le D\).
Each has zero higher cohomology by Kawamata--Viehweg vanishing:
\(L_2\) is nef and big, \(L\) is nef, and \(D-i\ge0\). The complex
is exact off degree zero. Indeed a generator is a unit locally away
from \(Z\), and at points of \(Z\) the equations are a regular system
of parameters. The cohomology of its global sections therefore computes
that of \((K_Y+L_2+DL)|_Z\). This gives the asserted lift \(b\).
Then \(\epsilon b\) restricts to an element of \(G_0=V_0\) whose
residue is \(v\). Equivalently one may first change the equations by
an invertible constant matrix to a sequence regular on successive
quotients, as in the source's restriction argument.

Let \(I\) be the ideal of \(Z\). A first-order vanishing section
\(g\in G_{l+1}\) with \(D\le l+1\) can be written
\[
                        g=\sum_{i=1}^D s_i g_i,\qquad g_i\in G_l .
\]
Its global extension restricts to zero on \(Z\), since \(Z\subset U\).
Apply the same Koszul complex tensored by \(K_Y+A+(l+1)L\).
Its terms have \(L\)-exponents \(l+1-i\ge0\), and the higher
cohomology again vanishes. Exactness of global sections at its
index-zero term identifies the restriction kernel with the image
of the \(D\) indicated multiplications. This proves the division.

We need the version for higher orders with an explicit lower bound
on the remaining degree. For \(j\ge D\), \(a\ge0\), every
\(g\in G_{j+a}\) vanishing to order \(a\) along \(Z\) admits
\[
                       g=\sum_{|\alpha|=a}s^\alpha g_\alpha,
                       \qquad g_\alpha\in G_j .
\]
For \(a=0\) this is immediate. For order \(a+1\), apply the
order-\(a\) result with coefficient degree \(j+1\).
Locally along \(Z\) the associated graded algebra of \(I\) is the
polynomial algebra on the conormal basis given by the \(s_i\);
this is the regular-sequence description
\(\operatorname{gr}_I\mathcal O_Y=\operatorname{Sym}_{\mathcal O_Z}
(I/I^2)\). Thus the degree-\(a\) symbol of the resulting expression
has a separate coefficient \(g_\alpha|_Z\) for each monomial,
in local frames. Since the section vanishes to order \(a+1\),
every such coefficient is zero. Apply the first-order division
above to each \(g_\alpha\in G_{j+1}\). Regrouping the terms with
the same multi-index \(\alpha+\mathbf e_i\) gives the assertion.

Write \(r=T-D\) and set
\[
 P_s=\sum_{|\alpha|=r}s^\alpha V_D,\qquad W_s=G_T+P_s\subset V_T,
 \qquad U_a=\{h\in W_s:h\text{ vanishes along }I^a\}.
\]
Here vanishing means membership in \(I^a\) in local bundle
coordinates along \(Z\). Stability under multiplication gives the
inclusion \(P_s\subset V_T\). Now \(P_s\subset U_r\) by construction.
Conversely, for \(h=g+p\in U_r\) with \(g\in G_T\), \(p\in P_s\),
the difference \(g=h-p\) vanishes to order \(r\). Higher-order
division with \(j=D\) expresses \(g\) in \(P_s\), proving \(U_r=P_s\).

For \(a<r\), the kernel of the symbol on \(U_a\) is \(U_{a+1}\)
by the same associated graded description. Its target after choosing
frames is \(\bigoplus_{|\alpha|=a} F'^{\,\delta}\) over the splitting
field \(F'\). To show surjectivity, choose \(\ell\) nonvanishing at
all the points, possible over the infinite field \(F'\).
For each \(\alpha\) evaluate the symbol of
\(s^\alpha\ell^{T-a}m\) with \(m\in V_0=G_0\). It vanishes in
every other monomial coordinate, and its \(\alpha\)-coordinate is the
residue vector of \(m\) multiplied at point \(P_\nu\) by
\(\ell(P_\nu)^{T-a}\). These nonzero factors and surjectivity of the
degree-zero residue prove that this block spans its entire target.
All these test sections lie in \(G_T\subset W_s\).
Writing \(y_s=\dim P_s\), summing graded dimensions now gives
\[
            \dim W_s=y_s+\delta\binom{T-1}{D},\qquad
            \dim(V_T/W_s)=d-y_s-\delta\binom{T-1}{D}.
\]
This identifies the exact dimension of the subspace requiring
additional top-basis columns in the determinant. The argument uses
ordinary adjoint sections on the compactification only in the
subspaces \(G_j\), and uses the actual \(V_D\) in \(P_s\).

Let \(R=\binom{T-1}{D}\). For \(0\le a<T-D\) the number of
monomials in \(D\) variables of degree \(a\) is
\(c_a=\binom{a+D-1}{D-1}\). Stars and bars and Pascal summation give
\(\sum c_a=R\). The total exponent of the last form on one residue
point in the determinant below is \(b=\sum_a c_a(T-a)\). Indeed
\(a c_a=D\binom{a+D-1}{D}\), so, taking the sum with \(a<T-D\),
\[
 b=TR-D\binom{T-1}{D+1},\qquad
 (D+1)b=R\bigl(T+D(D+1)\bigr).
\]
These are exact finite identities, including \(T=D+1\).

Choose a rational good \(s_0\) for which the deepest-step dimension
\(y\) is maximal. Such a point exists because \(F\) is infinite, \(Q\)
is nonzero, and the possible dimensions are integers between zero and
\(d\). After splitting the section at \(s_0\), select \(\delta\)
vectors of the short basis of \(V_0\) with independent evaluations.
Choose \(l\) avoiding all the points. For each of the above monomials
\(s^\alpha\) take the \(\delta\) columns \(s^\alpha l^{T-a}m_i\),
where \(m_i\) are the selected degree-zero sections. Their symbols are
independent: at each monomial the invertible evaluation matrix is
multiplied by a diagonal matrix with nonzero entries \(l(P_\nu)^{T-a}\).
The different symbols lie in independent summands or successive
filtration grades. Include \(y\) columns of the form \(s^\alpha m\),
\(|\alpha|=T-D,\ m\) a short degree-\(D\) basis vector, spanning the
deepest step. Complete with \(d-R\delta-y\) vectors from the short top
basis. These choices can be made by selecting from fixed rational
generating sets over the splitting field. After replacing \(s_0,l\)
by variables, their determinant \(P(s,l)\) is therefore a nonzero
polynomial over \(F\).

\paragraph{Divisibility before taking heights.}
At any other good \(s\), the dimension of the deepest step is at most
\(y\). If it is smaller, the \(y\) degree-\(D\) columns are dependent
and the determinant vanishes identically in \(l\). If the dimension
equals \(y\), either those columns are dependent, or they give a basis
of the deepest step. In the latter case take symbols successively
modulo the filtration. Only the \(R\delta\) columns from \(V_0\)
depend on the last block on the successive symbol spaces. The
determinant of their block for monomial \(\alpha\) is a scalar
independent of \(l\) times \(\prod_\nu l(P_\nu)^{T-a}\).
The remaining minors are independent of \(l\). This block determinant
identity holds for all \(l\) by polynomial identity. It includes the
possibility of zero scalar if a specialized residue matrix or final minor
loses rank. Consequently for every good \(s\),
\(P(s,l)=\gamma_s f(s,l)^b\), with \(\gamma_s\) in a splitting field.

To descend the divisibility, choose \(l_0\in F^N\) such that
\(g(s)=f(s,l_0)\) is a nonzero polynomial. Evaluation at \(l_0\)
and at \(l\) gives \(P(s,l)g(s)^b=f(s,l)^bP(s,l_0)\) on \(Q(s)\ne0\).
Multiplying the difference by \(Q\) gives a polynomial vanishing on
all rational tuples. The field is infinite, so the polynomial is zero;
the polynomial ring is a domain, so \(Q\) cancels. The polynomial \(f\)
has positive degree in \(l\) whereas \(g\) is independent of \(l\).
Irreducibility of \(f\) therefore gives \(\gcd(f,g)=1\).
In the polynomial unique factorization domain we obtain
\(P=f^bG\) for a nonzero multihomogeneous \(G\).

\paragraph{The lattice estimate.}
Let \(e\) be the sum of the multidegrees of \(G\). The sums of column
degrees on the two sides of this factorization give
\begin{equation}\label{eq:chow-column-degrees}
 b(D+1)\delta+e=R\delta T+y(T-D),\qquad R\delta+y\le d .
\end{equation}
Write \(H_f,H_G\) for the source's heights of the two forms in the
chosen coordinate frame and \(A_T\) for the arithmetic degree of
the Hermitian top lattice. The normalized multiform height is additive
on products. Its lower bound is minus \(c\) times the sum of block
degrees, with \(c=\log(n\Lambda)+(\log n+1)/2\), where \(n\ge1\)
counts the spanning coordinates and \(\Lambda\ge1\) bounds their
lattice data. Hence \(H_f\ge-c(D+1)\delta\), \(H_G\ge-ce\).
At each finite place the substituted polynomial vectors have integral
lattice coefficients. At each complex place the three types of columns
have norm bounded respectively by
\(\beta^{T/2}\Lambda_0,\ \beta^{(T-D)/2}\Lambda_D,\ \Lambda_T\)
for unit input vectors. The determinant inequality for polynomial
vectors in a Hermitian lattice gives
\begin{align}
 A_T\ \ge\ bH_f+H_G
   &-R\delta\log(\beta^{T/2}\Lambda_0)
       -y\log(\beta^{(T-D)/2}\Lambda_D) \notag\\[-2pt]
   &-(d-R\delta-y)\log\Lambda_T .\label{eq:chow-slope-raw}
\end{align}
This is a statement about the degree of the actual lattice, so possible
finite index of the columns in it is already included. The local
determinant estimate at infinity uses the exterior-product norm
inequality; summing the logarithms with the finite ideal contributions
uses the product formula. The stated multiform height bounds and this
lattice-height comparison are the metric inputs here.

Put \(h_f=H_f/((D+1)\delta)\) and \(z=1-R\delta/d\).
Equation \eqref{eq:chow-column-degrees} shows that \(d>0\),
\(0\le z\le1\) and \(e=y(T-D)-R\delta D(D+1)\le T(d-R\delta)\).
Furthermore
\(bH_f=R\delta (T+D(D+1))h_f
\ge R\delta(Th_f-D(D+1)c)\), because \(h_f\ge-c\).
All the logarithmic losses are nonnegative. In
\eqref{eq:chow-slope-raw} we may bound the loss on the final
\(d-R\delta\) columns by charging each
\(T\log\beta/2+\log\Lambda_D+\log\Lambda_T\). Division by \(d\) gives
\begin{align*}
 A_T/d \ \ge\
 &(1-z)\bigl(Th_f-D(D+1)c-\log\Lambda_0-(T/2)\log\beta\bigr)\\
 &-z\bigl(Tc+(T/2)\log\beta+\log\Lambda_D+\log\Lambda_T\bigr).
\end{align*}
This is the intermediate slope inequality in the source.

\paragraph{Comparison with the component height.}
The frame comparison gives \(h_f\ge h-(\log n+1)/2\), where \(h\)
is the geometric height used in Section~\ref{sec:hs-component}. Set
\[
 L=\log\Lambda+26\log\!\bigl(nD\delta(2+|h|+\log\Lambda)\bigr).
\]
Since \(nD\delta(2+|h|+\log\Lambda)\ge2n\), we have \(c\le L\).
In the first bracket of the slope inequality we may replace
\(Th_f\) by \(Th-T(\log n+1)/2\).
The three bounds
\[
 D(D+1)c\le26D^2(L+\log T),\qquad
 (T/2)\log\beta\le(T+D+1)\log\beta,
\]
\[
 \frac{T(\log n+1)}2
 \le (T+D^2)\left(\frac{(D+1)(\log n+1)}2+3\right)
\]
allow the additional subtraction \(3(D+1)|h|\) used by the component
formula. In the second bracket replace \(c\) by \(L+\log2\).
Both replacements weaken a lower bound since \(z,1-z\ge0\).
With \(J=\log\Lambda_0\) these are exactly the two terms of
\srcid{HSClaimS} at constant \(26\). The inputs on the geometric side are Bertini for the simple hyperplane
sections, adjoint vanishing on the smooth projective model and the
short-basis assertions above. Identifying the
birational image degree, its Chow height and the Hermitian lattice
degree transfers the inequality to the component of the window.

\paragraph{Source correspondences.}
The restriction and filtration are
\srcid{HilbertSamuelDegreeZeroResiduesRangeEqTop},
\srcid{FilMonoOfKerRes}, \srcid{FilSpaceVEqSpaceP} and
\srcid{HilbertSamuelVanishingFiltrationDeepestStepSpan}.
The selection of columns and the specialization identity are
\srcid{ExistsShortColumnsDetNeZeroHilbertSamuel} and
\srcid{DetRationalSectionsEqPowProductLinearForms}. The descent is
\srcid{PowDvdOfSpecialisations}; the weighted count and its numerical
use are \srcid{MonomialCountsDegreeAndWeightedSum} and
\srcid{ClaimSOfNumbers}. See also \srcid{SlopeLowerBoundTopLatticeOffHypersurface},
\srcid{ExplicitArithmeticHilbertSamuelBound}, and
\srcid{HilbertSamuelComponentShortBasisBound}.

\section{Determinant ideals and spherical heights}
\label{sec:chow-metric}
We prove the multiform estimates and the norm-to-degree step in
\eqref{eq:chow-slope-raw}. Let \(F\) be a number field, \(e_F=[F:\mathbb Q]\),
and \(w:F^n\twoheadrightarrow F^N\) a spanning coordinate matrix with
\(1\le N\le n\). At each embedding \(\sigma:F\hookrightarrow\mathbb C\)
let \(A_\sigma\) be an invertible frame, so the norm of \(x\in\mathbb C^N\)
in the coordinate space is \(\|A_\sigma^{-1}x\|_2\).
Assume every column of \(\sigma w\) has frame norm at most
\(\Lambda\ge1\). Form heights will use the mean over \emph{all} \(e_F\)
embeddings, counting a conjugate pair twice.

For a nonzero polynomial \(g\) with \(B\) blocks of \(N\) variables let
\(J(g)\) be the fractional ideal generated by the coefficients of the
polynomial obtained by substituting \(w\) in every block. The substituted
form is nonzero because \(w\) has a right inverse. Write \(m_\sigma(g)\)
for the expectation of \(\log|(\sigma g)(A_\sigma z_1,\ldots,A_\sigma z_B)|\),
where the \(z_i\) are independent uniform points on the unit sphere
in \(\mathbb C^N\), of real dimension \(2N-1\). Set
\[
 H(g)=-\frac{\log \mathrm N J(g)}{e_F}
               +\frac1{e_F}\sum_\sigma m_\sigma(g).
\]
This uses the norm on fractional ideals extended multiplicatively.
For a multihomogeneous form, its nonzero restriction to the product of
spheres has logarithm locally integrable and its zero set has measure
zero. One may use logarithmic integrability of a nonzero holomorphic
polynomial in complex affine charts: homogeneity reduces the statement
to the compact projective directions, and the radial and phase factors
are smooth in every unit chart. All expectations below are consequently
well defined. At a finite prime the minimum coefficient valuation is
additive on products: after dividing each factor by its Gauss norm
in a local uniformizer, the reductions have nonzero product over the
residue field. Hence \(J(gh)=J(g)J(h)\). At infinity the logarithmic
means add. The product formula gives \(H(cg)=H(g)\) for \(c\in F^\times\).

\paragraph{An evaluation bound including its constant.}
Let \(g\) have block degrees \(a_1,\ldots,a_B\), and put
\(a=\sum_i a_i\), \(H_{N-1}=\sum_{k=1}^{N-1}1/k\). For unit vectors
\(v_i\in\mathbb C^N\), unitary invariance and Jensen's inequality on
a circle give
\[
 \log|(\sigma g)(A_\sigma v_1,\ldots,A_\sigma v_B)|
      \le m_\sigma(g)+a H_{N-1}/2
\]
whenever the value on the left is nonzero. To see the constant, rotate
each \(v_i\) to the first coordinate, changing the polynomial by the
inverse unitary change of variables. In one block write a random vector
as \(z=(z_1,z')\). For fixed \(z_1,z'\), average
\(\log|g(z_1,e^{it}z')|\) over \(t\). The submean inequality bounds
this average below by \(\log|g(z_1,0)|\), allowing \(-\infty\) for
a vanishing evaluation. Repeat over all blocks. The left average
integrated over the spheres is unchanged by rotation, while the final
value equals the chosen evaluation times
\(\prod_i z_{i1}^{a_i}\). For \(N>1\), \(|z_1|^2\) has density
\((N-1)(1-u)^{N-2}\) on \(0<u<1\); this follows by normalizing
\(N\) independent standard complex Gaussians, whose squared radii
are independent exponentials. Its logarithmic loss is
\[
 \mathbb E[-\log |z_1|^2]
    =\int_0^1\frac{1-(1-t)^{N-1}}t\,dt
    =\sum_{j=0}^{N-2}\int_0^1(1-t)^j\,dt=H_{N-1}.
\]
For \(N=1\) the loss is zero. This proves the claim.

On the unit polydisk of the \(n\) substitution variables in each block,
the substitution \(\sigma w\) produces vectors of frame norm at most
\(n\Lambda\), by the triangle inequality. Homogeneity and the bound just
proved give \(|(\sigma g)(\sigma w z)|\le
\exp(m_\sigma(g)+a H_{N-1}/2)(n\Lambda)^a\).
By the coefficient formula on the unit torus the same bound holds for
each coefficient of \((\sigma g)(\sigma w z)\).
Take any nonzero coefficient \(\alpha\) of \(g(wz)\).
The inclusion \((\alpha)\subset J(g)\) gives
\(\log \mathrm N J(g)\le\sum_\sigma\log|\sigma\alpha|\).
Using the coefficient bounds and \(H_{N-1}\le1+\log n\) yields
\[
 H(g)\ \ge\ -a\left(\log(n\Lambda)+\frac{\log n+1}{2}\right).
\]
When \(g=f\) is a Chow form with \(D+1\) blocks, all of degree \(\delta\),
the conventional form height in these frames adds
\((D+1)\delta H_{N-1}/2\) to \(H(f)\).
Dividing by \((D+1)\delta\) gives precisely
\(h_f=h-H_{N-1}/2\ge h-(1+\log n)/2\), with the height \(h\) of the
image. Replacing its defining form by a scalar multiple or a positive
power leaves the normalized height unchanged when the degree is adjusted.

\paragraph{Lattice degree from polynomial columns.}
Let \(M\) be a full \(\mathcal O_F\)-lattice in a \(d\)-dimensional
\(F\)-space, \(b\) an \(F\)-basis and \(G_\sigma\) its positive Hermitian
Gram matrix. Define \(I_b(M)\) as the fractional ideal spanned by the
determinants of all \(d\)-tuples of lattice elements relative to \(b\).
With this convention
\[
 \widehat{\deg}(M)=-\frac{\log\mathrm N I_b(M)}{e_F}
                 -\frac1{2e_F}\sum_\sigma\log\det G_\sigma .
\]
Indeed, after localization at a finite prime the lattice is free, and
the local determinant ideal records the sum of its elementary-divisor
valuations relative to the basis lattice. Their sum weighted by residue
field logarithms gives the finite term. At infinity the norm of
\(b_1\wedge\cdots\wedge b_d\) is \(\sqrt{\det G_\sigma}\).
Under a change of basis by \(S\), the ideal is divided by \(\det S\)
and the Gram determinant is multiplied by \(|\sigma\det S|^2\);
the product formula cancels the two changes. This also covers a
projective lattice with nonprincipal determinant ideal.

Let \(P\) be a square matrix of polynomials, whose columns are vectors
in the \(b\)-coordinate space, and set \(g=\det P\ne0\).
After substituting \(w\), suppose every coefficient vector of every
column lies in \(M\). Expanding by multilinearity, the coefficient
of each monomial in the determinant is a sum of determinants of such
coefficient vectors. Hence \(J(g)\subset I_b(M)\).
At a complex place suppose column \(k\) has Hermitian norm at most
\(B_k>0\) on unit input vectors in every frame block. Hadamard's
inequality in the Hermitian inner product gives pointwise
\[
 |\det P|\,\sqrt{\det G_\sigma}\le\prod_{k=1}^{d}B_k.
\]
Taking the logarithm almost everywhere and integrating gives
\(m_\sigma(g)\le\sum_k\log B_k-\tfrac12\log\det G_\sigma\).
Since \(J(g)\subset I_b(M)\) implies
\(-\log\mathrm N J(g)\le-\log\mathrm N I_b(M)\),
summing these inequalities proves
\[
 H(\det P)-\sum_k\log B_k\le\widehat{\deg}(M).
\]
Taking the three families of columns and their bounds from
Section~\ref{sec:chow-determinant} proves \eqref{eq:chow-slope-raw},
including its finite-place contribution.

\paragraph{Source correspondences.}
The lower bound and harmonic correction are
\srcid{MultiformHeightLowerBound} and \srcid{FormHeightEqHt}.
The determinant and its norms are
\srcid{LatticeDegreeGeHeightPolynomialVectorsDet},
\srcid{CoeffDetMemSpanDetIntegralPolynomialVectors}, and
\srcid{FiniteHeightSubMeanLeDegreeFullLatticeDetIdeal}.

\section{An effective Hilbert--Samuel estimate on one component}
\label{sec:hs-component}
We derive the coefficient of \(\log(2+|h|)\) and the uniform additive
loss in \eqref{eq:cov-hs}. Fix a component of geometric degree
\(\delta\ge1\), height \(h\), and positive rank \(d\). Let its
unshifted and shifted arithmetic degrees be \(a_0,a_-\).
Write \(D\ge1\) for the flag dimension, \(T\ge3D+5\) for the length,
\(P\ge1\) for the size parameter and \(b\ge1\) for the error scale.
The hypotheses used below are
\begin{align*}
 &\log b\le P,\quad TP\le\sqrt b,\quad \rho,\kappa\ge0,\quad
 z=7(D^2/T+D\rho)\le\tfrac12,\quad zP\kappa\le b,\\
 &Q=\binom{T-1}{D}\delta,\qquad (1-z)d\le Q,\qquad
 a_-\ge a_0-dc,\quad c\ge0 .
\end{align*}
On this component the geometric comparison gives numbers
\(n,\Lambda,\beta,\Lambda_D,\Lambda_T\ge1\) and \(J\) such that
\begin{align*}
 &\log n,\log\beta,\log\delta\le P,\qquad
 \log\Lambda\le P(1+\kappa/T),\\
 &\log\Lambda_D,\log\Lambda_T\le P\kappa,\qquad
 J\le 4(DP(1+\kappa/T)+TP)+2DP\kappa/T .
\end{align*}
We spell out the geometric Hilbert--Samuel inequality used as input.
Set
\begin{align*}
 r&=\log\Lambda+26\log\big(nD\delta(2+|h|+\log\Lambda)\big),\\
 G&=(D+1)(\log n+1)/2+1,\\
 U&=Th-3(D+1)|h|-J-26D^2(r+\log T)
              -(T+D+1)\log\beta-(T+D^2)(G+2),\\
 V&=T(\log2+r+\tfrac12\log\beta)+\log\Lambda_D+\log\Lambda_T .
\end{align*}
The component assertion with constant \(26\) and initial multiplicity
one says \(0\le1-Q/d\le1\) and
\[
             a_0\ge Q U-(d-Q)V .
\]
It comes from an estimate for sections on the geometric component.
The following argument turns it into a degree bound with just one
height-dependent error and a constant independent of the component.

Put \(q=\kappa/T\) and \(l_0=27P(q+1)+65P\). Since \(T\le\sqrt b\)
we have \(\log D\le\log T\le P/2\). Also
\(\log(1+\log\Lambda)\le\log\Lambda\); hence
\[
 \log\Lambda+26\log\big(nD\delta(1+\log\Lambda)\big)\le l_0 .
\]
Indeed the three size logs contribute at most
\(26(P+P/2+P)=65P\), and the remaining terms at most
\(27\log\Lambda\). The inequality
\(2+|h|+\log\Lambda\le(2+|h|)(1+\log\Lambda)\) now yields
\(0\le r\le l_0+26\log(2+|h|)\).
From the rank share we have \(d\le2Q\) and \(d-Q\le zd\le2zQ\).
Since \(V\ge0\) the geometric comparison bounds \(a_0\) below by
\(Q(U-2zV)\). Substitution of the upper bound for \(r\) in its
negative terms gives
\begin{equation}
 a_-\ge Q\{Th-3(D+1)|h|-A\log(2+|h|)-H-2c\},\quad
       A=26(26D^2+2zT),                         \label{eq:hs-stage}
\end{equation}
where the complete constant \(H\) is
\begin{align*}
 H={}&J+26D^2(l_0+\log T)+(T+D+1)\log\beta\\
     &+(T+D^2)(G+2)
       +2z\{T(\log2+l_0+\tfrac12\log\beta)+\log\Lambda_D+\log\Lambda_T\}.
\end{align*}
In particular the rank defect multiplies the norm cost \(V\). Ignoring
it would lose the \(2zT\) term in \(A\).

Here are estimates for all five summands. The relations defining \(z\)
imply \(D^2\le zT/7\le T/14\), and therefore
\[
 D^2Pq\le b/7,\qquad D^2P\le b/14,\qquad
 DPq\le b/7,\qquad TP\le b,\qquad (TP)^2\le b .
\]
The first inequality follows by multiplying \(D^2\le zT/7\) by \(Pq\).
\begin{enumerate}
\item Expanding the bound for \(J\) gives \(6DPq+4DP+4TP\le9b\).
\item Since \(l_0+\log T\le27Pq+(185/2)P\), its contribution times
\(26D^2\) is at most \(26(27/7+185/28)b<300b\).
\item The logarithm of \(\beta\) contributes at most \(2TP\le2b\).
\item We have \(G+2\le5DP\) because \(P,D\ge1\) and \(\log n\le P\).
Also \(TD P\le (TP)^2\le b\) and \(D^2\le T\).
This gives at most \(10b\).
\item In the final term \(zTPq=zP\kappa\le b\). Both \(z\log\Lambda_D\)
and \(z\log\Lambda_T\) are bounded by \(b\). Moreover
\(zT\log2\le b/2\), \(zTP\le b/2\), and \(zT\log\beta/2\le b/4\).
Inserting \(l_0=27Pq+92P\) bounds the expression, including its
outer factor \(2\), by \(2(1/2+27+46+1/4+2)b<200b\).
\end{enumerate}
The five upper bounds give \(H\le4000b\).
Thus \eqref{eq:hs-stage} proves \eqref{eq:cov-hs}
with \(R=4000b+2c\).

There is also a cost for \(A\) in the subsequent component sum.
We have \(1\le A\le702T^2\), since \(D\le T\), \(2z\le1\) and \(T\ge1\).
Consequently \(\log A\le\log702+2\log T\le 8P\). Hence
\[
 A\log(\max(1,A))\le5616T^2P\le6000b .
\]
In the actual family the twist comparison supplies \(c=C_p Z_0^2\)
with \(C_p\ge0\) uniform and \(Z_0\) the base size. Its size conditions
give \(1\le Z_0,\ Z_0^{12}\le P,\ TP\le\sqrt b\), so \(Z_0^2\le b\).
As \(b\le E\), this proves
\(A\log(\max(1,A))+R\le(10000+2C_p)E\) with the same constant on the
common windows. The geometric component assertion and the twist
comparison, whose statements enter explicitly above, are the inputs
to this calculation.
\footnote{\href{https://github.com/Deicyde/riemann-hypothesis/blob/40e56bab/src/Formalism/LocalToGlobal/IndefiniteForms/IndefiniteTree/RealCubicMeans/LogMeanShares/LogMeanStatements/HSClaim.lean}{Component inequality};
\href{https://github.com/Deicyde/riemann-hypothesis/blob/40e56bab/src/Development/AbelianVarietyTheory/Proofs/DivisorProofs/DegBoundOfClaimNumbers.lean}{five errors and the twist};
\href{https://github.com/Deicyde/riemann-hypothesis/blob/40e56bab/src/Development/AbelianVarietyTheory/Proofs/DivisorProofs/HilbertSamuelALogALe.lean}{the logarithmic cost};
\href{https://github.com/Deicyde/riemann-hypothesis/blob/40e56bab/src/Development/AutomorphicArithmetic/ThetaCorrespondence/ThetaPairLifts/ThirdPlaceSeries/HilbertSamuelLowerBoundOfComponents.lean}{specialization at the common threshold}.}

\section{A lower bound for the slope across the components}
\label{sec:cov-heights}

The geometric argument compares heights of individual components with the
arithmetic slope of the space of sections on their union. Both the error in
their relative degrees and the error in their ranks enter this comparison.
We give the finite-sum argument used in the source, including the constant
with which it enters the covolume estimate.

Fix a nonempty finite set \(I\) of components. Write \(h_c\) for a component
height, \(s_c\) for its comparison slope, \(\delta_c\) for its geometric
degree, \(d_c\) for its rank and \(a_c\) for its arithmetic degree. Let
\(D\ge0,\ T\ge3D+5,\ A>0,\ B>0,\ \delta_*>0,\ E\ge1\). The
parameters \(u,v,X,R\) are nonnegative, and \(0\le z,\rho<1\).
Assume at every \(c\in I\) that
\begin{align}
 \delta_*\le\delta_c\le\delta_* /(1-\rho),\qquad
 B\delta_c\le d_c,\qquad (1-z)d_c\le B\delta_c, \label{eq:cov-shares}\\
 h_c+u\ge0,\qquad
 (1-\rho)\max(0,s_c+u)\le h_c+u+v, \label{eq:cov-shift}\\
 a_c\ge B\delta_c
 \bigl(Th_c-3(D+1)|h_c|-A\log(2+|h_c|)-R\bigr).\label{eq:cov-hs}
\end{align}
Assume also \(|I|^{-1}\sum s_c\ge-X\). If
\begin{align}
 TX\le C_1 E,\qquad Tv\le C_1 E,\qquad
 u\bigl(6D+10+T(z+2\rho)\bigr)\le C_1E, \nonumber\\
 A\log(\max(A,1))+R\le C_2 E,\qquad
 2\frac{\sum a_c}{\sum d_c}\le\mu,\label{eq:cov-budgets}
\end{align}
then
\begin{equation}
                  \mu\ge-2(3C_1+C_2+4)E. \label{eq:cov-u}
\end{equation}

\paragraph{Proof.}
Set \(m=\max(0,(1-\rho)u-X-v)\) and \(n=|I|\). Averaging
\eqref{eq:cov-shift} gives
\[
 \frac1n\sum(h_c+u)\ge (1-\rho)(u-X)-v
 \ge (1-\rho)u-X-v.
\]
The left side is also nonnegative, so it is at least \(m\).
For \(t\ge0\) one has
\[
 A\log(2+t)\le t+A\log(\max(A,1))+2 .
\]
If \(A\ge1\), apply \(\log y\le y-1\) to \(y=(2+t)/A\). If
\(0<A<1\), use \(\log(2+t)\le1+t\) and \(A(1+t)\le1+t\).
Put \(P=T-3D-4\), \(K=A\log(\max(A,1))+2\), and
\[
 N=Pu+2u(3D+4)+K+R,\qquad W=\sum B\delta_c,\qquad S=\sum d_c.
\]
Since \(h_c\ge-u\), one has \(|h_c|\le h_c+2u\). Substituting both
inequalities in \eqref{eq:cov-hs} gives
\[
 a_c\ge B\delta_c\bigl(P(h_c+u)-N\bigr).
\]
Here \(P\ge0\). In the nonnegative term replace \(\delta_c\) by
\(\delta_*\) and sum. The upper bound on \(\delta_c\) in
\eqref{eq:cov-shares} gives \(nB\delta_*\ge(1-\rho)W\). Hence
\[
 \sum a_c\ \ge\ B\delta_* Pnm-NW
       \ \ge\ W\bigl(P(1-\rho)m-N\bigr),\qquad 0<W\le S .
\]
As \(m\ge(1-\rho)u-X-v\),
\((1-\rho)m\ge(1-2\rho)u-X-v\). Subtract \(N\) and use
\(0\le P\le T\) on the terms with a minus sign. The result is
\[
 P(1-\rho)m-N\ \ge\ -T(X+v)-u(6D+10+T(z+2\rho))
              -A\log(\max(A,1))-4-R=:J .
\]
In fact the estimate before this last enlargement has \(6D+8\) in
place of \(6D+10\), no \(Tz\), and \(2\) in place of \(4\).
By construction \(J\le0\). If \(P(1-\rho)m-N\ge0\) then
\(\sum a_c/S\ge0\). Otherwise multiplication by \(W/S\le1\)
increases this negative quantity, so \(\sum a_c/S\ge J\) also in
that case. Substituting \eqref{eq:cov-budgets} and using \(4\le4E\)
proves \eqref{eq:cov-u}.
\footnote{\href{https://github.com/Deicyde/riemann-hypothesis/blob/40e56bab/src/Formalism/MeanHeightDescent/AverageDegreeBounds/SumComponents.lean}{Component sum};
\href{https://github.com/Deicyde/riemann-hypothesis/blob/40e56bab/src/Development/AbelianVarietyTheory/Proofs/DivisorProofs/UniformLowerBoundOfSums.lean}{specialization and common threshold}.}

At a given arithmetic point and a nonterminal lift index \(i\), the cover
supplies the finite set \(I\). For a quotient weight \(u_1\) the source
takes \(h_c\) and \(\delta_c\) at \(u_1\), and \(a_c,d_c\) at the
lower weight \(u_-\). In its notation, with \(P_0\) the size parameter,
\[
 T=\left\lceil P_0^2D^2 k/m_0\right\rceil,\qquad
 B=\binom{T-1}{D},\qquad z=7(D^2/T+D\rho),
\]
where \(D\) is the flag dimension, \(k\) is the weight parameter,
and \(m_0\) is the number of lifts. Here \(D,T\) are integers and
\(T\ge3D+5\) gives \(B>0\). One first chooses \(u_1\) among the
quotient weights guaranteed by the size condition. The degree and
rank-share estimates give \eqref{eq:cov-shares} for this same weight.
The component height estimate gives \eqref{eq:cov-shift} and the first
three budgets; the arithmetic Hilbert--Samuel estimate gives
\eqref{eq:cov-hs} and the fourth. The comparison for the sum of the
component degrees gives the last inequality of \eqref{eq:cov-budgets},
with \(\mu\) the cover slope at the weight indexed by \((u_-,0)\).
The seven uniform inputs each hold beyond a threshold. Their maximum
is a common threshold. Thus the constant \(C_U=2(3C_1+C_2+4)\)
works for every such cover on the same family.

\section{From the cover to the covolume of the theta lattice}
\label{sec:cov-lattice}

Fix a point, a lift index \(i\), and a cover at the required depth.
Write \(M\) for its main term, \(E\) for its error unit and \(b\) for
the scale in the norm bounds. Put \(K=P_0\kappa_i\), where \(\kappa_i\)
is the weight size. The size bounds include \(1\le b\le E,\ b\le K\)
and \(|M|\le K\). Write \(L_-\) for the cover lattice at the negative
shift and \(L_d\) for the integral lattice at the target weight
intersected with the prescribed type. Denote their squared-metric
slopes by \(\mu_-\) and \(\mu_d\); this normalization gives degree
equal to rank times slope, using squared norms at complex places.
Let \(\mu_0\) be the core slope at the original point.

The comparison for the flat line together with \eqref{eq:cov-u}
gives \(\mu_-\ge-M-(C_F+C_U)E\), with \(C_F\) the constant in the
flat-line comparison. The maps from \(L_-\) to \(L_d\) are integral,
have zero common kernel, and satisfy a sum-of-squares norm bound
with multiplier \(e^{2b}\). If there are \(j\) maps, the positive
shift bounds \(j\,\mathrm{rk}(L_d)\) above, and its rank is at most
\((1+b/K)\mathrm{rk}(L_-)\). The negative shift has positive rank.
The target is spanned over the field by integral sections of norm
at most \(e^K\) at each embedding. These are the hypotheses of the
integral lattice descent estimate in the source. With \(C=C_F+C_U\)
and \(t=-M-CE-2b\), that estimate gives
\[
                   \mu_d\ge t-(b/K)(|t|+2K).
\]
The comparison at the target ring gives \(\mu_0\ge\mu_d-2b\).
To see the exact loss, set \(\xi=b/K\). Then
\(\xi\le1,\ \xi K=b\), and
\[
 \xi |t|\le \xi|M|+\xi(CE+2b)\le b+CE+2b .
\]
Consequently \(\mu_0\ge-M-2CE-9b\ge-M-(2C+9)E\).
At the terminal lift there is a separate defect bound with constant
\(C_D\). Taking \(C_W=\max(C_D,2(C_F+C_U)+9)\) yields at every index
\begin{equation}                 \mu_0\ge-M-C_WE .\label{eq:cov-low}\end{equation}
\footnote{\href{https://github.com/Deicyde/riemann-hypothesis/blob/40e56bab/src/Development/AbelianVarietyTheory/Proofs/DivisorProofs/LowerBoundOfAllEstimates.lean}{Cover lattices and the numerical loss};
\href{https://github.com/Deicyde/riemann-hypothesis/blob/40e56bab/src/Development/AutomorphicArithmetic/ThetaCorrespondence/ThetaPairLifts/ThirdPlaceSeries/LowOfCoverHypotheses.lean}{last-index alternative}.}

\subsection{Removing the discarded subspace}
Let \(F\) be totally complex, and let \(L\) be a nonzero finitely
generated lattice in a hermitian \(F\)-space. Suppose \(e^2=e\), and at
each complex embedding squared norms add for \(ex\) and \((1-e)y\).
In the Gram realization write the inner product as conjugate-linear in
its first entry. The norm identity first gives
\(\Re\langle ex,(1-e)y\rangle=0\). At the chosen embedding pick
\(a\in F\) with nonreal image, which exists because the place is complex.
Applying the same identity with \(ax\) in place of \(x\) gives
\(\Re(\overline{a}\langle ex,(1-e)y\rangle)=0\). These two equations
force the inner product to vanish. Thus the image and kernel are
orthogonal and \(e\) is self-adjoint.

Let \(L_g=L\cap\mathrm{im}(e)\) and \(L_b=L\cap\ker(e)\), of ranks
\(d_g,d_b\). Suppose \(g\ge1\) is an integer such that \(geL\subset L\).
The map \(g(1-e)\) carries \(L\) into \(L_b\), its image \(J\)
contains \(gL_b\), and its kernel on \(L\) is \(L_g\). Hence
\(\operatorname{rk}L=d_g+d_b\) and \(J\) has rank \(d_b\).
The hermitian quotient by the saturated kernel identifies isometrically
at each embedding with the orthogonal projection onto \(\ker e\),
before multiplication by \(g\). Equivalently the operator from the
quotient to the image has norm at most \(g\). Taking determinants and
using \(J\subset L_b\) of full rank gives the cut inequality
\begin{equation}
 (d_g+d_b)\,\mu(L)\le d_g\,\mu(L_g)+d_b\{\mu(L_b)+2\log g\}.
 \label{eq:cov-cut}
\end{equation}
Indeed the squared archimedean determinant incurs at most \(g^{2d_b}\)
per place; finite index in \(L_b\) can only decrease the degree of
the image. Additivity of determinant degree in the saturated sequence
for \(L_g\subset L\) gives \eqref{eq:cov-cut}. Products of rank and
slope here mean degree, also on a rank-zero summand.

Fix \(0\le\delta\le1/12\) and suppose \(d_b\le3\delta(d_g+d_b)\).
Let \(Q\) bound the maximal slope of an ambient lattice containing \(L\),
as well as \(2\log g\) and \(-\mu(L)\). Then \(Q\ge0\),
\(d_g>0\), and \(d_b\le4\delta d_g\) since \(1-3\delta\ge3/4\).
In \eqref{eq:cov-cut} the bad degree is at most \(d_bQ\) by the
definition of maximal slope. Moving its terms to the left gives
\[
 d_g(\mu(L)-\mu(L_g))\le d_b(Q+2\log g-\mu(L))
                 \le 3d_b Q\le12\delta d_gQ .
\]
This proves \(\mu(L_g)\ge\mu(L)-12\delta Q\), including the case \(d_b=0\).
\footnote{\href{https://github.com/Deicyde/riemann-hypothesis/blob/40e56bab/src/Formalism/MeanHeightDescent/AverageDegreeBounds/GoodBadCut.lean}{Integral cut and height of the image};
\href{https://github.com/Deicyde/riemann-hypothesis/blob/40e56bab/src/Development/AbelianVarietyTheory/Proofs/DivisorProofs/MajAtAddRight.lean}{orthogonality and the factor twelve}.}

\subsection{The uniform conclusion}
In the application \(L\) is the lattice of the prescribed type and
\(L_g\) its retained part. Put \(\delta=\exp(-x^{\varepsilon_0/7})\),
where \(\varepsilon_0>0\) is fixed on the family. The truncation
provides the idempotent above, \(g\) the order of its finite group,
and \(2\log g\le b\). The discarded rank is at most \(3\delta\,\mathrm{rk}L\).
Suppose the maximal slope is bounded by \(x^A b\), \(|M|\le x^B b\),
and \(\mu(L)\ge\mu_0-C_KE\) for \(C_K\ge0\). By \eqref{eq:cov-low},
\(-\mu(L)\le M+(C_K+C_W)E\). Then
\[
 Q=(x^A+x^B+C_K+C_W+1)E
\]
dominates each of the three bounds required above. The condition
\(12\delta(x^A+x^B+C_K+C_W+1)\le1\) also ensures \(\delta\le1/12\).
It holds eventually because \(x^j\exp(-x^{\varepsilon_0/7})\to0\)
for each real \(j\). Therefore \(\mu(L_g)\ge\mu(L)-E\).
The preceding recursion gives the normalized log covolume at \(i\)
at most \(-\mu(L_g)\). Combining the three inequalities yields
\[
  \frac{\log\operatorname{covol}_i}{D_{\rm norm}}
      \le M+(C_K+C_W+1)E .
\]
Here the left side denotes precisely the lattice log covolume and
normalization in the recursion at that index. All inequalities hold
simultaneously above the maximum of the loss threshold and the
thresholds in the eight uniform inputs. For each larger constant the
conclusion persists because \(E\ge1\).
\footnote{\href{https://github.com/Deicyde/riemann-hypothesis/blob/40e56bab/src/Development/AbelianVarietyTheory/Proofs/DivisorProofs/CovolumeFromGramAndBounds.lean}{The pointwise bound};
\href{https://github.com/Deicyde/riemann-hypothesis/blob/40e56bab/src/Development/AutomorphicArithmetic/ThetaCorrespondence/ThetaPairLifts/ThirdPlaceSeries/CovolXOfSteps.lean}{the uniform covolume conclusion}.}

The cover degree, its rank share, its component height and
Hilbert--Samuel bounds, and the recursive covolume comparison were
inputs in this derivation. Their formal constructions enter through
\href{https://github.com/Deicyde/riemann-hypothesis/blob/40e56bab/src/Development/AutomorphicArithmetic/ThetaCorrespondence/ThetaPairLifts/ThirdPlaceSeries/TopCovolG.lean}{the geometric assembly theorem}.

\section{Descent through a family of integral maps}
\label{sec:map-descent}

Here is the descent inequality used in the comparison of \(L_-\)
with \(L_d\). Write \(F\) for a totally complex number field and
\(v\) for one archimedean place from each conjugate pair. Let \(m\)
be the number of such places. Give the finite-dimensional \(F\)-spaces
\(V,W\) positive definite hermitian norms at each place. We use
the degree \(\widehat d\) in the squared-norm normalization of the
preceding section. In particular the degree of a free line generated
by \(x\) is \(-m^{-1}\sum_v\log\|x\|_v^2\). For an integral lattice
\(C\) in its \(F\)-span put
\(\mu(C)=\widehat d(C)/\operatorname{rk}_F(C)\) when its rank is positive.

Let \(A\subset V,\ B\subset W\) be finitely generated
\(\mathcal O_F\)-submodules, \(A\ne0\). Suppose \(n\) linear maps
\(\phi_j:V\to W\) send \(A\) into \(B\), have zero common kernel and
obey, for all \(v\) and \(x\in V\),
\begin{equation}
 \sum_{j=1}^n\|\phi_jx\|_v^2\ \le\ e^{2P}\|x\|_v^2 .
 \label{eq:descent-map-bound}
\end{equation}
Suppose \(B\) is spanned over \(F\) by elements of \(B\) of norm at
most \(e^\Lambda\) at every place. Put \(a=\operatorname{rk}_F A\),
\(N=n\operatorname{rk}_F B\). If \(\Lambda,\epsilon\ge0\),
\(N\le(1+\epsilon)a\) and \(t\le\mu(A)-2P\), then
\begin{equation}
                  \mu(B)\ \ge\ t-\epsilon(|t|+2\Lambda).
 \label{eq:descent-result}
\end{equation}
Neither a positive \(t\) nor a free lattice \(A\) is required.

\paragraph{The lower bound on quotient slopes.}
We give the determinant argument underlying the estimate for the
minimal slope. Let \(C\) be spanned over \(F\) by integral vectors of
norm at most \(e^\Lambda\) at all places. For a saturated sublattice
\(D\subset C\), put \(Q=C/D\) with its quotient norms. The images of
these vectors span \(Q\) over \(F\), and the quotient norms do not
increase. If \(r=\operatorname{rk}_F Q>0\), select \(r\) images
\(q_1,\ldots,q_r\) forming an \(F\)-basis. Their integral span
\(Q_0\) is free and of finite index in \(Q\). At every \(v\) Hadamard's
inequality gives
\[
  \det(\langle q_i,q_j\rangle_v)_{i,j}
        \ \le\ \prod_i\|q_i\|_v^2 \ \le\ e^{2r\Lambda}.
\]
The determinant formula for the degree gives
\(\widehat d(Q_0)\ge-2r\Lambda\). Enlarging a full-rank lattice
by finite index increases its degree. Indeed in a common real
realization the volume of the smaller lattice is the index times
the volume of the larger one, and degree uses the negative logarithm
of volume with positive normalization. Therefore
\(\widehat d(Q)\ge-2r\Lambda\). This also holds at rank zero, with
degree zero. Equivalently all nonzero quotient slopes are at least
\(-2\Lambda\). In the formal development the same conclusion is
obtained by putting the short lines into a saturated step of the
Harder--Narasimhan filtration. The full-span property gives full rank
of that step. A saturated sublattice of full rank equals the ambient
lattice.\footnote{\href{https://github.com/Deicyde/riemann-hypothesis/blob/40e56bab/src/Formalism/MeanHeightDescent/AverageDegreeBounds/SpanMuMin.lean}{Short generators, line heights and the saturated filtration step}.}

\paragraph{The image lattice.}
Give \(B^n\) the orthogonal-sum metric and put
\(\Phi x=(\phi_1x,\ldots,\phi_nx)\).
The common-kernel hypothesis gives injectivity. Choosing a nonzero
element of \(A\) shows that \(n>0\) and \(B\ne0\). Write \(I=\Phi A\),
and \(D=B^n\cap (F I)\) for its saturation in \(B^n\).
Then \(a=\operatorname{rk}_F I=\operatorname{rk}_F D\le N\), and
\(\widehat d(D)\ge\widehat d(I)\). By
\eqref{eq:descent-map-bound} the operator norm of \(\Phi\) at each
place is at most \(e^P\). The bound extends to the completion at that
place by density of \(F\) and continuity of the linear map.
On the top exterior power of \(F A\) the norm multiplier is at most
\(e^{aP}\), since each singular value on this subspace is at most
\(e^P\). The induced map on the determinant modules of \(A\) and
\(I\) is an integral isomorphism, so no index arises there. The
change of squared determinant norms contributes at most \(2aP\)
on averaging their logarithms. Hence
\[
             \widehat d(D)\ge\widehat d(I)
             \ge\widehat d(A)-2aP \ge at .
\]
The short generators of \(B\), placed one at a time in each coordinate,
give integral short generators of \(B^n\) over \(F\). The preceding
quotient bound yields \(\widehat d(B^n/D)\ge-2(N-a)\Lambda\).
In a saturated exact sequence, determinant degrees add when the
quotient has its quotient norm. This follows by wedging a basis of
the subspace with orthogonal lifts of a basis of the quotient at each
place and using the integral determinant isomorphism. In particular
it applies even when \(D\) is only projective over \(\mathcal O_F\).
Since the degree of \(B^n\) is \(n\widehat d(B)\), we obtain
\[
                 N\mu(B)\ \ge\ at-2(N-a)\Lambda .
\]
Subtract the right side from \(Nt\). The difference is at most
\((N-a)(|t|+2\Lambda)\). The rank hypothesis gives
\(N-a\le\epsilon a\le\epsilon N\), and \(N>0\).
Division by \(N\) proves \eqref{eq:descent-result}.
\footnote{\href{https://github.com/Deicyde/riemann-hypothesis/blob/40e56bab/src/Formalism/MeanHeightDescent/AverageDegreeBounds/DescentSlope.lean}{Slope inequality for an injective map};
\href{https://github.com/Deicyde/riemann-hypothesis/blob/40e56bab/src/Development/AbelianVarietyTheory/Proofs/DivisorProofs/HtMapLeOfPoints.lean}{direct-sum norm, comparison of ranks and final inequality}.}
In the preceding section the source takes \(A=L_-\), \(B=L_d\),
\(P=b,\ \Lambda=K,\ \epsilon=b/K\). Thus the inequality holds for
the actual ranks of the lattices on each window once the listed maps,
short generators and rank estimates have been supplied there.

\section{A bound for the height error at each member of a window}
\label{sec:height-gap}

The finite Fourier calculation uses an upper bound for a difference
\(\mathrm{ASK}-\mathrm{Main}\). It is important that this bound holds at
each member of a finite family. An upper bound on its mean over the
family would leave individual prime sums uncontrolled. We translate
the pointwise slope calculation which produces the error before its
terms are estimated by powers of the scale.

Fix a member of the family and write \(W\ne\varnothing\) for its
finite set of words and \(d_u>0\) for the weight of \(u\in W\).
Write \(\langle f\rangle=\sum d_u f_u/\sum d_u\). For each word let
\(\nu_u,\nu_u^+,\mu_u,\mu_u^+\) denote respectively the slope and
maximal slope of its target and source lattices. Write \(h_u\) for its
height, \(T_u>0\) for its count, and \(E_u\) for its exponent sum.
Let \(G,F,\mu_C,\mu_C^{\min}\) be the quantities denoted
\(\mathrm{Ga},\mathrm{Fo},\mu_{\mathrm{CC}},\mu_{\min,\mathrm{CC}}\)
at this member. Set
\[
 s=\langle\nu_u-h_u\rangle,\qquad
 z=\max_{u\in W}(\nu_u^+-h_u)-s .
\]
Since \(\nu_u\le\nu_u^+\), we have \(z\ge0\).

Let there be \(n_T\) top indices \(t\) with weights \(a_t\ge0\).
Put \(A=\sum_t a_t\). Write \(m\) for the chain length, \(f>0\) for
the number of geometric components, \(v>0\) for the volume parameter,
\(w_{j\rho}>0\) for the local weights, and \(g\) for the exponent
parameter \(\mathrm{lg}\). The notation \(L_a,L,a,b,b',k\) and
\(B_0,N_c,X''\) is local to this section, with \(k,b>0\). The
constants denoted by \(c\) with subscripts below are those of the
window in the source. We keep separate the upper error in the count,
the cost of division, and the loss in passing to the quotient. Put
\begin{align*}
 C_T&=c_{\rm count}(L_a+m^2L_a/k),\\
 \lambda&=\sum_{j=1}^m f^{-1}\sum_{\rho=1}^f
            \log(w_{j\rho}v/\pi)+mg-C_T,\\
 E_D&=c_{\rm div}n_TaL_a+
          c_{\rm div}m(B_0+N_c)x^{e_{\rm div}\varepsilon},\\
 S_-&=c_{\rm slope}n_TaL_a(1+m/k),\qquad
 S_+=c_{\rm slope}n_TaL_a,\\
 P&=c_{\rm loss}n_Tmbx^{e_{\rm loss}\varepsilon},\\
 L_u&=c_1\mathrm{IO}+c_1 n_Tm(a/k+b')\log T_u+c_1 X''+1 .
\end{align*}
Here \(\mathrm{IO}\) is the local loss and \(c_1,n_T,m,a,b'\ge0\).
The two quotient-slope comparisons read
\begin{align*}
 \nu^+_u&\le\mu_u^++G+F-E_u+c_1mB_0,\\
 \nu_u&\ge\mu_u+G+F-E_u-L_u .
\end{align*}
The slope bounds give \(-S_-\le\mu_u\) and \(\mu^+_u\le S_+\).
The exponent, height and count comparisons used here give
\begin{align*}
 |E_u-mAg|&\le c_E m B_0,\\
 |(h_u+A\log T_u)-(h_{u'}+A\log T_{u'})|
       &\le c_H(n_T a L_a+mB_0),\\
 \lambda\le\log T_u&\le c_{\rm count}mL,\qquad
 |\log T_u-\log T_{u'}|\le c_{\rm count}\log k.
\end{align*}
Assume also \(\mathrm{ASK}-G\le P,\ \mathrm{IO}\le P\). These
pointwise hypotheses are quantitative inputs from earlier parts of
the construction.

\paragraph{Cancellation of the common exponent.}
Put \(\omega=\gamma+\log(\pi/v)-\mathcal L\), where
\(\mathcal L\) denotes the source's logarithmic-derivative parameter.
Let
\[
 D_0=m\sum_t(a_t\log a_t-a_t-a_t\omega).
\]
Terms with \(a_t=0\) are interpreted by continuity. The division
inequality is \(\mu_C+h_u-F\le D_0-A\log T_u+E_D\).
Insert it in the lower quotient comparison. For every \(u\) this
gives
\[
 \nu_u-h_u\ge G+\mu_C
          -S_- -m(Ag+c_EB_0)-D_0+A\lambda-E_D-L_u .
\]
Expanding \(\lambda\) and \(\omega\) yields the identity
\begin{align*}
 D_0-A\lambda+mAg
  &=\mathrm{Main}+AC_T,\\
 \mathrm{Main}
  &=\sum_{j=1}^m\sum_t a_t\!
  \left(\log a_t-1-\gamma+\mathcal L
               - f^{-1}\sum_\rho\log w_{j\rho}\right).
\end{align*}
Indeed the \(mg\) terms cancel first, and
\(\log(\pi/v)+\log(w_{j\rho}v/\pi)=\log w_{j\rho}\).
Taking the weighted mean proves
\begin{equation}
  G+\mu_C-\mathrm{Main}-K_0-\langle L\rangle \le s,\qquad
  K_0=AC_T+E_D+mc_EB_0+S_- .
  \label{eq:height-gap-lower}
\end{equation}

\paragraph{Variation from word to word.}
For two words \(u,u'\), subtracting the two exponent bounds and the
height comparisons gives
\[
  (-E_u-h_u)-(-E_{u'}-h_{u'})\le V,\quad
  V=2mc_EB_0+c_H(n_TaL_a+mB_0)+Ac_{\rm count}\log k.
\]
Subtract the lower quotient comparison at \(u'\) from the upper
one at \(u\). The terms \(G,F\) cancel. After subtracting the heights
the result is at most \(K_3+L_{u'}\), with
\(K_3=S_-+S_++V+c_1mB_0\). Taking the mean over \(u'\) and the
maximum over \(u\) yields \(z\le K_3+\langle L\rangle\). By the
upper bound on \(\log T_u\) and the bound on \(\mathrm{IO}\),
\[
 \langle L\rangle\le M_L
    :=c_1P+c_1 n_Tm(a/k+b')c_{\rm count}mL+c_1X''+1 .
\]

\paragraph{The logarithmic term.}
The slope inequality on the common lattice gives
\[
 s-\mu_C\le\log |W|+
 \frac{4z\log(1+H/z)+(\eta+4\sqrt e)H}{1-\delta}
\]
when \(s-\mu_C^{\min}\le H,\ H>0,\ 0\le\delta<1,\ 0\le e\le1/4\)
and \(\eta\ge0\). At \(z=0\) use the continuous value zero for the
product. Suppose \(H\le Xb\), \(X=x^{c_{HH}}>0\).
For \(z>0\), the tangent-line inequality for the logarithm gives
\[
 z\log(1+H/z)
 \le z\log(1+X)+\frac{H-Xz}{1+X}
 \le (z+b)\log(1+X).
\]
The last inequality uses \(X/(1+X)\le\log(1+X)\), by integrating
\((1+t)^{-1}\) from \(0\) to \(X\). This holds at \(z=0\) as well.
If the window gives uniform bounds
\(\delta\le\bar\delta<1,\ \eta\le\bar\eta,\ e\le\bar e\le1/4\),
the nonnegativity of all numerator terms gives
\begin{align}
 \mathrm{ASK}-\mathrm{Main}
  &\le K_0+\log|W|+P+M_L \notag\\
  &\quad+
    \frac{4\log(1+X)(K_3+M_L+b)
                     +(\bar\eta+4\sqrt{\bar e})Xb}{1-\bar\delta}.
    \label{eq:unit-gap}
\end{align}
This follows by adding \eqref{eq:height-gap-lower},
the bound for \(s-\mu_C\), and \(\mathrm{ASK}-G\le P\) with the
indicated signs. The right-hand side is the error denoted
\(\mathrm{ErrX5at}\) for this member. In particular \(\log|W|=0\)
is permitted. No averaging over different windows, scales or
members occurs in this derivation.
\footnote{\href{https://github.com/Deicyde/riemann-hypothesis/blob/40e56bab/src/Formalism/RankGapBounds/MeanBoundSecondForm.lean}{Slope, division and height hypotheses};
\href{https://github.com/Deicyde/riemann-hypothesis/blob/40e56bab/src/Formalism/SizeWindowBrackets/QuotientSlopeBoundsWith.lean}{count and quotient inequalities and the error formula};
\href{https://github.com/Deicyde/riemann-hypothesis/blob/40e56bab/src/Formalism/SizeWindowBrackets/ChainSlopeWindows/LevelOfCongruence.lean}{the pointwise deduction}.}

\subsection{Absorbing the height errors on the prescribed windows}
\label{sec:absorb-height}
We now estimate the right side of \eqref{eq:unit-gap}. Write
\[
 y=x^\varepsilon,\quad y_D=x^{e_{\rm div}\varepsilon},\quad
 y_L=x^{e_{\rm loss}\varepsilon},\quad
 h_1=1+c_{HH},\quad g_1=1+e_k ,
\]
The exponents \(e_{\rm div}\), \(e_{\rm loss}\), \(e_k\) and \(c_{HH}\)
are nonnegative. Assume that each of the coefficients
\[
   c_{\rm count},\quad c_{\rm div},\quad c_E,\quad c_H,\quad
   c_{\rm slope},\quad c_1,\quad c_{\rm loss}
\]
lies in \([0,K]\). The window hypotheses on the sizes are
\begin{align*}
 &x,b,n_T,k\ge1,\quad 0<\varepsilon\le1,\quad xb\le a,\qquad
 m\le x^{1/2+\varepsilon/2},\quad xm\le k\le \min(a,x^{e_k}),\\
 &0\le L,L_a\le2y,\quad 0\le kb'\le a,\quad 0\le X''\le a,\quad
 0\le a_t\le2a,\quad B_0,N_c\ge0,\\
 &\log|W|\le(k-3)\log2,\quad\bar\delta\le\tfrac12,\quad
 (\bar\eta+4\sqrt{\bar e})x^{c_{HH}}b\le Ky\,n_Ta .
\end{align*}
The last inequality is a joint bound on the height and the exceptional
mass. The pointwise argument uses their product. Increase the threshold
on \(x\) until
\[
      \log x\le y,\qquad 200(1+K)^2 h_1g_1\le y .
\]
For fixed positive \(\varepsilon\) and fixed other constants, such a
threshold exists by \(\log x=o(x^\varepsilon)\).
We claim that
\begin{equation}
  \mathrm{ASK}-\mathrm{Main}
  \le x^{(e_{\rm div}+e_{\rm loss}+3)\varepsilon}
      \big\{n_Ta(1+m^2/k)+m(B_0+N_c)\big\}.
 \label{eq:height-power}
\end{equation}

Set \(E=n_Ta(1+m^2/k)+m(B_0+N_c)\), \(Y=y_Dy_Ly\),
\(R=YE\). Each of \(1,\ a,\ n_Ta,\ n_Ta m^2/k,\ mB_0,\ mN_c\)
is bounded by \(E\); the stronger bound \(m(B_0+N_c)\le E\)
also holds. Since \(m\le x\), \(mb\le a\), and \(m/k\le1\).
The bound on each \(a_t\) gives \(A\le2n_Ta\). Positivity of the
word count and \(k\ge1\) give
\(\log|W|\le (k-3)\log2\le k\le a\). Also
\[
 \log k\le g_1 y,\qquad
 \log(1+x^{c_{HH}})\le \log 2+c_{HH}\log x\le h_1 y .
\]
In the first inequality one can use \(e_k y\) in place of \(g_1 y\).
All of \(1,y,y_D,y_L\) are bounded by \(Y\).

The contributions to \(K_0\) in \eqref{eq:height-gap-lower}, in units of
\(R\), have the following bounds.
\[
 \begin{array}{c|ccccc}
 \text{term}&AC_T& c_{\rm div}n_Ta L_a&
 c_{\rm div}m(B_0+N_c)y_D&mc_EB_0&S_-\\ \hline
 \text{bound divided by }R&4K&2K&K&K&4K .
 \end{array}
\]
For the first entry use
\(A(1+m^2/k)\le2E\) and \(L_a\le2Y\).
The last uses \(1+m/k\le2\). Thus \(K_0\le12KR\).
Similarly \(P\le KR\) because \(n_Tmb\le n_Ta\).
For the middle term in \(M_L\), note that
\[
 n_Tm^2(a/k+b')\le 2n_Tam^2/k\le2E.
\]
Multiplying by \(c_1 c_{\rm count}L\le2K^2Y\) bounds this term
by \(4K^2R\). The other terms in \(M_L\) have bounds \(K^2R,KR,R\);
hence \(M_L\le(5K^2+K+1)R\).

In \(K_3\) the two source-slope errors together contribute at most
\(6KR\). The exponent variation contributes \(2KR\), the height
variation \(3KR\), and the term \(Ac_{\rm count}\log k\) contributes
\(2K g_1R\). The remaining \(c_1mB_0\) is at most \(KR\). Thus
\[
 K_3+M_L+b\le (5K^2+13K+2K g_1+2)R .
\]
The exceptional-mass term is at most \(KR\), and
\((1-\bar\delta)^{-1}\le2\). Insert all these bounds in
\eqref{eq:unit-gap}. After division by \(R>0\) the upper bound becomes
\[
  5K^2+16K+2+8h_1y(5K^2+13K+2K g_1+2).
\]
The first polynomial is at most \(10(1+K)^2\). The one in parentheses
is at most \(8g_1(1+K)^2\), since \(g_1\ge1\). Consequently the last
display is at most \(200(1+K)^2h_1g_1y\le y^2\).
Multiplying back by \(R=y_Dy_LyE\) proves \eqref{eq:height-power}.
In the Fourier comparison \(B_0\) and \(N_c\) have additional explicit
upper bounds. They may now be substituted on the right of
\eqref{eq:height-power} since their common coefficient is nonnegative.
All inequalities were evaluated at the same scale and family member.
\footnote{\href{https://github.com/Deicyde/riemann-hypothesis/blob/40e56bab/src/Formalism/SizeWindowBrackets/ErrorFormsWeakWindowsFalse.lean}{Hypotheses on the sizes and the joint exceptional-mass bound};
\href{https://github.com/Deicyde/riemann-hypothesis/blob/40e56bab/src/Formalism/SizeWindowBrackets/ChainSlopeWindows/SubThreeMulLogTwoLe.lean}{termwise absorption};
\href{https://github.com/Deicyde/riemann-hypothesis/blob/40e56bab/src/Development/AutomorphicArithmetic/DedekindLFunctions/LandauTheorems/ChTop.lean}{substitution at the arithmetic member}.}

\section{From the height error to a one-sided prime-sum estimate}
\label{sec:primecomparison}

Fix a scale \(x\geq3\) and a tolerance \(\varepsilon>0\).
The comparison in this section takes place separately at each scale.
The finite groups, weights, sets and auxiliary parameters may vary
with \(x\). We display the constants that have to be fixed before the
scale is chosen.

\subsection{The finite Fourier identity, including ramified primes}

Let \(H\) be a finite abelian group of order \(n\), written additively,
and put \(E=(\mathbb Z/2)^2\times H\),
\(\Phi=\{(0,v,h)\}\subset E\), \(f=|\Phi|=2n\).
Elements denoted \(t\in H\) are identified with \((0,0,t)\) in \(E\).
Let \(\kappa=(1,1,0)\), \(e=(1,0,0)\),
\(\xi(v_1,v_2,h)=(-1)^{v_1}\), and
\(\rho_t(g)=\kappa+g+t\). The map \(\rho_t\) carries the
complement of \(\Phi\) bijectively onto \(\Phi\).

The arithmetic input at a prime \(p\) is an inertia subgroup \(I_p\)
of \(E\) and a Frobenius representative \(F_p\in E\). Only its
coset modulo \(I_p\) enters the sums. For a finite order character
\(\psi\in\widehat E\), its value at a prime power in this construction is
\[
 v_\psi(p^i)=\frac1{|I_p|}\sum_{h\in I_p}\psi(iF_p+h)
 \qquad(i\geq1).
\]
Indeed restriction of \(\psi\) to inertia is either trivial, in
which case the sum is \(\psi(iF_p)\), or has zero mean, in which
case the local character contributes zero. The host construction
identifies these values with those of its associated primitive
Dirichlet character. Orthogonality now gives an exact formula,
\begin{equation}
 \frac1f\sum_{\psi\in\widehat E}\overline{\psi(g)}v_\psi(p^i)
 =\frac{2}{|I_p|}{\bf1}_{\,g-iF_p\in I_p}.
 \label{eq:pc-inertia}
\end{equation}
To see this, insert the average defining \(v_\psi\). The inner
character sum is \(|E|=2f\) when \(iF_p+h=g\) and is zero otherwise.
There is at most one such \(h\).

Here the host provides a finite set \(S\) containing \(2,q_1\) and
a prime \(d\notin S\). At \(p\in S\) it gives \(I_p\subset\Phi\)
and \(F_p\in\Phi\), at \(d\) it gives \(I_d=\{0,e\}\), and at all
other primes it gives \(I_p=0\). Let \(w_\rho>0\) for \(\rho\in\Phi\),
\(A>0\), and \(N\geq A/w_\rho\) for every \(\rho\). Write
\(r_{m,g}=(A/m-w_{\rho_t(g)})_+\), where \(z_+=\max(z,0)\).
The arithmetic prime contribution for this \(A,t,w\) equals
\begin{align*}
 P(A,t,w)={}&\sum_{\substack{p\leq N,\ p\notin S\cup\{d\}\\ F_p\notin\Phi}}
  2\log p\sum_{\substack{1\leq i\leq N\\iF_p\notin\Phi}}r_{p^i,iF_p}\\
 &+\log d\sum_{1\leq i\leq N}
       \sum_{\substack{g\notin\Phi\\g-iF_d\in I_d}}r_{d^i,g}.
\end{align*}
The sums indexed by \(p\) in this section range over primes. When
\(p^i>N\) the ramp vanishes, so the second line is zero if \(d>N\).
In view of \eqref{eq:pc-inertia}, the same real number is
\begin{equation}
 P(A,t,w)=\frac1f\sum_{g\notin\Phi}\sum_{\psi\in\widehat E}
  \overline{\psi(g)}
  \sum_{m=1}^N\Lambda(m)v_\psi(m)(A/m-w_{\rho_t(g)})_+ .
 \label{eq:pc-prime}
\end{equation}
At \(p\in S\) the entire coset is in \(\Phi\). At an unramified
prime the surviving term has coefficient two and \(g=iF_p\).
At \(d\) the surviving coset terms have coefficient one. These are
all the terms, since \(\Lambda(m)\) vanishes away from prime powers.

Let \(G_\psi(y)\) be the ramp with coefficients \(\Lambda(m)v_\psi(m)\),
centered by \(y\ell_\psi\) for \(\psi\ne1\), where \(\ell_\psi\)
is the logarithmic derivative of the associated primitive \(L\)-function
at one. For \(\psi=1\) the centering is
\(-y\log y+(1+\gamma)y\). These formulas also define the functions
for \(0<y<1\) by taking the ramp sum to be empty. Put
\[
 M(A,w)=A\left(\log A-1-\gamma+\ell_\xi
                         -\frac1f\sum_{\rho\in\Phi}\log w_\rho\right).
\]
Then
\begin{equation}
 P(A,t,w)-M(A,w)
   =\frac1f\sum_{g\notin\Phi}\sum_{\psi\in\widehat E}
     \overline{\psi(g)}\,w_{\rho_t(g)}G_\psi(A/w_{\rho_t(g)}).
 \label{eq:pc-balance}
\end{equation}
For \(\psi\ne1\), multiplying the centered ramp by \(w\)
adds \(A\ell_\psi\) to the inner prime sum of \eqref{eq:pc-prime}.
For \(\psi=1\) it adds \(-A\log A+A\log w+(1+\gamma)A\).
The sum of \(\overline{\psi(g)}\) over \(g\notin\Phi\) is \(f\)
at \(1\), \(-f\) at \(\xi\), and zero otherwise, by character
orthogonality on the index-two subgroup. The bijection \(\rho_t\)
turns the logarithmic sum into the sum over \(\Phi\) in \(M\).
Moving the two corrections to the left gives \eqref{eq:pc-balance}.

\subsection{The weights and the surviving characters}

Let \(H_0\subset H\) be a subgroup, \(T\subset H\), \(T_2\subset H_0\),
\(n_T=|T|\), \(n_2=|T_2|\) and \(\Delta=n_2/n\).
For \(u\in H_0\), the thick positions are precisely
\(\rho=(0,1,s+u)\), \(s\in T_2\).
For each \(1\leq j\leq m\) take positive weights \(w_{u,j,\rho}\).
Take \(b,b',X'>0,\ \eta,\vartheta\geq0\) such that
\[
 b'X'\leq xb,\qquad b\leq w_{u,j,\rho}\leq b+\eta+1
 \ \ (\rho\text{ thick}),\qquad
 w_{u,j,\rho}\leq2b'\ \ (\rho\text{ otherwise}).
\]
At \(t\in T\) put \(A=a_t\), \(y_t=a_t/b\), and require
\(|y_t-x|\leq1\) and \(a_t/w_{u,j,\rho}\leq N\).
The arithmetic estimates on the thin scales and the window costs are
\begin{equation}
 |G_\psi(y)|\leq y\vartheta\quad(y\geq X'/3),\qquad
 2f(x+1)\vartheta\leq x^{-1},\qquad
 6f(x+1)(\eta+1)\leq b/x .
 \label{eq:pc-cost}
\end{equation}

The thick part after replacing each thick weight by \(b\) is
\[
 T_u(t)=\frac1f\sum_{s\in T_2}\sum_{\psi\in\widehat E}
          \psi\bigl(-(e+s+u-t)\bigr)bG_\psi(y_t).
\]
Indeed \(\rho_t(g)=(0,1,s+u)\) if and only if \(g=e+s+u-t\).
We claim
\begin{equation}
 \Re T_u(t)\leq P(a_t,t,w_{u,j})-\Re M(a_t,w_{u,j})+2b/x .
 \label{eq:pc-individual}
\end{equation}
For the thick terms the inequality used is
\[
 |wG_\psi(A/w)-bG_\psi(A/b)|
                  \leq3A\,|w-b|/\min(w,b).
\]
For a nontrivial \(\psi\) the two linear centering terms cancel.
The ramps are \(1\)-Lipschitz in \(w\) and vanish for
\(m>A/\min(w,b)\). The classical Chebyshev bound
\(\sum_{m\leq z}\Lambda(m)\leq3z\) proves the claim.
For the trivial \(\psi\), suppose first \(b\leq w\). The decrease
of the uncentered ramp sum lies between \(0\) and \(3A(w-b)/b\);
the change from the centering has the opposite sign and size
\(A\log(w/b)\leq A(w-b)/b\). The same bound therefore holds.
Interchanging \(w,b\) covers the other order. Thus the replacement
cost at a thick term is at most \(3(x+1)(\eta+1)\).
At a thin term \(w>0\), \(w\leq2b'\) and \(A\geq(x-1)b\).
Since \(x\geq3\), these give \(A/w\geq X'/3\).
Its absolute contribution is at most \(A\vartheta\leq(x+1)b\vartheta\).
In \eqref{eq:pc-balance} there are \(f\) choices for \(g\), \(2f\)
characters and the factor \(1/f\). Hence its change in absolute value
is bounded by \(2f[3(x+1)(\eta+1)+(x+1)b\vartheta]\leq2b/x\)
by \eqref{eq:pc-cost}. This proves \eqref{eq:pc-individual}.

Let \(\mathcal A\subset\widehat E\) consist of the characters trivial
on \(H_0\). For each \(s\in T_2\), orthogonality on \(H_0\) gives
\[
 \sum_{u\in H_0}\psi(-(e+s+u-t))=
 \begin{cases}|H_0|\psi(e+t),&\psi\in\mathcal A,\\0,&\psi\notin\mathcal A.\end{cases}
\]
Here \(e=-e\) and \(\psi(s)=1\) in the first case. It follows that
\[
 \frac1{|H_0|}\sum_{u\in H_0}T_u(t)
      =b\,\frac{\Delta}{2}\sum_{\psi\in\mathcal A}\psi(e+t)G_\psi(y_t).
\]
Write \(E_{\rm top}\) for a common upper bound, for every \(u\in H_0\),
on \(\sum_{t\in T}\sum_{j=1}^m [P(a_t,t,w_{u,j})-\Re M(a_t,w_{u,j})]\).
Average \eqref{eq:pc-individual} over \(u\) and sum over \(t,j\).
The result is
\begin{equation}
 mb\frac{\Delta}{2}\sum_{t\in T}\sum_{\psi\in\mathcal A}
       \Re[\psi(e+t)G_\psi(y_t)]
              \leq E_{\rm top}+2n_Tmb/x .
 \label{eq:pc-top}
\end{equation}

\subsection{Comparison with a fixed family of primitive characters}

Let \(K\subset\mathcal A\) contain \(1\). Suppose that \(\psi(e)=1\)
on \(K\). Set \(c_\psi=1\) if \(\psi\) is trivial on \(H\), and
\(c_\psi=c\) otherwise. Fix \(|c|\leq1\) and constants \(M_0,K_0,C_0\)
valid throughout the windows, with
\begin{gather*}
 |\mathcal A|\leq M_0,\qquad
 \left|\sum_{t\in T}\psi(t)-c_\psi n_T\right|
                  \leq M_0x^{-1+\varepsilon}n_T\quad(\psi\in K),\\
 |\ell_\psi|\leq K_0\quad(\psi\in K,\ \psi\ne1),\qquad
 |\ell_\psi|\leq C_0\log x\quad(\psi\in\mathcal A\setminus K).
\end{gather*}
For \(\psi\) trivial on \(H\) the first sum equals \(n_T\) exactly.
Define
\[
 G_K=\sum_{\psi\in K}c_\psi\Re G_\psi(x),\qquad
 \Sigma=\sum_{\psi\in\mathcal A\setminus K}|G_\psi(x)|,\qquad
 V=G_K-\Sigma .
\]
For \(x\) sufficiently large depending only on the fixed constants
and \(\varepsilon\), we have
\begin{equation}
 n_T(V-4x^{1/2+\varepsilon})\leq
  \sum_{t\in T}\sum_{\psi\in\mathcal A}\Re[\psi(e+t)G_\psi(y_t)] .
 \label{eq:pc-lower}
\end{equation}

For the proof, the ramp is locally Lipschitz. Off the integers the
derivative of the principal centered ramp is
\(\sum_{m\leq y}\Lambda(m)/m-\log y+\gamma\).
The elementary Mertens estimate
\(\sum_{m\leq y}\Lambda(m)/m=\log y+O(1)\) therefore gives an absolute
bound \(D\) for this derivative on \(y\geq1\). Continuity at integers
extends the Lipschitz estimate across them. For a nontrivial \(\psi\),
the absolute derivative is at most \(\log y+D+|\ell_\psi|\) by the
same estimate, increasing \(D\geq0\) once if necessary. Since
\(G_\psi(1)=\ell_\psi\) for \(\psi\ne1\), we obtain
\begin{align*}
 |G_\psi(y_t)-G_\psi(x)|&\leq\log(x+1)+D+|\ell_\psi|,\\
 |G_\psi(x)|&\leq x(\log x+D+|\ell_\psi|).
\end{align*}
For \(1\ne\psi\in K\), subtract \(n_Tc_\psi G_\psi(x)\) from its
sum in \eqref{eq:pc-lower} before taking real parts. The difference is
\[
  \sum_{t\in T}\psi(t)[G_\psi(y_t)-G_\psi(x)]
          +\left(\sum_{t\in T}\psi(t)-c_\psi n_T\right)G_\psi(x).
\]
In absolute value this is at most \(n_T\) times
\(\log(x+1)+D+|K_0|+|M_0|x^\varepsilon(\log x+D+|K_0|)\).
For \(\psi=1\) the lower error per term is at most \(D\).
For \(\psi\in\mathcal A\setminus K\) the real part of each term is
bounded below by \(-|G_\psi(x)|-\log(x+1)-D-|C_0|\log x\).
Thus all three errors per term are at most \(A_1x^{1/4+\varepsilon}\), where
\[
 A_1=5+D+|K_0|+|M_0|(4+D+|K_0|)+4|C_0|.
\]
Indeed \(\log x\leq4x^{1/4}\), \(\log(x+1)\leq5x^{1/4}\) for \(x\geq3\).
Enlarge the threshold so that \(|M_0|A_1\leq x^{1/4}\). The total
per \(t\) is then at most \(x^{1/2+\varepsilon}\), which gives the
stated bound with the source's constant four.

\subsection{Accounting for the size error}

The source supplies a geometric upper bound for \(E_{\rm top}\) at
this same \(x\). In numerical form, with \(C\geq0\), \(a_E\geq0\), it is
\begin{align}
 E_{\rm top}\leq x^{C\varepsilon}\bigg[
 n_T a_E(1+m^2/k)
 +m\left(2nmbx^\varepsilon+
       \frac{n_2 k b x^\varepsilon}{n(n x^{-\varepsilon})}\right)\bigg].
 \label{eq:pc-geometry}
\end{align}
The size hypotheses that enter division by the main coefficient are
\[
 \begin{gathered}
 0\leq a_E\leq2xb,\quad n_T\geq nx^{-\varepsilon/8}/2,
 \quad x^{-\varepsilon}\leq\Delta\leq1,\\
 x^{1/2-\varepsilon}\leq m\leq x^{1/2+\varepsilon/2},\qquad
 xm\leq k\leq n^2x^{1/2-\varepsilon},\qquad
 26\leq x^{3\varepsilon/8}.
 \end{gathered}
\]
Let \(Q=n_Tmb>0,\ Y=x^{1/2+13\varepsilon/8},\ P_0=x^{C\varepsilon}\).
Divide the four summands inside the bracket by \(Q\).
The first is \(a_E/(mb)\leq2x/m\leq2x^{1/2+\varepsilon}\leq2Y\).
The second is \(a_E m/(bk)\leq2xm/k\leq2\leq2Y\).
The third is \(2nm x^\varepsilon/n_T
      \leq4x^{1/2+13\varepsilon/8}=4Y\).
The fourth is \(n_2 k x^{2\varepsilon}/(n_Tn^2)
      \leq2x^{1/2+9\varepsilon/8}\leq2Y\).
Hence \(E_{\rm top}\leq10P_0YQ\). Combining
\eqref{eq:pc-lower} and \eqref{eq:pc-top}, dividing by \(Q\)
and using \(2/x\leq1\leq P_0Y\), gives
\[
 \Delta(V-4x^{1/2+\varepsilon})\leq22P_0Y .
\]
As \(\Delta^{-1}\leq x^\varepsilon\) and \(C\geq0\), it follows that
\begin{equation}
 V\leq26x^{1/2+(C+21/8)\varepsilon}
       \leq x^{1/2+(C+3)\varepsilon}.
 \label{eq:pc-core}
\end{equation}

For the cyclic level of order \(o\) used to extract a prescribed
primitive character, the host parameters are
\[
 c=(3o+1)^{-1},\qquad M_0=4o,\qquad
 K_0=\sum_{\chi\ \mathrm{retained}} |L'(1,\chi)/L(1,\chi)|
\]
with the formal value for the principal component in \(K_0\) understood
as the library's totalized logarithmic derivative. Equivalently one
may increase any bound for the nonprincipal summands by a fixed
nonnegative constant, since it is used only on \(\psi\ne1\).
The character group \(\mathcal A\) is the dual of \(E/H_0\),
so \(|\mathcal A|=4[H:H_0]=4o\).
The arithmetic host identifies the retained characters bijectively
with the primitive family of the level and identifies their ramps
and weights, hence \(G_K=G_{\mathcal L}(x)\).
The host's bounded-character estimate gives
\(|G_\psi(x)|\leq x^{1/2+\varepsilon}\) on
\(\mathcal A\setminus K\), whence \(\Sigma\leq4o x^{1/2+\varepsilon}\).
The height error exponent at this scale is at most \(e_D+e_L+3\)
with \(e_D+e_L\geq0\). Thus \eqref{eq:pc-geometry} may be used with
\(C=e_D+e_L+3\). Adding \(\Sigma\) in \eqref{eq:pc-core} and requiring
\(1+4o\leq x^\varepsilon\) gives
\[
 G_{\mathcal L}(x)\leq x^{1/2+(e_D+e_L+7)\varepsilon}.
\]
Every large-scale threshold here depends on the fixed level, its
error constants and \(\varepsilon\). Consequently, if the arithmetic
inputs hold at every point of one interval above those thresholds,
the conclusion holds at every point of that same interval.

\paragraph{Correspondence and inputs.}
The Fourier calculation translates \texttt{ch\_\allowbreak{}balance}.
The projection, replacement of the weights, and character comparison
translate \texttt{Character\allowbreak{}Balance\allowbreak{}Moves.\allowbreak{}sum\_\allowbreak{}Th}, \texttt{th\_\allowbreak{}le},
\texttt{one\_\allowbreak{}char}, \texttt{core} and \texttt{ch\_\allowbreak{}core}.
The four-term calculation translates \texttt{proposition\_\allowbreak{}size\_\allowbreak{}bound}.
The conversion at a fixed level translates \texttt{ch\_\allowbreak{}top},
\texttt{gm\allowbreak{}A\_\allowbreak{}eq}, \texttt{sig\allowbreak{}A\_\allowbreak{}le} and \texttt{core\_\allowbreak{}gm}.
The upstream inputs used here are the host with the specified
Frobenius and inertia data, the common-window character and weight
estimates, the inequalities on their sizes, and the geometric
inequality \eqref{eq:pc-geometry}. In \texttt{ch\_\allowbreak{}top} this last
inequality follows from the source's unitwise height bound
\texttt{each\_\allowbreak{}unit\_\allowbreak{}bound\_\allowbreak{}statement} and error comparison
\texttt{err\allowbreak{}X5at\_\allowbreak{}le} after substituting the upper bounds for its
two coefficients \(B_0,N_c\).

\section{Common windows for the prime pairs and the Selmer estimate}
\label{sec:rh-common-windows}

Fix a primitive character $\psi$ of conductor $m$, $0<\epsilon<1/52$,
and a scale factor $c\ge1$. Put $z=U^{\epsilon/5}$,
$z'=z(\log U)^2$, $P=\exp z$, and $D=\exp z'$. For every fixed $A\ge1$
all sufficiently large $U$ admit primes $q,d$ with $P^{1/2}<q\le P$,
$D^{1/2}<d\le D$, and $(m,qd)=1$. They satisfy $q\equiv3\bmod 8$,
$d\equiv7\bmod16$, $q\mid d+1$, and $p\mid d+1$ for every odd prime
$p\le z$. For every character $\xi$ modulo $m$, every nonzero
$(i,j)\in\{0,1\}^2$, and every primitive character $\chi$ inducing
$\xi(\frac{\cdot}{q})^i(\frac{\cdot}{d})^j$ on units prime to $2mqd$,
this same choice gives
\[
 |G_\chi(y)|\le y^{1/2+\epsilon/2},\qquad
 |L'(1,\chi)/L(1,\chi)|\le4\log y\qquad (U\le y\le U^A).
\]
Here $G_\chi$ is the complex centered prime sum. For any primitive real
$\rho$ with conductor dividing $8m$ in place of $\xi$, the corresponding
inducing character has no zero for $\Re s>1-1/400$ and
$|\Im s|\le8mPD^j$. These are the conclusions of \texttt{PairGood}.

\subsection{Simultaneous choice of the primes}

The first candidates are the primes in $(\sqrt P,P]$ congruent to
$3\bmod8$. For each candidate $q$, the second candidates are the primes
in $(\sqrt D,D]$ congruent to $7\bmod16$ and to $-1$ modulo $q$ and
all odd primes at most $z$. These congruences form a single reduced
residue class modulo $16q\prod_{2<p\le z}p$ for all sufficiently
large $U$, with the product taken over primes. For every fixed
$\eta>0$ the estimates \texttt{stage1} and \texttt{stage2} exclude
at most an $\eta$ fraction in each selection. The second estimate is
uniform in the first candidate. Both candidate sets are nonempty.
Taking $\eta=1/2$ leaves a $q$ meeting every first-stage requirement,
then a $d$ meeting every requirement for that $q$. In the first stage
the family comprises the primitive representatives of all characters
modulo $8m$. In the second these characters are multiplied by
$(\frac{\cdot}{q})^i$, $i=0,1$, before taking primitive representatives;
the number of characters to test is at most twice the number modulo
$8m$. For large $U$, $d>q>8m$.

Here are the exceptional-set bounds in \texttt{StagePow}. Write $V$
for the upper endpoint of the candidate progression, $\mathcal P$ for
the candidate set, and $F$ for the maximum allowed base conductor.
Then $1\le F\le\sqrt V$, each candidate exceeds $F$, and the
prime-progression lower bound in the admissible ranges gives
$|\mathcal P|\ge V^{9/10}$. For $Q\ge3$ the zero-density input
\texttt{CountZAll.zero\_density\_box} bounds the number of primitive
characters of conductor at most $Q$ with a zero in
$\Re s>1-1/400,\ |\Im s|\le Q$ by $C Q^{1/4}$.
For a failing prime $p$, choose one bad primitive twist. Its conductor
equals $fp$ with $1\le f\le F$. Distinct primes yield distinct
characters: equality gives $fp=f'p'$, hence $p\mid f'p'$ and
$p=p'$ since $0<f'<p$. Taking $Q=FV$ bounds the number discarded by
$CV^{3/8}$, at most $(\eta/2)|\mathcal P|$ once $V$ is large.
The zero-density theorem imports mollifier zero detection, a block
large-sieve estimate and a coefficient-square bound.

The second exceptional set concerns a bad value at some real point,
which can depend on the prime. With a fixed truncation $T$ put
\[
  S_p(y)=\sum_{n\le T}\Lambda(n)\min(1,y/n)\kappa(n)
                             (\tfrac np),\qquad |\kappa(n)|\le1 .
\]
The mean-square input \texttt{meanSqK} gives
$\sum_{p\in\mathcal P}|S_p(y)|^2\le C y\log y\,|\mathcal P|$
uniformly in the permitted $y$ and in $\kappa$. Let $D_p(y)$ denote
the piecewise derivative in $y$. Then $yD_p(y)$ equals $S_p(y)$ with
$\kappa(n)$ replaced by $\kappa(n)\mathbf 1_{n>y}$ away from its
finitely many corners. Thus $\sum_p|D_p(y)|^2\le C(\log y)|\mathcal
P|/y$. For $L\le W\le2L$, integration and Cauchy--Schwarz give
\[
  \sup_{L\le y\le W}|S_p(y)|^2
      \le 2|S_p(L)|^2+2L\int_L^W|D_p(v)|^2\,dv.
\]
Sum over $p$ and apply the two mean-square bounds. The result is
$O(CL\log(2L)|\mathcal P|)$. Markov's inequality bounds the fraction
with $|S_p(y)|>y^{1/2+\tau}/2$ at some point in the block by
$O(C\log(2L)L^{-2\tau})$. For any fixed $\tau>0$ the explicit
relaxation used in \texttt{BadShare} at large $L$ is
$(32C/\tau)L^{-\tau}$. Summing over blocks starting at $2^hU$ bounds
the total fraction by $(32C/\tau)U^{-\tau}/(1-2^{-\tau})$, uniformly
in the upper endpoint. For at most $R$ multipliers, take the last
fraction at most $\eta/(2(R+1))$ and use the union bound. Including
the zero-density exceptions gives at most $\eta|\mathcal P|$ in total.
\texttt{BoxOfS} supplies the explicit-formula comparison to the twists.
Take $\tau=\epsilon/2$. The constant four for the logarithmic derivative
comes from \texttt{kappa4} at heights $H=8mPD^j$ satisfying
$U^{\epsilon/6}\le\log H\le3U^{\epsilon/5}(\log U)^2$.
The prime-progression, mean-square, zero-density and explicit-formula
inputs cited here have separate proofs in the source.

\subsection{Transferring both estimates to one interval}

Take $a(u)$ with $u^{2+8\epsilon}\le2^{a(u)}\le2u^{2+8\epsilon}$ for
$u\ge1$. Apply the strict window theorem to obtain $J\ge1,X_0$, and
the good-pair theorem with $A=J$ to obtain $U_0$. Set
$k=c^{1/(2+8\epsilon)}$. Given $y_0$, choose
$M\ge\max(1,y_0,U_0,k^3,k^{400})$ and $X\ge\max(X_0,M^J)$.
Set $U=X^{1/J}$. Choose one good pair at $U$ through $U^J=X$.
Apply the strict window theorem to this $X,q,d$. It supplies $W\ge U$
with $W^2\le X$ and the strict estimates throughout $[W,W^2]$.

Put $Z=kW$. Since $k^3\le W$, we have $Z^{3/2}\le W^2$. For every
$x\in[Z,Z^{3/2}]$ the point $u=x/k$ satisfies $W\le u\le x\le W^2$.
Set $e=a(u)$. Then $x^{1/J}\le U\le x$ and the same $q,d$ work for
every argument from $U$ to $x$. Furthermore
\[
 x^{2+8\epsilon}\le c2^e\le2x^{2+8\epsilon},
 \qquad x^{400}=k^{400}u^{400}\le u^{401}.
\]
For every field $E$ realizing the specified fixing subgroup for
$(\psi,e,q,d)$, apply the strict estimate at $u$, shorten the range
of twists, and increase the right side using $u^\epsilon\le x^\epsilon$.
This gives $\mathrm{VVle}(E,\lfloor x^{400}\rfloor,
\sqrt{\exp(x^\epsilon)})$ as in \texttt{WinPairs.winPairs}.
Here $Z\ge y_0$ and $q,d$ stay fixed on the whole indicated interval.

\subsection{The family used in the analytic ending}

The level $\Lambda$ consists of the primitive representatives of the
powers of $\psi$, including $\psi$ and the primitive trivial character.
The weight is one for the trivial member and the distinguished member
if present, and $1/(3\Lambda.o+1)$ for every other member. Thus the
weighted centered sum $G_\Lambda$ is fixed once $\psi$ is fixed. The
host requirement $G_\psi$ asks for the character comparisons
$\mathrm{EAtH}(\Lambda,H,T,x,\epsilon,4)$ together with the strict
bound $\mathrm{VVle}(E,t,\sqrt b)$ on that same host. In
\texttt{supplyOfAdmOf} the admissible-rank construction yields the
host. The prime rows and twist estimates give the first conjunct.
The fixing subgroup identifies its field with the field of the strict
estimate, and $t\le x^{400}$ gives the second conjunct. The resulting
host supply \texttt{supplyChi} has window exponent $\lambda=3/2$.

In \texttt{topDesWin} the closed estimates for the level, divisibility,
auxiliary character, irreducibility, small-level rank, selected rank,
and the three comparisons $G,S,\mathrm{GLOBAL}$ are conjoined. Their
conjunction gives a positive $\epsilon_1$ and, for each sufficiently
small fixed $\epsilon$, one lower-bound tuple $K$ for all nine.
Design existence gives a positive $\epsilon_3$ and quantifies over
every such $K$. Take $\epsilon<\min(\epsilon_1,\epsilon_3)$ and
select $K$ before selecting the design family. On the resulting
family its two size bounds, level and character comparisons and
strict Selmer estimate imply the divisibility relation. Its rank
conclusions then give the required multiplicity relation and the
global comparison supplies the chain. In \texttt{shapeWin} the
specific exponents passed into this argument are $e_D=2,e_L=0$.
The nine comparison inputs are imported here.

Consequently for some $\epsilon_0>0$ and every $0<\epsilon<\epsilon_0$
there is a family with a scale set $S$ containing arbitrarily late
intervals $[Z,Z^{3/2}]$. Beyond its threshold every host index
satisfies the same character comparison with constant four. On its
sufficiently large scales in $S$ the compatible chain gives, by
\texttt{tiedCore\_of\_chainTied}, a host index witnessing the core
numerical comparison. The character comparison applies to that index
because it applies to all indices. Beyond the maximum of the three
thresholds for the chain, the character comparison and
\texttt{core\_gm}, the latter estimate yields
\[
        G_\Lambda(x)\le x^{1/2+(e_D+e_L+7)\epsilon}\qquad(x\in S).
\]
For $\Re s>1/2$, put $b=\tfrac12\min(\Re s-1/2,1/2)$ and
$C=e_D+e_L+7$. With $\epsilon=\min(\epsilon_0/2,b/(|C|+1))$ one
gets $C\epsilon\le b$ and $1/2<\theta=1/2+b<\min(1,\Re s)$.
Given any lower endpoint, choose an entire interval above it, above
one, and beyond all three thresholds. Throughout it $G_\Lambda(x)
\le x^\theta$. The analytic window criterion now gives
$L(s,\psi)\ne0$. This is the deduction in
\texttt{zeroFreeAll\_of\_topG} from \texttt{level\_of\_tiedWin}.

\section{The cyclotomic CM field on a common window}
\label{sec:host-field}

Fix a primitive Dirichlet character $\psi$ of conductor $M\geq 1$, and
write $\Xi=\langle\psi\rangle$, $o=|\Xi|$. The arithmetic object supplied
to the window argument is a number field $E$ with a specified Artin map,
two quadratic characters, a CM type, ranks and level primes. We construct
its field and its type, including the sign correction in the Artin map
which is needed when $\psi$ is odd.

\subsection{The tower and the factor in the degree}

Write $U_n=(\mathbb Z/n\mathbb Z)^\times$. For $e\geq 0$ put
\[
 B=U_M/\ker\psi,\qquad D_e=2^{e+2},\qquad
 L_e=\operatorname{lcm}(D_e,M),\qquad C_e=U_{D_e}/\{1,-1\}.
\]
The group $C_e$ has order $2^e$. The order of $B$ is $o$ because
$\operatorname{im}\psi$ is cyclic of order $o$. The annihilator of
$\Xi$ on $U_M$ is $\ker\psi$, since a unit on which $\psi$ takes value one has the same property for
every power.

The source first selects a modulus $Q$ as the least divisor of $L_e$
such that $D_e\mid Q$ and the quotient map $U_{L_e}\to B$ kills all
units congruent to $1$ modulo $Q$. The modulus $L_e$ itself is
admissible. In the primitive case used here, $Q=L_e$. To prove this,
take a unit $v$ modulo $M$ congruent to $1$ modulo $h=\gcd(Q,M)$.
The Chinese remainder theorem supplies $z$ congruent to $v$ modulo
$M$ and to $1$ modulo $Q$. It is coprime to $M$ and $Q$, hence to
$L_e$ since $D_e\mid Q$. The defining property of $Q$ gives
$\psi(v)=\psi(z)=1$. The reduction $U_M\to U_h$ is surjective.
Its kernel is thus killed by $\psi$ and $\psi$ factors through $h$.
Primitivity gives $M\le h$. Since $h\mid M$, equality follows,
so $M\mid Q$. Combining $D_e\mid Q$, $M\mid Q$ and $Q\mid L_e$
proves the assertion.

Let $\eta_e:U_{L_e}\to C_e\times B$ be reduction on both factors and
let $A_e=\operatorname{im}\eta_e$. This finite abelian group is the
multiplicative presentation of the frame. Its projection to $C_e$
is surjective. Its kernel maps bijectively by the second projection
onto
\[
 B_e=\bigl\{[u\bmod M]\in B: u\in U_{L_e},\
                                   u\equiv\pm1\bmod D_e\bigr\}.
\]
Indeed a kernel element has first coordinate one; conversely the unit
in the definition of $B_e$ supplies a preimage. Therefore
$|A_e|=2^e|B_e|$. Reduction of a unit from level $L_{e+1}$ to $L_e$
preserves its $B$ coordinate and takes each sign modulo $D_{e+1}$
to the same sign modulo $D_e$. Thus $B_{e+1}\subseteq B_e$. The
positive integers $|B_e|\le o$ attain a minimum $c$ at some $e_1$,
after which they are all $c$. Consequently
\[
          1\le c\le o,\qquad n_e:=|A_e|=c2^e \quad(e\ge e_1).
\]
The constant depends only on the initial character. This proves the
degree stabilization used to choose the scale with the correct factor.

\subsection{The field, including the sign of complex conjugation}

Take $e\ge\max(e_1,2,4M,64o)$ and distinct primes $q,d$ as in
Section~\ref{sec:rh-common-windows}, with $(M,qd)=1$.
Set $N=Qqd$, where $Q=L_e$. For $a\in U_N$ let $\lambda_q(a)$ and
$\lambda_d(a)$ be its quadratic residue signs. Identify each sign with
$1$ or $-1$ in $U_Q$ and put $w(s)=\eta_e(s)\in A_e$.
In the Galois coordinates used in the source, the Artin map is
\[
 \alpha(a)=\bigl(\lambda_q(a),\lambda_d(a),
           \eta_e(a\bmod Q)\,w(\lambda_q(a))\bigr)
       \ \in\ \{1,-1\}^2\times A_e .
\]
For any desired triple $(s,t,v)$, choose a unit modulo $Q$ mapping to
$v w(s)$, a unit modulo $q$ with sign $s$, and a unit modulo $d$ with
sign $t$. Each sign occurs at an odd prime. Since $Q,q,d$ are
pairwise coprime, one unit $a\bmod N$ realizes all three choices.
As $w(s)^2=1$, we have $\alpha(a)=(s,t,v)$. Both primes are
$3\bmod4$, whence $\alpha(-1)=(-1,-1,1)$, including for odd $\psi$.

Identify $\operatorname{Gal}(\mathbb Q(\zeta_N)/\mathbb Q)$ with
$U_N$ by $\sigma_a(\zeta_N)=\zeta_N^a$. Define
$E=\mathbb Q(\zeta_N)^{\ker\alpha}$. Galois correspondence and the
surjectivity just proved give $\operatorname{Gal}(E/\mathbb Q)
=\{1,-1\}^2\times A_e$ in these coordinates, and $[E:\mathbb Q]=4n_e$.
For every integer $a$ prime to $N$ the restriction $\sigma_a|_E$
is the identity precisely when
\begin{equation}
\label{eq:host-field-kernel}
 a\equiv\pm1\bmod D_e,\qquad \psi(a)=1,\qquad
 \lambda_q(a)=1,\qquad\lambda_d(a)=1.
\end{equation}
Indeed the sign coordinates must first be one, leaving
$\eta_e(a\bmod Q)=1$ as the third condition. Its two coordinates give
the first two conditions of \eqref{eq:host-field-kernel}.
In particular $E$ is the compositum inside $\mathbb Q(\zeta_N)$
of $\mathbb Q(\zeta_{D_e})^+$, the field cut out by $\psi$ in
$\mathbb Q(\zeta_M)$, and the two quadratic fields
$\mathbb Q(\sqrt{-q})$ and $\mathbb Q(\sqrt{-d})$. Indeed the
subgroup fixing a compositum is the intersection of the four fixing
subgroups. This description allows an intersection of the first
two fields and accounts for the factor $c$ in the degree.

For the quadratic identifications, let $p=q$ or $d$, select a
primitive $p$th root $\zeta$, and extend $\lambda_p$ by zero at zero.
The Gauss sum $g=\sum_{a\bmod p}\lambda_p(a)\zeta^a$ has
$\sigma_u(g)=\lambda_p(u)g$ by change of variables. In a complex
embedding,
\[
 |g|^2=\sum_{t\in\mathbb F_p^\times}\lambda_p(t)
               \sum_{b\in\mathbb F_p^\times}\zeta^{(t-1)b}
       =(p-1)-\sum_{t\ne0,1}\lambda_p(t)=p .
\]
Also $\overline g=\lambda_p(-1)g=-g$, so $g^2=-p$.
It belongs to $E$ by \eqref{eq:host-field-kernel}. This supplies
the root and the transformation rule invoked in
\texttt{CharData.exists\_root}. Write $s_d$ for the chosen root of
$-d$ and $K=\mathbb Q(s_d)\subset E$. Every embedding of $E$ into
$\mathbb C$ sends $s_d$ to $i\sqrt d$ or $-i\sqrt d$.
Cyclotomic complex conjugation induces the nontrivial involution
$(-1,-1,1)$ on $E$ and acts by ordinary conjugation in every embedding.
Its fixed field $F$ is totally real; every embedding of $F$ extends
to $E$, and its image is pointwise fixed by complex conjugation.
As $s_d$ is purely imaginary under every embedding, $E$ is totally
imaginary and $E/F$ is quadratic. Thus $[F:\mathbb Q]=2n_e$ and the
signature of $E$ is $(0,2n_e)$. The selected CM type is
\[
   \Phi=\{\tau:E\hookrightarrow\mathbb C:\tau(s_d)=i\sqrt d\}.
\]
Conjugation pairs the embeddings, with exactly one of each pair in
$\Phi$. Thus $|\Phi|=2n_e$. For $\tau\in\Phi$ and
$\sigma=(s,t,v)$, one has $\tau\circ\sigma\in\Phi$ exactly when
$t=1$, since $\sigma(s_d)=t s_d$.

\subsection{Which characters are tested on the same window}

Let $P_2:A_e\to B$ be the second projection and $A_0=\ker P_2$.
It is surjective by surjectivity of $U_{L_e}\to U_M$.
For every complex character $\theta$ of $\{1,-1\}^2\times A_e$
trivial on the subgroup $\{(1,1,v):v\in A_0\}$, there are
$\xi\in\Xi$ and $i,j\in\{0,1\}$ such that for all $a$ prime to $N$,
\[
   \theta(\alpha(a))=\xi(a)\lambda_q(a)^i\lambda_d(a)^j,
       \qquad \theta(1,-1,1)=(-1)^j.
\]
Indeed $v\mapsto\theta(1,1,v)$ descends to $B$. Every character
of the cyclic group $B$ is a power of the faithful character induced
by $\psi$. Its pullback to $U_M$ is the required $\xi$.
Each of the homomorphisms
$s\mapsto\theta(s,1,w(s))$ and $t\mapsto\theta(1,t,1)$ is either
trivial or the sign character. Let $i,j$ be their exponents. Multiplying
the three identities following from the factorization
\[
  \alpha(a)=(\lambda_q(a),1,w(\lambda_q(a)))\,
             (1,\lambda_d(a),1)\,(1,1,\eta_e(a\bmod Q))
\]
proves the formulas. This is the comparison \texttt{TwistAny}
which permits the previously selected prime pair to serve all these
characters, including an odd initial $\psi$.

\subsection{The decomposition condition at the level primes of the field}

We give the arithmetic behind \texttt{decomp\_typeI}. For every prime
$\ell\mid Qq$, every prime $\mathfrak P$ of $E$ above $\ell$, and
every automorphism stabilizing $\mathfrak P$, its second sign is $+1$.
For odd $\ell\mid Q$, we have $\ell\mid M$, and the common window
ensures $\ell\mid d+1$ once its splitting bound exceeds $4M$.
It also ensures $q\mid d+1$. Hence $-d$ is a nonzero square
modulo each odd $\ell\mid Qq$. It is $1\bmod8$ because $d=7\bmod16$.

Here is the simultaneous square-root calculation.
If $s^2+d=\ell^j h$ for $j\ge1$ at an odd prime $\ell\mid Qq$,
choose $v\bmod\ell$ with $h+2sv\equiv0\bmod\ell$. This is possible
since $\ell\nmid 2s$. Then $(s+\ell^jv)^2+d$ is divisible by
$\ell^{j+1}$. At two start with $s=1$ modulo $8$. If
$s^2+d=2^j h$, $j\ge3$, retain $s$ when $h$ is even and take
$s+2^{j-1}$ when $h$ is odd. In the latter case the quotient at $2^j$
becomes $h+s+2^{j-2}$, which is even. Chinese remaindering the
solutions at prime powers gives one odd $s$ with $Qq\mid s^2+d$.
Since $16\mid Q$, one has $4\ell\mid Qq$ for every $\ell\mid Qq$,
including $\ell=2$.

The element $\omega=(1+s_d)/2$ is integral, with polynomial
$X^2-X+(1+d)/4$. Put $v=(1+s)/2$. Then
\[
              (\omega-v)(\omega-(1-v))=-(s^2+d)/4\in\mathfrak P .
\]
At least one factor is in $\mathfrak P$. An automorphism stabilizing
$\mathfrak P$ with second sign $-1$ would exchange their vanishing
modulo $\mathfrak P$, as it sends $\omega$ to $1-\omega$. Both then
vanish, forcing $2v-1=s\in\mathfrak P$ and hence $\ell\mid s$.
This contradicts $s^2\equiv-d\bmod\ell$ and $\ell\nmid d$.
Thus the decomposition group lies in the subgroup preserving $\Phi$.

At $d$, the cyclotomic inertia group is the subgroup of $U_N$
congruent to $1$ modulo $Qq$. Its image in $\operatorname{Gal}(E/\mathbb Q)$
has exactly the elements $(1,1,1)$ and $(1,-1,1)$, by the Artin
formula and Chinese remaindering. Here we use the cyclotomic inertia
description and the restriction of inertia to a Galois subextension.
The same two facts at $q$ give the group
$\{(1,1,1),(-1,1,w(-1))\}$. Both ramification indices are two
and the residue characteristics are odd, so these two primes are tame.

\subsection{Compatibility with the width and the additional level}

In the interval of Section~\ref{sec:rh-common-windows}, a single pair
$q,d$ works for all $Z\le x\le Z^{3/2}$. At each such $x$ the
chosen exponent satisfies $x^{2+8\epsilon}\le c2^e\le
2x^{2+8\epsilon}$. For fixed $M,\epsilon$ and $J\ge1$, choose the
lower threshold so that $x>c2^{\max(e_1,2,4M,64o)}$ and
$x^{\epsilon/(5J)}\ge\max(5,4M)$. Then $e$ meets the construction
bound and $b=\lfloor Y^{\epsilon/5}\rfloor\ge4M$, since
$Y\ge x^{1/J}$. We have $Q\le4M\,2^e\le4M n_e$ and
$Q\le n_e^2$ since $4M\le2^e\le n_e$. Every odd prime divisor
$\ell$ of $Q$ divides $M$, so $4\ell\le b$. Finally
$o^{64}\le 2^{64o}\le2^e$. With $q>3$ at sufficiently large
scale, these are the size and splitting inequalities in
\texttt{FieldRows}. The first quadratic sign and the root of $-q$
were proved above. To verify exactness, fix a prime $\ell\mid Q$. We claim that
some $v\in U_Q$ congruent to one modulo $Q/\ell$ has
$\eta_e(v)\ne1$. If $D_e\mid Q/\ell$ and every such $v$ had
trivial image, $Q/\ell$ would meet all the defining requirements
for $Q$, contradicting minimality. If $D_e\nmid Q/\ell$, then
$\ell=2$ and $Q/D_e$ is odd. The integer $1+Q/2$ is coprime to
$Q$ and is $1+D_e/2$ modulo $D_e$, distinct there from both $1$
and $-1$ since $D_e\ge16$. It gives the desired $v$. Choose by
Chinese remaindering a unit $u\bmod N$ equal to $v\bmod Q$ and one
modulo $qd$. Then $u\equiv1\bmod N/\ell$, while
$\alpha(u)=(1,1,\eta_e(v))\ne1$. This proves the exact field row
of \texttt{AnyGrp.exact\_data} and \texttt{ModData.exact}.
Equation~\eqref{eq:host-field-kernel} identifies this very $E$ with
a field allowed by the universal quantifier in \texttt{WinDatum.bound}.
Thus the strict Selmer estimate applies to it, with the same $x,e,q,d$.

The additional split level primes enter through
\texttt{levelData} and \texttt{hostOfFieldOf}. With rank $k$ and the
proof's scale parameter $b_x=\exp(x^\epsilon)$, the former supplies
two distinct primes $g_i\in[(kb_x)^8,2(kb_x)^8]$ satisfying its
quadratic norm condition, and primes $p,p',q_0$ with
\[
 p,p'\equiv1\bmod 15N,\qquad q_0\equiv1\bmod N,\qquad
             p,p',q_0\le N^{C_8}.
\]
The three latter primes are pairwise distinct, and $p,p'$ differ
from each $g_i$. The norm condition gives $g_i\equiv1\bmod N$.
Therefore all the indicated level primes split completely in $E$,
since their cyclotomic Frobenius maps to the identity. The source
obtains the prime selection from \texttt{ShortNorm} and
\texttt{SmallOne}. The rank construction, the choices of thick
subsets, and the automorphic design are additional results. In
particular, when a regular algebraic cuspidal representation over $E$
is used later, regular algebraicity is a condition on that
representation. Computing the signature of $E$ supplies the
archimedean places over which that condition is formulated.

\paragraph{Source identifications.}
The conductor calculation and kernel appear in
\texttt{CharacterBoxes.KerChiAny}, the degree argument in
\texttt{CharacterBoxes.OrdAny}, and the field and rank bounds in
\texttt{CharacterBoxes.FieldChi} and \texttt{ShortNorms.supplyOfAdmOf}.
In the source directory \texttt{InertiaGaloisSymmetry},
\texttt{AnyFrame} gives the Artin map, \texttt{OfChar} takes its fixed
field, \texttt{CycBasic} and \texttt{CycAll} prove the degree and CM type,
and \texttt{CycGrp} proves the stabilizer and inertia statements.
The field identification uses elementary cyclotomic Galois theory,
finite abelian character duality, and the cyclotomic decomposition of
inertia. The common prime-selection and analytic window estimates are
used with the quantifiers in Section~\ref{sec:rh-common-windows}.

\section{Constructing the quadratic ray-class family}
\label{sec:quadratic-ray-family}

Let \(d\geq100\) be a prime with \(d\equiv3\bmod4\), let \(M\geq3\),
and set
\[
 K=\mathbb Q(\sqrt{-d}),\quad
 {\cal O}=\mathbb Z[\omega],\quad
 \omega=(1+\sqrt{-d})/2,\quad D=4dM^2.
\]
The discriminant is \(-d\), and the area of \({\cal O}\) in
\(\mathbb C\) is \(\sqrt d/2\).
The norm formula
\(N((A+B\sqrt{-d})/2)=(A^2+dB^2)/4\) shows that the units are
\(1,-1\). Since \(M\geq3\), the only unit congruent to one modulo
\(M{\cal O}\) is one.

Let \(G\) be the ray class group for \(M{\cal O}\), with no infinite
part. Its elements are fractional ideals prime to \(M\), modulo
principal ideals generated by elements congruent to one at the primes
of \(M\). Write its law additively and its trivial class as zero.
Every class has an integral representative prime to \(M\): given a
fractional representative, choose a denominator congruent to one at
the modulus and divisible by all its other denominator primes, using
the Chinese remainder theorem. Multiplication by that denominator
leaves the ray class unchanged.

The ray-class exact sequence gives
\[
 |G|\leq h_K|({\cal O}/M{\cal O})^\times|\leq h_KM^2.
\]
For completeness, the first inequality follows by mapping an ideal
to its ordinary ideal class. Its kernel is generated by principal
ideals prime to \(M\), whose ray classes are determined by the
generator's residue modulo \(M\), up to a unit. Reduction of a
fractional generator uses its local denominator inverse at \(M\).
Every ordinary ideal class has an ideal prime to \(M\) by the same
approximation argument.

Minkowski's ideal-class bound in this quadratic imaginary field is
\((2/\pi)\sqrt d\). The number of ideals of norm \(n\) is at most
the divisor function \(\tau(n)\). This follows prime by prime: a
split rational prime gives \(k+1\) ideals of norm \(p^k\), an inert
one gives zero or one, and a ramified one gives one. Hence, putting
\(Y=(2/\pi)\sqrt d\),
\[
 h_K\leq\sum_{n\leq Y}\tau(n)
       =\sum_{ab\leq Y}1
       \leq Y(1+\log Y)\leq4d .
\]
The last inequality follows from \(Y\leq\sqrt d\),
\(1+\log\sqrt d\leq4\sqrt d\), and \(d\geq100\).
Thus \(G\) is finite and
\begin{equation}
 \label{eq:ray-group-budget}
                            H:=|G|\leq D.
\end{equation}

\subsection{One lattice for each ray class}

Let \(a_g(n)\) count integral ideals prime to \(M\), in class \(g\),
of norm \(n\). Fix an integral ideal \(J\) prime to \(M\) representing
\(-g\). The Chinese remainder theorem supplies
\(\alpha_0\in J\) with \(\alpha_0\equiv1\bmod M{\cal O}\).
The map
\begin{equation}
 \label{eq:ray-lattice-bijection}
 \alpha\longmapsto(\alpha)J^{-1}
\end{equation}
is a bijection from \(\alpha_0+MJ\) to the ideals counted in
\(a_g\). An element in this coset lies in \(J\), so its image is
integral. It is prime to \(M\), has ray class \(g\), and is nonzero.
Conversely \(I\) in class \(g\) makes \(IJ\) ray principal, with
a generator \(\alpha\equiv1\bmod M\). This generator lies in \(J\).
Also
\(J\cap M{\cal O}=MJ\) by coprimality, so it lies in the indicated
coset. Two such generators of the same ideal differ by a unit
congruent to one modulo \(M\), which is one. This proves injectivity.
The assertion about an integral generator congruent to one follows
also with the localization definition of the ray relation. An integral
element of \(K\) congruent to one at every prime of \(M\) has difference
from one in \(M{\cal O}\), as is seen from the valuations at those primes.

Use the real quadratic form
\[
                         q(\alpha)=\frac{|\alpha|^2}{N(J)}
\]
on this coset. For \(\alpha\in J\), its value is the norm of
\((\alpha)J^{-1}\), and is therefore a positive integer on
\(\alpha_0+MJ\).

Normalize the lattice \(MJ\subset\mathbb C\) by \(N(J)^{-1/2}\).
Its area is \(M^2\sqrt d/2\).
Every nonzero vector is \(M\beta/N(J)^{1/2}\) with \(\beta\in J\),
and its squared length is at least \(M^2\), since
\(N((\beta)J^{-1})\) is a positive integer.
Choose a shortest nonzero vector \(u\). It is primitive, since an
integer multiple of another lattice vector would be longer.
Extend it to a basis \(u,v\), and replace \(v\) by \(v-ku\), with
\(k\) the nearest integer to its projection coefficient.
Then, writing
\[
 \alpha=|u|^2,\qquad
 l=\frac{\langle u,v\rangle}{|u|^2},\qquad
 \gamma=|v|^2-l^2|u|^2,
\]
we have
\[
 \alpha\geq M^2,\quad |l|\leq1/2,\quad
 \gamma\geq3\alpha/4\geq1,\quad
                   \alpha\gamma=dM^4/4 .
\]
Indeed \(|v|\geq|u|\) by minimality of \(u\).
In these coordinates the coset has form
\[
 Q(m,n)=\alpha\{m+\theta_1+l(n+\theta_2)\}^2
                           +\gamma(n+\theta_2)^2 .
\]
It follows from \eqref{eq:ray-lattice-bijection} that
\begin{equation}
 \label{eq:ray-fiber-count}
        a_g(k)=\#\{(m,n)\in\mathbb Z^2:Q(m,n)=k\}.
\end{equation}
All values \(Q(m,n)\) are positive integers. The constants
\(\theta_1,\theta_2,l,\alpha,\gamma\) may depend on \(g\), but
their product \(\alpha\gamma\) is the same for every class.

\subsection{Uniform counting and theta estimates}

For \(x\geq1\), a row \(n+\theta_2=y\) contributes integers \(m\)
in an interval of length
\[
                 f(y)=2\sqrt{(x-\gamma y^2)_+/\alpha}.
\]
For rows with \(|y|\leq\sqrt{x/\gamma}\), the discrepancy between
their integer count and that length is at most one.
There are at most \(2\sqrt{x/\gamma}+1\) such rows.
The function \(f\) is nonnegative and unimodal, with total variation
\(2\sup f=4\sqrt{x/\alpha}\).
Comparing its value at the left endpoint of each interval of length
one with its integral on that interval gives
\[
       \left|\sum_{n\in\mathbb Z}f(n+\theta_2)
                              -\int_{\mathbb R}f(y)\,dy\right|
                           \leq4\sqrt{x/\alpha}.
\]
To see this inequality directly, the discrepancy on one interval is
at most the variation of \(f\) on that interval. Sum these variations
over the partition. The function has compact support, so the sum is
its total variation.
The ellipse area is
\(\int f=\pi x/\sqrt{\alpha\gamma}\). Therefore
\[
 \left|\#\{(m,n):Q(m,n)\leq x\}
                         -\frac{\pi x}{\sqrt{\alpha\gamma}}\right|
 \leq2\sqrt{x/\gamma}+1+4\sqrt{x/\alpha}
                                  \leq10\sqrt x.
\]
Together with \eqref{eq:ray-fiber-count}, this proves
\begin{equation}
 \label{eq:ray-count-done}
 \left|\sum_{n\leq x}a_g(n)-\kappa x\right|\leq10\sqrt x,
                        \qquad \kappa=\frac{2\pi}{\sqrt d M^2}.
\end{equation}

We also prove the theta estimate uniformly in the shifts and shear.
For \(a>0\), the periodized Gaussian has Fourier series
\[
 \sum_{m\in\mathbb Z}e^{-a(m+c)^2}
       =\sqrt{\frac{\pi}{a}}\sum_{k\in\mathbb Z}
                         e^{-\pi^2k^2/a}e^{2\pi ikc}.
\]
Its coefficients are obtained by integrating on one period and
translating to the Gaussian integral on \(\mathbb R\).
Both the periodized sum and its Fourier series converge absolutely.
The Gaussian Fourier integral follows by differentiation in frequency
and the Gaussian integral at zero.
Since \(e^{-z}\leq2z^{-2}\) for \(z>0\), and
\(\sum_{k\geq1}k^{-4}\leq4/3\), this implies
\[
 \left|\sum_me^{-a(m+c)^2}-\sqrt{\pi/a}\right|
   \leq\sqrt{\pi/a}\frac{4a^2}{\pi^4}\sum_{k\geq1}k^{-4}
   \leq\sqrt{\pi/a}\,a^2 .
\]
The last inequality uses \(\pi>3\).

Apply this first in \(m\), at \(a=\alpha t\), uniformly in
\(c=\theta_1+l(n+\theta_2)\), and then in \(n\), at
\(a=\gamma t\). Positivity of the outer Gaussian gives
\[
 \left|\sum_{m,n}e^{-tQ(m,n)}
                      -\frac{\pi}{t\sqrt{\alpha\gamma}}\right|
 \leq\frac{\pi}{t\sqrt{\alpha\gamma}}
        \{(\alpha t)^2+(\gamma t)^2+(\alpha t)^2(\gamma t)^2\}.
\]
The same estimates establish absolute convergence of the double
sum. For \(0<t\leq1\), \(\alpha,\gamma\geq1\) imply
\[
 \alpha^2+\gamma^2+\alpha^2\gamma^2t^2
                                     \leq3(\alpha\gamma)^2.
\]
Since \(\alpha\gamma=dM^4/4\), the error is at most
\[
        3\pi(\alpha\gamma)^2t
             =\frac{3\pi}{16}d^2M^8t\leq D^{10}t .
\]
Grouping the positive terms by their positive integer norm
\(Q(m,n)\) gives the desired identity
\begin{equation}
 \label{eq:ray-theta-done}
       \left|\sum_{n\geq1}a_g(n)e^{-tn}-\frac{\kappa}{t}\right|
                                  \leq D^{10}t .
\end{equation}
The constants also satisfy
\(D^{-1}\leq\kappa\leq4\), by \(D=4dM^2\), \(d\geq100\),
\(M\geq3\) and \(3<\pi<4\).

\subsection{The prime weights and the positive convolution}

For an integral ideal \(I\) prime to \(M\), define
\(\Lambda_K(I)=\log N(P)\) when \(I=P^j\) for a prime ideal \(P\)
and \(j\geq1\), and zero otherwise. Define
\[
       \lambda_g(n)=
        \sum_{\substack{N(I)=n\\ [I]=g\\(I,M)=1}}\Lambda_K(I).
\]
These weights are nonnegative. The only nonzero terms occur at
rational prime powers \(n=p^k\).
For a prime \(P\) of residue degree \(f_P\) over \(p\), a contribution
at \(p^k\) is \(f_P\log p\), and occurs only if \(f_P\mid k\).
Since \(\sum_{P\mid p}e_Pf_P=2\), we get
\begin{equation}
 \label{eq:ray-lambda-budget}
                         \sum_g\lambda_g(n)\leq2\Lambda(n).
\end{equation}
Omitting the primes dividing the modulus can only decrease the sum.

Unique factorization of ideals gives
\(\log N(I)=\sum_{J\mid I}\Lambda_K(J)\).
Partitioning the divisors by norm and ray class proves
\[
 a_g(n)\log n
       =\sum_{h\in G}\sum_{ab=n}\lambda_h(a)a_{g-h}(b).
\]
All sums are finite. Divisors and complementary divisors of an ideal
prime to \(M\) are also prime to \(M\).

For a real character \(\chi\), the ideal coefficient of the
convolution of the trivial and \(\chi\) series at an ideal
\(C=\prod_PP^{e_P}\) is
\[
             r(C)=\sum_{J\mid C}\chi([J])
                      =\prod_P\sum_{j=0}^{e_P}\chi([P])^j\geq0.
\]
Each factor is \(e_P+1\) when \(\chi([P])=1\), and is zero or
one when \(\chi([P])=-1\). Grouping by \(N(C)=n\) proves
\[
                      \sum_{ab=n}a_1(a)a_\chi(b)\geq0,
\]
which is \eqref{eq:ray-convolution-positive}.
The second positivity property in the formal family follows too.
For every ideal \(I\), all exponents in \(I^2\) are even, so
\(r(I^2)\geq1\). The map \(I\mapsto I^2\) is injective by unique
factorization, and \(N(I^2)=N(I)^2\).
It follows that
\[
                  a_1(n)\leq\sum_{ab=n^2}a_1(a)a_\chi(b).
\]

Finally, if \(p\) is prime and \(\lambda_0(p)\ne0\), some ideal of
norm \(p\) is ray principal. It has a generator
\(\alpha\equiv1\bmod M{\cal O}\).
Write
\(\alpha=1+M(a+b\omega)=(A+B\sqrt{-d})/2\).
Then
\[
  A=2+2Ma+Mb,\qquad B=Mb,\qquad A^2+dB^2=4p.
\]
Thus \(M\mid B\) and \(2M\mid A-2-B\), which is precisely
\(\operatorname{NormOne}(d,M,p)\).

We have constructed every counting, theta, convolution and support
property of the family used in
Sections~\ref{sec:ray-real-gap} and~\ref{sec:level-prime-selection},
with the same parameter \(D=4dM^2\). The remaining analytic input
to the norm-prime count is the characterwise explicit formula
with its unique possible exceptional zero.

\paragraph{Source map.}
This argument follows
\srcid{RayClassIdealCountEqNcardShiftedBinaryForm},
\srcid{QFormCount}, \srcid{ThetaQForm},
\srcid{RayCountThetaFibreCount}, \srcid{FiniteIdealsIn},
\srcid{ConvRealChar}, and \srcid{ExistsLFamRay}.
Minkowski's ideal-class bound and the ray-class exact sequence are
used in their explicit standard forms stated above.
\footnote{\href{https://github.com/Deicyde/riemann-hypothesis/blob/40e56bab026552e089d7afcfeff08d2defb9c005/src/Development/AbelianVarietyTheory/Proofs/GlobalInterpolation/ExistsLFamRay.lean}{Assembly of the ray family};
\href{https://github.com/Deicyde/riemann-hypothesis/blob/40e56bab026552e089d7afcfeff08d2defb9c005/src/Development/AbelianVarietyTheory/Proofs/GlobalInterpolation/RayClassIdealCountEqNcardShiftedBinaryForm.lean}{the ideal and lattice bijection};
\href{https://github.com/Deicyde/riemann-hypothesis/blob/40e56bab026552e089d7afcfeff08d2defb9c005/src/Development/AbelianVarietyTheory/Proofs/GlobalInterpolation/ConvRealChar.lean}{the two real-character positivity assertions}.}

\section{A polynomial lower bound for real ray-class values}
\label{sec:ray-real-gap}

The real exceptional zero in
Section~\ref{sec:level-prime-selection} requires a polynomial separation
from one. We prove the quantitative implication from the ray-family
counting and theta identities. All constants below are independent of
the class group and its modulus.

Let \(G\) be a finite abelian group of order \(H\), let \(D\geq4\),
\(H\leq D\), and let \(a_g(n)\geq0\) satisfy
\[
 \begin{gathered}
 a_g(0)=0,\qquad a_0(1)=1,\qquad a_g(1)=0\quad(g\ne0),\\
 D^{-1}\leq\kappa\leq4,\qquad
 \left|\sum_{n\leq u}a_g(n)-\kappa u\right|
                                  \leq10\sqrt u\quad(u\geq1),\\
 \left|\sum_{n\geq1}a_g(n)e^{-tn}-\frac{\kappa}{t}\right|
                                   \leq D^{10}t\quad(0<t\leq1).
 \end{gathered}
\]
For every real character \(\chi:G\to\{1,-1\}\), put
\(a_\chi(n)=\sum_g\chi(g)a_g(n)\).
We require the positivity furnished by the ray-family Euler product,
in its precise coefficient form
\begin{equation}
 \label{eq:ray-convolution-positive}
          \sum_{mn=k}a_1(m)a_\chi(n)\geq0\qquad(k\geq1).
\end{equation}
The subscript \(1\) denotes the trivial character. These are exactly
the counting, theta, and real-convolution properties used from
\srcid{LFam}. This section proves their analytic consequence. Section~\ref{sec:quadratic-ray-family} constructs them from the ideals
of the quadratic field.

\subsection{Continuation and uniform bounds near one}

Nonnegativity and the counting bound give, for \(n\geq1\),
\[
       a_g(n)\leq\kappa n+10\sqrt n\leq14n,
                         \qquad |a_\chi(n)|\leq14Dn.
\]
Let \(\rho=\kappa H\), and write
\[
 A_\chi(u)=\sum_{n\leq u}a_\chi(n),\qquad
 R_\chi(u)=A_\chi(u)-{\bf1}_{\chi=1}\rho u .
\]
Orthogonality gives \(|R_\chi(u)|\leq10D\sqrt u\).
Consequently
\[
        I_\chi(z)=\int_1^\infty R_\chi(u)u^{-z-1}\,du
\]
is holomorphic for \(\Re z>1/2\), with
\[
 |I_\chi(z)|\leq\frac{10D}{\Re z-1/2},\qquad
 |I'_\chi(z)|\leq\frac{10D}{(\Re z-1/2)^2}.
\]
Both assertions follow by domination on every closed subregion of
that half-plane. Differentiation inserts \(-\log u\), and
\(\int_1^\infty u^{-a-1}\log u\,du=a^{-2}\) for \(a>0\).

For \(\Re z>1\), summation by parts gives
\[
       L_\chi(z)
          ={\bf1}_{\chi=1}\frac{\rho z}{z-1}
                              +zI_\chi(z).
\]
Thus the regularized function used in the source is
\[
 F_\chi(z)=
 \begin{cases}
   \rho z+(z-1)zI_\chi(z),&\chi=1,\\
   zI_\chi(z),&\chi\ne1.
 \end{cases}
\]
It is holomorphic for \(\Re z>1/2\), and \(F_1(1)=\rho\).

For \(|z-1|\leq1/4\), we have
\(\Re z\geq3/4\), \(|z|\leq5/4\),
\(|I_\chi(z)|\leq40D\), and \(|I'_\chi(z)|\leq160D\).
If \(\chi\ne1\), differentiation gives
\[
                   |F'_\chi(z)|\leq40D+(5/4)160D=240D.
\]
For the trivial character, use
\[
 F'_1(z)=\rho+(2z-1)I_1(z)+(z-1)zI'_1(z).
\]
Since \(\rho\leq4D\), its norm is at most
\(4D+(3/2)40D+(5/16)160D=114D\).
In particular the weaker uniform bound used by the formal argument holds:
\begin{equation}
 \label{eq:ray-derivative-bound}
                 |F'_\chi(z)|\leq18000D
                             \qquad(|z-1|\leq1/4).
\end{equation}

\subsection{The smoothed positive kernel}

For \(y>0\) define
\[
            W(y)=\int_0^\infty e^{-y/w-w}\,\frac{dw}{w}.
\]
It is positive and decreasing in \(y\).
The elementary inequality
\(e^{-y/w}\leq6w^3/y^3\) gives
\[
           0<W(y)\leq\frac6{y^3}
                       \int_0^\infty w^2e^{-w}\,dw
                     =\frac{12}{y^3}.
\]
In particular \(w_0=W(1)>0\) is an absolute finite constant.

Take a nontrivial real character \(\chi\), and set
\[
       P(t)=\sum_{n\geq1}a_1(n)e^{-tn},\qquad
       Q(t)=\sum_{n\geq1}a_\chi(n)e^{-tn}.
\]
For \(0<t\leq1\), summing the theta bounds and using character
orthogonality gives
\[
       |P(t)-\rho/t|\leq D^{11}t,\qquad
                          |Q(t)|\leq D^{11}t.
\]
For \(t\geq1\), the coefficient bound gives
\[
 |P(t)|,\ |Q(t)|
       \leq14D\sum_{n\geq1}ne^{-nt}
       =14D\frac{e^{-t}}{(1-e^{-t})^2}.
\]
There is therefore an absolute constant \(C\geq1\), for instance
\(C=100\), for which, with \(E=CD^{11}\),
\begin{equation}
 \label{eq:ray-theta-kernel-bounds}
 \begin{split}
       |P(t)-\rho/t|&\leq E\min(t,t^{-1}),\\
                    |Q(t)|&\leq E\min(t,t^{-1}),\\
       |Q(t)|&\leq CD e^{-t/2}\quad(t\geq1).
 \end{split}
\end{equation}
For the first two inequalities on \(t\geq1\), multiply the
preceding exponential estimate by \(t\) and use \(te^{-t}\leq1\),
together with \(\rho\leq4D\).
These bounds prove that \(Q\) is integrable at zero and infinity.

For \(v>0\), expansion under the integral gives
\begin{equation}
 \label{eq:ray-kernel-product}
 J(v):=\int_0^\infty P(v/w)Q(w)\,\frac{dw}{w}
           =\sum_{m,n\geq1}a_1(m)a_\chi(n)W(vmn).
\end{equation}
Indeed the substitution \(u=nw\) transforms
\(\int e^{-vm/w-nw}\,dw/w\) to \(W(vmn)\).
Every exchange is justified absolutely by
\[
 \sum_{m,n\geq1}|a_1(m)a_\chi(n)|W(vmn)
       \leq\frac{12(14D)^2}{v^3}
                              \sum_{m,n\geq1}\frac1{m^2n^2}<\infty .
\]
Grouping by \(mn=k\) and applying
\eqref{eq:ray-convolution-positive} yields
\[
                         J(v)\geq W(v),
\]
because the coefficient at \(k=1\) is one.

Let
\[
       \omega(w)=\frac{\min(w,w^{-1})}{w\sqrt w},\qquad
       K=\int_0^\infty\omega(w)\,dw=\frac83 .
\]
Write \(P(v/w)=\rho w/v+R_v(w)\).
The inequality \(\min(u,u^{-1})\leq\sqrt u\), for \(u>0\),
and \eqref{eq:ray-theta-kernel-bounds} give
\[
       \left|\frac{R_v(w)Q(w)}w\right|
                         \leq E^2\sqrt v\,\omega(w).
\]
The right side is integrable. Thus, putting
\(L=\int_0^\infty Q(w)\,dw\), we obtain
\begin{equation}
 \label{eq:ray-upper-kernel}
                  J(v)\leq\frac{\rho L}{v}+E^2\sqrt v K .
\end{equation}
This also gives an independent integrable majorant for the
decomposition of the integral.

Set
\[
 \delta=\min\left(1,\frac{w_0}{2(C^2K+1)}\right),\qquad
                         s=\delta D^{-22},\qquad v=s^2 .
\]
Then \(0<v\leq1\), so \(W(v)\geq w_0\), and
\[
            E^2\sqrt v K=C^2D^{22}sK=C^2\delta K\leq w_0/2 .
\]
Combining the lower and upper estimates for \(J(v)\) gives
\[
                \frac{w_0}{2}\leq\frac{\rho L}{s^2},
                    \qquad
                L\geq\frac{w_0\delta^2}{8D^{45}},
\]
where the second inequality uses \(0<\rho\leq4D\).

\subsection{Identifying the value at one}

For \(\Re z>2\), absolute convergence and the gamma integral give
\[
       \int_0^\infty Q(t)t^{z-1}\,dt
                   =\Gamma(z)L_\chi(z).
\]
The coefficient bound \(14Dn\) justifies the exchange of sum and
integral in that region. The small-\(t\) bound \(Q(t)=O(D^{11}t)\)
and the exponential bound at infinity show that the left side is
holomorphic for \(\Re z>-1\), by domination on compact subsets.
The gamma function is holomorphic and nonzero for \(\Re z>0\).
The identity theorem, applied in \(\Re z>1/2\), therefore identifies
this integral divided by \(\Gamma(z)\) with the continuation
\(F_\chi(z)\) there. At \(z=1\), since \(\Gamma(1)=1\), it gives
\[
                           F_\chi(1)=L.
\]
For the trivial character \(F_1(1)=\rho\geq D^{-1}\).
Consequently the positive absolute constant
\[
                  c_0=\min(1,w_0\delta^2/8)
\]
satisfies, for every real character including the trivial one,
\begin{equation}
 \label{eq:ray-real-value}
                         |F_\chi(1)|\geq c_0D^{-45}.
\end{equation}
If \(3/4\leq\beta<1\) and \(F_\chi(\beta)=0\), integrate
\(F'_\chi\) along the segment from \(\beta\) to \(1\) and use
\eqref{eq:ray-derivative-bound}. Then
\[
       c_0D^{-45}\leq|F_\chi(1)|
                         \leq18000D(1-\beta),
 \qquad
             1-\beta\geq\frac{c_0}{18000}D^{-46}.
\]
This supplies the required polynomial gap with an explicit
admissible exponent. The ray-family identities proved in Section~\ref{sec:quadratic-ray-family}
and the uniform characterwise explicit formula enter the norm-prime count. The lower bound at one and the
deduction of the real-zero gap have been proved here.

\paragraph{Source map.}
The exponential estimates correspond to
\srcid{ExpSumBounds}; the kernel and its positive convolution to
\srcid{NormRayClassFamilyFOneGeOfRealCharacter}; and the last integration
to \srcid{ExistsRayClassLFunctionRealZeroFarFromOne}.
\footnote{\href{https://github.com/Deicyde/riemann-hypothesis/blob/40e56bab026552e089d7afcfeff08d2defb9c005/src/Development/AbelianVarietyTheory/Definitions/GlobalTwisting/LamIdealOfNotPrimePow.lean}{The ray-family hypotheses};
\href{https://github.com/Deicyde/riemann-hypothesis/blob/40e56bab026552e089d7afcfeff08d2defb9c005/src/Development/AbelianVarietyTheory/Proofs/GlobalInterpolation/NormRayClassFamilyFOneGeOfRealCharacter.lean}{the positive-kernel proof};
\href{https://github.com/Deicyde/riemann-hypothesis/blob/40e56bab026552e089d7afcfeff08d2defb9c005/src/Development/AbelianVarietyTheory/Proofs/GlobalInterpolation/ExistsRayClassLFunctionRealZeroFarFromOne.lean}{the real-zero gap}.}

\section{The ray-family explicit formula and its exceptional zero}
\label{sec:ray-explicit}

Use the family constructed in Section~\ref{sec:quadratic-ray-family}.
For a character \(\chi\) write
\[
 b_\chi(n)=\sum_g\chi(g)\lambda_g(n),\qquad
 {\cal L}_\chi(z)=\sum_{n\geq1}b_\chi(n)n^{-z},\qquad
                \ell(t)=\log(D(|t|+2)).
\]
Equation~\eqref{eq:ray-lambda-budget} gives
\(|b_\chi(n)|\leq2\Lambda(n)\).
The ideal Euler product, absolutely convergent for \(\Re z>1\), gives
\[
 {\cal L}_\chi(z)=-\frac{L'_\chi(z)}{L_\chi(z)}
       ={\bf1}_{\chi=1}\frac1{z-1}-\frac{F'_\chi(z)}{F_\chi(z)}.
\]
To check convergence, the counting estimates give a total of
\(O(Du)\) ideals of norm at most \(u\).
The logarithm of the product is also dominated by
\(2\sum_{p,j}p^{-j\Re z}/j\).
Its exponential is nonzero, so every zero of \(F_\chi\) in
\(\Re z>1/2\) has real part at most one.

\subsection{A local formula with uniform constants}

There is an absolute \(C>0\) such that the zeros of \(F_\chi\)
in the disk
\[
             |z-(5/4+it)|\leq3/5
\]
have total multiplicity at most \(C\ell(t)\), and, for
\(|z-(5/4+it)|\leq1/2\) away from the zeros,
\begin{equation}
 \label{eq:ray-local-log}
 \left|\frac{F'_\chi(z)}{F_\chi(z)}
             -\sum_\rho\frac{m_\rho}{z-\rho}\right|
                                  \leq C\ell(t).
\end{equation}
We prove the assertion for the regularized principal function too.
On \(\Re z\geq11/20\), the summatory integral used above gives
\[
                    |F_\chi(z)|\leq300D(1+|z|)^2.
\]
For nontrivial characters use
\(|I_\chi(z)|\leq10D/(\Re z-1/2)\).
For the trivial character use
\(F_1(z)=\rho z+(z-1)zI_1(z)\), \(\rho\leq4D\).
At the center \(z_0=5/4+it\), the Euler product gives
\[
 |L_\chi(z_0)|\geq\zeta(5/4)^{-2},\qquad
                  |F_\chi(z_0)|\geq\frac14\zeta(5/4)^{-2}.
\]
Indeed the absolute value of the logarithm of the Euler product is
at most \(2\log\zeta(5/4)\); the extra factor in \(F_1\) has norm
at least \(1/4\).
The ratio of the upper bound on the disk of radius \(7/10\) to
this lower bound is \(M_0\) with
\(\log M_0\leq C_1\ell(t)\), for an absolute \(C_1\).

Here is the disk argument. Jensen's formula bounds the number \(k\)
of zeros in the disk of radius \(3/5\), with multiplicity, by
\(\log M_0/\log(7/6)\).
Boundary zeros can be handled by taking radii increasing to \(7/10\).
Let \(P\) be the product of the corresponding monic linear factors,
and set \(g=F_\chi/P\), extending across those zeros.
On the outer circle, \(|P|\geq(1/10)^k\), while
\(|P(z_0)|\leq(3/5)^k\). The maximum modulus principle gives
\[
                  |g(z)/g(z_0)|\leq M_0\,6^k .
\]
The function \(g\) has no zeros in the open disk of radius \(3/5\).
Its normalized logarithm \(h\) there has \(h(z_0)=0\) and
\(\Re h\leq A:=\log M_0+k\log6\).
For \(a>A\), Schwarz's lemma applied to \(h/(2a-h)\) gives
\[
                  |h(z)|\leq\frac{2a|z-z_0|}{3/5-|z-z_0|}.
\]
On the disk of radius \(11/20\) this is at most \(22a\).
Cauchy's derivative estimate on a circle of radius \(1/20\)
about any point in the disk of radius \(1/2\), followed by
\(a\downarrow A\), gives \(|h'|\leq440A\).
This proves \eqref{eq:ray-local-log} and the asserted zero count.

Since \(\Re\rho\leq1\), every zero term in
\eqref{eq:ray-local-log} has nonnegative real part at
\(s=1+a+it\), \(a>0\). Dropping unused zero terms gives, for
\(0<a\leq1/2\),
\begin{equation}
 \label{eq:ray-real-log-upper}
 \Re{\cal L}_\chi(1+a+it)
 \leq{\bf1}_{\chi=1}\frac{a}{a^2+t^2}
       +C\ell(t)-\sum_{\rho\in S}
                m_\rho\Re\frac1{1+a+it-\rho},
\end{equation}
for any selection \(S\) of the local zeros. Enlarge \(C\), when
necessary, so that the same constant works throughout.

\subsection{The zero-free strip}

Positivity of the prime weights gives
\begin{equation}
 \label{eq:ray-341}
 3{\cal L}_1(\sigma)+4\Re{\cal L}_\chi(\sigma+it)
                        +\Re{\cal L}_{\chi^2}(\sigma+2it)\geq0
                        \qquad(\sigma>1).
\end{equation}
Termwise this is \(3+4\cos u+\cos2u=2(1+\cos u)^2\geq0\).
For real \(\chi,\psi\) it also gives
\[
 {\cal L}_1(\sigma)\geq0,\qquad
 {\cal L}_1(\sigma)+{\cal L}_\chi(\sigma)\geq0,
\]
and
\[
 {\cal L}_1(\sigma)+{\cal L}_\chi(\sigma)
       +{\cal L}_\psi(\sigma)+{\cal L}_{\chi\psi}(\sigma)\geq0.
\]
The last two assertions follow respectively from \(1+\chi(g)\geq0\)
and \((1+\chi(g))(1+\psi(g))\geq0\).

Suppose \(\rho=1-\delta+i\gamma\) is a zero,
\(0\leq\delta\leq\epsilon\leq1/400\), with
\(\chi^2\ne1\) or \(\gamma\ne0\).
We show
\begin{equation}
 \label{eq:ray-epsilon}
                           \epsilon^{-1}\leq300C\ell(\gamma).
\end{equation}
First suppose that \(\chi^2\ne1\) or \(|\gamma|\geq30\epsilon\).
Use \(a=6\epsilon\) in
\eqref{eq:ray-real-log-upper} and \eqref{eq:ray-341}.
The selected zero contributes at least \(1/(7\epsilon)\).
Every possible principal pole at \(\gamma\) or \(2\gamma\) is at
most \(1/(156\epsilon)\), since
\[
          \frac{6\epsilon}{(6\epsilon)^2+\gamma^2}
                         \leq\frac1{156\epsilon}
                         \quad(|\gamma|\geq30\epsilon).
\]
When \(\chi^2\ne1\), both of those possible pole terms vanish.
Also \(\ell(0)\leq\ell(\gamma)\) and
\(\ell(2\gamma)\leq2\ell(\gamma)\).
Thus
\[
 0\leq
 \left(\frac12-\frac47+\frac5{156}\right)\epsilon^{-1}
                                      +9C\ell(\gamma)
       =-\frac{43}{1092}\epsilon^{-1}+9C\ell(\gamma),
\]
which implies \eqref{eq:ray-epsilon}.

In the other case \(\chi\) is real and
\(0<|\gamma|<30\epsilon\).
Conjugation gives the distinct second zero \(1-\delta-i\gamma\).
Both zeros lie in the local disk centered at \(5/4\).
Use \(a=60\epsilon\).
If \(\chi=1\), apply \({\cal L}_1(\sigma)\geq0\) directly.
If \(\chi\ne1\), apply
\({\cal L}_1(\sigma)+{\cal L}_\chi(\sigma)\geq0\), using the
empty zero selection for the first term.
In either case \eqref{eq:ray-real-log-upper} gives
\[
 2\frac{60\epsilon+\delta}{(60\epsilon+\delta)^2+\gamma^2}
                     \leq\frac1{60\epsilon}+2C\ell(0).
\]
The fraction on the left before multiplication by \(2\) is at
least \(60/(4621\epsilon)\). Subtracting \(1/(60\epsilon)\)
gives
\[
           \frac{2579}{277260}\epsilon^{-1}
                                      \leq2C\ell(0),
\]
which again implies \eqref{eq:ray-epsilon}.

Set \(K=300C\) and \(c_1=(4K+400)^{-1}\).
If such a zero had
\(\delta\leq c_1/\ell(\gamma)\), take
\(\epsilon=\{(2K+200)\ell(\gamma)\}^{-1}\).
Then \(0<\epsilon\leq1/400\), \(\delta\leq\epsilon\), and
\eqref{eq:ray-epsilon} would give \(2K+200\leq K\).
Consequently every zero in
\(\Re z\geq1-c_1/\ell(\Im z)\) is real and belongs to a real
character.

We also prove uniqueness and simplicity across that real family.
Choose
\[
                 c_2=\frac1{100(C+1)},\qquad
                    \sigma=1+\frac{2c_2}{\ell(0)}.
\]
A zero \(\beta\geq1-c_2/\ell(0)\), of multiplicity \(m\), contributes
at least \(m\ell(0)/(3c_2)\) to the zero sum at \(\sigma\).
If one character had two such zeros, or one multiple zero, choose
two zeros counting multiplicity.
Use \({\cal L}_1\geq0\) when that character is trivial, and
\({\cal L}_1+{\cal L}_\chi\geq0\) otherwise.
If two different real characters had such zeros, use
\({\cal L}_1+{\cal L}_\chi\geq0\) when one is trivial, and the
four-term positive expression when both are nontrivial.
In the last case \(\chi\psi\ne1\).
There is one principal pole in each expression. Its size is
\(\ell(0)/(2c_2)\), and the combined error is at most \(4C\ell(0)\).
Thus every case would imply
\[
             \frac{2}{3c_2}\leq\frac1{2c_2}+4C,
                         \qquad \frac1{6c_2}\leq4C ,
\]
contrary to the choice of \(c_2\).
There is therefore at most one pair \((\chi,\beta)\) in this
strip, and its zero is simple.

Fix
\[
                    c_0=\min(c_1,c_2,1/20)>0 .
\]
Let \(Z_\chi\) contain the real zero with
\(\beta\geq1-c_0/\ell(0)\) when it exists, and let it be empty
otherwise.
It lies in \([3/4,1)\).
Section~\ref{sec:ray-real-gap} gives
\(1-\beta\geq(c_*/18000)D^{-46}\), with an absolute \(c_*>0\).

\subsection{Removing the exceptional pole}

Uniformly for
\[
       1-\frac{c_0}{2\ell(\Im s)}\leq\Re s\leq3 ,
\]
the holomorphic continuation of
\[
     A_\chi(s)=-\frac{F'_\chi(s)}{F_\chi(s)}
                           +\sum_{\beta\in Z_\chi}\frac1{s-\beta}
\]
satisfies
\begin{equation}
 \label{eq:ray-corrected-log}
                         |A_\chi(s)|\leq C_2\ell(\Im s)^2 ,
\end{equation}
for an absolute \(C_2\), depending on the fixed \(c_0\).
The apparent singularity at the simple exceptional zero is removable.
Here are the distance estimates proving the bound.
When \(\Re s\leq7/4\), use \eqref{eq:ray-local-log}.
A local zero has \(|\Im\rho-\Im s|\leq3/5\), and
\[
       \ell(\Im\rho)\leq\frac43\ell(\Im s).
\]
Indeed
\(|\Im\rho|+2\leq(13/10)(|\Im s|+2)\), and
\(\log(13/10)\leq\ell(\Im s)/3\) since \(D\geq4\).
Every local zero other than the exceptional one either has real
part below \(3/4\), or satisfies
\[
        \Re\rho<1-\frac{3c_0}{4\ell(\Im s)}.
\]
Its distance from \(s\) is therefore at least
\(c_0/(4\ell(\Im s))\).
The local zero count is \(O(\ell(\Im s))\), so their total
logarithmic derivative contribution is \(O(\ell(\Im s)^2)\).
If the exceptional zero lies outside the local disk, its distance
from \(s\) is at least \(1/10\).
For \(\Re s\geq7/4\), absolute convergence of the Euler series
bounds \({\cal L}_\chi(s)\) by a constant, and
\(|s-1|^{-1}\), \(|s-\beta|^{-1}\leq4/3\).
This proves \eqref{eq:ray-corrected-log} everywhere stated.

\subsection{Perron inversion and the contour shift}

We recall the precise truncation estimate needed here, for arbitrary
coefficients with \(|b(n)|\leq2\Lambda(n)\).
For \(u\geq2\), \(2\leq T\leq u\), and \(c=1+1/\log u\),
\[
 \left|\sum_{n\leq u}b(n)
     -\frac1{2\pi i}\int_{c-iT}^{c+iT}
            \left(\sum_nb(n)n^{-s}\right)\frac{u^s}{s}\,ds\right|
                    \leq C_3\{u(\log u)^2/T+\log u\}.
\]
For clarity, the integer nearest the cutoff requires the last term.
The truncated step kernel
\((2\pi i)^{-1}\int_{c-iT}^{c+iT}y^s\,ds/s\) differs from
\({\bf1}_{y>1}\), for \(y\ne1\), by at most
\[
                  C y^c\min(1,(T|\log y|)^{-1}).
\]
The reciprocal bound follows by integration by parts in the
oscillatory tail. The bound by a constant follows by writing its
real integral in terms of \(\sin(t\log y)\), splitting at
\(t=|\log y|^{-1}\), and applying the Dirichlet bound thereafter.
At \(y=1\) the kernel is \(\pi^{-1}\arctan(T/c)\).
Summing the errors for \(n\leq u/2\) and \(n\geq2u\) uses
\(\sum(\log n)n^{-c}=O((c-1)^{-2})\).
In \([u/2,2u]\), remove the at most two integers nearest \(u\),
whose contribution is \(O(\log u)\).
The remaining errors are bounded by a harmonic sum times
\(C(u/T)\log u\). This gives the stated uniform estimate,
including integral and nonintegral \(u\).

Set \(L=\sqrt{\log u}\) and suppose \(\log D\leq L\).
Take \(a=c_0/12\), \(T=e^{aL}\), and
\[
                  \sigma=1-\frac{c_0}{2\log(D(T+2))}.
\]
For all sufficiently large \(u\), \(2\leq T\leq u\), and
\(\log(D(T+2))\leq3L\).
The corrected function \(A_\chi\) is holomorphic on and near the
rectangle from \(\sigma-iT\) to \(c+iT\).
Shift its Perron integral to the left boundary. Equation
\eqref{eq:ray-corrected-log} bounds the horizontal contributions
by \(C uL^2/T\), and the vertical contribution by
\[
                C u^\sigma L^2\log(T+2)
                         \leq C uL^3e^{-c_0L/6}.
\]
The integral of \(1/|\sigma+it|\) is \(O(\log(T+2))\);
also \(u^c=eu\), and \(\sigma>3/4\).
All these estimates are uniform in the character and modulus.

Finally, for \(3/4\leq b\leq1\) and \(c>b\), elementary inversion
gives
\[
 \frac1{2\pi i}\int_{c-iT}^{c+iT}
                         \frac{u^s}{s(s-b)}\,ds
               =\frac{u^b}{b}+O(1+u/T),
\]
uniformly in \(b\).
The full line integral equals \((u^b-1)/b\), by the residues at
\(b\) and zero. Its two tails are \(O(u^c/T)\), because
\(|s(s-b)|\geq t^2\).
Apply this once with \(b=1\) for the principal pole and, when
present, once with \(b=\beta\) for the exceptional pole.
Since
\({\cal L}_\chi=A_\chi+{\bf1}_{\chi=1}/(s-1)
-\sum_{\beta\in Z_\chi}1/(s-\beta)\),
Perron inversion and the rectangle estimate imply
\[
 \left|\sum_{n\leq u}b_\chi(n)-{\bf1}_{\chi=1}u
               +\sum_{\beta\in Z_\chi}\frac{u^\beta}{\beta}\right|
                         \leq C_4u e^{-\eta\sqrt{\log u}}
\]
for some absolute \(C_4,\eta>0\).
For example, after increasing \(C_4\), any fixed
\(0<\eta<a\) absorbs the powers of \(L\).
The bounded range where \(T<2\) is covered by increasing the same
constant, using \(|b_\chi(n)|\leq2\log n\).
This proves \eqref{eq:level-psi-input} with its required uniformity
and unique possible exceptional zero.

Together with the preceding two sections this establishes
\eqref{eq:level-short-norm}. Thus the analytic norm-prime input to
the simultaneous auxiliary-prime construction has now been derived.

\paragraph{Source map.}
The local disk formula is
\srcid{RayClassLFunctionLogDerivLocalFormula}.
The two zero arguments are \srcid{RayClassLFunctionZeroFreeRegion}
and \srcid{RayClassLFunctionUniqueSimpleExceptionalZero}.
The corrected derivative estimate and Perron shift are
\srcid{LogDerivSubExceptionalLe},
\srcid{TruncatedPerronOfLeTwoVonMangoldt}, and
\srcid{RayClassFamilyPsiWithExceptionalZero}.
\footnote{\href{https://github.com/Deicyde/riemann-hypothesis/blob/40e56bab026552e089d7afcfeff08d2defb9c005/src/Development/AbelianVarietyTheory/Proofs/GlobalInterpolation/RayClassLFunctionZeroFreeRegion.lean}{The zero-free argument};
\href{https://github.com/Deicyde/riemann-hypothesis/blob/40e56bab026552e089d7afcfeff08d2defb9c005/src/Development/AbelianVarietyTheory/Proofs/GlobalInterpolation/RayClassLFunctionUniqueSimpleExceptionalZero.lean}{uniqueness and simplicity};
\href{https://github.com/Deicyde/riemann-hypothesis/blob/40e56bab026552e089d7afcfeff08d2defb9c005/src/Development/AbelianVarietyTheory/Proofs/GlobalInterpolation/RayClassFamilyPsiWithExceptionalZero.lean}{the explicit formula}.}

\section{Choosing disjoint auxiliary primes on the common window}
\label{sec:level-prime-selection}

We give the elementary part of the prime selection used by
Section~\ref{sec:host-field}, retaining its modulus \(N\), prime \(d\),
rank \(r\), and scale \(x\).
For \(g\geq2\) and \(r\geq1\) define
\[
 {\cal B}_r(g)=\{g\}\cup
 \bigcup_{j=1}^{r}
 \{\,\ell\text{ prime}:g^j<\ell^{4v_\ell(g^j-1)}\,\}.
\]
These sets are the \srcid{BB} of the source. We first prove the
count which makes two such sets disjoint.

\subsection{Counting exceptional divisors}

For fixed \(g,j\), at most three primes occur in the \(j\)-th set.
Every such prime divides the positive integer \(n=g^j-1\).
If four distinct primes \(\ell_i\) occurred, their powers
\(\ell_i^{v_{\ell_i}(n)}\) would be pairwise coprime divisors of \(n\),
so
\[
 (g^j)^4<
 \prod_{i=1}^{4}\ell_i^{4v_{\ell_i}(n)}
 =\left(\prod_{i=1}^{4}\ell_i^{v_{\ell_i}(n)}\right)^4
 \leq n^4<(g^j)^4 .
\]
This is impossible. Each set is finite because all its elements
divide \(n\). Consequently
\begin{equation}
 \label{eq:level-bb-card}
       |{\cal B}_r(g)|\leq3r+1 .
\end{equation}

We also need the bound
\begin{equation}
 \label{eq:level-roots}
 \#\{a\bmod\ell^w:a^j=1\}\leq2j
       \qquad(w,j\geq1).
\end{equation}
Such residues are units. For odd \(\ell\) the unit group modulo
\(\ell^w\) is cyclic. Here is the usual lifting justification.
The multiplicative group of the residue field is cyclic: if \(m\)
is its exponent, the polynomial \(X^m-1\) has at most \(m\) roots,
while all its elements are roots. In a finite abelian group the
exponent is the order of an element, by taking commuting elements
with maximal order in each primary component. Thus \(m\) equals
the group order. Choose a primitive residue \(a\bmod\ell\).
One of the \(\ell\) lifts \(a+k\ell\bmod\ell^2\) satisfies
\[
          v_\ell((a+k\ell)^{\ell-1}-1)=1 ,
\]
because modulo \(\ell^2\) the coefficient of \(k\) is
\((\ell-1)a^{\ell-2}\ell\), which is nonzero after division by
\(\ell\). If \(z=1+\ell^b u\), \(\ell\nmid u\), \(b\geq1\),
the binomial expansion for odd \(\ell\) gives
\(v_\ell(z^\ell-1)=b+1\). Induction shows that this lift has order
\((\ell-1)\ell^{w-1}\) modulo \(\ell^w\), proving cyclicity.
In a cyclic group of order \(m\), the equation \(a^j=1\) has
\(\gcd(j,m)\leq j\) solutions.

For \(\ell=2\), the moduli \(2\) and \(4\) satisfy
\eqref{eq:level-roots} directly. When \(w\geq3\), repeated squaring
gives
\(v_2(5^{2^a}-1)=a+2\).
Indeed squaring an integer \(1+2^b u\), with \(b\geq2\) and \(u\)
odd, increases this valuation by one. Thus \(5\) generates the
\(2^{w-2}\) units which are \(1\bmod4\). Every odd residue is
uniquely \(5^a\) or \(-5^a\). Each of the two cosets contains at
most \(j\) solutions of \(a^j=1\), proving
\eqref{eq:level-roots} in all cases.

Fix now \(Y\geq2\), a prime \(\ell\), and \(j\geq1\).
Let \(S\) be any set of integers \(g\in[Y,2Y]\) satisfying
\(g^j<\ell^{4v_\ell(g^j-1)}\).
If \(S\) is nonempty, put
\[
            w=\min_{g\in S}v_\ell(g^j-1),\qquad Q=\ell^w .
\]
The defining inequality implies \(w\geq1\).
All \(g\in S\) satisfy \(g^j=1\bmod Q\).
For an element \(g_0\) attaining the minimum,
\[
                 Y\leq g_0\leq g_0^j<Q^4 ,
\]
so \(Q>Y^{1/4}\). In an interval of length \(Y\), each residue
class modulo \(Q\) contains at most \(Y/Q+1\) integers.
Equation~\eqref{eq:level-roots} therefore gives
\begin{equation}
 \label{eq:level-exception-fixed}
                 |S|\leq2j(Y^{3/4}+1).
\end{equation}
This is a bound on integers, and hence also on primes.

For any fixed integer \(\ell\), the number of \(g\in[Y,2Y]\) with
\(\ell\in{\cal B}_r(g)\) is at most
\[
 1+\sum_{j=1}^{r}2j(Y^{3/4}+1)
                  \leq1+2r^2(Y^{3/4}+1).
\]
The first term accounts for \(g=\ell\).
Combining this with \eqref{eq:level-bb-card}, the number of
\(g\in[Y,2Y]\) for which
\({\cal B}_r(g)\cap{\cal B}_r(g_1)\ne\varnothing\) is at most
\begin{equation}
 \label{eq:level-excluded}
       (3r+1)\{1+2r^2(Y^{3/4}+1)\}
                       \leq20r^3Y^{3/4}.
\end{equation}
For the last inequality use \(3r+1\leq4r\) and
\(1+2r^2(Y^{3/4}+1)\leq5r^2Y^{3/4}\).

It follows that a set \(G\subset[Y,2Y]\) with
\[
                  |G|\geq Y^{9/10},
       \qquad 20r^3<Y^{3/20}
\]
contains distinct elements \(g_1,g_2\) with disjoint exceptional
sets. Choose any \(g_1\in G\). Formula~\eqref{eq:level-excluded}
excludes strictly fewer than \(|G|\) choices for \(g_2\).
It excludes \(g_1\) itself since \(g_1\in{\cal B}_r(g_1)\), so
the two choices are distinct.

\subsection{The analytic prime input and its scale}

For \(d\equiv3\bmod4\) prime, define
\[
 \begin{split}
 \operatorname{NormOne}(d,N,p)\quad\Longleftrightarrow\quad&
 \exists A,B\in\mathbb Z,\quad A^2+dB^2=4p,\\
 &N\mid B,\qquad 2N\mid A-2-B .
 \end{split}
\]
The short norm-prime estimate used here asserts that there are
absolute constants \(A_0,c,c'>0\) such that, for
\(d\geq100\), \(N\geq3\), \(X\geq1\), and
\[
 \log(4dN^2)\leq c\sqrt{\log X},\qquad
 e^{-c\sqrt{\log X}}\leq h\leq1 ,
\]
the number of primes \(p\in(X,X(1+h)]\) satisfying
\(\operatorname{NormOne}(d,N,p)\) is at least
\begin{equation}
 \label{eq:level-short-norm}
           \frac{c'hX}{(4dN^2)^{A_0}\log X}.
\end{equation}
We will use \eqref{eq:level-short-norm} as a stated input until its
Hecke-family construction and exceptional-zero estimates are
translated. The implication from these family statements is proved below.
Its quantitative weakness is sufficient for all
the following uniform comparisons.
The exact statement is
\srcid{NormPrimesImagQuadShortIntervalLowerBoundStatement};
its subsequent deduction is given in
\srcid{NormPrimesImagQuadShortIntervalLowerBound}.

Suppose
\[
 d\mid N,\qquad
 \log N\leq3x^{\epsilon/5}(\log x)^2,\qquad \epsilon>0 .
\]
For all sufficiently large \(x\), uniformly in \(N,d\) satisfying
these inequalities and all \(Y\) with \(\log Y\geq8x^\epsilon\),
the bound \eqref{eq:level-short-norm} with \(h=1\) gives at least
\(Y^{9/10}\) primes in \((Y,2Y]\).
We spell out uniformity in \(Y\).
Put \(t=x^{\epsilon/5}\), \(L=\log Y\), and \(D=4dN^2\).
For large \(x\),
\((\log x)^2\leq t\), and, since \(d\leq N\),
\[
        \log D\leq\log4+3\log N\leq10t^2,
                  \qquad L\geq8t^5 .
\]
It follows that
\(\log D/\sqrt L\leq(10/\sqrt8)t^{-1/2}\), so the first
hypothesis of \eqref{eq:level-short-norm} holds uniformly.
The remaining size comparison is
\[
 \log\frac{c'Y}{D^{A_0}\log Y}
             -\frac9{10}\log Y
       =\frac L{10}-A_0\log D-\log L+\log c'\geq0 .
\]
The function \(L/10-\log L\) is increasing for \(L\geq10\).
Its minimum over \(L\geq8t^5\) is therefore attained at \(8t^5\),
where
\[
 \frac45t^5-10A_0t^2-\log(8t^5)+\log c'\longrightarrow+\infty .
\]
Thus a threshold depending only on \(\epsilon\) and the absolute
constants works for all allowed \(N,d,Y\).

Take now
\[
             b_x=e^{x^\epsilon},\qquad Y=(k b_x)^8,
             \qquad k,r\geq1,\quad r\leq x^5 .
\]
Then \(\log Y=8(\log k+x^\epsilon)\geq8x^\epsilon\), and
\[
       20r^3\leq20x^{15}<e^{x^\epsilon}
                       \leq(kb_x)^{6/5}=Y^{3/20}
\]
for sufficiently large \(x\). The preceding two subsections supply
two primes \(g_i\in[Y,2Y]\), each satisfying
\(\operatorname{NormOne}(d,N,g_i)\), with
\[
                g_1\ne g_2,\qquad
                  {\cal B}_r(g_1)\cap{\cal B}_r(g_2)=\varnothing .
\]
They also split in the cyclotomic host.
Indeed write \(B=Nb\) and \(A=2+B+2Na\). Substituting in
\(A^2+dB^2=4g_i\) and reducing modulo \(4N\) gives
\(4g_i\equiv4\bmod4N\), hence \(g_i\equiv1\bmod N\).

\subsection{The lower bound from a characterwise explicit formula}

We make the analytic content of \eqref{eq:level-short-norm} more
precise. The following implication accounts for the exceptional real
zero and the removal of prime powers. It can be checked with three
specified inputs about the ray-class family.

Let \(G\) be a finite abelian group of order \(H\), and let \(D\geq4\)
satisfy \(H\leq D\). Suppose nonnegative weights \(\lambda_g(n)\)
satisfy
\[
             \sum_{g\in G}\lambda_g(n)\leq2\Lambda(n).
\]
Write
\(\Psi_\chi(u)=\sum_{n\leq u}\sum_g\chi(g)\lambda_g(n)\).
Assume constants \(a,\eta,C,\gamma>0\), with \(\gamma\leq1\),
have the following uniform properties.
Across all characters there is at most one exceptional pair
\((\chi,\beta)\), where \(3/4\leq\beta<1\) and
\(1-\beta\geq\gamma D^{-a}\).
For \(u\geq2\), \(\log D\leq\sqrt{\log u}\), assume
\begin{equation}
 \label{eq:level-psi-input}
 \left|\Psi_\chi(u)-{\bf1}_{\chi=1}u+
               \sum_{\beta\in Z_\chi}\frac{u^\beta}{\beta}\right|
             \leq Cu e^{-\eta\sqrt{\log u}}.
\end{equation}
Here each \(Z_\chi\) is empty except for the possible exceptional
pair. Finally assume that
\(\lambda_0(p)\ne0\), for \(p\) prime, implies
\(\operatorname{NormOne}(d,N,p)\).

Put \(\delta=\gamma D^{-a}\).
Character orthogonality gives
\[
              \sum_\chi\Psi_\chi(u)
                        =H\sum_{n\leq u}\lambda_0(n).
\]
If \(V=X(1+h)\), \(0<h\leq1\), subtraction of
\eqref{eq:level-psi-input} at \(X\) and \(V\) therefore gives
\begin{equation}
 \label{eq:level-mass}
 H\sum_{X<n\leq V}\lambda_0(n)
 \geq hX-\sum_{\chi,\beta\in Z_\chi}
          \frac{V^\beta-X^\beta}{\beta}
                  -3DCX e^{-\eta\sqrt{\log X}}.
\end{equation}
The factor \(3\) bounds the two endpoint errors since \(V\leq2X\).

If an exceptional zero exists and \(\log X\geq1\), then
\[
 \frac{V^\beta-X^\beta}{\beta}
       \leq hX^\beta
       =hX e^{-(1-\beta)\log X}
       \leq hX e^{-\delta}
       \leq hX(1-\delta/2).
\]
The first inequality follows from the concavity of
\(t^\beta\) at \(1\). For the last one, use
\(e^{-\delta}\leq1/(1+\delta)\leq1-\delta/2\), valid for
\(0\leq\delta\leq1\). If there is no exceptional zero, the sum
in \eqref{eq:level-mass} is zero, and the same bound on that sum
still applies.

Choose \(c>0\) small enough that \(c\leq1\) and
\((a+2)c<\eta\). By decreasing \(c\) further, all
\(\log D\leq cL\), \(L=\sqrt{\log X}\), imply
\[
       3DCe^{-\eta L}
                     \leq\frac14 e^{-cL}\gamma D^{-a}.
\]
To justify the uniform choice, the ratio of the left side to the
right side is at most
\[
       \frac{12C}{\gamma}
             e^{-\{\eta-(a+2)c\}L}.
\]
Also \(L\geq\log4/c\). Take \(c\) sufficiently small that the last
display is at most one at this lower endpoint. It is decreasing
in \(L\). Hence, when \(h\geq e^{-cL}\),
\eqref{eq:level-mass} implies
\[
          \sum_{X<n\leq V}\lambda_0(n)
                       \geq\frac{\gamma hX}{4D^{a+1}}.
\]

Let \(P\) count the norm primes in \((X,V]\).
The prime terms in the last sum are at most \(2\log V\cdot P\).
The other terms are at most
\(2\{\psi(V)-\theta(V)\}\).
An elementary bound sufficient here is
\[
 \psi(V)-\theta(V)
       =\sum_{j=2}^{\lfloor\log V/\log2\rfloor}
                         \theta(V^{1/j})
       \leq\frac{\sqrt V(\log V)^2}{\log2}.
\]
Indeed there are at most \(\log V/\log2\) terms, and each is at
most \(\sqrt V\log V\), by counting all integers up to
\(V^{1/j}\). Since \(V\leq2X\), \(\log V\leq2\log X\) for
\(X\geq2\), the unwanted contribution is at most
\(C_1\sqrt X(\log X)^2\) for an absolute \(C_1\).
Its ratio to \(\gamma hX/(8D^{a+1})\) is bounded by
\[
 \frac{8C_1}{\gamma}L^4
        \exp\{-L^2/2+(a+2)cL\}.
\]
Decreasing \(c\) again makes this at most one for all
\(L\geq\log4/c\), because the negative quadratic term dominates
the linear term and \(\log L\). Thus at least
\(\gamma hX/(8D^{a+1})\) of the mass comes from primes.
Dividing by \(2\log V\leq4\log X\) gives
\[
             P\geq
               \frac{\gamma hX}{32D^{a+1}\log X}.
\]
This is \eqref{eq:level-short-norm} with
\(A_0=a+1\), \(c'=\gamma/32\), and \(D=4dN^2\).

The analytic inputs are now explicit.
The ray-class construction must supply the nonnegative weights and
their norm-prime support. The family explicit formula must give
\eqref{eq:level-psi-input} with a unique possible exceptional zero.
Section~\ref{sec:ray-real-gap} proves the real value and zero-gap estimates
from the counting, theta and positive-convolution identities of the family. In the release these are
\srcid{ExistsLFamRayStatement},
\srcid{RayClassFamilyPsiWithExceptionalZeroStatement}, and
\srcid{ExistsRayClassLFunctionRealZeroFarFromOneStatement}.
The last is itself obtained by combining a polynomial lower bound
for \(|F_\chi(1)|\) with
\[
             |F'_\chi(z)|\leq18000D
                           \qquad(|z-1|\leq1/4).
\]
If \(F_\chi(\beta)=0\), integration on \([\beta,1]\) gives
\(|F_\chi(1)|\leq18000D(1-\beta)\).
A bound \(c_0D^{-a_0}\leq|F_\chi(1)|\) therefore supplies
\(1-\beta\geq(c_0/18000)D^{-(a_0+1)}\).
This identifies the remaining analytic statements without suppressing
their dependence on \(D\).

\subsection{Three distinct small splitting primes}

For the second input, Linnik's theorem supplies a constant \(A\geq1\)
such that for every
integer \(M\geq2\) there is a prime \(p\equiv1\bmod M\) with
\(p\leq M^A\). Its formal name is \srcid{SmallOne}. For example, Xylouris's
Theorem~1.1 gives \(p\leq C M^{5.18}\) uniformly for each reduced
residue class. Replacing the exponent by
\(A=5.18+\max(0,\log C/\log2)\) absorbs the absolute constant for
all \(M\geq2\). This supplies the precise version used here.
\footnote{T.~Xylouris, \emph{On the least prime in an arithmetic
progression and estimates for the zeros of Dirichlet \(L\)-functions},
Acta Arithmetica 150 (2011), Theorem~1.1,
\href{https://doi.org/10.4064/aa150-1-4}{doi:10.4064/aa150-1-4}.}
For \(N\geq30\), first apply it to \(15N\), then to \(15Np\),
and finally to \(Npp'\). The resulting primes satisfy
\[
 \begin{gathered}
 p,p'\equiv1\bmod15N,\qquad q_0\equiv1\bmod N,\\
 p'\equiv1\bmod p,\qquad q_0\equiv1\bmod pp' .
 \end{gathered}
\]
They are pairwise distinct, as a prime cannot be both zero and one
modulo itself. Since \(15N\leq N^2\),
\[
 \begin{split}
 p&\leq N^{2A},\\
 p'&\leq (15Np)^A\leq N^{(2A+2)A},\\
 q_0&\leq (Npp')^A
             \leq N^{A\{1+2A+(2A+2)A\}}.
 \end{split}
\]
The common exponent
\[
                    C_8=A\{1+2A+(2A+2)A\}
\]
dominates all three. With \(N\) on the window above,
\[
 \log N^{C_8}
     \leq3C_8x^{\epsilon/5}(\log x)^2<x^\epsilon
\]
for all sufficiently large \(x\). Thus \(p,p',q_0<b_x\), whereas
\(g_i\geq(kb_x)^8\geq b_x\).
In particular \(p,p'\) differ from both \(g_i\), as required.

Every threshold used here depends on the fixed \(\epsilon\) and the
absolute analytic constants. They can be combined by taking their
maximum. Consequently the same choices of \(x,N,d,k,r\) are retained
throughout the construction. The remaining analytic dependencies are the construction of the ray-class
family with the stated identities and its characterwise explicit formula. The small splitting primes follow from the
stated classical least-prime theorem.

\paragraph{Source map.}
The definitions of the exceptional sets are in
\srcid{Formalism.CyclotomicNorms.StatementI}. The counting proofs
are \srcid{ExceptionalPrimesFiniteCardLe},
\srcid{CardPowEqOneLeOfIsCyclicUnits}, and
\srcid{DivLeRpowThreeQuarters}. Their use on the scale is
\srcid{ShortNorms.normPrimesN} and \srcid{ShortNorms.typePrimesN}.
The three-prime construction and its simultaneous application are
\srcid{ShortNorms.ExistsThreePrimesOneModLeRpow} and
\srcid{ShortNorms.levelData}.
\footnote{\href{https://github.com/Deicyde/riemann-hypothesis/tree/40e56bab026552e089d7afcfeff08d2defb9c005/src/Development/AutomorphicArithmetic/ShortNorms}{The prime-selection modules};
\href{https://github.com/Deicyde/riemann-hypothesis/blob/40e56bab026552e089d7afcfeff08d2defb9c005/src/Development/AbelianVarietyTheory/Proofs/FoundationProofs/DivLeRpowThreeQuarters.lean}{the fixed-prime residue count}.}

\section{Counting the exceptional characters and primes}
\label{sec:window-exceptions}

We prove the zero-density estimate and give the progression mean-square
reduction used in Section~\ref{sec:rh-common-windows}. The density
estimate counts primitive characters for which at least one zero
occurs. Each character contributes one to the count irrespective of
the multiplicity or number of such zeros.

\subsection{A detector for a single zero}

Put \(\sigma_0=399/400\). For \(Q\ge3\) consider primitive
characters \(\chi\) of conductor \(q\le Q\) with a zero \(\rho\)
satisfying \(\Re\rho>\sigma_0\), \(|\Im\rho|\le Q\).
There is an absolute constant \(C\) such that their number is at
most \(CQ^{1/4}\).

It suffices to treat large \(Q\), since a bounded conductor range
contains only finitely many characters. The character of conductor
one contributes at most one. For each of the others select one zero
in the box. These characters are nonprincipal, and the classical
nonvanishing theorem on \(\Re s\ge1\) gives \(\Re\rho<1\).
For \(Q\ge20\) set
\[
 X=\lceil Q\rceil,\quad Y=X^4,\quad
 J=8(\lfloor\log_2 X\rfloor+1),\quad N=2^J .
\]
Then \(X^8\le N\le256X^8\). Write \(P_\chi(s)=
\sum_{b\le N}\chi(b)b^{-s}\) and \(M_\chi(s)=
\sum_{a\le Y}\mu(a)\chi(a)a^{-s}\), with all sums starting at one.
Character orthogonality over a complete period bounds each prefix
sum of \(\chi\) by \(q\). In the half-plane \(\sigma=\Re s>0\)
the resulting Dirichlet series converges to \(L(s,\chi)\) by partial
summation and analytic continuation from \(\Re s>1\). More explicitly,
the tail is bounded by \(2(q+1)(1+|s|/\sigma)N^{-\sigma}\):
apply partial summation to \(t^{-s}\), use the prefix bound at the
endpoint and integrate \(|s|t^{-\sigma-1}\) to infinity.
At the selected zero this bounds \(|P_\chi(\rho)|\).
Now \(q+1\le2X\), \(1+|\rho|/\Re\rho\le3X\), and
\(N^{\Re\rho}\ge N^{15/16}\ge X^7\sqrt X\ge4X^7\).
As \(|M_\chi(\rho)|\le Y\), we have
\[
     |P_\chi(\rho)M_\chi(\rho)|
       \le\frac{2(2X)(3X)}{4X^7}X^4=\frac{3}{X}\le\tfrac12.
\]
The Möbius convolution makes this small product a detector. Define
\(a_n=\sum_{a\le Y,\,b\le N,\,ab=n}\mu(a)\).
Then \(a_1=1\), and for \(1<n\le Y\) the sum runs over every
divisor of \(n\), so \(a_n=0\). Consequently
\[
 P_\chi(\rho)M_\chi(\rho)=
     1+\sum_{Y<n\le YN}a_n\chi(n)n^{-\rho}.
\]
The last sum has magnitude at least \(1/2\). The \(J\) consecutive
blocks \((U,2U]\), \(U=Y2^j,\ 0\le j<J\), partition its range.
Select a block for which
\[
 \left|\sum_{U<n\le2U}a_n\chi(n)n^{-\rho}\right|
              \ge\frac{1}{2J}.
\]

\subsection{How many characters can use one block?}

The elementary divisor bound \(d(n)\le C_0 n^{1/400}\) applies
uniformly in \(Y,N\), since \(|a_n|\le d(n)\).
To obtain it, for sufficiently large primes \(p\) use
\(e+1\le2^e\le p^{e/400}\) for every \(e\ge0\).
For each of the finitely many smaller primes the supremum of
\((e+1)p^{-e/400}\) is finite. Their product gives \(C_0\).
Hence \(b_n=a_nn^{-\sigma_0}\) satisfies, for \(U\ge1\),
\[
 \sum_{U<n\le2U}|b_n|^2
   \le C_0^2 U^{1-2(199/200)}
   \le C_0^2 U^{-197/200}.
\]
Put \(C_1=C_0^2\). For the set \(T\) of characters assigned to this
block define \(v_\chi(n)=\chi(n)n^{-(\rho_\chi-\sigma_0)}\).
Each Gram matrix diagonal is at most \(U\). For distinct primitive
\(\chi,\chi'\) the character \(\chi\overline{\chi'}\) lifted to
modulus \(qq'\) is nonprincipal. Otherwise their inflations agree
on units and have the same primitive inducing character, forcing
the original characters, including conductors, to agree.
Its sum over a full period vanishes and every prefix has absolute
value at most \(qq'\). For \(w=\rho_\chi-\sigma_0+
\overline{\rho_{\chi'}-\sigma_0}\), partial summation gives
\[
 \left|\sum_{U<n\le2U}\chi(n)\overline{\chi'(n)}n^{-w}\right|
       \le 2qq'(1+|w|)\le8Q^3.
\]
Indeed \(\Re w>0\), \(|\Im w|\le2Q\), \(\Re w<2(1-\sigma_0)\);
on the interval the total variation of \(t^{-w}\) is at most
\(|w|\int_U^{2U}dt/t\le|w|\). Prefix differences cost at most
\(2qq'\). The sum of the absolute values in each Gram row is
therefore at most \(\Delta=U+8|T|Q^3\).

For clarity, put \(s_\chi=\sum_n b_n v_\chi(n)\),
\(Z=\sum_\chi|s_\chi|^2\), and \(w_n=\sum_\chi\overline{s_\chi}
v_\chi(n)\). Then \(Z=\sum_n b_n w_n\). By Cauchy--Schwarz,
\(Z^2\le\|b\|_2^2\sum_n|w_n|^2\). Expanding the second factor and
using \(2|s_\chi||s_{\chi'}|\le |s_\chi|^2+|s_{\chi'}|^2\)
bounds it by \(\Delta Z\), since the absolute Gram matrix is symmetric.
Division by \(Z\) when \(Z>0\), with the zero case immediate, gives the claimed bound.
The detector now gives
\[
 \frac{|T|}{4J^2}
 \le C_1 U^{-197/200}\{U+8|T|Q^3\}.
\]
Since \(U\ge X^4\ge Q^4\), the coefficient of \(|T|\) on the right
is at most \(8C_1Q^{-47/50}\), which is at most \(1/(8J^2)\)
for all sufficiently large \(Q\). Moving that term to the left yields
\[
                 |T|\le8C_1 J^2 U^{3/200}.
\]
There are \(J\) possible blocks, each with
\(U\le YN\le256X^{12}\le2^{20}Q^{12}\). For all sufficiently
large \(Q\), \(J\le32\log Q\) and \(J^3\le Q^{7/100}\).
Including the possible character of conductor one, the count is
therefore bounded by
\[
 1+8C_1 (2^{20})^{3/200}
           Q^{7/100+12(3/200)}\ll Q^{1/4}.
\]
The classical inputs in this argument are orthogonality of
Dirichlet characters, uniqueness of primitive induction, analytic
continuation and nonvanishing on the line \(\Re s=1\) for
nonprincipal Dirichlet characters. Each bound on the detector,
the block Gram matrix and the final count has been derived above.

In a prime stage with upper endpoint \(V\) and base conductor
at most \(F\le\sqrt V\), put \(Q=FV\). A candidate \(p>\sqrt V\)
failing the zero-free condition yields a primitive offending twist
of conductor \(fp\) for \(1\le f\le F\). If the twists for \(p,p'\)
coincide, then \(fp=f'p'\) and \(p>F\ge f'\) give \(p=p'\).
There are at most \(O(Q^{1/4})=O(V^{3/8})\) failures. Against the
source's separate prime-progression lower bound \(V^{9/10}\), this
is a fraction \(O(V^{-21/40})\).

\subsection{The quadratic mean square on the actual progressions}

For this part write \(Y\) for the lower endpoint denoted \(U\) in
Section~\ref{sec:rh-common-windows}, and \(V\) for the upper cutoff
on candidate primes. Choose \(0<\epsilon\le1\) and \(K\ge1\).
Set \(w=Y^{\epsilon/5}\), \(Z=(Y^K)^{400}\), \(N=\lfloor Z\rfloor\)
and \(L=\log Y\). The base progression has modulus \(m_0>0\)
divisible by \(8\), and a reduced class \(a_0\). Assume
\[
 \tfrac14 L^2\log m_0\le\log V,\qquad
 Y^{\epsilon/6}\le\log V\le Y^{\epsilon/5}L^2,\qquad
 p\mid m_0\ \Longrightarrow\ p\le w\ \text{or}\ p>Z
\]
for prime \(p\). Let \(\mathcal P=\{q:\sqrt V<q\le V,\,
q\text{ prime},\ q\equiv a_0\bmod m_0\}\).
For \(Y\) sufficiently large in terms of \(\epsilon,K\),
every \(q\in\mathcal P\) exceeds \(N^2\), since
\(1600K\log Y\le Y^{\epsilon/6}\le\log V\).
For \(Y\le y\le Y^K\) put
\[
 c(n)=\Lambda(n)\min(1,y/n),\qquad
 S_q=\sum_{1\le n\le N} c(n)\kappa(n)\left(\frac{n}{q}\right).
\]
We show that
\begin{equation}
 \sum_{q\in\mathcal P}|S_q|^2\le C y\log y\,|\mathcal P|
 \qquad (|\kappa(n)|\le1).
 \label{eq:window-progress-variance}
\end{equation}
The absolute constant \(C\) is independent of \(\epsilon,K\);
the lower threshold on \(Y\) may depend on them.

The input about primes in progressions is the following instance
of \texttt{twistK}. For all but at most one odd squarefree
\(s>1\) coprime to \(m_0\) and bounded by \(Y^{800K}\), one has
\begin{equation}
 \left|\sum_{q\in\mathcal P}\left(\frac{q}{s}\right)\right|
       \le\eta|\mathcal P|,\qquad \eta=Y^{-(K+1)} .
 \label{eq:window-jacobi}
\end{equation}
We give its classical inputs explicitly. The relative
prime number theorem in arithmetic progressions in the source
has main term \(\operatorname{li}(V)/\varphi(M)\), with a correction
\(-\chi_1(b)\operatorname{li}(V^{\beta_1})/\varphi(M)\)
when the real exceptional zero is present. Here
\(M=m_0s\), \(b\) is reduced, and the relative error is
\[
 O\!\left(e^{-c\log V/\log M}+e^{-\sqrt{c\log V}}\right)
 \quad\text{when}\quad C_0\log M\log(\log M+2)\le\log V .
\]
The zero-free region and the Landau--Page uniqueness bound are
applied across the moduli \(M\le H=m_0Y^{800K}\).
Summing the main term against the Jacobi character across the
classes \(b\equiv a_0\bmod m_0\) gives zero by orthogonality.
An exceptional term survives only if the exceptional character
agrees with the Jacobi character on units congruent to \(1\bmod m_0\).
Two such characters with \(1-\beta_1<c'/\log H\) induce the
same primitive character by Landau--Page. Hence the two Jacobi
characters coincide on those units. By the Chinese remainder
theorem a differing odd prime factor in their squarefree moduli
would allow a quadratic nonresidue there and ones at all other
factors and modulo \(m_0\), forcing opposite values.
There is thus at most one exceptional modulus \(s_*\).
When \(1-\beta_1\ge c'/\log H\) the surviving correction
relative to \(\operatorname{li}(V)\) is bounded by
\(4\log V\exp(-c'\log V/\log H)\). For fixed \(\epsilon,K\)
the hypotheses imply \(\log V/\log H\gg(\log Y)^2\) and
\(\log\log V=O(\log Y)\), while \(\sqrt{\log V}\) grows faster
than every fixed power of \(\log Y\). These errors are
\(o(Y^{-(K+1)})\). Removing primes up to \(\sqrt V\)
costs at most \(V^{1/2}\le V^{-2/5}|\mathcal P|\), by the
separate progression-count input \(|\mathcal P|\ge V^{9/10}\).
The relative-error estimate also applies to the unweighted count
of the classes, which is the denominator in \eqref{eq:window-jacobi}.
This is the role of the relative form of the prime number theorem.
The source proves the uniform comparison in \texttt{twistK}
and invokes the relative prime number theorem, the zero-free region
and the exceptional-zero bound there.

Expand the square of \(S_q\) and use \(|\kappa|\le1\).
For \(h=nn'\), define \(s(h)\) to be the product of primes
outside \(m_0\) having odd exponent in \(h\).
Quadratic reciprocity and its supplement for \(2\) give
\[
 \left|\sum_{q\in\mathcal P}\left(\frac{nn'}q\right)\right|
     =\left|\sum_{q\in\mathcal P}\left(\frac{q}{s(nn')}\right)\right|.
\]
Indeed every prime factor of \(nn'\) is less than \(q\);
square factors contribute one. The sign in reciprocity depends
only on \(q\bmod4\), the factor from \(2\) only on \(q\bmod8\),
and the factors from primes dividing \(m_0\) are fixed by the
progression. Their product is independent of \(q\).
The bound \(s(nn')\le N^2\le Y^{800K}\) permits
\eqref{eq:window-jacobi} at all pairs except those with
\(s(nn')=1\) or \(s(nn')=s_*\).

Write \(A_p=\sum_{p^j\le N,\,j\ge1}c(p^j)\),
\(B=\sup_p A_p\), and let \(R\) be the sum of \(c(p^j)\)
over powers for which \(j\) is even or \(p\mid m_0\).
Only prime powers have nonzero coefficient.
For \(s(nn')=1\), either both powers have the same base prime or
both have trivial squarefree part outside \(m_0\). Their total
weight is at most \(\sum_p A_p^2+R^2\).
For a specified \(s_*>1\), if it is prime, pairs with different
base primes have one factor with outside part \(s_*\) and the
other with trivial outside part, giving at most \(2BR\); pairs
with the same base prime have weight at most \(B^2\).
If \(s_*\) has two distinct prime factors the only possible
bases are those two, giving at most \(2B^2\).
Three or more prime factors give no pairs.
Use the sum \(2B^2+2BR\) in all cases. We have obtained
\begin{equation}
 \sum_{q\in\mathcal P}|S_q|^2\le|\mathcal P|
 \left\{\sum_p A_p^2+R^2+2B^2+2BR+
                      \eta\left(\sum_{n\le N}c(n)\right)^2\right\}.
 \label{eq:window-pair-split}
\end{equation}

Here are the size estimates that keep the constant absolute.
For \(y\ge3\), \(A_p\le3\log y\). If \(p\le y\) sum the terms
through \(\lfloor\log y/\log p\rfloor\) and then the decreasing
geometric tail, whose sum in units of \(\log p\) is at most \(2\).
If \(p\ge y\), use \(A_p\le 2y\log p/p\le2\log y\) by
monotonicity of \((\log t)/t\) for \(t\ge3\).
The prime diagonal satisfies
\(\sum_p A_p^2\le C_3 y\log y\). Indeed the primes up to
\(\sqrt y\) contribute at most \(9\sqrt y(\log y)^2\);
on \(\sqrt y<p\le y\) use
\(A_p^2\le18\log y\log p\) and
\(\sum_{p\le y}\log p\le y\log4\).
For \(p>y\) use \(A_p^2\le4y^2(\log p)^2/p^2\).
On \(t<p\le2t\), the sum of \((\log p)^2/p^2\) is at most
\(2(\log4)\log(2t)/t\), by the same Chebyshev estimate. Summing
the dyadic bounds \(t=2^j y\) gives \(O(\log y/y)\) and proves
the diagonal estimate. The Chebyshev bound follows, for example,
from the primorial bound \(\prod_{p\le n}p\le4^n\).

For \(R\), the total from even powers is at most
\(\sum_{a\le N}\Lambda(a)\min(1,y/a^2)\).
Use \(\sum_{a\le t}\Lambda(a)\ll t\), the elementary Chebyshev
bound, on \(a\le\sqrt y\) and on successive dyadic intervals above
\(\sqrt y\). The first sum is \(O(\sqrt y)\), and interval \(j\)
in the tail is \(O(2^{-j}\sqrt y)\).
All remaining base primes in \(R\) divide \(m_0\) and are at most
\(N\), hence at most \(w\). There are at most \(w\) of them,
each contributing at most \(3\log y\).
Thus \(R\le C_1(\sqrt y+w\log y)\).
Since \(w\le y^{1/5}\), the elementary inequality
\(\log y\le(10/3)y^{3/10}\) gives \(R\le5C_1\sqrt y\).
Finally, partial summation of the Chebyshev bound gives
\(\sum_{n\le N}c(n)\le y\sum_{n\le N}\Lambda(n)/n\le C_4 y\log Z\)
in the range at issue. For fixed \(K\) and all sufficiently large
\(Y\), \((\log Z)^2=(400K\log Y)^2\le Y\).
Then \(\eta y^2(\log Z)^2\le y\), since \(y\le Y^K\).
All five expressions in braces in \eqref{eq:window-pair-split}
are \(O(y\log y)\). More precisely with \(B\le C_2\log y\),
the source takes
\[
 C=C_3+25C_1^2+4C_2^2+20C_1C_2+C_4^2.
\]
This proves \eqref{eq:window-progress-variance}, including its
uniformity for arbitrary complex multipliers.

At the first selection take \(m_0=8,\ a_0=3,\ V=e^w\).
At the second take \(m_0=16q\prod_{2<p\le w}p\), the reduced
class of Section~\ref{sec:rh-common-windows}, and
\(V=e^{w(\log Y)^2}\). In this second stage \(q>\sqrt{e^w}>Z\)
for large \(Y\), so adjoining \(q\) to \(m_0\) preserves the
condition on prime divisors. The modulus bound follows from
\(\log q\le w\) and \(\sum_{p\le w}\log p\le w\log4\),
which give \(\log m_0\le4w\) for large \(Y\).
Thus the same \(K\) and lower endpoint \(Y\) apply in both stages.
The moving-point estimate and finite union over characters in
Section~\ref{sec:rh-common-windows} then give the required
exceptional fraction with \(\tau=\epsilon/2\) on all of
\([Y,Y^K]\).

\paragraph{Source correspondence.}
The density proof is
\texttt{CountZAll.zero\_density\_box}, with the exact detector in
\texttt{CountZDetect}, the coefficient square estimate in
\texttt{CountZMoll}, the interval character sums in
\texttt{CountZPair} and \texttt{CountZTwist}, the Gram bound in
\texttt{CountZBlock} and \texttt{CountZDual}, and the parameter
choice in \texttt{CountZNum} and \texttt{CountZFinal}. These are
the modules under
\srcid{Formalism.GL3MeanSquares.GenusMeanAsymptotics.JacobiCharacterSums}.
For the second estimate \srcid{QH.meanSqK} uses
\texttt{twistK}, \texttt{MeanSqAbs}, \texttt{SumMinLeDiv},
\texttt{SPsqLe}, \texttt{RRLe} and \texttt{SumCNLe} under
\srcid{Development.AutomorphicArithmetic.MeanSquareBounds}.
The relative prime number theorem is invoked as
\texttt{pnt\_progressions\_relative\_weak\_holds}, the zero-free
region as \texttt{zFRDirichlet}, the exceptional-zero estimate as
\texttt{excBound}, and the rational-prime count as
\texttt{count\_a}. The explicit formula that converts these
finite prime sums to centered sums of the twists is a separate
input in \texttt{BoxOfS}.

\section{A relative prime number theorem for the window moduli}
\label{sec:window-relative-pnt}

The denominator in Section~\ref{sec:window-exceptions} counts primes
in a single reduced residue class. An exceptional real zero can make
this count small. We prove the estimate with an error proportional to
the corrected main term, and then give the lower bound on the
denominator for the moduli used in the two selections of primes.

For \(z\ge2\) set
\(\ell(z)=\int_2^z dt/\log t\), and let \(\pi(X;M,b)\) count primes
\(p\le X\) congruent to \(b\bmod M\). Call \((\chi,\beta)\) eligible
at \(M\) if \(\chi\) is a nonprincipal real Dirichlet character modulo
\(M\), \(L(\beta,\chi)=0\), and
\[
       1-\frac{1}{8\log M}<\beta<1.
\]
There are absolute positive constants \(c,A,C\) with the following
property. For every positive multiple \(M\) of \(8\), select one
eligible pair if one exists, and select the empty option otherwise.
The selection can be made depending only on \(M\), in such a way that
for all integers \(b\) coprime to \(M\) and all real \(X\ge2\) with
\[
             \log X\ge C\log M\,\log(\log M+2)
\]
one has
\begin{align}
 B(X)&=\frac{\ell(X)-\chi(b)\ell(X^\beta)}{\varphi(M)},&
 E_c(X,M)&=\exp\!\left(-\frac{c\log X}{\log M}\right)
                   +\exp\!\left(-\sqrt{c\log X}\right),
 \nonumber\\
 |\pi(X;M,b)-B(X)|&\le A E_c(X,M) B(X).
 \label{eq:window-relative-pnt}
\end{align}
The character term is omitted for the empty option. The choice of
\(C\) below ensures \(X^\beta\ge2\) everywhere the estimate is used.

\paragraph{The weighted count.}
We use the unconditional progression theorem of J.~Thorner and
A.~Zaman. Their exceptional pair has \(\beta\ge1-1/(50\log M)\).
McCurley's zero-free region gives uniqueness in this interval,
simplicity, and a nonprincipal real character. With the exceptional
pair in their sense put \(\sigma=\chi(b)\) and
\(F(t)=t-\sigma t^\beta/\beta\). In its absence put \(F(t)=t\)
and \(\sigma=0\). Define
\(\vartheta(t;M,b)=\sum_{p\le t,\ p\equiv b\bmod M}\log p\).
When \(M\ge2\), \((b,M)=1\), \(t\ge3\), and
\(F(t)/\varphi(M)\ge t^{18/19}\), their Corollary~1.5 gives
\[
 \left|\vartheta(t;M,b)-\frac{F(t)}{\varphi(M)}\right|
 \ll \frac{F(t)}{\varphi(M)}
 \left\{\exp\!\left(-\frac{a\log t}{\log M}\right)
       +\exp\!\left(-a\frac{(\log t)^{3/5}}{(\log\log t)^{1/5}}\right)
 \right\}
\]
for an absolute \(a>0\). We use the weaker consequence
\(1-\beta\gg M^{-1}\) of the exceptional-zero estimate in their
equation~(1.2). Both results concern all characters at modulus
\(M\), which permits imprimitive characters. Equivalently, the
extra Euler factors of an induced character are nonzero at real
\(\beta>0\), since \(|\chi^*(p)p^{-\beta}|<1\).
See Thorner and Zaman, \emph{Refinements to the prime number theorem
for arithmetic progressions}, Math.~Z. 306 (2024), article~54,
\href{https://doi.org/10.1007/s00209-023-03414-3}{journal publication},
\href{https://arxiv.org/abs/2108.10878}{authors' text}.

Hereafter increase the absolute \(C\) so that
\(s=\sqrt X\ge M^{100}\) and \(s\) exceeds any required absolute
threshold. This is possible since \(M\ge8\).
For the sign \(\sigma=1\) write \(\delta=1-\beta\) and
\(w(t)=\min(1,\delta\log t)\). For \(\sigma=0,-1\) put \(w(t)=1\).
Uniformly on \(t\ge s\) we have \(F(t)\asymp t w(t)\).
For \(\sigma=1\), indeed, \(\beta>1/2\) yields
\(0\le-\log\beta\le2\delta\). Thus for \(\log t\ge4\),
\[
   1-e^{-\delta\log t/2}
      \le \frac{F(t)}{t}
      =1-\frac{e^{-\delta\log t}}{\beta}
      \le 1-e^{-\delta\log t}.
\]
Use \(1-e^{-v}\asymp\min(1,v)\) for \(v\ge0\).
For \(\sigma=-1\) the claimed comparison follows from
\(1\le F(t)/t\le3\). In all three cases \(w(t)\gg M^{-1}\).
Since \(\varphi(M)\le M\),
\[
           F(t)/\varphi(M)\gg t M^{-2}\ge t^{49/50}.
\]
This exceeds \(t^{18/19}\) once \(s\) exceeds an absolute constant.
The published hypothesis therefore holds at every \(t\in[s,X]\),
for the same class and the same exceptional option.

\paragraph{Removing the logarithmic weight.}
First form \(B\) with the authors' exceptional option.
For \(\sigma=1\),
\[
  \varphi(M) B(X)=\int_{X^\beta}^{X}\frac{dt}{\log t}
     \ge \frac{X(1-e^{-\delta\log X})}{\log X}.
\]
For \(\sigma=0,-1\), integration over \([X/2,X]\)
gives the corresponding lower bound \(X/(2\log X)\).
Together these inequalities yield
\begin{equation}
 B(X)\gg \frac{X w(X)}{\varphi(M)\log X}
          \gg\frac{X}{M^2\log X}.
 \label{eq:window-ap-main-lower}
\end{equation}
Write \(H(t)=\vartheta(t;M,b)-F(t)/\varphi(M)\).
Uniformly for \(s\le t\le X\), the weighted theorem gives
\[
 |H(t)|\ll \frac{t w(X)}{\varphi(M)} E_c(X,M)
\]
after choosing \(c>0\) sufficiently small. Indeed
\(\log t\ge\tfrac12\log X\), \(w(t)\le w(X)\), and
\(u^{3/5}/(\log u)^{1/5}\gg\sqrt u\) for large \(u\).
On this interval \(\varphi(M)B'(t)=F'(t)/\log t\), including the
factor \(1/\beta\) in \(F\). Partial summation with lower endpoint
\(s\) gives the exact identity
\[
 \pi(X;M,b)-B(X)=\pi(s;M,b)-B(s)
   +\frac{H(X)}{\log X}-\frac{H(s)}{\log s}
   +\int_s^X\frac{H(t)}{t(\log t)^2}\,dt.
\]
It follows either by integration by parts for the prime-mass step
function or by applying ordinary integration by parts separately to
\(F'(t)/\log t\). In particular the normalization of \(\ell\) at
\(2\) is retained at the lower endpoint. The two boundary terms
in \(H\) and the integral together are
\(O(X w(X) E_c(X,M)/(\varphi(M)\log X))\), since
\(\log s=\tfrac12\log X\) and
\(\int_s^Xdt/(\log t)^2\le4X/(\log X)^2\).
Also \(|\pi(s;M,b)-B(s)|\ll s\), by the elementary upper bounds
\(\ell(z)\ll z\) and \(\pi(s;M,b)\le s\).
As \(M\le X^{1/200}\), \eqref{eq:window-ap-main-lower} gives
\(s/B(X)\ll X^{-49/100}\log X\ll\exp(-\sqrt{c\log X})\).
This proves \eqref{eq:window-relative-pnt} for the authors' option.

\paragraph{The release's wider exceptional interval.}
If the authors' pair exists, it is eligible, so choose it. Otherwise,
if there is an eligible pair, fix any such \((\chi,\beta)\) once for
this \(M\). Its distance from \(1\) satisfies
\(1-\beta>1/(50\log M)\). Uniformly for \(\beta>1/2\) and
\(X^\beta\ge4\),
\[
 0\le\frac{\ell(X^\beta)}{\ell(X)}
     \ll X^{\beta-1}
     \le \exp\!\left(-\frac{\log X}{50\log M}\right).
\]
Here \(\ell(z)\asymp z/\log z\) for \(z\ge4\), obtained by splitting
the upper estimate at \(\sqrt z\) for large \(z\) and integrating
over \([z/2,z]\) for the lower estimate. Increase \(C\) so the last
bound, with its implied constant, is at most \(1/2\).
The newly selected main term is then between one half and twice
\(\ell(X)/\varphi(M)\), and their difference is bounded by a constant
times the new main term times \(E_c(X,M)\) on taking \(c\le1/50\).
If no eligible pair exists the empty option already works.
This proves the asserted choice for all \(M,X,b\) and also preserves
the lower bound \(B(X)\gg X/(M^2\log X)\). In particular the proof
requires uniqueness only for zeros in the narrower interval.

\paragraph{The actual prime interval.}
For every \(\epsilon>0\) there exists \(Y_0(\epsilon)\) such that
for every \(Y\ge Y_0(\epsilon)\), every \(m>0\) divisible by \(8\),
every \(a_0\) coprime to \(m\), and every \(V>0\) satisfying
\[
  \tfrac14(\log Y)^2\log m\le\log V,\qquad
  Y^{\epsilon/6}\le \log V\le Y^{\epsilon/5}(\log Y)^2,
\]
the set \(\mathcal P=\{p\text{ prime}:\sqrt V<p\le V,\
p\equiv a_0\bmod m\}\) satisfies
\begin{equation}
                       |\mathcal P|\ge V^{9/10}.
 \label{eq:window-ap-lower-count}
\end{equation}
Put \(L=\log Y\), \(u=\log V\).
Then \(\log m\le4u/L^2\) and \(\log m\le4Y^{\epsilon/5}\), whence
\(\log(\log m+2)=O_\epsilon(L)\).
Thus \(C\log m\log(\log m+2)\le u\) for sufficiently large \(Y\).
Both exponentials in \(E_c(V,m)\) tend to zero uniformly, because
\(u/\log m\ge L^2/4\) and \(u\ge Y^{\epsilon/6}\).
Consequently \(\pi(V;m,a_0)\ge B(V)/2\). For large \(Y\) one has
\(m\le V^{1/100}\). The preceding main-term lower bound gives
\(\pi(V;m,a_0)\gg V^{49/50}/\log V\). Subtracting at most
\(\sqrt V\) primes leaves at least \(V^{9/10}\) for sufficiently
large \(V\), uniformly in \(a_0,m\). The lower hypothesis on \(u\)
ensures that this last threshold follows from \(Y\ge Y_0(\epsilon)\).
The formal statement also quantifies over nonpositive real \(V\).
Under its real-power convention \(V^{9/10}=|V|^{9/10}\cos(9\pi/10)\)
for \(V<0\), and \(0^{9/10}=0\). These cases satisfy the inequality
by nonnegativity of the cardinality.

There is also uniformity over the enlarged moduli occurring in
the character sums. Given fixed \(D,B_0>0\), take \(M=ms\)
with any positive integer \(s\le Y^D\). Then
\(\log M\le4u/L^2+DL\le8u/L^2\) once \(Y\) is large, since
\(Y^{\epsilon/6}/L^3\) tends to infinity. The upper bound
\(\log M\le4Y^{\epsilon/5}+DL\) still gives
\(\log(\log M+2)=O_{\epsilon,D}(L)\).
Thus \eqref{eq:window-relative-pnt} applies with
\(E_c(V,M)=o(Y^{-B_0})\), uniformly over \(s\) and every reduced
class modulo \(M\). In particular one can use \(D=800K\) and
\(B_0=K+1\) for fixed \(K\ge1\) in Section~\ref{sec:window-exceptions}.

\paragraph{Exceptional pairs at two different moduli.}
For the use of Landau--Page uniqueness in
Section~\ref{sec:window-exceptions}, let \(H\ge3\), take integers
\(1\le M_1,M_2\le H\), and suppose nonprincipal real characters
\(\chi_i\bmod M_i\) have real zeros
\[
                  1-\frac{1}{100\log H}<\beta_i<1 .
\]
Put \(q=M_1M_2\). The induced character \(\widetilde\chi_i\bmod q\)
is nonprincipal, since the reduction map on unit groups is
surjective. This surjectivity follows by lifting a unit modulo
\(M_i\) to an integer congruent to \(1\) at each prime divisor of
\(q\) absent from \(M_i\), using the Chinese remainder theorem.
The identity
\[
 L(s,\widetilde\chi_i)=L(s,\chi_i)
   \prod_{\substack{p\mid q\\p\nmid M_i}}(1-\chi_i(p)p^{-s})
\]
preserves the zero at \(\beta_i\).
Since \(\log q\le2\log H\), both zeros lie in
\([1-1/(50\log q),1)\). Fixed-modulus uniqueness at \(q\) now yields
\(\widetilde\chi_1=\widetilde\chi_2\) and \(\beta_1=\beta_2\).
In particular \(\chi_1(n)=\chi_2(n)\) for all integers \(n\)
coprime to \(M_1M_2\). This supplies precisely the cross-modulus
conclusion of \texttt{MSqPage.page\_across} and verifies that the
same primitive character accounts for all such zeros. The constant
\(1/100\) is independent of the two moduli and the window. The
fixed-modulus theorem is due to K.~S.~McCurley,
\href{https://doi.org/10.1016/0022-314X(84)90089-1}{\emph{Explicit
zero-free regions for Dirichlet \(L\)-functions}}, J.~Number Theory
19 (1984), 7--32, in the form stated by Thorner and Zaman.

\paragraph{Source correspondence.}
In the tagged release the definitions of \(\ell,\pi,B\) and
\srcid{PNTAPRelW} are in \srcid{MeanSquareHypWIff}. The declaration
\texttt{pnt\_progressions\_relative\_weak\_holds} in
\srcid{PNTFinal} proves the estimate internally. Its
\srcid{PNTGlue} separates the ranges
\(\log M\le c_A\sqrt{\log t}\) and \(\log M>c_A\sqrt{\log t}\).
The latter uses the smoothed explicit formula, a differenced
explicit formula, the zero-density estimate, and the
Deuring--Heilbronn estimate. \srcid{PrimeCountOfPsi} converts the
combined estimate to the unweighted prime count. The proof above
uses the published weighted theorem to establish the same
\srcid{PNTAPRelW} directly. The window count
\eqref{eq:window-ap-lower-count} is \texttt{MSqCount.count\_a}.
These modules are in
\srcid{Formalism.GL3MeanSquares.GenusMeanAsymptotics.JacobiCharacterSums}.
The common-window application is
\srcid{Development.AutomorphicArithmetic.MeanSquareBounds.twistK}.

\section{The truncated twist and the logarithmic derivative at one}
\label{sec:window-explicit}

Let \(\chi\) be a nonprincipal Dirichlet character of modulus \(m\ge1\).
Write \(\ell_\chi(s)=-L'(s,\chi)/L(s,\chi)\) wherever defined,
and \(a_n=\Lambda(n)\chi(n)\). Write \(Q\ge2\) for the height
and conductor cutoff. The box hypothesis is
\[
 m\le Q,\qquad L(s,\chi)\ne0
 \quad\bigl(\Re s>399/400,\ |\Im s|\le Q\bigr).
\]
For \(v>0\) define
\[
 A_\chi(v)=\sum_{n\le v}a_n(v/n-1),\qquad
 G_\chi(v)=A_\chi(v)-v\ell_\chi(1),\qquad
 R_\chi(z)=\ell_\chi(1)-\sum_{n\le z}a_n/n .
\]
Positive-integer summation includes the floor of each real cutoff.
For \(v\le z\), direct separation at \(n=v\) gives the exact identity
\begin{equation}
 G_\chi(v)=-\sum_{n\le z}a_n\min(1,v/n)-vR_\chi(z).
 \label{eq:box-finite-split}
\end{equation}
The choice of sign records that the source centers by adding
\(vL'(1,\chi)/L(1,\chi)\).

\subsection{An explicit contour estimate}

For \(0<\delta\le1/800\), \(z\ge1\), \(0<\eta\le1\), and
\(h=z\eta\), we claim
\begin{equation}
 |A_\chi(z+h)-A_\chi(z)-h\ell_\chi(1)|
 \le C_\delta\log Q
   \left\{(\log Q)\eta z^{399/400+\delta}+z^2/Q\right\}.
 \label{eq:box-contour}
\end{equation}
Put \(\sigma=399/400+\delta\), \(T=Q/2\) and \(L_Q=\log Q\).
First fix the bound
\[
 |\ell_\chi(w)|\le B_\delta L_Q\quad
          (\Re w\ge \sigma,\ |\Im w|\le Q-1).
\]
In the range \(\Re w<11/10\), the source's local logarithmic
derivative lemma for \(L\) at \(w\) expands \(L'/L\) as a sum of
terms \(m_\rho/(w-\rho)\) for zeros in the closed disk centered at
\(21/20+i\Im w\) of radius \(11/60\), plus an error.
If \(C\) is the absolute polynomial-growth constant in
\(|L(s,\chi)|\le C m(1+|s|)\) on \(\Re s\ge1/2\), let
\(A=\log(42(|C|+1))/\log2+1\).
The local lemma bounds the error by
\((1440/11)A\log(m(|\Im w|+2))\), and the sum of the
multiplicities by \(A\log(m(|\Im w|+2))/\log2\).
All these zeros have imaginary part at most \(Q\) in absolute
value, so the box forces \(\Re\rho\le399/400\). Therefore
\(|w-\rho|\ge\delta\), and \(m(|\Im w|+2)\le Q^3\).
In the range \(\Re w\ge11/10\), the Euler-product derivative
bound \(|\ell_\chi(w)|\le 1/(\Re w-1)+C_R\) suffices. One may take
\[
 B_\delta=3A\left(\frac{1440}{11}
               +\frac{1}{\delta\log2}\right)
                  +\frac{10+C_R}{\log2}.
\]


Here is the disk argument underlying those two estimates. Suppose
\(f\) is holomorphic in the disk of radius \(r\) about \(z_0\),
continuous on its closure, \(f(z_0)\ne0\), and
\(|f|\le M|f(z_0)|\) there, \(M\ge1\).
Take every zero in the closed concentric disk of radius \(r/3\),
with its multiplicity, and let \(P\) be the product of the
corresponding monic linear factors. The set is finite by isolated
zeros and compactness. Write \(k=\deg P\) and \(g=f/P\), extending
across the zeros by their local power series. On the outer circle,
\(|P|\ge(2r/3)^k\), whereas \(|P(z_0)|\le(r/3)^k\). The maximum
modulus principle gives
\[
             2^k |g(z)|\le M |g(z_0)|\qquad (|z-z_0|\le r).
\]
At \(z_0\) this implies \(k\le\log M/\log2\). The quotient \(g\)
has no zeros in the disk of radius \(R=r/3\). Take its normalized
holomorphic logarithm \(u=\log(g/g(z_0))\) there. Then \(u(z_0)=0\)
and \(\Re u\le\log M\). For \(a>\log M\), the function
\(u/(2a-u)\) has modulus at most one and vanishes at \(z_0\).
Schwarz's lemma and the triangle inequality give
\[
 |u(z)|\le\frac{2a|z-z_0|}{R-|z-z_0|}\le6a
       \qquad (|z-z_0|\le3R/4).
\]
At any point within \(R/2\) of the center use the Cauchy estimate
on the circle of radius \(R/4\) about that point. Letting
\(a\) decrease to \(\log M\) gives \(|u'|\le24\log M/R\).
Logarithmic differentiation of \(f=Pg\) therefore gives
\[
 \left|\frac{f'(z)}{f(z)}
           -\sum_{|\rho-z_0|\le r/3}\frac{m_\rho}{z-\rho}\right|
     \le72\,\log M/r \qquad(|z-z_0|\le r/6,\ f(z)\ne0).
\]
For \(f=L(\cdot,\chi)\), use \(z_0=21/20+i\Im w\),
\(r=11/20\). The disk lies in \(\Re s\ge1/2\). On its closure
\(1+|s|\le2(|\Im w|+2)\); at the center the Euler-product inverse
bound gives \(1/|L(z_0,\chi)|\le21\). Thus one may use
\(M=42(|C|+1)m(|\Im w|+2)\).
Because \(m(|\Im w|+2)\ge2\), its logarithm satisfies
\(\log M\le A\log(m(|\Im w|+2))\).
The points in the range
\(\sigma\le\Re w<11/10\) satisfy \(|w-z_0|\le r/6\).
For completeness, on \(\Re s>0\) partial summation writes
\(L(s,\chi)=s\int_1^\infty P_\chi(t)t^{-s-1}\,dt\) with
\(P_\chi(t)=\sum_{n\le t}\chi(n)\). Full periods have zero sum,
so \(|P_\chi(t)|\le m\) and \(|L(s,\chi)|\le m|s|/\Re s\).
This supplies the asserted growth on the whole disk. For
\(\Re s>1\), the absolute Möbius series of the inverse Euler
product is bounded by \(\zeta(\Re s)\le1+1/(\Re s-1)\);
at \(z_0\) this is \(21\). Finally the only right-half-plane
bounds used above are on \(\Re s\ge11/10\) and \(\Re s=2\).
Both follow directly by majorizing the logarithmic-derivative
series by the convergent constant
\(\sum_{n\ge2}(\log n)n^{-11/10}\), increasing \(C_R\) if needed.
These computations justify the specializations of the local
expansion, including \(72/r=1440/11\).

For \(y\ge1\) the integrated Perron formula on the line \(2+it\)
gives \(A_\chi(y)\) as the integral of
\[
      F_y(w)=\ell_\chi(w)\frac{y^w}{w(w-1)}
\]
with measure \(dw/(2\pi i)\) upwards. The formula follows by the
absolutely convergent Dirichlet series for \(\ell_\chi\) on this
line and the elementary Mellin inversion of \((u-1)_+\), whose
transform is \(1/[w(w-1)]\). At \(u=1\) the summand is zero from
either side. Shift the finite part of the line between heights
\(-T,T\) to \(\Re w=\sigma\). The box excludes zeros on and
inside the rectangle. The pole of the kernel at \(w=1\)
contributes \(y\ell_\chi(1)\). The other pole, \(w=0\), lies outside the contour.
For the remaining error the two tails on the right are accompanied
by the horizontal segments at \(\pm T\) and the left vertical
segment. Apply this identity at \(y=z+h\) and \(y=z\), and
subtract before estimating the left segments.

On the right \(|\ell_\chi(2+it)|\le1+C_R\) and the integrand
is at most \((1+C_R)y^2/t^2\) for nonzero \(t\).
Since \(y\le2z\), each pair of omitted tails after division
by \(2\pi\) is bounded by \(8(1+C_R)z^2/Q\). On each horizontal
segment the denominator has modulus at least \(T^2\), the length
is at most \(2\), and the logarithmic derivative is bounded by
\(B_\delta L_Q\), since \(T\le Q-1\). Its unnormalized integral
is bounded by \(8B_\delta L_Q z^2/T^2\le16B_\delta L_Q z^2/Q\).

For \(w=\sigma+iu\), integrate the derivative of
\(\exp(tw\log(1+\eta))\), \(0\le t\le1\), to see that
\[
 |(z+h)^w-z^w|\le4\eta |w|z^\sigma .
\]
Thus the factor \(w\) in the kernel cancels in the magnitude of the
difference. Since \(1-\sigma\ge1/800\),
\[
 \int_{-T}^T\frac{du}{|\sigma-1+iu|}
 \le\int_{-T}^T\frac{2\,du}{1/800+|u|}
 =4\log(1+800T)\le40\log Q .
\]
The left difference before division by \(2\pi\) is at most
\(160 B_\delta L_Q^2\eta z^\sigma\). There are two pairs of
right tails and four horizontal segments. Using
\(L_Q\ge\log2\ge1/2\) and \(2\pi\ge1\), their bounds and
the vertical bound give \eqref{eq:box-contour} with
\(C_\delta=160B_\delta+32(1+C_R)\).

The independent error at the finite cutoff can be read directly
from the finite sum defining \(A_\chi\):
\[
 A_\chi(z+h)-A_\chi(z)-h\sum_{n\le z}a_n/n
     =\sum_{z<n\le z+h}a_n((z+h)/n-1).
\]
There are at most \(h+1\) summands, \(|a_n|\le\log(z+h)\),
and the factor is between zero and \(h/z\). Subtracting the
contour expression and dividing by \(h\) proves
\begin{equation}
 |R_\chi(z)|\le
 \frac{C_\delta L_Q\{L_Q\eta z^\sigma+z^2/Q\}
              +h(h+1)\log(z+h)/z}{h}.
 \label{eq:box-tail}
\end{equation}
This estimate applies to a general nonprincipal character obeying
the box and modulus conditions.

\subsection{Uniformity across a real interval}

Fix \(0<\epsilon\le1/4,\ \tau>0,\ K\ge1\) and take
\[
    Y^{\epsilon/6}\le \log Q\le Y^{\epsilon/4},\qquad
    x=Y^K,\quad z=x^{400},\quad \eta=x^{-2},\quad
    h=x^{398},\quad \delta=\epsilon/(1600K).
\]
In particular \(0<\delta\le1/800\).
Multiplication of \eqref{eq:box-tail} by \(x\) gives exactly
\[
 x|R_\chi(z)|\le C_\delta(\log Q)^2Y^{\epsilon/4}
   +C_\delta(\log Q)x^{403}/Q
   +\frac{(x^{398}+1)\log(x^{400}+x^{398})}{x^{399}} .
\]
Indeed \(z^\sigma=x^{399}Y^{\epsilon/4}\).
The first term is at most \(C_\delta Y^{3\epsilon/4}\), hence at
most \(Y^\epsilon/3\) for large \(Y\). Since \(\log Q\le x\)
and \(Q\ge\exp(Y^{\epsilon/6})\), the second is bounded by
\(C_\delta Y^{404K}\exp(-Y^{\epsilon/6})\), which is eventually
at most \(1/3\). For \(x\ge3\), \(x^{400}+x^{398}\le x^{401}\)
and \(x^{398}+1\le2x^{398}\). The last term is at most
\(802\log x/x\le802\), eventually at most \(Y^\epsilon/3\).
Increasing the common threshold accordingly gives
\[
                      x|R_\chi(x^{400})|\le Y^\epsilon .
\]
Now let \(q\) be the prime index and suppose
\(\chi(n)=\kappa(n)(\frac{n}{q})\) for every positive integer
\(n\). This is the primitive twist selected at the counting stage.
If throughout \(Y\le y\le Y^K\) the finite sum at
\(z=(Y^K)^{400}\) has absolute value at most
\(y^{1/2+\tau}/2\), then \eqref{eq:box-finite-split} yields
\[
 |G_\chi(y)|\le \tfrac12y^{1/2+\tau}+y|R_\chi(z)|
                    \le y^{1/2+\tau}.
\]
The last step uses \(y\le x\), and, for \(Y\ge16\),
\(Y^\epsilon\le Y^{1/4}\le Y^{1/2}/2\le y^{1/2+\tau}/2\).
The single choice of truncation applies even when the largest
finite sum is attained at a different \(y\) for each prime.

\subsection{The constant four in the logarithmic derivative bound}

The immediate bound in \texttt{box\_of\_S} is
\(|L'/L(1,\chi)|\le C\log Y\). Indeed apply the preceding tail
estimate at \(K=1\) to get \(|R_\chi(Y^{400})|\le1\), then use
\(\sum_{n\le t}\Lambda(n)/n\le C_0\log t\) with \(t=Y^{400}\). The constant
\(400C_0+1\) suffices for \(Y\ge3\).
To obtain \(4\), the source uses the sharper
\(\sum_{n\le t}\Lambda(n)/n\le\log t+2\) and changes the
truncation. For the latter Mertens estimate take \(N=\lfloor t\rfloor\).
The identity \(\sum_{d\mid n}\Lambda(d)=\log n\) yields
\[
 N\sum_{n\le N}\Lambda(n)/n
 \le\sum_{n\le N}\Lambda(n)\lfloor N/n\rfloor+\psi(N)
 =\log(N!)+\psi(N).
\]
The increasing-function integral bound gives
\(\log(N!)\le N\log N-N+1+\log N\).
The source's explicit elementary Chebyshev estimate
\(\psi(N)+1+\log N\le3N\) gives the claimed result after
division by \(N\).

Write \(l=\log Q,\ s=l^{1/4},\ x=s^{10}=l^{5/2}\). Now use
\(\delta=(2/5)/1600\) in \eqref{eq:box-tail}, with the same
\(\eta=x^{-2}\), \(z=x^{400}\). The numerical bound becomes
\[
 x|R_\chi(z)|\le C_\delta l^2 s+
                 C_\delta l x^{403}/Q+802 .
\]
If \(l\) is sufficiently large then \(l\ge81\), so \(x\ge3\),
and \(l^{1006}\le e^l=Q\). The first term without its constant
is \(s^9\le x\). The second satisfies \(lx^{403}=x l^{1006}\le xQ\).
Thus \(|R_\chi(z)|\le 2C_\delta+802\). Since
\(\log z=1000\log\log Q\), the Mertens estimate gives an absolute
constant \(C'\) such that
\[
       |L'(1,\chi)/L(1,\chi)|\le1000\log\log Q+C'.
\]
For the pair windows it suffices that \(0<\epsilon<1/52\) and,
for fixed \(A>0\),
\(
 Y^{\epsilon/6}\le\log Q\le A Y^{\epsilon/5}(\log Y)^2.
\)
The lower bound ensures the preceding threshold in \(l\) for
large \(Y\). Writing \(L=\log Y\), the upper bound gives
\[
 1000\log\log Q+C'\le200\epsilon L+2000\log L+1000\log A+C'.
\]
The first summand is at most \((50/13)L\), the second is at most
\(L/13\) eventually, and the two fixed constants together are at
most \(L/13\) eventually. Thus the expression is at most
\(4L\le4\log y\) for every \(y\in[Y,Y^K]\). In the actual
pair selection \(A=3\) suffices. The lower threshold depends
on the fixed parameters only. The constant four and the
prime-sum bound therefore hold throughout the same interval.

\paragraph{Source correspondence.}
The deductions are \texttt{BoxOfS.box\_of\_S}, \texttt{TailKExp.tailK},
\texttt{kappaC.kappaC} and \texttt{Kappa4Loglog.kappa4}, under
\srcid{Development.AutomorphicArithmetic.MeanSquareBounds}.
The rectangle, pointwise logarithmic-derivative bound, finite
difference and numerical substitutions are \texttt{TailContour},
\texttt{TailBox}, \texttt{TailDiff}, \texttt{TailKernel} and
\texttt{TailNum} under
\srcid{Formalism.GL3MeanSquares.GenusMeanAsymptotics.JacobiCharacterSums}.
The local zero expansion is proved above from its disk factorization.
The polynomial growth and Euler-product bounds proved above correspond to \texttt{LandauLemmaLfunctionZeros},
\texttt{LGrowth} and \texttt{LRight}. The exact elementary Mertens
bound and its explicit Chebyshev input occur in \texttt{mertensUpper}
under \srcid{Development.AutomorphicArithmetic.DedekindLFunctions.LandauTheorems}.

\section{From a one-sided prime-sum bound to zero-freeness}
\label{sec:window-analysis}

Let \(F=(\chi_i,a_i)_{i\in I}\) be a finite family of Dirichlet
characters, possibly of different positive moduli. Each weight \(a_i\)
is strictly positive. Each principal character in the family must have
modulus one. Write \(\Lambda(n)\) for the von Mangoldt function and
\(\gamma\) for Euler's constant. For \(y\geq1\) put
\[
 B_\chi(y)=\sum_{1\leq n\leq y}\Lambda(n)\chi(n)(y/n-1).
\]
Define
\[
 G_\chi(y)=
 \begin{cases}
 B_\chi(y)-y\log y+(1+\gamma)y,&\chi\text{ principal},\\
 B_\chi(y)+yL'(1,\chi)/L(1,\chi),&\chi\text{ nonprincipal}.
 \end{cases}
 \qquad G_F(y)=\sum_{i\in I}a_i\Re G_{\chi_i}(y).
\]
The second expression uses classical nonvanishing of \(L(1,\chi)\).
For definiteness, a zero of a meromorphic function is a point where its
order is positive. Consequently the pole of \(\zeta\) at one is excluded.
The formal library assigns a nonzero value to its total function at that
point, which gives the same assertion about zeros.

\paragraph{Window criterion.}
Fix \(0<\theta<1\), \(\delta>0\), and \(C\in\mathbb R\). Suppose that
for every \(Y_0\) there is \(Y\geq\max(Y_0,1)\) such that
\begin{equation}
 G_F(y)\leq C y^\theta\qquad(Y\leq y\leq Y^{1+\delta}).
 \label{eq:an-window}
\end{equation}
Then none of the \(L(s,\chi_i)\) has a zero with \(\Re s>\theta\).
The family, weights, exponent and constant are fixed across these
windows. Only their initial points need go to infinity.

\subsection{The transform and the sign of its residues}

Write \(g(u)=G_F(e^u)\) for \(u\geq0\). This is measurable. The individual
ramps are continuous across an integer because the new summand vanishes
there. Also \(|g(u)|\leq C_0e^{2u}\) for a constant depending on \(F\).
Indeed \(|\chi(n)|\leq1\) and \(\Lambda(n)\leq\log n\leq n\), so each
summand in \(B_\chi(y)\) has modulus at most \(y\), and there are at most
\(y\) summands. The centering terms have the same bound up to a fixed
constant. For \(\Re s>2\) consider the absolutely convergent transform
\[
 \Phi(s)=\int_0^\infty g(u)e^{-su}\,du .
\]
Changing variables to \(y=e^u\) gives the Mellin transform with measure
\(y^{-s-1}\,dy\). For \(n\geq1\),
\[
 \int_n^\infty (y/n-1)y^{-s-1}\,dy
       =\frac{n^{-s}}{s(s-1)} .
\]
The corresponding sum of absolute integrals converges because
\(\sum_{n\geq1}\Lambda(n)n^{-\Re s}\) converges for \(\Re s>1\).
Termwise integration followed by logarithmic differentiation of the
Euler product gives the following transforms \(\Phi_\chi\) of
\(G_\chi(e^u)\):
\begin{align*}
 \Phi_\chi(s)
 &=-\frac{L'(s,\chi)/L(s,\chi)}{s(s-1)}
          +\frac{L'(1,\chi)/L(1,\chi)}{s-1}
             &&(\chi\text{ nonprincipal}),\\
 \Phi_1(s)
 &=-\frac{\zeta'(s)/\zeta(s)}{s(s-1)}
          -\frac{1}{(s-1)^2}+\frac{1+\gamma}{s-1}.
\end{align*}
Here the transforms of \(y\) and \(y\log y\) are respectively
\((s-1)^{-1}\) and \((s-1)^{-2}\). The pole at one cancels in each
formula. For a nonprincipal character this follows by substitution
in the numerator after putting the terms over a common denominator.
For the principal character use
\(\zeta'(s)/\zeta(s)=-1/(s-1)+\gamma+O(s-1)\), which follows from the
Laurent expansion of \(\zeta\) at one. This is the divided-difference
removal of the singularity in the source.

On real \(u\) we have
\(\overline{G_\chi(e^u)}=G_{\chi^{-1}}(e^u)\). Indeed the Euler
coefficients are conjugate, and \(L(s,\chi^{-1})=
\overline{L(\overline{s},\chi)}\), first by the Dirichlet series and
then by continuation. Hence the meromorphic continuation of the transform is
\[
 \Phi(s)=\frac12\sum_i a_i\bigl(\Phi_{\chi_i}(s)
                                   +\Phi_{\chi_i^{-1}}(s)\bigr).
\]
Let \(Z\) be the set of zeros in \(\Re s>0\) of all the functions
\(L(s,\chi_i), L(s,\chi_i^{-1})\). Classical nonvanishing on
\(\Re s\geq1\) implies \(Z\subset\{0<\Re s<1\}\). A compact subset
of \(\Re s>0\) meets \(Z\) in finitely many points by isolatedness of
zeros and finiteness of the family. Away from \(Z\), the continuation is
holomorphic in that half-plane. For \(\rho\in Z\), write \(m_\chi(\rho)\)
for the zero multiplicity, taking it to be zero when the value is nonzero.
Locally \(L(s,\chi)=(s-\rho)^{m_\chi(\rho)}b(s)\) with \(b\)
holomorphic and \(b(\rho)\ne0\). Its logarithmic derivative is
\(m_\chi(\rho)/(s-\rho)+b'(s)/b(s)\). We therefore obtain
\begin{equation}
 \Phi(s)=\frac{c_\rho}{s-\rho}+\text{a holomorphic function},\qquad
 c_\rho=\frac{\sum_i a_i
             (m_{\chi_i}(\rho)+m_{\chi_i^{-1}}(\rho))}
                    {2\rho(1-\rho)} .
 \label{eq:an-res}
\end{equation}
The numerator is positive since at least one multiplicity is positive.
Thus \(c_\rho\ne0\). If \(\rho\) is real then \(0<\rho<1\), and
\(c_\rho\) is positive real. This last sign handles possible real zeros
in the same argument as the nonreal ones.

\subsection{What an upper bound on one window gives}

Take \(Y\) in \eqref{eq:an-window}, put \(U=\log Y\), and fix
\(\eta>0\). On \([U,(1+\delta)U]\) the function \(g\) has the nonnegative
majorant \(R(u)=|C|e^{\theta u}\). Its weighted integral is
\begin{equation}
 \int_U^{(1+\delta)U}R(u)e^{-\theta u}\,du
         =|C|\delta U\leq e^{\eta U}
 \label{eq:an-majorant}
\end{equation}
for all sufficiently large \(U\). We can take such \(U\) arbitrarily
large. The argument uses only \(g\leq R\), \(R\geq0\) and
\eqref{eq:an-majorant} on these intervals.

Put
\[
 A_k(z)=\frac{(-1)^k}{k!}\Phi^{(k)}(z)
       =\int_0^\infty g(u)\frac{u^k}{k!}e^{-zu}\,du
       \qquad(\Re z>2).
\]
For each fixed \(k\) and each closed neighborhood with
\(\Re z\geq2+b>2\), \(C_0u^ke^{-bu}\) is an integrable majorant;
the derivative formula follows by differentiation under the integral.
Since \(g\) is real, \(A_k(\sigma)\) is real for real \(\sigma>2\).
Taking absolute values on the vertical translate and adding the
untranslated integral yields, for every real \(t\),
\[
 |A_k(\sigma+it)|+A_k(\sigma)
  \leq 2\int_0^\infty \max(g(u),0)\frac{u^k}{k!}e^{-\sigma u}\,du .
\]
This follows pointwise from \(|g|+g=2\max(g,0)\). For \(u\geq0\),
the \(k\)-th term of the exponential series gives
\(u^ke^{-(\sigma-\theta)u}/k!\leq(\sigma-\theta)^{-k}\).
The contribution from the chosen interval is at most
\(2e^{\eta U}(\sigma-\theta)^{-k}\). Outside it we have the bound
\(2C_0u^ke^{-(\sigma-2)u}/k!\).

\paragraph{Concentration with the constants of the formal proof.}
For each \(\delta>0\) set
\[
 l=2+6/\delta,\quad M=l(1+\delta),\quad \sigma=M(2+l),\quad
 k=\left\lceil (1+\delta/2)\sigma U\right\rceil .
\]
Then \(\sigma>2\), and \(k\geq U\) when \(U\geq1\). Write
\(V=(1+\delta)U,\ a=\sigma-2-l,\ b=\sigma-2+l\); both \(a,b\) are positive.
The exact tail estimate used in the proof is
\begin{equation}
 \int_{(0,\infty)\setminus[U,V]}
 \frac{u^k}{k!}e^{-(\sigma-2)u}\,du
 \leq\left(\frac{e}{a}+\frac1b\right)\sigma^{-k}.
 \label{eq:an-tail}
\end{equation}
For \(u>V\) multiply the integrand by \(e^{l(u-V)}\geq1\), and for
\(u<U\) multiply it by \(e^{l(U-u)}\geq1\). Extending each resulting
nonnegative integral to \((0,\infty)\) gives
\(e^{-lV}a^{-k-1}+e^{lU}b^{-k-1}\), since repeated integration by parts
gives \(\int_0^\infty u^ke^{-zu}\,du=k!z^{-k-1}\) for \(z>0\).
For any \(z>0\), the inequality \(\log(\sigma/z)\leq(\sigma-z)/z\) yields
\[
 e^x z^{-k-1}\leq \frac{\sigma^{-k}}{z}
   \exp\!\left(x+k\frac{\sigma-z}{z}\right).
\]
For \(z=a\) the last ratio is \(1/(M-1)\). The identities
\((1+\delta/2)\sigma=M(M-1)\) and \(lV=MU\) show that
\(-lV+k/(M-1)\leq1\). For \(z=b\), the ratio is
\(-(l-2)/b\). Here \((1+\delta/2)(l-2)=l+1\) and \(lb\leq(l+1)\sigma\),
the latter following from \(l(l-2)\leq\sigma\). Thus \(k(l-2)/b\geq lU\),
so the second exponential is at most one. These two estimates prove
\eqref{eq:an-tail}.

Put \(E=\sigma-\theta\) and \(C_*=2+2C_0(e/a+1/b)\). We have proved
that for every \(\eta>0\) and every \(U_0\), some \(U\geq U_0\) and
integer \(k\geq U\) satisfy \emph{simultaneously for all} real \(t\)
\begin{equation}
 |A_k(\sigma+it)|+A_k(\sigma)\leq C_*e^{\eta U}E^{-k}.
 \label{eq:an-window-coeff}
\end{equation}
Indeed \(\sigma^{-k}\leq E^{-k}\) and we may choose \(U\geq1\).
The constants \(\sigma,C_*\) are independent of \(\eta,U_0,U,k,t\).

\subsection{Detecting a zero with high derivatives}

Suppose some \(\rho_0\in Z\) has \(\Re\rho_0>\theta\). The disk
\(\overline B(\sigma,E)\) is contained in the right half-plane since
\(\sigma-E=\theta>0\). If this disk contains a pole, let \(D\) be the
distance to the closest one; otherwise put \(D=E\). The open disk
\(B(\sigma,D)\) is pole-free.

First suppose some \(\rho_1\in Z\) satisfies \(\sigma-\Re\rho_1<D\).
At \(z_0=\sigma+i\Im\rho_1\) its distance is already \(<D\). Perturb
the imaginary part slightly, keeping this strict inequality, to obtain
\(z=\sigma+it\) with a unique nearest pole \(\rho\) and
\(d=|z-\rho|<D\). Here are the details. Choose \(\epsilon>0\) smaller
than both \(D-|z_0-\rho_1|\) and \(\sigma-D\). Only finitely many
poles lie in \(\overline B(z_0,D+\epsilon)\). For each pair from this
finite set with different imaginary parts, equality of squared distances
to \(\sigma+it\) is a linear equation in \(t\) with a nonzero coefficient.
For a pair of distinct poles with the same imaginary part, equality
would force equal real parts since both are to the left of \(\sigma\).
Choose \(t\) within \(\epsilon\) of \(\Im\rho_1\) outside the finite
exceptional set. All poles outside that disk have distance \(>D\)
from the chosen center. This proves the assertion.

Choose \(h\) strictly between \(d\) and the minimum of \(D\) and the
distance from \(z\) to any competing pole, using only poles in the
above compact disk. By \eqref{eq:an-res}, removing \(c_\rho/(s-\rho)\)
leaves a holomorphic function on a neighborhood of \(\overline B(z,h)\).
At the real center \(\sigma\), \(\Phi\) is holomorphic on a neighborhood
of the closed disk of radius \(h\), since \(h<D\). Cauchy's estimate on
these two circles gives \(M_1,M_2\), independent of \(k\), with
\[
 \left|A_k(z)-\frac{c_\rho}{(z-\rho)^{k+1}}\right|\leq M_1 h^{-k},
 \qquad |A_k(\sigma)|\leq M_2h^{-k}.
\]
Consequently the left side of \eqref{eq:an-window-coeff} is at
least \((|c_\rho|/d)d^{-k}-(M_1+M_2)h^{-k}\).

In the remaining case \(\sigma-\Re w\geq D\) for every pole \(w\).
The existence of \(\rho_0\) gives \(D<E\), so the defining disk
contains a closest pole. A pole at distance \(D\) has
\(D^2=(\sigma-\Re w)^2+(\Im w)^2\), forcing
\(w=\beta=\sigma-D\). Thus the closest pole is real and unique.
Take \(d=D\) and \(d<h<E\) small enough to exclude every other pole
from \(\overline B(\sigma,h)\). The same Cauchy estimate gives
\[
 A_k(\sigma)\geq(c_\beta/d)d^{-k}-M_1h^{-k}
\]
because \(c_\beta>0\) by \eqref{eq:an-res}. At \(t=0\) this also
bounds the left side of \eqref{eq:an-window-coeff} from below.

In either case there are \(c>0,\ 0<d<h\leq E\) and \(M'\geq0\)
such that \(c d^{-k}\leq M'h^{-k}+C_* e^{\eta U}E^{-k}\) on the
selected orders. Fix \(0<\eta<\log(h/d)\). Multiplication by \(d^k\)
and \(U\leq k\) bound the right side by
\(M'(d/h)^k+C_*(e^\eta d/h)^k\), which tends to zero as \(U_0\)
increases. This contradicts \(c>0\). All poles therefore satisfy
\(\Re\rho\leq\theta\). By construction each zero of a member of the
family with \(0<\Re s<1\) is one of these poles. The classical
nonvanishing theorem covers \(\Re s\geq1\), completing the window
criterion.

\subsection{The finite family attached to a primitive character}

Fix a primitive \(\psi\) of positive modulus \(N\). Let \(\Xi\) be the
finite cyclic group generated by \(\psi\) in the group of characters
modulo \(N\), and write \(o=|\Xi|\). Let \(S\) be the set of primitive
characters that induce the members of \(\Xi\), counted once.
It contains the modulus-one character, and it contains \(\psi\) because
\(\psi\) is already primitive. It is closed under inversion: induction
to a higher modulus commutes with inversion, and the conductors of
\(\chi\) and \(\chi^{-1}\) divide each other. Every principal member
of \(S\) therefore has modulus one, as also follows from the definition
of its conductor. Take the family indexed by \(S\), with weight \(1\)
on the principal character and \(1/(3o+1)\) on each other character.
For the cyclic level \(\mathcal L\) the formal construction chooses
no exceptional distinguished character. Its statistic is exactly
\[
 G_{\mathcal L}(y)=G_1(y)+\frac{1}{3o+1}
                  \sum_{\chi\in S\setminus\{1\}}\Re G_\chi(y)
                 =G_F(y).
\]
In this display \(G_1\) is real. Thus the preceding criterion applies
to every retained character with the arithmetic window held fixed.

Here is the precise arithmetic premise consumed by the analytic argument.
For this \(\mathcal L\) there exist \(e_D,e_L\in\mathbb R\) and
\(\varepsilon_1>0\) such that for every \(0<\varepsilon<\varepsilon_1\)
there are intervals \([Y,Y^{3/2}]\) with arbitrarily large \(Y\) on
which, simultaneously for all \(y\) in the interval,
\begin{equation}
 G_{\mathcal L}(y)\leq y^{1/2+(e_D+e_L+7)\varepsilon}.
 \label{eq:an-arith-input}
\end{equation}
In applying the formal comparison one first discards intervals lying
below the comparison threshold of the chosen family. The threshold
can depend on \(\varepsilon\), which is fixed before choosing \(Y\).

Take \(s\) with \(\Re s>1/2\), and put
\[
 t=\tfrac12\min(\Re s-1/2,1/2),\qquad
 \theta=1/2+t,\qquad
 \varepsilon=\min\!\left(\varepsilon_1/2,\,
                     \frac{t}{|e_D+e_L+7|+1}\right).
\]
Then \(1/2<\theta<1,\ \theta<\Re s\) and
\((e_D+e_L+7)\varepsilon\leq t\). On each of the selected windows
above one, \eqref{eq:an-arith-input} implies \eqref{eq:an-window} with
\(C=1,\delta=1/2\). The window criterion gives \(L(s,\psi)\ne0\),
subject to the pole convention already specified. This holds for
every primitive character because the construction can start with
any \(\psi\). The estimate \eqref{eq:an-arith-input} is the imported
arithmetic premise of this chapter; the analytic deduction does not
prove the arithmetic estimate.

\paragraph{Correspondence with the source.}
The ramp, positivity and admissibility conditions are the definitions
in \texttt{Deduction.\allowbreak{}Char\allowbreak{}Family}. The transform and residue calculation
translate \texttt{Deduction.\allowbreak{}pole\allowbreak{}Data}. The moment estimate and geometry
of the nearest pole translate \texttt{landau\_\allowbreak{}tails},
\texttt{landau\_\allowbreak{}window\_\allowbreak{}key}, \texttt{landau\_\allowbreak{}dichotomy},
\texttt{landau\_\allowbreak{}nearest} and \texttt{landau\_\allowbreak{}window}. The transition
between window hypotheses and zero-freeness is \texttt{lemma\allowbreak{}L} and
\texttt{zero\allowbreak{}Free\_\allowbreak{}of\_\allowbreak{}rpow\_\allowbreak{}window}. The cyclic family and the choice of
the exponent translate \texttt{Upper\allowbreak{}Theory.\allowbreak{}Comparison\ensuremath{{}_2}.\allowbreak{}level\_\allowbreak{}family},
\texttt{form\allowbreak{}B}, \texttt{Formalism.\allowbreak{}HC.\allowbreak{}level\_\allowbreak{}of\_\allowbreak{}tied\allowbreak{}Win} and
\texttt{zero\allowbreak{}Free\allowbreak{}On\_\allowbreak{}of\_\allowbreak{}tied\allowbreak{}Win\_\allowbreak{}cyc}. Arithmetic supplies the paired
inputs to the last theorem through \texttt{level\_\allowbreak{}gm\_\allowbreak{}le\_\allowbreak{}on\_\allowbreak{}family\_\allowbreak{}scales}
and \texttt{Upper\allowbreak{}Theory.\allowbreak{}Top.\allowbreak{}zero\allowbreak{}Free\allowbreak{}All\_\allowbreak{}of\_\allowbreak{}top\allowbreak{}G}.

\section{Reading the mathematical proof terms}
\label{sec:kernel-reading}
The accompanying literal edition prints the complete expressions of the
published solution export as numbered equations in dependent type theory.
This section specifies how to read a theorem value. It also relates the
short final line of the argument to the number-theoretic statement in the
first section. The edition contains shared subexpressions. Expanding each
occurrence into a separate tree would repeat the same mathematical
derivations many times.

\paragraph{Numbers, contexts and substitution.}
\(N_i\) is a name, \(U_i\) a universe level and \(E_i\) an expression.
These three numberings are independent of the printed record numbers.
\(N_0\) is the anonymous root and \(U_0=0\). A name is extended by
a string or by a natural number; these are distinct constructors.
Universe levels can be parameters, successors, maxima, and
\(\operatorname{imax}(u,v)\), which is \(0\) for \(v=0\) and
\(\max(u,v)\) otherwise. These formulas hold pointwise for all
assignments of natural numbers to the universe parameters.

Write \(\Gamma\vdash t:A\) for a typing judgment in the ordered context
\(\Gamma\). In an expression \(b_i\) is the variable \(i\) binders
outward, counting from zero. Thus both \(b_0\) in the body of
\(\lambda x.\,b_0\) and \(b_0\) in the body of
\(\lambda x.\lambda y.\,b_0\) can be represented by the
\emph{same} numbered expression. Their meanings are \(x\) and \(y\)
respectively. When extending a context, the indices referring to its older
variables increase by one. Substitution in a dependent body performs this
shift as it crosses internal binders and removes the outside binder that
has been supplied.

The basic rules, with bound variables written by name for legibility, are
\[
 \Gamma\vdash\operatorname{Sort}(u):\operatorname{Sort}(u+1),
 \qquad
 \frac{\Gamma\vdash A:\operatorname{Sort}(u)
       \quad\Gamma,x:A\vdash B:\operatorname{Sort}(v)}
      {\Gamma\vdash\prod_{x:A}B:\operatorname{Sort}(\operatorname{imax}(u,v))}.
\]
In particular \(\operatorname{Sort}(0)\) is the type of propositions.
A proposition universally quantified over a type in any universe is
again a proposition. The function introduction and elimination rules
read
\[
 \frac{\Gamma,x:A\vdash t:B}{\Gamma\vdash(\lambda x:A.\,t):\prod_{x:A}B},
 \qquad
 \frac{\Gamma\vdash f:\prod_{x:A}B\quad\Gamma\vdash a:A}
      {\Gamma\vdash f\,a:B[a/x]} .
\]
These rules presuppose well-formed contexts and types. A constant refers
to a declaration by \(N_i\), with the listed levels substituted in order
for its universe parameters. A let-expression checks its value against
the declared type and substitutes this value in the type of its body.
The binder brackets distinguish explicit, implicit, strict implicit and
type-class parameters. The expression contains the arguments actually
used, including those inferred in the source.

A type may be replaced by a definitionally equal one. In particular,
application of a lambda substitutes its argument, a let substitutes
its value, and an unfoldable definition substitutes its body at the
indicated universe parameters. A constructor followed by its projection
selects the indicated field. Applying an inductive recursor to a
constructor reduces by its listed rule, with parameters, fields and
recursive results in the indicated places. Functional extensional
\(\eta\)-conversion here identifies \(\lambda x.\,f\,x\) with \(f\)
when \(x\) is fresh. Proof irrelevance identifies two proofs of the
same proposition. The reference gives the full conversion and inductive
declaration rules; the mathematics in the proof terms uses this
underlying type theory.

\paragraph{A concrete piece of the export.}
In the first volume, record 60 completes the declaration of equality,
reflexivity and equality recursion. \(N_1\), \(N_9\) and \(N_{10}\)
are their respective names. Expressions \(E_4\) through \(E_6\)
assemble the type of equality:
\[
 \prod_{\{\alpha:\operatorname{Sort}(u)\}}\alpha\longrightarrow
       \alpha\longrightarrow\operatorname{Prop}.
\]
In the context \(\alpha:\operatorname{Sort}(u),\,a:\alpha\), \(E_7\) through \(E_{10}\)
assemble the expression \(a=a\).
\(E_{11},E_{12}\) bind \(a\) and \(\alpha\), giving the signature of
the reflexivity constructor. The equality recursor transports a proof of
a motive at \((a,\operatorname{refl})\) to the motive at
\((b,e)\) for \(e:a=b\). The row for its reflexivity reduction gives
the right-hand-side expression \(E_{41}\).

\paragraph{The theorem to look up.}
The named record \srcid{riemannHypothesis} supplies both a type and a
value. The type is the proposition that, for \(s\in\mathbb C\) with
\(\zeta(s)=0\), \(s\ne1\), and \(s\ne-2(n+1)\) for every \(n\ge0\),
one has \(\Re(s)=1/2\). The short value in the released source is the
declaration \srcid{rhOfAll}. Its value applies the general reflection
theorem \srcid{Formalism.VW.rh_of_every} to \srcid{allChi}.
The latter has value \srcid{UpperTheory.FinalTheorems.everyCharacter}
and quantifies over \emph{all} Dirichlet characters. This explains why
retrieving the handful of expressions in the final theorem record
leads back to the large arithmetic family. The source also exposes
separate theorem records for the primitive character statement, the
critical-strip statement for all characters, and right-half-plane
nonvanishing. In the literal publication, \srcid{edition-index.json}
locates the five endpoints by name and gives the type and value
expression numbers.

The external dependency audit for these endpoints reports precisely
propositional extensionality, classical choice, and soundness of
quotients. Propositional extensionality takes a logical equivalence of
propositions to their equality. Choice selects an element from a proof
that a type is nonempty. Quotient soundness takes related elements to
equal quotient classes. A definition or theorem with a body can be
followed to that body; a foundational axiom is read with its declared
type. The independent kernels checked the published formal export.
Coverage of the mathematical prose in this paper is recorded separately
below.
\footnote{\href{https://math.mit.edu/~drew/grh/}{Verification report
and dependency audit};
\href{https://lean-lang.org/doc/reference/latest/The-Type-System/}{rules
of the underlying type theory};
\href{https://github.com/leanprover/lean4export/blob/v4.33.0/format_ndjson.md}{format
of the published solution export}.}

\section*{Proofs still to be translated}
\addcontentsline{toc}{section}{Proofs still to be translated}
The proof of the reflection step, the determinant descent inequality,
the conversion of the pointwise slope error into a prime-sum bound, and
the passage from an admissible family of intervals to zero-freeness
are given above. The following inputs carry substantive arithmetic or
geometric content which the reader still needs for an independent proof.
The newer Part I proofs additionally supply the finite-level
Selmer estimates and auxiliary groups in
Section~\ref{sec:sr-arithmetic-selmer-producers}, the selected
curve regularity and distinguished local invariant calculation in
Section~\ref{sec:sr-curve-regular-cosets}, and the actual Case II
Hecke quotient in Subsection~\ref{subsec:sr-case-two-hecke}.
Sections~\ref{sec:sr-classical-doubling-inputs},
\ref{sec:sr-uniform-cusp-mean}, and
\ref{sec:sr-quadratic-period-witness} respectively match the classical
doubling inputs, prove the canonical mean bound, and construct the
quadratic period witness. The remaining clauses below are read with
these additions; they continue to require the compatible arithmetic
and geometric realizations explicitly identified in Part I.


The new Part I construction of the entry space and sharp ideal is in
Sections~\ref{sec:entry-space-flat} and~\ref{sec:entry-nice-sharp}.
Section~\ref{sec:sr-rank-four-actual-endpoint} gives the selected
rank-four tower and descent with the full local window predicate.
Section~\ref{sec:sr-rank-three-host-endpoints} proves the finite
rank-three host implication and distinguishes the actual support
output from the still-needed original-level clause.
Sections~\ref{sec:sr-ramified-matching-construction},
\ref{sec:sr-unramified-split-matching},
\ref{sec:sr-inert-matching-construction},
\ref{sec:sr-good-split-coinvariants}, and
\ref{subsec:sr-theta-decay-closed} supply all finite matching,
the exact coinvariant condition and the full ordinary decay clause.
The remaining global theta, rank-recursion and covolume statements
are independent of those local and convergence assertions.



\begin{itemize}
\item Sections~\ref{sec:selmer-generators}--\ref{sec:minus-derivatives}
give the generator count, the signed annihilator, the construction and
local relations of the derivative classes, and the simultaneous choice
of auxiliary primes. Proposition~\ref{prop:sr-PB} assembles the complete
prime-to-degree minus estimate from those printed arguments and the
Bernoulli integer calculation. Proposition~\ref{prop:sr-dyadic-generators}
proves the separate 24-generator statement for the unramified cyclotomic
module. Proposition~\ref{prop:sr-dyadic-descent} and
\eqref{eq:sr-dyadic-return} give its finite-twist transport and the
return to the CM field. Section~\ref{sec:sr-dyadic-bounds}
proves the cokernel growth, the fixed 24-element exceptional set and
the discriminant bound. Theorem~\ref{thm:sr-dyadic-window} proves the
uniform common-window estimate conditional on \(K_2(C_K)\).
Section~\ref{sec:sr-global-cohomology} proves that input by a separate
integral-homology argument, so Proposition~\ref{prop:sr-K2-proved}
discharges it with \(C_K=24\). On the rank-three side,
Sections~\ref{sec:common-trace-integer}--\ref{sec:deformation-euler}
derive the trace product and the full graded intersection inequality.
Sections~\ref{sec:kato-rigidity} and~\ref{sec:type-one-rigidity} derive
tangent rigidity at the identified modular points from the stated
Galois-theoretic hypotheses. The spectrum count is derived from the arithmetic stabilizer bound,
cohomological realization and cuspidal injection. Section~\ref{sec:ordinary-global-descent} identifies the twist and descends
the specified representation across the cyclic layers. The global deformation
model, the ordinary automorphy input over the auxiliary field, and the stated
local--global compatibility theorems still enter the comparison.
\item Sections~\ref{sec:sw-boundary} and~\ref{sec:theta-decay-proof}
give the two-cell argument, the normalization of the definite theta
value, the decay estimate from a nonzero minor, and the indefinite
chamber integral. Section~\ref{sec:sw-singular-vacuum} evaluates all the singular
vacuum tests at the distinguished parameter when \(m<n\). The differentiated
sections, the other singular terms, Fock-space generation, automorphic reduction and coefficient estimates
still enter these derivations.
\item Sections~\ref{sec:component-volume}--\ref{sec:hs-component}
derive the component degree and rank comparisons, the equality of
covolumes, the Chow-form determinant and its metric, and the uniform
height budget. The short period basis and the characteristic-zero harmonic
jet are derived in the preceding sections. Inputs include the toroidal
model and proportionality theorem, local isometry, strong approximation,
and the construction of the global integral model with its level structures.
Sections~\ref{sec:boundary-periods} and~\ref{sec:jet-connection}
identify the boundary periods and derive the connection and twist comparisons
on their respective domains. Section~\ref{sec:key-different} gives the fourth-root descent and the separate
different bound. Section~\ref{sec:key-rule-dvr} proves the integral comparison at residue
primes dividing an auxiliary denominator by changing invariants and comparing
their powers. The transverse intersections, determinant comparison and
connected-family input used there are specified in that section. Sections~\ref{sec:theta-translate-span}--\ref{sec:theta-unit-basis}
prove the integral spanning, the valuation statement and the primitive
frame from the specified local data. The horizontal and
mixed-characteristic comparisons use different local ideals.
\item Section~\ref{sec:window-exceptions} proves the character-count
bound and derives the required mean-square estimate on the two prime
progressions. Section~\ref{sec:window-explicit} derives the truncated explicit
formula and the logarithmic derivative bound \(4\log y\). Section~\ref{sec:host-field} constructs the cyclotomic CM field with its
conductor, local invariants and signature. Section~\ref{sec:window-relative-pnt} proves the required relative estimates,
including the progression denominator and comparison of possible exceptional
zeros. Section~\ref{sec:level-prime-selection} proves the exceptional-divisor
count and the simultaneous selection of the level primes. It also derives
the required norm-prime lower bound from the stated ray-class family
estimates. Section~\ref{sec:ray-real-gap} proves the polynomial value bound
and the real-zero gap from the counting, theta and positive-convolution
identities. Section~\ref{sec:quadratic-ray-family} constructs the family and its
shifted lattice, and proves the counting, theta, convolution and support
identities. Section~\ref{sec:ray-explicit} proves its zero-free strip,
uniqueness and simplicity of the possible exceptional zero, and the
uniform explicit formula used for the prime selection. The remaining
automorphic hypotheses and the passage from local geometric estimates to
the global model still enter the common-window construction.
\end{itemize}
These are outstanding translations from the formal development.
The independent verification of the released formal proof is recorded
in the source note at the start of Part II.
\end{document}
